\documentclass[11pt]{article}
\usepackage[a4paper,margin=25mm]{geometry}
\usepackage[T1]{fontenc}
\usepackage{lmodern,microtype}
\usepackage{amsmath,amssymb,amsthm,mathtools,bm,mathrsfs}
\usepackage{booktabs,array,tabularx,enumitem,aliascnt}
\usepackage{xurl}
\usepackage[hidelinks]{hyperref}
\usepackage[nameinlink,noabbrev]{cleveref}
\usepackage{etoolbox}
\makeatletter
\patchcmd{\l@section}{\setlength\@tempdima{1.5em}}{\setlength\@tempdima{2.4em}}{}{\PackageError{bridge}{TOC width patch failed}{}}
\makeatother
\newtheorem{theorem}{Theorem}[section]
\newaliascnt{proposition}{theorem}
\newtheorem{proposition}[proposition]{Proposition}
\aliascntresetthe{proposition}
\crefname{proposition}{proposition}{propositions}
\Crefname{proposition}{Proposition}{Propositions}
\newaliascnt{lemma}{theorem}
\newtheorem{lemma}[lemma]{Lemma}
\aliascntresetthe{lemma}
\crefname{lemma}{lemma}{lemmas}
\Crefname{lemma}{Lemma}{Lemmas}
\newaliascnt{corollary}{theorem}
\newtheorem{corollary}[corollary]{Corollary}
\aliascntresetthe{corollary}
\crefname{corollary}{corollary}{corollaries}
\Crefname{corollary}{Corollary}{Corollaries}
\theoremstyle{definition}
\newaliascnt{definition}{theorem}
\newtheorem{definition}[definition]{Definition}
\aliascntresetthe{definition}
\crefname{definition}{definition}{definitions}
\Crefname{definition}{Definition}{Definitions}
\theoremstyle{remark}
\newaliascnt{remark}{theorem}
\newtheorem{remark}[remark]{Remark}
\aliascntresetthe{remark}
\crefname{remark}{remark}{remarks}
\Crefname{remark}{Remark}{Remarks}
\newcommand{\R}{\mathbb R}
\newcommand{\C}{\mathbb C}
\newcommand{\T}{\mathbb T}
\newcommand{\Id}{I}
\newcommand{\tr}{\operatorname{tr}}
\newcommand{\rank}{\operatorname{rank}}
\newcommand{\Ran}{\operatorname{Ran}}
\newcommand{\spec}{\operatorname{spec}}
\newcommand{\sym}{\operatorname{sym}}
\newcommand{\diag}{\operatorname{diag}}
\newcommand{\curl}{\operatorname{curl}}
\newcommand{\diver}{\operatorname{div}}
\newcommand{\op}{\mathrm{op}}
\newcommand{\dd}{\mathrm d}
\newcommand{\ip}[2]{\langle #1,#2\rangle}
\newcommand{\norm}[1]{\lVert #1\rVert}
\newcommand{\src}[1]{\nolinkurl{#1}}
\newcolumntype{Y}{>{\raggedright\arraybackslash}X}
\title{\textbf{Exact-Action to Einstein: Source-Resolved Bridge}\\[0.4em]
\large Research notes, Versions 1--55\\
Slow-time critical sheet, moving principal,\\
and five-row compatibility gate}
\author{}
\date{}
\begin{document}
\maketitle
\begin{abstract}
V54 isolates the rapid and slow forward-expanding normal directions of the
raw exact action, but it does not prove that the tangential dynamics preserves
the small-rate Einstein-compatible regime on physical times of order $a$.
We now rescale the exact V51 constrained field on that longer time.  The
resulting weak equation contains an order-$a^{-2}$ frozen constraint and an
order-$a^{-1}$ state-dependent principal correction; the latter cannot be
classified as a small forcing.  The leading constraint has rank $26$, kernel
dimension $297$, and a three-dimensional cokernel consisting exactly of the
constant-shift means.  Any exact history that retains the existing $O(a)$
physical-rate scaling on a fixed slow-time interval must converge to this
297-dimensional critical sheet and satisfy five pointwise compatibility
conditions: two nonconstant weak-zero tangency equations and three mean
principal equations.  The first tangency test for the prepared branch is an
explicit five-component coefficient vector built from finite Taylor jets of
the original same-root action.  Its value is not inferred here.  Thus the
next $t=O(a)$ gate has been reduced from a 323-dimensional nonlinear
propagation problem to a finite source calculation followed, if it passes,
by a regular slow descriptor.  All fifty-four preceding mathematical bodies
are preserved byte for byte.  Every action-specific statement remains at
fixed spatial cutoff $N=3$, with $a$ denoting field amplitude.
\end{abstract}
\setcounter{tocdepth}{1}
\tableofcontents
\clearpage

% BEGIN IMMUTABLE EXACT-ACTION BRIDGE V1 BODY
\section{Context, target, and inherited results}
\label{v1:context}

There are two different dynamics in the supplied material.  The first is a
specified local harmonic evolution, together with an open coefficient family
and open, physically scaled sets of chronological transition rows.  Its
small-data nonlinear lifespan and Einstein limit are established in the
open-basin source \cite{OB25}.  The second is the canonical flow of the
\emph{unchanged reconstructed finite action}.  Its multiplier equations,
constraint preservation, and physical connection response have to be
analysed on their own.  A convergence theorem for the first dynamics is not
a convergence theorem for the second.

This note starts a separate proof lineage for the latter question.  Its
ultimate target is either a direct exact-action Einstein limit, or a proved
comparison with the stable family in the physical topology required by the
metric, original action, and literal curvature reads.  It does not change the
action, add a constraint projection to its flow, reset multiplier values,
rescale physical time, or identify a small field residual with a small
curvature response.

\paragraph{Source audit.}
The following are retained inputs, not claimed as results first proved here.
The labels in the middle column are literal labels in the supplied sources.
\begin{center}\small
\begin{tabularx}{\textwidth}{@{}>{\raggedright\arraybackslash}p{.20\textwidth}>{\raggedright\arraybackslash}p{.35\textwidth}Y@{}}
\toprule
Input & Source location & Role and boundary\\\midrule
Stationary reduction & \cite{OB25}, \src{obv19:reduced-response},
\src{obv24:velocity} & The exact multiplier Hessian factors through the
six-component connection tangent; its metric is indefinite.\\[3pt]
Original first jets & \cite{OB25}, \src{obv19:leading-response},
\src{obv20:coframe-jets}, \src{obv23:three-term} & Flat Lorentz conventions,
stationary connection, electric cubic term, and phase-product defects.\\[3pt]
Prepared free family & \cite{OB25}, \src{obv21:common-family},
\src{obv21:first-records} & One fixed smooth symmetric-gradient seed is
sampled at every odd cutoff.\\[3pt]
Stable law class & \cite{OB25}, \src{obv16:stability},
\src{obv16:law-theorem}, \src{obv17:Einstein} & Uniform stable-writer
propagation and action-prepared inputs; no raw-flow selection theorem.\\[3pt]
Static response criterion & \cite{OB25}, \src{obv25:angle-response},
\src{obv25:cofinal-criterion}; \cite{GT76},
\src{thm:GT76-Cartan-cofinal} & Exact source-weighted $2/3$ resonance
criterion, not a verification on the actual cofinal family.\\
\bottomrule
\end{tabularx}
\end{center}

In particular, the $O(h^2)$ weak drift, the static principal-angle formula,
the general forced shadowing estimate, and the fixed-cutoff exact-action
continuation theorem are not re-proved as new advances.  The new part begins
with an explicit all-cutoff elimination inside the \emph{actual first source},
then keeps the remaining source intact during the trace reduction and time
differentiation.  Schur elimination, projection calculus and subspace angles
are standard linear algebra; the subspace-angle context is \cite{BG}.
Energy-compatible finite differences are likewise established methodology
\cite{WAK}.  No novelty claim is made for those general methods.

\paragraph{Logical strength.}
A source-specific estimate may be substantially weaker than a uniform bound
on an entire inverse.  Nevertheless, a criterion for such an estimate is not
its verification.  Every statement below identifies whether it concerns the
actual retained coefficient family, an exact finite-dimensional reduction,
or a comparison witness.  The full inherited nonlinear logarithmic calculus
is an input, not independently validated by the finite checks accompanying
this note.

\section{Original coefficient conventions and the scalar-momentum branch}
\label{v1:notation}

Let $N=2m+1\ge3$, $h=N^{-1}$ and $V=N^3$.  The spatial pairing is
$\ip uv_h=h^3\sum_x\overline{u(x)}v(x)$, with the Frobenius pairing for
symmetric matrices.  Thus off-diagonal entries are counted twice.  Use the
\emph{original action} phase derivatives and literal shifts
\begin{equation}
 \widehat{\delta_i u}(k)=i\kappa_i(k)\widehat u(k),\quad
 \kappa_i(k)=2h^{-1}\sin(\pi h k_i),\quad
 \widehat{S_i u}(k)=z_i(k)\widehat u(k),\quad
 z_i(k)=e^{2\pi i h k_i}.
 \label{v1:phase}
\end{equation}
They commute, $S_i$ is unitary, and $\delta_i$ is skew-adjoint.  No Leibniz
rule is imposed on $\delta_i$.  Put $\Omega_h^2=-\sum_i\delta_i^2$.

The canonical field is $X=(\gamma-I,\pi)$; $a$ denotes field amplitude,
not time.  At the unit-lapse, zero-shift cut, retain
\begin{equation}
 \gamma=I+aU+O_h(a^2),\qquad
 \pi=-a\frac{\omega_h}{\sqrt2}I+O_h(a^2),\qquad
 \omega_h=2h^{-1}\sin(\pi h).
 \label{v1:first-data}
\end{equation}
Write $B_j,R_j$ for boost and rotation generators, with
\begin{equation}
 [R_i,R_j]=-\epsilon_{ijk}R_k,\quad
 [R_i,B_j]=-\epsilon_{ijk}B_k,\quad
 [B_i,B_j]=\epsilon_{ijk}R_k.
 \label{v1:Lorentz}
\end{equation}
The retained first stationary spatial connection is
\begin{equation}
 A_i^{[1]}=r_{ia}R_a+\varkappa_h B_i,\qquad
 r_{ia}=-\frac12\epsilon_{abc}\delta_b U_{ic},\qquad
 \varkappa_h=\frac{\omega_h}{2\sqrt2}.
 \label{v1:connection}
\end{equation}
The factors here use the $a$ normalization of \eqref{v1:first-data}.
In particular $\tr r=0$, although $r$ need not be diagonal or symmetric.
For a multiplier test $(n,\beta)$ put
\begin{equation}
 p_i=\delta_i n,\qquad w_i=-\frac12\epsilon_{ijk}\delta_j\beta_k.
 \label{v1:multiplier-test}
\end{equation}
The letter $w$ in this formula denotes the temporal-rotation test, not a
canonical metric velocity.

Let $J(a)$ be the reduced multiplier source in the exact factorization
$M=J^*K^{-1}J$, and write $J(a)=aJ_{1,h}+O_h(a^2)$ at the prepared cut.
The flat tangent consists of symmetric spatial rotational matrices and has
\begin{equation}
 P S=\frac13(\tr S)I,\qquad
 K_0=2I-6P,\qquad K_0^{-1}=\frac12I-\frac34P.
 \label{v1:K}
\end{equation}
The original flat electric inclusion $Z_0$ is isometric in this normalization.
All formulas for a physical norm below use this positive norm, not the signed
quadratic form of $K_0$.

\begin{lemma}[Two first-source identities on the scalar-momentum branch]
\label{v1:first-source}
For every symmetric $U$ in \eqref{v1:first-data}, the pure-shift part is
\begin{equation}
 (J_{1,h}(0,\beta))_{ii}=0,\qquad
 (J_{1,h}(0,\beta))_{ij}
   =-\frac{h\varkappa_h}{2}(S_iw_j+S_jw_i),\quad i\ne j.
 \label{v1:shift-source}
\end{equation}
The trace of the complete lapse source is
\begin{equation}
 \tr J_{1,h}(n,0)
   =-h\sum_i\sum_{a\ne i}r_{ia}S_i\delta_a n.
 \label{v1:trace-source}
\end{equation}
Both are identities of the first coefficient of the original reduced saddle.
Second preparation coefficients and higher logarithmic terms are not
replaced; they cannot enter this coefficient by degree counting.
\end{lemma}
\begin{proof}
We retain the original coefficient extraction rather than infer the result
from the continuum gauge algebra.  The flat coframe coefficients and their
first variations are
\[
 \Pi_i^0=2B_i,\quad \Sigma_{ij}^0=2\epsilon_{ijb}R_b,\quad
 \Pi_i^{[1]}=((\tr U)\delta_{ia}-U_{ia})B_a,\quad
 \Sigma_{ij}^{[1]}=\epsilon_{ijb}U_{ba}R_a.
\]
As in the differentiated saddle formula in \cite{OB25},
$J_1=Z_0^*(G_1-T_1^*R_0^*G_0)$: the auxiliary response is taken before
restricting to the symmetric rotational tangent.  For a pure shift, the
geometric spatial load and the velocity-load correction are boost-valued.
The magnetic Cartan contraction with the scalar boost matrix in
\eqref{v1:connection} is antisymmetric in the two rotational indices.
Their symmetric rotational projections are zero.  For a lapse, the
geometric part that survives this projection is
\[
 -\frac12\sum_{j,b}\bigl\{
 \epsilon_{jib}\mathfrak d_j(n,U_{\ell b})+
 \epsilon_{j\ell b}\mathfrak d_j(n,U_{ib})\bigr\},\qquad
 \mathfrak d_j(f,g)=\delta_j(fg)-(\delta_j f)g-f\delta_jg.
\]
This follows by substituting the displayed coframe coefficients in the
spatial load and auxiliary term.  Terms proportional to $\tr U$ cancel
under symmetric alternating contraction.  Taking $i=\ell$ and summing gives
zero because $U_{ib}=U_{bi}$, while $\epsilon_{jib}$ is antisymmetric.
These cancellations use the full product defect, not a continuum product
rule.

It remains to evaluate the electric cubic term
\[
 -\frac h2\sum_i\ip{2B_i}{[A_i,[A_i,S_iA_0]]}.
\]
At order $a$, its contribution to the differentiated auxiliary equation,
tested against a symmetric rotation $S$, is obtained by substituting
$A_i^{[1]}$ and $A_0^{[0]}=p_aB_a+w_aR_a$.  Let $s_i$ be row $i$ of $S$,
$r_i$ row $i$ of $r$, and $e_i$ the coordinate vector.  The boost coefficient
inside the double commutator is
\[
 s_i\mathbin\times(r_i\mathbin\times S_ip)
 +r_i\mathbin\times(s_i\mathbin\times S_ip)
 +\varkappa_h\{s_i\mathbin\times(e_i\mathbin\times S_iw)
             +e_i\mathbin\times(s_i\mathbin\times S_iw)\}.
\]
Its dot product with $e_i$ has coefficient, in $s_{ia}$,
\begin{equation}
 r_{ii}S_ip_a+\delta_{ia}\sum_b r_{ib}S_ip_b
 -2r_{ia}S_ip_i
 +\varkappa_h(S_iw_a-\delta_{ia}S_iw_i).
 \label{v1:electric-row-coefficient}
\end{equation}
Multiply by $-h$ and symmetrize in the Frobenius pairing.  At $p=0$ this
is exactly \eqref{v1:shift-source}.  On the diagonal its lapse part is
$-h\sum_{a\ne i}r_{ia}S_ip_a$, proving \eqref{v1:trace-source} after
adding the zero geometric trace.  The magnetic logarithmic term begins at
spatial connection degree three and contributes to $J_2$, not $J_1$.
Higher electric terms and the second prepared field coefficient also enter
only the higher amplitude orders.  Thus the stated identities retain every
term able to contribute.
\end{proof}

\section{An all-cutoff transverse-shift block}
\label{v1:transverse}

Define the real transverse shift space
\begin{equation}
 \mathcal T_h=\{\beta:\ \widehat\beta(0)=0,
                    \ \kappa(k)^T\widehat\beta(k)=0\text{ for }k\ne0\}.
 \label{v1:transverse-space}
\end{equation}
It is the positive orthogonal complement of the curl-free shifts and has
dimension $2(V-1)$.  Let $A_h=J_{1,h}(0,\cdot)|_{\mathcal T_h}$.

\begin{theorem}[Exact rank and uniform natural-order control]
\label{v1:transverse-rank}
On the actual scalar-momentum family, at every odd cutoff and independently
of $U$,
\begin{equation}
 \rank A_h=2(V-1),\qquad PA_h=0,
 \label{v1:rank}
\end{equation}
and, for every $\beta\in\mathcal T_h$,
\begin{equation}
 \frac{h^2\varkappa_h^2}{8}\norm{\Omega_h\beta}_h^2
 \le\norm{A_h\beta}_h^2
 \le\frac{h^2\varkappa_h^2}{2}\norm{\Omega_h\beta}_h^2.
 \label{v1:transverse-bounds}
\end{equation}
The same inequalities hold after multiplication by any scalar Fourier
Sobolev weight.  In particular
\begin{equation}
 A_h^*K_0^{-1}A_h=\frac12 A_h^*A_h\succ0.
 \label{v1:positive-shift}
\end{equation}
This is positivity of this specified block, not of the complete active
source or the full multiplier Hessian.
\end{theorem}
\begin{proof}
For arbitrary vector $w$ let $\mathcal L_hw$ be the symmetric matrix with
zero diagonal and off-diagonal entries $S_iw_j+S_jw_i$.
At a frequency $k$, put $t_i=z_i^{-1}\widehat w_i$ and multiply the
$(i,j)$ output by $z_i^{-1}z_j^{-1}$.  The off-diagonal output becomes
\[
 \begin{pmatrix}t_1+t_2\\t_1+t_3\\t_2+t_3\end{pmatrix}
 =\mathsf M t,\qquad
 \mathsf M=\begin{pmatrix}1&1&0\\1&0&1\\0&1&1\end{pmatrix}.
\]
All input and output phase multipliers are unitary.  The eigenvalues of
$\mathsf M^*\mathsf M$ are $4,1,1$.  Counting each off-diagonal matrix entry
twice therefore gives
\[
 2\norm w_h^2\le\norm{\mathcal L_hw}_h^2\le8\norm w_h^2.
\]
For $w=-\tfrac12\curl_\delta\beta$ and transverse $\beta$, Parseval gives
$\norm w_h^2=\tfrac14\norm{\Omega_h\beta}_h^2$.  Substitute
$A_h\beta=-h\varkappa_h\mathcal L_hw/2$ to obtain
\eqref{v1:transverse-bounds}.  The nonzero phase vector at every nonzero
odd-grid mode proves injectivity on $\mathcal T_h$.  Reality is preserved
by the conjugate frequency pairing, so the real rank is $2(V-1)$.
The diagonal is zero, proving $PA_h=0$ and then
\eqref{v1:positive-shift}.  A scalar Fourier weight commutes with every
matrix in the argument, proving the Sobolev assertion.
\end{proof}

\begin{corollary}[A phase-resolved projector without a small singular-value test]
\label{v1:shift-projector}
Let $\Pi_{A,h}$ be the positive orthogonal projector onto $\Ran A_h$.
At a nonzero mode use the off-diagonal coordinates $(12,13,23)$ and put
\[
 \mathsf L(k)=
 \begin{pmatrix}z_2&z_1&0\\z_3&0&z_1\\0&z_3&z_2\end{pmatrix},
 \qquad \nu(k)=\mathsf L(k)^{-*}\kappa(k).
\]
Then $\Pi_{A,h}$ is zero on diagonal coordinates, is zero at $k=0$, and
acts on the off-diagonal coordinates at $k\ne0$ as
\begin{equation}
 I_3-\frac{\nu(k)\nu(k)^*}{\norm{\nu(k)}^2},\qquad
 \frac12|\kappa(k)|\le\norm{\nu(k)}\le|\kappa(k)|.
 \label{v1:projector-symbol}
\end{equation}
Thus $\Pi_{A,h}$ is uniformly bounded on every Fourier Sobolev space and
commutes with the pointwise trace projection $P$; in fact both compositions
$P\Pi_{A,h}$ and $\Pi_{A,h}P$ vanish.
\end{corollary}
\begin{proof}
The singular values of $\mathsf L(k)$ are $2,1,1$ by the preceding unitary
factorization.  Its image of $\kappa(k)^\perp$ is exactly
$\nu(k)^\perp$.  Curl maps transverse shifts bijectively onto transverse
vectors, so this image is precisely the off-diagonal source range.  The
Frobenius pairing is twice the Euclidean pairing on these three coordinates;
its orthogonal projector is unchanged.  The remaining assertions follow
from the explicit formula and from zero diagonal output.
\end{proof}

The inverse of the map $w\mapsto\mathcal L_hw$ is also explicit.  For example,
\[
 w_1=\frac12\{S_2^{-1}Y_{12}+S_3^{-1}Y_{13}
                   -S_1S_2^{-1}S_3^{-1}Y_{23}\},
\]
with the other two components obtained by permutation.  On its transverse
range, $\beta=-2\Omega_h^{-2}\curl_\delta w$.  No claim is made that the
inverse elliptic factor is a local execution rule; it is an analytical
source-elimination device.

\section{A trace leak in the retained smooth free family}
\label{v1:trace-leak}

Use exactly the smooth field $\xi^\sharp$ and free parameter $\rho$ of the
source.  Its rational three-site table is reproduced in
\cref{v1:table} as input data.  Set
\begin{equation}
 U_{h,\rho}=U_0+\rho(\delta_i\xi_{h,j}+\delta_j\xi_{h,i}),
 \quad (U_0)_{23}=\cos2\pi x_1,
 \quad (U_0)_{13}=\cos2\pi x_2,
 \quad (U_0)_{12}=\cos2\pi x_3.
 \label{v1:common-seed}
\end{equation}
The diagonal of $U_0$ is zero and its off-diagonal entries are symmetric.
The scalar momentum is still \eqref{v1:first-data}.

\begin{theorem}[Nonzero active trace with a sharp smooth-cutoff order]
\label{v1:trace-leak-theorem}
For the lapse test $n(x)=\cos2\pi x_1$, $\beta=0$, the original first source
on \eqref{v1:common-seed} satisfies the exact identity
\begin{equation}
 \tr J_{1,h}(n,0)
 =\frac{\rho h}{2}(\delta_1n)
       (\delta_2+\delta_3)(\delta_2\xi_{h,3}-\delta_3\xi_{h,2}).
 \label{v1:trace-seed}
\end{equation}
There is a fixed nonzero trigonometric polynomial $\Psi$ such that, on every
sufficiently fine odd grid,
\begin{equation}
 \tr J_{1,h}(n,0)=\rho h\omega_h^3\,\mathcal S_h\Psi,
 \qquad
 \Psi=\frac{(\partial_1n)(\partial_2+\partial_3)
                (\partial_2\xi^\sharp_3-\partial_3\xi^\sharp_2)}
                 {2(2\pi)^3}.
 \label{v1:Psi}
\end{equation}
Consequently, for every fixed Sobolev index $r$, there are $c_r,C_r>0$
independent of $h,\rho$ with
\begin{equation}
 c_r|\rho|h\le\norm{\tr J_{1,h}(n,0)}_{H_h^r}
                   \le C_r|\rho|h.
 \label{v1:trace-rate}
\end{equation}
At $N=3$, $\rho=1$ and $x=(1/3,0,0)$, its exact value is $-27/320$.
This algebraic value is not asserted to be a finite-amplitude trajectory at
an admitted nonlinear radius.
\end{theorem}
\begin{proof}
For $U_0$ the rotational matrix $r$ is diagonal; its contribution to
\eqref{v1:trace-source} is therefore zero.  For the symmetric-gradient part,
commuting the original phase derivatives gives
\[
 r_{ia}^{\rm grad}=-\frac{\rho}{2}
                \delta_i(\curl_\delta\xi_h)_a.
\]
Only $p_1=\delta_1n$ is nonzero.  In \eqref{v1:trace-source} only $i=2,3$
remain, and $S_i p_1=p_1$ for those indices.  This proves
\eqref{v1:trace-seed} without a product-rule approximation.

Every frequency of $\xi^\sharp$ is in $\{-1,0,1\}^3$.  On this bank
$\delta_i=(\omega_h/(2\pi))\partial_i$, also after another phase derivative.
The same identity applies to $n$.  Hence \eqref{v1:Psi} is exact at the
sampled points, with the usual folded interpretation on the three-site grid.
There is no truncation of a high-frequency test in this assertion.

Nonvanishing can be checked in characteristic zero using the input table.
On $N=3$, $\delta_i=3(S_i-S_i^{-1})$.  At $x=(1/3,0,0)$ the table gives
\[
 \delta_1n=-\frac92,\qquad
 (\delta_2+\delta_3)(\delta_2\xi_3-\delta_3\xi_2)=\frac9{80}.
\]
Equation \eqref{v1:trace-seed} gives $-27/320$.  Since
$h\omega_h^3=27\sqrt3$ at that cutoff,
$\Psi(1/3,0,0)=-1/(320\sqrt3)\ne0$.  Thus the fixed trigonometric polynomial
is not zero.  On sufficiently fine grids its Sobolev norms are uniformly
comparable to its positive continuum norm, and $\omega_h\to2\pi$.
This proves both bounds in \eqref{v1:trace-rate}.
\end{proof}

\begin{corollary}[Smoothness does not make the normalized remaining source trace-free]
\label{v1:normalized-leak}
Let $G_{L,h}n=(I-\Pi_{A,h})J_{1,h}(n,0)$.  For the preceding test and a
fixed bounded interval of $\rho$,
\[
 \norm{G_{L,h}n}_{H_h^r}\le C_rh,\qquad
 \norm{P G_{L,h}n}_{H_h^r}\ge c_r|\rho|h.
\]
For each fixed nonzero admitted $\rho$, the ratio of these two norms is
therefore bounded below by a positive multiple of $|\rho|$.
\end{corollary}
\begin{proof}
The electric first-source coefficient on this fixed smooth bank is $O(h)$.
The geometric coefficient contains only phase-product defects, each $O(h^2)$
on the fixed bank by the Taylor expansion of \eqref{v1:phase}.  The projector
is bounded on $H_h^r$ by \cref{v1:shift-projector}.  It does not change the
trace, and $\norm{PS}=\norm{\tr S}/\sqrt3$.  Apply
\cref{v1:trace-leak-theorem}.
\end{proof}

This does not show that a $2/3$ resonance occurs on the actual cofinal
family.  It does rule out a shortcut: the trace-free property of the entire
unperturbed active source cannot be carried over to the fixed nonzero-$\rho$
family merely because its free field is smooth.  Only the transverse-shift
block has the unconditional trace-free property used in the next reduction.

\section{Exact elimination that keeps all remaining physical sources}
\label{v1:elimination}

On the mean-lapse-zero multiplier slice, use the orthogonal decomposition
into transverse shifts, mean-zero lapses, and curl-free shifts.  Their
respective dimensions are $2V-2$, $V-1$, and $V+2$.  The source's complete
leading graded matrix and target can then be written
\begin{equation}
 F_h=(A_h,G_h),\qquad
 G_h=(J_{1,h}|_{\text{mean-zero lapse}},\ J_{2,h}|_{\text{curl-free shift}}),
 \qquad b_h=(0,d_h).
 \label{v1:complete-F}
\end{equation}
The first $V-1$ entries of $d_h$ are also zero; its remaining entries are the
actual weak quadratic drift in the declared multiplier pairing.  This is
just a basis choice and the inherited amplitude grading.  No active equation
is discarded and $J_{2,h}$ is the full, prepared, original-action coefficient.

\begin{theorem}[Safe all-cutoff elimination of $2V-2$ source columns]
\label{v1:safe-elimination}
Put $\overline G_h=(I-\Pi_{A,h})G_h$.  Then
\begin{equation}
 \rank F_h=(2V-2)+\rank\overline G_h.
 \label{v1:rank-split}
\end{equation}
Moreover $F_h^*K_0^{-1}F_h$ is congruent to
\begin{equation}
 \diag\left(\tfrac12 A_h^*A_h,
                  \overline G_h^*K_0^{-1}\overline G_h\right).
 \label{v1:Gram-split}
\end{equation}
Thus its invertibility is exactly the invertibility of the remaining
$(2V+1)$-dimensional signed Gram.  On an invertible branch the complete
physical response is
\begin{equation}
 v_h=K_0^{-1}\overline G_h
       (\overline G_h^*K_0^{-1}\overline G_h)^{-1}d_h.
 \label{v1:reduced-response}
\end{equation}
The positive minimum source-reproduction vector is
\begin{equation}
 w_h=\overline G_h(\overline G_h^*\overline G_h)^{-1}d_h.
 \label{v1:minimum-reduced}
\end{equation}
Both vectors satisfy the original transverse equations $A_h^*v_h=A_h^*w_h=0$
and all remaining equations $G_h^*v_h=G_h^*w_h=d_h$.
\end{theorem}
\begin{proof}
Let $C_h=(A_h^*A_h)^{-1}A_h^*G_h$.  Right multiplication of $F_h$ by
$\left(\begin{smallmatrix}I&-C_h\\0&I\end{smallmatrix}\right)$ gives
$(A_h,\overline G_h)$ with orthogonal column ranges.  This proves
\eqref{v1:rank-split}.  Since $K_0^{-1}A_h=A_h/2$, orthogonality also gives
$A_h^*K_0^{-1}\overline G_h=0$.  The same invertible change of columns
therefore gives \eqref{v1:Gram-split}.  Its adjoint leaves the target
$(0,d_h)$ unchanged.  The first coefficient of the transformed signed solve
is zero, yielding \eqref{v1:reduced-response}.  The positive Gram gives
\eqref{v1:minimum-reduced} in the same way.  Finally $K_0^{-1}$ preserves
$(\Ran A_h)^\perp$, and $G_h=\overline G_h+A_hC_h$; these facts verify every
original source equation.
\end{proof}

The elimination is valid on every odd grid even before the remaining Gram
is known to be invertible.  It closes the rank and inverse problem for the
transverse block only.  For example, at $N=3$ it reduces a $107$-column graded
problem to $55$ remaining columns; at general $N$ the count is $2N^3+1$.
This finite-grid number $55$ has no relation to the independent
$55$-parameter stable-writer family.

\section{A scalar-trace solve for the complete remaining response}
\label{v1:scalar}

Work in the Hilbert space $(\Ran A_h)^\perp$, or in the full tangent space
without the preceding elimination.  Write $F$ for the remaining full-rank
source and $b$ for its complete target.  The discussion also applies to an
arbitrary full-rank finite source; no trace-free assumption on its columns
is imposed.  Let
\begin{equation}
 \Pi=F(F^*F)^{-1}F^*,\qquad
 w=F(F^*F)^{-1}b,\qquad
 E\tau=\frac{\tau}{\sqrt3}I,
 \label{v1:Pi-w}
\end{equation}
where $E$ is the isometric scalar-trace inclusion.  In the reduced Hilbert
space this inclusion is still available because $\Ran E\perp\Ran A_h$.
Define
\begin{equation}
 C=E^*\Pi E,\qquad H=2I-3C,\qquad \zeta=E^*w.
 \label{v1:trace-data}
\end{equation}
The scalar matrices $C,H$ have dimension $V$.

\begin{theorem}[Exact scalar-trace reconstruction]
\label{v1:scalar-reconstruction}
Assume $F$ has full column rank.  The signed Gram $F^*K_0^{-1}F$ is invertible
if and only if $H$ is invertible.  In that case set $\chi=H^{-1}\zeta$.
The \emph{unchanged} physical response obeys
\begin{align}
 v&=w-3(I-\Pi)E\chi,& E^*v&=-\chi,\label{v1:scalar-v}\\
 \norm v^2&=\norm w^2+9\ip\chi{(I-C)\chi}.
 \label{v1:scalar-norm}
\end{align}
In particular
\begin{equation}
 \norm w\le\norm v,\quad \norm\chi\le\norm v,\quad
 \norm v\le\norm w+3\norm\chi.
 \label{v1:scalar-two-sided}
\end{equation}
For any cofinal regular family, $v_h\to0$ is therefore equivalent to the two
\emph{separately measured} statements $w_h\to0$ and
$H_h^{-1}E^*w_h\to0$.
\end{theorem}
\begin{proof}
Write $F=QR$, where $Q^*Q=I$ and $R$ is invertible.  The signed Gram is
congruent to $(2I-3Q^*EE^*Q)/4$.  The nonzero eigenvalues of
$Q^*EE^*Q$ and $E^*QQ^*E=C$ agree: they are those of $L^*L$ and $LL^*$ for
$L=E^*Q$.  Since the dangerous eigenvalue $2/3$ is nonzero, the stated
invertibility equivalence follows.  This is also a direct finite Schur
complement argument, without computing a principal-angle basis.

The response is characterized by $F^*v=b$ and $K_0v\in\Ran F$.
The first equation gives $\Pi v=w$.  The second, with
$K_0=2I-6EE^*$, gives
\[
 (I-\Pi)v=3(I-\Pi)EE^*v.
\]
Let $\tau=E^*v$.  Then
$\tau=E^*w+3(I-C)\tau$, hence $H\tau=-\zeta$ and
$\tau=-\chi$.  Substitution proves \eqref{v1:scalar-v} and verifies both
original equations.  Orthogonality of $w$ and $(I-\Pi)E\chi$ proves
\eqref{v1:scalar-norm}.  The first two lower bounds follow respectively
from $w=\Pi v$ and $\chi=-E^*v$; the upper bound follows from the contraction
norms of the projections.  The cofinal equivalence is immediate.
\end{proof}

This is not a replacement of the source's static criterion by a different
physical rule.  It is an exact alternative way to compute the same vector,
using one positive source solve and a scalar-trace solve of dimension $V$.
The positive solve is still required: the inherited $O(h^2)$ unpreconditioned
drift alone does not bound $w$.  Nor does the small Fourier support of that
drift make $C$ a Fourier multiplier.

\begin{proposition}[Positive-space inverse on the source complement]
\label{v1:complement-inverse}
Let $\mathcal N=\ker F^*$ and $D=(I-\Pi)E$.  On a regular branch, the
restriction
\[
 L_0=(I-\Pi)K_0|_{\mathcal N}=2I_{\mathcal N}-6DD^*
\]
is invertible and
\begin{equation}
 L_0^{-1}=\frac12 I_{\mathcal N}-\frac32 D H^{-1}D^*.
 \label{v1:complement-inverse-formula}
\end{equation}
\end{proposition}
\begin{proof}
One has $D^*D=I-C$ and $I-3D^*D=-H$.  Multiply the proposed inverse by
$2(I-3DD^*)$.  The cross terms cancel using this identity, leaving
$I_{\mathcal N}$.  The multiplication in the reverse order is identical.
\end{proof}

\section{Nonlinear tangent Hessian: a physical correction estimate}
\label{v1:nonlinear}

At finite amplitude the original tangent Hessian need not equal $K_0$.
The following comparison states exactly what must be paid for that change;
it does not modify the actual Hessian.  The leading transverse block need
not remain invariant under that nonlinear Hessian.  Without a separate
commutation theorem, the finite-amplitude application must therefore use the
complete source $F_a$, not silently reuse the leading elimination of
$A_h$.  All tangent pairings in the proposition are positive physical
pairings.

\begin{proposition}[Source-resolved correction with the true Hessian]
\label{v1:K-perturbation}
Keep a full-rank $F$ and target $b$, and suppose the $K_0$ branch is regular.
Let $K=K_0+\Delta K$ be a self-adjoint invertible Hessian and set
\[
 q=\norm{L_0^{-1}(I-\Pi)\Delta K|_{\mathcal N}}_{\op}.
\]
If $q<1$, the true signed Gram $F^*K^{-1}F$ is invertible.  If $v_0$ denotes
the $K_0$ response and $v_K$ the true response, then
\begin{equation}
 v_K-v_0\in\mathcal N,\qquad
 \norm{v_K-v_0}
 \le\frac{\norm{L_0^{-1}(I-\Pi)\Delta K v_0}}{1-q}.
 \label{v1:K-correction}
\end{equation}
The numerator is the actual deformation acting on the actual comparison
response, not a worst-case source replacement.
\end{proposition}
\begin{proof}
For vectors $v$ satisfying $F^*v=b$, write $v=v_0+e$ with $e\in\mathcal N$.
The stationarity equation $(I-\Pi)Kv=0$ is
\[
 [I_{\mathcal N}+L_0^{-1}(I-\Pi)\Delta K|_{\mathcal N}]e
       =-L_0^{-1}(I-\Pi)\Delta K v_0.
\]
Its bracket has a Neumann inverse of norm at most $(1-q)^{-1}$.  This proves
existence, uniqueness, and \eqref{v1:K-correction} for every $b$.
If $F^*K^{-1}Fu=0$, then $e=K^{-1}Fu$ satisfies the homogeneous constrained
stationarity system, so $e=0$ and then $Fu=0$.  Full column rank gives $u=0$.
The square Gram is thus invertible, and its response is the unique one just
constructed.
\end{proof}

\begin{remark}[Changing $F$ and changing the physical inclusion]
The proposition can be used with the \emph{actual} finite-amplitude $F_a$
and $b_a$; their $K_0$ comparison must then be recomputed, not replaced
without a bound by the leading $F_0,b_0$.  Also the nonlinear inclusion $Z$
need not be isometric.  If $H_Z=Z^*Z$, use the physical coordinates
$\widetilde v=H_Z^{1/2}v$, $\widetilde K=H_Z^{-1/2}KH_Z^{-1/2}$ and
$\widetilde F=H_Z^{-1/2}F$.  Then
$\norm{Zv}=\norm{\widetilde v}$.  Apply \cref{v1:K-perturbation} to
$\widetilde K$ relative to the fixed flat reference $K_0$, charging the
\emph{entire} difference $\widetilde K-K_0$.  The congruence does not
preserve the special trace form of $K_0$ automatically.  Bounds for $H_Z$,
its derivatives, and these transformed coefficients remain necessary.
An $O_h(a)$ Taylor remainder with an unspecified $h$-dependent constant is
not a cofinal estimate for \eqref{v1:K-correction}.
\end{remark}

\section{Differentiated trace response: the extra temporal source}
\label{v1:time}

Pointwise smallness of the electric coefficient is not a time-regularity
estimate.  This section gives exact derivative formulas without choosing
smooth eigenvectors at crossings.  First consider a $C^r$ family $F(t),b(t)$
at a fixed finite cutoff, with full column rank and regular signed Gram,
in a fixed flat physical tangent space.  The trace inclusion $E$ and $K_0$
are fixed.  All projections and responses below are evaluated at $t$.

\begin{theorem}[Derivative hierarchy in scalar-trace coordinates]
\label{v1:time-hierarchy}
Let $\Pi,w,C,H,\zeta,\chi,v$ be as in
\cref{v1:scalar-reconstruction}.  Define $\chi_0=\chi$ and recursively
\begin{equation}
 \chi_j=H^{-1}\left[\zeta^{(j)}+
             3\sum_{\ell=1}^j\binom j\ell C^{(\ell)}\chi_{j-\ell}\right],
 \qquad 1\le j\le r.
 \label{v1:time-sources}
\end{equation}
Then $\chi_j=\chi^{(j)}=-E^*v^{(j)}$.  In particular
\begin{align}
 \chi'&=H^{-1}(\zeta'+3C'\chi),\label{v1:first-time-source}\\
 v'&=w'+3\Pi'E\chi-3(I-\Pi)E\chi'.
 \label{v1:v-derivative}
\end{align}
If $\max_{j\le r}\norm{\Pi^{(j)}}_{L^\infty_t\!,\op}\le L$ uniformly
in a cofinal family, then for every $1\le p\le\infty$ there is a constant
$C_{r,L}$ independent of the finite dimension such that
\begin{equation}
 C_{r,L}^{-1}\norm v_{W^{r,p}_t}
 \le \norm w_{W^{r,p}_t}+\norm\chi_{W^{r,p}_t}
 \le C_{r,L}\norm v_{W^{r,p}_t}.
 \label{v1:time-equivalence}
\end{equation}
Consequently a dynamical use of the trace reduction must control the
resolvent on the transported sources in \eqref{v1:time-sources}, not just on
$\zeta$.
\end{theorem}
\begin{proof}
Differentiate $H\chi=\zeta$ $j$ times.  Since
$H^{(\ell)}=-3C^{(\ell)}$ for $\ell\ge1$, moving the lower-order terms to
the right proves \eqref{v1:time-sources} by induction.  The identity
$E^*v=-\chi$ holds at each time and $E$ is fixed, proving the corresponding
identity for every derivative.  Differentiating \eqref{v1:scalar-v} gives
\eqref{v1:v-derivative}.  Repeated Leibniz differentiation of
$v=w-3(I-\Pi)E\chi$ bounds the left side of \eqref{v1:time-equivalence}
by its middle expression.  Conversely $w=\Pi v$ and $\chi=-E^*v$,
with the same derivative bounds on $\Pi$, give the other inequality.
Only finite binomial sums and contraction norms of $E$ are used, so no
constant depends on the number of sites.
\end{proof}

\begin{proposition}[The derivative of the positive reproduction vector]
\label{v1:w-derivative}
Let $F^\dagger=(F^*F)^{-1}F^*$ and
$\mathcal A=(I-\Pi)F'F^\dagger$.  Then
\begin{align}
 \Pi'&=\mathcal A+\mathcal A^*,\label{v1:Pi-derivative}\\
 w'&=\mathcal A w+F(F^*F)^{-1}(b'-F'^*w),\label{v1:w-prime}\\
 \norm{w'}^2&=\norm{\mathcal A w}^2+
       \norm{(F^*F)^{-1/2}(b'-F'^*w)}^2.
 \label{v1:w-prime-norm}
\end{align}
Moreover $\norm{\Pi'}_{\op}=\norm{\mathcal A}_{\op}$.
\end{proposition}
\begin{proof}
Differentiating $\Pi F=F$ gives $(I-\Pi)\Pi'F=(I-\Pi)F'$.
Differentiating $\Pi^2=\Pi$ shows that both diagonal blocks of $\Pi'$
relative to $\Ran\Pi\oplus\ker\Pi$ vanish.  Self-adjointness therefore gives
\eqref{v1:Pi-derivative} and the norm identity for its off-diagonal blocks.
Since $w=\Pi w$, its perpendicular derivative is $\mathcal A w$.
Differentiating $F^*w=b$ determines its parallel derivative as the second
term in \eqref{v1:w-prime}.  These two terms are orthogonal.  Polar
normalization of the parallel term proves \eqref{v1:w-prime-norm}.
\end{proof}

\begin{proposition}[Static suppression does not imply differentiated suppression]
\label{v1:time-witness}
There are smooth finite response families, with the same two eigenvalues
$-4,2$ as the original trace and trace-free Cartan sectors, for which on any
fixed interval $[0,T]$
\[
 \sup_t\norm{v_h(t)}=O(h),\qquad
 \sup_t\norm{\Pi_h'(t)}=O(1),\qquad
 \norm{w_h}_{W^{1,\infty}_t}=O(h^3),
\]
but $\norm{v_h'(0)}\ge c/h$.  Every signed Gram is invertible.
This is a matrix comparison witness, not a new original-action history.
\end{proposition}
\begin{proof}
Take $K_0=\diag(-4,2)$, $E=(1,0)^T$, and
\[
 c_h(t)=2+\sin(t/h^2),\quad
 \mu_h(t)=\frac23+h^2c_h(t),\quad
 F_h(t)=\binom{\sqrt{\mu_h(t)}}{\sqrt{1-\mu_h(t)}},\quad b_h=h^3.
\]
For sufficiently small $h$, $0<\mu_h<1$ and
$H_h=2-3\mu_h=-3h^2c_h$ never vanishes.  Since $F_h$ has unit norm,
$w_h=h^3F_h$, $\zeta_h=h^3\sqrt{\mu_h}$, and
\[
 \chi_h=-\frac{h\sqrt{\mu_h}}{3c_h}.
\]
Here $\mu_h'=\cos(t/h^2)$ is bounded, so $F_h'$, $\Pi_h'$ are bounded
and $w_h$ has the asserted bound.  The two-sided estimates of
\cref{v1:scalar-reconstruction} give $v_h=O(h)$.  Direct differentiation at
zero gives
\[
 \chi_h'(0)=-\frac{h}{12\sqrt{2/3+2h^2}}
             +\frac{\sqrt{2/3+2h^2}}{12h},\qquad
 \lim_{h\to0}h\chi_h'(0)=\frac{\sqrt{2/3}}{12}>0.
\]
Since $\chi_h'=-E^*v_h'$, the claimed lower bound follows.  The extra source
$3C_h'\chi_h$ in \eqref{v1:first-time-source} is therefore essential even
when the source subspace itself has bounded velocity.
\end{proof}

\subsection{The differentiated identity with the actual Hessian}
The flat reduction is useful for leading coefficients, but the unchanged
nonlinear saddle uses its true $K(t)$.  The following formula keeps that
Hessian throughout.

\begin{proposition}[Derivative of the exact constrained response]
\label{v1:true-derivative}
Let $K(t)$ be a differentiable, self-adjoint invertible Hessian, let $F(t)$
have full column rank, and suppose $\Gamma=F^*K^{-1}F$ is invertible.
Put $u=\Gamma^{-1}b$ and $v=K^{-1}Fu$.  Define
\[
 q=F'u-K'v,\qquad r=b'-F'^*v,\qquad
 \sigma=r-F^*K^{-1}q.
\]
Then the exact identity is
\begin{equation}
 v'=K^{-1}q+K^{-1}F\Gamma^{-1}\sigma.
 \label{v1:true-derivative-formula}
\end{equation}
If the measured response is $Z(t)v(t)$, its derivative is
$Z'v+Zv'$ and the $Z'$ term must be retained.
\end{proposition}
\begin{proof}
The two original equations are $Kv=Fu$ and $F^*v=b$.
Differentiating gives $Kv'=Fu'+q$ and $F^*v'=r$.
Substituting the first in the second yields
$\Gamma u'=r-F^*K^{-1}q=\sigma$, proving the formula.
The final assertion is the ordinary product rule in the fixed physical
comparison coordinates.  No derivative of an eigenbasis, pseudoinverse
threshold, or discarded null component occurs.
\end{proof}

\section{What the response estimates do and do not give the flow}
\label{v1:flow-interface}

The source's exact normalized canonical continuation has
\[
 X_t=\mathbb J\nabla_X\mathscr H_h,\qquad
 \lambda_t=-M_h^\dagger B_h+
                  \lambda\,\ell(M_h^\dagger B_h),\qquad \ell(\lambda)=1.
\]
At an exactly prepared cut its graded representation is
\begin{equation}
 \lambda_t=-U S_a^{-1}u_a,\qquad
 u_a=\Gamma_a^{-1}b_a,\qquad
 S_a=\diag(aI_{3V-3},a^2I_{V+2}).
 \label{v1:rate-grade}
\end{equation}
The electric response contains $K_a^{-1}F_a u_a$, not the unweighted
coefficient $u_a$.  These are different observations.

\begin{proposition}[An additional necessary initial test for harmonic shadowing]
\label{v1:clock-test}
Choose orthonormal multiplier coordinates $U$ on the mean-lapse-zero slice
at the unit-lapse, zero-shift cut.  Then
\begin{equation}
 \norm{\lambda_t}_h^2
 =a^{-2}\norm{(u_a)_A}_h^2+a^{-4}\norm{(u_a)_W}_h^2.
 \label{v1:rate-norm}
\end{equation}
Suppose $a(h)\to0$ and an exact-action metric is to shadow a small stable
harmonic metric in a topology controlling its initial first time derivative,
with the same lapse, shift, and spatial metric at the cut.  A necessary
condition is
\begin{equation}
 a(h)^{-1}\norm{(u_{a(h)})_A}_h\to0,\qquad
 a(h)^{-2}\norm{(u_{a(h)})_W}_h\to0.
 \label{v1:rate-necessary}
\end{equation}
This condition is additional to decay of the physical electric coefficient.
It uses the full finite-amplitude $u_a$, not just its first Taylor coefficient.
\end{proposition}
\begin{proof}
Orthogonality of $U$ and the two blocks of $S_a$ gives
\eqref{v1:rate-norm}.  In ADM coordinates
\[
 g_{00}=-N^2+\gamma_{ij}\beta^i\beta^j,\qquad
 g_{0i}=\gamma_{ij}\beta^j.
\]
At $N=1,\beta=0$ these imply
$\partial_tg_{00}=-2\partial_tN$ and
$\partial_tg_{0i}=\gamma_{ij}\partial_t\beta^j$.
Small compatible stable data have these first derivatives $O(a)$, while
$\gamma$ and its inverse have uniform positive margins.  Hence first-jet
shadowing forces $\norm{\lambda_t}\to0$.  Equation
\eqref{v1:rate-norm} then proves \eqref{v1:rate-necessary}.
No converse or common time interval is inferred.
\end{proof}

For a fixed small amplitude instead of $a(h)\to0$, the relevant test is the
difference from the \emph{actual stable} initial lapse--shift velocity, not
vanishing of that velocity.  A controlled change to harmonic coordinates
could provide another route, but its effect on the anchored operational
reads would have to be proved.  It is not performed here.

\begin{remark}[Why the static criterion is not an equivalent bridge theorem]
Even a one-column comparison displays the distinction.  With
$K_0=\diag(-4,2)$, $F_h=h(0,1)^T$ and $b_h=h^2$, one has
$v_h=h(0,1)^T$ but $u_h=2$.  If it occupies a weak amplitude slot, the
coefficient in \eqref{v1:rate-grade} is $2a^{-2}$.
This is only a linear-algebra witness, not a member of the prepared
original family.  Its purpose is to show that source-response decay alone
cannot imply the necessary clock test as an abstract theorem.
\end{remark}

The source's forced stable energy gives a comparison theorem once an actual
raw metric history has a common interval, a bounded physical chart, and a
vanishing acceleration residual with its required differentiated readouts.
Those comparison results are inherited.  The present work does not insert
those conclusions as hypotheses and call that a new exact-action closure.
In particular, a fixed-cutoff interval of length $O_h(a^3)$ cannot by itself
supply the required cutoff-independent interval.

\section{Results obtained and the next unresolved computation}
\label{v1:status}

The concrete original-action advances in this note are
\cref{v1:transverse-rank,v1:trace-leak-theorem,v1:safe-elimination}.
The transverse-shift rank and inverse no longer require a cutoff-by-cutoff
rank scan.  The remaining leading rank question can be formulated entirely
on $\overline G_h$, with $2V+1$ columns, and all signed trace amplification
can be evaluated using the $V$-dimensional matrix $H_h$.  The trace-leak
calculation explains why the larger active block cannot be eliminated by
an unproved trace-free argument.

The scalar and differentiated reductions are exact tools, not completed
cofinal estimates.  In particular, the newly exposed temporal sources and
clock test prevent a premature identification of three different targets:
\[
 \begin{gathered}
 \text{small initial electric coefficient},\\
 \text{controlled spacetime curvature reads},\\
 \text{shadowing in a full harmonic metric-jet chart}.
 \end{gathered}
\]
They need not follow from one another without additional estimates.

\paragraph{Immediate action-native target.}
Using the original prepared $J_{2,h}$ and the actual weak drift, form
\[
 \begin{aligned}
 \overline G_h&=(I-\Pi_{A,h})
        (J_{1,h}|_{\text{mean-zero lapse}},J_{2,h}|_{\text{curl-free shift}}),\\
 w_h&=\overline G_h(\overline G_h^*\overline G_h)^{-1}d_h,
 \end{aligned}
\]
then $C_h=E^*\Pi_{\overline G_h}E$ and $H_h=2I-3C_h$.
The next missing estimates are full rank of $\overline G_h$, regularity of
$H_h$, and decay of the \emph{actual} $w_h$ and $H_h^{-1}E^*w_h$ along a
cofinal family.  These matrices retain the original second preparation and
logarithmic coefficient terms.  Their Fourier support is not assumed finite
merely because the input drift has finite support.

For a physical-interval result one must additionally propagate the actual
finite-amplitude rank and chart margins, retain the positive inclusion
$Z$, bound the sources of \cref{v1:true-derivative}, and prove a common
physical lifespan.  Harmonic shadowing also requires the corresponding
lapse--shift control or a justified coordinate-comparison theorem.
The uniform nonconcentration estimate near $2/3$ on the actual source is
\emph{not} obtained in this iteration.  No direct raw-action Einstein
continuum theorem is claimed.

\paragraph{Verification and reproducibility.}
The accompanying Python file verifies the local double-commutator coefficient
symbolically, regenerates the retained rational seed, checks the exact trace
value and its rational squared norm, tests the shift-symbol singular values
and projector, and compares the complete signed response with both exact
reductions on finite matrices.  It also recovers, modulo $101$, the already
reported three-site leading ranks $77$ and $78$ at $\rho=0,1$ respectively;
these are consistency checks of inherited finite results, not new all-cutoff
rank assertions.  The true-Hessian correction and derivative identities are
tested separately.  Matrix comparison cases are explicitly
separated from original-seed calculations.  Floating checks support the
implementation; the all-cutoff statements rest on the proofs above.  No full
$J_2$ archive, cofinal signed inverse, or nonlinear exact-action spacetime
integration is represented as having been computed.

\appendix
\section{Retained seed and source fingerprints}
\label{v1:table}

The data below are copied from \cite{OB25},
\src{obv20:free-section}, as fixed input, not as a new construction.
For $0\le a,b,c\le2$ let $b_{abc}\in\R^3$ be the three-site vector array
whose rows are indexed by $(b,c)$ and whose three blocks give $80b_{abc}$
at $a=0,1,2$:
\begin{center}\small
\begin{tabular}{@{}c|rrr|rrr|rrr@{}}
\toprule
$(b,c)$&\multicolumn{3}{c|}{$a=0$}&\multicolumn{3}{c|}{$a=1$}&\multicolumn{3}{c}{$a=2$}\\\midrule
$(0,0)$&3&-3&-3&-2&-3&0&-1&2&2\\
$(0,1)$&1&-3&1&-3&-1&0&0&-3&0\\
$(0,2)$&1&1&-2&2&-1&3&0&2&-1\\
$(1,0)$&3&-1&-3&2&0&-2&-3&0&0\\
$(1,1)$&-1&3&3&0&0&3&2&2&0\\
$(1,2)$&-2&0&-2&1&-2&-1&0&1&2\\
$(2,0)$&3&-2&1&3&1&1&1&3&-3\\
$(2,1)$&3&-1&-1&-3&2&-2&0&3&3\\
$(2,2)$&0&-3&-1&-3&-2&-2&-2&-2&3\\\bottomrule
\end{tabular}
\end{center}
The refinement-compatible field is
\[
 L_a(x)=\frac{1+2\cos(2\pi(x-a/3))}{3},\qquad
 \xi^\sharp(x)=\sum_{a,b,c=0}^2b_{abc}L_a(x_1)L_b(x_2)L_c(x_3).
\]
In addition to the pointwise trace value in \cref{v1:trace-leak-theorem},
the exact normalized spatial-mean squared trace at $N=3,\rho=1$ is
$18981/25600$.  This is a three-site diagnostic, not a continuum norm.

The source files inspected for nonduplication have SHA-256 fingerprints:
\begin{description}[style=nextline,leftmargin=1.5em,labelindent=0pt]
\item[Open-basin source]
\path{Renewal_CMP_Open_Basin_3plus1_Research_Notes_V25(3).tex}\\
\path{0a7faaecf9866219c27bfa02b0d0a987b7c2fdfdf316ea38b235d0784eaf55ea}
\item[Grand Tensor source]
\path{Renewal_Grand_Tensor_Monograph_V076(6).tex}\\
\path{f4eed3332b5627bb6e19f95a7aba6695c806b8a1353c6376d299f9eb012c4ce3}
\end{description}

\begin{thebibliography}{9}
\raggedright
\bibitem{OB25}
\emph{Renewal Open-Basin $3+1$ Completion}, cumulative research source
V1--V25, supplied manuscript
\path{Renewal_CMP_Open_Basin_3plus1_Research_Notes_V25(3).tex}.
The literal labels cited in \cref{v1:context} distinguish inherited results
from the results of this note.
\bibitem{GT76}
\emph{Renewal Grand Tensor Monograph}, V076, supplied manuscript
\path{Renewal_Grand_Tensor_Monograph_V076(6).tex}, especially
\src{sec:GT76-Cartan-response} and
\src{thm:GT76-Cartan-cofinal}.
\bibitem{BG}
\AA. Bj\"orck and G. H. Golub,
Numerical methods for computing angles between linear subspaces,
\emph{Mathematics of Computation} \textbf{27} (1973), 579--594.
\href{https://doi.org/10.1090/S0025-5718-1973-0348991-3}{doi:10.1090/S0025-5718-1973-0348991-3}.
This is methodological context for subspace angles, not a proof of the
renewal source identities.
\bibitem{WAK}
S. Wang, D. Appel\"o and G. Kreiss,
An energy-based summation-by-parts finite difference method for the wave
equation in second order form, arXiv:2103.02006, version 2 (2022).
\url{https://arxiv.org/abs/2103.02006v2}.
This supplies context for energy-compatible discretization, not the
exact-action bridge claimed here.
\end{thebibliography}
% END IMMUTABLE EXACT-ACTION BRIDGE V1 BODY
% Future iterations append their research bodies here, before end{document}.

% BEGIN IMMUTABLE EXACT-ACTION BRIDGE V2 BODY
\clearpage
\renewcommand{\thesection}{\arabic{section}}
\renewcommand{\theHsection}{bridgeV2.\arabic{section}}
\setcounter{section}{10}
\section{An upstream change of initial records, not of the action}
\label{v2:context}

The preceding body leaves both the lapse part of the first source and the
complete second weak source unresolved.  On its fixed symmetric-gradient
seed, the lapse source also has a nonzero connection trace.  Rather than
postulate another bound on the resulting inverse, we return to the choice
of free initial records.  We retain the same canonical action, phase
derivatives, literal links, constraint map and physical pairing.  Only the
relative amplitudes of the three original axial configuration modes are
allowed to differ.  No new evolution equation is introduced.

This change yields a quantitative all-cutoff result for the \emph{entire
leading active source}.  Its proof uses an exact Jacobi representation of
the cosine product defect and a finite spectral-separation argument.  One
can select rational amplitude ratios differing from unity by $O(h)$, so the
new finite initial records have the same limiting initial data as the
unperturbed equal-amplitude family.  On this selected sequence the active
source has rank $3N^3-3$, is trace-free, and has a positive signed Gram with
an explicit polynomial lower bound.  Thus its elimination is legitimate at
every sufficiently fine odd cutoff and leaves only $N^3+2$ weak columns.

These statements concern leading amplitude coefficients.  They neither
assert full rank of the remaining weak source nor identify the nonlinear
finite-action flow with the stable writer.  In particular the finite-rank
certificates on the previous nonzero-$\rho$ family are \emph{not}
transferred to this different seed.  The actual second preparation and
second source must be evaluated on the new family.

\paragraph{Dependency and nonduplication audit.}
The source identities in \cref{v1:first-source}, and the original
unperturbed full first-source formula in \cite{OB25},
\src{obv19:J-geometric}--\src{obv19:J-electric}, are retained.  The
uniform initial range solve, analytic constraint map and mean calibration
method are inherited from \cite{OB25}, \src{obv17:range},
\src{obv17:quadratic-charge} and \src{obv17:exact-initializer}.
The following proof does not re-establish that general initializer or the
stable-writer Cauchy theorem.  The new ingredients are the cosine-defect
spectrum and its parity-dependent small eigenvalues, the tensor-product
lapse-rank calculation, the rational cofinal seed selection, and their
application to the complete leading active Cartan block.  Jacobi-matrix
spectral and bidiagonal methods are classical; see, for example,
\cite{v2:LST}.  All special matrix identities needed here are proved below.

\section{A three-amplitude family inside the original constraint chart}
\label{v2:family}

Continue to write $N=2m+1$, $h=N^{-1}$ and $V=N^3$; $N$ here is the number
of grid points in one direction, not the lapse.  Let $\mathsf T_i$ be the
symmetric configuration matrix whose two transverse off-diagonal entries
are one.  These configuration tensors are not the Lorentz rotation
generators.  For positive coefficients
$\alpha=(\alpha_1,\alpha_2,\alpha_3)$ set
\begin{equation}
 U_\alpha=\sum_{i=1}^3\alpha_i\mathsf T_i\cos(2\pi x_i),\qquad
 s_{h,\alpha}=\omega_h\sqrt{\frac{\alpha_1^2+\alpha_2^2+\alpha_3^2}{6}},
 \qquad p_{h,\alpha}=-s_{h,\alpha}I.
 \label{v2:seed}
\end{equation}
All amplitudes will remain in a fixed compact subset of $(0,\infty)^3$.
In particular $s_{h,\alpha}$ and its inverse are uniformly bounded.  At
$\alpha=(1,1,1)$ this is exactly the unperturbed $a$-normalized first jet
in \cref{v1:first-data,v1:common-seed}.

\begin{proposition}[Exact original-action preparation of the altered seed]
\label{v2:preparation}
For sufficiently small $|a|$, uniformly for $\alpha$ in a sufficiently
small fixed neighbourhood of $(1,1,1)$ and all sufficiently fine odd grids,
the original initial constraint map admits exact data
\begin{equation}
 \mathscr C_h(X_{h,\alpha}(a))=0,\qquad
 X_{h,\alpha}(a)=a(U_\alpha,p_{h,\alpha})+O(a^2)
 \label{v2:prepared}
\end{equation}
in the inherited Sobolev preparation spaces.  The remainder, and its fixed
parameter derivatives, have cutoff-independent bounds.  The full original
logarithmic action is used in the preparation.  Moreover
\begin{equation}
 \|X_{h,\alpha}(a)-X_{h,\widetilde\alpha}(a)\|_{\mathcal P_h^r}
       \le C_r|a|\,|\alpha-\widetilde\alpha|.
 \label{v2:seed-comparison}
\end{equation}
Consequently an $O(h)$ change of the three amplitudes changes the prepared
initial records by $O(|a|h)$, not their continuum initial limit.
\end{proposition}
\begin{proof}
The configuration is transverse and trace-free.  Allow, for the mean solve,
$\pi=-sI+\sum_i b_i\mathsf T_i\sin(2\pi x_i)$.  The inherited quadratic
canonical Hamiltonian, with constant lapse and shift kept independent,
gives the complete quadratic mean constraint vector
\begin{equation}
 Q_{h,\alpha}(s,b)=
 \begin{pmatrix}
 \frac32s^2-\sum_i b_i^2-\frac14\omega_h^2\sum_i\alpha_i^2\\
 -\omega_h\alpha_1 b_1\\-\omega_h\alpha_2 b_2\\-\omega_h\alpha_3 b_3
 \end{pmatrix}.
 \label{v2:mean-Q}
\end{equation}
Indeed $\|\mathsf T_i\|_F^2=2$, the mean sine and cosine squares are
$1/2$, and different axial modes are orthogonal on every odd grid.  Thus
$\|\pi\|^2-\|\tr\pi\|^2/2=-3s^2/2+\sum b_i^2$,
$\sum_i\|\delta_iU_\alpha\|^2=\omega_h^2\sum\alpha_i^2$, and
$\ip\pi{\delta_iU_\alpha}=-\omega_h\alpha_i b_i$.  This proves
\eqref{v2:mean-Q}, including its normalization.  It has the root
$(s,b)=(s_{h,\alpha},0)$ and derivative
\[
 \diag(3s_{h,\alpha},-\omega_h\alpha_1,
                         -\omega_h\alpha_2,-\omega_h\alpha_3),
\]
with a uniform inverse on the stated amplitude set.

Apply the original nonconstant-row contraction, retaining its four mean
rows.  Its range correction is $O(a^2)$ with uniform analytic derivatives.
After division of the remaining mean map by $a^2$, the integral Taylor
formula at $a=0$ gives $Q_{h,\alpha}+O(a)$ in $C^1$, uniformly in
$\alpha,h$.  The displayed inverse therefore gives the original four-mode
contraction with $s=s_{h,\alpha}+O(a)$ and $b=O(a)$.  This is the same
nonlinear constraint map and the same range solve as the inherited
initializer, with its quadratic starting root recalibrated, not a solve
of a replacement constraint.  Implicit differentiation in $\alpha$ is
uniform because the root derivative and range inverse have the margins
just described.  Multiplication by $a$ in the initial chart gives
\eqref{v2:seed-comparison}.  These assertions concern initial records only.
\end{proof}

For this family the first spatial rotation matrix is diagonal:
\begin{equation}
 r_i=\frac12(\alpha_{i+2}\delta_{i+2}\cos(2\pi x_{i+2})
                  -\alpha_{i+1}\delta_{i+1}\cos(2\pi x_{i+1})),
 \quad r_1+r_2+r_3=0,
 \label{v2:r}
\end{equation}
where indices are cyclic.  The scalar boost coefficient is
$\varkappa_{h,\alpha}=s_{h,\alpha}/2$.  The electric first-source
calculation in \cref{v1:first-source} is polynomial in this scalar; its
shift formula holds with $\varkappa_h$ replaced by
$\varkappa_{h,\alpha}$, while its lapse part is unchanged.  In particular
the full first source is trace-free on this entire three-amplitude family.

\section{The exact cosine product-defect operator}
\label{v2:Jacobi}

On the one-dimensional odd grid let $f(x)=\cos(2\pi x)$ and define
\begin{equation}
 D_hu=\delta(fu)-f\delta u-(\delta f)u.
 \label{v2:D}
\end{equation}
Multiplication is the original cyclic multiplication.  This is an operator
on lapse tests; it is neither a continuum derivative nor the Cartan
trace-compression matrix of \cref{v1:trace-data}.

\begin{theorem}[A one-point kernel and an explicit irreducible Jacobi block]
\label{v2:Jacobi-theorem}
The real matrix $D_h$ is self-adjoint, $D_h1=0$, and has $N$ distinct
real eigenvalues.  Its nonconstant part is unitarily equivalent to a
zero-diagonal Jacobi matrix of size $2m=N-1$.  Write $s=\pi h$ and
\begin{equation}
 \begin{aligned}
 t_j&=\frac4h\sin\frac{s}{2}\sin\frac{js}{2}\sin\frac{(j+1)s}{2}
                    &&(1\le j<m),\\
 t_m&=\frac2h\cos\frac{s}{2}+\frac1h\sin s,\\
 t_{2m-j}&=t_j &&(1\le j<m).
 \end{aligned}
 \label{v2:Jacobi-weights}
\end{equation}
These are its strictly positive successive off-diagonal entries.  The
nonzero eigenvalues occur in opposite pairs and
\begin{equation}
 \ker D_h=\R1,\qquad \rank D_h=N-1,
 \qquad c h^{-1}\le\|D_h\|_{\op}\le C h^{-1}.
 \label{v2:D-rank}
\end{equation}
Every eigenvalue is algebraic at a fixed integer $N$.
\end{theorem}
\begin{proof}
Since $\delta^*=-\delta$ and multiplication by a real field is
self-adjoint, $[\delta,M_f]-M_{\delta f}$ is self-adjoint.  Direct
substitution gives $D_h1=0$.  In the orthonormal signed Fourier basis its
only possible entries send mode $k$ to the folded modes $[k\pm1]$, with
coefficients
\[
 \frac i2\{\kappa([k+1])-\kappa(k)-\omega_h\},\qquad
 \frac i2\{\kappa([k-1])-\kappa(k)+\omega_h\}.
\]
The zero-mode row and column vanish.  Order the other modes as
$1,2,\ldots,m,-m,-m+1,\ldots,-1$.  They form one path.  For $1\le j<m$,
\[
 \frac{\kappa(j+1)-\kappa(j)-\omega_h}{2}
 =-\frac4h\sin\frac{s}{2}\sin\frac{js}{2}\sin\frac{(j+1)s}{2}.
\]
The folded central edge has coefficient
$-(2\kappa(m)+\omega_h)/2=-t_m$.  Reflection gives the other weights.
A diagonal unitary phase change makes all lower and upper off-diagonal
entries positive.  The fold has not been replaced by an unfurled symbol.

All $t_j$ are nonzero.  The first component of a Jacobi eigenvector
determines its successive components by the three-term recurrence; if the
first is zero, every component is zero.  Thus every eigenspace has
dimension at most one, and self-adjointness proves simplicity.  Expanding
the zero-diagonal determinant gives
$\det J_{2m}=(-1)^m(t_1t_3\cdots t_{2m-1})^2\ne0$.
Hence zero occurs only in the isolated constant mode.  Conjugation of
$J_{2m}$ by $\diag(1,-1,1,-1,\ldots)$ changes its sign, proving the paired
spectrum.  Its norm is at least its central entry and at most its largest
absolute row sum.  The interior weights are bounded and $t_m\asymp h^{-1}$,
which proves \eqref{v2:D-rank}.  Finally all sine and cosine values used at
integer $N$ are algebraic; the characteristic polynomial therefore has
algebraic coefficients and algebraic roots.
\end{proof}

\subsection{A sharp parity-dependent small eigenvalue}

The next estimate is used quantitatively in the seed selection.  It also
shows why rank alone must not be confused with an inverse margin.

\begin{theorem}[Fold-induced soft pair and inverse bounds]
\label{v2:soft-pair}
Let $D_h^\dagger$ be zero on constants and the inverse on their orthogonal
complement.  For sufficiently fine odd grids there are absolute positive
constants such that
\begin{equation}
 \begin{array}{ll}
 N\equiv1\pmod4:&c h^{-3}\le\|D_h^\dagger\|_{\op}\le C h^{-3},\\[2pt]
 N\equiv3\pmod4:&c h^{-2}\le\|D_h^\dagger\|_{\op}\le C h^{-2}.
 \end{array}
 \label{v2:inverse-parity}
\end{equation}
On $N\equiv1\pmod4$, precisely one opposite pair has magnitude smaller
than the order-$h^2$ lower bound for all other nonzero eigenvalues, and
\begin{equation}
 \lambda_{\mathrm{soft},\pm}(h)=\pm\bigl(2\pi^2h^3+O(h^4)\bigr).
 \label{v2:soft-asymptotic}
\end{equation}
In both parity classes the complete nonzero spectrum $\Sigma_h^\times$
satisfies
\begin{equation}
 \min_{\lambda\in\Sigma_h^\times}|\lambda|\ge c h^3,
 \qquad\sum_{\lambda\in\Sigma_h^\times}\frac1{|\lambda|}\le C h^{-3}.
 \label{v2:inverse-sum}
\end{equation}
\end{theorem}
\begin{proof}
Put $c_h=4h^{-1}\sin(s/2)$, $a_j=\sin(js/2)$ for $1\le j\le m$, and
\[
 (d_1,\ldots,d_{2m})=(a_1,\ldots,a_m,a_m,\ldots,a_1),\qquad
 b_h=\frac{t_m}{c_ha_m^2}.
\]
Here $c_h\asymp1$, $a_m\asymp1$, $b_h\asymp h^{-1}$, $b_h\ge1$, and
$\sum_{j=1}^{2m}d_j^{-2}\le C h^{-2}$.  The last assertion follows from
$a_j\ge jh$ and the convergent sum $\sum_{j\ge1}j^{-2}$.

Order the Jacobi path into odd and even vertices.  Its matrix is
$\left(\begin{smallmatrix}0&B_h\\B_h^*&0\end{smallmatrix}\right)$, where
$(B_h)_{ii}=t_{2i-1}$ and $(B_h)_{i+1,i}=t_{2i}$, with all other
entries zero.
It factors as
\[
 B_h=c_h\,\diag(d_{2i-1})\,\widehat B_h\,\diag(d_{2i}).
\]
All nonzero entries of $\widehat B_h$ are one except the central entry,
which is $b_h$.  It is on the diagonal when $m$ is odd and on the
subdiagonal when $m$ is even.  For any lower bidiagonal matrix with
nonzero diagonal, its inverse is zero above the diagonal and
\begin{equation}
 (B_h^{-1})_{ij}=(-1)^{i-j}
 \frac{t_{2j}t_{2j+2}\cdots t_{2i-2}}
      {t_{2j-1}t_{2j+1}\cdots t_{2i-1}},\qquad i\ge j.
 \label{v2:bidiagonal-inverse}
\end{equation}
Multiplication of a row of $B_h$ by this formula leaves one on its diagonal
and cancels every subdiagonal entry, proving the identity.

When $m$ is odd, the magnitudes of entries of $\widehat B_h^{-1}$ are
$0,1$ or $b_h^{-1}$.  The two diagonal weights and the sum just bounded
give $\|B_h^{-1}\|_{\mathrm{HS}}\le C h^{-2}$.  For $m\ge3$, its first
diagonal entry is $1/t_1\asymp h^{-2}$, giving the reverse norm bound.
Thus every nonzero Jacobi eigenvalue has magnitude at least $c h^2$.

When $m=2l$ is even, the entries crossing the central subdiagonal have
magnitude $b_h$, and all other nonzero entries have magnitude one.  Therefore
\begin{equation}
 B_h^{-1}=\frac{b_h}{c_h}\,u_h v_h^*+R_h,
 \qquad \|R_h\|_{\op}\le C h^{-2},
 \label{v2:rank-one-inverse}
\end{equation}
where $(u_h)_i=(-1)^i/d_{2i}$ for $i\ge l+1$ and zero otherwise, while
$(v_h)_j=(-1)^j/d_{2j-1}$ for $j\le l$ and zero otherwise.  The signs in
\eqref{v2:bidiagonal-inverse} prove this factorization entry by entry.
The nonzero singular value of its rank-one term is
\begin{equation}
 L_h=\frac{t_m}{c_h^2a_m^2}
            \sum_{j=1}^l\csc^2\!\left(\frac{(2j-1)\pi h}{2}\right).
 \label{v2:large-inverse-value}
\end{equation}
On $0<x\le\pi/4$, $\csc^2x-x^{-2}$ is bounded.  Subtracting this
bounded function term by term, and bounding the omitted tail of the odd
reciprocal-square sum, gives
\[
 \sum_{j=1}^l\csc^2\!\left(\frac{(2j-1)\pi h}{2}\right)
       =\frac1{2h^2}+O(h^{-1}).
\]
For completeness, $\sum_{j\ge1}(2j-1)^{-2}=\pi^2/8$ follows by
Parseval from the Fourier series of the unit square wave: its positive
odd sine coefficients are $4/(\pi(2j-1))$, and its squared $L^2$ mean is
one.  Also $t_m=2h^{-1}+O(1)$, $c_h=2\pi+O(h^2)$, and
$a_m^2=1/2+O(h)$.  Hence
\[
 L_h=\frac1{2\pi^2h^3}+O(h^{-2}).
\]
The largest singular value of $B_h^{-1}$ differs from $L_h$ by at most
$\|R_h\|$; its second singular value is at most $\|R_h\|$.
The first bound is the triangle inequality and its reverse, and the second
follows by restricting to the kernel of the rank-one term and applying the
min--max characterization.  The eigenvalues of the full bipartite Jacobi
matrix are the positive and negative singular values of $B_h$.  Reciprocation
therefore proves \eqref{v2:inverse-parity} and \eqref{v2:soft-asymptotic},
with all other magnitudes at least $c h^2$.

In the odd-$m$ case at most $N-1$ reciprocal eigenvalues are bounded by
$C h^{-2}$.  In the even-$m$ case two are bounded by $C h^{-3}$ and the
others by $C h^{-2}$.  Since $N=h^{-1}$, their sums are bounded by
$C h^{-3}$ in both cases.  Finitely many small grids can be absorbed by
changing constants.  This proves \eqref{v2:inverse-sum}.
\end{proof}

\section{The diagonal lapse source and its tensor spectrum}
\label{v2:diagonal}

Let $D_{i,h}$ act as \eqref{v2:D} in coordinate $i$ and as the identity
on the other two coordinates.  The original geometric first-source formula,
with $U_\alpha$ inserted, gives
\begin{equation}
 \mathscr D_{h,\alpha}n:=\diag\mathcal J_{h,\alpha}(n,\beta)
 =\begin{pmatrix}
  (\alpha_2D_{2,h}-\alpha_3D_{3,h})n\\
  (\alpha_3D_{3,h}-\alpha_1D_{1,h})n\\
  (\alpha_1D_{1,h}-\alpha_2D_{2,h})n
 \end{pmatrix}.
 \label{v2:diagonal-source}
\end{equation}
It is independent of $\beta$.  In fact its geometric off-diagonal entries
are zero: every such term differentiates an axial coefficient in a coordinate
on which that coefficient is constant.  The electric diagonal is zero by
\eqref{v2:r}.  Thus \eqref{v2:diagonal-source} is the diagonal of the
\emph{complete} original first source, not a geometric substitute for it.

\begin{theorem}[Exact tensor-product rank and kernel]
\label{v2:tensor-rank}
Let $\lambda_0=0,\lambda_1,\ldots,\lambda_{N-1}$ be the distinct
one-dimensional eigenvalues, with real orthonormal eigenvectors $\phi_j$
and $\phi_0=1$ in the normalized pairing.  In the product eigenbasis,
\begin{equation}
 \|\mathscr D_{h,\alpha}n\|_h^2
 =\sum_{a,b,c}|n_{abc}|^2
 \bigl[(\alpha_2\lambda_b-\alpha_3\lambda_c)^2
       +(\alpha_3\lambda_c-\alpha_1\lambda_a)^2
       +(\alpha_1\lambda_a-\alpha_2\lambda_b)^2\bigr].
 \label{v2:tensor-norm}
\end{equation}
For equal nonzero amplitudes, its kernel is exactly
$\operatorname{span}\{\phi_j\otimes\phi_j\otimes\phi_j:0\le j<N\}$,
and its rank is $V-N$.  More generally, if
\begin{equation}
 \alpha_1\Sigma_h\cap\alpha_2\Sigma_h=\{0\},
 \qquad \Sigma_h=\{\lambda_j:0\le j<N\},
 \label{v2:pair-separation}
\end{equation}
then its kernel is exactly the constants and its rank is $V-1$.
\end{theorem}
\begin{proof}
The three self-adjoint one-dimensional operators commute and have the
stated complete tensor eigenbasis.  Formula \eqref{v2:diagonal-source}
therefore diagonalizes to the three differences in
\eqref{v2:tensor-norm}.  The bracket is zero exactly when
$\alpha_1\lambda_a=\alpha_2\lambda_b=\alpha_3\lambda_c$.
At equal amplitudes, simplicity of the scalar spectrum forces $a=b=c$.
Under \eqref{v2:pair-separation}, their common value is zero, and nonzero
$\alpha_3$ forces $a=b=c=0$.  Dimension counts give the ranks.  Notice
that the equal-amplitude kernel described here is the kernel of the diagonal
observation only; its remaining off-diagonal source is not set to zero.
\end{proof}

\section{Quantitative cofinal separation by finite rational amplitudes}
\label{v2:selection}

A fixed nonalgebraic amplitude ratio would already exclude exact nonzero
spectral collisions, because \cref{v2:Jacobi-theorem} has an algebraic
spectrum at every cutoff.  That observation alone provides no rate and
need not serve as an operational initialization.  The following stronger
construction uses rational finite records approaching the original equal
amplitudes.

\begin{theorem}[Rational $O(h)$ seed detuning with an order-$h^6$ margin]
\label{v2:rational-selection}
For every sufficiently fine odd cutoff there is a rational number
$\theta_h$ with
\begin{equation}
 1<\theta_h<1+h,\qquad
 |\lambda_a-\theta_h\lambda_b|>h^6
       \quad(\lambda_a,\lambda_b\in\Sigma_h^\times).
 \label{v2:theta-margin}
\end{equation}
It may be chosen dyadic with denominator $2^{L_h}$, where
$L_h\le C+7\log_2N$.  A finite algorithm can select it by enumerating
those dyadic candidates and testing algebraic spectral inequalities.
For $\alpha_h=(1,\theta_h,1)$,
\begin{equation}
 \|\mathscr D_{h,\alpha_h}n\|_h\ge c h^6\|n\|_h
                    \qquad (\ip1n_h=0).
 \label{v2:diagonal-margin}
\end{equation}
The constant is independent of the selected candidate and the cutoff.
No floating rank threshold or infinite-precision transcendental input is
required at an individual cutoff.
\end{theorem}
\begin{proof}
For a tolerance $\eta>0$ the forbidden set of ratios is
\[
 \mathcal E_h(\eta)=\bigcup_{\lambda_a,\lambda_b\in\Sigma_h^\times}
          \{\theta:|\lambda_a-\theta\lambda_b|<\eta\}.
\]
Each interval has length $2\eta/|\lambda_b|$.  By
\eqref{v2:inverse-sum}, its total Lebesgue measure is at most
\begin{equation}
 |\mathcal E_h(\eta)|\le2\eta(N-1)
          \sum_{\lambda_b\ne0}|\lambda_b|^{-1}
          \le C\eta h^{-4}.
 \label{v2:bad-measure}
\end{equation}
Take $\eta=4h^6$.  For sufficiently small $h$, the forbidden measure is
less than the length $h/2$ of $(1+h/4,1+3h/4)$.  There is a point
$\overline\theta_h$ in that interval with every distance at least $4h^6$.
Since $\max|\lambda_b|\le C_0h^{-1}$, with $C_0=7$ sufficient from
$\|D_h\|\le4h^{-1}+2\pi$, a dyadic approximation with error
at most $h^7/C_0$ retains distances at least $3h^6$ and remains in
$(1,1+h)$.  Choosing the dyadic mesh smaller than this error gives the
stated bound on $L_h$ (with an enlarged absolute constant).

There are finitely many candidates at that precision.  Their coefficients
and the eigenvalues of the displayed Jacobi matrix are algebraic, so each
strict inequality against $h^6$ is decidable by real algebraic root
isolation and comparison.  A candidate with a strict larger margin was
just proved to exist, so enumeration terminates.  This is a finite
selection proof and a bound on amplitude-record bit length; no polynomial
bound on the total algebraic eigensolver cost or physical service deadline
is asserted.

For a nonconstant tensor mode, if $a,b$ are both nonzero then the last
squared difference in \eqref{v2:tensor-norm} is greater than $h^{12}$.
If exactly one is nonzero, its magnitude is at least $c h^3$ by
\eqref{v2:inverse-sum}, with $1<\theta_h<2$.  If $a=b=0$ but $c\ne0$,
the first two differences have magnitude $|\lambda_c|\ge c h^3$.
Taking their minimum and summing proves \eqref{v2:diagonal-margin}.
\end{proof}

\begin{corollary}[Fixed amplitude ratios also work on a full-measure set]
\label{v2:fixed-ratios}
For almost every fixed $\theta$ in any compact positive interval, and any
fixed exponent $p>5$, there is $c_\theta>0$ such that all odd cutoffs obey
\[
 |\lambda_a-\theta\lambda_b|\ge c_\theta h^p
                \quad(\lambda_a,\lambda_b\ne0).
\]
For these ratios the complete diagonal lapse kernel is constant at every
cutoff.  This is a statement about Lebesgue measure on an initial-parameter
interval, not a probabilistic statement about the physical renewal law.
\end{corollary}
\begin{proof}
The measure bound \eqref{v2:bad-measure}, with $\eta=h^p$, is summable
over odd $N$ when $p>5$.  The first Borel--Cantelli argument (equivalently,
subadditivity for the tail unions) excludes all but finitely many bad
cutoffs for almost every ratio.  Exact collisions form a countable union
of finite ratio sets and have measure zero.  The finitely many remaining
positive gaps can be absorbed into $c_\theta$.  Tensor separation follows
as before.  The full-measure assertion does not assert a cutoff-independent
open interval avoiding every resonance, or dynamical attraction to a ratio.
\end{proof}

\section{The complete leading active source is now controlled}
\label{v2:active}

Let $\mathcal T_h$ be the transverse shift space of
\eqref{v1:transverse-space}, and let $\mathcal L_h$ be the mean-zero lapse
space.  In this section $\alpha_h=(1,\theta_h,1)$ is selected by
\cref{v2:rational-selection} and
\[
 \mathcal A_h=\mathcal J_{h,\alpha_h}|_{\mathcal L_h\oplus\mathcal T_h}.
\]
The notation $\mathcal A_h$ denotes the whole leading active source, not
just the transverse-shift block $A_h$ used in V1.

\begin{theorem}[All-cutoff active rank, positive Cartan Gram, and a polynomial inverse]
\label{v2:active-theorem}
On the newly prepared family, at every sufficiently fine odd cutoff,
\begin{equation}
 \ker\mathcal J_{h,\alpha_h}
   =\R(1,0)\oplus\{(0,\beta):\curl_\delta\beta=0\},
 \qquad \rank\mathcal A_h=3V-3.
 \label{v2:full-active-rank}
\end{equation}
Moreover $P\mathcal A_h=0$ and, in the positive multiplier and Frobenius
pairings,
\begin{equation}
 \mathcal A_h^*K_0^{-1}\mathcal A_h
       =\frac12\mathcal A_h^*\mathcal A_h\succeq c h^{16}I.
 \label{v2:active-Gram}
\end{equation}
For every mean-zero lapse $n$ and transverse shift $\beta$ one has
\begin{equation}
 \|n\|_h\le C h^{-6}\|\mathcal A_h(n,\beta)\|_h,
 \qquad
 \|\Omega_h\beta\|_h\le C h^{-8}
                                      \|\mathcal A_h(n,\beta)\|_h.
 \label{v2:active-inverse}
\end{equation}
These exponents are sufficient bounds, not optimality claims.  They concern
$J_1$, not the complete nonlinear multiplier matrix.
\end{theorem}
\begin{proof}
The complete source diagonal is \eqref{v2:diagonal-source} and contains
no shift.  If the source is zero, \cref{v2:rational-selection} forces the
lapse to be constant.  A constant lapse has zero first source by the
original unit-lapse homogeneity identity; it also follows by putting $n=1$
in the product-defect and electric formulas.  The remaining pure-shift
source is exactly that in \cref{v1:first-source}, with the uniformly
positive scalar $\varkappa_{h,\alpha_h}$.  The proof of
\cref{v1:transverse-rank} therefore forces its transverse component to
vanish, and its kernel is precisely the curl-free shifts.  This proves
\eqref{v2:full-active-rank}.

The geometric trace is zero by the alternating contraction of a symmetric
configuration.  Here the rotational first connection is diagonal, so its
electric lapse diagonal and the pure-shift diagonal both vanish.  Thus the
whole first source is trace-free, including its remaining off-diagonal
entries.  This proves the identity in \eqref{v2:active-Gram}.

Write $y=\mathcal A_h(n,\beta)$.  The diagonal observation and
\eqref{v2:diagonal-margin} give the first inequality in
\eqref{v2:active-inverse}.  The lapse-only operator has the full finite-grid
bound $\|\mathcal J_{h,\alpha_h}(\cdot,0)\|_{L^2\to L^2}\le C h^{-1}$.
Indeed the geometric part contains at most one phase derivative on a
product or on a test, each of norm at most $2h^{-1}$, with uniformly
bounded axial coefficients.  In the electric part $r_i$ is uniformly
bounded and $hS_i\delta_j$ is bounded.  This argument includes all modes,
including folded ones.  Consequently
\[
 \|J_{1,h}(0,\beta)\|_h
 \le \|y\|_h+C h^{-1}\|n\|_h
 \le C h^{-7}\|y\|_h.
\]
The lower transverse estimate \eqref{v1:transverse-bounds} gives the
second inequality in \eqref{v2:active-inverse}.  The nonzero-mode
Poincar\'e bound for $\Omega_h$ gives $\|\beta\|_h\le C\|\Omega_h\beta\|_h$.
Thus the smallest singular value of $\mathcal A_h$ is at least $c h^8$;
squaring proves the last bound in \eqref{v2:active-Gram}.
\end{proof}

\begin{corollary}[Exact elimination down to the $V+2$ weak columns]
\label{v2:weak-only}
Let $\mathcal W_h=\{\beta:\curl_\delta\beta=0\}$, of dimension $V+2$,
and let $\mathcal W_h^{(2)}=J_{2,h}|_{\mathcal W_h}$ be the \emph{actual}
second source of the newly prepared family, including its second initial
coefficient and all contributing logarithmic terms.  Put
\[
 \Pi_{\mathcal A}=\mathcal A_h(\mathcal A_h^*\mathcal A_h)^{-1}
                     \mathcal A_h^*,\qquad
 \mathcal F_h=(I-\Pi_{\mathcal A})\mathcal W_h^{(2)}.
\]
The complete leading graded signed Gram is congruent to
\begin{equation}
 \diag\left(\tfrac12\mathcal A_h^*\mathcal A_h,
                         \mathcal F_h^*K_0^{-1}\mathcal F_h\right).
 \label{v2:weak-Schur}
\end{equation}
Thus its only remaining rank and signed-invertibility tests have size
$(V+2)\times(V+2)$.  When that Gram is invertible, its complete leading
physical response to the original weak target $b_{W,h}$ is
\begin{equation}
 v_h=K_0^{-1}\mathcal F_h
             (\mathcal F_h^*K_0^{-1}\mathcal F_h)^{-1}b_{W,h}.
 \label{v2:weak-response}
\end{equation}
All original active equations $\mathcal A_h^*v_h=0$ are preserved.
\end{corollary}
\begin{proof}
The proof of \cref{v1:safe-elimination} uses only injectivity of the
eliminated block and the identity $K_0^{-1}A=A/2$.  Both now hold for the
\emph{whole} $\mathcal A_h$, by \cref{v2:active-theorem}.  Its triangular
column change leaves the leading graded target $(0,b_{W,h})$ unchanged,
and gives \eqref{v2:weak-Schur} and \eqref{v2:weak-response} exactly.
No second-source column or original equation is removed by definition.
\end{proof}

\paragraph{A finite exact calibration.}
At $N=3$ the scalar characteristic polynomial is
$x(x^2-243/4)$.  Thus the rational choice $\alpha=(1,9/8,1)$ has
pairwise separated nonzero spectra in the first two coordinates.
Its diagonal lapse rank is $26$ and its full leading active rank is $78$.
The companion computation regenerates the complete first-source columns
and confirms both lower-rank certificates by exact elimination modulo
$101$; the displayed spectrum and kernel proof give the matching
characteristic-zero upper bounds.  This calculation uses a different seed
from the symmetric-gradient calibration in V1.

The scalar-trace reconstruction in \cref{v1:scalar-reconstruction} now
applies to $\mathcal F_h$ inside $\Ran(\mathcal A_h)^\perp$, since the
trace inclusion is orthogonal to the entire active range.  Its remaining
trace matrix still has size $V$, and the possible signed resonance is
still $2/3$.  The positive active Gram is not a proof that this last signed
Gram is regular.  Nor does a leading coefficient bound imply a common
finite-amplitude chart for the exact evolution.

\section{The detuning retains a nonzero original weak source}
\label{v2:weak-drift}

The altered seed must not obtain regularity merely by replacing the actual
drift by zero.  Its weak quadratic drift can be tracked explicitly.

\begin{proposition}[Weak drift on the three-amplitude scalar-momentum family]
\label{v2:drift}
Let $d_{h,\alpha}$ be the curl-free projection of the complete original
quadratic momentum drift on the first jet \eqref{v2:seed}.  In the same
$a$ normalization as V1 and the supplied source,
\begin{equation}
 d_{h,\alpha}=-4d_h^{\rm ax}
       (\alpha_i^2\sin(4\pi x_i))_{i=1}^3,
 \qquad d_h^{\rm ax}=\frac{\omega_h^3(1-\cos\pi h)}8.
 \label{v2:drift-vector}
\end{equation}
The doubled frequency is folded at $N=3$.  For sufficiently fine grids,
uniformly on compact positive amplitude sets and at each fixed Sobolev
index $r$,
\begin{equation}
 \mathcal I_hd_{h,\alpha}
       =-2\pi^5h^2(\alpha_i^2\sin(4\pi x_i))_{i=1}^3+O_{H^r}(h^4),
 \qquad c_rh^2\le\|d_{h,\alpha}\|_{H_h^r}\le C_rh^2.
 \label{v2:drift-size}
\end{equation}
For the rational sequence $\alpha_h=(1,\theta_h,1)$, the leading drift
therefore agrees with the equal-amplitude continuum coefficient up to
$O(h^3)$, while remaining nonzero of order $h^2$.
\end{proposition}
\begin{proof}
Use the inherited complete weak-shift quadratic identity and drift recipe
in \cite{OB25}, \src{obv23:weak-action}, \src{obv24:drift-recipe} and
\src{obv25:link-cancellation}.  The cancellation of the tested cubic
logarithm uses that a curl-free metric variation has zero flat stationary
connection.  It does not depend on equality of the three axial amplitudes
or on the value of the constant scalar momentum.  Thus it applies here
without changing the Hamiltonian or the multiplier source.

The remaining weak drift is quadratic in the first records.  Terms with
two constant scalar momentum factors give an integrated derivative in the
weak test and vanish.  The term pairing that momentum with the weak
configuration generator acting on the linear configuration velocity also
vanishes: both the momentum and that velocity are constant isotropic
tensors, and the exact trace-generator identity is a divergence.  The
spatial linear force is $-\omega_h^2 U_\alpha/2$.
Cross products of two distinct axial configuration modes have an exact
phase product rule in each coordinate, since one factor is constant in
that coordinate.  Their weak Ward cancellation is therefore the same
exact cancellation as for the inherited equal-amplitude vector formula.
A repeated mode in direction $i$ retains precisely its scalar factor
$\alpha_i^2$; its complete coefficient is the inherited repeated-mode
coefficient $-4d_h^{\rm ax}$.  These are all quadratic polarizations,
which proves \eqref{v2:drift-vector}.  The argument retains the folded
repeated mode and does not drop any term from the full second source.

Finally $\omega_h=2\pi+O(h^2)$ and
$1-\cos\pi h=\pi^2h^2/2+O(h^4)$ give the expansion.
The fixed Fourier bank has uniformly comparable discrete and continuum
Sobolev norms; positive amplitudes bounded away from zero give its lower
bound.  Since $\alpha_h-(1,1,1)=O(h)$, substitution changes the leading
$h^2$ coefficient by $O(h^3)$.  This is an application and extension of
the supplied weak-drift calculation to the new seed, not a claim to have
recomputed its complete $J_2$ inverse.
\end{proof}

\section{What is closed, and what must not be inferred}
\label{v2:status}

The upstream modification closes a specific previously unresolved block.
On rationally detuned original-action initial records, with an $O(h)$
change of the first mode amplitudes, the whole leading active source has
rank $3V-3$ and a positive signed Gram on every sufficiently fine odd grid.
The construction is compatible with the exact original initial constraints
and preserves the unperturbed continuum initial limit.  Its source inverse
has a proved, though deliberately conservative, polynomial bound.  These
are not deductions from a finite determinant scan.

The scalar spectral calculation also exposes a genuine refinement effect:
for $N\equiv1\pmod4$ the cosine product defect has an opposite pair of
eigenvalues $\pm(2\pi^2h^3+O(h^4))$.  In particular a putative uniform
order-$h^2$ lower bound for the diagonal lapse observation is false, even
after detuning.  Take a scalar soft eigenfunction in the first coordinate
and constants in the other two.  Formula \eqref{v2:tensor-norm} gives a
nonconstant unit lapse with diagonal source norm
$\sqrt2|\alpha_1|\,|\lambda_{\rm soft}|=O(h^3)$.
This statement concerns the diagonal observation only: the remaining
off-diagonal lapse source may improve the full active inverse.  The soft
Jacobi eigenvalue near zero is \emph{not} the Cartan trace-angle resonance
at $2/3$; the two spectral problems must not be conflated.

The exact-action bridge is still not closed.  The next source calculation
is now on
\[
 \mathcal F_h=(I-\Pi_{\mathcal A})J_{2,h}|_{\ker\curl_\delta},
 \qquad b_{W,h}=d_{h,\alpha_h}
\]
in the stated weak-coordinate pairing.  One must retain the second
preparation obtained from the new root \eqref{v2:mean-Q}; the second
coefficient from the old $\rho$-seed is not the same matrix.  Full rank of
$\mathcal F_h$, regularity of its signed Gram, decay of its positive
source-reproduction vector and its scalar-trace response have not been
proved here.  The numerator estimate \eqref{v2:drift-size} alone proves
none of them.  All nonlinear Hessian, differentiated-response, physical
inclusion and lapse--shift requirements in \cref{v1:nonlinear,v1:time,v1:flow-interface} remain in force, as does the common physical-lifespan
problem.

\paragraph{Verification scope.}
The companion program checks characteristic-zero lower-rank certificates for the three-site
diagonal and full first source, using exact modular elimination on a rationally unequal-amplitude seed, the Jacobi
Fourier representation including the folded edge, the bidiagonal inverse
and its central rank-one decomposition, the scalar soft-eigenvalue scaling,
and tensor-rank identities on finite grids.  The all-cutoff rank and
selection conclusions rely on the proofs above, not extrapolation from
those grids.  Symbolic checks verify the new mean calibration and drift
normalization.  The full second stationary source, a nonlinear exact-action
trajectory and a cofinal Cartan inverse are not reported as computed.
The predecessor body is preserved byte for byte, including its qualifications.

\begin{thebibliography}{9}
\raggedright
\setcounter{enumiv}{4}
\bibitem{v2:LST}
R. S. Leite, N. C. Saldanha and C. Tomei,
\emph{Reconstruction of tridiagonal matrices from spectral data},
arXiv:math/0508099 (2005).
\url{https://arxiv.org/abs/math/0508099}.
This is methodological context for Jacobi and bidiagonal coordinates,
not a source for the cosine-defect spectrum or the renewal source results.
\end{thebibliography}
% END IMMUTABLE EXACT-ACTION BRIDGE V2 BODY

% Future iterations append here.
% BEGIN IMMUTABLE EXACT-ACTION BRIDGE V3 BODY
\clearpage
\renewcommand{\thesection}{\arabic{section}}
\renewcommand{\theHsection}{bridgeV3.\arabic{section}}
\setcounter{section}{18}
\section{The weak coefficient is a first-jet problem}
\label{v3:context}

The leading active block on the detuned axial family is controlled by
\cref{v2:active-theorem}.  The next object is the actual weak coefficient
\(J_{2,h}|_{\ker\curl_\delta}\), not another estimate conditional on its
invertibility.  A direct use of the general second-source formula appears to
require the second prepared initial coefficient and the second stationary
connection.  We first examine whether that dependency survives the specified
weak tests.  It does not: a flat first-source identity annihilates every
curl-free shift for \emph{every symmetric first field variation}.  This
removes the second preparation from the complete weak coefficient, after
all its terms have been assembled.  A second cancellation removes the
undetermined tangent part of the second stationary connection.

The resulting coefficient construction uses only the first field jet, two
fixed pointwise linear maps, and the cubic electric and magnetic logarithmic
coefficients of the unchanged action.  It is valid at every odd cutoff and
for all weak test frequencies.  On the rational three-site seed already
used in \cref{v2:weak-only}, an exact finite-field computation then certifies
the full \(107\)-column graded signed Gram and its remaining \(29\)-column
weak Schur complement.  This supplies a new regular finite branch, not an
all-cutoff weak-rank theorem.

There is also a limitation on this particular upstream modification.  At
the same three-site cutoff, the inherited null/source obstruction at equal
amplitudes implies that the positive source-reproduction norm blows up at
least as \(|\theta-1|^{-1}\) along regular detunings.  Rank restoration is
therefore not a uniform response estimate even in this one-parameter family.
The parameter limit at fixed cutoff is kept distinct from refinement.

\paragraph{Inherited interfaces and nonduplication.}
We use the analytic stationary saddle and tangent inclusion of
\cite{OB25}, \src{obv19:split-saddle}--\src{obv19:KJ}; its general
second-source identity \src{obv20:J2-recipe}; and the exact symmetric-seed
null/source witness \src{obv20:quartic-obstruction}.  The latter is an
explicit inherited finite certificate.  The first-jet formulae, phase
calculus, initial preparation and active elimination in the preceding
bodies are not re-proved.  The new results concern the \emph{restricted}
second source, its exact evaluation, its detuned finite rank, and the response
as this detuning is removed.  Standard nullspace and saddle-point elimination
are background methods \cite{v3:BGL}; they are not claimed as new general
linear algebra.  No result below changes the original action, stationary
connection equation, physical norm, or canonical time.

\section{Universal flat weak-shift cancellation}
\label{v3:flat}

Fix an odd grid and the reference multiplier \(\lambda_*=(1,0)\).
Write \(X=(u,\pi)\), \(\gamma=I+u\).  In the inherited saddle use spatial
connections \(z=(A_i)_{i=1}^3\) and auxiliary variables
\(\eta=(A_0,w)\):
\begin{equation}
 \Phi_h(z,\eta;X,\lambda)
   =\ip{\eta}{\mathcal T_X(z)}_h+\mathcal V_{X,\lambda}(z).
 \label{v3:split}
\end{equation}
Here \(\mathcal T_X\) is independent of \(\lambda\) and
\(\mathcal V_{X,\lambda}\) is affine in \(\lambda\).  At its actual
stationary point set
\[
 T=D_z\mathcal T_X,\qquad
 G=D_z\nabla_\lambda\mathcal V_{X,\lambda},\qquad
 R_0=T_0^*(T_0T_0^*)^{-1}.
\]
The inclusion \(Z_0\) maps symmetric spatial matrices into pure rotational
spatial connections.  It parametrizes \(\ker T_0\).  Thus
\begin{equation}
 Z=Z_0-R_0(TR_0)^{-1}TZ_0,\qquad J_h(X)=Z^*G,
 \qquad J_h(0)=0.
 \label{v3:J}
\end{equation}
All adjoints here use the fixed positive coefficient pairings.  A different
fixed basis gives the corresponding congruent formulae.

For \(Y=(U,p)\) with \(U,p\) symmetric, the flat spatial stationary
connection differential is the inherited map
\begin{equation}
 (z_1)_i^{R_a}=-\frac12\epsilon_{abc}\delta_bU_{ic},\qquad
 (z_1)_i^{B_a}=p_{ia}-\frac12(\tr p)\delta_{ia}.
 \label{v3:z1}
\end{equation}
No divergence, trace, initial-constraint or scalar-momentum condition is
imposed on \(Y\) in this section.  Define
\(\mathcal W_h=\{\beta:\curl_\delta\beta=0\}\), including the three
constant shifts.  For \(\nu=(0,\beta)\) with \(\beta\in\mathcal W_h\),
the flat multiplier response has zero spatial and temporal connection
components and auxiliary velocity
\begin{equation}
 \eta_0(\beta)=(0,w_\beta),\qquad
 (w_\beta)_{ij}=\delta_i\beta_j+\delta_j\beta_i,
 \qquad G_0\nu+T_0^*\eta_0=0.
 \label{v3:eta0}
\end{equation}
These are identities of the original flat Cartan--Legendre saddle, including
the constant shift tests, not continuum gauge identifications.

\begin{lemma}[The first source annihilates all weak shifts for every first field jet]
\label{v3:universal-flat}
For every symmetric field variation \(Y=(U,p)\) and
\(\beta\in\mathcal W_h\),
\begin{equation}
 \boxed{D_XJ_h(0)[Y](0,\beta)=0.}
 \label{v3:flat-zero}
\end{equation}
The equality is exact at each cutoff and retains all spatial shifts.
\end{lemma}
\begin{proof}
Expand \(T=T_0+aT_1+O(a^2)\) and \(G=G_0+aG_1+O(a^2)\) on
\(X=aY\).  Equations \eqref{v3:J}--\eqref{v3:eta0} give
\[
 D_XJ_h(0)[Y]\nu=Z_0^*(G_1\nu+T_1^*\eta_0).
\]
We evaluate precisely this complete sum.  The first spatial load for a
shift, including derivatives of the products of \(\beta\) and coframe
coefficients, is boost-valued.  The velocity-load differential paired with
\(w_\beta\) is likewise boost-valued.  Their projections to symmetric
rotations vanish; no product rule is used to split their derivatives.
The first coframe variation of the shift Cartan coefficient has no
zeroth-order spatial connection on which to act.  The remaining flat
Cartan term pairs the shift's boost coefficient with the first spatial
connection and a symmetric rotational test.

Let \(S=S^T\) be this test and
\(B=p-\tfrac12(\tr p)I=B^T\).  Write their rows as \(S_i,B_i\).
The rotational part of \eqref{v3:z1} gives a rotation--rotation commutator,
annihilated by the shift's boost coefficient.  The possible boost part is,
up to the common normalization,
\begin{equation}
 \sum_{i<j}(\beta_j e_i-\beta_i e_j)\cdot
                    (-S_i\times B_j+S_j\times B_i)=0.
 \label{v3:Cartan-cancel}
\end{equation}
To verify the zero, rewrite the sum as
\(\sum_{i,j}\beta_j e_i\cdot(-S_i\times B_j+S_j\times B_i)\).
The first contraction vanishes by symmetry of \(S_{ia}\) and
antisymmetry of \(\epsilon_{iab}\); the second vanishes by symmetry of
\(B_{ib}\) and antisymmetry in \(i,b\).  This is a pointwise algebraic
cancellation for arbitrary symmetric \(B\), not only for scalar momentum.

The cubic electric remainder can enter \(T_1\), but only in the Gauss
rows dual to \(A_0\); the \(A_0\) component of \(\eta_0\) is zero.
The remainder has no explicit shift coefficient in \(G_1\).  A magnetic
cubic term has, after one spatial connection derivative, at least two
field factors, and cannot enter \(G_1\).  Higher electric and magnetic
terms have still higher degree.  These observations exhaust the terms
able to enter the displayed first coefficient, proving \eqref{v3:flat-zero}.
\end{proof}

\begin{theorem}[Independence of the second preparation on the weak block]
\label{v3:preparation-independent}
Let \(X(a)=aY+a^2Y_2+O(a^3)\) be any analytic field curve in the
original stationary chart.  It need not solve the initial constraints.
For every \(\beta\in\mathcal W_h\),
\begin{equation}
 \boxed{J_h(X(a))(0,\beta)=a^2\mathcal Q_h(Y;\beta)+O(a^3),
 \quad
 \mathcal Q_h(Y;\beta)=\frac12D_X^2J_h(0)[Y,Y](0,\beta).}
 \label{v3:weak-germ}
\end{equation}
In particular the complete \(J_2|_{\mathcal W_h}\) is independent of
\(Y_2\).  Changing a second free initial record while preserving the
first field jet cannot change this weak coefficient.
\end{theorem}
\begin{proof}
The inherited stationary inverse and tangent inclusion make \(J_h(X)\)
analytic at the fixed cutoff.  Its Taylor expansion on the stated curve
has second coefficient
\[
 D_XJ_h(0)[Y_2]\nu+\frac12D_X^2J_h(0)[Y,Y]\nu.
\]
The first coefficient and the first summand vanish by
\cref{v3:universal-flat}.  This proves the formula and the independence.
\end{proof}

\begin{remark}[The scope of the cancellation]
The general second-source formula in \cite{OB25} correctly retains
\(T_2\), \(G_2\) and the second preparation.  The theorem concerns their
\emph{assembled restriction to curl-free shifts}.  It does not say that
\(T_2\) itself is independent of preparation, or that one can delete its
term before doing the restriction.  Nor does it identify the weak
coefficients of two different \emph{first} jets.  The symmetric-gradient
seed and the detuned axial seed still have different weak sources.
\end{remark}

\section{Evaluation without a second stationary tangent solve}
\label{v3:compiler-section}

For \(\nu=(0,\beta)\), write
\(\mathcal V_{X,\nu}=\partial_\lambda\mathcal V_{X,\lambda}[\nu]\)
and define the scalar function
\begin{equation}
 \ell_\nu(X,z)=\mathcal V_{X,\nu}(z)
                   +\ip{\eta_0(\beta)}{\mathcal T_X(z)}_h.
 \label{v3:ell}
\end{equation}
For a power series, \([a^j]\) denotes its actual coefficient, without an
extra factorial.  On the straight field curve \(X=aY\), form
\begin{align}
 \tau_2&=[a^2]\mathcal T_{aY}(az_1),&
 z_2^\perp&=-R_0\tau_2,\label{v3:normal-second}\\
 g_1&=[a]D_z\ell_\nu(aY,az_1),&
 \eta_1&=-R_0^*g_1.\label{v3:aux-first}
\end{align}
The maps \(R_0,T_0,Z_0\) are the original pointwise maps, independent of
\(Y\).  In particular these equations do not invert a small weak Gram.

\begin{theorem}[A complete first-jet coefficient construction]
\label{v3:compiler}
For arbitrary symmetric \(Y\) and \(\beta\in\mathcal W_h\), the weak
coefficient of \cref{v3:preparation-independent} is exactly
\begin{equation}
 \boxed{\mathcal Q_h(Y;\beta)=
 Z_0^*[a^2]D_z\!
 \left(\mathcal V_{aY,\nu}(z)
       +\ip{\eta_0+a\eta_1}{\mathcal T_{aY}(z)}_h\right)
       \Big|_{z=az_1+a^2z_2^\perp}.}
 \label{v3:compiler-formula}
\end{equation}
There is no second initial constraint solve and no second stationary
\emph{tangent} solve in this formula.  The normal correction
\(z_2^\perp\) and auxiliary correction \(\eta_1\) must both be retained.
\end{theorem}
\begin{proof}
The exact stationary spatial connection on \(X=aY\) is
\(z(a)=az_1+a^2z_2+O(a^3)\).  Its constraint equation gives
\(T_0z_2+\tau_2=0\); hence
\(z_2=z_2^\perp+Z_0S_2\) for some symmetric field \(S_2\).
The universal cancellation gives \(Z_0^*g_1=0\).  Since
\(\Ran Z_0=\ker T_0\), \(g_1\in\Ran T_0^*\), and therefore
\(g_1+T_0^*\eta_1=0\).

Expanding the exact formula for \(Z^*G\) to second order gives
\[
 J_2\nu=Z_0^*(G_2\nu+T_2^*\eta_0+T_1^*\eta_1).
\]
Indeed \(\eta_1=-R_0^*(G_1\nu+T_1^*\eta_0)\), so this is precisely
\src{obv20:J2-recipe}, with its four terms grouped, not omitted.
The displayed coefficient is the right side of \eqref{v3:compiler-formula} when
the true \(z_2\) is used.

Replacing \(z_2\) by \(z_2^\perp\) changes that right side by
\[
 Z_0^*D_z^2\ell_\nu(0,0)Z_0S_2.
\]
This expression is zero.  At the flat point, the constraints paired with
\(w_\beta\) are linear in \(z\).  The remaining Hessian is the magnetic
shift Cartan coefficient.  Its two symmetric tangent variations are both
rotations, whose commutator is rotational and is annihilated by the boost
coefficient of a pure shift.  The electric Hessian is paired with the zero
\(A_0\) slot; logarithmic cubic terms have zero second spatial derivative
at \(z=0\).  Thus the omitted tangent correction contributes zero, as
claimed.  No assertion is made that the actual \(S_2\) is zero.
\end{proof}

\begin{proposition}[The quartic electric term cannot enter this restricted coefficient]
\label{v3:log-degree}
In \eqref{v3:compiler-formula}, the geometric coframe functions must be differentiated
to their required orders, and both the cubic electric and cubic magnetic
remainders must be retained.  The quartic electric remainder and all higher
logarithmic orders contribute zero to \(\mathcal Q_h(Y;\beta)\).
\end{proposition}
\begin{proof}
The quartic electric term is cubic in spatial connections and linear in
\(A_0\).  In a general second-source coefficient its spatial derivative
has degree two and can pair with a nonzero zeroth-order \(A_0\) multiplier
response.  For \(\beta\in\mathcal W_h\), that response is zero by
\eqref{v3:eta0}.  The first nonzero auxiliary \(A_0\) in
\eqref{v3:compiler-formula} has degree one, so this contribution has degree at
least three.  The normal constraint \(\tau_2\) also receives no quartic
electric term, since its \(A_0\) derivative has spatial degree three.
The cubic electric term instead has a spatial derivative of degree one;
paired with the first auxiliary response it can contribute at degree two.
It also contributes to \(\tau_2\).  It is not discarded.

A magnetic cubic term has a spatial derivative of degree two and a shift
coefficient of degree zero, and hence can contribute.  Magnetic terms of
degree at least four have spatial derivative of degree at least three.
All higher electric terms have greater spatial degree than the quartic
term just considered.  Coframe derivatives can lower the field degree
when paired with the auxiliary velocity and must still be retained.  These
are coefficient-level statements; the nonlinear remainder of the original
flow remains unchanged.
\end{proof}

\section{Polarization and finite Fourier bandwidth of the weak coefficient}
\label{v3:bands}

\begin{theorem}[Ten coefficient operators suffice for the axial first-jet family]
\label{v3:ten-banks}
Let \(Y_i=(\mathsf T_i\cos(2\pi x_i),0)\), \(1\le i\le3\), and
\(Y_4=(0,-I)\).  For \(t=(\alpha_1,\alpha_2,\alpha_3,s)\), set
\(Y(t)=\sum_{i=1}^4t_iY_i\).  Then, as an operator on the entire weak
space,
\begin{equation}
 \mathcal Q_h(Y(t);\cdot)=
 \sum_{i=1}^4t_i^2\mathcal Q_h(Y_i;\cdot)
 +\sum_{i<j}t_it_j
   \{\mathcal Q_h(Y_i+Y_j;\cdot)-\mathcal Q_h(Y_i;\cdot)
                                      -\mathcal Q_h(Y_j;\cdot)\}.
 \label{v3:polarization}
\end{equation}
These ten operator coefficients determine the actual weak source after
putting \(s=s_{h,\alpha}\).  Their computation does not assume that the
individual polarization directions satisfy nonlinear initial constraints.
\end{theorem}
\begin{proof}
Equation \eqref{v3:weak-germ} is a homogeneous quadratic function of
\(Y\), with values in finite linear operators on weak shifts.  Polarizing
that quadratic function gives \eqref{v3:polarization}.  The stationary
chart is defined for all sufficiently small fields, not just the initial
constraint zero set.  Its derivative at zero is therefore evaluable in
all the displayed directions.  Exact preparation is reintroduced only
when this coefficient is applied to an actual initial family.
\end{proof}

\begin{proposition}[Twenty-five possible frequency offsets, with aliases retained]
\label{v3:bandwidth}
For the axial first jet, the full unprojected weak source maps a Fourier
test at wave vector \(k\) into wave vectors
\begin{equation}
 \boxed{k+\ell\pmod N,\qquad
           \ell\in\mathbb Z^3,\quad |\ell|_1\le2.}
 \label{v3:offsets}
\end{equation}
There are \(25\) possible offsets before cyclic coincidences.  The claim
holds for all test frequencies, including the folded ones.  It is a
bandwidth statement for \(J_2|_{\mathcal W_h}\), not for its inverse or
for its projection by \(I-\Pi_{\mathcal A_h}\).
\end{proposition}
\begin{proof}
The first records have support in \(\{0,\pm e_1,\pm e_2,\pm e_3\}\).
The first stationary connection preserves these supports.  In
\eqref{v3:normal-second}, each coefficient is quadratic in these records;
phase derivatives and literal shifts preserve its Fourier support, and
\(R_0\) is pointwise constant.  In \eqref{v3:aux-first}, the first source
has one first-field factor and one test factor, with the same support
property.  Every term at degree two in \eqref{v3:compiler-formula} consequently
has one test and exactly two first-field factors, with the convention that
a degree-two normal factor already contains those two fields.  Thus its
output shift is a sum of two elements of the first-record support, giving
\(|\ell|_1\le2\).  This set has one zero, six unit, six doubled axial,
and twelve two-axis offsets.  Multiplication is cyclic throughout, so the
conclusion includes aliasing.  The active projector includes an inverse
and is not asserted to preserve this bandwidth.
\end{proof}

These results reduce the amount of coefficient calculation substantially,
but are not spectral estimates for the weak Schur complement.  In
particular finite frequency support of the \emph{drift} never justified
truncating the weak \emph{test space}; the latter is retained here.

\section{An exact detuned three-site signed-rank certificate}
\label{v3:certificate-section}

We now apply the coefficient construction, rather than only state it.
At \(N=3\), let
\begin{equation}
 \alpha=(1,9/8,1),\qquad s=\frac{3\sqrt{418}}{16},\qquad
 s^2=\frac{1881}{128}.
 \label{v3:finite-seed}
\end{equation}
This is exactly \eqref{v2:seed}, since \(\omega_h^2=27\).
Every phase derivative is the rational matrix \(3(S_i-S_i^{-1})\), and
the sampled cosine is \((1,-1/2,-1/2)\).  All coefficient operations
therefore take place in \(\mathbb Q(\sqrt{418})\).  The sign of the real
square root in \eqref{v3:finite-seed} fixes the scalar momentum branch.

Use lexicographic site order and raw symmetric coordinates
\((11,22,33,12,13,23)\).  Their Frobenius matrix is
\(\diag(1,1,1,2,2,2)\); an off-diagonal scalar variation therefore gives
a doubled covector component.  The inverse of the flat \emph{bilinear}
Cartan matrix in these raw covector coordinates is
\begin{equation}
 \mathsf K_{\rm raw}^{-1}
 =\frac14\diag(2I_3-\mathbf1\mathbf1^T,I_3).
 \label{v3:raw-K}
\end{equation}
This is the same physical Cartan form as \(K_0\), not a different trace
normalization.  A common spatial-volume factor can be omitted consistently
from the finite rank calculation.

Let \(J^{\rm all}_1\) be the \(162\times108\) first-source matrix with
columns ordered as all lapse point indicators followed by the three blocks
of shift point indicators.  Select the following zero-based column indices:
\begin{equation}
\begin{split}
 \mathcal P={}&\{0,\ldots,25\}\cup\{27,\ldots,34\}
 \cup\{36,\ldots,43\}\cup\{45,\ldots,52\}\\
 &\cup\{54,55,56,57,58,60,61,63,64,65,66,67,69,70,72,73,75,76,78,79\}\\
 &\cup\{81,84,87,90,93,96,99,102\}.
\end{split}
 \label{v3:pivots}
\end{equation}
There are \(78\) columns; denote their matrix by \(C\).
For the weak space use the three constant shifts, followed by
\(\beta=\delta(\mathbf1_{\{x\}}-\mathbf1_{\{x_{\rm last}\}})\) at the
other \(26\) sites.  These form a basis because the phase gradient is
injective on mean-zero scalars.  Let \(W\) be the \(162\times29\)
weak-source matrix from \eqref{v3:compiler-formula}, and put
\begin{equation}
 \begin{gathered}
 F=(C,W),\qquad K^{-1}=I_{27}\otimes\mathsf K_{\rm raw}^{-1},\\
 A=C^TK^{-1}C,\qquad L=C^TK^{-1}W,\\
 D=W^TK^{-1}W-L^TA^{-1}L.
 \end{gathered}
 \label{v3:finite-matrices}
\end{equation}

\begin{theorem}[The detuned seed has a full regular leading signed source]
\label{v3:finite-rank}
For the unchanged action on \eqref{v3:finite-seed},
\begin{equation}
 \rank C=78,\qquad \rank F=107,\qquad
 \det(F^TK^{-1}F)\ne0,\qquad \det D\ne0.
 \label{v3:rank-conclusion}
\end{equation}
In particular the remaining weak signed block has rank \(29\).
These are characteristic-zero assertions at this fixed cutoff.
\end{theorem}
\begin{proof}
We give the exact arithmetic certificate and its reconstruction rule.  Map
\(\sqrt{418}\) to \(30\) modulo \(241\); this is valid because
\(30^2\equiv418\equiv177\pmod{241}\).  Then
\(s\mapsto96\).  The coefficients of the construction have denominators
prime to \(241\).  In particular the fixed maps \(T_0,R_0,Z_0\), coframe
Taylor coefficients, literal \(h=1/3\) factors and Baker--Campbell--Hausdorff
coefficients are all defined in this residue field.

The companion coefficient program evaluates the finite differentiations
in \eqref{v3:normal-second}--\eqref{v3:compiler-formula} by exact polynomial
arithmetic, retaining amplitude degree at most two and first degree in
an independent variation variable.  No difference quotient or numerical
rank tolerance is used.  The explicit scalar functional and order
conventions used to regenerate its columns are given in
\cref{v3:implementation}.  Gaussian elimination on the resulting arrays
gives the following determinant residues:
\begin{equation}
 \boxed{\det A\equiv30,\qquad\det D\equiv227,\qquad
          \det(F^TK^{-1}F)\equiv62\pmod{241}.}
 \label{v3:det-residues}
\end{equation}
The Schur identity is checked separately:
\(30\cdot227\equiv62\pmod{241}\).  The full arrays and fixed column
list are included in the certificate, so these finite identities can be
verified without the source-search history.  Nonzero determinant images
imply nonzero determinants over \(\mathbb Q(\sqrt{418})\), and hence
under its specified real embedding.  A nonsingular signed Gram also
implies linear independence of its source columns.  This proves all
assertions.  The inherited all-cutoff active theorem is not used to
extrapolate the weak conclusion beyond this cutoff.
\end{proof}

\begin{corollary}[Regular detunings occur arbitrarily close to equal amplitudes at this cutoff]
\label{v3:fixed-parameter}
At \(N=3\), on a sufficiently small punctured interval about \(\theta=1\),
the first jets \(\alpha=(1,\theta,1)\),
\(s=\omega_h\sqrt{(2+\theta^2)/6}\), have a regular complete leading
graded source.  In particular arbitrarily small positive rational
detunings occur inside the admissible initial-parameter chart.  For each
such fixed parameter, exact initial preparation and the inherited
fixed-cutoff continuation theorem supply a local original-action branch
at sufficiently small nonzero field amplitude.
\end{corollary}
\begin{proof}
Use the fixed selected columns \eqref{v3:pivots} and the fixed weak basis.
Their coefficients are real analytic functions of \(\theta>0\), by
\eqref{v3:compiler-formula} and the positive analytic function \(s(\theta)\).
The complete signed determinant is nonzero at \(9/8\), so it is not
identically zero on this connected interval.  Its zero at \(1\), if
present, is isolated: otherwise the Taylor series there would have every
coefficient zero and analytic continuation would contradict the certified
nonzero value.  The same reasoning, or \cref{v2:tensor-rank} at \(N=3\),
ensures the active rank on a smaller punctured interval.  This proves the
parameter statement, without asserting a uniform inverse as \(\theta\to1\).

The initializer of \cref{v2:preparation} applies when \(\theta\) is
sufficiently close to one.  At fixed cutoff and parameter, the graded
finite-amplitude matrix is an analytic perturbation of its now nonsingular
leading Gram.  It is nonsingular for sufficiently small nonzero amplitude.
The original mean-lapse homogeneity then leaves its declared one-null
multiplier direction.  The continuation theorem in \cite{OB25},
\src{obv18:conditional-continuation}, applies on this local chart.  Its
existence interval, inverse bound and amplitude radius may depend on
\(\theta\) and on the cutoff.  No uniform lifespan follows from this
argument.
\end{proof}

\section{A positive-source pole when the detuning is removed}
\label{v3:detuning-pole-section}

The new finite-rank theorem must be interpreted together with the original
source, not merely as a successful determinant.  At \(N=3\) the
unperturbed equal-amplitude seed has an inherited weak test
\(\beta_*\in\mathcal W_h\) and an active multiplier \(a_*\) such that
\begin{equation}
 J_{1,1}a_*+\mathcal Q_h(Y(1);\beta_*)=0,
 \qquad \ip{\beta_*}{d_{h,1}}_h\ne0.
 \label{v3:inherited-null}
\end{equation}
The subscript \(1\) here denotes \(\theta=1\), not a new cutoff.
One may take the gradient of the point-pair scalar in
\src{obv20:null-witness}; the nonzero pairing and the exact active
companion are part of \src{obv20:quartic-obstruction}.  Its stated
\(\epsilon\)-normalization gives \(-81/16\); changing to the present
\(a\)-normalization rescales it by a fixed nonzero factor.  Only its
nonvanishing is used below.  The active companion can be chosen in the
mean-zero lapse and transverse-shift slice, since constants and curl-free
shifts are annihilated by the first source.

\begin{theorem}[Actual detuning causes a positive source-reproduction pole]
\label{v3:detuning-pole}
At \(N=3\), let \(v_\theta\) satisfy all the leading graded source
equations on the detuned family:
\[
 \mathcal A_\theta^*v_\theta=0,\qquad
 \mathcal Q_h(Y(\theta);\cdot)^*v_\theta=d_{h,\alpha(\theta)}
                  \quad\hbox{on }\mathcal W_h.
\]
There are constants \(c,\eta>0\), for this fixed cutoff, such that every
such solution with \(0<|\theta-1|<\eta\) obeys
\begin{equation}
 \boxed{\|v_\theta\|_h\ge\frac{c}{|\theta-1|}.}
 \label{v3:pole}
\end{equation}
This includes the positive minimum-norm reproduction vector and, on the
regular branches of \cref{v3:fixed-parameter}, the original signed Cartan
response.  The bound does not rely on a trace-angle resonance.
\end{theorem}
\begin{proof}
Keep the actual \(a_*,\beta_*\) of \eqref{v3:inherited-null} fixed and
set
\[
 r_\theta=\mathcal A_\theta a_*+
                            \mathcal Q_h(Y(\theta);\beta_*).
\]
Here \(\mathcal A_\theta a_*\) denotes the first source evaluated on its
active representative.  At \(\theta=1\), \(r_1=0\).  The coefficient
construction and \(s(\theta)\) are analytic, so on a fixed compact
neighbourhood \(\|r_\theta\|_h\le C|\theta-1|\).
The drift is the actual \eqref{v2:drift-vector}; continuity and
\eqref{v3:inherited-null} give
\(|\ip{\beta_*}{d_{h,\alpha(\theta)}}_h|\ge c_0>0\) on a smaller
neighbourhood.  Pair the two original source equations with the fixed
multiplier pair.  They yield exactly
\[
 \ip{r_\theta}{v_\theta}_h
                 =\ip{\beta_*}{d_{h,\alpha(\theta)}}_h.
\]
Cauchy--Schwarz gives \eqref{v3:pole}.  The argument applies to every
reproducing vector, so minimizing the positive norm cannot remove it.
The signed response is one such vector; no information about its signed
Gram eigenvalues was used.
\end{proof}

This is a new \emph{detuning}-parameter conclusion obtained from the
retained exact null/source witness; it is not a reassertion of that source's
field-amplitude multiplier-rate pole.  It is also not a cofinal divergence
theorem.  The displayed constants are for \(N=3\), whereas the sequence
\(\theta_h\to1\) in V2 changes \(N\) at the same time.  A cofinal use
requires a corresponding source/witness estimate with its actual
cutoff dependence.  This distinction prevents the finite rank restoration
from being promoted to a small-response conclusion.

\begin{corollary}[Unbounded initial electric read along exact prepared records at fixed cutoff]
\label{v3:physical-pole}
At \(N=3\), there are regular detunings \(\theta_j\to1\) and positive
field amplitudes \(a_j\to0\) such that the exact prepared records
\(X_{h,\alpha(\theta_j)}(a_j)\) tend to the flat canonical record, while
the original electric link read, evaluated on their own normalized exact
canonical continuations, satisfies
\[
 \|E_{h,\theta_j,a_j}(0)\|_h\longrightarrow\infty.
\]
Each continuation exists locally at its own parameter and amplitude.  No
common physical interval, and no corresponding cofinal divergence, is
asserted.
\end{corollary}
\begin{proof}
Fix a regular \(\theta\ne1\) in \cref{v3:fixed-parameter}.  In the
amplitude grading, the original source has columns of orders \(a\) on
the active space and \(a^2\) on the weak space.  Its normalized Gram
converges to the invertible leading Gram.  The exact stationary-response
identity used in \cite{OB25}, \src{obv24:velocity-exact}, therefore gives
at the initial unit-lapse, zero-shift cut
\[
 \dot z=Z_0v_\theta+O_\theta(a).
\]
Here the partial field term \(z_X F_h^{\rm can}\) is \(O(a)\); all
remaining factors are analytic after the fixed amplitude grading, with
constants allowed to depend on \(\theta\).  This argument uses the
original multiplier rate and does not prescribe the connection velocity.

The original read is
\[
 E_i=h^{-1}\{\dot U_iU_i^{-1}+A_0-U_iA_{0,+i}U_i^{-1}\},
 \qquad U_i=\exp(hA_i).
\]
At fixed \(h,\theta\), the initial stationary connections are \(O(a)\).
The differential of the exponential and the displayed formula give
\(E(0)=Z_0v_\theta+O_\theta(a)\).  The flat inclusion is isometric in
the physical coordinates.  Write the last error bound as
\(C_\theta a\) on its actual amplitude domain.

Choose regular \(\theta_j\to1\).  For each \(j\), choose \(a_j\)
inside that certified domain and the exact continuation domain, with
\(a_j<1/j\) and
\(C_{\theta_j}a_j\le c/(2|\theta_j-1|)\).
The reverse triangle inequality and \cref{v3:detuning-pole} give
\(\|E(0)\|_h\ge c/(2|\theta_j-1|)\to\infty\).
The uniform initial preparation estimate gives \(X_j\to0\).
This is a diagonal choice of actual initial histories at a fixed grid,
not an exchange of an uncontrolled pair of limits.  It proves no statement
about the size of their common-time domain.
\end{proof}

The unbounded read is a property of the \emph{exact-flow completion} of
canonical initial records; their full metric time jets need not remain
bounded.  It therefore does not contradict continuity of a geometric read
on a bounded complete metric-jet chart.  It does rule out a bounded extension
of this exact-flow electric read over the flat tip of this fixed-cutoff
prepared family.

\section{Reproducible coefficient functional and verification}
\label{v3:implementation}

The verification package contains the coefficient arithmetic, complete
residue matrices, determinant routine and preservation manifest.  The
following formulas fix the implementation independently of code-variable
names.  Write a Lorentz coefficient as \((r,b)\) and use
\begin{equation}
 [(r,b),(s,c)]=(-r\times s+b\times c,
                              -r\times c-b\times s).
 \label{v3:Lorentz}
\end{equation}
The invariant-action coefficients below are written as \emph{dual}
coefficient vectors, so the dot product with a bracket is the scalar
pairing in these coordinates.  Set
\[
 e=(I+u)^{1/2},\qquad \Pi_i=2\operatorname{cof}(e)_{i\cdot},\qquad
 \Sigma_{ij}=\bigl(2n\epsilon_{ij\ell}e_{\ell\cdot},
                                      \beta_j\Pi_i-\beta_i\Pi_j\bigr).
\]
The degree-extraction functional is the original geometric action and its
relevant logarithmic coefficients, minus the canonical momentum pairing:
\begin{align}
 \Phi={}&-\sum_i D\Pi_i[u][w]\cdot z_i^B-\pi:w
       +\sum_i(\delta_i\Pi_i)\cdot A_0^B
       +\sum_i\Pi_i\cdot[A_0,z_i]^B\nonumber\\
 &+\sum_{i<j}\Sigma_{ij}\cdot
             (\delta_i z_j-\delta_j z_i+[z_i,z_j])\nonumber\\
 &-\frac h2\sum_i\Pi_i\cdot[z_i,[z_i,S_iA_0]]^B
       +h\sum_{i<j}\Sigma_{ij}\cdot\mathcal L_3
                         (z_i,S_i z_j,-S_j z_i,-z_j).
 \label{v3:functional}
\end{align}
Spatial summation is understood.  Its derivative in \((A_0,w)\) is
\(\mathcal T_X\).  To compute \(\mathcal L_3\), start at
\(l_1=l_2=l_3=0\) and append each of the four factors \(x\) in the
displayed order, using the old values on the right:
\[
 l_3^{\rm new}=l_3+\tfrac12[l_2,x]
             +\tfrac1{12}([l_1,[l_1,x]]+[x,[x,l_1]]),\qquad
 l_2^{\rm new}=l_2+\tfrac12[l_1,x],\qquad l_1^{\rm new}=l_1+x.
\]
The coframe velocity derivative is important.  For extraction through
amplitude degree two, it is computed from
\begin{equation}
 \begin{split}
 e&=I+\tfrac12u-\tfrac18u^2+O(u^3),\\
 D e[u][w]&=\tfrac12w-\tfrac18(uw+wu)
           +\tfrac1{16}(u^2w+uwu+wu^2)+O(u^3w),
 \end{split}
 \label{v3:sqrt-jets}
\end{equation}
with \(D\Pi\) obtained by differentiating the cofactor.  Omitting the
last displayed velocity terms would omit a degree-two coefficient.
Equation \eqref{v3:functional} is used only to extract the coefficient
whose exactness was proved in \cref{v3:log-degree}, not to replace the
nonlinear action by a polynomial evolution.

All spatial tangent and auxiliary variations are taken before evaluating
their coefficients.  The program represents polynomials in the field
amplitude and an independent variation variable, retaining degrees
\((a^j,t^k)\) with \(j\le2\), \(k\le1\).  The coefficient of
\(a^jt\) is then the exact indicated first variation.  Products, literal
shifts, phase derivatives, cross products and matrix products are evaluated
in the same finite ring.  At the certified cutoff, \(\delta_i=3(S_i-S_i^{-1})\)
throughout, including its action on products; no continuum Leibniz rule is
substituted.

\paragraph{Tests actually performed.}
The finite checks verify the universal flat cancellation on arbitrary
symmetric first-field arrays, the equality after restoring the quartic
electric term, and invariance under an arbitrary added second tangent
connection.  They also verify the complete first-source formula against
an independently assembled coefficient differentiation for lapse and shift
point tests.  All \(29\) weak first coefficients vanish in the exact
residue computation.  The detuned matrices give
\eqref{v3:det-residues}.  A separate equal-amplitude residue calculation
recovers active rank \(77\), weak projected rank \(21\), and augmented
weak rank \(22\), consistent with the inherited source obstruction.
Those equal-amplitude ranks are cross-checks, not new characteristic-zero
rank proofs.  The symbolic symmetric-matrix contraction in
\eqref{v3:Cartan-cancel} is checked independently.

All modular arrays are reduced between operations; the checked integer
products fit within signed \(64\)-bit arithmetic.  The coefficient
calculation uses no finite-amplitude floating difference and the determinant
certificate uses no tolerance.  These checks are not proof-assistant
formalization, a nonlinear spacetime integration, or an all-cutoff
weak-response estimate.

\section{The remaining cofinal problem after the first-jet reduction}
\label{v3:status}

The new restricted identity removes an unnecessary second-preparation
calculation from the leading weak source.  The original general preparation
must still be supplied for exact finite initial histories and for nonlinear
remainders, but it need not be repeatedly differentiated to obtain
\(J_2|_{\mathcal W_h}\).  Likewise, the uncomputed second stationary
tangent can no longer be invoked as an extra unknown in this particular
coefficient.  The construction \eqref{v3:compiler-formula} retains the normal and
auxiliary responses that do matter.

The next cofinal calculation is therefore explicit:
\begin{equation}
 \mathcal F_h=(I-\Pi_{\mathcal A_h})\mathcal Q_h(Y_h;\cdot),
 \qquad
 b_{W,h}=d_{h,\alpha_h}.
 \label{v3:remaining}
\end{equation}
The first operator can be assembled using the ten polarization operators
and the \(25\)-offset coefficient support.  One must still prove full rank
of this \(V+2\)-column weak source, regularity of its signed Gram, and
source-resolved bounds for the positive reproduction and scalar-trace
response.  The single new rank certificate establishes none of these
cofinal bounds.  Its nearby detuning pole shows why those bounds must
retain the actual source, even before the signed trace amplification.

For the physical flow, the actual nonlinear Hessian and inclusion, the
transported derivative sources, the lapse--shift comparison and the
cutoff-independent lifespan requirements of V1 remain in force.  A future
refinement estimate cannot use the positive minimum-norm connection in
place of the original signed response, nor infer time-regularity from a
static coefficient bound.  No direct exact-action Einstein continuum
limit is claimed in this body.  The progress is an exact removal of two
coefficient dependencies, a computable action-native weak operator, and a
new regular detuned finite branch together with its quantitative
fixed-cutoff limitation.

\begin{thebibliography}{9}
\raggedright
\setcounter{enumiv}{5}
\bibitem{v3:BGL}
M. Benzi, G. H. Golub and J. Liesen,
\emph{Numerical solution of saddle point problems},
Acta Numerica \textbf{14} (2005), 1--137.
\href{https://doi.org/10.1017/S0962492904000212}{doi:10.1017/S0962492904000212}.
This is background for nullspace and saddle-point elimination, not a source
for the weak renewal coefficients or the detuned finite certificate.
\end{thebibliography}
% END IMMUTABLE EXACT-ACTION BRIDGE V3 BODY

% Future iterations append here.
% BEGIN IMMUTABLE EXACT-ACTION BRIDGE V4 BODY
\clearpage
\renewcommand{\thesection}{\arabic{section}}
\renewcommand{\theHsection}{bridgeV4.\arabic{section}}
\setcounter{section}{26}
\section{An upstream change from rank restoration to source cancellation}
\label{v4:context}

The detuning calculation in \cref{v3:detuning-pole} identifies a genuine
cost: at the certified three-site cutoff a nonzero drift meets an emerging
source cokernel, and even its positive minimum reproduction norm diverges.
It would be circular to replace this fact by another hypothesis on the
same inverse.  We therefore change the \emph{first initial field jet}, while
retaining the original action and every one of its field and multiplier
variations.  The new freedom is the second transverse tensor polarization
on each of the three axial configuration modes.

For arbitrary cosine and sine polarization tensors, the complete weak
quadratic drift is calculated below.  Its only contributions are doubled
axial modes; their coefficients are the unequal-norm and nonorthogonality
invariants of the two polarization tensors.  Equal norms and orthogonality
make this actual drift zero at every odd cutoff.  This is an exact
coefficient cancellation, including the three-site folded mode, not a small
source assumed after inversion.  The original nonlinear initial constraints
can be solved on these first jets by a recalibrated use of the inherited
range-and-mean preparation.  On one such balanced three-site seed, a new
exact signed-Gram certificate has full rank.  Thus the cancellation does not
merely put the source on an unverified singular branch.

The consequence is restricted but concrete.  On that regular finite branch
the leading normalized electric response is zero, and the literal initial
electric read is $O_h(a)$ as the field amplitude $a$ tends to zero.  The next
normalized source, the lapse--shift velocity and the physical lifespan are
not thereby controlled uniformly.  In particular this result is not an
exact-action Einstein continuum theorem.

\paragraph{Dependencies and nonduplication.}
The retained inputs are the original flat stationary map, the full
weak-shift quadratic action identity and trace test in \cite{OB25},
\src{obv23:weak-action}--\src{obv23:trace-test}; the stationary cubic
envelope; the nonconstant initial-constraint contraction in
\src{obv17:range}; and the exact normalized continuation on a verified
rank branch.  The weak-source compiler in \cref{v3:compiler} is reused, not
re-derived.  The new content is the polarization-complete drift, its exact
zero set, its original-action preparation, its regular finite source
certificate and its amplitude-response consequences.  Quadrature
polarizations themselves are standard linear algebra; polarization
conventions in gravitational-wave analysis are reviewed in \cite{v4:Isi}.
Here the quadrature is imposed on \emph{spatial initial profiles}, not
asserted to be a propagating circularly polarized wave.  The necessity of
checking nonlinear initial constraints rather than only their flat
linearization is also distinct from the new drift computation; see
\cite{v4:BD} for classical continuum context.

\section{The complete weak drift for arbitrary axial polarizations}
\label{v4:drift-section}

Keep $N=2m+1\ge3$, $h=N^{-1}$ and the phase derivative $\delta_{i,h}$ of
\cref{v1:phase}.  Put $c_i=\cos(2\pi x_i)$, $s_i=\sin(2\pi x_i)$ and
$\omega_h=2h^{-1}\sin(\pi h)$.  Thus $\delta_i c_i=-\omega_hs_i$ and
$\delta_i s_i=\omega_hc_i$ on every odd grid.  For each $i$ let
\[
 \mathcal V_i=\{A=A^T:\ A_{ij}=0\text{ for all }j,\ \tr A=0\}.
\]
This is the two-dimensional real space of transverse trace-free tensors
for that axis.  Take $A_i,B_i\in\mathcal V_i$ and set
\begin{equation}
 U=\sum_{i=1}^3U_i,\qquad U_i=A_ic_i+B_is_i,\qquad p=-sI.
 \label{v4:axial-jet}
\end{equation}
The scalar $s$ is spatially constant; its constraint calibration is supplied
later.  All products in this section remain the original cyclic products.
The first field jet is transverse and trace-free in its configuration slot.
Its flat unit-lapse Hamiltonian velocity is $F_{1,u}=sI$, and
$F_{1,p}=-\omega_h^2U/2$.

Let $b_{2,h}^{\beta}$ be the three momentum components of the actual
quadratic canonical drift at the unit-lapse, zero-shift cut.  Define
\[
 d_h=\mathsf P_h^{\rm cf}b_{2,h}^{\beta},\qquad
 \mathsf P_h^{\rm cf}=\text{positive orthogonal projection onto }
                         \ker\curl_\delta.
\]
This notation denotes the retained weak tests, not deletion of the other
canonical equations.  In particular $d_h=0$ will not mean that all active
quadratic drift components vanish.

\begin{lemma}[The full tested cubic envelope and its logarithmic cancellation]
\label{v4:full-drift}
For every curl-free shift $\beta$, put
$k_{ij}=\delta_i\beta_j+\delta_j\beta_i$.  On \eqref{v4:axial-jet},
\begin{equation}
 \ip\beta{d_h}_h
 =-D_U H_{3,h}^{\rm g}(U,p)[k]
                  -\frac{\omega_h^2}{2}\ip U{\mathcal G_{\beta,h}U}_h.
 \label{v4:drift-envelope}
\end{equation}
Here $H_{3,h}^{\rm g}$ is the cubic geometric stationary Hamiltonian and
$\mathcal G_{\beta,h}$ is the complete phase-product-defect generator of
\cite{OB25}.  The displayed equality is for the \emph{full original}
weak drift.  Cubic link terms have zero derivative in these tests; they
are not removed from the weak source $\mathcal Q_h$ or the finite action.
\end{lemma}
\begin{proof}
Expand the Poisson derivative of the original weak momentum row to second
field order.  The inherited quadratic action identity gives
\[
 \ip\beta{b_{2,h}^{\beta}}_h
 =-D_UH_{3,h}(U,p)[k]
  +\ip{F_{1,p}}{\mathcal G_{\beta,h}U}_h
  +\ip p{\mathcal G_{\beta,h}F_{1,u}}_h.
\]
The last term is zero by the retained trace identity, since both $p$ and
$F_{1,u}$ are constant scalar matrices.  The flat stationary spatial
connection differential of $(k,0)$ is zero: in each nonzero curl-free mode
$k_{ij}$ is a scalar multiple of $\kappa_i\kappa_j$, so the alternating
contraction in \eqref{v3:z1} vanishes.  Its velocity and temporal-connection
differentials at constant unit lapse and zero shift are also zero.  A
stationary cubic logarithmic Hamiltonian term is the original homogeneous
cubic remainder evaluated on this linear stationary connection.  Its
derivative in $(k,0)$ is consequently zero.  Coframe coefficients multiplying
that remainder cannot enter its cubic field coefficient.  The stationary
and Legendre envelopes cancel induced second-order saddle contributions.
This proves that $H_{3,h}$ can be replaced by $H_{3,h}^{\rm g}$ in this
particular derivative.  Substitution of $F_{1,p}$ gives
\eqref{v4:drift-envelope}.  This is an action-coefficient identity, not an
appeal to a future Einstein trajectory.
\end{proof}

For clarity the defect generator used in \eqref{v4:drift-envelope} can be
written, on $\tr U=0$, as
\begin{align}
 F_{ij}&=-\sum_k\mathfrak d_k(\beta_i,U_{kj})
                          +\sum_k\mathfrak d_k(\beta_k,U_{ij}),\notag\\
 (\mathcal G_{\beta,h}U)_{ij}
 &=\beta^k\delta_kU_{ij}+U_{kj}\delta_i\beta^k
       +U_{ik}\delta_j\beta^k
       +\tfrac12(F_{ij}+F_{ji})-\tfrac12\delta_{ij}\tr F,
 \label{v4:G}
\end{align}
where $\mathfrak d_k(f,g)=\delta_k(fg)-(\delta_k f)g-f\delta_kg$.
The $\delta_{ij}$ without an argument is the Kronecker symbol.

\begin{lemma}[Different axial modes do not contribute to the weak quadratic drift]
\label{v4:cross}
The polarization of \eqref{v4:drift-envelope} between $U_i$ and $U_j$,
$i\ne j$, is zero for every weak test.  A self-pair of $U_i$ contributes
only the $i$-component at the doubled $i$-axis frequency.  The constant
weak tests have zero pairing.
\end{lemma}
\begin{proof}
Use complex Fourier modes merely to evaluate the real multilinear
coefficient.  A cross-pair has field frequencies $\epsilon e_i$ and
$\epsilon'e_j$, $\epsilon,\epsilon'\in\{\pm1\}$; only a test at their
negative sum can have a nonzero spatial pairing.  In each coordinate the
nonzero frequencies in this trilinear pairing are then $1$ and $-1$.
Every derivative of a product occurring in a surviving term either keeps a
unit frequency or pairs these opposite frequencies into zero.  Its phase
symbol therefore obeys the exact additive rule.  No doubled coordinate
frequency occurs in that paired coefficient, including on $N=3$.
Consequently the complete paired geometric coefficient is the
ordinary-derivative coefficient with the common derivative scale
$\omega_h/(2\pi)$.  The continuum spatial Ward identity of the stationary
geometric Hamiltonian makes its canonical momentum drift an integrated
spatial divergence.  Its coefficient is zero.  This is exactly the
coefficient-level Ward identity retained in \cite{OB25},
\src{obv25:weak-source-expansion}, now applied only where the finite symbols
obey its product rule.  The logarithmic derivative has already been tested
and removed by \cref{v4:full-drift}; no link shift is approximated here.

A self-pair has frequencies $0,\pm2e_i$ modulo $N$.  At a nonzero such
frequency a curl-free test points along $e_i$.  For a constant shift,
$k=0$ and $\mathcal G_{\beta,h}U=\beta^i\delta_iU$ because all product
defects with a constant vanish.  Thus its pairing is
$-\omega_h^2\ip U{\beta^i\delta_iU}_h/2=0$ by skew-adjointness.
Only the asserted doubled axial tests remain.
\end{proof}

Define the \emph{folded} second-frequency symbol
\begin{equation}
 \widehat 2_N=\begin{cases}-1,&N=3,\\2,&N\ge5,\end{cases}
 \qquad \omega_{2,h}=2h^{-1}\sin(\pi h\widehat 2_N),
 \qquad \Delta_h=\omega_h-\frac12\omega_{2,h}>0.
 \label{v4:fold}
\end{equation}
The functions $\sin(4\pi x_i)$ and $\cos(4\pi x_i)$ below are sampled
literally.  In particular
$\delta_i\cos(4\pi x_i)=-\omega_{2,h}\sin(4\pi x_i)$, even at $N=3$.
Introduce the real polarization invariants
\begin{equation}
 e_i^c=\|A_i\|_F^2-\|B_i\|_F^2,\qquad
 e_i^s=2\ip{A_i}{B_i}_F.
 \label{v4:ellipticity}
\end{equation}

\begin{theorem}[Polarization-complete exact weak drift]
\label{v4:drift}
For \eqref{v4:axial-jet}, at every odd cutoff,
\begin{equation}
 \boxed{(d_h)_i=-\frac{\omega_h^2\Delta_h}{4}
             \bigl(e_i^c\sin(4\pi x_i)-e_i^s\cos(4\pi x_i)\bigr),
                  \qquad i=1,2,3.}
 \label{v4:drift-vector}
\end{equation}
This is the complete curl-free quadratic canonical drift.  It is independent
of the constant scalar momentum $s$.  Its positive spatial-mean norm is
\begin{equation}
 \boxed{\|d_h\|_h^2=\frac{\omega_h^4\Delta_h^2}{32}
                      \sum_{i=1}^3\bigl((e_i^c)^2+(e_i^s)^2\bigr).}
 \label{v4:drift-norm}
\end{equation}
For $N\ge5$, $\Delta_h=\omega_h(1-\cos\pi h)$; hence on every fixed
bounded polarization set and in every fixed Sobolev norm,
\begin{equation}
 (d_h)_i=-\pi^5h^2
       \bigl(e_i^c\sin(4\pi x_i)-e_i^s\cos(4\pi x_i)\bigr)+O(h^4).
 \label{v4:drift-refinement}
\end{equation}
\end{theorem}
\begin{proof}
By \cref{v4:cross}, it suffices to use axis one and a test
$\beta=(b(x_1),0,0)$.  Write the transverse two-by-two block of $U_1$ as
\[
 \begin{pmatrix}u&v\\v&-u\end{pmatrix},\qquad
 u=P\cos(2\pi x_1)+Q\sin(2\pi x_1),\quad
 v=R\cos(2\pi x_1)+S\sin(2\pi x_1).
\]
Primes in this paragraph mean $\delta_1$, not an assumed Leibniz
 derivative.  The first spatial rotation block is
\[
 (z_1)^R_{\{2,3\},\{2,3\}}
       =\tfrac12\begin{pmatrix}v'&-u'\\-u'&-v'\end{pmatrix},
       \qquad (z_1)^B=\tfrac{s}{2}I.
\]
For $k_{11}=2b'$ and all other $k_{ij}=0$, the derivative of the cubic
geometric \emph{Lagrangian--Legendre} functional in
\eqref{v3:functional} is
\begin{equation}
 -D_UH_{3,h}^{\rm g}[k]
       =\frac12\ip{b'}{(u')^2+(v')^2+s^2}_h.
 \label{v4:one-axis-action}
\end{equation}
Here are the individual coefficient checks.  The second coframe coefficient
is $-U^2/8$, whose derivative in $k$ is zero because $Uk+kU=0$.
The degree-one coframe coefficient contributes $k/2$ in the $23$ Cartan
row.  The rotational component of $[z_2,z_3]$ in that row is
$((u')^2+(v')^2+s^2)e_1/4$, giving the displayed term.  The degree-two
velocity-cofactor derivative paired with the scalar boost is zero in this
variation.  Explicitly, for any symmetric $K,W$, that derivative of
$2\operatorname{cof}\sqrt{I+U}$ is
\[
 \left(\tfrac12\tr K\tr W-\tr(KW)\right)I
 -\tfrac12(\tr W\,K+\tr K\,W)+\tfrac34(KW+WK);
\]
it vanishes at $K=\diag(2b',0,0)$ and $W=sI$.
These facts derive \eqref{v4:one-axis-action} from the finite scalar action,
not from an isotropy assumption about the answer.  The $s^2$ term integrates
to zero.

In \eqref{v4:G} the same axis restriction gives
$\mathcal G_\beta U=\delta_1(bU)-b'U$.  Thus its contribution in
\eqref{v4:drift-envelope}, after exact summation by parts, is
\[
 \omega_h^2\ip b{u u'+v v'}_h
                         +\omega_h^2\ip{b'}{u^2+v^2}_h.
\]
Put $q=u^2+v^2$.  The literal pointwise identity
$(u')^2+(v')^2=\omega_h^2(P^2+Q^2+R^2+S^2-q)$ gives
\begin{equation}
 \ip\beta{d_h}_h
 =-\frac{\omega_h^2}{2}
                 \ip b{\delta_1q-2(uu'+vv')}_h.
 \label{v4:one-axis-defect}
\end{equation}
No product rule has been applied to $q$.  Direct trigonometric multiplication
and the folded symbol give
\[
 \delta_1q-2(uu'+vv')
 =\Delta_h\bigl((P^2+R^2-Q^2-S^2)\sin4\pi x_1
                         -2(PQ+RS)\cos4\pi x_1\bigr).
\]
The full Frobenius invariants are twice the corresponding two-component
invariants.  This proves \eqref{v4:drift-vector}.  The resulting vector is
itself curl-free and has zero mean, so agreement on all weak tests identifies
$d_h$ without a missing projection or constant component.

The sampled doubled sine and cosine are orthogonal, with mean squares
$1/2$ on every odd grid.  This proves \eqref{v4:drift-norm}.
For $N\ge5$, the sine double-angle identity gives the stated $\Delta_h$;
$\omega_h^2\Delta_h/4=\pi^5h^2+O(h^4)$ gives
\eqref{v4:drift-refinement}.  All functions are in one fixed doubled Fourier
bank, so the estimate holds in the stated norms.  At $N=3$ the same exact
formula instead has $\Delta_h=3\omega_h/2$.  Replacing that folded value by
the nonfolded double-angle value would be incorrect.
\end{proof}

\begin{remark}[Explicit correction of the exceptional three-site prefactor]
For $B_i=0$ and $A_i=\alpha_i\mathsf T_i$, the new formula agrees with
\eqref{v2:drift-vector} for $N\ge5$ and leaves its refinement estimate
unchanged.  At $N=3$, folding the doubled sine alone is insufficient:
the scalar coefficient there must instead be
\[
 d_h^{\rm ax}=\frac{\omega_h^2\Delta_h}{8}
                         =\frac{3\omega_h^3}{16},
\]
not the nonfolded value $\omega_h^3(1-\cos\pi h)/8=\omega_h^3/16$.
Thus the literal three-site specialization of the displayed prefactor in
\eqref{v2:drift-vector} needs this factor-three correction.  For the
unperturbed first-axis shift $(0,-1,1)$ its spatial-mean pairing is then
$-243/4$, agreeing with the original coefficient witness in \cite{OB25},
\src{obv24:exact-witness}.  The cofinal statements in V2, the source-only
rank certificates and the nonzero-pairing pole argument of V3 are
unaffected.  The earlier bodies remain unchanged as records; the folded
coefficient used in this body is \eqref{v4:fold}.
\end{remark}

\section{The source-neutral first-jet class and its exact initial preparation}
\label{v4:neutral-section}

\begin{corollary}[Exact cancellation and its complete polarization conditions]
\label{v4:neutral}
For the axial first jets in \eqref{v4:axial-jet}, the following are equivalent:
\begin{enumerate}[label=\textup{(\roman*)},leftmargin=2.6em]
\item $d_h=0$ at one odd cutoff;
\item for each axis, $\|A_i\|_F=\|B_i\|_F$ and $\ip{A_i}{B_i}_F=0$;
\item $d_h=0$ at every odd cutoff.
\end{enumerate}
These conditions do not set the configuration to zero.  If every axis is
nonzero, the zero-drift set has codimension six in the twelve-dimensional
space $\prod_i(\mathcal V_i\times\mathcal V_i)$.
\end{corollary}
\begin{proof}
The prefactor in \eqref{v4:drift-norm} is strictly positive, so its sum of
squares vanishes exactly under (ii).  Those conditions are independent of
$h$.  At a nonzero orthogonal equal-norm pair, the differentials of
$\|A\|^2-\|B\|^2$ and $2\ip AB$ are independent: their gradient vectors
$(2A,-2B)$ and $(2B,2A)$ are nonzero and orthogonal.  Thus each two-dimensional
polarization plane contributes two independent scalar conditions, proving
the codimension assertion.  This is not an open set of arbitrary first jets.
\end{proof}

Let $\mathsf T_i$ be the off-diagonal tensor from \cref{v2:family}.  For the
ordered transverse pair $j<k$, let $\mathsf P_i$ have entries $+1$ at $jj$,
$-1$ at $kk$, and zero elsewhere.  These two tensors have squared Frobenius
norm two and are orthogonal.  A useful zero-drift family is
\begin{equation}
 U_i=r_i\bigl(\mathsf T_i\cos(2\pi x_i+\phi_i)
                 +\sigma_i\mathsf P_i\sin(2\pi x_i+\phi_i)\bigr),
 \quad r_i>0,\quad \sigma_i\in\{-1,1\}.
 \label{v4:balanced-family}
\end{equation}
This has six continuous parameters and discrete polarization orientations.
All three spatial coordinates occur.  If a constant translation vector
$\xi$ satisfies $\xi\cdot\nabla U=0$ for the continuum profiles, independence
of their axial Fourier supports gives $\xi_i=0$ for every $i$.  This simple
observation does not exclude all possible metric isometries or assert a
generic Cauchy-data class.

To prepare the original initial constraints, define
\[
 m_i=\tfrac12(\|A_i\|_F^2+\|B_i\|_F^2)>0,
 \qquad V_i=A_i s_i-B_i c_i,
 \qquad p(s,b)=-sI+\sum_i b_iV_i.
\]
For a fixed bounded polarization set with $m_i$ bounded below, put
\begin{equation}
 s_{h,\pm}=\pm\omega_h\sqrt{\frac16\sum_i m_i}.
 \label{v4:s-calibration}
\end{equation}
On \eqref{v4:balanced-family}, $m_i=2r_i^2$.

\begin{proposition}[Original-action mean balance for both polarizations]
\label{v4:mean}
The complete quadratic mean constraints on $(U,p(s,b))$ are
\begin{equation}
 \boxed{Q_h(s,b)=
 \begin{pmatrix}
  \frac32s^2-\sum_i m_i b_i^2-\frac{\omega_h^2}{4}\sum_i m_i\\
  -\omega_hm_1b_1\\-\omega_hm_2b_2\\-\omega_hm_3b_3
 \end{pmatrix}.}
 \label{v4:mean-vector}
\end{equation}
They have the roots $(s,b)=(s_{h,\pm},0)$ and a cutoff-uniform inverse for
the derivative in these four parameters on the stated polarization sets.
\end{proposition}
\begin{proof}
All configuration and nonconstant momentum modes are transverse and
trace-free.  Orthogonality gives
\[
 \|p\|_h^2-\tfrac12\|\tr p\|_h^2=-\tfrac32s^2+\sum_i m_i b_i^2,
 \qquad \sum_i\|\delta_iU\|_h^2=\omega_h^2\sum_i m_i,
 \qquad \ip p{\delta_iU}_h=-\omega_hm_i b_i.
\]
Insert these identities into the original constant-multiplier quadratic
Hamiltonian retained in \cite{OB25}, \src{obv17:quadratic-H}.  Its validity
for arbitrary transverse trace-free polarizations follows from its displayed
Frobenius energy, not from a single-polarization example.  The cubic
logarithmic remainder cannot alter these constant-multiplier quadratic
field coefficients.  This gives \eqref{v4:mean-vector}.  At a root, the
parameter derivative is
$\diag(3s_{h,\pm},-\omega_hm_1,-\omega_hm_2,-\omega_hm_3)$.
The positive lower bounds for $m_i$, $|s_{h,\pm}|$ and $\omega_h$ give the
uniform inverse.
\end{proof}

\begin{theorem}[Exact finite preparation of the source-neutral family]
\label{v4:preparation}
For sufficiently small $|a|$, each compact nonzero polarization subfamily
above admits exact original-action initial data
\begin{equation}
 \boxed{\mathscr C_h(X_h(a))=0,\qquad
 X_h(a)=a(U,p(s_{h,\pm},0))+O(a^2)}
 \label{v4:prepared}
\end{equation}
in the inherited preparation norm $\mathcal P_h^r$, uniformly at all
sufficiently fine odd cutoffs.  The remainder and fixed parameter
differentials have uniform bounds.  The first jet on
\eqref{v4:balanced-family} has $d_h=0$ exactly.  It remains nontrivial and
three-axis after the full nonlinear initial constraint solve.
\end{theorem}
\begin{proof}
The first jets solve the flat linear constraints: configuration divergence
and trace vanish, and the momenta have zero phase divergence.  Use the
inherited right inverse for the complete nonconstant initial rows.  Its
nonlinear contraction gives a range correction $O(a^2)$ with uniform
analytic coefficient bounds.  Retain its four mean equations.  Division by
$a^2$ is justified by the integral Taylor formula at the flat point, as in
\cref{v2:preparation}, and produces $Q_h(s,b)+O(a)$ in $C^1$.
\Cref{v4:mean} supplies the nondegenerate four-parameter root, so the
same finite contraction solves the remaining means with
$s=s_{h,\pm}+O(a)$ and $b=O(a)$.  All polarization profiles lie in a fixed
Fourier bank with uniform Sobolev bounds; the required coefficient estimates
are therefore uniform on the stated compact parameter set.  This argument
uses the \emph{original} constraint map, with its complete logarithmic
terms, in both contractions.  Only the starting mean calibration is
quadratic.

The correction is second order, so it does not replace the chosen first
jet.  The quadratic weak drift is therefore the coefficient in
\cref{v4:drift}, and is zero on the balanced family.  The three nonzero
fundamental axial coefficients persist at sufficiently small nonzero $a$.
No conclusion about the evolution of the polarization balance is implicit.
\end{proof}

The choice of sign in \eqref{v4:s-calibration} is a choice of initial trace
momentum, not a change of the protected orientation or a redefinition of
physical time.  At a fixed cutoff the same local preparation argument uses
the finite analytic range chart; the three-site application below needs
only that inherited fixed-cutoff initial chart.

\section{A regular balanced seed for the unchanged signed action}
\label{v4:certificate-section}

We now calculate the actual second weak source on a source-neutral seed;
it is not replaced by zero merely because its drift target is zero.  At
$N=3$, choose
\begin{equation}
 \boxed{U=\sum_{i=1}^3(\mathsf T_ic_i+\mathsf P_is_i),
                   \qquad p=3\sqrt3\,I.}
 \label{v4:finite-seed}
\end{equation}
Thus $s=-3\sqrt3$ in \eqref{v4:axial-jet}.  This is the negative root in
\eqref{v4:s-calibration}, since $m_1=m_2=m_3=2$ and $\omega_h^2=27$.
It is an exact mean-compatible first jet, and \cref{v4:preparation}
supplies its nonlinear initial records.

Use the raw symmetric coordinates, lexicographic sites, $K^{-1}$ and weak
basis of \cref{v3:certificate-section}.  In particular $K^{-1}$ is the
inverse \emph{bilinear} Cartan form, not the inverse Euclidean coefficient
Gram.  Form $J_1^{\rm all}$ from the complete first-source derivative.
Select exactly the 78 columns $\mathcal P$ in \eqref{v3:pivots}, obtaining
$C$.  Form the 29-column weak matrix $W$ by the full first-jet compiler
\eqref{v3:compiler-formula}, and put
\begin{equation}
 F=(C,W),\qquad A=C^TK^{-1}C,\qquad L=C^TK^{-1}W,
 \qquad D=W^TK^{-1}W-L^TA^{-1}L.
 \label{v4:certificate-matrices}
\end{equation}
The next proof verifies $A^{-1}$ at this cutoff.  It does not import the
trace-free active-block assertion from the different detuned family.

\begin{theorem}[A source-neutral seed with a nonsingular complete signed Gram]
\label{v4:finite-rank}
For \eqref{v4:finite-seed} in characteristic zero,
\begin{equation}
 \boxed{\rank J_1^{\rm all}=78,\qquad \rank F=107,
              \qquad\det(F^TK^{-1}F)\ne0,\qquad \det D\ne0.}
 \label{v4:rank}
\end{equation}
The original weak quadratic target is exactly zero, whereas the remaining
weak source has full rank 29.  Both statements concern the same first jet
and the same original action.
\end{theorem}
\begin{proof}
All coefficients lie in $\mathbb Q(\sqrt3)$.  The phase derivative is
$3(S_i-S_i^{-1})$, and the literal sampled trigonometric vectors are
\[
 c=(1,-1/2,-1/2),\qquad s=(0,\sqrt3/2,-\sqrt3/2).
\]
Specialize to $\mathbb F_{241}$ with $\sqrt3\mapsto56$, since
$56^2\equiv3\pmod{241}$.  Thus the momentum scalar is
$3\sqrt3\mapsto168$, while the convention $p=-sI$ has $s\mapsto73$.
This sign is part of the certificate; choosing an unrelated square-root
image for the momentum would not be the same algebraic first jet.
Every denominator of the pointwise saddle maps, coframe jets and
coefficient extraction is prime to 241.

The companion verifier constructs the actual $U,p,z_1$ in this ring.  It
evaluates the full first-source derivative and the compiler
\eqref{v3:compiler-formula}, retaining cubic electric and magnetic terms,
normal correction and first auxiliary response.  The complete arrays
$J_1^{\rm all},W,F,K^{-1}$ are supplied.  Direct modular elimination gives
\begin{equation}
 \boxed{\det A\equiv113,\qquad \det D\equiv56,\qquad
                      \det(F^TK^{-1}F)\equiv62\pmod{241}.}
 \label{v4:residues}
\end{equation}
Their independent Schur check is $113\cdot56\equiv62\pmod{241}$.
Nonzero specializations prove nonzero characteristic-zero determinants and
hence independence of the 107 source columns.  The first source has the
constant lapse and all 29 curl-free shifts in its kernel, by its mean
homogeneity and \cref{v3:universal-flat}.  Its rank is at most
$108-30=78$, while $\det A\ne0$ gives the reverse inequality.  This proves
all rank assertions without inferring a characteristic-zero upper rank
from a zero residue.

Finally \cref{v4:neutral} applies with $\|\mathsf T_i\|^2=
\|\mathsf P_i\|^2=2$ and $\ip{\mathsf T_i}{\mathsf P_i}=0$.
It gives the actual zero weak target.  A separate exact cubic-action check
pairs the full drift with all 29 weak basis elements and gives zero; its
role is to verify the implementation of the analytical identity.  The
nonzero Gram shows that this target cancellation does not come from
silently dropping weak source columns.
\end{proof}

\begin{corollary}[A regular neighborhood of balanced polarizations at fixed cutoff]
\label{v4:finite-neighborhood}
At $N=3$, there is a neighborhood of \eqref{v4:finite-seed} in the axial
polarization parameters, with the same choice of calibrated trace sign,
on which the complete leading signed Gram remains nonsingular.  Its
intersection with the nonzero balanced-polarization submanifold is
nonempty and relatively open there.  Every such balanced parameter has a
local exactly prepared original-action continuation at sufficiently small
nonzero $a$.
\end{corollary}
\begin{proof}
The first-jet compiler and the calibrated scalar square root are real
analytic in the polarization coefficients near the displayed point.  Keep
the same columns and weak basis.  Continuity of the nonzero determinant
preserves it on a smaller neighborhood, with a bounded inverse on compact
subsets.  The exact preparation remains analytic there.  Its normalized
finite-amplitude Gram is an analytic perturbation of this nonsingular
leading matrix; at fixed cutoff it is nonsingular for sufficiently small
nonzero amplitude.  Mean-lapse homogeneity retains the one null multiplier
direction.  The inherited exact normalized continuation theorem then
applies without a change of action or an added field force.  It supplies
only a local interval at these parameters.
\end{proof}

\section{The initial electric response loses its order-one amplitude term}
\label{v4:physical-section}

All statements in this section are at a fixed cutoff with a verified
regular leading source, as supplied by \cref{v4:finite-rank}.  They are not
claims that the displayed inverse or remainder constants are uniform in
$h$.  Use fixed multiplier coordinates on the mean-lapse-zero slice and
write $S_a=\diag(aI_A,a^2I_W)$ as in \eqref{v1:rate-grade}.  The normalized
initial consistency equation has
\begin{equation}
 M_h(a)=U S_a\Gamma_h(a)S_aU^*,\qquad
 \widehat b_h(a)=S_a^{-1}U^*\mathsf B_h(X_h(a),\lambda_*),\qquad
 \lambda_t=-U S_a^{-1}\Gamma_h(a)^{-1}\widehat b_h(a).
 \label{v4:grading}
\end{equation}
Signs agree with the retained normalized continuation.  Equivalent fixed
basis changes do not alter the orders below.  The source target at $a=0$
is $(0,d_h)$.

\begin{theorem}[Tame initial electric amplitude on the balanced regular branch]
\label{v4:initial-electric}
On the regular balanced family at $N=3$, uniformly on a sufficiently small
compact parameter neighborhood inside \cref{v4:finite-neighborhood},
\begin{equation}
 \boxed{\widehat b_h(a)=O_h(a),\quad
 \Gamma_h(a)^{-1}\widehat b_h(a)=O_h(a),\quad
                         \|E_h(a,0)\|_h\le C_h|a|.}
 \label{v4:small-electric}
\end{equation}
Here $E_h(a,0)$ is the literal initial electric link read evaluated with
its \emph{own normalized exact-action multiplier rate}.  No comparison
connection or prescribed time derivative is substituted.  The initial
spatial link curvature is also $O_h(a)$.
\end{theorem}
\begin{proof}
Exact preparation gives an analytic amplitude curve.  Its graded source
and normalized Gram are analytic at $a=0$ in the inherited stationary
chart.  The leading target is $(0,d_h)=0$ by \cref{v4:neutral}; hence
$\widehat b_h(a)=O_h(a)$.  The leading Gram is nonsingular, and its inverse
is uniformly bounded on the fixed-cutoff compact parameter set.  This
proves the first two assertions.  Notice that the active quadratic drift
is not set to zero: its first contribution to this normalized target has
an additional factor $a$ from the active grading.

The exact differentiated stationary equation gives, in the physical
spatial-connection coordinates,
\[
 \dot z=z_XF_h^{\rm can}
                 +Z_aK_a^{-1}F_a\Gamma_h(a)^{-1}\widehat b_h(a).
\]
At the initial unit-lapse, zero-shift cut,
$F_h^{\rm can}=O_h(a)$ and $z_X$ is bounded.  The remaining stationary
factors have analytic bounded limits at the fixed cutoff; the preceding
source estimate therefore gives $\dot z=O_h(a)$.  This is the exact
stationary-response identity used in \cref{v3:physical-pole}, now with a
vanishing rather than nonzero leading target.  It retains the true
finite-amplitude tangent Hessian $K_a$ and inclusion $Z_a$.

The stationary initial connections satisfy $A_i,A_0=O_h(a)$.  In the
literal read
\[
 E_i=h^{-1}\{\dot U_iU_i^{-1}+A_0-U_iA_{0,+i}U_i^{-1}\},\qquad
 U_i=\exp(hA_i),
\]
the differential of the exponential bounds the first term by $C_h|a|$;
the remaining difference and commutator have the same fixed-cutoff order.
This proves \eqref{v4:small-electric}.  The analytic spatial plaquette
read, zero at flat data, similarly has size $O_h(a)$.  Neither estimate
identifies that curvature with the Riemann tensor on a common spacetime
interval.
\end{proof}

\begin{proposition}[Near balance there is no detuning-type source pole at this cutoff]
\label{v4:near-neutral}
In a sufficiently small fixed-cutoff neighborhood from
\cref{v4:finite-neighborhood}, let $e=(e_i^c,e_i^s)_{i=1}^3$.  The leading
original physical response $v_0$ and the positive minimum reproduction
vector $w_0$ satisfy
\begin{equation}
 \|v_0\|_h+\|w_0\|_h\le C_h|e|,
 \qquad \|E_h(a,0)\|_h\le C_h(|a|+|e|).
 \label{v4:near-response}
\end{equation}
In particular a parameter perturbation with $|e|=O(|a|)$ retains the
$O_h(a)$ initial electric estimate.  This local parameter bound is not a
cofinal estimate and not attraction toward the balanced submanifold.
\end{proposition}
\begin{proof}
Both the positive source Gram and the complete signed Gram have uniformly
bounded inverses on a smaller fixed-cutoff compact neighborhood.  The
corresponding response maps are consequently bounded linear maps of their
actual targets, with coefficients analytic in the polarization parameters.
The leading target is $(0,d_h)$, and \eqref{v4:drift-norm} bounds its norm
by $C_h|e|$, including the fixed basis conversion.  The first inequality
follows.  The graded finite-amplitude response is an analytic perturbation
with error $O_h(a)$ on the same compact set.  The literal read comparison
in the previous proof then gives the second inequality.  Unlike the
singular equal-amplitude crossing in \cref{v3:detuning-pole}, the present
parameter center has a nonzero signed determinant.
\end{proof}

\section{The next source and the remaining clock obstruction}
\label{v4:next-source}

The improvement in \cref{v4:initial-electric} removes one particular
amplitude obstruction.  It does not make the normalized multiplier rate
small.  Write the actual initial drift as
\[
 U^*\mathsf B_h(X_h(a),\lambda_*)
              =a^2B_{2,h}+a^3B_{3,h}+O_h(a^4).
\]
Its weak quadratic component is zero on the balanced family.  Define
\begin{equation}
 b_{1,h}=\binom{(B_{2,h})_A}{(B_{3,h})_W},\qquad
 u_{1,h}=\Gamma_{0,h}^{-1}b_{1,h}.
 \label{v4:next-target}
\end{equation}
These are actual source coefficients, not newly chosen rates.

\begin{proposition}[The first surviving graded multiplier coefficient]
\label{v4:remaining-rate}
At the fixed regular initial branch,
\begin{equation}
 \widehat b_h(a)=a b_{1,h}+O_h(a^2),\qquad
 \Gamma_h(a)^{-1}\widehat b_h(a)=a u_{1,h}+O_h(a^2).
 \label{v4:next-response}
\end{equation}
In the fixed active/weak multiplier coordinates this gives
\begin{equation}
 (U^*\lambda_t)_A=-(u_{1,h})_A+O_h(a),\qquad
 (U^*\lambda_t)_W=-a^{-1}(u_{1,h})_W+O_h(1).
 \label{v4:remaining-pole}
\end{equation}
Thus the already proved initial electric bound can coexist with an
uncontrolled weak multiplier rate.  The coefficient $(u_{1,h})_W$ has
not been calculated or proved zero in this body.
\end{proposition}
\begin{proof}
Divide the active and weak components of the actual drift by $a$ and $a^2$,
respectively.  The weak $a^2$ coefficient is zero, yielding the first
formula.  Expand the inverse of the regular analytic Gram and use
\eqref{v4:grading}.  The different amplitude divisions give
\eqref{v4:remaining-pole}.  No inverse estimate uniform in $h$ is needed
or asserted for this fixed-cutoff calculation.
\end{proof}

In particular a full harmonic metric-jet shadowing statement must still
meet the initial lapse--shift test of \cref{v1:clock-test}, or supply an
independently controlled coordinate comparison.  Exact source neutrality
at the first jet is not shown to persist under the original flow.
The physical time interval remains that of the inherited local
continuation theorem.  A better initial coefficient alone cannot extend
that interval.

There is a useful distinction from \cref{v3:preparation-independent}:
the leading weak \emph{operator} $\mathcal Q_h$ does not depend on the
second preparation, but the next weak \emph{drift} generally does.
If $\mathsf B_{W,h}$ is the weak projection of the original drift map,
its retained flat first derivative is zero.  For
$X(a)=aY+a^2Y_2+a^3Y_3+\cdots$, Taylor expansion gives
\begin{equation}
 (B_{3,h})_W
 =D^2\mathsf B_{W,h}(0)[Y,Y_2]
                      +\tfrac16D^3\mathsf B_{W,h}(0)[Y,Y,Y].
 \label{v4:third-drift}
\end{equation}
The zero flat derivative follows from the original flat momentum
propagation identity, the same identity used to eliminate first-order
weak drift in \cite{OB25}, \src{obv24:finite-witness}.
Equation \eqref{v4:third-drift} identifies the next coefficient to compute;
it is not a solved preparation problem.  Altering $Y_2$ must retain the
second initial constraint equation and the next mean compatibility rows.
There is no license to choose $Y_2$ from the drift equation alone.

\section{Verification, scope, and the next action-native calculation}
\label{v4:status}

The all-cutoff conclusion is now stronger in one precise direction:
there is a nonzero three-axis first-jet family for which the original
weak quadratic drift vanishes identically, instead of having the sharp
nonzero $h^2$ size of the single-polarization axial family.  The proof uses
the finite defect, including folded frequencies, and does not set the
Cartan response equal to a positive-minimum surrogate.  On the same
source-neutral family an independently regenerated three-site coefficient
matrix has a regular complete signed Gram.  This gives an actual exact
finite initial branch with $O_h(a)$ literal electric read, rather than
only a conditional statement about a source-neutral matrix.

\paragraph{What was verified.}
The package supplies the complete first and second source residue arrays,
the raw Cartan inverse, the signed Gram and its active and weak Schur blocks.
The verifier checks every determinant in \eqref{v4:residues}, including the
quadratic-field embedding and momentum sign.  Its reconstruction option
re-evaluates all 108 first-source columns by direct coefficient
differentiation, and all 29 weak columns by \cref{v3:compiler}.

An independent cubic-action extractor retains one higher amplitude degree
and a first-variation variable.  It pairs the complete weak drift with all
29 weak basis vectors, including cubic link terms.  Every pairing is zero
on the balanced seed; nonbalanced sine amplitudes reproduce the nonzero
folded coefficient.  Symbolic checks verify the one-axis contractions and
mean calibration.  General-grid floating tests retain the phase symbols
and check balanced and nonbalanced profiles, but establish no all-cutoff
rank estimate.  The all-cutoff cancellation rests on the written proof.
No exact original-action spacetime integration has been performed.

\paragraph{Boundary of the advance.}
The balanced first-jet set has codimension six in the specified axial
polarization space.  It is not an open set of unrestricted microscopic
data and has not been dynamically selected.  Its lapse part need not be
trace-free; the all-cutoff positive active elimination for the different
family in V2 must not be transplanted to it.  Only the transverse-shift
elimination in V1 is automatically retained for scalar momentum.  The new
complete signed-rank certificate is at $N=3$ only.  In particular the zero
target does not resolve uniqueness or higher consistency conditions at
an unverified singular cutoff.

The next useful work has two concrete targets.  First, calculate the actual
first-jet source matrices on the balanced family along refinement, using
\cref{v3:compiler}, and establish their regularity or the additional
compatibility equations where they are singular.  Second, calculate
\eqref{v4:next-target} from exact second-order preparation and the full
cubic drift, rather than another determinant with the first source alone.
Cofinal control still requires the true nonlinear Hessian and physical
inclusion, differentiated sources, literal-curvature identification,
lapse--shift comparison where relevant, and a common physical lifespan.
The static $2/3$ resonance cannot amplify the \emph{zero leading target}
on a regular branch, but it may act on those next sources.  No raw-action
Einstein continuum limit is claimed.

\begin{thebibliography}{9}
\raggedright
\setcounter{enumiv}{6}
\bibitem{v4:BD}
D. R. Brill and S. Deser,
\emph{Instability of closed spaces in general relativity},
Communications in Mathematical Physics \textbf{32} (1973), 291--304.
\href{https://doi.org/10.1007/BF01645610}{doi:10.1007/BF01645610}.
This supplies context for nonlinear compatibility of compact initial data,
not the renewal drift or the finite source certificate.
\bibitem{v4:Isi}
M. Isi, \emph{Parametrizing gravitational-wave polarizations},
Classical and Quantum Gravity \textbf{40} (2023), 203001.
\href{https://doi.org/10.1088/1361-6382/acf28c}{doi:10.1088/1361-6382/acf28c};
\url{https://arxiv.org/abs/2208.03372v2}.
This supplies polarization terminology, not the spatial first-jet cancellation
or an evolution result for the present model.
\end{thebibliography}
% END IMMUTABLE EXACT-ACTION BRIDGE V4 BODY

% BEGIN IMMUTABLE EXACT-ACTION BRIDGE V5 BODY
\clearpage
\renewcommand{\thesection}{\arabic{section}}
\renewcommand{\theHsection}{bridgeV5.\arabic{section}}
\setcounter{section}{33}
\section{The next source: preparation freedom and a translation obstruction}
\label{v5:context}

The leading weak target is exactly zero on the balanced axial first jets of
\cref{v4:neutral}.  It does not follow that the next weak multiplier rate is
small.  The target in \eqref{v4:next-target} couples an active quadratic drift
to a weak cubic drift through the \emph{complete signed Gram}.  We investigate
which part of this target can actually be changed by the second initial
preparation.  The answer is not obtained by declaring a favourable remainder.

There are two complementary results.  On the balanced three-site seed, the
second configuration jet, subject to the linear initial constraint, controls
all 26 nonconstant weak drift directions.  Momentum corrections can restore
the four next mean-constraint equations without changing that control.  The
three constant-shift drift directions, however, are independent of every
second preparation.  Their cubic values vanish on the seed, but the active
Schur contribution to those same directions is nonzero.  Consequently an
$a^{-1}$ initial shift-velocity pole is unavoidable for every analytic exact
preparation with this first jet.  It coexists with the $O_h(a)$ initial
literal electric bound of \cref{v4:initial-electric}.

The positive preparation result is stronger than cancellation of a formal
coefficient: an analytic refinement of the exact initial preparation makes
the entire mean-zero curl-free part of the \emph{actual initial} shift
velocity zero.  Its remaining leading weak pole is spatially constant.
Nothing is subtracted from the equation of motion or from the measured
curvature.  Removing that constant pole by a change of coordinates would be
a different comparison theorem, not an operation performed here.

\paragraph{Inherited inputs and the scope of the advance.}
We retain the original analytic stationary saddle, canonical Hamiltonian and
initial-constraint map from \cite{OB25}, especially
\src{obv17:linear-C}, \src{obv17:range}, \src{obv23:weak-action}, and the
normalized exact continuation.  We use the first-jet weak-source theorem
\cref{v3:preparation-independent}, the balanced preparation
\cref{v4:preparation}, and its exact three-site Gram
\cref{v4:finite-rank}.  They are not reproved as new results.  The new
calculations concern the \emph{drift derivative}, not the derivative of the
quadratic momentum constraint studied in \cite{OB25},
\src{obv23:Ep}.  The two operators must not be identified.  Ordinary Schur
elimination and the analytic implicit-function theorem are tools, not new
abstract principles.  The source identities, finite minors, preparation
invariant and initial-rate conclusions below are specific to the unchanged
action.  Statements at $N=3$ are kept separate from the all-cutoff coefficient
identities.

\section{Original coefficients and the correction imposed by preparation}
\label{v5:coefficients}

Use the physical pairing and phase derivatives of \cref{v1:notation}.  Write
$Y=(u,p)$, $e=(1,0)$ for the unit-lapse, zero-shift multiplier, and
\[
 H_h(X)=\mathscr H_h(X,e),\qquad
 P_{h,\nu}(X)=\partial_\lambda\mathscr H_h(X,e)[\nu],
 \qquad \nu=(n,\beta).
\]
Here $P_{h,n}$ is the Hamiltonian lapse row, whereas the scalar component
of the source's initial map $\mathscr C_h$ is $-P_{h,n}$.  The common zero
set is unchanged.  With the convention
$\{P,H\}=DP[\mathbb J\nabla H]$, set
\[
 \mathsf B_{h,\nu}(X)=\{P_{h,\nu},H_h\}(X).
\]
At a fixed cutoff the homogeneous field coefficients are denoted by
$H_j$, $P_{\nu,j}$ and $\mathsf B_{\nu,j}$; for example
$H_h(aY)=a^2H_2(Y)+a^3H_3(Y)+O_h(a^4)$.  The unit-lapse integrated linear
Hamiltonian is zero.  No amplitude polynomial is used as an evolution law.

Put $t_u=\tr u$, $v_j=\delta_i u_{ij}$ and
$\Delta_{\delta,h}=\sum_i\delta_i^2$.  This phase Laplacian is
distinct from the scalar frequency defect in \eqref{v4:fold}.  The full quadratic
Hamiltonian is
\begin{align}
 H_2(u,p)={}&\|p\|_h^2-\tfrac12\|\tr p\|_h^2
   +\tfrac14\sum_i\|\delta_i u\|_h^2-\tfrac12\|v\|_h^2\notag\\
 &+\tfrac12\ip v{\delta t_u}_h
                         -\tfrac14\|\delta t_u\|_h^2.
 \label{v5:H2}
\end{align}
Its linear Hamiltonian vector field $F_1(Y)=(V(p),\mathcal F(u))$ has
\begin{align}
 V(p)&=2p-(\tr p)I,\notag\\
 \mathcal F(u)_{ij}
  &=\tfrac12\Delta_{\delta,h}u_{ij}
    -\tfrac12(\delta_i v_j+\delta_j v_i)
    +\tfrac12\delta_i\delta_jt_u
    +\tfrac12\delta_{ij}(\delta_kv_k-\Delta_{\delta,h}t_u).
 \label{v5:F1}
\end{align}
These are the original quadratic Fierz--Pauli coefficients.  The displayed
normalization follows either from the flat stationary envelope or by
varying \eqref{v5:H2} in the Frobenius pairing.  Every integration by parts
uses skew-adjointness, not a phase Leibniz rule.

The linear multiplier Hamiltonian and its vector field are
\begin{align}
 P_{\nu,1}(u,p)
  &=-\ip n{\delta_i\delta_j u_{ij}-\Delta_{\delta,h}\tr u}_h
            +\ip p{\delta_i\beta_j+\delta_j\beta_i}_h,\notag\\
 Y_\nu:=\mathbb J\nabla P_{\nu,1}
  &=(\delta_i\beta_j+\delta_j\beta_i,
                         \delta_i\delta_j n-\delta_{ij}\Delta_{\delta,h}n).
 \label{v5:linear-multiplier}
\end{align}

\begin{lemma}[The linear consistency row and the prepared quadratic drift]
\label{v5:prepared-B2}
For all symmetric finite records,
\[
 \mathsf B_{n,1}=-P_{(0,\delta n),1},\qquad
 \mathsf B_{(0,\beta),1}=0.
\]
If $X(a)=aY+a^2Y_2+O_h(a^3)$ satisfies the original initial constraints,
then its quadratic drift coefficient is independent of $Y_2$ and equals
\begin{equation}
 \boxed{\widehat B_{2,\nu}(Y)
 =-DH_3(Y)[Y_\nu]+DP_{\nu,2}(Y)[F_1Y]
                         +P_{(0,\delta n),2}(Y).}
 \label{v5:prepared-drift-formula}
\end{equation}
The last term is essential for the lapse tests.  It is zero for a pure
shift test.
\end{lemma}
\begin{proof}
Apply \eqref{v5:linear-multiplier} to $F_1Y$.  The scalar linear curvature
of $V(p)$ is $2\delta_i\delta_jp_{ij}$, and the phase divergence of
$\mathcal F(u)$ is zero by commutation of its derivatives.  These facts give
the two linear identities, including the sign in the first one.
At quadratic order on the straight curve, the Poisson expansion is
\[
 \mathsf B_{\nu,2}(Y)
 =-DH_3(Y)[Y_\nu]+DP_{\nu,2}(Y)[F_1Y].
\]
The actual prepared curve adds $\mathsf B_{\nu,1}(Y_2)$.
Exact preparation gives
$P_{\mu,1}(Y_2)=-P_{\mu,2}(Y)$ for every multiplier $\mu$.
Taking $\mu=(0,\delta n)$ gives the last term of
\eqref{v5:prepared-drift-formula}.  Thus it is obtained from the
\emph{original} second initial equation, not from a chosen value of $Y_2$.
\end{proof}

Here is a coefficient construction that fixes all the functionals in this
lemma independently of their implementation.  Let $\Phi_h$ be
\eqref{v3:functional}, let $z_1(Y)$ be \eqref{v3:z1}, and let
$\eta_1(Y)=(0,V(p))$ be the flat unit-multiplier auxiliary first jet.
For a general multiplier $\nu$, put
\[
 \eta_0(\nu)=
 \bigl((-\tfrac12\curl_\delta\beta,\delta n),
                         \delta_i\beta_j+\delta_j\beta_i\bigr).
\]
The ordered pair in the first slot lists rotation and boost coefficients
of $A_0$.  The flat equations give
$G_0\nu+T_0^*\eta_0(\nu)=0$.
Then
\begin{align}
 H_3(Y)&=-[a^3]\Phi_h(az_1(Y),a\eta_1(Y);aY,e),\notag\\
 P_{\nu,2}(Y)&=-[a^2]\Phi_h(az_1(Y),\eta_0(\nu);aY,\nu).
 \label{v5:coefficient-functional}
\end{align}
For the second identity the normal stationary second coefficient obeys
$T_0z_2=-\tau_2$.  Its contribution $G_0\nu\cdot z_2$ is
$\eta_0(\nu)\cdot\tau_2$, exactly the auxiliary term retained on the
right.  Its tangent part is annihilated by $G_0\nu$.  The first identity
is the stationary cubic envelope.  Through these degrees the geometric
coframe and velocity jets and the cubic electric and magnetic terms in
\eqref{v3:functional} are complete.  Quartic spatial terms have too high a
degree; quartic electric terms have degree four when $A_0$ is a first field
jet and degree three, not two, when it is $\eta_0$.  Thus they cannot enter
either displayed coefficient.  This degree check is not permission to
truncate the actual stationary or evolution equations.

\section{The drift derivative has a translation cokernel and free momentum slots}
\label{v5:drift-linearization}

Let $d_{2,h}(Y)$ denote the quadratic drift covector on the complete weak
shift space $\mathcal W_h=\ker\curl_\delta$.  For a fixed balanced first jet
$Y=(U,-sI)$ define
\begin{equation}
 \mathcal T_Y(Z):=Dd_{2,h}(Y)[Z]
                    =D^2\mathsf B_{W,h}(0)[Y,Z].
 \label{v5:T}
\end{equation}
This $\mathcal T_Y$ is not the stationary auxiliary constraint map
$\mathcal T_X$ of \eqref{v3:split}.
Write $\mathcal W_h=\mathcal W_c\oplus\mathcal W_g$, where
$\mathcal W_c$ is the three-dimensional constant-shift space and
$\mathcal W_g$ the mean-zero phase gradients.

\begin{lemma}[Constant weak quadratic drift is identically zero]
\label{v5:constant-B2}
For every odd cutoff, every symmetric $Y=(u,p)$ and every constant $c$,
\[
 d_{2,h}(Y)[c]=0.
\]
In particular $(\mathcal T_Y Z)[c]=0$ for every $Z$, whether or not it
satisfies a linear initial constraint.
\end{lemma}
\begin{proof}
A constant shift has $P_{c,1}=0$ and
$P_{c,2}=\ip p{c^i\delta_i u}_h$, by the full weak quadratic identity.
Its linear Hamiltonian vector field is
$(c^i\delta_i u,c^i\delta_i p)$.  Each constant-coefficient self-adjoint
operator in \eqref{v5:H2} commutes with the skew operator $c^i\delta_i$.
The derivative of its quadratic form in this direction is zero.  Hence
$\{P_{c,2},H_2\}=0$.  There is no $P_{c,1}$ contribution from $H_3$.
Differentiate the resulting identity to obtain the last assertion.
\end{proof}

\begin{lemma}[Momentum changes do not change the drift derivative at scalar momentum]
\label{v5:momentum-zero}
For $Y=(U,-sI)$ with arbitrary symmetric $U$ and constant $s$,
\begin{equation}
 \boxed{\mathcal T_Y(0,q)=0\quad\text{for every symmetric }q.}
 \label{v5:pzero}
\end{equation}
Consequently the preparation control acts through the configuration slot
alone.  The statement is exact at every odd cutoff.
\end{lemma}
\begin{proof}
For a weak test $\beta$ set $k=\delta_i\beta_j+\delta_j\beta_i$.
The full quadratic identity and the stationary cubic envelope give
\begin{equation}
 d_{2,h}(U,p)[\beta]
 =-D_UH_3(U,p)[k]
   +\ip{\mathcal F(U)}{\mathcal G_\beta U}_h
   +\ip p{\mathcal G_\beta V(p)}_h.
 \label{v5:general-weak-B2}
\end{equation}
The flat stationary connection variation in $(k,0)$ is zero for curl-free
$\beta$, so the tested cubic logarithmic derivative vanishes for arbitrary
$(U,p)$, just as in \cref{v4:full-drift}.  The only $p$-dependent geometric
cubic term is the ultralocal kinetic coefficient
\[
 H_{3,\mathrm{kin}}(U,p)
 =\ip U{2p^2-(\tr p)p}_h
   -\tfrac12\ip{\tr U}{\tr(p^2)-\tfrac12(\tr p)^2}_h.
\]
Its mixed derivative in $p=-sI$ and $q$ satisfies
\[
 D_pD_UH_{3,\mathrm{kin}}[k,q]
       =-s\ip kq_h+\tfrac s2\ip{\tr k}{\tr q}_h.
\]
On the other hand $V(-sI)=sI$, $\mathcal G_\beta(sI)=sk$, and the
inherited exact trace test gives
\[
 D_p\ip p{\mathcal G_\beta V(p)}_h[q]
       =-s\ip kq_h+\tfrac s2\ip{\tr k}{\tr q}_h.
\]
For example, obtain the latter by writing the derivative as
$\ip q{\mathcal G_\beta(sI)}_h+
\ip{-sI}{\mathcal G_\beta(2q-(\tr q)I)}_h$ and integrating its derivatives
by parts.  The two terms cancel with the opposite sign in
\eqref{v5:general-weak-B2}; its middle term is independent of $p$.
No phase product rule is used.
\end{proof}

For later reproduction, the configuration derivative is explicitly
\begin{equation}
 \begin{split}
 (\mathcal T_Y(Z,0))[\beta]
   ={}&-D_U^2H_3^{\mathrm g}(U,-sI)[Z,k]\\
     &+\ip{\mathcal F(Z)}{\mathcal G_\beta U}_h
                   +\ip{\mathcal F(U)}{\mathcal G_\beta Z}_h.
 \end{split}
 \label{v5:T-functional}
\end{equation}
The full-action extractor below retains the logarithms and independently
checks their cancellation in these tests.  The output vanishes on constant
shifts but need not vanish on nonconstant gradients.

\section{A full preparation-control minor on the balanced seed}
\label{v5:control}

This section uses the \emph{same} three-site first jet as
\eqref{v4:finite-seed}.  Put
\[
 \mathcal Z_u=\{Z:\delta_i\delta_jZ_{ij}
                          -\delta_i^2\tr Z=0\}.
\]
Its dimension is $6V-(V-1)=136$ at $V=27$.

\begin{theorem}[All nonconstant weak directions are controlled]
\label{v5:T-rank}
At the balanced three-site seed,
\begin{equation}
 \boxed{\mathcal T_Y:\mathcal Z_u\times\{0\}
             \longrightarrow\mathcal W_g^*\text{ is onto},\qquad
                  \rank\mathcal T_Y=26.}
 \label{v5:Trank}
\end{equation}
The rank remains 26 in a sufficiently small neighborhood within the nonzero
balanced polarization family, with its calibrated scalar momentum.
\end{theorem}
\begin{proof}
We give a finite characteristic-zero minor with its complete construction.
Use lexicographic sites and raw symmetric slots $(11,22,33,12,13,23)$,
retaining Frobenius weights in every scalar functional.  Finite residue
tables use unnormalized spatial sums; this common nonzero factor does not
affect a rank statement.  Let $C_1$ be the
$27\times162$ matrix of the scalar linear constraint.  Choose its 26 pivot
columns
\begin{equation}
 \begin{split}
 \mathcal I={}&(0,1,2,3,4,5,6,7,8,9,10,11,12,13,15,18,19,21,22,25,\\
              &\hspace{32mm}27,31,37,43,49,55).
 \end{split}
 \label{v5:Cpiv}
\end{equation}
Let $C_*$ be its first 26 rows on these columns.  The last row is minus
the sum of the others.  In the fixed specialization
$\sqrt3\mapsto56\in\mathbb F_{241}$,
\[
 \det C_*\equiv211\pmod{241}.
\]
This is also a rational, nonzero characteristic-zero minor.  List the
remaining 136 columns increasingly as $\mathcal J$.  Define the rational
kernel basis $Z_*$ by
\[
 (Z_*)_{\mathcal J,:}=I_{136},\qquad
 (Z_*)_{\mathcal I,:}=-C_*^{-1}(C_1)_{0:25,\mathcal J}.
\]
Then $C_1Z_*=0$ in characteristic zero.  This definition does not infer a
characteristic-zero nullspace from a modular zero.

Let $T$ be the $29\times162$ matrix obtained from
\eqref{v5:T-functional}, with the weak basis used in
\cref{v3:certificate-section}: three constants, followed by gradients of
point differences against the last site.  Select from $TZ_*$ the 26
nonconstant rows and the following 26 columns:
\[
 \mathcal K=(0,1,\ldots,23,25,28).
\]
The complete coefficient construction and residue arrays in the archive give
\begin{equation}
 \boxed{\det (TZ_*)_{g,\mathcal K}\equiv234\pmod{241}.}
 \label{v5:Tminor}
\end{equation}
All denominators of the original flat maps, coefficient extraction and the
specified kernel lift are units in this specialization.  Nonzero residue
therefore proves a nonzero characteristic-zero minor.  The upper rank 26
follows from \cref{v5:constant-B2}, not from a numerical or modular rank
threshold.  Continuity of this nonzero minor proves the neighborhood
assertion.  The calculation uses the drift derivative, not the source
$J_{2,h}$ or the derivative of $P_{\beta,2}$ alone.
\end{proof}

We next account for the initial constraints that a change of second jet
must preserve.  Let $\mathscr C_h=C_1^{\rm full}+C_2^{\rm full}+\cdots$ be
its homogeneous expansion.  Define
\begin{equation}
 \mathcal Z_Y^{(2)}=
 \{Z:\ D\mathscr C_h(0)Z=0,
                 \ P_0D C_2^{\rm full}(Y)[Z]=0\}.
 \label{v5:admissible-second}
\end{equation}
Here $P_0$ takes the four spatial means.

\begin{lemma}[Mean repair without loss of weak-drift control]
\label{v5:mean-repair}
Every $Z_u\in\mathcal Z_u$ has a momentum completion $Z_p$, using only
the four original balancing momentum directions, such that
$(Z_u,Z_p)\in\mathcal Z_Y^{(2)}$.  It satisfies
$\mathcal T_Y(Z_u,Z_p)=\mathcal T_Y(Z_u,0)$.
Consequently \cref{v5:T-rank} is still onto after imposing all the conditions
in \eqref{v5:admissible-second}.
\end{lemma}
\begin{proof}
Use the homogeneous scalar momentum and the three quadrature momentum
profiles $V_i$ of \cref{v4:mean}.  Each has zero phase divergence, so adding
them preserves the full linear constraints.  Their four mean derivatives
at the first jet form
\[
 \diag(3s,-\omega_hm_1,-\omega_hm_2,-\omega_hm_3),
\]
which is invertible.  It therefore uniquely cancels the four means caused
by $Z_u$.  At the displayed three-site specialization its determinant is
$144\pmod{241}$, independently checking the sign and normalization.
\Cref{v5:momentum-zero} shows that this momentum repair makes no change to
$\mathcal T_Y$.
\end{proof}

\begin{lemma}[Admissible second coefficients extend to exact preparations]
\label{v5:extend}
At the fixed three-site cutoff, let $X_0(a)=aY+a^2Y_2+O(a^3)$ be an
analytic exact preparation on the balanced chart.  For any fixed
$Z\in\mathcal Z_Y^{(2)}$, there is an analytic exact preparation with
second coefficient $Y_2+Z$ and the same first jet.  Finite-dimensional
families of such changes can be realized analytically in their parameters.
\end{lemma}
\begin{proof}
Use the retained range-and-mean initializer.  In its kernel coordinates,
change the free input by $a^2$ times the nonbalancing kernel part of $Z$.
The nonlinear range correction has derivative $O(a)$, so its change starts
at order $a^3$ and does not change this second kernel coefficient.
The normalized four mean equations have the invertible leading derivative
of \cref{v4:mean}.  Their order-$a$ equation restores exactly the balancing
momentum component of $Z$: its required condition is precisely
$P_0D C_2^{\rm full}(Y)[Z]=0$.
The analytic implicit-function theorem then solves the means and the range
at all higher orders using the original constraint map.  Any fixed finite
parameter value is admitted after reducing the amplitude interval; no
uniform radius for unbounded second coefficients is asserted.  The result
is an exact analytic curve, not merely a formal second-order solution.
\end{proof}

\section{The cubic translation drift does not depend on the second preparation}
\label{v5:translation-cubic}

For a constant shift $c$, the flat spatial response to that multiplier is
zero.  Its first auxiliary multiplier response at the unit cut is
\begin{equation}
 \eta_1^c(Y)=\left(c^iz_1(Y)_i,\ c^i\delta_i u\right).
 \label{v5:translation-aux}
\end{equation}
The first slot is the six-component temporal connection.  Substitution in
the differentiated flat saddle verifies this identity: constant $c$
commutes with every phase derivative, the spatial variation is zero by
\cref{v3:universal-flat}, and
$g_1+T_0^*\eta_1^c=0$.  In particular the temporal-connection term must
not be omitted when computing the next momentum coefficient.

The full cubic constant-shift row is therefore
\begin{equation}
 P_{c,3}(Y)=-[a^3\tau]\Phi_h
 \bigl(az_1(Y),a\eta_1(Y)+a\tau\eta_1^c(Y);\ aY,e+\tau(0,c)\bigr).
 \label{v5:P3}
\end{equation}
Stationarity removes the induced higher spatial connection coefficients
from this cubic envelope.  Both cubic logarithmic terms remain, including
the electric term paired with the nonzero \eqref{v5:translation-aux}.
Higher terms have field degree at least four in this expression.

\begin{theorem}[An exact formula for the preparation-independent cubic means]
\label{v5:cubic-means}
For every analytic exactly prepared curve with first jet $Y$, the cubic
constant-shift drift coefficient is
\begin{equation}
 \boxed{t_c(Y)=-DH_3(Y)[c^i\delta_iY]
                                +DP_{c,3}(Y)[F_1Y].}
 \label{v5:t}
\end{equation}
It is independent of $Y_2$ and of every higher preparation coefficient.
For the axial first jets of \eqref{v4:axial-jet}, at every odd $N\ge5$,
\begin{equation}
                    t_c(Y)=0\qquad\text{for all constant }c.
 \label{v5:t-fine}
\end{equation}
At $N=3$, \eqref{v5:t-fine} also holds for the balanced seed
\eqref{v4:finite-seed}.
\end{theorem}
\begin{proof}
The constant row has $P_{c,1}=0$, while its quadratic drift is identically
zero on all fields by \cref{v5:constant-B2}.  Its cubic Poisson coefficient
is $\{P_{c,2},H_3\}+\{P_{c,3},H_2\}$, which gives
\eqref{v5:t}.  The lower drift coefficients and all their derivatives are
zero.  Substitution of any prepared curve therefore adds no $Y_2$ or $Y_3$
term to this coefficient.

For completeness the fine-grid cancellation can be checked before any
inversion.  In the cubic unit Hamiltonian the rotation/boost grading and
\eqref{v3:functional} allow only monomials of types $u^3$ and $up^2$,
with the indicated derivatives and literal shifts on their factors.
The cubic constant-shift row \eqref{v5:P3} instead has types $u^2p$ and
$p^3$.  The same grading holds for the electric term after
\eqref{v5:translation-aux} is inserted; in particular it is not dropped.
These statements also follow term by term: spatial rotations are linear
in derivatives of $u$, boosts and the unit velocity are linear in $p$;
rotation action coefficients pair even numbers of boosts, and shift
coefficients pair odd numbers.

On the stated first jets the nonzero configuration frequencies are
$\{\pm e_1,\pm e_2,\pm e_3\}$ and $p$ is constant.  At degree three and
$N\ge5$, a spatially summed monomial can survive only with ordinary total
frequency zero, since no component of its total frequency can reach $N$.
The operator $c^i\delta_i$ on each input factor multiplies it by the
common scale $\omega_h$ times its signed unit coordinate; their sum is
zero on every surviving coefficient.  This proves the first term of
\eqref{v5:t} is zero even with its literal shifts.
For the second term, $F_{1,u}=sI$ and
$F_{1,p}=-\omega_h^2U/2$.  Differentiating the two permitted types of
$P_{c,3}$ produces $up^2$ or $u^3$.  A single axial configuration frequency
cannot sum to zero with two constants; three signed axial unit frequencies
cannot sum to zero either.  Thus that term also vanishes.  The argument is
about the actual finite coefficient and does not substitute a nonlinear
phase product rule.

The exceptional three-site aliases need a separate check.  The companion
characteristic-zero calculation evaluates the two terms of
\eqref{v5:t} using \eqref{v5:coefficient-functional} and \eqref{v5:P3}.
Replace $\sqrt3$ in the sampled sine and in $p=3\sqrt3 I$ by an indeterminate
$r$.  Each of the six resulting scalar expressions is a polynomial in $r$
of degree at most three.  Exact rational evaluations at
$r=0,1,-1,2$ give zero for each expression.  Cubic interpolation therefore
proves each polynomial is zero, before imposing $r^2=3$.
All array operations in this check are given by the displayed functional,
$\delta_i=3(S_i-S_i^{-1})$, the Lorentz bracket, and exact rational
arithmetic.  This is a characteristic-zero identity, not inference from
three modular zero entries.  The code and complete rational values are
included with the finite certificate.
\end{proof}

\section{The effective cubic source has a preparation invariant}
\label{v5:schur-invariant}

Fix a regular leading source and choose the active representatives and weak
basis of \cref{v4:certificate-section}.  A lapse representative may be
made mean-zero by subtracting its mean times $e$; this changes neither the
source nor the drift, since $J_he=0$ and
$\mathsf B_{h,e}=\{H_h,H_h\}=0$.  The active shift representatives need
not be orthogonal to the weak space.  Together with the weak basis they
still give a direct complement of the constant-lapse direction.
Write
\begin{equation}
 \Gamma_0=\begin{pmatrix}A&L\\L^*&W^*K_0^{-1}W\end{pmatrix},\qquad
 D=W^*K_0^{-1}W-L^*A^{-1}L.
 \label{v5:Schur}
\end{equation}
Here $A$ and $\Gamma_0$ are invertible on the verified branch; this is a
signed, not a positive, elimination.
Let $b_A$ be the active values of \eqref{v5:prepared-drift-formula} and
$b_3$ the weak cubic drift on a chosen exact preparation.
The first surviving normalized response of \cref{v4:remaining-rate} obeys
\begin{equation}
 (u_1)_W=D^{-1}\bigl(b_3-L^*A^{-1}b_A\bigr).
 \label{v5:uW}
\end{equation}
Let $R_c$ select the first three weak covector coordinates (constant tests),
and let $I_c$ insert the corresponding three multiplier coordinates.
Use $R_g,I_g$ for the other 26 coordinates at $N=3$.

\begin{theorem}[An invariant of all second preparations with the same first jet]
\label{v5:invariant}
The vector
\begin{equation}
 \boxed{\mathfrak c_Y
   :=R_c b_3-R_cL^*A^{-1}b_A
     =\bigl(t_{e_i}(Y)\bigr)_{i=1}^3-R_cL^*A^{-1}b_A}
 \label{v5:c}
\end{equation}
is independent of every second and higher exact preparation coefficient.
Every such preparation satisfies
\begin{equation}
                    R_cD(u_1)_W=\mathfrak c_Y.
 \label{v5:invariant-equation}
\end{equation}
Changing its second coefficient by
$Z\in\mathcal Z_Y^{(2)}$ changes $b_3$ by $\mathcal T_Y Z$ and leaves
$A,L,D,b_A$ unchanged.
\end{theorem}
\begin{proof}
The active coefficient is fixed by \cref{v5:prepared-B2}.
The leading active source depends only on the first jet.  The weak second
source does also, by \cref{v3:preparation-independent}.  Hence all the
blocks in \eqref{v5:Schur} are unchanged.
Taylor expansion of the weak drift on two prepared curves with the same
first jet gives a cubic difference
$D^2\mathsf B_{W,h}(0)[Y,Z]=\mathcal T_Y Z$; its first derivative is zero.
The constant part of this difference is zero by \cref{v5:constant-B2},
and its original cubic value is \cref{v5:cubic-means}.  This proves the
invariance and the change formula.  Eliminate the active coordinates from
$\Gamma_0u_1=(b_A,b_3)$ to obtain \eqref{v5:uW} and
\eqref{v5:invariant-equation}.  In particular zero \emph{raw} cubic means
need not give zero effective cubic means.
\end{proof}

\section{The actual balanced seed has a nonzero invariant}
\label{v5:invariant-certificate}

All data in this section again use \eqref{v4:finite-seed}.  Every polynomial
coefficient in the certificate is taken from the original scalar action.
The 108 prepared quadratic values are evaluated using
\eqref{v5:prepared-drift-formula}; its lapse correction is retained.  The
active values use exactly the 78 columns of \eqref{v3:pivots}.

\begin{theorem}[A characteristic-zero constant-shift obstruction]
\label{v5:nonzero-invariant}
At the balanced three-site seed, $\mathfrak c_Y\ne0$.
More precisely, with the fixed embedding $\sqrt3\mapsto56$ in
$\mathbb F_{241}$, the common spatial-sum convention gives
\begin{equation}
 \begin{split}
 (t_{e_i}(Y))_{i=1}^3&=0\quad\text{in characteristic zero},\\
 R_c L^*A^{-1}b_A&\equiv(161,57,199)^T\pmod{241},\\
 \boxed{\mathfrak c_Y\equiv(80,184,42)^T\pmod{241}.}
 \end{split}
 \label{v5:c-residue}
\end{equation}
The constant-constant weak Schur block is also invertible:
\begin{equation}
 \boxed{\det D_{cc}\equiv79\pmod{241},\qquad
   D_{cc}=R_cDI_c.}
 \label{v5:Dcc}
\end{equation}
These nonvanishing statements persist on a smaller balanced polarization
neighborhood at this fixed cutoff.
\end{theorem}
\begin{proof}
The full-action coefficient evaluator obtains $H_3$ and $P_{\nu,2}$ from
\eqref{v5:coefficient-functional}, then applies the exact polynomial
polarizations
\[
 DQ(Y)[Z]=\tfrac12\{Q(Y+Z)-Q(Y-Z)\},
\]
for homogeneous quadratics, and
\[
 DC(Y)[Z]=\tfrac12\{C(Y+Z)-C(Y-Z)\}-C(Z),
\]
for homogeneous cubics.  These are polynomial identities, not
finite-amplitude numerical difference quotients.  The actual canonical
linear field in \eqref{v5:F1} is used throughout.  It gives $b_A$, with
52 nonzero residue entries on the retained active basis.  The matrix
certificate of \cref{v4:finite-rank} gives
$\det A=113$ and $\det D=56$ modulo 241.  Solve $Ax=b_A$ in that field
and evaluate $R_cL^*x$.  The full vectors and matrices in the archive give
the second line of \eqref{v5:c-residue}; the first is the independent exact
calculation in \cref{v5:cubic-means}.  Direct elimination of the specified
$3\times3$ block gives \eqref{v5:Dcc}.

All denominators are units at 241.  Nonzero residues therefore prove the
stated nonzero characteristic-zero vector and determinant.  The calculation
uses one spatial-sum convention for both target and Gram.  Dividing both
by $V$ restores the physical mean convention and leaves every solution
coefficient unchanged; in particular the rate and its nonvanishing are not
normalization artifacts.  The data and the calibrated scalar momentum are
analytic in the polarization parameters.  Nonvanishing and invertibility
therefore persist on a sufficiently small neighborhood, with positive lower
bounds on compact subsets.
\end{proof}

\begin{corollary}[An unavoidable initial rate pole, despite small electric curvature]
\label{v5:actual-pole}
For every analytic exact initial preparation with first jet
\eqref{v4:finite-seed}, and its own normalized exact-action continuation,
\begin{equation}
 \boxed{\|\partial_t\beta_h(a,0)\|_h\ge\frac{c_h}{|a|}
                   \quad\text{for sufficiently small }0<|a|,}
 \label{v5:shift-pole}
\end{equation}
where the threshold may depend on the preparation.  The leading coefficient
has a positive lower bound uniform over changes of the second coefficient.
The same conclusion holds on a smaller compact regular balanced
polarization neighborhood.  At the same cuts the literal electric and
spatial curvature reads retain $O_h(a)$ size along each analytic preparation.
Thus full anchored first-metric-jet shadowing of a small harmonic history
cannot follow from the initial electric estimate.
\end{corollary}
\begin{proof}
Equation \eqref{v5:invariant-equation} gives
\[
 \|(u_1)_W\|\ge
      \frac{\|\mathfrak c_Y\|}{\|R_cD\|}>0.
\]
The right side is independent of the second preparation.  On the actual
initial curve, \cref{v4:remaining-rate} gives
\[
 \lambda_t=-a^{-1}\mathcal U_W(u_1)_W+O_h(1),
\]
where $\mathcal U_W$ is the fixed injective synthesis into curl-free shift
arrays; its lapse slot is zero.  Its least singular value is positive at
the fixed cutoff.  The active contribution and the mean-lapse normalization
are only $O_h(1)$, so they cannot cancel this nonzero $a^{-1}$ coefficient.
This proves \eqref{v5:shift-pole}.  Continuity gives the compact-parameter
statement.  Varying an unbounded second coefficient need not preserve a
common remainder or amplitude threshold.

The initial stationary source has zero leading target and a regular Gram.
The argument of \cref{v4:initial-electric} therefore applies to each of
these analytic preparations, giving the stated $O_h(a)$ reads.  At the
unit-lapse, zero-shift cut,
$\partial_tg_{0i}=\gamma_{ij}\partial_t\beta^j$.
The spatial metric has a uniform positive margin as $a\to0$, so the
first metric time jet is unbounded in these anchored coordinates even
though the literal electric read is small.  No contradiction with a
bounded metric-jet curvature theorem is implied.
\end{proof}

\section{An exact initial normal form: remove the nonconstant weak rate}
\label{v5:normal-form}

The invariant prevents suppression of the whole rate, but does not prevent
removing its nonconstant weak part by a legitimate choice of initial
records.  This gives a more useful next target than repeating the same
unsolvable second-preparation system.

\begin{proposition}[The attainable first weak-rate coefficients]
\label{v5:affine-rate}
At the three-site seed the first weak-rate coefficients obtained by exact
second preparations form precisely the affine plane
\begin{equation}
 \boxed{\{u\in\R^{29}:R_cDu=\mathfrak c_Y\}.}
 \label{v5:affine-plane}
\end{equation}
It contains a unique vector supported only on constant shifts,
\begin{equation}
 u=I_c\upsilon_Y,\qquad
 \upsilon_Y=D_{cc}^{-1}\mathfrak c_Y\ne0.
 \label{v5:constant-normal}
\end{equation}
For the stated certificate,
$\upsilon_Y\equiv(206,141,227)^T\pmod{241}$.
\end{proposition}
\begin{proof}
The invariant gives necessity.  Conversely \cref{v5:T-rank,v5:mean-repair}
allow every change in the 26 nonconstant components of $b_3$, with its
constant components fixed.  \Cref{v5:extend} realizes each finite such
change by an analytic exact preparation.  Applying the invertible matrix
$D^{-1}$ gives exactly \eqref{v5:affine-plane}.
For constant support the condition reduces to
$D_{cc}\upsilon_Y=\mathfrak c_Y$, which has the unique displayed nonzero
solution by \cref{v5:nonzero-invariant}.  The supplied exact residue solve
gives its three entries.
\end{proof}

Let $\mathbb P_{\nabla,0}$ be the physical orthogonal projection of a
shift array onto the mean-zero curl-free shifts.  It is an observation used
to specify initial data, not a projection imposed on the dynamical vector
field.

\begin{theorem}[Exact preparation with zero initial nonconstant weak shift velocity]
\label{v5:exact-rate-preparation}
At $N=3$, on a sufficiently small regular balanced first-jet neighborhood,
there are analytic exact initial preparations $X(a)$ such that their own
normalized original-action initial rate satisfies
\begin{equation}
 \boxed{\mathscr C_h(X(a))=0,\qquad
  X(a)=aY+O_h(a^2),\qquad
           \mathbb P_{\nabla,0}\partial_t\beta_h(a,0)=0.}
 \label{v5:exact-rate-zero}
\end{equation}
Their leading weak pole is the spatially constant vector
\begin{equation}
 \partial_t\beta_h(a,0)=-a^{-1}\upsilon_Y+O_h(1),
                         \qquad \upsilon_Y\ne0.
 \label{v5:exact-rate-form}
\end{equation}
The statements concern the initial cut.  They do not assert preservation
of the additional rate condition at later physical times.
\end{theorem}
\begin{proof}
Choose the 26 configuration controls in \eqref{v5:Tminor}, repair their
means as in \cref{v5:mean-repair}, and realize their affine second-jet
changes by the analytic original initializer.  Denote the resulting exact
initial data by $X(a,z)$, $z\in\R^{26}$, with common first jet $Y$.
Their graded matrix $\Gamma_h(a,z)$ has the same invertible limit
$\Gamma_0$.  Their normalized target has an analytic factor $a$ by
leading source neutrality.  Hence
\[
 \widetilde u(a,z)=a^{-1}\Gamma_h(a,z)^{-1}\widehat b_h(a,z)
\]
is analytic through $a=0$.  Its value at zero is the coefficient $u_1(z)$,
not an assumed small actual rate.

Use weak coordinates consisting first of constants, then of the specified
gradient basis.  Let $Q_g$ extract physical mean-zero gradient coordinates
from a full multiplier; thus $Q_g\mathcal U_g=I$,
$Q_g\mathcal U_c=0$.  It can have a nonzero value
$J_A=Q_g\mathcal U_A$ on the chosen nonorthogonal active representatives.
The exact initial rate gives the analytic equation
\begin{equation}
 aQ_g\lambda_t
   =-R_g\widetilde u_W(a,z)-aJ_A\widetilde u_A(a,z).
 \label{v5:scaled-gradient}
\end{equation}
The mean-lapse normalization term has zero $Q_g$ observation.
By \cref{v5:affine-rate}, there is a finite control $z_0$ at $a=0$ for
which $\widetilde u_W=I_c\upsilon_Y$; hence the right side is zero.
Its derivative in $z$ is
\[
 -R_gD^{-1}I_g\,(\mathcal T_YZ_*)_{g,\mathcal K}.
\]
The first factor is invertible because $D$ and $D_{cc}$ are invertible:
its inverse is the Schur complement of $D_{cc}$ in $D$.
The second factor is invertible by \eqref{v5:Tminor}.
In the common certificate its determinant, apart from the irrelevant
26-dimensional minus sign, is
\[
 (79/56)\,234\equiv201\pmod{241},
\]
which separately verifies this derivative.  The analytic implicit-function
theorem gives $z(a)$ through $z_0$ making
\eqref{v5:scaled-gradient} identically zero.  For $a\ne0$ this is
exactly the last equation in \eqref{v5:exact-rate-zero}.  All initial
constraints were already solved by the original initializer.
At $a=0$ the weak coefficient is the nonzero constant vector of
\eqref{v5:constant-normal}; expanding the exact rate gives
\eqref{v5:exact-rate-form}.  All inverses here have a verified nonzero
limit at the fixed cutoff.  No cutoff-independent radius or lifespan is
inferred from the implicit-function argument.
\end{proof}

The theorem chooses \emph{initial fields} so that their own exact rate has a
specified component zero.  It neither overwrites that rate after solving
nor adds a tangency projection to the Hamiltonian equations.  It also does
not set the transverse part of the $O_h(1)$ active rate to zero.
Although the surviving leading shift is spatially constant, a
cutoff-dependent time-dependent translation is not automatically an allowed
change of the anchored operational reads.  In particular its compatibility
with the finite cyclic products and literal-link observations must be
proved before it can be used in a continuum comparison.

\section{What is decided and what a cofinal continuation must now test}
\label{v5:status}

This body establishes an actual next-order obstruction and an actual
initial-data normal form on the balanced regular three-site branch.
Second preparation controls every nonconstant weak target but cannot
change the three effective translation components.  Those components are
nonzero for the unchanged action, even though their raw cubic drift values
are zero.  It is the signed active--weak coupling, not an invented external
forcing, that supplies them.  Suppression of the initial electric
coefficient therefore does not imply suppression of the original
lapse--shift first jet.  The matrix examples in V1 are not used as evidence
for this action-specific conclusion.

The positive result is an analytic exact preparation removing the complete
initial mean-zero curl-free shift velocity.  This isolates the remaining
leading pole into constant shifts and makes the next comparison problem
more specific.  It does not prove that those shifts are operationally
unobservable or that the prepared condition persists under evolution.
A rigorous coordinate-comparison theorem, or a different first-jet family
whose effective vector \eqref{v5:c} vanishes, would be an upstream route.
Further changes of $Y_2$ alone cannot remove \eqref{v5:shift-pole} on this
fixed first-jet neighborhood.

Several distinctions remain essential.  The universal quadratic
translation identity and the fine-grid cubic translation cancellation are
coefficient theorems at every stated cutoff.  The surjective preparation
minor, the complete regular Gram and the nonzero effective vector are
certified at $N=3$.  They do not establish an all-cutoff rank family, a
nonzero cofinal limit of that vector, or a divergence theorem as $h\to0$.
Likewise $O_h(a)$ electric curvature and an $a^{-1}$ fixed-cutoff rate
pole can coexist with a quite different joint amplitude/cutoff behavior.
The actual higher sources, nonlinear Hessian and inclusion, differentiated
response, curvature identification, and a common physical interval remain
necessary for a raw-action Einstein continuum theorem.  None is replaced
by an appeal to a small unpreconditioned drift or a successful finite
minor.

\paragraph{Reproducibility and checks.}
The package provides the complete drift derivative, the rationally defined
linear-constraint kernel, the selected 26-dimensional minor, all 108
prepared quadratic drift values, the active solve, the three cubic
translation coefficients, the Schur invariant and the constant-rate normal
form.  The coefficient functions retain the full action terms needed by
field degree.  Independent checks verify the quadratic Hamiltonian's
canonical gradient, the weak quadratic constraint on unrestricted fields,
the zero momentum derivative, new directional polarizations of the drift,
and the full signed-Gram solve for the normal form.  The exceptional
translation coefficients are also evaluated in exact rational arithmetic
and cubic interpolation, independently of modular zero tests.

The 26 nonconstant cubic drift values of an arbitrary baseline nonlinear
preparation have not been tabulated.  They enter the preparation theorem as
the actual analytic right-hand side of the original equations; the verified
onto derivative solves for them without assuming they vanish.  The theorem
is an existence and local uniqueness result for the selected controls, not
a reported numerical nonlinear initial root.

The modular arithmetic is reduced between operations.  The maximum
matrix dimension in these calculations is 162, the prime is 241, and the
implementation's products and finite sums stay within signed 64-bit integer
bounds; the rational translation check uses arbitrary-precision fractions.
The supplied scripts re-create the new arrays from the scalar coefficient
functional.  They do not perform a nonlinear exact-action spacetime
integration, an all-cutoff inverse enclosure, or a proof-assistant
formalization.  The four inherited research bodies are preserved byte for
byte, with fingerprints in the accompanying manifest.
% END IMMUTABLE EXACT-ACTION BRIDGE V5 BODY

% Future iterations append here.
% BEGIN IMMUTABLE EXACT-ACTION BRIDGE V6 BODY
\clearpage
\renewcommand{\thesection}{\arabic{section}}
\renewcommand{\theHsection}{bridgeV6.\arabic{section}}
\setcounter{section}{42}
\section{Uniform shifts: a change of frame is not a finite symmetry}
\label{v6:context}

The remaining inverse-amplitude rate of \cref{v5:exact-rate-preparation} is
spatially constant to leading order.  This makes a change of frame a natural
comparison, but not an automatic gauge identification.  The finite
Hamiltonian uses the original phase derivatives, cyclic products and ordered
plaquette logarithms.  None of these is silently replaced by its continuum
counterpart here.  In particular the constant-shift momentum has a cubic
term even on spatially constant canonical records.  A phase-product defect
alone cannot explain that term.

This body calculates the complete cubic constant-shift row on the balanced
axial family at every odd cutoff.  Its nonzero leading term is an ordered
magnetic-plaquette contribution of order $h$.  An exact time-dependent
canonical change of variables then isolates the surviving shift acceleration
without altering the action or its multiplier equations.  At the prepared
three-site cut it reduces the contribution of the $a^{-1}$ pole to a
nonzero contribution of order $a$ in the spatial canonical acceleration.
That phase-generated change of frame is distinguished from a genuine
translation of the interpolated spacetime metric.  Their second-jet
difference has an exact $|\omega_h-2\pi|$ coefficient; their third-jet
difference retains the as yet uncontrolled time derivative of the shift
acceleration.  Finally, transported products specify exactly how the
original literal-link observations are preserved, and a full-frequency
estimate bounds their departure from the untransported products.

\paragraph{Inputs, nonduplication and interpretation.}
The stationary coefficient functional, flat first connection, and cubic
constant-shift envelope are the inherited
\cref{v3:functional,v5:coefficient-functional,v5:P3}.  The original-action
initial preparation and its normal form are
\cref{v4:preparation,v5:exact-rate-preparation}.  Their finite-rank and
preparation-control proofs are not repeated.  The effective vector
$\mathfrak c_Y$ of \cref{v5:invariant} is not identified with the cubic
momentum row calculated below: the former includes the signed active--weak
elimination, whereas the latter is an off-shell action coefficient.
No first-jet root of $\mathfrak c_Y$ is asserted.

Time-dependent canonical transformations, finite algebra isomorphisms and
Fourier convolution estimates are established tools.  The new
original-action content is the explicit cubic coefficient and its origin,
the resulting nonvanishing test, and the quantitative use of the prepared
rate in the two distinct initial-frame comparisons.  The general issue of
broken finite diffeomorphism symmetry is not new \cite{v6:BD}; that reference
supplies context, not any of the source identities proved here.  All
finite-amplitude remainder constants below retain an $h$ subscript unless
uniformity is explicitly proved.  No new common physical lifespan is
inferred from a change of coordinates.

\section{The complete cubic constant-shift row on the balanced family}
\label{v6:cubic}

Use $N=2m+1\ge3$, $h=N^{-1}$ and the normalized spatial-mean pairing.
Let $T_j,P_j$ be the transverse tensors of \cref{v4:finite-seed}: on the two
coordinates $k<\ell$ different from $j$, $T_j$ has its $k\ell$ and $\ell k$
entries equal to one, while $P_j$ has entries $+1,-1$ on $kk,\ell\ell$.
Thus both have squared Frobenius norm two and they are orthogonal.  Put
\begin{equation}
 \begin{split}
 U&=\sum_{j=1}^3 r_j\{T_j\cos(2\pi x_j+\phi_j)
                 +\sigma_jP_j\sin(2\pi x_j+\phi_j)\},\\
 p&=p_0I,\qquad r_j\in\R,\quad \sigma_j\in\{-1,1\},\\
 \omega_h&=2h^{-1}\sin(\pi h),\quad c_h=\cos(2\pi h).
 \end{split}
 \label{v6:balanced}
\end{equation}
The scalar $p_0$ may initially be arbitrary.  For nonlinear compatible
preparation the retained mean calibration sets
$p_0^2=\omega_h^2\sum_jr_j^2/3$, with one of its nonzero sign branches.
The notation $p_0$ here is a scalar momentum, not a constant-mode projector.

Write
\[
 P_{h,c}(X)=\partial_\lambda\mathscr H_h(X,e)[(0,c)],\qquad
 G_{h,c}(u,p)=\ip p{c^i\delta_{i,h}u}_h
\]
for a spatially constant shift $c\in\R^3$.  The inherited quadratic
identity says $P_{h,c,2}=G_{h,c}$.  The following concerns the next
homogeneous field coefficient, not its Poisson drift.

\begin{theorem}[An exact all-cutoff cubic translation coefficient]
\label{v6:P3}
For the unchanged action and \eqref{v6:balanced}, at every odd cutoff,
\begin{equation}
 \boxed{
 P_{h,e_i,3}(U,p_0I)
 =-\frac{hp_0^3}{4}
  -\frac{hp_0\omega_h^2}{8}(1-c_h)\sum_{j\ne i}r_j^2 .}
 \label{v6:P3-formula}
\end{equation}
The result is independent of the phases and handedness signs.  On this
family the geometric and electric cubic contributions to this integrated
row each vanish.  The displayed value is the complete cubic magnetic
contribution.  This cancellation is not asserted for arbitrary field
records or other multiplier tests.

For the unit balanced seed with $r_j=1$ and $p_0=\omega_h$,
\begin{equation}
 \boxed{P_{h,e_i,3}(Y_h)
       =-\frac{h\omega_h^3}{4}(2-c_h)
       =-2\pi^3h+O(h^3),\qquad i=1,2,3.}
 \label{v6:P3-unit}
\end{equation}
In particular this translation coefficient is not $O(h^2)$.  For any fixed
nonzero balanced radii and either calibrated scalar-momentum sign,
$P_{h,e_i,3}$ has the opposite sign to $p_0$ and has magnitude bounded
below by $h|p_0|^3/4$.
\end{theorem}
\begin{proof}
We evaluate the original envelope \eqref{v5:P3}, with its auxiliary
response $\eta_1^c=(c^iz_{1,i},c^i\delta_iU)$ retained.  The Lorentz bracket
in the coefficient convention of \cref{v3:functional} is
\[
 [(r,b),(r',b')]
   =(-r\mathbin\times r'+b\mathbin\times b',
                    -r\mathbin\times b'-b\mathbin\times r').
\]
At a homogeneous scalar momentum the spatial first connection is
$z_{1,i}=-(p_0/2)B_i$.  For an ordered constant plaquette its cubic
logarithm is
\[
 \tfrac12[z_{1,i}+z_{1,j},[z_{1,i},z_{1,j}]].
\]
With $b=-p_0/2$, this is the boost
$b^3(e_j-e_i)/2$.  The shift bivector is
$2(c^je_i-c^ie_j)$, so their scalar product is
$-b^3(c^i+c^j)$.  The magnetic factor is $h$, and passage from the
stationary functional to the Hamiltonian changes the sign.  Summing over
the three ordered pairs gives
$2hb^3\sum_i c^i=-hp_0^3\sum_i c^i/4$.
The electric double commutator at this homogeneous record has zero
projection on its $i$th boost coefficient, and the geometric part has no
cubic constant-field term.  This proves the constant term in
\eqref{v6:P3-formula} directly from the original plaquettes.

For clarity, the remaining finite coefficient contractions are recorded
explicitly.  Keep one unit profile in coordinate $j$, let the shift test be
$e_i$, and subtract the homogeneous scalar-momentum value.  Write
$z=\exp(2\pi\mathrm i h)$.  The normalized zero-frequency contributions
of the three parts of the original cubic envelope are
\begin{equation}
\begin{array}{c|ccc}
 &\text{geometric}&\text{electric cubic}&\text{magnetic cubic}\\ \hline
 j=i&0&0&0\\
 j\ne i&0&0&\displaystyle
          \frac{hp_0\omega_h^2}{16}(z+z^{-1}-2).
\end{array}
 \label{v6:contraction-table}
\end{equation}
Here is a finite coefficient procedure for the table.  In
\eqref{v5:P3} substitute
\[
 z_{1,ka}^{R}=-\tfrac12\epsilon_{abc}\delta_bU_{kc},\qquad
 z_{1,ka}^{B}=-\tfrac{p_0}{2}\delta_{ka},\qquad
 V(p_0I)=-p_0I.
\]
Use $e=I+U/2-U^2/8$ and
$\dot e=w/2-(Uw+wU)/8+(U^2w+UwU+wU^2)/16$ at the required field
degrees.  A profile in axis $j$ has Fourier coefficients
$(T_j-\mathrm i\sigma_jP_j)e^{\mathrm i\phi_j}/2$ and its conjugate.
Its phase derivative multiplies these by $\pm\mathrm i\omega_h$;
$S_j$ multiplies them by $z^{\pm1}$; the other shifts leave them fixed.
Retain the coefficient of $a^3\tau$ and pair the two opposite profile
frequencies.  For the magnetic part compute the four-factor logarithm by
\[
 \begin{split}
 L_1^{+}&=L_1+x,\qquad
 L_2^{+}=L_2+\tfrac12[L_1,x],\\
 L_3^{+}&=L_3+\tfrac12[L_2,x]
       +\tfrac1{12}\{[L_1,[L_1,x]]+[x,[x,L_1]]\},
 \end{split}
\]
starting from zero and taking successively
$x=z_{1,i},S_i z_{1,j},-S_j z_{1,i},-z_{1,j}$.
These substitutions give the entries of \eqref{v6:contraction-table};
the factors $e^{\pm\mathrm i\phi_j}$ cancel in each pair, and the
terms odd in $\sigma_j$ cancel between the paired frequencies.
Equivalently the surviving transverse coefficient is
$hp_0\omega_h^2(z-1)^2/(16z)$.  Symbols for derivatives of the doubled
frequency cancel before any value for that symbol is inserted.  Thus no
unfolded second-frequency convention is needed for this calculation.
The accompanying symbolic coefficient script performs precisely these
contractions with the higher derivative symbols left independent.

Rotation--boost grading permits only $U^2p$ and $p^3$ in this row, as in
\cref{v5:cubic-means}.  Opposite frequencies on different axes cannot sum
to zero; therefore the cross-axis terms integrate to zero.  No nonzero
sum of two unit axial frequencies is a multiple of any odd $N\ge3$.
The same zero-frequency calculation consequently applies to every stated
cutoff, including $N=3$.  Multiply the single-profile contributions by
$r_j^2$, add the homogeneous term, and use
$z+z^{-1}-2=-2(1-c_h)$.  This proves the formula.
Finally $\omega_h=2\pi-\pi^3h^2/3+O(h^4)$ and $1-c_h=O(h^2)$ give
\eqref{v6:P3-unit} and the sign and lower bounds.
\end{proof}

\begin{corollary}[The balanced finite seed is not translation-covariant at cubic order]
\label{v6:finite-P3}
At $N=3$ on the first jet used in \cref{v5:nonzero-invariant}, the
spatial-mean cubic row is $-135\sqrt3/8$ for each constant coordinate
shift.  In the spatial-sum convention it is $-3645\sqrt3/8$, whose
residue is $31$ under $\sqrt3\mapsto56\in\mathbb F_{241}$.
For the nonzero constant-rate vector $\upsilon_Y$ of
\eqref{v5:constant-normal},
\begin{equation}
 P_{h,\upsilon_Y,3}(Y)\equiv201\pmod{241}
 \label{v6:Pnu}
\end{equation}
in that same spatial-sum convention, and hence its real value is nonzero.
\end{corollary}
\begin{proof}
Substitute $h=1/3$, $\omega_h=3\sqrt3$, $c_h=-1/2$ in
\eqref{v6:P3-unit}.  The source's exact rate coefficient has residues
$(206,141,227)$, whose sum is $92$ modulo 241.  The row is linear in the
constant test, so $31\cdot92\equiv201$.  All denominators in the
inherited constant-rate solve are units at this prime; a nonzero residue
therefore proves characteristic-zero nonvanishing.  The common
normalization of target and Gram leaves $\upsilon_Y$ unchanged when the
spatial mean replaces the spatial sum.
\end{proof}

The homogeneous scalar-momentum calculation is an off-shell test of an
action identity; it is not an additional exact vacuum initial datum.
The balanced first jet, in contrast, has the nonlinear compatible
preparations already proved.  Neither calculation supplies the effective
invariant $\mathfrak c_{Y,h}$ at a new cutoff.  In particular, the
nonvanishing of \eqref{v6:P3-formula} must not be substituted for an
all-cutoff nonvanishing or vanishing theorem about that different vector.

\section{Exact canonical transport with all multiplier equations retained}
\label{v6:canonical}

Let $D_{h,c}=c^i\delta_{i,h}$ act componentwise on the spatial canonical
pair.  It is skew-adjoint and its components commute.  For a real path
$q(t)\in\R^3$ put
\[
 U_{h,q}=\exp(q^i\delta_{i,h}),\qquad
 \mathcal U_{h,q}(u,p)=(U_{h,q}u,U_{h,q}p).
\]
This is a real orthogonal canonical transformation.  It preserves the
constant arrays and every Fourier Sobolev norm.  It is generally not a
permutation of sites, an order-preserving stochastic relabeling, or an
automorphism of the pointwise coefficient product.

\begin{proposition}[The exact moving-frame action]
\label{v6:canonical-action}
For a prescribed $C^1$ path $q$, set $X=\mathcal U_{h,q}\widetilde X$.
The original canonical action becomes exactly
\begin{equation}
 \int\{\ip{\widetilde p}{\dot{\widetilde u}}_h
              -\widetilde H_h(t,\widetilde X,\lambda)\}\,\dd t,
 \qquad
 \widetilde H_h
 =\mathscr H_h(\mathcal U_{h,q}\widetilde X,\lambda)
                         -G_{h,\dot q}(\widetilde X).
 \label{v6:Htilde}
\end{equation}
Original solutions and transformed solutions are in one-to-one
correspondence on their existing interval.  The transformed multiplier
rows are exactly the original rows evaluated at
$\mathcal U_{h,q}\widetilde X$; the same holds for every original observable
when that observable is pulled back as well.
\end{proposition}
\begin{proof}
Orthogonality and $\dot U_{h,q}=U_{h,q}D_{h,\dot q}$ give
\[
 \ip p{\dot u}_h
 =\ip{\widetilde p}{\dot{\widetilde u}}_h
                  +\ip{\widetilde p}{D_{h,\dot q}\widetilde u}_h.
\]
Substitution proves the action identity with no equation discarded.
The transformation is invertible and canonical, so variation in the
canonical pair gives the equivalent time-dependent Hamiltonian equations.
During every variation the prescribed path $q$ is held fixed.
Consequently $\partial_\lambda G_{h,\dot q}=0$ and the multiplier rows
are precisely the pulled-back original rows.  An observable
$\mathcal R_h(X,\lambda)$ is represented by
$\mathcal R_h(\mathcal U_{h,q}\widetilde X,\lambda)$, not by an
unjustified substitution of $\widetilde X$ into its old formula.
\end{proof}

For comparison with an already given exact history one may choose
$\dot q(t)=\overline\beta_h(t)$, its spatial-mean shift, and $q(0)=0$.
This uses the past and present history, not a future Einstein solution.
It does not make $q$ a varied function of $\lambda$ in
\eqref{v6:Htilde}.  Substituting $\dot q=\overline\beta$ \emph{before}
varying and subtracting the resulting momentum from the original multiplier
row would change the variational problem.

\begin{proposition}[The residual action is cubic, not identically zero]
\label{v6:residual-H}
On a fixed-cutoff analytic chart with a spatially constant shift $c$,
\begin{equation}
 \mathscr H_h(X,e+(0,c))
   =H_h(X)+G_{h,c}(X)+\mathcal R_h(X,c),\qquad
 \mathcal R_h(X,0)=0.
 \label{v6:Hdecomp}
\end{equation}
For $c$ in a sufficiently small fixed compact chart and $j=0,1,2$,
\[
 \|D_X^j\mathcal R_h(X,c)\|
             \le C_h|c|\|X\|^{3-j}.
\]
The derivative $D_c\mathcal R_h(X,0)[c]$ has cubic term
$P_{h,c,3}(X)$.  In particular the residual cannot be deleted, even on
homogeneous canonical records.  When $\dot q=c$, the right side of
\eqref{v6:Htilde} is
\[
 H_h(\mathcal U_{h,q}\widetilde X)
       +\mathcal R_h(\mathcal U_{h,q}\widetilde X,c),
\]
not in general $H_h(\widetilde X)$.
\end{proposition}
\begin{proof}
The flat stationary Legendre calculation with constant lapse one and shift
$c$ gives the quadratic Hamiltonian $H_2+G_{h,c}$: the velocity enters as
$w-c^i\delta_i u$, and the canonical Legendre transform restores
$G_{h,c}$.  The integrated constant and linear field terms vanish.  The
full analytic saddle has the same quadratic coefficient, since the
logarithmic remainder starts at degree three.  Subtracting those terms
defines $\mathcal R_h$.  Its first three field Taylor coefficients through
quadratic order vanish, and it vanishes for $c=0$.  Taylor's integral
formula in $X$ and then the fundamental theorem of calculus in $c$ give
the stated bounds.  Differentiation at $c=0$ gives $P_{h,c}-G_{h,c}$.
For $u=0$, $p=ap_0I$, phase transport fixes $X$ and $G_{h,c}(X)=0$, but
\cref{v6:P3} gives a nonzero cubic derivative whenever
$p_0\sum_i c^i\ne0$.  Thus even translation of a homogeneous array
cannot remove this off-shell action term.  Finally $G_{h,c}$ is invariant
under $\mathcal U_{h,q}$ by skew-adjointness and commutation; inserting
\eqref{v6:Hdecomp} in \eqref{v6:Htilde} proves the last identity.
\end{proof}

\section{The initial constant-rate pole in a phase-canonical frame}
\label{v6:initial}

Write $\mathsf F_h(X,\lambda)=\mathbb J\nabla_X\mathscr H_h(X,\lambda)$
for the actual canonical vector field in the physical Frobenius pairing,
and $\mathsf X_K=\mathbb J\nabla_XK$ for a Hamiltonian vector field.
Consider any existing exact continuation through one of the normal-form
initial records of \cref{v5:exact-rate-preparation}.  At the fixed certified
cutoff it satisfies
\begin{equation}
 \begin{gathered}
 X_a(0)=aY+O_h(a^2),\qquad \lambda_a(0)=e,\\
 c_a:=\overline{\dot\beta_a(0)}
       =-a^{-1}\upsilon_Y+O_h(1),\qquad
 \dot\lambda_a(0)-(0,c_a)=O_h(1).
 \end{gathered}
 \label{v6:normal-input}
\end{equation}
All the assertions below concern its initial derivatives.  Its local
existence interval need not be uniform in $a$ or $h$.

\begin{theorem}[The pole contributes only order amplitude to transformed acceleration]
\label{v6:initial-comparison}
Choose $q(0)=0$, $\dot q=\overline\beta_a$ and put
$\widetilde X_a=\mathcal U_{h,-q}X_a$.  At the initial cut define
\[
 \mathcal A_a^0
 =D_X\mathsf F_h(X_a,e)\mathsf F_h(X_a,e)
       +D_\lambda\mathsf F_h(X_a,e)[\dot\lambda_a-(0,c_a)].
\]
Then the exact identity and its amplitude expansion are
\begin{align}
 \ddot{\widetilde X}_a(0)-\mathcal A_a^0
   &=\mathsf X_{P_{h,c_a}-G_{h,c_a}}(X_a(0)),
       \label{v6:exact-acceleration}\\
   &=-a\,\mathsf X_{P_{h,\upsilon_Y,3}}(Y)+O_h(a^2).
       \label{v6:acceleration-order}
\end{align}
Consequently this constant-rate contribution is $O_h(a)$, although
$c_a=O_h(a^{-1})$.  The full transformed initial canonical acceleration
is bounded as $a\to0$.  It is not asserted to be $O_h(a)$: the retained
nonconstant multiplier rates in $\mathcal A_a^0$ can have nonzero bounded
limits.

On the actual three-site balanced normal form the coefficient in
\eqref{v6:acceleration-order} is nonzero.  Thus for sufficiently small
nonzero $a$,
\begin{equation}
 c_h'|a|\le
 \|\ddot{\widetilde X}_a(0)-\mathcal A_a^0\|
                         \le C_h'|a|.
 \label{v6:sharp-acceleration}
\end{equation}
\end{theorem}
\begin{proof}
Initially $\dot q=0$, so twice differentiating
$\widetilde X=\mathcal U_{h,-q}X$ gives
$\ddot{\widetilde X}=\ddot X-D_{h,c_a}X$.
Differentiate the actual canonical equation to obtain
$\ddot X=D_X\mathsf F_h\mathsf F_h+D_\lambda\mathsf F_h\dot\lambda$.
The uniform-shift part of the latter derivative is
$\mathsf X_{P_{h,c_a}}$.  Direct variation of
$G_{h,c_a}=\ip p{D_{h,c_a}u}_h$, using skew-adjointness, gives
$\mathsf X_{G_{h,c_a}}=D_{h,c_a}X$ on both canonical slots.  These facts
prove \eqref{v6:exact-acceleration}, without modifying the multiplier rate.

The analytic original row has
$P_{h,c}-G_{h,c}=P_{h,c,3}+O_h(|c|\|X\|^4)$.
Its canonical gradient at $aY+O_h(a^2)$ is
$a^2\mathsf X_{P_{h,c,3}}(Y)+O_h(|c||a|^3)$.
Substitute \eqref{v6:normal-input} to prove
\eqref{v6:acceleration-order}.  The subtracted multiplier rate in
$\mathcal A_a^0$ is bounded, the finite vector field is analytic, and
$\mathsf F_h(aY+O_h(a^2),e)=O_h(a)$; thus $\mathcal A_a^0$ is bounded.

For sharpness, Euler's identity for the cubic scalar gives
$DP_{h,\upsilon_Y,3}(Y)[Y]=3P_{h,\upsilon_Y,3}(Y)$.
The last number is nonzero by \cref{v6:finite-P3}; in the spatial-sum
certificate its residue is $3\cdot201\equiv121\pmod{241}$.
Thus its gradient and its canonical vector field are nonzero.  Taking
norms in \eqref{v6:acceleration-order} proves both bounds.  This is a
statement about one specified contribution to the transformed acceleration,
not an assertion that the whole acceleration cannot have a cancellation
with other terms.
\end{proof}

The result removes the initial inverse-amplitude divergence from this
canonical comparison; it does not remove \cref{v5:actual-pole} in the
anchored variables.  Nor does it turn the cubic action correction into an
exact symmetry.  The transformation is an exact re-expression of the same
history, with its original observations still evaluated by pullback.

\section{A genuine translated metric has a different second-jet comparison}
\label{v6:geometric}

Let $\mathcal I_h$ be trigonometric interpolation on the odd grid, and
let $\partial_{i,h}^{\rm sp}$ denote its exact differentiation on arrays:
its Fourier multiplier is $2\pi\mathrm i k_i$, not
$\mathrm i\kappa_{i,h}(k)$.  Reconstruct a smooth metric near the initial
cut from $\gamma=I+\mathcal I_hu$, $N$ and $\beta$ in the usual ADM
formula.  For the uniform path above make the genuine coordinate change
$x=y-q(t)$.  The transformed ADM fields are
\[
 \widetilde\gamma^{\rm geo}(t,y)=\gamma(t,y-q(t)),\quad
 \widetilde N^{\rm geo}(t,y)=N(t,y-q(t)),\quad
 \widetilde\beta^{\rm geo}(t,y)=\beta(t,y-q(t))-\dot q(t).
\]
These formulas are a coordinate calculation for the interpolated metric;
no covariance of the original finite action is being assumed.

\begin{theorem}[Exact phase-versus-geometric initial second-jet defect]
\label{v6:metric-jet}
At $q(0)=\dot q(0)=0$, $\ddot q(0)=c_a$, the constant part of the
initial shift velocity is removed exactly:
\[
 \partial_t\widetilde\beta^{\rm geo}(0)=\dot\beta(0)-c_a.
\]
The geometric and phase-canonical spatial second jets satisfy
\begin{equation}
 \boxed{
 \partial_t^2\widetilde\gamma^{\rm geo}(0)
       -\mathcal I_h\ddot{\widetilde u}_a(0)
   =\sum_i c_a^i\mathcal I_h
           (\delta_{i,h}-\partial_{i,h}^{\rm sp})u_a(0).}
 \label{v6:metric-exact}
\end{equation}
On \eqref{v6:balanced} and \eqref{v6:normal-input}, put
\[
 V_i=r_i\{-T_i\sin(2\pi x_i+\phi_i)
                       +\sigma_iP_i\cos(2\pi x_i+\phi_i)\}.
\]
Then the limit of the right side as $a\to0$ at fixed $h$ is
\begin{equation}
 -(\omega_h-2\pi)\sum_i\upsilon_Y^iV_i,
 \qquad
 \left\| (\omega_h-2\pi)\sum_i\upsilon_Y^iV_i\right\|_{L^2}^2
 =2(\omega_h-2\pi)^2\sum_i r_i^2|\upsilon_Y^i|^2.
 \label{v6:metric-limit}
\end{equation}
For the certified unit three-site seed this is nonzero.  Thus the two
frames do not become identical in the second-jet amplitude limit at a
fixed cutoff, even though neither retains the $a^{-1}$ constant-rate
pole in that jet.
\end{theorem}
\begin{proof}
Differentiate the displayed coordinate formulas at the cut.  Because
$\dot q(0)=0$,
$\partial_t^2\widetilde\gamma^{\rm geo}
=\partial_t^2\gamma-c_a^i\partial_i\gamma$.
The phase formula gives
$\ddot{\widetilde u}=\ddot u-c_a^i\delta_{i,h}u$.
Subtract them to prove \eqref{v6:metric-exact}.  Each unit axial profile
obeys $\delta_{i,h}U=\omega_hV_i$ and $\partial_iU=2\pi V_i$,
exactly on every odd grid.  Inserting $u_a=aU+O_h(a^2)$ and
$c_a=-\upsilon_Y/a+O_h(1)$ proves \eqref{v6:metric-limit} with an
$O_h(a)$ error.  The profiles $V_i$ are orthogonal at different axes and
have squared norm $2r_i^2$.  This proves the norm identity.  Finally
$\omega_h<2\pi$, and the unit seed has $\upsilon_Y\ne0$.
\end{proof}

\begin{corollary}[The actual scale needed for this cofinal frame comparison]
\label{v6:cofinal-frame}
On any \emph{supplied} cofinal family of such normal forms, with radii
bounded above and below by positive constants, suppose the remainder in
\eqref{v6:metric-limit} along an amplitude path is bounded by
$K_h|a(h)|\to0$.  Then the difference in
\eqref{v6:metric-exact} tends to zero if and only if
\begin{equation}
                       h^2\|\upsilon_{Y,h}\|\longrightarrow0.
 \label{v6:weighted-translation}
\end{equation}
In particular bounded $\upsilon_{Y,h}$ would suffice for this comparison.
Neither existence of that cofinal family nor the bound on its vector is
provided by the three-site certificate.
\end{corollary}
\begin{proof}
Use the positive norm identity in \eqref{v6:metric-limit} and
$|\omega_h-2\pi|\asymp h^2$.  The upper and lower bounds on the radii
and the vanishing remainder give necessity and sufficiency by the triangle
and reverse triangle inequalities.
\end{proof}

The first time jet of the transformed ADM metric is bounded on the fixed
normal-form cut: its shift derivative is $O_h(1)$, its lapse derivative is
$O_h(1)$ and its spatial derivative in time is $O_h(a)$.  This is weaker
than closeness to a small harmonic first jet, because the bounded residual
lapse and transverse-shift derivatives have not been removed.
Moreover a large coordinate acceleration reappears in higher comparison
identities unless its actual derivatives are controlled.  Specifically,
when the original history and $q$ are $C^3$, the same calculation gives
\begin{equation}
 \begin{split}
 &\partial_t^3\widetilde\gamma^{\rm geo}(0)
            -\mathcal I_h\partial_t^3\widetilde u(0)\\
 &\quad=\sum_i\mathcal I_h(\delta_{i,h}-\partial_{i,h}^{\rm sp})
          \{3c_a^i\dot u(0)+q^{(3),i}(0)u(0)\}.
 \end{split}
 \label{v6:third-jet}
\end{equation}
The quantity $q^{(3)}(0)=\overline{\ddot\beta(0)}$ has not been
estimated in V5 or here.  The leading spatial part of $\dot u(0)$ is
constant for the scalar-momentum first jet, but this observation supplies
no bound for the second term in \eqref{v6:third-jet}.  It would be
incorrect to infer the differentiated physical readouts from the
second-jet estimate alone.

\section{Transported products preserve literal finite observations}
\label{v6:products}

\begin{lemma}[Why a continuous finite relabeling cannot be assumed]
\label{v6:algebra-boundary}
Every unital real algebra automorphism of pointwise arrays on a finite
set is a site permutation.  A continuous one-parameter family of these
automorphisms starting at the identity is constant.  Also an orthogonal
unital map that preserves pointwise nonnegativity is a site permutation.
Consequently the nontrivial group $U_{h,q}$ is not such an automorphism or
positive stochastic relabeling for all small nonzero parameters.
\end{lemma}
\begin{proof}
The minimal nonzero idempotents of the array algebra are the single-site
indicators.  An algebra automorphism permutes them, proving the first
claim.  The permutation group is finite, proving the continuity claim.
For the positivity assertion, all matrix entries are nonnegative.
Orthogonal rows with nonnegative entries have disjoint supports.  In a
square invertible matrix this forces one nonzero entry per row and per
column; orthogonality makes it one.  The nonzero phase derivative is the
infinitesimal generator of $U_{h,q}$, so that group is not constant.
These are elementary finite-algebra facts, not a new obstruction specific
to gravity.
\end{proof}

There is nevertheless an exact algebraic pullback, which retains the
original observations rather than declaring them invariant.  For scalar
or matrix-valued arrays set
\begin{equation}
 f\star_qg=U_{h,-q}\bigl[(U_{h,q}f)(U_{h,q}g)\bigr].
 \label{v6:star}
\end{equation}
The product on the right is the original pointwise ordered product.

\begin{proposition}[An exact transport identity for finite words and logarithms]
\label{v6:word-transport}
For fixed $q$, $\star_q$ is associative and unital, with the transported
adjoint equal to the ordinary matrix adjoint.  The map $U_{h,q}$ is an
algebra isomorphism from $(\star_q)$ to the original array algebra.
It commutes with the literal shifts and the original phase derivatives,
and preserves the normalized spatial trace.  Hence every finite word,
inverse on its actual domain, and convergent exponential or logarithm
series satisfies the corresponding exact pullback identity.  In particular,
\begin{align}
 \exp_{\star_q}(U_{h,-q}A)&=U_{h,-q}\exp(A),\notag\\
 \log_{\star_q}(U_{h,-q}L)&=U_{h,-q}\log(L),
 \label{v6:log-transport}
\end{align}
where the logarithm uses the same convergent local branch.  The complete
ordered spatial-link word and its curvature read can therefore be
represented exactly in the transported algebra.
\end{proposition}
\begin{proof}
Apply $U_{h,q}$ to both bracketings of a triple $\star_q$ product; both
become the original ordered triple product.  The other algebraic claims
follow in the same way.  The map is real and fixes constants, so it
commutes with the adjoint and transports the unit to the unit.  All the
linear shifts and phase derivatives are diagonal Fourier multipliers,
so they commute with $U_{h,q}$.  Its zero-frequency multiplier is one,
proving preservation of the spatial trace.  Induction gives the identity
for every word and every power.  Conjugating an inverse gives the
inverse in $\star_q$, and convergence of the stated series at fixed
finite dimension gives \eqref{v6:log-transport} term by term.
\end{proof}

For example $U_{h,-q}L$ need not be a pointwise Lorentz-group-valued link
in the original multiplication; it is a link in the transported algebra.
Likewise positivity of a transported metric or rate refers to the
transported cone, not an unproved positivity-preserving action of $U_{h,-q}$
on its site values.  The original pointwise positive observations are
still recovered by $U_{h,q}$.  Temporal connection pullbacks and the
variation of the time-dependent product must also be retained for a
spacetime identity.  The proposition is not an assertion of exact finite
four-dimensional diffeomorphism covariance.

\section{Uniform full-frequency control of the product change}
\label{v6:uniform-products}

The preceding product is exact but nonlocal.  Its difference from the
original product admits a quantitative cofinal estimate with explicit
regularity loss.  Let $\|\cdot\|_{H_h^s}$ be the Fourier norm with weight
$\langle k\rangle^s$ on $\{-m,\ldots,m\}^3$ and normalized Parseval
pairing.  Fixed component dimensions are absorbed in the constants.
The norms are uniformly equivalent to the phase Sobolev norms at the
orders used here, since $|\kappa_{i,h}(k)|\asymp |k_i|$ on the full odd
frequency set.

\begin{theorem}[A full-frequency moving-product estimate]
\label{v6:product-bound}
For $s>3/2$ there is $C_s$, independent of the odd cutoff and of $q$, such
that
\begin{align}
 \|\delta_{i,h}(fg)-(\delta_{i,h}f)g-f\delta_{i,h}g\|_{H_h^s}
    &\le C_sh^2\|f\|_{H_h^{s+3}}\|g\|_{H_h^{s+3}},
       \label{v6:Leibniz-bound}\\
 \|f\star_qg-fg\|_{H_h^s}
    &\le C_s\min\{1,h^2|q|\}
                   \|f\|_{H_h^{s+3}}\|g\|_{H_h^{s+3}}.
       \label{v6:product-bound-eq}
\end{align}
For fixed arrays $f,g$ and a differentiable path $q(t)$,
\begin{equation}
 \left\|\frac{\dd}{\dd t}(f\star_{q(t)}g)\right\|_{H_h^s}
 \le C_sh^2|\dot q(t)|\|f\|_{H_h^{s+3}}\|g\|_{H_h^{s+3}}.
 \label{v6:product-velocity}
\end{equation}
At a time when $\dot q=0$, if $q$ is twice differentiable,
\begin{equation}
 \left\|\frac{\dd^2}{\dd t^2}(f\star_{q(t)}g)\right\|_{H_h^s}
 \le C_sh^2|\ddot q|\|f\|_{H_h^{s+3}}\|g\|_{H_h^{s+3}}.
 \label{v6:product-acceleration}
\end{equation}
These estimates include aliased modes; no low-frequency projection is
part of their assumptions.
\end{theorem}
\begin{proof}
For frequency representatives $p,r$ let $k$ represent $p+r$ modulo $N$.
The multiplier of the product defect is
$\mathrm i(\kappa_i(k)-\kappa_i(p)-\kappa_i(r))$.
If no folding occurs, the sine Taylor remainder after its linear term
gives a bound $C h^2(|p|+|r|)^3$.  If folding occurs in any coordinate,
$|p|+|r|\ge c/h$, while each of the three phase symbols is bounded by
$2/h$.  The same bound follows by increasing $C$.  Further,
$|k_i|\le|p_i+r_i|\le|p_i|+|r_i|$ for the least absolute odd-grid
representative.

The weighted convolution is therefore bounded by the sum of a cyclic
convolution with weights $\langle p\rangle^{s+3}\langle r\rangle^3$
and its transpose, times $Ch^2$.  Young's $\ell^2$--$\ell^1$ inequality
on the finite frequency group applies with no dimension factor.  The
uniform estimate
\[
 \sum_r\langle r\rangle^3|\widehat g(r)|
 \le \left(\sum_{r\in\mathbb Z^3}\langle r\rangle^{-2s}\right)^{1/2}
                         \|g\|_{H_h^{s+3}}
\]
then proves \eqref{v6:Leibniz-bound}.  The same argument without the cubic
symbol weight proves the usual uniform $H_h^s$ product bound.

Define the bilinear defect
$\mathcal L_i(f,g)=\delta_i(fg)-(\delta_i f)g-f\delta_i g$.
Exact differentiation of \eqref{v6:star} yields
\[
 \frac{\dd}{\dd t}(f\star_qg)
   =-\dot q^i U_{h,-q}\mathcal L_i(U_{h,q}f,U_{h,q}g).
\]
The transformations are isometries at all the displayed Sobolev orders.
This proves \eqref{v6:product-velocity}.  Integrate on the straight path
from zero to $q$ to obtain the $h^2|q|$ bound in
\eqref{v6:product-bound-eq}; the two uniform product bounds give its
constant alternative.  At $\dot q=0$ all terms in the next derivative
except $-\ddot q^iU_{h,-q}\mathcal L_i(U_{h,q}f,U_{h,q}g)$ vanish,
proving \eqref{v6:product-acceleration}.
\end{proof}

When both factors are perturbations of size $O(a)$ in the displayed
higher norms, the initial product-acceleration error caused by a rate
$|\ddot q|=O(a^{-1})$ is $O(h^2a)$, with a uniform bound only if those
higher-norm and rate constants are uniform.  Products with a constant
factor have zero defect.  This is a useful coefficient estimate; it is
not a bound on the inverse of the original nonlinear multiplier saddle.
In particular it cannot suppress the independent order-$h$ cubic
plaquette term in \cref{v6:P3}.

\section{Consequences, remaining targets and verification}
\label{v6:status}

There are three mathematically different facts.  First, the raw action has
an explicit cubic uniform-shift response of order $h$ on the balanced
first jets, even though their quadratic weak drift is zero.  This new
coefficient is entirely magnetic on that family.  Second, exact canonical
transport converts the surviving inverse-amplitude shift acceleration into
an order-amplitude contribution at the initial spatial canonical second
jet, without deleting an action term or changing any multiplier equation.
Third, a genuine geometric translation differs from that phase transport
by the explicitly measured $|\omega_h-2\pi|\upsilon_Y$ second-jet term.
Neither transport is an automatic equality of the original anchored
observations.

Thus the spatially constant pole need not be treated as proof of a
coordinate-independent singularity.  It also cannot be dismissed as an
exact finite symmetry.  The original-action comparison problem now has
concrete additional quantities: the effective rate vector
$\upsilon_{Y,h}$, its physical-time derivatives, the conjugated nonlinear
Hamiltonian remainder, and the transported literal-link reads.  The
three-site pole does not decide their behavior along refinement.

The next action-specific cofinal task is to form the balanced family's
complete active and weak sources at finer cutoffs and estimate their
actual effective vector, in particular the weighted quantity
$h^2\|\upsilon_{Y,h}\|$ appearing in \cref{v6:cofinal-frame}, alongside
the original response and its differentiated sources.  Alternatively one
may seek a different first-jet family with $\mathfrak c_{Y,h}=0$;
\cref{v6:P3} alone neither proves nor rules out that different cancellation.
The constant-rate normal form is proved only initially.  It has not been
propagated by the exact Euler flow, and no common physical interval,
cutoff-uniform nonlinear inverse, or raw-action Einstein limit is asserted.
A coordinate change cannot provide any of those estimates by definition.

\paragraph{Reproducibility.}
The symbolic Fourier verifier evaluates the unchanged cubic coefficient
functional over exact Laurent polynomials.  It keeps independent symbols
for the higher phase-derivative frequencies until they cancel.  It checks
the homogeneous plaquette coefficient, the three single-profile
contractions, both handedness signs and the phase independence.  Separate
exact rational evaluations of the three-site seed recover the cubic
polynomial in its algebraic root and the value $-3645\sqrt3/8$ in the
spatial-sum convention.  The inherited constant-rate residues then give
the nonzero radial derivative used in \cref{v6:initial-comparison}.
Floating general-grid tests use the actual sine phase symbols and full
cyclic products to check the closed formula; they are diagnostics, not
the all-cutoff proof.  Matrix and Fourier tests separately check the
canonical transformation identity, the transported ordered words and
logarithms, and the exact initial second- and third-jet difference formulas.
No nonlinear original-action trajectory, new rank certificate beyond V5,
or interval-certified cofinal estimate is represented as computed.

The source audit inspected the five preceding bridge bodies, the
open-basin source's constant-shift quadratic and cubic formulas, and the
Grand Tensor source's finite derivative and Cartan-response boundaries.
The elementary finite-algebra lemma and the general change-of-variables
identity are used as tools, not claimed as new abstract geometry.
All five preceding research bodies are preserved byte for byte in this
cumulative source.  The accompanying manifest records their hashes and
those of the new body and coefficient files.

\begin{thebibliography}{9}
\raggedright
\setcounter{enumiv}{8}
\bibitem{v6:BD}
B. Bahr and B. Dittrich,
\emph{Breaking and restoring of diffeomorphism symmetry in discrete gravity},
arXiv:0909.5688 (2009).
\url{https://arxiv.org/abs/0909.5688}.
This supplies primary-source context for broken finite diffeomorphism
symmetry, not the renewal coefficient, preparation, or frame-comparison
results proved here.
\end{thebibliography}
% END IMMUTABLE EXACT-ACTION BRIDGE V6 BODY

% Future iterations append here.
% BEGIN IMMUTABLE EXACT-ACTION BRIDGE V7 BODY
\clearpage
\renewcommand{\thesection}{\arabic{section}}
\renewcommand{\theHsection}{bridgeV7.\arabic{section}}
\setcounter{section}{50}
\section{A covector-order correction before the finer-cutoff calculation}
\label{v7:context}

This iteration retains the action and the balanced first-jet family of
\cref{v6:balanced}.  Its principal analytical result locates the effective
constant-shift source in an explicit Fourier-edge pairing.  The source is
not inferred from the order-$h$ cubic momentum value of \cref{v6:P3}.
A complete five-site calculation separately establishes signed regularity
and a nonzero effective source for that actual family.  These are two
different kinds of progress: an all-cutoff localization identity on every
regular active branch, and new characteristic-zero certificates at a
specified finer cutoff.

An audit needed for this calculation found a coordinate-order error in the
V5 target assembly.  It must be corrected before the inherited numerical
invariants are used.  The first-source columns in
\cref{v3:finite-matrices,v4:finite-rank} are ordered by multiplier component,
whereas the loop that assembled the last $3V$ values of V5's
\src{B2\_all} used site-major shift order.  The selected source columns were
therefore paired with some drift values belonging to different physical
tests.  The source matrices themselves and their determinant certificates
are unchanged.

\subsection{The precise permutation}
Fix lexicographic site labels $x=0,\ldots,V-1$ and shift components
$i=0,1,2$.  The physical column order used throughout this body is
\begin{equation}
 (n(x))_{x},\quad (\beta^1(x))_{x},\quad
 (\beta^2(x))_{x},\quad(\beta^3(x))_{x}.
 \label{v7:ordering}
\end{equation}
If $b^{\rm site}$ lists the lapse values followed by the values at
$(x,i)$, define the permutation $\mathcal P$ by
\begin{equation}
 b^{\rm comp}_x=b^{\rm site}_x,\qquad
 b^{\rm comp}_{V+iV+x}=b^{\rm site}_{V+3x+i}.
 \label{v7:permutation}
\end{equation}
The active covector is $b_A=b^{\rm comp}|_{\mathcal P_{\rm act}}$,
where $\mathcal P_{\rm act}$ is the selected column list; it is not the
restriction of $b^{\rm site}$ to that same list.  This change does not
rescale a field or a physical observation.

\begin{proposition}[Corrected three-site effective source]
\label{v7:corrected-three}
At the balanced three-site seed, with the original embedding
$\sqrt3\mapsto56\in\mathbb F_{241}$ and spatial-sum convention,
the corrected action-derived quantities are
\begin{equation}
 \begin{aligned}
 R_cL^*A^{-1}b_A&\equiv(226,87,166)^T,\\
 \mathfrak c_Y&\equiv(15,154,75)^T,\\
 \upsilon_Y=D_{cc}^{-1}\mathfrak c_Y&\equiv(217,147,105)^T.
 \end{aligned}
 \label{v7:corrected-values}
\end{equation}
The three determinants remain
$\det A\equiv113$, $\det D\equiv56$, and $\det D_{cc}\equiv79$.
In \cref{v6:finite-P3}, the corrected residue of
$P_{h,\upsilon_Y,3}(Y)$ is $79$, and its radial derivative has residue
$237$, rather than $201$ and $121$.
\end{proposition}
\begin{proof}
Regenerate the drift of \eqref{v5:prepared-drift-formula} on the actual
columns of \eqref{v7:ordering}, or apply \eqref{v7:permutation} to the
stored complete drift.  Solve the unchanged active Gram and multiply by
the unchanged constant weak columns.  The resulting vectors give
\eqref{v7:corrected-values}; the raw cubic translation drift is zero by
\cref{v5:cubic-means}.  In addition to the selected equations, the audit
verifies the \emph{complete} identity
\begin{equation}
 (J_1^{\rm all})^*K^{-1}C_A A^{-1}b_A=b^{\rm comp},
 \qquad C_A=J_1^{\rm all}|_{\mathcal P_{\rm act}}.
 \label{v7:all-active-check}
\end{equation}
Every one of its $108$ entries agrees in the residue field.  The unpermuted
calculation fails $28$ equations even against its own misordered complete
vector; against the correctly ordered physical vector, it differs in $79$
entries.  Solving its selected square subsystem therefore did not supply a
consistent physical covector.  This distinguishes the correction from a
harmless change of notation.

The inverse and all denominators are regular modulo $241$; the complete
arrays and physical column labels are provided in the certificate.  Hence
the nonzero vector in \eqref{v7:corrected-values} gives a valid
characteristic-zero nonvanishing certificate.  The unit cubic momentum
row has residue $31$ in spatial-sum normalization by
\cref{v6:finite-P3}.  Since $217+147+105\equiv228$, its contraction is
$31\cdot228\equiv79$; Euler's cubic identity multiplies this by three,
giving $237$.
\end{proof}

The specific vectors in \eqref{v5:c-residue} and
\eqref{v5:constant-normal}, and the numerical contractions in
\eqref{v6:Pnu}, are superseded by \cref{v7:corrected-three}.  Their general
formulas were not the error.  In particular the preparation-invariance
argument, the rank-26 control minor, the affine attainable-rate formula,
and the implicit-function argument are unchanged.  The corrected vector
and determinants are still nonzero, so the fixed-cutoff pole, its
constant-rate preparation, and the nonzero order-amplitude moving-frame
contribution retain their stated qualitative conclusions.  No claim is
made that the earlier incorrect numbers were equivalent representatives.
The original bodies remain verbatim, and this correction is part of the
append-only record.

\section{First-jet coefficients and a reproducible finer-grid audit}
\label{v7:coefficients}

The source construction is the inherited exact coefficient identity
\eqref{v3:compiler-formula}, with its normal correction and auxiliary
correction retained.  The drift is
\eqref{v5:prepared-drift-formula}, including the lapse correction imposed
by the second initial constraint equation.  Neither is replaced by a
finite-amplitude difference quotient.

For clarity, an efficient equivalent differentiation procedure is specified
here.  Let $\Phi_h$ be the scalar coefficient functional of
\cref{v3:implementation}, $z_1$ its flat first stationary connection,
$\eta=(A_0,w)$ its twelve auxiliary coordinates, and $R_0,Z_0$ the
retained flat maps.  Gradients below mean exact derivatives of the indicated
finite polynomial coefficients.  Set
\begin{equation}
 \begin{aligned}
 \tau_2&=\nabla_\eta[a^2]\Phi_h(az_1,\eta;aY,0)|_{\eta=0},
 &z_2^\perp&=-R_0\tau_2,\\
 \eta_1&=-\nabla_y[a]\Phi_h(az_1+R_0y,\eta_0;aY,\nu)|_{y=0},
 &\nu&=(0,\beta),
 \end{aligned}
 \label{v7:reverse-aux}
\end{equation}
where $\eta_0=(0,\delta_i\beta_j+\delta_j\beta_i)$ on a curl-free test.
The weak column is
\begin{equation}
 \mathcal Q_h(Y;\beta)=
 \nabla_s[a^2]\Phi_h(az_1+a^2z_2^\perp+Z_0s,
                         \eta_0+a\eta_1;aY,\nu)|_{s=0}.
 \label{v7:reverse-source}
\end{equation}
The gradient is in the raw symmetric dual coordinates of
\eqref{v3:raw-K}, so off-diagonal Frobenius weights are retained.
Equations \eqref{v7:reverse-aux}--\eqref{v7:reverse-source} are simply
reverse differentiation of the existing coefficient construction.  They
are a faster evaluation method, not a new field equation or a new
source identity.

All fine-grid coefficients can be evaluated in a cyclotomic number field.
For $N$ odd put $\zeta=\exp(\pi\mathrm i/(2N))$,
$t=\zeta^2$, $z=\zeta^4$ and $\mathrm i=\zeta^N$.  On the signed Fourier
representatives the derivative symbol is
\begin{equation}
 \mathrm i\kappa_h(k)=N(t^k-t^{-k}),\qquad
 \omega_h=-\mathrm iN(t-t^{-1}).
 \label{v7:cyclotomic}
\end{equation}
The sampled cosine and sine use
$(z^j+z^{-j})/2$ and $(z^j-z^{-j})/(2\mathrm i)$.  These formulas keep the
folded derivative, not an ordinary spectral derivative.

\paragraph{Exact arithmetic convention.}
The finite-field compiler uses addition, multiplication, linear coefficient
maps, and exact formal differentiation modulo a prime.  Reverse adjoints
are reduced modulo the same prime at each elementary step.  Its array
implementation uses exactly representable integer residues in binary64;
finite sums and products are bounded below the exact-integer threshold.
The signed-Gram verification itself uses integer modular arithmetic.
Independent forward coefficient tests and the unchanged three-site source
arrays check the differentiation implementation.  This is not numerical
differentiation of a discontinuous residue map: the derivative being
computed is the formal derivative of the underlying polynomial over the
residue field.

\section{The balanced five-site family has a regular complete signed source}
\label{v7:five-site}

Use the unit balanced profiles of \cref{v6:balanced}, all phases zero,
all handedness signs positive, and $p=\omega_h I$.  At $N=5$, this is the
same refinement-compatible physical family as the three-site balanced
seed, not a new detuning selected to make a minor nonzero.
There are $V=125$ sites, $372$ active coordinates, and $127$ weak coordinates.

\begin{theorem}[An original-action five-site signed certificate]
\label{v7:five-rank}
At this specified first jet,
\begin{equation}
 \rank J_1^{\rm all}=372,\qquad
 \rank(C_A,W)=499,\qquad
 \det\bigl((C_A,W)^*K^{-1}(C_A,W)\bigr)\ne0.
 \label{v7:five-rank-eq}
\end{equation}
The weak Schur complement $D$ and its constant-constant block $D_{cc}$
are both invertible.  The actual effective constant source is nonzero.
In the cyclotomic specialization $\zeta_{20}\mapsto235\in\mathbb F_{241}$,
the exact residues are
\begin{equation}
 \begin{aligned}
 \det A&\equiv190,& \det D&\equiv119,&
 \det\Gamma_0&\equiv197,& \det D_{cc}&\equiv73,\\
 \mathfrak c_{Y,1/5}&\equiv(1,160,213)^T,&
 D_{cc}^{-1}\mathfrak c_{Y,1/5}&\equiv(88,32,86)^T.
 \end{aligned}
 \label{v7:five-residues}
\end{equation}
These are fixed-cutoff characteristic-zero assertions, not bounds on the
real inverse norms or on a cofinal sequence.
\end{theorem}
\begin{proof}
The element $235$ has exact order $20$ modulo $241$, and satisfies the
cyclotomic polynomial $\Phi_{20}$.  Formula \eqref{v7:cyclotomic} gives
$\omega_h\mapsto77$.  All denominators in the flat inverses, coefficient
extraction, $h=1/5$, and the cyclotomic Fourier transform are units at this
prime.  The full source matrices are assembled in
\eqref{v7:ordering}.  For $C_A$, take the lexicographically first pivot
columns returned by exact elimination of $J_1^{\rm all}$; their full list
is included in the certificate.  Equivalently this is a deterministic
specification of one fixed integer column list, not a floating rank
threshold.  The weak columns are the three constants followed by gradients
of the $124$ point differences against the last site.

The raw inverse $K^{-1}$ is $I_{125}\otimes\mathsf K_{\rm raw}^{-1}$
from \eqref{v3:raw-K}.  Assemble $A,L,D,\Gamma_0$ as in
\eqref{v5:Schur}, with no term dropped.  Exact elimination gives the
four nonzero determinants in \eqref{v7:five-residues}.  The independent
Schur check is $190\cdot119\equiv197\pmod{241}$.  Thus the full
signed Gram is nonsingular over the original cyclotomic field and under
its physical complex embedding; its real coefficients give the asserted
real ranks.  The upper active rank $3V-3$ follows from the constant lapse
and the $V+2$ curl-free shift null directions of the first source.  The
certificate attains this bound.

The weak quadratic target vanishes by \cref{v4:neutral}.  Evaluate the
prepared active drift with the correct column convention and verify
\eqref{v7:all-active-check} on \emph{all $500$ multiplier columns}.
The cubic constant drift is zero by the fine-grid result
\cref{v5:cubic-means}.  Therefore
$\mathfrak c=-R_cL^*A^{-1}b_A$.  The exact active solve gives the last
line of \eqref{v7:five-residues}; the inverse of its $3\times3$ constant
block gives the displayed vector.  All determinants used in these solves
are nonzero modulo $241$, so the resulting nonzero residues prove
characteristic-zero nonvanishing, not just residue-field consistency.
\end{proof}

\begin{corollary}[A rate pole is not confined to the exceptional three-site grid]
\label{v7:five-pole}
On every analytic exact initial preparation with this five-site first jet,
the normalized original-action initial rate satisfies
\[
 \|\partial_t\beta_{1/5}(a,0)\|\ge c/|a|
\]
for all sufficiently small nonzero amplitudes, with a threshold that may
depend on the preparation.  Its initial literal electric and spatial
curvature reads are $O_{1/5}(a)$.
\end{corollary}
\begin{proof}
The mean calibration and original nonlinear initializer apply to the unit
balanced first jet.  The nonzero leading signed Gram of
\cref{v7:five-rank} persists at sufficiently small nonzero amplitude.
The normalized one-null-direction continuation of \cite{OB25} therefore
supplies the local exact-action history.  The algebraic identity
$R_cD(u_1)_W=\mathfrak c$ in \cref{v5:invariant} is independent of cutoff;
\cref{v7:five-rank} now supplies its previously missing nonzero vector
at $N=5$.  It bounds $\|(u_1)_W\|$ below.  The original amplitude grading
gives the displayed rate pole, exactly as in \cref{v5:actual-pole}.
Source neutrality and signed regularity give the $O_h(a)$ initial reads
by \cref{v4:initial-electric}.  Neither argument makes the local
interval or any constant uniform in $h$.
\end{proof}

\section{Parity of the actual active target and constant weak columns}
\label{v7:parity}

The finer certificate shows why another determinant is not the analytical
endpoint.  We now identify where its effective constant source can come
from.  Let $N=2m+1\ge5$ and define the real orthogonal Fourier involution
\begin{equation}
 \widehat{\mathscr R_h f}(k)=(-1)^{k_1+k_2+k_3}\widehat f(k),
 \qquad k\in\{-m,\ldots,m\}^3.
 \label{v7:R}
\end{equation}
It commutes with phase derivatives, literal shifts, the positive
Frobenius metric, $K_0$, and the constant and transverse Fourier
projections.  It is used only as an analytical parity test.  At odd
cutoff it is not assumed to be a site permutation, a positivity-preserving
map, or an automorphism of the original cyclic product.

Write the active multiplier Hilbert space as
\[
 \mathcal M_{A,h}=
 \{n:\overline n=0\}\oplus
 \{\beta:\overline\beta=0,\ \kappa(k)\cdot\widehat\beta(k)=0\}.
\]
Its physical inner product is the spatial mean of the ordinary scalar and
vector products.  Define
\begin{equation}
 \mathscr S_h(n,\beta)=(-\mathscr R_hn,\mathscr R_h\beta),\qquad
 \mathcal A_h=J_{1,h}|_{\mathcal M_{A,h}},\qquad
 W_{c,h}=\mathcal Q_h(Y_h;c),
 \label{v7:active-parity}
\end{equation}
where a source is identified with its Riesz representative in the
symmetric connection tangent space.  Adjoints in the following proofs
are physical Hilbert adjoints.  Raw-coordinate implementations use
\eqref{v3:raw-K} to represent these same operations.

\begin{lemma}[A time-reversal grading of the second constant source]
\label{v7:source-grading}
For a constant shift $c$, the quadratic source
$\mathcal Q_h((U,p);c)$ is odd in $p$.  In particular its polynomial types
are $Up$, with their indicated derivatives and literal shifts, and not
$U^2$ or $p^2$.  For the axial $U$ and spatially constant scalar $p$ of
\cref{v6:balanced},
\begin{equation}
 \operatorname{supp}\widehat W_{c,h}
       \subset\{\pm e_1,\pm e_2,\pm e_3\},\qquad
 \mathscr R_hW_{c,h}=-W_{c,h}.
 \label{v7:W-odd}
\end{equation}
\end{lemma}
\begin{proof}
In the Lorentz coefficient convention put
$\tau(r,b)=(r,-b)$.  This is an orthogonal Lie-algebra automorphism for
the bracket displayed in \cref{v6:P3}.  Apply to the coefficient functional
the replacements
\[
 (U,p,z,A_0,w,n,\beta)
 \longmapsto (U,-p,\tau z,-\tau A_0,-w,n,-\beta).
\]
The velocity and momentum terms are invariant.  The electric coefficient
is boost-valued and the nested bracket transforms as minus its $\tau$
image, leaving its boost pairing unchanged.  The magnetic field strength
transforms by $\tau$, and so does its bivector coefficient because its
shift part changes sign.  The spatial derivative and literal shift
operators commute with $\tau$.  Thus the full relevant coefficient
functional is invariant, including both cubic logarithmic terms.

The flat normal maps in \eqref{v7:reverse-aux} respect these rotation,
boost, and velocity gradings.  This can also be seen directly from their
Gauss and symmetric-velocity blocks in \cref{v3:compiler}: the kernel
inclusion consists of symmetric rotations and is fixed by $\tau$; the
orthogonal normal inverse intertwines the induced sign maps on the
auxiliary rows.  Consequently substituting the displayed replacements
in \eqref{v7:reverse-aux} and \eqref{v7:reverse-source} shows
\[
 \mathcal Q_h((U,-p);-c)=\mathcal Q_h((U,p);c).
\]
Linearity in $c$ makes the coefficient odd in $p$.  It is homogeneous of
total field degree two by \cref{v3:preparation-independent}; hence only
mixed $Up$ terms remain.  All flat inverses involved are pointwise
constant, and all other linear maps are Fourier multipliers.  With $p$
constant, no such term changes the unit axial support of $U$.  This
proves both statements in \eqref{v7:W-odd}.
\end{proof}

\begin{lemma}[The prepared active target has the opposite parity]
\label{v7:target-even}
Let $b_{A,h}$ be the prepared quadratic drift of
\eqref{v5:prepared-drift-formula} on the balanced axial scalar-momentum
first jet.  For every odd $N\ge5$,
\begin{equation}
                         \mathscr S_h b_{A,h}=b_{A,h}.
 \label{v7:b-even}
\end{equation}
The same assertion holds on the complete multiplier module before
restriction to the active slice.
\end{lemma}
\begin{proof}
The same grading gives types $U^3$ and $Up^2$ in $H_3$, types $U^2,p^2$
in the quadratic lapse row, and type $Up$ in the quadratic shift row.
The linear vector field is $(V(p),\mathcal F(U))$, with the original
constant-coefficient operators in \eqref{v5:F1}.  Substitution in
\eqref{v5:prepared-drift-formula}, with its last term retained, therefore
gives type $Up$ in the lapse drift and types $U^2,p^2$ in the shift drift.
After the derivatives on the multiplier test have been moved by exact
summation by parts, these support statements concern the actual Riesz
vectors, not merely their values on low-frequency tests.

The lapse target has only unit axial frequencies and is odd under
$\mathscr R_h$.  The shift target has only zero frequency and sums of two
signed axial unit frequencies.  When $N\ge5$ these lie in the retained
box without folding, and their total signed parity is even.  Derivatives
and literal shifts preserve the supports.  Thus the extra minus sign in
the lapse part of $\mathscr S_h$ gives \eqref{v7:b-even}.  Orthogonal
restriction to the active slice commutes with the two involutions.
The balanced weak quadratic drift is zero, so this restriction does not
discard an incompatible weak target.  At $N=3$ the doubled modes fold,
and the assertion is not made.
\end{proof}

\section{The active parity defect lives only on the folded Fourier faces}
\label{v7:seam}

Let $\mathscr P_{\partial,h}$ be the Fourier projection onto
\begin{equation}
 \Sigma_h=\{k\in\{-m,\ldots,m\}^3:
                  |k_i|=m\text{ for at least one }i\},\qquad
 r_h=|\Sigma_h|=N^3-(N-2)^3.
 \label{v7:seam-set}
\end{equation}
Use the same symbol componentwise on scalar and connection arrays.
This is a boundary of the \emph{retained frequency box}, not a physical
spatial boundary or a discarded set of states.  Define the actual defect
\begin{equation}
 \mathcal E_h=\mathscr R_h\mathcal A_h-
                                      \mathcal A_h\mathscr S_h.
 \label{v7:E}
\end{equation}
This $\mathcal E_h$ is not the electric curvature read of earlier bodies.

\begin{theorem}[An exact surface-rank defect for the original active source]
\label{v7:seam-theorem}
On the axial scalar-momentum family, $\mathcal E_h$ vanishes on the
transverse-shift slot and obeys the exact support identity
\begin{equation}
 \boxed{\mathcal E_h(n,\beta)=
 \mathscr P_{\partial,h}\mathcal E_h
                    (\mathscr P_{\partial,h}n,0).}
 \label{v7:seam-support}
\end{equation}
For radii in a fixed bounded set, uniformly in phase, handedness, and odd
$N\ge5$,
\begin{equation}
 \rank\mathcal E_h\le r_h=6N^2-12N+8,\qquad
 \|\mathcal E_h\|_{L^2\to L^2}\le C h^{-1}.
 \label{v7:seam-bound}
\end{equation}
No active signed inverse or full weak rank hypothesis is used here.
\end{theorem}
\begin{proof}
First let $f$ have unit axial Fourier support and let $M_f$ denote actual
cyclic multiplication.  In the matrix of
$\mathscr R_hM_f+M_f\mathscr R_h$, a nonzero entry connects an input
$p$ to $k=[p\pm e_i]_N$.  Without folding,
$(-1)^{k_1+k_2+k_3}=-(-1)^{p_1+p_2+p_3}$; the entry is zero.  Folding can
occur only from $p_i=\pm m$ to $k_i=\mp m$.  Both modes belong to
$\Sigma_h$.  Therefore
\begin{equation}
 \mathscr R_hM_f+M_f\mathscr R_h
 =\mathscr P_{\partial,h}
  (\mathscr R_hM_f+M_f\mathscr R_h)\mathscr P_{\partial,h}.
 \label{v7:scalar-seam}
\end{equation}
This statement also holds for each derivative and literal shift of $f$,
since neither changes its unit axial support.

The original lapse part of the first source is a finite sum of terms
$\delta_jM_{U_{lb}}-M_{U_{lb}}\delta_j-M_{\delta_jU_{lb}}$,
and terms $hM_{r_{ia}}S_i\delta_b$, with their fixed symmetric
alternating contractions; see \cref{v1:first-source} and its
explicit proof.  Each $r_{ia}$ is a first derivative of $U$.
All derivatives and shifts commute with $\mathscr R_h$ and
$\mathscr P_{\partial,h}$.  Applying \eqref{v7:scalar-seam} term by term
gives \eqref{v7:seam-support} on the lapse slot.  The transverse source
of \eqref{v1:shift-source} is a constant-coefficient derivative/shift map
when the momentum is scalar and constant.  It commutes with
$\mathscr R_h$, so its defect is zero.

Only the $r_h$ scalar lapse modes on the boundary can enter; this gives
the rank bound.  The unit-profile coefficients and their first phase
derivatives are uniformly bounded in supremum norm.  Literal shifts are
isometries, cyclic multiplication by a bounded array is bounded on
physical $L^2$, and $\|\delta_j\|\le2/h$.  Every lapse term has at most
one derivative of its input, and the electric terms carry an additional
factor $h$.  The finite displayed contractions therefore have norm at most
$C/h$.  Applying two isometric involutions gives the same bound for
$\mathcal E_h$.  No product-rule replacement or high-mode cutoff was made.
\end{proof}

\begin{corollary}[Only a surface-rank part of the active Gram mixes parity]
\label{v7:Gram-seam}
Let $\Gamma_{A,h}=\mathcal A_h^*K_0^{-1}\mathcal A_h$ on the active
Hilbert space.  Regardless of whether it is invertible,
\begin{equation}
 [\Gamma_{A,h},\mathscr S_h]
 =-\mathcal A_h^*K_0^{-1}\mathcal E_h
                         +\mathcal E_h^*K_0^{-1}\mathcal A_h.
 \label{v7:commutator}
\end{equation}
Its rank, and hence the rank of
$(\Gamma_{A,h}-\mathscr S_h\Gamma_{A,h}\mathscr S_h)/2$, is at most
$2r_h$.  This is $O(N^2)$ rather than the $O(N^3)$ active dimension.
\end{corollary}
\begin{proof}
Use $\mathcal A\mathscr S=\mathscr R\mathcal A-\mathcal E$,
its adjoint, and commutation of $K_0^{-1}$ with $\mathscr R$.
Their difference is \eqref{v7:commutator}.  Each summand has rank at
most $\rank\mathcal E$.  Multiplication by the involution converts
this commutator into the off-parity part without changing its rank.
This is an upper bound; it is not a claim that the parity-diagonal blocks
are themselves invertible.
\end{proof}

\section{The effective translation source is exactly a boundary pairing}
\label{v7:boundary-pairing}

Assume now that the \emph{active} signed Gram $\Gamma_{A,h}$ is invertible
at the retained cutoff.  The complete weak Gram need not be inverted in
this section.  For a constant shift test $c$ put
\begin{equation}
 \begin{aligned}
 x_h&=\Gamma_{A,h}^{-1}b_{A,h},&
 v_h&=K_0^{-1}\mathcal A_hx_h,\\
 k_{c,h}&=\mathcal A_h^*K_0^{-1}W_{c,h},&
 y_{c,h}&=\Gamma_{A,h}^{-1}k_{c,h},\qquad
 u_{c,h}=K_0^{-1}\mathcal A_hy_{c,h}.
 \end{aligned}
 \label{v7:primal-adjoint}
\end{equation}
These are the actual active solve and three actual adjoint active solves;
no worst-case inverse is substituted for their right sides.
The vectors $v_h,u_{c,h}$ are connection-tangent comparison vectors,
not the complete physical electric solution of the weak equations.
Since the raw cubic constant drift is zero at $N\ge5$,
\begin{equation}
 \mathfrak c_{Y,h}[c]=-\ip{W_{c,h}}{v_h}_h
                              =-\ip{k_{c,h}}{x_h}_h.
 \label{v7:c-again}
\end{equation}

\begin{theorem}[Exact source-resolved Fourier-boundary formula]
\label{v7:boundary-identity}
For every regular active cutoff on this family with $N\ge5$,
\begin{equation}
 \boxed{
 2\mathfrak c_{Y,h}[c]
 =\ip{\mathcal E_h y_{c,h}}{v_h}_h
                 -\ip{u_{c,h}}{\mathcal E_h x_h}_h.}
 \label{v7:boundary-exact}
\end{equation}
Consequently only the following four boundary tails can enter:
\begin{equation}
 \begin{split}
 2\mathfrak c_{Y,h}[c]
 ={}&\ip{\mathcal E_h(\mathscr P_{\partial,h}y_{c,h}^{n},0)}
                                     {\mathscr P_{\partial,h}v_h}_h\\
 &-\ip{\mathscr P_{\partial,h}u_{c,h}}
                   {\mathcal E_h(\mathscr P_{\partial,h}x_h^{n},0)}_h.
 \end{split}
 \label{v7:four-tails}
\end{equation}
Here $x^n,y^n$ are the lapse components of the active multiplier solutions.
\end{theorem}
\begin{proof}
The weak constant column has unit support by \cref{v7:source-grading},
whereas $\mathcal E_h$ has output only on $\Sigma_h$ by
\cref{v7:seam-theorem}.  Since $m\ge2$, these sets are disjoint.
The fixed Cartan map preserves frequency, so
$\mathcal E_h^*K_0^{-1}W_{c,h}=0$.
Taking the adjoint of the defining source defect gives
\[
 \mathscr S_h k_{c,h}
 =\mathcal A_h^*\mathscr R_hK_0^{-1}W_{c,h}
                    -\mathcal E_h^*K_0^{-1}W_{c,h}
 =-k_{c,h}.
\]
The target instead satisfies $\mathscr S_hb_{A,h}=b_{A,h}$ by
\cref{v7:target-even}.  From $\Gamma_Ax=b_A$,
\[
 \Gamma_A(I-\mathscr S)x=-[\Gamma_A,\mathscr S]x.
\]
Pair with $y_c=\Gamma_A^{-1}k_c$.  Self-adjointness of the signed Gram
and oddness of $k_c$ give
\[
 2\ip{k_c}{x}
 =-\ip{y_c}{[\Gamma_A,\mathscr S]x}
 =\ip{u_c}{\mathcal E x}-\ip{\mathcal E y_c}{v}.
\]
Equation \eqref{v7:c-again} proves \eqref{v7:boundary-exact}.
Finally apply the exact input and output support identity
\eqref{v7:seam-support}; orthogonality permits projection of the other
factor in each pairing.  This proves \eqref{v7:four-tails} with its
original signs and positive physical pairing.  No positivity of
$\Gamma_A$ is required.
\end{proof}

This separates two effects that should not be conflated.  The cubic
constant-shift \emph{momentum} in \cref{v6:P3} has an order-$h$ magnetic
term even on constant records.  The effective constant \emph{drift after
active elimination}, on the balanced family at $N\ge5$, has the exact
boundary representation \eqref{v7:boundary-exact}.  The second statement
does not cancel, modify, or dispute the first.  It identifies a more
specific mechanism for a different observable.

\section{Consequences for the cofinal estimate, without an inverse assumption in disguise}
\label{v7:tails}

\begin{corollary}[A four-tail estimate for the actual effective source]
\label{v7:tail-estimate}
Under the regular active hypotheses of \cref{v7:boundary-identity},
\begin{equation}
 \begin{split}
 |\mathfrak c_{Y,h}[c]|\le\frac{C}{h}\bigl(&
 \|\mathscr P_{\partial,h}y_{c,h}^{n}\|_h
                  \|\mathscr P_{\partial,h}v_h\|_h\\
 &+\|\mathscr P_{\partial,h}u_{c,h}\|_h
                  \|\mathscr P_{\partial,h}x_h^{n}\|_h\bigr).
 \end{split}
 \label{v7:tail-bound}
\end{equation}
For any $r\ge0$ this implies
\begin{equation}
 |\mathfrak c_{Y,h}[c]|\le C_rh^{2r-1}
 \left(\|y_{c,h}^{n}\|_{H_h^r}\|v_h\|_{H_h^r}
       +\|u_{c,h}\|_{H_h^r}\|x_h^{n}\|_{H_h^r}\right).
 \label{v7:Sobolev-tail}
\end{equation}
In particular, uniform bounds on these four \emph{actual} solution slots
at some $r>1/2$ would imply $\mathfrak c_{Y,h}\to0$.
No such uniform solution bounds are asserted here.
\end{corollary}
\begin{proof}
Apply Cauchy--Schwarz and \eqref{v7:seam-bound} to
\eqref{v7:four-tails}.  On $\Sigma_h$ one has
$|k|\ge m\ge (2/5)h^{-1}$ for $N\ge5$, so Parseval gives
$\|\mathscr P_{\partial,h}f\|_h\le C_rh^r\|f\|_{H_h^r}$.
Substitution proves both inequalities.  Apply them to the three coordinate
constant tests to obtain the vector statement.
\end{proof}

On a family for which $D_{cc,h}$ is separately invertible, the coefficient
$\upsilon_{Y,h}=D_{cc,h}^{-1}\mathfrak c_{Y,h}$ therefore obeys the
corresponding source-resolved bound multiplied by
$\|D_{cc,h}^{-1}\|$.  For example, if the four products in
\eqref{v7:Sobolev-tail} are bounded and this last inverse is $O(h^{-q})$,
then
\[
 h^2\|\upsilon_{Y,h}\|\le C h^{2r+1-q}.
\]
Thus $q<2r+1$ would suffice for the second-jet frame test of
\cref{v6:cofinal-frame}.  This is explicitly a consequence conditional
on those actual estimates, not a proof of them.  It still supplies neither
the finite-amplitude remainders nor the third-jet or lifespan bounds.

The noncircular part of this iteration is the localization itself: the
original active defect has only $O(N^2)$ input directions, and the
three effective outputs are exact pairings of specified boundary tails.
It is not enough to note that the input first jet is smooth.  Its active
solution and the adjoint active solutions may develop substantial Fourier
edge content when inverted; \eqref{v7:four-tails} retains precisely that
possibility.  Discarding the boundary would change the finite system and
is not used as a proof of convergence.

\section{Finite verification of the boundary mechanism and preparation interface}
\label{v7:finite-boundary}

At $N=5$, the component-major first-source matrix gives, in exact
cyclotomic residues, a parity defect of rank $98$.  Since
$125-3^3=98$, the upper bound of \cref{v7:seam-theorem} is attained
there.  Direct multiplication verifies its exact input and output
boundary support, oddness of all three constant weak columns, and
evenness of the complete prepared target.  Solving for the three adjoint
active columns in \eqref{v7:primal-adjoint} gives
\begin{equation}
 \left(\ip{\mathcal E y_{e_i}}v
                  -\ip{u_{e_i}}{\mathcal E x}\right)_{i=1}^3
 \equiv(2,79,185)^T
 \equiv2\mathfrak c_{Y,1/5}\pmod{241}.
 \label{v7:boundary-residue}
\end{equation}
All products in this check use the same spatial-sum convention.  It checks
the signs and basis interface of the written general identity, rather
than treating a modular zero as a proof of an all-cutoff cancellation.


\begin{theorem}[Exact five-site preparation isolates the remaining constant rate]
\label{v7:five-normal}
At the balanced five-site seed, the preparation derivative
$\mathcal T_Y$ of \cref{v5:drift-linearization} maps the configurations
satisfying the linear scalar constraint onto all $124$ nonconstant weak
drift directions.  After repairing the four mean constraints with the
retained momentum directions, there are analytic exact preparations such
that
\begin{equation}
 \boxed{\mathscr C_{1/5}(X(a))=0,\quad X(a)=aY+O_{1/5}(a^2),\quad
       \mathbb P_{\nabla,0}\partial_t\beta_{1/5}(a,0)=0.}
 \label{v7:five-normal-eq}
\end{equation}
Their leading weak rate is constant,
\[
 \partial_t\beta_{1/5}(a,0)=-a^{-1}\upsilon_{Y,1/5}+O_{1/5}(1),
 \qquad \upsilon_{Y,1/5}\ne0,
\]
with the coefficient in \eqref{v7:five-residues}.  The result persists on
a sufficiently small regular balanced parameter neighborhood at this
fixed cutoff.  It is an initial condition on the original fields, not a
projection of their dynamical vector field.
\end{theorem}
\begin{proof}
The scalar linear-constraint matrix $C_1$ has $125$ rows, $750$ raw
symmetric configuration columns, and rank $124$.  Its last row is minus
the sum of the others.  Choose the lexicographically first $124$ pivot
columns and invert the first $124$ rows on those columns.  This minor
has residue $60$ modulo $241$.  Put identity entries on the $626$ free
columns and minus that inverse times their constraint columns on the
pivot rows.  This defines a characteristic-zero kernel lift $Z_*$ with
$C_1Z_*=0$; it is not a nullspace guessed from a modular zero.

Evaluate the exact drift derivative \eqref{v5:T-functional} on these
columns and all weak tests, keeping the physical test order.  The three
constant rows are identically zero by \cref{v5:constant-B2}.  On the
other $124$ rows select the first independent $124$ columns.  Their
complete list is in the certificate and their determinant has residue
$158$ modulo $241$.  All denominators are units, so this proves
surjectivity in characteristic zero.  Re-evaluation at $N=3$ with the
same method reproduces the unchanged residue $234$ of
\eqref{v5:Tminor}; the ordering correction does not change this
configuration-to-weak map.

The four mean derivatives on the balancing momentum directions form
\[
 \diag(3s,-\omega_hm_1,-\omega_hm_2,-\omega_hm_3),
 \qquad s=-\omega_h,\quad m_i=2.
\]
This matrix is invertible.  The momentum-slot
cancellation of \cref{v5:momentum-zero} applies at every odd cutoff, so
this repair preserves the $124$-direction control.  The inherited
range-and-mean initializer then realizes any fixed such second-jet change
as an analytic exact preparation; its proof uses no special feature of
$N=3$ beyond the previously supplied control minor.

The leading weak coefficients consequently fill
$\{u:R_cDu=\mathfrak c\}$.  Since $D_{cc}$ is invertible by
\cref{v7:five-rank}, the unique constant-supported member is
$I_cD_{cc}^{-1}\mathfrak c$.  Choose the corresponding finite second
control.  The full scaled rate equation is analytic through $a=0$,
with derivative on the $124$ control slots
\[
 -R_gD^{-1}I_g\,(\mathcal T_YZ_*)_{g,\mathcal K}.
\]
Its first factor is invertible by the Schur complement of $D_{cc}$ in
$D$, and the second by the new minor.  The analytic implicit-function
theorem therefore solves the \emph{entire actual} initial gradient-rate
equation, not just its first coefficient, giving
\eqref{v7:five-normal-eq}.  Expanding the exact rate gives the constant
pole.  The original constraints have been solved at every step by the
unchanged initializer.  Nonzero finite determinants persist in a small
balanced parameter neighborhood; all constants and amplitude radii
remain cutoff-dependent.
\end{proof}

The five-site normal form supplies an additional actual branch to which
\cref{v6:initial-comparison,v6:metric-jet} apply.  In the same residue
convention, each spatial-sum cubic translation row has value $153$ and
its contraction with $(88,32,86)$ is $188\ne0$.  Thus the isolated
constant-rate contribution to the phase-frame initial canonical
acceleration is again nonzero of order $a$, rather than identically
removed.  This is an extension to a finer certified cutoff, not a
cutoff-uniform second- or third-jet theorem.


\section{Status and reproducibility}
\label{v7:status}

The source and target ordering is now checked against every original
multiplier equation before an active solve is used.  The corrected
three-site invariant remains nonzero; the numerical vector is different
from the one retained in V5.  A new full five-site signed certificate shows
that source neutrality, signed regularity and the next nonzero effective
constant source coexist away from the exceptional three-site aliases.
The all-cutoff analytical result is sharper than that finite observation:
for $N\ge5$, the effective source is an exact Fourier-boundary pairing,
and the active parity-mixing operator has surface rather than volume rank.

The next analytical target is therefore the boundary tails of the actual
active and three adjoint active solutions, together with the corresponding
constant Schur block.  The complete weak inverse and its trace-angle
response remain needed for full electric and nonlinear control; a bound
on $\mathfrak c$ alone would not replace them.  No all-cutoff regular rank
family, bound on the true nonlinear Hessian, propagated initial-rate
condition, common physical interval, or raw-action Einstein limit is
claimed.  The existing physical clock and literal observations are unchanged.

The package contains the full five-site source and drift matrices,
explicit selected columns, the four determinant certificates, a complete
covector-order test, and the primal/adjoint boundary arrays.
Its faster coefficient differentiation is checked against every inherited
three-site source column, the independently recomputed complete drift,
and forward coefficient evaluations.  Nonzero residue certificates use
good-prime maps of the explicitly specified cyclotomic fields; zero
structural claims rest on the written grading and support proofs.
Numerical linear algebra in a floating physical embedding is not used to
assert a real rank or a uniform inverse.

The six preceding bodies, including their historical numerical values,
remain byte-for-byte unchanged.  \Cref{v7:corrected-three} supplies the
explicit correction needed to use their downstream numerical tests.
The package records preservation hashes and verification scope.  These
computational checks do not constitute a proof-assistant formalization or
an independent reproof of every inherited analytic stationary estimate.
% END IMMUTABLE EXACT-ACTION BRIDGE V7 BODY

% Future iterations append here.
% BEGIN IMMUTABLE EXACT-ACTION BRIDGE V8 BODY
\clearpage
\renewcommand{\thesection}{\arabic{section}}
\renewcommand{\theHsection}{bridgeV8.\arabic{section}}
\setcounter{section}{59}
\section{An upstream phase selection rather than an assumed tail estimate}
\label{v8:context}

The Fourier-boundary formula of \cref{v7:boundary-identity} is an exact
localization, not a bound on the inverse boundary tails.  Supplying those
bounds as hypotheses would not close another part of the programme.  This
body takes a different upstream step: it changes the phases of the balanced
\emph{initial profiles}, retains their radii and scalar-momentum calibration,
and tests the actual effective constant source on that family.  Neither the
finite action nor its normalized Euler flow is changed.

The principal result is a computer-assisted fixed-cutoff existence theorem.
At $N=3$, a specified small rational phase-coordinate box contains a unique
zero of all three components of the original effective source.  Its phase
Jacobian, the complete leading signed Gram, and the nonconstant preparation
control remain nonsingular.  The zero is therefore not obtained by approaching
a singular source.  Combining these facts with the original analytic initial
constraint map gives exact preparations whose \emph{entire initial curl-free
shift velocity}, including its three means, is zero.  This removes the
inverse-amplitude initial weak-rate pole on a new first-jet family; it does not
propagate that condition or supply a cutoff-independent lifespan.

\paragraph{Inherited inputs and genuinely new work.}
The balanced weak quadratic cancellation and mean calibration are
\cref{v4:neutral,v4:preparation}.  The original weak coefficient construction,
prepared active drift, cubic translation drift and preparation derivative are
\cref{v3:compiler,v5:prepared-B2,v5:cubic-means,v5:T-functional}.
Their formulas are used with the corrected physical test order of
\cref{v7:ordering}.  They are not re-proved as advances.  The new ingredients
are the phase-parameter identities, a verified real zero of the complete
three-component effective source, regularity and preparation-control
enclosures on its neighborhood, and the joint analytic preparation theorem.
In particular, the computer calculation below evaluates the raw cubic means;
it does not silently assume that their special zero-phase cancellation holds
at another phase.

Interval inclusion and preconditioned contraction are established numerical
verification methods; see, for example, \cite{v8:Rump}.  A direct proof of the
specific inclusion test used here is provided.  Floating searches only propose
a center and preconditioners.  They are not used as existence or rank proofs.
All seven preceding research bodies, including the correction in V7, remain
verbatim.

\section{Balanced phases and exact parameter symmetries}
\label{v8:phases}

Use the transverse tensors $T_i,P_i$ of \cref{v6:balanced}, with squared
Frobenius norm two, $T_iP_i+P_iT_i=0$, and zero mutual Frobenius pairing.
For positive radii $r_i$ and phases $\phi_i$ put
\begin{equation}
 \begin{split}
 U_{r,\phi}(x)&=\sum_{i=1}^3r_i\{T_i\cos(2\pi x_i+\phi_i)
                                +P_i\sin(2\pi x_i+\phi_i)\},\\
 p_{r,h}&=p_0 I,\qquad
 p_0=\omega_h\sqrt{(r_1^2+r_2^2+r_3^2)/3},\qquad
 Y_{r,\phi,h}=(U_{r,\phi},p_{r,h}).
 \end{split}
 \label{v8:family}
\end{equation}
The sign of $p_0$ fixes the retained momentum branch.  The displayed positive
choice is used in the numerical certificate.  Every member has the original
zero weak quadratic drift and the same four-mean compatibility mechanism.
These properties follow from the existing balanced-family results and do not
require the phases to be close to zero.

At a regular active first jet define the effective three-vector in physical
spatial-mean normalization by
\begin{equation}
 \mathfrak c_h(r,\phi)[c]
 =t_c(Y_{r,\phi,h})
  -\ip{\mathcal A_h^*K_0^{-1}W_{c,h}}
         {\Gamma_{A,h}^{-1}b_{A,h}}_h,
 \label{v8:effective}
\end{equation}
where $t_c$ is the \emph{actual} preparation-independent cubic drift of
\eqref{v5:t}.  It is not the cubic momentum row of \cref{v6:P3}.

\begin{proposition}[Site-periodicity and a simultaneous phase reversal]
\label{v8:phase-symmetry}
Whenever the relevant active source is regular, for $\ell\in\mathbb Z^3$,
\begin{align}
 \mathfrak c_h(r,\phi+2\pi\ell/N)&=\mathfrak c_h(r,\phi),
                    \label{v8:site-periodicity}\\
 \mathfrak c_h(r,\phi+\pi(1,1,1))&=-\mathfrak c_h(r,\phi).
                    \label{v8:odd-phase}
\end{align}
Regularity is preserved by both transformations.  On an odd grid the second
transformation is equivalent, modulo the first, to translation of all three
phase parameters by $\pi/N$.
\end{proposition}
\begin{proof}
A phase increment $2\pi\ell/N$ is an actual permutation of the finite spatial
sites.  It commutes with the original derivatives and literal shifts and
preserves cyclic products, the spatial sum and the physical pairings.  Apply
that permutation to all field, source and multiplier slots in the defining
coefficient formulas.  The active Gram changes by an orthogonal conjugation;
its target and the constant weak source change by the corresponding
isometries.  A constant multiplier test itself is unchanged.  The scalar
pairing in \eqref{v8:effective} and its raw cubic term are consequently
unchanged.  This proves \eqref{v8:site-periodicity}, including regularity.

For the second transformation $U$ changes to $-U$, while the constant scalar
momentum is fixed.  Let $D(n,\beta)=(-n,\beta)$ on the physical active slice.
The original first-source formula has a lapse part linear in $U$ and a shift
part linear in the scalar momentum.  Thus $\mathcal A_-=\mathcal A D$.
The prepared quadratic target has types $Up$ in its lapse part and $U^2,p^2$
in its shift part, by the coefficient grading used in
\cref{v7:target-even}; hence $b_{A,-}=Db_A$.
The constant weak column has type $Up$, by
\cref{v7:source-grading}, so $W_{c,-}=-W_c$.
It follows that $\Gamma_{A,-}=D\Gamma_A D$ and that the actual active
connection response $K_0^{-1}\mathcal A\Gamma_A^{-1}b_A$ is unchanged.
Finally the formula \eqref{v5:t} has types $U^3$ and $Up^2$ on this family:
$H_3$ has these types, $P_{c,3}$ has types $U^2p,p^3$, and
$F_1=(V(p),\mathcal F(U))$.  It follows directly, before any frequency
cancellation, that $t_c(-U,p)=-t_c(U,p)$.  Substitution gives
\eqref{v8:odd-phase}.  Since $N$ is odd,
$\pi-\pi/N=2\pi(N-1)/(2N)$ is an integer number of site periods.
\end{proof}

These identities do not force a simultaneous zero of three components.
For example $(\cos N\phi_1,\sin N\phi_1,0)$ has both parameter symmetries
on an odd grid and is nowhere zero.  A sign change of a vector is not a
three-dimensional intermediate-value theorem.  The zero established below
therefore requires the actual source calculation and its validated Jacobian.

\begin{lemma}[Phase choices can be represented in a shrinking physical cell]
\label{v8:small-cell}
For any phases at an odd cutoff, a site translation puts their representatives
in $[-\pi/N,\pi/N]^3$.  For radii in a fixed compact set and every fixed
Sobolev order $s$, the associated first configuration differs from
$U_{r,0}$ by at most $C_s h$ in $H^s$, after trigonometric interpolation on
sufficiently fine grids.  The scalar momentum is unchanged.
\end{lemma}
\begin{proof}
Choose a nearest multiple of $2\pi/N$ independently in each coordinate.
The residual phases have the stated bound.  The configurations have only the
six fixed unit axial frequencies, and their coefficient derivatives with
respect to phase are bounded uniformly in the radii.  The fundamental theorem
of calculus on that fixed finite Fourier space gives the estimate.
\end{proof}

In particular, if a regular phase-selected family is eventually proved on
all relevant cutoffs, its phase selection need not change the limiting
first data.  The lemma does not assert that such zeros or uniform inverses
exist at those cutoffs.

\section{The original finite map whose zero is verified}
\label{v8:map}

For the certificate fix $N=3$, $h=1/3$, $r=(1,1,1)$ and
$p=3\sqrt3 I$.  Introduce rational phase coordinates
$\vartheta_i=\tan(\phi_i/2)$ and write
\begin{equation}
 C_i(\vartheta)=\frac{1-\vartheta_i^2}{1+\vartheta_i^2},\qquad
 S_i(\vartheta)=\frac{2\vartheta_i}{1+\vartheta_i^2}.
 \label{v8:rational-phase}
\end{equation}
The cosine and sine profile coefficients in \eqref{v8:family} are then
$C_iT_i+S_iP_i$ and $-S_iT_i+C_iP_i$, respectively.  Their equal norm and
orthogonality are exact rational identities.  The sampled base sine has
values $(0,\sqrt3/2,-\sqrt3/2)$, the base cosine has values
$(1,-1/2,-1/2)$, and the original derivative is precisely
$3(S_i-S_i^{-1})$.  No numerical trigonometric evaluation enters the
verified coefficient map.

Use the raw symmetric dual coordinates and inverse Cartan matrix in
\eqref{v3:raw-K}.  Let $C_A(\vartheta)$ be the $162\times78$ first-source
matrix on the fixed column list \eqref{v3:pivots}; all multiplier columns
have the physical component-major order \eqref{v7:ordering}.  Put
\begin{equation}
 \begin{gathered}
 A=C_A^TK^{-1}C_A,\qquad
 W_c=(\mathcal Q_h(Y;e_1),\mathcal Q_h(Y;e_2),\mathcal Q_h(Y;e_3)),\\
 L_c=C_A^TK^{-1}W_c,\qquad x=A^{-1}b_A,\qquad
 \boldsymbol f(\vartheta)=\frac1{27}\{\boldsymbol t(\vartheta)-L_c^Tx\}.
\end{gathered}
 \label{v8:f}
\end{equation}
Here $b_A$ consists of the \emph{prepared} quadratic drift on the same
78 physical columns, including the lapse correction in
\eqref{v5:prepared-drift-formula}.  The three entries of $\boldsymbol t$
are the spatial-sum versions of \eqref{v5:t}, computed from the full cubic
constant-shift envelope.  The factor $1/27$ converts the final covector to
spatial mean.  The source Gram and target use the same spatial-sum convention
before solving.  This removes no action term or source equation.

The map is rational in $\vartheta$ and $\sqrt3$ wherever $A$ is invertible.
The remaining weak columns use exactly the three constants followed by the
26 phase gradients of point differences against the last site.  With their
full matrix $W$, write
\begin{equation}
 F=(C_A,W),\qquad \Gamma=F^TK^{-1}F.
 \label{v8:full-Gram}
\end{equation}
The preparation-control minor is $M=(\mathcal T_Y Z_*)_{g,\mathcal K}$,
where $Z_*$ is the rational linear-constraint kernel lift of
\cref{v5:control} and $\mathcal K=(0,\ldots,23,25,28)$.
The configuration controls are the same fixed rational directions; their
values under $\mathcal T_Y$ are recomputed at the new phases.

\section{A verified regular zero of all three effective components}
\label{v8:root}

The center below is specified by exact dyadic rationals, not rounded phase
angles:
\begin{equation}
 \vartheta_0=
 \left(
 \frac{4886567754969805}{18014398509481984},\quad
 \frac{7638139748043125}{36028797018963968},\quad
 \frac{4319752166354399}{288230376151711744}
 \right).
 \label{v8:center}
\end{equation}
Set
\begin{equation}
 B=\{\vartheta:\|\vartheta-\vartheta_0\|_\infty\le2^{-35}\},\qquad
 B_+=\{\vartheta:\|\vartheta-\vartheta_0\|_\infty\le2^{-24}\}.
 \label{v8:boxes}
\end{equation}

\begin{theorem}[Computer-assisted phase root on the unchanged three-site action]
\label{v8:root-theorem}
The actual map \eqref{v8:f} has a unique zero $\vartheta_*$ in $B$.
Its derivative is nonsingular throughout $B$.  Throughout the larger box
$B_+$ the full signed Gram \eqref{v8:full-Gram} and the $26\times26$
preparation-control minor $M$ are nonsingular.  The active Gram is
nonsingular on $B$.  In particular the first jet at $\vartheta_*$ has
\begin{equation}
 \boxed{d_{2,h}(Y_*)=0,\qquad \mathfrak c_{Y_*}=0,\qquad
        \rank F=107,\qquad \rank(\mathcal T_{Y_*}|_{\mathcal Z_u})=26.}
 \label{v8:root-conclusions}
\end{equation}
The approximate phases, supplied only for interpretation, are
\[
 \phi_*\simeq(0.529769838,\ 0.417815742,\ 0.029972058)
 \quad\text{radians}.
\]
These statements are at $h=1/3$, not on an unproved cofinal rank family.
The exact phase is specified by its defining equations and isolating box;
a finite terminating binary expansion or an automatic microscopic selection
rule is not asserted.
\end{theorem}

\begin{proof}
We give the numerical inclusion conditions and their mathematical
justification.  Every bound in the following table is rounded \emph{upward}
from the interval certificate.  The files containing the complete enclosures
and the exact dyadic preconditioners are part of the accompanying package.
The symbol $R$ denotes the fixed $3\times3$ approximate inverse of the
phase Jacobian, encoded as exact binary rationals.  It is not assumed to be
an exact inverse.
\begin{center}\small
\begin{tabularx}{\textwidth}{@{}Yr@{}}
\toprule
Verified quantity & Strict upper bound\\\midrule
$\|R\boldsymbol f(\vartheta_0)\|_\infty$ & $9.0\,10^{-12}$\\
$\sup_{\vartheta\in B}\|I-RD\boldsymbol f(\vartheta)\|_\infty$ & $0.551$\\
$\|R\boldsymbol f(\vartheta_0)\|_\infty+
 2^{-35}\sup_B\|I-RD\boldsymbol f\|_\infty$ & $2.51\,10^{-11}$\\
Residual certifying the right source preconditioner on $F$ & $1.8\,10^{-14}$\\
$\sup_{B_+}\|I-R_\Gamma(S^T\Gamma S)\|_\infty$ & $0.083$\\
$\sup_{B_+}\|I-R_M M\|_\infty$ & $0.024$\\
\bottomrule
\end{tabularx}
\end{center}
Here $S$ is a fixed invertible right source preconditioner, whose
invertibility is itself verified by a residual against an invertible
triangular comparison.  Thus the Gram test concerns the original signed
source in changed coordinates, not a positive replacement.  The active
inverses used in evaluating \eqref{v8:f} are independently certified by
\cref{v8:linear-enclosure}.

Let $T(\vartheta)=\vartheta-R\boldsymbol f(\vartheta)$.  The second row
of the table is less than one, and the third is strictly less than
$2^{-35}$.  The mean-value integral therefore makes $T$ a strict contraction
of the complete closed cube $B$ into itself.  It has one fixed point there.
The derivative residual also implies that $RD\boldsymbol f$ is invertible;
since these are square matrices, $R$ and $D\boldsymbol f$ are invertible.
A fixed point is consequently a zero of $\boldsymbol f$, and every zero in
$B$ is that fixed point.  This proves existence, uniqueness and the derivative
assertion without assuming a zero exists.

The last two residual bounds are less than one.  Their Neumann-series
criterion proves that the full original signed Gram and the control minor
are nonsingular for \emph{every} parameter in $B_+$, hence also at the root.
A nonsingular signed Gram implies independent source columns.  The
preparation-control upper rank is 26 because constant quadratic translation
drifts vanish identically; the nonzero 26-dimensional minor attains it.
The leading weak drift vanishes by the balanced-family theorem, while the
remaining effective three-vector is zero by the root just proved.
\end{proof}

The table is an inclusion certificate, not a list of floating singular values
or residuals without inverse bounds.  No equality to zero is inferred from
roundoff.  In particular $\boldsymbol t$ is evaluated in the original
coefficient formula in the root enclosure; its contribution is not replaced
by a presumed phase symmetry.

\section{What the interval certificate encloses}
\label{v8:validation}

For reproducibility the coefficient and inverse steps are made explicit.
A real interval represents an enclosure of an exact quantity; it is not a
probability interval.  Each basic addition, subtraction, product, division
by an interval excluding zero, and finite sum is rounded outward.
The implementation uses extended-precision binary arithmetic with at least
64 significand bits and takes the next representable number in the outward
direction after each elementary operation.  Matrix contractions are finite
sums of these interval operations, not unchecked floating BLAS products.
The code checks the precision model, and the endpoints enclosing $\sqrt3$
are verified by rational squaring inequalities.  The proposed center and
preconditioners are exact dyadic input constants.

\begin{lemma}[A posteriori linear enclosure used in the certificate]
\label{v8:linear-enclosure}
Let $[A],[b]$ enclose finite real matrices and right sides, and let $R_0,x_0$
be fixed real point matrices of the required sizes.  Suppose
\[
 q\ge\sup_{A\in[A]}\|I-R_0A\|_\infty,\qquad q<1.
\]
Then every $A\in[A]$ is invertible, and every solution of $Ax=b$ satisfies
\begin{equation}
 \|x-x_0\|_\infty
 \le\frac{\sup_{A\in[A],\,b\in[b]}\|R_0(b-Ax_0)\|_\infty}{1-q}.
 \label{v8:linear-bound}
\end{equation}
Columnwise bounds apply for several right sides.  If $M$ bounds
$|I-R_0A|$ entrywise and $r$ bounds $|R_0(b-Ax_0)|$, an already valid error
radius vector $d$ may be replaced by $r+Md$.
\end{lemma}
\begin{proof}
The strict residual implies that $R_0A$ is invertible, so both square factors
are invertible.  With $e=x-x_0$, the exact equation is
$e=R_0(b-Ax_0)+(I-R_0A)e$.  Its convergent Neumann series gives
\eqref{v8:linear-bound}.  Taking entrywise absolute values proves the last
statement.  Dependence among the intervals can only make their product box
larger; it cannot invalidate this enclosure of the actual simultaneous data.
\end{proof}

The coefficient evaluator is a direct interval transcription of
\eqref{v7:reverse-aux}--\eqref{v7:reverse-source} and
\eqref{v5:coefficient-functional}.  It uses a finite formal polynomial in the
field amplitude, and retains the required cubic electric and magnetic
terms.  The first active source is also evaluated from the explicit
first-source formula.  Physical gradients are taken by the chain rule on
that finite polynomial; they are not differences of finite-amplitude
solutions.  A first-order three-parameter jet is propagated alongside every
interval value, so phase derivatives are enclosed before any nonlinear
solve is evaluated.  Induction over the finite arithmetic graph proves
inclusion of the exact value and these derivatives: forward operations use
the ordinary product and chain rules, and reverse adjoints use their exact
transpose rules with the same outward rounding.  Slicing, permutations and
literal shifts are exact index operations.  The raw off-diagonal dual slots
retain their Frobenius factors.

For the active solve, phase differentiation is implemented through
\begin{equation}
 A x=b_A,\qquad
 A\,\partial_jx=\partial_j b_A-(\partial_j A)x,\qquad
 \partial_j\boldsymbol f
 =\frac1{27}\{\partial_j\boldsymbol t-(\partial_j L_c)^Tx
                                    -L_c^T\partial_jx\}.
 \label{v8:derivative-solve}
\end{equation}
Every occurrence of a solve in \eqref{v8:derivative-solve} uses
\cref{v8:linear-enclosure}.  No derivative of an approximate inverse is
substituted for these equations.  The full signed-Gram test first changes
source coordinates by a fixed $S$, and encloses $(FS)^TK^{-1}(FS)$.
The independently checked right-preconditioner residual proves $S$ invertible.
This avoids unnecessary interval inflation from poorly scaled source columns
while preserving the exact signed problem.

The preparation minor is evaluated independently from
\eqref{v5:T-functional}, or equivalently by differentiating its complete
quadratic weak-drift polarization.  The kernel lift is defined by the exact
rational $26\times26$ linear-constraint minor in \cref{v5:control}; its
identity $C_1Z_*=0$ is checked in rational arithmetic.  Thus its columns are
actual constraint-compatible variations, not a numerical nullspace.

The following logical distinction is essential.  Cached enclosure files allow
a short independent check of all strict inclusion inequalities.  Regenerating
the enclosures additionally checks their incidence with the original scalar
coefficient functional.  Both procedures are supplied.  This finite
computer-assisted proof does not independently verify the entire inherited
analytic stationary-chart theory or the earlier continuum convergence results.

\section{Joint phase and second-preparation control removes every weak initial rate}
\label{v8:joint-preparation}

The verified root is useful because its phase Jacobian controls exactly the
three directions inaccessible to second preparation alone.  We now prove the
corresponding exact, not merely formal, initial-data statement.

\begin{theorem}[Exact original-action preparation with zero curl-free shift velocity]
\label{v8:exact-preparation}
At $h=1/3$ there is an analytic family of original exactly prepared records
$X(a)$, for sufficiently small nonzero $a$, such that
\begin{equation}
 \boxed{
 \mathscr C_h(X(a))=0,\qquad
 X(a)=aY_*+O_h(a^2),\qquad
 \mathbb P_{\mathrm{cf}}\partial_t\beta_h(a,0)=0.}
 \label{v8:full-weak-zero}
\end{equation}
The last observation includes the three constant shifts and every mean-zero
curl-free shift.  The derivative is the initial rate of each record's
\emph{own normalized original-action continuation}.  No multiplier rate is
prescribed after solving, and no field force or action term is changed.
The complete initial multiplier velocity is bounded as $a\to0$ on this
family.  Neither a common physical interval nor propagation of
\eqref{v8:full-weak-zero} is asserted.
\end{theorem}
\begin{proof}
For phases near $\vartheta_*$ take the original analytic range-and-mean
initializer with the balanced first jet $Y(\vartheta)$.  The four mean
calibration derivatives are nonzero on its homogeneous trace and three
quadrature momentum directions.  The construction therefore gives an
analytic exactly constrained family through each such first jet.

Choose the 26 rational configuration directions in the verified minor $M$.
For each phase, repair their four next mean constraints using the current
balancing momentum directions.  The momentum-slot cancellation of
\cref{v5:momentum-zero} leaves their weak-drift derivative unchanged.
The analytic extension argument of \cref{v5:extend} consequently supplies
exact records $X(a,\vartheta,z)$ with these second-jet controls
$z\in\mathbb R^{26}$.  Any fixed finite control is admitted after reducing
the amplitude interval; no radius for unbounded controls is asserted.

At $a=0$ the leading weak target is zero for every phase.  Let $D(\vartheta)$
be the full 29-dimensional weak Schur complement of the original signed
Gram and set
\[
 g_0(\vartheta,z)
 =b_3(\vartheta,z)-L(\vartheta)^TA(\vartheta)^{-1}b_A(\vartheta).
\]
Its three constant covector components are
$R_cg_0=27\boldsymbol f(\vartheta)$; the factor reflects the spatial-sum
coordinates.  They do not depend on $z$.  Its other components change
under $z$ by the invertible matrix $M(\vartheta)$, by the actual drift
Taylor identity in \cref{v5:invariant}.  Thus, at the verified phase
$\vartheta_*$, a finite $z_*$ makes \emph{all} components of $g_0$ zero.
We do not need to know or assume that the baseline nonconstant cubic drift
is zero: it is the finite right side of this invertible control equation.

On the normalized exact-action initial chart, write the true graded target
and matrix as $\widehat b(a,\vartheta,z)$ and $\Gamma(a,\vartheta,z)$.
Their inherited analytic expansions give
\[
 \widehat b=a\,b_1(\vartheta,z)+O_h(a^2),\qquad
 \Gamma=\Gamma_0(\vartheta)+O_h(a).
\]
The verified Gram is regular near $\vartheta_*$, so
$\widetilde u=a^{-1}\Gamma^{-1}\widehat b$ extends analytically through
$a=0$.  Its value there is the true first surviving coefficient, with
weak part $D^{-1}g_0$.

Let $Q_W$ extract coordinates of the physical projection onto the full
curl-free shift space.  Chosen active representatives need not be orthogonal
to that space.  Exactly as in \eqref{v5:scaled-gradient}, the actual rate
satisfies an analytic scaled equation
\begin{equation}
 \Psi(a,\vartheta,z):=aQ_W\lambda_t
       =-\widetilde u_W-aJ_A\widetilde u_A,
 \qquad \Psi(0,\vartheta,z)=-D^{-1}g_0(\vartheta,z).
 \label{v8:scaled-rate}
\end{equation}
The fixed synthesis matrices are retained in $J_A$.  The mean-lapse
normalization has zero shift observation.  At $(0,\vartheta_*,z_*)$,
$\Psi=0$.  Multiplying its derivative by $-D$ and ordering the three constant
rows before the 26 other rows gives a block triangular matrix
\begin{equation}
 \begin{pmatrix}
 27D\boldsymbol f(\vartheta_*)&0\\
 *&M(\vartheta_*)
 \end{pmatrix}.
 \label{v8:joint-Jacobian}
\end{equation}
There is no derivative of $D^{-1}$ multiplying a nonzero target at this
point, since $g_0=0$.  Both diagonal blocks are invertible by
\cref{v8:root-theorem}, and $D$ is invertible by the full Gram test.
The analytic implicit-function theorem yields $\vartheta(a),z(a)$ through
$(\vartheta_*,z_*)$ with $\Psi=0$.  For nonzero $a$ this is precisely the
last equation in \eqref{v8:full-weak-zero}.  The initial constraints have
already been solved identically.  Since $\vartheta(a)=\vartheta_*+O_h(a)$,
the first jet is $Y_*$.

Finally the full normalized solution has $u=a\widetilde u$.  Its active
rate, weighted by $a^{-1}$, is bounded.  At the selected root
$\widetilde u_W(0)=0$, and analyticity along the selected curve gives
$\widetilde u_W=O_h(a)$.  Its weak rate, weighted by $a^{-2}$ before the
factor $u=a\widetilde u$, is therefore bounded as well.  Invertibility of
the true finite-amplitude Gram persists for sufficiently small nonzero
amplitude, so the retained exact-action continuation theorem supplies
each local history.  This argument does not bound that interval uniformly.
\end{proof}

The phase adjustment in this theorem is an initial-data selection.  It is not
a new refresh rule, a projected Hamiltonian vector field, or a dynamical law
that drives arbitrary phases to the zero.  In particular it does not
contradict the preparation invariant at the original zero-phase first jet:
that first jet has a different, nonzero effective vector.  Higher preparation
alone still cannot change that vector.

\section{Persistence in the radii and the initial physical readouts}
\label{v8:persistence}

\begin{corollary}[A locally persistent balanced first-jet class]
\label{v8:radii}
There is a neighborhood of $r=(1,1,1)$ in which a unique nearby analytic
phase choice $\vartheta_*(r)$ makes the effective three-vector zero, using
the calibrated scalar momentum in \eqref{v8:family}.  The full signed
source and the preparation-control minor remain regular after reducing this
neighborhood.  Every such first jet has exact initial preparations of
\cref{v8:exact-preparation}.
\end{corollary}
\begin{proof}
The original coefficient map is analytic in the radii, rational phase
coordinates and the positive square-root momentum calibration.  At the
certified point its three phase derivatives are independent.  The analytic
implicit-function theorem gives the asserted phase graph.  The nonzero
signed determinant and control minor persist by continuity.  The preceding
joint initial-preparation argument then applies to each nearby parameter.
\end{proof}

This describes a three-dimensional phase-selected graph in the six-dimensional
balanced radius--phase family at a fixed cutoff, not an open ball of arbitrary
first fields or transition kernels.  Nonzero independent radii retain all
three spatial directions; they do not supply an all-cutoff or generic
microscopic selection theorem.

\begin{corollary}[No inverse-amplitude initial weak pole on the selected branch]
\label{v8:physical}
Along each analytic preparation obtained above, the original initial literal
electric and spatial link curvature reads are $O_h(a)$, and the complete
anchored first metric time jet is bounded as $a\to0$.  The initial spatial
metric velocity is $O_h(a)$; lapse and transverse-shift velocities are only
asserted $O_h(1)$.  For the comparison path $\dot q=\overline\beta$ with
$q(0)=0$, one has $\dot q(0)=\ddot q(0)=0$.  The initial phase-versus-geometric
second-jet discrepancy in \eqref{v6:metric-exact} is consequently zero.
\end{corollary}
\begin{proof}
The leading weak target is neutral, and the graded signed response is
regular near the selected first jet.  The same original response estimate
as in \cref{v4:initial-electric} therefore gives an order-$a$ initial
electric read; the analytic first stationary connection gives the spatial
bound.  No positive minimum response replaces the signed solution.
The final paragraph of \cref{v8:exact-preparation} bounds the actual full
multiplier rate.  At unit lapse and zero shift on the initial cut, the ADM identities
$\dot g_{00}=-2\dot N$ and $\dot g_{0i}=\gamma_{ij}\dot\beta^j$ then give
a bounded anchored first time jet.  The canonical spatial velocity is
$O_h(a)$ on its analytic chart.  The projection in
\eqref{v8:full-weak-zero} includes constant shifts, so
$\overline{\dot\beta}=0$ exactly.  Inserting this in the established
frame-comparison identity proves the last assertion.
\end{proof}

No smallness of the remaining bounded lapse or transverse rate,
no bound on $\ddot\beta$, and no differentiated spacetime-curvature estimate
is inferred.  The result removes an identified initial obstruction before
those dynamical questions are approached.

\section{Remaining cofinal work and reproducibility}
\label{v8:status}

The new result satisfies an actual source condition on an explicit regular
first-jet class.  It is therefore stronger than merely restating that a
vanishing effective vector would remove a leading pole.  The exact
three-component zero is enclosed, its Jacobian and the complete signed
source are verified regular, and their interaction with the original
constraint initializer gives exact full weak-rate cancellation at the cut.
The previously fixed zero-phase obstruction remains valid at its own first
jet.  The action, phase derivative, ordered plaquettes, physical clock and
literal observations have not changed.

The next cofinal question is whether regular zeros can be selected on all
relevant cutoffs with quantitative bounds on the phase derivative, signed
source and preparation controls.  The site-periodicity lemma shows how to
retain the same limiting first profiles if such zeros exist; it does not
establish their existence.  One must then control the true finite-amplitude
Hessian, differentiated sources, the remaining bounded initial rates and a
common physical interval.  Exact initial weak-rate cancellation is not
proved invariant under the raw Euler flow.  The V7 boundary mechanism may
still be useful for the cutoff dependence of the phase map, but its inverse
tails are not bounded here by assuming this new initial selection.
No raw-action Einstein continuum limit is claimed.

\paragraph{Verification scope.}
The package includes the outward-rounded interval arithmetic and reverse
coefficient differentiation, the exact dyadic center and preconditioners,
point and parameter-box enclosures of the original source map and its phase
Jacobian, all 29 weak columns on the larger box, and the independent
26-direction preparation-control enclosure.  The control lift is regenerated
and checked in rational arithmetic.  The physical multiplier ordering is
checked against the corrected convention, not inherited from an unlabeled
array.  Elementary interval operations are tested against rational endpoint
calculations, and the coefficient values and parameter differentials are
cross-checked against independent floating formal-coefficient evaluations;
those latter checks are diagnostics, not the inclusion proof.

The computer-assisted theorem uses the strict contraction and inverse
residual inequalities, not the exploratory floating roots.  All enclosures
can be regenerated from the supplied scalar coefficient program.  A shorter
verifier checks the supplied enclosure inequalities, exact kernel identity,
and preservation hashes.  The analytic nonlinear initializer and the
amplitude-dependent phase/control functions are established by the written
implicit-function proof; no numerical nonlinear initial root or spacetime
trajectory is reported.  No new all-cutoff rank certificate is claimed.

\begin{thebibliography}{9}
\raggedright
\setcounter{enumiv}{9}
\bibitem{v8:Rump}
S. M. Rump, \emph{Solving nonlinear systems with least significant bit
accuracy}, Computing \textbf{29} (1982), 183--200.
\url{https://doi.org/10.15480/882.305}.
This supplies primary-source context for verified nonlinear inclusions;
the particular renewal phase map, interval certificate and joint
initial-preparation theorem are established here.
\end{thebibliography}
% END IMMUTABLE EXACT-ACTION BRIDGE V8 BODY


% Future iterations append here.
% BEGIN IMMUTABLE EXACT-ACTION BRIDGE V9 BODY
\clearpage
\renewcommand{\thesection}{\arabic{section}}
\renewcommand{\theHsection}{bridgeV9.\arabic{section}}
\setcounter{section}{67}
\section{The active source left by phase selection}
\label{v9:context}

The exact preparations of \cref{v8:exact-preparation} have zero initial
curl-free shift velocity.  Their remaining lapse and transverse-shift rates
were proved bounded only at a fixed cutoff.  Before trying to propagate that
initial selection, one must determine whether these bounded rates have an
acceptable refinement scale.  This body returns to the original quadratic
consistency equation rather than postulating a uniform differentiated
inverse.

There is an explicit mixed-axis transverse test on which the prepared
quadratic drift is of order $h$, uniformly in the phases of the balanced
profiles.  The corresponding row of the leading multiplier Hessian is of
order $h^2$ on the full multiplier space.  Consequently any bounded
amplitude-leading rate solving the original consistency equation has norm at
least $c/h$ on every sufficiently fine cutoff at which it exists.  No signed
inverse bound is assumed in this lower estimate.  It does not establish a
cofinal existence family or rule out a geometrically controlled change of
coordinates.

We also examine a natural upstream enlargement: transverse tensor momenta
proportional to each circular spatial profile.  These preserve the complete
weak quadratic neutrality and admit exact nonlinear initial preparation.
They can cancel the three displayed mixed-axis tests, but then unavoidable
single-axis tests remain.  A uniform order-$h$ active target survives on this
entire specified class.  Finally, at the verified phase-selected three-site
root, the canonical initial acceleration has a nonzero amplitude limit.  An
exact affine canonical transport generated by the original linear constraint
rows reduces the \emph{whole} transformed initial canonical acceleration to
$O_h(a)$.  Its observations must be pulled back; it is not a finite
relabelling or an exact diffeomorphism symmetry of the action.

\paragraph{Inputs and nonduplication.}
The complete coefficient functional, prepared drift, and linear canonical
operators are \cref{v3:compiler,v5:prepared-B2} and
\eqref{v5:H2}--\eqref{v5:linear-multiplier}.  The first-source formula,
all-cutoff weak identities, nonlinear initial-constraint initializer, and V8
root are inherited.  The order-$h$ \emph{cubic momentum} in \cref{v6:P3} and
the effective \emph{cubic Schur source} in \cref{v7:boundary-identity} are not
the \emph{quadratic active drift} calculated here.  The lower bounds below
are not consequences of either of those scalar observations.  The new
coefficient tests use the corrected component-major convention throughout.

The symbolic calculations are finite equalities in a Fourier coefficient
ring, with explicit contraction rules below.  Floating finite-grid checks
are independent diagnostics, not the all-cutoff proof.  The V8 interval root
and analytic stationary-chart theory are used as inherited results, not
claimed to be independently revalidated here.  Broken finite gravitational
symmetry provides context \cite{v6:BD}; the particular coefficient and
comparison statements below follow from the unchanged renewal action.

\section{A mixed-axis test of the complete original drift}
\label{v9:mixed}

Let $\theta_i=2\pi x_i+\phi_i$ and use exactly the tensors $T_i,P_i$ of
\cref{v6:balanced}, with all handedness signs positive.  Set
\begin{equation}
 U_i=r_i(T_i\cos\theta_i+P_i\sin\theta_i),\qquad
 U=U_1+U_2+U_3,\qquad p=p_0I,
 \label{v9:scalar-family}
\end{equation}
where $r_i>0$.  For compatible scalar-momentum preparation,
$p_0^2=\omega_h^2\sum_i r_i^2/3$, with either nonzero sign, and
$\omega_h=2h^{-1}\sin(\pi h)$.  Define
\[
 \begin{aligned}
 \alpha_h&=2\pi h,&
 g_{12}(h)&=1+\sin\alpha_h+\sin(2\alpha_h),\\
 g_{13}(h)&=2\cos\alpha_h-1,&
 g_{23}(h)&=1-\sin\alpha_h-\sin(2\alpha_h).
 \end{aligned}
\]
For $i<j$, let $k$ be the remaining coordinate and put
\begin{equation}
 \beta^{ij}=e_k\cos(\theta_i+\theta_j).
 \label{v9:mixed-test}
\end{equation}
These are physical, mean-zero, transverse shift tests: their only vector
component is perpendicular to their two-coordinate wavevector.  Their
spatial-mean squared norms equal $1/2$.  They are not weak curl-free tests.

Denote by $\widehat B_{2,h}(Y)$ the prepared quadratic drift covector of
\eqref{v5:prepared-drift-formula}.  Its value on a pure shift agrees with the
unprepared quadratic Poisson coefficient, because the linear shift drift
vanishes identically.

\begin{theorem}[Exact phase-independent mixed-axis drift]
\label{v9:scalar-test}
For \eqref{v9:scalar-family}, on every odd $N\ge3$,
\begin{equation}
 \boxed{
 \widehat B_{2,h}(Y)[(0,\beta^{ij})]
    =\frac{h\omega_h^4 r_i r_j}{16}\,g_{ij}(h).}
 \label{v9:mixed-scalar}
\end{equation}
In particular the $12$ value is strictly positive at every odd $N\ge5$.
For radii in a fixed compact subset of $(0,\infty)^3$, uniformly in phase,
\begin{equation}
 \widehat B_{2,h}(Y)[(0,\beta^{12})]\ge c h,
 \qquad
 h^{-1}\widehat B_{2,h}(Y)[(0,\beta^{12})]
            =\pi^4r_1r_2+O(h).
 \label{v9:scalar-lower}
\end{equation}
The $12$ value comes entirely from the cubic magnetic part of the original
coefficient functional.  On the $13$ test the electric and magnetic
contributions are both retained and partially cancel.
\end{theorem}
\begin{proof}
Use the full shift coefficient from \cref{v5:prepared-B2}:
\begin{equation}
 \widehat B_{2,(0,\beta)}(Y)
   =-DH_3(Y)[(\delta_i\beta_j+\delta_j\beta_i,0)]
                +DP_{(0,\beta),2}(Y)[F_1Y].
 \label{v9:recipe}
\end{equation}
Here $H_3$ and $P_{\nu,2}$ are exactly the scalar extractions of
\eqref{v5:coefficient-functional}.  In particular the temporal connection
in the second functional is $-\curl_\delta\beta/2$, which is nonzero for
\eqref{v9:mixed-test}.  Reusing the cancellation valid for curl-free tests
would remove terms that contribute here.

The quadratic shift drift is even in momentum.  Its field types are
$U^2$ and $p^2$: this follows either from the reversal grading in
\cref{v7:source-grading}, applied to the Poisson coefficient, or directly
from \eqref{v9:recipe}.  Constant $p$ cannot pair with a nonzero mixed test.
Only the cross term involving $U_i$ and $U_j$ can contribute.  Two field
frequencies and the test frequency have coordinate totals in
$\{-2,0,2\}$ when both axes are present.  No nonzero such total is divisible
by an odd $N\ge3$.  Same-axis pairs miss one coordinate of the test and also
have zero sum.  Thus the required contractions are ordinary zero-frequency
contractions even on the three-site grid.  Their phase factors cancel
exactly.  Other multiplier tests are not being declared unaliased.

Here are the complete contractions for unit radii.  Write
$z=\exp(2\pi\mathrm i h)$ and retain the normalized spatial mean.  The table
lists the geometric, cubic electric, and cubic magnetic parts of
\eqref{v9:recipe}, each divided by $h\omega_h^4$:
\begin{equation}
\begin{array}{c|ccc}
 ij&\text{geometric}&\text{electric}&\text{magnetic}\\ \hline
 12&0&0&[2-\mathrm i(z+z^2-z^{-1}-z^{-2})]/32\\
 13&0&-1/8&(z+1+z^{-1})/16\\
 23&0&0&[2+\mathrm i(z+z^2-z^{-1}-z^{-2})]/32.
\end{array}
\label{v9:configuration-table}
\end{equation}
To reproduce the table, substitute
$\widehat U_i(\pm e_i)=(T_i\mp\mathrm iP_i)/2$,
$(z_1)_j^{R_a}=-\epsilon_{abc}\delta_b U_{jc}/2$ and
$(z_1)_j^{B_a}=-(p_0/2)\delta_{ja}$ in the coefficient functional.
Use $F_1Y=(-p_0I,-\omega_h^2U/2)$, the flat auxiliary test
$\eta_0=(-\curl_\delta\beta/2,\delta n,w_\beta)$ with $n=0$, and
$w_{\beta,ij}=\delta_i\beta_j+\delta_j\beta_i$.

More explicitly, if $k=w_\beta$, the first term of the recipe is the
$a^3\tau$ coefficient of $\Phi_h$ after substituting
\[
 z=a z_1(Y)+a\tau z_1(k,0),\quad
 u=aU+a\tau k,\quad p=ap_0I,\quad
 A_0=0,\quad w=-ap_0I,\quad (n,\beta)=(1,0).
\]
The second term is minus its $a^2\tau$ coefficient with
\[
 (u,p)=a(Y+\tau F_1Y),\quad z=a z_1(Y+\tau F_1Y),\quad
 (A_0,w)=\eta_0(0,\beta^{ij}),\quad(n,\beta)=(0,\beta^{ij}).
\]
All these substitutions are in the same original functional.  The
coframe and velocity coefficients are those displayed in the proof of
\cref{v6:P3}.  Form each magnetic logarithm by the three-coefficient BCH
recursion in that proof, on its four ordered factors.  Multiplying the
listed Fourier coefficients and selecting the zero character gives
\eqref{v9:configuration-table}.  This is a finite polynomial identity; the
companion symbolic script repeats precisely these substitutions with the
higher derivative symbols independent until cancellation.  Quartic and
higher action terms cannot enter the two coefficients just specified.

Finally $z+z^{-1}=2\cos\alpha_h$ and
$z^j-z^{-j}=2\mathrm i\sin(j\alpha_h)$.  Sum the three columns, restore
$r_ir_j$, and obtain \eqref{v9:mixed-scalar}.  For $N\ge5$, both sines
in $g_{12}$ are positive.  The lower bound uses a positive lower bound
on the radii and on $\omega_h$.  Its expansion follows from
$\omega_h=2\pi+O(h^2)$ and $g_{12}=1+O(h)$.
\end{proof}

The nonzero pairing holds for the actual phases selected in V8.  There is
no contradiction with $d_{2,h}=0$: that earlier equality concerns the
curl-free shift projection, whereas \eqref{v9:mixed-test} is transverse.
Nor is it the cubic effective vector made zero by the phase root.

\section{A source-row lower bound with no inverse assumption}
\label{v9:rate-scale}

Let $\mathcal J_{1,h}$ be the complete first multiplier source on the
mean-lapse-zero slice, including all shift directions, and define
\[
 \mathcal M_{2,h}=\mathcal J_{1,h}^*K_0^{-1}\mathcal J_{1,h}.
\]
This is the leading coefficient of the original multiplier Hessian in
ungraded physical coordinates.  Its curl-free shift kernel is retained.

\begin{lemma}[The actual smooth test row has order $h^2$]
\label{v9:test-row}
For the scalar-momentum family \eqref{v9:scalar-family}, uniformly on compact
sets of radii and all phases, there are $C,h_0>0$ such that
\begin{equation}
 \boxed{\|\mathcal M_{2,h}(0,\beta^{12})\|_h\le Ch^2
                \qquad (0<h<h_0).}
 \label{v9:row-bound}
\end{equation}
The norm is in the complete physical multiplier space.  No Fourier cutoff
is applied to the unknown rate.
\end{lemma}
\begin{proof}
The inherited transverse formula \eqref{v1:shift-source} is a
constant-coefficient derivative--shift map with prefactor $h\varkappa_h$.
On the two fixed frequencies of $\beta^{12}$ it therefore has norm $O(h)$
and has the same support.  Its trace is zero, so applying $K_0^{-1}$
multiplies this output by $1/2$.

Apply the \emph{full} adjoint first source to this output.  The shift part
has another factor $h$ and bounded fixed-frequency derivative.  The lapse
part is the sum of the geometric phase-product defects and the electric
terms displayed in the proof of \cref{v1:first-source}.  Moving their
phase derivatives and literal shifts by their exact adjoints leaves a
fixed Fourier bank: the input two-coordinate frequency is shifted by at
most one of the six axial field frequencies.  All indices have absolute
value at most two.  At sufficiently fine odd cutoffs no folding occurs
in this calculation.

On this bank $\delta_i=\partial_i+O(h^2)$, with a uniform coefficient
bound, and the geometric product defects and their adjoints are $O(h^2)$.
The electric terms and their adjoints are $O(h)$ because they carry an
explicit factor $h$ and bounded first-field coefficients.  Acting on the
already $O(h)$ connection output gives $O(h^2)$ in the lapse slot as
well.  The finite bank has fixed dimension; its Euclidean and physical
spatial-mean Fourier norms have the stated uniform bounds.  This proves
\eqref{v9:row-bound} for the complete resulting vector, not only for its
projection onto some tests.  Large-frequency components of the unknown
rate remain in its pairing with this exact row.
\end{proof}

\begin{theorem}[An unavoidable active refinement scale]
\label{v9:active-rate}
On every sufficiently fine cutoff in the scalar-momentum balanced family,
any vector $r_h$ satisfying the complete leading consistency equation
\begin{equation}
 \mathcal M_{2,h}r_h=-\widehat B_{2,h}(Y_h)
 \label{v9:leading-equation}
\end{equation}
obeys
\begin{equation}
                         \boxed{\|r_h\|_h\ge c/h.}
 \label{v9:rate-lower}
\end{equation}
The constants are uniform in the phases and compact positive radius ranges.
This is a necessary estimate for every solution; neither solvability nor
signed regularity is assumed as a conclusion.

In particular, if at a fixed cutoff an analytic exact initial preparation
has a finite limit $r_h=\lim_{a\to0}\partial_t\lambda_h(a,0)$, this limit
satisfies \eqref{v9:leading-equation}.  Thus a hypothetical cofinal extension
of V8's bounded-amplitude-rate preparations cannot have regulator-uniform
leading anchored multiplier velocities.
\end{theorem}
\begin{proof}
Pair \eqref{v9:leading-equation} with the actual transverse test.  Since
$\mathcal M_2$ is self-adjoint in the positive physical pairing,
\[
 ch\le |\widehat B_{2,h}[(0,\beta^{12})]|
   =|\ip{\mathcal M_{2,h}(0,\beta^{12})}{r_h}_h|
   \le Ch^2\|r_h\|_h.
\]
This proves the lower bound independently of every other eigenvalue or
source direction.  For the second assertion use the exact differentiated
constraint equation $\mathsf B_h+M_h\lambda_t=0$.  On an exactly prepared
curve its expansion is
\[
 M_h=a^2\mathcal M_{2,h}+O_h(a^3),\qquad
 \mathsf B_h=a^2\widehat B_{2,h}+O_h(a^3).
\]
The preparation correction in the lapse row is already included in
$\widehat B_2$.  Divide by $a^2$ and use the finite amplitude limit of the
rate.  The mean-lapse normalization adds only the fixed null direction,
which the displayed equation annihilates.
\end{proof}

\begin{proposition}[A finite-amplitude version with explicit remainder payments]
\label{v9:finite-rate}
At an actually supplied exactly prepared cut suppose
\[
 \left|a^{-2}\mathsf B_h[(0,\beta^{12})]
                -\widehat B_{2,h}[(0,\beta^{12})]\right|\le c_0h/2,
 \qquad
 \|a^{-2}M_h(0,\beta^{12})
                -\mathcal M_{2,h}(0,\beta^{12})\|_h\le h^2,
\]
where $c_0h$ is the lower bound from \cref{v9:scalar-test}.  Every actual
normalized exact rate at that cut satisfies
$\|\lambda_t\|_h\ge c_0/[2(C+1)h]$.
\end{proposition}
\begin{proof}
Pair the exact equation $\mathsf B_h+M_h\lambda_t=0$ with the same test
and apply the two displayed inequalities and \cref{v9:test-row}.
\end{proof}

At a fixed cutoff the analytic remainders permit sufficiently small
amplitudes satisfying these payments.  Their joint $a,h$ domains have not
been made uniform.  In particular, \cref{v9:active-rate} is not a theorem
about fixed nonzero amplitude as $h\to0$, nor a proof of invariant-curvature
divergence.  It excludes one proposed route to a uniform \emph{anchored
metric-jet} comparison and identifies its actual action source.

\section{An upstream momentum extension that preserves weak neutrality}
\label{v9:momentum-family}

To test whether scalar momentum is the cause of this difficulty, enlarge
the first data to
\begin{equation}
 Y=(U,p),\qquad U=\sum_iU_i,\qquad
 p=p_0I+Q,\qquad Q=\sum_i\ell_iU_i,
 \qquad \ell_i\in\R.
 \label{v9:locked-family}
\end{equation}
The tensor momentum on each axis follows the same spatial profile.  This is
an explicit class, not a description of all transverse momenta.  All fields
remain trace-free and divergence-free in their nonconstant slots.

\begin{lemma}[Exact kinetic part of the weak drift]
\label{v9:kinetic-defect}
For any trace-free, phase-divergence-free symmetric momentum $Q$, the weak
quadratic kinetic drift is represented on curl-free tests by the vector
\begin{equation}
 \mathcal K_i(Q)=\sum_{j,k}\mathfrak d_i(Q_{jk},Q_{jk})
                -2\sum_{k,j}\mathfrak d_k(Q_{ij},Q_{kj}),
 \label{v9:kinetic-vector}
\end{equation}
where $\mathfrak d$ is the original cyclic phase-product defect.  This is a
representative of the tested covector; a transverse part is not being
identified with the full drift by this equation.
\end{lemma}
\begin{proof}
The full weak identity in the inherited source gives
\[
 d_{2,h}(U,p)[\beta]
 =-D_UH_3^{\rm g}(U,p)[w_\beta]
   +\ip{\mathcal F(U)}{\mathcal G_\beta U}_h
   +\ip p{\mathcal G_\beta V(p)}_h.
\]
The tested cubic logarithmic derivative is zero for a curl-free test,
without being removed from other coefficient calculations.  For trace-free
$Q$, $V(Q)=2Q$, and the kinetic derivative is the pairing of $w_\beta$ with
$2Q^2-\tfrac12|Q|_F^2I$.  Substitute the complete product-defect expression
for $\mathcal G_\beta$ in \cite{OB25},
\src{obv23:weak-generator}.  Symmetry of $Q$ and $\operatorname{div}_\delta Q=0$
give, using exact summation by parts,
\[
 \begin{split}
 \ip Q{\mathcal G_\beta Q}_h
 ={}&\sum_{i,k,j}\ip{\beta_i Q_{kj}}{\delta_k Q_{ij}}_h
       +3\ip{Q^2}{\delta_i\beta_j}_h\\
   &-\sum_{k,i,j}\ip{\beta_k Q_{ij}}{\delta_k Q_{ij}}_h
       -\ip{|Q|_F^2}{\delta_k\beta_k}_h.
 \end{split}
\]
Twice this expression, minus the kinetic $w_\beta$ pairing, equals
\[
 2\sum_{i,k,j}\ip{\beta_i Q_{kj}}{\delta_kQ_{ij}}_h
   +2\ip{Q^2}{\delta_i\beta_j}_h
   -2\sum_{k,i,j}\ip{\beta_kQ_{ij}}{\delta_kQ_{ij}}_h
   -\ip{|Q|_F^2}{\delta_k\beta_k}_h.
\]
Move the derivatives on $\beta$ and use the definition of
$\mathfrak d_k$.  The ordinary product contributions involving
$\sum_k\delta_kQ_{kj}$ are zero.  The remaining terms are exactly
$\ip\beta{\mathcal K(Q)}_h$.  No Leibniz identity for $\delta$ was used.
\end{proof}

\begin{theorem}[Weak quadratic neutrality on the momentum extension]
\label{v9:extended-neutrality}
Every first jet \eqref{v9:locked-family}, with arbitrary real $p_0,\ell_i$,
has zero complete curl-free quadratic shift drift at every odd cutoff.
\end{theorem}
\begin{proof}
The weak drift splits into a configuration quadratic and a momentum
quadratic, as in the preceding proof.  The configuration part is neutral
by \cref{v4:neutral}.  The cross term of a scalar constant momentum with
any other momentum is zero by \cref{v5:momentum-zero}.  The constant term
itself is zero.  It remains to evaluate \eqref{v9:kinetic-vector} on $Q$.

Different-axis profile products have zero $\mathfrak d_k$ exactly, since in
each differentiated coordinate at least one factor is constant.  For a
single-axis profile $Q_i=\ell_iU_i$, its $i$th row is zero.  In the second
sum of \eqref{v9:kinetic-vector}, differentiation along that axis therefore
has a zero row factor; all other derivatives are zero.  In the first sum,
$|Q_i|_F^2=2\ell_i^2r_i^2$ is constant pointwise, while
$\ip{Q_i(x)}{\delta_iQ_i(x)}_F=0$ pointwise by equal polarization norms
and orthogonality.  Hence
$\sum_{j,k}\mathfrak d_i((Q_i)_{jk},(Q_i)_{jk})=0$.
All other coordinates give zero.  These statements include folded
frequencies and prove the claim before taking any curl-free projection.
\end{proof}

\begin{proposition}[Compatibility with the original nonlinear initializer]
\label{v9:extended-preparation}
Set $m_i=2r_i^2$ and
\begin{equation}
 p_0^2=\frac23\sum_i m_i(\ell_i^2+\omega_h^2/4).
 \label{v9:p0-calibration}
\end{equation}
Then the first jets \eqref{v9:locked-family} satisfy the four quadratic mean
compatibility equations.  On compact positive-radius and bounded-$\ell$
parameter sets, either fixed nonzero sign of $p_0$ admits the inherited
analytic exact original initial preparation with this first jet, after
reducing the amplitude radius.  This is an initial-constraint statement;
it supplies no new signed-rank or exact-evolution theorem.
\end{proposition}
\begin{proof}
Let $V_i=\omega_h^{-1}\delta_iU_i$, and add balancing momenta
$\sum_i b_iV_i$ while treating $p_0$ as the fourth balancing coordinate.
The four exact quadratic means, in the sign convention for
$\mathscr C_h$ used in \cref{v4:mean}, are
\begin{equation}
 \left(\frac32p_0^2-\sum_i m_i(\ell_i^2+b_i^2+\omega_h^2/4),\quad
                         (\omega_hm_ib_i)_{i=1}^3\right).
 \label{v9:mean-map}
\end{equation}
Indeed different axes have zero mean cross pairings, $U_i\perp V_i$,
$\|U_i\|_h^2=\|V_i\|_h^2=m_i$, and $\delta_iU_i=\omega_hV_i$.
The scalar formula follows from the original quadratic Hamiltonian
\eqref{v5:H2}.  The multiplier means see no cubic logarithmic contribution
at this quadratic field order.  At $b=0$ and \eqref{v9:p0-calibration}, the
four balancing derivatives form
$\operatorname{diag}(3p_0,\omega_hm_1,\omega_hm_2,\omega_hm_3)$.
They are uniformly nonzero on the stated compact parameter sets.

The nonconstant linear constraint inverse, the analytic nonlinear range
solve, and the normalized four-mean contraction are the already proved
original initializer.  Replace only its leading four-mode calibration by
\eqref{v9:mean-map}.  Its bounded fixed-frequency coefficients and the
nonzero four-dimensional derivative meet the same hypotheses.  The range
correction starts at order $a^2$, and the balancing correction changes no
first coefficient.  This yields exact preparations $X_h(a)=aY_h+O_h(a^2)$
with every original initial row zero.  No future acceleration is input to
this construction.
\end{proof}

\section{Mixed-channel cancellation creates an unavoidable single-axis source}
\label{v9:upstream-boundary}

For each axis let $j(i)$ be the smaller of its two transverse coordinates
and define the transverse unit-frequency test
$\gamma^i=e_{j(i)}\cos\theta_i$.

\begin{theorem}[Complete tested drift on the momentum extension]
\label{v9:extended-tests}
For \eqref{v9:locked-family}, on every sufficiently fine odd cutoff,
\begin{align}
 \widehat B_2[(0,\beta^{12})]
   &=\frac{h\omega_h^2r_1r_2}{16}g_{12}(h)
                            (\omega_h^2+4\ell_1\ell_2),\notag\\
 \widehat B_2[(0,\beta^{13})]
   &=\frac{h\omega_h^2r_1r_3}{16}g_{13}(h)
                            (\omega_h^2-4\ell_1\ell_3),\notag\\
 \widehat B_2[(0,\beta^{23})]
   &=\frac{h\omega_h^2r_2r_3}{16}g_{23}(h)
                            (\omega_h^2+4\ell_2\ell_3),
                       \label{v9:extended-mixed}\\
 \widehat B_2[(0,\gamma^i)]
   &=\frac{h\omega_h^2p_0\ell_ir_i}{8}
                         (\cos\alpha_h+\sin\alpha_h).
                       \label{v9:single-axis}
\end{align}
The mixed formulas hold at odd $N\ge3$; the single-axis formula is stated
at $N\ge5$ so that same-axis doubled frequencies cannot alias into its
test.  With positive radii, calibrated $p_0$, and sufficiently small $h$,
the full prepared active quadratic drift cannot vanish anywhere in this
class.  Uniformly on compact positive-radius ranges and bounded $\ell$,
\begin{equation}
                         \boxed{\|\widehat B_{2,A,h}\|_h\ge c h.}
 \label{v9:extended-lower}
\end{equation}
\end{theorem}
\begin{proof}
The configuration contributions are \cref{v9:scalar-test}.  Evenness in
momentum excludes mixed $Up$ terms.  A mixed test can see only the product
of its two axial momentum profiles, and not $p_0$ or a single profile.
Using the same coefficient recipe, with $U=0$ and
$p=U_i+U_j$, gives the following table divided by $h\omega_h^2$:
\begin{equation}
\begin{array}{c|ccc}
 ij&\text{geometric}&\text{electric}&\text{magnetic}\\ \hline
 12&0&0&[2-\mathrm i(z+z^2-z^{-1}-z^{-2})]/8\\
 13&0&1/2&-(z+1+z^{-1})/4\\
 23&0&0&[2+\mathrm i(z+z^2-z^{-1}-z^{-2})]/8.
\end{array}
\label{v9:momentum-table}
\end{equation}
In this substitution $F_1(0,p)=(2p-(\tr p)I,0)$; it must not be replaced
by the scalar-momentum velocity.  Coefficients on a single unit test use
$U=0$, $p=p_0I+U_i$.  Their geometric and electric cosine contributions
are zero; the magnetic contribution is
\[
 \frac{h\omega_h^2p_0}{16}(1-\mathrm i)(z+\mathrm i z^{-1})
       =\frac{h\omega_h^2p_0}{8}(\cos\alpha_h+\sin\alpha_h).
\]
This identity holds for all three choices of $(i,j(i))$.  The original
ordered-pair convention, not an assumed rotational covariance, is used in
these three finite contractions.  The substitutions specifying
\eqref{v9:configuration-table} apply with the stated momenta and tests and
produce \eqref{v9:momentum-table} and this single-axis coefficient.  The
package supplies their separate symbolic extractions.  Field-degree and
frequency counting exclude all remaining momentum monomials.  Restoring
$r_i,\ell_i$ proves the formulas.

For all sufficiently small $h$ each $g_{ij}$ and
$\cos\alpha_h+\sin\alpha_h$ is bounded below by a fixed positive number.
If all single-axis values vanished, calibrated $p_0\ne0$ and $r_i>0$
would force every $\ell_i=0$.  The $12$ mixed value would then be nonzero.
This proves impossibility of complete active neutrality.

For a uniform lower bound, if some $|\ell_i|\ge\omega_h/4$, the
corresponding single-axis value has magnitude at least $c h$, since
$|p_0|$ has a positive uniform lower bound by
\eqref{v9:p0-calibration}.  Otherwise
$|\omega_h^2+4\ell_1\ell_2|\ge3\omega_h^2/4$, and the $12$ test has
that lower bound.  Each test has fixed $L_h^2$ norm $1/\sqrt2$.
Cauchy--Schwarz proves \eqref{v9:extended-lower}.  All tested directions
are active transverse shifts, so weak neutrality does not remove them.
\end{proof}

\begin{corollary}[The tempting three-channel repair is incomplete]
\label{v9:three-cancel}
At sufficiently fine cutoffs the three mixed values vanish simultaneously
exactly for
\[
 (\ell_1,\ell_2,\ell_3)
       =\pm(\omega_h/2,-\omega_h/2,\omega_h/2).
\]
For these choices the calibrated scalar momentum is
$p_0^2=2\omega_h^2\sum_i r_i^2/3$, and all three single-axis values are
nonzero of order $h$ on compact positive-radius sets.
\end{corollary}
\begin{proof}
The three nonzero prefactors in \eqref{v9:extended-mixed} reduce the
conditions to $\ell_1\ell_2=-\omega_h^2/4$,
$\ell_1\ell_3=\omega_h^2/4$ and
$\ell_2\ell_3=-\omega_h^2/4$.  None of the variables is zero.
Taking their ratios gives the two stated solutions.  Substitute them in
\eqref{v9:p0-calibration} and \eqref{v9:single-axis}.
\end{proof}

This is not a no-go theorem for arbitrary transverse momenta, additional
Fourier modes, or other original-action first data.  It rules out a specified
and exactly preparable extension that might otherwise be mistaken for a
complete source cure.  In particular it does not alter the action to cancel
one of these coefficients.

\section{The phase-selected three-site cut has a nonzero canonical acceleration}
\label{v9:initial-acceleration}

Return to the actual V8 family, with scalar momentum and its exact
amplitude-dependent phase and preparation controls.  Let
$\mathcal G_h\nu=\mathbb J\nabla P_{\nu,1}$, so by
\eqref{v5:linear-multiplier},
\begin{equation}
 \mathcal G_h(n,\beta)
   =(\delta_i\beta_j+\delta_j\beta_i,
                 \delta_i\delta_jn-\delta_{ij}\Delta_{\delta,h}n).
 \label{v9:gauge-map}
\end{equation}

\begin{theorem}[A fixed-cutoff second-jet obstruction after weak-rate cancellation]
\label{v9:canonical-acceleration}
At the verified phase-selected root of V8, every selected analytic
preparation in \cref{v8:exact-preparation} has a uniquely determined
amplitude-leading multiplier rate $r_h$, independent of its higher
preparation controls, with
\[
 r_h\ne0,\qquad \overline{r_h^N}=0,\qquad
 \mathbb P_{\mathrm{cf}}r_h^\beta=0.
\]
Along its own exact continuation,
\begin{equation}
 \boxed{\ddot X_h(a,0)=\mathcal G_hr_h+O_h(a),\qquad
                         \mathcal G_hr_h\ne0.}
 \label{v9:nonzero-acceleration}
\end{equation}
Thus $X_h(a,0)=O_h(a)$ and $\dot X_h(a,0)=O_h(a)$ do not imply
$\ddot X_h(a,0)=O_h(a)$ in the original canonical coordinates.
\end{theorem}
\begin{proof}
The graded analytic solution in V8 gives
$\lambda_t(a,0)=r_h+O_h(a)$.  The weak observation and mean-lapse
normalization place $r_h$ on the physical active slice.  Dividing the
original consistency equation by $a^2$ gives
\eqref{v9:leading-equation}.  On that slice the active signed Gram is
invertible at the V8 root.  Its right side and first source depend only on
$Y_*$, so $r_h$ is unique and independent of higher preparation.

At $N=3$ the $12$ test in \cref{v9:scalar-test} has
$g_{12}=1$ and, for unit radii, value
$h\omega_h^4/16=243/16\ne0$ in spatial-mean normalization.
The leading rate therefore cannot be zero at any phase, including the
certified root.  This nonvanishing uses an exact action coefficient, not a
rounded value of the phase root.

Let $\mathsf F_h=\mathbb J\nabla_X\mathscr H_h$.  At the cut,
$\mathsf F_h(aY_*+O_h(a^2),e)=O_h(a)$ and its $X$ derivative is bounded.
Its multiplier derivative at $X=0,e$ is exactly $\mathcal G_h$.
Differentiation of $X_t=\mathsf F_h(X,\lambda)$ consequently gives
\eqref{v9:nonzero-acceleration}.  For mean-zero lapse and mean-zero
transverse shift, Parseval gives
\begin{equation}
 \|\mathcal G_h(n,\beta)\|_h^2
        =2\|\Omega_h\beta\|_h^2+2\|\Omega_h^2n\|_h^2.
 \label{v9:gauge-injective}
\end{equation}
Indeed the two slots in \eqref{v9:gauge-map} have Fourier squared norms
$2|\kappa|^2|\widehat\beta|^2$ and $2|\kappa|^4|\widehat n|^2$.
Every nonzero odd-grid frequency has $|\kappa|>0$.  Thus this map is
injective on the active slice and $\mathcal G_hr_h\ne0$.
\end{proof}

The leading acceleration is in the \emph{linear constraint-generated}
directions.  The theorem is not a curvature singularity theorem, and does
not contradict the $O_h(a)$ initial literal curvature in V8.  In particular,
passing to a coordinate-independent or appropriately transported comparison
requires another argument, not a reinterpretation of
\eqref{v9:nonzero-acceleration} as physical curvature.

\section{Exact affine canonical transport and the entire initial second jet}
\label{v9:affine}

The preceding acceleration suggests a larger comparison than V6's uniform
spatial translation.  This comparison is an exact change of variables, not
a symmetry assertion.  The linear original rows satisfy
\[
 \{P_{\nu,1},P_{\mu,1}\}=0.
\]
For completeness, their mixed bracket is the scalar linear curvature of a
symmetric phase gradient, equivalently the divergence of
$\delta_i\delta_jn-\delta_{ij}\Delta_\delta n$; both are zero by commuting
phase derivatives.  Pure lapse and pure shift brackets are zero directly.

\begin{proposition}[An exact affine moving-frame Hamiltonian]
\label{v9:affine-action}
For a prescribed twice differentiable multiplier-valued path $q(t)$, set
$X=\widetilde X+\mathcal G_hq(t)$.  Up to an exact time derivative and a
function of time alone, the original canonical action has transformed
Hamiltonian
\begin{equation}
 \boxed{\widetilde H_h(t,\widetilde X,\lambda)
   =\mathscr H_h(\widetilde X+\mathcal G_hq,\lambda)
                       -P_{\dot q,1}(\widetilde X).}
 \label{v9:affine-H}
\end{equation}
Every multiplier equation is the original equation evaluated at the shifted
canonical record.  Every literal observation is retained by the same
pullback.  The transformation is considered only where that record belongs
to the original stationary chart.
\end{proposition}
\begin{proof}
Write $\mathcal G_hq=(u_q,p_q)$.  Expanding the canonical one-form gives
\[
 \ip{\widetilde p+p_q}{\dot{\widetilde u}+\dot u_q}
 =\ip{\widetilde p}{\dot{\widetilde u}}
    +\ip{\widetilde p}{\dot u_q}-\ip{\dot p_q}{\widetilde u}
    +\frac{d}{dt}\ip{p_q}{\widetilde u}+\ip{p_q}{\dot u_q}.
\]
The two middle field-dependent terms are $P_{\dot q,1}(\widetilde X)$
because its Hamiltonian vector is $(\dot u_q,\dot p_q)$.
The remaining extra scalar depends only on the prescribed path.
This proves \eqref{v9:affine-H}.  Affine translations are symplectic;
substitution or direct differentiation gives
$\dot{\widetilde X}=\mathsf F_h(\widetilde X+\mathcal G_hq,\lambda)
-\mathcal G_h\dot q$.  During variation $q$ is fixed, so the subtraction
has no multiplier derivative.  Thus no original constraint equation or
observable is dropped.
\end{proof}

\begin{theorem}[Order-amplitude entire initial canonical acceleration in the affine frame]
\label{v9:affine-acceleration}
For an existing V8 exact history choose $q(0)=0$ and
$\dot q(t)=\lambda(t)-e$ as a prescribed comparison path for that history.
Then at its initial cut
\begin{equation}
 \widetilde X(a,0)=O_h(a),\qquad
 \dot{\widetilde X}(a,0)=O_h(a),\qquad
 \boxed{\ddot{\widetilde X}(a,0)=O_h(a).}
 \label{v9:affine-small}
\end{equation}
This bound concerns the whole canonical acceleration, not merely the
constant-shift contribution isolated in V6.
\end{theorem}
\begin{proof}
Initially $\lambda=e$, so $\dot q(0)=0$ and $\ddot q(0)=\lambda_t(0)$.
The exact identity is
\[
 \ddot{\widetilde X}(0)
   =D_X\mathsf F_h\,\mathsf F_h
       +(D_\lambda\mathsf F_h-\mathcal G_h)\lambda_t
                 \quad\text{at }(X(a),e).
\]
The first factor pair is $O_h(a)$ by analyticity and
$\mathsf F_h(X(a),e)=O_h(a)$.  The bracketed map is $O_h(a)$ because
$D_\lambda\mathsf F_h(0,e)=\mathcal G_h$.  V8 bounds the actual
$\lambda_t$ at the fixed cutoff.  This proves the acceleration statement;
the two lower jets follow directly from $q(0)=\dot q(0)=0$.
\end{proof}

Choosing the path from a particular existing history is legitimate for
comparison.  Substituting $\dot q=\lambda-e$ \emph{as a varied multiplier}
in \eqref{v9:affine-H} would subtract its linear constraint row and change
the theory.  Likewise a metric or link read in this frame is
$\mathcal R_h(\widetilde X+\mathcal G_hq,\lambda)$, not generally
$\mathcal R_h(\widetilde X,\lambda)$.  The affine transformation is not
asserted to be a site permutation, a positive stochastic map, or a nonlinear
ADM coordinate transformation.  Its third derivative retains
$\lambda_{tt}$ and differentiated coefficients.  No physical lifespan,
propagated rate condition, or cofinal observable bound follows from
\eqref{v9:affine-small} alone.

\section{Status, research direction, and verification}
\label{v9:status}

The main action-specific finding is the transverse order-$h$ quadratic
source with its order-$h^2$ smooth Hessian row.  It is present at every phase
and is different from the weak effective source eliminated in V8.  Therefore
establishing cofinal roots of the latter alone would not give the uniform
anchored leading-rate estimate.  The enlarged exactly preparable,
weak-neutral momentum class exposes the next upstream constraint: a
plausible cancellation of the mixed channels generates nonzero single-axis
channels.  Within this class the complete active quadratic source cannot be
zero.  This is a restriction on first data, not a redesign of the action.

There is also a precise positive comparison result.  The complete initial
canonical second jet of the verified phase-selected history is small after
an exact affine symplectic transport by the original linear constraints.
A useful continuation must now either find a more general first-jet family
that suppresses the \emph{full} active source at its natural Hessian scale,
or control this transported dynamics together with its original literal
observations.  More refinements of the same three phase equations do not
resolve that additional incidence.

No new all-cutoff full-rank or phase-root family is asserted.  No exact
nonlinear spacetime history is numerically integrated.  The scalar-family
rate lower bound keeps the finite-amplitude remainder payments and the
order of limits explicit.  It is not promoted to a fixed-amplitude
continuum failure theorem or to curvature divergence.  Higher time jets,
the true nonlinear Hessian and source, the physical reconstruction and a
common exact-action interval remain unresolved.

\paragraph{Verification.}
The supplied code regenerates the displayed electric, magnetic and geometric
contractions in exact symbolic Fourier arithmetic from the original scalar
coefficient functional.  It checks the closed Laurent identities separately
from numerical trigonometric substitution.  Independent floating AD tests
use the full cyclic coefficient program at several odd cutoffs, phases,
radii and tensor-momentum values.  They check all six selected channels,
the proposed mixed-channel cancellation and the surviving single-axis
source.  Additional tests verify the exact kinetic product-defect formula,
weak neutrality, the mean calibration, the full smooth Hessian-row scale
on finite grids, and the affine symplectic identity.  These diagnostics are
not proofs of uniform inverse or higher-jet estimates.  The earlier V8
phase-root theorem is an inherited input to the physical initial-cut
application, not a new interval computation in this body.

The first eight marked research bodies are unchanged byte for byte.  The
preservation manifest records their SHA-256 values, the new body, and the
coefficient programs.  The source audit distinguishes the new quadratic
active coefficient from the already proved cubic momentum, weak drift,
and constant-source invariants.
% END IMMUTABLE EXACT-ACTION BRIDGE V9 BODY


% Future iterations append here.
% BEGIN IMMUTABLE EXACT-ACTION BRIDGE V10 BODY
\clearpage
\renewcommand{\thesection}{\arabic{section}}
\renewcommand{\theHsection}{bridgeV10.\arabic{section}}
\setcounter{section}{75}
\section{The geometric observation after canonical transport}
\label{v10:context}

The affine transport of \cref{v9:affine-acceleration} makes the complete
initial canonical acceleration small at a fixed cutoff. It does not identify
the original metric or its curvature with an untransported function of those
new canonical variables. Repeating the canonical cancellation without
calculating this observation would leave the bridge at the same point.
We therefore return to the ordinary interpolated ADM metric already used in
\cref{v6:geometric}, and evaluate its full Riemann tensor at the actual
phase-selected initial cut. No new action, evolution, interpolation, spatial
scale, or literal-link convention is introduced.

The new conclusion is stronger than a coordinate-rate diagnostic. The
initial metric curvature has an explicitly computable amplitude limit,
although the original literal link curvature tends to zero. On the certified
three-site phase root this limit is nonzero, as proved below using exact
rational postprocessing of the inherited source enclosures. Its spatial
Einstein tensor and its Riemann-squared scalar are also nonzero. A genuine
spacetime coordinate change can remove all initial lapse and shift
velocities, but the same curvature then remains in the spatial metric's
second time derivative. It cannot be removed by a coordinate re-expression
of the same initial slice.

There is also a positive quantitative reduction. At every odd cutoff the
leading curvature is equivalent, in a positive norm, to the transverse
shift rate with the precise weight $h^2$ and three spatial derivatives.
Thus uniform boundedness of all anchored multiplier rates is not the
necessary geometric target. The appropriate initial-curvature target is a
source-dependent weighted rate estimate, together with the actual
finite-amplitude remainder. These estimates have not yet been obtained
along a cofinal exact-action family.

\paragraph{Inherited inputs and scope.}
We use the original stationary chart and exact initial preparations, the
verified phase root and its coefficient enclosures in
\cref{v8:root-theorem,v8:exact-preparation}, the actual active-rate limit and
canonical acceleration in \cref{v9:canonical-acceleration}, and the
$O_h(a)$ literal-curvature readouts of \cref{v8:physical}. Those statements
are inherited, not independently re-proved by the new arithmetic. In
particular, the large root-enclosure calculation is not repeated.
The $N=3$ calculations below are new rigorous consequences of its stored
interval matrices. General hypersurface and Riemann identities are standard
\cite{v10:G}; their particular specialization to the original-action jets,
the sharp frequency comparison, and the nonzero certified geometric
observation are proved here. The earlier mixed rotational link defect in
\cite{OB25}, \src{obv24:identification-defect}, concerns a different initial
family and the opposite ordering of small and non-small reads.

\section{The metric, derivative conventions, and actual initial jets}
\label{v10:metric}

Use the odd spatial grid $N=2m+1$, $h=N^{-1}$, and the trigonometric
interpolation $\mathcal I_h$ on the unit three-torus. On arrays let
$D_{i,h}$ denote differentiation of this interpolant, not the original
phase derivative:
\begin{equation}
 \widehat{D_{i,h} f}(k)=2\pi\mathrm i k_i\widehat f(k),\qquad
 \widehat{\delta_{i,h}f}(k)=\mathrm i\kappa_{i,h}(k)\widehat f(k),\qquad
 \kappa_{i,h}(k)=2h^{-1}\sin(\pi h k_i).
 \label{v10:two-derivatives}
\end{equation}
The first derivative is the $\partial^{\rm sp}$ of \cref{v6:geometric}.
Define $\mathcal Q_{i,h}=D_{i,h}-\delta_{i,h}$. These are analytical
comparison operators; none replaces $\delta_{i,h}$ in the finite action.

For an existing original-action history, reconstruct the same metric as in
\cref{v6:geometric}:
\begin{equation}
 \gamma=I+\mathcal I_hu,\quad N=\mathcal I_h\lambda^N,\quad
 \beta=\mathcal I_h\lambda^\beta,\qquad
 g=-N^2\,dt^2+\gamma_{ij}(dx^i+\beta^i dt)(dx^j+\beta^jdt).
 \label{v10:ADM}
\end{equation}
On every fixed sufficiently small initial chart this is a smooth Lorentzian
metric near its initial slice. Its interval is only the existing interval
of the history. Write $a$ for preparation amplitude and $t$ for physical
time. At the selected cut $\lambda(a,0)=e=(1,0)$ exactly.

The geometric calculation applies to any family with the following actual
jets at a fixed cutoff:
\begin{equation}
 \begin{gathered}
 X_h(a,0)=aY_h+O_h(a^2),\qquad X_{h,t}(a,0)=O_h(a),\\
 \lambda_{h,t}(a,0)=r_h+O_h(a),\qquad
 X_{h,tt}(a,0)=\mathcal G_hr_h+O_h(a).
 \label{v10:actual-jets}
 \end{gathered}
\end{equation}
Here $r_h=(r_h^N,r_h^\beta)$ and $\mathcal G_h$ is the unchanged map
\eqref{v9:gauge-map}. All these estimates include any fixed spatial
Sobolev order at a fixed grid, by equivalence of finite-dimensional norms.
They hold on the actual V8 preparations by
\cref{v8:exact-preparation,v9:canonical-acceleration}. No cofinal family
with these bounds uniform in $h$ is being assumed to exist.

Use the Riemann sign convention for which, at a flat metric,
\[
 R^{[1]}_{0i0j}
 =\tfrac12(\partial_i\partial_t g_{0j}
             +\partial_j\partial_t g_{0i}
             -\partial_i\partial_j g_{00}-\partial_t^2g_{ij}).
\]
The Ricci contraction is $R_{bd}=g^{ac}R_{abcd}$. At the initial slice
$N=1,\beta=0$, so the coordinate time vector is the unit normal. A spatial
orthonormalization of $\gamma$ changes the limiting formulas only by an
$O_h(a)$ term.

\section{The complete leading initial metric curvature}
\label{v10:curvature}

Define the real symmetric tensor
\begin{equation}
 \boxed{\mathfrak E_h(r^\beta)_{ij}
   =\tfrac12\bigl(\mathcal Q_{i,h}r^\beta_j
                         +\mathcal Q_{j,h}r^\beta_i\bigr).}
 \label{v10:E}
\end{equation}
It is different from the literal temporal-link electric read and from the
source-space operators previously denoted by electric or parity symbols.

\begin{theorem}[Actual initial curvature left by the finite derivative]
\label{v10:metric-curvature}
For the metric \eqref{v10:ADM} and the actual jets
\eqref{v10:actual-jets}, at each fixed cutoff,
\begin{equation}
 \boxed{R_{0i0j}(g_h)(a,0)=\mathcal I_h\mathfrak E_h(r_h^\beta)_{ij}
                              +O_h(a),\quad
 R_{0ijk}(g_h)(a,0)=O_h(a),\quad R_{ijkl}(g_h)(a,0)=O_h(a).}
 \label{v10:R-limit}
\end{equation}
The initial spatial Einstein tensor satisfies
\begin{equation}
 G_{ij}(g_h)(a,0)
  =-\mathcal I_h\mathfrak E_{h,ij}
       +\delta_{ij}\mathcal I_h\tr\mathfrak E_h+O_h(a),\qquad
 G_{00}(a,0)=G_{0i}(a,0)=O_h(a).
 \label{v10:G-limit}
\end{equation}
In particular no $\lambda_{tt}$ estimate is needed for these initial
curvatures. The leading lapse velocity $r_h^N$ cancels from the limit;
it is not asserted absent from the remainder constants.
\end{theorem}
\begin{proof}
At the slice, the exact ADM identities give
\[
 g_{00}=-1,\quad g_{0i}=0,\quad g_{ij}=\gamma_{ij},\quad
 \partial_t g_{00}=-2\partial_t N,\quad
 \partial_tg_{0i}=\gamma_{ij}\partial_t\beta^j.
\]
All spatial derivatives of $g_{00}$ and $g_{0i}$ vanish there.
The spatial derivatives of $\gamma-I$ and its first time derivative are
$O_h(a)$. By \eqref{v10:actual-jets},
\[
 \begin{split}
 \partial_t^2\gamma_{ij}(a,0)
 &=\mathcal I_h(\delta_i r^\beta_j+\delta_jr^\beta_i)+O_h(a),\\
 \partial_i\partial_tg_{0j}(a,0)
 &=\mathcal I_hD_i r^\beta_j+O_h(a).
 \end{split}
\]
Substitution in the second-derivative part of $R_{0i0j}$ gives
\eqref{v10:E}. The remaining Christoffel products are $O_h(a)$:
in a limiting jet with $\gamma=I$, $\gamma_t=0$, and vanishing spatial
first derivatives, the only possibly nonzero Christoffel symbols are
$\Gamma^0_{00}=r^N$ and $\Gamma^i_{00}=r^\beta_i$.
Every product in $R_{0i0j}$ pairs such a symbol with a symbol having
spatial lower indices, which is $O_h(a)$. Products involving two
uncontrolled pure-time second derivatives do not occur. This also proves
why a possibly large $\lambda_{tt}$ does not enter the calculation.

The intrinsic spatial curvature is $O_h(a)$ and the second fundamental
form is $O_h(a)$. The Gauss and Codazzi formulas therefore give the last
two estimates of \eqref{v10:R-limit}. Alternatively their coordinate
second-derivative formulas contain only spatial derivatives and at most
one time derivative of the small spatial data. There is no multiplier
velocity term without such a small factor.

At the limiting jet the only nonzero Riemann block is
$R_{0i0j}=\mathfrak E_{ij}$. Contract with $\eta^{00}=-1$ and
$\eta^{ij}=\delta^{ij}$ to obtain
\[
 R_{00}=\tr\mathfrak E,\quad R_{ij}=-\mathfrak E_{ij},\quad
 R_{0i}=0,\quad R=-2\tr\mathfrak E.
\]
These contractions prove \eqref{v10:G-limit}. The change from $g^{-1}$ to
$\eta^{-1}$, and from coordinate to spatial orthonormal components, costs
$O_h(a)$ at fixed cutoff.
\end{proof}

\begin{corollary}[A positive scalar test, not only a large coordinate component]
\label{v10:scalar-invariant}
Let $\mathscr K(g)=R_{abcd}(g)R^{abcd}(g)$ be the Riemann-squared scalar.
Then in the setting of \cref{v10:metric-curvature},
\begin{equation}
 \mathscr K(g_h)(a,0,x)
   =4|\mathcal I_h\mathfrak E_h(x)|_F^2+O_h(a),\qquad
 \int_{\T^3}\mathscr K(g_h)(a,0,x)\,dx
   =4\|\mathfrak E_h\|_h^2+O_h(a).
 \label{v10:Kretschmann}
\end{equation}
Here the integral is the fixed unit-torus spatial normalization. Including
$\sqrt{\det\gamma}$ gives the same limit. If the original temporal-link
curvature is $O_h(a)$ in its original frame, its boost component cannot
converge in this fixed-cutoff amplitude limit to the metric block
\eqref{v10:R-limit} unless $\mathfrak E_h=0$.
\end{corollary}
\begin{proof}
Each nonzero limiting Riemann entry has two temporal indices, so raising
all four indices has positive sign. The curvature symmetries give four
copies of each $\mathfrak E_{ij}^2$, proving the pointwise identity. There
are no other order-one blocks by \eqref{v10:R-limit}.
Normalized discrete Parseval equals the integral of the squared
trigonometric tensor. Also $\sqrt{\det\gamma}=1+O_h(a)$.
The original frame conversion tends to its fixed flat value and is bounded;
an $O_h(a)$ literal link vector still tends to zero after that conversion.
The asserted identification requirement follows without changing the
connection or replacing its signed response.
\end{proof}

These statements concern the specified interpolated metric, not a theorem
that every conceivable continuum reconstruction must have the same
finite-cutoff curvature. Rescaling component formulas without the tensorial
pullback of the metric would change that reconstruction. A genuine coordinate
change of the same metric does not change \eqref{v10:Kretschmann} at the same
event. The fixed spatial and temporal calibrations used here are precisely
the ones retained in \cref{v6:geometric}.

\section{A sharp full-frequency criterion for this geometric observation}
\label{v10:frequency}

The comparison does not require the full multiplier rate to remain bounded.
It requires the output of the particular operator \eqref{v10:E} to be small.
This operator can be characterized without a frequency cutoff.
Write
\[
 d_h(k)=2\pi k-\kappa_h(k),\qquad
 \|b\|_{\dot H_h^{s+3}}^2
 :=\sum_{k\ne0}\langle k\rangle^{2s}|k|^6|\widehat b(k)|^2.
\]
The dot here removes the constant mode and uses the displayed norm; it is
uniformly equivalent to $\|b-\overline b\|_{H_h^{s+3}}$.

\begin{theorem}[Exact Fourier norm and uniform two-sided derivative weight]
\label{v10:symbol}
For every real vector array $b$ and every real $s$,
\begin{equation}
 \boxed{\|\mathfrak E_h(b)\|_{H_h^s}^2
 =\frac12\sum_k\langle k\rangle^{2s}
       \left(|d_h(k)|^2|\widehat b(k)|^2
                           +|d_h(k)\cdot\widehat b(k)|^2\right).}
 \label{v10:exact-symbol}
\end{equation}
There are positive constants $c,C$, independent of the odd cutoff and of
$b$, such that
\begin{equation}
 \boxed{c h^2\|b\|_{\dot H_h^{s+3}}
      \le\|\mathfrak E_h(b)\|_{H_h^s}
      \le C h^2\|b\|_{\dot H_h^{s+3}}.}
 \label{v10:sharp-symbol}
\end{equation}
Thus the kernel consists exactly of spatially constant vectors.
No phase-transversality assumption is needed for either bound.
\end{theorem}
\begin{proof}
At frequency $k$ the symmetric matrix is
$\mathrm i(d_i\widehat b_j+d_j\widehat b_i)/2$.
Squaring its Frobenius norm gives
$(|d|^2|\widehat b|^2+|d\cdot\widehat b|^2)/2$, proving
\eqref{v10:exact-symbol} by Parseval.
For $0\le x\le\pi/2$, the alternating sine expansion with its remainder
sign gives
\[
 \left(\frac16-\frac{\pi^2}{480}\right)x^3
       \le x-\sin x\le\frac{x^3}{6}.
\]
The lower coefficient is positive. Applying these inequalities to
$x=\pi h|k_i|$, which is less than $\pi/2$ on every retained mode, gives
\[
 2\pi^3\left(\frac16-\frac{\pi^2}{480}\right)h^2|k_i|^3
 \le |d_{h,i}(k)|\le\frac{\pi^3}{3}h^2|k_i|^3.
\]
Use $|k|^6/9\le\sum_i|k_i|^6\le|k|^6$, followed by
$|d|^2|\widehat b|^2/2\le|\widehat{\mathfrak E}|_F^2
\le|d|^2|\widehat b|^2$. This proves both bounds on the entire frequency
box, including its boundary, with constants independent of $s$ in these
particular norm conventions. For $k\ne0$, $d_h(k)\ne0$; the lower bound
then gives the kernel statement.
\end{proof}

\begin{corollary}[The weighted cofinal rate test for the actual initial metric]
\label{v10:cofinal}
On an actually supplied cofinal family satisfying
\eqref{v10:actual-jets}, its leading initial metric curvature tends to zero
in $H^s$ if and only if
\begin{equation}
             h^2\|r_h^\beta\|_{\dot H_h^{s+3}}\longrightarrow0.
 \label{v10:weighted-rate}
\end{equation}
For a joint path $a=a(h)$, assume in addition that the full comparison
remainder in \eqref{v10:R-limit} has norm at most $C_h|a(h)|\to0$.
Then \eqref{v10:weighted-rate} is necessary and sufficient for all initial
metric Riemann components to tend to zero. If the literal-link curvature
also tends to zero on that path, it is likewise necessary and sufficient
for this initial metric/link-curvature discrepancy to vanish.
\end{corollary}
\begin{proof}
Apply \cref{v10:symbol} to \eqref{v10:R-limit}, followed by the triangle
and reverse triangle inequalities. The other metric components are covered
by the stated full remainder estimate. If the link tends to zero,
subtracting it does not change vanishing of the metric tensor.
\end{proof}

A bound $\|r_h^\beta\|_{H_h^{s+3}}=O(h^{-q})$ with $q<2$ would suffice
for this leading observation. In particular an $O(h^{-1})$ upper bound
with the displayed spatial regularity would be compatible with convergence.
The lower bound on the \emph{whole} rate in \cref{v9:active-rate} supplies
neither this upper bound nor a lower bound on the shift part alone. Lapse
rates can still enter $C_h$ and higher physical derivatives.
No bound in \eqref{v10:weighted-rate}, no cofinal phase-root existence, and
no common time interval is inferred from this corollary.

\section{An original-action curvature certificate at the verified phase root}
\label{v10:certificate}

Return to the exact V8 root at $N=3$, in the isolating cube
\eqref{v8:boxes}. Let $C_A$ be its original $162\times78$ selected
first-source matrix, $A=C_A^TK^{-1}C_A$, and $b_A$ its correctly ordered
prepared drift. Put $x=A^{-1}b_A$. The entries of $x$ are coordinates in
the selected source basis, not the physical transverse shift rate.

Let $S:\R^{78}\to\R^{108}$ insert these coordinates into the original
component-major multiplier columns \eqref{v3:pivots}; all its nonzero
entries equal one. On the $27$-site scalar grid let $D_i=3(S_i-S_i^{-1})$,
$P_0=\mathbf1\mathbf1^T/27$, and set
\begin{equation}
 \begin{aligned}
 \Omega^2&=-\sum_iD_i^2,&
 L&=(\Omega^2+P_0)^{-1}-P_0,\\
 \nabla&=(D_1,D_2,D_3)^T,&
 P_T&=I_{81}-\diag(P_0,P_0,P_0)-\nabla L\nabla^T.
 \end{aligned}
 \label{v10:rational-Hodge}
\end{equation}
These matrices are rational. The scalar inverse acts only on mean-zero
arrays. $P_T$ is the orthogonal projection onto the complete mean-zero
phase-transverse shift space.

\begin{lemma}[The physical leading rate from the actual active solve]
\label{v10:rate-reconstruction}
At the V8 root the actual amplitude-leading rate satisfies
\begin{equation}
 r_h^N=-(I-P_0)(Sx)^N,\qquad
 \boxed{r_h^\beta=-P_T(Sx)^\beta.}
 \label{v10:r-from-x}
\end{equation}
At this cutoff $D_{i,h}=(2\pi/(3\sqrt3))\delta_{i,h}$ on \emph{every}
array. Thus
\begin{equation}
 \boxed{\mathfrak E_{h,ij}
   =\frac12\left(\frac{2\pi}{3\sqrt3}-1\right)
                 (D_ir_h^\beta{}_j+D_jr_h^\beta{}_i),\qquad
                  \tr\mathfrak E_h=0.}
 \label{v10:three-E}
\end{equation}
\end{lemma}
\begin{proof}
The source annihilates constant lapse and curl-free shift arrays. Its
selected columns span the quotient by that kernel. Since
$Ax=b_A$, the representative $-Sx$ solves the leading multiplier equation
on every selected test; the original drift annihilates the same kernel,
so it solves the complete equation. Projecting off that kernel preserves
it. The V8 preparations have exactly zero curl-free shift velocity and
normalized mean lapse, so their limit is precisely this projected
representative. Uniqueness on the regular physical active slice proves
\eqref{v10:r-from-x}. This accounts for the weak correction needed when
the selected active representatives themselves are not orthogonal to the
weak subspace.

The inverse in \eqref{v10:rational-Hodge} is the usual exact finite Hodge
inverse because $\nabla^T\nabla=\Omega^2$. It gives
$P_T^2=P_T=P_T^T$ and $\nabla^TP_T=0$.
At $N=3$ each frequency component is $0,1$, or $-1$, so the two derivative
symbols have the same constant ratio. Substitution in \eqref{v10:E}
gives \eqref{v10:three-E}, and its trace vanishes by phase-transversality.
\end{proof}

\begin{theorem}[Nonzero geometric curvature at the original-action phase root]
\label{v10:verified-E}
At the exact V8 root, with its original physical normalization,
\begin{equation}
 \boxed{
  76.8<\mathfrak E_{23}(2/3,1/3,0)<76.9,\qquad
  78.9<\|\mathfrak E_h\|_h<78.92.}
 \label{v10:verified-values}
\end{equation}
Consequently along each of its selected exact initial preparations,
\begin{equation}
 \begin{split}
 \lim_{a\to0}\|G_{ij}(g_h)(a,0)\|_{L^2}&=\|\mathfrak E_h\|_h>78.9,\\
 24900<\lim_{a\to0}\int_{\T^3}\mathscr K(g_h)(a,0,x)\,dx&<24910.
 \end{split}
 \label{v10:verified-invariants}
\end{equation}
The original literal temporal and spatial link curvatures nevertheless
remain $O_h(a)$ on these same histories. The metric/link identification
therefore fails in this fixed-cutoff amplitude limit. These bounds are not
claims about a cutoff limit or a curvature in newly calibrated physical
units.
\end{theorem}
\begin{proof}
We specify the finite certificate and separate inherited inputs from new
arithmetic. The inherited V8 files enclose $A(\vartheta)$ and the complete
prepared $b(\vartheta)$ for every $\vartheta$ in the exact cube
\eqref{v8:boxes}. The factor-$27$ conversion is not made inside the solve:
its source Gram and target both have the same spatial-sum convention, so
$x$ and the physical rate are unchanged by conversion to spatial mean.
The new verifier restricts $b$ to the actual component-major pivot list,
not to site-major shift indices.

A fixed dyadic approximate inverse $R$ and a fixed dyadic vector $x_0$
are only proposals. Exact integer arithmetic on a binary grid of denominator
$2^{160}$ gives the verified residual
\[
 \sup\|I-RA\|_\infty<4.520\,10^{-6}<1.
\]
For every source matrix and target in the input enclosures,
\[
 x-x_0=R(b_A-Ax_0)+(I-RA)(x-x_0).
\]
The geometric-series bound followed by four componentwise enclosure
refinements bounds the maximum coordinate error by $0.002379$.
The complete componentwise radii, rather than that single maximum, are
used in the following linear image. All new additions, products,
roundings to the binary grid, and inequalities use arbitrary-precision
integers or rational numbers. No floating residual is accepted as a proof.

Construct \eqref{v10:rational-Hodge} over $\mathbb Q$, and verify its
idempotence, symmetry, and transverse identity exactly. Form all $243$
entries of the symmetric phase gradient in \eqref{v10:three-E} as rational
linear forms of $x$. Their exact interval images give, in particular,
\[
 734.625<(D_2r^\beta_3+D_3r^\beta_2)(2/3,1/3,0)<734.657.
\]
For the remaining scalar factor, square rational lower and upper bounds
on $\sqrt3$ to certify them. Bound $\pi$ by the alternating rational
series for $16\arctan(1/5)-4\arctan(1/239)$ with its explicit next-term
remainders. Multiplying the resulting positive scalar interval by each
linear-form interval proves the first inequality of
\eqref{v10:verified-values}. Summing the outward rational bounds on all
$243$ squared entries, divided by $27$, proves the norm inequalities.
The code checks the stronger squared inequalities directly; taking square
roots in a floating library is not part of the proof.

The same rational sum certifies
$24900<4\|\mathfrak E_h\|_h^2<24910$.
Its trace is exactly zero by \cref{v10:rate-reconstruction}, not merely
small numerically. Equations \eqref{v10:G-limit} and
\eqref{v10:Kretschmann} therefore prove \eqref{v10:verified-invariants}.
The literal link bounds are the inherited \cref{v8:physical}. The exact
root lies in the certified cube; thus every stated enclosure applies at
that root, without replacing it by its rounded phases.
\end{proof}

\paragraph{Reproducibility of the inclusion input.}
The package retains the inherited point and box source enclosures, phase
preconditioner, and original pivot list with their hashes. The V8 root
self-mapping and contraction inequalities are rechecked here in exact
rational arithmetic; the bounds are $0.550511$ for the contraction and
$2.49937\,10^{-11}<2^{-35}$ for the image radius. The V8 full-source and
preparation regularity theorem remains an inherited input. Rechecking an
inclusion does not independently regenerate the large original coefficient
enclosures. Their original regeneration code is supplied in the included,
unchanged V8 proof archive. This distinction prevents a cached array or a
hash from being mistaken for a new proof of its source incidence.

\section{A genuine spacetime normal form and the curvature that it retains}
\label{v10:diffeomorphism}

The preceding tensor is not a lapse--shift coordinate artifact. One can
remove all initial multiplier velocities by an actual coordinate change,
but the spatial acceleration then carries exactly the same curvature.
This gives a geometric counterpart to, rather than a repetition of, the
affine canonical transport.

\begin{theorem}[Exact initial geometric normalization]
\label{v10:geometric-normal}
For any existing smooth metric \eqref{v10:ADM} with $N(0)=1$, $\beta(0)=0$,
let
\[
 c_N(y)=\partial_tN(0,y),\qquad c^i(y)=\partial_t\beta^i(0,y).
\]
In a neighborhood of its initial slice the map
\begin{equation}
 \boxed{\Psi_a(t,y)=
       \left(t-\tfrac12t^2c_N(y),\quad y-\tfrac12t^2c(y)\right)}
 \label{v10:coordinate-map}
\end{equation}
is a local diffeomorphism. It fixes the slice pointwise and has differential
$I$ there. For $\widetilde g=\Psi_a^*g$ its initial ADM data satisfy
\begin{equation}
 \widetilde N=1,\quad\widetilde\beta=0,\quad
 \partial_t\widetilde N=\partial_t\widetilde\beta=0,\quad
 \widetilde\gamma=\gamma,\quad\partial_t\widetilde\gamma=\partial_t\gamma,
 \label{v10:normal-jets}
\end{equation}
and the exact spatial second-jet identity is
\begin{equation}
 \boxed{\partial_t^2\widetilde\gamma(0)
      =\partial_t^2\gamma(0)-c_N\partial_t\gamma(0)
                              -\mathcal L_c\gamma(0).}
 \label{v10:normal-acceleration}
\end{equation}
On \eqref{v10:actual-jets} this gives
\begin{equation}
 \widetilde g(a,0)-\eta=O_h(a),\quad
 \partial_t\widetilde g(a,0)=O_h(a),\qquad
 \boxed{\partial_t^2\widetilde\gamma(a,0)
        =-2\mathcal I_h\mathfrak E_h(r_h^\beta)+O_h(a).}
 \label{v10:normal-limit}
\end{equation}
The domain of the coordinate map does not supply an existence interval for
the original action history.
\end{theorem}
\begin{proof}
At $t=0$ the map and its differential are identity, so the inverse-function
theorem gives a local diffeomorphism near the slice. The coefficient fields
are smooth and the slice is compact; one may use a sufficiently small
neighborhood uniformly in $y$ at this fixed metric. The restriction must
also remain inside the metric's existing domain.

Differentiate
$\widetilde g_{\mu\nu}
 =g_{\alpha\beta}(\Psi)\partial_\mu\Psi^\alpha
                         \partial_\nu\Psi^\beta$ at $t=0$.
The only second derivatives of $\Psi$ with no factor of $t$ are
$\partial_t^2\Psi^0=-c_N$ and
$\partial_t^2\Psi^i=-c^i$.
Since $g_{00}=-1,g_{0i}=0$ there,
\[
 \partial_t\widetilde g_{00}
   =-2c_N+2c_N=0,\qquad
 \partial_t\widetilde g_{0i}
   =\gamma_{ij}c^j-\gamma_{ij}c^j=0.
\]
These identities and $\widetilde\beta(0)=0$ give the lapse and shift
parts of \eqref{v10:normal-jets}. The spatial first jet is unchanged.
Differentiating its spatial block a second time gives
\[
 \begin{split}
 \partial_t^2\widetilde\gamma_{ij}
 ={}&\partial_t^2\gamma_{ij}-c_N\partial_t\gamma_{ij}
       -c^k\partial_k\gamma_{ij}\\
    &-\gamma_{kj}\partial_i c^k-\gamma_{ik}\partial_jc^k.
 \end{split}
\]
The terms involving spatial derivatives of $c_N$ multiply $g_{0i}(0)=0$;
products of first mixed derivatives of $\Psi$ vanish at the slice.
This proves \eqref{v10:normal-acceleration} exactly.

For \eqref{v10:actual-jets}, $c_N$ and $c$ are bounded at fixed $h$,
$\gamma_t=O_h(a)$ and $\gamma=I+O_h(a)$ in every needed spatial norm.
Thus the limiting second jet is
$\mathcal I_h(\delta_i r^\beta_j+\delta_jr^\beta_i
                 -D_i r^\beta_j-D_jr^\beta_i)$.
This is $-2\mathcal I_h\mathfrak E_h$. The other assertions follow from
\eqref{v10:normal-jets}.
\end{proof}

\begin{corollary}[No coordinate removal of the certified scalar observation]
\label{v10:no-geometric-removal}
At the actual phase-selected three-site family, any coordinate
re-expression fixing the same initial events leaves the scalar limit
\eqref{v10:Kretschmann} unchanged. In particular the genuine normalization
\eqref{v10:coordinate-map} cannot make the complete initial metric
curvature $O_h(a)$. With identity differential at the slice its electric
components are also unchanged there.
\end{corollary}
\begin{proof}
The Riemann tensor is a tensor of the reconstructed metric and
$\mathscr K(\Psi^*g)=\mathscr K(g)\circ\Psi$.
At the same initial events the scalar values coincide, independently of
the second or higher derivatives of the coordinates. The strictly positive
limit from \cref{v10:verified-E} cannot tend to zero. For identity
differential the component transformation is itself identity at the cut.
\end{proof}

This corollary does not assert that the finite action is diffeomorphism
covariant. The coordinate change is applied to the already reconstructed
metric, while original literal observations remain evaluations on the
same original history. The map fixes the slice, so no source value is
reassigned to another initial event. V9's affine canonical action remains
an exact re-expression with its prescribed path held fixed under variation.
Its pulled-back metric observation recovers the curvature above. The small
canonical acceleration cannot instead be inserted into an unchanged
geometric read while discarding the transport terms.

\section{The specific remaining source-dependent inverse and verification}
\label{v10:status}

The sharper geometric test distinguishes the controlled canonical jets
from the metric curvature on the actual fixed-cutoff phase-selected family. The complete initial weak shift velocity is zero, and an
exact canonical transport makes the initial canonical acceleration small.
Nevertheless the ordinary interpolated metric has nonzero initial Ricci
and Riemann-squared limits while all original link-curvature coefficients
vanish. This is a proved observation mismatch on that fixed grid, not a
rejection of the separately retained stable Einstein writer, nor a cofinal
failure theorem for the raw action.

The new positive comparison identifies the rate component and derivative
weight that a cofinal geometric argument must control. In any regular
active chart where full weak-rate initial selection is actually available,
it is the composite output
\begin{equation}
 \boxed{-\tfrac12\operatorname{SymGrad}_{D-\delta}
       P_T S_\beta\,\Gamma_{A,h}^{-1}b_{A,h},}
 \label{v10:actual-composite}
\end{equation}
where $\operatorname{SymGrad}_Q b$ has entries $Q_i b_j+Q_j b_i$ and
$S_\beta$ is the shift synthesis from the selected active representatives.
Changing that representative basis changes neither the physical projected
rate nor this observation. The output includes the full inverse before
applying its frequency weight; smoothness of the un-inverted source does
not bound it. A source-specific estimate for this composite could succeed
even when the unweighted multiplier inverse or lapse rate is large.

The finite-amplitude remainder, the nonlinear stationary Hessian, the
literal-curvature synthesis, and later physical times remain distinct
requirements. In particular \eqref{v10:weighted-rate} is only an initial
metric-curvature condition. It does not propagate the phase preparation,
bound $\lambda_{tt}$, or establish a common physical lifespan. These
statements keep the fixed-cutoff amplitude limit separate from a joint
limit and prevent another coordinate cancellation from being mistaken for
the completed exact-action bridge.

\paragraph{Verification.}
The new verifier derives the active-solve enclosure by exact integer
approximate-inverse residuals, constructs the $27$-site Hodge matrices over
$\mathbb Q$, and evaluates the full symmetric-gradient observation with
rational bounds on $\pi$ and $\sqrt3$. It checks the stated point value,
positive norm, trace, and Riemann-squared bounds by rational inequalities.
The Riemann calculation is separately checked symbolically on all $256$
components of the limiting initial jet, retaining arbitrary pure-time
second lapse and shift derivatives until they cancel. The Ricci and
Einstein contractions and the positive scalar identity are checked from
those components. A separate exact rational polynomial-metric calculation
checks the initial geometric pullback with nonconstant rates and spatial
metric derivatives. General-grid Fourier calculations at $N=3,5,9,17,33$
are explicitly labelled floating diagnostics of the all-frequency identity;
the uniform estimates rest on the sine inequality proof.

The source audit inspected V9's canonical and physical boundaries, V6's
specified metric interpolation, V8's source-enclosure conventions, and the
older source's Gauss--Codazzi comparison. It does not claim an exhaustive
literature-priority result. The first nine marked bodies remain byte for
byte unchanged, including their corrections and limitations. Preservation
hashes and compilation verify the artifact, not the entire inherited
nonlinear mathematical theory. No nonlinear original-action trajectory or
new cofinal signed-rank family has been computed.

\begin{thebibliography}{9}
\raggedright
\setcounter{enumiv}{10}
\bibitem{v10:G}
E. Gourgoulhon, \emph{3+1 Formalism and Bases of Numerical Relativity},
lecture notes (2007), arXiv:gr-qc/0703035.
\url{https://arxiv.org/abs/gr-qc/0703035}.
This supplies primary-source context for ordinary hypersurface curvature
and coordinate transformations, not the original-action source,
phase-root certificate, or finite-symbol observation proved here.
\end{thebibliography}
% END IMMUTABLE EXACT-ACTION BRIDGE V10 BODY


% Future iterations append here.
% BEGIN IMMUTABLE EXACT-ACTION BRIDGE V11 BODY
\clearpage
\renewcommand{\thesection}{\arabic{section}}
\renewcommand{\theHsection}{bridgeV11.\arabic{section}}
\setcounter{section}{82}
\section{The weighted observation requires more than a small multiplier residual}
\label{v11:context}

The observation isolated in \cref{v10:actual-composite} contains the complete
active inverse before the difference-of-derivatives weight is applied.
Its three-derivative weight does not justify replacing that inverse by a
uniform elliptic estimate. This body separates the part of that observation
which can already be bounded from the scalar lapse problem which remains,
and then tests the principal form of the latter on the actual balanced
background. The resulting characteristic is present at interior Fourier
scales; it is distinct from the folded-face effect in V7.

There are three original-source conclusions. First, exact elimination of the
positive transverse-shift block writes the entire active geometric output as
an explicitly controlled direct term and one scalar-lapse-mediated term.
The direct term is $O(h)$ on the balanced prepared source. Second, the lapse
transfer has a sharp $h^{-2}$ amplification on a specified Fourier-edge
input, so its edge contribution cannot be hidden in an ordinary smooth
multiplier bound. Third, even away from that edge the frozen signed active
form reaches the trace fraction $2/3$ at an explicitly isolated real
characteristic. Its source columns remain independent and its shift output
is nonzero. Localized trigonometric packets turn this characteristic into
actual finite multiplier tests with small $H_h^1$ Hessian images and
nonvanishing geometric observations.

The last tests are not the prepared source $b_{A,h}$. They exclude a uniform
estimate for arbitrary source errors in that norm; they do not prove that
the original prepared source excites the characteristic or that the
raw-action Einstein limit fails. No field, multiplier equation, literal
product, or physical clock is altered.

\paragraph{Dependencies and nonduplication.}
We use the original first-source coefficient and its transverse block from
\cref{v1:first-source,v1:transverse-rank,v1:shift-projector}, the balanced
first jets and their prepared quadratic target from
\cref{v4:neutral,v5:prepared-B2,v7:target-even,v9:scalar-test}, and the
geometric observation \eqref{v10:E}. The ordinary block elimination used
below is not a new abstract theorem; V1 already supplies its homogeneous-target
special case. The new content is the nonzero-target observable split, its
uniform direct estimate and sharp edge behavior, and the exact characteristic
and finite-packet obstruction for this same action coefficient. No higher
source $J_2$, phase-root enclosure, or nonlinear existence theorem is
recomputed or inferred in this body. Localized high-frequency tests and
Schur complements are standard methods; the finite proofs below do not
invoke an external microlocal existence theorem.

\section{An exact scalar-lapse reduction of the geometric output}
\label{v11:split}

Work on the physical active Hilbert space
\[
 \mathcal M_{A,h}=\{n:\overline n=0\}\oplus\mathcal T_h,
 \qquad
 \mathcal T_h=\{\beta:\overline\beta=0,
                  \ \kappa_h(k)\cdot\widehat\beta(k)=0\}.
\]
Let $L_hn=J_{1,h}(n,0)$ and $T_h\beta=J_{1,h}(0,\beta)$.
Here $T_h$ is a transverse source operator, not a trace-angle compression.
Use the positive Frobenius pairing on symmetric connection matrices and
retain $K_0^{-1}=\tfrac12I-\tfrac34P_{\rm tr}$. Set
\begin{equation}
 B_h=T_h^*T_h,\quad T_h^\dagger=B_h^{-1}T_h^*,\quad
 \Pi_{\beta,h}=T_hT_h^\dagger,\quad
 F_{L,h}=(I-\Pi_{\beta,h})L_h,\quad
 \mathscr S_{L,h}=F_{L,h}^*K_0^{-1}F_{L,h}.
 \label{v11:scalar-data}
\end{equation}
The inverse in $B_h^{-1}$ is on $\mathcal T_h$ only. It exists at every
odd cutoff by \cref{v1:transverse-rank}; no lapse or full-source regularity
is needed to define it. The remaining scalar matrix has dimension $V-1$,
where $V=N^3$.

Write $b_h=(b_{N,h},b_{\beta,h})$ for the correctly ordered physical
prepared active target, and use $\mathfrak E_h$ for the geometric operator
of \eqref{v10:E}.

\begin{proposition}[Complete target and observation retained under shift elimination]
\label{v11:Schur}
The complete leading equation
\[
 (L_h,T_h)^*K_0^{-1}(L_h,T_h)(n,\beta)=-b_h
\]
is equivalent to
\begin{align}
 \beta_0&=-2B_h^{-1}b_{\beta,h},\qquad
 \beta=\beta_0-T_h^\dagger L_hn,\label{v11:beta-split}\\
 \mathscr S_{L,h}n
   &=-b_{N,h}+L_h^*T_hB_h^{-1}b_{\beta,h}.
                                                    \label{v11:lapse-equation}
\end{align}
The full active signed Gram is invertible if and only if
$\mathscr S_{L,h}$ is invertible. Every solution has the exact observation
\begin{equation}
 \boxed{\mathfrak E_h(\beta)
   =-2\mathfrak E_h B_h^{-1}b_{\beta,h}
                         -\mathfrak E_h T_h^\dagger L_hn.}
 \label{v11:output-split}
\end{equation}
\end{proposition}
\begin{proof}
The transverse source is trace-free, so $K_0^{-1}T_h=T_h/2$ and
$T_h^*K_0^{-1}T_h=B_h/2$. The transverse row of the original equation
is therefore $\tfrac12T_h^*L_hn+\tfrac12B_h\beta=-b_{\beta,h}$,
which is exactly \eqref{v11:beta-split}. Substitute it in the lapse row.
Since $\Pi_{\beta,h}$ is orthogonal and its range is trace-free,
$K_0^{-1}$ commutes with it and equals $I/2$ there. Hence
\[
 L_h^*K_0^{-1}L_h-\tfrac12L_h^*\Pi_{\beta,h}L_h
      =F_{L,h}^*K_0^{-1}F_{L,h}.
\]
This proves \eqref{v11:lapse-equation}, the invertibility equivalence by
block elimination, and the converse by substitution into both original
rows. Linearity of the unchanged geometric observation gives
\eqref{v11:output-split}.
\end{proof}

\begin{theorem}[The direct transverse contribution is already controlled]
\label{v11:direct}
For the balanced scalar-momentum family with radii in a fixed compact
positive set, and every fixed Sobolev index $s$, there is a constant independent
of the odd cutoff such that for every $f\in\mathcal T_h$,
\begin{equation}
 \|\mathfrak E_h B_h^{-1}f\|_{H_h^s}
                           \le C_s\|f\|_{H_h^{s+1}}.
 \label{v11:direct-operator}
\end{equation}
For the actual prepared target on that family, on sufficiently fine grids,
\begin{equation}
 \|b_{\beta,h}\|_{H_h^{s+1}}\le C_s h,
 \qquad
 \boxed{\|2\mathfrak E_h B_h^{-1}b_{\beta,h}\|_{H_h^s}\le C_s h.}
 \label{v11:direct-target}
\end{equation}
These estimates do not include the lapse-mediated term in
\eqref{v11:output-split}.
\end{theorem}
\begin{proof}
The exact transverse estimate in \cref{v1:transverse-rank}, with its
scalar momentum bounded away from zero, gives on each nonzero Fourier mode
\[
 \|B_h(k)^{-1}\|_{\mathcal T_h(k)\to\mathcal T_h(k)}
                      \le C h^{-2}|k|^{-2}.
\]
The exact symbol estimate of \cref{v10:symbol} bounds the geometric
operator on that mode by $C h^2|k|^3$. Their product is bounded by
$C|k|$, proving \eqref{v11:direct-operator} by Parseval without a
frequency restriction on $f$.

For the second assertion, the shift target is quadratic in the first records,
with types $U^2,p^2$. Its Riesz vector therefore has fixed support
$|k|_1\le2$ once the grid is sufficiently fine; restriction to the transverse
space preserves support. The geometric coefficients on this fixed bank
are analytic in $h^2$. Their value at $h=0$ vanishes: the
ordinary-derivative limit of the original stationary canonical action has
$\{P[\beta],H[1]\}=0$ by spatial pullback of the integrated scalar
Hamiltonian density. This is the same off-shell Ward input used in the
retained weak-source expansion, now applied before restricting the shift
test. Equivalently, expanding the complete prepared coefficient
\eqref{v5:prepared-drift-formula} with ordinary derivatives gives this
identity. A pure shift has no linear-drift preparation correction. Thus
its geometric part is $O(h^2)$ on this fixed bank. Every contributing
cubic electric or magnetic remainder has its explicit factor $h$ and
bounded fixed-bank coefficients, so their total is $O(h)$. There are
only finitely many coefficients, uniformly bounded in the phases and
compact radii. This proves the first estimate in \eqref{v11:direct-target};
\eqref{v11:direct-operator} proves the second. Weak neutrality removes
constant and curl-free shift components but was not used to delete any
transverse logarithmic term.
\end{proof}

Thus one part of the requested cofinal geometric estimate is a consequence
of existing source structure, not a new assumption on its inverse. The
remaining difficulty is the specific scalar solve
\eqref{v11:lapse-equation} observed through
$\mathfrak E_hT_h^\dagger L_h$.

\section{The lapse transfer has a genuine Fourier-edge loss}
\label{v11:edge}

For a source at a Fourier mode let its off-diagonal coordinates be
$(12,13,23)$; their Frobenius norm is twice the ordinary coordinate norm
squared. Put
\begin{equation}
 \mathsf L_z(k)=
 \begin{pmatrix}z_2&z_1&0\\z_3&0&z_1\\0&z_3&z_2\end{pmatrix},
 \quad z_i=e^{2\pi\mathrm i h k_i},\quad
 \widehat w=-\tfrac{\mathrm i}{2}\kappa_h(k)\times\widehat\beta.
 \label{v11:shift-symbol}
\end{equation}
The exact off-diagonal source is
$(hp_0/4)\mathsf L_z(k)\widehat w$, where $p=p_0I$.
This is just \eqref{v1:shift-source} with $\varkappa_h=-p_0/2$.
The singular values $2,1,1$ of $\mathsf L_z$ are inherited from V1.

Let $L_{\partial,h}$ retain only those entries of the \emph{geometric}
lapse-source product defects whose multiplication by a unit axial first-field
frequency folds across a Fourier face. Define
$L_{{\rm reg},h}=L_h-L_{\partial,h}$. This is an analytical decomposition,
not a truncation of the action. Electric terms, including their folded
products, remain in $L_{{\rm reg},h}$.

\begin{theorem}[Weighted transfer with its edge term retained]
\label{v11:edge-bound}
On the balanced scalar-momentum family, uniformly in phase and compact
positive radii, for every fixed real $s$,
\begin{equation}
 \boxed{\|\mathfrak E_hT_h^\dagger L_h n\|_{H_h^s}
 \le C_s h^2\|n\|_{H_h^{s+3}}
       +C_s h^{-2}\|\mathscr P_{\partial,h}n\|_{H_h^s}.}
 \label{v11:edge-estimate}
\end{equation}
The second term cannot in general be omitted or assigned the first term's
scale. On the unit zero-phase family there are real mean-zero scalar arrays
$n_h$ with $\|n_h\|_h=1$, supported on the two modes $(\pm m,0,0)$, such that
\begin{equation}
 \|\mathfrak E_hT_h^\dagger L_hn_h\|_h\ge c h^{-2}
                       \qquad(N=2m+1\text{ sufficiently large}).
 \label{v11:edge-sharp}
\end{equation}
No lapse inverse or full-rank hypothesis is needed in either statement.
\end{theorem}
\begin{proof}
A nonfolded geometric coefficient contains
$\kappa_j(k+\ell)-\kappa_j(k)-\kappa_j(\ell)$ with
$\ell=\pm e_i$. It vanishes for $j\ne i$. For $j=i$, the sine addition
formula gives the bound $C h^2\langle k\rangle^2$ over the entire
nonfolded range. An electric coefficient is bounded by
$C h\langle k\rangle$. All its outputs differ from the input frequency
by at most one unit; folded input and output magnitudes are still comparable.
The transverse pseudoinverse at a nonzero output has norm at most
$C/(h|k_{\rm out}|)$, and the observation has norm at most
$C h^2|k_{\rm out}|^3$. Their product is at most
$C h\langle k\rangle^2$. Multiplication by the two preceding source
bounds gives $C(h^2\langle k\rangle^3+
 h^3\langle k\rangle^4)\le C h^2\langle k\rangle^3$,
since $h\langle k\rangle$ is uniformly bounded on the grid. Zero output
modes have zero transverse inverse. Summing the finite number of offsets
proves the first term of \eqref{v11:edge-estimate}.

A folded geometric entry has input and output on $\Sigma_h$ and size at
most $C/h$. There $|k_{\rm out}|\asymp h^{-1}$, the transverse inverse
is $O(1)$ and the geometric observation is $O(h^{-1})$. This gives the
second term, with the actual Fourier-face input retained.

For sharpness first use the complex scalar $n=e^{2\pi\mathrm i m x_1}$.
At output $(-m,0,0)$ only the $+e_1$ first profile can fold this input.
Write
\[
 \kappa_m=2h^{-1}\sin(\pi hm),\quad
 d_m=2\pi m-\kappa_m,\quad
 \Delta_m=-2\kappa_m-\omega_h.
\]
Substitution in the original geometric first-source formula gives the
$(23)$ entry $\Delta_m/2$, the diagonal entries
$(0,-\mathrm i\Delta_m/2,\mathrm i\Delta_m/2)$, and zero other
off-diagonal entries. The transverse inverse ignores the diagonal entries.
For a target $q$ in off-diagonal slot $23$ at this output, direct inversion
of \eqref{v11:shift-symbol} on $\beta_1=0$ gives
\[
 \|\mathfrak E_hT_h^\dagger(qe_{23}^{\rm off})\|_F
                     =\frac{8|d_m q|}{3h|p_0|\kappa_m}.
\]
Indeed the two-column least-squares matrix has Gram
$\left(\begin{smallmatrix}2&-1\\-1&2\end{smallmatrix}\right)$,
up to the common nonzero factor, and its inverse maps the $23$ load to
opposite coefficients of equal size. The positive Frobenius factors cancel
in the inverse and are retained in the final observation. Setting
$q=\Delta_m/2$ gives
\begin{equation}
 \frac{4|d_m\Delta_m|}{3h|p_0|\kappa_m},\qquad
 h^2\frac{4|d_m\Delta_m|}{3h p_0\kappa_m}
               \longrightarrow\frac{4(\pi-2)}{3\pi}>0,
 \label{v11:edge-limit}
\end{equation}
where $p_0=\omega_h\to2\pi$ on the unit family. The electric source at
this output is $O(1)$, so its observed contribution is only $O(h^{-1})$.
It cannot cancel the displayed leading term. The real array
$n_h=\sqrt2\cos(2\pi m x_1)$ has two conjugate folded outputs, whose
combined squared norm equals the complex calculation. This proves
\eqref{v11:edge-sharp}. It is an operator witness, not a claim that the
actual scalar lapse solve has this Fourier content.
\end{proof}

\section{The actual interior principal form is not uniformly elliptic}
\label{v11:principal}

The edge loss is not the only difficulty. Work now with the unit balanced
family at arbitrary phases and with all handedness signs positive. Let
$r_{ia,h}(x)=-\epsilon_{abc}\delta_{b,h}U_{ic}(x)/2$ be the rotational
first connection. Its continuum coefficient is $r_0(x)=\pi R(x)$,
where $R$ is real symmetric and trace-free. In the regime
$1\ll\ell\ll h^{-1}$, at wave vector $\ell\xi$, the leading source
has scale $2\pi^2h\ell$ and real symbol
\begin{equation}
 \mathsf B(x,\xi)(n,\beta)=\ell_{R(x)}(\xi)n+	au_\xi\beta,
 \qquad \xi\cdot\beta=0,
 \label{v11:B}
\end{equation}
with
\begin{align}
 [\ell_R(\xi)]_{ia}
  &=-\tfrac12\bigl(R_{ii}\xi_a+R_{aa}\xi_i
       +\delta_{ia}[(R\xi)_i+(R\xi)_a]
                      -2R_{ia}(\xi_i+\xi_a)\bigr),\label{v11:ell}\\
 (\tau_\xi\beta)_{ii}&=0,\qquad
 (\tau_\xi\beta)_{ij}
          =-\tfrac14\bigl[(\xi\times\beta)_i+
                          (\xi\times\beta)_j\bigr]\quad(i\ne j).
                                                    \label{v11:tau}
\end{align}
The original source includes a common factor $\mathrm i$ when applied to
an oscillatory mode. It cancels from its Hermitian signed Gram.
These formulas describe a limit of the \emph{given} source, not a replacement
finite evolution. A finite-packet statement justifying their use is proved
in the next section.

At a point where the three profile angles satisfy
\begin{equation}
 \theta_1=0,\quad\theta_2=3\pi/2,\quad
        (\cos\theta_3,\sin\theta_3)=(3/5,-4/5),
 \label{v11:point}
\end{equation}
we obtain the exact matrix
\begin{equation}
 R_*=
 \begin{pmatrix}-1/5&-3/5&0\\-3/5&-4/5&-1\\0&-1&1\end{pmatrix}.
 \label{v11:Rstar}
\end{equation}
Such a spatial point exists for \emph{every} choice of the three initial
phases: set $x_i=(\theta_i-\phi_i)/(2\pi)$ modulo one. Thus phase selection
cannot remove this local coefficient pattern.

\begin{theorem}[A simple interior characteristic with independent source columns]
\label{v11:characteristic}
Let $\xi=(1,t,0)$ and $R=R_*$. After eliminating the two positive transverse
columns from $\mathsf B^*K_0^{-1}\mathsf B$, its scalar lapse Schur form is
\begin{equation}
 \boxed{q(t)=
 \frac{103t^4-168t^3+146t^2-86t+19}
                    {25(3t^2-2t+3)}.}
 \label{v11:q}
\end{equation}
It is positive at $t=0$, negative at $t=1/2$, and has a unique simple zero
$t_*$ in $(39/100,2/5)$. At this zero the three source columns of
$\mathsf B$ remain independent, but the signed Gram has a one-dimensional
kernel. Its source/trace compression has an eigenvalue exactly $2/3$.
A kernel multiplier can be chosen $(1,\beta_*)$, with $\beta_*\ne0$ and
$\xi_*\cdot\beta_*=0$.
\end{theorem}
\begin{proof}
The source \eqref{v11:ell} at this point is
\[
 \ell_{R_*}(1,t,0)=
 \begin{pmatrix}
 3t/5&-t/2-1/5&-1/2\\
 -t/2-1/5&3/5&-3t/2\\
 -1/2&-3t/2&t
 \end{pmatrix}.
\]
Write $d$ for its three diagonal entries and $o$ for its off-diagonal
coordinates $(12,13,23)$. Put
\[
 M=\begin{pmatrix}1&1&0\\1&0&1\\0&1&1\end{pmatrix},\qquad
 \nu=M^{-T}\xi
       =\tfrac12(1+t,1-t,t-1)^T.
\]
The transverse off-diagonal range is $M\xi^\perp=\nu^\perp$.
Eliminating it leaves off-diagonal vector
$\nu(\nu\cdot o)/|\nu|^2$. The Cartan inverse on diagonal matrices
is $I/2-\mathbf1\mathbf1^T/4$, and on off-diagonal matrices is $I/2$
in the Frobenius pairing. Therefore the exact Schur scalar is
\begin{equation}
 q=\tfrac12|d|^2-\tfrac14(\mathbf1\cdot d)^2
                             +\frac{(\nu\cdot o)^2}{|\nu|^2}.
 \label{v11:scalar-principal}
\end{equation}
Direct rational multiplication gives \eqref{v11:q}. In particular
$q(0)=19/75$ and $q(1/2)=-3/100$.
For $p(t)=103t^4-168t^3+146t^2-86t+19$,
\[
 p(39/100)=\frac{8385223}{10^8}>0,\qquad
 p(2/5)=-\frac{97}{625}<0.
\]
Throughout this interval,
\[
 p'(t)=412t^3-504t^2+292t-86
             \le-\frac{24363}{1250}<0.
\]
The intermediate-value theorem and strict monotonicity give the unique
simple root; the approximate value $0.3934666194$ is not its definition.

The transverse source is injective on $\xi^\perp$, since both the cross
product there and $M$ on $\mathbb R^3$ are invertible. The lapse source
has the diagonal entry $3/5$, whereas every transverse column has zero
diagonal. Thus the complete source has rank three at the characteristic.
The transverse Gram is positive definite and the scalar Schur complement
vanishes simply; hence the full signed Gram has nullity one.

For an explicit kernel multiplier define
\[
 w=M^{-1}\left(o-\nu\frac{\nu\cdot o}{|\nu|^2}\right),\qquad
 \beta(t)=-\frac{4\xi\times w}{|\xi|^2}.
\]
Then $w\perp\xi$ and $\tau_\xi\beta$ cancels the transverse part of
$o$. This proves the kernel assertion at $t_*$. In coordinates,
\begin{equation}
 \begin{aligned}
 \beta_1(t)&=\frac{2t(25t^3-2t^2+9t+8)}
                {5(1+t^2)(3t^2-2t+3)},\\
 \beta_2(t)&=-\frac{2(25t^3-2t^2+9t+8)}
                {5(1+t^2)(3t^2-2t+3)},\\
 \beta_3(t)&=\frac{2(5t^2+16t-1)}{5(3t^2-2t+3)}.
 \end{aligned}
 \label{v11:beta-star}
\end{equation}
The numerator of $-\beta_2$ is positive on the isolating interval, so the
shift is nonzero. Finally orthonormalize the three source columns as
$\mathsf B=QR$. Since
\[
 \mathsf B^*K_0^{-1}\mathsf B
          =\tfrac14R^*(2I-3Q^*P_{\rm tr}Q)R,
\]
its singularity is precisely a $2/3$ eigenvalue of the compression. At a
single spatial point the trace space has dimension one, so the other two
compression eigenvalues are zero. No global finite-grid spectral gap is
inferred from this frozen calculation.
\end{proof}

The sign change, column independence, nonzero shift, and simple root persist
under sufficiently small changes of the three radii, with their calibrated
scalar momentum, by continuity and the ordinary implicit-function theorem.
The point \eqref{v11:point} moves with the phases. These are assertions about
the local source symbol, not cofinal phase roots or the global source rank.

\section{Interior finite packets with a nonvanishing geometric output}
\label{v11:packets}

The preceding characteristic is not used merely to suggest a bad numerical
condition number. We now construct tests of the actual finite matrix
$\mathcal M_{2,h}=J_{1,h}^*K_0^{-1}J_{1,h}$ and keep the geometric
observation. All norms in this section use the same Fourier weights as V10;
$H_h^1$ of a multiplier pair includes the lapse and all three shift slots.

\begin{lemma}[Localized trigonometric packets with exact discrete normalization]
\label{v11:packet-lemma}
Let $\ell\to\infty$, $M=\lfloor\ell/8\rfloor$, and suppose
$h\ell\to0$. Center at any $x_*$ and put
\[
 A_M(x)=\prod_{j=1}^3\left(\frac{1+\cos(2\pi(x_j-x_{*,j}))}{2}\right)^M,
 \qquad a_M=A_M/\|A_M\|_{L^2}.
\]
Let $k_\ell=(\ell,\operatorname{round}(t_*\ell),0)$ and
$\psi_h=\mathcal S_h(a_Me^{2\pi\mathrm i k_\ell\cdot x})$.
For sufficiently large $\ell$ and fine odd grids, $\|\psi_h\|_h=1$;
its Fourier support is contained in
$k_\ell+\{-M,\ldots,M\}^3$, disjoint from zero and from the grid faces.
Moreover, uniformly in the center,
\begin{equation}
 \|\operatorname{dist}_{\T^3}(x,x_*)\psi_h\|_h\le C\ell^{-1/2},
 \quad
 \|(D_h-2\pi\mathrm i\ell\xi_*)\psi_h\|_h\le C\ell^{1/2}.
 \label{v11:packet-bounds}
\end{equation}
Multiplication by a uniformly Lipschitz coefficient $f$ therefore replaces
$f(x)$ by $f(x_*)$ on this packet at an $O(\ell^{-1/2})$ relative error.
\end{lemma}
\begin{proof}
Each envelope factor has the exact Fourier expansion
\[
 \cos^{2M}(\pi y)=4^{-M}\sum_{j=-M}^{M}
                    \binom{2M}{M+j}e^{2\pi\mathrm i j y}.
\]
This proves the support assertion. Its squared modulus has coordinate degree
at most $2M$, so its discrete mean equals its integral when $2M<N$,
independently of the center and carrier. Hence the normalization is exact.
For $I_M=\int_{-1/2}^{1/2}\cos^{4M}(\pi y)\,dy$, the elementary beta-integral
recursion, or integration by parts, gives
\[
 \frac{\int\sin^2(\pi y)\cos^{4M}(\pi y)\,dy}{I_M}
       =\frac1{4M+2},\qquad
 \frac{\int|\partial_y\cos^{2M}(\pi y)|^2\,dy}{I_M}
       =\frac{4\pi^2M^2}{4M-1}.
\]
These identities also hold for the sampled polynomial expressions when
$2M+1<N$. On $[-1/2,1/2]$, $|y|\le|\sin\pi y|/2$, proving the first
bound, coordinate by coordinate. The derivative identity gives
$\|D_ha_M\|_h\le C\sqrt M$. Differentiation of the carrier adds
$2\pi\mathrm i k_\ell$, and $|k_\ell-\ell\xi_*|\le1$, giving the
second bound. The stated coefficient replacement follows from the first
bound and the Lipschitz inequality. These are finite quadrature identities,
not an assumption about concentration of the unknown active solution.
\end{proof}

\begin{theorem}[Geometrically visible interior quasimodes of the original active Hessian]
\label{v11:quasimodes}
On the unit balanced scalar-momentum family, for every phase sequence and
all sufficiently fine odd cutoffs, there are real physical active multiplier
arrays $\nu_h=(n_h,\beta_h)$ with
\begin{equation}
 \|\nu_h\|_h=1,\qquad
 \boxed{\|\mathcal M_{2,h}\nu_h\|_{H_h^1}\le C h^{1/3},\qquad
        \|\mathfrak E_h(\beta_h)\|_h\ge c>0.}
 \label{v11:quasimode-result}
\end{equation}
Their spectra, and the spectra of $J_{1,h}\nu_h$ and
$\mathcal M_{2,h}\nu_h$, lie strictly inside the retained frequency box.
Also $\|J_{1,h}\nu_h\|_h\asymp h^{1/3}$. The constants are uniform in
the phases. Neither global source regularity nor an actual nonlinear
history is assumed in this construction.
\end{theorem}
\begin{proof}
Take $\ell=\lfloor h^{-2/3}\rfloor$ and the packet of
\cref{v11:packet-lemma}, centered at the point \eqref{v11:point} for the
current phases. Put $\xi_*=(1,t_*,0)$, $\beta_*=\beta(t_*)$, and
$a_*=(\beta_*\times\xi_*)/|\xi_*|^2$, so
$\xi_*\times a_*=\beta_*$. Define first the complex active pair
\begin{equation}
 \nu_h^{\C}=
 \left(\psi_h,\;\frac1{2\pi\mathrm i\ell}
                         \curl_\delta(a_*\psi_h)\right).
 \label{v11:packet-rate}
\end{equation}
The lapse mean is zero by support. The shift mean is zero and its phase
divergence vanishes exactly, since the phase derivatives commute. Its
norm tends to $\sqrt{1+|\beta_*|^2}$. Indeed the exact derivative symbols
and \eqref{v11:packet-bounds} give
\[
 \beta_h^{\C}=\beta_*\psi_h+
ho_h,\qquad
 \|\rho_h\|_h\le C(\ell^{-1/2}+(h\ell)^2).
\]
All these fields have support in a box of radius $M$ about $k_\ell$.

We give the coefficient error bounds needed to apply the frozen symbol.
On any such box, expanded by at most two unit offsets, the original first
source and its adjoint have operator norm at most $Ch\ell$ for small
$h\ell$. The electric and transverse terms have one phase derivative
and a factor $h$. The geometric terms have a unit-profile product defect;
there is no folding in this box, so the sine addition estimate used in
\cref{v11:edge-bound} bounds them by $Ch^2\ell^2$. There are finitely
many such terms and the shifts are isometric. The same bound holds for
the adjoints by their finite Fourier matrices.

The electric terms on the packet have the following replacements: the
rotational coefficient $r_h(x)$ can be replaced by $\pi R_*$ with error
$C(\ell^{-1/2}+h^2)$ on its normalized envelope; its differentiated
carrier can be replaced by $2\pi\mathrm i\ell\xi_*$ with relative
error $C(\ell^{-1/2}+(h\ell)^2)$; and every unit shift of the carrier
and envelope differs from identity by $Ch\ell$ in norm. Derivatives
of coefficient factors in the adjoint add at most $C/\ell$ after this
normalization. The scalar momentum is $2\pi+O(h^2)$.
The geometric term itself contributes only relative $Ch\ell$ at
scale $h\ell$. Therefore, writing
$\epsilon_h=\ell^{-1/2}+h\ell$ and
$v_*=(1,\beta_*)$, we obtain
\begin{align}
 J_{1,h}\nu_h^{\C}
    &=2\pi^2\mathrm i h\ell\,
                 \mathsf B(x_*,\xi_*)v_*\,\psi_h+e_h,
             &\|e_h\|_h&\le Ch\ell\epsilon_h,\label{v11:packet-source}\\
 J_{1,h}^*(S\psi_h)
    &=-2\pi^2\mathrm i h\ell\,
                 \mathsf B(x_*,\xi_*)^*S\,\psi_h+e'_h,
             &\|e'_h\|_h&\le Ch\ell\epsilon_h|S|
                                                        \label{v11:packet-adjoint}
\end{align}
for every fixed symmetric tensor $S$. The norm pairings in these formulas
are the physical scalar/vector and Frobenius pairings. The estimates just
listed follow term by term from \eqref{v11:ell} and
\eqref{v11:shift-symbol}; no pseudodifferential inverse or cutoff-uniform
stationary assumption is involved.

Apply \eqref{v11:packet-adjoint} to the main term of
\eqref{v11:packet-source}, with $K_0^{-1}$ retained. The signed symbol
annihilates $v_*$ by \cref{v11:characteristic}. The adjoint acting on
$e_h$ has norm at most $Ch\ell$ because that error still has the same
bounded-offset Fourier support. Consequently
\[
 \|\mathcal M_{2,h}\nu_h^{\C}\|_h
           \le Ch^2\ell^2\epsilon_h,\qquad
 \|\mathcal M_{2,h}\nu_h^{\C}\|_{H_h^1}
           \le Ch^2\ell^3\epsilon_h.
\]
The last inequality uses the exact support bound $|k|\le C\ell$.
The nonzero diagonal entry of $\mathsf Bv_*$ and
\eqref{v11:packet-source} also give
$\|J_{1,h}\nu_h^{\C}\|_h\asymp h\ell$.

For the geometric observation, the sine expansion on the same support gives
\[
 d_{h,i}(k)=\frac{\pi^3}{3}h^2k_i^3+O(h^4|k|^5).
\]
The first derivative envelope bound and its finite bandwidth show that
replacing $k_i^3$ by $\ell^3\xi_{*,i}^3$ costs relative
$C\ell^{-1/2}$: factor the cubic difference and use the bound on
$(D-2\pi\mathrm i\ell\xi_*)\psi_h$. The same argument applies to
$\beta_h^{\C}$, using its explicit curl formula. Hence
\begin{equation}
 \frac{\mathfrak E_h(\beta_h^{\C})}{h^2\ell^3}
  =\frac{\pi^3\mathrm i}{6}
    \bigl(\xi_{*,i}^3\beta_{*,j}+\xi_{*,j}^3\beta_{*,i}\bigr)_{ij}
                         \psi_h+o_{L_h^2}(1).
 \label{v11:visible-symbol}
\end{equation}
The fixed tensor is nonzero: for real nonzero $d$ and $b$,
$\|\sym(d\otimes b)\|_F^2=(|d|^2|b|^2+|d\cdot b|^2)/2>0$.
Here $d=(\xi_{*,i}^3)_i\ne0$ and $b=\beta_*\ne0$.
Since $h^2\ell^3\to1$ and $\epsilon_h=O(h^{1/3})$, the asserted
bounds follow after normalization. Finally take the real part and renormalize.
The positive and negative carrier boxes are disjoint, even after the two
unit source offsets, so their squared norms split equally. All source and
observation operators are real; the same estimates hold for this real pair.
The support remains at frequency $O(h^{-2/3})\ll h^{-1}$ throughout.
\end{proof}

\begin{corollary}[No uniform observation estimate from an arbitrary $H^1$ source residual]
\label{v11:no-uniform}
There is no cutoff-independent $C$ such that, on these active spaces,
\begin{equation}
 \|\mathfrak E_h(\beta)\|_h
       \le C\|\mathcal M_{2,h}(n,\beta)\|_{H_h^1}
 \label{v11:false-estimate}
\end{equation}
for all real multiplier tests and all sufficiently fine cutoffs.
At any regular cutoffs, the composite inverse
$\mathfrak E_hP_\beta\mathcal M_{2,h}^{-1}:H_h^1\to L_h^2$
therefore has norm at least $c h^{-1/3}$ along this construction.
The same obstruction remains if one restricts the tests and their matrix
images away from the folded Fourier faces.
\end{corollary}
\begin{proof}
Insert \cref{v11:quasimodes} into \eqref{v11:false-estimate}.
At a regular cutoff choose the exact source residual
$f_h=\mathcal M_{2,h}\nu_h$; its inverse image is precisely $\nu_h$.
The ratio gives the stated norm lower bound. The support conclusion is
part of the packet theorem, not a deletion of modes from the matrix.
\end{proof}

\section{The actual prepared source has an exact finite interaction separation}
\label{v11:source-separation}

The visible packets do not provide a counterexample for the actual prepared
source. There is a stronger source-specific reason than merely noting that
the packet was chosen separately: the original low-frequency source has an
exact finite interaction separation from it.

\begin{proposition}[Vanishing short-power pairings for the actual target]
\label{v11:short-powers}
On the balanced scalar-momentum family, for the actual prepared active
Riesz target $b_{A,h}$ and the packets in \cref{v11:quasimodes},
\begin{equation}
 \boxed{\ip{\nu_h}{\mathcal M_{2,h}^{\,j} b_{A,h}}_h=0
 \quad\text{if}\quad 2+2j<\ell-M.}
 \label{v11:short-power-zero}
\end{equation}
The same conclusion holds with $\mathcal M_{2,h}^{\,j}$ replaced by any
polynomial of degree satisfying this bound. There are order $h^{-2/3}$
consecutive vanishing pairings. This is not a conclusion about the inverse.
\end{proposition}
\begin{proof}
The complete correctly ordered target has Fourier support in $|k|_1\le2$:
the lapse part has type $Up$, and the shift part has types $U^2,p^2$, with
all derivatives and literal shifts preserving the support of their factors.
The first source changes a frequency by zero or one signed coordinate unit,
so its signed Gram changes it by an offset of $\ell^1$ size at most two.
The physical active projections are Fourier multipliers and add no offsets.
Induction therefore places $\mathcal M_{2,h}^{\,j}b_{A,h}$ in
$|k|_1\le2+2j$ as long as that ball does not reach a folded face. The stated
bound implies this for sufficiently fine grids, because $\ell\ll N$.
Both carrier boxes of the real packet have $|k_1|\ge\ell-M$.
The two supports are disjoint, and exact physical Parseval proves
\eqref{v11:short-power-zero}. Linearity proves the polynomial statement.
No approximation, projection of the action, or spectral-gap argument was
used.
\end{proof}

Thus the bad arbitrary-residual estimate and the original-source problem
must remain distinct. The actual source cannot reach these localized tests
through a bounded number of applications of the finite leading Hessian.
An inverse need not be approximable by such a polynomial with a uniformly
controlled remainder near a signed characteristic, however. Estimating the
longer interaction chains, or using an actual transport cancellation of the
prepared source, is additional mathematics; the finite separation is not
promoted to a small characteristic spectral mass.

\section{What this establishes for the original source, and what it does not}
\label{v11:implications}

The three conclusions have different logical roles. The direct transverse
part of the \emph{actual prepared target} has the proved geometric decay
\eqref{v11:direct-target}. The sharp Fourier-edge transfer is an exact
property of the actual lapse-to-shift operator, but its witness is not the
unknown actual lapse solution. The interior packets are likewise actual
finite multiplier tests of the original Hessian, but their small matrix
images are constructed residuals, not the prepared drift $b_{A,h}$.
It would be incorrect to substitute either witness for that drift and claim
a cofinal failure of the Einstein limit.

Conversely, a proposed bridge proof cannot now appeal to a uniform estimate
of the form \eqref{v11:false-estimate}, or to a uniformly elliptic scalar
lapse Schur form on this balanced family. The failure is not solely a
large-frequency aliasing effect: the signed characteristic is realized by
packets strictly below the Fourier faces and has a nonzero geometric output.
At the frozen symbol it is exactly the same trace fraction $2/3$ as in the
original Cartan analysis. This statement is not a proof that a particular
finite-grid trace-compression eigenvalue converges to $2/3$; passing from
localized signed quasimodes to such a normalized global spectral statement
would require additional control of the positive source normalization.

The source-dependent target can be written without a worst inverse:
\begin{equation}
 \boxed{\mathfrak E_hT_h^\dagger L_h
  \mathscr S_{L,h}^{-1}
       \bigl(-b_{N,h}+L_h^*T_hB_h^{-1}b_{\beta,h}\bigr).}
 \label{v11:remaining-composite}
\end{equation}
The first term of \eqref{v11:output-split} has already been controlled.
For \eqref{v11:remaining-composite}, the genuinely new research question is
whether the \emph{specific} prepared source has enough cancellation or
transport structure along the characteristic, while its solved lapse has
sufficiently small folded-face tails. The complete active inverse remains
inside this expression. Neither cancellation is assumed merely because the
un-inverted source has finite Fourier support.

No all-cutoff signed rank, cofinal phase-root selection, nonlinear source
bound, differentiated curvature estimate, or common physical lifespan is
established here. In particular the exact initial geometric criterion from
V10 is not a spacetime theorem. The successful stable-writer population and
its law-space neighborhoods remain separate inherited results.

\section{Verification and preservation}
\label{v11:verification}

The symbolic verifier constructs \eqref{v11:ell} from the original electric
first-source row, eliminates the transverse block in the positive Frobenius
coordinates, and checks the rational polynomial \eqref{v11:q}. It verifies
the isolating-interval signs, strict derivative bound, explicit nonzero
shift, the full signed null equation modulo the defining quartic, and the
trace-compression fraction $2/3$. These are exact rational/symbolic identities;
no rounded root is used in the proofs.

A separate Fourier implementation evaluates the full original first source
with every folded mode retained. At $N=3,5$ its values are compared against
the independently retained real-space coefficient program on random full
multiplier arrays. The adjoint is checked by the physical Parseval pairing.
The scalar-lapse elimination is checked against every row of the original
leading equation, not only a selected square subsystem. These floating
comparisons test the implementation; they do not supply new signed-rank
certificates. The actual-target direct output, the sharp folded-output
formula, and the constructed interior packets are evaluated separately and
labelled as diagnostics. Exact integer frequency-support checks verify the
short-power separation from the actual target. The packet inequalities under refinement are
proved by \cref{v11:packet-lemma,v11:quasimodes}, not by extrapolating those
finite values.

The new code does not solve $J_2$, regenerate the V8 interval root, integrate
an original-action spacetime, or estimate the unknown prepared source's
characteristic projections. Its comparison evaluator and source files are
fingerprinted. The cumulative manuscript preserves all ten preceding
marked bodies verbatim, including their explicit corrections and scope
statements. Compilation and hashes verify the artifact, not every inherited
mathematical assertion.
% END IMMUTABLE EXACT-ACTION BRIDGE V11 BODY

% Future iterations append here.
% BEGIN IMMUTABLE EXACT-ACTION BRIDGE V12 BODY
\clearpage
\renewcommand{\thesection}{\arabic{section}}
\renewcommand{\theHsection}{bridgeV12.\arabic{section}}
\setcounter{section}{90}
\section{Complementary momenta before estimating a signed inverse}
\label{v12:context}

The signed characteristic in \cref{v11:characteristic} belongs to the
circular-configuration, scalar-momentum first jets used there. It is not a
theorem about every initial record of the unchanged action. A second
possibility is to keep the linearly polarized configurations for which the
first-source trace vanishes, and cancel their weak drift with a different
momentum polarization. This body implements that possibility. It changes
initial data, not the action, phase derivative, multiplier continuation,
physical clock, or literal curvature observation.

The resulting family has two properties simultaneously. Its complete weak
quadratic drift is zero, and its entire leading active source is trace-free
with a positive Cartan Gram. After the rational amplitude selection already
proved in \cref{v2:rational-selection}, that active Gram is nonsingular at
every sufficiently fine odd cutoff, with a stated polynomial lower bound.
The transverse estimate now uses a spatially variable diagonal boost matrix;
its inverse is established below, rather than borrowed from the
scalar-momentum formula. The initial mean constraints are balanced by an
explicit homogeneous anisotropic momentum, and the original nonlinear
initializer supplies exact initial records on a common small-amplitude
interval. This is an initial-data interval, not a physical evolution interval.

A separate characteristic-zero calculation establishes full leading graded
signed rank at a three-site member of this new family. At that member the
original action consequently has local exact continuations and initial
literal curvature of order amplitude. The weak second source is not
trace-free, and no cofinal full signed-rank or curvature theorem is inferred
from active positivity.

\paragraph{Inputs and nonduplication.}
The configuration weak-drift formula is \cref{v4:drift}; the momentum
product-defect identity and scalar cross cancellation are
\cref{v9:kinetic-defect,v5:momentum-zero}. The scalar lapse-spectrum and
rational detuning results are \cref{v2:tensor-rank,v2:rational-selection}.
The original uniform initial constraint calculus, range inverse, and nonlinear
range solve are the retained \cite{OB25},
\src{obv17:C-calculus}, \src{obv17:linear-inverse}, and \src{obv17:range}.
They are not proved again as new results. The contribution is a different
configuration--momentum cancellation, its explicit anisotropic mean
calibration, the variable diagonal transverse inverse, and the simultaneous
application to the actual first source. The exact finite certificate uses
\cref{v3:compiler}, not a new approximate action. The quartic root and
quasimodes of V11 remain valid on their own, different first-jet family.

\section{Configuration--momentum cancellation of the complete weak drift}
\label{v12:neutrality}

Use the original odd grid and the folded symbols of \eqref{v4:fold}:
\[
 \omega_h=2h^{-1}\sin(\pi h),\qquad
 \Delta_h=\omega_h-\tfrac12\omega_{2,h}>0.
\]
Let $A_i,B_i,C_i,D_i$ be real symmetric transverse trace-free tensors for
axis $i$: their $i$th rows and columns are zero. Let $P_c$ be any constant
symmetric momentum, not necessarily scalar. Consider the first records
\begin{equation}
 \begin{split}
 U&=\sum_i\{A_i\cos(2\pi x_i)+B_i\sin(2\pi x_i)\},\\
 p&=P_c+Q,\qquad
 Q=\sum_i\{C_i\cos(2\pi x_i)+D_i\sin(2\pi x_i)\}.
 \end{split}
 \label{v12:general-data}
\end{equation}
Both nonconstant slots are phase-divergence-free and trace-free. Define
\[
 e_i^c(A,B)=\|A_i\|_F^2-\|B_i\|_F^2,\qquad
 e_i^s(A,B)=2\ip{A_i}{B_i}_F,
\]
and define $e_i^c(C,D),e_i^s(C,D)$ in the same way.

\begin{theorem}[An exact configuration--momentum weak-neutrality law]
\label{v12:matching}
For \eqref{v12:general-data}, the complete curl-free quadratic canonical
shift drift is represented by
\begin{equation}
 \boxed{
 (d_{2,h})_i=\Delta_h\left\{
   \left[e_i^c(C,D)-\frac{\omega_h^2}{4}e_i^c(A,B)\right]\sin(4\pi x_i)
  -\left[e_i^s(C,D)-\frac{\omega_h^2}{4}e_i^s(A,B)\right]\cos(4\pi x_i)
 \right\}.}
 \label{v12:matching-formula}
\end{equation}
This holds at every odd cutoff, with the folded second derivative at $N=3$.
In particular it is zero if and only if, for each axis,
\begin{equation}
 e_i^c(C,D)=\frac{\omega_h^2}{4}e_i^c(A,B),\qquad
 e_i^s(C,D)=\frac{\omega_h^2}{4}e_i^s(A,B).
 \label{v12:matching-conditions}
\end{equation}
The constant symmetric momentum $P_c$ does not enter these conditions.
The configuration and momentum contributions need not vanish separately.
\end{theorem}
\begin{proof}
The original weak quadratic drift separates into its configuration and
momentum quadratic terms; its formula is \eqref{v5:general-weak-B2}.
The configuration term at zero momentum is exactly \eqref{v4:drift-vector}.
That coefficient calculation retains the full tested cubic envelope and its
proved logarithmic cancellation.

Write $P_c=p_cI+Q_c$, with $Q_c$ constant and trace-free. The scalar
momentum term and its cross term with any other momentum vanish by
\cref{v5:momentum-zero}. The remaining momentum is the trace-free,
phase-divergence-free field $Q_c+Q$. On that space its full tested kinetic
contribution is the inherited vector \eqref{v9:kinetic-vector}:
\[
 \mathcal K_i=\sum_{j,k}\mathfrak d_i(Q_{jk},Q_{jk})
      -2\sum_{k,j}\mathfrak d_k(Q_{ij},Q_{kj}),\qquad
 \mathfrak d_k(f,g)=\delta_k(fg)-f\delta_kg-g\delta_kf.
\]
Adding $Q_c$ changes no term in this expression, because a product defect
with a constant factor is exactly zero. Cross products of different axial
profiles also have zero defect: in each differentiated coordinate at least
one of the factors is constant. For a self-pair on axis $i$, the second sum
is zero. Its only potentially nonzero derivative has $k=i$, but the $i$th
row of that profile is zero. The first sum therefore reduces to
\[
 \delta_i\|C_i\cos(2\pi x_i)+D_i\sin(2\pi x_i)\|_F^2
 -2\ip{C_i\cos(2\pi x_i)+D_i\sin(2\pi x_i)}
          {\delta_i(C_i\cos(2\pi x_i)+D_i\sin(2\pi x_i))}_F.
\]
Using the literal doubled trigonometric functions and their folded symbol,
this is
\[
 \Delta_h\{e_i^c(C,D)\sin(4\pi x_i)-e_i^s(C,D)\cos(4\pi x_i)\}.
\]
Adding the configuration term proves \eqref{v12:matching-formula}.
The represented vector has zero mean and is curl-free, so agreement on all
weak tests identifies the entire weak covector. At every odd $N\ge3$ the
doubled sine and cosine have mean square $1/2$ and zero mutual pairing.
Since $\Delta_h>0$, their independent coefficients give
\eqref{v12:matching-conditions}. No phase Leibniz rule or changed link term
has been used.
\end{proof}

Here is the particular realization used below. Let $T_i$ have unit entries
in the two transverse off-diagonal slots, and let $P_i$ have transverse
diagonal entries $+1,-1$, as in \cref{v6:balanced}. Then
$\|T_i\|_F^2=\|P_i\|_F^2=2$ and $\ip{T_i}{P_i}_F=0$. Put
\begin{equation}
 \boxed{
 U_\alpha=\sum_i\alpha_iT_i\cos(2\pi x_i),\qquad
 Q_{h,\alpha}=\frac{\omega_h}{2}
                    \sum_i\alpha_iP_i\cos(2\pi x_i).}
 \label{v12:complementary}
\end{equation}
The two fields use complementary transverse tensor polarizations at the
same phase. Neither is individually circularly polarized. The matching
conditions hold exactly, for arbitrary real $\alpha_i$ and every constant
symmetric $P_c$. When some $\alpha_i\ne0$, the configuration weak drift
alone is nonzero; the momentum term is its negative. These are spatial
initial profiles, not an assertion about traveling-wave solutions.

\section{An anisotropic mean calibration and exact initial records}
\label{v12:preparation}

Restrict $1/2\le\alpha_i\le3/2$ and write
$\Sigma_\alpha=\sum_i\alpha_i^2$. Set
\begin{equation}
 \begin{gathered}
 \tau_h=8\omega_h,\qquad
 s_{h,\alpha}=\omega_h\left(4-\frac{\Sigma_\alpha}{64}\right),\\
 B_c=\diag(s_{h,\alpha},-\tau_h,-\tau_h),\qquad
 P_c=B_c-(\tr B_c)I
     =\omega_h\diag\left(16,4+\frac{\Sigma_\alpha}{64},
                                 4+\frac{\Sigma_\alpha}{64}\right),\\
 Y_{h,\alpha}=(U_\alpha,p_{h,\alpha}),\qquad
 p_{h,\alpha}=P_c+Q_{h,\alpha}.
 \end{gathered}
 \label{v12:family}
\end{equation}
The constant tensor is anisotropic. Its numerical coefficients are bounded
independently of $h$; the actual canonical record is $aY_{h,\alpha}+O(a^2)$
and is small when the preparation amplitude $a$ is small. No background
finite metric is changed to accommodate these coefficients.

\begin{lemma}[The original quadratic means have a uniform simple zero]
\label{v12:means}
The first records \eqref{v12:family} satisfy every linear initial constraint
and all four quadratic mean compatibility equations. To calculate the
latter, replace $s_{h,\alpha}$ temporarily by $s$ in $P_c$ and add momenta
$\sum_i z_i\alpha_iT_i\sin(2\pi x_i)$. In the original sign convention for
$\mathscr C_h$, the complete quadratic mean map is
\begin{equation}
 \boxed{
 Q_h(s,z)=\left(
 2\tau_h^2-4s\tau_h-\frac{\omega_h^2}{2}\Sigma_\alpha
                         -\sum_i\alpha_i^2z_i^2,
                 \;(-\omega_h\alpha_i^2z_i)_{i=1}^3\right).}
 \label{v12:means-formula}
\end{equation}
Its value at $(s,z)=(s_{h,\alpha},0)$ is zero, and its derivative there is
\[
 \diag(-32\omega_h,-\omega_h\alpha_1^2,
                      -\omega_h\alpha_2^2,-\omega_h\alpha_3^2).
\]
This derivative has a cutoff-independent inverse bound on the stated
amplitude box.
\end{lemma}
\begin{proof}
Each nonconstant tensor is transverse to the coordinate on which it depends;
its divergence is zero. The configuration is also trace-free. The constant
momentum has zero divergence. These facts verify the full linear initial
map retained in \src{obv17:linear-C}.

Use the original constant-multiplier quadratic Hamiltonian, equivalently
\eqref{v5:H2} on these transverse records. For a symmetric matrix
$P_c=B_c-(\tr B_c)I$,
\[
 \|P_c\|_F^2-\tfrac12(\tr P_c)^2
      =\|B_c\|_F^2-(\tr B_c)^2=4s\tau_h-2\tau_h^2.
\]
Different axial modes have zero spatial cross averages. Their squared
configuration norms are $\alpha_i^2$, their squared momentum norms in
\eqref{v12:complementary} are $\omega_h^2\alpha_i^2/4$, and their spatial
Hamiltonian contribution is also $\omega_h^2\alpha_i^2/4$.
The added sine momenta contribute $\alpha_i^2z_i^2$ to the kinetic mean.
They are off-diagonal and hence have zero cross pairing with the diagonal
constant and oscillatory momenta. Constant-shift quadratic rows are exactly
$\ip p{\delta_iU}_h$, and their added sine contribution is
$-\omega_h\alpha_i^2z_i$. All other such pairings vanish.
The scalar constraint is the negative lapse derivative of the Hamiltonian;
these computations give \eqref{v12:means-formula} with its signs.
Cubic link terms do not enter a quadratic field coefficient with a spatially
constant multiplier. Substitution of $\tau_h=8\omega_h$ and
$s_{h,\alpha}$ proves the zero. The derivative and its inverse bound follow
from $\alpha_i\ge1/2$ and $3\sqrt3\le\omega_h<2\pi$.
\end{proof}

\begin{theorem}[Exact nonlinear preparation of the complementary family]
\label{v12:exact-preparation}
For every fixed retained integer Sobolev index $r\ge3$, there are
$a_r,h_r,C_r>0$, independent of the cutoff and of
$\alpha\in[1/2,3/2]^3$, such that, for $h<h_r$ and $|a|<a_r$, the original
initial constraint map has analytic preparations
\begin{equation}
 \boxed{
 \mathscr C_h(X_{h,\alpha}(a))=0,\qquad
 X_{h,\alpha}(a)=aY_{h,\alpha}+O_{\mathcal P_h^r}(a^2),\qquad
 \|X_{h,\alpha}(a)\|_{\mathcal P_h^r}\le C_r|a|.}
 \label{v12:exact-initial}
\end{equation}
Here $\mathcal P_h^r=H_h^{r+2}(\operatorname{Sym}_3)\times
H_h^{r+1}(\operatorname{Sym}_3)$, as in the retained initializer. The
remainder bound is uniform on this initial family. At any fixed additional
odd cutoff, including $N=3$, the same construction holds on its finite
analytic stationary chart, with cutoff-dependent constants if needed.
This theorem concerns initial records, not their exact-action physical
lifespan.
\end{theorem}
\begin{proof}
All first fields and the four balancing directions used in
\cref{v12:means} belong to the kernel of the linear initial constraint map.
Their indicated Sobolev norms are uniformly bounded, since only constants
and unit axial frequencies occur. Apply the retained nonlinear range solve
\src{obv17:range} to $z=aY(s,z_1,z_2,z_3)$, writing its exactly range-solved
record as $z+R_hy_h(z)$. The correction is $O(\|z\|^2)$ uniformly. Its
remaining mean map is analytic, has zero constant and linear coefficients,
and has quadratic coefficient \eqref{v12:means-formula}.

Divide its value at $aY$ by $a^2$, using the analytic integral Taylor
extension at $a=0$. The resulting four-vector is
$Q_h(s,z)+O_r(a)$ in $C^1$ on a fixed neighborhood of the root. Its derivative
at $a=0$ has the uniform inverse of \cref{v12:means}. A preconditioned
contraction, first choosing the balancing neighborhood and then the common
$a_r$, solves these four exact mean equations. It changes $s$ and the sine
coefficients by $O_r(a)$, so it changes no first field coefficient. The
range equations were already zero; all original initial rows now vanish.
The same uniform derivative inverse gives analytic parameter dependence and
the bounds in \eqref{v12:exact-initial}. The retained chart guarantees
pointwise positive $I+u$ after decreasing $a_r$.
At a fixed cutoff the identical finite-dimensional analytic argument uses
the original flat stationary inverse and linear range inverse. This explains
the separate assertion without promoting a local physical continuation
interval to a uniform one. No initial constraint or its mean is removed.
\end{proof}

\section{The original first source on diagonal boost records}
\label{v12:first-source}

The first spatial connection is the retained linear stationary map
\[
 r_{ia}=-\tfrac12\epsilon_{abc}\delta_bU_{ic},\qquad
 B=p-\tfrac12(\tr p)I.
\]
For \eqref{v12:family}, both are diagonal: write $r=\diag(r_i)$ and
$B=\diag(b_i)$. More explicitly, with $f_i=\cos(2\pi x_i)$,
\begin{equation}
 r_1=\tfrac12(\alpha_3\delta_3f_3-\alpha_2\delta_2f_2),\quad
 r_2=\tfrac12(\alpha_1\delta_1f_1-\alpha_3\delta_3f_3),\quad
 r_3=\tfrac12(\alpha_2\delta_2f_2-\alpha_1\delta_1f_1).
 \label{v12:rotation-diagonal}
\end{equation}
Also $b_i=(B_c)_{ii}+(Q_{h,\alpha})_{ii}$, and
\begin{equation}
 S_i b_i=b_i,\qquad
 b_1>2\omega_h,\quad b_2,b_3<-2\omega_h,\qquad
 |b_i|<10\omega_h.
 \label{v12:boost-margins}
\end{equation}
The shift identity follows because the $i$th diagonal momentum entry has
no $i$-axis oscillation. The inequalities follow from
$|(Q_{h,\alpha})_{ii}|\le3\omega_h/2$ and
$3/4\le\Sigma_\alpha\le27/4$; for example
$b_1/\omega_h\ge4-27/256-3/2>2$.

\begin{lemma}[Complete first-source formula for the new family]
\label{v12:J1}
For a multiplier $(n,\beta)$ set $w=-\curl_\delta\beta/2$.
The actual first source $J_{1,h}(n,\beta)$ has diagonal
$\mathscr D_{h,\alpha}n$ from \eqref{v2:diagonal-source}, and off-diagonal
entries
\begin{equation}
 \boxed{
 (J_{1,h}(n,\beta))_{ij}
 =-\frac h2\left(
 r_iS_i\delta_jn+r_jS_j\delta_in+
 b_iS_iw_j+b_jS_jw_i\right),\qquad i\ne j.}
 \label{v12:J1-off}
\end{equation}
In particular its entire range is trace-free. Its pure-lapse part is
independent of the diagonal momentum, while its pure-shift part has zero
diagonal. These are identities of the original first coefficient, including
the electric logarithmic contribution.
\end{lemma}
\begin{proof}
Use the complete first differential
$Z_0^*(G_1+T_1^*\eta_0)$ from \cref{v1:first-source,v3:universal-flat}.
The geometric lapse coefficient is the symmetric alternating contraction of
the original product defects in $n$ and $U$. Substituting $U_\alpha$ gives
exactly \eqref{v2:diagonal-source}; its geometric off-diagonal entries are
zero. This substitution involves no momentum.

The non-electric pure-shift terms have zero symmetric rotational projection
also for a general symmetric boost matrix. The only nontrivial pointwise
Cartan contraction is
\[
 \sum_{i,j}\beta_j e_i\cdot
                    (-s_i\mathbin\times B_j+s_j\mathbin\times B_i)=0
\]
for symmetric matrices with rows $s_i,B_i$. The first term vanishes by
symmetry of $s_{ia}$ and the alternating tensor in $i,a$; the second by
symmetry of $B_{ib}$ and alternation in $i,b$. This identity does not require
$\beta$ to be curl-free. The remaining geometric shift loads and the
velocity correction are boost-valued and have zero rotational projection.
For a lapse, a boost variation cannot pair with its rotational magnetic
coefficient, so only the rotational first connection contributes there.

It remains to contract the original cubic electric term. For a symmetric
rotational test with row $s_i$, diagonal $r_i e_i$ and $b_i e_i$, and flat
temporal components $p^n=\delta n$ and $w$, its vector triple products give
\[
 \begin{aligned}
 e_i\cdot\bigl\{&s_i\times(r_i e_i\times S_ip^n)
       +r_i e_i\times(s_i\times S_ip^n)\\
       &+s_i\times(b_i e_i\times S_iw)
       +b_i e_i\times(s_i\times S_iw)\bigr\}
 =\sum_{j\ne i}s_{ij}(r_iS_i\delta_jn+b_iS_iw_j).
 \end{aligned}
\]
Multiply by the original factor $-h$ and symmetrize in the Frobenius pairing.
This gives \eqref{v12:J1-off} and zero electric diagonal. The magnetic cubic
term begins at second source order, as in the original degree audit; higher
logarithmic terms cannot enter $J_1$. These observations account for every
term at this coefficient. The geometric diagonal sums to zero, and all other
diagonal terms are zero, proving trace-freeness.
\end{proof}

\section{Uniform inversion of the variable diagonal transverse block}
\label{v12:transverse}

Let $T_h^{\rm new}\beta=J_{1,h}(0,\beta)$, restricted to mean-zero
phase-transverse shifts. This notation distinguishes it from the
constant-momentum Fourier multiplier in V11. Introduce the off-diagonal map
\[
 (\mathcal L_b w)_{ij}=b_iS_iw_j+b_jS_jw_i,\qquad i<j.
\]
The three output entries use the Euclidean norm before Frobenius weights
are restored.

\begin{theorem}[An all-frequency positive transverse estimate]
\label{v12:T-inverse}
For every odd cutoff and every first jet in \eqref{v12:family},
$\mathcal L_b$ is invertible on arbitrary three-component arrays. Its
inverse is a conjugate of a pointwise $3\times3$ inverse. On transverse
shifts it gives
\begin{equation}
 \boxed{
 \frac{h^2\omega_h^2}{5000}\|\Omega_h\beta\|_h^2
 \le\|T_h^{\rm new}\beta\|_h^2
 \le50h^2\omega_h^2\|\Omega_h\beta\|_h^2.}
 \label{v12:T-bounds}
\end{equation}
Its rank is exactly $2(V-1)$. Its kernel on all shifts is precisely the
curl-free shifts, and its Cartan Gram is
$(T_h^{\rm new})^*K_0^{-1}T_h^{\rm new}
=\tfrac12(T_h^{\rm new})^*T_h^{\rm new}$.
\end{theorem}
\begin{proof}
Set $t_i=S_i^{-1}w_i$ and replace the $ij$ output by
$S_i^{-1}S_j^{-1}(\mathcal L_b w)_{ij}$. Both are isometries of their
respective positive array spaces. Since $S_i b_i=b_i$, the transformed
map is pointwise multiplication by
\begin{equation}
 M_b(x)=\begin{pmatrix}
 S_1^{-1}b_2&S_2^{-1}b_1&0\\
 S_1^{-1}b_3&0&S_3^{-1}b_1\\
 0&S_2^{-1}b_3&S_3^{-1}b_2
 \end{pmatrix}.
 \label{v12:pointwise-M}
\end{equation}
Write $a_{ij}=S_j^{-1}b_i$. Its determinant is
\[
 \det M_b=-a_{21}a_{13}a_{32}-a_{12}a_{31}a_{23}.
\]
Each product contains one positive first diagonal boost and two negative
ones, so the two terms have the same sign. By \eqref{v12:boost-margins},
$|\det M_b|\ge16\omega_h^3$. Each row and each column has at most two
entries of magnitude at most $10\omega_h$, whence
$\|M_b\|_2\le20\omega_h$. The smallest singular value is at least the
absolute determinant divided by $\|M_b\|_2^2$, giving
$\sigma_{\min}(M_b)\ge\omega_h/25$.
Thus, for arbitrary $w$,
\[
 \frac{\omega_h}{25}\|w\|_h
       \le\|\mathcal L_b w\|_h\le20\omega_h\|w\|_h.
\]
This is an exact shifted-array argument; it includes all folded modes.

Equation \eqref{v12:J1-off} gives $T_h^{\rm new}\beta=-h\mathcal L_b w/2$
in off-diagonal entries. Their Frobenius norm squared is twice the squared
three-entry norm. For $w=-\curl_\delta\beta/2$ and transverse $\beta$,
$\|w\|_h^2=\|\Omega_h\beta\|_h^2/4$ by Parseval. Combining these facts
gives \eqref{v12:T-bounds}. The curl maps the transverse space bijectively
to itself on every nonzero odd-grid mode, so the rank and full-shift kernel
follow. The range has zero diagonal, which proves the stated Cartan identity.
\end{proof}

This inverse is an analytical source operation, not an asserted local rule
for the multiplier evolution. On the range of $T_h^{\rm new}$ it recovers
$w$ using \eqref{v12:pointwise-M}, then
$\beta=-2\Omega_h^{-2}\curl_\delta w$. On arbitrary off-diagonal outputs
this procedure is not being identified with the positive least-squares
pseudoinverse. Also, unlike V11's scalar-momentum operator,
$T_h^{\rm new}$ is not a Fourier multiplier. Its variable coefficients must
be retained in any later source-dependent estimate.

\section{Simultaneous cofinal active regularity and weak neutrality}
\label{v12:cofinal-active}

Now choose the same \emph{configuration} detuning supplied by the inherited
\cref{v2:rational-selection}:
\[
 \alpha_h=(1,\theta_h,1),\qquad
 1<\theta_h<1+h,\qquad
 \|\mathscr D_{h,\alpha_h}n\|_h\ge ch^6\|n\|_h
 \quad(\overline n=0).
\]
The new momentum is always recalibrated by \eqref{v12:family}; it is not the
scalar momentum of V2. Denote the full active restriction of its first source
by $\mathcal A_h^{\rm new}$.

\begin{theorem}[A weak-neutral family with a positive regular active source]
\label{v12:main}
At every sufficiently fine odd cutoff, the actual first jets
$Y_{h,\alpha_h}$ satisfy
\begin{equation}
 \begin{gathered}
 d_{2,h}(Y_{h,\alpha_h})=0,\qquad
 \rank\mathcal A_h^{\rm new}=3V-3,\qquad
 P_{\rm tr}\mathcal A_h^{\rm new}=0,\\
 \boxed{
 (\mathcal A_h^{\rm new})^*K_0^{-1}\mathcal A_h^{\rm new}
 =\tfrac12(\mathcal A_h^{\rm new})^*\mathcal A_h^{\rm new}
                         \succeq c h^{16}I.}
 \end{gathered}
 \label{v12:main-eq}
\end{equation}
The constant $c>0$ is independent of the cutoff and the selected rational
candidate. The complete first-source kernel is exactly constant lapse plus
curl-free shifts. Quantitatively, for $y=\mathcal A_h^{\rm new}(n,\beta)$,
\[
 \|n\|_h\le Ch^{-6}\|y\|_h,\qquad
 \|\Omega_h\beta\|_h\le Ch^{-8}\|y\|_h.
\]
These powers are sufficient, not claimed optimal. The same jets have the
exact original initial preparations of \cref{v12:exact-preparation}.
\end{theorem}
\begin{proof}
Weak neutrality follows from \cref{v12:matching} with
$A_i=\alpha_iT_i$, $B_i=0$, $C_i=\omega_h\alpha_iP_i/2$, and $D_i=0$.
By \cref{v12:J1} the source diagonal has not changed from the scalar
operator to which the inherited detuning theorem applies. It bounds the
mean-zero lapse by $Ch^{-6}\|y\|$. A constant lapse has zero first source:
both the product defect with that constant and every $\delta_jn$ vanish.
The remaining pure-shift source has exactly the kernel and rank established
in \cref{v12:T-inverse}. This proves the full kernel and active rank.

The complete lapse operator has $L_h^2\to L_h^2$ norm at most $Ch^{-1}$,
by its finite sum of original product defects and shifted derivatives with
bounded unit-profile coefficients. This is the elementary bound used in
V2; its lapse formula is unchanged. Hence
\[
 \|T_h^{\rm new}\beta\|_h
    \le\|y\|_h+Ch^{-1}\|n\|_h\le Ch^{-7}\|y\|_h.
\]
The lower estimate in \eqref{v12:T-bounds}, and the uniform positive lower
bound on $\omega_h$, give the asserted $h^{-8}$ estimate. The nonzero-mode
Poincare bound for $\Omega_h$ controls $\|\beta\|_h$ as well. Thus the
smallest active source singular value is at least $ch^8$.
Trace-freeness was proved for the whole source in \cref{v12:J1}; on that
range $K_0^{-1}=I/2$. Squaring the singular-value estimate proves
\eqref{v12:main-eq}. The original preparation theorem was proved independently
of this rank calculation and applies to the uniformly bounded selected
first jets.
\end{proof}

\begin{corollary}[No signed trace amplification in the leading active solve]
\label{v12:no-active-resonance}
Every orthonormal basis $Q_A$ of the new active source range satisfies
$Q_A^*P_{\rm tr}Q_A=0$. Thus the active range has no trace fraction $2/3$.
For any active target $b$, its original signed response is exactly the positive
minimum-norm reproduction vector
\[
 K_0^{-1}\mathcal A_h^{\rm new}
 [ (\mathcal A_h^{\rm new})^*K_0^{-1}\mathcal A_h^{\rm new}]^{-1}b
 =\mathcal A_h^{\rm new}
 [ (\mathcal A_h^{\rm new})^*\mathcal A_h^{\rm new}]^{-1}b.
\]
This equality follows from the source itself; it does not replace a signed
inverse by a positive inverse in a problem where they differ.
\end{corollary}
\begin{proof}
Use $P_{\rm tr}\mathcal A_h^{\rm new}=0$ and
$K_0^{-1}\mathcal A_h^{\rm new}=\mathcal A_h^{\rm new}/2$.
The characterization as a minimum-norm vector is the usual orthogonal
source decomposition, already established in the inherited Cartan analysis.
\end{proof}

The scalar detuning is therefore reused for a genuinely different purpose:
V2's positive active family had a nonzero weak quadratic target, while the
circular scalar-momentum family made that target zero but introduced the
signed active characteristic analyzed in V11. Here both desired properties
hold at once. This does not give a uniform elliptic estimate. Positive
source singular values can still become small, and the bound $ch^{16}$
does not prove decay of the weighted geometric output.

The first jets also have a fixed refinement limit. The rational choice
$\theta_h=1+O(h)$ and $\omega_h=2\pi+O(h^2)$ give, in every fixed Sobolev
norm on their finite Fourier support,
\[
 \mathcal I_hY_{h,\alpha_h}=Y_*+O(h),
\]
where
\[
 U_*=\sum_iT_i\cos(2\pi x_i),\qquad
 p_*=2\pi\diag(16,4+3/64,4+3/64)
                         +\pi\sum_iP_i\cos(2\pi x_i).
\]
This is a first-data consistency statement, not a spacetime convergence claim.

\section{A complete three-site signed certificate on the new family}
\label{v12:certificate}

Fix $N=3$ and $\alpha=(1,9/8,1)$. Then $\Sigma_\alpha=209/64$,
$\omega_h=3\sqrt3$, and
\begin{equation}
 P_c=3\sqrt3\,\diag(16,16593/4096,16593/4096),\qquad
 Q=\frac{3\sqrt3}{2}\sum_i\alpha_iP_i\cos(2\pi x_i).
 \label{v12:three-seed}
\end{equation}
These exact data, with $U_\alpha$ as above, are in $\mathbb Q(\sqrt3)$.
All spatial tests retain $\delta_i=3(S_i-S_i^{-1})$.
Use the original component-major multiplier order and the same 78 active
columns listed in \eqref{v3:pivots}. Let $C$ be the resulting $162\times78$
first-source matrix, and let $W$ be the $162\times29$ full weak second-source
matrix given by \cref{v3:compiler}, on the three constant shifts followed by
the 26 gradients of point differences. Put
\[
 A=C^TK^{-1}C,\quad L=C^TK^{-1}W,\quad
 D=W^TK^{-1}W-L^TA^{-1}L,\quad
 \Gamma=(C,W)^TK^{-1}(C,W),
\]
with the original raw dual Cartan matrix \eqref{v3:raw-K}. No weak column
from the circular or scalar-momentum families is reused.

\begin{theorem}[Regular full leading source at a complementary-polarization seed]
\label{v12:finite-rank}
At \eqref{v12:three-seed},
\[
 \rank C=78,\qquad \rank(C,W)=107,\qquad
 \det\Gamma\ne0,\qquad \det D\ne0.
\]
With $\sqrt3\mapsto56\in\mathbb F_{241}$, the exact residues are
\begin{equation}
 \boxed{\det A\equiv131,\quad \det D\equiv208,\quad
              \det\Gamma\equiv15,\quad \det D_{cc}\equiv133.}
 \label{v12:certificate-values}
\end{equation}
Here $D_{cc}$ is the three-dimensional constant--constant weak Schur block.
These are characteristic-zero nonvanishing assertions at this fixed cutoff.
They do not bound the real inverse norms or establish the same weak rank
at other cutoffs.
\end{theorem}
\begin{proof}
The companion program first evaluates all 108 original first-source columns
from \eqref{v12:J1-off} and the retained diagonal formula. It independently
regenerates all 108 columns by formal differentiation of the original scalar
coefficient functional \eqref{v3:functional}, retaining the auxiliary flat
response. The two complete arrays agree in the stated residue field.
All diagonal traces are zero. Every curl-free shift column and the constant
lapse column is checked against the full source, not a selected subsystem.
The complete prepared drift, in the corrected component-major convention,
has zero pairing with all 29 weak basis tests, consistently with the
all-cutoff proof in \cref{v12:matching}.

For the higher source the program recomputes the normal correction and first
auxiliary response in \eqref{v3:compiler-formula} on this first jet, then
regenerates all 29 weak columns, including the required cubic electric and
magnetic terms. All coefficients and denominators are in
$\mathbb Q(\sqrt3)$ and are regular at the displayed specialization.
The implementation equivalently uses $\zeta_{12}\mapsto237$, which gives
$\omega_h\mapsto168=3\cdot56$. Exact elimination gives
\eqref{v12:certificate-values}. The separate Schur determinant check is
$131\cdot208\equiv15\pmod{241}$. Nonzero residues prove nonzero determinants
in the coefficient field and hence in its physical real embedding. A
nonsingular signed Gram implies independence of its source columns.
The original first-source null directions give the upper active rank 78,
which the certificate attains. Full matrices, column labels and regeneration
code are supplied with the note.
\end{proof}

The higher source is not trace-free. For example, the trace of its first
constant-shift column at the origin has residue $17\pmod{241}$ in the same
raw diagonal convention, and is therefore nonzero in characteristic zero.
The other two constant columns have trace residues $78,102$ there. This
explicitly prevents extending \cref{v12:no-active-resonance} from the active
block to the full graded source without additional mathematics.

\section{The exact-action branch supplied by the finite certificate}
\label{v12:physical}

\begin{corollary}[A nonempty exact branch with a neutral leading physical target]
\label{v12:local-branch}
At the three-site first jet \eqref{v12:three-seed}, the exact initial
preparations have a normalized local original-action continuation for every
sufficiently small nonzero amplitude. On each such initial curve the original
literal temporal and spatial link curvatures obey
\[
 \|E_h(a,0)\|_h+\|F_h^{\mathrm{sp}}(a,0)\|_h\le C_h|a|.
\]
The same conclusions persist in a sufficiently small parameter neighborhood
inside the complementary family at this fixed cutoff. No bound on the
initial weak multiplier velocity, no phase-selection equation, and no common
physical time interval is asserted.
\end{corollary}
\begin{proof}
The original finite analytic initializer applies by
\cref{v12:exact-preparation}. Its graded source has the invertible leading
Gram of \cref{v12:finite-rank}, which persists at sufficiently small nonzero
amplitude. The original homogeneity leaves precisely the normalized
constant-lapse null direction; the retained exact-action local continuation
theorem applies. It is the original continuation, not a stable-writer
replacement.

In the inherited active/weak amplitude grading, the leading normalized
source target is $(0,d_{2,h})$. The latter is zero by
\cref{v12:matching}, so the true normalized target is $O_h(a)$. The inverse
of its true normalized signed Gram is bounded at this fixed cutoff and
first jet. The physical tangent response is consequently $O_h(a)$; the
remaining analytic stationary and canonical contributions to the initial
literal electric read have the same size. This is the original response
argument of \cref{v4:initial-electric}, now with both hypotheses established
on the new first jet. The spatial stationary connection is $O_h(a)$ and
its original spatial link read is $O_h(a)$. No positive surrogate is
substituted for the higher signed response. Continuity of the nonzero finite
determinants proves the parameter-neighborhood assertion. The amplitude
weights in the multiplier rate are stronger than those in this physical
response, so these estimates do not bound that rate or its derivatives.
\end{proof}

On the cofinally active-regular family, the exact analytic elimination down
to the weak source is still available. Let
\[
 \mathcal W_h^{\rm new}=\mathcal Q_h(Y_{h,\alpha_h};\cdot)
                   \big|_{\ker\curl_\delta},\qquad
 \Pi_A=\mathcal A_h^{\rm new}
       [ (\mathcal A_h^{\rm new})^*\mathcal A_h^{\rm new}]^{-1}
                         (\mathcal A_h^{\rm new})^*,\qquad
 \overline W_h=(I-\Pi_A)\mathcal W_h^{\rm new}.
\]
The leading signed Gram is congruent to
\begin{equation}
 \diag\left(\tfrac12(\mathcal A_h^{\rm new})^*\mathcal A_h^{\rm new},
                       \overline W_h^*K_0^{-1}\overline W_h\right).
 \label{v12:remaining-weak}
\end{equation}
This is the inherited exact block identity, applied only after its new
trace-free active hypotheses have been proved. The second block has $V+2$
columns and retains every actual weak second-source term. The zero leading
weak target prevents a leading response on a regular branch, but the next
normalized drift and the complete nonlinear Hessian need not share that
cancellation.

\section{What this changes in the programme and what remains}
\label{v12:status}

This upstream family separates an avoidable signed obstruction from the
unresolved positive and higher-source estimates. The interior $2/3$ active
characteristic established for circular configurations and scalar momentum
is not forced by the finite action on all weak-neutral data. Complementary
configuration and momentum polarizations, together with the explicit
homogeneous anisotropy, give exact weak neutrality and a trace-free,
cofinally regular leading active range at once. Both features are properties
of the unchanged coefficients, not conditions imposed on an unknown inverse.

The construction does not make the exact-action bridge unconditional.
The active lower bound deteriorates polynomially and has not been improved
to the source-dependent weighted estimate required by V10. In particular the
variable diagonal transverse block is not the Fourier multiplier used to
prove V11's direct $O(h)$ estimate, so that estimate cannot simply be copied
to this family. The full weak block in \eqref{v12:remaining-weak} has been
certified only at $N=3$. Its higher source has nonzero trace; the signed
response at subsequent amplitude orders still needs its actual source
weights. No complete next-drift cancellation, cofinal weak rank, nonlinear
rank propagation, differentiated observation estimate, or cutoff-independent
physical lifespan has been proved here. Initial preparation and leading
active regularity are different statements from those evolution properties.

The next action-specific calculation is therefore the actual
$\overline W_h$ and the next normalized drift for these complementary
first jets, together with a positive-source estimate for the new active
lapse response. The verified three-site member supplies a genuine regular
starting branch, rather than a formal zero-source example with an unknown
finite continuation. It does not supply a generic microscopic attraction
mechanism or an open basin of unrestricted first data.

\paragraph{Verification and source scope.}
The proof package contains the new source formula, the unchanged inherited
coefficient evaluator and flat maps, all first and weak columns on the
three-site member, the complete prepared drift, signed matrices and selected
columns. All 108 first columns are cross-checked by independent formal
coefficient differentiation. The active trace identity, the positive active
Gram in physical dual coordinates, all first-source null directions, the 29
weak-drift pairings and the four determinant residues have exact modular
checks. Nonzero modular determinants are used only for characteristic-zero
nonvanishing; the zero identities rest on the written polarization and
source proofs. The finite arithmetic uses reduced residues; its largest
matrix contraction has dimension 162 and prime 241, keeping unreduced
integer products and sums well below exact binary64 integer capacity.
An instrumented first-source pass additionally records forward and adjoint
integer bounds.

Symbolic checks verify the full folded polarization identity, the constant
momentum mean balance, the shifted determinant, strict boost sign margins,
and the original electric triple-product contraction. Separate floating
comparisons at $N=3,5,9,13$ test the complete source, weak drift and pointwise
shift inverse; they are diagnostics, not uniform inverse certificates.
The all-cutoff active theorem explicitly inherits V2's scalar-spectrum and
rational-selection proofs instead of presenting them as new discoveries.
The initial range-and-mean calculus is inherited from the original-action
notes, with the new mean map proved here. Its role should not be confused
with the familiar continuum issue that linearized compact initial constraints
need not be integrable without their nonlinear compatibility conditions
\cite{v4:BD}. No nonlinear spacetime integration or proof-assistant
verification is represented as performed. All eleven earlier research bodies
remain byte-for-byte unchanged.
% END IMMUTABLE EXACT-ACTION BRIDGE V12 BODY


% Future iterations append here.
% BEGIN IMMUTABLE EXACT-ACTION BRIDGE V13 BODY
\clearpage
\renewcommand{\thesection}{\arabic{section}}
\renewcommand{\theHsection}{bridgeV13.\arabic{section}}
\setcounter{section}{99}
\section{The true transverse inverse before the remaining lapse solve}
\label{v13:context}

The complementary first jets of \cref{v12:family} remove the signed trace
characteristic from the leading active source. They do not make its transverse
block a Fourier multiplier. The pointwise inverse of the shifted three-component
map in \cref{v12:T-inverse} is not, on arbitrary outputs, the least-squares
inverse on phase-transverse shifts. This distinction matters both for the
original Schur equation and for its geometric output. We resolve this inverse
rather than assuming that the constant-coefficient result of V11 still applies.

The result has three consequences for the actual new family. After its explicit
curl scale is factored out, the transverse normal matrix has a uniform positive
spectral interval and Fourier bandwidth two. Its inverse therefore has
cutoff-independent exponential off-diagonal decay in cyclic Fourier distance,
and uniform Sobolev bounds of every fixed order. These statements give an
$O(h)$ bound for the direct transverse curvature output of the original
prepared source. The diagonal geometric lapse source is annihilated exactly by
this pseudoinverse; consequently the additional folded-face transfer loss of
V11 is absent from the complementary family's lapse-to-shift observation.
Finally the full active lower bound improves from $ch^{16}$ to $ch^{12}$ in
the unweighted physical multiplier norm.

The remaining scalar lapse inverse and the higher weak source are not replaced
by this normalized transverse matrix. Their source-dependent estimates and the
nonlinear physical lifespan remain separate. None of the inverse formulas
below changes the action, the phase derivative, its cyclic product, the initial
family, or a literal curvature observation.

\paragraph{Inputs and novelty boundary.}
The original first source is \cref{v12:J1}, the exact initial preparation is
\cref{v12:exact-preparation}, and the diagonal lapse detuning is the inherited
\cref{v2:rational-selection}. The geometric comparison operator is
\cref{v10:symbol}. We use the actual prepared drift formula
\eqref{v5:prepared-drift-formula}, not a residual fitted to an evolution scheme.
Exponential inverse decay from polynomial approximation is a classical method
\cite{v13:DMS}; the argument is included to specify its constants and cyclic
frequency geometry. The new application is to the true variable-coefficient
transverse source and to its actual targets, together with the observation
cancellation and improved full-active estimate. These are not new claims about
the unchanged higher weak matrices or their previously certified ranks.

\section{A uniformly conditioned normal operator on all Fourier modes}
\label{v13:normal}

Retain $\alpha\in[1/2,3/2]^3$ and the first data of \cref{v12:family}, at
any odd $N\ge3$, $h=N^{-1}$. The diagonal boost entries are $b_i$ and the
diagonal rotational entries are $r_i$ of \cref{v12:first-source}. Let
$\mathcal V_h$ be all three-component arrays with positive spatial-mean pairing,
and let $P_T$ be the orthogonal projection onto its mean-zero phase-transverse
subspace $\mathcal T_h$. All adjoints in this body use these physical pairings.
Set
\[
 \Omega_h=(-\textstyle\sum_i\delta_{i,h}^2)^{1/2},\qquad
 C_h=\curl_\delta,\qquad R_h=C_h\Omega_h^{-1},\qquad Q_T=I-P_T,
\]
where $\Omega_h^{-1}$ is zero on constants. Then
\begin{equation}
 C_h^*=C_h,\quad C_h^2=\Omega_h^2P_T,\quad R_h^*=R_h,
 \quad R_h^2=P_T.                                      \label{v13:curl}
\end{equation}
These identities follow at each nonzero mode from the cross-product symbol
$\mathrm i\kappa_h(k)\times$. They include the folded phase symbols. On
$\mathcal T_h$, $R_h$ is an isometry, not a derivative with an unbounded norm.

Write $\widehat{\mathcal L}_h=\omega_h^{-1}\mathcal L_b$, where
\[
 (\mathcal L_b w)_{ij}=b_iS_iw_j+b_jS_jw_i,\qquad (ij)=(12,13,23).
\]
Both sides here have the ordinary three-component positive pairing. Let
$\iota_o:\mathcal V_h\to\operatorname{Sym}_3$ be the isometric inclusion
with off-diagonal entries $z_{ij}/\sqrt2$ and zero diagonal. In particular,
$\iota_o^*Y=\sqrt2(Y_{12},Y_{13},Y_{23})$. The original pure-shift source,
now denoted simply $T_h$, factors as
\begin{equation}
 \boxed{T_h=\frac{h\omega_h}{2\sqrt2}\,
          \iota_o\widehat{\mathcal L}_hR_h\Omega_h
              \quad\hbox{on }\mathcal T_h.}            \label{v13:T-factor}
\end{equation}
This keeps the two off-diagonal Frobenius occurrences: it follows directly from
$w=-C_h\beta/2$ in \eqref{v12:J1-off}.

\begin{lemma}[A quantitative improvement of the shifted coefficient bounds]
\label{v13:coefficient-bounds}
For the actual coefficient family and every complex or real vector array,
\begin{equation}
 \frac34\|w\|_h\le\|\widehat{\mathcal L}_hw\|_h
                           \le18\|w\|_h.              \label{v13:Lhat-bounds}
\end{equation}
The constants are independent of $N$ and of the admitted $\alpha$.
\end{lemma}
\begin{proof}
Use the input and output shift isometries of \cref{v12:T-inverse}. They
conjugate $\widehat{\mathcal L}_h$ to the pointwise matrix $M_b/\omega_h$.
Compare that matrix to the constant matrix
\[
 M_0=\begin{pmatrix}-8&4&0\\-8&0&4\\0&-8&-8\end{pmatrix}.
\]
The eigenvalues of $M_0^*M_0$ are $16$ and $136\pm8\sqrt{33}$. Thus its
smallest singular value is $4$ and its norm is less than $14$.
Each nonzero entry of $M_b/\omega_h-M_0$ has magnitude at most
\[
 \frac32+\frac{27}{256}=\frac{411}{256}.
\]
Indeed the oscillatory boost entries have size at most $3/2$, and the first
constant boost differs from $4$ by at most $\Sigma_\alpha/64\le27/256$.
Shifting those entries changes neither bound. Each row and column contains at
most two entries, so the induced $2$-norm of the difference is at most
$411/128<13/4$. The triangle and reverse triangle inequalities give the
lower bound $4-13/4=3/4$ and an upper bound less than $14+13/4<18$.
Integrate the pointwise bounds and undo the two isometries. This proof uses the
specific anisotropic calibration, not a free condition-number assumption.
\end{proof}

Define on \emph{all} of $\mathcal V_h$ the analytical normal matrix
\begin{equation}
 \boxed{\mathsf H_h
  =R_h\widehat{\mathcal L}_h^*\widehat{\mathcal L}_hR_h+Q_T.}
                                                               \label{v13:H}
\end{equation}
The identity on $\ker P_T$ only extends an inverse off the physical transverse
slice. It is not a mass, action term, or new equation on longitudinal shifts.

\begin{theorem}[Uniform spectral interval and exact Fourier bandwidth]
\label{v13:H-spectrum}
The matrix \eqref{v13:H} is self-adjoint, commutes with $P_T$, and obeys
\begin{equation}
 \boxed{\frac9{16}I\preceq\mathsf H_h\preceq324I.}   \label{v13:H-bounds}
\end{equation}
In the three-component unitary Fourier representation, its block entry between
frequencies $k,l$ is zero when their cyclic $\ell^1$ distance exceeds two.
No active lapse or higher weak rank assumption is used.
\end{theorem}
\begin{proof}
Decompose $v=P_Tv+Q_Tv$. The quadratic form is
$\|\widehat{\mathcal L}_hR_hP_Tv\|_h^2+\|Q_Tv\|_h^2$.
Equation \eqref{v13:curl} and \cref{v13:coefficient-bounds} give
\eqref{v13:H-bounds} on both orthogonal summands. The formula also proves
self-adjointness and commutation.
The coefficients $b_i/\omega_h$ have only constant and unit axial Fourier
frequencies. Each term of $\widehat{\mathcal L}_h$ consequently changes
frequency by at most one coordinate unit, modulo $N$; its literal shift is a
Fourier multiplier. Its adjoint has the same bandwidth. The operators $R_h$
and $Q_T$ are Fourier multipliers even at their zero modes. Composition gives
bandwidth two on the complete cyclic frequency group. No Fourier-face entry
has been deleted.
\end{proof}

\section{Exact pseudoinverse and the failure of an inverse-then-project shortcut}
\label{v13:pseudoinverse}

Let $B_h=T_h^*T_h$ on $\mathcal T_h$, and write $T_h^\dagger=B_h^{-1}T_h^*$
for its physical least-squares inverse. The following formulas distinguish this
operator from the unrestricted pointwise inverse in V12.

\begin{theorem}[The original transverse normal solve]
\label{v13:inverse-formulas}
With $\Omega_h^{-1}$ extended by zero on constants, the exact formulas are
\begin{align}
 B_h^{-1}f
   &=\frac8{h^2\omega_h^2}\,
       \Omega_h^{-1}\mathsf H_h^{-1}\Omega_h^{-1}f,
                    &&f\in\mathcal T_h,               \label{v13:B-inverse}\\
 T_h^\dagger Y
   &=\frac{2\sqrt2}{h\omega_h}\,
       \Omega_h^{-1}\mathsf H_h^{-1}R_h
       \widehat{\mathcal L}_h^*\iota_o^*Y,             \label{v13:T-dagger}\\
 \Pi_{T,h}:=T_hT_h^\dagger
   &=\iota_o\widehat{\mathcal L}_hR_h\mathsf H_h^{-1}R_h
                    \widehat{\mathcal L}_h^*\iota_o^* .\label{v13:PiT}
\end{align}
In particular $T_h^\dagger$ annihilates \emph{every} diagonal symmetric array.
The projection \eqref{v13:PiT} is the original positive orthogonal projection
onto $\Ran T_h$, not a prescribed projection of the dynamics.
\end{theorem}
\begin{proof}
Equation \eqref{v13:T-factor} gives
$B_h=(h^2\omega_h^2/8)\Omega_h\mathsf H_h\Omega_h$ on $\mathcal T_h$.
All factors are invertible on that slice. Inverting them in their actual
order proves \eqref{v13:B-inverse}; no commutation of $\Omega_h$ with the
variable normal matrix is assumed. Multiplication by the exact adjoint of
\eqref{v13:T-factor} proves \eqref{v13:T-dagger}, using that $\mathsf H_h^{-1}$
preserves $\mathcal T_h$ and $R_h$ lands there. Multiplication on the left by
$T_h$ yields \eqref{v13:PiT}. The characterization as an orthogonal projection
follows from $T_h^\dagger T_h=I$ and self-adjointness of $B_h^{-1}$.
Finally $\iota_o^*$ annihilates a diagonal array, exactly at every frequency.
\end{proof}

\begin{proposition}[The pointwise inverse followed by Hodge projection is different]
\label{v13:naive}
For an invertible map $M$ on a positive finite Hilbert space and an orthogonal
projection $P$, the minimizer of $\|Mw-y\|$ over $w\in\Ran P$ equals
$PM^{-1}y$ for every $y$ if and only if $[P,M^*M]=0$.
For $M=\widehat{\mathcal L}_h$, $P=P_T$, this commutation fails at the actual
three-site seed of \cref{v12:three-seed}.
\end{proposition}
\begin{proof}
Write $z=M^{-1}y$ and $A=M^*M$. The correct normal equation is
$PAPw=PAz$ on $\Ran P$. Substituting $w=Pz$ solves it for every $z$ exactly
when $PA(I-P)=0$. Self-adjointness of $A$ makes this equivalent to the
commutator condition.
For the stated action family all entries of $\widehat{\mathcal L}_h$ and
$P_T$ are rational at $N=3$. Construct $P_T$ using the original rational
phase derivatives as in \eqref{v10:rational-Hodge}, and construct
$\widehat{\mathcal L}_h$ from \eqref{v12:pointwise-M} with the literal input
and output shifts undone. In component-major vector order, zero-based entry
$(0,1)$ of the commutator has residue $87$ modulo $241$. The companion exact
verifier gives the full construction, checks all $6561$ pure-shift source
entries against the retained action certificate, and uses only denominators
that are units at that prime. This nonzero residue proves noncommutation in
characteristic zero. The general zero identities used to define the projection
follow from the displayed exact Hodge algebra, not from modular zeros alone.
\end{proof}

\section{Exponential localization in cyclic Fourier distance}
\label{v13:decay}

Represent frequencies by $\{-m,\ldots,m\}^3$, and define
\[
 d_N(k,l)=\sum_{i=1}^3\min\{|k_i-l_i|,N-|k_i-l_i|\}.
\]
This is a distance on the finite frequency group. In particular
$d_N(k,0)=|k|_1$ for these representatives. Let $P_E$ project onto any set
$E$ of frequencies, retaining all three components there. Distances between
sets mean the minimum of $d_N$.

\begin{theorem}[Uniform off-diagonal inverse estimate]
\label{v13:decay-theorem}
Put $q=23/25$. For any nonempty frequency sets $E,F$ at distance $d>0$,
\begin{equation}
 \boxed{\|P_E\mathsf H_h^{-1}P_F\|_{2\to2}
                 \le\frac{32}{9}\,q^{\lceil d/2\rceil}.} \label{v13:set-decay}
\end{equation}
The diagonal bound is $\|\mathsf H_h^{-1}\|\le16/9$.
All constants are uniform in the cutoff and in the admitted first-jet
parameters. The same inequality bounds an individual $3\times3$ inverse
block. This is a Fourier localization of an analytical inverse; it is not
finite spatial execution locality.
\end{theorem}
\begin{proof}
Let $a_0=9/16$ and $b_0=324$. For $j\ge0$ let $\mathrm T_{j+1}$ be the
Chebyshev polynomial and define
\[
 r_{j+1}(x)=
 \frac{\mathrm T_{j+1}((a_0+b_0-2x)/(b_0-a_0))}
      {\mathrm T_{j+1}((a_0+b_0)/(b_0-a_0))},\qquad
 p_j(x)=\frac{1-r_{j+1}(x)}x.
\]
Since $r_{j+1}(0)=1$, $p_j$ is a polynomial of degree $j$. On $[a_0,b_0]$
the numerator of $r_{j+1}$ has absolute value at most one. The denominator
is $(q^{-(j+1)}+q^{j+1})/2$, where
\[
 q=\frac{\sqrt{b_0}-\sqrt{a_0}}{\sqrt{b_0}+\sqrt{a_0}}=\frac{23}{25}.
\]
The spectral theorem and \cref{v13:H-spectrum} give
\[
 \|\mathsf H_h^{-1}-p_j(\mathsf H_h)\|
                    \le\frac2{a_0}q^{j+1}.
\]
Exact bandwidth two implies that $p_j(\mathsf H_h)$ connects no frequencies
at distance greater than $2j$. Choose $j=\lceil d/2\rceil-1$. Then
$P_Ep_j(\mathsf H_h)P_F=0$ and the last estimate proves
\eqref{v13:set-decay}. This is a direct polynomial proof on the original
cyclic graph; it does not replace that graph by an infinite one or discard
folded paths.
\end{proof}

\begin{corollary}[Uniform Sobolev and analytic-weight bounds]
\label{v13:weighted}
For every real $s$ there is $C_s$ independent of $h,\alpha$ such that
\begin{equation}
 \|\mathsf H_h^{-1}f\|_{H_h^s}\le C_s\|f\|_{H_h^s}.
                                                         \label{v13:Hs}
\end{equation}
More generally, if $0\le\sigma<\frac12\log(25/23)$, the same assertion
holds uniformly with the norm
\[
 \|f\|_{s,\sigma,h}^2
   =\sum_k\langle k\rangle^{2s}e^{2\sigma|k|_1}|\widehat f(k)|^2.
\]
The coefficients and their original shifts act boundedly on all these
spaces, with constants depending only on $s,\sigma$ and the fixed parameter
box. Every fixed parameter differential of the inverse has the same type
of bound, on a compact parameter neighborhood where the coefficient bounds
hold.
\end{corollary}
\begin{proof}
For the chosen frequency representatives, the triangle inequality in $d_N$
gives
\[
 \frac{\langle k\rangle^s}{\langle l\rangle^s}
   \le C_s(1+d_N(k,l))^{|s|},\qquad
 |k|_1-|l|_1\le d_N(k,l).
\]
An inverse block in \cref{v13:decay-theorem} is at most
$C\exp[-\gamma d_N(k,l)]$ in norm, with $\gamma=\tfrac12\log(25/23)$.
The number of frequency points at integer distance $d$ is bounded by
$C(1+d)^2$, uniformly in the finite torus. Weighting the blocks therefore
leaves a summable row and column majorant
$C_{s}(1+d)^{|s|+2}e^{-(\gamma-\sigma)d}$. The elementary Schur row-column
estimate proves both norm bounds, including negative $s$.
Coefficient multiplication changes frequency by at most one, and its shifted
coefficients are uniformly bounded, so the corresponding estimates follow
from the same weight ratios. For parameter derivatives differentiate
$\mathsf H_h\mathsf H_h^{-1}=I$, starting with
$\partial\mathsf H_h^{-1}=-\mathsf H_h^{-1}(\partial\mathsf H_h)\mathsf H_h^{-1}$.
All fixed coefficient differentials have bounded bandwidth and norms; repeated
product differentiation proves the last statement. No differentiation of a
floating inverse is part of the proof.
\end{proof}

\section{The actual prepared source has a controlled transverse output}
\label{v13:direct}

Let $b_h=(b_{N,h},b_{\beta,h})$ be the prepared quadratic drift in the
original physical active multiplier coordinates for \cref{v12:family}.
The sign convention for its leading equation is
$\mathcal A_h^*K_0^{-1}\mathcal A_h(n,\beta)=-b_h$.
Weak neutrality says that the shift target already lies in $\mathcal T_h$;
no incompatible target is removed by that restriction.

\begin{lemma}[Fixed-bank size of the complete prepared target]
\label{v13:target}
On sufficiently fine odd cutoffs, the complete prepared target has Fourier
support in $|k|_1\le2$, and for every fixed $s$,
\begin{equation}
 \|b_{N,h}\|_{H_h^s}+\|b_{\beta,h}\|_{H_h^s}\le C_s h. \label{v13:target-size}
\end{equation}
The bounds are uniform on the stated $\alpha$ box. They use the full prepared
lapse correction as well as the transverse logarithmic terms.
\end{lemma}
\begin{proof}
The original extracted drift is quadratic in the first fields, whose only
frequencies are zero and the six unit axial modes. The flat connection and
auxiliary maps are constant coefficient, and phase derivatives and shifts
preserve the frequency of a factor. After exact summation by parts in the
multiplier test, this proves the fixed-bank support, before any inverse.

The geometric part on this bank is analytic in $h^2$, since its only spatial
derivative symbols are $2h^{-1}\sin(\pi hk_i)$. Its constant coefficient is
zero after the actual preparation correction. We spell out the inherited
continuum Ward incidence used for this claim. The ordinary-derivative limit
of the extracted stationary Palatini coefficients is the canonical
gravitational Hamiltonian. Its momentum functional generates spatial pullback,
so $\{P_\beta,H_1\}=0$ for a fixed unit lapse. Its lapse bracket is
$\{P_n,H_1\}=-P_{\gamma^{-1}\nabla n}$, in the sign convention of
\cref{v5:prepared-B2}. At quadratic field order the latter gives
$-P_{\nabla n,2}(Y)+P_{U\nabla n,1}(Y)$.
Every first jet here satisfies the linear momentum constraint. Hence
$P_{\xi,1}(Y)=0$ for every vector test $\xi$, including $\xi=U\nabla n$.
The added $+P_{\nabla n,2}(Y)$ in
\eqref{v5:prepared-drift-formula} cancels the remaining term. This is an
initial coefficient identity, not a future Einstein-solution assumption.
It uses the same ordinary canonical envelope and Ward identities as the
retained fixed-bank consistency argument. Thus the geometric contribution is
$O(h^2)$ on a bank of fixed dimension.

Every cubic electric or magnetic term able to contribute has its explicit
factor $h$. Its remaining coefficients, including $\omega_h$, the calibrated
constant momentum, and all derivatives at fixed frequencies, stay bounded
uniformly on the compact $\alpha$ box. These terms are consequently $O(h)$.
Higher link terms cannot enter this quadratic drift extraction. The exact
constant and curl-free shift pairings are zero by \cref{v12:matching} and
the constant-shift identity. Combining the estimates gives
\eqref{v13:target-size}, with equivalent norms on the fixed bank. In
particular the proof does not apply the curl-free logarithmic cancellation
to transverse tests.
\end{proof}

\begin{theorem}[Uniform direct curvature decay and exponentially small face tails]
\label{v13:direct-theorem}
For the actual direct transverse solution
\[
 \beta_{0,h}=-2B_h^{-1}b_{\beta,h},
\]
and every fixed real $s$, one has
\begin{equation}
 \boxed{\|\beta_{0,h}\|_{H_h^{s+3}}\le C_s h^{-1},\qquad
        \|\mathfrak E_h(\beta_{0,h})\|_{H_h^s}\le C_s h.}
                                                        \label{v13:direct-result}
\end{equation}
There are $c_*>0$ and $C_s$, independent of the cutoff, such that
\begin{equation}
 \boxed{\|\mathscr P_{\partial,h}\beta_{0,h}\|_{H_h^s}
                          \le C_s h^{-1}e^{-c_*/h}.}     \label{v13:direct-tail}
\end{equation}
This concerns the actual direct source response, not a test vector or a
prescribed low-frequency projection of the solution.
\end{theorem}
\begin{proof}
On nonzero odd-grid modes $|\kappa_h(k)|\asymp|k|$, with uniform constants.
Thus $\Omega_h^{-1}$ gains one Sobolev derivative on mean-zero arrays.
Equation \eqref{v13:B-inverse} and \cref{v13:weighted} imply, for any $f\in
\mathcal T_h$ and real $s$,
\begin{equation}
 \|B_h^{-1}f\|_{H_h^{s+2}}
                             \le C_s h^{-2}\|f\|_{H_h^s}.\label{v13:B-Sobolev}
\end{equation}
Here $3\sqrt3\le\omega_h<2\pi$ is harmless, and all operator orders in the
factorization are retained. Apply this estimate at $s+1$ to
\cref{v13:target}; then apply the upper bound of \cref{v10:symbol} to obtain
\eqref{v13:direct-result}.

The same proof works with the exponential weights of \cref{v13:weighted}.
A fixed-bank target satisfies $\|b_{\beta,h}\|_{s,\sigma,h}\le C_{s,\sigma}h$.
Choose a fixed $0<\sigma<\gamma$ in that corollary. It follows that
$\|\beta_{0,h}\|_{s,\sigma,h}\le C_s h^{-1}$. On the folded faces
$|k|_1\ge m=(N-1)/2$, so restriction costs $e^{-\sigma m}$. For all
sufficiently fine grids this is at most $e^{-c_*/h}$, with a fixed smaller
$c_*>0$. This proves \eqref{v13:direct-tail}. The argument controls the tail
of the full inverse rather than deleting the face before solving.
\end{proof}

\section{The lapse observation has no separate folded-face loss on this family}
\label{v13:lapse-transfer}

Decompose the original lapse source exactly into its diagonal and off-diagonal
parts,
\[
 L_hn=\mathscr D_{h,\alpha}n+V_hn,\qquad
 (V_hn)_{ij}=-\frac h2(r_iS_i\delta_jn+r_jS_j\delta_in)\quad(i\ne j),
\]
with $V_h$ having zero diagonal. The geometric product defects, including their
large folded-frequency entries, are \emph{all} in $\mathscr D_{h,\alpha}$
by \cref{v12:J1}. They are retained in the lapse equation.

\begin{theorem}[All-frequency zero-order lapse-to-shift transfer]
\label{v13:transfer}
For every real $s$ and every scalar array $n$,
\begin{align}
 T_h^\dagger L_hn&=T_h^\dagger V_hn,\label{v13:diagonal-killed}\\
 \|T_h^\dagger V_hn\|_{H_h^s}&\le C_s\|n\|_{H_h^s},\label{v13:transfer-zero}\\
 \boxed{\|\mathfrak E_hT_h^\dagger L_hn\|_{H_h^s}
                         \le C_s h^2\|n\|_{H_h^{s+3}}.}\label{v13:transfer-output}
\end{align}
These estimates include every folded mode. There is no additional term
$h^{-2}\|\mathscr P_{\partial,h}n\|$ in this family's bound.
\end{theorem}
\begin{proof}
Equation \eqref{v13:T-dagger} annihilates any diagonal output, which proves
\eqref{v13:diagonal-killed} exactly. The same formula and
\cref{v13:weighted}, with the one-derivative gain of $\Omega_h^{-1}$, give
\[
 \|T_h^\dagger Y\|_{H_h^s}\le C_s h^{-1}\|\iota_o^*Y\|_{H_h^{s-1}}
\]
for all real $s$. Each coefficient $r_i$ has fixed unit axial support and
bounded coefficients; a literal shift is an isometry of every Fourier
Sobolev space. A phase derivative maps $H_h^s$ to $H_h^{s-1}$ uniformly,
including negative $s$. Multiplication by the fixed-support coefficient is
bounded on $H_h^{s-1}$ uniformly even across a cyclic face, by the frequency
weight ratio used in \cref{v13:weighted}. Therefore
$\|V_hn\|_{H_h^{s-1}}\le C_s h\|n\|_{H_h^s}$.
The two estimates prove \eqref{v13:transfer-zero}; their $h$ factors cancel
in the actual order. Apply \cref{v10:symbol} at index $s$ and then
\eqref{v13:transfer-zero} at $s+3$ to prove \eqref{v13:transfer-output}.
No phase Leibniz rule or unproved commutation of the pseudoinverse with
coefficients is used.
\end{proof}

The distinction from V11 is structural, not notational. On its circular
configuration family the folded geometric lapse output had an off-diagonal
component visible to the transverse inverse. Here that geometric output is
diagonal, so it is invisible to this particular observation before any
estimate is taken. It remains visible to the scalar lapse energy below.
The new bound does not assert that the unknown solved lapse has a bounded
$H_h^{s+3}$ norm, nor that an arbitrary Fourier-edge input has small output.

\section{A positive scalar lapse equation with the complete target retained}
\label{v13:scalar}

Since the entire leading source is trace-free, its signed active equation is
exactly the normal equation for its positive source energy. Using the true
projection \eqref{v13:PiT}, define
\begin{equation}
 \begin{split}
 \mathscr S_h&=\frac12\left\{
       \mathscr D_{h,\alpha}^*\mathscr D_{h,\alpha}
                      +V_h^*(I-\Pi_{T,h})V_h\right\},\\
 g_h&=-b_{N,h}+V_h^*T_hB_h^{-1}b_{\beta,h}.
 \end{split}                                             \label{v13:S-g}
\end{equation}
These operators act on the mean-zero lapse. Diagonal and off-diagonal
symmetric arrays are orthogonal in the original Frobenius pairing.

\begin{proposition}[Exact active reduction and the unresolved observation]
\label{v13:Schur}
Every solution of the actual leading active equation is equivalently given by
\begin{equation}
 \boxed{\mathscr S_hn_h=g_h,\qquad
        \beta_h=\beta_{0,h}-T_h^\dagger V_hn_h.}          \label{v13:active-solve}
\end{equation}
Its scalar quadratic form is the positive identity
\begin{equation}
 2\ip n{\mathscr S_hn}_h
     =\|\mathscr D_{h,\alpha}n\|_h^2
                          +\|(I-\Pi_{T,h})V_hn\|_h^2.  \label{v13:scalar-energy}
\end{equation}
The complete geometric output consequently satisfies
\begin{equation}
 \boxed{\|\mathfrak E_h(\beta_h)\|_{H_h^s}
              \le C_s h+C_s h^2\|n_h\|_{H_h^{s+3}}.}   \label{v13:final-output}
\end{equation}
The actual scalar right side $g_h$ has $\|g_h\|_{H_h^s}\le C_s h$ at every
fixed order. It also has a uniform exponentially weighted bound of this
size for a fixed sufficiently small positive weight. Its solved lapse is
not asserted to inherit that bound.
\end{proposition}
\begin{proof}
The transverse row is
$\tfrac12T_h^*V_hn+\tfrac12B_h\beta=-b_{\beta,h}$; its exact elimination
gives the second equation of \eqref{v13:active-solve}. Substituting in the
lapse row gives \eqref{v13:S-g}. The diagonal is orthogonal to $V_h$ and
to $\Ran T_h$, which explains why it contributes just
$\mathscr D^*\mathscr D/2$. The remaining off-diagonal term is
$V_h^*(I-\Pi_{T,h})V_h/2$. This proves both original rows and the converse.
Since $I-\Pi_{T,h}$ is an orthogonal projection,
\eqref{v13:scalar-energy} follows. The two actual output bounds are
\cref{v13:direct-theorem,v13:transfer}, proving \eqref{v13:final-output}.

For the last assertion, \cref{v13:target} bounds $b_N$ by $C_sh$.
The direct inverse has $\|B_h^{-1}b_\beta\|_{H_h^{s+2}}\le C_sh^{-1}$.
Both $T_h$ and $V_h^*$ are first-order finite coefficient maps with explicit
factor $h$, bounded from $H_h^{j+1}$ to $H_h^j$ by $C_jh$.
Their composition therefore gives
$\|V_h^*T_hB_h^{-1}b_\beta\|_{H_h^s}\le C_sh$.
The same argument in exponential weights uses the fixed-band coefficients
and \cref{v13:weighted}. It estimates this \emph{constructed right side}, not
$\mathscr S_h^{-1}$.
\end{proof}

This reduction neither identifies a pure-active solution with the full raw
multiplier rate nor suppresses the next weak source. Such an identification
needs the appropriate weak-rate preparation or control of its coupling.
The higher block in \eqref{v12:remaining-weak} is still signed and retains its
actual nonzero trace. The active subsystem and its observed output are the
objects estimated here.

\section{A sharper cofinal active lower bound}
\label{v13:active-bound}

The uniform zero-order estimate also improves the coarse active bound without
solving the remaining source-specific lapse problem.

\begin{theorem}[Four powers gained in the unweighted active Gram estimate]
\label{v13:active-theorem}
On the same rationally selected configuration family $\alpha_h$ used in
\cref{v12:main}, at every sufficiently fine odd cutoff,
\begin{equation}
 \boxed{\mathcal A_h^*K_0^{-1}\mathcal A_h
      =\tfrac12\mathcal A_h^*\mathcal A_h\succeq c h^{12}I.}
                                                        \label{v13:h12}
\end{equation}
The constant is uniform over the candidates supplied by the inherited
selection. This improves the sufficient $ch^{16}$ bound in V12; optimality
is not asserted. For the actual leading active target there is a unique
active response, with the coarse estimate
\[
 \|(n_h,\beta_h)\|_h\le C h^{-11}.
\]
This last bound is not strong enough to imply \eqref{v13:final-output}
tends to zero.
\end{theorem}
\begin{proof}
For $y=\mathcal A_h(n,\beta)$ its diagonal is
$\mathscr D_{h,\alpha_h}n$. The inherited detuning gives
$\|n\|_h\le Ch^{-6}\|y\|_h$. Apply the \emph{true} transverse pseudoinverse
to its off-diagonal equation. It gives
\[
 \beta=T_h^\dagger y-T_h^\dagger V_hn.
\]
The smoothing estimate proved before \eqref{v13:transfer-zero}, at $s=0$,
and \cref{v13:transfer} give
\[
 \|\beta\|_h\le Ch^{-1}\|\iota_o^*y\|_{H_h^{-1}}+C\|n\|_h
                  \le Ch^{-6}\|y\|_h.
\]
In contrast to the earlier elementary estimate, no full-grid derivative
bound $\|\delta n\|\le Ch^{-1}\|n\|$ is spent. Thus the least active
source singular value is at least $ch^6$. Trace-freeness of the original
range gives \eqref{v13:h12}. In finite dimension this also proves
invertibility; its inverse norm is at most $Ch^{-12}$.
Finally \cref{v13:target} gives $\|b_h\|_h\le Ch$, proving the coarse
response bound. Neither this estimate nor positive definiteness supplies
the high-order source-specific bound for the solved lapse.
\end{proof}

\section{Status, source audit, and reproducible checks}
\label{v13:status}

The transverse problem is now closed at the required operator level on the
complementary family: its normalized normal matrix is uniformly positive,
its inverse has uniform full-frequency localization and Sobolev bounds, and
its original least-squares projection is explicit. The direct prepared-source
curvature contribution is $O(h)$, with exponentially small Fourier-face
tails. The transfer from an arbitrary lapse has no separate folded-face loss
because the corresponding geometric source is exactly diagonal. These are
actual properties of the new coefficients, not assumptions about an unknown
inverse response. The cofinal unweighted active lower bound is also stronger.

What remains is the specific positive scalar inverse in \eqref{v13:active-solve}
observed through $\mathfrak E_hT_h^\dagger V_h$, and the independently signed
higher weak system. Exponential localization of $\mathsf H_h^{-1}$ is not
exponential localization of $\mathscr S_h^{-1}$: the latter still contains
the degenerate diagonal lapse operator and a different coefficient geometry.
Although $g_h$ itself is exponentially localized and $O(h)$, the present
polynomial bound does not control its inverse image at the required weighted
orders. No nonlinear rank propagation, initial-rate cancellation on this
family, differentiated curvature bound, common physical lifespan, or
raw-action Einstein continuum theorem follows from this body.

\paragraph{Verification scope.}
The exact verifier checks the constant comparison spectrum, all rational
singular-value margins, the Chebyshev residual construction through six
explicit degrees, and the exact factor $23/25$. It reconstructs the original
three-site pure-shift source and matches all $6561$ entries against the
retained coefficient certificate. The nonzero commutator entry in
\cref{v13:naive} is a new good-prime characteristic-zero witness; structural
zero identities have the written algebraic proofs above.

Separate floating checks use the unchanged complete coefficient evaluator
at $N=3,5,9,13,17$. They evaluate the actual prepared targets and the variable
transverse least-squares normal equations, check the physical adjoints and
full-bandwidth identities, and measure the direct observation and its face
tails. They include arbitrary full-frequency right sides, not only smooth
trials. Krylov solves there are diagnostics with recorded residuals, not
interval certificates or the proof of the uniform bounds. A Fourier-band
mask uses a stated tolerance only to classify mathematically integral mode
labels in floating output; no tolerance defines a rank or changes a retained
mode. The analytic proofs are independent of those diagnostic thresholds.

No new higher weak columns, nonlinear initial roots, or original-action
spacetime trajectories are represented as computed. The inherited initial
constraint calculus, scalar-detuning theorem and source coefficient identities
are explicit inputs. A targeted audit of the supplied bridge bodies and
relevant monograph ranges found the earlier warnings about the variable
pseudoinverse, but not this normalized full-frequency solve. No exhaustive
literature-priority assertion is made. All twelve preceding marked research
bodies are retained byte for byte, with their corrections and boundaries;
the accompanying manifests record preservation and current verification.

\begin{thebibliography}{9}
\raggedright
\setcounter{enumiv}{11}
\bibitem{v13:DMS}
S. Demko, W. F. Moss and P. W. Smith,
\emph{Decay rates for inverses of band matrices},
Mathematics of Computation \textbf{43} (1984), 491--499.
\href{https://doi.org/10.1090/S0025-5718-1984-0758197-9}
{doi:10.1090/S0025-5718-1984-0758197-9}.
This supplies primary-source context for polynomial inverse-decay methods,
not the particular normalized renewal source, its calibration, or its
geometric-output identities. The needed finite cyclic estimate is proved here.
\end{thebibliography}
% END IMMUTABLE EXACT-ACTION BRIDGE V13 BODY


% Future iterations append here.
% BEGIN IMMUTABLE EXACT-ACTION BRIDGE V14 BODY
\clearpage
\renewcommand{\thesection}{\arabic{section}}
\renewcommand{\theHsection}{bridgeV14.\arabic{section}}
\setcounter{section}{108}
\section{The positive lapse problem has a flux structure}
\label{v14:context}

The complementary family has a positive leading active Gram, and its true
transverse inverse has already been controlled in \cref{v13:inverse-formulas}.
Neither fact gives uniform ellipticity of the remaining scalar lapse equation.
We first express that equation using an explicit flux and a uniformly
conditioned metric on its longitudinal and constant components. We then test
the flux symbol on the actual complementary first data. It has a simple real
characteristic at which a nonzero transverse shift cancels the leading lapse
source. This is a loss of rank of the principal \emph{source}, not cancellation
by an indefinite pairing. The finite active Gram may remain positive and
invertible, as proved in V12--V13, while this lower-frequency characteristic
persists.

The results give an exact scalar flux formulation, uniform control of every
auxiliary metric in that formulation, and actual finite scalar tests on which
small residual does not imply a small geometric output. The actual right side
has an exponentially small pairing with those localized tests. It is therefore
not legitimate to substitute them for the prepared source and conclude failure
of its continuum limit. The remaining estimate is still source-specific, but
its differential structure is now explicit.

\paragraph{Inputs and scope.}
The first-source formula is \cref{v12:J1}; its uniformly conditioned shifted
map, true transverse inverse, and scalar reduction are
\cref{v13:H-spectrum,v13:inverse-formulas,v13:Schur}. The actual right side and
its analytic Fourier bounds are retained from \cref{v13:Schur}. We use the
geometric output \cref{v10:symbol} and the trigonometric localization estimates
of \cref{v11:packet-lemma}. Those estimates are not proved again as new work.
The new characteristic below is different from \cref{v11:characteristic}:
there the frozen source has three independent columns and an indefinite Gram
becomes singular; here the frozen source itself has rank two and the Cartan
form is positive on its range. No new higher weak source, phase root, exact
nonlinear trajectory, or physical lifespan is asserted.

\section{An exact longitudinal-flux reconstruction with the means retained}
\label{v14:flux}

Work on V12's complementary family, with
$\alpha\in[1/2,3/2]^3$. All finite adjoints are for the positive spatial-mean
pairings. Write
\[
 M_h=\widehat{\mathcal L}_h,\qquad P_h=P_T,\qquad Q_h=I-P_h,
 \qquad A_h=M_h^*M_h.
\]
Here $M_h$ is the invertible three-component coefficient map of V13, not the
multiplier Hessian. The notation $Q_h$ in this body is a vector Hodge
projection, not a geometric difference of derivatives. Its range consists of
constant vectors and phase gradients. Let $V_h$ be the off-diagonal lapse
source in \cref{v13:lapse-transfer} and define
\begin{equation}
 J_h^L=h^{-1}\iota_o^*V_h,\qquad
 z_h(n)=M_h^{-1}J_h^L n,\qquad \sigma_h(n)=Q_hz_h(n).
 \label{v14:z}
\end{equation}
Thus, in raw off-diagonal coordinates,
\[
 (J_h^Ln)_{ij}=-\frac1{\sqrt2}
     (r_iS_i\delta_jn+r_jS_j\delta_in),\qquad i<j.
\]
The factors $\sqrt2$ are from the original Frobenius pairing. Define on
$\Ran Q_h$ the positive metric
\begin{equation}
 K_{Q,h}=Q_hM_h^{-1}M_h^{-*}Q_h\big|_{\Ran Q_h}.
 \label{v14:KQ}
\end{equation}

\begin{theorem}[Exact flux projection and positive scalar energy]
\label{v14:flux-theorem}
At every odd cutoff, $K_{Q,h}$ is invertible and
\begin{equation}
 \frac1{324}I\preceq K_{Q,h}\preceq\frac{16}{9}I,
 \qquad \frac9{16}I\preceq K_{Q,h}^{-1}\preceq324I.
 \label{v14:K-bounds}
\end{equation}
The original off-diagonal source residual satisfies
\begin{equation}
 (I-\Pi_{T,h})V_hn
   =h\,\iota_o M_h^{-*}Q_hK_{Q,h}^{-1}\sigma_h(n).
 \label{v14:residual}
\end{equation}
Consequently the \emph{unchanged} scalar operator of \eqref{v13:S-g} is
\begin{equation}
 \boxed{\mathscr S_h
 =\tfrac12\mathscr D_{h,\alpha}^*\mathscr D_{h,\alpha}
       +\tfrac{h^2}{2}(\sigma_h)^*K_{Q,h}^{-1}\sigma_h,}
 \label{v14:S-flux}
\end{equation}
and its exact energy is
\begin{equation}
 2\ip n{\mathscr S_h n}_h
 =\|\mathscr D_{h,\alpha}n\|_h^2
                   +h^2\ip{\sigma_h(n)}{K_{Q,h}^{-1}\sigma_h(n)}_h.
 \label{v14:energy}
\end{equation}
The second summand is uniformly comparable to the squared flux norm,
with the explicit constants in \eqref{v14:K-bounds}. This comparison does
not control the lapse norm by its flux.
\end{theorem}
\begin{proof}
V13 proves $3\|v\|/4\le\|M_hv\|\le18\|v\|$. Hence
$M_h^{-1}M_h^{-*}$ has spectrum in $[1/324,16/9]$. Compression to the
positive Hilbert space $\Ran Q_h$ preserves these quadratic-form bounds,
proving \eqref{v14:K-bounds}.

Suppress $h$ and consider $B=QM^{-1}$ as a map from all off-diagonal
coordinates to $\Ran Q$. It is onto, since $B(Mq)=q$ for $q\in\Ran Q$,
and its kernel is $M\Ran P$. Thus
$B^*(BB^*)^{-1}B=M^{-*}QK_Q^{-1}QM^{-1}$ is the positive orthogonal
projection onto $(M\Ran P)^\perp$. For example, it is self-adjoint and
idempotent, is the identity on $\Ran B^*$, and has kernel $\ker B$.
The physical transverse range is $\iota_o M\Ran P$ by
\eqref{v13:T-factor}. Substituting $\iota_o^*Vn=hJ^Ln$ proves
\eqref{v14:residual}. Squaring its positive norm gives
\[
 \|(I-\Pi_T)Vn\|^2
   =h^2\ip{\sigma(n)}{K_Q^{-1}\sigma(n)}.
\]
Insert this in \eqref{v13:scalar-energy}; polarization gives the operator
identity as well as the energy. No source equation is replaced by a
minimizing rule on the physical evolution.
\end{proof}

The flux can be represented by one scalar array and three means. On
$\mathcal Z_h=\R^3\oplus\{\varphi:\overline\varphi=0\}$ let
\begin{equation}
 U_h(c,\varphi)=c+\nabla_\delta\Omega_h^{-1}\varphi.
 \label{v14:U}
\end{equation}
This is an isometry onto $\Ran Q_h$: the two summands are orthogonal and
$\nabla_\delta^*\nabla_\delta=\Omega_h^2$. In these coordinates the actual
flux is
\begin{equation}
 \boxed{\mathcal F_h n=U_h^*\sigma_h(n)
   =\left(\overline{z_h(n)},\ -\Omega_h^{-1}\diver_\delta z_h(n)\right).}
 \label{v14:F}
\end{equation}
The metric is $\mathcal K_h=U_h^*K_{Q,h}U_h$. Thus
\eqref{v14:S-flux} also holds with $\sigma_h,K_{Q,h}$ replaced by
$\mathcal F_h,\mathcal K_h$. The three means are retained; deleting them
would change the scalar energy. All maps use the original phase operators
and literal shifted coefficients, including their folded frequencies.

\section{There is no further singular inverse hidden in the flux metric}
\label{v14:metric-control}

\begin{proposition}[An exact inverse formula using the existing normalized solve]
\label{v14:metric-inverse}
Let $R_h=\curl_\delta\Omega_h^{-1}$ and $\mathsf H_h$ be as in V13. On
$\Ran Q_h$,
\begin{equation}
 \boxed{K_{Q,h}^{-1}
   =Q_h\{A_h-A_hR_h\mathsf H_h^{-1}R_hA_h\}Q_h.}
 \label{v14:inverse-Schur}
\end{equation}
Both $K_{Q,h}^{-1}$ and $M_h^{-1}$ have cutoff-independent bounds on every
fixed Fourier Sobolev space. They also have such bounds in a sufficiently
small fixed positive exponential Fourier weight. In particular
\begin{equation}
 \|z_h(n)\|_{H_h^{s-1}}+\|\mathcal F_hn\|_{H_h^{s-1}}
                      \le C_s\|n\|_{H_h^s}
 \label{v14:flux-order}
\end{equation}
for every real $s$.
\end{proposition}
\begin{proof}
Relative to $\Ran P\oplus\Ran Q$, the compression $QA^{-1}Q$ is
$K_Q$. Its inverse is the Schur complement
$QAQ-QAP(PAP)^{-1}PAQ$. One can verify this without an inverse formula:
solve $Ax=q$, eliminate $Px$ using the positive $PAP$ block, and substitute
in the $Q$ row. On $\Ran P$ the inverse of $PAP$ is
$R\mathsf H^{-1}R$, since $R^2=P$ and $\mathsf H=RAR+Q$.
This proves \eqref{v14:inverse-Schur} in the correct operator order.
The fixed-band maps $A,R,Q$ and V13's inverse estimate give its Sobolev
and exponential-weight bounds directly.

For clarity, the full coefficient inverse in \eqref{v14:z} also admits an
explicit uniform argument. Undo the input and output shift isometries of
\cref{v12:T-inverse}. In those coordinates $M=M_0+E_h$, where $M_0$ is
V13's constant comparison matrix and
\[
 \|M_0^{-1}\|=1/4,\qquad \|E_h\|\le411/128,
 \qquad \rho:=\|M_0^{-1}\|\,411/128=411/512<1.
\]
The perturbation changes cyclic Fourier frequency by at most one. The
Neumann series for $(I+M_0^{-1}E_h)^{-1}M_0^{-1}$ therefore gives, for
frequency sets at positive cyclic distance $d$,
\begin{equation}
 \|P_E M_h^{-1}P_F\|\le\frac{128}{101}\rho^d.
 \label{v14:M-decay}
\end{equation}
Indeed its terms through degree $d-1$ have zero such block, and the tail
is at most $\rho^d/[4(1-\rho)]$. Restoring the shift isometries changes
no Fourier support. The weighted row-column argument already proved in
\cref{v13:weighted} applies to this exponential majorant, with any
exponential weight smaller than $-\log\rho$. This proves the stated
bounds for $M_h^{-1}$; its adjoint has the same bounds by duality.
The flux metric inverse may use a weight smaller than both this constant
and the V13 weight.

Finally $J_h^L$ is a first-order map with fixed-support bounded coefficients
and literal shifts. Thus it maps $H_h^s$ to $H_h^{s-1}$ uniformly.
$Q_h$, the means, and the normalized gradient/divergence in $U_h,U_h^*$
are uniformly bounded Fourier maps of order zero. This proves
\eqref{v14:flux-order}. None of these statements is a bound for
$\mathscr S_h^{-1}$.
\end{proof}

Equations \eqref{v14:energy} and \eqref{v14:F} locate the remaining issue:
it is the ability of the particular lapse to be recovered from a diagonal
product-defect observation and a weighted divergence flux. It is not an
unknown condition number in an auxiliary Hodge projection. In particular a
pointwise coefficient inverse alone does not discard the divergence condition.

\section{The frozen flux on the actual complementary profiles}
\label{v14:principal}

Let $\alpha$ stay in the admitted positive box and first take the
ordinary smooth limit of its coefficient fields. At the spatial point
\begin{equation}
                        x_*=(1/4,0,0),
 \label{v14:point}
\end{equation}
the diagonal rotation and boost coefficients satisfy
\begin{equation}
 r_0(x_*)=\pi R,\quad R=\alpha_1(0,-1,1),\qquad
 b_0(x_*)=2\pi B,
 \label{v14:RB}
\end{equation}
where
\begin{equation}
 B_1=4-\frac{\sum_i\alpha_i^2}{64}
                       +\frac{\alpha_2+\alpha_3}{2},\quad
 B_2=-8-\frac{\alpha_3}{2},\quad B_3=-8-\frac{\alpha_2}{2}.
 \label{v14:B-values}
\end{equation}
Thus $B_1>0$, $B_2,B_3<0$, uniformly away from zero. This is a point of
the original trigonometric profiles; it need not itself be a grid site.
Localized samples near it will be used below.

For any diagonal real $R,B$ define
\begin{equation}
 M_B=\begin{pmatrix}B_2&B_1&0\\B_3&0&B_1\\0&B_3&B_2\end{pmatrix},
 \quad
 l_R(\xi)=
 \begin{pmatrix}R_1\xi_2+R_2\xi_1\\
 R_1\xi_3+R_3\xi_1\\R_2\xi_3+R_3\xi_2\end{pmatrix}.
 \label{v14:M-l}
\end{equation}
On an oscillatory mode of wavevector $\ell\xi$ with
$1\ll\ell\ll h^{-1}$, the leading raw off-diagonal source is
\begin{equation}
 \mathrm i\pi^2h\ell\,
 \mathcal B(R,B,\xi)(n,\beta),\qquad
 \mathcal B(R,B,\xi)(n,\beta)
       =-l_R(\xi)n+M_B(\xi\times\beta),\quad\xi\cdot\beta=0.
 \label{v14:source-symbol}
\end{equation}
The geometric diagonal source is of smaller order $O(h^2\ell^2)$ in this
regime. These statements follow by substituting
$r=\pi R+O(h^2)$, $b=2\pi B+O(h^2)$,
$\delta=2\pi\mathrm i\ell\xi+o(\ell)$, and $S_i=I+O(h\ell)$ into
\eqref{v12:J1-off}. A finite norm version is proved below; no inverse is
inferred from a frozen symbol.

\begin{lemma}[Scalar principal flux and positive Schur form]
\label{v14:flux-symbol}
Suppose $B_1B_2B_3\ne0$ and $\xi\ne0$. Set
\[
 \chi(R,B,\xi)=\xi^TM_B^{-1}l_R(\xi),\qquad
                 \nu_B(\xi)=M_B^{-*}\xi.
\]
After eliminating the two transverse columns of the positive Gram of
\eqref{v14:source-symbol}, its scalar Schur form, with the common
$\pi^4h^2\ell^2$ factor omitted, is
\begin{equation}
           \boxed{q_+(R,B,\xi)=
                   \frac{\chi(R,B,\xi)^2}{|\nu_B(\xi)|^2}.}
 \label{v14:qplus}
\end{equation}
This normalization uses the original Frobenius and Cartan pairings.
The flux characteristic is $\chi=0$, not a change in the sign of $q_+$.
\end{lemma}
\begin{proof}
$\det M_B=-2B_1B_2B_3\ne0$. The range of the transverse off-diagonal
columns is $M_B\xi^\perp$, whose Euclidean orthogonal complement is
spanned by $M_B^{-*}\xi$. The residual of $l_R$ after orthogonal
projection onto that range is
$\nu_B(\nu_B\cdot l_R)/|\nu_B|^2$. Its squared Euclidean norm is
\eqref{v14:qplus}. On raw off-diagonal entries the Frobenius norm counts
each entry twice and the Cartan inverse is $I/2$. These factors cancel
in the signed Gram, which here equals the positive Euclidean Gram.
This proves the stated normalization and the formula.
\end{proof}

\section{A visible source-null characteristic throughout the diagonal calibration}
\label{v14:characteristic}

\begin{theorem}[A simple flux zero and a nonzero shift without signed resonance]
\label{v14:characteristic-theorem}
For \eqref{v14:RB} and $\xi=(0,1,t)$,
\begin{equation}
 \boxed{\chi(t)=
       \frac{\alpha_1(1-t)(B_2+B_3t)}{2B_2B_3}.}
 \label{v14:factor}
\end{equation}
It has the two distinct simple roots $1$ and
\begin{equation}
                         t_*=-B_2/B_3<0.
 \label{v14:tstar}
\end{equation}
At the latter root the frozen complete active source has rank two, with
kernel spanned by $(1,\beta_*)$, where
\begin{equation}
 w_*=M_B^{-1}l_R(\xi_*),\qquad
 \boxed{\beta_*=-\frac{\xi_*\times w_*}{|\xi_*|^2}\ne0,
             \qquad\xi_*=(0,1,t_*).}
 \label{v14:betastar}
\end{equation}
All source columns are off-diagonal and trace-free. Their positive Gram has
nullity one; this is not a $2/3$ trace-compression eigenvalue.
For $\alpha=(1,1,1)$ one has the exact values
\begin{equation}
 B=(317/64,-17/2,-17/2),\quad t_*=-1,\quad
 \boxed{\beta_*=(2/17,-317/9248,-317/9248),}
 \label{v14:unit}
\end{equation}
and
\begin{equation}
 q_+(t)=
 \frac{100489(t-1)^2(t+1)^2}
                  {692361t^2-982766t+692361}.
 \label{v14:unit-q}
\end{equation}
The simple flux zero and nonzero shift are uniform on the stated compact
$\alpha$ box. In particular V12's cutoff-dependent rational detunings do
not remove this principal characteristic.
\end{theorem}
\begin{proof}
At $x_*$, only the derivative of the first configuration profile is
nonzero. Equation \eqref{v12:rotation-diagonal} gives
$R=(0,-\alpha_1,\alpha_1)$; \eqref{v12:family} gives
\eqref{v14:B-values}. For $\xi=(0,1,t)$, the lapse vector is
$l_R(\xi)=(0,0,\alpha_1(1-t))^T$. Solving the three pointwise equations
$M_Bw=l_R$ gives
\[
 w=\alpha_1(1-t)
       \left(-\frac{B_1}{2B_2B_3},\frac1{2B_3},\frac1{2B_2}\right).
\]
Pairing with $\xi$ proves \eqref{v14:factor}.
Since $B_2,B_3$ have the same nonzero sign, $t_*<0$ differs from one.
The derivative at that root is
$\alpha_1(B_2+B_3)/(2B_2B_3)\ne0$.
At the root $w_*\perp\xi_*$ and $w_*\ne0$, because
$\alpha_1(1-t_*)\ne0$ and $M_B$ is invertible. The vector identity
$\xi\times(\xi\times w)=-|\xi|^2w$ on $\xi^\perp$ proves that
$\xi_*\times\beta_*=w_*$. Thus \eqref{v14:source-symbol} annihilates
$(1,\beta_*)$. The two transverse columns are independent: curl on
$\xi^\perp$ and $M_B$ are both invertible there. Hence the complete rank
is exactly two and its positive Gram has nullity one. The entire frozen
range is off-diagonal, so every orthonormal range basis has zero trace
compression. There is no third independent column on which an indefinite
null cancellation could be claimed.

At unit amplitudes the displayed solution gives \eqref{v14:unit}.
Direct rational substitution in \eqref{v14:qplus} gives
\eqref{v14:unit-q}. Its denominator is a positive squared norm times a
nonzero positive constant, so the two zeros are double zeros of this
\emph{positive} Schur form. The derivative, the inverse coefficient
margins, and $|\beta_*|$ are continuous and nonzero on the compact
parameter box, giving uniform bounds. All the detuned amplitudes remain
in that box. This proves the last statement without a global finite-rank
or inverse assumption.
\end{proof}

The obstruction is not tied to the numerical constants $8$ and $4$ used
in the homogeneous boost calibration. The factorization
\eqref{v14:factor} holds for any nonzero diagonal $B$ with $B_2B_3>0$ at
this coefficient pattern. Altering those homogeneous diagonal values while
keeping their signs cannot remove the visible flux root. It can move the
root and change its shift coefficient. This is a pointwise principal
statement, not a no-go theorem for arbitrary configurations, nondiagonal
momenta, or the evolution of the original source.

\section{Actual finite positive quasimodes with nonzero geometric output}
\label{v14:packets}

The preceding calculation is not used as a claim about a finite-grid exact
kernel. The following theorem retains the diagonal product defects and every
literal shift, and converts it into actual tests of the positive finite
matrix. Denote by $\mathcal M_{2,h}=J_{1,h}^*K_0^{-1}J_{1,h}$ the leading
active multiplier Hessian on the complementary family.

\begin{theorem}[Interior source-null packets for the original positive active system]
\label{v14:packet-theorem}
For every sequence of admitted complementary amplitudes, and all sufficiently
fine odd cutoffs, there are real physical active pairs
$\nu_h=(n_h,\beta_h)$ with
\begin{equation}
 \boxed{\|\nu_h\|_h=1,\qquad
 \|J_{1,h}\nu_h\|_h\le Ch^{2/3},\qquad
 \|\mathcal M_{2,h}\nu_h\|_{H_h^1}\le Ch^{1/3},\qquad
 \|\mathfrak E_h(\beta_h)\|_h\ge c>0.}
 \label{v14:packet-bounds}
\end{equation}
The spectra of these pairs, their source images, and their Hessian images
stay strictly inside the retained Fourier box. The constants are uniform
on the parameter box, and no full higher-source rank or exact history is
required. On the rationally selected active-regular family the positive
finite matrices remain invertible; \eqref{v14:packet-bounds} does not
contradict that fact.
\end{theorem}
\begin{proof}
Put $\ell=\lfloor h^{-2/3}\rfloor$, $m_e=\lfloor\ell/8\rfloor$ and
$k_\ell=(0,\ell,\operatorname{round}(t_*\ell))$. Use the normalized
trigonometric envelope of \cref{v11:packet-lemma}, centered at
$x_*=(1/4,0,0)$, with degree $m_e$ in each coordinate. Its product with
the carrier $e^{2\pi\mathrm i k_\ell\cdot x}$ is denoted $\psi_h$.
The same finite quadrature and envelope estimates apply to this carrier:
$\|\psi_h\|_h=1$, coefficient replacement at $x_*$ costs
$O(\ell^{-1/2})$, and
\[
 \|(D_h-2\pi\mathrm i\ell\xi_*)\psi_h\|_h\le C\ell^{1/2}.
\]
The carrier and its finitely expanded box are away from zero and the
folded faces since $|t_*|$ is uniformly bounded and $\ell/N\to0$.
All these facts are uniform in the current amplitudes.

Set $a_*=(\beta_*\times\xi_*)/|\xi_*|^2$ and define a complex active pair
\begin{equation}
 \nu_h^{\C}=
 \left(\psi_h,\;\frac1{2\pi\mathrm i\ell}
                         \curl_\delta(a_*\psi_h)\right).
 \label{v14:packet}
\end{equation}
Its shift is exactly phase-transverse and mean-zero, while its lapse has
zero mean. By the phase-symbol expansion and the envelope estimate, the
shift is $\beta_*\psi_h+O_{L_h^2}(\ell^{-1/2}+(h\ell)^2)$.
The norm of the pair consequently tends to
$\sqrt{1+|\beta_*|^2}$, uniformly bounded above and below.

Let $\varepsilon_h=\ell^{-1/2}+h\ell$. On the carrier box expanded by two
unit offsets, the original source and its adjoint have norm at most
$Ch\ell$. Indeed the off-diagonal lapse and shift terms each have one
phase derivative, explicit factor $h$, bounded fixed-frequency coefficients,
and literal shifts. The nonfolded geometric product defects have norm at
most $Ch^2\ell^2$ by the exact sine addition bound of V11; this is at most
$Ch\ell$ here. The same bound applies to the adjoint Fourier blocks.

Replacing the coefficients by \eqref{v14:RB}, the derivative by its
carrier value, and the literal shifts by the identity in the off-diagonal
terms has total error at most $Ch\ell\varepsilon_h$. Derivatives of the
envelope cost relative $\ell^{-1/2}$; a one-site shift costs $Ch\ell$;
the difference $\omega_h-2\pi$ is $O(h^2)$; and the phase derivative
error has relative size $O((h\ell)^2)$. The entire diagonal source is
of order $h^2\ell^2$, and is included in that error. Thus, unlike V11's
signed characteristic, the \emph{complete source} here satisfies
\[
 J_{1,h}\nu_h^{\C}
 =\mathrm i\pi^2h\ell\,
   \iota_{\rm raw}\mathcal B(R,B,\xi_*)(1,\beta_*)\psi_h+e_h
 =e_h,\qquad \|e_h\|_h\le Ch\ell\varepsilon_h.
\]
The fixed raw inclusion $\iota_{\rm raw}$ simply duplicates symmetric
off-diagonal entries. The principal term is zero by
\cref{v14:characteristic-theorem}, not by a signed Gram cancellation.
The error still has the same bounded-offset Fourier support. Applying the
actual bounded adjoint and $K_0^{-1}$ gives
\[
 \|\mathcal M_{2,h}\nu_h^{\C}\|_h
       \le Ch^2\ell^2\varepsilon_h,\qquad
 \|\mathcal M_{2,h}\nu_h^{\C}\|_{H_h^1}
       \le Ch^2\ell^3\varepsilon_h.
\]
Since $h\ell\asymp h^{1/3}$ and
$\varepsilon_h=O(h^{1/3})$, these give the first two small bounds in
\eqref{v14:packet-bounds}.

For the observation, the exact full-frequency symbol in V10 and its sine
expansion on this interior box give
\begin{equation}
 \frac{\mathfrak E_h(\beta_h^{\C})}{h^2\ell^3}
 =\frac{\pi^3\mathrm i}{6}
       (\xi_{*,i}^3\beta_{*,j}+\xi_{*,j}^3\beta_{*,i})_{ij}\psi_h
          +o_{L_h^2}(1).
 \label{v14:visible}
\end{equation}
To bound the cubic multiplier replacement, factor
$k_i^3-\ell^3\xi_{*,i}^3$ and use the first derivative envelope bound
and the exact finite bandwidth. Its relative error is $O(\ell^{-1/2})$;
the remaining phase-symbol error is $O((h\ell)^2)$. The explicit curl
in \eqref{v14:packet} obeys the same estimates. The tensor in
\eqref{v14:visible} is nonzero because both
$(\xi_{*,i}^3)_i$ and $\beta_*$ are nonzero and
$\|\sym(d\otimes b)\|_F^2=(|d|^2|b|^2+|d\cdot b|^2)/2$.
Its norm has a uniform positive lower bound on the compact parameter box.
Finally $h^2\ell^3\to1$. Take the real part and normalize the full pair.
The positive and negative carrier boxes are disjoint, also after the two
source offsets, so the squared norms of real and imaginary parts split
as in the inherited packet construction. This proves all bounds for real
arrays with every claimed support property.
\end{proof}

The estimate $\|J_{1,h}\nu_h\|=O(h^{2/3})$ is the diagnostic difference
from V11's $\|J_{1,h}\nu_h\|\asymp h^{1/3}$ signed example. In the
present family the principal source loses rank before its positive Gram
is formed. A zero trace compression therefore does not remove this
mechanism. Nor does this theorem identify any constructed test with the
actual prepared target.

\section{The characteristic survives exact scalar minimization}
\label{v14:scalar-tests}

\begin{theorem}[Geometrically visible quasimodes of the positive scalar lapse equation]
\label{v14:scalar-theorem}
On the same family there are real mean-zero lapse arrays $\widehat n_h$
with $\|\widehat n_h\|_h=1$ such that
\begin{equation}
 \boxed{\|\mathscr S_h\widehat n_h\|_{H_h^1}\le Ch^{1/3},\qquad
 \|\mathfrak E_hT_h^\dagger V_h\widehat n_h\|_h\ge c>0.}
 \label{v14:scalar-tests-eq}
\end{equation}
The lapse itself has the strictly interior carrier support above. The
minimizing shift is the actual full-frequency least-squares solution and
is not asserted to have compact Fourier support. In particular there is
no cutoff-independent estimate
\begin{equation}
 \|\mathfrak E_hT_h^\dagger V_hn\|_h
                 \le C\|\mathscr S_hn\|_{H_h^1}
 \label{v14:false-bound}
\end{equation}
for arbitrary mean-zero scalar inputs. This conclusion applies even on
the cofinally positive, invertible family selected in V12.
\end{theorem}
\begin{proof}
Start with the real pair $(n,\beta)$ of
\cref{v14:packet-theorem}, and put
$y=J_{1,h}(n,\beta)$ and $\beta_{\min}=-T_h^\dagger V_hn$.
Since the diagonal part of $y$ is annihilated by the pseudoinverse,
\[
 \beta_{\min}-\beta=-T_h^\dagger y.
\]
The source error has $\|y\|_{H_h^2}\le C\ell^2h\ell\varepsilon_h$ by
its exact carrier bandwidth. The full-frequency smoothing estimate in
V13 therefore gives
\[
 \|\mathfrak E_h(\beta_{\min}-\beta)\|_h
 \le Ch^2\|T_h^\dagger y\|_{H_h^3}
 \le Ch\|\iota_o^*y\|_{H_h^2}
 \le Ch^2\ell^3\varepsilon_h=o(1).
\]
Hence the actual minimizing shift has the nonvanishing observation in
\eqref{v14:scalar-tests-eq}.

For the residual, let $(m_N,m_\beta)=\mathcal M_{2,h}(n,\beta)$.
Exact elimination of the shift block, without requiring a lapse inverse,
gives
\begin{equation}
 \mathscr S_hn=m_N-(T_h^\dagger L_h)^*m_\beta.
 \label{v14:scalar-residual}
\end{equation}
Here $T_h^\dagger L_h=T_h^\dagger V_h$ is uniformly bounded on every
$H_h^s$ by V13. Its adjoint is bounded on $H_h^1$ by the same assertion
at $s=-1$ and physical duality. Thus
$\|\mathscr S_hn\|_{H_h^1}\le C\|\mathcal M_{2,h}(n,\beta)\|_{H_h^1}$,
which is $O(h^{1/3})$. The lapse norm tends to a positive bounded value;
rescale it, together with its minimizing shift, to one. This proves the
stated theorem and the contradiction to \eqref{v14:false-bound}.
\end{proof}

Positivity has not been lost. For example exact minimization also gives
$2\ip{\widehat n_h}{\mathscr S_h\widehat n_h}\le Ch^{4/3}$.
The inherited lower bound is of order $h^{12}$, which is compatible with
this test. The theorem rules out a particular uniform observation estimate;
it does not negate finite uniqueness, imply a negative mode, or prove
failure for the actual source.

\section{The actual scalar right side is separated from these localized tests}
\label{v14:actual-source}

\begin{proposition}[Exponential source separation for the scalar flux characteristic]
\label{v14:source-separation}
Let $g_h$ be the \emph{actual} right side in \eqref{v13:S-g}. For the
normalized lapse tests of \cref{v14:scalar-theorem}, there are $c_0,C>0$
independent of the cutoff and selected amplitudes such that
\begin{equation}
             \boxed{|\ip{\widehat n_h}{g_h}_h|
                           \le Ch\exp(-c_0h^{-2/3}).}
 \label{v14:source-pairing}
\end{equation}
No spectral projection onto the small eigenvalues of $\mathscr S_h$ is
replaced by these packet tests in this assertion.
\end{proposition}
\begin{proof}
By \cref{v13:Schur}, the actual constructed right side has a bound
$\|g_h\|_{0,\sigma,h}\le C_\sigma h$ for one fixed $\sigma>0$.
Both carrier boxes of $\widehat n_h$ have
$|k|_1\ge|k_2|\ge\ell-m_e\ge c\ell$.
Physical Parseval and Cauchy--Schwarz, with the exponential weight on
$g_h$, give
$|\ip{\widehat n_h}{g_h}|\le C h e^{-\sigma(\ell-m_e)}$.
Here the lapse, not the nonlocal minimizing shift, is paired with the scalar
source, so its compact Fourier support is exact. Since
$\ell\asymp h^{-2/3}$, this is \eqref{v14:source-pairing}.
\end{proof}

This is a positive fact about the original source, not merely a disclaimer
that the packet was chosen independently. It remains insufficient for
$\mathfrak E_hT_h^\dagger V_h\mathscr S_h^{-1}g_h$: packets are not an
orthogonal eigenbasis, and the inverse may mix them with distant frequencies.
The polynomial lower bound on $\mathscr S_h$ does not, without another
argument, turn \eqref{v14:source-pairing} into a spectral nonconcentration
estimate or a bound for the full inverse output.

\section{Consequences and reproducibility}
\label{v14:status}

The positive scalar problem is now written in an exact flux form. Its
longitudinal metric is uniformly controlled and uses only the already solved
normalized transverse inverse. The remaining degeneracy is in the actual
map from lapse to weighted divergence flux and diagonal product defects.
On the original complementary profiles it has a simple visible flux
characteristic, with an explicit rational specialization. This characteristic
persists throughout the admitted diagonal calibration and therefore cannot
be removed by another small change of its homogeneous diagonal boost values.
The finite packet and scalar minimization theorems show why a general
uniform elliptic or residual-to-observation argument is unavailable even
though the active finite matrices are positive.

The result does not reopen the signed active $2/3$ problem resolved by V12.
It distinguishes that issue from positive source-rank loss. Nor does it
establish a new full higher-weak-source branch. The actual scalar right side
has a proved exponentially small pairing with the localized characteristic
tests; its full observed inverse remains the next target. A useful upstream
extension would have to change this flux geometry or prove a source-specific
range/transport property, not merely enlarge the positive spectral interval
of the already controlled transverse metric. All nonlinear response,
physical reconstruction, differentiated-time estimates, and common-lifespan
requirements remain separate. No raw-action Einstein continuum theorem or
cofinal counterexample to it is asserted.

\paragraph{Verification scope.}
The accompanying symbolic verifier proves the general characteristic
factorization, its simple root, the exact unit coefficients, source rank two,
and the nonzero shift. It checks the positive Schur polynomial, a rational
noncommuting example of the exact flux projection and metric inverse, the
Neumann constants, and integer support margins. A separate full-frequency
implementation of the actual first source verifies physical adjoints,
pointwise coefficient inversion, and both routes to the flux residual on
random arrays. It is also compared at small grids with the unchanged original
scalar coefficient extractor. Floating packet diagnostics retain the full
phase symbols and literal shifts, and test exact scalar minimization at a
specified finite grid. They are not rank certificates, bounds on the actual
solved lapse, or extrapolations used in the analytical proofs.

The source audit inspected all thirteen preceding bridge bodies and the
relevant original-action and monograph ranges. V13's normalized transverse
calculus and V11's localization lemma are explicit inherited inputs. No
large inherited phase-root enclosure, higher weak matrix, or nonlinear
spacetime history is represented as recomputed. The first thirteen marked
bodies remain byte-for-byte unchanged, including their corrections and
limitations; the preservation manifest records their hashes.
% END IMMUTABLE EXACT-ACTION BRIDGE V14 BODY

% Future iterations append here.
% BEGIN IMMUTABLE EXACT-ACTION BRIDGE V15 BODY
\clearpage
\renewcommand{\thesection}{\arabic{section}}
\renewcommand{\theHsection}{bridgeV15.\arabic{section}}
\setcounter{section}{117}
\section{The actual scalar forcing, rather than another characteristic test}
\label{v15:context}

The visible tests of \cref{v14:scalar-theorem} do not decide the inverse
response of the original prepared source. This body returns to that source.
On the complementary family approaching unit configuration amplitudes, its
raw lapse row has no term of order $h$, but its effective scalar row does.
We compute the complete leading transverse coefficient and prove that its
feedback through the actual transverse inverse produces a nonzero scalar
forcing. A rational Fourier residual certificate establishes that nonvanishing
for the limiting coefficient, without solving or assuming invertibility of
the characteristic scalar operator. Uniform consistency of the already
controlled transverse inverse transfers it to every sufficiently fine cutoff.

This gives a source-specific lower bound of order $h^{-1}$ for the lapse
component of the original leading \emph{active} solve, even in any fixed
negative Sobolev norm. It is not a claim about an unprepared full multiplier
trajectory. There is also a positive inverse estimate: the exponentially
small high-frequency part of the actual right side remains negligible after
the \emph{entire} scalar inverse and geometric observation. Only
$O((\log N)^3)$ low-frequency source coordinates are needed to approximate
that output to any prescribed algebraic accuracy. The solution is not
frequency-truncated, and the remaining low-source response is not bounded
by calling the source finite-dimensional.

\paragraph{Dependencies and what is new.}
We retain the complementary data and original initializer of
\cref{v12:family,v12:exact-preparation}, the correctly ordered prepared
coefficient \eqref{v5:prepared-drift-formula}, the normalized transverse
inverse and weighted bounds of V13, and the geometric observation of V10.
The scalar equation is precisely \eqref{v13:S-g}; its flux interpretation
is V14. The new items are the complete leading-source table, its certified
nonzero scalar feedback, the resulting cofinal lower bound for this actual
source, and the inverse-level source approximation. No higher weak coefficient,
nonlinear rate selection, or physical lifespan is inferred. The continuum
operators below are limits of specified coefficients, not replacements for
the finite action or the metric reconstruction.

\section{A complete leading coefficient for the prepared source}
\label{v15:leading-source}

Let $\alpha_h=(1,\theta_h,1)$ with $1<\theta_h<1+h$, and use the
\emph{recalibrated} complementary first jets of \cref{v12:family}.
The source results in this section need no rank hypothesis and hold for
any such choices; the inverse applications later use the rational choices
of \cref{v2:rational-selection}. Put $\lambda_*=2\pi$ and
$\vartheta_i=2\pi x_i$. In this body $\partial_i^\circ$ denotes the
ordinary angle derivative, with Fourier symbol $\mathrm i k_i$ on
$e^{2\pi\mathrm i k\cdot x}$. Thus ordinary spatial differentiation is
$\lambda_*\partial_i^\circ$. This normalization is only for displaying
coefficient formulas. The original physical coordinates are unchanged.

For a smooth periodic function $f$, write $\mathcal S_h f$ for its samples,
and $\mathcal I_h$ for the original trigonometric interpolation. The
physical, component-major prepared target is $b_h=(b_{N,h},b_{\beta,h})$.
Define a real transverse trigonometric vector $q$ by
\begin{equation}
 q(x)=2\sum_{k\in\mathcal B_+} q_k\cos(2\pi k\cdot x),
 \label{v15:q-definition}
\end{equation}
where the nonzero coefficients are
\begin{equation}
\begin{array}{c|rrr}
 k&q_{k,1}&q_{k,2}&q_{k,3}\\\hline
 (1,0,0)&0&1&-1\\
 (0,1,0)&-253/512&0&-1\\
 (0,0,1)&-253/512&-1&0\\
 (2,0,0)&0&-1/8&-1/8\\
 (0,2,0)&-1/8&0&-1/8\\
 (0,0,2)&-1/8&-1/8&0\\
 (1,-1,0)&0&0&-1/8\\
 (1,0,1)&0&1/8&0\\
 (0,1,-1)&-1/8&0&0
\end{array}
\label{v15:source-table}
\end{equation}
All other Fourier coefficients are zero. In particular $k\cdot q_k=0$,
and $q$ has zero mean. It is a coefficient of the prepared drift, not a
chosen substitute forcing.

\begin{theorem}[The lapse term starts one order later than the transverse term]
\label{v15:target-expansion}
For every fixed real Sobolev order $s$, uniformly over the stated detunings,
\begin{equation}
 \boxed{\|b_{N,h}\|_{H_h^s}\le C_s h^2,\qquad
 \|b_{\beta,h}-h\lambda_*^4\mathcal S_hq\|_{H_h^s}
                                      \le C_s h^2.}
 \label{v15:target-estimate}
\end{equation}
These are estimates of the complete prepared coefficients of the unchanged
action on every sufficiently fine odd cutoff. In particular
$\|b_{\beta,h}\|_{H_h^s}\asymp h$.
\end{theorem}
\begin{proof}
The complete target has the fixed Fourier bank $|k|_1\le2$, as established
in \cref{v13:target}. Its geometric term is analytic in $h^2$ on this
bank and its constant coefficient vanishes after the original lapse
preparation correction. Its electric and magnetic cubic terms have a factor
$h$. Their coefficient at $h=0$ is therefore obtained with the literal
shifts replaced by their \emph{leading coefficient} $I$ and the phase
derivatives replaced by $\lambda_*\partial^\circ$. This is Taylor
extraction, not a change to a finite-cutoff equation. The dependence on
$\alpha$ is analytic on a fixed neighborhood, including the actual
anisotropic momentum calibration. Replacing $\alpha_h$ by $(1,1,1)$ in
this leading coefficient costs $O(h^2)$ in the complete target.

We specify the finite rational calculation of that coefficient. In unit
angle normalization the first records are
\[
 U=\sum_iT_i\cos\vartheta_i,\qquad
 p=\diag(16,259/64,259/64)+\tfrac12\sum_iP_i\cos\vartheta_i.
\]
Use the original first connection $z_1$, quadratic field $F_1=(V(p),-U/2)$,
and multiplier field $Y_\nu$ in \eqref{v5:linear-multiplier}, with
$\delta=\partial^\circ$. For each test $\nu$, take the three original
scalar coefficients
\begin{equation}
 -DH_3(Y)[Y_\nu],\qquad DP_{\nu,2}(Y)[F_1Y],\qquad
                                  P_{(0,\partial^\circ n),2}(Y).
 \label{v15:three-terms}
\end{equation}
Retain the coefficient of an independent symbol $\varepsilon$ multiplying
the two original cubic logarithmic remainders, with their original signs
and Frobenius factors. At this order the shifts are $I$, so a magnetic
plaquette uses exactly
$[z_i+z_j,[z_i,z_j]]/2$. The electric term is the retained
$-\varepsilon\sum_i\langle\Pi_i,[z_i,[z_i,A_0]]\rangle/2$.
The coefficient functions are those of \cref{v3:implementation} and
\eqref{v5:coefficient-functional}; the temporal auxiliary response is
retained for every multiplier test.

Test each of the four multiplier slots against $e^{-\mathrm i k\cdot
\vartheta}$, for $k=0$ and one representative of every pair $\{k,-k\}$
with $0<|k|_1\le2$. There are 13 such representatives and hence 52 scalar
calculations. Exact Fourier contraction gives zero in the lapse slot in
each of the three terms of \eqref{v15:three-terms}. In the shift slots,
the three entries are respectively $2q_k,-q_k,0$ on the rows of
\eqref{v15:source-table}; they are all zero at $k=0$ and at the three
remaining pairs $(1,1,0),(1,0,-1),(0,1,1)$. Negative representatives are
conjugates. The package supplies the full three-term rational table and
its regeneration from the original scalar functional. No numerical root,
finite-amplitude difference quotient, or modular zero proves these
identities. Every operation in this finite table is addition,
multiplication, and ordinary differentiation in a rational Fourier ring.

The support bound exhausts all possible target coefficients, so these
tests identify the complete vector. Restoring physical derivatives and
the momentum scale contributes $\lambda_*^4$: spatial connections have
scale $\lambda_*$, a shift-generated configuration variation contributes
one derivative, and a lapse-generated momentum variation contributes two.
Equivalently the simultaneous scaling $\delta\mapsto\lambda_*\partial^\circ$,
$p\mapsto\lambda_*p$ in \eqref{v15:three-terms} gives that factor in
both multiplier slots. Taylor's theorem on the fixed bank, including the
$O(h)$ detuning, proves \eqref{v15:target-estimate}. Norms on this bank
are uniformly equivalent. Since the table is nonzero, the two-sided
transverse size follows. The assertion is deliberately for sufficiently
fine grids, not a redefinition of the exceptional three-site folding.
\end{proof}

The cancellation in the first estimate does not say that the effective
scalar source $g_h$ is $O(h^2)$. That source includes the original
transverse elimination, which is evaluated next.

\section{The limiting transverse problem fixes the effective forcing}
\label{v15:limiting-transverse}

All operators in this section act on the ordinary unit three-torus with
normalized positive pairing, using angle derivatives. Define the real
coefficient vectors
\begin{align}
 b^\circ&=(253/64+(\cos\vartheta_2+\cos\vartheta_3)/2,
          -8+(\cos\vartheta_1-\cos\vartheta_3)/2,
          -8-(\cos\vartheta_1+\cos\vartheta_2)/2),\notag\\
 r^\circ&=\tfrac12(\sin\vartheta_2-\sin\vartheta_3,
                   \sin\vartheta_3-\sin\vartheta_1,
                   \sin\vartheta_1-\sin\vartheta_2).
 \label{v15:limit-coefficients}
\end{align}
Let $C^\circ=\curl_{\partial^\circ}$,
$\Omega^\circ=(-\sum_i(\partial_i^\circ)^2)^{1/2}$, and let $P^\circ$
be the mean-zero transverse Fourier projection. Put $Q^\circ=I-P^\circ$.
The matrix coefficient $M^\circ$ and the two raw off-diagonal sources are
\begin{equation}
 \begin{split}
 (M^\circ w)_{ij}&=b_i^\circ w_j+b_j^\circ w_i,\\
 (\overline T\beta)_{ij}&=\tfrac14(M^\circ C^\circ\beta)_{ij},\\
 (\overline Vn)_{ij}&=-\tfrac12(r_i^\circ\partial_j^\circ n+
                              r_j^\circ\partial_i^\circ n),\\
 \overline B&=\overline T^*\overline T
                    \quad\hbox{on }\Ran P^\circ.
 \end{split}
 \label{v15:bar-operators}
\end{equation}
Adjoints of the raw off-diagonal sources count both Frobenius occurrences;
$M^\circ$ itself uses the ordinary three-vector pairing.

Define the bounded positive operator
\begin{equation}
 H^\circ=P^\circ(M^\circ)^*M^\circ P^\circ+Q^\circ.
 \label{v15:Hcirc}
\end{equation}
The explicit coefficient comparison of V13 applies at $h=0$ as well:
\begin{equation}
             \tfrac9{16}I\preceq H^\circ\preceq324I.
 \label{v15:Hcirc-bounds}
\end{equation}
It is invertible on the Hilbert space, for example by the norm-convergent
Neumann series after scaling by the midpoint of this interval. Its Fourier
bandwidth is two. The polynomial inverse-decay proof of V13 also applies
on $\mathbb Z^3$; it uses only a bounded self-adjoint operator with that
spectral interval. The same summable weighted block bounds therefore give
all fixed Sobolev and sufficiently small exponential Fourier bounds.
These facts concern a controlled \emph{transverse} normal operator,
not the characteristic scalar inverse.

Set
\begin{equation}
 f^\circ=C^\circ(\Omega^\circ)^{-2}q,\quad
 w^\circ=(H^\circ)^{-1}f^\circ,\quad
 \overline\beta=8C^\circ(\Omega^\circ)^{-2}w^\circ,\quad
 g^\circ=2\overline V^*M^\circ w^\circ.
 \label{v15:g-definition}
\end{equation}
Inverse powers are zero on constants. The source $f^\circ$ is transverse;
so is $w^\circ$, because $H^\circ$ commutes with $P^\circ$. These
functions have an exponential Fourier bound. Exact multiplication gives
\begin{equation}
 \overline\beta=\overline B^{-1}q,\qquad
 \overline T\overline\beta=2M^\circ w^\circ,\qquad
                  g^\circ=\overline V^*\overline T\overline B^{-1}q.
 \label{v15:g-feedback}
\end{equation}
For example, $\overline B=(C^\circ(M^\circ)^*M^\circ C^\circ)/8$
on the transverse slice; invert its two curl factors in their actual order.
This proves \eqref{v15:g-feedback} without assuming they commute with $M^\circ$.

\begin{theorem}[Full effective scalar source expansion]
\label{v15:g-expansion}
For $g_h$ defined by the original \eqref{v13:S-g}, and the detuned first
jets used in \cref{v15:target-expansion}, at every fixed Sobolev order,
\begin{equation}
 \boxed{\|g_h-h\lambda_*^4\mathcal S_hg^\circ\|_{H_h^s}
                                                  \le C_s h^2.}
 \label{v15:g-limit}
\end{equation}
There is no scalar inverse in this assertion.
\end{theorem}
\begin{proof}
Normalize the original maps by
$\overline T_h=T_h/(h\lambda_*^2)$ and
$\overline V_h=V_h/(h\lambda_*^2)$, and set
$\overline B_h=\overline T_h^*\overline T_h$ on the original transverse
slice. V13 gives uniform bounds for $\overline B_h^{-1}:H_h^j\to H_h^{j+2}$.
The right side
$q_h=b_{\beta,h}/(h\lambda_*^4)$ differs from $\mathcal S_hq$ by $O(h)$
in every fixed Sobolev norm, by \cref{v15:target-expansion}.

We detail the consistency needed through this inverse. On sampled smooth
functions, $\delta_h/\lambda_*$ differs from $\partial^\circ$ by $O(h^2)$
with three derivatives charged. A one-site shift differs from the identity
by $O(h)$ with one derivative charged. The coefficient vectors converge
to \eqref{v15:limit-coefficients} at order $h$, including the detuning.
The original transverse Fourier projection differs from the ordinary one
on smooth samples at order $h^2$, with a fixed finite derivative loss. To
check the latter, at a nonzero frequency use
$P_T(k)=I-\kappa_h(k)\kappa_h(k)^T/|\kappa_h(k)|^2$ and
$P^\circ(k)=I-kk^T/|k|^2$; sine Taylor expansion gives an $O(h^2|k|^2)$
bound in the interior. The complementary frequencies and sampling aliases
are bounded by the smooth tail, exponentially small for the functions in
\eqref{v15:g-definition}. No projection is applied to an unknown solution.

Insert $P_T\mathcal S_h\overline\beta$ in the exact normalized transverse
equation. The preceding coefficient and projection estimates show that its
residual from $q_h$ is $O(h)$ in every fixed order. The proved uniform
inverse of $\overline B_h$ then gives
\[
 hB_h^{-1}b_{\beta,h}
          =P_T\mathcal S_h\overline\beta+O_{H_h^j}(h)
          =\mathcal S_h\overline\beta+O_{H_h^j}(h).
\]
Adjoints of $\overline T_h,\overline V_h$ obey the same smooth consistency
and uniform first-order Sobolev bounds, by their exact shifted-coefficient
formulas. Consequently
\[
 \frac{V_h^*T_hB_h^{-1}b_{\beta,h}}{h\lambda_*^4}
              =\mathcal S_h(\overline V^*\overline T\overline\beta)
                    +O_{H_h^s}(h).
\]
The omitted term here is the \emph{retained} $-b_{N,h}$ in the definition
of $g_h$; its normalized size is $O(h)$ by
\eqref{v15:target-estimate}. This proves \eqref{v15:g-limit}. The argument
uses the already controlled transverse inverse only; it assumes neither
solvability nor regularity for a limiting scalar equation.
\end{proof}

\section{A rational certificate of nonzero scalar feedback}
\label{v15:feedback-certificate}

Let $\psi(x)=\sin(\vartheta_2+\vartheta_3)$. In raw off-diagonal
coordinates put
\begin{equation}
                     j^\circ=4(M^\circ)^*\overline V\psi.
 \label{v15:j}
\end{equation}
The factor four comprises the factor two in \eqref{v15:g-definition}
and the two Frobenius occurrences. Thus
$\langle\psi,g^\circ\rangle=\langle j^\circ,w^\circ\rangle$ in the
ordinary vector pairing.

\begin{theorem}[Verified nonzero leading effective source]
\label{v15:nonzero-g}
For the actual coefficient functions in \eqref{v15:g-definition},
\begin{equation}
       \boxed{\frac{151}{1000}
          <\ip\psi{g^\circ}<\frac{153}{1000}.}
 \label{v15:g-pairing}
\end{equation}
In particular the original scalar right side satisfies
$\|g_h\|_{H_h^s}\asymp h$ for every fixed real $s$, on all sufficiently
fine cutoffs of the stated family. The leading forcing comes from transverse
feedback, even though the raw lapse drift is $O(h^2)$.
\end{theorem}
\begin{proof}
We give an a posteriori rational certificate on the \emph{continuum
coefficient} Hilbert space. It does not infer a continuum statement from a
finite-grid inverse. The supplied file \src{force_trial_L8.json} specifies
an odd, purely imaginary Fourier vector with dyadic coefficients of
denominator $2^{44}$ and support $0<|k|_1\le8$. Apply the exact rational
matrix $P^\circ(k)=I-kk^T/|k|^2$ at each of these frequencies to obtain
a real mean-zero transverse trigonometric trial $\widehat w$.
The trial was proposed by floating arithmetic; only the following exact
residual verifies it.

The coefficient vectors in \eqref{v15:limit-coefficients} are rational
Fourier polynomials. The vector $f^\circ$ has coefficients
$\mathrm i(k\times q_k)/|k|^2$, also rational multiples of $\mathrm i$.
Compute, over noncyclic finite Fourier support,
\[
 e=f^\circ-P^\circ(M^\circ)^*M^\circ\widehat w.
\]
This residual has support $|k|_1\le10$. The exact rational computation
in the package gives
\begin{equation}
 \begin{gathered}
 \ip{j^\circ}{\widehat w}
       =\frac{1026718915432535}{6755399441055744},\\
 \|j^\circ\|^2=\frac{1112585}{4096},\\
 \|e\|^2<\frac{49}{10^{16}},\qquad
 \frac{256}{81}\|j^\circ\|^2\|e\|^2<\frac1{400000^2}.
 \end{gathered}
 \label{v15:rational-certificate}
\end{equation}
The last inequality uses the computed exact residual sum, not only the
preceding rounded bound. All Fourier convolution, Hodge projection,
pairing, and squared-norm operations use arbitrary-precision rational
numbers. The files contain the trial coefficients and complete rational
residual sums. No numerical quadrature, interval eigenvalue threshold,
or floating residual is accepted as an inclusion proof.

Since $\widehat w$ is transverse, $H^\circ\widehat w=f^\circ-e$.
The lower bound \eqref{v15:Hcirc-bounds} gives
\[
 \left|\ip{j^\circ}{w^\circ-\widehat w}\right|
       \le\frac{16}{9}\|j^\circ\|\|e\|<\frac1{400000}.
\]
The resulting rational interval around the first line of
\eqref{v15:rational-certificate} is strictly inside
$(151/1000,153/1000)$. This proves \eqref{v15:g-pairing}.
The certificate validates a residual in the bounded positive transverse
problem on the full torus; it does not invert $\mathscr S_h$ or its
characteristic limit. Finally use \cref{v15:g-expansion}, the fixed norm
of $\psi$, and duality at any fixed Sobolev order for the lower bound on
$g_h$. The upper bound follows from the same expansion.
\end{proof}

\section{The solved lapse has a necessary scale visible in smooth tests}
\label{v15:lapse-lower}

For the rational detunings of V12, the original scalar matrix is positive
and invertible. Denote its unique original-source solution by
\begin{equation}
                         n_h=\mathscr S_h^{-1}g_h.
 \label{v15:n}
\end{equation}
This is the lapse \emph{component of the leading active multiplier solve}.
It is not the background lapse, which is one at the prepared cut. Nor is
it automatically the amplitude-leading rate of the full exact continuation:
that identification still requires the higher weak rate to be controlled.

\begin{theorem}[Source-specific lower bound in every fixed negative order]
\label{v15:n-lower}
For every fixed $s\ge0$, there are cutoff-independent positive constants
$c_s,C_s$ on the selected family such that, for sufficiently small $h$,
\begin{equation}
              \boxed{\|n_h\|_{H_h^{-s}}\ge c_s h^{-1}.}
 \label{v15:negative-lower}
\end{equation}
The known coarse upper estimate is $\|n_h\|_h\le C h^{-11}$.
No matching $O(h^{-1})$ upper bound, weighted curvature bound, or full
nonlinear multiplier estimate is asserted.
\end{theorem}
\begin{proof}
For the fixed sampled smooth test $\psi_h=\mathcal S_h\psi$, the exact
scalar formula \eqref{v13:S-g} gives, for every fixed $s$,
\begin{equation}
                        \|\mathscr S_h\psi_h\|_{H_h^s}\le C_s h^2.
 \label{v15:test-S}
\end{equation}
Indeed $\mathscr D_h\psi_h=O(h^2)$ on a fixed Fourier bank, and its
adjoint on that bank gives another factor $h^2$ on fine grids. The
other term is $V_h^*(I-\Pi_{T,h})V_h\psi_h/2$. The inner source is
$O(h)$ in every fixed Sobolev order; V13 bounds $\Pi_{T,h}$ on each
of those spaces, and $V_h^*$ is first order with factor $h$.
This gives \eqref{v15:test-S} for the full resulting vector, including
any nonlocal Fourier tail of the projection.

By \cref{v15:g-expansion,v15:nonzero-g},
$\langle\psi_h,g_h\rangle\ge c h$. Self-adjointness and physical
Sobolev duality therefore imply
\[
 ch\le |\ip{\mathscr S_h\psi_h}{n_h}_h|
        \le C_sh^2\|n_h\|_{H_h^{-s}},
\]
which proves \eqref{v15:negative-lower}. The existing diagonal detuning
and the positive scalar energy give
$\mathscr S_h\succeq c h^{12}I$. Applying that inverse bound to
$\|g_h\|=O(h)$ gives the stated coarse upper bound. These are estimates
of the unique source-dependent solve; no constructed packet is substituted
for $g_h$.
\end{proof}

For completeness, the natural normalized coefficient equation is fixed by
these formulas. Define $\overline\Pi$ as the positive projection onto
$\Ran\overline T$ and
\[
 \mathscr S^\circ=\tfrac12\overline V^*(I-\overline\Pi)\overline V.
\]
On each fixed smooth test,
$\mathscr S_h\mathcal S_h\varphi/(h^2\lambda_*^4)
=\mathcal S_h\mathscr S^\circ\varphi+O(h)$ in every fixed order.
The diagonal term is $O(h^4)$ before this normalization; the remaining
consistency follows through the already controlled transverse projection.
Consequently, if $\mathcal I_h(hn_h)$ is bounded and converges weakly in
some fixed negative Sobolev space along a subsequence, its limit $n^\circ$
obeys, in the distributional test sense,
\begin{equation}
                         \mathscr S^\circ n^\circ=g^\circ.
 \label{v15:limiting-equation}
\end{equation}
One may prove this by testing the exact equation against a fixed smooth
$\varphi$, moving $\mathscr S_h$ by self-adjointness, and using its
smooth-test consistency and \eqref{v15:g-limit}. The bounded negative norm
pays the consistency remainder. Equation \eqref{v15:g-pairing} makes any
such limit nonzero. This is a necessary limiting equation, not a proof
that a bounded sequence or a regular solution of this characteristic
coefficient problem exists.

\section{High source frequencies remain small through the full scalar inverse}
\label{v15:source-tail}

There is a useful positive estimate which V14's individual packet pairings
did not provide. Let $P_{\le R}$ be the scalar Fourier projection onto
$|k|_1\le R$ and $P_{>R}=I-P_{\le R}$, on the mean-zero scalar space.
Let the actual lapse-to-curvature operator be
\[
                         \mathcal O_h=\mathfrak E_hT_h^\dagger V_h.
\]
It retains the full variable-coefficient pseudoinverse from V13.

\begin{theorem}[Inverse-level approximation using only logarithmically many source modes]
\label{v15:inverse-tail}
Choose one fixed positive analytic weight $\sigma$ for which V13 bounds
the actual $g_h$. For every fixed $s\ge0$ and $R\ge0$, on the original
selected positive scalar family,
\begin{align}
 \|\mathscr S_h^{-1}P_{>R}g_h\|_h
                       &\le C h^{-11}e^{-\sigma R},\label{v15:state-tail}\\
 \boxed{\|\mathcal O_h\mathscr S_h^{-1}P_{>R}g_h\|_{H_h^s}
                       \le C_s h^{-12-s}e^{-\sigma R}.}
                                                       \label{v15:observed-tail}
\end{align}
In particular, for any prescribed $p>0$ choose
\begin{equation}
 R_h=\left\lceil\frac{p+12+s}{\sigma}\log(1/h)\right\rceil.
 \label{v15:R}
\end{equation}
Then
\begin{equation}
 \boxed{\|\mathcal O_hn_h-
     \mathcal O_h\mathscr S_h^{-1}P_{\le R_h}g_h\|_{H_h^s}
                                                        \le C_s h^p.}
 \label{v15:compressed-output}
\end{equation}
For all sufficiently fine cutoffs the retained real source dimension is
\begin{equation}
                  m_{R_h}=\frac{4R_h^3+6R_h^2+8R_h}{3}
                               =O((\log N)^3).
 \label{v15:dimension}
\end{equation}
No solution mode, operator entry, or physical observation is discarded.
\end{theorem}
\begin{proof}
The actual source has $\|g_h\|_{0,\sigma,h}\le Ch$ by V13. Thus
$\|P_{>R}g_h\|_h\le Ch e^{-\sigma R}$. Positivity and the already
proved scalar lower bound give
$\|\mathscr S_h^{-1}\|_{2\to2}\le Ch^{-12}$; this establishes
\eqref{v15:state-tail} \emph{after the full inverse}. From V13 and the
full-frequency observation estimate of V10,
\[
 \|\mathcal O_hu\|_{H_h^s}
       \le C_sh^2\|u\|_{H_h^{s+3}}
       \le C_sh^{-1-s}\|u\|_h.
\]
The last step is the explicit finite-frequency bound $|k|\le C/h$;
it is an upper estimate, not an assumption that $u$ is smooth. Combining
it with \eqref{v15:state-tail} proves \eqref{v15:observed-tail}.
Substitution of \eqref{v15:R} yields \eqref{v15:compressed-output} by
linearity of the \emph{unchanged} inverse.

Eventually $R_h<(N-1)/2$, so its frequency ball is unwrapped. The number
of lattice points in the three-dimensional integer $\ell^1$ ball of
radius $R$ is $(4R^3+6R^2+8R+3)/3$, by summing its coordinate slices;
subtract the constant mode. Conjugate pairs give the same real dimension.
This proves \eqref{v15:dimension}. Before that asymptotic regime one may
retain the whole available source space; no claim of useful small-grid
computational savings follows from the asymptotic count.
\end{proof}

This theorem does more than pair $g_h$ with a particular localized
characteristic test: it controls the inverse image of its complete high
\emph{source} tail. It does not control the high \emph{solution} tail of
$\mathscr S_h^{-1}P_{\le R_h}g_h$. In particular it does not justify
replacing the scalar matrix by its low-frequency principal block. All
resonant scattering from the retained low source is still present.

\section{An exact low-source interface that does not hide a high-mode deletion}
\label{v15:low-interface}

The finite source approximation has an exact Schur representation. This
ordinary algebra is recorded only to state precisely what a subsequent
verified reduced calculation would have to keep. Put $P=P_{\le R}$,
$Q=I-P$ on the mean-zero scalar space and write
\[
 A=P\mathscr S_hP,\quad C=Q\mathscr S_hP,\quad D=Q\mathscr S_hQ.
\]
Since $\mathscr S_h\succeq ch^{12}I$, the block $D$ on $\Ran Q$ is
invertible. Define
\begin{equation}
 K_R=A-C^*D^{-1}C,\qquad
 Z_Rx=x-D^{-1}Cx\quad(x\in\Ran P).
 \label{v15:low-Schur}
\end{equation}
Both spaces in this formula are the actual finite Fourier spaces, and
$K_R\succeq ch^{12}I$ on $\Ran P$. Direct substitution in the two
block equations proves
\begin{equation}
 \boxed{\mathcal O_h\mathscr S_h^{-1}P g_h
                   =\mathcal O_h Z_R K_R^{-1}P g_h.}
 \label{v15:low-exact}
\end{equation}
For its positivity, $\langle x,K_Rx\rangle$ is the minimum of
$\langle x+y,\mathscr S_h(x+y)\rangle$ over $y\in\Ran Q$; the
minimum is attained at $y=-D^{-1}Cx$ and is at least
$ch^{12}\|x\|^2$. This also verifies the formula without any choice of
small-eigenvalue threshold.

In \eqref{v15:low-exact} the input and final scalar Schur matrix have only
$m_R$ coordinates, but $D^{-1}$ and the observation of its lifted vectors
are retained. They may be expensive and remain mathematically nontrivial.
A low-dimensional source is not a proof of a well-conditioned response.
Conversely a verified calculation of these low-source responses, together
with \cref{v15:inverse-tail}, would control the original output without
requiring a uniform estimate for every scalar residual, which V14 excludes.
No complexity or physical execution-deadline claim is made here.

\section{Consequences, verification, and the remaining range problem}
\label{v15:status}

The source cannot be declared one order smaller merely because its raw lapse
row is. Its leading scalar coefficient is now specified through a uniformly
controlled transverse problem and certified nonzero in rational arithmetic.
The cofinal positive scalar solve therefore has a necessary $h^{-1}$ scale
in every fixed negative Sobolev norm. This lower bound is compatible with an
$O(h)$ geometric output if a suitable $O(h^{-1})$ \emph{upper} estimate
with the required spatial regularity can be proved; it supplies no such
upper estimate. The bound concerns the active subsystem and is not a new
pole theorem for the full raw trajectory before its higher weak rates have
been controlled.

There is also a completed inverse-level estimate for the actual high source
tail. It reduces the outstanding geometric response, to arbitrary algebraic
accuracy, to the image of an explicit logarithmic source bank through the
full scalar inverse. It does not eliminate propagation from that bank into
the characteristic directions. The unresolved question is now whether the
specific nonzero $g^\circ$, and its finite corrections, lie in a sufficiently
regular range of the positive flux operator with the diagonal defects
retained. A uniform elliptic assertion cannot replace that range argument.
No new cofinal higher-weak rank, weak-rate cancellation, finite-amplitude
inverse bound, propagated initial condition, or common physical lifespan
is established. The exact-action Einstein continuum bridge is not closed.

\paragraph{Verification scope.}
The symbolic source program regenerates all 52 leading test coefficients
from the retained scalar functional in an exact rational Fourier ring. It
keeps the prepared lapse correction and both cubic logarithmic terms until
after their contraction. A separate arbitrary-precision rational verifier
reads a dyadic trigonometric trial, projects it exactly, and computes its
full noncyclic residual in $H^\circ$. Its norm certificate proves
\eqref{v15:g-pairing} on the continuum coefficient Hilbert space, rather
than treating a sampled numerical inverse as that proof. The finite trial
and every rational inclusion inequality are supplied.

Separately labeled floating diagnostics use the unchanged original prepared
coefficient evaluator at $N=9,17,25,33$, with $\alpha=(1,1+h/2,1)$.
They compare the leading source table, its $O(h^2)$ lapse remainder, and
the normalized scalar feedback. These amplitudes are probes of the source
formula, not claimed to satisfy the special diagonal detuning test. They
do not solve the characteristic scalar equation or supply a new signed-rank
certificate. Exact rational tests also check the source-tail exponent and
the block elimination, and integer enumeration verifies the source dimension.
The fourteen preceding research bodies remain byte-for-byte unchanged.
Their analytical initial preparation and the V13 inverse estimates are
inherited inputs, not independently formalized or reproved by these checks.
% END IMMUTABLE EXACT-ACTION BRIDGE V15 BODY


% Future iterations append here.
% BEGIN IMMUTABLE EXACT-ACTION BRIDGE V16 BODY
\clearpage
\renewcommand{\thesection}{\arabic{section}}
\renewcommand{\theHsection}{bridgeV16.\arabic{section}}
\setcounter{section}{125}
\section{A source-excited singular sector of the diagonal regulator}
\label{v16:context}

The scalar forcing in V15 is not zero, and its high source tail has already
been controlled after the full inverse. Repeating either estimate would not
resolve the response of the remaining low source. This body instead examines
an upstream feature of the positive scalar energy: the small eigenvalues of
its \emph{diagonal product-defect term}. Their magnitudes were determined in
V2, but their limiting functions, their incidence with the actual forcing,
and the effect of the retained flux on those functions were not.

The results distinguish three statements. First, the folded scalar inverse
has a rank-two operator limit whose range consists of a square wave and a
logarithmic function. These are genuine distributional null solutions of the
limiting scalar differential expression, although neither is smooth. We give
the exact smooth-source compatibility conditions. Second, an exact rational
certificate proves that the original effective forcing has a nonzero pairing
with a tensor product of these modes. If the scalar flux term were deleted,
that component of the inverse would have order $h^{-5}$, with an explicit
nonzero normalized limit. Third, the \emph{retained} flux excludes every
nonzero vector of this four-dimensional limiting sector from its finite-energy
domain. Its normalized energy on the corresponding finite sector tends to
infinity. Thus neither deleting the flux nor identifying the bare diagonal
soft modes with eigenmodes of the full scalar operator is legitimate.

These are results for the actual scalar suboperator and the actual prepared
forcing, not a substituted field equation. The diagonal-only inverse below
is explicitly a comparison used to test an approximation; the original scalar
equation always retains both terms. No bound on its full low-source inverse,
higher weak source, nonlinear evolution, or physical lifespan is claimed.

\paragraph{Dependencies and novelty boundary.}
We use the exact cosine Jacobi representation, inverse decomposition and
soft eigenvalue asymptotic of \cref{v2:Jacobi-theorem,v2:soft-pair}; those
spectral magnitudes are not new results. The new conclusions identify their
physical Fourier profiles and source-dependent operator limits. The original
complementary data, coefficient inverses and scalar energy are
\cref{v12:family,v13:Schur,v14:flux-theorem}. The complete leading target,
its uniform consistency and its dyadic transverse trial are inherited from
\cref{v15:target-expansion,v15:g-expansion,v15:nonzero-g}. We recompute a
different, nonsmooth source pairing with a new exact residual estimate. The
source table itself is an inherited coefficient calculation, not represented
as independently regenerated here. Singular domains of Jacobi operators
provide general context \cite{v16:Yafaev}; the finite-fold limits and all
source-specific statements used here have the direct proofs below. No
classification of self-adjoint extensions is invoked or claimed.

\section{The limiting differential expression and its smooth range}
\label{v16:range}

In this body $D_h$ denotes the \emph{one-dimensional cosine product-defect}
operator of \eqref{v2:D}, not the ordinary spectral derivative of V10. Use
$x\in\T=\R/\mathbb Z$, $\theta=2\pi x$, and normalized measure
$dx=d\theta/(2\pi)$. Put
\[
 f(x)=\cos(2\pi x),\qquad
 D_hu=\delta_h(fu)-f\delta_hu-(\delta_hf)u.
\]
Let $\mathcal I_h$ be trigonometric interpolation from the odd grid and let
$\mathcal I_h^*$ be its Hilbert adjoint. Thus
$\mathcal I_h\mathcal I_h^*=P_m$, the orthogonal Fourier projection onto
$|k|\le m=(N-1)/2$. For an arbitrary $L^2$ function this adjoint is Fourier
restriction followed by synthesis, \emph{not} evaluation at singular points.
It is used to define tests, not to alter any initial-record read.

\begin{lemma}[Smooth-test limit of the original diagonal coefficient]
\label{v16:smooth-limit}
For each smooth periodic $u$,
\begin{equation}
 h^{-2}\mathcal I_hD_h\mathcal I_h^*u\longrightarrow\mathcal L u,
 \qquad
 \boxed{\mathcal L u=\tfrac18\partial_x(f'\partial_xu)
       =-\pi^3\partial_\theta(\sin\theta\,\partial_\theta u).}
 \label{v16:L}
\end{equation}
The error in any fixed Sobolev norm is $O_u(h^2)$. All cyclic face terms
of the finite operator are retained in this assertion.
\end{lemma}
\begin{proof}
On nonfolded modes the sine expansion gives
$\delta_h=\partial_x+h^2\partial_x^3/24+O(h^4\partial_x^5)$.
The first-order product rule cancels its order-zero defect. The next
coefficient is
\[
 \tfrac1{24}\{\partial_x^3(fu)-f\partial_x^3u-f'''u\}
  =\tfrac18(f'u''+f''u')=\tfrac18\partial_x(f'u').
\]
Changing to $\theta$ gives \eqref{v16:L}. The remaining nonfolded symbols
have a polynomial fifth-order bound times $h^2$ after normalization.
A folded entry of $D_h/h^2$ is at most $Ch^{-3}$ and can use only an outer
Fourier coefficient of $u$, since $f$ has unit frequency. Smooth Fourier
tails absorb this factor and any fixed output weight to arbitrary algebraic
order. The omitted ordinary output modes obey the same tail estimate. For
example sufficiently many fixed derivatives of $u$ give an $O(h^2)$ bound
without discarding a face entry. This proves the assertion on each fixed
smooth test.
\end{proof}

Define two real $L^2$ functions, with their values at the singular points
irrelevant, by
\begin{equation}
 \boxed{s(\theta)=\operatorname{sgn}(\sin\theta),\qquad
 l(\theta)=\frac2\pi\log\left|\tan\frac\theta2\right|.}
 \label{v16:sl}
\end{equation}
Their mean-zero Fourier expansions are
\begin{equation}
 s=\frac4\pi\sum_{\substack{k\ge1\\k\ \mathrm{odd}}}\frac{\sin k\theta}{k},
 \qquad
 l=-\frac4\pi\sum_{\substack{k\ge1\\k\ \mathrm{odd}}}\frac{\cos k\theta}{k}.
 \label{v16:sl-Fourier}
\end{equation}
They are orthogonal and both have norm one. These equalities follow from
Parseval and $\sum_{k>0,\,k\ \mathrm{odd}}k^{-2}=\pi^2/8$. In particular
neither function belongs to $H^{1/2}(\T)$.

\begin{theorem}[Two additional distributional null modes and exact smooth compatibility]
\label{v16:range-theorem}
The $L^2$ distributional nullspace of \eqref{v16:L} is
\begin{equation}
             \ker_{L^2,\mathrm{dist}}\mathcal L
                        =\operatorname{span}\{1,s,l\}.
 \label{v16:dist-kernel}
\end{equation}
The smooth nullspace consists only of constants. For a smooth mean-zero
$F$, there is a smooth mean-zero periodic solution of $\mathcal Lu=F$ if
and only if
\begin{equation}
                \boxed{\ip{s}{F}=0,\qquad\ip{l}{F}=0.}
 \label{v16:compat}
\end{equation}
When it exists this smooth solution is unique.
\end{theorem}
\begin{proof}
The distributional equation $\mathcal Lu=0$ says that
$\sin\theta\,u'$ is constant. On each of the two open intervals between
$0$ and $\pi$, integration gives a multiple of
$\log|\tan(\theta/2)|$ plus a constant. The flux constant is the same on
both intervals. Their two constants may differ: a jump contributes a delta
to $u'$, killed by the simple zero of $\sin\theta$. The logarithm has
principal-value derivative $\operatorname{pv}(1/\sin\theta)$, whose product
with $\sin\theta$ is one. These functions are in $L^2$. An $L^2$ solution
cannot have an additional distribution supported at either point. Thus all
solutions have the form in \eqref{v16:dist-kernel}. For a smooth solution the
flux constant must be zero at a zero of $\sin\theta$, and continuity joins
the two interval constants; only a global constant remains.

For the inhomogeneous equation put $A(\theta)=\int_0^\theta F(t)\,dt$.
Mean zero gives $A(2\pi)=A(0)=0$. Integration yields
\[
 \sin\theta\,u'=-\pi^{-3}A(\theta)+c.
\]
Smoothness at zero forces $c=0$; smoothness at $\pi$ then requires
$A(\pi)=0$. Conversely those zeros make $A/\sin\theta$ smooth and
periodic. The remaining condition for its antiderivative to be periodic is
$\int_0^{2\pi}A/\sin\theta=0$. For a mean-zero $F$,
$\ip{s}{F}=A(\pi)/\pi$. When $A(\pi)=0$, integration by parts gives
\[
 \ip{l}{F}=-\frac1{\pi^2}\int_0^{2\pi}\frac{A(\theta)}{\sin\theta}\,d\theta.
\]
All boundary terms are zero because $A$ vanishes there and the logarithm
is locally integrable. Hence \eqref{v16:compat} is necessary and sufficient.
Integrate $u'=-\pi^{-3}A/\sin\theta$ and fix its additive mean to prove
existence and uniqueness.
\end{proof}

This is a domain statement for a specific differential expression. It does
not claim that a nonsmooth null function belongs to every self-adjoint
realization, nor that smooth-test convergence determines a unique inverse.
The finite inverse below identifies what the retained fold actually selects.

\section{A rank-two operator limit of the folded inverse}
\label{v16:green}

Let $D_h^\dagger$ vanish on constants and invert $D_h$ on their orthogonal
complement. On $L^2(\T)$ define
\[
 \mathbb G_h=h^3\mathcal I_hD_h^\dagger\mathcal I_h^*,\qquad
 E=\operatorname{proj}_{\operatorname{span}\{s,l\}},\qquad
 K F=s\ip{l}{F}+l\ip{s}{F}.
\]
Thus $K$ is self-adjoint, $K^2=E$, and its two nonzero eigenvalues are
$+1,-1$.

\begin{theorem}[The actual fold selects the square-wave/logarithmic pair]
\label{v16:inverse-limit}
Along $N\equiv1\pmod4$,
\begin{equation}
       \boxed{\left\|\mathbb G_h-\frac1{2\pi^2}K\right\|_{L^2\to L^2}
                                                     \le C h^{1/2}.}
 \label{v16:rank-two-limit}
\end{equation}
Let $E_h$ project onto the two scalar soft eigenvalues of
\cref{v2:soft-pair}, and write their common absolute value as
$\varepsilon_h=2\pi^2h^3+O(h^4)$. With interpolation and zero extension,
\begin{equation}
 \mathcal I_hE_h\mathcal I_h^*\longrightarrow E,\qquad
 \varepsilon_h^{-1}\mathcal I_hD_hE_h\mathcal I_h^*\longrightarrow K
                       \quad\hbox{in operator norm}.       \label{v16:projectors}
\end{equation}
These limits identify physical Fourier modes, not only eigenvalue magnitudes.
\end{theorem}
\begin{proof}
Write $N=2m+1$ with $m=2j$ even. Use the inherited exact weights
\[
 c_h=4h^{-1}\sin(\pi h/2),\quad a_k=\sin(k\pi h/2),\quad
 t_m=2h^{-1}\cos(\pi h/2)+h^{-1}\sin(\pi h).
\]
Set
\[
 Z_h=\sum_{\substack{1\le k<m\\k\ \mathrm{odd}}}a_k^{-2},\qquad
 s_h=\sqrt{2/Z_h}\sum_{\substack{1\le k<m\\k\ \mathrm{odd}}}a_k^{-1}\sin k\theta,
 \qquad
 l_h=-\sqrt{2/Z_h}\sum_{\substack{1\le k<m\\k\ \mathrm{odd}}}a_k^{-1}\cos k\theta.
\]
These are exactly orthonormal finite trigonometric functions. Undo the
Jacobi phases in V2's inverse decomposition. In the ordered physical modes
$1,\ldots,m,-m,\ldots,-1$, the lower adjacent entries of $D_h$ are
$-\mathrm i t_j$, and the conjugating phases are
$(-\mathrm i)^{j-1}$. Its exceptional inverse entries connect positive
odd modes $k$ to negative odd modes $-r$ and are
\[
       \frac{\mathrm i t_m}{c_h^2a_m^2a_ka_r}.
\]
The reverse entries are their conjugates. The real form of this rank-two
part is therefore exactly
\begin{equation}
 \mathcal I_hD_h^\dagger\mathcal I_h^*
   =\Lambda_h(s_h\otimes l_h+l_h\otimes s_h)+\mathcal R_h,
 \quad \Lambda_h=\frac{t_mZ_h}{c_h^2a_m^2},\quad
                   \|\mathcal R_h\|\le Ch^{-2}.           \label{v16:finite-rank-two}
\end{equation}
Here $u\otimes v$ denotes $F\mapsto u\ip vF$. The remainder bound is the
already proved bidiagonal inverse remainder of V2, after an isometry; the
new sign and physical profiles follow from the displayed entries. A direct
recurrence verification is also given in the next section.

The inherited sums give $Z_h=(2h^2)^{-1}+O(h^{-1})$ and
$h^3\Lambda_h=(2\pi^2)^{-1}+O(h)$. On $0<x\le\pi/4$,
$|1/\sin x-1/x|\le Cx$. Consequently the difference between
$\sqrt{2/Z_h}/a_k$ and $4/(\pi k)$ is at most
$C(h/k+h^2k)$ for the retained odd modes. The sum of their squared errors
is $O(h)$; the omitted squared Fourier tail is also $O(h)$. Parseval gives
\[
                      \|s_h-s\|+\|l_h-l\|\le Ch^{1/2}.
\]
Insert these bounds in \eqref{v16:finite-rank-two}; its normalized remainder
is $O(h)$. This proves \eqref{v16:rank-two-limit} on the common Hilbert
space, including its complement outside the finite frequency box.

For clarity, the projection conclusion uses a gap in this \emph{normalized
inverse}, not an assumed gap for the original source. The two soft inverse
eigenvalues tend to $\pm(2\pi^2)^{-1}$, while every other eigenvalue of
$\mathbb G_h$ is $O(h)$ by V2. If $v\in\Ran E_h$, its normalized inverse
is invertible there with uniformly bounded inverse. Thus
$(I-E)v=(I-E)(\mathbb G_h-(2\pi^2)^{-1}K)
(\mathbb G_h|_{\Ran E_h})^{-1}v$. The reverse subspace bound follows in
the same way, or by equality of the two finite ranks. Hence the projections
converge in norm. On this pair $D_h^2=\varepsilon_h^2I$, so
$D_hE_h/\varepsilon_h=\varepsilon_hD_h^\dagger E_h$.
Use $\varepsilon_h/h^3\to2\pi^2$ and the first limit to prove the second
formula in \eqref{v16:projectors}.
\end{proof}

The finite inverses are still those of the original cyclic operator. The
rough functions in the limit are analytical outputs; they are not sampled
as admissible metric records or substituted into the initial-data family.

\section{Fixed smooth sources detect the fold, and compatibility removes it}
\label{v16:fixed-sources}

\begin{theorem}[A complete compatibility test for bounded normalized scalar inverses]
\label{v16:bounded-iff}
For a smooth mean-zero $F$, the following are equivalent:
\begin{enumerate}[label=(\roman*),itemsep=2pt]
 \item $F$ obeys \eqref{v16:compat}.
 \item $F=\mathcal Lu$ for a smooth mean-zero periodic $u$.
 \item $h^2\|D_h^\dagger\mathcal I_h^*F\|_h$ is bounded over all
 sufficiently fine odd cutoffs.
\end{enumerate}
When they hold,
\begin{equation}
 h^2\mathcal I_hD_h^\dagger\mathcal I_h^*F\longrightarrow u
          \quad\hbox{in }L^2,                           \label{v16:compatible-convergence}
\end{equation}
with error $O_u(h)$, and $O_u(h^2)$ on $N\equiv3\pmod4$.
If either compatibility pairing is nonzero, the norm in (iii) diverges
like $h^{-1}$ along $N\equiv1\pmod4$.
\end{theorem}
\begin{proof}
The first equivalence is \cref{v16:range-theorem}. For its smooth solution,
\cref{v16:smooth-limit} gives
$D_h\mathcal I_h^*u-h^2\mathcal I_h^*F=O_h(h^4)$ in the positive grid
norm, with an $h$-independent constant depending on finitely many derivatives
of $u$. Apply the inherited bound $\|D_h^\dagger\|\le Ch^{-3}$,
or $Ch^{-2}$ on the other parity class. The means of both finite vectors
are zero, so this proves the stated errors relative to $P_mu$; its omitted
smooth Fourier tail is smaller. This implies (iii).
Conversely boundedness in (iii) would give $\mathbb G_hF\to0$. Along
$N\equiv1\pmod4$, \cref{v16:inverse-limit} identifies that limit as
$KF/(2\pi^2)$. Orthogonality of $s,l$ forces both pairings to vanish.
If one is nonzero, the same limit gives a strictly positive limiting norm
for $h^3D_h^\dagger\mathcal I_h^*F$, which proves the final assertion.
\end{proof}

There is an explicit contrast even for the first harmonic. For a real source
with Fourier coefficients $F_1=z$, $F_{-1}=\overline z$ and all others zero,
let $u_k$ be the positive Fourier coefficients of $D_h^\dagger F$. The
physical three-term recurrence and its folded last equation give
\begin{equation}
 u_k=\begin{cases}
 -\mathrm i z/(c_ha_1a_k),&k\text{ even},\\
 \mathrm i\overline z\,t_m/(c_h^2a_1a_m^2a_k),
                    &k\text{ odd},\ m\text{ even},\\
 \mathrm i\overline z\,a_m^2/(t_ma_1a_k),
                    &k\text{ odd},\ m\text{ odd}.
 \end{cases}                                             \label{v16:recurrence}
\end{equation}
Negative coefficients are conjugates. To verify the formulas, the interior
row is $\mathrm i(t_ku_{k+1}-t_{k-1}u_{k-1})=F_k$; the first row gives
the even sequence and the last row gives the displayed odd sequence.
Products of $t_k=c_ha_ka_{k+1}$ telescope. This is an exact formula, not
an asymptotic boundary condition.

\begin{corollary}[Two explicitly different low-source limits]
\label{v16:examples}
Along $N\equiv1\pmod4$,
\[
 h^3\mathcal I_hD_h^\dagger\mathcal I_h^*\cos\theta\to-\pi^{-3}s,
 \qquad
 h^3\mathcal I_hD_h^\dagger\mathcal I_h^*\sin\theta\to\pi^{-3}l.
\]
Along $N\equiv3\pmod4$, the same sources instead have the $L^2$ limits
\begin{align}
 h^2\mathcal I_hD_h^\dagger\mathcal I_h^*\cos\theta
    &\longrightarrow\frac1{\pi^3}\sum_{j\ge1}\frac{\sin(2j\theta)}j,
                                                   \label{v16:cos-rough}\\
 h^2\mathcal I_hD_h^\dagger\mathcal I_h^*\sin\theta
    &\longrightarrow\frac1{\pi^3}\log(2|\sin\theta|).
                                                   \label{v16:sin-rough}
\end{align}
The latter limits are mean-zero distributional particular solutions of
$\mathcal Lu=\cos\theta,\sin\theta$, but are not smooth.
\end{corollary}
\begin{proof}
The first limits follow from \cref{v16:inverse-limit} and
$\ip l{\cos\theta}=-2/\pi$, $\ip s{\sin\theta}=2/\pi$.
For odd $m$, \eqref{v16:recurrence} shows that the normalized odd output
has norm $O(h)$ and that the even coefficient tends to
$-2\mathrm iz/(\pi^3 k)$. The same $1/\sin x$ bound and square-summable
tail argument as above prove $L^2$ convergence. With $z=1/2$ or
$z=-\mathrm i/2$ these are the two displayed Fourier series. The sine
series is the period-$\pi$ sawtooth with value
$(\pi-2\theta)/(2\pi^3)$ on $0<\theta<\pi$. Its jumps occur where
$\sin\theta=0$, so multiplying its distributional derivative by
$\sin\theta$ removes the delta terms. The logarithm has derivative
$\operatorname{pv}(\cot\theta)/\pi^3$ and hence flux
$\cos\theta/\pi^3$. Applying \eqref{v16:L} verifies both particular
solutions, including across the singular points.
\end{proof}

For comparison, the compatible even-frequency sources have exact finite
solutions already visible without an inverse estimate. At all sufficiently
fine grids,
\[
 D_h\cos\theta=\Delta_h\sin(2\theta),\qquad
 D_h\sin\theta=-\Delta_h\cos(2\theta),\qquad
                   \Delta_h=\pi^3h^2+O(h^4).
\]
Thus the normalized inverses of $\pi^3\sin(2\theta)$ and
$-\pi^3\cos(2\theta)$ converge to $\cos\theta$ and $\sin\theta$ on
both parity classes. The identities are obtained by the original doubled
trigonometric product and its actual derivative symbol; no continuum
product rule is imposed on the finite operator.

An exponentially localized source is therefore not, by itself, protected
from low-frequency consequences of the folded diagonal inverse. This is
stronger than a worst-case eigenvalue statement, but still concerns the
diagonal coefficient, not the full scalar operator of V15.

\section{The actual prepared source couples to a tensor soft mode}
\label{v16:actual-source}

Return to three dimensions and the actual limiting forcing $g^\circ$ in
\eqref{v15:g-definition}. All derivatives in this section are angle
 derivatives $\partial_i^\circ$, and the coefficient matrices are exactly
$M^\circ,r^\circ,H^\circ$ in V15. Define the real unit-norm test
\begin{equation}
                        \Phi(\theta)=s(\theta_2)l(\theta_3).
 \label{v16:Phi}
\end{equation}
It is independent of $\theta_1$. This is a legitimate $L^2$ pairing with
the analytic forcing, not pointwise evaluation of a singular initial field.

\begin{theorem}[A nonzero original-source pairing with the rough sector]
\label{v16:rough-pairing}
For the original coefficient forcing,
\begin{equation}
                  \boxed{-\frac18<\ip\Phi{g^\circ}<-\frac3{25}.}
 \label{v16:rough-interval}
\end{equation}
The assertion is certified by exact rational Fourier arithmetic and the
proved transverse inverse. It contains no inversion of the characteristic
scalar operator.
\end{theorem}
\begin{proof}
Retain the exact projected dyadic trial $\widehat w$ from V15 and recompute
its full noncyclic residual
$e=f^\circ-H^\circ\widehat w$. Put
$\widehat g=2\overline V^*M^\circ\widehat w$ and
$u=w^\circ-\widehat w=(H^\circ)^{-1}e$.
We first bound the entire $L^2$ error in $g$, rather than differentiating
the rough test. V13--V15 give $\|(H^\circ)^{-1}\|\le16/9$ and
$\|M^\circ\|\le18$. Each angle derivative of $M^\circ$ has pointwise
Frobenius norm at most one: it has four entries of magnitude at most $1/2$.
Since the ordinary transverse projection commutes with differentiation,
\begin{equation}
 \|\nabla u\|\le\frac{16}{9}\bigl(\|\nabla e\|+63\|u\|\bigr),
 \qquad \|u\|\le\frac{16}{9}\|e\|.                 \label{v16:H1-residual}
\end{equation}
Indeed each differentiated normal coefficient has norm at most $36$, and
$36\sqrt3<63$. Differentiation is justified by the inherited Sobolev bounds
for this uniformly positive transverse inverse, not by a scalar estimate.

For a raw off-diagonal vector $y$, the physical Frobenius adjoint is
\[
 2\overline V^*y
 =2\sum_{i<j}\{\partial_j^\circ(r_i^\circ y_{ij})
                         +\partial_i^\circ(r_j^\circ y_{ij})\}.
\]
Here $|r_i^\circ|\le1$ and $|\partial_j^\circ r_i^\circ|\le1/2$.
Cauchy--Schwarz gives
$\|2\overline V^*y\|\le2\sqrt3\|y\|+2\sqrt6\|\nabla y\|$.
Apply this to $y=M^\circ u$, using
$\|\nabla(M^\circ u)\|\le\sqrt3\|u\|+18\|\nabla u\|$.
The rational constants
\begin{equation}
                   \|g^\circ-\widehat g\|
                           \le72\|u\|+89\|\nabla u\|       \label{v16:g-error}
\end{equation}
are therefore valid upper bounds.

The new rational computation retains every residual frequency through
$|k|_1\le10$ and proves
\[
 \|e\|<\frac{661}{10^{10}},\qquad
 \|\nabla e\|<\frac{364}{10^9}.
\]
Set $d_0=(16/9)(661/10^{10})$ and
$d_1=(16/9)(364/10^9+63d_0)$. Direct rational comparison gives
$72d_0+89d_1<1/800$. Thus \eqref{v16:g-error} proves
$\|g^\circ-\widehat g\|<1/800$ on the \emph{full} coefficient Hilbert
space.

Write $s_o(\theta)=\sum_{k>0,\,k\ \mathrm{odd}}\sin(k\theta)/k$ and
$c_o(\theta)=\sum_{k>0,\,k\ \mathrm{odd}}\cos(k\theta)/k$. Their product
has rational imaginary Fourier coefficients. The exact finite Fourier
pairing with $\widehat g$ is
\begin{equation}
 \ip{s_o(\theta_2)c_o(\theta_3)}{\widehat g}
   =\frac{199947659541650831082275417}
           {2636573470272607903285248000}.               \label{v16:rational-pair}
\end{equation}
The actual test is $\Phi=-16s_oc_o/\pi^2$ and has norm one. Rational
upper and lower bounds for $\pi$, obtained by the alternating Machin
arctangent series, and \eqref{v16:rational-pair} therefore enclose
$\ip\Phi{\widehat g}$. Adding the error bound from
\eqref{v16:H1-residual}--\eqref{v16:g-error} gives, for orientation, the
outward interval $(-0.124179,-0.121703)$; the verified statements are the
strict rational inequalities in \eqref{v16:rough-interval}.
The code compares squared residual norms and all endpoint inequalities
using arbitrary-precision fractions. It also recovers V15's distinct smooth
test pairing from the same new $\widehat g$ image, checking the adjoint
normalization and the sign. No equality or sign is inferred from a small
floating residual.
\end{proof}

The proof reuses a trigonometric trial, not an assumed exact trigonometric
solution. The error estimate is strengthened to $H^1$ before applying the
first-order output. A mere $L^2$ bound on $w^\circ-\widehat w$ would not
justify this rough-source certificate.

\section{What the actual source would do without the flux term}
\label{v16:diagonal-comparison}

Use the original rationally selected complementary family of V12--V15, and
restrict now to $N\equiv1\pmod4$. Let
\[
 A_h^{\rm diag}=\tfrac12\mathscr D_{h,\alpha_h}^*
                                      \mathscr D_{h,\alpha_h},
 \qquad z_h=(A_h^{\rm diag})^{-1}g_h.
\]
The inverse exists on mean-zero scalars by the inherited detuning theorem.
\emph{This $z_h$ is a comparison solve, not the original lapse solution}
$n_h=\mathscr S_h^{-1}g_h$. It tests the effect of omitting a positive term
from an actual inverse; that omission is not performed in the evolution.

Let $P_{0,1}$ be the constant projection in the first coordinate and define
\[
 P_h^{\rm soft}=P_{0,1}\otimes E_h\otimes E_h,\qquad
 P^{\rm soft}=P_{0,1}\otimes E\otimes E,\qquad
 J=P_{0,1}\otimes K\otimes K\quad\hbox{on }\Ran P^{\rm soft}.
\]
This real subspace has dimension four. It reduces $A_h^{\rm diag}$ because
the scalar coefficient is a polynomial in the three commuting copies of
$D_h$. The projections are analytical tests of the inverse, not a truncation
of the original scalar equation.

\begin{theorem}[An actual low-source component of order $h^{-5}$ in the diagonal comparison]
\label{v16:diag-response}
In the physical trigonometric embedding,
\begin{equation}
 \boxed{h^5\mathcal I_hP_h^{\rm soft}z_h
   \longrightarrow\frac43(2I+J)P^{\rm soft}g^\circ
                         \quad\hbox{in }L^2(\T^3).}     \label{v16:diag-limit}
\end{equation}
The limit is nonzero and satisfies
\begin{equation}
           \boxed{\lim_{h\to0}h^5\|P_h^{\rm soft}z_h\|_h>\frac4{25}.}
 \label{v16:diag-lower}
\end{equation}
The four eigenvalues of the comparison compression have order $h^6$.
This does not give a lower bound of order $h^{-5}$ for the full solved lapse.
\end{theorem}
\begin{proof}
On the four-dimensional subspace the first scalar $D_h$ is zero and the
other two have eigenvalues $\pm\varepsilon_h$. Write
$K_h=D_hE_h/\varepsilon_h$ on the scalar soft space and
$J_h=I\otimes K_h\otimes K_h$ there. Since
$\alpha_h=(1,\theta_h,1)$, the exact tensor formula
\eqref{v2:tensor-norm} gives
\begin{equation}
 A_h^{\rm diag}\big|_{\Ran P_h^{\rm soft}}
   =\varepsilon_h^2\{(\theta_h^2+1)I-\theta_hJ_h\},
 \qquad J_h^2=I.                                       \label{v16:four-matrix}
\end{equation}
Its inverse bracket is
$\{(\theta_h^2+1)I+\theta_hJ_h\}/(\theta_h^4+\theta_h^2+1)$.
It is uniformly bounded above and below on the admitted detunings. Thus its
four eigenvalues have the stated scale, without using the much smaller
worst-case diagonal margin.

The embedded projections and involutions converge by
\cref{v16:inverse-limit}; $\theta_h\to1$ and
$\varepsilon_h^2/h^6\to4\pi^4$.
V15 proves $\mathcal I_hg_h/h\to(2\pi)^4g^\circ$ in $L^2$ for the
actual source. Apply the exact inverse of \eqref{v16:four-matrix} to that
source and use these convergences. The ratio of the two physical constants
is $(2\pi)^4/(4\pi^4)=4$, giving
$4(2I-J)^{-1}P^{\rm soft}g^\circ$. Since $J^2=I$,
$(2I-J)^{-1}=(2I+J)/3$, which is \eqref{v16:diag-limit}.

The test $\Phi$ of \eqref{v16:Phi} belongs to this limiting subspace and
has norm one. Hence \cref{v16:rough-pairing} gives
$\|P^{\rm soft}g^\circ\|>3/25$. The smallest eigenvalue of
$4(2I-J)^{-1}$ is $4/3$. This proves \eqref{v16:diag-lower}.
Every inverse and projection in the argument is the specified comparison;
adding the flux can change its inverse response and its invariant subspaces.
\end{proof}

This gives an actual-source example of the danger that V15's individual
high-frequency packet estimates did not address. The limiting rough modes
have nonzero low Fourier coefficients, so exponential localization of the
\emph{source} does not make their coupling exponentially small. Their
role in the full operator must nevertheless be evaluated, rather than
inferred from the diagonal comparison.

\section{The retained flux excludes the entire limiting soft sector}
\label{v16:flux-domain}

Keep the unit limiting coefficients $b^\circ,r^\circ$ of
\eqref{v15:limit-coefficients}. Let $M_b=M^\circ$ and let $M_r$ be the
same off-diagonal map with $b^\circ$ replaced by $r^\circ$:
\[
 M_b=\begin{pmatrix}b_2&b_1&0\\b_3&0&b_1\\0&b_3&b_2\end{pmatrix},\qquad
 M_r=\begin{pmatrix}r_2&r_1&0\\r_3&0&r_1\\0&r_3&r_2\end{pmatrix},\qquad
 A(\theta)=M_b(\theta)^{-1}M_r(\theta).
\]
The inverse coefficient is smooth, since its determinant never vanishes.
For distributions define the unscaled flux
\begin{equation}
                 \mathcal Z u=Q^\circ A\nabla^\circ u,
       \qquad Q^\circ=I-P^\circ.                        \label{v16:Z}
\end{equation}
It includes the constant vector components. The limit of the physical
$\sigma_h$ in \eqref{v14:z}, tested on smooth functions, is
$-(2\pi)^2\mathcal Z/\sqrt2$. This constant comes from the original
Frobenius inclusion and the two physical derivatives; it is not a change
of spatial or temporal units.

\begin{lemma}[A local obstruction to an $L^2$ flux]
\label{v16:local-singularity}
Let $a,b$ be smooth functions of tangential coordinates and let $x$ be a
normal coordinate. If they do not both vanish identically on a patch, a
distribution with leading normal terms
\[
                 a\delta'_0+b\operatorname{fp}(x^{-2})
\]
and remainder in $H^{-1}_{\rm loc}$ does not belong to $H^{-1}_{\rm loc}$,
provided the leading coefficients are not identically zero after tangential
localization. Consequently, if $\diver z$ has those terms, the longitudinal
projection of $z$ cannot be in $L^2$.
\end{lemma}
\begin{proof}
Choose a smooth compact tangential test so that at least one resulting
coefficient is nonzero. Tangential pairing is bounded from local $H^{-1}$
to one-dimensional local $H^{-1}$. The Fourier transforms of the localized
normal terms have leading terms $\mathrm i c\xi$ and $-\pi d|\xi|$,
respectively. The latter identity follows by differentiating the principal
value $1/x$; its sine-transform value can be obtained from
$\int_0^\infty e^{-\epsilon t}\sin t/t\,dt=\arctan(1/\epsilon)$
and then letting $\epsilon$ tend to zero. Smooth localization adds only
lower-order terms. Considering both signs of $\xi$ shows that their
squared magnitude divided by $1+\xi^2$ has a nonintegrable positive
leading value unless $c=d=0$. This remains true for complex coefficients,
since the sum over the two signs is a positive multiple of
$|c|^2+|d|^2$. Lower $H^{-1}$ terms cannot cancel that failure of
membership. Finally $\diver z=\diver(Q^\circ z)$, and the divergence
of an $L^2$ vector is in $H^{-1}$. This proves the last assertion.
\end{proof}

\begin{theorem}[No nonzero tensor soft mode has finite limiting flux energy]
\label{v16:no-soft-flux}
For every nonzero $u\in\Ran P^{\rm soft}$,
\begin{equation}
                           \boxed{\mathcal Z u\notin L^2(\T^3).}
 \label{v16:flux-not-L2}
\end{equation}
This is a statement about the actual coefficient flux, not an arbitrary
positive perturbation of the diagonal regulator.
\end{theorem}
\begin{proof}
Write
\[
 u=s(\theta_2)\{a_1s(\theta_3)+a_2l(\theta_3)\}
    +l(\theta_2)\{b_1s(\theta_3)+b_2l(\theta_3)\}.
\]
Temporarily denote the two braced functions by $a(\theta_3),b(\theta_3)$;
the constants here are not the boost coefficients. On an interval in
$0<\theta_3<\pi$ they are smooth and cannot both vanish identically
unless all four constants vanish, since $s=1$ there and $l$ is strictly
varying. Near $\theta_2=0$,
\[
 u=a(\theta_3)\operatorname{sgn}(\theta_2)
       +\frac2\pi b(\theta_3)\log|\theta_2|
       +\hbox{terms smooth in the normal coordinate}.
\]
The logarithmic correction to this expansion is smooth on the chosen
patch. Direct inversion of $M_b$ gives
\begin{equation}
                 A_{22}=\tfrac12\left(\frac{r_1}{b_1}
                                         +\frac{r_3}{b_3}\right).
 \label{v16:A22}
\end{equation}
At $\theta_1=\theta_2=0$ the actual coefficients give
\begin{equation}
  A_{22}=-\frac{\sin\theta_3}{4b_1},\qquad
          b_1=\frac{253}{64}+\frac{1+\cos\theta_3}{2}>0.
 \label{v16:normal-nonzero}
\end{equation}
Thus $A_{22}$ is nonzero on a sufficiently small tangential patch where
at least one of $a,b$ is nonzero.

The distribution $\diver(A\nabla^\circ u)$ on this patch has leading
normal terms
\[
 2A_{22}|_{\theta_2=0}\,a(\theta_3)\delta'_0
   -\frac2\pi A_{22}|_{\theta_2=0}\,b(\theta_3)
                               \operatorname{fp}(\theta_2^{-2}).
\]
All remaining terms have at most a delta or principal-value $1/\theta_2$
normal singularity, or a logarithm, with smooth localized tangential
coefficients. They belong to local $H^{-1}$. Derivatives of coefficients
multiplying the two displayed terms also produce only these lower
singularities. By \cref{v16:local-singularity}, the divergence is not in
$H^{-1}$. If $\mathcal Zu=Q^\circ A\nabla^\circ u$ were in $L^2$,
its divergence would be in $H^{-1}$ and would equal that same distribution.
This is a contradiction. No eigenbasis of the full scalar operator has
been substituted for the tensor modes.
\end{proof}

\section{The full finite scalar energy lifts these modes}
\label{v16:finite-lift}

\begin{theorem}[Uniform divergence of the normalized energy on the bare soft sector]
\label{v16:lift-theorem}
On the actual detuned complementary family, along $N\equiv1\pmod4$,
\begin{equation}
 \boxed{\inf_{u_h\in\Ran P_h^{\rm soft},\ \|u_h\|_h=1}
       h^{-2}\ip{u_h}{\mathscr S_hu_h}_h
                                      \longrightarrow+\infty.}
 \label{v16:lift}
\end{equation}
The diagonal contribution alone is $O(h^6)$ on this subspace, while the
retained flux has the stronger property \eqref{v16:lift}. No quantitative
rate of this divergence is asserted. This is a compression estimate, not a
lower bound on the scalar Schur complement after eliminating its complement.
\end{theorem}
\begin{proof}
Suppose the infimum stayed bounded along a subsequence, and choose unit
vectors there with bounded normalized energy. The embedded rank-four
projections converge in operator norm by \cref{v16:inverse-limit} and their
tensor products. Compactness in the fixed limiting four-dimensional space
therefore gives, after a subsequence,
$\mathcal I_hu_h\to u\in\Ran P^{\rm soft}$ strongly in $L^2$, with
$\|u\|=1$.

The original positive identity \eqref{v14:energy} and the fixed lower
bound for its flux metric give
\[
 h^{-2}\ip{u_h}{\mathscr S_hu_h}_h
          \ge\frac9{32}\|\sigma_h(u_h)\|_h^2.
\]
Hence the interpolated fluxes have a weakly convergent $L^2$ subsequence.
Their distributional limit is
$-(2\pi)^2\mathcal Zu/\sqrt2$. We justify this statement without assuming
regularity of $u_h$. For a smooth vector test, move the exact finite flux
by its adjoint. V14's uniform coefficient inverse and V15's smooth-test
consistency show that this adjoint converges in $L^2$ to the adjoint of
the displayed limiting first-order flux. Phase Hodge projections, literal
shifts and coefficient products are all retained in that consistency.
Pair with the strongly convergent $\mathcal I_hu_h$. Thus the weak $L^2$
limit must equal the indicated distribution. This would put
$\mathcal Zu$ in $L^2$, contradicting \cref{v16:no-soft-flux}.
This proves \eqref{v16:lift}. The diagonal order follows directly from
\eqref{v16:four-matrix}.
\end{proof}

The conclusion is useful precisely because it goes in the opposite
direction from \cref{v16:diag-response}. The actual source excites the
bare diagonal soft sector, but the full scalar energy penalizes it much
more strongly than the diagonal floor. The inverse can still mix this
sector with its complement to reduce the flux. A compression bound does
not control that mixing or justify using the compressed inverse in place
of the original inverse.

\begin{corollary}[The domain required of a bounded normalized actual lapse]
\label{v16:actual-domain}
Let $n_h=\mathscr S_h^{-1}g_h$ be the original leading active lapse
component on the selected family. If $\mathcal I_h(hn_h)$ is bounded in
$L^2$ along a sequence, every weak limit $n^\circ$ has
\[
        \mathcal Z n^\circ\in L^2,\qquad
                     \mathscr S^\circ n^\circ=g^\circ
                          \quad\hbox{distributionally}.
\]
Such a limit cannot be a nonzero vector in $\Ran P^{\rm soft}$, and it
cannot be zero. The boundedness hypothesis itself is not proved here.
\end{corollary}
\begin{proof}
Put $u_h=hn_h$. Its exact equation gives
\[
 h^{-2}\ip{u_h}{\mathscr S_hu_h}_h
                             =\ip{u_h}{g_h/h}_h.
\]
The right side is bounded by the assumed $L^2$ bound and V15's actual-source
estimate. The same positive flux bound and smooth adjoint consistency used
in \cref{v16:lift-theorem} give $\mathcal Zn^\circ\in L^2$, now with
weak rather than strong convergence of $u_h$; the tested adjoints converge
strongly. The limiting scalar equation is exactly the consistency statement
already proved in \eqref{v15:limiting-equation}. Its right side is nonzero
by V15, or by \cref{v16:rough-pairing}, so the limit is not zero.
\Cref{v16:no-soft-flux} excludes the other possibility stated.
\end{proof}

\section{What this resolves and the remaining full response}
\label{v16:status}

The fold-induced eigenvalue scale from V2 now has an identified two-function
range, an operator-norm inverse limit, and exact compatibility conditions on
fixed smooth sources. The actual V15 scalar forcing fails a corresponding
tensor-sector orthogonality test by a certified margin. Consequently a
diagonal-only approximation would predict an order-$h^{-5}$ response in
that sector. This is not a bound for the original solve. The original
flux instead excludes every nonzero limiting vector of that sector from
its finite-energy domain, and its normalized compression energy diverges.
Both sides of this distinction are proved for the same coefficients and
source.

The relevant question is therefore not whether an analytic source can excite
a bare diagonal soft mode; here it does. It is whether the \emph{full}
positive flux equation permits a controlled response after all complementary
components are allowed to cancel its flux. The low-source interface in
\eqref{v15:low-exact} retains precisely those components. Neither its high
block nor its lifted observation can be replaced by the bare diagonal
compression. The next substantive estimate remains
\[
 \mathfrak E_hT_h^\dagger V_h\mathscr S_h^{-1}P_{\le R_h}g_h,
\]
with V15's controlled high-source remainder. The new limiting energy-domain
condition is more restrictive than distributional solvability alone, but
is not an existence or uniform inverse theorem for this range problem.

No finite action, source pairing, derivative symbol, initial calibration,
physical clock, or literal link is changed. No new higher weak signed-rank
certificate, actual initial multiplier cancellation, nonlinear rank
propagation, differentiated curvature estimate, or common exact-action
lifespan is claimed. In particular the scalar active analysis remains
separate from identifying it with a full nonlinear multiplier rate. The
raw-action Einstein continuum bridge is not closed by a comparison inverse.

\section{Verification, sources, and append-only preservation}
\label{v16:verification}

The package supplies the new rational source-pairing computation, including
its full residual derivative, the $H^1$ inverse error bound, exact bounds
for $\pi$, and all final inequalities. The unchanged V15 dyadic trial and
its rational coefficient module are fingerprinted inherited inputs. The
new calculation independently recovers the prior smooth pairing from its
new full output polynomial; it does not claim to regenerate the 52
original-action source coefficients or the entire earlier stationary chart.

Symbolic checks verify the physical inverse recurrences on formal rational
Jacobi coefficients of both parities, the differential limit, the two
rough particular solutions away from their singular points, the exact
normal flux coefficient, and the four-dimensional inverse polynomial.
Distributional behavior at the singular points and the uniform asymptotic
bounds have the written proofs above. Floating Jacobi diagnostics at
$N=9,13,17,25,33,49,65,97,129,193,257$ separately compare the normalized
inverse and spectral projector with their stated limits; they are not
interval eigenvalue certificates. Separate full-flux diagnostics at
$N=9,17,33,65$ evaluate the four-dimensional compression with probe
detunings $\theta_h=1+h/2$; they do not certify the special rational
selection or invert the full scalar equation. No arbitrary small floating
residual is used as a rank, vanishing, or domain proof.

A targeted audit of the preceding bridge and original-action source found
the inherited soft eigenvalue asymptotic and its reciprocal-square sum,
but not these physical rough-mode limits, range conditions or actual-source
tensor pairing. No exhaustive literature-priority claim is made. The
fifteen preceding marked research bodies remain byte-for-byte unchanged,
including their corrections and qualifications. The package records
preservation hashes and compilation checks separately from mathematical
verification. No nonlinear trajectory or proof-assistant formalization is
represented as having been performed.

\begin{thebibliography}{9}
\raggedright
\setcounter{enumiv}{12}
\bibitem{v16:Yafaev}
D. R. Yafaev, \emph{Self-adjoint Jacobi operators in the limit circle case},
arXiv:2104.13609 (2021).
\url{https://arxiv.org/abs/2104.13609}.
This supplies primary-source context for the role of domains and boundary
conditions in Jacobi inverses, not a theorem about the present finite fold,
its selected modes, or the prepared renewal source. The required facts in
this body are proved directly.
\end{thebibliography}
% END IMMUTABLE EXACT-ACTION BRIDGE V16 BODY


% Future iterations append here.
% BEGIN IMMUTABLE EXACT-ACTION BRIDGE V17 BODY
\clearpage
\renewcommand{\thesection}{\arabic{section}}
\renewcommand{\theHsection}{bridgeV17.\arabic{section}}
\setcounter{section}{135}
\section{The domain of the full scalar energy, not its isolated compression}
\label{v17:context}

The exclusion proved in \cref{v16:no-soft-flux} concerns a particular
four-dimensional space of coordinate-aligned square-wave and logarithmic
modes. It does not exclude every discontinuous scalar from the domain of the
original flux. Nor does the divergent energy of that compression determine
the inverse after all complementary components are retained. This body
therefore identifies the closed domain of the \emph{full limiting scalar
energy} and proves convergence of the complete finite scalar forms to it.
An explicit family of characteristic jump profiles makes this domain larger
than a classical Sobolev domain. Their jump parts are annihilated by the
actual off-diagonal lapse source. The limiting scalar operator maps them to
analytic right sides. The actual prepared forcing has a certified nonzero
pairing with one of them.

There is also a quantitative finite result. A specified sequence of
trigonometric approximations to that jump profile has convergent normalized
\emph{full} scalar energy. When used at the lapse scale $h^{-1}$, its actual
lapse-to-shift geometric output tends to zero, although the sufficient
three-derivative lapse bound from V13 diverges. No inverse is replaced by a
compression in these statements. They give a larger viable observation domain,
not a claim that the actual solved lapse has been shown to enter it.

\paragraph{Dependencies and interpretation.}
We retain the exact flux identity and coefficient bounds of V14, the
normalized transverse inverse of V13, the smooth consistency and actual
forcing of V15, and V16's rational trial and output-error estimate. The
calculation of the 52 original source coefficients, the soft scalar
spectrum, and the nonlinear initializer are inherited, not repeated as new
results. First-order graph approximation and weak/strong convergence of
quadratic forms are classical methods \cite{v17:Ruppenthal,v17:Mosco}.
The proofs below specify their application to the actual projected flux,
including its constant components and the diagonal finite regulator.
The resolvents used below are analytical spectral probes. No positive mass
is added to the action or proposed as a physical evolution law.

\section{The maximal flux graph and its smooth core}
\label{v17:domain}

Work on the real mean-zero Hilbert space $\mathcal H=L_0^2(\T^3)$, with
normalized measure and angle derivatives $\partial_i^\circ$ as in V15.
All coefficient fields in this section are the \emph{unit limiting
complementary fields} $b^\circ,r^\circ$ in
\eqref{v15:limit-coefficients}. Write
\[
 M=M_b,\qquad A=M_b^{-1}M_r,\qquad P=P^\circ,\qquad Q=I-P,
 \qquad K=QM^{-1}M^{-*}Q\big|_{\Ran Q}.
\]
Here $M_r$ is the same off-diagonal coefficient map with $b$ replaced by $r$.
The coefficient inverse is smooth and real analytic. The range of $Q$
contains both gradients and the three constant vector fields. Its retained
metric satisfies $324^{-1}I\preceq K\preceq16I/9$. Define, in distributions,
\begin{equation}
 Z u=QA\nabla^\circ u,\qquad
 \mathcal D=\{u\in\mathcal H:Zu\in L^2(\T^3;\R^3)\}.
 \label{v17:Z-domain}
\end{equation}
Neither $A\nabla u$ nor $u$ is required separately to have a full derivative
in $L^2$. The projection in \eqref{v17:Z-domain} is the original flux
projection, not an additional dynamical projection.

\begin{theorem}[Closed maximal graph and approximation without a hidden boundary condition]
\label{v17:graph}
The operator $Z:\mathcal D\subset\mathcal H\to\Ran Q$ is closed and
densely defined. Smooth mean-zero periodic functions are dense in its graph
norm $\|u\|+\|Zu\|$. Consequently
\begin{equation}
 a(u,v)=\tfrac14\ip{Zu}{K^{-1}Zv},\qquad u,v\in\mathcal D,
 \label{v17:form}
\end{equation}
is a densely defined closed nonnegative form. Its associated nonnegative
self-adjoint operator $\mathscr A$ agrees with
$\mathscr S^\circ=\overline V^*(I-\overline\Pi)\overline V/2$ on smooth
functions, and its distributional expression on its operator domain is
$Z^*K^{-1}Z/4$.
\end{theorem}
\begin{proof}
If $u_j\to u$ in $L^2$ and $Zu_j\to v$ in $L^2$, multiplication by the
smooth $A$, differentiation, and the order-zero Fourier projection $Q$
pass to distributions. Thus $Zu=v$, proving closedness. Smooth functions
are in the domain and dense in $\mathcal H$.

Let $J_\epsilon$ be convolution with a smooth periodic approximate identity.
It commutes with $Q$ and every ordinary derivative. For any $u\in\mathcal D$,
\[
 ZJ_\epsilon u-J_\epsilon Zu
       =Q[A\nabla^\circ,J_\epsilon]u.
\]
The commutator on the right is uniformly bounded $L^2\to L^2$ and tends
strongly to zero. Here is a direct verification, also valid when
$A\nabla u$ is only a distribution. For one entry its kernel formula is
\begin{align*}
 [A_{ij}\partial_j,J_\epsilon]u(x)
  ={}&\int(A_{ij}(x)-A_{ij}(y))\partial_j\rho_\epsilon(x-y)u(y)\,dy\\
     &+\int\rho_\epsilon(x-y)(\partial_jA_{ij})(y)u(y)\,dy.
\end{align*}
The Lipschitz bound and $\int |z||\nabla\rho_\epsilon(z)|\,dz\le C$
give a cutoff-independent $L^2$ operator bound by Young's inequality. On a
smooth $u$ the commutator tends to zero by ordinary approximation of both
terms in the original commutator. Density and the uniform bound give
strong convergence for every $L^2$ input. Therefore
$J_\epsilon u\to u$ and $ZJ_\epsilon u\to Zu$ in $L^2$.
The mollifier preserves the scalar mean. This proves graph density, not
merely weak distributional approximation.

The fixed upper and lower bounds on $K^{-1}$ make the form norm equivalent
to the graph norm, so the form is closed. The usual variational construction
of its operator is unambiguous: $u\in D(\mathscr A)$ exactly when there is
$f\in\mathcal H$ with $a(u,v)=\ip fv$ for every $v\in\mathcal D$, and
then $\mathscr Au=f$. Existence of its resolvent follows by completing
$\mathcal D$ in $a+\rho\ip\cdot\cdot$ and applying Riesz representation,
for each $\rho>0$. This is the closed nonnegative self-adjoint realization.
On smooth functions the exact flux projection identity in V14, with
$\overline V=-M_r\nabla^\circ/2$, gives
\[
 \tfrac12\|(I-\overline\Pi)\overline V u\|_F^2
       =\tfrac14\ip{Zu}{K^{-1}Zu}.
\]
The two off-diagonal Frobenius occurrences are included here. Polarization
and the smooth core prove the asserted expression and agreement.
\end{proof}

This theorem does not impose a choice among the self-adjoint realizations
of V16's \emph{diagonal-only} differential expression. The operator here is
fixed by a different, complete positive flux form. We do not assert that
its kernel on mean-zero functions is zero or that its range is closed.

\section{The complete finite scalar forms have one cofinal limit}
\label{v17:mosco}

Let $\lambda_*=2\pi$. Use the exact finite scalar matrix $\mathscr S_h$
of \eqref{v13:S-g}, including both terms, and any complementary choices
$\alpha_h=(1,\theta_h,1)$ with $1<\theta_h<1+h$. Special rational spectral
selection is not required in this section. Define
\begin{equation}
 \mathscr A_h=\frac{\mathscr S_h}{h^2\lambda_*^4},\qquad
 a_h(u_h)=\ip{u_h}{\mathscr A_hu_h}_h.
 \label{v17:Ah}
\end{equation}
Embed the mean-zero grid space isometrically in $\mathcal H$ by the original
trigonometric interpolation $\mathcal I_h$, and extend $a_h$ by $+\infty$
off that finite-dimensional space. Write $\sigma_h$ for the \emph{actual}
flux in \eqref{v14:z} and put
\[
 Z_h=-\frac{\sqrt2}{\lambda_*^2}\sigma_h.
\]
Its sign has no effect on energy, but fixes convergence to $Z$ rather than
$-Z$. The exact finite identity is
\begin{equation}
 a_h(u_h)=\frac{\|\mathscr D_{h,\alpha_h}u_h\|_h^2}
                       {2h^2\lambda_*^4}
               +\tfrac14\ip{Z_hu_h}{K_{Q,h}^{-1}Z_hu_h}_h.
 \label{v17:finite-form}
\end{equation}

\begin{lemma}[Smooth consistency including adjoints and the metric]
\label{v17:consistency}
On each fixed smooth periodic scalar or vector test, the interpolated $Z_h$
and its adjoint converge in every fixed Sobolev order to $Z$ and $Z^*$.
The finite $Q_h$ and $K_{Q,h}^{\pm1}$ likewise converge on smooth tests to
$Q$ and $K^{\pm1}$, after projection onto their respective ranges.
Also
\[
 \|\mathscr D_{h,\alpha_h}\mathcal I_h^*u\|_h=O_u(h^2).
\]
All statements retain the means and all finite folded products.
\end{lemma}
\begin{proof}
For a fixed smooth test, $\delta_h/\lambda_*$ converges to
$\partial^\circ$ with three derivatives charged and an $O(h^2)$ error.
The literal shifts differ from identity by $O(h)$ with one derivative
charged. The complementary coefficients and their fixed derivatives converge
at order $h$, including the detuning. Smooth Fourier tails absorb the
sampling or restriction difference and every folded term. The Hodge symbols
$\kappa_i\kappa_j/|\kappa|^2$ converge to $k_ik_j/|k|^2$ on smooth tests;
the zero mode is retained in $Q_h$ and $Q$ in exactly the same way.

The coefficient inverses and the normalized transverse inverse have the
uniform Sobolev bounds of V13--V14. Inserting a sampled smooth limiting
inverse into its finite equation gives an $O(h)$ smooth residual; the
uniform inverse bounds transfer this to the exact finite solution. This
proves consistency through $M_h^{-1}$ and through the inverse in the
Schur formula for $K_{Q,h}^{-1}$. The formula for $K_{Q,h}$ itself only
uses coefficient inverses and the projections. Exact adjoints reverse the
same finite products and shifts and use $\delta_h^*=-\delta_h$, so the
identical test argument applies to them, without an inverse for
$\mathscr S_h$. The physical constants in V16 give
$\sigma_h\to-\lambda_*^2Z/\sqrt2$, as normalized above.
Finally the diagonal source is the three differences of the one-dimensional
product defects. The smooth-test expansion in \cref{v16:smooth-limit}
gives its stated $O(h^2)$ bound. These are consistency assertions on fixed
tests, not uniform estimates on unknown nonsmooth solutions.
\end{proof}

\begin{theorem}[Weak lower bound and strong recovery for the full scalar regulator]
\label{v17:form-limit}
The forms $a_h$ converge to $a$ in the following two senses:
\begin{enumerate}[label=(\roman*),itemsep=3pt]
\item If $\mathcal I_hu_h\rightharpoonup u$ in $\mathcal H$, then
$\liminf_h a_h(u_h)\ge a(u,u)$, where the right side is $+\infty$ for
$u\notin\mathcal D$.
\item For every $u\in\mathcal D$ there are finite mean-zero $u_h$ with
$\mathcal I_hu_h\to u$ strongly and $a_h(u_h)\to a(u,u)$.
\end{enumerate}
Thus this is convergence of the \emph{full} forms, usually called Mosco
convergence. It holds along all fine odd grids and has the same limit on
both congruence classes modulo four.
\end{theorem}
\begin{proof}
For (i) pass to a subsequence with bounded energies realizing the lower
limit. The uniform metric margin in \eqref{v17:finite-form} bounds
$\mathcal I_hZ_hu_h$ in $L^2$. Extract a weak limit $z$. Test against a
smooth vector and move $Z_h$ by its exact adjoint. The strong adjoint
consistency of \cref{v17:consistency}, paired with weak convergence of
$\mathcal I_hu_h$, identifies $z=Zu$ distributionally. In particular
$u\in\mathcal D$ and $z\in\Ran Q$.

For every smooth $v\in\Ran Q$ use $v_h=Q_h\mathcal I_h^*v$. The positive
finite metric has the dual inequality
\[
 \ip{Z_hu_h}{K_{Q,h}^{-1}Z_hu_h}_h
  \ge2\ip{Z_hu_h}{v_h}_h-\ip{v_h}{K_{Q,h}v_h}_h.
\]
Both terms on the right converge. Taking the supremum over smooth $v$ in
$\Ran Q$, which are dense there, gives
$\liminf\ip{Z_hu_h}{K_{Q,h}^{-1}Z_hu_h}\ge\ip{Zu}{K^{-1}Zu}$.
The diagonal contribution is nonnegative and cannot invalidate this lower
bound. This proves (i), including its domain assertion.

For smooth $u$, take $u_h=\mathcal I_h^*u$. Its diagonal contribution is
$O(h^2)$ by the lemma. Its flux and metric converge strongly, so
$a_h(u_h)\to a(u,u)$ and interpolation converges strongly to $u$.
For arbitrary $u\in\mathcal D$, the preceding graph theorem supplies
smooth mean-zero $u^{(j)}$ with $u^{(j)}\to u$ and $Zu^{(j)}\to Zu$.
Their energies converge. For each $j$ choose a threshold below which the
finite consistency error for $u^{(j)}$ is less than $1/j$; choose these
thresholds decreasing. Using $u^{(j)}$ on the associated successive ranges
of $h$ gives the required diagonal recovery sequence. This construction
changes no finite operator or action coefficient. Its purpose is comparison
of their quadratic forms, not prescription of a physical history.
Neither argument uses the parity-dependent inverse of the diagonal term.
\end{proof}

\section{A convergent full resolvent, and the zero-parameter boundary}
\label{v17:resolvents}

The actual normalized right side is
\begin{equation}
 f_h=\frac{g_h}{h\lambda_*^4},\qquad
          \mathcal I_h f_h\longrightarrow g^\circ\quad\hbox{in }L^2,
 \label{v17:forcing}
\end{equation}
by V15. The original unregularized scalar response $n_h$ obeys
$\mathscr A_h(hn_h)=f_h$ on the rationally selected invertible family.

\begin{theorem}[Cofinal resolvent convergence for the actual normalized source]
\label{v17:resolvent}
For every fixed $\rho>0$ set
\[
 u_{h,\rho}=(\mathscr A_h+\rho I)^{-1}f_h,
 \qquad u_\rho=(\mathscr A+\rho I)^{-1}g^\circ.
\]
Then along the full odd refinement,
\begin{equation}
 \mathcal I_hu_{h,\rho}\longrightarrow u_\rho\ \hbox{strongly in }L^2,
 \qquad a_h(u_{h,\rho})\longrightarrow a(u_\rho,u_\rho).
 \label{v17:resolvent-limit}
\end{equation}
In addition the interpolated fluxes converge strongly,
$\mathcal I_hZ_hu_{h,\rho}\to Zu_\rho$, and
$\|\mathscr D_{h,\alpha_h}u_{h,\rho}\|_h/h\to0$.
The same statements hold for any strongly convergent sequence of right
sides in place of \eqref{v17:forcing}.
\end{theorem}
\begin{proof}
The resolvent is the unique minimizer of
\[
 F_h(v)=a_h(v)+\rho\|v\|^2-2\ip{f_h}{v}.
\]
Comparison with zero and Cauchy--Schwarz bound its $L^2$ norm and energy
uniformly. Weak compactness and the weak lower inequality in
\cref{v17:form-limit} identify every cluster point as the unique minimizer
of $F(v)=a(v,v)+\rho\|v\|^2-2\ip{g^\circ}{v}$, using a strong recovery
sequence for each competitor. Thus the minimum values converge. If $v_h$
is a strong recovery sequence for that limiting minimizer, the exact
quadratic identity at the finite minimizer gives
\[
 F_h(v_h)-F_h(u_{h,\rho})
  =a_h(v_h-u_{h,\rho})+\rho\|v_h-u_{h,\rho}\|^2.
\]
The left side tends to zero; this proves strong convergence. The Euler
identity $a_h(u_{h,\rho})+\rho\|u_{h,\rho}\|^2
=\ip{f_h}{u_{h,\rho}}$ then gives energy convergence.

The flux lower bound in the previous proof is attained in the limit; the
nonnegative diagonal term must therefore tend to zero. To see strong flux
convergence rather than just convergence of total energies, approximate
$Zu_\rho$ in $L^2$ by smooth vectors in $\Ran Q$, use their strongly
convergent finite metric images in the polarized squared norm, and then
use the uniform upper and lower metric margins. Weak convergence of the
fluxes and convergence of their metric squared norms give strong $L^2$
convergence. The diagonal assertion follows from its exact coefficient in
\eqref{v17:finite-form}.
\end{proof}

This is a resolvent statement about the original scalar matrix at a
\emph{fixed spectral parameter}. It is not a modified action, a new lapse
mass, or the unproved conclusion at $\rho=0$.

\begin{proposition}[What the limit proves about the actual unregularized susceptibility]
\label{v17:susceptibility}
On the selected invertible family define
\[
 E_h=\ip{f_h}{\mathscr A_h^{-1}f_h}_h,
 \qquad
 E^\circ=\sup_{u\in\mathcal D}
          \{2\ip{g^\circ}{u}-a(u,u)\}\in[0,+\infty].
\]
Then
\begin{equation}
                     \liminf_h E_h\ge E^\circ.
 \label{v17:susceptibility-lower}
\end{equation}
Moreover
\[
 E^\circ=\sup_{\rho>0}\ip{g^\circ}{(\mathscr A+\rho I)^{-1}g^\circ}.
\]
No finite value of $E^\circ$ or matching upper bound for $E_h$ is asserted.
\end{proposition}
\begin{proof}
For any positive finite $\mathscr A_h$, adding $\rho I$ decreases the
quadratic inverse, as is seen by diagonalizing that finite matrix. The
previous theorem therefore gives
$\liminf E_h\ge\ip{g^\circ}{(\mathscr A+\rho I)^{-1}g^\circ}$ for every
fixed $\rho$. The latter number is the supremum over $u\in\mathcal D$ of
$2\ip{g^\circ}{u}-a(u,u)-\rho\|u\|^2$. For each fixed competitor its
last term tends to zero as $\rho\downarrow0$. Taking the two suprema proves
the identity, with $+\infty$ allowed, and then the inequality.
\end{proof}

In particular a future proof that this \emph{actual} limiting source has
infinite energy-dual susceptibility would imply divergence of the original
normalized scalar energy. Conversely finiteness alone does not control
small finite eigenvalues invisible to fixed-$\rho$ convergence. The
interchange of zero spectral parameter and refinement has not been made.

\section{Characteristic jumps belonging to the full operator domain}
\label{v17:jumps}

The domain above has explicit non-Sobolev elements, different from V16's
excluded tensor modes. Let $w$ be the $2\pi$-periodic sawtooth equal to
$\varphi$ on $-\pi<\varphi<\pi$. Its values at the jump are immaterial.
For $\{i,j,k\}=\{1,2,3\}$ set
\begin{equation}
 W_i(\theta)=w(\theta_j+\theta_k),\qquad
 \xi_i=e_j+e_k.
 \label{v17:saw}
\end{equation}
These are mean-zero $L^2$ functions, with squared norm $\pi^2/3$, and
none belongs to $H^{1/2}$. Their Fourier series is
$W_i=2\sum_{m\ge1}(-1)^{m+1}\sin(m(\theta_j+\theta_k))/m$.

\begin{theorem}[An exact source-null jump with analytic image]
\label{v17:jump-domain}
Define $v_1=(1,1,-1)$, $v_2=(1,-1,1)$ and $v_3=(-1,1,1)$ in raw
$(12,13,23)$ coordinates. For the actual limiting coefficients,
\begin{equation}
 M_r\nabla^\circ W_i=r_i v_i
                 \quad\hbox{as distributions},
 \qquad ZW_i=QM^{-1}r_i v_i.
 \label{v17:jump-cancel}
\end{equation}
In particular $W_i\in\mathcal D$ and in fact $W_i\in D(\mathscr A)$.
Its operator image $\mathscr A W_i$ is a real analytic periodic mean-zero
function. Thus analytic forcing of this limiting scalar operator does not,
without a separate condition, force every solution to be in $H^{1/2}$.
For $i=1$ the flux and energy are nonzero.
\end{theorem}
\begin{proof}
Since $r_1+r_2+r_3=0$, direct multiplication gives
$M_r\xi_i=r_iv_i$. Distributionally,
$\nabla^\circ W_i=\xi_i(1-2\pi\delta_{\{\theta_j+\theta_k=\pi\}})$,
where the delta is normalized in the summed angle coordinate. On that
surface $\sin\theta_j=\sin\theta_k$, so $r_i$ vanishes there. Multiplication
by this smooth vanishing coefficient annihilates the entire delta term.
This proves \eqref{v17:jump-cancel}, including across the jump, not merely
away from it.

The right side is a trigonometric vector. Multiplication by $M^{-1}$ is
real analytic, and the ordinary Hodge projection preserves an exponential
Fourier bound. Hence $ZW_i$ is real analytic. The same is true after
applying $K^{-1}$: its formula through the uniformly conditioned transverse
inverse has the analytic-weight bounds of V13--V15. Finally $Z^*$ is the
adjoint of an order-one operator with analytic coefficients and an
order-zero Hodge projection; it preserves analyticity after reducing the
positive strip width. Thus $f_i=Z^*K^{-1}ZW_i/4$ is analytic and mean-zero.
Integration by parts against smooth tests gives
$a(W_i,v)=\ip{f_i}{v}$, and the graph core extends it to every
$v\in\mathcal D$. This proves membership in the operator domain.

For a nonzero check at $i=1$, away from any coordinate ambiguity
$A\nabla W_1=A\xi_1$. Direct rational evaluation of its divergence at
$(\theta_1,\theta_2,\theta_3)=(0,\pi/2,0)$ gives
\begin{equation}
 \diver^\circ(A\xi_1)=-\frac{21594937}{751168800}\ne0.
 \label{v17:jump-nonzero}
\end{equation}
One can reproduce it by setting
$b=(253/64+(c_2+c_3)/2,-8+(c_1-c_3)/2,-8-(c_1+c_2)/2)$,
$r=(s_2-s_3,s_3-s_1,s_1-s_2)/2$, forming $M_b^{-1}M_r\xi_1$, and applying
$-s_a\partial_{c_a}+c_a\partial_{s_a}$ to its $a$th component before
substituting $c=(1,0,1)$, $s=(0,1,0)$. Since $\diver(Qz)=\diver z$,
\eqref{v17:jump-nonzero} proves $ZW_1\ne0$. Its form energy is therefore
strictly positive by the metric margin.
\end{proof}

The surfaces in \eqref{v17:saw} are characteristic sheets of the actual
source, on which the whole normal source vector vanishes. They are not
V16's coordinate-aligned jump surfaces with a nonzero normal coefficient.
Their opposite domain behavior is consequently consistent. The functions
are analytical scalar tests and possible limiting lapse components, not
nonsmooth initial metric records inserted into the nonlinear action.

\section{The original forcing has a verified nonzero jump pairing}
\label{v17:source-jump}

\begin{theorem}[Actual-source coupling to an admissible jump mode]
\label{v17:pairing}
For the original $g^\circ$ defined in \eqref{v15:g-definition},
\begin{equation}
                 \boxed{\frac3{10}<\ip{W_1}{g^\circ}<\frac{31}{100}.}
 \label{v17:pairing-interval}
\end{equation}
This is not the assertion that $g^\circ=\mathscr A W_1$, or that the
unregularized solution is a multiple of $W_1$.
\end{theorem}
\begin{proof}
Use the inherited dyadic transverse trial $\widehat w$ and recompute its
full noncyclic residual, derivative residual and output
$\widehat g=2\overline V^*M\widehat w$, exactly as specified by the
rational coefficient operations of V16. Those operations and their two
residual bounds verify again
$\|g^\circ-\widehat g\|<1/800$. The input source table and inverse margin
are inherited; the new pairing is evaluated afresh.

The Fourier coefficient of $W_1$ at $(0,m,m)$ is
$\mathrm i(-1)^m/m$ for every nonzero integer $m$, and its other
coefficients vanish. Pair this with the complete rational Fourier
polynomial $\widehat g$. Exact rational arithmetic gives
\begin{equation}
               \ip{W_1}{\widehat g}
                  =\frac{171018477772789}{562949953421312}.
 \label{v17:pairing-rational}
\end{equation}
Although $W_1$ has infinitely many coefficients, only finitely many enter
this particular polynomial pairing. It is not a truncation of the true
forcing error. Since $\|W_1\|=\pi/\sqrt3<2$, that error in the pairing is
less than $1/400$. Adding and subtracting $1/400$ from the rational number
in \eqref{v17:pairing-rational} gives an interval strictly contained in
$(3/10,31/100)$. Every comparison is rational; for the norm bound it
suffices to use $\pi<22/7$ and $(22/7)^2<12$. The verifier independently
recovers the preceding smooth-source pairing from the same output polynomial.
\end{proof}

This test has finite full flux energy and nonzero coupling to the actual
source, unlike the bare tensor sector excluded in V16. It follows, for
example, that $E^\circ>0$: optimize the variational expression over scalar
multiples of $W_1$. This does not decide whether $E^\circ$ is finite.

\section{An explicit full-regulator recovery sequence for the jump}
\label{v17:recovery}

The abstract recovery argument can be made quantitative on this original-source
coupled direction. Let
\begin{equation}
 w_L(\varphi)=2\sum_{m=1}^L
       (-1)^{m+1}\left(1-\frac m{L+1}\right)\frac{\sin(m\varphi)}m,
 \quad U_L(\theta)=w_L(\theta_2+\theta_3).
 \label{v17:Fejer}
\end{equation}
These are genuine trigonometric polynomials, so the finite arrays below
require no sampling at a discontinuity. Write $F_L$ for the Fejer kernel
with coefficients $a_m=\max(1-|m|/(L+1),0)$. Then
$w_L'=1-F_L(\varphi-\pi)$.

\begin{lemma}[Exact cancellation norm and polynomial derivative bounds]
\label{v17:Fejer-bounds}
One has $\|U_L-W_1\|\le CL^{-1/2}$. For every fixed integer $j\ge1$,
\[
 c_jL^{j-1/2}\le\|U_L\|_{\dot H^j}\le C_jL^{j-1/2}.
\]
The actual rotation coefficient satisfies the exact identity
\begin{equation}
 \boxed{\|r_1F_L(\theta_2+\theta_3-\pi)\|^2
                              =\frac1{4(L+1)}.}
 \label{v17:Fejer-exact}
\end{equation}
Consequently $M_r\nabla U_L\to r_1v_1$ in $L^2$ with error
$\sqrt{3}/(2\sqrt{L+1})$, and $ZU_L\to ZW_1$ in $L^2$.
\end{lemma}
\begin{proof}
The sawtooth coefficients are proportional to $m^{-1}$. On $m\le L$
the coefficient difference is $O((L+1)^{-1})$, and the omitted squared
tail is $O(L^{-1})$. Parseval proves the first claim. Squaring the
$j$th derivative coefficients gives a sum comparable to
$\sum_{m\le L}m^{2j-2}(1-m/(L+1))^2\asymp L^{2j-1}$; restricting to
$L/4\le m\le L/2$ gives the lower bound.

At fixed $\varphi=\theta_2+\theta_3$, averaging over the other angle gives
$\mathbb E(r_1^2\mid\varphi)=(1+\cos\varphi)/4$.
With $s=\varphi-\pi$, the desired squared norm is therefore
$\int(1-\cos s)F_L(s)^2\,ds/(8\pi)$.
By Parseval this equals
\[
 \frac14\left(\sum_m a_m^2-\sum_m a_ma_{m+1}\right)
  =\frac18\sum_m(a_{m+1}-a_m)^2=\frac1{4(L+1)}.
\]
There are exactly $2(L+1)$ nonzero differences, each of magnitude
$1/(L+1)$. Since $M_r\xi_1=r_1v_1$ and $|v_1|^2=3$, differentiation
of \eqref{v17:Fejer} proves the vector norm assertion. Multiplication
by the bounded $M^{-1}$ and the orthogonal $Q$ proves graph convergence.
The error vector has Fourier degree at most $L+1$ in each coordinate;
its $H^j$ norm is at most $C_jL^{j-1/2}$ by the same exact $L^2$ identity
and the finite bandwidth.
\end{proof}

\begin{theorem}[Quantitative recovery through the unchanged finite energy]
\label{v17:quantitative-recovery}
Set $L_h=\lfloor h^{-1/2}\rfloor$ and let $u_h$ be the samples of
$U_{L_h}$ on the original odd grid. For every sufficiently small $h$
these arrays are mean-zero and their interpolants equal $U_{L_h}$ exactly.
On all complementary detunings $1<\theta_h<1+h$,
\begin{equation}
 \begin{split}
  \|\mathcal I_hu_h-W_1\|&\le Ch^{1/4},\\
  \|\mathcal I_hZ_hu_h-ZW_1\|&\le Ch^{1/4},\\
  \left|a_h(u_h)-a(W_1,W_1)\right|&\le Ch^{1/4}.
 \end{split}
 \label{v17:explicit-recovery}
\end{equation}
In particular the original full normalized energy is bounded on this
explicit sequence, despite its loss of $H^{1/2}$ regularity in the limit.
Its actual source pairing satisfies
$\ip{u_h}{f_h}_h\to\ip{W_1}{g^\circ}>3/10$.
\end{theorem}
\begin{proof}
For small $h$, $L_h+2<(N-1)/2$, so products with the unit coefficient bank
are unfurled in the operations on $u_h$. No matrix entry is removed by this
fact. The nonfolded cosine-defect estimate gives
\[
 \|\mathscr D_{h,\alpha_h}u_h\|\le Ch^2\|U_{L_h}\|_{H^2}
                                      \le Ch^2 L_h^{3/2}.
\]
Its contribution to \eqref{v17:finite-form} is at most $Ch^2L_h^3$.
For the off-diagonal source, sine Taylor expansion, the exact shift
estimate, and coefficient detuning give
\begin{equation}
 \left\|\frac{\mathcal I_hV_hu_h}{h\lambda_*^2}
                         -\overline VU_{L_h}\right\|_{H^j}
 \le C_j\left(h L_h^{j+3/2}+h^2L_h^{j+5/2}\right)
 \quad(j=0,2).
 \label{v17:source-recovery-estimate}
\end{equation}
Indeed shifting a differentiated input spends one further derivative,
$\|U_L\|_{H^{j+2}}\le C_jL^{j+3/2}$; the phase Taylor error spends
three and has factor $h^2$; and the $O(h)$ coefficient error spends
only one. Finite fixed-bank multiplication has uniform bounds. These
facts prove the displayed estimate directly on the exact polynomial
inputs. Since $hL_h<1$, the second term is absorbed by the first.

By \cref{v17:Fejer-bounds}, $\overline VU_L$ differs from the fixed smooth
vector $-r_1v_1/2$ in $L^2$ by $O(L^{-1/2})$.
The finite coefficient inverse and Hodge projection are uniformly $L^2$
bounded, and on that \emph{fixed smooth vector} their consistency is
$O(h)$ by \cref{v17:consistency}. Applying them gives
\[
 \|\mathcal I_hZ_hu_h-ZW_1\|
             \le C(L_h^{-1/2}+hL_h^{3/2}+h).
\]
The same argument for the uniformly bounded flux metric inverse, with
its smooth-test consistency on $ZW_1$, bounds the difference of flux
energies by this quantity. Add the diagonal contribution $Ch^2L_h^3$.
With $L_h\asymp h^{-1/2}$ the errors are at most $Ch^{1/4}$, proving
\eqref{v17:explicit-recovery}. The scalar $L^2$ error is the first part
of the preceding lemma. Strong convergence of $f_h$ in
\eqref{v17:forcing} and the new certified pairing prove the last assertion.
\end{proof}

\section{The genuine geometric output allows this rough recovery}
\label{v17:observation}

Let $\mathcal O_h=\mathfrak E_hT_h^\dagger V_h$ be the unchanged
lapse-to-curvature observation. The following applies it \emph{after}
the true least-squares transverse inverse, not after the pointwise inverse.

\begin{theorem}[A vanishing output despite a divergent sufficient lapse norm]
\label{v17:visible-recovery}
For the explicit $u_h$ of \cref{v17:quantitative-recovery}, set
$n_h^{\rm test}=h^{-1}u_h$. Then
\begin{equation}
 \boxed{\|\mathcal O_h n_h^{\rm test}\|_h\le Ch^{1/4},\qquad
 h^2\|n_h^{\rm test}\|_{\dot H_h^3}\asymp h^{-1/4}.}
 \label{v17:geometric-recovery}
\end{equation}
Thus the sufficient lapse bound in \eqref{v13:final-output} is not a
necessary condition for this output to vanish. The exact shift observation
criterion of V10 remains unchanged.
\end{theorem}
\begin{proof}
The full-frequency estimates of V10 and V13 give
$\|\mathfrak E_hT_h^\dagger Y\|\le Ch\|\iota_o^*Y\|_{H_h^2}$.
For $Y=V_hu_h/h$, this bounds the desired norm by
$C\|V_hu_h\|_{H_h^2}$. The limiting source has a fixed smooth part and
an error bounded in $H^2$ by $CL^{3/2}$, from the last sentence of
\cref{v17:Fejer-bounds}. Equation \eqref{v17:source-recovery-estimate}
at $j=2$ therefore yields
\[
 \|V_hu_h\|_{H_h^2}
       \le C\left(h(1+L_h^{3/2})+h^2L_h^{7/2}\right)
       \le Ch^{1/4}.
\]
No cancellation was assumed inside an unknown inverse; it is the exact
source cancellation before applying the already controlled transverse
pseudoinverse. For the second assertion, the two-sided polynomial norm
estimate gives
$h^2\|n_h^{\rm test}\|_{\dot H_h^3}
=h\|U_{L_h}\|_{\dot H^3}\asymp hL_h^{5/2}\asymp h^{-1/4}$.
\end{proof}

These are tests and a recovery sequence, \emph{not} a solution of
$\mathscr S_hn_h=g_h$. One may add V13's actual direct transverse response
$\beta_{0,h}$ and set
$\beta_h^{\rm test}=\beta_{0,h}-T_h^\dagger V_hn_h^{\rm test}$.
This pair solves the original leading \emph{transverse} row exactly and
has total geometric output $O(h^{1/4})$. Its scalar row has not been
proved to equal the actual target. No nonlinear initial datum, trajectory,
or physical time interval is obtained from this comparison construction.

\section{What is now fixed, and what remains at zero spectral parameter}
\label{v17:status}

The full normalized scalar regulator now has a proved closed-form limit,
with a smooth core and all-cutoff weak lower/strong recovery convergence.
At the level of forms and fixed positive-parameter resolvents, the limit is
independent of the two odd-grid parity classes; parity-independent behavior
of the zero-parameter inverse is not asserted. Its fixed positive-parameter
resolvents converge
strongly on the actual normalized source, with convergent energies and
fluxes. The zero-parameter inverse is a separate limit; the lower
susceptibility statement retains that distinction explicitly.

The new source-specific domain example also changes the regularity target.
A characteristic jump belongs to the domain of the full scalar operator,
has an analytic operator image, and has a certified nonzero pairing with
the actual forcing. Its explicit original-regulator recovery has bounded
normalized energy and a vanishing geometric output at the $h^{-1}$ lapse
scale, even though the simple sufficient three-derivative lapse norm
diverges. This does not assert that the actual solution is that recovery,
but it rules out treating all nonsmooth limiting lapse components as the
forbidden coordinate-aligned tensor modes of V16.

The substantive unresolved step is a zero-spectral-parameter estimate for
the actual $g^\circ$ and its finite counterparts, in this identified flux
and observation domain. Strong resolvent convergence at fixed $\rho$ is
not a bound allowing $\rho\downarrow0$ uniformly in $h$. Nor does form
convergence imply convergence of the third-order geometric observation on
arbitrary energy-bounded sequences; V14 already supplies the relevant
warning. The higher weak signed system, actual finite-amplitude rate
selection, nonlinear propagation, differentiated curvature, and a common
physical lifespan remain separate requirements. No raw-action Einstein
continuum theorem or cofinal failure theorem is asserted.

\paragraph{Verification and source scope.}
The exact verifier checks the source-null sheet contractions, the rational
nonzero divergence, the Fejer difference identity, all derivative exponent
payments, and the actual-source jump pairing through a recomputed rational
residual and output polynomial. It retains the inherited original source
table rather than claiming a new regeneration of that table. Independent
floating diagnostics use the complete original phase operators, variable
coefficient transverse solve and diagonal regulator on the explicit recovery
sequence; they are not estimates for the unregularized scalar inverse.
The closed-graph, form-limit and resolvent arguments are written analytic
proofs, not conclusions drawn from these finite diagnostics. A targeted
source audit found the earlier one-dimensional sawtooth example but not
these characteristic sheets, complete form limit or observed recoveries;
no exhaustive priority claim is made. All sixteen earlier research bodies
remain byte-for-byte unchanged, with their corrections and limitations.

\begin{thebibliography}{9}
\raggedright
\setcounter{enumiv}{13}
\bibitem{v17:Ruppenthal}
J. Ruppenthal, \emph{Friedrichs' extension lemma with boundary values and
applications in complex analysis}, arXiv:0803.0092 (2008).
\url{https://arxiv.org/abs/0803.0092}.
This supplies context for minimal/maximal graph approximation. The projected
periodic flux and its commutator argument are proved here.
\bibitem{v17:Mosco}
U. Mosco, \emph{Convergence of convex sets and of solutions of variational
inequalities}, Advances in Mathematics \textbf{3} (1969), 510--585.
\href{https://doi.org/10.1016/0001-8708(69)90009-7}
{doi:10.1016/0001-8708(69)90009-7}.
This supplies context for weak-lower/strong-recovery convergence; the full
finite scalar convergence and its limits of interpretation are proved above.
\end{thebibliography}
% END IMMUTABLE EXACT-ACTION BRIDGE V17 BODY


% Future iterations append here.
% BEGIN IMMUTABLE EXACT-ACTION BRIDGE V18 BODY
\clearpage
\renewcommand{\thesection}{\arabic{section}}
\renewcommand{\theHsection}{bridgeV18.\arabic{section}}
\setcounter{section}{144}
\section{Cross-programme transfer and the correct zero-parameter question}
\label{v18:context}

The four supplied parallel lineages suggest an upstream distinction rather
than an imported gravitational estimate. Exact conditional centering in the
Yang--Mills notes separates a controlled innovation norm from the cost of
returning to the original variables. The shell calculation retains every
source and measure term before taking its limit. The regenerative notes
replace a failing configuration graph by a full phase-space comparison,
without equating that comparison with an unproved nonlinear evolution.
The source-selection notes likewise distinguish an available algebraic
construction from the operation supplied by the original source.
For the present problem these distinctions point to the \emph{complete
flux domain and the actual source}, not another assertion of uniform
coercivity in a lapse-coordinate norm.

There is a concrete new finding in that direction. The limiting original
scalar form has a nonconstant zero-energy mode. It is a two-chamber sign
function, annihilated by the unprojected lapse source before the longitudinal
projection is taken. The actual limiting forcing is exactly orthogonal to
this mode, chamber by chamber. We classify the kernel of the unprojected
source, construct a recovery sequence through the full finite scalar
energy, and show that its geometric output also tends to zero at the
natural inverse-cutoff scale. A separate source decomposition isolates
an already energy-bounded part of the forcing and reduces its remaining
energy-dual question to an explicit equation with a finite Fourier stress.
None of these statements assumes that the zero-parameter inverse exists.

\paragraph{What transfers, and what does not.}
The source audit uses the actual supplied versions, with literal labels:
\begin{center}\small
\begin{tabularx}{\textwidth}{@{}>{\raggedright\arraybackslash}p{.19\textwidth}X@{}}
\toprule
Source & Transfer used and scope retained\\\midrule
Yang--Mills V43 \cite{v18:YM43} &
\src{thm:v43-triangular}, \src{thm:v43-free-history},
\src{prop:v43-norm-price}: exact centering retains the action, the old
measure factors, and the norm of the inverse coordinate map. Its
innovation floor is not a floor for the present scalar operator.\\[.3em]
Shell mixing V16 \cite{v18:Shell16} &
\src{prop:v14-resolvent}, \src{thm:v16-flat}, \src{thm:v16-limit}:
the stated coercive transverse inverse and complete measured coefficient
have different domains from this characteristic lapse form. Their
source-accounting rule is useful; their coercivity is not borrowed.\\[.3em]
Regenerative V21 \cite{v18:Regen21} &
\src{v21:thm:caustic}, \src{v21:thm:quadratic},
\src{v21:thm:all-modes}: control of a full positive observable can survive
a singular coordinate graph. These are Wilson/phase-space results, not
an exact-action gravitational lifespan theorem.\\[.3em]
SM selection V17 \cite{v18:SM17} &
\src{sel17:thm:marginals}, \src{sel17:cor:old-covariance}:
restricted observations or an invariant algebra do not identify the complete
original operation. Here the full forcing, constants, and output remain
in the same source equations.\\
\bottomrule
\end{tabularx}
\end{center}
The review is a targeted dependency audit of the relevant results and their
terminal scope statements, not an independent verification of all four
cumulative theories. No source supplies the missing zero-parameter scalar
estimate for this action. The new results below are proved using its own
coefficients; no Yang--Mills gap, Hamilton--Jacobi evolution, quantum
instrument, or modified gravitational flow is substituted.

The inherited mathematical inputs are V15's complete prepared source table
and effective-source definition, V13--V14's original transverse inverse and
flux metric, and V17's closed maximal graph and finite-form convergence.
The new kernel and source identities do not repeat the nonzero sawtooth
energy in V17. The new sign function has \emph{zero} raw source and energy.

\section{The remaining energy-dual source is an explicit finite-stress equation}
\label{v18:source}

Retain angle derivatives $\partial_i^\circ$ on the unit three-torus,
normalized Haar measure, and the unit limiting complementary coefficients
$b^\circ,r^\circ$ of \eqref{v15:limit-coefficients}. Suppress their superscripts
in this body and set
\[
 \begin{gathered}
 M=M_b,\qquad N_r=M_r,\qquad A=M^{-1}N_r,\\
 P=P^\circ,\quad Q=I-P,\qquad K=QM^{-1}M^{-*}Q|_{\Ran Q}.
 \end{gathered}
\]
The letter $A$ here is the smooth \emph{coefficient} matrix, not the scalar
operator $\mathscr A$. Both gradients and all three constant vectors belong
to $\Ran Q$. Use
\[
 D u=A\nabla^\circ u,\qquad Z=QD,\qquad
 \mathcal D=\{u\in L_0^2:Zu\in L^2\},\qquad
 a(u,v)=\tfrac14\ip{Zu}{K^{-1}Zv}.
\]
The raw derivative $D$ may be only a distribution on $\mathcal D$; this is
allowed. V17 proves that $Z$ is closed, has a smooth core, and defines the
nonnegative self-adjoint scalar operator $\mathscr A$ through this form.

Let $q,f^\circ,w^\circ,g^\circ$ be exactly the source objects in
\eqref{v15:q-definition}--\eqref{v15:g-feedback}. In particular
\[
 f^\circ=C^\circ(\Omega^\circ)^{-2}q,\qquad
 H^\circ=PM^*MP+Q,\qquad w^\circ=(H^\circ)^{-1}f^\circ\in\Ran P.
\]
The finite source table gives, by ordinary Parseval,
\begin{equation}
 \|f^\circ\|^2=\frac{594441}{65536}.
 \label{v18:f-norm}
\end{equation}
Indeed $q$ has the nine positive-frequency rows in
\eqref{v15:source-table}, $k\cdot q_k=0$, and the squared norm is
$2\sum_k |q_k|^2/|k|^2$. This is a check of the retained table, not a
recalculation of its original action coefficients.

\begin{theorem}[Exact splitting of the actual limiting forcing]
\label{v18:source-split}
Put $y=M^*Mw^\circ$. The original limiting source splits as
\begin{equation}
 \boxed{g^\circ=g_{\mathrm b}+g_{\mathrm r},\qquad
 g_{\mathrm b}=-2Z^*Qy,\qquad
 g_{\mathrm r}=-2D^*f^\circ
       =2\diver^\circ(A^*f^\circ).}
 \label{v18:g-split}
\end{equation}
These are smooth mean-zero functions. The first term satisfies the proved
energy-dual bound
\begin{equation}
 |\ip{g_{\mathrm b}}u|
 \le4\|K^{1/2}Qy\|\sqrt{a(u,u)}
 \qquad(u\in\mathcal D).
 \label{v18:bounded-part}
\end{equation}
Thus the unresolved energy-dual part is determined by the explicit finite
Fourier field $f^\circ$ and the explicit analytic matrix $A$, with no unknown
transverse inverse in its formula. The solution of the transverse inverse
occurs only in the already bounded part.
\end{theorem}
\begin{proof}
On $\Ran P$ the equation $H^\circ w^\circ=f^\circ$ says $Py=f^\circ$;
therefore $y=f^\circ+Qy$. Track the physical off-diagonal pairing: it counts
both symmetric entries, whereas the vectors acted on by $M$ have the ordinary
three-component pairing. Since $\overline V=-N_r\nabla^\circ/2$, its physical
adjoint satisfies
\[
 \overline V^*v=\diver^\circ(N_r^*v).
\]
Consequently the original $g^\circ=2\overline V^*Mw^\circ$ is
$-2D^*M^*Mw^\circ=-2D^*y$. Substituting the decomposition of $y$, and
using $D^*Q=Z^*$, proves \eqref{v18:g-split} with its signs and factors.
Analyticity follows from the inherited analytic-weight bound for $w^\circ$
and the analytic coefficients. The divergence expressions have zero mean.
For a smooth $u$,
\[
 |\ip{g_{\mathrm b}}u|
 =2|\ip{Qy}{Zu}|
 \le2\|K^{1/2}Qy\|\|K^{-1/2}Zu\|.
\]
The last factor is $2\sqrt{a(u,u)}$. V17's smooth graph core extends the
identity to $\mathcal D$, even if $Du$ itself is not $L^2$.
\end{proof}

For precision, the exact remaining range question can be stated without a
coordinate-norm inverse. For $g\in L_0^2$ define
\[
 \mathcal E^*(g)=\sup_{u\in\mathcal D}
                      \{2\ip gu-a(u,u)\}\in[0,+\infty].
\]
\begin{proposition}[Energy-dual range and the source-specific stress equation]
\label{v18:dual}
One has $\mathcal E^*(g)<\infty$ if and only if $g=Z^*j$ for some
$j\in\Ran Q\cap L^2$, in the distributional sense. When finite,
\begin{equation}
 \mathcal E^*(g)=4\min_{Z^*j=g}\|K^{1/2}j\|^2.
 \label{v18:dual-norm}
\end{equation}
For the \emph{actual} source $g^\circ$, this property is equivalent to the
existence of $j\in\Ran Q\cap L^2$ solving
\begin{equation}
 \boxed{\diver^\circ\{A^*(j+2f^\circ)\}=0.}
 \label{v18:stress-equation}
\end{equation}
This last equation is a remaining equation, not a solution asserted here.
\end{proposition}
\begin{proof}
Put $T=K^{-1/2}Z/2$, so $a(u,u)=\|Tu\|^2$. Finiteness of the supremum,
by optimizing over every real multiple of $u$, is equivalent to a bound
$|\ip gu|\le C\|Tu\|$. It implies that $\ip gu$ is a well-defined bounded
functional of $Tu$, also when $T$ has a kernel. Extend it continuously to
$\overline{\Ran T}$ and apply Hilbert-space Riesz representation. The
minimum-norm representative $v\in\overline{\Ran T}$ obeys
$g=T^*v=Z^*K^{-1/2}v/2$, and the supremum equals $\|v\|^2$ by density
of the range. Set $j=K^{-1/2}v/2$. Conversely any such $j$ gives the required
bound, and orthogonal projection onto $\overline{\Ran T}$ minimizes its
weighted norm. This proves \eqref{v18:dual-norm}, including attainment.
A distributional equality on smooth tests is sufficient: the smooth core
extends it to the graph domain and identifies the adjoint there.

The class of energy-bounded functionals is a vector space, and
\eqref{v18:bounded-part} already puts $g_{\mathrm b}$ in it. Thus
$g^\circ$ is in that class exactly when $g_{\mathrm r}$ is. The condition
$Z^*j=g_{\mathrm r}=-2D^*f^\circ$, with $Qj=j$, is
$-\diver^\circ(A^*j)=2\diver^\circ(A^*f^\circ)$, exactly
\eqref{v18:stress-equation}. If it is solved, the stress for the full source
is $j-2Qy$. This is an elementary closed-range-functional argument; it is not
claimed as a new abstract factorization theorem or a proof that the range
is closed in $L^2$.
\end{proof}

\section{A nonconstant kernel of the actual unprojected source}
\label{v18:kernel}

Define the real analytic scalar
\begin{equation}
 \begin{split}
 F(\theta)&=\cos\theta_1+\cos\theta_2+\cos\theta_3+
                                  \cos(\theta_1+\theta_2+\theta_3)\\
 &=4\cos\frac{\theta_2+\theta_3}{2}
       \cos\frac{\theta_1+\theta_3}{2}
       \cos\frac{\theta_1+\theta_2}{2},\qquad
 \kappa(\theta)=\operatorname{sgn}F(\theta).
 \end{split}
 \label{v18:kappa}
\end{equation}
The values on $F=0$ are immaterial for $L^2$ and distributions. The product
of the three half-angle signs is periodic in each original angle: adding
$2\pi$ to one angle reverses two factors. This $\kappa$ is an analytical
lapse comparison function. It is not a new protected orientation Read,
microscopic binary axiom, or nonsmooth physical initial metric.

\begin{theorem}[An exact nonconstant zero-energy mode]
\label{v18:kernel-theorem}
The function $\kappa$ is real, mean-zero and unit-norm. For the actual
limiting coefficient of the complementary family,
\begin{equation}
 \boxed{M_r\nabla^\circ\kappa=0\quad\hbox{in distributions},\qquad
 Z\kappa=0,\qquad
 \kappa\in D(\mathscr A),\quad \mathscr A\kappa=0.}
 \label{v18:zero-source}
\end{equation}
It follows that the normalized limiting scalar form has no positive
$L_0^2$ spectral floor. The complete projected-flux kernel is not asserted
to be exhausted by this one mode.
\end{theorem}
\begin{proof}
The simultaneous translation $\theta\mapsto\theta+\pi(1,1,1)$ reverses
$F$ and $\kappa$. It preserves Haar measure. Since the zero set of the
nonzero trigonometric polynomial $F$ has measure zero, $|\kappa|=1$
almost everywhere, proving its mean and norm assertions.

There is an exact smooth coefficient identity
\begin{equation}
 M_r\nabla^\circ F=F a_c,\qquad
 a_c=\tfrac12(\cos\theta_2-\cos\theta_1,
              \cos\theta_1-\cos\theta_3,
              \cos\theta_3-\cos\theta_2)^T.
 \label{v18:smooth-identity}
\end{equation}
For example its first row is
$r_2\partial_1^\circ F+r_1\partial_2^\circ F$; substitute
$r_2=(\sin\theta_3-\sin\theta_1)/2$,
$r_1=(\sin\theta_2-\sin\theta_3)/2$ and
$\partial_j^\circ F=-\sin\theta_j-\sin(\theta_1+\theta_2+\theta_3)$.
The sine sum and difference identities give
$F(\cos\theta_2-\cos\theta_1)/2$. The other two rows give the other
entries of \eqref{v18:smooth-identity}. This also verifies the full vector,
not only its divergence or a longitudinal projection.

Set $\kappa_\epsilon=F/(F^2+\epsilon^2)^{1/2}$. Then
$\kappa_\epsilon\to\kappa$ in $L^2$ and
\[
 M_r\nabla^\circ\kappa_\epsilon
   =a_c\frac{F\epsilon^2}{(F^2+\epsilon^2)^{3/2}}\longrightarrow0
                                      \quad\hbox{in }L^2.
\]
The scalar factor is bounded uniformly and tends to zero wherever $F\ne0$;
dominated convergence proves the claim. Differentiation and smooth
multiplication are continuous into distributions, giving the first identity
of \eqref{v18:zero-source}. Multiplying by the smooth $M^{-1}$ and applying
$Q$ gives $Z\kappa=0$, so $\kappa\in\mathcal D$. The form pairing
$a(\kappa,v)$ is zero for every $v\in\mathcal D$; the defining operator
criterion in V17 therefore puts $\kappa$ in $D(\mathscr A)$ with zero
image. Since its mean is zero and its norm is one, no positive mean-zero
floor is possible. This argument does not replace the finite scalar matrices
by a singular one, nor classify additional solutions of $QD u=0$ whose raw
flux $Du$ may be nonzero and transverse.
\end{proof}

\section{Two chambers and a complete kernel classification before projection}
\label{v18:chambers}

Let
\[
 H_1=\{\theta_2+\theta_3=\pi\bmod 2\pi\},\quad
 H_2=\{\theta_1+\theta_3=\pi\bmod 2\pi\},\quad
 H_3=\{\theta_1+\theta_2=\pi\bmod 2\pi\}.
\]
Their complement is denoted $\mathcal C$. These are the characteristic
sheets on which the sign in \eqref{v18:kappa} can jump.

\begin{theorem}[Exactly two raw-source chambers]
\label{v18:raw-kernel}
The set $\mathcal C$ has exactly two connected components, on which
$\kappa$ takes its two signs. Moreover
\begin{equation}
 \boxed{\{u\in L^2(\T^3):M_r\nabla^\circ u=0
                     \hbox{ in distributions}\}
                   =\operatorname{span}_{\R}\{1,\kappa\}.}
 \label{v18:raw-classification}
\end{equation}
In particular the smooth raw-source kernel consists only of constants.
This classifies the unprojected source, not the larger possible kernel of $Z$.
\end{theorem}
\begin{proof}
Consider the integer torus map with matrix
\[
 B_c=\begin{pmatrix}0&1&1\\1&0&1\\1&1&0\end{pmatrix},\qquad \det B_c=2.
\]
It is a two-sheet covering of the torus by itself. Its nontrivial deck
translation is $\pi(1,1,1)$. Removing the three coordinate sheets at angle
$\pi$ from the target leaves a simply connected open cube. The preimage of
this cube is $\mathcal C$, so it has exactly two connected components,
each mapped bijectively to the cube. The deck translation interchanges
them and reverses $F$. Thus the two components are precisely its two signs.

The matrix determinant is $\det M_r=-2r_1r_2r_3$. In addition to the
three $H_i$, its zero set consists of the equality sheets
$E_1=\{\theta_2=\theta_3\}$, $E_2=\{\theta_1=\theta_3\}$,
$E_3=\{\theta_1=\theta_2\}$, all modulo $2\pi$.
Off this finite sheet arrangement, multiplication by the smooth inverse
of $M_r$ gives $\nabla^\circ u=0$ in distributions. Thus an $L^2$
solution is locally constant on every connected open cell.

Across a regular patch of $E_1$ away from all other sheets, its jump
contributes $(u_+-u_-)(e_2-e_3)\delta_{E_1}$ to the gradient, up to the
common positive normalization of the surface measure. There $r_1=0$
and $r_3=-r_2\ne0$, so
\[
 M_r(e_2-e_3)=(0,0,-2r_2)^T\ne0.
\]
The distributional equation forces $u_+=u_-$. The same direct computation
applies to the other equality sheets. A path in either connected component
of $\mathcal C$ can be slightly perturbed to cross this finite arrangement
only at such regular patches; intersections have codimension at least two.
It follows that all constants within each component agree. A nonzero
$L^2$ function cannot be supported only on the measure-zero sheets.
Hence every solution has one constant on each of the two components.
Conversely $1$ and $\kappa$ are raw-source null by
\cref{v18:kernel-theorem}. This proves \eqref{v18:raw-classification}.
A smooth function cannot have different constant values on the two sides
of a regular $H_i$ patch, so only constants are smooth solutions.
\end{proof}

An equivalent direct check of the allowed normals is useful below. Set
\[
 \begin{array}{lll}
 \xi_1=(0,1,1),&\xi_2=(1,0,1),&\xi_3=(1,1,0),\\
 v_1=(1,1,-1),&v_2=(1,-1,1),&v_3=(-1,1,1).
 \end{array}
\]
Then the actual matrix satisfies
\begin{equation}
                M_r\xi_i=r_iv_i,
 \qquad r_i|_{H_i}=0.                         \label{v18:normals}
\end{equation}
The vanishing is before projection and holds on whole sheet patches, not
only at the single characteristic point used in V14.

\section{An explicit Fourier kernel and exact source compatibility}
\label{v18:Fourier}

Let $\widehat\kappa(k)$ denote the coefficient of $e^{\mathrm i k\cdot\theta}$.
For $k\in\mathbb Z^3$ put
\[
 n(k)=(-k_1+k_2+k_3,\ k_1-k_2+k_3,\ k_1+k_2-k_3).
\]
\begin{proposition}[All kernel coefficients and their regularity]
\label{v18:Fourier-formula}
One has
\begin{equation}
 \boxed{\widehat\kappa(k)=
 \begin{cases}
 \displaystyle\frac8{\pi^3}\prod_{i=1}^3
                  \frac{\sin(\pi n_i(k)/2)}{n_i(k)},
                      & k_1+k_2+k_3\hbox{ odd},\\[2mm]
 0,                    & k_1+k_2+k_3\hbox{ even}.
 \end{cases}}
 \label{v18:Fourier-coefficients}
\end{equation}
For $s\ge0$, $\kappa\in H^s$ exactly when $s<1/2$.
\end{proposition}
\begin{proof}
Set $t_1=(\theta_2+\theta_3)/2$ and cyclically, and use the integer
matrix $A_c=2B_c^{-1}$ to map $t\mapsto\theta=A_ct$. Its determinant
is four and it preserves normalized Haar measure. Under this map the
three half-angle factors become $\sigma(t_i)$, where
$\sigma(t)=\operatorname{sgn}(\cos t)$. Also
$e^{-\mathrm i k\cdot\theta}=e^{-\mathrm i(A_c^Tk)\cdot t}$ and
$A_c^T=A_c$, giving $n(k)$ above. The one-dimensional coefficient is
$\widehat\sigma(n)=2\sin(\pi n/2)/(\pi n)$ for odd $n$ and zero for
even $n$, including zero. The product integral proves the formula.
All $n_i(k)$ have the parity of $k_1+k_2+k_3$. The odd case cannot have a
zero denominator. Conversely every triple of odd integers $n$ gives an
integer $k=B_cn/2$, so the parametrization has no missing odd triples.

The square sum with weight $(1+|k|^2)^s$ is bounded for $s<1/2$ by a
constant times the sum of three products of the convergent one-dimensional
series $\sum_{n\ {
m odd}}|n|^{-2+2s}$ and $\sum|n|^{-2}$.
For failure at $1/2$, choose $k=(0,m,m+1)$ with $m\ge1$.
Then $n(k)=(2m+1,1,-1)$ and
$|\widehat\kappa(k)|=8/[\pi^3(2m+1)]$. The $H^{1/2}$ square sum over
these modes diverges as a harmonic series. Higher orders also fail.
\end{proof}

\begin{theorem}[The complete limiting forcing obeys both chamber balances]
\label{v18:orthogonality}
For the actual source $g^\circ$ and both terms of \eqref{v18:g-split},
\begin{equation}
 \boxed{\ip\kappa{g^\circ}=0,\qquad
 \ip\kappa{g_{\mathrm b}}=\ip\kappa{g_{\mathrm r}}=0.}
 \label{v18:orthogonality-eq}
\end{equation}
In particular, with $\mathcal C_\pm$ the two chambers,
\begin{equation}
                  \int_{\mathcal C_+}g^\circ
                =\int_{\mathcal C_-}g^\circ=0.        \label{v18:chamber-balance}
\end{equation}
For every fixed $\rho>0$,
\begin{equation}
        \ip\kappa{(\mathscr A+\rho I)^{-1}g^\circ}=0.
 \label{v18:resolvent-kernel}
\end{equation}
\end{theorem}
\begin{proof}
The raw identity implies $D\kappa=0$ as well as $Z\kappa=0$.
All stresses $f^\circ,y,Qy$ in \eqref{v18:g-split} are smooth, so
integration by parts in distributions gives each of
\eqref{v18:orthogonality-eq} exactly. There is no limiting numerical
cancellation or assumption about a source symmetry. Mean-zero forcing and
$\kappa=\mathbf1_{\mathcal C_+}-\mathbf1_{\mathcal C_-}$ give the two
linear equations whose solution is \eqref{v18:chamber-balance}.
Finally $\mathscr A\kappa=0$ implies
$(\mathscr A+\rho I)^{-1}\kappa=\kappa/\rho$. Self-adjointness and
\eqref{v18:orthogonality-eq} prove \eqref{v18:resolvent-kernel}.
\end{proof}

The two balances are necessary conditions, not a sufficient Fredholm
alternative: the full projected flux may have additional null modes and
its range need not be closed. A source may annihilate this known kernel
and still fail the energy-dual bound of \cref{v18:dual}. Conversely this
nonconstant kernel is not evidence that the actual forcing has a pole on
it: the forcing is exactly orthogonal. Analytically separating
$\operatorname{span}\{\kappa\}$ in the self-adjoint space does not declare
it a physical gauge mode or remove an original finite equation.

\section{A finite recovery of the kernel through the original scalar form}
\label{v18:recovery}

The singular sign must not be inserted as a sampled discontinuous physical
record. We instead give trigonometric comparison arrays on the original
finite scalar space. For odd $L\ge1$ define the Fejer polynomial
\begin{equation}
 \sigma_L(t)=\frac4\pi\sum_{\substack{1\le n\le L\\n\ {\rm odd}}}
       \left(1-\frac n{L+1}\right)
       \frac{\sin(\pi n/2)}n\cos(nt),\qquad
 U_L(\theta)=\prod_{i=1}^3\sigma_L(t_i).
 \label{v18:UL}
\end{equation}
The half-angles are those of \cref{v18:Fourier-formula}. The product is
$2\pi$ periodic in each original coordinate, with only integer frequencies
$|k_i|\le L$. Its mean is exactly zero: the deck translation reverses each
factor and hence their product.

\begin{lemma}[Quantitative source-null approximation]
\label{v18:Fejer}
With normalized one-dimensional measure,
\begin{align}
 \|\cos t\,\sigma_L'(t)\|^2
      &=\frac{4L-2}{\pi^2(L+1)^2},
                     &\|\sigma_L'\|^2
      &=\frac{4L(L+2)}{3\pi^2(L+1)},                 \label{v18:oneD}\\
 \|U_L-\kappa\|&\le CL^{-1/2},
                     &\|M_r\nabla^\circ U_L\|&\le CL^{-1/2}.
                                                            \label{v18:raw-recovery}
\end{align}
For every integer $j\ge1$, $\|U_L\|_{H^j}\le C_jL^{j-1/2}$.
For every integer $j\ge0$,
$\|M_r\nabla^\circ U_L\|_{H^j}\le C_jL^{j-1/2}$.
Writing
\[
 N_L=\frac8{\pi^2}\sum_{\substack{1\le n\le L\\n\ {\rm odd}}}
       \frac{(1-n/(L+1))^2}{n^2},\qquad
 C_L=\frac8{\pi^2}\sum_{\substack{1\le n\le L\\n\ {\rm odd}}}
       \frac{1-n/(L+1)}{n^2},
\]
one also has the exact identity
\begin{equation}
 \|U_L-\kappa\|^2=1+N_L^3-2C_L^3.                  \label{v18:exact-error}
\end{equation}
\end{lemma}
\begin{proof}
The Fejer kernel is nonnegative and has integral one. Its convolution with
$\sigma=\operatorname{sgn}\cos$ gives \eqref{v18:UL}, so
$|\sigma_L|\le1$. The Fourier formula shows
$\sigma_L(t+\pi)=-\sigma_L(t)$ and
$\|\sigma_L-\sigma\|^2=O(L^{-1})$: the squared error on the retained odd
modes is $O(L^{-1})$, and the omitted reciprocal-square tail is also
$O(L^{-1})$.

For the first exact derivative identity, write $L=2m+1$ and differentiate
the displayed cosine series. Multiplication by $\cos t$ gives adjacent
sine frequencies. The nonzero coefficients of
$\pi\cos t\,\sigma_L'$ are
$4(-1)^j/(L+1)$ at $2j$, $1\le j\le m$, and
$2(-1)^{m+1}/(L+1)$ at $2m+2$. Orthogonality gives
$(8m+2)/(\pi^2(L+1)^2)=(4L-2)/(\pi^2(L+1)^2)$.
The other derivative has squared norm
$8\sum_{n\ {\rm odd}\le L}(1-n/(L+1))^2/\pi^2$; summing the squares
of the first $m+1$ odd integers in reverse order gives the second identity.

The integer covering change of variables used earlier preserves Haar
measure, and turns $U_L$ and $\kappa$ into products of independent
one-dimensional factors. Thus $\|U_L\|^2=N_L^3$ and
$\ip\kappa{U_L}=C_L^3$, proving \eqref{v18:exact-error}. Telescoping the
three product differences and using $|\sigma_L|,|\sigma|\le1$ proves
its stated order without a logarithmic loss.

By \eqref{v18:normals},
\[
 M_r\nabla^\circ U_L
  =\tfrac12\sum_i r_iv_i\,\sigma_L'(t_i)
                                \prod_{j\ne i}\sigma_L(t_j).
\]
Each $r_i$ equals $\cos t_i$ times a sine of a half-angle difference, with
the appropriate cyclic sign. Therefore the first identity in
\eqref{v18:oneD}, the covering change of variables, and $|v_i|=\sqrt3$
give the raw-source bound in \eqref{v18:raw-recovery}. In particular no
distributional cancellation at intersecting sheets is being inferred from a
sampled pointwise product. The derivative norm of $U_L$ is $O(\sqrt L)$,
by the second identity in \eqref{v18:oneD} and the product rule. Its Fourier
bandwidth is $L$ in each coordinate, so higher derivatives give the stated
$L^{j-1/2}$ bounds. The raw source has bandwidth at most $L+1$; apply the
same finite Fourier inequality to its $O(L^{-1/2})$ norm to prove its
$H^j$ bounds.
\end{proof}

\begin{theorem}[Vanishing full finite energy and vanishing observed output]
\label{v18:finite-recovery}
Let $L_h$ be the largest positive odd integer not exceeding $h^{-1/2}$,
and put $u_h=\mathcal S_hU_{L_h}$. On the original complementary family
with $\alpha_h=(1,\theta_h,1)$ and $1<\theta_h<1+h$, for all sufficiently
fine odd grids,
\begin{equation}
 \boxed{\|\mathcal I_hu_h-\kappa\|\le Ch^{1/4},\qquad
 \|Z_hu_h\|_h\le Ch^{1/4},\qquad
 a_h(u_h)\le Ch^{1/2}.}
 \label{v18:finite-energy}
\end{equation}
Here $a_h$ and $Z_h$ are the complete normalized finite scalar form and
flux of V17, not just their smooth limits. The actual lapse-to-curvature
observation $\mathcal O_h=\mathfrak E_hT_h^\dagger V_h$ satisfies
\begin{equation}
               \boxed{\|\mathcal O_h(h^{-1}u_h)\|_h\le Ch^{1/4}.}
 \label{v18:finite-output}
\end{equation}
On the rationally selected invertible subfamily,
$\lambda_{\min}(\mathscr A_h)\le Ch^{1/2}$, although every such finite
matrix is positive. These arrays are comparison tests, not asserted
solutions of the actual scalar equation or nonlinear initial metric data.
\end{theorem}
\begin{proof}
Eventually the support box of $U_{L_h}$, expanded by the fixed source
coefficients, lies strictly within the odd-grid Fourier box. Sampling is
therefore exact on the input polynomial and on the finite products used
before the controlled inverses. All finite operators still act on their
full spaces. The first bound of \eqref{v18:finite-energy} follows from
\cref{v18:Fejer} and $L_h\asymp h^{-1/2}$.

We record the consistency costs needed here, to avoid applying smooth-test
convergence to an $h$-dependent function without a bound. For a trigonometric
polynomial of coordinate degree $L$ in this range, the original phase-symbol
estimate is $\|\delta_h/\lambda_* -\partial^\circ\|_{H^{j+3}\to H^j}
\le C_jh^2$. A literal shift costs at most $C_jh$ with one derivative.
The calibrated coefficients differ from their unit limits by $O(h)$ in
all fixed orders, including $\theta_h-1$. It follows, in raw off-diagonal
norms equivalent to the physical ones, that
\begin{equation}
 \left\|\frac{\mathcal I_hV_hu_h}{h\lambda_*^2}
                       -\overline VU_L\right\|_{H^j}
 \le C_j\{hL^{j+3/2}+h^2L^{j+5/2}\},\quad j=0,2.
 \label{v18:source-consistency}
\end{equation}
For example $V_h/h$ is a bounded unit-profile coefficient times one shifted
phase derivative; the shift spends two derivatives of the input in this
comparison, while the phase error spends three. Their input bounds are
exactly those in \cref{v18:Fejer}. The $O(h)$ coefficient error spends
only one and is covered by the displayed bound. These are coefficient
estimates, not a replacement of $V_h$ by $\overline V$.
Also the diagonal product-defect source on this unfurled input obeys
\begin{equation}
              \|\mathscr D_{h,\alpha_h}u_h\|_h
                         \le Ch^2L^{3/2}.             \label{v18:diagonal}
\end{equation}
This is the sine-addition bound for a unit-profile defect, with its two
input derivatives charged.

The exact finite energy of V17 is
\[
 a_h(u)=\frac{\|\mathscr D_hu\|_h^2}{2h^2\lambda_*^4}
                 +\tfrac14\ip{Z_hu}{K_{Q,h}^{-1}Z_hu}_h.
\]
Its normalized flux is a uniformly bounded coefficient inverse and Hodge
projection applied to $V_hu/(h\lambda_*^2)$, with the fixed Frobenius
factor retained. Thus \eqref{v18:raw-recovery} and
\eqref{v18:source-consistency} give
\[
 \|Z_hu_h\|_h\le C\{L^{-1/2}+hL^{3/2}+h^2L^{5/2}\}.
\]
The metric $K_{Q,h}^{-1}$ is uniformly bounded, and the normalized diagonal
energy is at most $Ch^2L^3$. Substituting $L=L_h\asymp h^{-1/2}$ proves
the other bounds in \eqref{v18:finite-energy}.

For the output use the actual estimate
$\|\mathfrak E_h T_h^\dagger Y\|_h\le Ch\|\iota_o^*Y\|_{H_h^2}$
from V13 and V10. With $Y=V_hu_h/h$, \eqref{v18:source-consistency} and
the $H^2$ raw-source bound in \cref{v18:Fejer} give
\[
 \|\mathcal O_h(h^{-1}u_h)\|_h
 \le C h\{L^{3/2}+hL^{7/2}+h^2L^{9/2}\}\le Ch^{1/4}.
\]
This retains the nonlocal true pseudoinverse. Finally $\|u_h\|_h\to1$,
so its Rayleigh quotient proves the stated eigenvalue upper bound. No
scalar inverse is used in this construction.
\end{proof}

\section{A quantitative source-dependent resolvent test on the new mode}
\label{v18:resolvent}

Keep the actual $f_h=g_h/(h\lambda_*^4)$, its inherited $O(h)$ convergence
to $g^\circ$, and the normalized original scalar operator $\mathscr A_h$.
For $\rho>0$ write
\[
             v_{h,\rho}=(\mathscr A_h+\rho I)^{-1}f_h.
\]
The added $\rho$ is an analytical resolvent parameter, not a physical mass
or a change to the original action.

\begin{theorem}[The known chamber mode is not excited at a controlled moving parameter]
\label{v18:resolvent-bound}
The comparison sequence of \cref{v18:finite-recovery} satisfies
\begin{equation}
                  |\ip{u_h}{f_h}_h|\le Ch^{1/4}.
 \label{v18:finite-pairing}
\end{equation}
For $0<\rho\le1$, independently of the cutoff,
\begin{equation}
 \boxed{|\ip\kappa{\mathcal I_hv_{h,\rho}}|
       \le Ch^{1/4}(\rho^{-1}+\rho^{-3/2}).}
 \label{v18:rho-bound}
\end{equation}
In particular this component tends to zero if $\rho_h\downarrow0$ and
$\rho_h/h^{1/6}\to\infty$. This is a result about one known null mode,
not convergence of the entire moving-parameter resolvent or the inverse at
$\rho=0$.
\end{theorem}
\begin{proof}
The fixed source has bounded $L^2$ norm and the inherited smooth expansion
gives $\|\mathcal I_hf_h-g^\circ\|\le Ch$, including the smooth sampling
tail. Pair with $\mathcal I_hu_h$, use
\eqref{v18:orthogonality-eq} and the $O(h^{1/4})$ recovery error, and obtain
\eqref{v18:finite-pairing}.

The finite spectral theorem for a nonnegative operator gives
\[
 \|v_{h,\rho}\|_h\le\rho^{-1}\|f_h\|_h,\qquad
 a_h(v_{h,\rho})\le\frac1{4\rho}\|f_h\|_h^2,
\]
since $t/(t+\rho)^2\le1/(4\rho)$ for $t\ge0$.
Test the \emph{complete} resolvent equation against $u_h$. The form
Cauchy--Schwarz inequality and \eqref{v18:finite-energy} give
\[
 \rho|\ip{u_h}{v_{h,\rho}}_h|
 \le |\ip{u_h}{f_h}_h|
        +\sqrt{a_h(u_h)a_h(v_{h,\rho})}
 \le Ch^{1/4}(1+\rho^{-1/2}).
\]
Replacing $\mathcal I_hu_h$ by $\kappa$ costs at most
$Ch^{1/4}\|v_{h,\rho}\|_h\le Ch^{1/4}/\rho$. This proves
\eqref{v18:rho-bound}. Finally
$h^{1/4}\rho_h^{-3/2}=(h^{1/6}/\rho_h)^{3/2}$ tends to zero under the
stated scale condition, and controls the other term when $\rho_h\le1$.
No uniform scalar spectral gap has been assumed.
\end{proof}

Source compatibility and small compressed energy do not license simply
erasing $u_h$ from the finite solve. It is not an exact finite eigenvector,
and the complementary components can couple to it. V15's exact low-source
Schur formula already records the needed coupling and lifted observation;
we do not count the same Schur identity as a new theorem here. In particular,
\eqref{v18:rho-bound} is not a bound for
$\mathcal O_h\mathscr S_h^{-1}g_h$. The source-specific stress equation
\eqref{v18:stress-equation}, the rest of the projected kernel, and the
zero spectral tail remain distinct questions.

\section{Direction after the transfer audit and verification scope}
\label{v18:status}

The parallel programmes help sharpen the target, but do not close it by
analogy. The limiting lapse form is not coercive even on mean-zero $L^2$:
a specific two-chamber mode is in its actual kernel. Borrowing an auxiliary
Yang--Mills innovation floor would contradict this calculation. At the same
time the original source exactly annihilates that mode, including the
chamber means; a singular operator is not by itself a source pole. The
explicit finite recovery has vanishing \emph{full} normalized energy and
vanishing observed geometric output at the $h^{-1}$ scale. This makes it
possible to distinguish a physically visible obstruction from an analytical
coordinate null direction, without declaring a new gauge quotient.

The source split is a useful next reduction. Its transverse-normal solution
appears only in an already bounded energy-dual term. The remaining finite
stress $f^\circ$ is explicit, and the equation
$\diver^\circ[A^*(j+2f^\circ)]=0$, $Qj=j\in L^2$, asks for the missing
longitudinal/constant compensation. It has not been solved in this body.
Likewise the raw-source kernel classification is not a classification of all
projected-flux null modes or a proof of a closed range. Either a further
actual-source compensation or a zero-spectral source estimate is still needed
before the unregularized response can be controlled.

No new higher weak signed rank, weak-rate preparation, nonlinear exact-action
trajectory, differentiated spacetime bound, or common physical lifespan is
asserted. The original finite action and its literal observations remain
unchanged. In particular this is neither an exact-action Einstein continuum
theorem nor a counterexample to one.

\paragraph{Reproducibility.}
The exact verifier checks the trigonometric source factorization, every normal
contraction, the integer covering and Fourier parity, the Fejer derivative
identities, complete rational Fourier convolutions on finite test banks, and
the explicit stress norm. A separate rational matrix regression checks the
adjoint and energy-dual normalization in the source split; it is labelled an
algebraic regression rather than original-action data. The all-mode identities,
raw-kernel classification, domain, and recovery estimates have the written
proofs above, not finite samples as substitutes.

Floating diagnostics at stated odd grids use the retained original finite
source and true transverse inverse to evaluate the recovery, full energy,
flux, and geometric output. They do not solve the scalar inverse, certify a
new rank, or approximate a nonlinear spacetime. The supplied source audit
lists the inspected ranges of all four related files, their relevant theorem
labels and what was not imported. No external programme is presented as an
independent validation of the predecessor's full stationary or continuum
theory. The seventeen earlier marked bodies remain byte-for-byte unchanged,
with source and output hashes in the preservation manifest.

\begin{thebibliography}{9}
\raggedright
\setcounter{enumiv}{12}
\bibitem{v18:YM43}
\emph{Renewal Yang--Mills Mass-Gap Research Notes}, supplied cumulative V43,
Section~\src{sec:v43}, particularly \src{thm:v43-triangular},
\src{thm:v43-free-history}, and \src{prop:v43-norm-price}.
This is a source for the exact-centering comparison and its boundaries, not a
gravity coercivity or nonlinear continuation input.
\bibitem{v18:Shell16}
\emph{Renewal Yang--Mills Shell Mixing Research Notes}, supplied cumulative V16,
\src{prop:v14-resolvent}, \src{thm:v16-flat}, and \src{thm:v16-limit}.
Its complete measured response and controlled transverse norm are kept in
their stated source spaces.
\bibitem{v18:Regen21}
\emph{Renewal Native Regenerative Stationarity Closure Notes}, supplied
cumulative V21, \src{v21:thm:caustic}, \src{v21:thm:quadratic}, and
\src{v21:thm:all-modes}. Its Wilson phase-space calculation is not identified
with the exact gravitational evolution considered here.
\bibitem{v18:SM17}
\emph{Renewal SM Source Selection Research Notes}, supplied cumulative V17,
\src{sel17:thm:marginals} and \src{sel17:cor:old-covariance}. Used for the
distinction between an actual complete source and a comparison construction;
no CP instrument is inserted into the gravitational law.
\end{thebibliography}
% END IMMUTABLE EXACT-ACTION BRIDGE V18 BODY


% Future iterations append here.
% BEGIN IMMUTABLE EXACT-ACTION BRIDGE V19 BODY
\clearpage
\renewcommand{\thesection}{\arabic{section}}
\renewcommand{\theHsection}{bridgeV19.\arabic{section}}
\setcounter{section}{152}
\section{The remaining stress equation and an essential zero threshold}
\label{v19:context}

The chamber balance in \cref{v18:orthogonality} removes the known raw-kernel
obstruction for the actual source. It does not imply that deleting that one
mode leaves a coercive operator. This body tests that possible next step on
the actual coefficients. The limiting scalar form has zero in its essential
spectrum even in the inversion-odd sector containing the forcing. The
chamber mode is inversion-even. No bounded finite-rank penalty opens a
positive floor in the sourced sector.

The mechanism is a quadratic zero of the raw coefficient at two specified
points, not a repetition of the interior high-frequency characteristics in
V11 or V14. It also gives a quantitative finite conclusion. A growing space
of original scalar arrays has uniformly small normalized energy, small
original geometric output at the inverse-cutoff lapse scale, and small
projection of the actual normalized source. The dimension is at least of
order $h^{-3/10}$, its energy is at most order $h^{2/5}$, and its observed
output is at most order $h^{2/5}$. These arrays are comparison tests, not
solutions of the scalar equation. Their source projection is not a spectral
projection. The distinction is retained throughout.

First we put V18's unsolved compensation equation into an equivalent
curl--divergence system with explicit trigonometric coefficients. Inversion
parity shows that, for this particular limiting forcing, compensation can
be sought with a mean-zero longitudinal stress. This is a derived reduction
of the unknown compensation, not deletion of the constant flux components
from the action or its scalar form. The reduced equation is not solved here.

\paragraph{Inherited inputs and verification boundary.}
The complete source table and smooth effective-source consistency are
\cref{v15:target-expansion,v15:g-expansion}. The original finite transverse
pseudoinverse and its all-frequency observation bound are V13; the complete
normalized scalar forms and their closed limiting realization are V17.
V18 supplies the raw chamber kernel, the actual source split, and its
energy-dual range equivalence. None is reproved as a new acquisition.
The spectral dimension argument is standard; a short proof in the present
form domain is included, with background in \cite{v19:Teschl}. The new
results concern the actual limiting coefficients and the original finite
operators, not a substitute elliptic action. No nonlinear exact history,
new higher weak rank, or zero-parameter inverse bound is inferred.

\section{An explicit curl--divergence compensation system and its parity}
\label{v19:stress}

Use angle derivatives and normalized Haar measure as in V18. Let
$M=M_b$, $N_r=M_r$, $A=M^{-1}N_r$, $D=A\nabla^\circ$, $Z=QD$ and
$K=QM^{-1}M^{-*}Q|_{\Ran Q}$. The known trigonometric source is $q$ in
\eqref{v15:source-table}, and $f^\circ=C^\circ(\Omega^\circ)^{-2}q$.
In particular
\[
 \curl^\circ f^\circ=q,\qquad \diver^\circ f^\circ=0,
 \qquad \overline{f^\circ}=0.
\]
Every derivative and equality in this section is distributional when its
input is only $L^2$. The fixed analytic matrix $M$ and its inverse are
bounded. The three components on which $M$ acts have the ordinary vector
pairing used in V18; no extra off-diagonal Frobenius factor is inserted.

\begin{proposition}[An equivalent first-order stress system]
\label{v19:mixed-system}
The actual source $g^\circ$ has finite energy-dual susceptibility if and
only if there is a real three-component field $\tau\in L^2$ such that
\begin{equation}
 \boxed{\curl^\circ(M^*\tau)=2q,\qquad
                    \diver^\circ(N_r^*\tau)=0.}
 \label{v19:mixed}
\end{equation}
Its compensating stress for the remainder in V18 is then
\begin{equation}
                  j=M^*\tau-2f^\circ,\qquad Qj=j.
 \label{v19:j-from-tau}
\end{equation}
Conversely $\tau=M^{-*}(j+2f^\circ)$. No differential inverse or
unspecified source term appears in the coefficients of \eqref{v19:mixed}.
\end{proposition}
\begin{proof}
V18 proves that finite susceptibility is equivalent to
$Qj=j\in L^2$ and $\diver^\circ[A^*(j+2f^\circ)]=0$.
On the torus an $L^2$ vector is in $\Ran Q$ exactly when its curl vanishes:
at a nonzero Fourier mode this means its coefficient is parallel to $k$;
its three constant coefficients are unrestricted. Substituting
\eqref{v19:j-from-tau} and $\curl f^\circ=q$ gives the first equation.
Since $A^*=N_r^*M^{-*}$, the divergence equation gives the second. The
bounded coefficient inverse proves equivalence of the $L^2$ requirements.
This retains the harmonic components rather than silently setting them to
zero in a general source problem.
\end{proof}

Write $\mathcal Iu(\theta)=u(-\theta)$, also componentwise on vectors.
This is an analytical involution on coefficient functions, not a physical
relabeling of a finite chronological source. Its odd scalar subspace is
$\mathcal H_-:=\{u\in L_0^2:\mathcal Iu=-u\}$.

\begin{theorem}[The forcing is odd, and compensation needs no harmonic stress]
\label{v19:parity}
For the unit limiting complementary coefficients,
\begin{equation}
 \mathcal I M=M,\quad \mathcal I N_r=-N_r,\quad
 \mathcal I f^\circ=-f^\circ,\quad
 \boxed{\mathcal I g^\circ=-g^\circ.}
 \label{v19:parities}
\end{equation}
The form domain is invariant under $\mathcal I$, and
$a(\mathcal Iu,\mathcal Iv)=a(u,v)$. Thus $\mathcal H_-$ reduces
$\mathscr A$ and contains its actual forced resolvents. The known chamber
mode $\kappa$ is even and is orthogonal to this whole sector.

If \eqref{v19:mixed} has an $L^2$ solution, it has an odd one. Equivalently,
the actual compensation can be chosen
\begin{equation}
 \boxed{j=\nabla^\circ\varphi,\quad
  \varphi\in H^1\cap L_0^2,\quad \mathcal I\varphi=\varphi,
 \qquad
 \diver^\circ\{A^*(\nabla^\circ\varphi+2f^\circ)\}=0.}
 \label{v19:potential}
\end{equation}
This is an equivalence for the actual limiting source, not an assertion of
existence of $\varphi$.
\end{theorem}
\begin{proof}
The entries of $b$ use only constants and cosines; those of $r$ use sines.
The table for $q$ is a real cosine polynomial. Curl reverses inversion
parity and $(\Omega^\circ)^{-2}$ preserves it, so $f^\circ$ is odd.
The Hodge maps $P,Q$ commute with componentwise inversion. The positive
transverse operator $H^\circ=PM^*MP+Q$ also commutes with it. Consequently
$w^\circ=(H^\circ)^{-1}f^\circ$ and $y=M^*Mw^\circ$ are odd.
Both $A$ and ordinary differentiation reverse parity, so $D$ and $Z$
commute with $\mathcal I$; the same holds for $K$ on its invariant range.
These facts prove domain and form invariance. Applying the source identity
$g^\circ=-2D^*y$ proves its oddness. The associated self-adjoint operator
and each resolvent commute with $\mathcal I$, by uniqueness of the form
resolvent. The cosine definition of $\kappa$ makes it even.

Given any $\tau$ satisfying \eqref{v19:mixed}, replace it by
$\tau_-=(\tau-\mathcal I\tau)/2$. Since $M$ is even, curl of its odd
part is exactly the even part of $\curl(M^*\tau)$, namely $2q$.
Since $N_r$ is odd, divergence of $N_r^*\tau_-$ is the odd part of the
zero divergence. The resulting $j$ is odd, curl-free, and hence has zero
mean. Fourier Hodge decomposition gives its unique mean-zero potential
$\varphi\in H^1$, with $\|\varphi\|_{H^1}\le C\|j\|_2$.
Oddness of the gradient makes this potential even. Conversely any solution
of \eqref{v19:potential} gives the compensator of V18. Constants are absent
from this chosen stress by the proved parity, while they remain in $Q$,
$K$, the scalar form and every original finite equation.
\end{proof}

An additional exact coefficient identity is useful for assembling this
system. For $v=(1,-1,1)$,
\begin{equation}
      (v\cdot\nabla^\circ)b=r,\qquad
      (v\cdot\nabla^\circ)M=N_r,\qquad
      A=M^{-1}(v\cdot\nabla^\circ M).                 \label{v19:direction}
\end{equation}
It follows by differentiating the three displayed cosine coefficients in
\eqref{v15:limit-coefficients}. It is a logarithmic derivative in an
analytical torus direction, not a physical-time equation or an invertibility
argument for \eqref{v19:potential}. In particular it does not replace the
curl and divergence conditions by pointwise inversion.

\section{Quadratic coefficient zeros and vanishing of the actual forcing}
\label{v19:nodes}

In angle coordinates put $p_+=(\pi/2,\pi/2,\pi/2)$ and $p_-=-p_+$,
with torus representatives understood. In physical unit-torus coordinates
these points are $x_+=(1/4,1/4,1/4)$ and $x_-=-x_+$.

\begin{lemma}[The raw coefficient vanishes to second order at both points]
\label{v19:quadratic-zero}
At $p_+$, in local angle increments $z$,
\begin{equation}
 \begin{split}
 r(p_++z)&=\frac14(z_3^2-z_2^2,\ z_1^2-z_3^2,\ z_2^2-z_1^2)
                                  +O(|z|^4),\\
 b(p_+)&=(253/64,-8,-8).
 \end{split}                                           \label{v19:node-jets}
\end{equation}
At $p_-$ the quadratic part of $r$ has the opposite sign and $b$ has the
same value. There is a fixed neighborhood on which, for $j=0,1$,
\[
 |\nabla^jN_r(\theta)|\le C\,\operatorname{dist}(\theta,\{p_+,p_-\})^{2-j}.
\]
All higher fixed derivatives and the derivatives of $M^{-1}$ are bounded.
The actual limiting scalar forcing satisfies
\begin{equation}
 \boxed{g^\circ(p_+)=g^\circ(p_-)=0,\qquad
 |g^\circ(\theta)|\le C\,\operatorname{dist}(\theta,\{p_+,p_-\})}
 \label{v19:source-zero}
\end{equation}
near these points. The same vanishing holds for both terms in V18's split.
\end{lemma}
\begin{proof}
Substitute $\sin(\pi/2+z_i)=\cos z_i=1-z_i^2/2+O(z_i^4)$ in the
original $r$. The cosine values of $b$ at the point are zero except for
its constants. Inversion proves the formula at $p_-$. The derivative
bounds follow by Taylor's formula, and coefficient invertibility is the
retained explicit boost-margin estimate, not a new assumption.

The exact source is $g^\circ=2\diver^\circ(N_r^*Mw^\circ)$ by V15--V18.
The stress $Mw^\circ$ is smooth with bounded fixed derivatives by the
proved transverse inverse. Both $N_r$ and its first derivatives vanish
at the two points. This gives \eqref{v19:source-zero}. The expressions
$g_{\rm r}=2\diver^\circ(A^*f^\circ)$ and
$g_{\rm b}=2\diver^\circ(A^*Qy)$ give the same conclusion separately:
$A=M^{-1}N_r$ and its first derivatives vanish at those points. No estimate
of the unknown scalar inverse is used.
\end{proof}

\section{Zero is essential even in the sector containing the actual source}
\label{v19:essential}

\begin{theorem}[No finite-dimensional coercivity repair in the sourced sector]
\label{v19:essential-theorem}
Let $\mathscr A_-$ be the restriction of the closed scalar operator of V17
to $\mathcal H_-$. Then
\begin{equation}
          \boxed{0\in\sigma_{\rm ess}(\mathscr A_-).} \label{v19:essential-zero}
\end{equation}
More explicitly, for every $\eta>0$ its spectral projection onto
$[0,\eta)$ has infinite rank. There are real smooth odd functions
$u_\epsilon$ with $\|u_\epsilon\|_2=1$ such that
\begin{equation}
 u_\epsilon\rightharpoonup0,\qquad
 \|Zu_\epsilon\|_2\le C\epsilon,
 \qquad a(u_\epsilon,u_\epsilon)\le C\epsilon^2.
 \label{v19:singular-form-sequence}
\end{equation}
For every bounded finite-rank map $B$, the form
$a(u,u)+\|Bu\|^2$ has no positive $L^2$ lower bound on $\mathcal H_-$.
The same conclusion holds for a bounded compact penalty. In particular,
removing or penalizing the single chamber mode cannot repair the estimate.
\end{theorem}
\begin{proof}
Choose a real even smooth compactly supported bump in a sufficiently small
Euclidean ball, and normalize it in $L^2$. Scale it to width $\epsilon$
around $x_+$ and subtract its reflected copy around $x_-$. For small
$\epsilon$ the supports are disjoint; normalize the pair to one and call
it $u_\epsilon$. It is exactly odd and mean-zero. Its gradient has norm
at most $C/\epsilon$, while $N_r=O(\epsilon^2)$ on its support. The
coefficient inverse and $Q$ are bounded, so
$\|Zu_\epsilon\|\le C\epsilon$. The bounds on $K^{-1}$ give the form
estimate. Support measure tends to zero, and Cauchy--Schwarz against the
restriction of any fixed $L^2$ function proves weak convergence to zero.
It is also orthogonal to the even $\kappa$, without any change to its fields.

Suppose the spectral projection $E_-([0,\eta))$ had finite rank. Its image
of $u_\epsilon$ would tend strongly to zero by weak convergence. The spectral
theorem for the nonnegative form would give
\[
 a(u_\epsilon,u_\epsilon)\ge
 \eta\|(I-E_-([0,\eta)))u_\epsilon\|^2\longrightarrow\eta,
\]
a contradiction. Infinite rank for every such interval is the asserted
essential-spectrum statement. This proof uses form vectors and does not
claim an unproved $L^2$ operator-residual bound for $\mathscr A u_\epsilon$.
Finally a bounded compact map sends a bounded weakly null sequence to a
strongly null sequence, so the same tests exclude either penalty estimate.
\end{proof}

This theorem does not classify the projected-flux kernel. Zero could be an
infinite-multiplicity eigenvalue, an accumulation threshold, or both.
In particular, a gap after quotienting the \emph{entire possibly infinite}
kernel has not been ruled out by a finite-rank assertion. What has been ruled
out is the finite-dimensional Fredholm repair suggested by the one known
chamber mode. The source is in the odd sector used in the theorem, but a
spectral threshold is not by itself a pole for that particular source.

\section{Two-scale estimates with all finite modes and sampling errors retained}
\label{v19:finite-calculus}

Return to the original complementary grids, with any
$\alpha_h=(1,\theta_h,1)$ satisfying $1<\theta_h<1+h$; the inverse
applications may use the already selected rational subfamily. Let
\begin{equation}
 \epsilon_h=h^{3/10},\qquad k_h=h^{-2/5},\qquad
 \lambda_*=2\pi,\qquad
 \mathscr A_h=\frac{\mathscr S_h}{h^2\lambda_*^4},\qquad
 \mathcal O_h=\mathfrak E_hT_h^\dagger V_h.
 \label{v19:scales}
\end{equation}
The real scale $k_h$ describes derivative size, not a deleted Fourier band.
For integer $0\le j\le30$, suppose a smooth function $U$ is supported in
the two cubes of side $C\epsilon_h$ about $x_\pm$, and
\begin{equation}
       \|U\|_{H^j}\le C_j k_h^j\|U\|_2.              \label{v19:family-derivatives}
\end{equation}
The constants in this hypothesis will be provided by disjoint rescaled
copies of one fixed bump, uniformly in the dimension of their span.

\begin{lemma}[Finite source bounds for the localized family]
\label{v19:localized-calculus}
For $u_h=\mathcal S_hU$, uniformly over such families,
\begin{align}
 \|\mathcal I_hu_h-U\|_2&\le Ch^{18}\|U\|_2,
                                                \label{v19:sampling}\\
 \|Z_hu_h\|_h&\le C\{\epsilon_h^2k_h+hk_h^2+h^2k_h^3+h^4\}\|U\|_2,
                                                \label{v19:flux-est}\\
 \frac{\|\mathscr D_{h,\alpha_h}u_h\|_h}{h}
       &\le C(hk_h^2+h^4)\|U\|_2,               \label{v19:diag-est}\\
 \|\mathcal O_h(h^{-1}u_h)\|_h
       &\le C h\{\epsilon_h^2k_h^3+hk_h^4+h^2k_h^5+h^4\}\|U\|_2.
                                                \label{v19:output-est}
\end{align}
Here $Z_h$ is the complete normalized flux of V17. Its constant components
and every original folded mode are retained.
\end{lemma}
\begin{proof}
For completeness, the sampling estimate used here can be obtained directly
from Fourier series. For a smooth $U$, the coefficient of its interpolated
samples at a retained frequency $k$ is $\sum_{m\in\mathbb Z^3}\widehat
U(k+Nm)$. Weighted Cauchy--Schwarz on the nonzero alias indices, together
with the high-frequency Fourier tail, gives
\begin{equation}
 \|\mathcal I_h\mathcal S_hU-U\|_{H^j}
          \le C_{J,j}h^{J-j}\|U\|_{H^J},\qquad J>j+3/2.
 \label{v19:alias}
\end{equation}
Indeed the sum of $(1+|k+Nm|^2)^{-J}$ for $m\ne0$ is at most
$C_JN^{-2J}$ on the retained cube; the output weight is at most $C_jN^{2j}$.
The nonalias high tail has the same bound. Sum first over $k$, then over
its disjoint alias classes. With $J=30$, \eqref{v19:family-derivatives}
gives $h^{30-j}k_h^{30}=h^{18-j}$. This proves
\eqref{v19:sampling} and controls every needed alias contribution below.

On the support of $U$, \cref{v19:quadratic-zero} and Leibniz differentiation
give, for $j=0,2$,
\begin{equation}
 \|M_r\nabla^\circ U\|_{H^j}
                \le C_j\epsilon_h^2 k_h^{j+1}\|U\|_2.
 \label{v19:raw-local}
\end{equation}
To check the powers, no coefficient derivative costs more than needed:
undifferentiated $r$ is $O(\epsilon_h^2)$, one derivative is
$O(\epsilon_h)$, and higher derivatives are bounded. Since
$\epsilon_h k_h\ge1$, each differentiated product is bounded by the
right side. The constants depend on the fixed derivative order, not the
number of bumps.

The original off-diagonal lapse formula, with its literal shifts, has the
smooth consistency bound
\begin{equation}
 \left\|\frac{V_h\mathcal S_hU}{h\lambda_*^2}
                 -\mathcal S_h(\overline VU)\right\|_{H_h^j}
 \le C_j\{h k_h^{j+2}+h^2k_h^{j+3}+h^4\}\|U\|_2,
                   \qquad j=0,2.                    \label{v19:V-consistency}
\end{equation}
This is a quantified use of the original coefficient formula, not smooth
consistency applied without a derivative budget. On unaliased inputs the
phase-derivative error is at most $Ch^2$ with three derivatives charged;
a one-site shift costs at most $Ch$ with one derivative; and the coefficient
error from $\alpha_h-1$ and $\omega_h-\lambda_*$ is $O(h)$.
The shift in a first-order map therefore charges two input derivatives.
Apply \eqref{v19:family-derivatives}. The high retained input modes and
aliases are bounded using \eqref{v19:alias} and the full operator bound
$\|\delta_h\|\le2/h$. Fixed-support coefficient multiplication and its
shifts spend only fixed powers of $h^{-1}$. With $J=30$ these errors are
smaller than the stated $h^4$ remainder, also at $j=2$.
The sampled smooth term in \eqref{v19:V-consistency} obeys
\eqref{v19:raw-local}; applying \eqref{v19:alias} to that product controls
its own sampling error.

The normalized flux is a fixed multiple of the bounded coefficient inverse
and Hodge projection applied to $V_hu_h/h$, by V14--V17. Its metrics have
uniform bounds. Taking $j=0$ proves \eqref{v19:flux-est}.
For the diagonal source, every nonfolded cosine product-defect entry is
bounded by $Ch^2\langle k\rangle^2$. Its output has a fixed unit offset.
The folded entries can have size $Ch^{-1}$, but their inputs belong to the
high tail just bounded. Thus
$\|\mathscr D_hu_h\|\le Ch^2\|U\|_{H^2}+Ch^5\|U\|$, proving
\eqref{v19:diag-est}. No boundary coefficient is set to zero.
Finally the true pseudoinverse and geometric symbol bounds in V13 give
\[
 \|\mathfrak E_hT_h^\dagger Y\|_h
                      \le Ch\|\iota_o^*Y\|_{H_h^2}.
\]
Use $Y=V_hu_h/h$ and the $j=2$ source estimate to prove
\eqref{v19:output-est}. The nonlocal pseudoinverse remains in its actual order.
\end{proof}

\section{A growing low-energy space with a small original geometric observation}
\label{v19:subspace}

\begin{theorem}[Many finite low-energy directions on the unchanged regulator]
\label{v19:subspace-theorem}
For all sufficiently fine odd cutoffs there is a real subspace
$\mathcal V_h$ of the original mean-zero scalar array space such that
\begin{equation}
 \boxed{\dim\mathcal V_h\ge c h^{-3/10},\qquad
        u_h(-x)=-u_h(x)\quad(u_h\in\mathcal V_h).}
 \label{v19:dimension}
\end{equation}
Every vector in it satisfies
\begin{equation}
 \boxed{
 a_h(u_h)\le Ch^{2/5}\|u_h\|_h^2,
 \qquad
 \|\mathcal O_h(h^{-1}u_h)\|_h
                      \le Ch^{2/5}\|u_h\|_h.}
 \label{v19:subspace-bounds}
\end{equation}
Its interpolants are exactly orthogonal to the even chamber mode.
In particular the complete normalized finite scalar matrix obeys
\begin{equation}
 \boxed{\rank\mathbf1_{[0,Ch^{2/5}]}(\mathscr A_h)
                              \ge c h^{-3/10}.}       \label{v19:counting}
\end{equation}
The constants are independent of the allowed detuning. On the rationally
selected family all these matrices are positive and invertible, as before.
The subspace is not asserted invariant or equal to a spectral subspace.
\end{theorem}
\begin{proof}
Choose a fixed real $C_c^\infty$ bump supported in a cube about zero,
normalize it in $L^2$, and place disjoint translates of its $k_h$-rescaling
inside a cube of side $\epsilon_h$ about $x_+$. Use lattice spacing a
fixed multiple of $k_h^{-1}$. For small $h$ their number is at least
$c(\epsilon_h k_h)^3=c h^{-3/10}$. The reflected cube about $x_-$ is
disjoint. For each bump subtract its reflection and divide by $\sqrt2$.
The resulting smooth functions are odd, mutually orthonormal, and supported
in the two cubes. Every linear combination $U$ satisfies
\eqref{v19:family-derivatives}, because derivatives of the different bumps
have disjoint supports; integer derivative norms add in squares. Passing
to equivalent Fourier Sobolev norms changes only order-dependent constants.

Let $\mathcal V_h$ be the space of their sampled arrays. By
\eqref{v19:sampling}, sampling is injective and its norms differ from the
continuum norm by $O(h^{18})$, uniformly on the whole span. Alternatively
injectivity follows from these same estimates applied to a putative zero
sample. Inversion permutes the original odd-grid sites, so oddness and
zero discrete mean hold exactly. Trigonometric interpolation preserves
that inversion parity, proving the orthogonality to $\kappa$.

The exact normalized form is
\[
 a_h(u)=\frac{\|\mathscr D_hu\|_h^2}{2h^2\lambda_*^4}
            +\tfrac14\ip{Z_hu}{K_{Q,h}^{-1}Z_hu}_h.
\]
Insert the proved bounds of \cref{v19:localized-calculus}. The relevant
powers are
\[
 \epsilon_h^2 k_h=h^{1/5},\quad hk_h^2=h^{1/5},\quad
 h^2k_h^3=h^{4/5}.
\]
The bounded flux metric and sampling norm equivalence give the first
estimate of \eqref{v19:subspace-bounds}. For the observation they are
\[
 h\epsilon_h^2k_h^3=h^{2/5},\quad
 h^2k_h^4=h^{2/5},\quad h^3k_h^5=h.
\]
This proves the second estimate on the entire subspace, not just on each
basis vector separately.

For the counting statement, choose the constant in its spectral interval
larger than the energy-bound constant. If that projection had rank smaller
than $\dim\mathcal V_h$, a nonzero vector in $\mathcal V_h$ would be
orthogonal to its range. The finite spectral theorem would bound its energy
strictly above the value just proved. This contradiction proves
\eqref{v19:counting}. This is the elementary variational dimension principle;
no numerical eigenvalue sampling establishes the asymptotic count.
\end{proof}

These are lapse comparison arrays in a linear source problem, not new
nonsmooth nonlinear metric records. The observation scale $h^{-1}u_h$ is
the natural active-lapse scale from V15. The estimate does not say that the
original inverse selects this space, nor that the actual low spectral
projection has a small observation. A variational counting subspace and the
eigenspace counted by it are not interchangeable. Also the full finite
operator need not preserve inversion parity; only the continuum symmetry
was proved in \cref{v19:parity}.

\section{The actual source has a small projection onto the entire constructed space}
\label{v19:source-projection}

Let $P_{\mathcal V_h}$ be the positive physical orthogonal projection onto
$\mathcal V_h$. Retain the actual normalized forcing
$f_h=g_h/(h\lambda_*^4)$, not a chosen residual.

\begin{theorem}[Dimension-independent actual-source suppression]
\label{v19:source-bound}
On the same family,
\begin{equation}
          \boxed{\|P_{\mathcal V_h}f_h\|_h\le Ch^{3/4}.}
 \label{v19:source-projection-bound}
\end{equation}
This is a bound on the whole growing subspace, not a sum of individually
small overlaps with an uncontrolled dimension factor.
\end{theorem}
\begin{proof}
V15 gives $f_h-\mathcal S_hg^\circ=O(h)$ in every fixed Sobolev norm.
The uniform discrete Sobolev embedding, or its Fourier Cauchy--Schwarz proof,
therefore gives the same $O(h)$ supremum bound. On the two support cubes,
\cref{v19:quadratic-zero} bounds $|g^\circ|$ by $C\epsilon_h$.
Their discrete normalized volume is at most $C(\epsilon_h+h)^3\le
C\epsilon_h^3$, since $h/\epsilon_h\to0$. Every vector in
$\mathcal V_h$ is supported on those sites. Consequently, for any unit
vector $u_h$ in this space,
\[
 |\ip{u_h}{f_h}_h|
 \le\|\mathbf1_{\rm cubes}f_h\|_h
 \le C(\epsilon_h+h)\epsilon_h^{3/2}
 \le C\epsilon_h^{5/2}=Ch^{3/4}.
\]
Take the supremum over unit vectors in the whole space. This equals the
projection norm and proves \eqref{v19:source-projection-bound}, with no
basis-dependent factor. This argument uses the actual source's nodal
vanishing; boundedness of an arbitrary forcing would give a weaker power.
\end{proof}

This theorem does not replace $P_{\mathcal V_h}$ by a spectral projection.
In particular it is not yet an inverse estimate for
$\mathscr A_h^{-1}f_h$. The original positive inverse can mix this space
with its complement. The spectral count and source bound have distinct
logical roles and are not composed through an unproved identification.

\section{A controlled observed part of the actual moving resolvent}
\label{v19:resolvent}

For $0<\rho\le1$, put
$v_{h,\rho}=(\mathscr A_h+\rho I)^{-1}f_h$. As before, $\rho$ is an
analytical resolvent parameter and is not added to the gravitational action.

\begin{theorem}[Projected source-resolvent and geometric bounds]
\label{v19:resolvent-theorem}
Uniformly at all sufficiently fine odd cutoffs,
\begin{align}
 \|P_{\mathcal V_h}v_{h,\rho}\|_h
       &\le C\left(h^{3/4}\rho^{-1}
                              +h^{1/5}\rho^{-3/2}\right),
                                                     \label{v19:rho-state}\\
 \boxed{\|\mathcal O_h(h^{-1}P_{\mathcal V_h}v_{h,\rho})\|_h
       \le C\left(h^{23/20}\rho^{-1}
                              +h^{3/5}\rho^{-3/2}\right).}
                                                     \label{v19:rho-output}
\end{align}
In particular the observed part in \eqref{v19:rho-output} tends to zero
for any $\rho_h\downarrow0$ satisfying $\rho_h/h^{2/5}\to\infty$.
The claim is not convergence of the entire moving resolvent or the inverse
at zero.
\end{theorem}
\begin{proof}
The actual normalized forcing has bounded $L_h^2$ norm. Nonnegativity and
the spectral theorem give
\[
 \|v_{h,\rho}\|_h\le C/\rho,\qquad
 a_h(v_{h,\rho})\le C/(4\rho).
\]
Test the full resolvent equation against an arbitrary unit
$u_h\in\mathcal V_h$. Form Cauchy--Schwarz and
\cref{v19:subspace-theorem,v19:source-bound} give
\[
 \rho|\ip{u_h}{v_{h,\rho}}_h|
 \le Ch^{3/4}+\sqrt{a_h(u_h)a_h(v_{h,\rho})}
 \le Ch^{3/4}+Ch^{1/5}\rho^{-1/2}.
\]
Taking the supremum proves \eqref{v19:rho-state}. Apply the uniform
observation bound on the same whole subspace to its orthogonal projection;
this yields \eqref{v19:rho-output}. The more restrictive term there is
$h^{3/5}\rho_h^{-3/2}=(h^{2/5}/\rho_h)^{3/2}$. The other tends to zero
under the same condition. No commutation of $P_{\mathcal V_h}$ with the
scalar matrix has been used.
\end{proof}

The analytical projection decomposes an existing resolvent for the estimate.
It is not an imposed projection in the field equation or a declaration
that its complementary observation vanishes. The state estimate itself would
require the stronger condition $\rho_h/h^{2/15}\to\infty$ to vanish;
\eqref{v19:rho-output} controls the stated geometric output on a wider
moving-parameter range. The zero-parameter problem remains different.

\section{Direction, nonduplication, and verification}
\label{v19:status}

The limiting sourced sector cannot be made coercive by removing a fixed
finite set of scalar modes. The known chamber mode was even; the new
low-energy concentration occurs in the odd sector containing the actual
forcing. The quantitative finite theorem also supplies a growing variational
space and a corresponding lower count of small normalized eigenvalues.
Neither result requires signed active amplification, which V12 already
removed on this family.

At the same time, the original forcing vanishes at the quadratic coefficient
zeros. It has a small projection onto the entire constructed finite space,
and that space has a small original geometric observation at the natural
lapse scale. The observed component of an actual moving resolvent is
controlled accordingly. These are positive source-and-observation results,
not a proof of divergence caused by the essential threshold.

The compensation question is equivalently the trigonometric
curl--divergence system \eqref{v19:mixed}, or the even-potential equation
\eqref{v19:potential}. No $L^2$ compensator has been constructed here. The
initial small Fourier ansatz tests did not supply one; their numerical
residuals are not a proof of nonexistence at larger support or in the flux
domain. A further source-specific compensation, zero-spectral estimate, or
control of the complementary observed response is needed. In particular the
new finite spaces cannot be substituted for spectral subspaces in the
unregularized inverse.

The full projected kernel, the zero-parameter scalar response, higher signed
weak equations, finite-amplitude remainders, differentiated spacetime
observations and common physical lifespan remain unresolved. No original
nonlinear evolution or physical clock has been changed. This body neither
proves nor disproves a raw-action Einstein continuum limit.

\paragraph{Verification scope.}
The exact script checks the directional coefficient identity, inversion
parities, all nodal coefficient values and their quadratic Taylor terms,
the complete inherited stress curl and divergence identities, and the
algebraic powers in the sampling, packing, form, source and resolvent bounds.
The spectral arguments and the dimension-independent smooth localization
estimates are the written proofs, not conclusions from a finite eigenvalue
sample. Separate floating diagnostics apply the retained full original
source and true transverse inverse to shrinking odd comparison bumps at
specified odd grids. They check means, parity, the normal-equation residual
and the separate energy and geometric outputs; they do not certify the
asymptotic counting constant, solve the scalar inverse, or integrate an
exact-action spacetime. All eighteen earlier marked bodies are preserved
byte for byte, with the new source and artifact hashes recorded.

\begin{thebibliography}{9}
\raggedright
\setcounter{enumiv}{16}
\bibitem{v19:Teschl}
G. Teschl, \emph{Mathematical Methods in Quantum Mechanics},
Graduate Studies in Mathematics \textbf{99}, American Mathematical Society,
2009, Sections 4.3--4.4 and 6.4, in particular Theorem 4.12 and the spectral
projection characterization of essential spectrum.
\url{https://www.mat.univie.ac.at/~gerald/ftp/book-schroe/schroe.pdf}.
This is background for the standard variational spectral argument; the
quadratic nodes, actual source parities, finite localization and observed
bounds are proved here. No external spectral floor is imported.
\end{thebibliography}
% END IMMUTABLE EXACT-ACTION BRIDGE V19 BODY


% Future iterations append here.
% BEGIN IMMUTABLE EXACT-ACTION BRIDGE V20 BODY
\clearpage
\renewcommand{\thesection}{\arabic{section}}
\renewcommand{\theHsection}{bridgeV20.\arabic{section}}
\setcounter{section}{161}
\section{Initial multiplier values before another singular inverse}
\label{v20:context}

The stress equation of \cref{v19:mixed-system} is still an unsolved range
problem on the unit-lapse, zero-shift initialization used to derive it.
Neither its essential zero threshold nor an unsuccessful finite Fourier
trial decides the actual zero-parameter response. This body does not claim
to solve that equation. It makes a different upstream change: allow the
\emph{initial values} of lapse and shift to differ from $(1,0)$ by order
field amplitude, while keeping their original equations and the mean-lapse
normalization. These are different initial records of the same action.
They are not a projected vector field, a prescribed multiplier velocity,
or a coordinate transformation of an already selected history.

The effect is governed by the first \emph{Poisson constraint bracket},
not the multiplier Hessian repeatedly inverted in the preceding bodies.
We derive its full contribution to the prepared quadratic drift. On the
actual complementary three-site first jet, a characteristic-zero rational
calculation then finds a normalized initial tilt cancelling all 108
quadratic consistency rows. The bracket has rank 90, so this is a genuine
range test with 18 null directions, not the inversion of a presumed
nonsingular matrix. After the four mean normalizations the admissible
quadratic-cancelling tilts form an affine space of dimension 14.

The original analytic initializer supplies exactly constrained fields at
these tilted multiplier values. Their complete initial drift is order
$a^3$ instead of order $a^2$. The leading active and higher weak source
coefficients remain those of the same first jet, so its inherited regular
signed certificate supplies a local original-action branch. This does not
cancel the next weak drift or bound the full multiplier velocity. In
particular no common physical lifespan or cofinal tilted family is inferred.

\paragraph{Inputs and nonduplication.}
The stationary split and flat weak cancellation are
\cref{v3:split,v3:universal-flat,v3:compiler}; the prepared drift is
\cref{v5:prepared-B2}; and the complementary first jet, mean calibration,
initializer, and full three-site signed certificate are
\cref{v12:family,v12:means,v12:exact-preparation,v12:finite-rank}.
The original Poisson matrix and normalized continuation appear in
\cite{OB25}, \src{obv21:Dirac-matrix} and
\src{obv18:conditional-continuation}. Their general constrained-system
interpretation is standard \cite{v20:Dirac}. The new objects here are the
initial-tilt transformation law for this action, the full original bracket
range certificate, its exactly normalized representative, and the
connection to the original higher-source preparation. The monograph's
separate diagonal one-direction lapse-selection comparison is not a
certificate for this shared-link action and is not used as one.

\section{The flat germ and the constraint bracket controlling an initial tilt}
\label{v20:germ}

Fix a finite odd cutoff in the original analytic stationary chart. Write
\[
 H(X,\lambda)=\mathscr H_h(X,\lambda),\quad
 P(X,\lambda)=D_\lambda H(X,\lambda),\quad
 F(X,\lambda)=\mathbb J\nabla_X H(X,\lambda),
 \quad e=(1,0).
\]
All multiplier covectors are for the physical positive pairing. The initial
map \(\mathscr C_h\) differs from \(P\) only by its retained scalar-row
sign. Their zero sets are identical. Define the two different matrices
\begin{equation}
 M(X,\lambda)=D_\lambda P(X,\lambda),\qquad
 \mathbb K(X,\lambda)=D_XP(X,\lambda)\mathbb J D_XP(X,\lambda)^*.
 \label{v20:two-matrices}
\end{equation}
The first is symmetric; the second is skew-adjoint. The actual constraint
drift is \(\mathsf B=D_XP\,F\), and every exact continuation satisfies
\(\mathsf B+M\lambda_t=0\).

Let \(L=D_XP(0,e)\). For a multiplier test \(\nu=(n,\beta)\), the flat
Hamiltonian vector of its original linear row is
\begin{equation}
 G\nu=\mathbb J L^*\nu
 =\left(\delta_i\beta_j+\delta_j\beta_i,
        \delta_i\delta_jn-\delta_{ij}\Delta_\delta n\right).
 \label{v20:G}
\end{equation}
The original commuting phase derivatives give \(L\mathbb J L^*=0\).
This is the inherited linear identity, not a nonlinear gauge algebra.

\begin{lemma}[The multiplier dependence of the flat field germ]
\label{v20:flat}
For \(\lambda\) in a sufficiently small neighborhood of \(e\),
\begin{equation}
 H(0,\lambda)=0,\quad
 D_XH(0,\lambda)=L^*\lambda,\quad
 P(0,\lambda)=0,\quad D_XP(0,\lambda)=L.
 \label{v20:flat-identities}
\end{equation}
Consequently, uniformly on a smaller fixed-cutoff neighborhood,
\begin{equation}
 M(X,\lambda)=O_h(\|X\|^2),\qquad
 D_XM(X,\lambda)=O_h(\|X\|),\qquad
 F(0,\lambda)=G(\lambda-e).
 \label{v20:flat-orders}
\end{equation}
\end{lemma}
\begin{proof}
Use the original stationary split
\(\Phi=\langle\eta,\mathcal T_X(z)\rangle+
\mathcal V_{X,\lambda}(z)\), with \(\eta=(A_0,w)\).
At \(X=0\), the spatial stationary connection is exactly \(z=0\).
The flat auxiliary solution has temporal rotation
\(-\curl_\delta\beta/2\), temporal boost \(\nabla_\delta N\), and
symmetric velocity \(\delta_i\beta_j+\delta_j\beta_i\). Substitution
in the flat saddle equations verifies this solution for every admitted
\(\lambda\): the spatial source is cancelled by the corresponding
flat auxiliary row. The logarithmic terms have at least two spatial
connection factors in the electric variation and at least three in the
magnetic remainder, so contribute nothing at this zero connection.
The original flat saddle inverse and analytic uniqueness identify this
solution with the retained branch.

Its stationary value is zero. The envelope derivative in \(X\), at this
solution, is just the retained linear lapse--shift row
\(\langle\lambda,LX\rangle\). Its gradient is therefore \(L^*\lambda\),
with no nonlinear multiplier term at first field order. Differentiating
these identities in \(\lambda\) proves the last two equalities of
\eqref{v20:flat-identities}. In particular \(M\) and its first field
derivative vanish at \(X=0\) throughout this multiplier neighborhood.
Analytic Taylor expansion proves the first two bounds in
\eqref{v20:flat-orders}. Finally \(Ge=0\), because the constant lapse
has zero flat phase gradient, giving its last equality.
\end{proof}

Let \(P_{\nu,2}(Y)\) denote the actual homogeneous quadratic field
coefficient of \(\langle\nu,P(X,e)\rangle\). Define
\begin{equation}
 \boxed{\mathbb K_1(Y)[\nu,\mu]
 =DP_{\nu,2}(Y)[G\mu]-DP_{\mu,2}(Y)[G\nu].}
 \label{v20:K1}
\end{equation}
Expansion of \eqref{v20:two-matrices} gives
\(\mathbb K(aY+O(a^2),e)=a\mathbb K_1(Y)+O_h(a^2)\).
All entries of this coefficient are linear in the first field jet. It is
not the quadratic symmetric matrix \(J_1^*K_0^{-1}J_1\).

\begin{lemma}[An exact derivative identity for the original drift]
\label{v20:drift-derivative}
For every fixed multiplier direction \(\mu\),
\begin{equation}
 \boxed{D_\lambda\mathsf B[\mu]
   =\mathbb K\mu+D_X(M\mu)[F].}
 \label{v20:derivative}
\end{equation}
\end{lemma}
\begin{proof}
Differentiate \(\mathsf B=D_XP\,F\) and commute its finite analytic
partial derivatives. Since \(D_\lambda F[\mu]=\mathbb J D_XP^*\mu\),
the differentiated second factor gives \(\mathbb K\mu\). The first
factor is
\[
 D_X(D_\lambda P[\mu])[F]=D_X(M\mu)[F].
\]
This is an identity of the full nonlinear stationary Hamiltonian; no term
is omitted.
\end{proof}

\section{Exactly constrained initialization at tilted multiplier values}
\label{v20:preparation}

Let \(Y\) be a complementary first jet with its calibrated four quadratic
mean equations, as in \cref{v12:means}. Fix a finite multiplier field
\(\ell=(\ell^N,\ell^\beta)\), with \(\overline{\ell^N}=0\), and set
\begin{equation}
                         \lambda_\ell(a)=e+a\ell.
 \label{v20:initial-values}
\end{equation}
The mean lapse is exactly one. For sufficiently small amplitude its lapse
is positive. Setting the three shift means to zero is optional for the
initializer; the certified representative below has that normalization too.

\begin{theorem}[The original initializer permits an order-amplitude multiplier tilt]
\label{v20:tilted-initializer}
At each fixed admitted cutoff there are analytic original initial records
\(X_\ell(a)\), for sufficiently small \(a\), satisfying
\begin{equation}
 \boxed{P(X_\ell(a),e+a\ell)=0,\qquad
 X_\ell(a)=aY+O_h(a^2).}
 \label{v20:exact-initial}
\end{equation}
The construction is analytic also in \(\ell\) on a bounded finite parameter
neighborhood. Its nonconstant linear equation and its four quadratic mean
calibration equations are unchanged. Neither a uniform bound for a
cutoff-dependent \(\ell_h\) nor a common physical interval is asserted.
\end{theorem}
\begin{proof}
Retain the original nonconstant range inverse and the four balancing field
directions of V12. By \cref{v20:flat}, the full constraint map at any
\(\lambda\) near \(e\) has the same flat derivative \(L\). Its
nonlinear range equation therefore has the same contraction construction,
with analytic multiplier parameters, on a smaller fixed-cutoff ball.

More precisely, for \(X=O_h(a)\), integrating the first bound in
\eqref{v20:flat-orders} along the segment from \(e\) to \(e+a\ell\)
gives
\begin{equation}
               P(X,e+a\ell)-P(X,e)=O_h(a^3).
 \label{v20:constraint-cost}
\end{equation}
Thus the linear and quadratic initial equations are unaltered. The range
correction starts at order \(a^2\). Divide the four remaining mean rows
by \(a^2\), using the analytic Taylor extension at \(a=0\). Their
value is the calibrated quadratic mean map of \cref{v12:means} plus
\(O_h(a)\). Its four balancing derivatives are nonsingular. The analytic
implicit-function theorem therefore solves the exact mean equations,
changing the balancing coefficients by \(O_h(a)\) and leaving \(Y\)
unchanged. Every original constraint row is now zero. The same inverses
prove local analytic dependence on \(\ell\). Positivity of the spatial
metric and lapse follows by reducing the amplitude neighborhood. This
changes initial multiplier \emph{values}; it supplies no multiplier rate.
\end{proof}

\begin{theorem}[The complete prepared quadratic drift transforms by the first Poisson bracket]
\label{v20:tilt-law}
For every exactly prepared curve in \eqref{v20:exact-initial},
\begin{equation}
 \boxed{\mathsf B(X_\ell(a),e+a\ell)
  =a^2\{\widehat B_2(Y)+\mathbb K_1(Y)\ell\}+O_h(a^3).}
 \label{v20:law}
\end{equation}
Here \(\widehat B_2(Y)\) is the original \emph{prepared} coefficient
in \eqref{v5:prepared-drift-formula}, not the coefficient obtained by
forgetting its lapse correction. In particular
\begin{equation}
                  \mathbb K_1(Y)\ell=-\widehat B_2(Y)
 \label{v20:range-equation}
\end{equation}
is precisely the initial-value condition removing every quadratic drift row.
\end{theorem}
\begin{proof}
The exact preparation and \eqref{v20:constraint-cost} imply
\(P(X_\ell(a),e)=O_h(a^3)\). Hence its first two field coefficients
satisfy exactly the same linear and second constraint equations used in
the inherited prepared-drift calculation. That calculation gives
\(\mathsf B(X_\ell(a),e)=a^2\widehat B_2(Y)+O_h(a^3)\); the
linear constraint violation at order \(a^3\) changes no displayed
coefficient.

In \eqref{v20:derivative}, on
\((X_\ell(a),e+ta\ell)\), \(0\le t\le1\), one has
\(F=O_h(a)\) and \(D_XM=O_h(a)\) by \cref{v20:flat}.
Their product is \(O_h(a^2)\). Also
\(\mathbb K=a\mathbb K_1(Y)+O_h(a^2)\) there. To check its
multiplier dependence, its constant field coefficient is
\(L\mathbb J L^*=0\) for every \(\lambda\); differentiating its
first field coefficient costs \(O_h(a)\), and the multiplier change
is itself order \(a\). Integrate \eqref{v20:derivative} against
\(a\ell\) over that segment. The extra drift is
\(a^2\mathbb K_1(Y)\ell+O_h(a^3)\), proving \eqref{v20:law}.
The last statement is equality of its complete quadratic covector, not a
selected subset of equations.
\end{proof}

This law does not contradict V9--V19's unit-initial-multiplier results.
Those hold at \(\lambda=e\). Here the intrinsic first canonical jet
\(Y\) is preserved, but the full initial spacetime record has different
order-\(a\) lapse and shift components. They are not silently identified
with the earlier record or reset during its evolution.

\section{A reproducible rational construction of the original bracket}
\label{v20:coefficient}

Fix the actual first jet of \cref{v12:three-seed}: \(N=3\), \(h=1/3\),
\(\alpha=(1,9/8,1)\), and
\begin{equation}
 \begin{split}
 U&=\sum_i\alpha_iT_i\cos(2\pi x_i),\\
 p&=3\sqrt3\,\diag(16,16593/4096,16593/4096)
       +\frac{3\sqrt3}{2}\sum_i\alpha_iP_i\cos(2\pi x_i).
 \end{split}
 \label{v20:seed}
\end{equation}
The full multiplier order is 27 lapse point indicators followed by the
three component-major groups of 27 shift indicators. The original phase
derivative is the rational matrix \(3(S_i-S_i^{-1})\).

The certificate uses exact homogeneous coefficient polarization. We give
the scalar extraction to fix signs and its incidence with the action.
For \(Y=(U,p)\), put
\[
 z_i=(R_i,B_i),\quad
 R_{ia}=-\tfrac12\epsilon_{abc}\delta_bU_{ic},\quad
 B=p-\tfrac12(\tr p)I,
\]
\[
 \Pi_1=(\tr U)I-U,\qquad
 \Pi_2=\tfrac34U^2-\tfrac12(\tr U)U+
       \left\{\tfrac14(\tr U)^2-\tfrac12\tr(U^2)\right\}I.
\]
For \(\nu=(n,\beta)\), use the complete flat auxiliary response
\(a_R=-\curl_\delta\beta/2\), \(a_B=\nabla_\delta n\),
\(w=w_\beta\). For \(i<j\), let \(\epsilon_{ijk}=\epsilon\),
with \(k\) the remaining coordinate, and set
\[
 \Sigma_{0,ij}=(2\epsilon n e_k,\,2\beta_j e_i-2\beta_i e_j),\qquad
 \Sigma_{1,ij}=(\epsilon n U_k,\,\beta_j(\Pi_1)_i-
                                             \beta_i(\Pi_1)_j).
\]
Rows of a symmetric tensor carry their original Frobenius occurrences.
The retained Lorentz bracket is
\[
 [(r,b),(s,c)]=(-r\times s+b\times c,\,-r\times c-b\times s).
\]
In the spatial-sum convention, the full quadratic scalar coefficient is
\begin{equation}
\begin{split}
 -P_{\nu,2}(Y)={}&-\langle D\Pi_2(U)[w],B\rangle
       +\sum_i\langle\delta_i(\Pi_2)_i,a_B\rangle
       +\sum_i\langle(\Pi_1)_i,[a,z_i]_B\rangle\\
 &-\frac h2\sum_i\langle2e_i,[z_i,[z_i,S_i a]]_B\rangle\\
 &+\sum_{i<j}\left\{
  \langle\Sigma_{1,ij},\delta_i z_j-\delta_jz_i\rangle
              +\langle\Sigma_{0,ij},[z_i,z_j]\rangle\right\}.
\end{split}
\label{v20:P2}
\end{equation}
This is the original scalar functional's degree-two coefficient with its
flat auxiliary test retained. The cubic electric term contributes here;
the cubic magnetic remainder has too high a field degree to enter this
particular coefficient.

For the cubic constant-multiplier coefficient \(\Phi_3=-H_3\), write
\(V(p)=2p-(\tr p)I\) and use
\begin{equation}
\begin{split}
 \Phi_3(Y)={}&-\langle D\Pi_2(U)[V(p)],B\rangle\\
 &+\sum_{i<j}\epsilon\left\{
 -\tfrac14\langle(U^2)_k,\delta_iR_j-\delta_jR_i\rangle
 +\langle U_k,[z_i,z_j]_R\rangle
 +2h\langle e_k,\mathcal B_{3,ij}(z)_R\rangle\right\}.
\end{split}
\label{v20:Phi3}
\end{equation}
Here \(\mathcal B_{3,ij}\) is the cubic BCH coefficient of the original
ordered factors \((z_i,S_i z_j,-S_jz_i,-z_j)\), not of an unshifted
comparison plaquette. It is reproduced by starting \(l_1=l_2=l_3=0\)
and, for each next factor \(v\), making the simultaneous updates
\[
 l_3\leftarrow l_3+\tfrac12[l_2,v]
       +\tfrac1{12}([l_1,[l_1,v]]+[v,[v,l_1]]),\quad
 l_2\leftarrow l_2+\tfrac12[l_1,v],\quad l_1\leftarrow l_1+v.
\]
The derivative in these formulas is explicitly
\[
 D\Pi_2(U)[w]=\tfrac34(Uw+wU)-\tfrac12\{(\tr w)U+(\tr U)w\}
                +\{\tfrac12(\tr U)(\tr w)-\tr(Uw)\}I.
\]

The exact polarization identities
\begin{align}
 DP_{\nu,2}(Y)[Z]
   &=\tfrac12\{P_{\nu,2}(Y+Z)-P_{\nu,2}(Y-Z)\},\notag\\
 D\Phi_3(Y)[Z]
   &=\tfrac12\{\Phi_3(Y+Z)-\Phi_3(Y-Z)\}-\Phi_3(Z)
 \label{v20:polarization}
\end{align}
therefore generate every entry of \eqref{v20:K1} and of
\begin{equation}
 \widehat B_2(Y)[\nu]
 =D\Phi_3(Y)[G\nu]+DP_{\nu,2}(Y)[F_1Y]
                           +P_{(0,\delta n),2}(Y).
 \label{v20:B2}
\end{equation}
The last term is the prepared lapse correction. These are polynomial
identities with no small differencing parameter or numerical derivative.
The package evaluates them using arbitrary-precision rational arithmetic.

\section{All quadratic consistency rows admit an exact normalized tilt}
\label{v20:certificate}

The arithmetic may be rationalized without approximating \(\sqrt3\).
Write \(p=\sqrt3\,p_R\) in \eqref{v20:seed}, and evaluate the bracket
at \((U,p_R)\). Denote its rational blocks by \(K_{NN},K_{N\beta},
K_{\beta N},K_{\beta\beta}\). Its same-type blocks are linear in
\(p\), whereas its mixed blocks depend only on \(U\). The prepared
lapse target has type \(Up\), and its shift target has types \(U^2,p^2\).
These statements follow directly from \eqref{v20:P2}--\eqref{v20:B2},
or from the retained reversal grading. Thus, with
\(Q_3=\diag(I_{27},\sqrt3 I_{81})\), the physical equation
\eqref{v20:range-equation} for \(\ell=Q_3x\) is exactly
\begin{equation}
 \widehat Kx=-\widehat b,\qquad
 \widehat K=\frac1{\sqrt3}Q_3\mathbb K_1(Y)Q_3
       =\begin{pmatrix}K_{NN}&K_{N\beta}\\
                       K_{\beta N}&3K_{\beta\beta}\end{pmatrix},
 \quad \widehat b=\frac1{\sqrt3}Q_3\widehat B_2(Y).
 \label{v20:rational-system}
\end{equation}
All these arrays are rational. In particular their scaled target is recovered
from two rational evaluations by
\[
 \widehat b_N=\widehat B_{2,N}(U,p_R),\qquad
 \widehat b_\beta=\widehat B_{2,\beta}(U,p_R)
                              +2\widehat B_{2,\beta}(0,p_R).
\]
Common factors converting spatial sum to spatial mean cancel from the
linear equation. The physical squared norm of the tilt remains
\(x^TW_3x/27\), where \(W_3=\diag(I_{27},3I_{81})\).

\begin{theorem}[An exact full-row range certificate for the actual action]
\label{v20:range-certificate}
For \eqref{v20:seed}, \(\widehat K\) has rank 90 and
\(\widehat b\in\Ran\widehat K\). There is a unique solution
\(x_*\) minimizing \(x^TW_3x\), and the physical tilt \(\ell_*=Q_3x_*\)
satisfies
\begin{equation}
 \boxed{\mathbb K_1(Y)\ell_*=-\widehat B_2(Y),\qquad
     \overline{\ell_*^N}=0,\qquad\overline{\ell_*^\beta}=0.}
 \label{v20:tilt-cancel}
\end{equation}
Every one of its 108 equations holds in characteristic zero. Quantitatively,
\begin{equation}
        356.60<\|\ell_*\|_h<356.61,\qquad
          \max_{x,\mu}|\ell_*^\mu(x)|<519.
 \label{v20:tilt-size}
\end{equation}
The normalized solutions form a 14-dimensional affine space. This is a
fixed-three-site statement, not a bound under spatial refinement.
\end{theorem}
\begin{proof}
The files in the proof package specify \(\widehat K,\widehat b,x_*\)
as exact rational strings and give the independent 18-column null matrix
\(N_*\). To specify the rank witness compactly, the omitted indices from
the selected principal 90-column set \(\mathcal J\) are, in zero-based
component-major order,
\begin{equation}
 \mathcal F=(8,17,20,23,24,25,26,35,44,53,74,77,80,89,98,105,106,107),
 \qquad \mathcal J=\{0,\ldots,107\}\setminus\mathcal F.
 \label{v20:minor-indices}
\end{equation}
The exact checks are
\begin{equation}
 \boxed{\widehat K^T=-\widehat K,\quad
 \widehat K N_*=0,\quad (N_*)_{\mathcal F,:}=I_{18},\quad
 \det\widehat K_{\mathcal J,\mathcal J}=25\pmod{241}.}
 \label{v20:rank-certificate}
\end{equation}
All rational denominators in this minor are units modulo 241. Its nonzero
residue proves rank at least 90 over \(\mathbb Q\); the 18 explicit
independent null columns prove the reverse bound. This does not infer an
exact rank from a floating threshold.

Exact elimination of the full augmented rational matrix supplies a
particular solution \(x_0\). The verifier multiplies it against
\emph{all} 108 rows. The minimum physical-norm representative is then
specified without an approximate inverse by
\begin{equation}
 x_*=x_0-N_*(N_*^TW_3N_*)^{-1}N_*^TW_3x_0.
 \label{v20:minimizer}
\end{equation}
The middle matrix is positive definite. Exact rational multiplication
verifies \(\widehat Kx_*=-\widehat b\) and
\(N_*^TW_3x_*=0\). Every other solution differs by a null vector;
Pythagoras in this physical weight proves uniqueness and minimality.
The four component-constant multiplier arrays are verified null vectors.
Orthogonality to them proves all four zero means in \eqref{v20:tilt-cancel}.
Their mean map has rank four on the 18-dimensional kernel, so imposing
these four means leaves dimension 14.

The complete rational sum \(x_*^TW_3x_*/27\) lies strictly between
\((35660/100)^2\) and \((35661/100)^2\). Each rational square
\(x_{*,j}^2\), multiplied by 1 on lapse slots or 3 on shift slots, is
strictly below \(519^2\). These rational comparisons prove
\eqref{v20:tilt-size}; no floating square root is used as a bound.
The exact entries of \(\ell_*\) are defined by the rational arrays and
\eqref{v20:minimizer}, not by rounded displayed coordinates.

Finally, a fresh independent differentiation of the inherited complete
scalar-action evaluator at \(\sqrt3\mapsto56\in\mathbb F_{241}\)
reproduces every one of the 11,664 bracket entries and all 108 prepared
target entries after the congruence \eqref{v20:rational-system}. It
also checks the cancellation on every original row. This is a regression
against the original functional; the exact rational range proof is given
by the identities above and their coefficient regeneration.
\end{proof}

The norm in \eqref{v20:tilt-size} is not small. It is finite at this
cutoff, so \(a\ell_*\to0\) and the initializer admits it after reducing
\(|a|\). For example \(|a|<1/1038\) keeps the lapse between \(1/2\)
and \(3/2\); other stationary-chart restrictions may require a smaller
amplitude. No universal amplitude radius or cofinal tilt bound is claimed.

\section{Why the existing leading signed certificate survives the initial tilt}
\label{v20:source-stability}

Initial multiplier values enter the stationary connection, so reusing the
higher weak matrix needs a proof. The needed cancellation is supplied by
its original saddle structure, not by assuming that every second source
is multiplier-independent.

\begin{lemma}[The leading graded source is unchanged by an order-amplitude tilt]
\label{v20:graded-source}
Let \(X(a)=aY+a^2Y_2+O_h(a^3)\),
\(\lambda(a)=e+a\ell+O_h(a^2)\), on the original stationary branch.
Then the leading active coefficient of the reduced source is
\(aJ_1(Y)|_A\), and its restriction to every curl-free shift is
\begin{equation}
           J(X(a),\lambda(a))|_W
              =a^2\mathcal Q_h(Y;\cdot)|_W+O_h(a^3),
 \label{v20:same-weak}
\end{equation}
with the same \(\mathcal Q_h\) as \cref{v3:compiler}. The leading
stationary tangent Hessian is still \(K_0\). These equalities hold for
fixed finite \(Y_2,\ell\); no uniform nonlinear remainder is asserted.
\end{lemma}
\begin{proof}
Retain \(\Phi=\langle\eta,\mathcal T_X(z)\rangle+
\mathcal V_{X,\lambda}(z)\). The first summand is independent of
\(\lambda\), and the second is affine in it. The flat spatial connection
\(z(0,\lambda)=0\) from \cref{v20:flat} implies that its first
coefficient on the stated curve is precisely \(z_1(Y)\). Multiplier
tilt can enter its second coefficient, but not its first. Analyticity also
gives the same limiting tangent Hessian \(K_0\).

The second normal connection coefficient is fixed by expanding
\(\mathcal T_{X(a)}(z(a))=0\), which contains no multiplier. At fixed
\(Y,Y_2\), two initial tilts can therefore change this second coefficient
only by \(Z_0S_2\), a flat tangent vector. The reduced source itself is
\(Z(X,z)^*D_z\partial_\lambda\mathcal V_X(z)\); it has no additional
explicit multiplier dependence. For a weak test \(\nu=(0,\beta)\),
\(\curl_\delta\beta=0\), its flat normal correction pairs with
\(\eta_0=(0,w_\beta)\). The change caused by \(Z_0S_2\) is exactly
\[
 Z_0^*D_z^2\{\mathcal V_{0,\nu}(z)+
                    \langle\eta_0,\mathcal T_0(z)\rangle\}_{z=0}Z_0S_2.
\]
This map is zero. The flat constraint paired with \(w_\beta\) is linear
in \(z\); both tangent variations in the magnetic term are rotations,
whose commutator is annihilated by the boost coefficient of a pure shift.
The electric term has zero flat temporal response on this weak test, and
higher logarithms have too high a spatial degree. This is exactly the
flat tangent-Hessian cancellation proved in \cref{v3:compiler}.
Thus the tilt contributes nothing to the weak second coefficient.

The possible \(Y_2\) contribution is
\(D_XJ(0,e)[Y_2]\nu=0\), by \cref{v3:universal-flat}; its general
symmetric-field statement is used, not just the scalar-momentum case.
The remaining second coefficient is consequently the full construction
\(\mathcal Q_h(Y;\nu)\), including its normal and first auxiliary
corrections. The leading active statement follows directly by the
first-coefficient expansion. No higher weak coefficient is declared
unchanged by this argument.
\end{proof}

\begin{theorem}[An actual tilted local branch with cubic initial drift]
\label{v20:local-branch}
At the first jet \eqref{v20:seed} and the exact tilt
\eqref{v20:tilt-cancel}, the original action admits analytic exactly
constrained initial records with
\begin{equation}
 \boxed{\begin{gathered}
 \lambda(a,0)=e+a\ell_*,\qquad X(a,0)=aY+O_h(a^2),\\
 P(X(a,0),\lambda(a,0))=0,\\
 \mathsf B(X(a,0),\lambda(a,0))=O_h(a^3).
 \end{gathered}}
 \label{v20:main}
\end{equation}
Each sufficiently small nonzero amplitude has a normalized local exact
original-action continuation. Its initial literal temporal and spatial
link curvatures are \(O_h(a)\). No bounded initial multiplier velocity
or corresponding metric-curvature limit is asserted.
\end{theorem}
\begin{proof}
The exact initializer is \cref{v20:tilted-initializer};
\cref{v20:tilt-law,v20:range-certificate} give the cubic drift.
By \cref{v20:graded-source}, the leading graded source is the same
107-column matrix already certified at this first jet in
\cref{v12:finite-rank}. Its determinant residues remain
\(131\) for the active Gram, \(208\) for the weak Schur block,
and \(15\) for the full Gram modulo 241. These inherited arrays and
determinants are rechecked by the package; new higher weak columns are
not represented as independently regenerated in this body.

Use the fixed mean-lapse-zero multiplier complement, of dimension 107.
When using the inherited point-column certificate, subtract each lapse
column's mean multiple of \(e\). Since \(J_1e=0\), this changes no
leading active source column and supplies exactly that fixed complement.
In active/weak coordinates the source columns have weights \(a,a^2\).
After division by these weights, the full signed matrix tends to that
invertible leading Gram. It stays invertible for sufficiently small
nonzero \(a\). The exact radial null direction is \(\lambda(a)\),
which is not in the complement because its mean lapse is one. Thus the
complete multiplier Hessian has exactly this one null direction.
The inherited normalized original-action continuation theorem in
\cite{OB25}, \src{obv18:conditional-continuation}, supplies each local
history, preserving every original constraint. It does not provide a
common interval as \(a\to0\) or \(h\to0\).

For the original electric response, write \(D_a=\diag(aI_A,a^2I_W)\).
Its normalized source target is \(D_a^{-1}\mathsf B=O_h(a)\) by
\eqref{v20:main}, and its normalized signed Gram has a bounded inverse
at this fixed first jet. The spatial connection's multiplier contribution
is therefore \(O_h(a)\), by the unchanged stationary response formula
\(D_\lambda z=-ZK^{-1}J\). The remaining field contribution is also
\(O_h(a)\): \(X_t=F(X,e+a\ell_*)=O_h(a)\), and all fixed-cutoff
stationary derivatives are bounded on the small chart. The initial
stationary spatial and temporal connections are \(O_h(a)\). Their
literal ordered curvature formulas give the stated temporal and spatial
bounds. This uses the original signed solution, not its positive minimum
surrogate. The stronger amplitude weights of the multiplier \emph{rate}
do not follow from this physical response estimate.
\end{proof}

\section{The next weak coefficient is simplified, not proved zero}
\label{v20:next}

For clarity, keep a fixed basis of the mean-lapse-zero multiplier complement,
with active and curl-free shift coordinates as in the preceding proof. Let
\(\Gamma(a)\) be its exact graded Hessian and let
\(\mathsf B(a)=a^3B_3+O_h(a^4)\) on one of the tilted prepared curves.
Writing \(R_A,R_W\) for coordinate restriction, define
\[
 b_3^W=R_WB_3,\qquad
 \Gamma(0)=\begin{pmatrix}A&L\\L^*&C\end{pmatrix},\qquad
 D=C-L^*A^{-1}L.
\]
The physical active representatives need not be orthogonal to the weak
shift subspace; their fixed synthesis maps must be retained when converting
these coordinates back to a rate.

\begin{proposition}[The first surviving weak target has no quadratic active feedback]
\label{v20:weak-order}
In the normalized complement coordinates the exact multiplier rate has
\begin{equation}
 \begin{split}
 \lambda_t^A&=-v_{1,A}+O_h(a),\qquad
 \lambda_t^W=-a^{-1}v_{1,W}+O_h(1),\\
 v_1&=\Gamma(0)^{-1}\binom0{b_3^W},\qquad
 \boxed{v_{1,W}=D^{-1}b_3^W,\quad
           v_{1,A}=-A^{-1}LD^{-1}b_3^W.}
 \end{split}
 \label{v20:weak-expansion}
\end{equation}
Thus the initial-tilt cancellation removes the former quadratic active
feedback from this first surviving weak target. Absence of the
inverse-amplitude weak pole is equivalent to \(b_3^W=0\), on this
regular fixed-cutoff branch. That last equality is not established here.
\end{proposition}
\begin{proof}
Restriction of the exact consistency equation to the fixed complement gives
\(D_a\Gamma(a)D_a\lambda_t=-\mathsf B(a)\), since the normalization
puts \(\lambda_t\) in that complement. Its normalized right side is
\[
 D_a^{-1}\mathsf B(a)
    =a\binom0{b_3^W}+O_h(a^2).
\]
Analytic inversion gives
\(\Gamma(a)^{-1}D_a^{-1}\mathsf B=a v_1+O_h(a^2)\).
Undoing the two column weights proves the first line of
\eqref{v20:weak-expansion}. Block elimination of the invertible leading
Gram gives its last line. A physical curl-free observation of active
representatives contributes only a bounded term from \(\lambda_t^A\),
so it cannot cancel a nonzero inverse-amplitude weak coefficient.
Invertibility of \(D\) proves the equivalence. The coefficient \(B_3\)
includes the new initial tilt, the actual second preparation, and the full
next action terms; it is not borrowed from the zero-tilt family.
\end{proof}

There is a useful but explicitly conditional consequence. If further original
initial preparation on this first-jet class made the \emph{entire actual}
initial curl-free shift velocity zero, the complete initial multiplier
velocity would be \(O_h(a)\), not just bounded. Indeed the velocity would
then lie on the physical active slice; its restricted leading matrix is
\(a^2\) times the invertible active Gram, whereas its target is
\(O_h(a^3)\). This gives the estimate directly. At those cuts,
\(\lambda-e=O_h(a)\), \(X=O_h(a)\), \(X_t=O_h(a)\), and
\(X_{tt}=D_XF F+D_\lambda F\lambda_t=O_h(a)\). The original
interpolated ADM metric would then have initial Riemann components
\(O_h(a)\), using the same initial curvature identities as V10 with
the small nonconstant initial lapse and shift retained. This is not a new
weak-rate preparation theorem, nor a fixed-amplitude refinement statement.

\section{Scope, next calculation, and reproducible checks}
\label{v20:status}

The positive acquisition is actual finite data: the unchanged action has
initial multiplier values \(e+a\ell_*\), exact original initial constraints,
and zero complete quadratic drift at a verified regular complementary first
jet. The first Poisson matrix is singular, yet the actual prepared target
passes its full range test. The inherited higher signed source remains
regular after this initial-value change, which is why this is more than a
formal cancellation on records with unknown local continuation.

The earlier stress-compensation equation remains relevant to its own fixed
unit-initial-multiplier branch. No compensator for it is claimed by this
construction. The next concrete calculation on the new branch is instead
its full cubic weak covector \(b_3^W\), including its dependence on the
14 normalized null-tilt directions and on admissible higher field preparation.
The current result neither sets that covector to zero nor proves that its
controls are onto. It also leaves cofinal tilt existence and size, higher
weak signed rank, nonlinear remainder budgets, differentiated observations,
and the common physical lifespan separate. A further claimed bridge must
retain the changed complete initial records and the original action; it
cannot simply reuse the old source \(g_h\) while changing its initial
multiplier values.

\paragraph{Verification and preservation.}
The exact coefficient program regenerates every entry of the rational first
Poisson matrix and both rational prepared-target evaluations used in
\eqref{v20:rational-system}. The short verifier checks all 108 cancellation
rows over \(\mathbb Q\), skewness, the 90-dimensional nonzero principal
minor, all 18 exact null vectors, the weighted minimum-norm conditions, all
four means, and the rational norm and pointwise bounds. A separate run of
the unchanged original scalar-action differentiation verifies all bracket
and target entries modulo a good prime; its integer operations and modular
adjoints are checked below the exact arithmetic threshold. The quadratic
and cubic polarization rules and the flat constraint-vector identities are
also tested in exact arithmetic.

The existing full leading source determinant is rechecked from the inherited
V12 arrays, with its new applicability proved by \cref{v20:graded-source};
those arrays are not described as freshly regenerated weak coefficients.
No numerical nonlinear initial root, spacetime integration, all-cutoff tilt
certificate, or proof-assistant formalization is reported. The previous
19 marked research bodies are preserved byte-for-byte, including their
corrections and scope statements. New labels use the prefix \texttt{v20:}.

\begin{thebibliography}{9}
\raggedright
\setcounter{enumiv}{17}
\bibitem{v20:Dirac}
P. A. M. Dirac, \emph{Generalized Hamiltonian Dynamics},
Canadian Journal of Mathematics \textbf{2} (1950), 129--148.
\href{https://doi.org/10.4153/CJM-1950-012-1}{doi:10.4153/CJM-1950-012-1}.
This is primary-source background for constraint Poisson brackets; the
initial-tilt law, original-action coefficient certificate, and its local
preparation consequences are established in this body.
\end{thebibliography}
% END IMMUTABLE EXACT-ACTION BRIDGE V20 BODY


% Future iterations append here.
% BEGIN IMMUTABLE EXACT-ACTION BRIDGE V21 BODY
\clearpage
\renewcommand{\thesection}{\arabic{section}}
\renewcommand{\theHsection}{bridgeV21.\arabic{section}}
\setcounter{section}{169}
\section{Cubic weak preparation after the initial-value selection}
\label{v21:context}

The initial multiplier selection in \cref{v20:main} cancels every quadratic
constraint-drift row. Its cubic weak row was not evaluated or proved
controllable. Since the weak multiplier Hessian begins two amplitude orders
later than the active one, cubic drift alone does not bound the multiplier
velocity. This body addresses that remaining initial-cut obstruction on the
same complementary three-site first jet. Neither the intrinsic first canonical
record nor the original action is changed.

There are two distinct controls. Three of the previously unused normalized
first-tilt directions change the cubic constant-shift rows. Admissible second
canonical records control the remaining mean-zero weak rows. Their interaction
is triangular because the cubic mean row, although it depends on the first
multiplier tilt, is independent of second and higher canonical preparation.
The first statement is certified by an exact rational source calculation; the
second by a good-prime nonzero minor of the complete constrained control map.
Together they give analytic original initial records whose \emph{entire actual}
initial curl-free shift velocity is zero and whose complete multiplier
velocity is of order amplitude. The original initial metric and literal-link
curvatures consequently tend to zero at that fixed cutoff.

\paragraph{Inherited inputs and nonduplication.}
The stationary chart, weak flat-source cancellation and cubic constant-shift
envelope are \cref{v3:universal-flat,v5:P3,v5:constant-B2}. The complete
quadratic drift and first Poisson bracket are \cref{v20:tilt-law} and
\eqref{v20:K1}. The complementary first jet, mean calibration and full signed
source certificate are \cref{v12:means,v12:finite-rank}. Their applicability
under an order-amplitude initial tilt is \cref{v20:graded-source}. We retain
the original analytic range-and-mean initializer and fixed-cutoff continuation,
not a new prescribed flow. The new work consists of the cubic-mean tilt map,
its actual three-dimensional range certificate, the constrained second-jet
control and their nonlinear preparation consequence. The analytic implicit
function theorem and the constraint-algorithm interpretation are standard;
see the already retained \cite{v20:Dirac}. No cofinal inverse or physical
lifespan estimate is imported from that reference.

All numerical certificates in this body concern $N=3$, $h=1/3$, unless a
statement explicitly says otherwise. In particular the new family is not
silently identified with the old unit-initial-multiplier source of V15--V19.
No zero of a cubic coefficient is inferred from a small floating residual.

\section{The cubic mean row depends only on the first tilt and first field}
\label{v21:mean-law}

Use $H,P,F,M,\mathbb K,L,G$ from \cref{v20:germ}, and put
$\mathcal W=\ker\curl_\delta=\mathcal W_c\oplus\mathcal W_g$, with
$\mathcal W_c$ the three constant shifts and $\mathcal W_g$ the mean-zero
phase gradients. Superscripts on a physical multiplier denote lapse and shift,
not powers. Fix $Y\in\ker L$ and an exactly constrained analytic initial curve
\begin{equation}
 X(a)=aY+a^2Y_2+O_h(a^3),\qquad
 \lambda(a)=e+a\ell+a^2m+O_h(a^3),\qquad P(X(a),\lambda(a))=0.
 \label{v21:initial-jet}
\end{equation}
The mean lapse of every multiplier correction is zero. Additional shift-mean
normalizations will be imposed on the explicit construction below.

For a constant shift $c$, let $\delta_c=c^i\delta_i$. Retain the complete
stationary cubic coefficient $P_{c,3}$ of \eqref{v5:P3}, including its temporal
auxiliary response. Define
\begin{align}
 t_c(Y)&=D\Phi_3(Y)[\delta_cY]+DP_{c,3}(Y)[F_1Y],
                                                    \label{v21:t}\\
 \mathcal R_Y[c,\ell]
 &=DP_{c,3}(Y)[G\ell]-DP_{\ell,2}(Y)[\delta_cY]
                      -\ip{\delta_c\ell}{P_2(Y)}_h . \label{v21:R}
\end{align}
Here $\Phi_3=-H_3$, $P_2(Y)$ is the complete multiplier covector and
$P_{\ell,2}=\ip\ell{P_2}_h$. All three entries are supplied by the original
quadratic and cubic scalar envelopes. This $\mathcal R_Y$ is not the
stationary connection response or the first Poisson bracket.

\begin{theorem}[Complete cubic constant drift under initial multiplier tilt]
\label{v21:constant-law}
On \eqref{v21:initial-jet}, at every fixed admitted odd cutoff,
\begin{equation}
 \boxed{\mathsf B(X(a),\lambda(a))[(0,c)]
   =a^3\{t_c(Y)+\mathcal R_Y[c,\ell]\}+O_h(a^4).}
 \label{v21:constant-law-eq}
\end{equation}
The displayed coefficient is independent of $Y_2$, of all higher field
coefficients, and of $m$ and all higher multiplier coefficients. Its dependence
on $\ell$ is affine, not quadratic. No vanishing of this coefficient is assumed.
\end{theorem}
\begin{proof}
For a constant shift the original $P_{c,1}$ is zero and
$P_{c,2}(X)=\ip p{\delta_c u}_h$. Its Hamiltonian vector is
$\delta_c X$, and $\{P_{c,2},H_2\}=0$ on arbitrary fields. Also the
universal first-source cancellation gives $J_1(X)c=0$ for every symmetric
first field. The leading multiplier Hessian is $J_1^*K_0^{-1}J_1$, so
$D_\lambda P_{c,2}(Y)[\mu]=0$ at $e$ for every $Y,\mu$, together with
its field derivatives. Thus tilting the multiplier introduces no new term
through the cubic \emph{value} of $P_c$, and none through the quadratic field
coefficient of its field gradient.

Write $X_2=Y_2$ for this proof and let $X_{P_{\ell,2}}$ be the canonical
vector of its homogeneous quadratic functional. The first two coefficients
of the actual canonical vector field are
\[
 F(X(a),\lambda(a))
 =a(F_1Y+G\ell)
 +a^2\{F_1X_2+F_2(Y)+X_{P_{\ell,2}}(Y)+Gm\}+O_h(a^3).
\]
This follows by Taylor expansion of the original stationary Hamiltonian;
$H_1(X,\lambda)=\ip\lambda{LX}$ is affine in $\lambda$, so no quadratic
$\ell$ term occurs at this vector-field order. Terms of that type can occur
at the next order and are not removed from the action.

Pair the field gradient of $P_c$ with this expansion. Terms involving
$F_1X_2$ and $DP_{c,2}(X_2)F_1Y$ cancel by polarization of
$\{P_{c,2},H_2\}=0$. The coefficient without tilt is exactly
\eqref{v21:t}. The remaining terms at order three are
\[
 DP_{c,2}(X_2)[G\ell]+DP_{c,2}(Y)[X_{P_{\ell,2}}(Y)]
                      +DP_{c,3}(Y)[G\ell]+DP_{c,2}(Y)[Gm].
\]
For any symmetric $X$ and multiplier $\mu$, commuting $L$ with $\delta_c$
and using skew-adjointness gives the exact identity
\begin{equation}
 DP_{c,2}(X)[G\mu]=\ip{\delta_c\mu}{LX}_h.
 \label{v21:translation-incidence}
\end{equation}
Indeed the Poisson bracket in question is
$-DP_{\mu,1}(X)[\delta_cX]=-\ip\mu{\delta_c LX}_h$.
Since $LY=0$, the $m$ term is zero. Exact initial preparation and V20's
$O(a^3)$ constraint cost imply $LX_2=-P_2(Y)$, so the first term becomes
$-\ip{\delta_c\ell}{P_2(Y)}_h$, with no free second coefficient left.
Finally $DP_{c,2}(Y)[X_{P_{\ell,2}}(Y)]
=-DP_{\ell,2}(Y)[\delta_cY]$. These are exactly \eqref{v21:R}.
The quadratic constant drift also vanishes, using $LY=0$ and
\eqref{v21:translation-incidence}. This proves the complete expansion and
all the stated independence assertions.
\end{proof}

The last preparation term in \eqref{v21:R} matters. Leaving out
$-\ip{\delta_c\ell}{P_2(Y)}$ would differentiate at an unconstrained second
field coefficient and produce a different cubic mean map. Conversely this
proof does not allow arbitrary cubic weak rows to be computed from cubic
constant-shift envelopes: a nonconstant weak row has a nonzero linear
constraint, and its absolute cubic value can require the quartic Hamiltonian.

\section{An exact tilt cancelling all quadratic rows and all three cubic means}
\label{v21:mean-certificate}

Use precisely $Y$ from \eqref{v20:seed}. Its canonical entries are in
$\mathbb Q(\sqrt3)$, with $p=\sqrt3p_R$. Let $Q_3$ and $W_3$ be the
multiplier synthesis and positive physical weight of \eqref{v20:rational-system}.
Write $x_{20}$ for the exact weighted minimum-norm tilt in V20, so
$\ell_{20}=Q_3x_{20}$. Let $N_*$ be its supplied $108\times18$ rational
null matrix, and let $\mathsf m$ take the four component means of a multiplier.
In the rational coordinates define
\[
 E_0=\diag(\mathbf1_{27},\mathbf1_{27},\mathbf1_{27},\mathbf1_{27}),
 \qquad N_0=N_*-E_0\mathsf m N_*.
\]
Here $E_0$ is a $108\times4$ constant-column matrix. The matrix $N_0$ has
range equal to the normalized 14-dimensional null space; it need not have
independent columns.

For the three unit constant shifts form the rational vector $t$ and rational
$3\times108$ matrix $R$ by
\[
 t_i=t_{e_i}(Y),\qquad (Rx)_i=\mathcal R_Y[e_i,Q_3x].
\]
Use spatial sums consistently in this certificate; multiplying its three
rows by $1/27$ would change none of their zeros or ranks. The full original
$\widehat K,\widehat b$ from V20 retain that same sum convention.

\begin{theorem}[A characteristic-zero joint first-tilt certificate]
\label{v21:joint-tilt}
For this actual first jet, $t=0$. If $E$ consists of columns $0,1,2$ of
$N_0$, then $H_c=RE$ is invertible and
\begin{equation}
                         \det H_c=37\pmod{241}.
 \label{v21:mean-minor}
\end{equation}
Define the exact rational vector and its physical tilt by
\begin{equation}
 x_{21}=x_{20}-E H_c^{-1}(t+Rx_{20}),\qquad \ell_{21}=Q_3x_{21}.
 \label{v21:x21}
\end{equation}
They satisfy, in characteristic zero,
\begin{equation}
 \boxed{\widehat Kx_{21}=-\widehat b,\qquad t+Rx_{21}=0,
                     \qquad\mathsf m x_{21}=0.}
 \label{v21:joint-zero}
\end{equation}
Thus every original quadratic drift row and each cubic constant-shift row
vanishes on exactly prepared records with this first tilt. The normalized
first tilts satisfying both conditions form an affine space of dimension 11.
The certificate also gives
\begin{equation}
          361.82<\|\ell_{21}\|_h<361.84,
                  \qquad \max_{x,\mu}|\ell_{21}^\mu(x)|<520.
 \label{v21:tilt-bounds}
\end{equation}
No minimum-norm assertion is made for this new representative.
\end{theorem}
\begin{proof}
For reproducibility the constant-shift envelope is evaluated independently
as follows. Set $z=z_1(Y)$, use the original Lorentz bracket and ordered BCH
terms, and substitute
\[
 (z,A_0,w,u,p,N,\beta)
 =(az_1(Y),az_1(Y)_i,a\delta_iU,aU,ap,0,e_i)
\]
in the full scalar coefficient functional. The negative coefficient of
$a^3$ is $P_{e_i,3}(Y)$, by \eqref{v5:P3} and linearity in the auxiliary
and multiplier slots. In particular its cubic electric term is retained:
$A_0$ is not zero in this substitution. Quadratic and cubic homogeneous
polarization as in \eqref{v20:polarization} supplies all derivatives in
\eqref{v21:t}--\eqref{v21:R}. For the last term use exact discrete
summation by parts, or equivalently the component-major covector
$\delta_iP_2(Y)$. These prescriptions define every entry of the new matrix,
not only a selected test.

The field grading is also checked directly in those polynomials.
$t_c$ has types $U^3$ and $Up^2$; the lapse columns of $R$ have types
$U^2,p^2$, and its shift columns before $Q_3$ have type $Up$.
Accordingly, evaluation at $(U,p_R)$ and $(U,0)$ gives
\[
 t=3t(U,p_R)-2t(U,0),\qquad
 R_N=3R_N(U,p_R)-2R_N(U,0),\qquad R_\beta=3R_\beta(U,p_R).
\]
All these calculations use arbitrary-precision fractions on the original
three-site arrays and $\delta_i=3(S_i-S_i^{-1})$. They give exactly three
zero entries for $t$, and all 324 rational entries of $R$. The two unscaled
constant vectors are zero separately. These zeros are not inferred from
reduction modulo a prime.

Subtracting the means from the null columns preserves their null identity,
because all four constant columns are null in V20. It gives
$\widehat K E=0$ and $\mathsf m E=0$. The determinant of $RE$ has the
nonzero residue \eqref{v21:mean-minor}, with every denominator a unit at
241. Hence it is nonzero over $\mathbb Q$. Substitution in \eqref{v21:x21}
proves all equations in \eqref{v21:joint-zero}; the short verifier additionally
multiplies the rational arrays against every one of the 108 original rows
and all three new rows. Rank three on the 14-dimensional normalized kernel
gives dimension $14-3=11$ for the affine solution set.

Finally the exact rational value $x_{21}^TW_3x_{21}/27$ lies between the
squares of the two endpoints in \eqref{v21:tilt-bounds}. Every physical
component square, namely $x_j^2$ or $3x_j^2$, is less than $520^2$.
This proves the bounds without a floating square root. An independent formal
AD evaluation of the unchanged scalar action at the good specialization
$\sqrt3\mapsto56\pmod{241}$ matches all 324 new entries and all three
constant terms. That regression checks the new explicit envelope against the
original functional; it is separate from the characteristic-zero cancellations.
\end{proof}

The finite norm of $\ell_{21}$ permits the original chart after reducing the
amplitude. The records have $\lambda=e+a\ell_{21}+O_h(a^2)$, not a large
fixed lapse perturbation. No refinement-uniform bound on the tilt or its
spatial derivatives follows from this fixed-cutoff certificate.

\section{Which second initial records change the remaining cubic weak rows}
\label{v21:second-control}

Fix a first tilt $\ell$ satisfying the quadratic range equation of V20.
Let $d_{2,W}(Y)$ be the complete un-tilted quadratic weak drift, on all weak
shift tests. For a canonical variation $Z$ and a second multiplier variation
$m$, define the finite control covector
\begin{equation}
 \boxed{\mathcal C_{Y,\ell}(Z,m)[\nu]
 =D d_{2,W}(Y)[Z,\nu]+\mathbb K_1(Z)[\nu,\ell]
                           +\mathbb K_1(Y)[\nu,m],\quad\nu\in\mathcal W.}
 \label{v21:C}
\end{equation}
The notation $D d_{2,W}(Y)[Z,\nu]$ means the field derivative applied to
$Z$, then evaluated on the multiplier test $\nu$.

Admissible homogeneous second-field changes at fixed first tilt form
\begin{equation}
 \mathcal Z_Y^{(2)}=\{Z:LZ=0,\quad P_0 DP_2(Y)[Z]=0\},
 \label{v21:Z2}
\end{equation}
where $P_0$ is spatial mean in the four original constraint rows. The second
condition includes the actual anisotropic momentum and both complementary
polarizations; a scalar-momentum cancellation is not used to remove it.

\begin{lemma}[Exact cubic variation and extension of admissible second records]
\label{v21:second-law}
At fixed $Y,\ell$, the change of the complete cubic weak drift under a change
$Y_2\mapsto Y_2+Z$ and $\lambda_2\mapsto\lambda_2+m$ is exactly
$\mathcal C_{Y,\ell}(Z,m)$. It is affine in these two second coefficients
and independent of all higher initial coefficients. For $Z\in\mathcal Z_Y^{(2)}$
and a normalized $m$, these changes extend to analytic exact original
preparations, locally in their finite parameters. Furthermore
\begin{equation}
        \mathcal C_{Y,\ell}(Z,m)[(0,c)]=0
                 \quad(Z\in\mathcal Z_Y^{(2)},\ c\hbox{ constant}).
 \label{v21:C-means}
\end{equation}
\end{lemma}
\begin{proof}
For a pure shift the un-tilted linear drift is identically zero on all field
records. Expand the full finite analytic identity of V20,
$D_\lambda\mathsf B[\mu]=\mathbb K\mu+D_X(M\mu)[F]$.
At $X=aY+a^2Y_2+\cdots$, $\lambda=e+a\ell+\cdots$, the change of the
$a^3$ coefficient caused by $Z$ in the $a^2$ field slot is
$D d_{2,W}(Y)[Z]+\mathbb K_1(Z)\ell$. Terms involving $D_XM\,F$
first occur at order $a^3$ in the drift tilt, but their value at that order
uses only $Y,\ell$, so varying $Y_2$ affects them only at the next order.
Changing the second multiplier adds $\mathbb K_1(Y)m$, because
$D_\lambda\mathsf B=a\mathbb K_1(Y)+O_h(a^2)$ on this curve.
Higher field coefficients enter the zero linear shift drift and higher
multiplier coefficients enter a derivative of order $a$, so they do not
change this cubic weak coefficient. This proves \eqref{v21:C} without
deleting the higher action terms from the \emph{baseline} cubic value.

For another direct check, expand the canonical field through degree three.
The terms linear in $Z$ are precisely the field derivative of the un-tilted
quadratic weak Poisson coefficient plus the two polarizations defining
$\mathbb K_1(Z)[\nu,\ell]$. The terms linear in $m$ are the analogous
first-bracket polarization at $Y$. The possible quartic Hamiltonian term
multiplies the linear weak row but is fixed by $Y$ at this coefficient;
it does not enter the displayed derivative in the second controls.

The second initial equation is $LY_2=-P_2(Y)$. Differences must therefore
have $LZ=0$. The four next mean equations change by
$P_0DP_2(Y)[Z]$, giving the second condition in \eqref{v21:Z2}. A second
multiplier perturbation changes the constraints only at order $a^4$, since
$M=O_h(\|X\|^2)$; it does not alter this condition. The original kernel
range solve and four mean balancing directions therefore extend each such
second change exactly as in \cref{v5:extend}, now with the tilted initial
multiplier and the V12 mean calibration. Their leading balancing derivative
is unchanged and nonsingular. Their analytic dependence gives families for
bounded neighborhoods of the finite parameters. The higher corrections do
not change the second field coefficient just selected.

Finally the derivative of $d_{2,W}$ on a constant test is zero on all fields
by \cref{v5:constant-B2}; \eqref{v21:translation-incidence} gives
$\mathbb K_1(Z)[c,\ell]=\ip{\delta_c\ell}{LZ}_h=0$ and
$\mathbb K_1(Y)[c,m]=0$. This proves \eqref{v21:C-means}.
\end{proof}

The absolute cubic weak vector need not be explicitly tabulated to prove
controllability. It is a finite analytic right side on the original stationary
chart. An onto \emph{actual} derivative on admissible second records solves
for that right side without assuming it is zero. The next section establishes
this onto property at the new, exactly specified tilt; it is not taken as an
additional hypothesis.

\section{A full-rank control on exactly admissible second field records}
\label{v21:control-certificate}

Use raw independent symmetric coordinates at each of the 27 sites, with the
independent entries ordered as $(11,22,33,12,13,23)$. The first 162 coordinates are configuration;
the next 162 are momentum. If $z\in\R^{324}$, let $\mathsf J_f z$ be the
physical variation with those configuration coordinates and momentum
coordinates multiplied by $\sqrt3$. This is an invertible real coordinate
change, not a restriction on possible variations.

Let $\mathsf L$ contain 26 of the 27 point rows of each original linear
constraint, omitting the last site in each group. Divide its shift rows by
$\sqrt3$ after composing with $\mathsf J_f$. They are rational. The omitted
rows are recovered because every linear constraint has zero spatial mean.
Append the four rows $P_0DP_2(Y)\mathsf J_f$, again dividing the three shift
rows by $\sqrt3$; use spatial sums rather than means for these four rows.
Call the resulting $108\times324$ rational matrix $\mathsf A_c$. Its kernel
is exactly $\mathsf J_f^{-1}\mathcal Z_Y^{(2)}$.

On the weak space use the three constant shifts followed by
\[
 \beta_j=\nabla_\delta(e_j-e_{26}),\qquad 0\le j<26,
\]
where $e_j$ is a scalar point indicator in lexicographic site order. Let
$\mathsf C_g$ be the $26\times324$ matrix
\begin{equation}
 (\mathsf C_gz)_j=
       \mathcal C_{Y,\ell_{21}}(\mathsf J_fz,0)[(0,\beta_j)],
 \qquad \mathsf B_c=\begin{pmatrix}\mathsf A_c\\\mathsf C_g\end{pmatrix}.
 \label{v21:stacked-control}
\end{equation}
Its entries are rational: the weak drift has field types $U^2,p^2$, and
the same-type and mixed first-bracket blocks have precisely the grading in
\cref{v20:certificate}. The $\sqrt3$ factors in the field and tilt synthesis
therefore occur in pairs in these weak rows.

\begin{theorem}[Actual constrained second-field control of every nonconstant weak row]
\label{v21:control-theorem}
At the unchanged first jet \eqref{v20:seed} and the exact new tilt
\eqref{v21:x21},
\begin{equation}
 \boxed{\rank\mathsf A_c=108,\qquad\rank\mathsf B_c=134,
 \qquad\rank\bigl(\mathcal C_{Y,\ell_{21}}(\cdot,0)\big|_{
                  \mathcal Z_Y^{(2)}}\bigr)=26.}
 \label{v21:control-ranks}
\end{equation}
The last map has image the entire nonconstant weak covector space and zero
constant components. Second multiplier-value changes are not needed to obtain
this rank.
\end{theorem}
\begin{proof}
The original linear-constraint rank is $4V-4=104$. The four quadratic mean
calibration derivatives of V12 make the additional mean rows independent;
the certificate below independently checks the combined rank. There are
only 108 rows in $\mathsf A_c$ and only 134 in $\mathsf B_c$.

Here is a compact specification of a nonzero 134-column minor. An interval
$a:b$ denotes every integer from $a$ to $b$, inclusive, and all indices
are zero based. In their increasing order take
\begin{equation}
 \begin{split}
 \mathcal J={}&(0:47,\ 49:50,\ 53:58,\ 162:175,\ 177:193,\ 195:198,\\
 &201:202,\ 204,\ 207:208,\ 210,\ 213:214,\ 216:229,\\
 &231:247,\ 252,\ 255,\ 258,\ 264).
 \end{split}
 \label{v21:control-indices}
\end{equation}
The corresponding 108-column minor for $\mathsf A_c$ replaces the first
three intervals by
\[
 (0:13,\ 15:16,\ 18:23,\ 25,\ 27,\ 31,\ 37,\ 43,\ 49,\ 55:56),
\]
and retains the rest from $162:175$ onward. Under the same good-prime
specialization as V20, the exact values are
\begin{equation}
 \boxed{\det(\mathsf A_c)_{:,\mathcal J_c}=25\pmod{241},
              \qquad\det(\mathsf B_c)_{:,\mathcal J}=132\pmod{241}.}
 \label{v21:control-minors}
\end{equation}
All denominators, including those of the exactly solved tilt, are units
at this prime. Consequently both determinants are nonzero in characteristic
zero and the two row ranks are maximal.

For clarity, the coefficient evaluation uses no quartic baseline substitution.
It differentiates the exact polynomial expression
\[
 d_{2,W}(Y)[\nu]
 =D\Phi_3(Y)[G\nu]+DP_{\nu,2}(Y)[F_1Y]
       \qquad(\nu\hbox{ a pure weak shift}),
\]
and adds the derivative of the exact polarization defining
$\mathbb K_1(Y)[\nu,\ell_{21}]$. Thus the resulting field derivative is
exactly \eqref{v21:C} with $m=0$. The flat vector fields, both off-diagonal
Frobenius occurrences, the constant auxiliary terms, and the complete electric
and magnetic coefficients are retained. The matrix is regenerated from the
original scalar polynomial functional using modular formal adjoints. All
forward and reverse arithmetic is integer-valued; the verification checks its
intermediate bounds below the exact binary arithmetic threshold. An independent
polynomial-value polarization of the original prepared drift checks the whole
29-row weak image on two specified full-field variations.

For a finite linear map $A$ and an additional row map $C$,
$\rank\left(\begin{smallmatrix}A\\C\end{smallmatrix}\right)
-\rank A=\rank(C|_{\ker A})$: split the domain into $\ker A$ and a
complement, and eliminate $C$ on that complement. Applying this identity
to the two certified ranks proves the final rank in \eqref{v21:control-ranks}.
Its constant output is zero by \cref{v21:second-law}. This proves actual
surjectivity on the admissible second records, not only on unconstrained
field directions.
\end{proof}

An exact right inverse, if desired for reconstruction of the controls, is
specified without a numerical nullspace. Let $E_{\mathcal J}$ insert the
columns in \eqref{v21:control-indices}. Then
\begin{equation}
 \mathsf Z_*=
 \mathsf J_f E_{\mathcal J}(\mathsf B_cE_{\mathcal J})^{-1}
                  \begin{pmatrix}0_{108\times26}\\ I_{26}\end{pmatrix}
 \label{v21:right-inverse}
\end{equation}
defines 26 real physical second-field directions. Every column satisfies
both original homogeneous second constraints, and its complete weak control
has zero means and the indicated unit nonconstant covector. This is a
characteristic-zero definition using the proved nonsingular matrix. The
large rational inverse need not be rounded or printed to define it.

\section{Joint analytic preparation with zero entire weak velocity}
\label{v21:joint-preparation}

The mean and second-field controls can now be combined without assuming that
the still unlisted baseline cubic weak target is zero. Fix the physical
active slice consisting of mean-zero lapse and mean-zero transverse shift,
and the physical curl-free shift slice. The old active point representatives
can be made to lie in this active slice by subtracting their constant lapse
and curl-free shift parts. Their leading $J_1$ columns are unchanged.
Together with the same weak basis they form a fixed direct complement to the
constant-lapse direction. The leading graded Gram is the same invertible
one certified in V12, up to the corresponding invertible coordinate changes.

\begin{theorem}[Exact original-action initial preparation with small complete rate]
\label{v21:main}
There are analytic original initial records at $h=1/3$, for sufficiently
small nonzero $a$, with
\begin{equation}
 \boxed{\begin{gathered}
 X(a,0)=aY+O_h(a^2),\qquad
 \lambda(a,0)=e+a\ell_{21}+O_h(a^2),\qquad\overline{\lambda^N(a,0)}=1,\\
 P(X(a,0),\lambda(a,0))=0,\qquad
 \mathsf B(X(a,0),\lambda(a,0))=O_h(a^3),\\
 \mathbb P_{\rm cf}\lambda_t^\beta(a,0)=0,\qquad
                         \lambda_t(a,0)=O_h(a).
 \end{gathered}}
 \label{v21:main-eq}
\end{equation}
The last two observations are the initial velocity of each record's own
normalized original-action continuation. The initial shift means can also
be kept zero. Each sufficiently small nonzero amplitude has a unique local
exact continuation on its original stationary chart. No multiplier velocity,
field force, or projected evolution is substituted after preparation.
\end{theorem}
\begin{proof}
Let $E$ be the three normalized null-tilt columns in \cref{v21:joint-tilt}
and put $\ell(s)=\ell_{21}+Q_3Es$ for $s\in\R^3$. For every $s$ near
zero the quadratic drift is still zero, because
$\widehat K E=0$, and all four initial multiplier means have the retained
normalization. Take the original analytic initializer with initial values
$e+a\ell(s)$. Its exact constrained fields have first coefficient $Y$.

Add second-field controls $z\in\R^{26}$ using \eqref{v21:right-inverse}.
The extension in \cref{v21:second-law} gives an analytic exactly constrained
family $X(a,s,z)$ whose second coefficient changes by $\mathsf Z_*z$.
All finite $z$ in a neighborhood of any fixed value are allowed after reducing
the amplitude. No bound uniform in unbounded control values is asserted.
The constraints are solved identically, not added later to a drift-zero curve.

Let $B_3(s,z)$ be the actual complete cubic drift coefficient of this family
and restrict it to the 29 weak tests. At $s=0$ its three mean components are
zero by \cref{v21:joint-tilt}; the 26 other components depend affinely on $z$
with derivative $I_{26}$, by \eqref{v21:right-inverse}. Hence there is a
finite $z_*$ making \emph{all} weak components of $B_3(0,z_*)$ zero.
This solves for the actual analytic right side, whatever its baseline value;
no truncation of the quartic action entering that baseline is used.

On the fixed normalized multiplier complement write the true Hessian as
$D_a\Gamma(a,s,z)D_a$, where
$D_a=\diag(aI_A,a^2I_W)$. V20's graded-source result shows that
$\Gamma(0,s,z)=\Gamma_0$ is independent of $s,z$ here: the canonical first
jet $Y$ is unchanged. This is the original signed Gram, not a positive
replacement of its higher weak block. Its nonzero inherited determinant
makes it invertible near $(0,0,z_*)$. The exact normalized drift satisfies
\[
 D_a^{-1}\mathsf B(a,s,z)
       =a\binom0{B_{3,W}(s,z)}+O_h(a^2).
\]
Thus the scaled actual weak velocity
\[
 \Psi(a,s,z)=a\,[\mathbb P_{\rm cf}\lambda_t^\beta(a,0;s,z)]_W
\]
has an analytic extension through $a=0$. Brackets denote its coordinates
in the fixed weak basis. If $D_0$ is the weak Schur complement of
$\Gamma_0$, the extension is
\[
                   \Psi(0,s,z)=-D_0^{-1}B_{3,W}(s,z).
\]
In particular $\Psi(0,0,z_*)=0$. Multiply the derivative at this point by
$-D_0$, order the three mean rows before the other weak rows, and order
$s$ before $z$. Its matrix is
\begin{equation}
                       \begin{pmatrix}H_c&0\\ *&I_{26}\end{pmatrix}.
 \label{v21:joint-derivative}
\end{equation}
Both diagonal blocks are nonsingular. With spatial-mean rather than
spatial-sum coordinates the first block has a harmless nonzero factor;
the same factor appears in its source definition. There is no derivative of
$D_0^{-1}$ multiplying a nonzero target at the base point. The analytic
implicit function theorem therefore gives $s(a),z(a)$ with
$s(0)=0,z(0)=z_*$ and $\Psi(a,s(a),z(a))=0$. For nonzero $a$ this is
exactly the claimed entire curl-free velocity condition. Since
$s(a)=O_h(a)$, the first initial tilt is still $\ell_{21}$.

Every resulting rate lies in the physical active slice: its lapse mean is
zero by normalization, and its curl-free shift part has just been made zero.
Restrict its own exact consistency equation
$M\lambda_t=-\mathsf B$ to that active slice. The restricted matrix is
$a^2$ times a matrix tending to the invertible positive active Gram of V12;
the target is $O_h(a^3)$. Thus the whole actual rate is $O_h(a)$.
This step uses the exact weak-rate cancellation, not merely vanishing of its
leading pole.

The original graded signed Gram remains regular at sufficiently small
nonzero amplitude. Its only full null direction is the radial multiplier,
which is transverse to the fixed mean-lapse-zero complement. The inherited
normalized continuation theorem therefore supplies each local exact history,
preserving every original constraint. Reducing the amplitude keeps the spatial
metric and lapse positive. No estimate of the common physical time of these
histories has been used or proved.
\end{proof}

The condition in \eqref{v21:main-eq} is imposed by selecting the initial
records on which the original rate has the stated value. It is not an
instruction to reset the multiplier during a history. The actual cubic
nonconstant target and the analytic nonlinear roots are not represented as
numerically integrated trajectories; the onto certificate and the analytic
implicit function proof give their existence.

\section{Initial metric and literal curvature now have the same small-amplitude scale}
\label{v21:geometry}

\begin{corollary}[Removal of the initial metric-jet obstruction on the new branch]
\label{v21:initial-geometry}
On the initial records of \cref{v21:main}, for every fixed spatial Sobolev
order at this cutoff,
\begin{equation}
 \begin{gathered}
 X=O_h(a),\quad X_t=O_h(a),\quad X_{tt}=O_h(a),\qquad
       \lambda-e=O_h(a),\quad\lambda_t=O_h(a),\\
 \boxed{\|\operatorname{Riem}(g_h)(a,0)\|=O_h(|a|),\qquad
 \|E_h^{\rm link}(a,0)\|+\|F_h^{{\rm sp},{\rm link}}(a,0)\|
                                                       =O_h(|a|).}
 \label{v21:curvature-small}
 \end{gathered}
\end{equation}
Here $g_h$ is the same ordinary interpolated ADM metric as V10 and the link
reads are the original ordered reads on the same history. The initial
Einstein tensor is $O_h(a)$ and the metric Riemann-squared scalar is
$O_h(a^2)$. In fixed compatible frames, their initial metric/link curvature
difference is also $O_h(a)$. No equality of the coefficients divided by $a$
is asserted.
\end{corollary}
\begin{proof}
At fixed $h$, the original analytic field vector is
$F(X,e+a\ell+O(a^2))=O_h(a)$ by \cref{v20:flat}. Its field and multiplier
differentials stay bounded on the chosen stationary neighborhood.
The exact identity $X_{tt}=D_XF\,F+D_\lambda F\,\lambda_t$ and
\cref{v21:main} give the second field derivative bound. Finite-dimensional
spatial norm equivalence gives each stated fixed-order version.

The initial lapse and shift are now small nonconstant fields; they must not
be replaced by exactly one and zero in the metric calculation. Nevertheless
$\gamma-I,N-1,\beta$, their needed spatial derivatives, $\gamma_t,N_t,
\beta_t$ and $\gamma_{tt}$ are all $O_h(a)$ on the original interpolated
record. The initial ADM curvature uses precisely these jets, together with
spatial derivatives of $\gamma,N,\beta$ and at most first spatial derivatives
of their first time jets. In coordinates this is seen from
\[
 R_{0i0j}=\tfrac12(g_{0j,it}+g_{i0,tj}-g_{00,ij}-g_{ij,tt})
                    +\hbox{Christoffel products};
\]
the other independent blocks use no more time derivatives. Neither $N_{tt}$
nor $\beta_{tt}$ occurs in these expressions. The second-derivative terms
are $O_h(a)$, while Christoffel products are $O_h(a^2)$. This proves the
metric estimate without assuming a bound for the next multiplier derivative.
The inverse metric and frame transformations are uniformly bounded on this
fixed chart, so the Einstein tensor has the same order and the curvature
square has order $a^2$.

The initial stationary spatial and temporal connections are $O_h(a)$.
Their velocity is obtained by differentiating the original stationary map:
$z_t=D_Xz\,X_t+D_\lambda z\,\lambda_t=O_h(a)$ on this fixed analytic
chart. Every original temporal or spatial ordered link curvature therefore
has order $a$ by its analytic flat expansion; all spatial $h$ factors are
fixed here. This also follows from the retained graded signed-response
argument, without changing the link formulas. Putting both observations in
compatible fixed frames and subtracting gives the last bound.
\end{proof}

No affine canonical transport or coordinate renormalization is needed for
\eqref{v21:curvature-small}. It concerns a different exactly prepared initial
record of the same action, not a removal of V10's invariant curvature at its
old circular phase root. Both curvature observations tend to the flat value
here. This is an initial-cut statement, not a proof that their first-order
coefficients agree, that their later difference remains small, or that the
nonlinear development satisfies the Einstein equation.

\section{Persistence within the normalized first-tilt solution set}
\label{v21:persistence}

\begin{corollary}[An eleven-parameter family of admissible leading tilts]
\label{v21:eleven}
Let
\[
 \mathcal N_{21}=\ker\widehat K\cap\ker\mathsf m\cap\ker R.
\]
It has dimension 11. For every $\eta$ in a sufficiently small neighborhood
of zero in $\mathcal N_{21}$, there are analytic exact initial preparations
with the same canonical first coefficient $Y$, initial multiplier first
coefficient $Q_3(x_{21}+\eta)$, and every conclusion of
\cref{v21:main,v21:initial-geometry}. Their dependence on $\eta$ is analytic
on a common small fixed-cutoff neighborhood.
\end{corollary}
\begin{proof}
The dimension follows from \cref{v21:joint-tilt}. These leading tilts satisfy
both complete quadratic cancellation and the cubic mean equations identically.
The constrained second-field control is polynomial in the first tilt, so its
certified nonzero minor remains nonzero nearby. Alternatively retain the 26
fixed second directions at $\eta=0$; their control matrix is $I+O(\eta)$
and stays invertible on a smaller neighborhood. The three mean controls
$Q_3Es$ still have derivative $H_c$. Repeat the joint analytic implicit
function argument with $\eta$ as an external parameter. Exact initial
constraints and the rate condition persist simultaneously. On a compact
smaller parameter neighborhood all finite matrices, base controls and analytic
chart radii have common fixed-cutoff bounds. This supplies the stated local
parameter family without claiming a cutoff-independent one.
\end{proof}

This is a family of initial multiplier coefficients and higher preparations
at a specified canonical first jet. It is not an open basin of arbitrary
microscopic laws or generic canonical Cauchy data. The original choice
$\lambda=e$ is not included merely by a change of terminology.

\section{Next dynamical incidence and reproducibility}
\label{v21:status}

The new preparation removes all previously identified initial multiplier and
metric-jet poles on this fixed complementary branch. It is achieved with the
original action by three first-tilt controls and 26 exactly admissible second
field controls. The original quadratic drift vanishes, the complete actual
initial weak velocity is zero, and the remaining active velocity is of order
amplitude. The literal and reconstructed metric curvatures have the same
vanishing initial scale. In particular no bounded-rate assumption is left
unverified at this initial point.

The next dynamical question is different. The prepared zero set need not be
invariant under the exact action vector field. A concrete next observation is
$\frac{d}{dt}(\mathbb P_{\rm cf}\lambda_t^\beta)$ along that vector field,
with the initial records just constructed and all differentiated Hessian and
source terms retained. Re-solving the initial selection at successive times
would define a new law and is not a continuation theorem for the old one.
Nor do finite analytic initial cancellations give uniform bounds on
$\lambda_{tt}$, the size of a rank neighborhood, or a physical lifespan.

The second independent requirement is refinement: the first Poisson range
certificate, cubic mean controls, constrained second-control rank and their
norms have here been established only at $N=3$. No cofinal tilted family or
uniform amplitude radius is asserted. The source $g_h$ in V15--V19 belongs
to the old unit-initial-value branch and cannot be reused without recalculating
its changed initial data. The complete higher weak signed system also remains
in the nonlinear continuation. Thus the raw-action Einstein continuum bridge
is not closed, although its initial obstruction is removed on an actual
nonempty branch rather than merely in a comparison scheme.

\paragraph{Verification scope.}
The new exact rational program regenerates all three constant cubic values
and every entry of their $3\times108$ tilt map. Its chosen tilt is checked
against all 108 inherited quadratic rows, the three new mean rows and four
mean normalizations over $\mathbb Q$. A separate full scalar-envelope AD
evaluation reproduces the entire new mean map at a good prime. The control
program regenerates the complete $29\times324$ weak field derivative and
the original 108 combined homogeneous constraint rows, certifies the stated
nonzero minors, and verifies the constant-row dependence on the linear
constraints. Independent value polarizations of the original prepared drift
check every weak row on two further full-field directions. These are finite
coefficient checks, not nonlinear spacetime integrations.

The proof of exact nonlinear preparation uses the retained analytic stationary
and initial-constraint theorems with the newly verified controls. Its baseline
nonconstant cubic weak target and the resulting nonlinear initial roots are
not represented as numerically tabulated. No proof-assistant formalization
or independent verification of all inherited results is claimed. All twenty
preceding marked bodies are preserved byte for byte. The package records
input/output hashes, the exact certificates, regeneration code and the
actually executed verification stages.
% END IMMUTABLE EXACT-ACTION BRIDGE V21 BODY


% Future iterations append here.
% BEGIN IMMUTABLE EXACT-ACTION BRIDGE V22 BODY
\clearpage
\renewcommand{\thesection}{\arabic{section}}
\renewcommand{\theHsection}{bridgeV22.\arabic{section}}
\setcounter{section}{178}
\section{Transport of the prepared records, not repeated preparation}
\label{v22:context}

The added Yang--Mills source \cite{v22:YM86} proves local transport estimates
for its complete restricted Wilson law and uses a literal plaquette second
difference to obtain nonconcentration at a specified scale. Its final
conclusion keeps that scale distinct from a thermal normalization and from
a physical continuum limit. Those probability estimates do not apply to the
deterministic Palatini continuation here: no corresponding gravitational
probability law, translation estimate, or action-density comparison has been
supplied. The useful transfer is methodological and more limited. One must
transport the original object, retain its complete defining equations, and
measure its change at the normalization at which a conclusion is sought.
A new initial root at every time is not such a transport.

We therefore test the initial set constructed in \cref{v21:main} against
its own exact action vector field. A new null test of the first Poisson
bracket detects an initial tangency defect independent of every unlisted
higher preparation coefficient. At the actual rationally specified first
tilt, that defect is nonzero. The entire weak velocity is zero initially,
but its derivative has magnitude at least a constant times $a^{-2}$ as the
positive preparation amplitude $a$ tends to zero, with $h=1/3$ fixed.
This is a statement about the next multiplier derivative, not a reversal of
V21's proved small initial velocity or small initial curvature.

There is also a constructive dynamical conclusion. The original histories
exist on every fixed interval $0\le t\le a^3S$, for all sufficiently small
$a$ depending on $S$. Their canonical and multiplier displacements on that
interval are $O_h(a^4)$, and their complete rates and both curvature reads
remain $O_h(a)$. An explicit finite linear system is the limit of their
rescaled displacement. A particular scalar observation of the actual weak
velocity, divided by $a$, converges to a nonzero linear function of $s=t/a^3$.
A constant-rate-free test also detects the layer in the original metric
curvature and excludes an amplitude-linear curvature bound on a fixed physical
interval for this family. Thus the initial zero set is not invariant, although
unscaled observations remain small throughout the proved, shrinking interval.

\paragraph{Audit and dependencies.}
The new input was reviewed through its V85--V86 transport and fluctuation
results and their stated normalization dependencies; this is not an independent
verification of its entire cumulative Yang--Mills programme. We retain the
original gravitational stationary chart, the universal weak first-source
cancellation, V20's tilted germ and full Poisson matrix, and V21's exact
initial preparation and signed-source regularity. In particular the absolute
cubic nonconstant preparation baseline is still not computed here. The new
coefficient test does not require it. A targeted search of the preceding
bridge bodies found the proposed tangency test in \cref{v21:status}, not its
evaluation or the physical-time layer proved below. Consistent multiplier
selection and preservation of an additional chosen condition are distinct
questions also in the general constrained-discretization literature
\cite{v22:DBGP}; no external estimate supplies the coefficient or lifespan
claimed in this body.

\section{A first-jet tangency covector with a preparation-independent test}
\label{v22:drift}

Keep $H,P,F,M,\mathbb K,L,G$ of V20, the pure weak shift space
$\mathcal W=\ker\curl_\delta$, and the physical active complement
$\mathcal A$ of V21. The mean-lapse-zero multiplier space is the fixed direct
sum $\mathcal A\oplus\mathcal W$. Fix real bases there; covectors below
are evaluated on those same test columns. The first Poisson matrix is
$\mathbb K_1(Y)$, and it is not the symmetric multiplier Hessian.
Let $\ell=\ell_{21}$ initially. For any analytic exact preparation supplied
by \cref{v21:main}, write
\begin{equation}
 \begin{gathered}
 X_a=aY+O_h(a^2),\qquad \lambda_a=e+a\ell+O_h(a^2),\qquad P(X_a,\lambda_a)=0,\\
 r_a:=\lambda_t(a,0)=(r_{A,a},0),\qquad
 r_{A,a}=a v_A+O_h(a^2),\qquad V=F_1Y+G\ell .
 \end{gathered}
 \label{v22:base}
\end{equation}
The vector $v_A$ is the actual coefficient of this preparation's rate. It
need not be known numerically. Its existence follows from the analytic
active inverse in V21 after the exact weak-rate cancellation. All constants
in this body can depend on the fixed cutoff and the chosen analytic
preparation; none is asserted uniform under refinement.

Define the weak covector-valued linear maps
\begin{equation}
 \begin{split}
 \mathcal C_\ell Z[\nu]
   &=Dd_{2,W}(Y)[Z,\nu]+\mathbb K_1(Z)[\nu,\ell],\\
 \mathcal K_W\eta[\nu]&=\mathbb K_1(Y)[\nu,\eta],\qquad \nu\in\mathcal W,
 \end{split}
 \label{v22:C}
\end{equation}
and the weak Poisson-null space
\begin{equation}
 \mathcal N_W=\{\nu\in\mathcal W:
               \mathbb K_1(Y)[\nu,\mu]=0\ \hbox{for every multiplier }\mu\}.
 \label{v22:NW}
\end{equation}
For $\nu\in\mathcal N_W$ the candidate tangency observation is
\begin{equation}
 \boxed{\Delta_\nu(Y,\ell)=\mathcal C_\ell(F_1Y+G\ell)[\nu].}
 \label{v22:Delta}
\end{equation}
It depends only on the canonical first jet and first multiplier tilt.

\begin{lemma}[The actual weak drift and its first derivatives]
\label{v22:weak-Taylor}
In a neighborhood of $(X,\lambda)=(0,e)$, the weak restriction of the
original canonical drift has the joint Taylor expansion
\begin{equation}
 \mathsf B_W(X,e+\zeta)
   =d_{2,W}(X)+\mathbb K_1(X)|_W\zeta
                           +O_h((\|X\|+\|\zeta\|)^3).
 \label{v22:joint-Taylor}
\end{equation}
At \eqref{v22:base}, its differentials are
\begin{equation}
 D_X\mathsf B_W=a\mathcal C_\ell+O_h(a^2),\qquad
 D_\lambda\mathsf B_W=a\mathcal K_W+O_h(a^2).
 \label{v22:DB}
\end{equation}
Moreover $LV=0$. Along the actual exact continuation,
\begin{equation}
 \left.\frac{d}{dt}\mathsf B_W\right|_{t=0}
       =a^2\{\mathcal C_\ell V+\mathcal K_W v_A\}+O_h(a^3),
 \label{v22:dotB}
\end{equation}
where $v_A$ is inserted as an active multiplier. Consequently
\begin{equation}
 \boxed{\left.\frac{d}{dt}\mathsf B[\nu]\right|_{t=0}
               =a^2\Delta_\nu(Y,\ell)+O_h(a^3)
                       \quad(\nu\in\mathcal N_W).}
 \label{v22:invariant-derivative}
\end{equation}
\end{lemma}
\begin{proof}
At $\lambda=e$ the weak linear drift vanishes identically on every field,
not just on the initial constraint set, by \cref{v5:prepared-B2}.
The quadratic term is $d_{2,W}(X)$. V20 proves
$\mathsf B(0,\lambda)=0$ and
$D_\lambda\mathsf B[\mu]=\mathbb K\mu+D_X(M\mu)[F]$.
Near the joint origin, $\mathbb K=\mathbb K_1(X)+O_h((\|X\|+\|\zeta\|)^2)$,
$D_XM=O_h(\|X\|)$, and $F=O_h(\|X\|+\|\zeta\|)$.
Integrating this identity in $\zeta$ gives \eqref{v22:joint-Taylor}; its
extra $D_XM\,F$ contribution is of total degree at least three. The exact
analytic Taylor formula also controls fixed differentials of the remainder.
Differentiation at $X_a=aY+O(a^2)$ and $\zeta_a=a\ell+O(a^2)$ yields
\eqref{v22:DB}. No constraint or higher coefficient was dropped to obtain it.

The original quadratic field satisfies
$L F_1Y=0$ whenever $LY=0$: its momentum divergence is zero and its scalar
constraint derivative is a derivative of the initial momentum constraint,
as in the proof of \cref{v5:prepared-B2}. Also $LG=L\mathbb J L^*=0$.
Thus $LV=0$. The actual initial field velocity is $aV+O_h(a^2)$, and
\eqref{v22:base} gives its multiplier velocity. Apply the chain rule to
\eqref{v22:DB} to get \eqref{v22:dotB}. The definition of $\mathcal N_W$
annihilates its unknown $v_A$ term. This proves \eqref{v22:invariant-derivative}
for every analytic preparation with these first jets.
\end{proof}

For the active rows let $\mathsf B_1=D_X\mathsf B(0,e)$. The inherited
flat identity is
\begin{equation}
 \mathsf B_1Z[ (n,\beta)]=-P_{(0,\nabla_\delta n),1}(Z).
 \label{v22:B1}
\end{equation}
It factors through $LZ$, vanishes on weak tests, and satisfies
$\mathsf B_1V=0$. It follows from the joint Taylor expansion in those rows
that $\dot{\mathsf B}_A(a,0)=O_h(a^2)$ as well. This observation prevents
an unretained linear scalar-constraint term in the elimination below.

\section{A new exact original-source obstruction at the selected tilt}
\label{v22:certificate}

Use the exact $N=3$ data \eqref{v20:seed}, the rationally defined
$\ell_{21}=Q_3x_{21}$ from \eqref{v21:x21}, and lexicographic point
indicators $e_j$, $0\le j<27$. Put
\begin{equation}
 \varphi_*=-e_1+e_2-e_4+e_5-e_7+e_8,
 \qquad \nu_*=(0,\nabla_\delta\varphi_*).
 \label{v22:test}
\end{equation}
The derivatives are the original rational matrices
$\delta_i=3(S_i-S_i^{-1})$. In particular this is an actual nonconstant
curl-free test, not a continuum test substituted for a finite one.

\begin{theorem}[A nonzero tangency defect invisible to all higher initial controls]
\label{v22:nonzero}
For these original coefficients,
\begin{equation}
 \boxed{\nu_*\in\mathcal N_W,\qquad
         \|\nu_*\|_h^2=10,\qquad
         \Delta_*:=\Delta_{\nu_*}(Y,\ell_{21})\ne0.}
 \label{v22:nonzero-eq}
\end{equation}
With normalized spatial-mean pairing the fresh rational calculation gives
\begin{equation}
           -6304458<\Delta_*/\sqrt3<-6304457.
 \label{v22:delta-interval}
\end{equation}
Equivalently the spatial-sum coefficient divided by $\sqrt3$ is rational
and has residue $143$ modulo $241$. Under $\sqrt3\mapsto56$, the complete
spatial-sum coefficient has residue $55$. The size in
\eqref{v22:delta-interval} uses the inherited unit-torus and amplitude
normalization; no new physical unit or field normalization is implied.
\end{theorem}
\begin{proof}
We specify both the original-source incidence and the arithmetic. Let
$W$ be the $108\times29$ rational weak-test matrix with its three constant
shift columns followed by $\nabla_\delta(e_j-e_{26})$, $0\le j<26$.
In these coordinates \eqref{v22:test} has the six nonzero entries
\[
 (w_*)_4=-1,\quad(w_*)_5=1,\quad(w_*)_7=-1,\quad
 (w_*)_8=1,\quad(w_*)_{10}=-1,\quad(w_*)_{11}=1.
\]
Use zero-based indices throughout. The rationalized V20 matrix
$\widehat K$ satisfies $\widehat K Ww_*=0$ on all 108 rows, checked over
$\mathbb Q$. A pure shift has the same $\sqrt3$ factor in every component
under $Q_3$, so \eqref{v20:rational-system} implies the original physical
null identity, not only a null identity for a changed pairing. Direct
rational differentiation of $\varphi_*$ gives $\|\nu_*\|_h^2=10$.

For the independent scalar calculation write $p=\sqrt3p_R$ and
$\ell_{21}=(n,\sqrt3\beta_R)$. Define rational fields
\[
 v_U=2p_R-(\tr p_R)I+\operatorname{SymGrad}_\delta\beta_R,
 \qquad v_p=\mathcal F(U)+\nabla_\delta^2n-I\Delta_\delta n.
\]
Then the actual first field velocity is $V=(\sqrt3v_U,v_p)$.
Every row of $LV=0$ is checked over $\mathbb Q$ after this exact grading.
The weak quadratic drift has field types $U^2,p^2$; the mixed first-bracket
block is linear in $U$ and the same-type shift block is linear in $p$.
Consequently
\begin{equation}
 \begin{split}
 \frac{27\Delta_*}{\sqrt3}
 ={}&Dd_{2,W}(U,p_R)[(v_U,v_p),\nu_*]\\
   &+\mathbb K_1(v_U,v_p)[\nu_*,(n,\beta_R)]
                  \quad\hbox{in spatial-sum convention}.
 \end{split}
 \label{v22:rational-delta}
\end{equation}
The factor $27$ converts the mean pairing to the sum. The right side is
evaluated by the full $P_{\mu,2}$ and $\Phi_3$ in
\eqref{v20:P2}--\eqref{v20:Phi3}, including the literal shifts and electric
and magnetic coefficients. For example
\[
 d_{2,W}(U,p_R)[\nu_*]
   =D\Phi_3(U,p_R)[G\nu_*]
                   +DP_{\nu_*,2}(U,p_R)[F_1(U,p_R)].
\]
Its directional derivative is obtained by exact quadratic polarization.
The derivative of $\Phi_3$ uses the exact cubic polarization subtracting
$\Phi_3(G\nu_*)$, not a finite-amplitude difference quotient.

The file \src{exact_tangency_defect.json} records the two rational summands
of \eqref{v22:rational-delta}, their complete rational sum, and its mean
conversion. Comparing its numerator and denominator with the two integer
bounds proves \eqref{v22:delta-interval}. Every denominator is nonzero and
a unit at $241$. There are two separate regressions. First the result equals
$w_*^T\mathsf C\mathsf J_f^{-1}V$ under the good-prime specialization of
V21's \emph{complete} weak field-control derivative. Its stored source row
is not used to compute the rational scalar value. Second, a fresh evaluation
of the unchanged scalar stationary coefficient functional reproduces the
derivative and bracket residues $181$ and $115$, respectively; their sum is
$55$ modulo $241$. All forward modular intermediates in that run are
integers below $2^{48}$. The rational calculation, not a modular zero or
floating threshold, establishes the nonzero value.
\end{proof}

Thus the new calculation has not used the unknown first active rate $v_A$,
the second-field solution $z_*$ from V21, or its absolute quartic-action
baseline. This independence is a consequence of the actual null test.
It is not permission to omit those quantities from other observations.

\section{The initial weak acceleration has an inverse-square amplitude pole}
\label{v22:acceleration}

On the fixed physical multiplier complement write the exact Hessian as
\begin{equation}
 M_{\mathrm{comp}}=D_a\Gamma_aD_a,\qquad
 D_a=\diag(aI_A,a^2I_W),\qquad
 \Gamma_0=\begin{pmatrix}A_0&J_0\\J_0^*&C_0\end{pmatrix},
 \qquad D_0=C_0-J_0^*A_0^{-1}J_0.
 \label{v22:Gamma}
\end{equation}
The off-diagonal block $J_0$ is distinct from the flat initial constraint
map $L$ and the field-dependent stationary source $J$. Both $A_0$ and the
signed $D_0$ are invertible at the inherited seed.
Neither is replaced by an unsigned weak matrix. Denote coordinate
restrictions by $R_A,R_W$ and synthesis by $I_A,I_W$.

\begin{theorem}[Loss of tangency of the exact weak-velocity zero set]
\label{v22:pole}
For every analytic preparation in \eqref{v22:base} with
$\ell=\ell_{21}$, define
\[
 t_W=\mathcal C_\ell V+\mathcal K_W I_Av_A.
\]
Its own exact continuation satisfies
\begin{equation}
 \boxed{R_W\lambda_{tt}(a,0)
       =-a^{-2}D_0^{-1}t_W+O_h(a^{-1}),\qquad
       w_*^Tt_W=\Delta_*\ne0.}
 \label{v22:weak-accel}
\end{equation}
In particular there are $a_0,c_h>0$ for which
\begin{equation}
 \boxed{\|\mathbb P_{\rm cf}\partial_t^2\lambda^\beta(a,0)\|_h
                    \ge c_h a^{-2}\quad(0<a<a_0).}
 \label{v22:accel-lower}
\end{equation}
The leading nonzero test is independent of the higher preparation. The
threshold and remainders can depend on the chosen analytic curve.
\end{theorem}
\begin{proof}
Differentiate the original consistency equation along its own field:
\begin{equation}
 M\lambda_{tt}=-\dot{\mathsf B}-\dot M\lambda_t,
 \qquad
 \dot M=D_XM[F]+D_\lambda M[\lambda_t].
 \label{v22:differentiated}
\end{equation}
The normalization keeps $\lambda_t,\lambda_{tt}$ in the fixed mean-lapse-zero
complement. Restricting to it is therefore exact. On the initial family,
$D_XM=O_h(a)$, $D_\lambda M=O_h(a^2)$, $F=O_h(a)$ and
$\lambda_t=O_h(a)$; hence $\dot M\lambda_t=O_h(a^3)$.
By \cref{v22:weak-Taylor} and \eqref{v22:B1}, the other term has active size
$O_h(a^2)$ and weak part $a^2t_W+O_h(a^3)$.

The blocks of the restricted matrix are
$a^2(A_0+O(a))$, $a^3(J_0+O(a))$, and $a^4(C_0+O(a))$.
Eliminate its active equation without changing either target. The active-target
feedback to the weak row has an additional factor $a$, and is $O_h(a^3)$.
The weak Schur matrix is $a^4(D_0+O(a))$. Analytic inversion proves the first
formula of \eqref{v22:weak-accel}. The Poisson-null identity removes the
$v_A$ term, giving its second formula by \cref{v22:nonzero}.
Thus $t_W$ and $D_0^{-1}t_W$ cannot vanish. Fixed-dimensional equivalence
of the chosen coordinates with the physical weak norm, followed by the
reverse triangle inequality, proves \eqref{v22:accel-lower}. This is a weak
acceleration statement; an active response coupled to it may have its own
lower amplitude order and is not discarded.
\end{proof}

The zero initial weak velocity in V21 is unchanged. Its derivative is now
known to be nonzero, so the selected zero set is not invariant under the
unchanged flow near these initial points. This does not prove divergence
of a coordinate-invariant curvature: the initial Riemann tensor uses
$\lambda_t$ but not $\lambda_{tt}$, as proved in V21. Nor does it exclude a
small unscaled rate on an amplitude-dependent interval.

\section{An analytic displacement scaling of the actual vector field}
\label{v22:layer-field}

The last distinction can be made quantitative. Fix one analytic preparation
$(X_a,\lambda_a)$ of V21 and keep its actual original physical time $t$.
For $a>0$ introduce comparison variables
\begin{equation}
 t=a^3s,\qquad X(t)=X_a+a^4\xi(s),\qquad
                   \lambda(t)=\lambda_a+a^4\eta(s).
 \label{v22:scales}
\end{equation}
This is a description of original histories at the displayed physical times,
not a different normalization of the physical clock. The field $\eta$ has
mean-zero lapse and splits into its fixed active and weak coordinates.

\begin{lemma}[An exact analytic extension of the scaled system]
\label{v22:extension}
On every bounded set of $(\xi,\eta)$, the original normalized vector field
in \eqref{v22:scales} extends analytically to $a=0$. Its limit is
\begin{align}
 \xi'&=V,&\eta_A'&=v_A,\notag\\
 \eta_W'&=-D_0^{-1}\left\{
    \mathcal C_\ell\xi+\mathcal K_W(I_A\eta_A+I_W\eta_W)
                      -J_0^*A_0^{-1}R_A\mathsf B_1\xi\right\}.
 \label{v22:limit-system}
\end{align}
The convergence of these vector fields is $O_h(a)$ in each fixed finite
$C^k$ norm on a compact set. No full physical history is needed to define
this extension, and no higher action term is removed from the finite system.
\end{lemma}
\begin{proof}
Let $\delta X=a^4\xi$, $\delta\lambda=a^4\eta$ and use analytic Taylor
expansion at the actual initial record. The canonical field obeys
$F(X_a+\delta X,\lambda_a+\delta\lambda)/a=V+O_h(a)$ uniformly with fixed
differentials. The weak drift derivative in \eqref{v22:DB} gives
\[
 \delta\mathsf B_W
   =a^5(\mathcal C_\ell\xi+\mathcal K_W\eta)+O_h(a^6),
 \qquad
 \delta\mathsf B_A=a^4R_A\mathsf B_1\xi+O_h(a^5).
\]
All displayed remainders are analytic with the stated powers, on bounded
sets. In the second formula a multiplier variation starts at order $a^5$.

For completeness, the graded Hessian has the estimate required for inversion.
The actual source $J_A$ is of joint order one near $(0,e)$; the weak source
$J_W$ is of joint order at least two. The latter follows from the universal
first-field cancellation and the flat spatial-source identity throughout
the multiplier neighborhood, as in \cref{v20:graded-source}. Consequently
$\delta J_A=O_h(a^4)$ and $\delta J_W=O_h(a^5)$ on this scaled tube.
The stationary tangent Hessian is analytic and invertible there. In
$M=J^*\mathsf K^{-1}J$, its three block changes are therefore of orders
$a^5,a^6,a^7$, respectively. Equivalently
\begin{equation}
 \Gamma(a,\xi,\eta)-\Gamma(a,0,0)=O_h(a^3).
 \label{v22:Gamma-change}
\end{equation}
The existing nonzero leading signed determinant makes this inverse analytic
uniformly on each fixed scaled tube for all sufficiently small $a$.

At the base the rate is $(a v_A+O(a^2),0)$. Also
$D_a^{-1}\mathsf B(X_a,\lambda_a)=O_h(a^2)$, since
$\mathsf B=-M\lambda_t$ there. Multiplying the normalized drift difference
by the graded inverse, the inverse change in \eqref{v22:Gamma-change}
contributes only $O_h(a^3)$ to the unscaled weak rate. The leading change
is thus
\[
 \begin{split}
 \delta r_W&=-a\,R_W\Gamma_0^{-1}
       \binom{R_A\mathsf B_1\xi}{\mathcal C_\ell\xi+\mathcal K_W\eta}
                                                     +O_h(a^2),\\
 \delta r_A&=O_h(a^2).
 \end{split}
\]
Exact block inversion yields the weak expression in
\eqref{v22:limit-system}. Since $d\lambda/dt=a\eta'$ and
$dX/dt=a\xi'$, these are precisely the scaled equations. The apparent
negative powers have removable singularities by the analytic divisibilities
just proved. Differentiating the finite Taylor formulas proves the $C^k$
claim. This argument retains the full original vector field before taking
its comparison limit.
\end{proof}

The scalar linear-drift term in \eqref{v22:limit-system} has not been omitted
from the off-constraint extension. Along the limiting trajectory from zero,
$\xi(s)=sV$ and $\mathsf B_1V=0$, so it vanishes by the original flat
constraint identities. This is different from deleting it for arbitrary
scaled field inputs.

\section{A proved original-action initial layer and its nonzero rate observation}
\label{v22:layer}

\begin{theorem}[Uniform existence and convergence on the rescaled interval]
\label{v22:layer-theorem}
For every fixed $S<\infty$ and each analytic initial family under consideration,
there is $a_S>0$ such that the original normalized exact histories exist on
$0\le t\le a^3S$ for $0<a<a_S$. They preserve all original constraints.
In the variables \eqref{v22:scales} they converge, uniformly together with
their first two $s$ derivatives, to the solution of
\eqref{v22:limit-system} with initial value zero. In particular
\begin{equation}
 \xi_0(s)=sV,\qquad \eta_{A,0}(s)=s v_A,\qquad
 \eta_{W,0}'=-D_0^{-1}\{s t_W+K_{WW}\eta_{W,0}\},
 \label{v22:triangular-limit}
\end{equation}
where $K_{WW}=\mathcal K_W I_W$, and $\eta_{W,0}(0)=0$.
The full original displacements on this interval are $O_{h,S}(a^4)$.
\end{theorem}
\begin{proof}
The limiting system is finite-dimensional affine linear, so its zero-initial
solution exists and is bounded, with its derivatives, on every fixed $[0,S]$.
Choose a compact tube containing that trajectory in its interior. On this
tube \cref{v22:extension} gives common Lipschitz constants and $O_h(a)$
vector-field differences for small $a$. The integral equations and Gronwall's
inequality bound the difference of their solutions by $C_{h,S}a$ while both
stay in the tube. Taking $a$ still smaller prevents exit through its boundary.
Local existence and continuation in the tube prove existence on all of
$[0,S]$. The $C^1$ vector-field comparison and differentiation of the equations
give convergence of the first two derivatives. This argument neither assumes
that the unscaled singular vector field has a common neighborhood nor uses
a physical-time existence interval not yet established.

Undo \eqref{v22:scales}. Its states remain within $O(a^4)$ of the original
initial state. The analytic stationary chart, positive lapse and spatial
metric, and graded rank persist on that tube, by the same fixed-cutoff
margins used in \cref{v22:extension}. Thus the constructed solution is the
original normalized continuation, not an equation obtained by re-solving
initial preparation. Its exact identity $\dot P=\mathsf B+M\lambda_t=0$
and its initial $P=0$ preserve every constraint. Substitution of
$\xi_0=sV$ and $\eta_{A,0}=sv_A$ gives \eqref{v22:triangular-limit}.
\end{proof}

\begin{corollary}[An explicitly forced weak-rate observation on actual histories]
\label{v22:linear-output}
Uniformly for $0\le s\le S$,
\begin{equation}
 \boxed{
 a^{-1}w_*^TD_0 R_W\lambda_t(a,a^3s)
                     \longrightarrow -\Delta_* s.}
 \label{v22:actual-profile}
\end{equation}
For every fixed $s>0$ in such an interval,
$\|\mathbb P_{\rm cf}\lambda_t^\beta(a,a^3s)\|_h\ge c_{h,s}a$
for all sufficiently small $a$. Initially this quantity is exactly zero.
The scalar in \eqref{v22:actual-profile} uses the fixed leading weak Schur
form as a dual test; it is a diagnostic of the original rate, not a new
curvature observable or an altered equation.
\end{corollary}
\begin{proof}
The rate relation in \eqref{v22:scales} gives
$a^{-1}R_W\lambda_t=\eta_W'$. Left-multiply
\eqref{v22:triangular-limit} by $w_*^TD_0$.
The exact physical Poisson-null identity gives
$w_*^TK_{WW}=0$ and $w_*^T\mathcal K_WI_A=0$.
Also $w_*^Tt_W=\Delta_*$. Therefore
\[
                 w_*^TD_0\eta_{W,0}'(s)=-s\Delta_*
\]
for the entire limiting interval, not just to first order in $s$.
The convergence theorem proves \eqref{v22:actual-profile}. At fixed positive
$s$ its limit is nonzero. Cauchy--Schwarz in the fixed weak coefficient norm,
and equivalence to the original physical norm, give the lower bound. No
unknown value of $v_A$ occurs in this scalar conclusion.
\end{proof}

The estimate at $s=0$ differentiated once also recovers the inverse-square
acceleration pole. The existence statement implies a lower lifespan of
order $a^3$ at this cutoff. It does not imply that $a^3$ is the maximal
lifespan, or that the histories cannot extend much farther. For a fixed
$a$, the allowed upper $S$ is not asserted arbitrary independently of $a$.

\section{Both curvature observations stay small throughout the proved layer}
\label{v22:geometry}

\begin{theorem}[A dynamical extension of the initial curvature bounds]
\label{v22:curvature-layer}
On the original physical interval of \cref{v22:layer-theorem}, for every
fixed spatial Sobolev order at $h=1/3$,
\begin{equation}
 \begin{gathered}
 \sup_{0\le t\le a^3S}
 (\|X\|+\|X_t\|+\|X_{tt}\|+\|\lambda-e\|+\|\lambda_t\|)
                                                  \le C_{h,S}a,\\
 \boxed{\sup_{0\le t\le a^3S}
  \{\|\operatorname{Riem}(g_h)\|
           +\|E_h^{\rm link}\|+\|F_h^{{\rm sp},{\rm link}}\|\}
                                                  \le C_{h,S}a.}
 \end{gathered}
 \label{v22:layer-curvature}
\end{equation}
The initial weak-velocity condition is not propagated, but its departure
is compatible with these small unscaled observations. No equality of the
first curvature coefficients divided by $a$ is asserted.
\end{theorem}
\begin{proof}
The base record is $O_h(a)$ from flat values, its displacement is $O_{h,S}(a^4)$,
and both scaled velocities are bounded by the preceding theorem. Thus
$X,\lambda-e,X_t,\lambda_t$ have the displayed bounds. On the fixed
stationary neighborhood the analytic derivatives $D_XF,D_\lambda F$ are
bounded, giving
$X_{tt}=D_XF X_t+D_\lambda F\lambda_t=O_{h,S}(a)$ throughout the interval.
All needed spatial derivatives obey the same estimate at this fixed grid.

For the original interpolated ADM metric, use exactly the coordinates and
reconstruction of V10 and V21. Its curvature is a smooth function on a
positive small chart of $\gamma,N,\beta$, their spatial derivatives,
$\gamma_t,\gamma_{tt},N_t,\beta_t$, and the required first spatial derivatives
of first time jets. Its independent Riemann blocks do not involve
$N_{tt}$ or $\beta_{tt}$, as the coordinate identities in
\cref{v21:initial-geometry} show at any time, not only at a zero-shift cut.
All differences from the flat values are $O_{h,S}(a)$, proving the metric
bound. This argument does not assert a bounded time derivative of curvature.

The original stationary connection, including its temporal auxiliary slot,
is $O_{h,S}(a)$. Its spatial connection velocity is obtained by the actual
chain rule $z_t=D_Xz X_t+D_\lambda z\lambda_t=O_{h,S}(a)$.
The original analytic ordered temporal and spatial link-curvature formulas
therefore give the same bound. At fixed compatible frames their difference
is also $O_{h,S}(a)$. Neither a transported surrogate metric nor a positive
replacement of the signed connection response has been used.
\end{proof}

This theorem is a positive original-flow result on a shrinking interval.
It cannot be turned into a fixed physical interval by calling $s$ physical
time. Similarly a nonzero limit in \eqref{v22:actual-profile} is at the
$a^{-1}$ rate normalization; it does not negate the unscaled $O(a)$ bounds.
These distinctions are the useful lesson from the source-normalized
nonconcentration calculation in the supplied Yang--Mills notes.

\section{The rate defect is visible in the original metric curvature}
\label{v22:observed-layer}

The first scalar rate test could in principle have detected only a spatially
constant shift. Such a conclusion would not establish a geometric response:
\cref{v10:symbol} shows that the difference-of-derivatives observation
annihilates precisely constant vectors. The following additional projection
is a change of analytical test, not a deletion of those physical modes.

Let $I_c:\R^3\to\mathcal W$ insert the three constant-shift coordinates and
let $D_{cc}=I_c^*D_0I_c$. This actual weak Schur compression is invertible
at the inherited seed: \cref{v12:finite-rank} certifies determinant residue
$133$ modulo $241$, and a common mean-versus-sum conversion does not change
its nonvanishing. Define
\begin{equation}
 \boxed{\overline w=w_*-I_cD_{cc}^{-1}I_c^*D_0w_* .}
 \label{v22:corrected-test}
\end{equation}
The original weak coordinate synthesis gives $\overline\nu=I_W\overline w$.

\begin{lemma}[A Poisson-null test that detects no constant rate]
\label{v22:geometric-test}
One has
\begin{equation}
 \overline\nu\in\mathcal N_W,\qquad
 \overline w^*D_0I_c=0,\qquad
 \Delta_{\overline\nu}(Y,\ell_{21})=\Delta_*.
 \label{v22:corrected-identities}
\end{equation}
Consequently there is a bounded real linear functional $\Lambda_h$ on the
finite symmetric-tensor observation space for which
\begin{equation}
 \Lambda_h\bigl(\mathfrak E_h((I_W z)^\beta)\bigr)
                      =\overline w^*D_0z\quad(z\in\mathcal W).
 \label{v22:Lambda-temp}
\end{equation}
In this formula $\mathfrak E_h$ acts on the shift component of $I_Wz$;
more explicitly its argument is $(I_Wz)^\beta$.
\end{lemma}
\begin{proof}
Each constant shift belongs to $\mathcal N_W$, by the original translation
identity and $LY=0$. Its tangency covector is zero: the constant quadratic
drift is identically zero, and
$\mathbb K_1(V)[c,\ell]=\ip{\delta_c\ell}{LV}_h=0$.
Subtracting the three constant tests therefore preserves the Poisson-null
identity and the value $\Delta_*$. Symmetry of $D_0$ and invertibility of
$D_{cc}$ give the middle identity in \eqref{v22:corrected-identities}.

On weak shifts the kernel of $\mathfrak E_h$ is exactly the constants,
by the exact all-frequency formula of V10. The linear functional
$z\mapsto\overline w^*D_0z$ annihilates that same kernel, so it defines a
well-defined linear functional on the image of $\mathfrak E_h I_W$.
Finite dimension makes it bounded. Extend it to the full observation space
by its positive orthogonal projection onto that image. This constructs
$\Lambda_h$ without removing a mode from the original curvature. Its norm
is finite at this cutoff; it is not asserted uniformly bounded as $h\to0$.
\end{proof}

\begin{theorem}[A nonzero normalized metric response and its initial derivative]
\label{v22:metric-profile}
Let $\mathcal R_a(s)$ denote the symmetric coordinate block
$(R_{0i0j}(g_h)(a,a^3s))_{i,j}$ of the same original interpolated metric.
Uniformly on every fixed $s$ interval of the layer,
\begin{equation}
 \boxed{\frac{\mathcal R_a(s)-\mathcal R_a(0)}a
     \longrightarrow
        \mathcal I_h\mathfrak E_h\bigl((I_W\eta_{W,0}'(s))^\beta\bigr).}
 \label{v22:metric-change}
\end{equation}
The convergence holds also after one $s$ derivative. Thus, with the functional extended to continuous tensor fields by
$\Lambda_h\mathcal I_h^*$ (Fourier restriction followed by the fixed finite
test),
\begin{equation}
 \boxed{\Lambda_h\!
       \left(\frac{\mathcal R_a(s)-\mathcal R_a(0)}a\right)
          \longrightarrow-\Delta_*s,\qquad
       \|\partial_t(R_{0i0j})(a,0)\|\ge c_h a^{-2}.}
 \label{v22:metric-slope}
\end{equation}
This nonzero observation is independent of the unknown first active rate.
The initial rate and curvature bounds in V21 remain valid.
\end{theorem}
\begin{proof}
The state displacements in the layer are $a^4$ times bounded smooth
functions of $s$, whereas the original initial fields are $aY+O(a^2)$ and
$\lambda=e+a\ell+O(a^2)$. Hence their first amplitude coefficients are
independent of $s$. Also $X_t/a\to V$ and
\[
 X_{tt}/a\longrightarrow F_1V+G(I_Av_A+I_W\eta_{W,0}'(s)).
\]
This follows directly by differentiating the original canonical equation,
using its analytic field derivatives at $(0,e)$ and the rate limit already
proved. In the spatial-configuration component the only $s$-dependent term
is $\operatorname{SymGrad}_\delta(I_W\eta_{W,0}')^\beta$.

In the same metric coordinates, $\partial_tg_{0i}/a$ has the corresponding
$s$-dependent coefficient $(I_W\eta_{W,0}')^\beta_i$. All purely spatial
metric terms at first amplitude order, and the terms involving $X_t/a$,
are independent of $s$. Products of first metric derivatives are $O(a^2)$
on the layer. Subtract the coordinate Riemann identities at $s$ and zero:
\[
 R_{0i0j}=\tfrac12(g_{0j,it}+g_{i0,tj}-g_{00,ij}-g_{ij,tt})
                                      +\hbox{Christoffel products}.
\]
The remaining first amplitude term is exactly one half of the symmetric
gradient with the geometric derivative minus the phase derivative, giving
\eqref{v22:metric-change}. This calculation keeps the initial nonconstant
lapse and shift in the common, subtracted terms; it does not set them to zero.

For the differentiated assertion, express the Riemann block as the analytic
finite function of $X,\lambda,F$ and the actual rate used in
\cref{v22:curvature-layer}. After substituting the scaled variables, its
division by $a$ extends analytically on the bounded layer tube. The $C^1$
convergence of the scaled flow therefore justifies differentiating
\eqref{v22:metric-change} in $s$. Apply \cref{v22:geometric-test} and the
Poisson-null calculation from \cref{v22:linear-output}. It gives the first
limit of \eqref{v22:metric-slope}. Since
$\partial_s(\mathcal R_a/a)=a^2\partial_t\mathcal R_a$, its derivative
at zero has nonzero scalar limit $-\Delta_*$. Boundedness of $\Lambda_h$
and the reverse triangle inequality prove the last bound. The derivative
is taken in the declared physical time, not the rescaled comparison time.
\end{proof}

\begin{corollary}[No amplitude-linear curvature bound on a fixed physical interval]
\label{v22:no-linear-fixed-time}
For these analytic initial families there do not exist $T,C,a_0>0$ such that
all histories with $0<a<a_0$ both exist on $[0,T]$ and satisfy
\[
                 \sup_{0\le t\le T}\|R_{0i0j}(g_h)(a,t)\|\le Ca.
\]
This is not a nonexistence theorem for continuation on $[0,T]$, and does not
exclude a weaker vanishing bound, a different initial family, or a joint
amplitude--cutoff limit.
\end{corollary}
\begin{proof}
Such a bound would make the absolute value of the scalar in
\eqref{v22:metric-slope} at most $2\|\Lambda_h\|C$ whenever $a^3s\le T$.
Choose a fixed finite $s>2\|\Lambda_h\|C/|\Delta_*|$. The layer theorem
applies to this particular $s$ for all sufficiently small $a$, and then
$a^3s\le T$. Passing to the limit contradicts
$|\Delta_*|s>2\|\Lambda_h\|C$. The argument takes $s$ fixed first; it does
not assume estimates on an unbounded rescaled interval.
\end{proof}

\section{The next upstream gate and the verification record}
\label{v22:status}

The first-tilt freedom left in V21 is not automatically a way to propagate
its chosen initial condition. With $Y$ fixed, let
$\ell(\eta)=\ell_{21}+Q_3\eta$, where
$\eta\in\mathcal N_{21}$ is one of the eleven normalized null-tilt
parameters. Define for any $\nu\in\mathcal N_W$
\begin{equation}
 \mathcal T_\nu(\eta)=
 Dd_{2,W}(Y)[F_1Y+G\ell(\eta),\nu]
 +\mathbb K_1(F_1Y+G\ell(\eta))[\nu,\ell(\eta)].
 \label{v22:tilt-gate}
\end{equation}
This is an explicit polynomial of degree at most two in the eleven parameters.
It uses only the original already specified quadratic and cubic envelopes.
Vanishing of every such component is a necessary next first-jet tangency
condition for small weak acceleration. It is not inferred from initial
weak-rate preparation.

\begin{proposition}[Higher preparation cannot repair this defect near the current tilt]
\label{v22:local-obstruction}
There is a neighborhood of $0$ in $\mathcal N_{21}$ on which
$\mathcal T_{\nu_*}(\eta)\ne0$. For every analytic family of V21 in that
neighborhood, the same tangency obstruction and a nonzero scalar initial-layer
profile hold, with its own corresponding nonzero coefficient. No change of
higher field or multiplier coefficients at a fixed first tilt eliminates
that leading observation.
\end{proposition}
\begin{proof}
Equation \eqref{v22:tilt-gate} is polynomial, hence continuous, and its value
at zero is the nonzero number in \cref{v22:nonzero}. It stays away from zero
on a smaller neighborhood. V21's initial controllability and leading signed
regularity persist there. The Taylor, null-test, and scaled-flow arguments
apply to each such family. Their invariant coefficient contains no higher
preparation parameter by \cref{v22:weak-Taylor}.
\end{proof}

We have not solved the polynomial gate elsewhere on the eleven-parameter
space, or proved that all its components are independent. Its vanishing,
if achieved, would still be a necessary next condition, not an invariant
manifold theorem. An alternative is to retain and estimate the actual rapid
weak evolution rather than require its exact zero. Either route must keep
the full signed higher source, metric observations, and physical-time scale.
The present layer supplies no cofinal tilted family, cutoff-independent
constants, differentiated-curvature estimate, or common positive physical
lifespan. The raw-action Einstein continuum bridge remains unclosed.

\paragraph{Executed checks and limitations.}
The new verifier constructs the weak test and checks its full Poisson-null
identity over $\mathbb Q$, its physical norm and the flat-velocity constraints.
A fresh rational value-polarization calculation evaluates
\eqref{v22:rational-delta} from the original $P_2,\Phi_3$ formulas and proves
its strict interval. A separate evaluation of the retained full scalar
stationary coefficient functional checks the same two summands at the good
prime, independently of the stored weak-control derivative. The V21 control
matrix gives a third incidence check. No unused reverse-AD audit is counted
as an executed derivative: the new good-prime regression uses forward
polynomial values only.

Exact finite block checks retain the signed weak Schur matrix and verify the
amplitude orders in the displaced rate. Original-coefficient recurrences test
the linear scalar observation in the limiting layer without assigning a value
to the unknown active rate. These are algebraic regressions, not nonlinear
spacetime integrations or numerical constructions of V21's implicit initial
root. The layer convergence and curvature bounds are the written analytic
proofs based on the named inherited stationary and preparation results.
All twenty-one preceding marked bodies are preserved byte-for-byte.

\begin{thebibliography}{9}
\raggedright
\setcounter{enumiv}{18}
\bibitem{v22:YM86}
\emph{Renewal Yang--Mills Mass-Gap Research Notes V86}, supplied cumulative
source, V85--V86, especially \src{lem:v86-three-events},
\src{thm:v86-joint}, and \src{lem:v86-doubling}. Used for the audited
source-transport and scale distinction, not as a gravitational estimate.
\bibitem{v22:DBGP}
C. Di Bartolo, R. Gambini and J. Pullin,
\emph{Consistent and mimetic discretizations in general relativity},
Journal of Mathematical Physics \textbf{46} (2005), 032501;
\url{https://arxiv.org/abs/gr-qc/0404052}.
Primary-source context for the difference between solving discrete constraints
with selected multipliers and preserving an additional condition. The present
coefficient test and original-flow layer are derived above, not imported.
\end{thebibliography}
% END IMMUTABLE EXACT-ACTION BRIDGE V22 BODY


% Future iterations append here.
% BEGIN IMMUTABLE EXACT-ACTION BRIDGE V23 BODY
\clearpage
\renewcommand{\thesection}{\arabic{section}}
\renewcommand{\theHsection}{bridgeV23.\arabic{section}}
\setcounter{section}{187}
\section{The full leading-tilt gate before changing the canonical first jet}
\label{v23:context}

V22 exhibits one weak Poisson-null test whose tangency coefficient is nonzero at
the V21 leading tilt.  The remaining eleven normalized leading-tilt parameters
were not tested there.  It would be circular to move to another higher
preparation before deciding whether those parameters can remove the defect:
V22 already proves that its selected tangency observable is independent of all
higher field and multiplier preparation coefficients.  This body therefore
computes the tangency polynomial on the \emph{entire} admissible leading-tilt
affine space at the same canonical first jet.

The result is a no-go for that route.  A particularly simple weak null test,
generated by
\[
 \phi_\diamond=\frac23\{\cos(2\pi x_1)-\cos(2\pi x_2)\},
 \qquad \nu_\diamond=(0,\nabla_\delta\phi_\diamond),
\]
has a tangency coefficient which is exactly constant on all eleven admissible
leading-tilt directions and is nonzero.  Thus no change of leading lapse--shift
values inside the V21 lower-incidence solution set can make the prepared
zero-weak-rate condition tangent to the original action flow.  Since the
canonical first coefficient has remained fixed throughout V20--V22, the next
upstream object is the canonical first jet itself.

We consequently identify an explicit five-dimensional family of complementary,
weak-neutral, mean-compatible first jets which contains the current seed and
admits exact original initial preparation.  No tangency root is claimed on this
larger family in the present body.  Its role is to specify the next finite
source problem without silently changing the action or recycling the now
excluded eleven-tilt route.

\paragraph{Inherited inputs and nonduplication.}
The first Poisson bracket, its exact 18-dimensional nullspace, and the complete
quadratic drift are \cref{v20:K1,v20:range-certificate}.  The cubic constant-row
map and its three independent controls are \cref{v21:constant-law,v21:joint-tilt};
the resulting eleven-dimensional leading-tilt space is
\cref{v21:eleven}.  V22 supplies the preparation-independent tangency formula
and the implication from a nonzero weak-null tangency coefficient to the
inverse-square initial weak acceleration.  None of those statements is
reproved as new work.  The new items are an exact basis of the complete
leading-tilt solution space, the simple null test above, the characteristic-zero
constancy certificate for its full tangency polynomial, the resulting global
leading-tilt no-go, and the upstream five-parameter first-jet chart.  Every
statement about exact evolution remains fixed-cutoff unless explicitly stated
otherwise.

\section{The exact eleven-dimensional leading-tilt space}
\label{v23:tilt-space}

Work at the complementary three-site first jet \(Y\) of \eqref{v20:seed}.
Let \(N_*\) denote the exact \(108\times18\) rational null matrix of the first
Poisson bracket from V20, in the scaled multiplier coordinates used there.  Let
\(\mathsf m\) be the four component-mean map and let \(R\) be the three-row
cubic constant-shift tilt map of V21.  Define
\begin{equation}
 \mathsf C_{21}=\begin{pmatrix}\mathsf mN_*\\ RN_*\end{pmatrix}:
                 \R^{18}\longrightarrow\R^7 .
 \label{v23:C21}
\end{equation}
The affine set of leading tilts satisfying the complete quadratic cancellation,
the four normalizations, and all three cubic constant-shift equations is
\begin{equation}
 \mathfrak L_{21}=\ell_{21}+\mathcal N_{21},\qquad
 \mathcal N_{21}=\Ran N_*\cap\ker\mathsf m\cap\ker R .
 \label{v23:L21}
\end{equation}

\begin{lemma}[An exact basis of every admissible leading-tilt direction]
\label{v23:E11}
The matrix \(\mathsf C_{21}\) has rank seven and
\[
                  \dim\mathcal N_{21}=11.
\]
If \(C_0\) is any rational basis matrix for \(\ker\mathsf C_{21}\) and
\(E=N_*C_0=(E_1,\ldots,E_{11})\), then
\begin{equation}
 \boxed{\mathbb K_1(Y)E_j=0,\qquad
        \mathsf mE_j=0,\qquad RE_j=0\quad(1\le j\le11).}
 \label{v23:E-properties}
\end{equation}
For the supplied exact basis a seven-column minor of \(\mathsf C_{21}\), using
columns \(0,1,2,3,7,10,15\) of the V20 null coordinates, has determinant
\begin{equation}
 -\frac{859417883246151146653209}{268435456}
e0,
 \qquad \det\equiv126\pmod{241}.
 \label{v23:C-minor}
\end{equation}
\end{lemma}
\begin{proof}
V20 gives \(\rank N_*=18\) and \(\mathbb K_1(Y)N_*=0\).
The displayed minor proves \(\rank\mathsf C_{21}=7\) in characteristic zero,
so its kernel has dimension eleven.  Multiplication by the injective map
\(N_*\) preserves that dimension.  The three identities in
\eqref{v23:E-properties} are then the definitions of the three rows of
incidence in \eqref{v23:C21}, together with the original Poisson-null identity.
No singular-value threshold is used.
\end{proof}

\section{A simple exact weak-null test}
\label{v23:diamond-test}

Let scalar grid sites use the same zero-based lexicographic ordering as V21.
On the three-site torus the cosine values are rational.  Define
\begin{equation}
 \phi_\diamond(x)=\frac23\{\cos(2\pi x_1)-\cos(2\pi x_2)\},\qquad
 \nu_\diamond=(0,\nabla_\delta\phi_\diamond).
 \label{v23:diamond}
\end{equation}
It has no constant-shift component and has zero scalar mean.

\begin{lemma}[The diamond test belongs to the exact weak Poisson kernel]
\label{v23:diamond-null}
At the first jet \(Y\),
\begin{equation}
                      \boxed{\mathbb K_1(Y)\nu_\diamond=0.}
 \label{v23:diamond-null-eq}
\end{equation}
In the 11-column weak-null basis used in V22, whose free weak coordinates are
\[
 (0,1,2,11,18,20,23,25,26,27,28),
\]
the coordinate vector of \(\nu_\diamond\) is
\begin{equation}
                 (0,0,0,1,0,0,-1,0,0,0,0)^T .
 \label{v23:diamond-coordinates}
\end{equation}
\end{lemma}
\begin{proof}
Represent a mean-zero weak scalar by the V21 basis
\(q_j=e_j-e_{26}\), \(0\le j<26\), and apply the original phase gradient.
The scalar values of \eqref{v23:diamond} are rational and their first 26
entries are precisely its coefficients in this basis.  Exact multiplication by
the V20 first Poisson matrix gives zero in all 108 rows.  Row reduction of the
same exact \(108\times29\) weak restriction gives
\eqref{v23:diamond-coordinates}.  Thus the asserted null identity is over
\(\mathbb Q\), not a modular or floating rank inference.
\end{proof}

\section{The tangency polynomial is constant on all eleven leading-tilt directions}
\label{v23:constant-gate}

For a weak Poisson-null test \(\nu\), V22's leading tangency functional can be
written entirely in terms of the original quadratic weak drift and first
Poisson bracket as
\begin{equation}
 \mathcal T_\nu(\ell)=
 Dd_{2,W}(Y)[F_1Y+G\ell,\nu]
 +\mathbb K_1(F_1Y+G\ell)[\nu,\ell].
 \label{v23:T}
\end{equation}
Here \(G\ell\) is the original flat multiplier Hamiltonian vector.
The first term is affine in \(\ell\), and the second is at most quadratic;
therefore \(\mathcal T_\nu\) is a polynomial of degree at most two in the
leading tilt.

Put \(\ell=\ell_{21}+\sum_jc_jE_j\).  For later reproducibility, the exact
coefficient formulas are
\begin{align}
 L_j={}&Dd_{2,W}(Y)[GE_j,\nu]
 +\mathbb K_1(F_1Y)[\nu,E_j]
 +\mathbb K_1(GE_j)[\nu,\ell_{21}]
 +\mathbb K_1(G\ell_{21})[\nu,E_j],\label{v23:Lcoef}\\
 Q_{jj}={}&\mathbb K_1(GE_j)[\nu,E_j],\label{v23:Qdiag}\\
 Q_{jk}={}&\mathbb K_1(GE_j)[\nu,E_k]
             +\mathbb K_1(GE_k)[\nu,E_j]\quad(j<k).\label{v23:Qcross}
\end{align}
No finite-amplitude difference quotient is involved in these formulas.

\begin{theorem}[A characteristic-zero global leading-tilt no-go]
\label{v23:global-no-go}
For the exact weak-null test \(\nu_\diamond\), every one of the 77 nonconstant
coefficients in \eqref{v23:Lcoef}--\eqref{v23:Qcross} vanishes:
\begin{equation}
 \boxed{L_j=0,\qquad Q_{jj}=0,\qquad Q_{jk}=0
        \quad(1\le j<k\le11).}
 \label{v23:all-zero}
\end{equation}
The remaining constant is nonzero.  Writing
\begin{equation}
 \Delta_\diamond:=\mathcal T_{\nu_\diamond}(\ell_{21}),
 \label{v23:Delta}
\end{equation}
one has
\begin{equation}
 \boxed{\Delta_\diamond\ne0,\qquad
        \Delta_\diamond\equiv148\pmod{241},\qquad
        \Delta_\diamond\simeq-7.908635358866945\times10^6}
 \label{v23:Delta-value}
\end{equation}
in the rational spatial-sum scaling used throughout V20--V22.  Hence
\begin{equation}
 \boxed{\mathcal T_{\nu_\diamond}(\ell)=\Delta_\diamond
        \quad\hbox{for every }\ell\in\mathfrak L_{21}.}
 \label{v23:constant-T}
\end{equation}
\end{theorem}
\begin{proof}
The coefficient formulas \eqref{v23:Lcoef}--\eqref{v23:Qcross} reduce the claim
to finitely many homogeneous polynomial contractions of the original
\(P_2\) and \(\Phi_3\) coefficients.  The supplied verifier evaluates them over
\(\mathbb Q\), using the literal three-site phase derivative, original shifted
BCH coefficients, both off-diagonal Frobenius occurrences and the exact V21
tilt.  It obtains zero for all 11 linear coefficients, all 11 diagonal
quadratic coefficients and all 55 cross coefficients.  These are exact
fraction equalities, not modular zeros.

The same independent scalar extraction gives \(\Delta_\diamond\) as one exact
rational number.  Its denominator is a unit modulo 241 and its numerator divided
by that denominator has residue 148.  Therefore it is nonzero in characteristic
zero.  Since \(\mathcal T_{\nu_\diamond}\) has degree at most two, the complete
coefficient list proves \eqref{v23:constant-T} on the entire affine space, not
only in a neighborhood of \(\ell_{21}\).
\end{proof}

The full exact fraction in \eqref{v23:Delta} is retained in the companion
certificate rather than printed across several manuscript lines.  Its decimal
value in \eqref{v23:Delta-value} is explanatory only; nonvanishing is certified
by the rational residue and the exact stored numerator and denominator.

\section{Consequences for prepared histories}
\label{v23:consequence}

The algebraic no-go concerns the leading tilt on the fixed canonical first jet.
To translate it into an exact-flow statement, one must still remain on a
regular initial branch for which the V21 higher preparation exists.  This is
a local requirement on the finite parameters, not an extra cancellation
condition.

\begin{corollary}[Higher preparation cannot repair the V22 tangency defect at this first jet]
\label{v23:no-higher-repair}
Fix any \(\ell\in\mathfrak L_{21}\) for which the original analytic initial
constraint preparation, the V21 zero-initial-weak-rate preparation and the
regular graded source chart exist.  Then the resulting original-action history
cannot make that zero-weak-rate condition tangent at the prepared cut.  More
precisely, the \(\nu_\diamond\) component of its differentiated weak target has
leading coefficient \(\Delta_\diamond\), independent of all higher initial
field and multiplier coefficients.  Consequently
\begin{equation}
 \boxed{\|\mathbb P_{\rm cf}\lambda_{tt}^{\beta}(a,0)\|_h
                         \ge c_h a^{-2}}
 \label{v23:acceleration}
\end{equation}
for sufficiently small nonzero amplitude, with \(c_h>0\) on the fixed
three-site branch.
\end{corollary}
\begin{proof}
V22 derives the differentiated weak consistency equation precisely on a
Poisson-null test and proves that its leading coefficient depends only on the
canonical first jet and leading multiplier tilt.  Higher preparation therefore
cannot alter \(\mathcal T_{\nu_\diamond}\).  Apply
\cref{v23:global-no-go}.  The leading weak Schur matrix is the same invertible
matrix fixed by the unchanged canonical first jet, so a nonzero tested target
produces the inverse-square rate lower bound exactly as in V22.
\end{proof}

Thus another optimization of the eleven V21 leading-tilt parameters is a proved
dead end.  This conclusion is stronger than the local statement in V22 and is
why the remaining attack moves upstream below.

\section{A five-dimensional upstream family of canonical first jets}
\label{v23:first-jet-family}

The complementary weak-neutral construction itself has unused first-jet
freedom.  We now record it explicitly so that the next tangency calculation
has a well-defined source family rather than an informal parameter count.
Retain the transverse tensors \(T_i,P_i\) of V12 and put
\begin{equation}
 \begin{split}
 U_\alpha&=\sum_{i=1}^3\alpha_iT_i\cos(2\pi x_i),\\
 Q_{h,\alpha}&=\frac{\omega_h}{2}
          \sum_{i=1}^3\alpha_iP_i\cos(2\pi x_i),\\
 B_c(b)&=\omega_h\diag(b_1,b_2,b_3),\qquad
 P_c(b)=B_c(b)-(\tr B_c(b))I,\\
 Y_{h}(\alpha,b)&=(U_\alpha,\,P_c(b)+Q_{h,\alpha}).
 \end{split}
 \label{v23:Y5}
\end{equation}
Let \(\Sigma_\alpha=\sum_i\alpha_i^2\).  The quadratic scalar-mean condition is
\begin{equation}
 \boxed{b_1b_2+b_1b_3+b_2b_3=\frac14\Sigma_\alpha.}
 \label{v23:mean-surface}
\end{equation}
Near the present seed, \(b_2+b_3\ne0\), so
\begin{equation}
             b_1(\alpha,b_2,b_3)
       =\frac{\Sigma_\alpha/4-b_2b_3}{b_2+b_3}.
 \label{v23:b1}
\end{equation}
The base point is
\begin{equation}
 \alpha_*=(1,9/8,1),\qquad
 b_*=(16175/4096,-8,-8).
 \label{v23:base-first}
\end{equation}

\begin{theorem}[A genuine five-parameter weak-neutral initial-data chart]
\label{v23:five-chart}
In a neighborhood of \eqref{v23:base-first}, the variables
\begin{equation}
                 (\alpha_1,\alpha_2,\alpha_3,b_2,b_3)
 \label{v23:five-parameters}
\end{equation}
define a real-analytic five-dimensional family of canonical first jets
\eqref{v23:Y5} with the following properties.
\begin{enumerate}[label=\textup{(\roman*)},leftmargin=2.8em]
\item Every linear original initial constraint vanishes.
\item The complete curl-free quadratic shift drift vanishes exactly at every
      odd cutoff:
      \[
                         d_{2,h}(Y_h(\alpha,b))=0.
      \]
\item All four quadratic mean constraints vanish at the uncorrected first-jet
      level when \eqref{v23:mean-surface} holds.
\item The original nonlinear range-and-mean initializer extends these first
      jets to analytic exactly constrained initial records for sufficiently
      small amplitude and parameters near the base point.
\end{enumerate}
No tangency equation or exact-action evolution estimate is asserted by this
theorem.
\end{theorem}
\begin{proof}
The configuration and oscillatory momentum are transverse and trace-free in
each axial mode; the homogeneous momentum is constant.  This proves the linear
constraint rows.  In the notation of \cref{v12:matching}, take
\(A_i=\alpha_iT_i,B_i=0,C_i=\omega_h\alpha_iP_i/2,D_i=0\).
Because \(\|T_i\|_F=\|P_i\|_F\) and \(T_i\perp P_i\), the two matching
conditions in \eqref{v12:matching-conditions} hold identically.  The constant
symmetric momentum does not enter them.  This proves (ii).

For the scalar mean, write \(B_c=\omega_h\diag(b)\).  The constant momentum
satisfies
\[
 \|P_c\|_F^2-\tfrac12(\tr P_c)^2
 =\omega_h^2\{\|b\|^2-(b_1+b_2+b_3)^2\}
 =-2\omega_h^2(b_1b_2+b_1b_3+b_2b_3).
\]
The matched oscillatory configuration and momentum contribute
\(\omega_h^2\Sigma_\alpha/2\) with the opposite sign in the original scalar
constraint convention.  Hence the scalar row is zero exactly when
\eqref{v23:mean-surface} holds.  The three constant-shift mean rows vanish by
mode orthogonality and transversality, proving (iii).

It remains to justify nonlinear preparation rather than merely its quadratic
starting point.  Varying \(b_1\) changes the dimensionless scalar mean by
\(2(b_2+b_3)\); at the base this derivative is \(-32\ne0\).  The same three
sine momentum balancing directions used in V12 have shift-mean derivatives
\(-\omega_h\alpha_i^2\), all nonzero near the base.  Thus the four-dimensional
mean balancing Jacobian remains nonsingular.  The nonconstant range inverse is
the unchanged original one.  The analytic range solve followed by the four
mean implicit equations therefore gives exact initial records, proving (iv).
Finally \(\partial_{b_1}\) of the left side of
\eqref{v23:mean-surface} is \(b_2+b_3=-16\) at the base, so the ordinary
implicit-function theorem gives the analytic parametrization \eqref{v23:b1}
and its five-dimensionality.
\end{proof}

The special V12 calibration \(b_2=b_3=-8\),
\(b_1=4-\Sigma_\alpha/64\), is one two-codimension slice of this family.  The
new chart does not assert that all five directions preserve the leading source
rank or its V21 control minors; those are open finite conditions to be retested
on whichever upstream candidate is selected.

\section{The next upstream finite incidence problem}
\label{v23:next}

The no-go in \cref{v23:global-no-go} changes the next calculation.  At the
fixed first jet, neither higher preparation nor any of the eleven admissible
leading-tilt directions can remove the V22 tangency source.  The appropriate
finite problem is therefore to vary \(Y\) inside \cref{v23:five-chart} and solve
all lower incidences and tangency conditions simultaneously.

For a first-jet parameter \(p\in\R^5\), let \(Y(p)\) be
\eqref{v23:Y5}.  On a branch where the first Poisson rank is constant, define
\begin{align}
 \mathfrak F_2(p,\ell)&=\widehat B_2(Y(p))+
                               \mathbb K_1(Y(p))\ell,\label{v23:F2}\\
 \mathfrak F_c(p,\ell)&=t(Y(p))+\mathcal R_{Y(p)}\ell,\label{v23:Fc}
\end{align}
and let \(\mathfrak F_T\) denote the complete weak Poisson-null tangency
covector obtained from \eqref{v23:T}, with the null bundle transported on that
constant-rank chart.  The upstream target is
\begin{equation}
 \boxed{\mathfrak F_2=0,\qquad \mathsf m\ell=0,\qquad
        \mathfrak F_c=0,\qquad \mathfrak F_T=0.}
 \label{v23:upstream-system}
\end{equation}
This is a finite algebraic/analytic incidence problem at each cutoff.  It must
be solved before repeating V21's higher weak preparation and V22's propagation
test on a new first jet.

The present body does not assert a root of \eqref{v23:upstream-system}.  In
particular, changing \(Y\) changes the Poisson null bundle, the admissible
leading-tilt affine space and the cubic mean map simultaneously.  Reusing the
old eleven-dimensional basis after changing the first jet would be another
unjustified substitution.  The next useful computation is the derivative of
the \emph{tilt-invariant tangency quotient} with respect to the five parameters
in \eqref{v23:five-parameters}, followed, if it has the required rank, by a
verified finite root and a fresh regular-source certificate.

\section{Verification, preservation, and scope}
\label{v23:status}

The exact certificate constructs \(\mathcal N_{21}\) from the supplied V20
null basis rather than a floating nullspace.  It verifies the rank-seven lower
incidence by \eqref{v23:C-minor}, constructs \(\nu_\diamond\) directly from
its rational scalar potential and checks all 108 Poisson-null equations.
The tangency verifier uses the original rational \(P_2\) and \(\Phi_3\)
coefficient functions.  It evaluates the coefficient formulas
\eqref{v23:Lcoef}--\eqref{v23:Qcross}; all 77 nonconstant coefficients are
exactly zero over \(\mathbb Q\), while the stored constant has residue 148
modulo 241 with a nonzero denominator.  No rank, zero, or nonzero statement in
\cref{v23:global-no-go} is determined by a floating tolerance.

A separate symbolic check verifies the five-parameter mean identity
\eqref{v23:mean-surface}, its base point and the nonzero balancing derivative.
This check does not constitute a tangency-root calculation on the enlarged
family.  No cofinal first-jet family, higher weak signed-rank theorem,
cutoff-uniform control norm, propagated nonlinear estimate, differentiated
curvature theorem or common physical lifespan is claimed.  The exact-action
Einstein bridge therefore remains open, but optimization of the old V21
leading-tilt family is now closed as a route.

All twenty-two preceding marked research bodies are preserved byte for byte.
The package records their hashes and the new exact certificates.  The finite
coefficient scripts verify the displayed algebra; they are not a proof-assistant
formalization or a nonlinear spacetime integration.
% END IMMUTABLE EXACT-ACTION BRIDGE V23 BODY


% BEGIN IMMUTABLE EXACT-ACTION BRIDGE V24 BODY
\clearpage
\renewcommand{\thesection}{\arabic{section}}
\renewcommand{\theHsection}{bridgeV24.\arabic{section}}
\setcounter{section}{195}

\section{Upstream first-jet incidence after the global tilt no-go}
\label{v24:context}

Version V23 proves that, at the fixed complementary three-site canonical first
jet, the V22 weak-null tangency defect is constant and nonzero on the entire
11-dimensional affine family of leading multiplier tilts satisfying the lower
quadratic, normalization and cubic-mean conditions.  That result closes a
pure leading-tilt repair at the old first jet.  The next possible repair is to
change the canonical first jet itself inside the five-parameter weak-neutral
chart of \cref{v23:five-chart}, while simultaneously re-solving the lower
incidences.

This body performs the first characteristic-zero calculation on that enlarged
problem.  It records one reproducibility correction to the V23 companion
calculation, constructs an exact lower-incidence tangent in the
$\alpha_2$ direction, and proves that the V22 tangency quotient changes
nontrivially along that tangent.  Thus the V23 obstruction is attached to the
chosen first jet rather than being invariant under all admissible canonical
first-jet changes.  We do not yet claim that this formal tangent integrates to
an actual nearby root of the full upstream system.

\paragraph{Inherited inputs and nonduplication.}
The first Poisson bracket, prepared quadratic target and rational scaling are
\cref{v20:K1,v20:rational-system}.  The normalized lower tilt and cubic mean
map are \cref{v21:joint-tilt}; the exact 11-dimensional fixed-first-jet null
space and the weak test $\nu_\diamond$ are
\cref{v23:E11,v23:diamond-null}.  The five-parameter first-jet chart is
\cref{v23:five-chart}.  No rank, control or tangency statement from those
bodies is repeated as new work.  All finite calculations below use the same
three-site phase derivative, ordered logarithmic coefficients and physical
multiplier pairing.

\section{A reproducibility correction that does not change the V23 no-go}
\label{v24:v23-correction}

The theorem statement \eqref{v23:T} uses the complete V22 tangency functional
\[
 \mathcal T_\nu(Y,\ell)
 =Dd_{2,W}(Y)[F_1Y+G\ell,\nu]
   +\mathbb K_1(F_1Y+G\ell)[\nu,\ell].
\]
The auxiliary V23 recomputation script used to print the base constant wrote
its first term as $Dd_{2,W}(Y)[F_1Y,\nu]$, while its nonconstant coefficient
calculation did include the variation generated by $G\ell$.  The omitted base
term must therefore be checked rather than silently ignored.

\begin{lemma}[The omitted base term is exactly zero]
\label{v24:missing-zero}
At the V21 selected first tilt and the V23 diamond test,
\begin{equation}
 \boxed{Dd_{2,W}(Y)[G\ell_{21},\nu_\diamond]=0.}
 \label{v24:missing-zero-eq}
\end{equation}
Consequently the exact rational constant, its residue $148\pmod{241}$ and the
nonzero conclusion of \cref{v23:global-no-go} are unchanged.
\end{lemma}
\begin{proof}
Use the original homogeneous quadratic weak drift, with both the cubic
Hamiltonian polarization and the quadratic constraint polarization retained,
and insert the exact rational V21 tilt.  The coefficient in
\eqref{v24:missing-zero-eq} is one rational contraction of those arrays.  The
companion exact calculation gives numerator zero before reduction modulo any
prime.  The corrected base formula is therefore equal to the rational number
already stored in V23.  This is a correction of the reproducibility expression,
not of the theorem's mathematical value or its no-go conclusion.
\end{proof}

\section{The $\alpha_2$ canonical tangent and the lower incidence}
\label{v24:alpha2}

Retain the V23 chart with $b_2=b_3=-8$ and vary only $\alpha_2$ at the base
point
\[
 \alpha_*=(1,9/8,1),\qquad b_*=(16175/4096,-8,-8).
\]
The mean equation \eqref{v23:mean-surface} gives
\begin{equation}
 \dot b_1=-\frac9{256},\qquad \dot b_2=\dot b_3=0.
 \label{v24:b1dot}
\end{equation}
Let $Z_2=\partial_{\alpha_2}Y|_*$ be the resulting canonical first-jet
variation.  In physical variables it consists of the second axial
configuration profile, its complementary oscillatory momentum, and the
homogeneous diagonal correction induced by \eqref{v24:b1dot}.  In the V20
rationalized coordinates the same variation is the exact fraction-valued
array supplied with the certificate.

For the lower incidences put
\begin{equation}
 \mathcal L(Y,\ell)=
 \begin{pmatrix}
  \widehat B_2(Y)+\mathbb K_1(Y)\ell\\
  \mathsf m\ell\\
  t(Y)+R(Y)\ell
 \end{pmatrix}.
 \label{v24:lower-map}
\end{equation}
At $(Y,\ell_{21})$ this vanishes.  Its multiplier derivative is
\begin{equation}
 A:=D_\ell\mathcal L
   =\begin{pmatrix}\widehat K\\ \mathsf m\\R\end{pmatrix},
 \qquad \rank A=97,
 \label{v24:A}
\end{equation}
where the rank is the inherited rank-90 Poisson block plus the exact rank-seven
mean/cubic incidence of V23.

Define the canonical derivative covector
\begin{equation}
 c_2:=D_{\alpha_2}\mathcal L(Y,\ell_{21})
 =\begin{pmatrix}
   \dot{\widehat K}_2\ell_{21}+\dot{\widehat B}_{2,2}\\
   0\\
   \dot t_2+\dot R_2\ell_{21}
  \end{pmatrix}.
 \label{v24:c2}
\end{equation}
Every dot in \eqref{v24:c2} is obtained by exact homogeneous polarization of
the original coefficient functionals; it is not a finite-amplitude numerical
difference.

\begin{theorem}[A characteristic-zero upstream tangent of the complete lower incidence]
\label{v24:lower-tangent}
There is an exact rationalized multiplier tangent $m_2\in\R^{108}$ such that
\begin{equation}
 \boxed{A m_2=-c_2.}
 \label{v24:lower-tangent-eq}
\end{equation}
Equivalently, the formal curve
\[
 Y_s=Y+sZ_2+O(s^2),\qquad
 \ell_s=\ell_{21}+s m_2+O(s^2)
\]
satisfies every quadratic drift row, all four multiplier normalizations and
all three cubic constant-shift equations to first order in $s$.
\end{theorem}
\begin{proof}
The construction avoids a numerical nullspace.  Let $J$ be the 90 principal
indices of the V20 Poisson certificate.  Solve the nonsingular characteristic-
zero system
\[
 \widehat K_{J,J}(m_2^{(0)})_J
   =-\bigl(\dot{\widehat K}_2\ell_{21}
                 +\dot{\widehat B}_{2,2}\bigr)_J,
 \qquad (m_2^{(0)})_{J^c}=0.
\]
Exact multiplication verifies the remaining 18 Poisson rows as well.  Thus
$m_2^{(0)}$ solves the first block of \eqref{v24:lower-tangent-eq}.

Let $N_*$ be the exact 18-column V20 Poisson-null matrix.  The seven lower
rows on this nullspace are
\[
 C_7=\begin{pmatrix}\mathsf mN_*\\RN_*\end{pmatrix}.
\]
V23 already certifies that the columns
\[
             (0,1,2,3,7,10,15)
\]
form a nonsingular $7\times7$ minor; its displayed determinant is nonzero in
characteristic zero and has residue $126\pmod{241}$.  Solve on these seven
columns for $y$ so that
\[
 \mathsf m(m_2^{(0)}+N_*y)=0,
 \qquad
 R(m_2^{(0)}+N_*y)+\dot t_2+\dot R_2\ell_{21}=0.
\]
Put $m_2=m_2^{(0)}+N_*y$.  The companion certificate multiplies the resulting
large rational vector against every original row and obtains exact zero.  This
proves \eqref{v24:lower-tangent-eq} without a floating rank threshold.
\end{proof}

The theorem is a first-order accessibility statement.  It does not yet prove
that an exact nonlinear curve of lower-incidence zeros exists for nonzero
$s$; higher-order range conditions can in principle obstruct such an
integration.

\section{The diamond null test is stationary under all five first-jet directions}
\label{v24:diamond-stability}

The fixed test $\nu_\diamond$ can be used for the upstream differential only
if its Poisson-null property is not immediately lost when the first jet moves.
At the base the first Poisson matrix is linear in the canonical first jet, so
its derivative in a first-jet direction $Z$ is simply
$\mathbb K_1(Z)$.

\begin{lemma}[First-order persistence of the weak-null test]
\label{v24:diamond-first-order}
Let $Z_j$, $1\le j\le5$, be the five coordinate tangents of the first-jet chart
\eqref{v23:five-parameters}, with $b_1$ differentiated through
\eqref{v23:b1}.  Then
\begin{equation}
 \boxed{\mathbb K_1(Z_j)\nu_\diamond=0,
             \qquad j=1,\ldots,5.}
 \label{v24:diamond-five}
\end{equation}
Hence the derivative of the transported weak-null condition has no correction
term at the base point in any of the five canonical directions.
\end{lemma}
\begin{proof}
For each $Z_j$ the complete $108\times108$ first-Poisson derivative is
regenerated from the original quadratic constraint functional.  Exact rational
multiplication by the 108-component vector of \eqref{v23:diamond} gives the
zero vector in all five cases.  No statement about the null bundle away from
the base is inferred from these first derivatives.
\end{proof}

\section{The V23 obstruction moves under a lower-incidence canonical deformation}
\label{v24:tangency-derivative}

Let
\[
 V_s=F_1Y_s+G\ell_s,
 \qquad
 \Delta_\diamond(s)
 =Dd_{2,W}(Y_s)[V_s,\nu_\diamond]
   +\mathbb K_1(V_s)[\nu_\diamond,\ell_s].
\]
Because the weak drift is quadratic and the first Poisson bracket is linear in
its canonical field argument, $\Delta_\diamond$ is a polynomial of degree at
most two in the simultaneous formal increment $(Z_2,m_2)$.  Its directional
derivative can therefore be evaluated by exact polarization,
\begin{equation}
 \dot\Delta_2
 =\frac12\left\{
   \mathcal T_{\nu_\diamond}(Y+Z_2,\ell_{21}+m_2)
  -\mathcal T_{\nu_\diamond}(Y-Z_2,\ell_{21}-m_2)
 \right\}.
 \label{v24:Delta-polarization}
\end{equation}
This is a polynomial identity, not a finite-difference approximation.

\begin{theorem}[Nonzero characteristic-zero first-jet derivative of the tangency obstruction]
\label{v24:nonzero-upstream}
For the exact lower-incidence tangent of \cref{v24:lower-tangent},
\begin{equation}
 \boxed{\dot\Delta_2\ne0.}
 \label{v24:nonzero-upstream-eq}
\end{equation}
In the rationalized coefficient convention used by the V20--V23 exact
certificates,
\begin{equation}
 \boxed{\dot\Delta_2\equiv174\pmod{241},\qquad
        \dot\Delta_2\simeq-5.272804285668177\times10^7.}
 \label{v24:derivative-value}
\end{equation}
The exact numerator and denominator are retained in the proof package.  The
invertible V20 physical-coordinate congruence gives the same nonvanishing
conclusion for the original physical tangency covector.
\end{theorem}
\begin{proof}
Use the exact $m_2$ constructed in \cref{v24:lower-tangent} and evaluate
\eqref{v24:Delta-polarization} with the original rational $P_2$ and $\Phi_3$
coefficient functionals.  The base value is the same nonzero rational
$\Delta_\diamond$ as in V23, after the zero correction of
\cref{v24:missing-zero}.  The resulting derivative is one exact fraction.  Its
denominator is nonzero modulo 241 and its residue is 174, proving nonvanishing
in characteristic zero.  The displayed decimal is explanatory only.

As an independent scaling check, the good-prime evaluator run directly in the
physical variables gives residue $40$.  This equals $3\cdot174\pmod{241}$,
precisely the nonzero congruence factor induced by the retained
$\sqrt3$ momentum/shift rationalization.  Neither computation uses a floating
zero criterion.
\end{proof}

The combined conclusion can be summarized as
\begin{equation}
 \boxed{
 D_\ell\Delta_\diamond\big|_{\mathcal N_{21}}=0,
 \qquad
 D_{(Y,\ell)}\Delta_\diamond[(Z_2,m_2)]\ne0.}
 \label{v24:upstream-summary}
\end{equation}
Thus the global V23 no-go is a fixed-first-jet statement.  It is not an
invariant obstruction of the unchanged finite action under all canonical
first-jet deformations.

\section{Why the one-parameter Newton direction is not yet a root theorem}
\label{v24:newton}

The ratio of the exact base value to the derivative in
\cref{v24:nonzero-upstream} gives the formal linear Newton displacement
\[
                  -\Delta_\diamond/\dot\Delta_2
                       \simeq -0.14998917.
\]
Starting from $\alpha_2=9/8$, that linear prediction lies below
$\alpha_2=1$.  The equal-amplitude surface is precisely where the inherited
V2 diagonal lapse observation acquires its enlarged kernel; it is therefore a
structurally singular surface for the source-rank problem, not a point through
which a regular one-dimensional continuation may be assumed.

This numerical ratio is only a diagnostic.  No root is inferred from a first
Taylor coefficient, and the lower-incidence tangent of
\cref{v24:lower-tangent} has not been integrated through that singular surface.
The correct next search should use the multi-parameter regular chart rather
than force the $\alpha_2$ line through the known equal-amplitude degeneracy.

\section{Status and next exact calculation}
\label{v24:status}

V24 supplies the first characteristic-zero evidence that going upstream in the
canonical first jet is mathematically effective.  One complete canonical
direction can be accompanied by an exact leading-tilt correction preserving
all lower incidences to first order, and the V23 tilt-invariant tangency defect
then changes with a nonzero exact derivative.  The same weak test remains
Poisson-null to first order in every coordinate direction of the five-parameter
first-jet chart.

The result is deliberately weaker than an upstream root theorem.  A formal
lower-incidence tangent need not integrate; the transported Poisson rank and
the V21 higher weak-control minor must remain regular; and the full tangency
quotient contains more weak-null components than the single diamond test.  A
good-prime exploratory calculation indicates a four-dimensional canonical
projection of the linearized lower-incidence kernel, but that broader rank
statement is not promoted here because only the $\alpha_2$ direction has been
fully reconstructed over characteristic zero.

The next action-specific calculation is therefore an exact multi-parameter
continuation of the lower incidence on the regular side of the equal-amplitude
surface, followed by a verified zero or no-go for the complete transported
weak-null tangency quotient.  Only after such a new first jet is found should
V21's constrained second-field control and the full signed-source rank be
recomputed there.  No cofinal first-jet family, nonlinear propagation theorem,
differentiated-curvature estimate or common physical lifespan is proved in
this body.  The raw-action Einstein continuum bridge remains open.

\paragraph{Verification and preservation.}
The package contains the exact five first-Poisson derivatives, the corrected
$\alpha_2$ prepared-target derivative, the complete cubic-mean derivative,
the exact lower tangent and the tangency derivative fraction.  The verifier
checks \eqref{v24:missing-zero-eq}, all five identities
\eqref{v24:diamond-five}, every row of \eqref{v24:lower-tangent-eq}, the
nonzero residue in \eqref{v24:derivative-value}, and the physical good-prime
scaling regression.  A cached derivative which omitted the homogeneous
$b_1$ contribution to the $p$-only source term is excluded from the delivered
certificate; the corrected source derivative is regenerated and checked
against the original physical evaluator.

All twenty-three preceding marked research bodies are preserved byte for byte.
The finite coefficient checks certify the displayed algebra; they are not a
proof-assistant formalization, an exact nonlinear root of the upstream system,
or a spacetime integration.
% END IMMUTABLE EXACT-ACTION BRIDGE V24 BODY


% Future iterations append here.
% BEGIN IMMUTABLE EXACT-ACTION BRIDGE V25 BODY
\clearpage
\renewcommand{\thesection}{\arabic{section}}
\renewcommand{\theHsection}{bridgeV25.\arabic{section}}
\setcounter{section}{202}

\section{Complete canonical tangent certification before a nonlinear crossing}
\label{v25:context}

The canonical first-jet calculation in \cref{v24:lower-tangent} supplies one
exact tangent, but does not establish a manifold of lower-incidence zeros.
The next step is to resolve the complete five-parameter differential before
using it to search for a new tangency root. This body completes that
characteristic-zero calculation. Four independent canonical directions admit
full lower-incidence lifts; the remaining canonical direction is excluded by
an explicit nonzero Poisson-cokernel pairing. All four lifted directions
change the diamond tangency observation. A further homogeneity calculation
removes one common-scale freedom from the search.

There is an important distinction between these positive differential results
and the nonlinear problem. We give an exact local elimination of 98 lower
rows. Its 17 remaining scalar equations have zero value and zero first
differential at the base, but are not proved identically zero. Consequently
the four-dimensional canonical projection below is a projection of the
\emph{linearized} solution space, not a four-dimensional manifold assertion.
After fixing common scale, three canonical shape coordinates and eleven
multiplier coordinates remain in this exact residual formulation.

\paragraph{Inherited inputs and scope.}
The original coefficient functions are \eqref{v20:P2}, \eqref{v20:Phi3},
\eqref{v20:B2}, and \eqref{v21:R}. We retain the V20 rational Poisson matrix
and its 18-column null matrix, the exact V21 tilt and cubic-mean map, and
V23's fixed-first-jet no-go. V24's five Poisson differentials and its
corrected $\alpha_2$ calculation are inputs and independent regressions,
not new discoveries here. The new source calculation regenerates all five
prepared-target differentials and their complete cubic-mean differentials,
then solves the original lower rows exactly. No old uncorrected derivative
cache is used. All coefficient statements in this body concern $N=3$.
No nonlinear root, time propagation, or refinement-uniform estimate is
inferred from these finite calculations.

\section{The exact mean-surface differential and the full lower map}
\label{v25:coordinates}

Write the canonical parameter as
\[
 p=(\alpha_1,\alpha_2,\alpha_3,b_2,b_3),\qquad
 b_1(p)=\frac{\Sigma_\alpha/4-b_2b_3}{b_2+b_3},
 \qquad \Sigma_\alpha=\sum_i\alpha_i^2.
\]
The first field $Y(p)$ is precisely \cref{v23:five-chart}: the configuration
and oscillatory momentum use complementary polarizations, and the
homogeneous boost is $\omega_h\diag(b_1,b_2,b_3)$. Work near
\[
 p_*=(1,9/8,1,-8,-8),\qquad b_1(p_*)=16175/4096.
\]
The denominator is nonzero there. In V20's rationalized variables, write
$x$ for the multiplier coordinates, $\ell=Q_3x$, and let $x_*$ be the exact
V21 representative. Use the following unambiguous rational map:
\begin{equation}
 \mathcal L(p,x)=
 \begin{pmatrix}
  \widehat K(p)x+\widehat b(p)\\
  \mathsf m x\\
  t(p)+R(p)x
 \end{pmatrix}\in\R^{115}.
 \label{v25:lower}
\end{equation}
Here $\mathsf m$ uses component sums; its zero equations are identical to
component means. The first block uses all 108 original rows, and the last
block uses the same spatial-sum convention as V21. In particular
$\mathcal L(p_*,x_*)=0$. Put
\[
 A=D_x\mathcal L(p_*,x_*),\qquad C=D_p\mathcal L(p_*,x_*).
\]
These symbols refer only to the lower-incidence differential, not to the
Cartan Hessian or the scalar operator of V13--V19.

The five derivatives of the homogeneous mean-surface coefficient are
\begin{equation}
 D_pb_1\big|_*=
 \left(-\frac1{32},-\frac9{256},-\frac1{32},
                -\frac{16593}{65536},-\frac{16593}{65536}\right).
 \label{v25:b1-derivatives}
\end{equation}
They follow by differentiating
$b_1(b_2+b_3)+b_2b_3=\Sigma_\alpha/4$. The two last terms, as well as
the first three, enter the actual momentum variation. In particular the
$p$-only part of the prepared quadratic target is differentiated too.

\paragraph{Exact coefficient convention.}
Let $Z_j=\partial_{p_j}Y|_*$. In the rational scaling of V20 the $j$th
column of $C$ has the form
\begin{equation}
 C_j=\begin{pmatrix}c_j\\0_4\\d_j\end{pmatrix},\quad
 c_j=(\partial_j\widehat K)x_*+\partial_j\widehat b,\quad
 d_j=\partial_jt+(\partial_jR)x_*.
 \label{v25:C}
\end{equation}
All derivatives are derivatives of the original homogeneous coefficient
polynomials. For example a polynomial $F$ of degree at most three obeys
\begin{equation}
 DF(Y)[Z]=\frac{8\{F(Y+Z/2)-F(Y-Z/2)\}
                       -\{F(Y+Z)-F(Y-Z)\}}6.
 \label{v25:polarization}
\end{equation}
The identity follows by expanding in the scalar increment; its cubic odd
part cancels exactly. This is not a small-step difference quotient.
The evaluations in this body use rational multiplication by $1/2$ and $1/6$
and exact integer/fraction arrays throughout. The full prepared source
includes the original lapse correction. The cubic-mean functional includes
all three terms of \eqref{v21:R}. The regenerated $j=2$ entries agree
exactly, before reduction modulo a prime, with the corrected V24 certificate.

\section{The canonical projection of the linearized lower incidence}
\label{v25:projection}

Let $N_*$ be the original $108\times18$ rational Poisson-null matrix. Its
rows at the 18 free indices form $I_{18}$, as in V20. Define
\[
 \Omega=N_*^T(c_1,c_2,c_3,c_4,c_5).
\]
This operation tests the complete quadratic target against the actual
Poisson cokernel. It is not a nullspace of the positive multiplier Hessian.

\begin{theorem}[An exact one-row canonical compatibility condition]
\label{v25:canonical-theorem}
There is a nonzero rational vector $v\in\R^{18}$ such that
\begin{equation}
 \boxed{\Omega=v(1,0,-1,0,0).}
 \label{v25:Omega}
\end{equation}
Consequently
\begin{equation}
 \boxed{
 \pi_p\ker(A\ C)=
       \{\dot p\in\R^5:\dot\alpha_1=\dot\alpha_3\}.}
 \label{v25:canonical-space}
\end{equation}
Moreover
\[
 \rank A=97,\qquad \rank(A\ C)=98,\qquad
 \dim\ker(A\ C)=15.
\]
The kernel with $\dot p=0$ has dimension 11. These are exact
characteristic-zero ranks at the stated base point.
\end{theorem}
\begin{proof}
The source program evaluates all entries of \eqref{v25:C} by the displayed
coefficient formulas and \eqref{v25:b1-derivatives}. Multiplication over
$\mathbb Q$ gives
\[
 N_*^Tc_2=N_*^Tc_4=N_*^Tc_5=0,\qquad
 N_*^T(c_1+c_3)=0.
\]
Its first nonzero remaining component has zero-based index 13 and value
$v_{13}$, with
\begin{equation}
 v_{13}\ne0,\qquad v_{13}\equiv230\pmod{241}.
 \label{v25:obstruction-residue}
\end{equation}
The exact fraction is stored in \src{canonical_obstruction.json}; its
denominator is a unit at 241. The full 18-by-5 product, not just this
component, is supplied. Exact rational zeros prove \eqref{v25:Omega}; the
nonzero residue proves its rank is one in characteristic zero.

The Poisson block has rank 90 and kernel $\Ran N_*$. The seven remaining
rows on this kernel are
\[
 C_7=\begin{pmatrix}\mathsf mN_*\\RN_*\end{pmatrix}.
\]
The columns $(0,1,2,3,7,10,15)$ have the inherited nonzero determinant
$126\pmod{241}$. Thus these seven rows are onto on the Poisson kernel,
and $\rank A=90+7=97$.

If $A\dot x+C\dot p=0$, multiplication of its top block by $N_*^T$
forces $\Omega\dot p=0$, hence $\dot\alpha_1=\dot\alpha_3$.
Conversely that equality makes the entire top target orthogonal to
$\ker\widehat K$. Since $\widehat K$ is a real skew matrix, its range
is that orthogonal complement. Solve the top rows, then use the onto
$C_7$ to correct all seven remaining rows without changing the top solution.
This proves equality in \eqref{v25:canonical-space}. Its dimension is four;
rank-nullity and the 11-dimensional vertical kernel give total kernel
dimension 15 and augmented rank $113-15=98$.
\end{proof}

The four lifts are explicitly available, rather than only supplied by the
range argument. Use the parameter directions
\begin{equation}
 q_2=(0,1,0,0,0),\quad q_{13}=(1,0,1,0,0),\quad
 q_{b2}=(0,0,0,1,0),\quad q_{b3}=(0,0,0,0,1).
 \label{v25:four-directions}
\end{equation}
Let $J$ be V20's 90 principal indices and let $U=(0,1,2,3,7,10,15)$ be
the seven null-column indices. For $c=(Cq)_{1:108}$ and
$d=(Cq)_{113:115}$, first set
\[
 (m^{(0)})_J=-\widehat K_{J,J}^{-1}c_J,\qquad
 (m^{(0)})_{J^c}=0.
\]
Then set
\begin{equation}
 m=m^{(0)}-N_{*,U}(C_{7,:,U})^{-1}
       \begin{pmatrix}\mathsf m m^{(0)}\\Rm^{(0)}+d\end{pmatrix}.
 \label{v25:lift}
\end{equation}
The certificate multiplies each of these four rational vectors against
\emph{all 115 rows} and obtains
\begin{equation}
                  A m_q+Cq=0.
 \label{v25:lift-zero}
\end{equation}
The $q_2$ vector reproduces the V24 lift byte for byte as a list of reduced
fractions. Different valid lifts differ by the exact V21 vertical kernel;
no choice of minimum norm is required here.

\section{A structural diamond null test on the whole first-jet family}
\label{v25:diamond}

V24 verified the five tangent null identities at the base. To use a fixed
diamond test at a finite canonical deformation, that alone is insufficient.
The remaining independent linear field direction can be checked explicitly.

\begin{theorem}[Global coefficient nullity at the three-site cutoff]
\label{v25:diamond-theorem}
At $N=3$, for every first field assembled from the three complementary axial
amplitudes and an arbitrary homogeneous diagonal boost, one has
\begin{equation}
 \boxed{\mathbb K_1(Y)\nu_\diamond=0,\qquad
 \nu_\diamond=(0,\nabla_\delta\phi_\diamond),\quad
 \phi_\diamond=\tfrac23(\cos2\pi x_1-\cos2\pi x_2).}
 \label{v25:global-diamond}
\end{equation}
In particular this holds throughout the five-parameter mean-calibrated
chart, not merely to first order at its base.
\end{theorem}
\begin{proof}
The first Poisson coefficient is linear in the canonical first field.
The linear field family in question has six coefficient coordinates
$(\alpha_1,\alpha_2,\alpha_3,b_1,b_2,b_3)$. The five coordinate tangents
of \eqref{v25:b1-derivatives}, together with the independent normal direction
$(0,0,0,1,0,0)$ in these six coordinates, form a matrix of determinant one.
The first five annihilation identities are the exact inherited V24 products.
For the sixth direction, synthesize the original field
\[
 U=0,\qquad p_R=\diag(0,-3,-3),\qquad p=\sqrt3p_R.
\]
Polarize the original $P_2$ against all 108 multiplier tests to construct its
complete first Poisson coefficient. The exact product of that rational
matrix with the diamond vector has all 108 entries zero. The full matrix
and zero product are retained in \src{structural_diamond.json}. The diagonal
physical scaling changes no pure-shift zero product: each output block has
its common nonzero grading factor. Linearity and the six-coordinate spanning
identity prove \eqref{v25:global-diamond} for every member of the linear
family. The mean relation is a restriction of that family, so inherits the
identity. This is a polynomial coefficient proof at $N=3$, not a statement
about every nonlinear finite Poisson matrix or other cutoffs.
\end{proof}

\section{All four admissible canonical directions move the obstruction}
\label{v25:derivatives}

Use the complete tangency polynomial of V22--V24,
\[
 \mathcal T_\nu(Y,\ell)
   =Dd_{2,W}(Y)[F_1Y+G\ell,\nu]
            +\mathbb K_1(F_1Y+G\ell)[\nu,\ell].
\]
Let $\widehat\Delta(p,x)$ denote its diamond evaluation in the retained
V24 rational coefficient convention. Its physical version differs by the
unchanged nonzero grading/normalization factors; the table below does not
introduce new physical units. At the base its exact value is $d_*$, the
nonzero V23 value, with $d_*\equiv148\pmod{241}$.

\begin{theorem}[A well-defined nonzero canonical derivative in four directions]
\label{v25:derivative-theorem}
For the four complete lower-incidence lifts \eqref{v25:lift}, the derivative
$D\widehat\Delta[(q,m_q)]$ is independent of the choice of lift. In the
inherited rational convention its values have the following exact residues
and explanatory approximations:
\begin{center}
\begin{tabular}{@{}lrr@{}}
\toprule
Canonical direction & Residue modulo 241 & Approximate derivative\\
\midrule
$q_2$    & 174 & $-5.2728043\times10^7$\\
$q_{13}$ & 224 & $\phantom{-}6.0559478\times10^7$\\
$q_{b2}$ & 56  & $-1.262752940\times10^8$\\
$q_{b3}$ & 208 & $\phantom{-}1.284075066\times10^8$\\
\bottomrule
\end{tabular}
\end{center}
Every derivative is nonzero in characteristic zero. No zero of
$\widehat\Delta$, or nonlinear lower-incidence curve, is inferred.
\end{theorem}
\begin{proof}
Two lifts of the same canonical direction differ by $\ker A$, precisely the
11-dimensional fixed-first-jet admissible tilt space. V23 proves that the
diamond polynomial is constant on that entire affine space. Its differential
therefore annihilates their difference, proving independence.

The weak drift is homogeneous quadratic in its canonical argument, and
$\mathbb K_1$ and $F_1$ are linear in theirs. Consequently $\mathcal T$
is homogeneous quadratic in the simultaneous pair $(Y,\ell)$. For
$Z=D_pY[q]$, its derivative is exactly
\[
 \frac12\{\mathcal T_{\nu_\diamond}(Y+Z,\ell_*+Q_3m_q)
         -\mathcal T_{\nu_\diamond}(Y-Z,\ell_*-Q_3m_q)\}.
\]
The rational implementation retains every term in this graded formula.
Each result is stored as a reduced fraction in \src{tangency_derivatives.json}. Its denominator is
nonzero modulo 241 and its residue is the displayed nonzero entry. This
proves nonvanishing over $\mathbb Q$ and in the physical real embedding.
The $q_2$ fraction is identical to V24's certified fraction, not just to its
decimal approximation. The other three fractions use the newly reconstructed
complete lower lifts. No modular-only tangent or finite differencing step
is substituted into the calculation.
\end{proof}

\section{Common scaling is an exact redundancy of the coefficient search}
\label{v25:scaling}

The mean relation in the five-parameter family is homogeneous of degree two
in $(\alpha,b)$. Hence simultaneous scaling of those six coefficients by
$c>0$ sends $Y$ to $cY$ and preserves mean compatibility.

\begin{proposition}[Scale covariance and three independent canonical shape coordinates]
\label{v25:scale-theorem}
If $(Y,\ell)$ satisfies the complete lower coefficient equations, then
$(cY,c\ell)$ does too. The diamond tangency polynomial satisfies
\begin{equation}
 \boxed{\mathcal T_{\nu_\diamond}(cY,c\ell)
                         =c^2\mathcal T_{\nu_\diamond}(Y,\ell).}
 \label{v25:scale}
\end{equation}
Thus nonzero common scaling cannot turn a nonzero tangency value into zero.
At the base the radial canonical tangent is
\[
 q_{\rm rad}=q_{13}+\tfrac98q_2-8q_{b2}-8q_{b3}.
\]
The exact derivative fractions obey
\begin{equation}
 \boxed{
 \dot\Delta_{13}+\tfrac98\dot\Delta_2
                -8\dot\Delta_{b2}-8\dot\Delta_{b3}=2d_*.}
 \label{v25:Euler}
\end{equation}
Modulo this radial freedom, the canonical projection of the linearized
lower incidence has dimension three.
\end{proposition}
\begin{proof}
The quadratic target has degree two in $Y$, the first Poisson matrix degree
one, the cubic mean target degree three, and its tilt map degree two.
Therefore the first, mean-normalization, and last blocks of the lower map
scale respectively by $c^2,c,c^3$. Each remains zero. The two terms of
$\mathcal T$ both have simultaneous degree two, proving \eqref{v25:scale}.
Differentiation at $c=1$ gives derivative $2d_*$ along $(Y,\ell_*)$.
The specified radial parameter vector synthesizes exactly $Y$, including
its homogeneous boost. Its multiplier lift may be chosen $x_*$ itself.
The linear combination of the four certified lifts differs from this lift
by $\ker A$, which the tangency differential annihilates. This gives
\eqref{v25:Euler}; the verifier additionally checks it as an equality of
the computed fractions. The radial vector is nonzero in the four-dimensional
canonical projection, so the quotient has dimension three.
\end{proof}

This scaling is redundancy of the homogeneous coefficient search, not a
change of the physical clock. It can also be viewed as a reparametrization
of the small initial amplitude. On $\alpha_3>0$ one may set $\alpha_3=1$
by dividing $(\alpha,b,\ell)$ by $\alpha_3$. The remaining canonical
linearized shape directions at the base are represented by
$\alpha_2,b_2,b_3$, with $\dot\alpha_1=\dot\alpha_3=0$. Both boost
directions have nonzero diamond derivatives without changing the configuration
detuning to first order. This avoids the \emph{linearized} necessity of
crossing the equal-amplitude value on the old $\alpha_2$-only trial, but does
not prove a regular finite crossing in either boost direction.

\section{An exact nonlinear reduction, and the equations not removed by rank}
\label{v25:nonlinear}

We now give a local nonlinear representation of the lower zero set. This
representation makes explicit why the tangent dimension cannot yet be called
a manifold dimension. The original maps are rational analytic functions of
$p$ on $b_2+b_3\ne0$ and affine in $x$.

Let $I_J$ insert V20's 90 principal multiplier coordinates. Keep the
\emph{fixed base} matrix $N_*$ as an analytical coordinate basis, not as an
asserted nullspace of $\widehat K(p)$ away from the base. Split its 18 columns
into $U=(0,1,2,3,7,10,15)$ and its 11-column complement $W$. Write
\begin{equation}
 x=x_*+I_Jz+N_{*,U}y+N_{*,W}\eta,
 \quad z\in\R^{90},\ y\in\R^7,\ \eta\in\R^{11}.
 \label{v25:multiplier-chart}
\end{equation}
This is an invertible affine coordinate change, since $N_{*,J^c,:}=I_{18}$.

\begin{theorem}[Local lower incidence with seventeen unresolved analytic residuals]
\label{v25:residual-theorem}
Near $(p_*,x_*)$, 97 original lower equations determine $z,y$ uniquely and
analytically as functions of $(p,\eta)$. One further original row determines
\begin{equation}
 \alpha_1=\chi(\alpha_2,\alpha_3,b_2,b_3,\eta),\qquad
 \chi=\alpha_3+O(\|(\alpha_2,\alpha_3,b_2,b_3,\eta)
                     -(9/8,1,-8,-8,0)\|^2).
 \label{v25:chi}
\end{equation}
All remaining lower equations are exactly equivalent to a specified analytic
vector $\mathcal R\in\R^{17}$ on these 15 coordinates. At the base,
\begin{equation}
                         \mathcal R=0,\qquad D\mathcal R=0.
 \label{v25:residual-double}
\end{equation}
No assertion that $\mathcal R$ vanishes identically is made. After the exact
common-scale normalization $\alpha_3=1$, the same lower zero set modulo scale
is described by 17 residual equations on three canonical shape coordinates
and eleven multiplier coordinates.
\end{theorem}
\begin{proof}
Select the 90 top rows with indices $J$, and all four normalization and three
cubic-mean rows. Their derivative in $(z,y)$ at the base is
\[
 \begin{pmatrix}
  \widehat K_{J,J}&0\\
  (\mathsf m;R)I_J&C_{7,:,U}
 \end{pmatrix}.
\]
Its determinant is nonzero; in the fixed ordering its residue is
$25\cdot126=17\pmod{241}$. At every nearby $p$ the selected equations remain
affine in $(z,y)$ with an invertible coefficient matrix. Exact inversion
therefore gives analytic (indeed rational in these coefficient coordinates)
functions $z(p,\eta),y(p,\eta)$. This solves these 97 equations exactly,
not only their differential.

Let $\rho\in\R^{18}$ be the 18 unselected top rows after that substitution,
in the order of $J^c$. At the base the selected top differentials vanish.
Since $N_{*,J^c,:}=I_{18}$, the full Poisson-cokernel product identifies
\[
 D_p\rho=\Omega=v(1,0,-1,0,0),\qquad D_\eta\rho=0.
\]
The 13th zero-based component has nonzero $\alpha_1$ derivative
$v_{13}$. This is the original multiplier row with zero-based index 89.
The analytic implicit function theorem solves that one row for $\alpha_1$,
giving \eqref{v25:chi}. All its other first derivatives are zero except
$\partial_{\alpha_3}\chi=1$. Substitute this solution into the other 17
components of $\rho$ to define $\mathcal R$. Their entire first differential
is a multiple of $d\alpha_1-d\alpha_3$, which vanishes after substitution.
This proves \eqref{v25:residual-double} and exact equivalence of the remaining
equations. No original row has been dropped.

Finally the exact scaling action of \cref{v25:scale-theorem} has one unique
intersection with $\alpha_3=1$ for every sufficiently nearby positive
$\alpha_3$. Restrict the preceding equations to this slice. It removes one
coordinate and no lower-incidence solution orbit. This gives the final
three-plus-eleven-coordinate description. A zero first differential of the
17 residuals is not proof that they vanish as functions or that their zero
set has that ambient dimension. The statement is only their exact
representation and vanishing through first order.
\end{proof}

For the reduced diamond function on this scale-fixed analytical chart,
write $\zeta=(\alpha_2-9/8,b_2+8,b_3+8,\eta)$. The preceding results give
\begin{equation}
 \widehat\Delta_{\rm red}(\zeta)
   =d_*+\dot\Delta_2(\alpha_2-9/8)
       +\dot\Delta_{b2}(b_2+8)+\dot\Delta_{b3}(b_3+8)
       +O(\|\zeta\|^2).
 \label{v25:reduced-Taylor}
\end{equation}
There is no linear $\eta$ term, by the global fixed-first-jet constancy.
Equation \eqref{v25:reduced-Taylor} specifies genuine sensitivity on the
selected-row chart, but a valid new first jet must also solve
$\mathcal R(\zeta)=0$ and the \emph{complete} weak-null tangency system.
Only the diamond test has here been proved fixed throughout the family;
other tests cannot be transported by reusing their old null coordinates.

\section{Verification record and the new frontier}
\label{v25:status}

The unfinished characteristic-zero tangent calculation is now complete.
The actual lower system admits exactly four canonical directions to first
order, each has an explicit rational full-row lift, and each changes the
formerly tilt-invariant diamond observation. The diamond is a structural
first-Poisson null test on the entire six-coefficient linear first-field
family at this cutoff. These conclusions use the true mean-surface momentum
derivatives, not their earlier modular indications or a stale scalar source.

The next search is also smaller and better specified. One of the four
canonical directions is an exact common-scale redundancy. After removing
it, the lower equations have an exact residual representation in three
canonical shape coordinates and eleven multiplier coordinates, with 17
analytic residuals starting at second order. Those residuals must be evaluated
or proved dependent before declaring a nonlinear lower-incidence manifold.
The diamond equation is only one tangency condition, and its nonzero derivative
is not a root certificate. No full weak-null root, new higher-source rank,
new exactly prepared dynamical branch, cofinal tilt bound, or common physical
lifespan is claimed in this body. The original action and clock are unchanged.

\paragraph{Executed checks.}
The package regenerates all five 108-component prepared-target derivatives
and all five three-component cubic-mean derivatives with exact fractions.
The corrected V24 $\alpha_2$ outputs are recovered exactly. It supplies and
verifies the full cokernel product, four 115-row lower lifts, the additional
normal Poisson-null calculation, and all four reduced tangency fractions.
The exact Euler identity \eqref{v25:Euler} is checked independently after
those computations. Inherited coefficient matrices are identified as inputs;
no numerical zero, sampled eigenvalue, unverified Gröbner calculation, or
proof-assistant formalization is substituted for these certificates. The
analytic residual theorem uses the stated nonsingular blocks and the ordinary
analytic implicit function theorem. Its 17 higher-order residuals have not
been numerically or symbolically expanded in this iteration.

The complete V1--V24 marked bodies are preserved byte for byte. The new
verifier records their hashes and checks the resulting source references and
labels separately from the mathematical certificates. Compilation and visual
inspection are artifact checks, not independent validation of all inherited
mathematics. Background on first-order tangent data as distinct from an
actual parameterized solution family is given in \cite{v25:Stacks}; all finite
source-specific ranks and constructions used here are provided above.

\begin{thebibliography}{9}
\raggedright
\setcounter{enumiv}{18}
\bibitem{v25:Stacks}
The Stacks Project Authors, \emph{The Stacks Project}, Section 33.16,
Tangent spaces, Tag 0B28.
\url{https://stacks.math.columbia.edu/tag/0B28}.
Background for the meaning of first-order solution data, not a source for the
present original-action coefficients or a claim of nonlinear integrability.
\end{thebibliography}
% END IMMUTABLE EXACT-ACTION BRIDGE V25 BODY


% Future iterations append here.
% BEGIN IMMUTABLE EXACT-ACTION BRIDGE V26 BODY
\clearpage
\renewcommand{\thesection}{\arabic{section}}
\renewcommand{\theHsection}{bridgeV26.\arabic{section}}
\setcounter{section}{210}

\section{Integrate the lower incidence before imposing tangency}
\label{v26:context}

The seventeen residuals retained in \cref{v25:status} are a genuine logical
obligation: their zero first differential does not by itself make them
identities. We discharge that obligation on the actual complementary
three-site family. Its symmetric-amplitude slice has a fixed eighteen-dimensional
Poisson kernel, and the complete prepared quadratic target annihilates that
kernel identically. These two coefficient identities imply that all seventeen
residuals vanish, and that the lower solution set really is an analytic
manifold with four canonical parameters and eleven multiplier parameters.
The complete weak-null frame, not only its diamond member, can be kept fixed
on this manifold.

We then evaluate a finite, rather than infinitesimal, change of the homogeneous
boost. A short rational interval contains a zero of the diamond tangency
component on an exactly solved lower-incidence curve. Exact endpoint values
and inverse bounds throughout the interval exclude a sign change caused by a
pole. This is a one-component tangency zero, not a zero of the complete eight-row
tangency system. In particular no new full higher weak-source certificate or
propagating Einstein history is inferred at that point.

\paragraph{Inputs and nonduplication.}
The scalar coefficient functions are the unchanged \eqref{v20:P2},
\eqref{v20:Phi3}, \eqref{v20:B2}, and \eqref{v21:R}. We retain their
rationalization, V20's eighteen-column Poisson null matrix, V21's seven
normalization/mean controls, and V25's transverse canonical derivative.
The initial constraint and weak-rate preparation theorems remain inherited
analytic inputs. What is new is the exhaustive coefficient identity on a
linear first-field family, the resulting \emph{nonlinear} solution manifold,
its complete fixed weak-null frame, and the interval-certified scalar crossing.
Every new rank and crossing assertion in this body concerns $N=3$, $h=1/3$.
Neither a finite-action term nor its time normalization is changed.

\section{A linear first-field family and its exact common Poisson kernel}
\label{v26:kernel}

Let $a,c,d_1,d_2,d_3$ be real coefficients and set
\begin{equation}
 \begin{gathered}
 \alpha=(a,c,a),\qquad B_c=\omega_h\diag(d_1,d_2,d_3),\qquad
 P_c=B_c-(\tr B_c)I,\\
 Y=\left(\sum_i\alpha_iT_i\cos(2\pi x_i),\;
 P_c+\frac{\omega_h}{2}\sum_i\alpha_iP_i\cos(2\pi x_i)\right).
 \end{gathered}
 \label{v26:linear-family}
\end{equation}
Here $T_i,P_i$ are precisely the complementary transverse tensors of V12,
not Lorentz generators. At $N=3$, $\omega_h=3\sqrt3$.
Before imposing the mean equation, \eqref{v26:linear-family} is a
five-dimensional \emph{linear} space of canonical first fields. Off the mean
equation it is used only to establish polynomial identities; its elements are
not all asserted to be compatible initial records.

Use V20's rational coordinates $p=\sqrt3p_R$, $\ell=Q_3x$, and write
$\widehat K(Y)$, $\widehat b(Y)$ for its rationalized Poisson matrix and
complete prepared quadratic target. In these coordinates the first field is
linear in $q=(a,c,d_1,d_2,d_3)$, $\widehat K$ is homogeneous linear, and
$\widehat b$ is homogeneous quadratic. The last claim concerns the
\emph{prepared} polynomial, including its lapse correction, not the
unprepared canonical drift.

For explicit descriptions of grid functions put
\[
 f_i(x)=\cos(2\pi x_i),\qquad
 s_i(x)=\frac2{\sqrt3}\sin(2\pi x_i).
\]
Their one-dimensional value lists are $(1,-1/2,-1/2)$ and $(0,1,-1)$.
Define the following eighteen multiplier directions, with spans unaffected
by the nonzero uniform shift scaling in $Q_3$:
\begin{equation}
 \begin{split}
 \mathcal N={\rm span}\{\;&(1,0),\ (f_i,0),\ (s_i,0),\ (0,e_i),
                         \ (0,\nabla_\delta f_i),\ (0,\nabla_\delta s_i)
                         \quad (1\le i\le3),\\
                      &(0,\nabla_\delta(f_1s_3)),\
                        (0,\nabla_\delta(s_1f_3))\;\}.
 \end{split}
 \label{v26:explicit-kernel}
\end{equation}
The constant vectors $e_i$ in this display are spatial shift vectors.

\begin{theorem}[A fixed complete Poisson kernel on the symmetric-amplitude slice]
\label{v26:common-kernel}
For every field in \eqref{v26:linear-family}, the original first Poisson
coefficient annihilates $\mathcal N$. This space has dimension eighteen.
On a neighborhood of the retained base
\[
 (a,c,d_1,d_2,d_3)=(1,9/8,16175/4096,-8,-8)
\]
it is the complete kernel and the Poisson rank is exactly ninety. In the
rational coordinates it is the same fixed kernel as the matrix $N_*$ of V20.
\end{theorem}
\begin{proof}
The seven lapse functions in \eqref{v26:explicit-kernel} are independent by
one-dimensional grid orthogonality. The three constant shifts are independent
and orthogonal to gradients. The eight scalar potentials in its gradient
part are mutually orthogonal, nonzero, and mean-zero. The original phase
gradient is injective on mean-zero arrays, because its Fourier symbol has no
nonconstant zero on an odd grid. This proves dimension eighteen.

Here is a finite coefficient verification specifying the entire identity,
not an inference from randomly selected parameters. Synthesize the five
unit fields $Y_j$ of \eqref{v26:linear-family}. V24's five evaluated first-field
directions, together with V25's additional homogeneous boost direction,
span the full six-coefficient linear first-field space, with exact spanning
determinant one. Linear combination of those retained original Poisson
coefficients therefore supplies each $\widehat K(Y_j)$ exactly. The companion
certificate verifies
\begin{equation}
             \widehat K(Y_j)N_*=0,\qquad j=1,\ldots,5,
 \label{v26:five-kernel-products}
\end{equation}
as $108\times18$ rational matrix identities. It also builds the rational
column matrix $B_{\mathcal N}$ from the eighteen displayed functions and verifies
\[
 B_{\mathcal N}=N_*C_{\mathcal N},\qquad
 C_{\mathcal N}=(B_{\mathcal N})_{J^c,:},\qquad
 \det C_{\mathcal N}=-1162261467/2\ne0.
\]
Here $J$ is V20's ninety-column principal set and $N_{*,J^c,:}=I_{18}$.
Thus these are the same kernel coordinates, with an explicit invertible
change of basis. Linearity in the five field coefficients proves annihilation
at every real parameter, including those not on the mean surface.

The ninety-dimensional principal minor at the base is the retained nonzero
minor of residue $25$ modulo $241$. It remains nonzero on an open real
neighborhood. There the rank is at least ninety, while
\eqref{v26:five-kernel-products} gives rank at most $108-18=90$.
Hence the kernel and rank are exactly as stated. This conclusion concerns the
first Poisson coefficient; it is not a classification of the full nonlinear
constraint bracket.
\end{proof}

\section{The complete quadratic target belongs to that range identically}
\label{v26:target}

\begin{theorem}[An exhaustive source--cokernel identity]
\label{v26:range-identity}
Throughout the linear family \eqref{v26:linear-family},
\begin{equation}
                    \boxed{N_*^T\widehat b(Y)=0.}
 \label{v26:range-zero}
\end{equation}
Consequently, wherever its ninety-dimensional minor remains nonzero, the
complete equation $\widehat K(Y)x=-\widehat b(Y)$ is solvable on all
108 original rows.
\end{theorem}
\begin{proof}
Write $Y(q)=\sum_{j=1}^5q_jY_j$. Because the complete target is quadratic,
\[
 \widehat b(Y(q))=\sum_jq_j^2 b_{jj}+
                         \sum_{j<k}q_jq_k b_{jk},\qquad
 b_{jj}=\widehat b(Y_j),\quad
 b_{jk}=\widehat b(Y_j+Y_k)-\widehat b(Y_j)-\widehat b(Y_k).
\]
The certificate evaluates all five diagonal and all ten mixed contractions
with $N_*$. Each of their eighteen entries is exactly zero over $\mathbb Q$.
This exhausts all fifteen monomials; it is not testing a continuum of
parameters by a finite sample without a degree bound.

For reproducibility, the contractions use \eqref{v20:B2} directly. For a
rational test column $\nu$ they retain the cubic polarization
\[
 \tfrac12\{\Phi_3(Y+G\nu)-\Phi_3(Y-G\nu)\}-\Phi_3(G\nu),
\]
the quadratic constraint polarization in $F_1Y$, and the prepared lapse row
$P_{(0,\delta n),2}(Y)$. The physically rationalized shift target is evaluated
as $B_{2,\beta}(U,p_R)+2B_{2,\beta}(0,p_R)$, exactly as in V20;
the lapse target has its original scaling. Contracting before extraction
reduces arithmetic but changes none of these rows or terms. The files give
the original coefficient functions, all fifteen eighteen-component results,
and a regeneration program using arbitrary-precision fractions.

The resulting polynomial identity proves \eqref{v26:range-zero} everywhere
on the linear space. On the regular neighborhood,
$\ker\widehat K=\operatorname{Ran}N_*$ and
$\widehat K^T=-\widehat K$. In a finite positive pairing the range of this
skew matrix is the orthogonal complement of its kernel. This proves the
full-row solvability claim.
\end{proof}

The already established weak quadratic neutrality accounts for the
curl-free part of this compatibility; it does not by itself account for the
additional lapse-null tests. The exhaustive calculation above includes those
lapse rows and the original lapse preparation correction.

\section{The nonlinear lower solution set is a genuine manifold}
\label{v26:manifold}

Now impose the original mean surface
\begin{equation}
 d_1(d_2+d_3)+d_2d_3=\frac{2a^2+c^2}{4},\qquad
 d_1=\frac{(2a^2+c^2)/4-d_2d_3}{d_2+d_3}.
 \label{v26:mean}
\end{equation}
Its denominator is nonzero near the base. Let
$\mathsf m$ take the four component sums, and let $t_3(Y),R(Y)$ be
V21's complete cubic-mean constant and tilt matrix. Set
\[
 C_7(Y)=\begin{pmatrix}\mathsf mN_*\\R(Y)N_*\end{pmatrix},\qquad
 U=(0,1,2,3,7,10,15),\qquad V=\{0,\ldots,17\}\setminus U.
\]
The seven-by-seven matrix $H(Y)=C_{7,:,U}(Y)$ has the retained nonzero
base determinant of residue $126$ modulo $241$.

\begin{theorem}[Exact integration of the four canonical directions]
\label{v26:lower-manifold}
Near the V25 base, the complete lower-incidence zero set is exactly an
analytic manifold with
\begin{equation}
                 \boxed{\alpha_1=\alpha_3,\qquad
                  \dim\mathcal L^{-1}(0)=4+11=15.}
 \label{v26:manifold-dimension}
\end{equation}
Its projection to canonical parameters is a relatively open neighborhood in
that four-dimensional mean-compatible hypersurface. The multiplier fibre
has dimension eleven. In V25's selected-row coordinates,
\begin{equation}
                    \boxed{\chi=\alpha_3,\qquad
                              \mathcal R\equiv0.}
 \label{v26:residuals-zero}
\end{equation}
Thus all seventeen residuals retained there vanish identically on a sufficiently
small neighborhood. After fixing common scale by $\alpha_3=1$, the quotient
has three canonical shape parameters and eleven multiplier parameters.
\end{theorem}
\begin{proof}
We first give a direct rational formula on \eqref{v26:mean}. Define a particular
quadratic-drift solution by
\[
 (x^{(0)})_J=-\widehat K_{J,J}^{-1}\widehat b_J,\qquad
 (x^{(0)})_{J^c}=0.
\]
The source--cokernel identity and the unchanged kernel imply that this solves
all top rows: its residual vanishes on $J$, lies in
$(\operatorname{Ran}N_*)^\perp$, and therefore has zero remaining coordinates,
since $N_{*,J^c,:}=I$. For arbitrary $\eta\in\mathbb R^{11}$ put
\begin{equation}
 \begin{split}
 x(Y,\eta)&=x^{(0)}+N_{*,V}\eta+N_{*,U}y,\\
 y&=-H(Y)^{-1}\left\{
       \begin{pmatrix}\mathsf m x^{(0)}\\t_3(Y)+R(Y)x^{(0)}\end{pmatrix}
                    +C_{7,:,V}(Y)\eta\right\}.
 \end{split}
 \label{v26:exact-chart}
\end{equation}
Every added column is an actual Poisson-null column, so all 108 top rows remain
zero. The last formula solves all four normalization and three cubic-mean
rows exactly. Conversely every top solution is $x^{(0)}+N_*z$, and its last
seven rows solve uniquely for $z_U$ in terms of $z_V$. This proves that
\eqref{v26:exact-chart} describes every lower solution on this hypersurface.
All its functions are rational and analytic where the two named minors and
$d_2+d_3$ are nonzero. Center $\eta$ at the coordinates of $x_*$ to obtain
an analytic chart through the retained solution. These fibre coordinates and
V25's eleven free coordinates differ by an invertible linear change at the
base and hence by an analytic coordinate change nearby.

It remains to exclude a nearby lower branch off $\alpha_1=\alpha_3$.
V25 eliminated the same 97 selected rows and used one unselected original
row, with nonzero transverse derivative $v_{13}$, to solve uniquely for
$\alpha_1=\chi(\alpha_2,\alpha_3,b_2,b_3,\eta)$.
On $\alpha_1=\alpha_3$, \eqref{v26:exact-chart} solves every row for each of
these nearby free parameters. Uniqueness in that implicit-function chart
therefore forces $\chi=\alpha_3$. The remaining seventeen rows are already
zero on that chart, which proves \eqref{v26:residuals-zero}. Thus there is no
hidden second-order compatibility condition left in the lower system locally.

The ambient dimension is $5+108=113$; the selected 98-row derivative remains
of rank 98 after reducing the neighborhood. The explicit fifteen-dimensional
solution chart implies rank at most 98 on it, so its rank is exactly 98.
This proves the manifold assertion rather than merely its tangent dimension.
The inherited common scaling has one intersection with $\alpha_3=1$ near
the positive base and removes one parameter without deleting a solution orbit.
\end{proof}

\section{All weak-null tests are fixed, and the prepared branches extend locally}
\label{v26:weak-frame}

\begin{corollary}[The complete eight-component nonconstant tangency gate]
\label{v26:weak-gate}
On the neighborhood in \cref{v26:lower-manifold}, the weak Poisson-null space
is exactly the three constant shifts together with gradients of
\begin{equation}
             f_1,s_1,f_2,s_2,f_3,s_3,f_1s_3,s_1f_3.
 \label{v26:eight-potentials}
\end{equation}
All eleven tests are independent of the canonical parameters. For the V22
tangency functional, the three constant-shift components vanish identically
on the compatible first fields. Hence its complete weak-null condition is
exactly eight scalar equations in the fixed frame \eqref{v26:eight-potentials}.
\end{corollary}
\begin{proof}
Intersect the complete kernel \eqref{v26:explicit-kernel} with the space of
pure curl-free shifts. Independence of the seven lapse functions leaves
precisely the eleven displayed shift columns. For the constant tangency
rows, $Dd_{2,W}(Y)[Z,c]=0$ for every $Z$ by the inherited constant-drift
identity. Also
$\mathbb K_1(Z)[c,\ell]=\langle\delta_c\ell,LZ\rangle_h$ by
\eqref{v21:translation-incidence}. Here $Z=F_1Y+G\ell$ has $LZ=0$:
$LG=0$ is the flat bracket identity, and the inherited linear consistency
formula gives $LF_1Y=0$ whenever $LY=0$. This proves the constant rows vanish
without assuming tangency of a nonlinear history. All remaining independent
weak-null tests are precisely \eqref{v26:eight-potentials}.
\end{proof}

\begin{corollary}[Local canonical-parameter extension of the exact initial preparations]
\label{v26:initial-extension}
After reducing the neighborhood of the lower manifold, its canonical and
multiplier parameters admit the exact initial preparations of V21, with
\[
 \begin{gathered}
 X(a_0,0)=a_0Y+O(a_0^2),\qquad
 \lambda(a_0,0)=e+a_0\ell+O(a_0^2),\qquad P=0,\\
 \mathbb P_{\rm cf}\lambda_t^\beta(a_0,0)=0,\qquad
 \lambda_t(a_0,0)=O(a_0).
 \end{gathered}
\]
Here $a_0$ is the preparation amplitude, distinct from the coefficient $a$
in \eqref{v26:linear-family}. The original initial metric and link curvatures
are $O(|a_0|)$. On a compact smaller parameter neighborhood the constants
and initial amplitude radius can be chosen uniformly at this fixed cutoff.
No physical evolution interval independent of $a_0$ or $h$ is asserted.
\end{corollary}
\begin{proof}
The actual mean-calibration derivative, full leading signed-source determinant,
and constrained second-field control minor are nonzero at the V21 base.
They are analytic in the canonical and tilt parameters and remain nonzero
nearby. The common null matrix proved here gives the quadratic-preserving
first-tilt directions; the seven normalization/mean controls remain onto.
The admissible second-field spaces vary analytically by their original
constraint minors. Choose the same nonzero minors on a smaller neighborhood
to obtain analytic local right inverses, rather than freezing an inadmissible
second-field variation.

The triangular implicit-function argument of \cref{v21:main} now applies
with these lower-manifold parameters as external parameters. Its absolute
cubic weak baseline remains the original analytic right side; no value for
that baseline or a quartic action truncation is inserted. The existing onto
controls solve it and then make the entire actual initial weak velocity zero.
The restricted active equation then gives the $O(a_0)$ rate, and the original
metric/link argument gives the curvature estimates. Compactness of a smaller
finite parameter box supplies common initial bounds. This is an extension
of the inherited initial theorem to actual canonical variations, not a
new claim that its prepared set is invariant under the exact flow.
\end{proof}

\section{A finite boost deformation crosses the diamond equation}
\label{v26:crossing}

The preceding chart permits an actual one-parameter lower-incidence curve.
Keep $\alpha=(1,9/8,1)$ and set
\begin{equation}
 d_2=-8+z,\qquad d_3=-8-z,\qquad
 d_1=\frac{16175}{4096}-\frac{z^2}{16}.
 \label{v26:line}
\end{equation}
The mean equation is exact for all $z$. Write $K(z),b(z),R(z),t_3(z)$ for
the original rationalized coefficients on this line. Let $x_*$ be the retained
V21 tilt and use the same $J,N_{*,U}$ as above. Define
\begin{align}
 u(z)&=x_*-I_JK_{J,J}(z)^{-1}(K(z)x_*+b(z))_J,\notag\\
 H(z)&=\begin{pmatrix}\mathsf mN_{*,U}\\R(z)N_{*,U}\end{pmatrix},\notag\\
 x(z)&=u(z)-N_{*,U}H(z)^{-1}
                \begin{pmatrix}\mathsf m u(z)\\t_3(z)+R(z)u(z)\end{pmatrix}.
 \label{v26:line-tilt}
\end{align}
This fixes the unused null coordinates; it is not a minimum-norm rule.
Where its two inverses exist, the common kernel and source identities show
that it solves all 115 lower equations. In particular $x(0)=x_*$.

Let $\widehat\Delta(z)$ be the diamond tangency functional evaluated by the
original rational coefficient functions at this field and tilt. The physical
spatial-sum functional on the unscaled weak diamond test is
$\sqrt3\,\widehat\Delta(z)$. To verify this conversion, the un-tilted weak
quadratic drift has only field types $U^2$ and $p^2$. Under
$p\mapsto\sqrt3p_R$, $\ell^\beta\mapsto\sqrt3x^\beta$, its configuration
velocity scales by $\sqrt3$, while its momentum velocity does not. Its
polarized derivative therefore scales by $\sqrt3$. In the first bracket,
the shift--lapse block is linear in the configuration velocity and the
shift--shift block is linear in the momentum velocity and multiplied by the
shift tilt. Both terms have the same factor. Thus the zero tested here is a
zero of the actual physical component, not of an unrelated rescaled forcing.

\begin{theorem}[A regular lower-incidence zero of one physical tangency component]
\label{v26:diamond-crossing}
There exists
\begin{equation}
            z_\diamond\in
       \left(-\frac{220629}{10^7},-\frac{220628}{10^7}\right)
 \label{v26:root-interval}
\end{equation}
for which \eqref{v26:line}--\eqref{v26:line-tilt} solve every lower-incidence
equation and the physical diamond tangency component is zero.
The Poisson rank is ninety and the seven lower null controls are regular
throughout this interval. The leading active source there remains
trace-free of rank 78 with positive Cartan Gram. No uniqueness of this
scalar zero and no vanishing of the other seven independent tangency
components is asserted.
\end{theorem}
\begin{proof}
The exact endpoint and interval certificate is described in the next section.
It establishes that both matrices inverted in \eqref{v26:line-tilt} are
nonsingular on the full closed interval and that
\begin{equation}
    100<\widehat\Delta(-220629/10^7)<110,\qquad
    -45<\widehat\Delta(-220628/10^7)<-35.
 \label{v26:endpoint-signs}
\end{equation}
The formula is therefore continuous on the interval, and the intermediate
value theorem gives \eqref{v26:root-interval}. The complete lower equations
hold at every point by the proved source--cokernel identity and exact
seven-row solve. Nonzero Poisson minor and the common kernel give rank ninety.
The seven-control determinant is separately retained.

For the active-source claim, the configuration remains exactly the previously
certified $\alpha=(1,9/8,1)$. Its mean-zero diagonal lapse observation
therefore has rank 26. The new boost is still diagonal and its $i$th entry
has no oscillation in coordinate $i$. Throughout the interval its normalized
first diagonal entry is greater than two and the other two are less than
minus two, with magnitudes less than ten, including the oscillatory
complementary momenta. These bounds follow directly from
$|z|<3/100$, $\max_i\alpha_i=9/8$ and \eqref{v26:line}.
The shifted pointwise determinant proof of \cref{v12:T-inverse} applies
unchanged under precisely those bounds and gives transverse rank 52.
The diagonal and off-diagonal source arguments then give rank $26+52=78$.
Trace-freeness uses the same diagonal-rotation/diagonal-boost coefficient
identity, so the original Cartan form is positive on this active range.
No conclusion about the higher 29-column weak source follows from that
active calculation.
\end{proof}

\section{Exact endpoint values and a no-pole certificate}
\label{v26:interval-certificate}

The floating crossing is used only to propose the rational interval
\eqref{v26:root-interval}. All statements of \cref{v26:diamond-crossing}
use the following exact checks. At each rational endpoint the original
$P_2,\Phi_3,P_{c,3}$ coefficient routines regenerate $b$, all needed cubic
means and controls, and the actual diamond polynomial. Exact rational
linear solves define $x(z)$. The complete 108 quadratic rows, all four
normalizations and all three original cubic mean evaluations are checked
as rational zeros. The two resulting fractions obey
\eqref{v26:endpoint-signs}; the full fractions, rather than their decimal
roundings, are stored in the package.

Here is the uniform nonsingularity argument. Put $z_0=-220629/10^7$ and
$\delta=10^{-7}$. The ninety-dimensional principal block has the exact
quadratic form
\[
 K_J(z)=K_J(z_0)+(z-z_0)K'_J(z_0)+(z-z_0)^2K_{J,2}.
\]
With its exact rational inverse $P_J=K_J(z_0)^{-1}$, calculate the induced
infinity norms of $P_JK'_J(z_0)$ and $P_JK_{J,2}$. The certified inequality is
\begin{equation}
 \rho_J:=\delta\|P_JK'_J(z_0)\|_\infty+
                 \delta^2\|P_JK_{J,2}\|_\infty<\tfrac12.
 \label{v26:rho-J}
\end{equation}
Thus $P_JK_J(z)$ stays within distance $\rho_J$ of the identity throughout
the interval and is invertible by its norm-convergent Neumann series.

The seven-dimensional matrix is a polynomial
$H(z)=\sum_{j=0}^4H_jz^j$. The degree bound is exact: $R(Y)$ is homogeneous
quadratic in the first fields, and those fields have degree two along
\eqref{v26:line}. Its coefficients are reconstructed over $\mathbb Q$ by
five evaluations at $z=-2,-1,0,1,2$. These evaluations are polynomial
coefficient tests only; no physical preparation is claimed at those distant
values. The complete cubic baseline is subtracted before forming the control
matrix. Independent evaluations at additional rational line points reproduce
the polynomial entry by entry.

Let $P_H=H(z_0)^{-1}$ and $m_0=220629/10^7$. The factorization of
$z^j-z_0^j$ gives the interval bound
\begin{equation}
 \rho_H:=\delta\sum_{j=1}^4j m_0^{j-1}\|P_HH_j\|_\infty<\tfrac12.
 \label{v26:rho-H}
\end{equation}
This proves nonsingularity of $H(z)$ on the same interval. The matrices,
exact row-sum norms and rational comparisons in
\eqref{v26:rho-J}--\eqref{v26:rho-H} are supplied and rechecked by the
portable verifier. There is no inferred spectral gap from floating singular
values and no assumption that two nonsingular endpoints exclude an intervening
pole. No assertion is made that the root interval lies inside the unspecified
full signed-source neighborhood used in \cref{v26:initial-extension}; the
minors certified here have their separately stated lower/active meanings.

\section{Consequences and the remaining full tangency problem}
\label{v26:status}

The nonlinear lower compatibility problem left in V25 is now locally closed:
its seventeen putative higher-order residuals are identities, not hidden
constraints. The common Poisson kernel has an explicit fixed basis on that
branch, so the complete eight-component nonconstant tangency gate can be
assembled without a changing null-bundle calculation. All four canonical
directions of V25 integrate; after removing scale there are three canonical
shape variables and eleven remaining tilt variables.

The finite crossing proves more than a nonzero derivative of the diamond
component. It supplies actual mean-compatible canonical coefficients and
exact lower tilts at which that component vanishes, with no pole in the
certified crossing interval. This does not solve the complete gate. The other
seven independent tangency rows must be imposed, and any proposed full root
needs its own higher weak signed-source and preparation-control verification.
The local preparations near the old base and the scalar root interval have
been kept as distinct domains rather than identified without a bound.

Nothing here proves preservation of the initial weak-rate condition over a
fixed physical interval, or supplies refinement-uniform first-jet selection,
higher weak rank, differentiated geometric bounds, or a common lifespan.
The raw-action Einstein continuum theorem remains unresolved. The next
well-defined task is the complete eight-row tangency system on the actual
scale-fixed lower manifold, not another expansion of the now eliminated
seventeen residuals or another tilt-only search at the excluded old first jet.

\paragraph{Verification and preservation.}
The exhaustive source certificate consists of five original Poisson
coefficient matrices and fifteen complete source--cokernel monomials over
$\mathbb Q$. The explicit kernel functions are checked against the retained
point-coordinate kernel, and its independence is proved above. The line
certificate regenerates original endpoint source values, exact lower solves,
cubic means and diamond contractions; the polynomial control matrix and
both no-pole inequalities are verified with arbitrary-precision fractions.
Floating probes are labelled as discovery calculations only. The packet
contains a portable exact verifier and separate regeneration commands.
All twenty-five preceding marked mathematical bodies, including their
corrections and qualifications, are retained byte for byte. These checks are
not proof-assistant formalization, an independent validation of every
inherited analytic theorem, or numerical integration of an exact-action
spacetime.
% END IMMUTABLE EXACT-ACTION BRIDGE V26 BODY


% Future iterations append here.
% BEGIN IMMUTABLE EXACT-ACTION BRIDGE V27 BODY
\clearpage
\renewcommand{\thesection}{\arabic{section}}
\renewcommand{\theHsection}{bridgeV27.\arabic{section}}
\setcounter{section}{218}
\section{The complete tangency system, rather than another scalar crossing}
\label{v27:context}

The lower-incidence manifold in V26 is an actual analytic solution manifold,
not merely a tangent approximation. Its scale-fixed dimension is fourteen:
three canonical shape coordinates and eleven multiplier coordinates. That
count does not say that fourteen independent controls act on the eight
nonconstant weak-null tangency equations. The present body calculates their
actual action. Six equations are completely independent of the multiplier
fibre. On the original lower equations the other two change in only one
common direction. Consequently seven canonical compatibility functions,
not eight freely controllable equations, govern this local first-tangency
selection problem.

This is an exact coefficient result at the retained three-site cutoff. Its
proof uses every original prepared quadratic row and the original cubic
mean equations. We neither project those equations nor choose a new
Hamiltonian. The eight quantities are still the differentiated weak targets
of V22, and not a claim of full invariance of an initial preparation under
the nonlinear evolution. The one-component zero proved in V26 remains a
one-component zero; its status in the full system is tested separately below.

\paragraph{Inputs and novelty boundary.}
We retain the original polynomial envelopes \eqref{v20:P2} and
\eqref{v20:Phi3}, their physical rationalization in
\eqref{v20:rational-system}, the complete weak derivative of V21, and the
V22 preparation-independent tangency incidence. V26 supplies the fixed
Poisson kernel, the nonlinear lower solution manifold, its explicit weak
frame, and the certified narrow diamond-root interval. These results are
inputs, not new claims here. The new results are the full fibre transformation
law, the two original-row identities that reduce its rank, and the resulting
seven-function quotient. Exact rational coefficient identities are supplied
in the accompanying package. Floating searches are not used to decide their
zeros, ranks, or solution sets.

\section{The coefficient family and its eight original tangency rows}
\label{v27:setup}

Throughout this body $N=3$, $h=1/3$, and $V=27$. Put
\[
 f_i=\cos(2\pi x_i),\qquad s_i=\frac{2}{\sqrt3}\sin(2\pi x_i).
\]
On the three sites in coordinate $i$, their values are respectively
$(1,-1/2,-1/2)$ and $(0,1,-1)$. Products are the original cyclic products,
and $\delta_i=3(S_i-S_i^{-1})$. In particular
\begin{equation}
 \delta_i f_i=-\frac92s_i,\qquad \delta_i s_i=6f_i.
 \label{v27:trig}
\end{equation}
Write the five independent linear field coefficients as
\[
 q=(a,c,d_1,d_2,d_3),\qquad \alpha=(a,c,a),
\]
with the original complementary fields
\begin{equation}
 \begin{split}
 U(q)&=\sum_i\alpha_iT_if_i,\\
 p(q)&=\sqrt3\,p_R(q),\qquad
 p_R(q)=3\{\diag(d)- (d_1+d_2+d_3)I\}
                         +\frac32\sum_i\alpha_iP_if_i.
 \end{split}
 \label{v27:fields}
\end{equation}
For exactly preparable first jets impose the original mean hypersurface
\begin{equation}
 d_1(d_2+d_3)+d_2d_3=(2a^2+c^2)/4.
 \label{v27:mean}
\end{equation}
The coefficient calculations below initially allow the five $q$'s to vary
independently. This larger linear carrier is only an algebraic device for
exhausting the polynomial identities; its off-hypersurface points are not
called exact initial records.

A rational multiplier vector $x\in\R^{108}$ has component-major order,
and synthesizes the physical tilt $\ell=Q_3x$,
$Q_3=\diag(I_{27},\sqrt3I_{81})$. Let $\widehat K(q)$ and
$\widehat b(q)$ be exactly the rationalized original first Poisson matrix
and prepared target of V20--V26. Define the eight tests by
\begin{equation}
 \nu_i=(0,\nabla_\delta\psi_i),\qquad
 (\psi_1,\ldots,\psi_8)
       =(f_1,s_1,f_2,s_2,f_3,s_3,f_1s_3,s_1f_3).
 \label{v27:weak-frame}
\end{equation}
The three constant weak tests have identically zero tangency as in V26.
The eight tests in \eqref{v27:weak-frame} and those constants exhaust the
common weak Poisson kernel on the regular lower branch.

Use $\mathsf T_i(q,x)$ for the spatial-\emph{sum} tangency functional on
$\nu_i$, divided by the inherited nonzero physical factor $\sqrt3$:
\begin{equation}
 \sqrt3\,\mathsf T_i(q,x)=
 Dd_{2,W}(Y(q))[F_1Y(q)+GQ_3x,\nu_i]
 +\mathbb K_1(F_1Y(q)+GQ_3x)[\nu_i,Q_3x].
 \label{v27:T}
\end{equation}
The physical spatial-mean value is therefore $\sqrt3\mathsf T_i/27$.
Only this fixed normalization is removed; no equation is changed by it.
Each $\mathsf T_i$ is a homogeneous polynomial of total degree two in
$(q,x)$. Indeed the weak drift is quadratic in $Y$, its field differential
is linear, the first Poisson coefficient is linear in its field slot, and
$F_1Y+G\ell$ is linear in $(Y,\ell)$. The field-type calculation in the
inherited rationalization makes \eqref{v27:T} a polynomial over $\mathbb Q$.

Let $\mathsf N$ denote the following fixed $108$-by-$18$ rational kernel
frame, with the indicated column order:
\begin{equation}
 \begin{split}
 &(1,0),\ (f_1,0),(s_1,0),(f_2,0),(s_2,0),(f_3,0),(s_3,0),\\
 &(0,e_1),(0,e_2),(0,e_3),\ \nu_1,\ldots,\nu_8.
 \end{split}
 \label{v27:N}
\end{equation}
The notation here lists rational multiplier coordinates; physical synthesis
still uses $Q_3$. V26 proves $\widehat K(q)\mathsf N=0$ throughout
\eqref{v27:fields}. On its regular region this is the entire Poisson kernel.
Write a kernel coefficient as $\eta\in\R^{18}$, with lapse coordinates
$\eta_{f_i},\eta_{s_i}$, and define
\begin{equation}
                         \Lambda(\eta)=\eta_{f_1}+\eta_{f_3}.
 \label{v27:Lambda}
\end{equation}

\section{An exact fibre law for every one of the eight equations}
\label{v27:fibre}

In the following four linear functionals, $\sum_x$ is the unnormalized sum
over the original $27$ sites and $x^{\beta j}$ is the $j$th rationalized
shift component of $x$. Set
\begin{align}
 L_1(x)&=\sum_x\{f_1s_3x^{\beta1}-s_1f_3x^{\beta3}\},\notag\\
 L_2(x)&=\sum_x\{s_1f_3x^{\beta1}-f_1s_3x^{\beta3}\},\label{v27:Lmoments}\\
 A(x)&=\sum_x\{\tfrac34s_1s_3x^{\beta1}+f_1f_3x^{\beta3}\},\notag\\
 B(x)&=\sum_x\{f_1f_3x^{\beta1}+\tfrac34s_1s_3x^{\beta3}\}.
 \label{v27:ABmoments}
\end{align}

\begin{theorem}[Complete original tangency transformation on the common kernel]
\label{v27:fibre-theorem}
For every real $q$ in the linear family \eqref{v27:fields}, every multiplier
$x$, and every $\eta\in\R^{18}$, the exact identities are
\begin{align}
 \mathsf T_i(q,x+\mathsf N\eta)&=\mathsf T_i(q,x),
                                      &&1\le i\le6,\label{v27:six}\\
 \mathsf T_7(q,x+\mathsf N\eta)-\mathsf T_7(q,x)
 &=6561\{\Lambda(\eta)A(x)+\eta_{s_1}L_1(x)
                                  +\eta_{s_3}L_2(x)\},\label{v27:seven}\\
 \mathsf T_8(q,x+\mathsf N\eta)-\mathsf T_8(q,x)
 &=6561\{\Lambda(\eta)B(x)-\eta_{s_1}L_2(x)
                                  -\eta_{s_3}L_1(x)\}.\label{v27:eight}
\end{align}
Each of $L_1,L_2,A,B$ annihilates $\Ran\mathsf N$. In particular there
is no quadratic term in $\eta$ in these formulas. Their validity does not
assume that $x$ has yet solved the lower equations.
\end{theorem}
\begin{proof}
We describe the exhaustive rational coefficient proof, including its
normalization and degree bound. The underlying scalar coefficient is
exactly \eqref{v20:P2}. Work first with rational field and multiplier
coordinates; the field-type scaling leading to \eqref{v27:T} puts all
physical factors into the same rational formula. Let $\mathbf e_k$ be
the $108$ standard multiplier columns. For the $j$th column $n_j$ of
\eqref{v27:N}, form the matrix
\[
 (Q_j)_{\mu k}=DP_{\mathbf e_\mu,2}(Gn_j)[G\mathbf e_k].
\]
Here $G$ is evaluated on the rational coordinates by its original commuting
phase derivatives. Every matrix entry is computed as
\[
 \frac{P_{\mathbf e_\mu,2}(Gn_j+G\mathbf e_k)
              -P_{\mathbf e_\mu,2}(Gn_j-G\mathbf e_k)}2.
\]
This is exact polarization of a quadratic polynomial, not a numerical
difference quotient. For $w_i=\nu_i$ in these coordinates, put
\begin{equation}
 (C_i)_{j,:}=2w_i^TQ_j-(Q_{w_i}n_j)^T-n_j^TQ_{w_i}.
 \label{v27:C-algebra}
\end{equation}
Here $Q_{w_i}$ is defined by the same displayed derivative with the kernel column $n_j$ replaced by $w_i$.

For clarity about why this is the needed coefficient, at $q=0$ the tangency
polynomial is
\[
 \mathsf T_i(0,x)=DP_{w_i,2}(Gx)[Gx]-DP_{x,2}(Gx)[Gw_i].
\]
Its symmetric cross term in $x,n_j$ is exactly \eqref{v27:C-algebra}.
The coefficient with one canonical slot $Y$ and one kernel tilt is
\[
 Dd_{2,W}(Y)[Gn_j,w_i]+\mathbb K_1(F_1Y)[w_i,n_j].
\]
It is linear in $Y$, so five field basis vectors exhaust it. The exact
coefficient files verify all $5\cdot8\cdot18=720$ such entries are zero.
The cubic scalar envelope \eqref{v20:Phi3} is retained in this evaluation.

The remaining $C_i$ matrices are computed in full, with $18\cdot108$
entries each. For the first six tests every entry is zero. For test seven
the only nonzero rows are
\[
 (C_7)_{f_1}=(C_7)_{f_3}=6561A,\qquad
 (C_7)_{s_1}=6561L_1,\qquad(C_7)_{s_3}=6561L_2,
\]
where the symbols on the right denote the row vectors of the explicitly
displayed functionals. For test eight they are
\[
 (C_8)_{f_1}=(C_8)_{f_3}=6561B,\qquad
 (C_8)_{s_1}=-6561L_2,\qquad(C_8)_{s_3}=-6561L_1.
\]
These are exact identities over $\mathbb Q$. They can also be checked in
the complete tensor Fourier basis with one-dimensional factors $1,f,s$:
\eqref{v27:trig}, finite cyclic multiplication, and orthogonality exhaust
all $27$ scalar array coordinates. The package supplies all eighteen $Q_j$
matrices, the cross matrices, and the mixed-field arrays, together with
regeneration from the original coefficient functional. It uses no modular
zero test to infer any of these vanishing identities.

Direct pairing of the four displayed rows with \eqref{v27:N} gives zero.
For example their possible nontrivial pairings with the two mixed gradients
cancel using $\sum f^2=3/2$, $\sum s^2=2$, and
\eqref{v27:trig}; all other terms have a zero independent-coordinate mean.
Thus $C_i\mathsf N=0$ for all eight tests. The pure quadratic restriction
$\mathsf T_i(0,\mathsf N\eta)$ is consequently zero as well. Combining
these facts with the vanishing mixed-field coefficients and the proved
total degree two gives \eqref{v27:six}--\eqref{v27:eight} for arbitrary
real coefficients, not only rational sample points.
\end{proof}

\section{Two original weak rows remove two apparent tilt controls}
\label{v27:range-moments}

The last two equations in \cref{v27:fibre-theorem} initially involve three
lapse combinations. Two disappear on the original quadratic solution set,
not by assuming a favorable first-tilt choice.

\begin{lemma}[Exact source incidence for the two shift moments]
\label{v27:moment-rows}
Let $w_1,w_2\in\R^{108}$ represent the rows $L_1,L_2$ in
\eqref{v27:Lmoments}. Define the two original curl-free multiplier tests
\[
 u=(0,\nabla_\delta(f_1f_3)),\qquad
 v=(0,\nabla_\delta(s_1s_3)).
\]
Throughout the five-dimensional linear field family,
\begin{equation}
 \boxed{\widehat K(q)u=-\frac{6561}{4}a\,w_2,
        \qquad \widehat K(q)v=-2187a\,w_1.}
 \label{v27:row-incidence}
\end{equation}
For $a\ne0$, every solution of the \emph{complete original} equation
$\widehat K(q)x=-\widehat b(q)$ consequently obeys
\begin{equation}
                         \boxed{L_1(x)=L_2(x)=0.}
 \label{v27:moment-zero}
\end{equation}
\end{lemma}
\begin{proof}
The first Poisson coefficient is linear in \eqref{v27:fields}. Substitution
of the two tests in \eqref{v20:K1}, using the original quadratic envelope,
gives the two nonzero $a$ coefficients displayed in
\eqref{v27:row-incidence}. Its other four linear coefficients are zero.
The package checks both complete $108$-component identities against the
five original Poisson coefficient matrices. Equivalently they follow by
\eqref{v27:trig} and the same finite tensor Fourier contraction used above.
The tests are not asserted to be Poisson-null; their nonzero images are
precisely what is needed.

The complete weak quadratic target is zero on the complementary family,
independently of its constant momentum, by \cref{v12:matching}. Hence
$u^T\widehat b=v^T\widehat b=0$. Skewness gives, for example,
\[
 -\tfrac{6561}{4}a\,L_2(x)
       =(\widehat Ku)^Tx=-u^T\widehat Kx=u^T\widehat b=0.
\]
The other identity gives $L_1(x)=0$. This uses all original constraint
rows, not an orthogonal projection of the target. A separate exact check
in the package evaluates these zero source pairings on all fifteen
quadratic field monomials.
\end{proof}

Set $\mathsf A_7=6561A(x)$ and $\mathsf A_8=6561B(x)$. Since the four
moments annihilate the common kernel, these numbers depend only on the
canonical parameters whenever the quadratic range equation is soluble and
its full kernel is \eqref{v27:N}. On that set the entire transformation is
\begin{equation}
 \boxed{\mathsf T(q,x+\mathsf N\eta)
   =\mathsf T(q,x)+\Lambda(\eta)
                  (0,0,0,0,0,0,\mathsf A_7,\mathsf A_8)^T.}
 \label{v27:rank-one-law}
\end{equation}
This conclusion is stronger than a numerical Jacobian rank at the base
point: it is a finite difference identity for arbitrary kernel displacement.

\section{Seven canonical compatibility functions and one effective fibre coordinate}
\label{v27:quotient}

On the regular lower manifold retain the original cubic mean map
$t(q)+R(q)x=0$ and the four multiplier means $\mathsf m x=0$. Its fibre
at a fixed canonical point is an affine translate of
\begin{equation}
 \mathcal F_q=\{\mathsf N\eta:\ \mathsf m\mathsf N\eta=0,
                                           \ R(q)\mathsf N\eta=0\}.
 \label{v27:Fq}
\end{equation}
It has dimension eleven near the base by V26. The next certificate shows
that \eqref{v27:rank-one-law} has rank exactly one there, not merely at most
one.

\begin{lemma}[A normalized, cubic-mean-preserving unit gain]
\label{v27:unit-gain}
At the base
\[q_*=(1,9/8,16175/4096,-8,-8),\]
set all kernel coefficients to zero except
\begin{equation}
 \eta_{f_1}=1,\quad \eta_{s_1}=\frac{69355}{382422},\quad
 \eta_{f_2}=\frac{6602752}{46464273},\quad
 \eta_{s_2}=\frac{16572416}{46464273}.
 \label{v27:eta-unit}
\end{equation}
Then $\mathsf N\eta\in\mathcal F_{q_*}$ and $\Lambda(\eta)=1$.
The same unit-gain condition has an analytic solution on a neighborhood
of $q_*$ on the canonical mean hypersurface. Moreover
$\mathsf A_7(q_*)\ne0$ and $\mathsf A_8(q_*)\ne0$.
\end{lemma}
\begin{proof}
The four selected columns have zero means. On the six nonconstant lapse
columns of \eqref{v27:N}, stack the three actual rows $R(q_*)$ and the
row $\Lambda$. The minor on $(f_1,s_1,f_2,s_2)$ is
\begin{equation}
 \begin{pmatrix}
 0&-19683/2&-819068679/262144&819068679/131072\\
 111537/16&-111537/8&0&-1594323/128\\
 -98415/16&98415/8&-7026831/1024&7026831/512\\
 1&0&0&0
 \end{pmatrix}.
 \label{v27:gain-minor}
\end{equation}
Its exact determinant is
$-354317843330549369451/268435456$, with nonzero residue $87$ modulo
$241$. Multiplication of this matrix by the four entries in
\eqref{v27:eta-unit} gives $(0,0,0,1)^T$ over $\mathbb Q$.
The remaining columns of the coefficient vector are zero, so this verifies
the complete assertion. The inverse persists analytically near the base,
and its right side remains $(0,0,0,1)^T$.

Using the exactly specified $x_{21}$ as one lower solution at the base,
the package computes both gain entries as exact fractions. Their good-prime
residues are
\begin{equation}
              \mathsf A_7(q_*)=56\pmod{241},\qquad
              \mathsf A_8(q_*)=129\pmod{241}.
 \label{v27:gain-residues}
\end{equation}
Every denominator is nonzero at that prime. These residues prove
characteristic-zero nonvanishing; zeros used elsewhere in the argument were
instead checked over $\mathbb Q$. Continuity keeps both entries nonzero
on a smaller canonical neighborhood.
\end{proof}

Choose any analytic local section $x_0(q)$ of V26's lower manifold and define
\begin{equation}
 \begin{split}
 I_i(q)&=\mathsf T_i(q,x_0(q)),\qquad 1\le i\le6,\\
 I_7(q)&=\mathsf A_8(q)\mathsf T_7(q,x_0(q))
                         -\mathsf A_7(q)\mathsf T_8(q,x_0(q)).
 \end{split}
 \label{v27:invariants}
\end{equation}
All seven functions are independent of the section by
\eqref{v27:rank-one-law}. They are analytic on the regular lower domain.
In fact they may be evaluated on any analytic section solving only
$\widehat Kx=-\widehat b$. Two such sections differ by the same complete
Poisson kernel. The two forced moment identities still hold, and
\eqref{v27:rank-one-law} leaves all seven expressions invariant. The lower
mean and normalization equations are not discarded from the physical
selection: V26 supplies their separate solutions. This observation merely
avoids solving them again when evaluating their already invariant canonical
compatibility functions.

\begin{theorem}[Exact local reduction of all eight first-tangency conditions]
\label{v27:quotient-theorem}
On a sufficiently small neighborhood of the base canonical point, the complete
lower equations and all eight first-tangency equations admit a multiplier
solution at $q$ if and only if
\begin{equation}
                         \boxed{I_1(q)=\cdots=I_7(q)=0.}
 \label{v27:seven-conditions}
\end{equation}
If these conditions hold, the admissible multipliers form an affine space
of dimension ten in the eleven-dimensional lower fibre. In particular six
of the original tangency equations are completely canonical; the two
remaining equations use only one effective multiplier coordinate.

After fixing common scale, \eqref{v27:seven-conditions} is a system of
seven analytic functions of three canonical shape variables. This statement
does not assert independence of those seven functions, the existence of a
common zero, or the nonexistence of such a zero away from the base.
\end{theorem}
\begin{proof}
The six equations are independent of the fibre by \eqref{v27:six}.
Their vanishing and $I_7=0$ are necessary if all eight rows vanish.
For sufficiency let $\eta_*(q)$ be the analytic unit-gain coefficient from
\cref{v27:unit-gain}. On the smaller neighborhood
$\mathsf A_7\ne0$, set
\[
 x=x_0-\frac{\mathsf T_7(q,x_0)}{\mathsf A_7(q)}
                                             \mathsf N\eta_*(q).
\]
It remains an exact lower solution. Formula \eqref{v27:rank-one-law}
sets row seven to zero, and the residual in row eight is zero exactly when
$I_7=0$. The first six rows were not changed. Any other solution differs
by a fibre vector in $\ker\Lambda$. Since \cref{v27:unit-gain} makes
$\Lambda$ onto on the eleven-dimensional fibre, this kernel has dimension
ten. This proves both directions without a full-rank assumption on the
eight-by-fourteen tangency Jacobian.

Common scaling $Y\mapsto c_0Y$, $\ell\mapsto c_0\ell$, $c_0>0$, preserves
the homogeneous coefficient equations as proved in V25. The six $I_i$ scale
by $c_0^2$, while the gains scale by $c_0$ and $I_7$ by $c_0^3$.
Their zero set is therefore scale invariant. Locally take $a=1$ on
\eqref{v27:mean}; the three shape variables are $c,d_2,d_3$ and $d_1$
is its displayed rational mean solve. No physical-time scaling is made.
\end{proof}

As an additional exact base regression, the seventh expression in
\eqref{v27:invariants} has residue $44$ modulo $241$ at $x_{21}$, with a
unit denominator. Its full rational fraction is recorded. The base is not
a common zero, consistently with the earlier tangency obstruction. This
new invariant does not turn equation counting into a nonexistence theorem
for other canonical shapes.

\section{The certified diamond crossing is not a zero of the full system}
\label{v27:root-window}

The invariant axial components permit an exact full-system test of the
one-component crossing in V26. Retain its unchanged canonical curve
\begin{equation}
 a=1,\quad c=9/8,\quad d_2=-8+z,\quad d_3=-8-z,\quad
 d_1=\frac{16175}{4096}-\frac{z^2}{16},
 \qquad z\in I,
 \label{v27:curve}
\end{equation}
where
\[
 I=\left[-\frac{220629}{10^7},-\frac{220628}{10^7}\right].
\]
V26 proves that the mean hypersurface is satisfied, that the full lower
incidence has an analytic solution throughout $I$, and that the diamond
component has at least one zero there. Its Poisson and seven-control
inverses were certified separately on that entire interval.

\begin{theorem}[Full tangency is excluded on the complete scalar-root interval]
\label{v27:root-exclusion}
For every $z\in I$ and every solution $x$ of the complete lower equations
at \eqref{v27:curve}, the component associated with
$\nu_2=(0,\nabla_\delta s_1)$ satisfies
\begin{equation}
            \boxed{-24\,272\,317<\mathsf T_2(q(z),x)<-24\,271\,262<0.}
 \label{v27:root-exclusion-bound}
\end{equation}
The bounds use the rational spatial-sum convention of \eqref{v27:T}.
Thus no point of this canonical interval, with any admissible lower-fibre
tilt, solves the complete eight-component tangency system. In particular
the diamond zero certified in V26 cannot be repaired by adjusting the
remaining multiplier fibre at that canonical point.
\end{theorem}
\begin{proof}
By \eqref{v27:six}, this component is independent of the entire common
Poisson kernel, even before imposing the cubic mean equations. It can
therefore be enclosed using a Poisson-only representative. The existing
lower solutions ensure that this invariant value is also their value;
no initial normalization or mean row is removed from the selection.

Here is the complete verified calculation. Put $z_0=-220629/10^7$,
$t=z-z_0\in[0,10^{-7}]$, and let $x_0$ be V26's exact left-endpoint
lower solution. Let $E_J$ insert the original ninety principal Poisson
coordinates, and $P=\widehat K_{JJ}(z_0)^{-1}$ be the supplied exact
rational inverse. Expand
\[
 \widehat K(z_0+t)=K_0+tK_1+t^2K_2,\qquad
 \widehat b(z_0+t)=\sum_{j=0}^4t^jb_j.
\]
The first expansion follows from the five original linear coefficient
matrices. The second has degree at most four because the prepared target
is quadratic in the canonical first fields, whose dependence on $t$ is
quadratic. Five exact rational evaluations of the original full target
supply its complete coefficients. Its value at the right endpoint also
matches all 108 independently retained endpoint entries.

Define the quadratic multiplier trial $x^{\rm tr}(t)=x_0+tx_1+t^2x_2$ by
\begin{align*}
 x_1&=-E_JP(K_1x_0+b_1)_J,\\
 x_2&=-E_JP(K_1x_1+K_2x_0+b_2)_J.
\end{align*}
Its full residual $r(t)=\widehat K(z_0+t)x^{\rm tr}(t)+\widehat b(z_0+t)$
is a rational polynomial of degree at most four with an exact triple zero
at $t=0$. This is verified on all 108 rows, not only on $J$. It also
follows from the fixed common kernel and full source-range identity.
The exact Poisson solution obtained by adding $E_Jv(t)$ has
\[
 v(t)=-\widehat K_{JJ}(z_0+t)^{-1}r_J(t).
\]
The residual and correction stay in the correct full Poisson range, so
all original quadratic equations are solved by this representative.

V26's exact preconditioned bound gives
\[
 \sup_{t\in[0,10^{-7}]}\|P(\widehat K_{JJ}(z_0+t)-\widehat K_{JJ}(z_0))\|_\infty
                       \le\rho<1/100.
\]
If $v_0(t)=-Pr_J(t)$, the convergent Neumann series gives the rigorous
componentwise enclosure
\[
 v(t)\in v_0([0,10^{-7}])+
 \left[-\frac\rho{1-\rho}\|v_0([0,10^{-7}])\|_\infty,
        \frac\rho{1-\rho}\|v_0([0,10^{-7}])\|_\infty\right]^{90}.
\]
The polynomial-centred residual makes this correction much smaller than
an enclosure obtained by directly interval-evaluating two large nearly
equal source values. The certified result is
\begin{equation}
                       \sup_{t\in[0,10^{-7}]}\|v(t)\|_\infty<4\cdot10^{-10}.
 \label{v27:trial-correction}
\end{equation}

The trial value $T^{\rm tr}(t)=\mathsf T_2(q(z_0+t),x^{\rm tr}(t))$
is again a polynomial of degree at most four, now by total degree two
of \eqref{v27:T}. It is constructed in exact rational arithmetic from
the original scalar coefficient functions. An extra-node evaluation
checks the resulting interpolation independently; the degree bound is
algebraic, not inferred from these evaluations. For the correction use
the exact quadratic identity
\[
 \begin{split}
 \mathsf T_2(q,x^{\rm tr}+E_Jv)
   ={}&T^{\rm tr}
        +D_x\mathsf T_2(q,x^{\rm tr})[E_Jv]
        +\mathsf T_2(0,E_Jv).
 \end{split}
\]
The middle term is evaluated by a first directional jet of the original
quadratic functional, with interval values and interval directions. This
avoids subtracting two dependent large interval values. All endpoints use
100-bit dyadics with Python-integer outward rounding. Every addition,
multiplication, reciprocal, and Taylor-coefficient operation encloses its
exact rational or real operation. No floating root or eigenvalue threshold
enters the argument. The sum of the three resulting intervals is strictly
inside the rational bounds in \eqref{v27:root-exclusion-bound}; all endpoint
fractions, trial coefficients, full residual coefficients and regeneration
code are supplied.

Finally any other solution of the quadratic Poisson equation differs by
$\mathsf N\eta$ on this regular interval. Equation \eqref{v27:six} leaves
its second component unchanged. This includes all solutions satisfying the
additional cubic means and normalizations, proving the stated quantifier
over the entire lower fibre. The interval does contain a diamond-component
zero by the inherited theorem, but its second component is nonzero by this
new calculation. There is no contradiction: the two are different original
tangency observations.
\end{proof}

The theorem excludes this particular canonical interval, not the whole
three-shape domain. It neither requires nor supplies a new higher weak
signed-source rank there. A hypothetical full exact-action branch satisfying
all first-tangency conditions would already have to pass the coefficient
test that fails here. No conclusion about arbitrary initial fields or the
cofinal Einstein limit is inferred.


\section{Direction of continuation and verification boundary}
\label{v27:status}

The new reduction changes how the upstream search should be organized.
The fourteen local lower parameters are not fourteen independent controls
on the eight differentiated weak conditions. Ten multiplier directions are
invisible to all of them; a single further multiplier coordinate moves two
rows in a fixed common direction. The seven section-independent canonical
functions in \eqref{v27:invariants} are the next actual selection problem.
It is not productive to repeat an unconstrained search in the ten invisible
directions, or to treat a zero of one linear combination as a full root.

The exact lower-incidence manifold and its Poisson regularity have not been
lost. Nor does the reduction prove that all seven canonical functions are
independent, that they have no common zero on the existing mean hypersurface,
or that a new polarization is already required by a global theorem. Those
possibilities remain distinct. Exploratory floating solves of the original
coefficient equations did not find a simultaneous zero in the tested regions;
that outcome is not used as a no-go certificate. Candidate points outside a
previously verified rank or preparation chart are not promoted to original
nonlinear branches.

The second independent task also remains: a simultaneous first-tangency zero,
if found, would still need the higher weak signed-source and admissible
preparation-control certificates, subsequent tangency or nonlinear evolution
estimates, and a common physical interval. Every result in this body is at
$N=3$. No refinement-uniform tilt, all-cutoff higher weak rank, bounded nonlinear
spacetime interval, or raw-action Einstein continuum theorem follows from
these coefficient identities.

\paragraph{Reproducibility.}
The package retains the original scalar coefficient evaluator. It provides
all eighteen complete gauge-second matrices, all eight cross matrices, the
720 mixed-field entries, the four sparse moment rows, the two complete
Poisson-image identities, the normalized unit-gain certificate, and the
seven-function reconstruction. Vanishing identities are verified over
arbitrary-precision rationals; modular residues are used only to certify
nonzero exact fractions. An independent polynomial-jet implementation
recomputes all eight base tangency values without subtracting large nearby
floating values. Its ordinary Taylor coefficients use the original cubic
and quadratic envelopes, including their auxiliary and logarithmic terms.

The short checker, the exhaustive coefficient-regeneration driver, and the
interval inclusion driver have separately stated roles. Neither a floating
search nor stored success flags replace the exact identity checks. The
written analytic arguments are not represented as proof-assistant
formalizations or independent verification of the entire inherited programme.
The source audit directly inspected the V26 body, its original coefficient
interfaces and exact lower-solve certificates. All twenty-six preceding
marked mathematical bodies are preserved byte for byte, including earlier
corrections and scope restrictions. The preservation manifest records their
individual hashes and the cumulative input/output hashes.
% END IMMUTABLE EXACT-ACTION BRIDGE V27 BODY


% Future iterations append here.
% BEGIN IMMUTABLE EXACT-ACTION BRIDGE V28 BODY
\clearpage
\renewcommand{\thesection}{\arabic{section}}
\renewcommand{\theHsection}{bridgeV28.\arabic{section}}
\setcounter{section}{225}
\section{Leave the axial Poisson stratum before solving another compatibility quotient}
\label{v28:context}

The seven compatibility functions of V27 belong to a special first-field family
on which the first Poisson bracket has eighteen null directions. They have not
been proved incompatible on that entire family. There is nevertheless a more
upstream question: are its fourteen nonconstant null directions unavoidable
for the original action? The answer at the retained three-site cutoff is no.
Four explicit mixed-frequency tensor pairs preserve the original linear and
quadratic mean constraints and raise the first Poisson rank from ninety to its
maximum possible value, one hundred and four. Every nonconstant member of the
old weak-null frame then ceases to be a null test.

The change is to the canonical initial first field, not to the action or its
Poisson definition. It supplies a relatively open dense maximal-rank subset of
a connected local chart of compatible first fields. The complete prepared
quadratic drift is automatically in the Poisson range there; a unique
four-mean-zero initial multiplier tilt cancels it. At the next order, a second
initial multiplier coefficient controls all one hundred and four nonconstant
cubic drift rows. Three cubic constant-shift conditions remain and are not
silently suppressed. At a separately evaluated rational witness they are all
nonzero. Thus maximal Poisson rank removes the old nonconstant null obstruction,
but does not itself give a bounded raw multiplier rate or invariant preparation.

\paragraph{Retained inputs and scope.}
The original coefficient maps are \eqref{v20:P2}, \eqref{v20:Phi3},
\eqref{v20:K1} and \eqref{v20:B2}. The initial-tilt law is
\cref{v20:tilt-law}; the actual constant cubic source is
\cref{v21:constant-law}. We use the nonlinear range-and-mean initializer and
its parameter version in \cref{v20:tilted-initializer}, and the original
signed-source stability statement \cref{v20:graded-source}. The full leading
signed Gram at the old complementary first jet is the inherited
\cref{v12:finite-rank}, not a new determinant calculated here. V27's fibre
identities and interval exclusion remain valid on their original smaller
family. We do not reuse its constant null frame after changing the first field.
All new finite rank statements below concern $N=3$, $h=1/3$. No statement
about arbitrarily fine cutoffs follows from the integer $104$.

\section{Four mixed modes satisfying the original first-field constraints}
\label{v28:family}

Use the three-site torus $x\in\{0,1/3,2/3\}^3$, the original
$\delta_i=3(S_i-S_i^{-1})$, and $\omega=3\sqrt3$. For the axial tensors
$T_i,P_i$ retain the definitions in \cref{v12:family}; in particular
$\|T_i\|_F^2=\|P_i\|_F^2=2$ and $T_i:P_i=0$. Put
\[
 \alpha=(1,9/8,1),\qquad f_i(x)=\cos(2\pi x_i),\qquad
 \Sigma_\alpha=\sum_i\alpha_i^2=209/64.
\]
The four new frequency representatives are
\begin{equation}
 \mathcal K=\{(1,1,0),(1,0,1),(0,1,1),(1,-1,0)\}.
 \label{v28:frequencies}
\end{equation}
For $k\in\mathcal K$, let $i<j$ be its nonzero coordinate indices and let
$m$ be the remaining index. Define
\begin{equation}
 v_k=k_j e_i-k_i e_j,\quad w_k=e_m,\quad
 T_k=v_kw_k^T+w_kv_k^T,\quad Q_k=v_kv_k^T-2w_kw_k^T,
 \label{v28:tensors}
\end{equation}
and
\begin{equation}
 \phi_k(x)=\cos(2\pi k\cdot x)+\frac1{\sqrt3}\sin(2\pi k\cdot x).
 \label{v28:profiles}
\end{equation}
The values of this last profile are exactly $(1,0,-1)$ according as
$k\cdot(3x)$ is $0,1,2$ modulo three. They are finite rational records;
no rounded trigonometric evaluations define the source.

For a real parameter $\varepsilon$ set
\begin{equation}
 d_2=d_3=-8,\qquad d_1(\varepsilon)=\frac{16175}{4096}-\frac{\varepsilon^2}{3}.
 \label{v28:d}
\end{equation}
The new canonical first field is
\begin{align}
 U_\varepsilon&=\sum_i\alpha_iT_if_i+
                       \varepsilon\sum_{k\in\mathcal K}T_k\phi_k,\notag\\
 p_\varepsilon&=\omega\left\{
       \diag(d)-(d_1+d_2+d_3)I
       +\frac12\sum_i\alpha_iP_if_i
       +\frac\varepsilon2\sum_{k\in\mathcal K}Q_k\phi_k\right\},
                                                        \label{v28:Y}\\
 Y_\varepsilon&=(U_\varepsilon,p_\varepsilon).\notag
\end{align}
The physical amplitude of a prepared record will be denoted $a_0$, separately
from this first-field parameter: $X(a_0)=a_0Y_\varepsilon+O(a_0^2)$.
The profile change is not a physical-time rescaling.

\begin{lemma}[Exact linear and mean compatibility]
\label{v28:compatibility}
For every real $\varepsilon$, the first field \eqref{v28:Y} satisfies
\begin{equation}
 LY_\varepsilon=0,\qquad
 P_0P_2(Y_\varepsilon)=0,
 \label{v28:constraints}
\end{equation}
where $P_0$ takes the four original spatial means. At $\varepsilon=0$ it is
exactly the first jet \eqref{v20:seed}. Near zero the original nonlinear
initializer extends these first fields, with any fixed admitted bounded
initial tilt, to analytic exact initial constraint zeros.
\end{lemma}
\begin{proof}
Each $T_k,Q_k$ is symmetric and trace-free and annihilates $k$. Since its
entries in $k$ are $0,\pm1$, the original phase wavevector at that mode is
$\omega k$. Thus each mixed mode is exactly phase-transverse. Its tensor
identities are
\[
 |v_k|^2=2,\quad |w_k|^2=1,\quad v_k\cdot w_k=0,\quad
 \|T_k\|_F^2=4,\quad\|Q_k\|_F^2=8,\quad T_k:Q_k=0.
\]
The nonzero axial and mixed frequency pairs are distinct in the original
finite frequency group. Their profiles and derivatives are orthogonal across
different pairs. One has $\overline{\phi_k}=0$ and
$\overline{\phi_k^2}=2/3$. These assertions follow either by the exact three
values or by finite Fourier orthogonality; they do not apply an unwrapped
product rule to an aliased product.

The linear constraints are the double-divergence/trace configuration row and
the momentum divergence. All inhomogeneous tensors above are transverse and
trace-free and the homogeneous momentum is constant. This proves $LY=0$.
The three original quadratic momentum means are $\langle p,\delta_iU\rangle_h$.
They vanish because same-frequency polarization contractions are zero and
different pairs are orthogonal. The homogeneous part pairs with a mean-zero
derivative.

The remaining mean is the full quadratic Hamiltonian \eqref{v5:H2}.
On these fields its divergence and trace-derivative terms vanish. Direct
orthogonal summation, retaining the physical momentum factor $\omega$, gives
\begin{equation}
 \frac{H_2(Y_\varepsilon)}{\omega^2}
 =-2(d_1d_2+d_1d_3+d_2d_3)+\frac{\Sigma_\alpha}{2}
                                      +\frac{32}{3}\varepsilon^2.
 \label{v28:mean-energy}
\end{equation}
Indeed one new configuration mode contributes
$\frac14\omega^2|k|^2\|T_k\|_F^2\overline{\phi_k^2}
 =4\omega^2/3$ before multiplication by $\varepsilon^2$. Its momentum
partner contributes the same amount. Four such pairs give $32\omega^2/3$.
The constant momentum contributes $-2\omega^2\sum_{i<j}d_id_j$ and the
old axial pairs contribute $\omega^2\Sigma_\alpha/2$. Substitution of
\eqref{v28:d} makes \eqref{v28:mean-energy} exactly zero.

For preparation, the derivative of this scalar mean in the homogeneous
$d_1$ direction is $-2\omega^2(d_2+d_3)=32\omega^2\ne0$.
The three axial transverse momentum sine controls retained in the original
initializer have independent shift-mean derivatives, since every $\alpha_i$
is nonzero. The mixed frequencies do not alter these diagonal pairings.
Together with the original nonconstant right inverse they give the same
range-and-mean implicit-function argument, locally in $\varepsilon$ and the
initial tilt. This proves exact nonlinear initial preparation, not just the
four quadratic equalities. It supplies no time interval for evolution.
\end{proof}

The four mixed pairs are a sufficient explicit choice. Their minimal number is
not asserted. Nor is the old all-cutoff trace-free-source theorem extended to
these generally nondiagonal connection coefficients. At the witness below the
\emph{un-tilted} quadratic weak source is nonzero; it is cancelled by the
subsequently solved initial tilt, not by importing axial weak neutrality.

\section{The four constant Poisson null directions and the complete target}
\label{v28:four}

Let $\mathcal E\subset\R^{108}$ be the space of the four component-constant
multiplier fields, and let $\mathcal E^\perp$ impose zero mean in each slot.
The original Poisson pairing is the physical spatial mean. Constant blockwise
rescalings used for rational arithmetic preserve both these subspaces.

\begin{lemma}[Maximum possible first Poisson rank]
\label{v28:constants}
For every symmetric first field $Y$ satisfying $LY=0$,
\begin{equation}
 \mathcal E\subseteq\ker\mathbb K_1(Y),\qquad
 \widehat B_2(Y)\perp\mathcal E,\qquad
                       \rank\mathbb K_1(Y)\le104.
 \label{v28:universal-constants}
\end{equation}
If the rank is $104$, the complete quadratic drift equation has a unique
solution with all four means zero.
\end{lemma}
\begin{proof}
For a constant shift $c$, $Gc=0$ and
$P_{c,2}(Y)=\langle p,\delta_cu\rangle_h$. The exact translation incidence
\eqref{v21:translation-incidence} gives
\[
 \mathbb K_1(Y)[c,\mu]=DP_{c,2}(Y)[G\mu]
                                  =\langle\delta_c\mu,LY\rangle_h=0.
\]
For the unit constant lapse $e$, $Ge=0$ and $P_{e,2}=H_2$. Thus
$\mathbb K_1(Y)[e,\mu]=-\langle\mu,LF_1Y\rangle_h=0$.
Here $LF_1Y=0$ follows from the inherited exact flat constraint-propagation
identity when $LY=0$. It uses commuting skew phase derivatives, not a
nonlinear lattice Leibniz rule. The four independent constants give the rank
upper bound.

The prepared quadratic drift on a constant shift is identically zero by
\cref{v5:constant-B2}. Its value on $e$ is
$DH_2(Y)[F_1Y]=0$: both the $Ge$ term and the prepared lapse correction in
\eqref{v20:B2} vanish for this test, and a Hamiltonian is constant along its
own Hamiltonian vector. Hence all four target means vanish. A real
skew-adjoint matrix has range equal to the orthogonal complement of its
kernel. At rank $104$ that complement is exactly $\mathcal E^\perp$.
Its restriction there is bijective, giving the final assertion.
\end{proof}

This is a property of the \emph{first Poisson bracket}. The multiplier
Hessian governing the exact rate is symmetric and has different amplitude
weights. No positivity, regularity or rate estimate for that Hessian follows
from \eqref{v28:universal-constants} alone.

\section{Exact maximal-rank and active-control certificates}
\label{v28:certificate}

Write $p_\varepsilon=\sqrt3\,p_{R,\varepsilon}$ and set
\[
 Q_3=\diag(I_{27},\sqrt3I_{81}),\qquad
 \widehat K(\varepsilon)=\frac1{\sqrt3}Q_3\mathbb K_1(Y_\varepsilon)Q_3.
\]
As in \eqref{v20:rational-system}, this matrix is rational at rational
$\varepsilon$. The physical tilt is $\ell=Q_3x$ and its rational target is
$\widehat b=Q_3\widehat B_2/\sqrt3$. The first Poisson coefficient is linear
in the first field, so
\begin{equation}
 \widehat K(\varepsilon)=K^{[0]}+\varepsilon K^{[1]}
                                   +\varepsilon^2K^{[2]}.
 \label{v28:Kpoly}
\end{equation}
The first term is the inherited complementary matrix; the second is generated
from the four mixed pairs; the third is the homogeneous boost correction
$-K_{d_1}/3$. In rational momentum coordinates that correction is simply
$\varepsilon^2\diag(0,1,1)$.

All indices below are zero based in the original component-major multiplier
order. Put
\[
 \mathcal J=\{0,\ldots,107\}\setminus\{26,53,80,107\}.
\]
Let $F$ be the $27$-by-$26$ scalar matrix with columns $e_j-e_{26}$,
$0\le j<26$. In rational multiplier coordinates let $E_A$ consist first of
these twenty-six lapse columns and then of the eighty-one vector columns
$\curl_\delta e_{i,x}$. It spans the physical mean-zero lapse and mean-zero
transverse-shift slice, of dimension $78$. Let $E_W$ contain the twenty-six
pure shifts $\nabla_\delta(e_j-e_{26})$. Its range is the full nonconstant
curl-free shift slice. The fixed congruence $Q_3$ preserves both slices.

\begin{theorem}[Original first Poisson rank and complete nonconstant weak control]
\label{v28:rank}
At $\varepsilon_*=1/16$, for which $d_1=48509/12288$,
\begin{equation}
 \boxed{\det\widehat K(\varepsilon_*)_{\mathcal J,\mathcal J}
           =36\pmod{241},\qquad
        \rank\mathbb K_1(Y_{\varepsilon_*})=104.}
 \label{v28:104}
\end{equation}
The active-to-nonconstant-weak map has full row rank:
\begin{equation}
 \boxed{\rank(E_W^T\widehat K(\varepsilon_*)E_A)=26.}
 \label{v28:26}
\end{equation}
A nonzero $26$-column minor is given by
\begin{equation}
 \mathcal J_A=(0{:}10,\ 12{:}16,\ 26{:}30,\ 36,\ 53,\ 54,\ 80,\ 83),
 \qquad
 \det(E_W^T\widehat K E_A)_{:,\mathcal J_A}=223\pmod{241}.
 \label{v28:active-minor}
\end{equation}
An interval $i{:}j$ in this formula includes both endpoints. All denominators
in the displayed matrices are units at the prime $241$.
\end{theorem}
\begin{proof}
The certificate is regenerated directly from the original quadratic scalar
coefficient \eqref{v20:P2}. For the 108 original point multipliers $\mu$, form
$G\mu$ with the commuting original phase derivatives. Quadratic polarization
is the exact identity
\[
 DP_{\nu,2}(Y)[G\mu]
 =\tfrac12\{P_{\nu,2}(Y+G\mu)-P_{\nu,2}(Y-G\mu)\}.
\]
Subtract its transpose and apply the fixed rational congruence. Every entry
of $K^{[1]}$ is thereby a specified rational number. The original $K^{[0]}$
and homogeneous coefficient give \eqref{v28:Kpoly}; direct regeneration at
the witness verifies that polynomial against all $108^2$ entries. The package
contains these coefficients, the full target and tilt, and a portable checker.
An independent forward scalar contraction of \eqref{v20:P2}, before the
adjoint assembly used for the matrix, checks all 108 quadratic rows on the
new field and additional variations.

Rational reduction and Gaussian elimination modulo $241$ give the first
residue in \eqref{v28:104}. Since the denominators are nonzero at that prime,
the determinant is nonzero over $\mathbb Q$. This proves rank at least $104$;
\cref{v28:constants} proves the reverse bound. The four constant columns are
also checked exactly against each coefficient in \eqref{v28:Kpoly}.

The columns defining $E_W$ and $E_A$ contain only integers, with the gradient
and curl matrices built on the same original sites. The second rational
matrix product gives the nonzero minor \eqref{v28:active-minor}, establishing
row rank at least $26$; it has only $26$ rows. Blockwise rational/physical
rescalings do not change either range, so this is precisely surjectivity onto
the nonconstant weak covectors from the physical active slice. No floating
rank tolerance, frequency deletion, or modified source operator is used.
\end{proof}

The second statement is stronger than surjectivity from arbitrary normalized
multipliers. It permits all nonconstant weak cubic controls below to be chosen
among active \emph{initial multiplier values}. It does not make the actual
active multiplier \emph{velocity} an independently selectable control.

\begin{proposition}[Every old nonconstant weak-null test is released]
\label{v28:old-null}
For the eight tests $W_i$ of V26--V27 one has
$K^{[0]}W_i=K^{[2]}W_i=0$, whereas $K^{[1]}W_i\ne0$. Thus none is
Poisson-null on \eqref{v28:Y} at any $\varepsilon\ne0$. In particular, for
$\nu_\diamond=\frac23(W_1-W_3)$,
\begin{equation}
 \boxed{\|\widehat K(\varepsilon)\nu_\diamond\|_{\ell^2(108)}^2
                         =344373768\,\varepsilon^2.}
 \label{v28:diamond}
\end{equation}
This norm is a rational coordinate certificate, not a physical curvature norm.
\end{proposition}
\begin{proof}
The old kernel identity gives the first two zeros, also verified by exact
multiplication. The exact squared coordinate norms of the eight $K^{[1]}W_i$,
in their original order, are the successive entries in the two rows
\[
 \begin{array}{rrrr}
 1937102445/4,&645700815,&1937102445/4,&645700815,\\[2pt]
 645700815/2,&430467210,&8193225897/8,&8193225897/8.
 \end{array}
\]
They are all positive rational numbers. The same multiplication on the
specified difference gives \eqref{v28:diamond}. Since the higher homogeneous
coefficient also annihilates these tests, no other power of $\varepsilon$
occurs in their images. These are identities of the full original Poisson
matrix, not a change to the tests after projecting away their images.
\end{proof}

\section{Arbitrarily small enlargement and a relatively open dense first-field class}
\label{v28:generic}

\begin{theorem}[A maximal first Poisson stratum with a retained regular signed chart]
\label{v28:generic-theorem}
There is $\varepsilon_0>0$ such that, for every
$0<|\varepsilon|<\varepsilon_0$, both ranks in \cref{v28:rank} hold, the
original leading active symmetric Gram remains positive, and the complete
leading graded signed Gram remains invertible.

More generally, take a sufficiently small connected analytic chart of first
fields around $Y_0$ satisfying $LY=0$ and $P_0P_2(Y)=0$. It has dimension
$216$. In this chart there is a relatively open dense subset on which the
first Poisson rank is $104$ and the active-to-nonconstant-weak rank is $26$.
The chart may be chosen so that its leading graded signed Gram is invertible
throughout. These assertions are at the fixed cutoff, not uniform in $h$.
\end{theorem}
\begin{proof}
The determinant in \eqref{v28:104} is a polynomial in $\varepsilon$ of degree
at most $208$; that in \eqref{v28:active-minor} has degree at most $52$.
Both are nonzero polynomials by their certified values at $1/16$. A nonzero
polynomial has only finitely many roots. Factoring out its possible zero at
$\varepsilon=0$ therefore gives a punctured interval on which each is
nonzero. The first-field compatibility holds on that whole interval by
\cref{v28:compatibility}. This argument gives existence of a punctured interval;
no numerical lower bound on its radius is asserted.

At $Y_0$ the physical active Gram is positive and the complete leading graded
signed Gram is invertible by the inherited complementary certificate. Their
coefficients are analytic in $Y$. The leading weak source coefficient is
independent of the second field and the order-amplitude initial multiplier
tilt by \cref{v20:graded-source}. Consequently the same matrices vary
continuously with the new first field alone. Positive definiteness of the
active block and nonvanishing of the complete determinant persist after
shrinking the interval. Trace-freeness is not needed for this fixed-cutoff
continuity argument and is not claimed for the mixed family.

There are $324$ independent real canonical first-field entries and the original
linear map $L$ has rank $104$. Its kernel has dimension $220$. On that kernel
the four quadratic mean functions have independent derivatives at $Y_0$, as
certified by the original mean calibration. The analytic implicit-function
theorem supplies a local $216$-dimensional manifold. Choose a connected
coordinate ball in it and shrink it so that the two symmetric-Gram properties
just used hold throughout. The two Poisson minors are analytic functions on
this chart. Neither is identically zero, since the mixed curve has arbitrarily
small nonzero points with both minors nonzero. The zero set of a nonzero
real-analytic function on a connected chart has empty interior. Thus their
simultaneous nonzero set is open and dense. On it the rank upper bounds make
both ranks exactly the stated values. This is relative density of a specified
finite first-field class, not a probability law or universality theorem for
microscopic renewal kernels.
\end{proof}

The explicit witness $\varepsilon=1/16$ certifies the two Poisson minors. It
is \emph{not} asserted to lie in the possibly smaller signed-source
neighborhood obtained by continuity. The existence statement near zero and
the numerical witness have different roles. The latter is not silently used
as a new $107$-column signed-source certificate.

\section{Complete quadratic cancellation on original exact initial records}
\label{v28:tilt}

\begin{theorem}[Automatic quadratic range and unique normalized initial tilt]
\label{v28:quadratic}
On the maximal-rank set in \cref{v28:generic-theorem}, the equation
\begin{equation}
 \mathbb K_1(Y)\ell(Y)=-\widehat B_2(Y),\qquad
              \overline{\ell^N}=\overline{\ell^\beta}=0
 \label{v28:tilt-eq}
\end{equation}
has a unique solution, analytic locally in $Y$. The original action has analytic
exact initial preparations at $\lambda=e+a_0\ell(Y)$ with
\begin{equation}
 \boxed{P(X(a_0),\lambda(a_0))=0,\quad
 X(a_0)=a_0Y+O(a_0^2),\quad
 \mathsf B(X(a_0),\lambda(a_0))=O(a_0^3).}
 \label{v28:prepared}
\end{equation}
For each such $Y$ and sufficiently small nonzero $a_0$, the regular signed
chart gives its own local normalized original-action continuation. Its initial
literal electric and spatial link-curvature reads are $O_Y(a_0)$. A small
complete multiplier velocity or metric-curvature bound is not inferred.
\end{theorem}
\begin{proof}
\Cref{v28:constants} proves range compatibility, so no additional test on an
unspecified target is needed. Inversion on the fixed four-mean-zero space gives
uniqueness and local analyticity. Equivalently use the nonzero principal
$104$-column system and subtract the four component means. Because the
matrix and target have four zero sums, solving those $104$ rows recovers all
$108$ rows. This is an inverse of the actual skew Poisson matrix, not of a
positive substitute.

The original initializer applies to $Y$ by the construction of its constraint
chart and permits the order-amplitude multiplier values by
\cref{v20:tilted-initializer}. The exact transformation law
\cref{v20:tilt-law} gives \eqref{v28:prepared}. At fixed $Y$ the tilt is
finite, so positivity of the lapse and the original stationary chart hold
for sufficiently small amplitude. The leading signed Gram is regular by the
choice of chart and unchanged by that tilt. The same graded argument as
\cref{v20:local-branch} yields one radial multiplier-Hessian null direction,
and the inherited normalized exact-action continuation theorem applies.
The normalized stationary-response target is $O_Y(a_0)$ under the original
column weights $(a_0,a_0^2)$, giving the stated literal-curvature estimate as
in that theorem. This does not undo those weights in the multiplier rate.
No common physical interval or uniform bound as $\varepsilon\to0$ is claimed.
\end{proof}

At $\varepsilon_*=1/16$ the exact rational solution is supplied in
\src{new_tilt.json}, with the complete actual target in
\src{new_target.json}. The checker verifies all $108$ equations, not just a
square subsystem. Its physical norm and component bounds are
\begin{equation}
 83.64<\|\ell_*\|_h<83.65,\qquad
                   \max_{x,\mu}|\ell_*^\mu(x)|<122.
 \label{v28:witness-norm}
\end{equation}
The proof compares the exact rational quadratic form
$(\|x_N\|^2+3\|x_\beta\|^2)/27$ with the squares of those endpoints, and
each component square with $122^2$. This witness tilt minimizes the physical
norm among all solutions, since the only freedom is addition of a component
constant and its four means are zero. Its source and tilt values are valid
coefficient calculations independently of whether that particular
$\varepsilon_*$ is in the smaller signed-source chart.

\section{All nonconstant cubic rows are controlled by second initial multiplier values}
\label{v28:cubic}

The remaining difficulty should not be hidden by the maximal rank. Fix a
first field in the set of \cref{v28:generic-theorem}, its unique normalized
$\ell$, and an exactly constrained initial curve with those first coefficients.
Change only its second initial multiplier coefficient by $m$, keeping the
canonical free records fixed in the original analytic initializer:
\[
 \lambda_m(a_0)=e+a_0\ell+a_0^2m+O(a_0^3),\qquad
                         \overline{m^N}=\overline{m^\beta}=0.
\]
The initial canonical second coefficient is retained. Any third-coefficient
change required by exact re-preparation belongs to $\ker L$ and does not
change the cubic drift, as proved below. This is selection of initial values,
not control during an evolution.

\begin{theorem}[Full nonconstant cubic controllability and its three remaining means]
\label{v28:cubic-theorem}
Write the original complete drift of one such preparation as
$\mathsf B(a_0)=a_0^3B_3+O(a_0^4)$. Its change under $m$ is exactly
\begin{equation}
                         B_3\longmapsto B_3+\mathbb K_1(Y)m.
 \label{v28:cubic-control}
\end{equation}
Consequently every prescribed zero-mean target for the four nonconstant cubic
rows has a unique normalized second-multiplier correction. In particular,
there are exact initial preparations for which the entire nonconstant part of
$B_3$ vanishes. The only possible remaining cubic drift is its three
constant-shift components
\begin{equation}
                         C_c(Y)=t_c(Y)+\mathcal R_Y[c,\ell(Y)].
 \label{v28:C}
\end{equation}
These are independent of all higher initial coefficients. All cubic drift
rows vanish precisely when the three quantities in \eqref{v28:C} vanish.
For control of only the $26$ nonconstant weak rows, $m$ can be chosen on the
physical active slice, by \cref{v28:26}.
\end{theorem}
\begin{proof}
The exact identity \eqref{v20:derivative} is
$D_\lambda\mathsf B[\mu]=\mathbb K\mu+D_X(M\mu)[F]$.
On the present curves, $\mathbb K=a_0\mathbb K_1(Y)+O(a_0^2)$,
$D_XM=O(a_0)$, and $F=O(a_0)$. Integrating against the value change
$a_0^2m$ gives the difference $a_0^3\mathbb K_1(Y)m+O(a_0^4)$
at fixed canonical fields. The exact constraint difference is
$O(a_0^4)$ because $M=O(a_0^2)$. Keeping the free canonical parameters
fixed, the original initializer has the same canonical coefficients through
order $a_0^2$: its nonconstant range change starts at order $a_0^4$,
and its divided four-mean equation changes by $O(a_0^2)$, so the mean
balancing change in the full field can start at order $a_0^3$.
Subtracting the two exactly zero constraint expansions at order $a_0^3$
gives $L\Delta X_3=0$. The only possible resulting change to the cubic
drift is its linear coefficient on $\Delta X_3$. That coefficient is
$\mathsf B_{n,1}(\Delta X_3)=-P_{(0,\delta n),1}(\Delta X_3)$ on a
lapse test and zero on a shift test, so it vanishes when $L\Delta X_3=0$.
Thus no such correction changes the cubic drift, and
\eqref{v28:cubic-control} holds for the exactly prepared curves.

The restriction of $\mathbb K_1$ to $\mathcal E^\perp$ is invertible,
so it cancels the nonconstant projection of $B_3$ or assigns it any desired
value. Its four mean components cannot change. The cubic constant-shift law
\cref{v21:constant-law} identifies three of them with \eqref{v28:C}.
The constant-lapse component is zero: the exact homogeneity identity is
$\langle\lambda,\mathsf B\rangle=0$, and $\mathsf B=O(a_0^3)$,
$\lambda=e+O(a_0)$ imply $\langle e,B_3\rangle=0$.
This proves both necessity and sufficiency at cubic order.
Finally \cref{v28:26} gives a right inverse from the active slice to the
nonconstant weak covectors; inserting that $m$ in \eqref{v28:cubic-control}
proves the last claim without an arbitrary weak multiplier control.
\end{proof}

\begin{remark}[The mean-order bookkeeping]
An order-$a_0^2$
change of multiplier produces an order-$a_0^4$ mean-constraint change, and a
mean-balancing field change can therefore occur at order $a_0^3$.
It is in $\ker L$ and its linear drift vanishes. The exact statement used
above is cubic-drift invariance under that necessary re-preparation, not a
false prohibition on all third-field changes.
\end{remark}

There is an actual nonvanishing check on the new witness. Its complete cubic
constant-source vector, evaluated with the original constant-shift cubic
functional and the uniquely normalized tilt, obeys
\begin{equation}
 \begin{split}
 -163753&<C_{e_1}(Y_{1/16})<-163752,\\
 -61857&<C_{e_2}(Y_{1/16})<-61856,\\
 12903&<C_{e_3}(Y_{1/16})<12904.
 \end{split}
 \label{v28:means-witness}
\end{equation}
These bounds use the physical spatial-\emph{sum} coefficient convention of
V21; dividing all three rows by $27$ gives their mean convention. Their
rational residues are respectively $172,44,29$ modulo $241$. The package
regenerates the full prepared functional, including its last preparation
term, and checks the rational inequalities. Thus the explicit rank witness
is not being claimed as a simultaneous cubic-mean zero. No higher initial
multiplier coefficient can repair those three values while its $Y,\ell$
remain fixed.

\section{The quadratic slow coefficient and the new tangency distinction}
\label{v28:tangency}

The previous weak-null tests no longer give preparation-independent
nonconstant tangency obstructions. This can be expressed without asserting a
new physical evolution. On the maximal-rank compatible first-field chart let
\[
 \mathcal V(Y)=F_1Y+G\ell(Y).
\]

\begin{proposition}[A unique differential lift of quadratic cancellation]
\label{v28:lift}
The vector $\mathcal V(Y)$ is tangent to $LY=0$, and to the four quadratic
mean constraints. The first-tilt map has a unique four-mean-zero differential
along it, satisfying
\begin{equation}
 \boxed{\mathbb K_1(Y)D\ell(Y)[\mathcal V]
 =-D\widehat B_2(Y)[\mathcal V]
                   -\mathbb K_1(\mathcal V)\ell(Y).}
 \label{v28:differential-lift}
\end{equation}
For the old nonconstant weak tests the right side is no longer required to
vanish by a null pairing; their Poisson images are nonzero on the mixed curve.
This coefficient lift does not identify $D\ell[\mathcal V]$ with the actual
normalized original-action multiplier rate.
\end{proposition}
\begin{proof}
The identities $LF_1Y=0$ and $LG=0$ give linear tangency. For the constant
lapse mean $H_2$, its derivative on $F_1Y$ is zero and its derivative on
$G\ell$ is $\mathbb K_1(Y)[e,\ell]=0$. For a constant shift mean
$P_{c,2}$, its derivative on $F_1Y$ is the identically zero quadratic
translation drift, and its derivative on $G\ell$ is
$\langle\delta_c\ell,LY\rangle_h=0$. Thus $\mathcal V$ is tangent to the
compatible first-field manifold. Differentiate the identity
$\widehat B_2(Y)+\mathbb K_1(Y)\ell(Y)=0$ along that tangent. Linearity of
$\mathbb K_1$ in $Y$ gives \eqref{v28:differential-lift} and the invertible
normalized restriction gives uniqueness. This is an analytic consequence
of the already constructed coefficient map; it does not change any finite
evolution equation or make an independent rate prescription. The actual
rate still solves its symmetric-Hessian consistency equation and must also
meet the cubic mean and subsequent differentiated conditions.
\end{proof}

\section{Direction and verification}
\label{v28:status}

The first Poisson degeneracy responsible for the old eight nonconstant weak
null tests is now removed on a nonempty, relatively open dense compatible
first-field class at the same cutoff. It is compatible, on a punctured small
mixed-mode interval, with the inherited regular full signed chart. The actual
complete quadratic source then has a unique normalized initial-value
cancellation without another Poisson-cokernel condition. At the next order,
all nonconstant cubic source rows are controllable by second initial values;
three actual cubic means remain. Their nonvanishing at the explicit witness
is recorded rather than suppressed. The old seven canonical compatibility
functions are not tests for this larger Poisson stratum: reusing their null
frame there would test a different problem.

The next finite target is the three-component function \eqref{v28:C} on the
maximal-rank first-field chart, retaining its uniquely solved tilt and the
original mean constraints. A zero would remove all cubic drift rows by the
proved second-value control, but would not alone propagate a zero weak rate
or supply a physical-time interval. No common zero, zero weak velocity on
this enlarged class, differentiated-curvature estimate, refinement-uniform
rank theorem or cofinal exact-action Einstein limit is established here.
The finite genericity assertion is not a microscopic law-space universality
claim. No action, lapse normalization, phase derivative or ordered link read
has been changed.

\paragraph{Executed certificate scope.}
The exact checker verifies the four tensor pairs, their grid means and
orthogonality, the original linear and quadratic mean constraints, all four
constant Poisson-null identities, the $104$-column principal determinant,
the independent $26$-column active-control determinant, and every released
old weak-null image. The new mixed Poisson coefficient and full witness
matrix are regenerated from the retained rational action coefficients.
Independent forward scalar contractions test the adjoint-assembled quadratic
source. The complete original prepared target, its exact normalized tilt,
and its three cubic mean values are retained as rational arrays and checked
against all original rows. Bounds on the displayed norms and scalar values
use exact rational comparisons, not floating tolerances.

The open-dense and punctured-interval statements are analytical consequences
of the exact finite nonvanishing certificates and the named inherited
initialization and signed-chart results. No finite list of matrix samples is
represented as a proof under spatial refinement. Earlier exploratory searches
on the smaller canonical quotient supplied neither a common root nor a
nonexistence theorem and are not used as proofs. The complete V1--V27 marked
bodies are retained byte for byte. The package records the preservation
hashes, coefficient files, checker and the stages actually executed; it does
not assert proof-assistant formalization or nonlinear spacetime integration.
% END IMMUTABLE EXACT-ACTION BRIDGE V28 BODY


% Future iterations append here.
% BEGIN IMMUTABLE EXACT-ACTION BRIDGE V29 BODY
\clearpage
\renewcommand{\thesection}{\arabic{section}}
\renewcommand{\theHsection}{bridgeV29.\arabic{section}}
\setcounter{section}{234}
\section{Simultaneous cubic means on the mixed first-field stratum}
\label{v29:context}

The maximal first-Poisson-rank family of V28 removes the additional axial
Poisson null tests. It does not remove the three cubic constant-shift means:
those means were nonzero at its explicit rational rank witness. The next
question is therefore their simultaneous zero, with the uniquely determined
quadratic-cancelling tilt retained. This body answers that finite coefficient
question positively. Independent cosine and sine amplitudes on the four
mixed modes give a regular zero of all three original cubic means, without
losing maximal Poisson rank. A rational contraction certificate, rather than
a floating residual, establishes the zero.

The new point is not claimed to lie in the unspecified neighborhood where
V28 retained its higher signed-source certificate by continuity. The first
Poisson matrix and the symmetric multiplier Hessian are different objects.
Accordingly, the result gives exactly constrained original initial records
with full drift of order four in field amplitude, but does not yet assert a
normalized exact-action history through this particular zero. That requires
the original higher signed chart to be checked at the selected fields.

\paragraph{Related-source audit and change of direction.}
The supplied regenerative V38 excludes a nearby restoration of the missing
axial profile on its own boundary branch and quantifies the loss in mean
conditioning \cite[\src{v38:thm:isolation}, \src{v38:prop:mean}]{v29:Native}.
This supports keeping all three axial profiles nonzero and continuing with
mixed frequencies; it is not a theorem about the new family below.
The SM source-selection examples distinguish a successful restricted reversal
panel from symmetry of the complete original instrument
\cite[\src{sel35:prop:mixed}, \src{sel35:prop:scalar}]{v29:SM}.
The perturbative ESM zero-transfer result is specific to its completed action
and constant-metric sector; it does not assert a general parity of an oriented
cubic vertex \cite[\src{thm:v23-zero-transfer}, \src{sec:v23-remaining}]{v29:ESM}.
Finally the YM differentiated Laplace estimate retains its original local
normalization and leading potential \cite[\src{r32:derivatives},
\src{r32:kernels}]{v29:YM}. No symmetry, inverse floor, or physical-time bound
from these different models is substituted here. Their relevant methods and
scope restrictions are used to organize a calculation of the complete
original three-component source. The accompanying audit records exact source
filenames, labels, inspected ranges, and hashes.

\paragraph{Inherited inputs.}
The coefficient definitions remain \eqref{v20:P2}, \eqref{v20:Phi3},
\eqref{v20:K1}, \eqref{v20:B2}, and \eqref{v21:R}. The last includes its
original second-preparation term. The flat stationary chart, range-and-mean
initializer, first-tilt law, and V28's second-value control are retained as
inputs. They are not independently re-proved by this coefficient certificate.
All new ranks and zeros in this body concern $N=3$, $h=1/3$. There is no
change to the action, physical clock, phase derivative, or ordered link read.

\section{Independent mixed profiles and exact mean calibration}
\label{v29:family}

Use $\omega=3\sqrt3$, the axial amplitudes
$\alpha=(1,9/8,1)$, and V28's four frequency vectors
\[
 k_1=(1,1,0),\quad k_2=(1,0,1),\quad k_3=(0,1,1),\quad k_4=(1,-1,0).
\]
For each $k$, let $i<j$ be its nonzero coordinate indices, let $m$ be the
remaining index, and set
\[
 v_k=k_j e_i-k_i e_j,\quad w_k=e_m,\quad
 T_k=v_kw_k^T+w_kv_k^T,\qquad Q_k=v_kv_k^T-2w_kw_k^T.
\]
Both tensors are transverse and trace-free; their squared Frobenius norms
are four and eight, respectively, and $\langle T_k,Q_k\rangle_F=0$.
The original signed phase vector at $k$ is $\omega k$, so this is also
exact phase transversality on the retained grid.

Let $c_k(x)=\cos(2\pi k\cdot x)$ and
$s_k(x)=2\sin(2\pi k\cdot x)/\sqrt3$. Their grid values, ordered by
$k\cdot(3x)\bmod3$, are
\[
 c_k=(1,-1/2,-1/2),\qquad s_k=(0,1,-1).
\]
For eight real coefficients $a=(a_1,\ldots,a_4)$ and
$b=(b_1,\ldots,b_4)$ define $\phi_j=a_jc_{k_j}+b_js_{k_j}$. The $b_j$
here are profile coefficients, not the homogeneous boosts $d_i$ below.
Let $T_i,P_i$ be the original complementary axial tensors. Define
\begin{align}
 U_{a,b}&=\sum_{i=1}^3\alpha_iT_i\cos(2\pi x_i)
                         +\sum_{j=1}^4T_{k_j}\phi_j,\notag\\
 p_{a,b}&=\omega\{\diag(d_1,-8,-8)-(d_1-16)I\}
       +\frac\omega2\left\{\sum_{i=1}^3\alpha_iP_i\cos(2\pi x_i)
                                +\sum_{j=1}^4Q_{k_j}\phi_j\right\},\notag\\
 d_1&=\frac{16175}{4096}
              -\frac1{16}\sum_{j=1}^4a_j^2-\frac1{12}\sum_{j=1}^4b_j^2.
 \label{v29:fields}
\end{align}
In rational coordinates $p=\sqrt3\,p_R$, the homogeneous correction to
$p_R$ is $\diag(0,q,q)$, where
$q=3\sum a_j^2/16+\sum b_j^2/4$. The former V28 curve is
$a_j=\varepsilon$, $b_j=\varepsilon/2$ for every $j$.

\begin{theorem}[Exact compatibility of the enlarged first fields]
\label{v29:compatibility}
Every first field $Y(a,b)=(U_{a,b},p_{a,b})$ satisfies the original linear
constraint equations and all four quadratic mean equations:
\begin{equation}
                  LY(a,b)=0,\qquad P_0P_2(Y(a,b))=0.
 \label{v29:compatible}
\end{equation}
The four mean balancing directions remain independent near the selected
profiles below. Each such first field extends to analytic exact initial
constraint zeros of the original action at sufficiently small field amplitude.
The amplitude neighborhood is a fixed-cutoff one.
\end{theorem}
\begin{proof}
The mixed tensors are transverse and trace-free, and the homogeneous
momentum is divergence-free. Thus $LY=0$ directly with the original phase
derivatives. The seven nonzero frequency pairs (three axial and four mixed)
are distinct on the cyclic three-site group, so different pairs are
orthogonal in the original spatial-mean pairing. Also
\[
 \overline{\phi_j^2}=a_j^2/2+2b_j^2/3,
 \qquad \overline{c_{k_j}s_{k_j}}=0.
\]
The mixed contribution of one tensor pair to the original quadratic
Hamiltonian is $4\omega^2\overline{\phi_j^2}$. Its configuration part is
$\omega^2|k_j|^2\|T_{k_j}\|_F^2\overline{\phi_j^2}/4
=2\omega^2\overline{\phi_j^2}$; its momentum part has the same value.
The homogeneous contribution is
$-2\omega^2(d_1d_2+d_1d_3+d_2d_3)$, with $d_2=d_3=-8$.
Changing $d_1$ therefore changes it by $32\omega^2\Delta d_1$.
The correction in \eqref{v29:fields} cancels the sum of the mixed energies
exactly. The original axial calibration already cancels the remaining energy.
This proves the scalar mean equation.

Each constant-shift quadratic constraint is
$\langle p,\delta_c U\rangle_h$. The homogeneous term integrates to zero,
different frequency pairs are orthogonal, and the configuration and momentum
tensors at each common frequency are orthogonal. Thus all three momentum
means vanish. No product rule for the phase derivative is used.

For the nonlinear preparation, retain a free change of $d_1$ and the three
axial sine-momentum directions $T_i\sin(2\pi x_i)$. At the displayed first
fields their mean derivative is diagonal: the energy derivative in $d_1$
is $32\omega^2\ne0$, and the shift derivatives in the respective sine
momenta are $-\omega\alpha_i\ne0$, up to the unchanged common sign
convention for the constraint rows. The other entries vanish by the same
frequency and tensor orthogonality. The original nonconstant linear inverse
and analytic stationary chart are unchanged. Solve their range equation,
divide the four remaining means by the square of field amplitude, and apply
the analytic implicit-function theorem with this nonsingular derivative.
The result is an exact zero of every original initial constraint, not only
of its quadratic approximation. Analytic parameter dependence is local at
this fixed cutoff and follows from the same two inverses.
\end{proof}

\section{The complete cubic-mean zero as a polynomial system}
\label{v29:system}

Retain the original rational synthesis $Q_3=\diag(I_{27},\sqrt3 I_{81})$.
For a rationalized multiplier $x\in\R^{108}$ the quadratic cancellation is
\[
 \widehat K(Y)x=-\widehat b(Y),\qquad
 \widehat K(Y)=\frac1{\sqrt3}Q_3\mathbb K_1(Y)Q_3,
 \qquad \widehat b(Y)=\frac1{\sqrt3}Q_3\widehat B_2(Y).
\]
The complete cubic constant-shift vector is
\begin{equation}
 C(Y,x)=\bigl(t_{e_i}(Y)+\mathcal R_Y[e_i,Q_3x]\bigr)_{i=1}^3.
 \label{v29:C}
\end{equation}
It is the original physical spatial-sum coefficient. Its rational evaluation is
\[
 3C_{\mathrm{basic}}(U,p_R,x)-2C_{\mathrm{basic}}(U,0,x),
\]
as in the retained coefficient functions. In particular the preparation term
$-\langle\delta_i\ell,P_2(Y)\rangle$ is included.

Let $\mathcal J$ omit indices $26,53,80,107$ from the component-major
multiplier order $0,\ldots,107$. Write $E_{\mathcal J}$ for insertion of
104 coordinates, with zero values at those four omitted indices.

\begin{lemma}[Point representatives and physical normalization]
\label{v29:point-gauge}
For the compatible fields \eqref{v29:fields}, the four component-constant
vectors lie in $\ker\widehat K$, and $\widehat b$ has zero component sums.
If $\widehat K_{\mathcal J\mathcal J}$ is invertible, its 104 equations
recover all 108 original quadratic equations. Moreover $C(Y,x)$ is unchanged
by adding any component-constant multiplier. Hence
\[
 x^{\rm n}=x-E_0\mathsf m x,\qquad \ell=Q_3x^{\rm n}
\]
is the unique four-mean-zero tilt whenever $x=E_{\mathcal J}
\widehat K_{\mathcal J\mathcal J}^{-1}(-\widehat b_{\mathcal J})$.
\end{lemma}
\begin{proof}
The first two assertions are the original constant-row identities of
\cref{v28:constants}. They also follow coefficientwise from the linear
constraint $LY=0$, the flat translation incidence, and the integrated
quadratic Hamiltonian. Each omitted row is minus the sum of the other 26
rows in its component, so no equation is deleted. The same constants are
four independent kernel vectors, giving rank at most $104$; the invertible
principal block gives equality and uniqueness modulo those constants.

For the cubic means, $Ge=G(0,c')=0$ for constant $c'$. Formula
\eqref{v21:R} gives
$\mathcal R_Y[c,e]=-DH_2(Y)[\delta_cY]=0$, since its constant-coefficient
quadratic form commutes with the skew phase derivative.
For a constant shift $c'$ it gives
$-DP_{c',2}(Y)[\delta_cY]=0$, by skew-adjointness and the commutation of
$\delta_c$ and $\delta_{c'}$ in
$P_{c',2}=\langle p,\delta_{c'}U\rangle_h$.
The last preparation term is zero for both constant choices. Thus $C$ is
unchanged by normalization. This argument uses the full cubic-mean formula;
it is not a parity assumption about an oriented plaquette.
\end{proof}

Fix five profile coefficients as the following exact dyadic numbers:
\begin{equation}
 (a_2,a_3,a_4,b_1,b_4)
   =2^{-20}(43835,43443,91967,-2807,70091).
 \label{v29:fixed}
\end{equation}
The three free coefficients are $z=(a_1,b_2,b_3)$. For
$y=(z,x_{\mathcal J})\in\R^{107}$ define the polynomial map
\begin{equation}
 \mathcal F(y)=
 \begin{pmatrix}
  [\widehat K(Y(z))E_{\mathcal J}x_{\mathcal J}+\widehat b(Y(z))]_{\mathcal J}\\
  C(Y(z),E_{\mathcal J}x_{\mathcal J})
 \end{pmatrix}.
 \label{v29:F}
\end{equation}
It has rational coefficients and total degree at most six. The system contains
no inverse. Its first 104 components are the original Poisson-range equation;
its last three are all the remaining cubic means.

For explicit coefficient reconstruction write
\[
 Y(z)=Y_0+z_1Y_1+z_2Y_2+z_3Y_3+
       \left(\frac3{16}z_1^2+\frac14z_2^2+\frac14z_3^2\right)Y_4
\]
in rational canonical coordinates, where $Y_4$ has zero configuration and
momentum $\diag(0,1,1)$. The fields $Y_0,\ldots,Y_3$ follow literally
from \eqref{v29:fields} and \eqref{v29:fixed}. Linearity of $\widehat K$
and quadratic homogeneity of $\widehat b$ require five complete Poisson
matrices and fifteen complete target vectors. The mixed quadratic vectors
are evaluated by exact polarization. All are regenerated from the retained
$P_2/\Phi_3$ formulas. The three final components of \eqref{v29:F} are
evaluated by the original $P_{c,3}$ envelope and its prepared derivatives.
The rational coefficient files and their hashes accompany the note.

\section{A regular simultaneous zero certified over the rationals}
\label{v29:certificate}

The exact center $y_0$ is the dyadic vector with denominator $2^{350}$ in
\src{candidate_point.json}; its complete 107 integer numerators are supplied.
The unique root below is defined by this rational center and the following
radius, not by the rounded decimal display:
\begin{equation}
 \mathcal B=\{y\in\R^{107}:\|y-y_0\|_\infty\le2^{-240}\}.
 \label{v29:box}
\end{equation}
The three selected profile coefficients are approximately
\begin{equation}
 (a_1,b_2,b_3)\simeq
 (0.0515308896122,\ 0.0071497379600,\ 0.0400678398141).
 \label{v29:decimal}
\end{equation}
These decimals only identify the location. They are not interval endpoints.

\begin{theorem}[An exact zero of the three original cubic means]
\label{v29:root}
The box \eqref{v29:box} contains a unique zero $y_*$ of
\eqref{v29:F}. Throughout that box the principal first-Poisson matrix
$\widehat K_{\mathcal J\mathcal J}$ is invertible, so the complete first
Poisson rank is $104$. The full derivative $D\mathcal F$ is invertible on
the box. At $Y_*=Y(z_*)$, the normalized physical tilt consequently satisfies
\begin{equation}
 \boxed{\mathbb K_1(Y_*)\ell_*=-\widehat B_2(Y_*),\qquad
        \overline{\ell_*^N}=\overline{\ell_*^\beta}=0,\qquad
        t_c(Y_*)+\mathcal R_{Y_*}[c,\ell_*]=0\quad(c=e_1,e_2,e_3).}
 \label{v29:joint-root}
\end{equation}
No small floating residual or modular zero is used to infer any of these zeros.
\end{theorem}
\begin{proof}
Here are the explicit verification inequalities. Let $R$ be the supplied
$107\times107$ dyadic preconditioner, with denominator $2^{70}$, and let
$y_c$ be the componentwise downward dyadic rounding of $y_0$ to128 bits.
The coefficient evaluator gives exact rational values for
$\mathcal F(y_0)$ and $D\mathcal F(y_c)$. It also bounds every component
of the holomorphic polynomial \eqref{v29:F} by $M$ on the complex unit
polydisc about $y_0$. The executed bounds are recorded in the following
certificate table.
\begin{center}\small
\begin{tabular}{@{}ll@{}}
\toprule
Quantity & Certified bound\\\midrule
$\sup_{\|y-y_0\|_\infty\le1}|\mathcal F_i(y)|$ & $M\le 2^{44}$\\
$\|R\|_\infty$ & $<21$\\
$\|I-RD\mathcal F(y_c)\|_\infty$ & $<10^{-8}$\\
$\|R\mathcal F(y_0)\|_\infty/2^{-240}$ & $<3\times10^{-29}$\\
$\sup_{\mathcal B}\|I-RD\mathcal F\|_\infty$ & $<10^{-8}$\\
Principal Poisson preconditioned defect on $\mathcal B$ & $<2\times10^{-10}$\\
\bottomrule
\end{tabular}
\end{center}
Every comparison is made as an inequality of integer cross-products or exact
fractions. The actual fractions, center, preconditioners, values and derivative
arrays are part of the certificate, not just the displayed upper bounds.

For the derivative, reverse accumulation in the rational arithmetic circuit
uses only the identities $D(u+v)=Du+Dv$ and $D(uv)=vDu+uDv$; divisions are
by fixed nonzero rational constants. Its three output derivatives are thus
exact derivatives of the original polynomial, not finite differences of an
inverted matrix. Independent degree-six polarization checks and exact
coefficient regressions are recorded in the verification report.

The polynomial majorant is obtained from the same scalar formulas, replacing
each addition or subtraction by the sum of absolute bounds, each product by
the product of its bounds, and every fixed division by division by the
absolute constant. Bounds for field and tilt entries follow directly from
\eqref{v29:fields} on the unit polydisc. No intermediate cancellation is
needed for this majorant. The full rational target bank supplies the first
104 majorants; the complete cubic-mean circuit supplies the last three.

Cauchy's one- and two-variable integral estimates imply
$|\partial_j\partial_k\mathcal F_i|\le8M$ on the concentric half polydisc.
Put $n=107$, $r=2^{-240}$ and $d=\|y_c-y_0\|_\infty$. Thus
\[
 q:=\|I-RD\mathcal F(y_c)\|_\infty
       +8n^2\|R\|_\infty M(r+d)
       \ \ge\ \sup_{\mathcal B}\|I-RD\mathcal F\|_\infty.
\]
The computed fractions give $q<1/2$ and
$\beta+qr<r$, where $\beta=\|R\mathcal F(y_0)\|_\infty$.
Consequently $y\mapsto y-R\mathcal F(y)$ is a contraction of the closed
real cube strictly into itself. Completeness gives a unique fixed point.
The bound on $RD\mathcal F$ makes $R$ and $D\mathcal F$ nonsingular, so
its fixed point is a zero of $\mathcal F$ and the derivative stays nonsingular
throughout the cube. This supplies the root assertion without a numerical
rank decision.

For completeness the Poisson block has its own dyadic preconditioner. Its
center defect is bounded by exact integer products. Since $\widehat K(Y(z))$
is a quadratic polynomial in $z$, the norms of its five exact coefficient
matrices give a bound for its entire variation on the root cube. The last
line of the table is their sum after preconditioning and is strictly below
one. The Neumann series therefore proves invertibility everywhere on the
cube, independently of the cubic-mean derivative. The constant-kernel and
row-completion identities in \cref{v29:point-gauge} then recover all $108$
quadratic equations. Finally normalization leaves all three means unchanged,
proving \eqref{v29:joint-root}.
\end{proof}

The computational proof uses the standard distinction between proposing a
root and verifying one \cite{v29:Rump}. Floating arithmetic proposed the
center and preconditioners. The final assertions use only their declared
rational values and the proved majorants; no floating success flag is an
input to the contraction inequalities.

\section{Five-dimensional persistence without an unverified rank transfer}
\label{v29:persistence}

\begin{theorem}[The three means are transverse on the eight-profile family]
\label{v29:five}
Near the certified point, the normalized first-tilt map $\ell(Y(a,b))$ is
analytic. The map
\[
 (a,b)\longmapsto
 \bigl(t_{e_i}(Y(a,b))+\mathcal R_{Y(a,b)}[e_i,\ell(Y(a,b))]\bigr)_{i=1}^3
\]
is a submersion. Its zero set in the eight-dimensional mixed-profile family
is locally an analytic manifold of dimension five. All three axial profiles
remain nonzero and the first Poisson rank remains $104$ there.
\end{theorem}
\begin{proof}
The principal inverse just proved gives the analytic tilt. In the derivative
of \eqref{v29:F}, eliminate its 104 multiplier columns by that same inverse.
The remaining $3\times3$ Schur matrix is precisely the derivative of the
reduced cubic means in $(a_1,b_2,b_3)$, with the other five profiles fixed.
The complete derivative and the Poisson block are both invertible by
\cref{v29:root}, so this Schur matrix is invertible. Treat the other five
profile coefficients as parameters and apply the analytic implicit-function
theorem. It gives an analytic local graph of simultaneous zeros. The field
synthesis has independent nonzero Fourier profiles, so this is a genuine
five-dimensional first-field family, not a redundancy of coordinates.
Continuity preserves maximal Poisson rank on a smaller neighborhood. The
axial amplitudes were held fixed at $(1,9/8,1)$ throughout.
\end{proof}

One can also state the local dimension in the full compatible first-field
chart, using its declared unquotiented canonical coordinates. Its linear constraint rank is $104$ on $324$ real canonical coordinates.
The four mean derivatives in \cref{v29:compatibility} are independent on
that kernel, so the compatible chart has dimension $216$ near $Y_*$. The
three certified mixed-profile tangent directions lie in this chart and
make its cubic-mean derivative onto. Hence the same zero set in the full
maximal-Poisson-rank compatible chart is locally a codimension-three
analytic submanifold, of dimension $213$. This is a first-coefficient selection
statement, not an open class of microscopic laws or generic Cauchy data.

\section{Exact initial constraints with all cubic drift rows cancelled}
\label{v29:preparation}

\begin{theorem}[Quartic initial drift on the unchanged action]
\label{v29:quartic}
At each first field in a sufficiently small neighborhood of $Y_*$ within
the zero set of \cref{v29:five}, the original action admits analytic initial
records, for sufficiently small amplitude $a_0$, satisfying
\begin{equation}
 \boxed{\begin{gathered}
 X(a_0)=a_0Y+O(a_0^2),\qquad
 \lambda(a_0)=e+a_0\ell(Y)+a_0^2m(Y)+O(a_0^3),\\
 P(X(a_0),\lambda(a_0))=0,\qquad
 \overline{\lambda^N}=1,\quad\overline{\lambda^\beta}=0,\\
 \mathsf B(X(a_0),\lambda(a_0))=O(a_0^4).
 \end{gathered}}
 \label{v29:quartic-eq}
\end{equation}
The spatial metric and lapse are positive for small enough $|a_0|$. The
bounds and analytic parameter dependence are local at this fixed cutoff.
This is an initial-record theorem, not an assertion that the selected
records have a normalized exact continuation on a regular higher signed chart.
\end{theorem}
\begin{proof}
Use the original analytic initializer from \cref{v29:compatibility} at
$\lambda=e+a_0\ell(Y)$. The first-tilt law and
\eqref{v29:joint-root} cancel the entire quadratic drift, giving
$\mathsf B=a_0^3B_3+O(a_0^4)$. Its three constant-shift components vanish
by the certified cubic-mean equations. The constant-lapse component vanishes
by the exact homogeneous identity $\langle\lambda,\mathsf B\rangle=0$
used in \cref{v28:cubic-theorem}. Thus $B_3$ has zero means in all four
components.

The normalized first Poisson matrix is invertible. Choose
$m(Y)$ to solve $\mathbb K_1(Y)m=-B_3$ on the four-mean-zero space.
The absolute baseline $B_3$ is the original analytic cubic drift, including
higher action terms required by its nonconstant rows. It need not be assumed
zero or replaced by the already computed three means. V28's exact
second-value law, with the corresponding re-preparation of the original
constraints, changes that coefficient by precisely $\mathbb K_1(Y)m$.
Its necessary third-field correction lies in $\ker L$ and has zero linear
drift, so it does not alter this cancellation. This proves every equation
in \eqref{v29:quartic-eq}.

The inverse and all fixed-cutoff initial coefficient maps depend analytically
on $Y$ in a smaller neighborhood. Their actual values, including the
unlisted baseline $B_3$, determine $m(Y)$ analytically. The initializer then
supplies the exact, not merely formal, constrained curve. Small amplitude
keeps the original stationary chart, spatial metric, and lapse positive.
Neither an eigenvalue of the higher signed multiplier Hessian nor a
physical-time existence estimate is inferred from this argument.
\end{proof}

The constant components of the initial multiplier are fixed once. No mean
source is discarded and no value is reset during an evolution. The four-order
bound refers to the original constraint drift at these prepared cuts. Since
the active and weak symmetric-Hessian columns have different amplitude
weights, it does not by itself give a bounded full multiplier velocity.

\section{Remaining dynamical and refinement requirements}
\label{v29:status}

The finite target left by V28 is now populated: there is a regular
simultaneous zero of all three original cubic means with maximal first
Poisson rank. It belongs to a five-dimensional local family within an
explicit eight-profile class and removes every cubic drift row after the
already proved second-value control. This is not the old diamond-component
crossing, and it is not an inference from a count of adjustable parameters.

The next rank test is the original complete higher weak signed source at
$Y_*$. The root certificate verifies the first Poisson matrix and the
three-mean submersion only. V28's continuity neighborhood about its axial
point has no quantitative radius placing the present zero inside it.
A new higher signed-source certificate or an explicit overlap estimate is
therefore required before invoking normalized exact continuation at this
root. Further propagation of the selected condition, differentiated physical
curvature bounds, cofinal selection, and a common physical lifespan remain
separate. Neither a raw-action Einstein continuum theorem nor a no-go theorem
for that limit is proved here.

The audit of the supplied related notes does not change these distinctions.
Their axial obstruction is not re-proved on another notation, their local
coefficient or measure estimates are not imported as a gravitational inverse,
and a partial symmetry panel is not used to replace an actual source.
The positive acquisition is the independently verified simultaneous cubic
source zero of the unchanged gravitational action.

\paragraph{Verification and preservation.}
The package regenerates five complete first-Poisson matrices and fifteen
complete quadratic target vectors for the exact slice, evaluates the three
original cubic means, differentiates their complete rational circuit, and
checks the root and Poisson-invertibility inequalities by integer arithmetic.
The explicit mean and tensor identities, source normalization, all omitted-row
relations, and comparisons with the retained V28 family are separately checked.
A rounded numerical solution is used only to propose rational inputs to the
certificate. The full nonlinear stationary initial root, the baseline
nonconstant cubic vector, an exact-action time history, and the higher weak
signed determinant at the new zero are not represented as computed.

All twenty-eight preceding marked bodies are preserved byte for byte.
The mathematical proofs rely on their named stationary and initialization
inputs; the finite checks are not an independent reproof of the entire
programme or a proof-assistant formalization. The report identifies which
regenerations, rational checks, and rendering checks were actually executed.

\begin{thebibliography}{9}
\raggedright
\bibitem{v29:Native} \emph{Renewal Native Regenerative Stationarity Closure
Research Notes V38}, supplied file
\src{Renewal_Native_Regenerative_Stationarity_Closure_Notes_V38(1).tex}.
The named axial and conditioning results are used only for the source audit.
\bibitem{v29:SM} \emph{Renewal SM Source Selection Research Notes V35},
supplied file \src{Renewal_SM_Source_Selection_Research_Notes_V35(1).tex}.
\bibitem{v29:ESM} \emph{Renewal Perturbative ESM Continuum Research Notes V23},
supplied file \src{Renewal_Perturbative_ESM_Continuum_Research_Notes_V23(1).tex}.
\bibitem{v29:YM} \emph{Renewal Yang--Mills Mass-Gap Research Notes V32},
supplied file \src{Renewal_Yang_Mills_Mass_Gap_Research_Notes_V32.tex}.
Version numbers in these distinct lineages are not chronological equivalences.
\bibitem{v29:Rump} S. M. Rump, \emph{Verification methods: Rigorous results
using floating-point arithmetic}, Acta Numerica \textbf{19} (2010), 287--449.
\href{https://doi.org/10.1017/S096249291000005X}
{doi:10.1017/S096249291000005X}. Background for verification methodology;
the finite polynomial majorant and contraction proof used here are explicit.
\end{thebibliography}
% END IMMUTABLE EXACT-ACTION BRIDGE V29 BODY


% Future iterations append here.
% BEGIN IMMUTABLE EXACT-ACTION BRIDGE V30 BODY
\clearpage
\renewcommand{\thesection}{\arabic{section}}
\renewcommand{\theHsection}{bridgeV30.\arabic{section}}
\setcounter{section}{241}
\section{The signed chart at the cubic-mean root and the rate to select}
\label{v30:context}

The simultaneous source zero in \cref{v29:root} concerns the first Poisson
matrix and three cubic means. It does not certify the symmetric multiplier
Hessian. This body evaluates the complete leading signed source at that
\emph{same} root. A full-source interval enclosure, followed by independent
integer matrix arithmetic, proves its nonsingularity. No continuity radius
from the old axial seed is substituted for this calculation.

There is then a stronger preparation consequence than merely making the
initial drift quartic. The three mean controls and the 104 second-value
controls realize any prescribed finite order-amplitude initial multiplier
rate, by selecting initial records whose \emph{own} exact rate has that value.
In particular the entire initial canonical constraint drift can be made
exactly zero. This does not make its zero set invariant. An original-source
test detects a nonzero derivative at such a cut. The useful alternative is
the differential of the already determined first-tilt map along its canonical
first-field motion. Selecting that rate cancels the entire leading fast
forcing and gives an actual short-time tracking theorem for the unchanged
action. The physical interval still has length proportional to the cube of
field amplitude, and the cutoff here is fixed.

\paragraph{Inherited inputs and scope.}
The input point is exactly \cref{v29:root}, with the eight-profile fields
\eqref{v29:fields}. We retain its source definitions, original analytic
stationary chart and range-and-mean initializer. The first-tilt law, the
independence of the graded source from higher preparation, and the full
second-value control are \cref{v20:tilt-law,v20:graded-source,v28:cubic-theorem}.
The differentiated first-tilt identity is \cref{v28:lift}; it is not claimed
as a new identity here. The new use of it follows a verified signed-source
calculation and an exact initial-rate realization. The initial-layer method
of V22 is reused with the new point and rate; its old secular null-test
conclusion is not transferred. All assertions involving numerical margins
in this body have $N=3$, $h=1/3$. There is no change to the finite action,
phase derivatives, time normalization, or literal link observations.

\section{The complete source is evaluated on one and the same root box}
\label{v30:source}

Let $Y_*=Y(z_*)$ be the first field in \cref{v29:root}. Denote the original
first source by $J_1(Y)$ and its complete weak second coefficient by
$\mathcal Q_h(Y;\cdot)$, as defined in \cref{v3:compiler}. Their target is the
162-dimensional space of symmetric connection covectors, in site order and
independent-entry order $(11,22,33,12,13,23)$. In these covector coordinates
the physical inverse Cartan matrix is the direct sum over the 27 sites of
\begin{equation}
 K_{\rm site}^{-1}
 =\diag(1/2,1/2,1/2,1/4,1/4,1/4)
       -\frac14\begin{pmatrix}\mathbf1_{3\times3}&0\\0&0\end{pmatrix}.
 \label{v30:Ksite}
\end{equation}
This keeps the two off-diagonal Frobenius occurrences. A common conversion
from spatial sum to spatial mean changes no rank, inertia, or zero equation.

Let $e_j$ now denote scalar point indicators, and choose the 29 weak tests
\[
 (0,e_i^{\rm spatial}),\quad i=1,2,3,
 \qquad (0,\nabla_\delta(e_j-e_{26})),\quad 0\le j<26.
\]
The first three are constant shifts. Let $W(Y)$ have the corresponding
$\mathcal Q_h$ columns. For the active first-source columns use the following
78 indices of the original component-major 108-multiplier list:
\begin{equation}
\begin{split}
\mathcal I_A=( {}&54,99,96,63,100,101,60,84,73,57,92,90,91,\\
 &74,55,81,95,78,102,61,58,106,94,64,71,104,\\
 &56,66,88,38,107,68,44,85,28,33,52,45,32,\\
 &46,30,76,41,53,49,23,22,5,34,14,21,7,\\
 &17,36,19,16,77,15,11,25,2,8,62,86,3,\\
 &6,12,0,1,4,18,20,9,43,67,24,26,13).
\end{split}
\label{v30:indices}
\end{equation}
Set
\begin{equation}
 C(Y)=J_1(Y)_{:,\mathcal I_A},\quad
 F_*(Y)=(C(Y)\ W(Y)),\qquad
 \Gamma_*(Y)=F_*(Y)^*K_0^{-1}F_*(Y).
 \label{v30:full-source}
\end{equation}
The notation $F_*$ denotes this source matrix, not the canonical vector field.

The first source is evaluated by the degree-one tangent variation of the
original scalar saddle functional, including the complete flat auxiliary
multiplier response. In particular the diagonal-boost-only first-source
shortcut of V12 is \emph{not} used for the mixed tensor momenta. The second
source uses the normal second connection correction, its first auxiliary
response, and the tangent test in \cref{v3:compiler}. The cubic electric and
magnetic terms remain where that coefficient requires them. The degree
selection here is extraction of a coefficient of the original action, not
an alternative evolution equation.

\begin{lemma}[A same-root full-source enclosure]
\label{v30:enclosure}
The dyadic endpoint array in \src{source_dyadic_bounds.json}, with common
denominator $2^{128}$, encloses every entry of $F_*(Y(z))$ for every $z$
occurring in the V29 root cube. All 108 first columns and all 29 weak
columns are regenerated before selecting \eqref{v30:indices}.
\end{lemma}
\begin{proof}
Use the exact five-field polynomial synthesis of \cref{v29:system}, with
its five fixed profile coefficients \eqref{v29:fixed}. Each free coordinate
is enclosed by outward dyadic endpoints of precision 56 which contain the
350-bit V29 center plus and minus $2^{-240}$. Those inclusions are checked
as rational inequalities. All fixed field coefficients are enclosed by
exactly checked 50-bit dyadics. With
\[
 k=\left\lfloor\sqrt{3\,2^{124}}\right\rfloor,
 \qquad k^2<3\,2^{124}<(k+1)^2,
\]
the physical momentum factor $\sqrt3$ is enclosed by
$[k2^{-62},(k+1)2^{-62}]$. Its endpoints are exactly representable in the
specified extended binary arithmetic. Thus both the source point and the
physical algebraic scaling are enclosed, not replaced by their decimal
approximations.

The coefficient circuit is the interval transcription of the original
saddle functional retained from V8. Its additions, products, fixed rational
divisions, literal permutations, and all contraction sums are rounded
outwards at each operation. Matrix contractions are accumulated by explicit
interval sums, not by an unchecked floating BLAS reduction. Reverse
accumulation applies the exact sum and product derivatives with the same
outward operations. Induction over this finite arithmetic graph encloses
the original coefficient and its tangent gradient. The coefficient normal
solve uses the exact rational flat right inverse; it is not a nonlinear
approximate saddle solve.

The implementation checks the extended-precision format and finite input
ranges. No overflow or unbounded denominator occurs in the circuit.
Converting its binary endpoints outwards to denominator $2^{128}$ produces
the supplied platform-independent integer arrays. The final matrix
inequalities below are then checked independently using only Python
integers. The finite circuit, input inclusions, and all endpoint arrays
are supplied for regeneration. The argument uses the inherited
source-extraction identity, rather than purporting to reprove that identity
from numerical samples.
\end{proof}

\section{The full signed Gram is regular, and its signature is retained}
\label{v30:regularity}

\begin{theorem}[Original signed-source regularity at the simultaneous mean zero]
\label{v30:signed}
Throughout the V29 root cube, \eqref{v30:full-source} is nonsingular. The
first source has rank 78 and the complete graded source has rank 107.
The active Gram $C^*K_0^{-1}C$ is positive definite. In the convention
listing positive and negative dimensions, the inertias are
\begin{equation}
 \boxed{\operatorname{inertia}(\Gamma_*)=(96,11),\qquad
 \operatorname{inertia}(C^*K_0^{-1}C)=(78,0),\qquad
 \operatorname{inertia}(D_0)=(18,11),}
 \label{v30:inertias}
\end{equation}
where $D_0$ is the original weak Schur complement. These are inertias of
auxiliary source-response forms, not spacetime signatures or assertions
about the physical Hamiltonian energy.
\end{theorem}
\begin{proof}
The supplied binary rational matrices $S,R_C,R_G$ are analytical
preconditioners. Their floating construction proposes them; their stated
values define exact rational matrices. Put
\[
 \widetilde F=F_*S,\qquad
 \widetilde\Gamma=\widetilde F^*K_0^{-1}\widetilde F.
\]
The integer endpoint calculation verifies, over the entire source enclosure,
\begin{equation}
 \begin{gathered}
 \|I-R_CS\|_\infty<3\,10^{-14},\
 \boxed{\|I-R_G\widetilde\Gamma\|_\infty<2\,10^{-11},\qquad
                  \|\widetilde\Gamma^{-1}\|_\infty<417.}
 \label{v30:margins}
 \end{gathered}
\end{equation}
Here $R_C$ is the fixed triangular QR factor used to propose $S$, not an
assumed exact factorization of the unknown source. The first inequality
proves that $S$ is invertible. The second proves that
$R_G\widetilde\Gamma$ is invertible by its Neumann series, hence that
$\widetilde\Gamma$ and $\Gamma_*$ are invertible. The displayed inverse
bound is $\|R_G\|_\infty/(1-q)$ with the complete verified $q$, rather
than the rounded threshold $2\,10^{-11}$. Its computed upper bound is less
than $416.320$. It is in these declared preconditioned coordinates, not a
cutoff-uniform estimate in an unmentioned norm.

For clarity, the independent integer multiplication uses interval centers
and radii. If dyadic numerator endpoints are $L_A,U_A$, put
$C_A=L_A+U_A$ and $E_A=U_A-L_A$, and similarly for $B$. The exact product
center has numerator $C_AC_B$ and its radius is bounded by
\[
 |C_A|E_B+E_A|C_B|+E_AE_B
\]
at denominator $4\,2^{256}$. Round the lower numerator down and the upper
numerator up to the common output denominator. This proves each matrix
product enclosure by the triangle inequality, with no floating product in
the acceptance test. Every reported norm comparison is an integer
cross-product inequality.

To certify the inertia, use the supplied fixed binary rational eigenvector
proposals $Q_G,Q_A$. Their exact orthogonality defects are less than
$10^{-12}$, so both are invertible. In the full congruence
$Q_G^*\widetilde\Gamma Q_G$, every diagonal interval is separated from the
sum of the absolute off-diagonal intervals by more than $4/1000$; 96
diagonals are positive and 11 negative. In the active congruence the
corresponding margin exceeds $249/1000$ and all 78 diagonals are positive.
The first 78 columns of $S$ have zero entries in its last 29 rows, verified
exactly, so that active congruence concerns precisely $C^*K_0^{-1}C$.

These diagonal-dominance tests determine inertia without an eigenvalue
rounding threshold. Multiply the off-diagonal entries of an actual symmetric
matrix by $t\in[0,1]$. Strict absolute diagonal dominance precludes a null
vector: its component of largest absolute value contradicts the equation
in that row. Thus no eigenvalue crosses zero along this homotopy, and the
inertia is that of the signed diagonal. Invertible congruence preserves it.
Ordinary block elimination then gives the weak inertia by subtraction of
the 78 positive active dimensions.

Invertibility of the full Gram implies independence of its 107 source
columns. The universal first-source cancellation gives kernel containing
constant lapse and every curl-free shift, so the first-source rank is at
most $108-30=78$. Its 78 independent selected columns attain this bound.
They therefore span the active quotient, and together with the stated weak
tests supply the full normalized multiplier complement. No negative Cartan
direction has been projected out.
\end{proof}

\begin{corollary}[The missing higher signed chart is populated]
\label{v30:chart}
For any analytic exactly prepared curve with first canonical field $Y_*$,
initial multiplier $e+O(a)$, and the retained higher source calculus, its
normalized multiplier Hessian has exactly one radial null direction for
all sufficiently small nonzero $a$. In a fixed active/weak complement,
\begin{equation}
 M_{\mathcal H}(a)=D_a\Gamma(a)D_a,\qquad
 D_a=\diag(aI_{78},a^2I_{29}).
 \label{v30:graded}
\end{equation}
Here $\Gamma(0)$ equals $\Gamma_*$ up to a fixed invertible congruence.
The same conclusion holds on a sufficiently small neighborhood of $Y_*$
within the first-field families used below.
\end{corollary}
\begin{proof}
Project the selected active multiplier representatives to the physical
active slice by subtracting their constant-lapse and curl-free parts.
Their leading $J_1$ columns do not change. The full weak basis is retained.
The universal weak first-source identity and \cref{v20:graded-source}
give exactly the weights in \eqref{v30:graded}, independently of the
second canonical coefficient and first multiplier tilt. The leading signed
Gram is nonsingular by \cref{v30:signed}; analytic continuity proves
invertibility for small nonzero amplitude. Homogeneity supplies the radial
null vector $\lambda$, which is transverse to the mean-lapse-zero
complement since $\overline{\lambda^N}=1$. This is the sole null direction.
Continuity also gives the stated first-field neighborhood, without a claim
that a specified large parameter region or another cutoff lies in it.
\end{proof}

\section{Exact initial-rate realization, not only a quartic drift coefficient}
\label{v30:rate-selection}

Let $\mathcal E_0$ be the 104-dimensional multiplier space with all four
component means zero, and let $\mathcal H_0$ be the 107-dimensional space
with only the lapse mean zero. Its complement is the constant-lapse line.
Write $\Pi_0$ for the physical projection onto $\mathcal E_0$ and $R_c$
for the three constant-shift covector tests. Throughout this section a
specified $r\in\mathcal H_0$ is finite at the fixed cutoff. It is an
initial-rate target used to select records, not a substituted vector field.

\begin{theorem}[Every finite order-amplitude rate is attained on exact initial records]
\label{v30:rate-theorem}
Fix $r\in\mathcal H_0$. There are analytic original initial records, for
sufficiently small nonzero amplitude $a$, with first field $Y_*$ and first
tilt $\ell_*$ from V29, satisfying
\begin{equation}
 \boxed{\begin{gathered}
 X(a,0)=aY_*+O(a^2),\qquad
 \lambda(a,0)=e+a\ell_*+O(a^2),\\
 \overline{\lambda^N}=1,\quad\overline{\lambda^\beta}=0,\qquad
 P(X(a,0),\lambda(a,0))=0,\\
 \mathsf B(X(a,0),\lambda(a,0))+aM(X(a,0),\lambda(a,0))r=0,\\
 \lambda_t(a,0)=ar.
 \end{gathered}}
 \label{v30:rate-main}
\end{equation}
The last equality is for the normalized original-action continuation.
The construction persists locally with the five free profile parameters of
\cref{v29:five}, and locally in the target $r$. In particular, $r=0$
gives $\mathsf B=0$ and the entire actual initial multiplier velocity zero,
not merely its weak component or its leading coefficient.
\end{theorem}
\begin{proof}
Write the eight profile parameters as $(z,t)$, with three selected coordinates
$z=(a_1,b_2,b_3)$ and five remaining coordinates $t$. The verified first
Poisson block determines the analytic four-mean-zero map $\ell(z,t)$.
For $m\in\mathcal E_0$ apply the original initializer at
\[
 \lambda=e+a\ell(z,t)+a^2m
\]
while retaining its fixed canonical free-record convention. This gives an
analytic exactly constrained family $X(a,z,t,m)$ with first field $Y(z,t)$.
The first-tilt law cancels every quadratic drift row identically in these
parameters. Thus $\mathsf B/a^3$ has an analytic extension through $a=0$.
Also $M/a^2$ extends analytically there, with leading value
$M_2(Y)=J_1(Y)^*K_0^{-1}J_1(Y)$ in the full multiplier coordinates.

Define the 107-component analytic map
\begin{equation}
 \mathcal G(a,z,t,m;r)=
 \begin{pmatrix}R_c\\ \Pi_0\end{pmatrix}
       \frac{\mathsf B+aMr}{a^3}.
 \label{v30:G}
\end{equation}
Coordinates for $\Pi_0$ mean a fixed basis of its 104-dimensional range.
At $a=0$ the complete second-value law gives
\[
 \frac{\mathsf B+aMr}{a^3}\Big|_{a=0}
  =B_3^{\rm base}(z,t)+\mathbb K_1(Y(z,t))m+M_2(Y(z,t))r.
\]
The raw cubic baseline is the original analytic coefficient, including its
unlisted higher action contributions. It is not replaced by its three means.
The constant-shift rows of both $\mathbb K_1m$ and $M_2r$ are zero.
The first identity follows from the compatible first-Poisson constants;
the second follows from the universal weak first-source cancellation.
Therefore the top three components of $\mathcal G(0,\cdot)$ are exactly
V29's $C(z,t)$, independently of $m$ and $r$.

At $t=t_*$ and $z=z_*$, $C=0$ and $D_zC$ is invertible. The restriction
$\mathbb K_1:\mathcal E_0\to\mathcal E_0$ is invertible. There is
therefore a unique finite $m_0(r)$ cancelling the lower 104 components.
The constant-lapse row of this leading coefficient is zero by homogeneity
and $M_2e=0$. At $(0,z_*,t_*,m_0(r))$ the derivative in $(z,m)$ is
\begin{equation}
 \begin{pmatrix}D_zC&0\\ *&\mathbb K_1|_{\mathcal E_0}\end{pmatrix},
 \label{v30:triangular}
\end{equation}
Both diagonal blocks are nonsingular. The analytic implicit-function theorem
supplies $z(a,t;r)$ and $m(a,t;r)$ with $\mathcal G=0$.
At $t_*$, $z(a)-z_*=O(a)$, so the canonical first field and first tilt
remain $Y_*,\ell_*$. All controls are finite locally; their numerical
nonlinear roots are not needed for this existence proof.

Vanishing of the 107 tested rows means that
$\mathsf B+aMr$ is a multiple of the constant-lapse covector. The exact
identities $\langle\lambda,\mathsf B\rangle=0$ and $M\lambda=0$, and
$\overline{\lambda^N}=1$, force that multiple to be zero. Thus the
untested row is recovered without dropping a mean constraint.

By \cref{v30:chart}, these records have their own unique local normalized
exact-action continuation. Its rate obeys $M\lambda_t=-\mathsf B$.
Subtract the identity just proved: $M(\lambda_t-ar)=0$. The only kernel
is $\R\lambda$, and both rates have lapse mean zero, so $\lambda_t=ar$.
This proves attainment by initial selection. The evolution law is never
replaced by the target $ar$ after the initial cut. Analytic dependence on
$t,r$ follows from the same nonsingular derivative, and continuity supplies
a smaller common fixed-cutoff signed chart.
\end{proof}

The case $r=0$ strengthens V29's quartic initial drift to an exact zero.
For nonzero $r$, the canonical drift need not remain quartic: the actual
identity is $\mathsf B=-aMr$, whose active rows can begin at order $a^3$.
Those rows are balanced by the original multiplier motion, not deleted.
Small actual rates, rather than vanishing canonical drift by itself, are
the relevant dynamical quantity.

\begin{corollary}[Small original initial geometry]
\label{v30:initial-geometry}
For every fixed target in \cref{v30:rate-theorem}, all original initial
metric curvature components and literal temporal/spatial link curvatures
are $O_r(a)$ at this cutoff. The canonical fields and their first two time
jets obey $X=O(a)$, $X_t=O(a)$, $X_{tt}=O_r(a)$.
\end{corollary}
\begin{proof}
The original field vector has $F(X,e+O(a))=O(a)$ and bounded derivatives
on the fixed stationary chart. Its acceleration is
$D_XF\,F+D_\lambda F\,ar=O_r(a)$. The same original ADM and link
identities used in \cref{v21:initial-geometry} then apply with the small
nonconstant lapse and shift retained. No $\lambda_{tt}$ enters the
initial Riemann tensor. The stationary spatial connection has velocity
$D_Xz\,F+D_\lambda z\,ar=O_r(a)$. Its original ordered curvatures have
the same bound. This gives small unnormalized observations as the fields
approach the flat cut; it does not identify their coefficients divided by $a$.
\end{proof}

\section{A zero initial rate is not the appropriate transport condition}
\label{v30:departure}

Use the actual first canonical motion
\[
 V_*=F_1Y_*+G\ell_*.
\]
For a weak shift test $\nu$, define
\begin{equation}
 T_W[\nu]=Dd_{2,W}(Y_*)[V_*,\nu]
                   +\mathbb K_1(V_*)[\nu,\ell_*].
 \label{v30:T}
\end{equation}
The complete counterpart $T$ uses $D\widehat B_2$ in place of $Dd_{2,W}$.
Its weak restriction is exactly \eqref{v30:T}.

\begin{proposition}[A certified nonzero departure source]
\label{v30:T-nonzero}
For $\nu_\sharp=(0,\nabla_\delta\cos(2\pi x_1))$, at the actual V29 root,
\begin{equation}
 \boxed{-55872<T_W[\nu_\sharp]<-55870.}
 \label{v30:T-value}
\end{equation}
This is the physical spatial-sum coefficient. The squared physical
spatial-mean norm of the test is $27/2$. In particular $T_W\ne0$.
\end{proposition}
\begin{proof}
Enclose all 104 point-gauge tilt coordinates using their exact V29 root
cube, subtract their four means, and restore their original $Q_3$ factors.
Use precisely the same enclosed $Y$ as in \cref{v30:enclosure}. Compute
$V=F_1Y+G\ell$ with the original commuting phase derivatives. The complete
coefficient is obtained from the following homogeneous polynomial rules:
\[
 DP_{\nu,2}(Y)[Z]=\frac{P_{\nu,2}(Y+Z)-P_{\nu,2}(Y-Z)}2,
\]
\[
 \begin{aligned}
 D\Phi_3(Y)[Z]&=\frac{\Phi_3(Y+Z)-\Phi_3(Y-Z)}2-\Phi_3(Z),\\
 d_{2,W}(Y)[\nu]&=D\Phi_3(Y)[G\nu]+DP_{\nu,2}(Y)[F_1Y].
 \end{aligned}
\]
Differentiate this quadratic $d_{2,W}$ by the first rule and add the two
$P_2$ polarizations defining $\mathbb K_1(V)[\nu,\ell]$.
The implementation retains the temporal auxiliary and both logarithmic
contributions of those original envelopes. No derivative of an approximate
inverse is taken. Outward interval evaluation, followed by exact conversion
of the binary endpoints, gives
\[
 \frac{-15726351721733464465}{281474976710656}
 \le T_W[\nu_\sharp]\le
 \frac{-15726351711603531071}{281474976710656}.
\]
These rational endpoints lie strictly between the integers in
\eqref{v30:T-value}. The entire root cube is enclosed. Independently,
the original Torch coefficient evaluator gives the same scalar within its
floating diagnostic accuracy; that comparison is not the sign proof.
Finally the sampled cosine derivative is $(0,-9/2,9/2)$ along the first
coordinate and zero in the other components, giving squared mean $27/2$.
\end{proof}

\begin{lemma}[The leading differentiated consistency target]
\label{v30:derivative}
On the records of \cref{v30:rate-theorem}, with fixed $r$, the weak part of
the differentiated consistency source has expansion
\begin{equation}
 R_W\{\dot{\mathsf B}+\dot M\lambda_t\}(0)
   =a^2R_W\{T+\mathbb K_1(Y_*)r\}+O_r(a^3).
 \label{v30:derivative-source}
\end{equation}
It depends at this order only on $Y_*,\ell_*,r$, not on higher initial
preparation. For the exact-zero-rate choice $r=0$,
\begin{equation}
 \boxed{\|\mathbb P_{\rm cf}\lambda_{tt}^{\beta}(a,0)\|_h
                    \ge c a^{-2}}
 \label{v30:zero-rate-acceleration}
\end{equation}
for all sufficiently small positive $a$.
\end{lemma}
\begin{proof}
For a pure shift the linear un-tilted drift vanishes identically. Its field
derivative at $(aY+O(a^2),e+a\ell+O(a^2))$, divided by $a$, is
$Dd_{2,W}(Y)+[Z\mapsto\mathbb K_1(Z)\ell]_W$.
The exact identity \eqref{v20:derivative} gives
$D_\lambda\mathsf B=a\mathbb K_1(Y)+O(a^2)$ there; its additional
$D_X(M\mu)[F]$ term has order $a^2$. Since $X_t=aV+O(a^2)$ and
$\lambda_t=ar$, the first two contributions give \eqref{v30:derivative-source}.
Also $D_XM=O(a)$, $D_\lambda M=O(a^2)$, so $\dot M=O(a^2)$ and
$\dot M\lambda_t=O(a^3)$. This proves the order assertion with all
terms retained.

Differentiate $M\lambda_t=-\mathsf B$ and restrict to the normalized
complement. For $r=0$, \eqref{v30:graded} gives
\[
 (\lambda_{tt})_W=-a^{-2}D_0^{-1}T_W+O(a^{-1}).
\]
The active component is at most $O(a^{-1})$ and cannot cancel this physical
weak leading component when active representatives are projected as in
\cref{v30:chart}. Invertibility of $D_0$ and
\cref{v30:T-nonzero} give the lower bound. This is a derivative at the
initial cut, not a finite-time curvature-divergence theorem.
\end{proof}

\section{Realizing the canonical differential lift cancels the leading fast forcing}
\label{v30:tracking}

The map $\ell(Y)$ is defined not only on the chosen eight-profile slice:
on the full compatible maximal-first-Poisson-rank chart it is the unique
four-mean-zero solution of the complete quadratic cancellation. The vector
$\mathcal V(Y)=F_1Y+G\ell(Y)$ is tangent to that compatible chart by
\cref{v28:lift}. Its differential is therefore defined even though
$\mathcal V$ need not be tangent to the eight-profile slice or to its
three-mean zero set.

\begin{theorem}[A populated first-jet transport rate]
\label{v30:transport-rate}
The finite multiplier
\begin{equation}
 r_*:=D\ell(Y_*)[V_*]\in\mathcal E_0
 \label{v30:rstar}
\end{equation}
is well-defined by the original coefficient maps and the verified first
Poisson inverse. It satisfies
\begin{equation}
 \boxed{\mathbb K_1(Y_*)r_*=-T.}
 \label{v30:tracking-cancel}
\end{equation}
Selecting the initial records of \cref{v30:rate-theorem} with this target
makes their actual rate exactly $\lambda_t(a,0)=ar_*$. On these records
\begin{equation}
 \boxed{\lambda_{tt}(a,0)=O(a^{-1}),\qquad
                   a^2\lambda_{tt}(a,0)\longrightarrow0.}
 \label{v30:acceleration-improved}
\end{equation}
In particular the entire order-$a^{-2}$ weak acceleration coefficient is
cancelled, not only the test \eqref{v30:T-value}.
\end{theorem}
\begin{proof}
Differentiate the full quadratic cancellation along its compatible tangent
$V_*$. This is exactly the identity of \cref{v28:lift}, giving
\eqref{v30:tracking-cancel}. The verified normalized Poisson inverse proves
existence, uniqueness and finiteness of $r_*$, and differentiation of the
four mean normalizations gives $r_*\in\mathcal E_0$. Thus it is an actual
specified coefficient, not a still unsolved rate assumption. The initial
realization theorem is applicable to it.

By \cref{v30:derivative}, the weak order-two differentiated source now
vanishes. On exactly constrained histories the same complete order-two
identity in all components follows either from the prepared-drift formula,
or by adding the linear constraint completion used in the next proof.
Consequently $\dot{\mathsf B}+\dot M\lambda_t=O(a^3)$ at the cut.
Applying the original two graded inverses in \eqref{v30:graded} gives
$(\lambda_{tt})_A=O(1)$ and $(\lambda_{tt})_W=O(a^{-1})$.
This proves \eqref{v30:acceleration-improved}. No $O(1)$ or $O(a)$
upper bound for the full second multiplier derivative is inferred.
\end{proof}

This selection deliberately permits a small nonzero multiplier velocity.
It does not solve the unproductive requirement that the exact-zero-drift
set from $r=0$ should be invariant. Nor does \eqref{v30:tracking-cancel}
prove that the three-mean surface is invariant under $\mathcal V$ or that
a slow manifold of the nonlinear action has been constructed.

\section{Original-action tracking on a populated initial layer}
\label{v30:layer}

\begin{theorem}[A physical $a^3$ interval with no leading fast departure]
\label{v30:layer-theorem}
Use the transport-selected records of \cref{v30:transport-rate} and
$a>0$. For every fixed finite $S$, their original normalized exact-action
histories exist on $0\le t\le a^3S$ for all sufficiently small $a$.
Uniformly for $0\le s\le S$,
\begin{equation}
 \boxed{\begin{aligned}
 X(a,a^3s)-X(a,0)&=a^4\{sV_*+O_S(a)\},\\
 \lambda(a,a^3s)-\lambda(a,0)&=a^4\{sr_*+O_S(a)\},\\
 \lambda_t(a,a^3s)&=ar_*+O_S(a^2),\\
 \lambda_{tt}(a,a^3s)&=O_S(a^{-1}).
 \end{aligned}}
 \label{v30:layer-main}
\end{equation}
The first two error estimates hold with two $s$ derivatives. Every original
constraint is preserved, and the original metric and literal link
curvatures are $O_S(a)$ on this physical interval. All constants here are
fixed-cutoff constants. No common positive interval as $a\to0$ or $h\to0$
is asserted.
\end{theorem}
\begin{proof}
We spell out the use of the inherited initial-layer method so that no
changed off-constraint equation is mistaken for a new physical law.
Define the fixed linear covector map $\mathscr R$ by
\[
 \langle(n,\beta),\mathscr R P\rangle_h
       =\langle(0,\nabla_\delta n),P\rangle_h,
 \qquad \overline{\mathsf B}=\mathsf B+\mathscr R P.
\]
The bar denotes only this analytical completion, not an average. The flat
identities in \cref{v5:prepared-B2,v20:flat} imply
\[
 \overline{\mathsf B}(0,\lambda)=0,\qquad
 D_{(X,\lambda)}\overline{\mathsf B}(0,e)=0.
\]
On $P=0$, the completed source is precisely $\mathsf B$.

In the mean-lapse-zero complement solve its symmetric matrix equation with
$-\overline{\mathsf B}$ instead of $-\mathsf B$, retaining the original
canonical field vector. Off the constraint surface the exact evolution of
$P$ for this extension is
\begin{equation}
 \dot P=-\mathscr R P+e\langle\lambda,\mathscr R P\rangle_h.
 \label{v30:P-extension}
\end{equation}
Indeed the projected equation implies that
$M\lambda_t+\overline{\mathsf B}$ is a multiple of $e$; pairing with
$\lambda$ gives that multiple, using $\overline{\lambda^N}=1$ and
$\langle\lambda,\mathsf B\rangle=0$. Equation \eqref{v30:P-extension}
is homogeneous linear in $P$, so initial $P=0$ stays zero and the extension
agrees with the original normalized action flow throughout the history.
It is only a device for a bounded coordinate existence proof, not an
operational constraint projection or a replacement field equation.

Write $X_a^0,\lambda_a^0$ for the selected initial records and set
\[
 t=a^3s,\qquad X=X_a^0+a^4\xi,\quad
 \lambda=\lambda_a^0+a^4\eta,\quad\eta\in\mathcal H_0.
\]
On bounded $(\xi,\eta)$ sets, the completed-source difference from its
initial value is $a^5$ times an analytic bounded function: its first
ordinary differential is $O(a)$ and the displacement is $a^4$. Moreover
$[M(X,\lambda)-M(X_a^0,\lambda_a^0)]ar_*=O(a^6)$.
The initial rate is exactly $ar_*$. Subtract its exact initial matrix
identity before inverting. The graded factorization \eqref{v30:graded}
then makes $\lambda_t/a$ analytic through $a=0$; its active correction
has an additional factor $a$. The canonical equation $\xi'=F/a$ likewise
has an analytic extension with limit $V_*$. Thus the scaled vector fields
and their fixed derivatives converge on every bounded set, with error
$O(a)$.

For precision, the limiting scaled equations are
\begin{equation}
 \begin{split}
 \xi'&=V_*,\qquad \eta_A'=r_{*,A},\\
 \eta_W'&=r_{*,W}-D_0^{-1}
  \{sT_W+K_{WA}\eta_A+K_{WW}\eta_W\},\qquad
       \xi(0)=\eta(0)=0,
 \end{split}
 \label{v30:limit-layer}
\end{equation}
where the blocks of $\mathbb K_1(Y_*)$ are in the same fixed physical
active/weak coordinate convention as $D_0$. The linearized completed weak
source on $\xi=sV_*$ is $sT_W$, and its multiplier differential is
$\mathbb K_1(Y_*)$, by \eqref{v20:derivative}. Terms from the moving
symmetric matrix against the initial rate are one order higher, as the
$a^6$ bound shows. This derives every term in \eqref{v30:limit-layer}
from the original source.

Identity \eqref{v30:tracking-cancel} makes its exact solution
$\xi=sV_*,\eta=sr_*$. It exists for every finite $s$. On a compact tube
about this solution up to a fixed $S$, ordinary finite-dimensional ODE
existence and continuous dependence for the convergent analytic vector
fields give the original scaled solution, with $C^2_s$ error $O_S(a)$.
For example subtract the two integral equations, apply the common local
Lipschitz bound and Gronwall, and then differentiate their vector fields
once to obtain the second derivative estimate. The tube stays inside the
stationary and signed charts for sufficiently small $a$, because its
physical displacements are $a^4$ and its leading normalized Gram remains
nonsingular. Equation \eqref{v30:P-extension} returns the solution to the
unchanged constraint flow.

Undoing $t=a^3s$ proves all four lines of \eqref{v30:layer-main}; in
particular $\eta''=O_S(a)$ gives $\lambda_{tt}=a^{-2}\eta''=O_S(a^{-1})$.
Throughout the layer $X,\lambda-e,X_t,\lambda_t$ are $O_S(a)$ and
$X_{tt}=D_XF F+D_\lambda F\lambda_t=O_S(a)$. The original geometric
argument of \cref{v30:initial-geometry} applies at every such time. It
uses no bound on the unneeded second lapse or shift derivatives in the
Riemann tensor. All curvature and time normalizations remain the original
ones.
\end{proof}

The homogeneous rapid operator $-D_0^{-1}K_{WW}$ is retained in
\eqref{v30:limit-layer}. Cancelling its leading forcing is not a stability
estimate for its long rescaled evolution or its nonlinear perturbations.
The result cannot be extended to physical times of order $a^2$, let alone
order one, merely by letting $S$ depend on $a$ in a fixed-$S$ theorem.

\section{Status and reproducible checks}
\label{v30:status}

The missing higher signed source at V29's simultaneous cubic-mean zero is
now verified at that same point. The original initial-rate equation is
therefore populated, not conditional on an untested symmetric Hessian.
A triangular analytic preparation selects initial records with any finite
order-amplitude target rate, including exactly zero complete canonical drift.
The latter choice is not dynamically invariant: a newly certified original
source test detects its leading departure. Realizing the canonical
differential lift instead cancels all leading fast forcing and gives the
actual short-time tracking estimate \eqref{v30:layer-main}.

The next questions concern the rapid homogeneous operator, higher forcing,
and motion of the cubic-mean surface. Instantaneous preparation does not
control them. The full signed matrix is not positive; its negative auxiliary
directions are not identified as physical instabilities. No cutoff beyond
$N=3$, refinement-uniform source estimate, or common positive physical lifespan
is certified here. Nonlinear observation matching and the exact-action Einstein
continuum limit remain unresolved.

\paragraph{Executed verification.}
The 108 original first-source columns and all 29 complete weak columns are
freshly rebuilt over the V29 root enclosure. Independent integer arithmetic
verifies the coordinate inverse, full signed inverse, and two inertia
certificates. The rerun V29 checker evaluates the source, checks an independent
polynomial derivative, and verifies its majorant and root inequalities. Its
stored full Jacobian remains an inherited certificate input, not a fresh
full-Jacobian regeneration. The departure coefficient is rebuilt by original
polynomial polarizations; its enclosed endpoints pass exact rational sign tests.

Additional exact calibrations check the triangular control, radial recovery,
graded orders and layer cancellation; they do not replace the analytical
implicit-function and ODE proofs. Floating coefficients remain diagnostics.
No nonlinear initial root, baseline full cubic vector, physical spacetime
trajectory, or proof-assistant formalization is represented as computed.
All 29 preceding bodies are preserved byte for byte. The package records
source fingerprints, endpoint arrays and executed checks. Verification uses
standard methods \cite{v29:Rump} on the specific original source proved here.
% END IMMUTABLE EXACT-ACTION BRIDGE V30 BODY


% Future iterations append here.
% BEGIN IMMUTABLE EXACT-ACTION BRIDGE V31 BODY
\clearpage
\renewcommand{\thesection}{\arabic{section}}
\renewcommand{\theHsection}{bridgeV31.\arabic{section}}
\setcounter{section}{249}
\section{The rapid homogeneous equation must be tested before extending time}
\label{v31:context}

The populated transport theorem of \cref{v30:layer-theorem} cancels the
leading fast forcing at one and the same verified mixed-profile root. It
retains the homogeneous rapid matrix, and does not prove that its evolution
is bounded. This body evaluates that matrix with the original Poisson and
signed-source coefficients. There is a simple positive real eigenvalue,
paired with a negative one. Its positive eigenvector has a quantitatively
nonconstant shift and is visible to the original reconstructed metric
curvature. The assertion is not inferred from the negative inertia of an
auxiliary form.

This calculation changes the next stability question. A positive conserved
quadratic norm for the entire rapid system is impossible at this root.
Transport selection can nevertheless control the original histories on an
interval longer than any fixed multiple of $a^3$: a direct remainder estimate
proves an $a^3\log(1/a)$ physical interval with uniform $O(a)$ curvature.
This estimate does not replace the physical clock or remove the growing mode.
It pays the exponential amplification of the higher-order remainder.

\paragraph{Inputs, new conclusions, and nonduplication.}
The root and the first Poisson coefficient bank are those of V29. The
complete original higher signed source, its enclosure, the initial-rate
realization, and the transport rate $r_*$ are V30 inputs. All five first
Poisson coefficient matrices are freshly reconstructed from the original
quadratic scalar functional; the higher signed-source enclosure is retained
and its integer matrix checks are rerun, not represented as a new extraction
of all weak source columns. V10 supplies the geometric symbol. The new
objects are a verified real rapid eigenpair, its application to actual
rate-selected histories and their original metric observation, a quantitative
extension of the scaled remainder estimate, and the resulting logarithmic
physical interval. The elementary contraction method is proved below; the
verification methodology is contextualized by the already cited
\cite{v29:Rump}. Every numerical certificate here has $N=3$, $h=1/3$.

\section{The original signed Poisson pencil and its exact coordinate changes}
\label{v31:pencil}

Let $Y_*$ denote exactly the V29 root, and let $\Gamma_*$ be the complete
107-dimensional source Gram of \eqref{v30:full-source}. On its fixed
active/weak coefficient spaces write
\begin{equation}
 \Gamma_* =\begin{pmatrix}A&L\\L^*&C\end{pmatrix},\qquad
 D=C-L^*A^{-1}L,\qquad
 K=I_W^*\mathbb K_1(Y_*)I_W,\qquad
 \mathscr B=-D^{-1}K.                         \label{v31:BK}
\end{equation}
Here $I_W$ is the actual 108-by-29 multiplier synthesis: three constant shifts,
then $\nabla_\delta(e_j-e_{26})$, $0\le j<26$, all with zero lapse.
The physical shift gradient uses $\delta_i=3(S_i-S_i^{-1})$. All bilinear
forms use the same spatial-sum convention; changing all of them to physical
spatial means multiplies $D$ and $K$ by the same factor and leaves
$\mathscr B$ unchanged. V30 proves that $A$ and $D$ are invertible, with
inertias $(78,0)$ and $(18,11)$ respectively. The 29-dimensional matrix in
\eqref{v31:BK} is precisely the homogeneous matrix in
\eqref{v30:limit-layer}, not the inverse of the full first Poisson matrix.

Let $S$ be the retained upper triangular binary rational whitening matrix,
partitioned as
\[
 S=\begin{pmatrix}S_A&S_{AW}\\0&S_W\end{pmatrix},\qquad
 \widetilde\Gamma=S^*\Gamma_*S.
\]
The Schur complement and skew block transform as
\begin{equation}
 \widetilde D=S_W^*DS_W,\qquad
 \widetilde K=S_W^*KS_W,\qquad
 -\widetilde D^{-1}\widetilde K=S_W^{-1}\mathscr B S_W.
                                                        \label{v31:congruence}
\end{equation}
For the first identity, either expand the three blocks, or minimize the
active quadratic expression at a fixed weak coefficient. The triangular
active translation $S_{AW}$ cancels from its Schur complement. The last
identity then follows by multiplication in the displayed order. This is an
analytical similarity, not a transformation of the action or a projection of
negative directions.

\begin{lemma}[Signed conservation and spectral pairing]
\label{v31:signed-law}
The original rapid matrix obeys
\begin{equation}
 \mathscr B^*D+D\mathscr B=0,\qquad
 e^{s\mathscr B^*}D e^{s\mathscr B}=D.                 \label{v31:D-law}
\end{equation}
Its spectrum, with algebraic multiplicities, is invariant under
$\lambda\mapsto-\lambda$ and complex conjugation. A real eigenvector with
nonzero real eigenvalue is $D$-null. The three constant-shift columns are in
$\ker K$; they are not discarded.
\end{lemma}
\begin{proof}
$K^*=-K$ and $D^*=D$ give \eqref{v31:D-law} by direct multiplication and
ordinary differentiation of the exponential. In particular
$\mathscr B^*=-D\mathscr B D^{-1}$, which proves the opposite-spectrum
identity. Real coefficients give conjugate pairing. If $\mathscr Bv=\gamma v$
with $\gamma\in\R\setminus\{0\}$, the first identity gives
$2\gamma v^*Dv=0$. Constant shifts are first-Poisson null on compatible
first fields by the original translation incidence. Their exact null identity
is also checked in each of the five newly regenerated coefficient matrices.
Conservation of this indefinite form is not a positive stability estimate.
\end{proof}

\section{A certified growing eigenpair at the same mixed-profile root}
\label{v31:eigenpair}

\begin{theorem}[A simple real expanding pair in the original rapid matrix]
\label{v31:eigen-theorem}
At the exact V29 root, $\mathscr B$ has a simple real eigenvalue $\gamma_*$
with
\begin{equation}
 \boxed{0.09479542352<\gamma_*<0.09479542359.}          \label{v31:gamma}
\end{equation}
It also has the simple eigenvalue $-\gamma_*$. Let $v_+$ be the positive
mode's physical weak multiplier, normalized by
$\|v_+^\beta\|_h=1$. Then
\begin{equation}
 \boxed{\|v_+^\beta-\overline{v_+^\beta}\|_h^2>\frac{21}{25}.}
                                                        \label{v31:visible-fraction}
\end{equation}
No assertion about the number or location of all other eigenvalues is needed
or made by this certificate.
\end{theorem}
\begin{proof}
The inputs and acceptance inequalities are specified sufficiently to permit
independent rational reproduction. V29 gives the root cube with radius
$2^{-240}$ and its 350-bit dyadic center. The three free profile coordinates
are $z_1,z_2,z_3$. Its exact first-Poisson bank gives
\[
 \widehat K(z)=\widehat K_0+z_1\widehat K_1+z_2\widehat K_2+
 z_3\widehat K_3+
 \{3z_1^2/16+(z_2^2+z_3^2)/4\}\widehat K_4.
\]
This uses the same $\widehat K$ rational scaling as V20--V29. If $E$ denotes
the bottom 81 rows of $I_W$, its physical weak block is exactly
\begin{equation}
 K(z)=\frac{\sqrt3}{3}E^*\widehat K(z)_{\beta\beta}E.
                                                        \label{v31:weak-bank}
\end{equation}
Indeed the physical bracket is
$\sqrt3 Q_3^{-1}\widehat K Q_3^{-1}$, with
$Q_3=\diag(I_{27},\sqrt3 I_{81})$. This proves the factor in
\eqref{v31:weak-bank}; replacing it by an unscaled bracket would change the
time exponent. Each of the five full rational 108-by-108 bank matrices is
regenerated from $P_2$, and all 58,320 entries agree with the retained bank.
The synthesis is checked against the original phase derivative, including
its orientation and the subtraction of the 26th scalar indicator.

Evaluate \eqref{v31:weak-bank} by outward dyadic arithmetic. The input
$\sqrt3$ is enclosed by integer square bounds, and every rational coefficient
and root coordinate is rounded outwards. For the signed part, use the unchanged
V30 source endpoint array, form $\widetilde\Gamma$, and partition it into
$\widetilde A,\widetilde L,\widetilde C$. A fixed binary rational matrix
$R_A$ satisfies
\[
 q_A:=\|I-R_A\widetilde A\|_\infty<7\,10^{-14}.
\]
It follows, for every admitted matrix, that
\[
 \|\widetilde A^{-1}-R_A\|_\infty
 \le\frac{q_A}{1-q_A}\|R_A\|_\infty.
\]
This yields an entrywise enclosure of its true inverse. Substitution in
$\widetilde C-\widetilde L^*\widetilde A^{-1}\widetilde L$ encloses
$\widetilde D$ in its actual order. In particular no floating Schur
complement is accepted as the exact matrix.

For the eigenpair use the real system, with $j=22$ in zero-based indexing,
\begin{equation}
 \mathcal F(\gamma,v)=
 \begin{pmatrix}(\widetilde K+\gamma\widetilde D)v\\v_j-1\end{pmatrix}=0.
                                                        \label{v31:eigen-system}
\end{equation}
The binary rational center $(\gamma_0,v_0)$, a binary rational 30-by-30
preconditioner $R$, and all endpoint arrays are in the package. For
$r=10^{-6}$, exact integer arithmetic verifies
\begin{equation}
 \begin{split}
 q&:=\sup_{\|(\gamma,v)-(\gamma_0,v_0)\|_\infty\le r}
       \|I-RD\mathcal F(\gamma,v)\|_\infty<0.004468,\\
 \eta&:=\|R\mathcal F(\gamma_0,v_0)\|_\infty<2.172\,10^{-11},\\
 \eta+qr&<r,\qquad
 \frac{\eta}{1-q}<2.182\,10^{-11}.                  \label{v31:eigen-bounds}
 \end{split}
\end{equation}
These bounds are uniform over the original root cube and source enclosures.
For clarity, the Jacobian used here has the exact form
\[
 D\mathcal F(\gamma,v)=
 \begin{pmatrix}\widetilde Dv&\widetilde K+\gamma\widetilde D\\
                 0&e_j^*\end{pmatrix}.
\]
All finite products in \eqref{v31:eigen-bounds} are sums of integer
center--radius products at denominator $2^{128}$, with outward integer
division. Floating eigensolvers propose the center and preconditioner only.

For every actual $(\widetilde D,\widetilde K)$, the map
$(\gamma,v)\mapsto(\gamma,v)-R\mathcal F(\gamma,v)$ maps the displayed
closed cube strictly into itself and is a contraction there. Completeness
therefore supplies its unique real fixed point. The defect bound also makes
$RD\mathcal F$ invertible, so $R$ is invertible and a fixed point solves
\eqref{v31:eigen-system}. Its distance from the center is at most
$\eta/(1-q)$. The resulting eigenvalue enclosure, whose exact rational
endpoints are stored, is contained in \eqref{v31:gamma}.

The invertible bordered Jacobian also proves simplicity, not just existence.
It forces the kernel of $\widetilde K+\gamma_*\widetilde D$ to be
one-dimensional. A generalized eigenvector would solve
$(\widetilde K+\gamma_*\widetilde D)w=-\widetilde Dv$.
Subtracting a multiple of $v$ makes $w_j=0$, giving a nonzero kernel vector
of that bordered Jacobian, a contradiction. Thus the algebraic multiplicity
is one. \Cref{v31:signed-law} supplies the simple negative mate.

Finally the physical weak vector is $I_WS_Wv$. Its lapse is zero, and its
shift norm is positive. Reconstruct the shift and its three means by exact
interval matrix products, using the sharpened radius in
\eqref{v31:eigen-bounds}. If $b= (I_WS_Wv)^\beta$, the complete integer
sums of squares prove
\[
 \frac{\|b-\overline b\|_h^2}{\|b\|_h^2}>0.84395>\frac{21}{25}.
\]
The lower sum uses zero whenever a scalar interval straddles zero; otherwise
it uses the smaller endpoint square. The upper sum uses the larger endpoint
square, and both are divided by the physical spatial volume count 27.
These operations produce rational lower and upper bounds, not rounded
norms. Normalizing $b$ proves \eqref{v31:visible-fraction}.
\end{proof}

\begin{corollary}[A positive full-sector Lyapunov floor is unavailable]
\label{v31:no-positive-metric}
There is no positive definite real matrix $H$ satisfying
$\mathscr B^*H+H\mathscr B\preceq0$. In particular the complete rapid
semigroup cannot be uniformly bounded for all $s\ge0$ in any fixed norm.
This conclusion does not identify a physical Hamiltonian instability.
\end{corollary}
\begin{proof}
Pair the proposed inequality with $v_+$ to obtain
$2\gamma_*v_+^*Hv_+\le0$, which contradicts positivity. Also
$e^{s\mathscr B}v_+=e^{\gamma_*s}v_+$, so any fixed norm grows on this
vector. The statement concerns the auxiliary rapid multiplier equation of
\eqref{v30:limit-layer}; transferring it to actual histories requires the
next argument, not the inertia count alone.
\end{proof}

\section{The growing mode is detected by original metric curvature}
\label{v31:geometry}

\begin{lemma}[A quantitative physical observation of the positive mode]
\label{v31:observation}
With the normalization of \cref{v31:eigen-theorem}, the original geometric
comparison operator of V10 obeys
\begin{equation}
                \boxed{\|\mathfrak E_h(v_+^\beta)\|_h>\frac7{10}.}
                                                        \label{v31:E-lower}
\end{equation}
\end{lemma}
\begin{proof}
On the three-site grid, every nonzero scalar wave component has ordinary
symbol $2\pi k_i$ and phase symbol $3\sqrt3 k_i$, with $k_i\in\{-1,0,1\}$.
V10's exact Fourier identity consequently gives
\[
 \|\mathfrak E_h(b)\|_h^2
 \ge\frac{(2\pi-3\sqrt3)^2}{2}\|b-\overline b\|_h^2.
\]
Use \eqref{v31:visible-fraction}. The elementary rational bounds
$\pi>31415/10000$ and $\sqrt3<17321/10000$ give
\[
 \frac{(2\pi-3\sqrt3)^2}{2}\frac{21}{25}
 >\frac{2479925469}{5000000000}>\frac{49}{100}.
\]
The checker verifies the square bound for $\sqrt3$ and the stated $\pi$
bound by the alternating remainders in Machin's arctangent identity. Thus no
floating transcendental constant enters the lower bound. This also shows why
a spatially constant multiplier mode would not suffice for this conclusion.
\end{proof}

\begin{theorem}[Exponential sensitivity on populated exact-action initial layers]
\label{v31:populated-instability}
Fix a sufficiently small real $\delta$. Apply \cref{v30:rate-theorem} at the
same $Y_*,\ell_*$ with target
\[
 r_\delta=r_*+\delta v_+.
\]
Write $(X^\delta_a,\lambda^\delta_a)$ for the resulting original histories,
and use $\delta=0$ for the transport-selected reference. For every fixed
finite $S$, these histories exist on $0\le t\le a^3S$ for sufficiently
small positive $a$, and uniformly for $0\le s\le S$,
\begin{equation}
 \boxed{a^{-1}\lambda^\delta_{a,t}(a^3s)
       \longrightarrow r_*+\delta e^{\gamma_*s}v_+.}    \label{v31:rate-growth}
\end{equation}
In the original common interpolated coordinate frame their curvature blocks
satisfy
\begin{equation}
 \boxed{\frac{(R_{0i0j}(g^\delta_a)(a^3s)-R_{0i0j}(g^0_a)(a^3s))_{ij}}{a}
 \longrightarrow\delta e^{\gamma_*s}\mathfrak E_h(v_+^\beta).}
                                                        \label{v31:metric-growth}
\end{equation}
The limiting difference has norm greater than
$7|\delta|e^{\gamma_*s}/10$ when $\delta\ne0$. This is a sensitivity result
between different exactly prepared initial records, not a claim that the
unperturbed transport-selected history has a nonzero leading growing component.
\end{theorem}
\begin{proof}
The analytic rate-selection theorem applies to $r_\delta$, and its first
canonical coefficient and first initial tilt are independent of $\delta$.
Reuse the original scaled-equation derivation of V30 with initial rate
$ar_\delta$. Its limiting canonical and active motions are
$\xi=sV_*$ and $\eta_A=sr_{*,A}$, because $v_+$ is purely weak. Write
$\eta_W=sr_{*,W}+w$. The transport identity
$T+\mathbb K_1r_*=0$ gives
\[
 w'=\delta v_++\mathscr B w,\qquad w(0)=0.
\]
Here and below vectors in the weak equation mean their fixed weak
coordinates. Its exact solution is
\[
 w(s)=\delta\frac{e^{\gamma_*s}-1}{\gamma_*}v_+,
 \qquad \eta'_W=r_{*,W}+\delta e^{\gamma_*s}v_+.
\]
This limiting solution is bounded on every fixed finite interval. The same
analytic scaled ODE and constraint-completion argument as V30 therefore
gives the actual histories and convergence with two $s$ derivatives. In
particular $\lambda_t/a=\eta'$ proves \eqref{v31:rate-growth}.

To check the observation, all these histories have, in those intervals,
$X/a\to Y_*$, $(\lambda-e)/a\to\ell_*$ and $X_t/a\to V_*$.
The original canonical equation implies
\[
 X_{tt}/a\longrightarrow F_1V_*+G\{r_*+\delta e^{\gamma_*s}v_+\}.
\]
Indeed $D_XF(0,e)=F_1$ and $D_\lambda F(0,e)=G$; all remaining products
cost another amplitude. In the linearized expression
\[
 R_{0i0j}=\tfrac12(g_{0j,it}+g_{i0,tj}-g_{00,ij}-g_{ij,tt})
                        +\text{Christoffel products},
\]
the spatial and initial-multiplier terms have the same limit for all
$\delta$. The difference in $g_{0i,t}/a$ is
$\delta e^{\gamma_*s}v_{+,i}^\beta$, while the difference in
$g_{ij,tt}/a$ is its original phase symmetric gradient. Christoffel products
are $O_S(a^2)$. Their difference is exactly
\eqref{v31:metric-growth}, with the ordinary spatial derivative retained on
the metric side. \Cref{v31:observation} proves visibility. No affine canonical
change is substituted for the metric read, and no assertion about a
physical invariant Hamiltonian energy is needed.
\end{proof}

For fixed $\delta\ne0$, these compact-$s$ limits also exclude a
simultaneous amplitude-independent bound $C a$ on the original metric
curvature on a fixed positive physical interval for this perturbed family,
with all histories defined there. The reference leading curvature on a fixed
layer is independent of $s$, by the same formula with $\delta=0$.
Choose a finite $S$ with the displayed exponential term larger than that
reference bound and the proposed $C$, then take $a\to0$, so $a^3S<T$.
This contradicts the proposed bound. It is not a nonexistence theorem for
longer continuation, and does not exclude a weaker vanishing estimate, the
unperturbed selected family, or another joint amplitude--cutoff regime.

\section{A quantitative scaled remainder beyond bounded rescaled intervals}
\label{v31:remainder}

The growing mode does not justify stopping at a fixed-$S$ theorem. We can
extend the actual transport-selected solution for a logarithmically growing
rescaled interval, provided the remainder is estimated before taking that
interval. Work henceforth with $r=r_*$. All norms below are fixed
finite-dimensional norms at this cutoff; their constants may depend on the
verified root and the declared coordinate bases.

Let $X_a^0=aY_a$, $\lambda_a^0=e+a\ell_a$ be the analytic exactly prepared
initial records of V30. The functions $Y_a,\ell_a$ are analytic through
$a=0$, with values $Y_*,\ell_*$. Set $z=(\xi,\eta)$ in
\[
 t=a^3s,\qquad X=X_a^0+a^4\xi,\qquad
 \lambda=\lambda_a^0+a^4\eta,\quad\eta\in\mathcal H_0.
\]
Use the same completed source $\overline{\mathsf B}=\mathsf B+\mathscr R P$
and the same homogeneous constraint evolution \eqref{v30:P-extension}.
The resulting scaled vector field is denoted $\mathscr F_a(z)$.

\begin{lemma}[A polynomial growth bound from the original analytic germ]
\label{v31:uniform-remainder}
There are $a_0,d,C>0$, a fixed vector $b=(V_*,r_*)$, and a fixed real matrix
$J_*$ such that, for $0<a<a_0$ and $a^3\|z\|\le d$,
\begin{equation}
 \mathscr F_a(z)=b+J_*z+\mathcal R_a(z),\qquad J_*b=0,
                                                        \label{v31:scaled-form}
\end{equation}
\begin{equation}
 \boxed{\begin{aligned}
 \|\mathcal R_a(z)\|&\le C\{a(1+\|z\|)+a^3\|z\|^2\},\\
 \|D_z\mathcal R_a(z)\|&\le C(a+a^3\|z\|).
 \end{aligned}}                                         \label{v31:R-bounds}
\end{equation}
The canonical and active rows of $J_*$ are zero. On a purely weak multiplier
vector its weak row is $\mathscr B$, so the certified positive mode remains
in this full scaled linearization.
\end{lemma}
\begin{proof}
Introduce the bounded first-field variables $q=(U,L)$ by
$X=aU$, $\lambda=e+aL$. The original flat identities give analytic maps
\[
 F(aU,e+aL)=a\mathcal F(a,q),\qquad
 \overline{\mathsf B}(aU,e+aL)=a^2\mathcal B(a,q).
\]
The second equality uses vanishing of the completed source and its first
joint differential at $(0,e)$, not vanishing of the uncompleted lapse row.
On the fixed normalized complement the original source factorization is
\[
 M_{\mathcal H}(aU,e+aL)=D_a\Gamma(a,q)D_a,
 \qquad D_a=\diag(aI_A,a^2I_W).
\]
The universal weak first-source cancellation supplies this analytic
factorization also off the exact constraint surface. V30's higher-source
independence gives an analytic $\Gamma$ near $(0,Y_*,\ell_*)$ with invertible
value there. Hence $\Gamma$ and its inverse and fixed derivatives are bounded
on a fixed small $q$-neighborhood. This is not a bound for the ungraded
physical inverse as $a\to0$.

Write $q_a=(Y_a,\ell_a)$ and $q=q_a+a^3z$. Define the analytic difference
quotients, retaining their vector arguments,
\[
 B_a[z]=\int_0^1D_q\mathcal B(a,q_a+ta^3z)[z]dt,
 \quad G_a[z]=\int_0^1D_q\Gamma(a,q_a+ta^3z)[z]dt.
\]
Subtract the \emph{exact} initial identity
$\overline{\mathsf B}(a q_a)+aM(a q_a)r_*=0$ before inverting. Then
\[
 \overline{\mathsf B}(a q)-\overline{\mathsf B}(a q_a)=a^5B_a[z],
 \qquad M(a q)-M(a q_a)=a^3D_aG_a[z]D_a.
\]
Here $a q$ is shorthand for $(aU,e+aL)$. With
$E_a=\diag(aI_A,I_W)$ and $R_a=\diag(I_A,aI_W)$, the exact scaled multiplier
row is
\begin{equation}
 \eta'=r_*-E_a\Gamma(a,q)^{-1}E_aB_a[z]
             -a^2E_a\Gamma(a,q)^{-1}G_a[z]R_ar_*.
                                                        \label{v31:exact-scaled-rate}
\end{equation}
This follows from $D_a^{-1}=a^{-2}E_a$ and $D_ar_*=aR_ar_*$.
The formula keeps the change of the symmetric matrix against the initial
rate; it is not just a source-difference approximation.

Let $H_*=D_q\mathcal B(0,Y_*,\ell_*)$.
On the specified neighborhood, Taylor's theorem gives
\[
 B_a[z]=H_*z+O((a+a^3\|z\|)\|z\|),\quad
 \Gamma(a,q)^{-1}=\Gamma_*^{-1}+O(a+a^3\|z\|),
 \quad \|G_a[z]\|\le C\|z\|.
\]
Use these bounds in \eqref{v31:exact-scaled-rate}; its limiting correction is
$-I_WD^{-1}R_WH_*z$. The canonical row is
$\xi'=\mathcal F(a,q)=V_*+O(a+a^3\|z\|)$.
This proves the first bound in \eqref{v31:R-bounds}. Differentiating the
integral formulas and using their bounded second derivatives gives the
second bound; the factor $a^3$ is retained whenever the evaluation point is
differentiated. Terms such as $a^6\|z\|^2$ are bounded by a fixed multiple
of $a^3\|z\|$ under $a^3\|z\|\le d$.

The limiting weak differential applied to $b=(V_*,r_*)$ is
$(T+\mathbb K_1r_*)_W=0$ by the original transport identity. Thus $J_*b=0$.
Its restriction to purely weak shifts is exactly $-D^{-1}K$. All estimates
hold on a fixed bounded $q$-neighborhood; they do not infer a growing-domain
bound from compact-$z$ convergence alone.
\end{proof}

\section{A logarithmically longer physical interval for the selected original flow}
\label{v31:log-layer}

\begin{theorem}[Quantitative logarithmic transport window]
\label{v31:log-theorem}
Let $L=1+\|J_*\|$ in a fixed submultiplicative Euclidean matrix norm, and
fix any $0<\theta<1$. Put
\begin{equation}
 S_a=\frac{1-\theta}{L}\log(1/a),\qquad T_a=a^3S_a.
                                                        \label{v31:times}
\end{equation}
For all sufficiently small positive $a$, the transport-selected original
histories exist on $0\le t\le T_a$, with every original constraint zero.
In scaled coordinates, for $0\le s\le S_a$,
\begin{equation}
 \boxed{\|z_a(s)-sb\|+\|z_a'(s)-b\|+\|z_a''(s)\|
       \le C a e^{Ls}(1+s)^2.}                         \label{v31:log-error}
\end{equation}
In particular, throughout this physical interval,
\begin{equation}
 \begin{split}
 X(a,a^3s)&=X(a,0)+a^4sV_*+O(a^{4+\theta}\log^2(1/a)),\\
 \lambda(a,a^3s)&=\lambda(a,0)+a^4sr_*+
                                  O(a^{4+\theta}\log^2(1/a)),\\
 \lambda_t(a,a^3s)&=ar_*+O(a^{1+\theta}\log^2(1/a)),\\
 \lambda_{tt}(a,a^3s)&=O(a^{-2+\theta}\log^2(1/a)).
 \end{split}                                          \label{v31:physical-log}
\end{equation}
Moreover, with constants independent of $a$ on these intervals,
\begin{equation}
 \boxed{\|\operatorname{Riem}(g_h)\|+
        \|E_h^{\rm link}\|+\|F_h^{{\rm sp},{\rm link}}\|\le Ca.}
                                                        \label{v31:log-curvature}
\end{equation}
These are fixed-cutoff estimates. Although $T_a/a^3\to\infty$, $T_a\to0$;
no positive amplitude-independent or refinement-independent lifespan follows.
\end{theorem}
\begin{proof}
Let $w=z-sb$. Since $J_*b=0$, \eqref{v31:scaled-form} gives
$w'=J_*w+\mathcal R_a(sb+w)$, $w(0)=0$. Bootstrap $\|w\|\le1$.
For $s\le S_a$, $a^3(\|sb\|+1)<d$ when $a$ is small. The first remainder
bound then gives $\|\mathcal R_a(sb+w)\|\le Ca(1+s)^2$.
The variation-of-constants formula and
$\|e^{tJ_*}\|\le e^{Lt}$ imply
\[
 \|w(s)\|\le Ca\int_0^s e^{L(s-u)}(1+u)^2du
                  \le C'a e^{Ls}(1+s)^2.
\]
Its supremum is at most $C''a^\theta\log^2(1/a)\to0$; hence the bootstrap
cannot exit through $\|w\|=1$. It also cannot leave the stationary or
normalized signed chart: $q=q_a+a^3z$ stays in the fixed coefficient
neighborhood, the actual field displacement is $a^4z$, and the normalized
Gram stays invertible. Ordinary finite-dimensional continuation therefore
reaches $S_a$.

The equation gives the same bound for $w'$. Differentiate it once:
\[
 w''=J_*w'+D_z\mathcal R_a(sb+w)(b+w').
\]
The derivative estimate in \eqref{v31:R-bounds}, the boundedness of $b$, and
the already small $w,w'$ prove the same bound for $w''$. This establishes
\eqref{v31:log-error}, before undoing the clock factors. Since
$X=X_a^0+a^4\xi$, $\lambda=\lambda_a^0+a^4\eta$,
$\lambda_t=a\eta'$ and $\lambda_{tt}=a^{-2}\eta''$,
\eqref{v31:physical-log} follows.

The extension used here preserves $P=0$ by the exact homogeneous equation
\eqref{v30:P-extension}. Thus the continued histories are original-action
histories, not trajectories of a modified on-constraint law. Throughout them
$X,\lambda-e,X_t,\lambda_t$ are $O(a)$, and
$X_{tt}=D_XF F+D_\lambda F\lambda_t=O(a)$. The original ADM formulas
use these jets and bounded fixed-cutoff spatial derivatives, but not
$\lambda_{tt}$ in the Riemann tensor. They give the first observation in
\eqref{v31:log-curvature}. The stationary connection and its first time
derivative are $O(a)$ by their original analytic maps, giving the literal
link estimates as in V30. This proof pays a growing rescaled interval
explicitly and never replaces a fixed-$S$ constant by an assumed uniform one.
\end{proof}

\section{The unstable source projection is the next object to control}
\label{v31:next}

The transport identity cancels the leading forcing, not the homogeneous
rapid matrix. The certified mode identifies a necessary component of any
longer-time selection or remainder estimate. For precision, choose right
modes $v_\pm$ for $\pm\gamma_*$. Simplicity and
\eqref{v31:D-law} imply $v_-^*Dv_+\ne0$. Thus
\[
 \ell_+=\frac{v_-^*D}{v_-^*Dv_+},\qquad
 \ell_+\mathscr B=\gamma_*\ell_+,\qquad \ell_+v_+=1
\]
is an actual signed dual observation of the expanding mode. For the linear
inhomogeneous weak equation $w'=\mathscr Bw+f(s)$, its exact component is
\begin{equation}
 \ell_+w(s)=e^{\gamma_*s}
 \left[\ell_+w(0)+\int_0^s e^{-\gamma_*u}\ell_+f(u)du\right].
                                                        \label{v31:unstable-source}
\end{equation}
If $f$ is bounded on $[0,\infty)$, this component is bounded exactly when
\begin{equation}
 \ell_+w(0)=-\int_0^\infty e^{-\gamma_*u}\ell_+f(u)du. \label{v31:compatibility}
\end{equation}
Indeed subtract the full integral in \eqref{v31:unstable-source}; the tail
is bounded by $\|\ell_+f\|_\infty/\gamma_*$. Conversely a nonzero full
bracket produces exponential growth. Nonvanishing of the dual pairing used
above is the usual elementary simple-eigenvalue fact: if a left eigenvector
annihilated its right eigenvector, the latter would lie in the range of the
shifted matrix and have a generalized preimage, contradicting simplicity.

Equation \eqref{v31:compatibility} is an analytical boundary condition for
one mode, not an operational permission to choose future data, add damping,
or reset the multiplier. For the original nonlinear problem the higher
forcing is not supplied as an arbitrary external $f$. Its exact projection,
and any admissible initial preparation realizing the required cancellation,
remain to be calculated. The other rapid spectral directions are not
classified by the single-pair certificate. Neither a bounded full rapid
semigroup nor the needed nonlinear slow manifold has been established.

\paragraph{Status.}
Proved here are a simple real expanding/contracting pair at the same verified
root; a quantitative nonconstant geometric observation of its positive mode;
its excitation by actual nearby rate-selected initial records; the explicit
polynomial remainder bound; and original transport with $O(a)$ curvatures on
an $a^3\log(1/a)$ physical interval. The first-source bank is freshly checked;
the higher signed-source enclosure, original initializer and stationary-chart
calculus remain named inputs. Open requirements include higher unstable-source
compensation, the complete rapid spectrum and its nonlinear coupling,
refinement-uniform selection and signed response, differentiated geometric
identification, and a common positive physical lifespan. No raw-action Einstein
continuum theorem, finite-amplitude blow-up or escape on a fixed physical
interval, or proof-assistant formalization is asserted.

\paragraph{Reproducibility.}
The portable checker encloses the original weak pencil and solves the bordered
real eigenpair system with integer interval products and rational acceptance
inequalities. It reruns the V30 integer signed-matrix certificate. A separate
fresh exact reconstruction checks all five original first-Poisson bank
matrices and the weak synthesis. The symbolic checks retain noncommuting
signed block algebra, the rate and time factors, and the geometric constant.
Floating eigenvalues are exploratory only; only the one enclosed real pair is
promoted to a theorem. No nonlinear initial root beyond the inherited V29
certificate, no full higher-source extraction, and no spacetime integration is
represented as newly performed. All thirty preceding marked mathematical
bodies are retained byte for byte, with input/output hashes and executed
verification reports in the package.
% END IMMUTABLE EXACT-ACTION BRIDGE V31 BODY


% Future iterations append here.
% BEGIN IMMUTABLE EXACT-ACTION BRIDGE V32 BODY
\clearpage
\renewcommand{\thesection}{\arabic{section}}
\renewcommand{\theHsection}{bridgeV32.\arabic{section}}
\setcounter{section}{256}
\section{Select the expanding initial components rather than suppress the original flow}
\label{v32:context}

The single real growing pair of V31 rules out a positive nonincreasing norm
on the whole rapid multiplier system. It does not exhaust that system's
expanding directions. We first certify their complete number at the same
mixed-profile root, retaining the original signed pencil. There are eleven
real expanding dimensions, eleven contracting dimensions, and seven bounded
center dimensions. The center is positive for the original weak Schur form;
its boundedness is proved by a spectral and inertia argument, not by assuming
that an unverified near-zero eigenvalue is exactly zero.

The resulting splitting permits a different use of the initial-rate theorem
of V30. Eleven higher initial-data parameters cancel the expanding part of
an endpoint displacement. The entire subsequent curve still solves the
original normalized action equations. A forward/backward integral equation
bounds its complete nonlinear remainder without an exponentially growing
Green operator. This supplies a physical interval of length $c a^{5/2}$,
with original metric and literal-link curvatures $O(a)$, at the fixed
three-site cutoff. The constant $c>0$ is an existence constant from the
fixed analytic charts, not a newly computed numerical lifespan.

This is a boundary-selected family of initial records. It is neither an
attracting population nor an operational theorem that a microscopic source
performs the selection. In particular the endpoint condition is used only
to select initial data once. No damping, feedback force, multiplier reset,
or new field equation is applied during the history.

\paragraph{Inherited inputs and scope of the addition.}
The original root, full signed chart and analytic initial-rate family are
\cref{v30:chart,v30:rate-theorem}. The completed-source extension,
its exact constraint propagation, and the analytic graded factorization are
\eqref{v30:P-extension} and \cref{v31:uniform-remainder}. The new spectral
certificate uses the complete original $D,K$ enclosures constructed in V31;
it is not a re-extraction of the higher weak columns. The preceding
single-eigenpair certificate and its observation remain unchanged. The new
items are the complete hyperbolic count, the positive central complement,
a parameter-dependent remainder estimate, a uniformly bounded mixed
initial/terminal Green map, and its application to actual prepared records.
All finite matrix data and all evolution claims here have $N=3$, $h=1/3$.
The general matrix-verification methodology is the already cited
\cite{v29:Rump}; the estimates needed below are proved explicitly.

\section{A complete stable--unstable count for the original rapid pencil}
\label{v32:spectrum}

Retain $D,K,\mathscr B=-D^{-1}K$ and the exact physical synthesis of
\eqref{v31:BK}. Thus $D^*=D$ is nonsingular with inertia $(18,11)$,
$K^*=-K$, and
\[
                   \mathscr B^*D+D\mathscr B=0.
\]
Use V31's triangular whitening only as an analytical similarity. Denote the
resulting pencil by $\widetilde D,\widetilde K$. The theorem concerns the
original physical matrix as well as this similar one.

\begin{theorem}[Eleven expanding and eleven contracting dimensions]
\label{v32:spectral-theorem}
The original real rapid space has an invariant direct sum
\begin{equation}
 \R^{29}=E_{\rm u}\oplus E_{\rm s}\oplus E_{\rm c},\qquad
       \dim E_{\rm u}=11,\quad\dim E_{\rm s}=11,\quad\dim E_{\rm c}=7.
 \label{v32:split}
\end{equation}
Every eigenvalue on $E_{\rm u}$ has real part greater than $1/500$;
every eigenvalue on $E_{\rm s}$ has real part smaller than $-1/500$.
All twenty-two hyperbolic eigenvalues are simple over $\C$.
The restriction of $D$ to $E_{\rm c}$ is positive definite. In particular
the center restriction is diagonalizable over $\C$, has purely imaginary
spectrum, and has a bounded two-sided real semigroup. For some finite $C$,
\begin{equation}
 \|e^{-s\mathscr B}|_{E_{\rm u}}\|+
 \|e^{s\mathscr B}|_{E_{\rm s}}\|\le C e^{-s/500},\qquad
 \|e^{s\mathscr B}|_{E_{\rm c}}\|\le C\quad(s\in\R)
 \label{v32:semigroups}
\end{equation}
with the first inequality interpreted for $s\ge0$. No assertion that all
five eigenvalues in the certified near-zero cluster equal zero is needed.
\end{theorem}

\begin{proof}
We give the complete enclosure and then the exact inertia argument. The
package supplies a real binary-rational matrix $V$, an inverse proposal
$R_V$, a block-diagonal real matrix $\Lambda$, and a proposal $R_D$ for
$\widetilde D^{-1}$. One-by-one blocks of $\Lambda$ represent real
centers; two-by-two blocks have the form
$\left(\begin{smallmatrix}x&y\\-y&x\end{smallmatrix}\right)$.
There is also a five-by-five zero block. All entries are exact binary
rationals once written to the certificate, regardless of how proposed.

With original dyadic interval entries for the pencil, integer matrix products
and rational comparisons prove
\begin{equation}
 \begin{split}
 q_V&:=\|I-R_VV\|_\infty<2\,10^{-14},\qquad
 q_D:=\|I-R_D\widetilde D\|_\infty<1.5\,10^{-10},\\
 E&=-\widetilde K V-\widetilde D V\Lambda,\\
 \delta&:=2\frac{\|R_V\|_\infty}{1-q_V}
                 \frac{\|R_D\|_\infty}{1-q_D}\|E\|_\infty
                              <3\,10^{-9}.
 \end{split}                                                    \label{v32:certificate}
\end{equation}
The actual computed upper value of $\delta$ is between
$2.7599\,10^{-9}$ and $2.7600\,10^{-9}$. All comparisons in
\eqref{v32:certificate} use stored integer numerators and denominators.
The inverse bounds follow from the convergent Neumann series, so
\[
 V^{-1}(-\widetilde D^{-1}\widetilde K)V
          =\Lambda+V^{-1}\widetilde D^{-1}E.
\]
Diagonalize each displayed real two-by-two block by the fixed complex
matrix with columns $(1,i)^T,(1,-i)^T$. Its infinity norm is $2$ and
its inverse has infinity norm $1$. Therefore the final error row norm
is at most $\delta$.

The positive-real-part diagonal centers are shown for readability below.
The certificate uses their exact binary-rational values and the complete
negative and conjugate partners, not the rounded table.
\begin{center}\small
\begin{tabular}{rrc}
\toprule
Real part & Absolute imaginary part & Real expanding dimension\\\midrule
$0.094795423552$&$0$&1\\
$0.055265503613$&$0$&1\\
$0.041642597705$&$0.026239295895$&2\\
$0.014307090390$&$0$&1\\
$0.008514231990$&$0.005045457969$&2\\
$0.005352363125$&$0.009018403688$&2\\
$0.002158359497$&$0.003526597911$&2\\\bottomrule
\end{tabular}
\end{center}
There are two further isolated disks about imaginary centers near
$\pm0.133964634i$, and five disks centered at zero. The twenty-four
nonzero disks of radius $\delta$ are pairwise disjoint and disjoint from
the zero cluster, as checked by exact squared-distance comparisons. Each
isolated disk contains exactly one eigenvalue, and the zero cluster contains
five counted algebraically. Here is the elementary counting argument:
scale the error from zero to its actual value. The eigenvalues remain in
the same disjoint disk groups by the row-dominance estimate, while the roots
of the finite characteristic polynomial vary continuously as a multiset.
Thus the number in each group stays equal to its number at the diagonal
matrix. This also proves simplicity of all isolated eigenvalues. The
hyperbolic disks give eleven real dimensions on each side with the stated
$1/500$ margin. The other seven algebraic dimensions have
$|\operatorname{Re}\lambda|<10^{-5}$.

It remains to prove that the last seven directions are truly central, not
seven unverified small instabilities. Let $E_0$ initially denote their
real generalized spectral space. Conservation of the $D$ pairing implies
that $E_{\rm s}$ is totally isotropic, by sending time to $+\infty$;
$E_{\rm u}$ is totally isotropic by sending time to $-\infty$.
The same conservation proves that each is $D$-orthogonal to $E_0$:
the possible polynomial growth on $E_0$ can be bounded by
$C\exp(|s|/1000)$, strictly slower than the hyperbolic decay used in
the appropriate time direction. The pairing between $E_{\rm s}$ and
$E_{\rm u}$ is nondegenerate. Otherwise a nonzero vector in one of these
spaces would be $D$-orthogonal to all three spectral spaces, contradicting
nonsingularity of $D$. Consequently $D$ has inertia $(11,11)$ on their
sum: in paired bases its matrix is
$\left(\begin{smallmatrix}0&H\\H^*&0\end{smallmatrix}\right)$
with $H$ invertible. The orthogonal complementary restriction is therefore
positive definite, of inertia $(7,0)$, because the full inertia is $(18,11)$.
Set $E_{\rm c}=E_0$.

On this space $D^{1/2}\mathscr B D^{-1/2}$ is skew-adjoint in the
corresponding restricted positive metric. It is diagonalizable with imaginary
eigenvalues and its exponential is unitary there. Norm equivalence gives
the center estimate. The simple hyperbolic spectral decomposition gives the
other two estimates. The three exact constant-shift null vectors remain;
the argument neither deletes them nor requires an exact classification of
the other members of the near-zero cluster.
\end{proof}

\section{The full scaled center grows at most linearly}
\label{v32:full-center}

Use the full scaled displacement space
$\mathcal X=\mathcal P_h\oplus\mathcal H_0$, of dimension $324+107=431$,
and the matrix $J_*$ of \cref{v31:uniform-remainder}. Its rows for canonical
variables and active multipliers are zero. Grouping those 402 coordinates
as $x$, and the 29 weak coordinates as $y$, its form is
\begin{equation}
            J_*\binom{x}{y}=\binom{0}{C_*x+\mathscr By}.
 \label{v32:full-J}
\end{equation}
The symbol $C_*$ here is just this fixed linear map, not a constraint map
or the cubic-mean selection.

\begin{lemma}[A full linear splitting with controlled central derivatives]
\label{v32:full-splitting}
There are real invariant projections $P_{\rm u},P_{\rm s},P_{\rm c}$ on
$\mathcal X$, summing to the identity, with ranks $11,11,409$.
The first two spaces are exactly the weak embeddings of
$E_{\rm u},E_{\rm s}$. Their semigroups satisfy the bounds in
\eqref{v32:semigroups}. On the full center space,
\begin{equation}
 \|e^{sJ_{\rm c}}\|\le C(1+|s|),\qquad
                  \|J_{\rm c}e^{sJ_{\rm c}}\|\le C
                  \quad(s\in\R).
 \label{v32:center-bound}
\end{equation}
In particular $b=(V_*,r_*)$, for which $J_*b=0$, has $P_{\rm u}b=P_{\rm s}b=0$.
\end{lemma}
\begin{proof}
The stable and unstable parts of $\mathscr B$ are invertible. A fixed
triangular change subtracts their responses to $C_*x$ and removes that
coupling. On the remaining weak center, the positive-metric diagonalization
of \cref{v32:spectral-theorem} splits it into its zero eigenspace and its
nonzero imaginary eigenspaces. Subtracting the latter's fixed inverse
responses to $C_*x$ leaves rotations. The zero block obeys
$x'=0$, $y_0'=C_0x$, so its nilpotent part squares to zero. This proves
\eqref{v32:center-bound} in fixed equivalent norms. No inverse is taken
on a zero eigenspace. Every eigenvector of $J_*$ at a nonzero hyperbolic
eigenvalue has $x=0$ by \eqref{v32:full-J}, proving the assertion about the
embedded control spaces and their dimensions. The projections are the
associated invariant spectral projections, generally not orthogonal physical
coordinate projections. The last assertion follows since the hyperbolic
restrictions are invertible.
\end{proof}

This lemma does not bound all rapid motion forward in time. It only separates
the parts for which forward, backward, and polynomial estimates are valid.
The full projection $P_{\rm u}$ can involve canonical and active components
through its left covectors; it must not be replaced on arbitrary displacements
by a projection of just the weak coordinate vector.

\section{The exact initial-rate family is available uniformly in the control}
\label{v32:family}

Identify a vector $d\in E_{\rm u}$ with its physical purely weak multiplier
synthesis, and let $\iota d=(0,d)$ be its full scaled embedding. Apply
\cref{v30:rate-theorem} with
\[
                         r(d)=r_*+d.
\]
Denote its analytic initial records by
$X^0_{a,d}=aY_{a,d}$ and $\lambda^0_{a,d}=e+a\ell_{a,d}$.
They satisfy the original constraints and have their own exact initial
rate $a(r_*+d)$. On a fixed small $d$ ball the maps are jointly analytic,
with
\begin{equation}
 q_{a,d}:=(Y_{a,d},\ell_{a,d})=q_*+O(a),\qquad
 q_*=(Y_*,\ell_*),\qquad D_dq_{a,d}=O(a).
 \label{v32:initial-dependence}
\end{equation}
The last statement follows because $q_{0,d}=q_*$ is independent of $d$.

For the completed-source extension already used in V30, put
\[
 X=X^0_{a,d}+a^4\xi,\quad \lambda=\lambda^0_{a,d}+a^4\eta,
 \qquad t=a^3s,\quad z=(\xi,\eta).
\]

\begin{lemma}[The original nonlinear remainder with all rate parameters retained]
\label{v32:parameter-remainder}
There are fixed $d_0,d_1,a_0,C>0$ such that, for
$\|d\|\le d_1$, $0<a<a_0$, and $a^3\|z\|\le d_0$, the exact scaled
extension is
\begin{equation}
 z'=b+\iota d+J_*z+R(a,d,z),\qquad z(0)=0,
 \label{v32:scaled}
\end{equation}
where
\begin{align}
 \|R(a,d,z)\|&\le C\{a(1+\|z\|)+a^3\|z\|^2\},\notag\\
 \|D_zR(a,d,z)\|&\le C(a+a^3\|z\|),\notag\\
 \|D_dR(a,d,z)\|&\le C\{a(1+\|z\|)+a^3\|z\|^2\}.
 \label{v32:R}
\end{align}
Every solution of this extension from its exact initial records preserves
$P=0$ and hence agrees with the original normalized action flow.
\end{lemma}
\begin{proof}
Repeat the exact subtraction in \eqref{v31:exact-scaled-rate}, now with
$q_{a,d}$ and $r_*+d$. In the same notation its multiplier row is
\[
 \eta'=r_*+d-E_a\Gamma(a,q)^{-1}E_aB_{a,d}[z]
       -a^2E_a\Gamma(a,q)^{-1}G_{a,d}[z]R_a(r_*+d),
 \qquad q=q_{a,d}+a^3z.
\]
This is an equality of the original symmetric-Hessian equation after its
exact initial identity has been subtracted. In particular the change of
$M$ against the initial rate is retained. The canonical row remains
$\xi'=\mathcal F(a,q)$. Analyticity of $\Gamma^{-1},\mathcal B,\mathcal F$
on one fixed coefficient neighborhood, together with
\eqref{v32:initial-dependence}, gives the first two estimates precisely as
in V31, uniformly on the fixed $d$ ball. Their limiting constant and linear
parts are $b+\iota d$ and $J_*$: neither $q_*$ nor the leading graded
matrix depends on $d$.

For the last estimate differentiate these same integral difference quotients
with respect to $d$. Differentiation of an evaluation at $q_{a,d}$ brings
$D_dq_{a,d}=O(a)$. Differentiation of $r_*+d$ in the last displayed term
retains its factor $a^2\|z\|$. No ungraded inverse is differentiated by
itself. The bounded coefficient derivatives therefore give
$C a(1+\|z\|)+C a^3\|z\|^2$ as a sufficient bound. The derivatives
and bounds also hold on line segments in a smaller such neighborhood, which
will be used for the Lipschitz estimates below. Finally the completed-source
extension obeys the unchanged homogeneous constraint equation
\eqref{v30:P-extension}. Its initial constraints are exactly zero for every
$d$, so uniqueness of that homogeneous linear equation proves the last claim.
\end{proof}

\section{A bounded mixed initial--terminal Green map}
\label{v32:Green}

Write $J_{\rm u},J_{\rm s},J_{\rm c}$ for the restrictions in
\cref{v32:full-splitting}. We identify $d$ with $\iota d$ in formulas
confined to the full unstable space. For a continuous forcing $f$ on $[0,S]$,
$S\ge1$, consider
\begin{equation}
 w'=J_*w+\iota d+f(s),\qquad
               w(0)=0,\qquad P_{\rm u}w(S)=0.
 \label{v32:mixed-linear}
\end{equation}
Define $\|f\|_{1,S}=\sup_{0\le s\le S}\|f(s)\|/(1+s)$.

\begin{lemma}[Uniform inverse without forward exponential amplification]
\label{v32:Green-lemma}
Problem \eqref{v32:mixed-linear} has a unique pair $(d,w)$. Set
\begin{equation}
 A_S=\int_0^S e^{-uJ_{\rm u}}du
       =J_{\rm u}^{-1}(I-e^{-SJ_{\rm u}}).
\end{equation}
Its control and solution satisfy the exact formulas
\begin{align}
 d&=-A_S^{-1}\int_0^S e^{-uJ_{\rm u}}P_{\rm u}f(u)du,
                                                    \label{v32:control-formula}\\
 P_{\rm u}w(s)&=-\int_s^S e^{-(u-s)J_{\rm u}}
                                      (\iota d+P_{\rm u}f(u))du,\notag\\
 (P_{\rm s}+P_{\rm c})w(s)&=\int_0^s
 e^{(s-u)J_*}(P_{\rm s}+P_{\rm c})f(u)du.             \label{v32:Green-formulas}
\end{align}
With one constant independent of $S\ge1$,
\begin{equation}
 \|d\|+\sup_{s\le S}\frac{\|w(s)\|}{(1+s)^3}
       +\sup_{s\le S}\frac{\|w'(s)\|}{(1+s)^2}
                                      \le C\|f\|_{1,S}. \label{v32:Green-bound}
\end{equation}
The unstable and stable solution components satisfy the sharper bound
$C\|f\|_{1,S}(1+s)$.
\end{lemma}
\begin{proof}
No eigenvalue of $e^{-SJ_{\rm u}}$ is one, so $A_S$ is invertible.
For large $S$ its inverse is uniformly bounded by the convergent Neumann
series and exponential decay; on a remaining compact interval $1\le S\le S_0$
continuity and nonsingularity give a finite common bound. Multiplying the
forward unstable solution at $S$ by $e^{-SJ_{\rm u}}$ gives
\eqref{v32:control-formula}. Backward integration then proves its displayed
solution and its initial value. Uniqueness follows from these two formulas.

Exponential decay gives a uniform bound on
$\int_0^S\|e^{-uJ_{\rm u}}\|(1+u)du$. Hence
$\|d\|\le C\|f\|_{1,S}$. In the backward integral use
$1+u\le(1+s)(1+u-s)$; the same exponential moment proves its linear
weight bound. Forward stable integration has that bound as well. For the
center, \eqref{v32:center-bound} gives
\[
 \int_0^s C(1+s-u)(1+u)du\le C'(1+s)^3.
\]
For its first derivative use $\|J_{\rm c}e^{tJ_{\rm c}}\|\le C$
instead, and retain the endpoint $P_{\rm c}f(s)$. This gives the square
weight. On the hyperbolic components the equation and the linear weight
bound give the stronger linear weight for $w'$. Summing proves
\eqref{v32:Green-bound}. These are finite-dimensional identities for a
linear comparison problem; their use with the actual nonlinear $f$ is the
next step, not an assumption of an externally chosen physical force.
\end{proof}

\section{Exact nonlinear selection on an amplitude-to-the-five-halves interval}
\label{v32:selection}

\begin{theorem}[A longer original-action interval from eleven initial controls]
\label{v32:main}
At the same certified three-site root there are constants $c,a_0,C>0$
with the following property. For every $0<a<a_0$ set
\begin{equation}
                 S_a=c a^{-1/2},\qquad T_a=a^3S_a=c a^{5/2}.
 \label{v32:time}
\end{equation}
There is $d_a\in E_{\rm u}$, $\|d_a\|\le Ca$, and a unique solution
within the controlled ball specified in the proof, such that the initial
records of \cref{v32:family} with rate target $r_*+d_a$ have a normalized
original-action history on $0\le t\le T_a$. With its scaled displacement
$z_a(s)=sb+w_a(s)$,
\begin{equation}
 \boxed{w_a(0)=0,\qquad P_{\rm u}w_a(S_a)=0,\qquad
 \|w_a(s)\|\le Ca(1+s)^3,\quad
 \|w_a'(s)\|\le Ca(1+s)^2.}
 \label{v32:w-bounds}
\end{equation}
Also $\|w_a''(s)\|\le Ca(1+s)$ throughout this interval.
Every original initial and evolved constraint is zero. The original
multiplier rate satisfies
\begin{equation}
 \boxed{\lambda_t(a,0)=ar_*+O(a^2),\qquad
 \lambda_t(a,a^3s)=ar_*+O\bigl(a^2(1+s)^2\bigr)=O(a).}
 \label{v32:rate-bound}
\end{equation}
The selected initial canonical and multiplier records differ from the
$r_*$-selected records by $O(a^3)$, so their first canonical coefficient
and first initial tilt are still exactly $Y_*,\ell_*$.
\end{theorem}
\begin{proof}
Put $z=sb+w$ and use $J_*b=0$. Then the original scaled equation is
\eqref{v32:mixed-linear} with its own forcing
\[
                     f(s)=R(a,d,sb+w(s)).
\]
We solve this as a fixed point using the Green map of
\cref{v32:Green-lemma}. On pairs consisting of $d\in E_{\rm u}$ and a
continuous $w:[0,S]\to\mathcal X$, use the norm
\[
 \|(d,w)\|_*=
    \max\left\{\|d\|,\sup_{s\le S}\frac{\|w(s)\|}{(1+s)^3}\right\}.
\]
The closed ball $\|(d,w)\|_*\le K_0a$ is complete. Choose $K_0$
large enough for the forthcoming Green bound, then choose $c>0$ small.
For $S=c a^{-1/2}$ and sufficiently small $a$, this ball obeys
\[
 \|w(s)\|\le K_0a(1+S)^2(1+s)\le(1+s),\qquad
                         \|sb+w(s)\|\le C_b(1+s).
\]
The first inequality follows by making $K_0c^2$ sufficiently small and
then $a$ small. It keeps $a^3\|z\|<d_0$ and $\|d\|<d_1$.
The first estimate of \eqref{v32:R} gives, with constants independent of $a,S$,
\begin{equation}
 \|f\|_{1,S}\le C\{a+a^3(1+S)\}\le C'a.
 \label{v32:self-map-source}
\end{equation}
The Green bound makes the ball invariant if $K_0$ is chosen large enough.
This choice precedes the small choice of $c$.

For two pairs in the ball, the last two derivative bounds of
\eqref{v32:R} on their connecting segments imply
\begin{align*}
 \|f-\widetilde f\|_{1,S}
 &\le C\{(a+a^3(1+S))(1+S)^2
                   +a+a^3(1+S)\}\|(d,w)-(\widetilde d,\widetilde w)\|_*\\
 &\le C' a(1+S)^2\|(d,w)-(\widetilde d,\widetilde w)\|_*.
\end{align*}
For small $a$, the last factor is at most a fixed multiple of $c^2$.
Choose $c$ so the Green map times this constant is smaller than $1/2$.
The contraction theorem gives the unique fixed point in the stated ball.
It is a solution of the exact scaled ODE, not only of its Taylor polynomial,
and is continuously differentiable by the integral formulas. Equation
\eqref{v32:Green-bound} gives the bound for $w'$.

For the second derivative, $z'=b+w'$ stays uniformly bounded because
$a(1+S)^2$ is small. Consequently
$\|f'(s)\|=\|D_zR(a,d,z)z'\|\le Ca$. On the center, differentiation
of the convolution after one integration by parts gives
\[
 w_{\rm c}''(s)=J_{\rm c}e^{sJ_{\rm c}}P_{\rm c}f(0)
    +\int_0^s J_{\rm c}e^{(s-u)J_{\rm c}}P_{\rm c}f'(u)du
    +P_{\rm c}f'(s).
\]
Its norm is at most $Ca(1+s)$. On the stable and unstable spaces the
linear-weight bounds for $w,w'$, and the equation
$w''=J_*w'+f'$, give that bound too. This proves the asserted second
derivative estimate without a forward exponential majorant.

The fixed point remains in the original stationary and graded signed charts:
$q=q_{a,d}+a^3z$ differs from $q_*$ by $O(a)+O(a^3(1+S))$.
Its physical initial data were exactly constrained by V30. The homogeneous
constraint equation \eqref{v30:P-extension} therefore gives $P=0$ for the
whole constructed solution. Thus it is an original-action history, with
positive lapse and spatial metric after reducing $a_0$.
Undoing the time factor gives
$\lambda_t=a(r_*+w'_\eta)$ and the second estimate in
\eqref{v32:rate-bound}. At $s=0$ the exact rate-selection identity gives
$a(r_*+d_a)$, proving its first estimate. Finally
\eqref{v32:initial-dependence} and $\|d_a\|\le Ca$ give
\[
 \|X^0_{a,d_a}-X^0_{a,0}\|+
 \|\lambda^0_{a,d_a}-\lambda^0_{a,0}\|
                     \le Ca^2\|d_a\|\le Ca^3.
\]
No assertion of uniqueness outside the controlled ball, or of analyticity
at $a=0$ for the endpoint-selected family, is needed.
\end{proof}

\section{The leading correction is determined by the original higher forcing}
\label{v32:correction}

The terminal selection is not an unspecified cancellation of an arbitrary
external source. Its first nontrivial correction can be identified from the
actual analytic germ, although that germ's higher coefficient is not
numerically tabulated here. On every bounded $z$ set define
\[
       R_1(z)=\partial_aR(0,0,z).
\]
The exact formula in the proof of \cref{v32:parameter-remainder} shows
\begin{equation}
                 R_1(z)=c_1+J_1z.                    \label{v32:first-remainder}
\end{equation}
Indeed the nonlinear dependence on the displacement point first brings
$a^3z$, whereas all first-order-in-$a$ changes of the base record and
coefficients give a constant or a linear function of $z$. The term involving
the change of the symmetric matrix against the initial rate starts with
$a^2$ in that scaled formula and is retained at its actual order.
The notation $J_1$ in this section is a derivative of the full scaled vector
field, not the first connection source denoted $J_{1,h}$ earlier.

\begin{proposition}[First initial correction with no expanding exponential]
\label{v32:first-correction}
For the family of \cref{v32:main},
\begin{equation}
 \boxed{\frac{\iota d_a}{a}\longrightarrow
 d_1=-P_{\rm u}c_1-J_{\rm u}^{-1}P_{\rm u}J_1b.}
 \label{v32:d1}
\end{equation}
On every fixed finite rescaled interval, $w_a/a$ converges to the solution
of
\[
 w_1'=J_*w_1+d_1+c_1+sJ_1b,\qquad w_1(0)=0.
\]
Its expanding component is the polynomial
\begin{equation}
                  P_{\rm u}w_1(s)=-sJ_{\rm u}^{-1}P_{\rm u}J_1b.
 \label{v32:no-exponential-first}
\end{equation}
No numerical value or nonvanishing of $d_1$ is asserted.
\end{proposition}
\begin{proof}
At a fixed $s$, the bounds for the fixed point imply $d_a=O(a)$,
$w_a(s)=O(a)$, and hence
$a^{-1}R(a,d_a,sb+w_a(s))\to R_1(sb)$.
The uniform bound \eqref{v32:self-map-source}, divided by $a$, is
$C(1+s)$ on its whole increasing interval. In
\eqref{v32:control-formula} exponential decay therefore permits dominated
convergence as $S_a\to\infty$. Also $A_{S_a}^{-1}\to J_{\rm u}$.
Thus
\[
 \lim a^{-1}\iota d_a
  =-J_{\rm u}\int_0^\infty e^{-sJ_{\rm u}}
                    P_{\rm u}(c_1+sJ_1b)ds,
\]
which is \eqref{v32:d1}. Here
$\int_0^\infty e^{-sJ_{\rm u}}ds=J_{\rm u}^{-1}$ and
$\int_0^\infty s e^{-sJ_{\rm u}}ds=J_{\rm u}^{-2}$, by integration
by parts with exponentially decaying boundary terms. On fixed intervals,
the integral equation and its uniform bounds give the claimed limit for
$w_a/a$. Substitution verifies \eqref{v32:no-exponential-first}, including
its zero initial value. This derives the required first correction from
the original higher forcing and its original preparation chart. It does
not claim that $d_1=0$ for the uncorrected V31 family.
\end{proof}

\section{Original observations on the longer physical interval}
\label{v32:geometry}

\begin{corollary}[Order-amplitude metric and link curvature on the selected histories]
\label{v32:geometry-theorem}
Throughout $0\le t\le c a^{5/2}$, the histories of \cref{v32:main} obey
\begin{equation}
 \boxed{\|\operatorname{Riem}(g_h)\|+
       \|E_h^{\rm link}\|+\|F_h^{{\rm sp},{\rm link}}\|\le Ca.}
 \label{v32:curvature}
\end{equation}
The norms may be taken at any fixed spatial order at this cutoff. Moreover,
\begin{align}
 X(a,a^3s)&=X^0_{a,d_a}+a^4sV_*+O\bigl(a^5(1+s)^3\bigr),\notag\\
 \lambda(a,a^3s)&=\lambda^0_{a,d_a}+a^4sr_*+
                                      O\bigl(a^5(1+s)^3\bigr),\notag\\
 \|\lambda_{tt}(a,a^3s)\|&\le Ca^{-1}(1+s).
 \label{v32:physical-bounds}
\end{align}
Both total physical displacements at the endpoint are $O(a^{7/2})$.
For every fixed $0<p<1/2$, restriction to $s\le a^{-p}$ gives the
stronger rate comparison $\lambda_t=ar_*+o(a)$, uniformly on that interval.
\end{corollary}
\begin{proof}
The displacement bounds and rate bound follow from
\eqref{v32:w-bounds}, since $X,\lambda$ displacements have factor $a^4$
and $\lambda_{tt}=a^{-2}\eta''$. The endpoint sizes follow on inserting
$S_a=c a^{-1/2}$. The initial records have $X=O(a)$,
$\lambda-e=O(a)$, and the displacements are smaller. The original analytic
canonical vector is therefore $F=O(a)$, while \eqref{v32:rate-bound}
gives $\lambda_t=O(a)$. Its exact differentiated equation gives
\[
                X_{tt}=D_XF\,F+D_\lambda F\,\lambda_t=O(a).
\]
This is stronger for the canonical acceleration than a crude bound obtained
by using the whole $w''$ norm. On the fixed grid, every required spatial
derivative has a fixed finite bound. The original interpolated ADM curvature
uses precisely these canonical and multiplier jets, but not
$\lambda_{tt}$. The same explicit curvature calculation as
\cref{v30:initial-geometry} applies at every time, proving its $O(a)$ bound.
The original stationary connection and its first time derivative are $O(a)$
by the unchanged analytic stationary map. Its literal ordered link curvatures
are consequently $O(a)$ as well. All frames and physical normalizations are
retained. Finally $a^2(1+a^{-p})^2=o(a)$ for $p<1/2$, which proves the
last statement. No equality of the curvature coefficients divided by $a$
is inferred.
\end{proof}

\section{What has been selected, and what remains unresolved}
\label{v32:status}

\paragraph{Proved.}
The full original rapid pencil has eleven expanding and eleven contracting
real dimensions and a seven-dimensional positive, bounded center. The
indefinite conservation law has not been replaced by a positive full-sector
law. The analytic original initial-rate family has enough independent
parameters to implement a mixed initial/terminal selection on all expanding
directions. The exact nonlinear fixed-point argument supplies initial
corrections $d_a=O(a)$, initial-record changes $O(a^3)$, and original-action
histories on $[0,c a^{5/2}]$ with all original constraints and $O(a)$ metric
and literal-link curvatures. The leading correction is a specified expression
in the original higher-order germ; its numerical coefficient remains unlisted.

\paragraph{Inherited inputs.}
The root, stationary and initial-constraint maps, analytic rate realization,
and source enclosures are the explicitly identified V29--V31 inputs. The
fresh spectral calculation uses the original dyadic pencil enclosures and
rechecks them through the inherited integer certificates. The nonlinear
proof uses bounded derivatives on their fixed finite chart; it does not
silently assert any bound under spatial refinement.

\paragraph{Remaining requirements.}
The new selected family is not the uncorrected $r_*$ family, is not proved
attracting, and is not an invariant center-stable manifold. Its endpoint
condition is a boundary-value method for choosing initial records, not a
new operational source law. Sensitivity of a physical preparation to finite
errors is not controlled on the entire longer interval by the present
existence theorem. The eleven hyperbolic directions remain present in the
original dynamics. No estimate for the full nonlinear flow at order-one
physical time, no all-cutoff family of these root and spectral certificates,
and no differentiated metric/link matching is proved. In particular
$c a^{5/2}\to0$ and neither a nontrivial fixed-amplitude Einstein continuum
nor the raw-action Einstein bridge is closed here. The next useful problem
is a compatible longer-time or invariant selection in the original
constraint geometry, together with the refinement estimates; it is not
another attempt to make a positive full-sector rapid energy.

\paragraph{Executed verification.}
The new integer/rational checker verifies both inverse preconditioners, the
complete pencil residual, all disk separations and the hyperbolic counts.
The bounded-center conclusion is the written signed-inertia proof, not a
floating decision about near-zero eigenvalues. A separate exact algebra
checker verifies noncommuting pencil similarities, the mixed-endpoint Green
identities, the polynomial first correction, and all amplitude/time powers.
These tests do not replace the analytic fixed-point proof or compute a
nonlinear shooting solution. The retained V31 eigenpair and V30 signed-matrix
checks are rerun, whereas original higher-source extraction is explicitly
inherited. All thirty-one preceding marked mathematical bodies are unchanged,
with hashes in the preservation manifest. No proof-assistant formalization
or numerical original-action spacetime integration is claimed.
% END IMMUTABLE EXACT-ACTION BRIDGE V32 BODY


% Future iterations append here.
% BEGIN IMMUTABLE EXACT-ACTION BRIDGE V33 BODY
\clearpage
\renewcommand{\thesection}{\arabic{section}}
\renewcommand{\theHsection}{bridgeV33.\arabic{section}}
\setcounter{section}{265}
\section{Retain the block structure before estimating central growth}
\label{v33:context}

The interval in \cref{v32:main} is obtained by applying one remainder norm
to all 431 displacement coordinates. The full limiting center may grow
linearly, although its 29-dimensional rapid subsystem has a bounded center.
An order-$a$ remainder treated indiscriminately in all rows then costs two
powers of the rescaled time. The original matrix equation has more structure:
its apparent order-$a$ active--weak response is the off-diagonal part of one
symmetric block inverse. A fixed, amplitude-dependent triangular shear cancels
that part. The remaining slow-row differential is order $a^2$, not order $a$.
This identity is established below before any longer interval is claimed.

Using the different row estimates, the same eleven-dimensional initial-rate
freedom supplies a new boundary-selected family on
\[
                         0\le t\le c a^2.
\]
Every curve solves the original normalized action equations and preserves all
original constraints. Its metric and literal-link curvatures remain $O(a)$.
The positive constant $c$ is independent of the field amplitude but is still a
fixed-three-site constant. The initial correction is $O(a^2)$ in the rate
and $O(a^3)$ in the records. Selection is performed once; the endpoint
condition is not imposed by a feedback force during the history.

\paragraph{Inputs and nonduplication.}
We retain the root and signed chart of \cref{v30:chart}, the analytic rate
realization in \cref{v30:rate-theorem}, the exact scaled matrix equation
\eqref{v31:exact-scaled-rate}, and the complete rapid splitting of
\cref{v32:spectral-theorem}. Their integer certificates are rerun, not counted
as new coefficient extractions. The full-source completion and its homogeneous
constraint equation are exactly \eqref{v30:P-extension}. The new points are
the amplitude-graded Schur shear, its sharper row and derivative bounds, an
anisotropic mixed Green estimate, and a longer-time contraction using the
actual nonlinear source. Boundary-value methods for saddle-type slow motion
are established methodology \cite{v33:GK}; no normally hyperbolic slow manifold
or external lifespan theorem is imported here. The retained rapid center is
not uniformly invertible, and no inverse on it will be used.

All action-specific results in this body have $N=3$, $h=1/3$. The argument
neither constructs an invariant manifold nor extends the unchanged initial
records selected in V32. It selects possibly different records from the same
original analytic rate family for its new endpoint.

\section{An exact amplitude-graded Schur cancellation}
\label{v33:shear}

Keep the physical active multiplier slice of dimension 78 and weak slice of
dimension 29. In the normalized first-field variables $q=(U,L)$, where
$X=aU$ and $\lambda=e+aL$, write the original analytic graded matrix as
\[
 \Gamma(a,q)=\begin{pmatrix}A(a,q)&L(a,q)\\L(a,q)^*&C(a,q)\end{pmatrix},
 \qquad H(a,q)=A(a,q)^{-1}L(a,q).
\]
The use of $L(a,q)$ for a matrix block in this section is distinct from the
multiplier component of $q$. Both $A$ and $\Gamma$ are invertible on the
fixed small coefficient neighborhood supplied by V30. Let
\[
 A_*=A(0,q_*),\quad L_*=L(0,q_*),\quad H_*=A_*^{-1}L_*,\qquad
 E_a=\operatorname{diag}(aI_A,I_W),\quad
 R_a=\operatorname{diag}(I_A,aI_W).
\]
The star on $A_*,L_*,H_*$ denotes evaluation at the reference point, not an
adjoint. All adjoints and block decompositions retain the physical pairings.

\begin{lemma}[The weak inverse cancels from one combined row]
\label{v33:block-identity}
For every admitted $(a,q)$,
\begin{align}
 [I_A\ \ aH(a,q)]E_a\Gamma(a,q)^{-1}E_a
       &=[a^2 A(a,q)^{-1}\ \ 0],\label{v33:identity-one}\\
 [I_A\ \ aH(a,q)]E_a\Gamma(a,q)^{-1}
       &=a[A(a,q)^{-1}\ \ 0].\label{v33:identity-two}
\end{align}
No commutation between the blocks is assumed. The fixed shear
\begin{equation}
 \mathcal T_a(\xi,\eta_A,\eta_W)
       =(\xi,\eta_A+aH_*\eta_W,\eta_W)                 \label{v33:shear-map}
\end{equation}
is invertible, has determinant one, and obeys
$\|\mathcal T_a\|,\|\mathcal T_a^{-1}\|\le1+a\|H_*\|$ in the
corresponding direct-sum norm, with an equivalent harmless fixed constant
for another chosen norm.
\end{lemma}
\begin{proof}
The first block row of $\Gamma$ is $[A\ L]$, so
$[I\ H]\Gamma^{-1}=[A^{-1}\ 0]$. Also
$[I\ aH]E_a=a[I\ H]$. These two equalities prove both identities in
their original matrix order. The off-diagonal part of
$\mathcal T_a-I$ squares to zero. Therefore its inverse changes the plus
sign to a minus sign, and the determinant and bounds follow. No signed block
has been replaced by a positive one.
\end{proof}

The shear is constant along a history. We do not interpret the state-dependent
row $d\eta_A+aH(a,q)d\eta_W$ as an exact differential; it need not be one.
Freezing at $H_*$ supplies an actual coordinate map, and all errors from
$H(a,q)-H_*$ will be estimated. The map is analytical only. The physical
multiplier is always recovered by
\begin{equation}
          \eta_A=\zeta_A-aH_*y,\qquad\eta_W=y.         \label{v33:undo}
\end{equation}
In particular $\zeta_A$ is not substituted for the active lapse--shift
observation.

We record the exact application to the source equation. Use V32's initial
rate parameter $d\in E_{\rm u}$, with purely weak physical synthesis,
$r_d=r_*+d$, and its exactly constrained initial point $q_{a,d}$.
Set $\widehat z=(\xi,\eta_A,\eta_W)$ and
$q=q_{a,d}+a^3\widehat z$. Let $B_{a,d}[\widehat z]$ and
$G_{a,d}[\widehat z]$ be the integral difference quotients of
$\mathcal B$ and $\Gamma$ in V31. The full, exact multiplier row is
\begin{equation}
 \eta'=r_d-E_a\Gamma^{-1}E_a B_{a,d}[\widehat z]
          -a^2E_a\Gamma^{-1}G_{a,d}[\widehat z]R_ar_d. \label{v33:exact-rate}
\end{equation}
The inverse and both difference quotients are evaluated at their actual
arguments. In particular the final term, coming from the change of the
symmetric matrix against the initial rate, is retained.

Write $\delta H=H_*-H(a,q)$ and suppress these arguments only within the
following display. Applying \cref{v33:block-identity} to \eqref{v33:exact-rate}
gives the exact centered active equation
\begin{align}
 \zeta_A'={}&r_{*,A}+aH_*(r_{*,W}+d)
       -a^2A^{-1}R_AB_{a,d}[\widehat z]\notag\\
 &-a\delta H\,R_W\Gamma^{-1}E_aB_{a,d}[\widehat z]
       -a^3A^{-1}R_AG_{a,d}[\widehat z]R_ar_d\notag\\
 &-a^3\delta H\,R_W\Gamma^{-1}G_{a,d}[\widehat z]R_ar_d.
                                                               \label{v33:slow-exact}
\end{align}
Here $R_A,R_W$ restrict a multiplier vector to its block; they are not the
weight matrix $R_a$. Every term in this identity belongs to the original
matrix equation. In particular its last two terms are not assumed zero.

\section{Different bounds for the slow and rapid rows}
\label{v33:row-bounds}

Group the new coordinates as
\[
 x=(\xi,\zeta_A)\in\R^{402},\qquad y=\eta_W\in\R^{29},\qquad z=(x,y).
\]
Let $b_x=(V_*,r_{*,A})$ and $b_y=r_{*,W}$, so $b=(b_x,b_y)$ is the
same limiting transport vector as before. The limiting linearization is
\[
 J_*z=(0,C_*x+\mathscr B y),\qquad J_*b=0.
\]
The matrix $C_*$ is the original fixed coupling of canonical/active
coordinates into the weak row. The constant shear is identity at $a=0$, so
this limiting matrix is unchanged.

\begin{theorem}[A quadratic-amplitude slow differential]
\label{v33:row-theorem}
On fixed small $d$ and coefficient neighborhoods, for $0<a<a_0$ and
$a^3\|z\|\le\rho_0$, the exact transformed extension is
\begin{equation}
 x'=b_x+F_{\rm s}(a,d,z),\qquad
 y'=b_y+d+C_*x+\mathscr B y+F_{\rm f}(a,d,z),\qquad z(0)=0.
                                                        \label{v33:rows}
\end{equation}
Its analytic remainders obey
\begin{align}
 \|F_{\rm s}\|&\le C(a+a^2\|z\|+a^4\|z\|^2),&
 \|D_zF_{\rm s}\|&\le C(a^2+a^4\|z\|),\label{v33:slow-bounds}\\
 \|F_{\rm f}\|&\le C\{a(1+\|z\|)+a^3\|z\|^2\},&
 \|D_zF_{\rm f}\|&\le C(a+a^3\|z\|),\label{v33:fast-bounds}\\
 \|D_z^2F_{\rm f}\|&\le Ca^3,&
 \|D_dD_zF_{\rm f}\|&\le C(a+a^3\|z\|).\label{v33:second-bounds}
\end{align}
One also has
\begin{align}
 \|D_dF_{\rm s}\|&\le C(a+a^2\|z\|+a^4\|z\|^2),\notag\\
 \|D_dF_{\rm f}\|&\le C\{a(1+\|z\|)+a^3\|z\|^2\}. \label{v33:parameter-bounds}
\end{align}
The constants do not depend on $a,d,z$ in the stated ranges. The canonical
part of $D_zF_{\rm s}$ in fact has norm at most $Ca^3$.
\end{theorem}
\begin{proof}
The initial coefficient map has $q_{a,d}=q_*+O(a)$ and
$D_dq_{a,d}=O(a)$ by \eqref{v32:initial-dependence}. The inverse shear is
uniformly bounded. Thus
\[
 \|H(a,q)-H_*\|\le C(a+a^3\|z\|),\quad
 \|B_{a,d}[\widehat z]\|+\|G_{a,d}[\widehat z]\|\le C\|z\|.
\]
These are bounds for analytic maps on one fixed neighborhood, not for the
ungraded multiplier inverse. Each $z$ derivative of an evaluation at
$q_{a,d}+ta^3\mathcal T_a^{-1}z$ has its explicit factor $a^3$.
The two difference quotients are linear in their displayed vector argument
before differentiating that evaluation point. Consequently their first
$z$ derivatives are bounded by $C(1+a^3\|z\|)$ and their second ones
by $Ca^3(1+a^3\|z\|)$. All coefficient and inverse derivatives of fixed
order are bounded on the smaller neighborhood.

In \eqref{v33:slow-exact}, the first varying-source term is $O(a^2\|z\|)$.
The next is bounded by
$a\,C(a+a^3\|z\|)\|z\|$. The last two are bounded by
$Ca^3\|z\|$ and $Ca^3(a+a^3\|z\|)\|z\|$.
Together with the constant $aH_*(r_{*,W}+d)$ these prove the active part
of \eqref{v33:slow-bounds}. The canonical row is
$\mathcal F(a,q_{a,d}+a^3\widehat z)$, whose difference from $V_*$ is
$O(a)+O(a^3\|z\|)$ and whose $z$ derivative is $O(a^3)$.
Differentiating the exact active terms proves the slow derivative bound;
factors such as $a^7\|z\|^2$ are absorbed using $a^3\|z\|\le\rho_0$.
This is the cancellation that is lost in a single row-independent bound.

For the weak row use \eqref{v33:exact-rate} and the limiting differential
$H_0=D_q\mathcal B(0,q_*)$. The leading term is
$-R_W\Gamma_*^{-1}I_WH_{0,W}\widehat z$ in the notation of V31.
It is precisely $C_*x+\mathscr B y$ at $a=0$. Replacing
$\widehat z$ by $\mathcal T_a^{-1}z$ differs from that limiting argument
by a fixed linear $O(a\|z\|)$ term. All other errors have size
$C(a\|z\|+a^3\|z\|^2)$, plus a harmless $Ca$ allowance as displayed.
The exact weak remainder at $z=0$ is zero. The same evaluation-factor
accounting proves both \eqref{v33:fast-bounds} and
$\|D_z^2F_{\rm f}\|\le Ca^3$. There is no $z$ dependence in the fixed
shear itself and hence no missing derivative of a moving coordinate map.

Finally differentiate the same formulas in $d$. Every differentiated base
point brings $D_dq_{a,d}=O(a)$, and differentiating $r_*+d$ only affects
the explicitly retained terms of orders $a$ or $a^2$ and higher. The
limiting matrix $J_*$ is independent of $d$. These facts give
\eqref{v33:parameter-bounds} and the mixed derivative in
\eqref{v33:second-bounds}. The estimates also hold along parameter and
state line segments in a smaller neighborhood, as required for the
contraction below. No new estimate on a cofinal grid family has been used.
\end{proof}

\section{A mixed Green estimate that uses the two different source orders}
\label{v33:Green}

Let $p_{\rm u},p_{\rm s},p_{\rm c}$ be the invariant projections for
$\mathscr B$ supplied by V32, and write $B_j=\mathscr B|_{E_j}$ and
$C_j=p_jC_*$. Only $B_{\rm u}$ and $B_{\rm s}$ will be inverted.
For the full limiting matrix, its unstable projection is
\begin{equation}
 P_{\rm u}(x,y)=(0,p_{\rm u}y+B_{\rm u}^{-1}C_{\rm u}x).
                                                        \label{v33:full-projection}
\end{equation}
It is not $p_{\rm u}y$ on a general displacement.

For $S\ge1$, a continuous $f_{\rm s}$ and a continuously differentiable
$f_{\rm f}$, put
\[
 \|f\|_{\mathcal F_S}=\sup_{s\le S}\|f_{\rm s}(s)\|
       +\sup_{s\le S}\frac{\|f_{\rm f}(s)\|}{1+s}
       +\sup_{s\le S}\|f_{\rm f}'(s)\|.
\]
For $w=(u,v)$ define
\begin{equation}
 \|(d,w)\|_{\mathcal W_S}=\|d\|
 +\sup_{s\le S}\left\{
 \frac{\|u(s)\|}{1+s}+\|u'(s)\|
 +\frac{\|v(s)\|}{(1+s)^2}+\frac{\|v'(s)\|}{1+s}\right\}.
                                                        \label{v33:Wnorm}
\end{equation}
The real initial control $d$ lies in the original $E_{\rm u}$.

\begin{lemma}[Only one secular integration is paid in the slow forcing]
\label{v33:Green-theorem}
The problem
\begin{equation}
 u'=f_{\rm s},\qquad
 v'=C_*u+\mathscr Bv+d+f_{\rm f},\qquad
 u(0)=v(0)=0,\quad P_{\rm u}(u(S),v(S))=0             \label{v33:linear}
\end{equation}
has a unique solution and satisfies
\begin{equation}
                   \|(d,w)\|_{\mathcal W_S}
                         \le C_G\|f\|_{\mathcal F_S}, \label{v33:Green-bound}
\end{equation}
with $C_G$ independent of $S\ge1$.
\end{lemma}
\begin{proof}
First $u(s)=\int_0^sf_{\rm s}(r)dr$. Set
\[
 \widetilde v_j=p_jv+B_j^{-1}C_ju,\qquad
 g_j=p_jf_{\rm f}+B_j^{-1}C_jf_{\rm s},\qquad j={\rm u},{\rm s}.
\]
Then $\widetilde v_{\rm u}'=B_{\rm u}\widetilde v_{\rm u}+d+g_{\rm u}$,
with zero initial and terminal values, while
$\widetilde v_{\rm s}'=B_{\rm s}\widetilde v_{\rm s}+g_{\rm s}$ has zero
initial value. The exact formulas are
\begin{align}
 d&=-\left(\int_0^S e^{-rB_{\rm u}}dr\right)^{-1}
                 \int_0^S e^{-rB_{\rm u}}g_{\rm u}(r)dr,\notag\\
 \widetilde v_{\rm u}(s)&=-\int_s^S e^{-(r-s)B_{\rm u}}
                                           (d+g_{\rm u}(r))dr,\notag\\
 \widetilde v_{\rm s}(s)&=\int_0^s e^{(s-r)B_{\rm s}}g_{\rm s}(r)dr.
                                                        \label{v33:hyperbolic-Green}
\end{align}
V32 already proves the uniform inverse bound for the first parenthesis.
Its stable and unstable exponential estimates give
$\|d\|\le C\|f\|_{\mathcal F_S}$ and
$\|\widetilde v_j(s)\|\le C\|f\|_{\mathcal F_S}(1+s)$.
Undoing the displayed triangular coordinates and using
$\|u(s)\|\le s\|f\|_{\mathcal F_S}$ gives that same bound for
$p_jv$. Their differential equations give the corresponding first-derivative
bound, without differentiating $f_{\rm s}$.

On the actual seven-dimensional weak center,
\[
 v_{\rm c}(s)=\int_0^s e^{(s-r)B_{\rm c}}
                                  (C_{\rm c}u(r)+p_{\rm c}f_{\rm f}(r))dr.
\]
The bounded two-sided semigroup gives its quadratic weight. The needed
improvement for its derivative is obtained by integration by parts:
\begin{equation}
 v_{\rm c}'(s)=e^{sB_{\rm c}}p_{\rm c}f_{\rm f}(0)
   +\int_0^s e^{(s-r)B_{\rm c}}
                   (C_{\rm c}f_{\rm s}(r)+p_{\rm c}f_{\rm f}'(r))dr.
                                                        \label{v33:center-derivative}
\end{equation}
It is bounded by $C\|f\|_{\mathcal F_S}(1+s)$. This formula is valid also
on the zero eigenspace; it takes no inverse there. Combining the components
proves \eqref{v33:Green-bound}. The explicit formulas prove uniqueness.
They retain the canonical/active term in the full endpoint projection
\eqref{v33:full-projection}; replacing it by a projection of $v$ alone would
be a different boundary condition.
\end{proof}

The estimate uses a uniformly bounded slow forcing and a linearly weighted,
differentiable fast forcing. It is not the false assertion that an arbitrary
forcing with weight $(1+s)$ in every row has a quadratic rather than cubic
response. The different hypotheses will now be verified for the original
nonlinear remainder.

\section{An original-action interval of order amplitude squared}
\label{v33:selection}

\begin{theorem}[Row-sensitive nonlinear initial selection]
\label{v33:main}
At the same certified root and fixed cutoff there are $c,a_0,C>0$ such
that, for $0<a<a_0$ and
\begin{equation}
                      S_a=c/a,\qquad T_a=a^3S_a=c a^2,\label{v33:time}
\end{equation}
one can select $d_a\in E_{\rm u}$ with $\|d_a\|\le Ca$ in the original
initial-rate family. Its normalized exact-action history exists on $[0,T_a]$.
In the centered displacement coordinates,
\[
                 z_a(s)=sb+w_a(s),\qquad w_a=(u_a,v_a),
\]
where $w_a(0)=0$ and $P_{\rm u}w_a(S_a)=0$. Uniformly on $[0,S_a]$,
\begin{equation}
 \boxed{\begin{aligned}
 \|u_a(s)\|&\le Ca(1+s),&\|u_a'(s)\|&\le Ca,\\
 \|v_a(s)\|&\le Ca(1+s)^2,&\|v_a'(s)\|&\le Ca(1+s).
 \end{aligned}}                                             \label{v33:main-bounds}
\end{equation}
Every original constraint is preserved. The exact initial physical rate is
$a(r_*+d_a)$; the initial canonical and multiplier records differ from
those with target $ar_*$ by $O(a^3)$. No common positive physical interval
or cutoff-uniform constant is asserted.
\end{theorem}
\begin{proof}
Use the Banach space of $d\in E_{\rm u}$ and $C^1$ functions $w$ with
zero initial value, with the norm \eqref{v33:Wnorm}. The endpoint condition
will be imposed by the linear solution map. Substitute $z=sb+w$ into
\eqref{v33:rows}. Since $J_*b=0$, the equation for $w$ is exactly
\eqref{v33:linear} with
\[
 f_{\rm s}(s)=F_{\rm s}(a,d,sb+w(s)),\qquad
 f_{\rm f}(s)=F_{\rm f}(a,d,sb+w(s)).
\]
These are the original nonlinear terms, not prescribed forcing. Compose
this evaluation with the Green map of \cref{v33:Green-theorem}.

Take a ball $\|(d,w)\|_{\mathcal W_S}\le Ma$ and require
$a(1+S)\le\kappa_0$, with $\kappa_0$ to be chosen after $M$.
If $M\kappa_0$ is small, then
\[
 \|sb+w(s)\|\le C_b(1+s),\qquad
 \|b+w'(s)\|\le C_b,
\]
with $C_b$ fixed independently of $S,a$. The coefficient displacement is
$O(a^3(1+S))=O(a^2)$, so all evaluations lie in the proved coefficient
neighborhood for small $a$.
The slow bound in \eqref{v33:slow-bounds} now gives
\[
 \sup_{s\le S}\|f_{\rm s}(s)\|
 \le C\{a+a^2(1+S)+a^4(1+S)^2\}\le C'a.
\]
The fast value and chain rule give
\[
 \sup_{s\le S}\frac{\|f_{\rm f}(s)\|}{1+s}\le C'a,
 \qquad
 \sup_{s\le S}\|f_{\rm f}'(s)\|
 \le C\{a+a^3(1+S)\}C_b\le C'a.
\]
Thus $\|f\|_{\mathcal F_S}\le C_0a$. Choose $M>2C_GC_0$, and then
reduce $\kappa_0$ so the preceding bounds hold with those fixed constants.
The map sends the ball into itself.

We give the Lipschitz budget explicitly. For two inputs in the ball, let
$\delta$ be their distance in \eqref{v33:Wnorm}. Their point and velocity
differences are bounded respectively by $C(1+s)^2\delta$ and
$C(1+s)\delta$. Applying the row and parameter derivative bounds gives
\begin{equation}
 \|\Delta f\|_{\mathcal F_S}
 \le C\left\{a(1+S)+a^2(1+S)^2
                 +a^3(1+S)^2+a^4(1+S)^3\right\}\delta. \label{v33:contraction}
\end{equation}
For the only term requiring two derivatives, write
\[
 \Delta f_{\rm f}'
 =(D_zF_{{\rm f},1}-D_zF_{{\rm f},2})(b+w_1')
              +D_zF_{{\rm f},2}(w_1'-w_2').
\]
The first difference is bounded by
$Ca^3(1+S)^2\delta+C(a+a^3(1+S))\delta$ using
\eqref{v33:second-bounds}; the second term is bounded by
$C(a+a^3(1+S))(1+S)\delta$. These are the fast terms of
\eqref{v33:contraction}. The slow source difference uses
$D_zF_{\rm s}=O(a^2+a^4(1+S))$, giving its square and cubic terms,
and $D_dF_{\rm s}$ gives smaller covered terms. No derivative of an
uncontrolled input acceleration is needed.

With $S=c/a$, the bracket in \eqref{v33:contraction} is bounded by
$C(c+c^2+a)$ after decreasing $a_0$. Choose $c>0$ and then $a_0$ so
$C_G$ times that quantity is below $1/2$ and $a(1+S)\le\kappa_0$.
The contraction theorem supplies the unique fixed point in this weighted
ball, proving \eqref{v33:main-bounds} and the endpoint condition.
Its right side is analytic in the state, so the solution is as differentiable
in $s$ as needed on this interval. In particular
\[
 \|u_a''(s)\|\le Ca^2,\qquad
 \|v_a''(s)\|\le Ca(1+s),
\]
by differentiating the actual slow row and using
$v_a''=C_*u_a'+\mathscr Bv_a'+f_{\rm f}'$.

Undo the fixed shear. The normalized coefficient point is
$q=q_{a,d_a}+a^3\mathcal T_a^{-1}z_a$, at distance $O(a)+O(a^2)$
from $q_*$. It remains in the stationary and signed charts, and the physical
metric and lapse stay positive after reducing $a_0$. The initial constraints
are exactly zero by V30. The same homogeneous equation
\eqref{v30:P-extension} proves $P=0$ throughout the constructed curve.
The curve therefore solves the unchanged normalized exact-action flow.
Finally $D_dq_{a,d}=O(a)$ and $d_a=O(a)$ give an $O(a^3)$ change of
physical initial records. Their initial rate equals $a(r_*+d_a)$ by the
original rate-realization theorem, not by assigning a new vector field.
\end{proof}

The constants in this proof are actual fixed-chart bounds on the original
analytic maps and the verified spectral spaces. Their numerical values are
not evaluated. The proof neither requires nor asserts that the five
near-zero rapid eigenvalues vanish exactly.

\section{The physical observations and the price of undoing the shear}
\label{v33:geometry}

\begin{corollary}[Original curvature bounds on the longer interval]
\label{v33:geometry-theorem}
The histories of \cref{v33:main} obey, uniformly on $0\le t\le c a^2$,
\begin{equation}
 \boxed{\|\operatorname{Riem}(g_h)\|+
       \|E_h^{\rm link}\|+\|F_h^{{\rm sp},{\rm link}}\|\le Ca.}
                                                        \label{v33:curvature}
\end{equation}
In the original, unsheared variables one has
\begin{align}
 \|\lambda_t(a,a^3s)-ar_*\|&\le Ca^2(1+s),\notag\\
 \|\lambda_t^A(a,a^3s)-ar_{*,A}\|&\le Ca^2,\qquad
 \|X_t(a,a^3s)-aV_*\|\le Ca^2,\label{v33:rates}\\
 \|\lambda_{tt}(a,a^3s)\|&\le Ca^{-1}(1+s).\notag
\end{align}
The total physical displacements are $O(a^3)$ at the endpoint. More precisely,
\begin{align}
 X(a,a^3s)&=X^0_{a,d_a}+a^4sV_*+O(a^5(1+s)),\notag\\
 \lambda_W(a,a^3s)&=\lambda^0_{W,a,d_a}+a^4sr_{*,W}
                                      +O(a^5(1+s)^2),\label{v33:displacements}\\
 \lambda_A(a,a^3s)&=\lambda^0_{A,a,d_a}+a^4sr_{*,A}
                  +O(a^5(1+s)+a^6(1+s)^2).\notag
\end{align}
For each fixed $0<p<1$, restriction to $s\le a^{-p}$ gives
$\lambda_t=ar_*+o(a)$ uniformly. This is a rate comparison, not equality of
metric and link curvature coefficients divided by $a$.
\end{corollary}
\begin{proof}
The physical velocity is $a$ times the scaled displacement derivative.
Undoing \eqref{v33:undo} gives
\[
 \lambda_{t,W}=a(r_{*,W}+v_a'),\qquad
 \lambda_{t,A}=a\{r_{*,A}+u_{a,A}'-aH_*(r_{*,W}+v_a')\}.
\]
The bounds in \cref{v33:main} and $a(1+S_a)=O(1)$ prove
\eqref{v33:rates}. The canonical formula and the second-derivative bounds
at the end of that proof give the other rate estimates. Integrating, or
undoing the state shear directly, proves \eqref{v33:displacements}; its
endpoint powers follow from $S_a=c/a$. The sharper statement for $p<1$
follows from $a^2(1+a^{-p})=o(a)$.

Throughout the interval, $X,\lambda-e,X_t,\lambda_t$ are $O(a)$.
The original canonical equation, rather than the combined slow-row bound,
gives
\[
                     X_{tt}=D_XF\,F+D_\lambda F\,\lambda_t=O(a).
\]
All required spatial norms are equivalent at the fixed cutoff. The original
interpolated ADM curvature depends on precisely these canonical and first
multiplier time jets, not on $\lambda_{tt}$, as checked in V21 and V30.
Its second-derivative terms are $O(a)$ and its Christoffel products $O(a^2)$.
The original stationary connection and its first time derivative are likewise
$O(a)$ by their analytic maps and the displayed physical velocities. The
literal temporal and spatial ordered-link curvatures are therefore $O(a)$.
Inverse metrics and frame comparisons stay bounded on the retained chart.
This proves \eqref{v33:curvature} without regarding $\zeta_A$ as an actual
lapse or shift and without dropping its $-aH_*y$ pullback.
\end{proof}

In particular $T_a/a^{5/2}=c a^{-1/2}\to\infty$. The improvement is for
newly endpoint-selected initial records. It is not obtained by allowing V32's
endpoint to grow in a theorem stated only on a smaller interval, nor by
changing the physical time variable: $t=a^3s$ still denotes the original
clock at the specified physical times.

\section{Why the row estimates alone do not give an order-one interval}
\label{v33:boundary}

The preceding cancellation does not prove that the full center has become
semisimple or that all its higher perturbations are harmless. The following
separate two-coordinate example identifies a limitation of the available
bounds, not a coefficient calculation for the gravitational action.

\begin{proposition}[A sharp time scale for inference from the displayed row orders]
\label{v33:comparison}
Consider the real analytic comparison system
\[
 x'=a^2y,\qquad y'=1+x,\qquad x(0)=y(0)=0.
\]
It has $b=(0,1)$, $J(x,y)=(0,x)$, $Jb=0$, and exactly the strengthened
slow order $D_zF_{\rm s}=O(a^2)$ with zero fast remainder. Nevertheless
\begin{equation}
 x(s)=\cosh(as)-1,\qquad y(s)=a^{-1}\sinh(as),\qquad
                               y'(s)=\cosh(as).       \label{v33:comparison-solution}
\end{equation}
Thus $y'$ is bounded uniformly when $s\le c/a$, but not on any family with
$aS_a\to\infty$. At $S_a=a^{-1}\log(1/a)$ the physical-rate scaling
$a y'(S_a)$ tends to $1/2$, not to zero. The example can be supplemented
with decoupled contracting, expanding and bounded oscillatory blocks without
changing these conclusions.
\end{proposition}
\begin{proof}
Differentiate $y'=1+x$ and use $x'=a^2y$, so $y''=a^2y$ with
$y(0)=0$, $y'(0)=1$. Solving this scalar equation gives all the formulas.
The last rate is $(1+a^2)/2$. Additional decoupled blocks leave the two
coordinates unchanged.
\end{proof}

This example does not prove that $c a^2$ is optimal for the original action.
It proves only that the proved row orders, a bounded rapid center, and
initial control of the already identified fast expanding sector cannot by
themselves imply the next lifespan improvement. A further argument must use
the actual finite-$a$ center coupling, additional cancellations of its
secular source, or a suitable transported nonlinear constraint geometry.
No external positive energy or phase reset supplies that information.

\section{Status, source audit, and reproducible verification}
\label{v33:status}

\paragraph{Proved.}
The original amplitude-graded symmetric equation has the exact combined-row
identity \eqref{v33:identity-one}. Freezing its Schur coefficient gives a
uniformly invertible near-identity coordinate shear; the actual remaining
slow differential is order $a^2$. The row-sensitive mixed Green map retains
the full unstable projection and the un-inverted rapid center. Its nonlinear
application selects eleven higher initial-data components and proves original
histories on $[0,c a^2]$ with exact constraints, $O(a)$ physical rates and
metric/link curvatures, and $O(a^3)$ changes of initial records. The weak rate
comparison is $O(a^2(1+s))$, rather than the previous quadratic time weight.

\paragraph{Inherited inputs.}
The actual root and source enclosures, the analytic stationary and initial
preparation maps, the rate-realization theorem, and the rapid splitting are
the named V29--V32 inputs. The portable check reruns V32's exact spectral acceptance,
its inherited V31/V30 matrix certificates, and its algebra checks. No new
higher-source extraction, numerical value of the lifespan constant, or
nonlinear original-action trajectory is represented as computed here.

\paragraph{Remaining requirements.}
The new family is boundary selected, not an operationally selected attracting
population or an invariant manifold. Arbitrary preparation errors can still
excite the retained expanding modes. The physical interval $c a^2$ still
shrinks, every constant is fixed-cutoff, and no refinement theorem or
nontrivial fixed-amplitude Einstein development follows. Higher physical-time
curvature derivatives and equality of normalized metric/link observations
remain separate. The next useful upstream calculation is the actual slow
center coupling after the Schur cancellation; the comparison in
\cref{v33:comparison} prevents claiming more from the present bounds alone.

\paragraph{Checks and preservation.}
The new exact checker verifies the graded block identity with noncommuting
blocks, the fixed-versus-moving Schur remainder, the near-identity inverse,
the full unstable projection, the mixed initial/terminal formulas, the
central derivative integration, and the time/observation exponent arithmetic.
A separate exact comparison checks \eqref{v33:comparison-solution}. These
are finite algebra checks of the written analytic proof, not a
proof-assistant formalization. All 32 previous marked mathematical bodies
are preserved byte-for-byte. Source and artifact hashes and the results of
checks actually executed are recorded in the package. The search index did
not expose the latest lineage reliably; the supplied V32 LaTeX and archive
were therefore read directly. No conclusion is based on unrelated older
search matches, and no exhaustive literature-priority claim is made.

\begin{thebibliography}{9}
\raggedright
\bibitem{v33:GK}
J. Guckenheimer and C. Kuehn,
\emph{Computing slow manifolds of saddle type},
SIAM Journal on Applied Dynamical Systems \textbf{8} (2009), 854--879.
\href{https://doi.org/10.1137/080741999}{doi:10.1137/080741999}.
This is primary-source context for boundary-value selection in saddle-type
slow--fast problems. Its hypotheses and numerical algorithm are not substituted
for the original-action block identities or the nonlinear proof above.
\end{thebibliography}
% END IMMUTABLE EXACT-ACTION BRIDGE V33 BODY


% Future iterations append here.

% BEGIN IMMUTABLE EXACT-ACTION BRIDGE V34 BODY
\clearpage
\renewcommand{\thesection}{\arabic{section}}
\renewcommand{\theHsection}{bridgeV34.\arabic{section}}
\setcounter{section}{273}
\section{Separate multiplier transport from canonical feedback}
\label{v34:context}

The last comparison in \cref{v33:comparison} shows that the row orders alone
cannot extend the physical interval. It does not identify the coefficient
responsible for a further central growth in the original action. Here we
return to the original, ungraded consistency equation before estimating its
linearization. The derivative of its symmetric matrix cannot be omitted.
It has an exact covector-transport identity. After retaining that derivative,
the full 107-dimensional multiplier block has a uniformly controlled
stable--center--unstable boundary response in its amplitude-weighted norm.
The canonical input to that response remains a separate, explicitly displayed
term.

There is an action-specific algebraic reason for this improvement. The eleven
rapid expanding and eleven contracting dimensions already account for all
eleven negative directions of the original leading signed Gram. In balanced
multiplier coordinates, the remaining 85-dimensional space is positive,
uniformly on a small coefficient neighborhood and down to zero amplitude.
This statement is for the complete multiplier Poisson factor, not just the
29-dimensional weak block of V32. In particular no new hyperbolicity can
appear in its central complement by restoring the active multiplier rows.
It does not assert that the complete 431-dimensional variational system is
stable: its canonical feedback is retained below.

\paragraph{Dependence and nonduplication.}
The analytic stationary chart, normalized original continuation, and initial
rate family are the named inputs of V30. The graded Hessian is the one in
V31--V33, not a chosen positive replacement. The first Poisson rank 104 at
the same mixed-profile root is V29; the signed inertia $(96,11)$ and active
positivity are V30; the rapid splitting and its positive seven-dimensional
center are V32. Those certificates are rerun and retained as inputs. The
new conclusions are the exact moving-Hessian identities, the extension to
all 107 multiplier directions, the uniform nonautonomous block estimate,
and the exact remaining canonical source. The finite inertia method is
standard Hamiltonian linear algebra; \cite{v34:KKS} is contextual primary
literature, not an imported estimate on the present coefficients.

All action-specific constants in this body remain at $N=3$, $h=1/3$.
No new nonlinear lifespan, invariant manifold, cofinal rank family, or
zero-amplitude Einstein development is claimed. An auxiliary boundary
condition in a linear response problem is not a physical feedback rule.

\section{An exact covector identity for the original normalized rate}
\label{v34:covector}

Use a fixed orthonormal synthesis from $\R^{107}$ onto the physical
mean-lapse-zero multiplier complement. Write the normalized multiplier as
$e+\iota u$ and, only in this body, put
\[
 H(X,u)=\mathscr H_h(X,e+\iota u),\quad
 P=H_u,\quad M=H_{uu},\quad F=\mathbb J H_X,
 \quad Q_X=H_{XX},\quad A_X=H_{uX}.
\]
Thus $P$ consists of the 107 restricted original covectors, not a new
constraint map; $X$ has the original 324 canonical coordinates. Set
\begin{equation}
 K=A_X\mathbb J A_X^*,\qquad B=A_XF,\qquad r=-M^{-1}B.
 \label{v34:definitions}
\end{equation}
On the retained one-radial-null chart $M$ is nonsingular. The normalized
original equations are $X_t=F$, $u_t=r$. Their omitted radial equation is
still the original one and is supplied by homogeneity, as in V30. A dot on
a coefficient denotes its total derivative along this actual vector field.
In particular
\[
 \dot M=D_XM[F]+D_uM[r],\qquad
 \dot A_X=D_XA_X[F]+D_uA_X[r].
\]
No stationary or preparation coefficient is frozen in these definitions.

\begin{theorem}[The full multiplier derivative and its covector transport]
\label{v34:derivative}
Throughout this fixed analytic chart,
\begin{align}
 \boxed{M D_ur=-K-\dot M,}\label{v34:rate-derivative}\\
 \boxed{M D_Xr=-\dot A_X-A_X\mathbb J Q_X.}\label{v34:canonical-source}
\end{align}
For a variation $(\xi,\eta)$ of an actual normalized history, let
$p=M\eta$. Its complete linearized equations are exactly
\begin{align}
 \xi_t&=\mathbb J Q_X\xi+\mathbb J A_X^*M^{-1}p,\notag\\
 \boxed{p_t=-K M^{-1}p-(\dot A_X+A_X\mathbb J Q_X)\xi.}
 \label{v34:covector-system}
\end{align}
In particular
\begin{equation}
                  \boxed{\frac{d}{dt}(A_X\xi+p)=0.}
 \label{v34:linear-constraints}
\end{equation}
The last identity is conservation of the variation of the original restricted
constraint, not a projection imposed on the variation.
\end{theorem}
\begin{proof}
Differentiate $Mr=-B$. For a fixed multiplier direction $\eta$,
\[
 D_uB[\eta]=K\eta+D_X(M\eta)[F].
\]
The first term comes from $D_uF=\mathbb J A_X^*$; the second from the
other factor of $B=A_XF$. This is V20's exact drift derivative on the fixed
normalized complement. Symmetry of the three multiplier derivatives of
$H$ gives $(D_uM[\eta])r=(D_uM[r])\eta$. Consequently
\[
 M D_ur[\eta]=-K\eta-D_XM[F]\eta-D_uM[r]\eta,
\]
which proves \eqref{v34:rate-derivative}. In particular the instantaneous
Jacobian is not generally $-M^{-1}K$.

For a fixed canonical direction $\xi$, symmetry of the two canonical
slots of $H_{uXX}$ gives
\[
 D_XB[\xi]=(D_XA_X[F])\xi+A_X\mathbb J Q_X\xi.
\]
Symmetry of the mixed partials also gives
$(D_XM[\xi])r=(D_uA_X[r])\xi$. Substitution in the differentiated
consistency equation proves \eqref{v34:canonical-source}. The original
canonical variational equation is
$\xi_t=\mathbb J Q_X\xi+\mathbb J A_X^*\eta$.
Differentiate $p=M\eta$ and use both identities just established. The
$\dot M\eta$ terms cancel exactly, proving \eqref{v34:covector-system}.
Finally differentiate $A_X\xi+p$ and substitute the two equations. Its
remaining coefficient on $M^{-1}p$ is
$A_X\mathbb J A_X^*-K=0$. This proves
\eqref{v34:linear-constraints} with all terms retained.
\end{proof}

For example, for the multiplier block with canonical input set to zero,
\[
 \eta_t=-M^{-1}(K+\dot M)\eta,
 \qquad p_t=-K M^{-1}p.
\]
The latter equation is an exact change of response coordinates, not an
assumption that $\dot M=0$. Its signed quadratic quantity satisfies
\begin{equation}
        \frac{d}{dt}\langle p,M^{-1}p\rangle
                   =-\langle\eta,\dot M\eta\rangle.
 \label{v34:energy-derivative}
\end{equation}
Indeed the two terms containing $K$ cancel by skew-adjointness and
$\frac{d}{dt}M^{-1}=-M^{-1}\dot M M^{-1}$. Neither signed conservation
nor a positive full-sector energy is inferred from this identity.

\section{Amplitude-balanced coordinates and the complete Poisson factor}
\label{v34:balanced}

Write $X=aU$, $u=aL$ and $q=(U,L)$, with $a>0$ held constant along a
history. On the fixed active/weak decomposition put
\[
 E_a=\diag(aI_{78},I_{29}),\qquad
 D_a=\diag(aI_{78},a^2I_{29})=a^2E_a^{-1}.
\]
The inherited exact factorization and the first Poisson expansion give
\begin{equation}
 M(a,q)=D_a\Gamma(a,q)D_a,
 \qquad K(a,q)=a\mathcal K(a,q).
 \label{v34:graded}
\end{equation}
Both $\Gamma$ and $\mathcal K$ extend analytically through $a=0$ near
$q_*=(Y_*,\ell_*)$. Their symmetries are $\Gamma^*=\Gamma$ and
$\mathcal K^*=-\mathcal K$; at $a=0$ the latter is the original first
Poisson coefficient. The factor of $a$ follows from the flat commuting
constraint identity, with its multiplier dependence retained.

For a physical multiplier variation set $y=E_a^{-1}\eta$. Equations
\eqref{v34:rate-derivative} and \eqref{v34:graded} give the exact block
\begin{equation}
 \boxed{y_t=a^{-3}\mathcal L(a,q(t))y
                  -\Gamma^{-1}\dot\Gamma\,y+f,\qquad
 \mathcal L=-\Gamma^{-1}E_a\mathcal K E_a,\quad
 f=E_a^{-1}D_Xr\,\xi.}
 \label{v34:balanced-equation}
\end{equation}
For an independently supplied canonical input the same formula defines the
corresponding multiplier response. Its rapid Poisson factor obeys
\begin{equation}
                    \boxed{\mathcal L^*\Gamma+\Gamma\mathcal L=0.}
 \label{v34:signed-identity}
\end{equation}
The transformation weights active variations by $a^{-1}$ and weak variations
by one. It is not uniformly equivalent to their unweighted norm.

For clarity, this $\mathcal L$ is a 107-dimensional multiplier coefficient,
not V32's 431-dimensional full displacement matrix $J_*$. The canonical
input $f$ in \eqref{v34:balanced-equation} is not presumed small. Also the
Poisson factor alone is not the complete multiplier Jacobian: the moving
metric term $-\Gamma^{-1}\dot\Gamma$ is part of the original derivative.

\section{An 85-dimensional positive complement for the original multiplier pencil}
\label{v34:spectrum}

\begin{theorem}[Uniform complete-multiplier splitting near the certified root]
\label{v34:complete-splitting}
There are a coefficient neighborhood $\mathcal U$ of $q_*$ and $a_0>0$
such that, for $0\le a\le a_0$ and $q\in\mathcal U$, the original
factor $\mathcal L(a,q)$ has real invariant subspaces
\begin{equation}
 \R^{107}=\mathcal E_{\mathrm s}(a,q)\oplus
          \mathcal E_{\mathrm u}(a,q)\oplus
          \mathcal E_{\mathrm c}(a,q),\qquad
                (\dim\mathcal E_{\mathrm s},\dim\mathcal E_{\mathrm u},
                 \dim\mathcal E_{\mathrm c})=(11,11,85).
 \label{v34:dimensions}
\end{equation}
They have analytic group projections with uniformly bounded fixed-order
coefficient derivatives on a smaller neighborhood. The first two spectra
have real parts respectively less than $-\gamma$ and greater than $\gamma$
for some fixed $\gamma>0$. The restriction of $\Gamma$ to the central
space is uniformly positive:
\begin{equation}
       c_0\|v\|^2\le\langle v,\Gamma(a,q)v\rangle
                              \le C_0\|v\|^2,
                \qquad v\in\mathcal E_{\mathrm c}(a,q).
 \label{v34:positive-center}
\end{equation}
Every central eigenvalue is purely imaginary and semisimple. With uniform
constants on that neighborhood,
\begin{align}
 \|e^{\tau\mathcal L}\mathcal P_{\mathrm s}\|
 +\|e^{-\tau\mathcal L}\mathcal P_{\mathrm u}\|
       &\le C e^{-\gamma_0\tau}\quad(\tau\ge0),\notag\\
 \boxed{\|e^{\tau\mathcal L}\mathcal P_{\mathrm c}\|\le C
                                      \quad(\tau\in\R).}
 \label{v34:frozen-bounds}
\end{align}
These are frozen-coefficient bounds for the actual Poisson factor; they do
not freeze the coefficients in \eqref{v34:balanced-equation}.
\end{theorem}
\begin{proof}
At $(0,q_*)$, write the original graded matrix as
$\Gamma_*=(\begin{smallmatrix}A&L\ L^*&C\end{smallmatrix})$ and set
$H=A^{-1}L$, $D=C-L^*A^{-1}L$. The change
\[
 \mathcal R(y_A,y_W)=(y_A+Hy_W,y_W)
\]
puts its quadratic form into $\diag(A,D)$. Since $E_0\mathcal K_*E_0$
has only its weak--weak block $K_{WW}$ nonzero, direct block inversion gives
\begin{equation}
 \mathcal R\mathcal L(0,q_*)\mathcal R^{-1}
              =\diag(0_{78},-D^{-1}K_{WW}).
 \label{v34:zero-amplitude}
\end{equation}
The weak block is exactly the original rapid pencil certified in V32.
Its eleven stable and eleven unstable dimensions have a separated spectrum;
its seven-dimensional center is positive for $D$. V30 makes $A$ positive.
Consequently the complete central space at this base point has dimension
$78+7=85$ and is positive for $\Gamma_*$. No assertion about five exact
zero weak eigenvalues is needed.

Choose disjoint finite contours enclosing the stable and unstable spectral
groups and a bounded contour enclosing their complement. Matrix inverse
continuity on these contours gives their spectral projections on one common
neighborhood, analytically in the coefficients. This can also be proved
entry by entry using the adjugate divided by the determinant under each
contour integral. Their ranks are constant and their derivatives are bounded
on a compactly contained neighborhood. The stable and unstable spectra stay
a fixed positive distance from the imaginary axis.

The central projection remains close to its base projection, and $\Gamma$
remains close to $\Gamma_*$. Positive definiteness of the base restricted
form therefore proves \eqref{v34:positive-center} uniformly. Identity
\eqref{v34:signed-identity} restricted to that positive space conjugates
the central generator to a skew-adjoint matrix. Thus it is diagonalizable
with imaginary eigenvalues and preserves its positive energy, proving the
central semigroup bound. The base hyperbolic eigenvalues are simple; by
shrinking the neighborhood their isolated spectral disks and bounded
projectors persist, yielding the stated hyperbolic bounds, with a possibly
smaller $\gamma_0$. Alternatively their uniformly stable Lyapunov matrices
give the same estimates.

This is also an inertia check. Each hyperbolic spectral space is isotropic
for $\Gamma$; they pair nondegenerately and together have inertia $(11,11)$.
The complete inertia is $(96,11)$, so their orthogonal complementary
restriction has inertia $(85,0)$. The argument retains, rather than repairs,
all eleven negative directions.
\end{proof}

An equivalent central norm in physical multiplier coordinates is obtained by
undoing $y=E_a^{-1}\eta$ and completing the Schur square:
\begin{equation}
 \langle y,\Gamma y\rangle
   =a^{-2}\langle\eta_A+aA^{-1}L\eta_W,
                    A(\eta_A+aA^{-1}L\eta_W)\rangle
                       +\langle\eta_W,D\eta_W\rangle.
 \label{v34:physical-metric}
\end{equation}
It is positive on the transported central space, not on the whole multiplier
space. The $a^{-1}$ active-coordinate cost cannot be omitted when using the
uniform estimate. Pointwise Schur completion in this display is an equality
of quadratic forms, not an assumption that a moving differential integrates
to coordinates without derivative terms.

\begin{corollary}[Exact count for the leading-coefficient pencil]
\label{v34:frozen-leading}
Retain the original $\Gamma_*$ and $\mathcal K_*=\mathbb K_1(Y_*)$,
without their higher coefficients, and consider the analytical pencil
\[
                 \mathcal L_*^{[0]}(a)
                  =-\Gamma_*^{-1}E_a\mathcal K_*E_a.
\]
For every sufficiently small $a>0$ it has eleven stable dimensions, eleven
unstable dimensions, exactly three semisimple zero eigenvalues, and 41 pairs
of nonzero imaginary eigenvalues, counted with algebraic multiplicity.
This is a statement about the named leading-coefficient pencil, not a claim
that the complete nonlinear action has exactly three zero modes.
\end{corollary}
\begin{proof}
The preceding proof applies to this particular analytic pencil as well.
V29 certifies rank 104 of the first Poisson coefficient on the full
108-dimensional multiplier space, with its four constant null directions.
Removing the constant lapse direction removes one radical vector and leaves
rank 104 and nullity three on the fixed 107-dimensional complement.
Congruence by the nonsingular $E_a$ and multiplication by $\Gamma_*^{-1}$
preserve that nullity for $a>0$. A zero eigenvalue lies in the positive
center and is semisimple by \cref{v34:complete-splitting}. Of its remaining
$85-3=82$ dimensions, reality and the nonzero imaginary spectrum give
41 conjugate pairs. No rate at which those frequencies approach zero, or
simplicity of the additional imaginary pairs, is asserted.
\end{proof}

\section{A nonautonomous multiplier response with the moving Hessian retained}
\label{v34:transport}

The frozen result is not substituted for a time-dependent theorem. Let
$q:[0,T]\to\mathcal U$ be a $C^1$ path with $T\le T_0$ and
$\sup_t\|q_t\|\le L_0$. Keep $a$ constant. All norms below are balanced
multiplier norms. The constants may depend on $T_0,L_0$ and the fixed
coefficient chart, but not on $a$ or on a lower bound for a nonzero central
frequency.

Let $\mathcal P_j(t)=\mathcal P_j(a,q(t))$ be the three instantaneous
projections. For $f$ continuous set $f_j=\mathcal P_jf$ and
$f_{\mathrm f}=f_{\mathrm s}+f_{\mathrm u}$. Prescribe stable and central
initial components $s_0,c_0$ and an unstable terminal component $u_T$ in the
corresponding ranges at their endpoints.

\begin{theorem}[Uniform mixed response of the actual multiplier block]
\label{v34:green}
For all sufficiently small $a>0$, the equation
\begin{equation}
 y_t=a^{-3}\mathcal L(a,q(t))y-\Gamma^{-1}\dot\Gamma\,y+f(t)
 \label{v34:nonautonomous}
\end{equation}
with boundary data
\[
 \mathcal P_{\mathrm s}(0)y(0)=s_0,\quad
 \mathcal P_{\mathrm c}(0)y(0)=c_0,\quad
 \mathcal P_{\mathrm u}(T)y(T)=u_T
\]
has a unique solution. Writing $y_j=\mathcal P_j(t)y(t)$ and
$y_{\mathrm f}=y_{\mathrm s}+y_{\mathrm u}$, it satisfies
\begin{align}
 \|y_{\mathrm c}\|_\infty
  &\le C_{T_0,L_0}\bigl(\|c_0\|+\|f_{\mathrm c}\|_{L^1}\notag\\
  &\hspace{18mm}+a^3(\|s_0\|+\|u_T\|+\|f_{\mathrm f}\|_\infty)\bigr),
 \label{v34:center-estimate}\\
 \|y_{\mathrm f}\|_\infty
  &\le C_{T_0,L_0}(\|s_0\|+\|u_T\|)\notag\\
  &\hspace{7mm}+C_{T_0,L_0}a^3
       (\|f_{\mathrm f}\|_\infty+\|c_0\|+\|f_{\mathrm c}\|_{L^1}),
 \label{v34:fast-estimate}\\
 \|y_{\mathrm f}\|_{L^1}
  &\le C_{T_0,L_0}a^3
       (\|s_0\|+\|u_T\|+\|f_{\mathrm f}\|_\infty
                         +\|c_0\|+\|f_{\mathrm c}\|_{L^1}).
 \label{v34:fast-L1}
\end{align}
The term $-\Gamma^{-1}\dot\Gamma$ and the derivatives of all three
projections are included in these estimates. No center inverse is used.
The boundary condition is a response selection, not a change to the action.
\end{theorem}
\begin{proof}
Fix the base group projections $P_j^*$. On a sufficiently small common
neighborhood the map
\[
                 V(a,q)=\sum_j\mathcal P_j(a,q)P_j^*
\]
is invertible, with uniformly bounded inverse and fixed-order derivatives.
It maps each base range onto its current range. The uniform closeness to
identity proves invertibility by a Neumann series. Thus $y=Vw$ puts
\eqref{v34:nonautonomous} into a fixed direct sum:
\begin{equation}
 w_j'=a^{-3}L_j(t)w_j+\sum_k C_{jk}(t)w_k+\widetilde f_j(t),
 \qquad \|C(t)\|\le C_{L_0}.
 \label{v34:fixed-bundle}
\end{equation}
Here $V^{-1}\mathcal L V$ is block diagonal and
\[
 C=-V^{-1}\Gamma^{-1}\dot\Gamma V-V^{-1}\dot V,
             \qquad \widetilde f=V^{-1}f.
\]
Both derivative contributions are retained. Norms of the component sources
are uniformly equivalent to the stated projected-source norms.

On the central block, $G_c(t)$, the restriction of $V^*\Gamma V$, has
uniform positive upper and lower bounds and bounded derivative. Its principal
matrix obeys $L_c^*G_c+G_cL_c=0$. The homogeneous equation including
$C_{cc}$ therefore has a propagator bounded by
$C\exp(C_{L_0}|t-u|)$, by differentiating $\langle w_c,G_cw_c\rangle$.
In particular no factor $a^{-3}$ enters its energy growth.

For the stable block, let $G_s(a,q)$ be its positive Lyapunov matrix solving
$L_s^*G_s+G_sL_s=-I$ in the fixed coordinates. It exists by the convergent
integral of $e^{\tau L_s^*}e^{\tau L_s}$, has uniform upper and lower bounds,
and has bounded coefficient derivatives. The latter also follow from the
invertible finite Sylvester equation, whose sums of stable eigenvalues stay
away from zero. Along $q(t)$ its derivative is bounded. Including the
bounded $C_{ss}$ and reducing $a_0$ gives a propagator bound
\[
        \|\Phi_s(t,u)\|\le C e^{-\gamma_1(t-u)/a^3}\quad(t\ge u).
\]
Apply the same argument to $-L_u$ backward in time. The central propagator
and these two exponential bounds are valid for the time-dependent diagonal
blocks, not merely for their instantaneous spectra.

First solve the stable and unstable integral equations with a proposed center
path. Their off-diagonal mutual coupling has supremum-norm bound $Ca^3$,
since each exponential Green kernel has integral at most $Ca^3$. For small
$a$ its Neumann series converges. For the resulting fast path one obtains
\begin{align*}
 \|w_{\mathrm f}\|_\infty
   &\le C(\|s_0\|+\|u_T\|)
                   +Ca^3(\|w_c\|_\infty+\|f_{\mathrm f}\|_\infty),\\
 \|w_{\mathrm f}\|_{L^1}
   &\le Ca^3(\|s_0\|+\|u_T\|)
             +C_{T_0}a^3(\|w_c\|_\infty+\|f_{\mathrm f}\|_\infty).
\end{align*}
The second inequality follows by integrating the boundary exponential and
using Fubini on each Volterra term. It is not obtained by bounding the
boundary contribution by its supremum over an interval of length $T_0$.
Differences between two proposed center paths give the same bounds without
boundary or source terms.

Insert this solution into the center variation-of-constants formula. Its
supremum norm is at most
\[
 C_{T_0,L_0}(\|c_0\|+\|f_c\|_{L^1}+\|w_{\mathrm f}\|_{L^1}).
\]
The dependence of this formula on the proposed center has norm at most
$C_{T_0,L_0}a^3$, so another Neumann series solves it uniquely. Substitution
in the two fast estimates proves \eqref{v34:center-estimate}--
\eqref{v34:fast-L1}. Existence and uniqueness follow from these two
invertible integral systems; equivalence with the differential equation and
its endpoints follows by differentiating their finite-dimensional kernels.
The estimates are uniform also when $T$ shrinks to zero. No value of a
central frequency was divided out at any step.
\end{proof}

\section{Application to the populated histories and the canonical term left over}
\label{v34:application}

\begin{corollary}[Uniform multiplier-block response along V33's actual histories]
\label{v34:populated}
On the endpoint-selected original histories of \cref{v33:main}, the coefficient
path $q(t)=(X(t)/a,u(t)/a)$ obeys $\|q_t\|\le C$ and remains in
$\mathcal U$ for all $0\le t\le c a^2$ and sufficiently small $a$.
Therefore \cref{v34:green} applies on that entire populated interval with
constants independent of $a$. For every actual full variational solution,
its multiplier part is the solution of that block response with its own
endpoint data and the original canonical source
\begin{equation}
 \boxed{f_\xi
 =-a^{-4}\Gamma^{-1}E_a
              (\dot A_X+A_X\mathbb J Q_X)\xi.}
 \label{v34:remaining-source}
\end{equation}
No bound for this source in terms of only $\|\xi\|$ is asserted.
\end{corollary}
\begin{proof}
V33 proves $X_t,u_t=O(a)$ on the entire stated interval, so $q_t=O(1)$.
Its selected initial coefficients are $q(0)=q_*+O(a)$; its total physical
displacements are $O(a^3)$, so $q(t)-q(0)=O(a^2)$. These bounds put the
whole path in the fixed coefficient neighborhood. The derivative bounds of
its analytic $\Gamma$ and $\mathcal K$ are consequently uniform.
The exact multiplier variational equation is
\eqref{v34:balanced-equation}. Substitution of
\eqref{v34:canonical-source} and
$E_a^{-1}M^{-1}=a^{-4}\Gamma^{-1}E_a$ gives
\eqref{v34:remaining-source}. The linear response theorem now applies with
this actual source and the actual endpoint components. It does not bound
those components or the canonical input on its own.
\end{proof}

This corollary is an estimate for a block of the actual linearization along
already supplied nonlinear histories. It is not a nonlinear continuation
beyond $c a^2$, nor a statement that an arbitrary purely multiplier variation
is tangent to the original constrained manifold. In fact, when the restricted
linearized constraints vanish, \eqref{v34:linear-constraints} gives
\[
                    M\eta=-A_X\xi.
\]
Thus a tangent variation with $\xi=0$ also has $\eta=0$. The independent
block boundary problem is a way of measuring its response to canonical
input, not an assertion of independent physical multiplier data.

For comparison, eliminating the multiplier variation on this tangent space
gives exactly
\begin{equation}
 \boxed{\xi_t=\mathbb J\bigl(Q_X-A_X^*M^{-1}A_X\bigr)\xi.}
 \label{v34:reduced-hessian}
\end{equation}
This is the usual constrained-Hessian Schur identity, already present in the
programme's exact multiplier elimination. It is recorded here to identify the
remaining object, not counted as a new abstract reduction theorem. Neither
the positivity of the 85-dimensional multiplier center nor the block Green
bound proves positivity or stability of this different canonical Hessian.
The remaining radial energy-tangency condition is still retained for a fully
constrained physical variation.

The auxiliary comparison $x'=a^2y$, $y'=1+x$ in V33 therefore cannot be
assigned to the multiplier Poisson factor merely from the row orders. A
new growing central loop, if present, must involve the actual canonical
feedback, a transport effect outside the stated coefficient-speed budget,
or a change of coefficient region. Within the stated budget the moving
Hessian alone is included in the controlled block estimate.

\section{Verification, scope, and the next noncircular calculation}
\label{v34:status}

\paragraph{Proved.}
The full original multiplier derivative contains $-M^{-1}\dot M$ in addition
to $-M^{-1}K$. The covector $M\eta$ removes precisely that derivative,
retaining the canonical source and the exact linearized constraint identity.
The full balanced 107-dimensional Poisson factor has eleven stable, eleven
unstable, and 85 positive central dimensions on a uniform coefficient chart.
Its central semigroup is bounded in the stated amplitude-weighted norm.
A moving-coefficient boundary response includes the Hessian and spectral-frame
derivatives, has $a^3$ fast Green gain, and requires no central inverse.
It applies with amplitude-independent constants along the actual V33
histories. The leading-coefficient pencil additionally has exactly three
semisimple zeros and 41 nonzero imaginary pairs at positive small amplitude.

\paragraph{Inherited and actually rerun.}
The V29 root and first Poisson rank, V30 signed enclosure, and V32 rapid
splitting are explicit inputs. The supplied source arrays are not represented
as a newly extracted higher-action coefficient. The portable tests rerun the
inherited exact acceptance and check the new identities symbolically on an
independent analytic Hamiltonian, as well as in noncommuting finite block
examples. A separate floating diagnostic reconstructs the complete leading
pencil from the original root bank, projects the original active representatives
onto the physical active slice, and retains all 107 multiplier coordinates.
It checks the conserved form and spectral pattern at five positive amplitudes.
Those numerical counts are diagnostics only; the all-small-amplitude counts
are proved above from the certified base matrices and continuity.

\paragraph{Remaining requirements.}
No additional nonlinear lifespan exponent is claimed in this body. The next
calculation is the coupled canonical source in \eqref{v34:remaining-source},
or equivalently the constrained canonical Hessian and its actual transport.
The block estimate has not bounded that feedback or supplied its independent
stable/unstable endpoint data. Its positive norm charges an $a^{-1}$ cost to
active multiplier variations. None of these costs may be removed by referring
to the bounded central semigroup. Refinement-uniform preparation and signed
response, a physical invariant or operational selection mechanism, higher
curvature derivatives, normalized metric/link matching, and a common positive
physical lifespan remain unresolved. The raw-action Einstein bridge is not
closed by this linear block result.

All 33 earlier marked mathematical bodies remain byte-for-byte unchanged.
The accompanying records distinguish exact algebra, inherited certificate
acceptance, floating diagnostics, and analytical arguments. No nonlinear
spacetime integration, proof-assistant formalization, or independent audit of
every earlier analytic theorem is claimed.

\begin{thebibliography}{9}
\raggedright
\bibitem{v34:KKS}
T. Kapitula, P. G. Kevrekidis, and B. Sandstede,
\emph{Counting eigenvalues via the Krein signature in infinite-dimensional
Hamiltonian systems}, Physica D \textbf{195} (2004), 263--282,
\href{https://doi.org/10.1016/j.physd.2004.03.018}
{doi:10.1016/j.physd.2004.03.018}.
This gives primary-source context for signed-energy spectral counting. The
finite proof above uses the actual certified matrices and is included in full;
no infinite-dimensional index theorem is applied to the present source.
\end{thebibliography}
% END IMMUTABLE EXACT-ACTION BRIDGE V34 BODY

% Future iterations append here.

% BEGIN IMMUTABLE EXACT-ACTION BRIDGE V35 BODY
\clearpage
\renewcommand{\thesection}{\arabic{section}}
\renewcommand{\theHsection}{bridgeV35.\arabic{section}}
\setcounter{section}{280}
\section{Put the canonical feedback into the spectral equation}
\label{v35:context}

The positive multiplier center of V34 does not estimate the canonical
input to that block. Rather than bound that input independently by its
worst norm, we include it in the original constrained canonical operator.
An exact finite Schur transfer then identifies its contribution at every
nonzero rapid spectral parameter. The original flat subsidiary identity
makes this contribution three amplitude orders smaller than the balanced
multiplier Poisson factor. In physical eigenvalue units the correction to
each of the twenty-two separated hyperbolic rates is bounded, whereas the
rates themselves are of order $a^{-3}$.

This is a statement about the full canonical feedback in a frozen
linearization of the original action, not a substitution of the multiplier
block for that linearization. The canonical matrix has 324 coordinates;
the multiplier complement has 107. We also retain the remaining radial
energy constraint. The fast eigenvectors need not be exactly tangent to its
level set. A separately proved correction supplies actual tangent
variations, with a quantified error, and shows that one of the fast modes
has a nonzero differential in the original initial metric-curvature read.
Neither a frozen spectrum nor this instantaneous sensitivity is promoted
to a longer-time stability or instability theorem.

The inputs are the original flat coefficient identities in V5 and V20,
the simultaneous root and signed chart in V29--V30, the separated simple
rapid spectrum in V31--V32, and V34's exact normalized reduction. These
are inherited results, not newly extracted higher action jets. The
matrix-valued spectral reduction is standard methodology
\cite{v35:GT}; its finite determinant and multiplicity argument are given
below. The new source-specific step is the flat subsidiary cancellation
inside that reduction, and its consequences for the actual canonical
pencil and radial-tangent read. All physical constants remain at $N=3$.

\section{Exact canonical elimination and the flat subsidiary cancellation}
\label{v35:flat}

Use the fixed physical canonical and normalized multiplier coordinates of
\cref{v34:covector}. To avoid confusing the canonical source with a
spectral parameter, write
\[
 C=H_{uX},\qquad Q=H_{XX},\qquad M=H_{uu},\qquad
 K=C\mathbb J C^*,\qquad G=\mathbb J L^*,\qquad L=C(0,0).
\]
Here $u$ is the 107-dimensional mean-lapse-zero multiplier correction;
the complete multiplier is $e+\iota u$. The star is the positive physical
adjoint, extended as the Hermitian adjoint on complex vectors. Transposes
in analytic matrix polynomials mean the complexification of the real
transpose, with real coefficient parameters held fixed.

At each fixed sufficiently small nonzero amplitude, $M$ is invertible.
Solving $H_u(X,u)=0$ locally gives $u=\mathfrak u(X)$. The ordinary
reduced Hamiltonian is
\[
 \mathscr H^{\rm red}(X)=H(X,\mathfrak u(X)).
\]
The inherited envelope and Schur identities give
\begin{equation}
 \nabla\mathscr H^{\rm red}=H_X,\qquad
 S:=D^2\mathscr H^{\rm red}=Q-C^*M^{-1}C,\qquad
 T:=\mathbb J S .                                      \label{v35:T}
\end{equation}
These are identities of the original action. Its radial equation remains
$\mathscr H^{\rm red}=0$; it is not automatically included by eliminating
only the 107 multiplier equations.

Let $F_0=\mathbb J Q_0$ be the original flat quadratic field operator,
where $Q_0=Q(0,0)$. There is a fixed subsidiary row operator $B_0$ with
\begin{equation}
       LF_0=B_0L,\qquad L\mathbb J L^*=0.              \label{v35:flat-subsidiary}
\end{equation}
In unreduced point-row notation $B_0(c_N,c_\beta)=
(\diver_\delta c_\beta,0)$. This preserves the normalized row complement.
The first identity is exactly the linear drift calculation in V5; the
second is the original commuting-derivative identity, not a nonlinear
first-class constraint assertion.

\begin{lemma}[The entire flat canonical-to-constraint transfer vanishes]
\label{v35:flat-transfer}
For every integer $j\ge0$,
\begin{equation}
                LF_0^jG=0.                           \label{v35:all-flat-moments}
\end{equation}
In particular $L\mathbb J Q_0\mathbb J L^*=0$. Whenever the inverse exists,
\begin{equation}
             L(zI-F_0)^{-1}G=0.                       \label{v35:flat-resolvent}
\end{equation}
\end{lemma}
\begin{proof}
Induction in \eqref{v35:flat-subsidiary} gives $LF_0^j=B_0^jL$.
Multiplication by $G$ gives \eqref{v35:all-flat-moments}. For large
$|z|$, the norm-convergent geometric resolvent series proves
\eqref{v35:flat-resolvent}. Each matrix entry is a rational function of
$z$, so its identically zero numerator gives the same equality on the
whole resolvent set. Alternatively the finite Cayley--Hamilton identity
expresses the resolvent as a polynomial in $F_0$ divided by its
characteristic polynomial. No phase product rule is used.
\end{proof}

Write $X=aU$, $u=a\ell$, $q=(U,\ell)$ as in V34. On a small fixed
coefficient neighborhood of $q_*=(Y_*,\ell_*)$,
\begin{equation}
 \begin{gathered}
 C(a,q)=L+O(a),\qquad Q(a,q)=Q_0+O(a),\qquad K=a\mathcal K(a,q),\\
 M=a^4E_a^{-1}\Gamma(a,q)E_a^{-1},\qquad
 E_a=\diag(aI_{78},I_{29}).
 \end{gathered}                                      \label{v35:grading}
\end{equation}
All displayed remainders and fixed coefficient derivatives are uniform on
a smaller compact neighborhood. Their analytic extensions follow from the
original stationary chart. $\Gamma$ is symmetric and uniformly invertible;
$\mathcal K$ is skew. The relations in
\eqref{v35:grading} describe different matrices and may not be interchanged.

\section{A spectral transfer retaining the full canonical feedback}
\label{v35:transfer}

For $\sigma\ne0$ in a fixed compact complex set and small $a$, put
\[
 R_a(\sigma)=(I-a^3\mathbb JQ/\sigma)^{-1},\qquad
 \Theta_a(\sigma)=C R_a(\sigma)\mathbb JQ\mathbb J C^*.
\]
All coefficients below are evaluated at the same $(a,q)$.

\begin{theorem}[Exact nonlinear spectral matrix and an extra amplitude factor]
\label{v35:transfer-theorem}
There is an analytic matrix $\Theta_1(a,q,\sigma)$ on every such compact
spectral set with
\begin{equation}
                 \Theta_a(\sigma)=a\Theta_1(a,q,\sigma),
 \qquad \|\Theta_1\|+\|\partial_\sigma\Theta_1\|\le C. \label{v35:Theta}
\end{equation}
Define the $107$-dimensional matrix-valued function
\begin{equation}
 \boxed{\mathcal F_a(\sigma)=
   \sigma\Gamma+E_a\mathcal K E_a
       +\frac{a^3}{\sigma}E_a\Theta_1(a,q,\sigma)E_a.} \label{v35:F}
\end{equation}
For $z=\sigma/a^3$ outside the bounded free canonical spectrum, $z$ is an
eigenvalue of the original $T$ if and only if $\mathcal F_a(\sigma)$ is
singular. The correspondence preserves algebraic multiplicity. More
precisely,
\begin{equation}
 \boxed{\det(\sigma I_{324}-a^3T)=
 \frac{\det(\sigma I_{324}-a^3\mathbb JQ)}
      {\sigma^{107}\det\Gamma}\det\mathcal F_a(\sigma).} \label{v35:det}
\end{equation}
The geometric factor is not discarded: it is the last term of
\eqref{v35:F}, bounded by $Ca^3$ away from $\sigma=0$.
\end{theorem}
\begin{proof}
The inverse $R_a(\sigma)$ exists uniformly by a geometric series, since
$Q$ is bounded and the spectral set is separated from zero. Analyticity
and \cref{v35:flat-transfer} give
$\Theta_0=L\mathbb JQ_0\mathbb J L^*=0$. Division by $a$ is therefore
analytic: it equals the integral of the $a$ derivative along the segment
from zero to $a$, with $q,\sigma$ fixed. This proves \eqref{v35:Theta},
including its derivative bound. Without the original flat identity this
step would leave a term of order $a^2$, not $a^3$.

Let $R(z)=(zI-\mathbb JQ)^{-1}$. The resolvent identity, in its actual
matrix order, gives
\[
 C R(z)\mathbb J C^*
       =\frac{a^3}{\sigma}K+\frac{a^6}{\sigma^2}\Theta_a(\sigma).
\]
Indeed $R(z)=(a^3/\sigma)R_a(\sigma)$ and
$R_a-I=(a^3/\sigma)R_a\mathbb JQ$.
Apply the finite determinant identity
$\det(I+UV)=\det(I+VU)$ to \eqref{v35:T}:
\[
 \det(zI-T)=\det(zI-\mathbb JQ)
          \frac{\det(M+C R(z)\mathbb J C^*)}{\det M}.
\]
Multiplying the smaller matrix on the left and right by $E_a$, using
\eqref{v35:grading}, and cancelling the determinant factors yields
\eqref{v35:det}. The prefactor there is analytic and nonzero locally at
any nonzero $\sigma$ under consideration. Orders of zeros of determinants
therefore agree. Those orders are precisely algebraic multiplicities
for the ordinary canonical characteristic polynomial and for the
matrix-valued spectral equation, respectively.

For the eigenvectors, if $\mathcal F_a(\sigma)y=0$, set
\begin{equation}
 \eta=E_a y,\qquad
 \xi=\frac{a^3}{\sigma}R_a(\sigma)\mathbb J C^*E_a y.
                                                               \label{v35:eigen-lift}
\end{equation}
Then $C\xi+M\eta=0$ and $T\xi=z\xi$. The vector $\xi$ cannot be zero:
otherwise the first equality and nonsingularity of $M$ would force
$\eta=0$ and then $y=0$. Conversely $T\xi=z\xi$ implies
\eqref{v35:eigen-lift} with $\eta=-M^{-1}C\xi$ and
$y=E_a^{-1}\eta$. This also proves the claimed kernel correspondence.
\end{proof}

This calculation is a frequency-dependent elimination, not a freezing of
$Q$ to $Q_0$. The complete $Q,C,\Gamma,\mathcal K$ remain in
\eqref{v35:F}. The estimate is local in coefficient space and is not
uniform as $\sigma\to0$.

\section{The rapid canonical spectrum and the bounded correction to its rates}
\label{v35:spectrum}

Let $q_a=q_*+O(a)$ be any of the analytic initial families supplied in V30;
the transport-selected initial family is one such choice. The same statements
apply pointwise along V33 histories while their normalized coefficients
remain in the retained neighborhood. Write
\[
 \mathcal L_a=-\Gamma(a,q_a)^{-1}E_a\mathcal K(a,q_a)E_a,\qquad
 \mathscr B_*=-D_*^{-1}K_{WW,*}.
\]
The first is the complete multiplier Poisson factor of V34. It is not the
complete multiplier Jacobian, which also contains its moving metric term.

\begin{theorem}[Full canonical rapid limit and hyperbolic persistence]
\label{v35:canonical-spectrum}
On the actual branch,
\begin{equation}
 \boxed{\det(\sigma I_{324}-a^3T_a)
     \longrightarrow \sigma^{295}\det(\sigma I_{29}-\mathscr B_*)}
                                                               \label{v35:char-limit}
\end{equation}
coefficientwise and hence uniformly on compact complex sets. Its roots
converge as multisets with multiplicity. In particular the original
canonical matrix has eleven simple expanding and eleven simple contracting
rapid eigenvalues associated with the twenty-two separated hyperbolic
roots of V32. For some $c,C>0$,
\begin{equation}
 c a^{-3}\le\max_{z\in\spec T_a}|z|\le C a^{-3}.       \label{v35:radius}
\end{equation}
For each of these separated roots, let $\mu_j(a)$ be the nearby simple
eigenvalue of the \emph{same finite-amplitude} multiplier factor
$\mathcal L_a$, and $z_j(a)$ the canonical eigenvalue. Then
\begin{equation}
                  \boxed{z_j(a)=a^{-3}\mu_j(a)+O(1).} \label{v35:rate-comparison}
\end{equation}
The constants are uniform on a sufficiently small coefficient neighborhood.
No assertion is made that these are the only expanding canonical modes:
additional rates with $a^3z\to0$ are not excluded.
\end{theorem}
\begin{proof}
At $a=0$, $E_0$ vanishes on the 78 active coordinates and is the identity
on the 29 weak coordinates. Block elimination in the invertible active
part of $\Gamma_*$ gives, for $\sigma\ne0$,
\[
 \frac{\det(\sigma\Gamma_*+E_0\mathcal K_*E_0)}{\det\Gamma_*}
       =\sigma^{78}\det(\sigma I_{29}+D_*^{-1}K_{WW,*}).
\]
Combine this with \eqref{v35:det}. The free canonical factor tends to
$\sigma^{324}$, giving \eqref{v35:char-limit} away from zero. To justify
coefficientwise convergence without assuming a bounded canonical matrix,
evaluate these degree-324 monic polynomials at 325 distinct fixed nonzero
complex numbers. The evaluation map has a fixed invertible Vandermonde
matrix. Convergence of all values proves convergence of every coefficient.
The limiting degree is 324, so roots remain bounded and converge, with
multiplicity; this also follows by small disjoint contour circles around
the finite roots and polynomial convergence. It proves the upper bound in
\eqref{v35:radius}. The inherited real root $\gamma_*>0$ proves the lower
bound and persistence of a positive rapid rate.

For the sharper comparison fix one of the simple hyperbolic roots and a
small disk containing no other root of $\mathcal L_0$. The unperturbed
matrix function is
$\sigma\Gamma+E_a\mathcal K E_a=\Gamma(\sigma I-\mathcal L_a)$.
Its determinant has one simple zero $\mu_j(a)$ there, and its derivative
at that zero is bounded away from zero on a smaller coefficient
neighborhood. By \eqref{v35:Theta} the added term in
\eqref{v35:F}, and its $\sigma$ derivative, are $O(a^3)$. On a sufficiently
large fixed multiple of the radius $a^3$ about $\mu_j(a)$ the linear
term of the unperturbed determinant dominates both its quadratic remainder
and this perturbation. The argument principle, or the elementary analytic
implicit-function theorem with uniform derivative, gives a unique zero
$\sigma_j(a)=\mu_j(a)+O(a^3)$. Formula \eqref{v35:det} identifies it with
$a^3z_j(a)$, proving \eqref{v35:rate-comparison}. For the real root,
conjugation and uniqueness make $z_j(a)$ real. Since $S_a$ is real symmetric,
$\mathbb J S_a$ has the usual $z,-z,\bar z,-\bar z$ symmetry: indeed
$(\mathbb J S_a)^T$ is similar to $-\mathbb J S_a$. Thus the contracting
partners are retained as well. The other separated roots are treated on
their disjoint disks.
\end{proof}

The leading real value is the already certified
$0.09479542352<\gamma_*<0.09479542359$. V35 does not recompute that
number. It proves how the actual full canonical eigenvalue inherits it.
The $O(1)$ in \eqref{v35:rate-comparison} compares to the finite-amplitude
Poisson factor. Comparing instead to $a^{-3}\gamma_*$ generally leaves
an $O(a^{-2})$ correction from the variation of $\Gamma,\mathcal K$.
The direct operator-norm estimate is only $\|T_a\|\le Ca^{-4}$.
The spectral-radius bound does not improve that norm estimate or control
nonnormal transients.

\begin{corollary}[No additional unbounded growth rate in a fixed nonzero rapid annulus]
\label{v35:annulus}
Fix $0<\delta<R<\infty$. Every canonical eigenvalue $z$ with
$\delta\le|a^3z|\le R$ satisfies
\[
 \operatorname{dist}(a^3z,\spec\mathcal L_a)\le C_{\delta,R}a^3.
\]
Outside the separated hyperbolic groups already identified, it follows that
$|\operatorname{Re}z|\le C_{\delta,R}$. This does not assert that the
remaining roots are exactly imaginary, nor that their semigroups are bounded.
\end{corollary}
\begin{proof}
V34's invariant projections are uniformly bounded. The central restriction
of $\Gamma$ is uniformly positive, so conjugation by its positive square
root makes the central restriction of $\mathcal L_a$ skew-Hermitian.
It therefore has a uniformly conditioned unitary spectral representation.
The separated hyperbolic eigenvalues are simple, so their eigenvector
matrices also have uniformly bounded condition numbers on a smaller
coefficient neighborhood. Combining the invariant summands gives
\[
 \|(\sigma I-\mathcal L_a)^{-1}\|
       \le C\,\operatorname{dist}(\sigma,\spec\mathcal L_a)^{-1}
\]
outside the spectrum, without dividing by a particular central frequency.
At a zero of \eqref{v35:F}, a unit kernel vector satisfies
$\| (\sigma I-\mathcal L_a)y\|\le C_{\delta,R}a^3$.
The resolvent inequality proves the stated distance bound; if $\sigma$
already belongs to the multiplier spectrum it is immediate. In the
remaining groups the nearby multiplier eigenvalue is purely imaginary
by V34. Divide the bound for its real part by $a^3$ to obtain the result.
\end{proof}

\section{The leading correction is an identified original coefficient}
\label{v35:correction}

The spectral comparison also states what would have to be computed for a
numerical correction. At fixed $q=q_*$, let
\[
 C=L+aC_1+O(a^2),\qquad Q=Q_0+aQ_1+O(a^2).
\]
These are derivatives of the original stationary Hamiltonian, not trial
matrices. The analytic coefficient in \eqref{v35:Theta} has limit
\begin{equation}
 \Theta_{1,*}=
 C_1\mathbb JQ_0\mathbb J L^*
 +L\mathbb JQ_1\mathbb J L^*
 +L\mathbb JQ_0\mathbb J C_1^*.                       \label{v35:theta-first}
\end{equation}
The resolvent itself contributes first at order $a^3$, and hence not in
this derivative. A path $q_a=q_*+O(a)$ has the same limiting expression,
because the $a=0$ matrices $L,Q_0$ do not depend on $q$.

\begin{proposition}[A bounded physical shift of a separated rapid eigenvalue]
\label{v35:first-shift}
Let $y_j,w_j$ be right and left null vectors of
$\mu_{j,0}\Gamma_*+E_0\mathcal K_*E_0$, with
$w_j^*\Gamma_*y_j\ne0$. Then
\begin{equation}
 \boxed{z_j(a)-a^{-3}\mu_j(a)\longrightarrow
 -\frac{w_j^*E_0\Theta_{1,*}E_0y_j}
       {\mu_{j,0}\,w_j^*\Gamma_*y_j}.}                \label{v35:shift}
\end{equation}
Here $\mu_{j,0}$ denotes the limiting eigenvalue. The left-vector pairing
is the Hermitian pairing shown.
\end{proposition}
\begin{proof}
Choose analytic right and left vectors for the simple root of the
unperturbed matrix function. Its derivative in $\sigma$ is $\Gamma$.
Subtract its zero equation from \eqref{v35:F} at the perturbed root,
pair with the left vector, divide by $a^3$, and use
\eqref{v35:rate-comparison}. The eigenvector correction is $O(a^3)$;
its leading contribution is killed by the left null vector, and all other
products with it vanish after this division. The surviving perturbation is
$(\mu_{j,0})^{-1}E_0\Theta_{1,*}E_0$. Division by the nonzero simple-root
pairing proves \eqref{v35:shift}.
\end{proof}

The coefficient \eqref{v35:theta-first} is specified but not numerically
extracted here. In particular its sign is not guessed from the signed
inertia, and it is not replaced by zero. This section gives a formula for
a finite correction to a separated rapid rate, not an estimate of the
central modes near $\sigma=0$.

\section{Retain the radial constraint in the canonical fast directions}
\label{v35:radial}

Work now at the initial points of an actual V30 transport-prepared curve.
Then $P=0$, $\mathscr H^{\rm red}=0$, $X=aY_*+O(a^2)$,
$u=a\ell_*+O(a^2)$ and the actual rates are $O(a)$. Let $\sigma_a>0$
be the canonical scaled eigenvalue converging to $\gamma_*$, and normalize
its kernel vector $y_a$ in \eqref{v35:eigen-lift} so that
\[
 y_a\longrightarrow y_*=
 \binom{-A_*^{-1}L_*v_+}{v_+},\qquad
 \|v_+^\beta\|_h=1.
\]
The symbols $A_*,L_*$ in this formula are the active and mixed blocks of
$\Gamma_*$, not the flat row map $L$. The inherited $v_+$ has a nonzero
nonconstant weak shift. Put $\nu_+=E_0y_*=(0,v_+)$ in physical synthesis.
The eigenvector lift gives
\begin{equation}
 \xi_a=a^3\gamma_*^{-1}G\nu_++O(a^4),\qquad
 \eta_a=E_a y_a=\nu_++O(a).                           \label{v35:leading-vectors}
\end{equation}
In particular $\|\xi_a\|\asymp a^3$; the symmetric phase gradient
$G\nu_+$ is nonzero because a periodic curl-free nonconstant shift cannot
have zero symmetric gradient.

Let $g_a=\nabla\mathscr H^{\rm red}(X_a)$. Its first coefficient is
\[
 g_a=a g_0+O(a^2),\qquad g_0=Q_0Y_*+L^*\ell_*.
\]
There is an explicit original canonical normal direction
\begin{equation}
 W=(0,\omega_h\diag(0,-1,-1)),\qquad
 LW=0,\qquad \langle g_0,W\rangle=32\omega_h^2=864.
                                                               \label{v35:W}
\end{equation}
Indeed $W$ differentiates the homogeneous boost $d_1$ with the profiles
and $d_2=d_3=-8$ held fixed. The homogeneous momentum part of $H_2$ is
$-2\omega_h^2(d_1d_2+d_1d_3+d_2d_3)$; its derivative is the displayed
number. All mixed profiles have zero spatial mean, so their momentum
cross term with $W$ integrates to zero. This use of $W$ is a transverse
energy test, not a change of the prepared first jet.

\begin{theorem}[Exact radial-tangent representatives of the rapid mode]
\label{v35:radial-theorem}
Set
\begin{equation}
 \widetilde\xi_a=\xi_a-
       \frac{\langle g_a,\xi_a\rangle}{\langle g_a,W\rangle}W,
 \qquad
 \widetilde\eta_a=-M^{-1}C\widetilde\xi_a.              \label{v35:radial-project}
\end{equation}
These variations satisfy \emph{every} original linearized constraint,
including the radial one. They are derivatives of local curves of exactly
constrained original initial data. Quantitatively,
\[
 \widetilde\xi_a-\xi_a=O(a^4),\qquad
 E_a^{-1}(\widetilde\eta_a-\eta_a)=O(a).
\]
\begin{equation}
 \boxed{\|(a^3T_a-\sigma_a I)(a^{-3}\widetilde\xi_a)\|\le Ca.}
 \label{v35:radial-bounds}
\end{equation}
No exact frozen eigenvector on the energy tangent space is asserted.
\end{theorem}
\begin{proof}
The subsidiary identity and $LY_*=0$ give
\[
 \langle g_0,G\nu\rangle
       =-\langle LF_0Y_*,\nu\rangle
           +\langle\ell_*,L\mathbb JL^*\nu\rangle=0
\]
for every normalized multiplier test. Thus \eqref{v35:leading-vectors}
implies $\langle g_a,\xi_a\rangle=O(a^5)$. By \eqref{v35:W},
$\langle g_a,W\rangle=864a+O(a^2)$, so the coefficient of $W$ in
\eqref{v35:radial-project} is $O(a^4)$. Since $CW=O(a)$ and
$E_a^{-1}M^{-1}=a^{-4}\Gamma^{-1}E_a$, its multiplier correction in
balanced norm is $O(a)$. Also
$\|T_aW\|\le C a^{-3}$, from \eqref{v35:T}, bounded $C,Q$,
$CW=O(a)$ and $\|M^{-1}\|\le Ca^{-4}$. Therefore
\[
 (T_a-\sigma_a/a^3)\widetilde\xi_a
 =-\frac{\langle g_a,\xi_a\rangle}{\langle g_a,W\rangle}
                   (T_a-\sigma_a/a^3)W=O(a),
\]
which is exactly \eqref{v35:radial-bounds}.

The restricted variation is zero by construction of
$\widetilde\eta_a$. The radial variation is
$\langle g_a,\widetilde\xi_a\rangle=0$. Since $M$ is invertible and
$g_a\ne0$, the normalized constraint graph has the smooth energy
hypersurface $\mathscr H^{\rm red}=0$. Every vector tangent to this
hypersurface is the velocity of a local smooth curve in it, by a local
implicit-function chart using $W$ as the transverse coordinate. Lifting
by $\mathfrak u(X)$ proves the stated realization in exactly constrained
initial records, with positive lapse and the original rank chart retained
for a sufficiently small curve parameter at each fixed $a$.
\end{proof}

\section{The full initial geometric differential sees the rapid canonical direction}
\label{v35:geometry}

For each fixed nonzero $a$, vary initial data along the just constructed
constraint curve and take its own local normalized exact-action histories.
Let $\delta$ mean differentiation in that curve parameter at zero. All
observations in the next statement are evaluated at its initial cut.

\begin{theorem}[A populated curvature differential with an explicit rapid scale]
\label{v35:geometric-differential}
For the tangent direction of \cref{v35:radial-theorem},
\begin{equation}
 \boxed{a^3\delta\lambda_t\longrightarrow\gamma_*\nu_+,\qquad
 a^3\delta(R_{0i0j}(g_h))_{ij}
           \longrightarrow\gamma_*\mathfrak E_h(v_+^\beta).}
                                                               \label{v35:geometric-limit}
\end{equation}
The second limit is nonzero: the inherited original-read certificate has
$\|\mathfrak E_h(v_+^\beta)\|_h>7/10$. Consequently
\begin{equation}
       \|\delta(R_{0i0j}(g_h))_{ij}\|_h\ge c a^{-3}    \label{v35:sensitivity}
\end{equation}
for sufficiently small $a$. These are differentials of actual initial
observations, not curvature divergence of the reference histories.
\end{theorem}
\begin{proof}
First use the uncorrected eigenvector. V34's exact derivative identities
and $\eta_a=-M^{-1}C\xi_a$, $T_a\xi_a=z_a\xi_a$ give
\[
 \delta r=z_a\eta_a-M^{-1}(\dot C\xi_a+\dot M\eta_a),
                  \qquad z_a=\sigma_a/a^3.
\]
Along the prepared histories $\dot C=O(a)$, because the original fixed-chart
partial derivatives are bounded and both $X_t,u_t$ are $O(a)$. Also
$M=D_a\Gamma D_a$, with $a$ held constant along the history and
$\dot\Gamma=O(1)$. Since $\xi_a=O(a^3)$ and $\eta_a=E_a y_a$,
\[
 \|M^{-1}\dot C\xi_a\|=O(1),\qquad
 M^{-1}\dot M E_a y_a=E_a\Gamma^{-1}\dot\Gamma y_a=O(1).
\]
Hence $a^3\delta r\to\gamma_*\nu_+$.

We check rather than assume that the radial correction preserves this
limit. It changes the canonical vector by $O(a^4)W$.
Since $CW=O(a)$, \eqref{v35:flat-subsidiary} also gives
$(\dot C+C\mathbb JQ)W=O(a)$, using $LF_0W=B_0LW=0$.
V34 therefore gives $D_Xr[W]=O(a^{-3})$.
The multiplier correction is $E_a O(a)$; the exact balanced derivative
has $\|D_ur\,E_a\|\le Ca^{-3}$. The changes of $\delta r$ are thus
$O(a)$ and $O(a^{-2})$, respectively. Multiplication by $a^3$ sends
both to zero. This proves the first limit for the actual radial-tangent
variation.

The original canonical equation gives $\delta X_t=O(1)$ and
\[
 \delta X_{tt}=G\,\delta r+O(a\|\delta r\|)+O(1).
\]
This follows by differentiating
$X_{tt}=D_XF\,F+D_uF\,r$; all partial derivatives of the stationary
canonical field $F$ are bounded, the reference $F,r$ are $O(a)$, and
$D_uF=G+O(a)$. In particular
$a^3\delta\gamma_{tt}$ tends to the symmetric \emph{phase} gradient of
$\gamma_*v_+^\beta$; the lapse part of $G$ occupies only the momentum
slot. The other leading contribution to the metric curvature is the
ordinary spatial gradient of $\delta\beta_t$. At this fixed cutoff all
spatial derivatives are bounded maps. Variations of the undifferentiated
metric coefficients, lapse, shift, and first time jets contribute at most
$O(1)$ or $O(a\|\delta r\|)$ besides those two terms. The reference
connection products contribute the same smaller orders. In the exact
coordinate formula
\[
 R_{0i0j}=\tfrac12(g_{0j,it}+g_{i0,tj}-g_{00,ij}-g_{ij,tt})
                     +\hbox{Christoffel products},
\]
no second lapse or shift time derivative occurs. Multiplication by $a^3$
therefore leaves exactly
$\gamma_*\sym((D-\delta)\otimes v_+^\beta)
=\gamma_*\mathfrak E_h(v_+^\beta)$, with the V10 convention for the
symmetric gradient. This proves \eqref{v35:geometric-limit}. Its nonzero
norm and convergence imply \eqref{v35:sensitivity}. The original V31
normalization supplies the cited visibility bound, not a new evaluation
of a nonlinear metric trajectory.
\end{proof}

The canonical component of this direction is $O(a^3)$, while its multiplier
component has a nonzero $O(1)$ limit. Thus the theorem is not a small
perturbation of the complete initial record in the same norm as the
canonical coordinates alone. For each $a$ it is an infinitesimal tangent,
and its parameter may be reduced to stay in the original chart. The
reference selected histories still obey their $O(a)$ curvature bounds.
No finite-time amplification factor, attracting basin, or divergence of a
coordinate-invariant scalar is inferred solely from this differential.

\section{Direction, source audit, and verification}
\label{v35:status}

\paragraph{Proved.}
The full canonical spectral equation has an exact 107-dimensional transfer
with all canonical feedback retained. The original flat subsidiary identity
removes its order-$a^2$ correction, leaving order $a^3$ at fixed nonzero
rapid spectral parameter. The entire scaled canonical characteristic
polynomial has the limit \eqref{v35:char-limit}; every separated
hyperbolic canonical rate differs by only $O(1)$ from the corresponding
same-amplitude multiplier Poisson rate. A real growing direction can be
corrected to satisfy the radial energy constraint, and it gives a nonzero
$a^{-3}$ differential of the original initial metric-curvature observation.

\paragraph{Inherited inputs and executed checks.}
V29's root, V30's signed regularity, and the V31--V34 spectral acceptance
are inherited. Their numerical enclosures are not described as newly
extracted canonical Hessian coefficients. The new verifier reconstructs the
entire original flat $324$-coordinate canonical operator and 108 original
linear rows at $N=3$ in exact rational arithmetic, and checks the subsidiary,
commuting-bracket, adjoint, and transfer-zero identities. Independent small
rational matrix families check the full determinant reduction, eigenvector
lifting, the extra amplitude factor, and the radial normalization algebra.
Those families are explicitly labeled algebra tests, not alternative
physical sources. The conditional-to-source-specific passage in the theorems
uses the written analytic argument and the named original inputs.

\paragraph{What remains.}
The bounds for \eqref{v35:F} are separated from $\sigma=0$. The limit
therefore does not classify central rates with $a^3z\to0$, bound the
nonautonomous coupled propagator, or estimate the unresolved canonical
source of V34 over longer intervals. Frozen eigenvectors are not solutions
of a varying-coefficient linear equation; the radial-tangent correction is
stated as a quasimode, not as an invariant eigenspace of the energy
hypersurface. No new nonlinear lifespan exponent is claimed. The next
productive calculation is the small-$\sigma$ canonical transfer, or a
source-resolved bound on its time-dependent action on the selected
histories, not another high-frequency positivity argument. Refinement,
higher observation derivatives, normalized metric/link matching and a
common positive physical lifespan remain open requirements. The exact-action
Einstein bridge is not closed by the present spectral and differential
results.

All 34 prior marked bodies are preserved byte-for-byte. The package records
which inherited checks were rerun, exact new finite checks, and the analytic
scope of each theorem. No numerical nonlinear spacetime integration,
proof-assistant formalization, or independent reproof of every inherited
analytic theorem is claimed.

\begin{thebibliography}{9}
\raggedright
\bibitem{v35:GT}
S. G\"uttel and F. Tisseur,
\emph{The nonlinear eigenvalue problem}, Acta Numerica \textbf{26}
(2017), 1--94, \href{https://doi.org/10.1017/S0962492917000034}
{doi:10.1017/S0962492917000034}.
The publisher's article record was checked for the established matrix-valued
spectral viewpoint. No numerical algorithm or external convergence theorem
from this survey replaces the finite determinant proof or original-source
cancellation above.
\end{thebibliography}
% END IMMUTABLE EXACT-ACTION BRIDGE V35 BODY

% Future iterations append here.

% BEGIN IMMUTABLE EXACT-ACTION BRIDGE V36 BODY
\clearpage
\renewcommand{\thesection}{\arabic{section}}
\renewcommand{\theHsection}{bridgeV36.\arabic{section}}
\setcounter{section}{288}
\section{The smaller spectral scale contains an additional real pair}
\label{v36:context}

V35 separates the canonical rates of order $a^{-3}$ but deliberately leaves
$a^3z\to0$ unresolved. The original constant-shift rows determine a different
scale inside that region. Their linear field gradients vanish at the flat
point, their first Poisson rows vanish at the compatible first field, and
their mutual Poisson bracket vanishes through quadratic field order. Keeping
these cancellations, rather than dividing V35's estimate by a small rapid
parameter, gives a finite limiting pencil at $z=\zeta/a$.

The new pencil has a three-dimensional negative stiffness term. Its
coefficient is the positive quadratic canonical energy of the constant phase
translations of the actual first field. The original signed form is positive
on the active-plus-constant-shift carrier, and the first Poisson bracket is
invertible on the remaining nonconstant multiplier carrier. These two facts
force a change of determinant sign on the positive real $\zeta$ axis. Hence
the full canonical matrix has an additional real expanding/contracting pair
of order $a^{-1}$. This is separate from the eleven rapid pairs of V35.
A second theorem bounds the real parts of \emph{all} the remaining canonical
rates by $C/a$, making this slower growth order sharp.

\paragraph{Inputs and scope.}
All statements concern the same $N=3$ root $Y_*$ of V29, its signed chart in
V30, and exactly constrained initial families with
$X_a=aY_*+O(a^2)$ and $\lambda_a=e+a\ell_*+O(a^2)$.
The full canonical matrix is precisely \eqref{v35:T}; the radial energy
constraint has not been deleted from the physical problem. V35's rapid
counts and V34's positive multiplier center are inherited inputs.
The constant-shift coefficient $P_{c,2}$ and the complete cubic envelope
$P_{c,3}$ are retained from V5--V6, and the prepared cubic-mean map is
\eqref{v21:R}. The new work is their incidence in the small-parameter
canonical transfer, its real spectral consequence, and the bound on every
other growth rate. The polynomial-pencil viewpoint is standard
\cite{v36:TM}; all finite reductions and sign arguments needed here are proved.
No higher-action Hessian is numerically substituted, and no new physical
lifespan is asserted.

\section{Constant translations fix the missing orders}
\label{v36:translations}

Use physical orthonormal multiplier coordinates in the fixed mean-lapse-zero
complement. Decompose that space as
\[
 \mathcal E=\mathcal A\oplus\mathcal W_g\oplus\mathcal H,
 \qquad (\dim\mathcal A,\dim\mathcal W_g,\dim\mathcal H)=(78,26,3),
 \qquad \mathcal R=\mathcal A\oplus\mathcal W_g.
\]
Here $\mathcal H$ consists of the three spatially constant shifts, not the
radial lapse. Denote the orthogonal projections by $P_A,P_g,P_H,P_R$ and
the constant-shift inclusion by $I_H:\R^3\to\mathcal E$.
Keep $C,Q,M,K,F=\mathbb JQ,L,F_0$ from V35. Thus $K=C\mathbb JC^*$,
$E_a=\diag(aI_A,I_g,I_H)$ and
$M=a^4E_a^{-1}\Gamma E_a^{-1}$.

For a constant $c$, write $\delta_c=c^i\delta_i$ and
$P_c(X,\lambda)=\langle I_Hc,P(X,\lambda)\rangle_h$.
The original quadratic identity, on arbitrary canonical fields, is
\[
 P_{c,2}(u,p)=\langle p,\delta_cu\rangle_h,
 \qquad \mathbb J\nabla P_{c,2}(Y)=\delta_cY.
\]
Define the actual translation-energy matrix
\begin{equation}
 \mathcal V c=\delta_cY_*,\qquad
 \mathsf T_{cd}=\langle\delta_cY_*,Q_0\delta_dY_*\rangle_h,
 \qquad \mathsf T=\mathcal V^*Q_0\mathcal V.
 \label{v36:Ttranslation}
\end{equation}
The symbol $\mathsf T$ in this body is a $3$-by-$3$ energy matrix, not the
$324$-dimensional canonical generator $T_a$.

\begin{lemma}[Prepared Poisson orders involving all three translation rows]
\label{v36:Poisson-orders}
Put $K_1=\mathbb K_1(Y_*)$. Its kernel in $\mathcal E$ is exactly
$\mathcal H$, so $K_{1,R}=P_RK_1P_R|_{\mathcal R}$ is invertible. On an
exactly prepared family as above,
\begin{align}
 K_{RR}&=aK_{1,R}+O(a^2),\\
 K_{HR}&=a^2\mathsf R+O(a^3),\\
 K_{HH}&=a^3\mathsf S+O(a^4),                         \label{v36:Korders}
\end{align}
where the coefficients are specified by the original scalar envelopes:
\begin{align}
 \mathsf R[c,\nu]
 &=DP_{c,3}(Y_*)[G\nu]-DP_{\nu,2}(Y_*)[\delta_cY_*]
       -\langle\delta_c\nu,P_2(Y_*)\rangle_h,
          &&\nu\in\mathcal R,\label{v36:R}\\
 \mathsf S[c,d]
 &=DP_{c,3}(Y_*)[\delta_dY_*]
                -DP_{d,3}(Y_*)[\delta_cY_*].        \label{v36:S}
\end{align}
In particular $\mathsf S^*=-\mathsf S$. These three leading blocks do not
depend on the second or higher preparation coefficients.
\end{lemma}
\begin{proof}
The maximal first Poisson rank at the actual V29 root is $104$ on this
$107$-dimensional complement. The three constant shifts are null by the
original compatible-first-field identity, so they exhaust the kernel.
A real skew matrix maps the orthogonal complement of its kernel into itself
and is invertible there. This proves the first assertion.

The universal weak first-source identity gives $J_1(Y)c=0$ for every
symmetric $Y$. The leading multiplier Hessian is $J_1^*K_0^{-1}J_1$.
Thus $D_\lambda P_{c,2}(Y,e)=0$, with its field derivatives, for every
constant $c$. Consequently, through second order in its field gradient,
\[
 D_XP_c(X_a,\lambda_a)
 =aDP_{c,2}(Y_*)
   +a^2\{DP_{c,2}(Y_2)+DP_{c,3}(Y_*)\}+O(a^3).
\]
No first-tilt contribution occurs at these orders. The other rows have
$D_XP_\nu=L_\nu+aDP_{\nu,2}(Y_*)+O(a^2)$.
Exact initial preparation supplies $LY_2=-P_2(Y_*)$; the
order-amplitude multiplier tilt changes this equation only at the next
order, as proved in V20.

The order-$a$ constant row is
$DP_{c,2}(Y_*)[G\nu]=\langle\delta_c\nu,LY_*\rangle_h=0$.
Its order-$a^2$ coefficient is
\[
 DP_{c,2}(Y_2)[G\nu]+DP_{c,3}(Y_*)[G\nu]
                 -DP_{\nu,2}(Y_*)[\delta_cY_*].
\]
Substitution of the second constraint equation gives \eqref{v36:R},
including the preparation term. This is the same complete coefficient as
V21's cubic-mean tilt map, now used as a mixed Poisson coefficient.

The quadratic translation generators commute exactly:
$\{P_{c,2},P_{d,2}\}=0$, because the constant-coefficient skew phase
translations commute on both canonical slots. This identity holds on all
fields, so its polarization eliminates every $Y_2$ contribution to the
next mutual bracket. The first remaining term is exactly
\eqref{v36:S}. Higher field or multiplier terms enter only at order $a^4$.
This proves all displayed expansions with the original envelopes retained.
\end{proof}

\section{A uniform transfer before the rapid parameter is kept away from zero}
\label{v36:uniform-transfer}

Choose $z_0$ larger than twice a bound for $\|F\|$ on the retained chart.
For every complex $|z|\ge z_0$, set
\begin{equation}
 R_z=(I-F/z)^{-1},\qquad
 H_a(z)=C R_z F\mathbb JC^*.
 \label{v36:H}
\end{equation}
This uses the physical spectral variable $z$, not a new time coordinate.

\begin{lemma}[Uniform flat cancellation and a stronger translation column bound]
\label{v36:uniform-H}
Uniformly for $|z|\ge z_0$ on the specified initial families,
\begin{equation}
 H_a(z)=a\Theta_a(z),\qquad
 \|\Theta_a(z)\|\le C,\qquad
 \|P_H\Theta_a(z)\|+\|\Theta_a(z)P_H\|\le Ca.
 \label{v36:Hbounds}
\end{equation}
For $z=\zeta/a$ on a compact set separated from $\zeta=0$,
\begin{equation}
 a^{-2}H_{HH}(\zeta/a)=-\mathsf T+O(a).
 \label{v36:HHlimit}
\end{equation}
The limits and fixed $\zeta$ derivatives are uniform on such compact sets.
\end{lemma}
\begin{proof}
The Neumann series gives uniform resolvents and coefficient derivatives on
$|z|\ge z_0$. At the flat point,
$L(I-F_0/z)^{-1}F_0G=0$ by the entire flat-transfer identity of V35.
Taylor expansion in amplitude therefore gives the first two assertions.

In a constant-shift column $C_H=aC_{1,H}+O(a^2)$, with
$\mathbb JC_{1,H}^*c=\delta_cY_*$. The possible first-order term of
$H_{RH}$ is
\[
 L_R(I-F_0/z)^{-1}F_0\delta_cY_*.
\]
It is zero: $F_0$ commutes with constant phase translations and
$LF_0^j\delta_cY_*=B_0^j\delta_cLY_*=0$ for every $j$.
For the other mixed block use $F_0G=-GB_0^*$, which follows by taking
the adjoint of $LF_0=B_0L$ and using the Hamiltonian identity
$F_0\mathbb J=-\mathbb JF_0^*$. Also
$C_{1,H}G=0$ by $LY_*=0$. Hence
$C_{1,H}(I-F_0/z)^{-1}F_0G_R=0$.
Both mixed blocks are therefore $O(a^2)$, uniformly in the stated region;
the $HH$ block already has two factors $C_H=O(a)$.

At $z=\zeta/a$, $R_z=I+O(a)$. Thus its $HH$ coefficient is
$C_{1,H}\mathbb JQ_0\mathbb JC_{1,H}^*=-\mathcal V^*Q_0\mathcal V$.
This proves \eqref{v36:HHlimit}, including its sign. Uniform analytic
Taylor estimates give the differentiated versions. No statement here
crosses a bounded free canonical eigenvalue or divides by $z=0$.
\end{proof}

For later use the exact spectral equation of V35 can now be written, for
$\sigma=a^3z$ and $|z|\ge z_0$, as
\begin{equation}
 \left\{\sigma\Gamma+E_a\mathcal K E_a
            +z^{-1}E_a\Theta_a(z)E_a\right\}y=0,
 \qquad \mathcal K=K/a.
 \label{v36:exact-balanced}
\end{equation}
Its kernel is in bijection with the canonical eigenvectors, since the free
canonical resolvent exists in this region. This is the same determinant
reduction as V35 with a stronger range of estimates, not a new physical
operator.

\section{Every nonrapid exponential rate is at most of order inverse amplitude}
\label{v36:ceiling}

\begin{theorem}[A sharp upper scale for every additional real part]
\label{v36:growth-ceiling}
Remove the twenty-two separated rapid hyperbolic eigenvalues of V35, counted
with their algebraic multiplicities. Every one of the other $302$ canonical
eigenvalues obeys
\begin{equation}
                  \boxed{|\operatorname{Re}z|\le C a^{-1}.}
 \label{v36:ceiling-eq}
\end{equation}
The constant is independent of small positive amplitude on the stated
bounded preparation family. This does not bound their imaginary parts,
their Jordan structure, or the nonautonomous propagator.
\end{theorem}
\begin{proof}
V35's characteristic-polynomial convergence places the remaining scaled
roots $\sigma=a^3z$ in a fixed strip separated from the hyperbolic spectra
of the full multiplier factor
$\mathcal L_a=-\Gamma^{-1}E_a\mathcal K E_a$. Its hyperbolic resolvents
are uniformly bounded there. Let $y$ be a unit kernel vector in
\eqref{v36:exact-balanced}, with $|z|$ sufficiently large. Projection onto
the hyperbolic groups and \eqref{v36:Hbounds} give
$\|y_{\rm hyp}\|\le C/|z|$.

The hyperbolic and central spectral spaces are $\Gamma$-orthogonal.
The central restriction is uniformly positive by V34. Increasing $z_0$ if
necessary therefore gives $\langle y,\Gamma y\rangle\ge c>0$.
Set $\nu=E_ay$. Since $E_a\mathcal K E_a$ is skew-Hermitian,
taking real parts of the exact equation yields
\begin{equation}
 |\operatorname{Re}\sigma|
             \le \frac C{|z|}\,|\nu^*\Theta_a(z)\nu|.
 \label{v36:energy-eigen}
\end{equation}
In particular the right side is at most $C/|z|$.

If $|\sigma|\ge a$, then $|z|\ge a^{-2}$, and this last bound gives
$|\operatorname{Re}z|\le C/(a^3|z|)\le C/a$.
If $|z|\le a^{-1}$, the conclusion is immediate from
$|\operatorname{Re}z|\le|z|$. It remains to treat
$a^{-1}\le|z|\le a^{-2}$.

Multiply \eqref{v36:exact-balanced} by $E_a^{-1}$ and use
$\mathcal K=K_1+O(a)$ and $\ker K_1=\mathcal H$. The inverse of
$K_{1,R}$ gives
\begin{equation}
 \|P_R\nu\|\le C\left(\frac{|\sigma|}{a}+a+|z|^{-1}\right)
              \le C\left(\frac{|\sigma|}{a}+a\right).
 \label{v36:Rnu}
\end{equation}
Here $\|\nu\|\le1$ and $\|E_a^{-1}\|=a^{-1}$; no inverse of a
small finite-amplitude Poisson eigenvalue is used.
The translation row and column bounds in \eqref{v36:Hbounds} now imply
\[
 |\nu^*\Theta_a(z)\nu|
 \le C\{\|P_R\nu\|^2+a\}
 \le C\left(\frac{|\sigma|^2}{a^2}+a\right).
\]
Insert this in \eqref{v36:energy-eigen} and divide by $a^3$:
\[
 |\operatorname{Re}z|
       \le C\left(a|z|+\frac1{a^2|z|}\right)\le C/a
       \qquad(a^{-1}\le|z|\le a^{-2}).
\]
These three regions cover every root; bounded $z$ are absorbed by the same
constant. Spectral projections were used only in the frozen eigenvalue
identity. The argument is not a positive energy estimate for the coupled
physical dynamics.
\end{proof}

\section{A finite pencil at the inverse-amplitude scale}
\label{v36:slow-pencil}

Define another, explicitly singular analytical scaling
\[
 B_a=\diag(I_R,aI_H),\qquad
 F_a=aE_a^{-1}B_a^{-1}=P_A+aP_g+P_H,
 \qquad F_*=P_A+P_H.
\]
This $B_a$ is a coordinate scaling, not a drift or a Poisson matrix.
At $z=\zeta/a$, put
\begin{equation}
 \mathcal P_a(\zeta)=B_a^{-1}
 \left\{\frac Ka+a^2\zeta E_a^{-1}\Gamma E_a^{-1}
                       +\frac1\zeta H_a(\zeta/a)\right\}B_a^{-1}.
 \label{v36:Pa}
\end{equation}
Multiplying the physical resolvent constraint
$M+C(zI-F)^{-1}\mathbb JC^*$ by $z/a$, and then applying this congruence,
gives exactly \eqref{v36:Pa}. Thus its determinant is zero exactly at the
original canonical eigenvalues in this spectral region.

\begin{theorem}[A preparation-independent leading inverse-amplitude pencil]
\label{v36:pencil-limit}
Uniformly on compact $\zeta$ sets separated from zero,
\begin{equation}
 \boxed{\mathcal P_a(\zeta)
    =\zeta\mathcal G+\mathcal J-\zeta^{-1}\mathcal H+O(a),}
 \label{v36:pencil}
\end{equation}
where
\begin{equation}
 \mathcal G=F_*\Gamma_*F_*,\qquad
 \mathcal J=\begin{pmatrix}K_{1,R}&-\mathsf R^*\\
                           \mathsf R&\mathsf S\end{pmatrix},\qquad
 \mathcal H=I_H\mathsf T I_H^*.
 \label{v36:pencil-coefficients}
\end{equation}
In particular $\mathcal G^*=\mathcal G$, $\mathcal J^*=-\mathcal J$,
and $\mathcal H^*=\mathcal H$. The associated quadratic polynomial is
\begin{equation}
             \mathcal Q_0(\zeta)=\zeta^2\mathcal G+\zeta\mathcal J-\mathcal H.
 \label{v36:quadratic-pencil}
\end{equation}
Its three coefficients depend on the original first field $Y_*$, its known
leading signed source, and quadratic/cubic scalar envelopes only. Neither
the unlisted second initial field nor a fourth-order action coefficient is
needed for this limiting pencil.
\end{theorem}
\begin{proof}
For the mass term the identity $aE_a^{-1}B_a^{-1}=F_a$ gives exactly
$\zeta F_a\Gamma F_a\to\zeta\mathcal G$. The Poisson block limits are
\eqref{v36:Korders}, with $K_{RH}=-K_{HR}^*$.
By \cref{v36:uniform-H}, the feedback has $RR$ block $O(a)$, $RH$ and
$HR$ blocks $O(a^2)$, and $HH$ block $-a^2\mathsf T+O(a^3)$ at
$z=\zeta/a$. The two factors $B_a^{-1}$ therefore leave only
$-\mathsf T/\zeta$ in its limit. All mixed feedback blocks tend to zero.
This proves \eqref{v36:pencil} with its stated uniform error. The coefficient
formulas and the preparation independence proved in
\cref{v36:Poisson-orders} establish the final assertion.

Multiplication by $\zeta$ in \eqref{v36:quadratic-pencil} does not give a
claim about eigenvalues at zero. The convergence above is used only on
compact nonzero sets; possible zero or infinite eigenvalues of a singular
mass pencil are not silently counted as canonical physical modes.
\end{proof}

\section{Both decisive signs follow from the actual source}
\label{v36:signs}

\begin{lemma}[Positive retained mass and strictly positive translation energy]
\label{v36:positive-data}
The matrix $\mathcal G$ is positive semidefinite, with
\begin{equation}
 \ker\mathcal G=\mathcal W_g,\qquad
 \mathcal G\big|_{\mathcal A\oplus\mathcal H}\succ0,
 \qquad \rank\mathcal G=81.
 \label{v36:mass-positive}
\end{equation}
For the actual root's fixed axial amplitudes,
\begin{equation}
 \boxed{\mathsf T\succeq
       729\diag(1,81/64,1)\succeq729I_3.}
 \label{v36:translation-positive}
\end{equation}
In particular $\mathcal H$ is positive semidefinite of rank three.
\end{lemma}
\begin{proof}
At zero amplitude the complete multiplier factor has only the weak Poisson
block. Every vector in $\mathcal A\oplus\mathcal H$ is in its kernel:
$E_0$ kills the active coordinates and the constant shifts are Poisson-null.
This kernel belongs to the positive central space established in V34.
The restriction of $\Gamma_*$ to this $81$-dimensional subspace is therefore
positive definite. Compression by $F_*$ proves \eqref{v36:mass-positive}.
This uses the certified signed form, not positivity on the entire weak carrier.

Constant phase differentiation removes the homogeneous momentum. All remaining
configuration and momentum profiles in V29 are transverse and trace-free.
On such nonconstant records the original quadratic Hamiltonian has Hessian
pairing
\[
 \langle(\dot u,\dot p),Q_0(\dot u,\dot p)\rangle_h
     =2\|\dot p\|_h^2+\tfrac12\sum_j\|\delta_j\dot u\|_h^2.
\]
Different axial and mixed wavevectors are orthogonal on the three-site grid,
including the folded representatives. Hence the mixed profiles contribute a
positive semidefinite addition to the translation-energy matrix.
For the retained axial pair at axis $i$,
$U_i=\alpha_iT_i\cos(2\pi x_i)$ and
$p_i=(\omega_h/2)\alpha_iP_i\cos(2\pi x_i)$, with both tensor squared
norms equal to two. Its contribution is
$\omega_h^4\alpha_i^2c_i^2$.
At $N=3$, $\omega_h=3\sqrt3$ and
$\alpha=(1,9/8,1)$, giving \eqref{v36:translation-positive}.
These are the first canonical profiles of the verified root, not a new
homogeneous approximation to it.
\end{proof}

\section{An unavoidable real pair and an exact three-row reduction}
\label{v36:real-pair}

Let $\mathcal P_0(\zeta)=\zeta\mathcal G+\mathcal J-\zeta^{-1}\mathcal H$.
For real $\zeta\ge0$, its $RR$ block is
\[
 B_R(\zeta)=K_{1,R}+\zeta\mathcal G_{RR}.
\]

\begin{lemma}[The nonconstant block never has a positive-real pole]
\label{v36:Rinvertible}
For every $\zeta\ge0$, $B_R(\zeta)$ is invertible with positive determinant.
For $\zeta>0$, the limiting spectral equation is exactly the three-row
condition
\begin{equation}
 \det\mathcal Z_H(\zeta)=0,\qquad
 \mathcal Z_H(\zeta)=\zeta\mathcal G_{HH}+\mathsf S-\frac{\mathsf T}{\zeta}
  -(\mathsf R+\zeta\mathcal G_{HR})B_R(\zeta)^{-1}
                (-\mathsf R^*+\zeta\mathcal G_{RH}).
 \label{v36:three-rows}
\end{equation}
The inverse and all couplings in this formula are retained.
\end{lemma}
\begin{proof}
At zero, $K_{1,R}$ is an invertible skew matrix of even dimension, so its
determinant is the positive product of its nonzero singular-value pairs.
For $\zeta>0$, if $B_Rx=0$, its real energy pairing gives
$\langle x,\mathcal G_{RR}x\rangle=0$.
By \eqref{v36:mass-positive}, $x$ then lies in $\mathcal W_g$, and
$K_{1,R}x=0$ forces $x=0$. Thus no positive-real singularity occurs, and
continuity preserves the positive determinant. The exact Schur determinant
identity gives \eqref{v36:three-rows}; no condition at $\zeta=0$ is inferred
from its displayed reciprocal factor.
\end{proof}

\begin{theorem}[Additional canonical expansion of order $a^{-1}$]
\label{v36:slow-real}
There exist fixed $0<c<C<\infty$ such that, for every sufficiently small
positive amplitude in the stated original prepared family, the full canonical
matrix has a real eigenvalue $z_a$ and its negative with
\begin{equation}
                         \boxed{c/a\le z_a\le C/a.}
 \label{v36:slow-rate}
\end{equation}
They are distinct from the twenty-two rapid hyperbolic modes of V35.
Consequently at least twelve expanding and twelve contracting canonical
eigenvalues are present, counted algebraically. The maximum absolute real
part outside the twenty-two rapid modes has the sharp order
\begin{equation}
 \boxed{\max_{z\notin\text{the rapid hyperbolic groups}}
                  |\operatorname{Re}z|\asymp a^{-1}.}
 \label{v36:sharp-growth}
\end{equation}
No simplicity, unique slow positive root, numerical value of $c,C$, or
nonautonomous growth estimate is asserted.
\end{theorem}
\begin{proof}
At small positive $\zeta$, eliminate the nonconstant block, which tends to
$K_{1,R}$. The harmonic Schur matrix is
$-\mathsf T/\zeta+O(1)$, so
\begin{equation}
 \det\mathcal P_0(\zeta)
      =-\det(K_{1,R})\det(\mathsf T)\zeta^{-3}(1+O(\zeta))<0.
 \label{v36:small-sign}
\end{equation}
The minus sign is important: there are three retained constant shifts.

For sufficiently large real $\zeta$, the symmetric part
$S_\zeta=\zeta\mathcal G-\mathcal H/\zeta$ is positive on
$\mathcal A\oplus\mathcal H$ and zero on $\mathcal W_g$.
If $(S_\zeta+\mathcal J)x=0$, taking the real part puts $x$ in
$\mathcal W_g$. Its $R$ row then gives $K_{1,R}x=0$, hence $x=0$.
The matrix is invertible. Its determinant is positive: adding $\epsilon I$
makes its Hermitian part strictly positive, so all real eigenvalues are
positive and all nonreal eigenvalues occur in conjugate pairs of positive
product. Let $\epsilon\downarrow0$ and use the just proved nonsingularity.
Thus $\det\mathcal P_0(\zeta)>0$ at every sufficiently large positive
$\zeta$.

Choose two fixed finite positive endpoints where these signs are strict.
Uniform convergence in \eqref{v36:pencil} preserves the signs for
$\mathcal P_a$ when $a$ is small. Its entries are real and continuous on
the intervening interval, where $|\zeta/a|$ lies beyond the bounded free
canonical spectrum. The intermediate value theorem therefore gives
$\det\mathcal P_a(\zeta_a)=0$ for some endpoint-bounded $\zeta_a>0$.
The exact determinant correspondence in \eqref{v36:Pa} gives
$z_a=\zeta_a/a$ as an actual canonical eigenvalue. Hamiltonian spectral
symmetry gives $-z_a$. Their scaled rapid values $a^3z_a=a^2\zeta_a$
tend to zero, so they are not the separated rapid hyperbolic roots.
Combining their lower bound with \cref{v36:growth-ceiling} proves
\eqref{v36:sharp-growth}. No spectral sampling is used in this argument.
\end{proof}

The sign change does not require values for the skew cubic coefficients
$\mathsf R$ and $\mathsf S$. They remain essential for locating the roots
and their eigenvectors. Nor does positive translation energy alone suffice:
the proof also uses the actual maximal Poisson rank and positivity of the
retained active-plus-translation mass. Removing the constant shifts would
remove the very rows producing \eqref{v36:small-sign}, and would change the
original problem.

\section{Consequences for preparation and the next propagation calculation}
\label{v36:direction}

The limiting pencil depends only on the canonical first field and its
original leading source coefficients. Thus changing higher preparation
within the V30--V33 class, including selecting a different small initial rate,
does not erase this $a^{-1}$ \emph{frozen spectral} pair. Its excitation can
still depend on that selection. A spectral mode can be present without being
excited by a particular exact history; the present result is not an argument
that every selected history grows.

In particular V33's $ca^2$ physical interval is not contradicted. A frozen
rate of order $a^{-1}$ acts by a bounded factor $\exp(O(a))$ on that short
interval. Extending selection to longer intervals must now retain both the
rapid $a^{-3}$ hyperbolic sector and this additional inverse-amplitude scale.
The exact three-row pencil \eqref{v36:three-rows} is a concrete next coefficient
calculation: evaluate the already identified cubic envelopes there, locate
and separate all of its nonzero roots, and compare the resulting subspaces
with the original time-dependent constraint-tangent system.

Two distinctions remain compulsory. First, the radial energy level is a moving
physical hypersurface. As in V35, a frozen eigenvector is not automatically
tangent to it. V35's quantitatively controlled radial correction was proved
for the rapid branch only and is not silently applied to the present slower
modes. Second, no original metric-curvature visibility certificate for the
new slow eigenvector is asserted here. The translation-energy coefficient
is an actual canonical calculation, but it is not itself that nonlinear
geometric readout. These checks are separate from the canonical spectral
existence proved above.

\section{Verification and status}
\label{v36:status}

\paragraph{Proved in this body.}
The constant-shift Poisson blocks have the exact additional preparation
orders in \eqref{v36:Korders}. Their canonical feedback has the uniform
translation suppression \eqref{v36:Hbounds}. Every canonical real part
outside the inherited rapid hyperbolic groups is $O(a^{-1})$. The actual
inverse-amplitude pencil is specified by quadratic/cubic envelopes and the
known signed source. Its mass and translation energy have the proved signs,
and its nonconstant block is invertible on the full positive-real axis.
A determinant sign change forces at least one additional real pair of order
$a^{-1}$, making the nonrapid growth bound sharp.

\paragraph{Inherited inputs and executed checks.}
The nonlinear root, maximal first Poisson rank, signed Gram and rapid center
are inherited. The new checker reconstructs the original flat operators,
verifies their translation identities, and checks the physical axial
translation-energy matrix. Independent rational Hamiltonian germs test the
small-scale determinant and its mixed translation brackets; a separate
positive/skew example retains massless coordinates in the three-row Schur
reduction. These are algebra checks, not alternative physical sources.
The V35 verifier, including its nested certificate acceptance, was rerun.
No full nonlinear action Hessian was newly extracted.

\paragraph{Remaining.}
The new slow roots and their physical tangent representatives have not been
numerically enclosed or completely classified. There is no new propagator
bound, longer nonlinear lifespan, spatial-refinement theorem, or normalized
metric/link matching result. The first-field choice may be changed upstream,
but it must retain the original action and reverify its source assumptions.
The raw-action Einstein continuum bridge remains unclosed. All thirty-five
prior marked mathematical bodies are preserved byte-for-byte. No nonlinear
spacetime integration, proof-assistant formalization, or independent reproof
of every inherited analytic statement is claimed.

\begin{thebibliography}{9}
\raggedright
\bibitem{v36:TM}
F. Tisseur and K. Meerbergen,
\emph{The quadratic eigenvalue problem}, SIAM Review \textbf{43} (2001),
235--286, \href{https://doi.org/10.1137/S0036144500381988}
{doi:10.1137/S0036144500381988}.
The primary publisher and author-repository records were checked for the
standard polynomial-pencil framework. No external spectral estimate or
algorithm from that survey replaces the original-source derivation here.
\end{thebibliography}
% END IMMUTABLE EXACT-ACTION BRIDGE V36 BODY

% Future iterations append here.
% BEGIN IMMUTABLE EXACT-ACTION BRIDGE V37 BODY
\clearpage
\renewcommand{\thesection}{\arabic{section}}
\renewcommand{\theHsection}{bridgeV37.\arabic{section}}
\setcounter{section}{297}
\section{The full slow index before another propagation estimate}
\label{v37:context}

The sign argument in \cref{v36:slow-real} supplies a real pair at the
inverse-amplitude scale, but does not count the other roots of its singular
pencil. Computing only a positive root of the three-row determinant would
leave that issue unresolved. We instead linearize the \emph{reciprocal}
spectral parameter. The retained constant-shift stiffness makes this
linearization regular. After quotienting only its exact algebraic mass
kernel, it has dimension 84 and a nonsingular metric with inertia $(81,3)$.
This is a spectral calculation, not removal of original constraint rows.

Fresh evaluation of the original cubic envelopes at the entire certified
first-field root encloses one real reciprocal pair and one imaginary pair
whose two-dimensional metric is negative. Those two invariant planes use
all three negative directions. Their orthogonal complement is positive.
This proves that the \emph{complete limiting slow pencil has exactly one
expanding pair}, rather than merely finding one of several possible pairs.
No numerical decision about a small reciprocal zero cluster is needed.

There is a second improvement. A correction of a frozen eigenvector to the
radial energy tangent should not be allowed to create an uncontrolled
multiplier change. The original four mean-calibration directions permit an
analytic canonical normal which annihilates all 107 normalized multiplier
rows exactly. This supplies constraint-tangent representatives of the new
slow mode without changing its multiplier component. It does not, by
itself, evaluate the time-dependent curvature variation of that mode.

\paragraph{Dependencies and scope.}
We retain the actual root and compatible initial chart of
\cref{v29:root,v29:compatibility}, the full leading signed source and
initial-rate preparation of V30, and the exact canonical transfer of
V35--V36. In particular the finite-amplitude convergence of the slow pencil
is \cref{v36:pencil-limit}, not an assumption added here. All calculations
are at $N=3$ and concern $a\downarrow0$ at that fixed spatial cutoff.
The original higher signed-source enclosure is inherited and rechecked;
its cubic mixed Poisson coefficients are freshly evaluated below. Index
methods for quadratic pencils are established tools
\cite{v36:TM,v37:BJK}; the finite reciprocal and inertia arguments used
here are proved directly. No external instability theorem is substituted
for the original source calculation.

\section{A regular reciprocal descriptor with only three negative directions}
\label{v37:descriptor}

Use the real multiplier order
$\mathcal A\oplus\mathcal W_g\oplus\mathcal H$ of V36, with dimensions
$78,26,3$, and let $I_H:\R^3\to\mathcal E$ synthesize the three original
constant shifts. Write
\[
 \mathcal P_0(\zeta)=\zeta\mathcal G+\mathcal J-\zeta^{-1}\mathcal H,
 \qquad \mathcal H=I_H\mathsf T I_H^*.
\]
The coefficients here are \eqref{v36:pencil-coefficients}, including
$\mathsf R$ and $\mathsf S$. The actual $\mathsf T$ is positive definite;
$\mathcal G$ is positive on $\mathcal A\oplus\mathcal H$ and its kernel
is exactly $\mathcal W_g$. The nonconstant skew block $K_{1,R}$ is invertible.

\begin{theorem}[Regular descriptor and exact reciprocal quotient]
\label{v37:reciprocal}
On $\mathcal E\oplus\R^3$, define
\begin{equation}
 \mathbf S=\begin{pmatrix}
 \mathcal J&I_H\mathsf T\\-\mathsf T I_H^*&0
 \end{pmatrix},\qquad
 \mathbf E=\begin{pmatrix}\mathcal G&0\\0&-\mathsf T\end{pmatrix}.
 \label{v37:descriptor-matrices}
\end{equation}
Then $\mathbf S$ is invertible and skew, and
\begin{equation}
 \det\mathbf S=\det K_{1,R}(\det\mathsf T)^2>0,\qquad
 \det(\mathbf S+\zeta\mathbf E)
       =(-\zeta)^3\det\mathsf T\,\det\mathcal P_0(\zeta)
       \quad(\zeta\ne0).
 \label{v37:descriptor-det}
\end{equation}
Put $\mathbf C=-\mathbf S^{-1}\mathbf E$. Its exact kernel is
$\ker\mathbf E=\mathcal W_g\oplus\{0\}$, of dimension 26. On the
84-dimensional quotient by this kernel there is a nonsingular symmetric
metric $E_o$ with inertia $(81,3)$ and an induced map $C_o$ satisfying
\begin{equation}
                 E_oC_o+C_o^*E_o=0.                 \label{v37:skew-metric}
\end{equation}
The complete nonzero finite roots of the original slow pencil are the
reciprocals of the nonzero eigenvalues of $C_o$, with algebraic multiplicity:
\begin{equation}
 \boxed{\det(\mathbf S+\zeta\mathbf E)
        =\det\mathbf S\det(I_{84}-\zeta C_o).}
 \label{v37:quotient-det}
\end{equation}
Zero eigenvalues of $C_o$ correspond to the singular pencil's infinite
spectral part, not to additional finite slow eigenvalues.
\end{theorem}
\begin{proof}
If $\mathbf S(x,w)=0$, its bottom row gives $x_H=0$. The top nonconstant
rows then give $K_{1,R}x_R=0$, so $x_R=0$. The remaining three rows give
$\mathsf T w=0$. This proves invertibility without an assumed inverse
for the complete skew $\mathcal J$. Eliminating the invertible $K_{1,R}$
block in $\mathbf S$ leaves
$\left(\begin{smallmatrix}S'&\mathsf T\\-\mathsf T&0\end{smallmatrix}\right)$,
whose determinant is $(\det\mathsf T)^2$, independently of the skew $S'$.
A real nonsingular skew matrix has positive determinant, proving the first
formula. Eliminating the bottom block $-\zeta\mathsf T$ in
$\mathbf S+\zeta\mathbf E$ leaves exactly $\mathcal P_0(\zeta)$,
with the displayed minus sign.

Since $\mathbf S$ is invertible, $\ker\mathbf C=\ker\mathbf E$.
Moreover $\mathbf E\mathbf C=-\mathbf E\mathbf S^{-1}\mathbf E$ is
skew. Choose any real complement $O$ to this kernel and order the coordinates
as $O\oplus\ker\mathbf E$. In those coordinates,
\[
 \mathbf E=\begin{pmatrix}E_o&0\\0&0\end{pmatrix},\qquad
 \mathbf C=\begin{pmatrix}C_o&0\\ *&0\end{pmatrix}.
\]
The nondegenerate metric has the inertia of $\mathcal G$ on its
81-dimensional support, together with that of $-\mathsf T$: it is
$(81,3)$. Taking the upper block of the skew identity proves
\eqref{v37:skew-metric}. Taking the determinant of
$\mathbf S(I-\zeta\mathbf C)$ proves \eqref{v37:quotient-det}.
At finite nonzero $\zeta$, the prefactors in
\eqref{v37:descriptor-det} are nonzero, so root multiplicities agree.
This quotient removes an exact kernel of an auxiliary reciprocal map.
It does not delete the 26 gradient equations from $\mathcal P_0$.
\end{proof}

The construction distinguishes two possible zeros: the 26-dimensional
kernel of $\mathbf C$ is exact by construction, whereas the induced $C_o$
may have further zero eigenvalues. Their number is not determined by that
kernel statement.

\begin{proposition}[An a priori bound on slow hyperbolicity]
\label{v37:index-bound}
The algebraic dimension of the right-half-plane generalized eigenspace of
$C_o$, and hence the number of right-half-plane finite slow roots, is at
most three. The same bound holds in the left half-plane.
\end{proposition}
\begin{proof}
Complexify the real spaces and metric. Equation \eqref{v37:skew-metric}
implies $e^{tC_o^*}E_oe^{tC_o}=E_o$. On the right-half-plane generalized
space, $e^{-tC_o}$ tends to zero as $t\to+\infty$, including its finite
Jordan polynomial factors. Applying the metric identity to two such
vectors proves that their pairing is zero. This space is therefore totally
isotropic for a Hermitian form of inertia $(81,3)$ and has dimension at
most three. For completeness, projection onto the negative coordinate
space is injective on any isotropic space: a vector with zero negative
component has nonnegative squared form equal to zero, hence vanishes.
This proves the dimension bound without diagonalizability. Reciprocal
inversion preserves the sign of the real part away from zero. The left
space is handled by $t\to-\infty$.
\end{proof}

\section{The actual mixed and translation coefficients are now enclosed}
\label{v37:source}

The numerical data used below describe the original V29 root, not a nearby
rational seed substituted for it. The certificate retains the dyadic root
center of \cref{v29:root} and its complete cube of radius $2^{-240}$.
All coefficient operations are performed with outward-rounded dyadic
\emph{integer} intervals at 128 binary places. A cube narrower than that
working scale is enlarged, never contracted, by this arithmetic.

The spatial-sum convention is used for every Gram and covector matrix in
the certificate. It multiplies $\mathcal G,\mathcal J,\mathcal H$ by the
same factor 27 relative to the spatial-mean convention, and therefore
changes none of the roots. The literal phase derivatives are
$\delta_i=3(S_i-S_i^{-1})$, and the physical $\sqrt3$ is enclosed by
integer square-root bounds rather than rounded to a selected decimal.

The following are computed freshly from the retained original scalar
coefficient functions:
\begin{equation}
 \begin{split}
 \mathsf R[c,\nu]={}&DP_{c,3}(Y_*)[G\nu]
     -DP_{\nu,2}(Y_*)[\delta_cY_*]
     -\langle\delta_c\nu,P_2(Y_*)\rangle,\\
 \mathsf S[c,d]={}&DP_{c,3}(Y_*)[\delta_dY_*]
                    -DP_{d,3}(Y_*)[\delta_cY_*],\\
 \mathsf T_{cd}={}&2\langle\delta_cp_*,\delta_dp_*\rangle
       +\tfrac12\sum_j\langle\delta_j\delta_cU_*,
                                       \delta_j\delta_dU_*\rangle.
 \end{split}                                                    \label{v37:source-formulas}
\end{equation}
The last equality uses the actual transverse trace-free nonconstant fields;
the differentiated homogeneous momentum is zero. The original Frobenius
pairing counts both occurrences of each off-diagonal entry.

\begin{proposition}[Complete source enclosure at the specified root]
\label{v37:source-enclosure}
The delivered coefficient bank encloses all 324 entries of $\mathsf R$
against the complete 108 original multiplier tests, all nine entries of
$\mathsf S$, and all nine entries of $\mathsf T$, at every point of the
retained root cube. The largest entry widths are less than
\[
 8\,10^{-32},\qquad 11\,10^{-31},\qquad 3\,10^{-33},
\]
respectively. The physical first Poisson matrix is enclosed from all five
original coefficient matrices of the already established polynomial bank.
The resulting exact $\mathcal J$ and $\mathsf T$ have the definitions
in V36, not comparison values inferred from numerical eigenvectors.
\end{proposition}
\begin{proof}
The archive supplies the five exact first-field basis records and the
original rational $P_2$, $\Phi_3$, and constant-shift $P_{c,3}$ routines.
Their function bodies are retained unchanged. For a quadratic functional,
its derivative is obtained from the half difference at $Y+Z,Y-Z$; for a
cubic it is that half difference minus its value at $Z$. Apply these exact
polarization rules to \eqref{v37:source-formulas}. Discrete summation by
parts turns its last preparation term into the same component-major
$\delta_cP_2$ contribution used in V21. The constant-shift envelope retains
its temporal auxiliary response and original electric and magnetic cubic
terms. No vanishing is inferred from interval proximity to zero.

Every evaluation is a finite composition of additions, products, literal
permutations, contractions, and the specified square-root enclosure.
The supplied interval primitives bound each such operation with Python
integers; their denominator is $2^{128}$. Induction through the unmodified
functional therefore encloses its exact value. The complete output arrays,
not just their displayed maximum widths, are serialized in the package.
The first Poisson bank is converted from its rational congruence by
$K=\sqrt3 Q_3^{-1}\widehat K Q_3^{-1}$, with both $Q_3$ factors retained.
All 108 tests are evaluated before restriction to the normalized physical
basis. This proves the enclosure statement by a finite, reproducible
certificate.
\end{proof}

The known 162-by-107 reduced source enclosure and its tangent metric $K_0^{-1}$
are inherited from V30. We form their full Gram in the stated physical
basis, then apply the \emph{prescribed limiting mass} $F_*\Gamma_*F_*$.
This last operation belongs to V36's spectral limit; it is not imposed on
the finite multiplier Hessian. Active representative columns are first
projected by the exact rational phase-Hodge matrix; their leading source
is unchanged by the universal weak first-source cancellation. The proof
package records both this synthesis and the literal input hashes.

\section{A rational spectral certificate with a negative oscillatory plane}
\label{v37:certificate}

The finite certificate specifies two proposed real coordinate matrices. The
first orders the 84-dimensional support of $\mathbf E$ before its 26-dimensional
kernel and approximately normalizes its metric. The second approximately
block-diagonalizes the induced reciprocal matrix. Each proposed binary
floating entry is interpreted as its \emph{exact binary rational value}.
No floating inverse, eigensolver decision, or rank threshold is accepted.

Let $E_o$ be the supported metric after the first congruence and let
$C_o$ be its induced reciprocal operator. The second rational similarity
$P$ gives
\[
                 B=P^{-1}C_oP=B_0+\mathcal E.
\]
The reference $B_0$ is a real normal matrix: real scalar blocks, real
rotation blocks, and a zero reference cluster. A zero block of $B_0$ is
\emph{not} a declaration that the corresponding eigenvalues of $B$ vanish.
All proposed matrices, exact residual bounds and interval endpoints are
stored in \src{results/slow_proposals.json} and
\src{results/slow_spectral_certificate.json}.

\begin{theorem}[Two isolated pairs and their metric signs]
\label{v37:certified-pairs}
For the actual root, the supported reciprocal operator has one simple
real pair $\pm\mu_+$, with $\mu_+>0$, and a simple imaginary pair
$\pm i\omega_-$ whose real invariant plane is negative definite for
$E_o$. Their reciprocals satisfy the strict rational decimal bounds
\begin{equation}
 \boxed{0.1585835166<\gamma:=\mu_+^{-1}<0.1585835168,}
 \label{v37:gamma}
\end{equation}
\begin{equation}
 \boxed{0.03017120680<\varpi:=\omega_-^{-1}<0.03017120690.}
 \label{v37:negative-frequency}
\end{equation}
The original slow pencil therefore has the simple roots $\pm\gamma$ and
$\pm i\varpi$.
\end{theorem}
\begin{proof}
We state the acceptance inequalities and their implications explicitly.
For a proposed rational inverse $R$ of an interval matrix $A$, put
$E=I-RA$ and $q=\|E\|_\infty<1$. Then
\begin{equation}
 A^{-1}=R+ER+E^2R+\mathcal T,\qquad
 \|\mathcal T\|_\infty\le\frac{q^3\|R\|_\infty}{1-q}.
 \label{v37:inverse-enclosure}
\end{equation}
The identity follows by summing $(I-E)^{-1}R$ in its actual order.
An entrywise enlargement by the last bound encloses the inverse. It is
used for both the descriptor inverse and $P^{-1}$.

The exact integer/rational checks give the strict bounds in this table.
A norm of an interval matrix means the supremum bounded entrywise, not
its midpoint norm.
\begin{center}\small
\begin{tabular}{@{}p{.61\textwidth}r@{}}
\toprule
Acceptance quantity & Upper bound\\\midrule
Descriptor inverse residual $\|I-R_S S_X\|_\infty$ & $6.22\,10^{-10}$\\
Supported metric error from $\diag(I_{81},-I_3)$ & $6.20\,10^{-13}$\\
Similarity inverse residual $\|I-R_PP\|_\infty$ & $5.45\,10^{-14}$\\
$\delta=\max(\|B-B_0\|_1,\|B-B_0\|_\infty)$ & $1.109\,10^{-9}$\\
Two-disk spectral projection distance $\epsilon$ & $2.218\,10^{-8}$\\
\bottomrule
\end{tabular}
\end{center}
In particular $\|B-B_0\|_2\le\delta$. The metric error proves inertia
$(81,3)$ directly by continuity from the nonsingular diagonal metric.

The four reference centers selected in the certificate are a positive and
negative real center and a positive and negative imaginary center. Each is
more than $1/5$ from every other reference center. Take the circle of radius
$r=1/10$ about each. On those circles the resolvent of normal $B_0$ has
norm at most $1/r$. The Neumann estimate for $B_0+t\mathcal E$ is valid
uniformly for $0\le t\le1$ because $\delta<r$. Each circle therefore
contains exactly one eigenvalue, with algebraic multiplicity one, by
continuity of its spectral projection. The sharper normal-resolvent bound
places that eigenvalue within $\delta$ of its reference center.
Conjugation and uniqueness make the real roots real.
The symmetry $\mu\mapsto-\overline\mu$ from
\eqref{v37:skew-metric}, with uniqueness in each imaginary disk, makes the
other two roots exactly imaginary. Inverting the corresponding exact
rational center intervals gives \eqref{v37:gamma}--\eqref{v37:negative-frequency}.

The imaginary pair's sign is also checked, not inferred from its frequency.
Let $\Pi_0$ be the orthogonal projection onto its two reference coordinates,
and let $\Pi$ be the union of the two exact Riesz projections. The resolvent
identity integrated around the two circles gives
\begin{equation}
                \|\Pi-\Pi_0\|_2\le
                  \epsilon:=\frac{2\delta}{r-\delta}.
 \label{v37:projector-error}
\end{equation}
Set $Q=P^*E_oP$. On $\Ran\Pi_0$, the exact interval Gershgorin upper
bound is strictly negative; write its absolute margin as $m_0>0$.
For a unit $v$ in this reference plane,
\[
 (\Pi v)^*Q(\Pi v)
       \le-m_0+\|Q\|_2(2\epsilon+\epsilon^2).
\]
The delivered rational calculation verifies
\begin{equation}
 m_0-\max(\|Q\|_1,\|Q\|_\infty)
                          (2\epsilon+\epsilon^2)>0.1498>0.
 \label{v37:negative-margin}
\end{equation}
As $\epsilon<1$, $\Pi$ maps the reference plane injectively onto the
exact two-dimensional spectral plane. This proves negativity on that
whole plane. All displayed acceptance inequalities are comparisons of
complete rational fractions; decimals in the table are strict bounds
checked against those fractions.
\end{proof}

\section{The complete finite slow spectrum has exactly one expanding pair}
\label{v37:full-index}

\begin{theorem}[Inertia exhaustion and a bounded remaining reciprocal center]
\label{v37:full-count}
At the original certified root, the limiting pencil $\mathcal P_0$ has
exactly one root in the open right half-plane and one in the open left
half-plane, with algebraic multiplicity. They are the simple real roots
$\pm\gamma$ in \eqref{v37:gamma}. Every other finite nonzero slow root
is purely imaginary and semisimple.

On the 84-dimensional reciprocal quotient the invariant splitting is
\begin{equation}
 \R^{84}=E_{\rm hyp}\mathbin{\perp_{E_o}}E_{\rm neg}
                         \mathbin{\perp_{E_o}}E_{\rm pos},
 \quad \dim E_{\rm hyp}=2,\quad\dim E_{\rm neg}=2,\quad\dim E_{\rm pos}=80.
 \label{v37:inertia-splitting}
\end{equation}
Its respective metric inertias are $(1,1),(0,2),(80,0)$. The center
$E_{\rm neg}\oplus E_{\rm pos}$ admits a positive conserved quadratic
form and a bounded reciprocal semigroup for all real times. This does
not make the original $E_o$ positive on that center.
\end{theorem}
\begin{proof}
For an $E_o$-skew map, generalized spectral spaces whose eigenvalues do
not sum to zero after conjugation are $E_o$-orthogonal. One proof uses
the Sylvester equation for the restricted pairing; its homogeneous equation
has only zero solution when the two finite spectra are disjoint after
that reflection. Equivalently the exponential metric identity and the
Jordan expansion give the same pairings recursively.

The two real eigenlines at $\pm\mu_+$ are individually isotropic. Their
mutual pairing is nonzero: otherwise one line would be orthogonal to its
partner and all other generalized eigenspaces, contradicting nonsingularity
of $E_o$. Thus their real plane is nondegenerate of inertia $(1,1)$.
The isolated imaginary plane is negative definite by
\cref{v37:certified-pairs} and is orthogonal to that real plane. The
orthogonal complement of their direct sum is invariant, nondegenerate,
and has inertia
\[
                   (81,3)-(1,1)-(0,2)=(80,0).
\]
On this complement, conjugating $C_o$ by the positive square root of its
metric gives a real skew-adjoint operator. Its eigenvalues are imaginary,
including possible semisimple zeros, and its exponential is bounded.
The same argument with the negative of the metric applies to the negative
imaginary plane. Their positive metric direct sum is conserved on the
remaining center. This proves the complete eigenvalue count, not merely
a statement about the isolated disks. Reciprocal inversion and
\eqref{v37:quotient-det} transfer it, with multiplicities and semisimplicity,
to every nonzero finite root of $\mathcal P_0$.
\end{proof}

The negative oscillatory plane is the information needed to sharpen the
three-pair upper bound of \cref{v37:index-bound} to one pair. Its sign is
not a physical spacetime signature. Nor is the reciprocal semigroup a
physical-time propagation operator. On the nonzero spectral part one may
invert $C_o$ to obtain a generator whose eigenvalues are the finite slow
roots; its remaining center is again bounded in a positive metric. No
inverse is taken on any reciprocal zero space.

The certificate does not decide the dimension of that zero space from the
small numerical reference cluster. Accordingly no exact total count of
finite imaginary pairs, no count of infinite pencil roots, and no exact
zero-mode count for the nonlinear action is asserted.

\section{What the slow count implies for the original canonical matrix}
\label{v37:canonical}

\begin{theorem}[A located original slow pair and bounded correction in slow annuli]
\label{v37:canonical-roots}
On the analytic prepared families of V30--V36, for sufficiently small
positive $a$, the original canonical matrix has a simple real pair
\begin{equation}
                  \boxed{z_+(a)=\gamma/a+O(1),\qquad z_-(a)=-z_+(a).}
 \label{v37:physical-rate}
\end{equation}
The coefficient $\gamma$ is the actual enclosed number in
\eqref{v37:gamma}, independent of unlisted higher preparation coefficients.

For any fixed compact annulus $0<r\le|\zeta|\le R<\infty$, every other
original canonical eigenvalue with $\zeta=az$ in this annulus has
\begin{equation}
                         |\operatorname{Re}z|\le C_{r,R},
 \label{v37:slow-annulus}
\end{equation}
after removing the two small neighborhoods of the real limiting roots.
Multiplicity is retained at any imaginary cluster. This statement neither
covers $az\to0$ nor covers $|az|\to\infty$.
\end{theorem}
\begin{proof}
The exact original transfer satisfies
$\mathcal P_a(\zeta)=\mathcal P_0(\zeta)+O(a)$ with its fixed spectral
derivatives uniformly on compact sets separated from zero by
\cref{v36:pencil-limit}. The real roots of $\mathcal P_0$ are simple by
\cref{v37:full-count}; their determinant derivatives are nonzero. The
analytic implicit-function theorem therefore gives unique real roots
$\zeta_\pm(a)=\pm\gamma+O(a)$. The transfer correspondence identifies
them with $az_\pm(a)$, proving \eqref{v37:physical-rate}. The original
Hamiltonian matrix has opposite-pair symmetry, so the pairing is exact.

For the annulus assertion it is useful to state the resolvent estimate
rather than infer a rate from polynomial root continuity. The quotient
$C_o$ is similar, with a fixed bounded condition number, to real scalar
hyperbolic blocks and skew-adjoint blocks on its two definite centers.
Thus, away from its spectrum, its resolvent is bounded by a constant
inverse distance to the spectrum. The block triangular lift back to
$\mathbf C$ adds a bounded factor on compact nonzero $\zeta$ sets, since
its other diagonal block in $I-\zeta\mathbf C$ is the identity.
The bottom Schur elimination in \eqref{v37:descriptor-det} is also bounded
on those compact sets. Consequently
\[
 \|\mathcal P_0(\zeta)^{-1}\|
     \le C_{r,R}\operatorname{dist}(\zeta,\mathcal Z_0)^{-1},
\]
locally near its finitely many roots $\mathcal Z_0$ in an enlarged annulus;
elsewhere the inverse is uniformly bounded. This includes repeated
semisimple imaginary roots and does not divide by an individual imaginary
frequency approaching zero.

At a zero of the original $\mathcal P_a$, a unit kernel vector has
$\|\mathcal P_0(\zeta)v\|\le Ca$. The resolvent estimate puts that
root within $O(a)$ of $\mathcal Z_0$. Outside the real root groups the
limiting roots are imaginary, so $|\operatorname{Re}\zeta|\le Ca$.
Divide by $a$ to obtain \eqref{v37:slow-annulus}. The exact determinant
correspondence from V36 preserves its interpretation as a canonical
statement.
\end{proof}

The extra positive slow rate and the eleven positive rapid rates occupy
different spectral scales. The preceding theorem does not say that they
exhaust every expanding eigenvalue of the complete finite-amplitude
324-dimensional matrix. In particular, smaller or intermediate scales and
nonnormal transient behavior remain distinct questions. It also does not
convert frozen spectral information into a time-dependent growth theorem.

\section{A radial normal which preserves the multiplier component exactly}
\label{v37:radial}

The energy constraint is not discarded by the spectral elimination. Let
$C_{\rm full}=D_XP$ denote the derivative of all 108 original constraint
rows at a selected initial point, let $C$ be its restriction to the fixed
107-dimensional mean-lapse-zero multiplier complement, and put
$g_a=\nabla\mathscr H^{\rm red}(X_a)$. Homogeneity gives
$g_a=C_{\rm full}^*\lambda_a$ there. Every norm and adjoint in this
section uses the original physical pairings.

\begin{lemma}[A uniformly bounded canonical normal in the kernel of all normalized rows]
\label{v37:lossless-normal}
There is a real analytic bounded vector $W_a$ for small $a\ge0$ such that
\begin{equation}
             \boxed{CW_a=0,\qquad \langle g_a,W_a\rangle=a,\qquad
                    T_aW_a=\mathbb JQW_a=O(1).}
 \label{v37:normal}
\end{equation}
These are exact identities at each nonzero $a$, not leading-order
annihilations only.
\end{lemma}
\begin{proof}
Split the original multiplier rows into their 104 nonconstant Fourier
coordinates and their four means. Divide the latter rows of
$C_{\rm full}$ by $a$. On the analytic prepared initial family this
row-scaled matrix extends analytically to $a=0$. Its top part has limit
the nonconstant flat derivative $L$, and its mean part has limit
$P_0DP_2(Y_*)$.

The latter four rows are independent on $\ker L$. This is the original
mean-calibration submersion in \cref{v29:compatibility}, not a new rank
assumption: varying the homogeneous diagonal boost supplies scalar mean
derivative $32\omega_h^2=864$, while the three retained axial sine
momentum directions supply the three nonzero shift-mean derivatives.
Their other mean components vanish, and all lie in $\ker L$.
The 108-row limiting matrix is therefore onto. Choose a fixed nonzero
108-column minor. Its inverse gives a bounded analytic right inverse of
the row-scaled matrix for small $a$.

Apply that right inverse to the row vector with scalar mean one and all
other coordinates zero. Restoring the mean scaling yields a bounded
analytic $W_a$ with
\[
       \langle\nu,C_{\rm full}W_a\rangle=a\overline{\nu^N}
       \quad\hbox{for every original multiplier }\nu.
\]
It annihilates every normalized complementary row exactly. Pairing with
$\lambda_a$, whose mean lapse is one, gives
$\langle g_a,W_a\rangle=a$. Finally the canonical formula for $T_a$
has term $-\mathbb JC^*M^{-1}CW_a=0$ identically; its remaining $Q$
is bounded on the original stationary chart. This proves the last identity.
The original scalar-row sign can be accommodated in the target of the
same right inverse; the displayed pairing fixes its convention.
\end{proof}

Unlike the fixed homogeneous normal used for the rapid mode in V35, this
normal lies in $\ker C$ \emph{at the actual finite-amplitude point}.
Its construction costs only the four already verified mean derivatives.
It avoids multiplying an $O(a)$ normalized constraint residual by the
singular multiplier inverse.

\begin{theorem}[An admissible slow canonical direction with unchanged multiplier lift]
\label{v37:slow-tangent}
Let $v=(v_R,c)$ be a real null vector of $\mathcal P_0(\gamma)$,
normalized in any fixed nonzero convention. Then $c\ne0$. For the real
canonical root in \eqref{v37:physical-rate}, choose corresponding real
null vectors $v_a=v+O(a)$ of $\mathcal P_a(\zeta_a)$, where
$\zeta_a=az_+(a)$. An eigenvector lift may be normalized as
\begin{equation}
 \begin{split}
  \eta_a&=a B_a^{-1}v_a=(a v_{R,a},c_a),\\
  \xi_a&=(z_+(a)I-F)^{-1}\mathbb JC^*\eta_a
     =\frac{a^2}{\gamma}\{Gv_R+\delta_cY_*\}+O(a^3).
 \end{split}                                                    \label{v37:slow-lift}
\end{equation}
In particular $\|\xi_a\|\asymp a^2$. Define
\begin{equation}
 \widetilde\xi_a=\xi_a-\frac{\langle g_a,\xi_a\rangle}{a}W_a,
                       \qquad\widetilde\eta_a=\eta_a.
 \label{v37:tangent-correction}
\end{equation}
Then
\begin{equation}
 \boxed{C\widetilde\xi_a+M\widetilde\eta_a=0,
        \quad\langle g_a,\widetilde\xi_a\rangle=0,
        \quad\widetilde\xi_a-\xi_a=O(a^3),}
 \label{v37:tangent}
\end{equation}
and
\begin{equation}
 \boxed{\|(aT_a-\zeta_a I)(a^{-2}\widetilde\xi_a)\|\le Ca.}
 \label{v37:quasimode}
\end{equation}
These vectors are tangents to local curves of exactly constrained original
initial records. They are not asserted to form an invariant subspace of
the moving physical energy hypersurface.
\end{theorem}
\begin{proof}
If $c=0$ in a positive-real pencil null vector, its nonconstant equation
is $B_R(\gamma)v_R=0$. That block is invertible by
\cref{v36:Rinvertible}, so the entire vector would vanish. Simplicity of
the enclosed root supplies $v_a=v+O(a)$ by an analytic null-vector chart.
The exact congruence in \eqref{v36:Pa} maps $v_a$ to the multiplier
vector $B_a^{-1}v_a$; multiplying both lifted vectors by $a$ gives the
chosen normalization. V35's resolvent lift then gives
$C\xi_a+M\eta_a=0$ and $T_a\xi_a=z_+(a)\xi_a$ exactly.

The nonconstant rows of $C$ tend to $L$, whereas
$\mathbb JC_H^*c=a\delta_cY_*+O(a^2)$ by the original quadratic
translation coefficient. Since
$(z_+(a)I-F)^{-1}=a\gamma^{-1}I+O(a^2)$, these identities prove
\eqref{v37:slow-lift}. The vector in braces is nonzero: $Gv_R$ has zero
transverse-trace-free projection on each nonzero phase mode in both
canonical slots, whereas $\delta_cY_*$ is transverse and trace-free.
It cannot vanish for $c\ne0$, since its quadratic energy is
$\langle c,\mathsf T c\rangle>0$. This proves the norm comparison.

The first radial gradient is
$g_a=a g_0+O(a^2)$, with $g_0=Q_0Y_*+L^*\ell_*$.
The inherited flat subsidiary identities give
$\langle g_0,Gv_R\rangle=0$. Also
\[
 \langle Q_0Y_*,\delta_cY_*\rangle=0,\qquad
 \langle L^*\ell_*,\delta_cY_*\rangle=0:
\]
the first is skew phase-translation invariance of the original quadratic
Hamiltonian, and the second follows from $LY_*=0$ and commutation with
constant phase translations. Thus
$\langle g_a,\xi_a\rangle=O(a^4)$. The coefficient of $W_a$ in
\eqref{v37:tangent-correction} is $O(a^3)$. Using
\cref{v37:lossless-normal} proves all identities in
\eqref{v37:tangent}, including the \emph{unchanged} multiplier vector.
The residual of the corrected eigenvector is that coefficient times
$(aT_a-\zeta_a I)W_a=O(1)$; division by $a^2$ proves
\eqref{v37:quasimode}.

For each nonzero $a$, solving the 107 normalized constraints gives the
original smooth multiplier graph because $M$ is invertible. Its remaining
radial equation is the regular energy hypersurface, since
$\langle g_a,W_a\rangle=a\ne0$. Every tangent vector to that
hypersurface is realized by a local smooth curve in it; lifting by the
same multiplier graph gives exactly constrained original records with
tangent \eqref{v37:tangent}. The curve parameter may be restricted at
fixed $a$ to retain positive lapse and the verified stationary rank chart.
No uniform parameter radius or original-time growth assertion is used.
\end{proof}

\section{The remaining physical observation is not determined by the frozen root}
\label{v37:physical-scope}

The new tangent theorem closes the radial constraint issue for the slow
pair, but does not justify importing V35's rapid curvature-visibility
argument unchanged. Its leading multiplier component in
\eqref{v37:slow-lift} is a \emph{constant shift}, on which
$\mathfrak E_h$ vanishes. Also the exact differential of the original
rate on an uncorrected canonical eigenvector is
\begin{equation}
 \delta r=z\eta-M^{-1}(\dot C\xi+\dot M\eta).
 \label{v37:actual-rate}
\end{equation}
At the slow scale, $\xi=O(a^2)$ and $\dot C=O(a)$ only give
$M^{-1}\dot C\xi=O(a^{-1})$. That has the same possible order as
$z\eta$. It was lower order for the rapid $a^{-3}$ mode, but cannot be
discarded here. The exact constraint-normal correction does not authorize
setting this source to zero. An evaluation of the actual moving coefficient
and the original metric read is still required for slow geometric visibility.

The limit has now been classified without requiring a numerical decision on
all small eigenvalues. Only one slow expanding scalar channel remains after
the eleven rapid expanding channels have been distinguished, within these
two separated spectral scales. The remaining limiting slow directions are
oscillatory, with one negative-energy plane whose sign was explicitly paid
for in the index. A subsequent longer-time construction would need the
actual time-dependent spectral transport, the relevant source couplings,
and compatible initial controls on these different scales. The bounded
reciprocal center alone is not that construction.

All action-specific conclusions remain at the fixed three-site cutoff.
There is no new nonlinear lifespan exponent: V33's $ca^2$ physical interval
is retained. No estimate uniform in spatial refinement, no equality of
normalized metric and link curvature coefficients, and no exact-action
Einstein continuum limit is asserted.

\section{Verification, provenance, and append-only status}
\label{v37:verification}

The original coefficient extractor is rerun over integer dyadic intervals
on the entire inherited V29 root cube. It supplies every mixed and mutual
constant-shift Poisson coefficient and the translated quadratic energy.
The complete physical basis, including the 26 phase gradients and three
constant shifts, is synthesized using exact rational discrete Hodge maps.
The inherited higher-source enclosure enters the full Gram without a
numerical rank truncation. The noncommuting descriptor inverse and quotient
are then enclosed with \eqref{v37:inverse-enclosure}.

The spectral checker verifies the matrix residuals, all four isolation
circles, the two-dimensional negative metric bound, and both reciprocal
root intervals in rational arithmetic. The proof of the complete count
uses that certificate and the direct inertia argument; it does not count
unverified eigenvalue samples. Separate exact algebra tests verify the
interval primitives, descriptor determinants, quotient and metric identities,
original flat translation and transverse-trace-free incidences, and the
lossless radial-normal formulas. The small independent matrix tests are
marked as algebra examples, not replacement gravitational sources.

The original V36 verifier and its nested source/spectral acceptance were
rerun successfully. Those acceptance checks do not constitute a new full
nonlinear Hessian extraction, an independent reproof of every inherited
analytic theorem, or a proof-assistant formalization. The new portable
verifier, source audit, exact certificate files and preservation manifest
are supplied. All thirty-six earlier marked mathematical bodies are retained
byte-for-byte; corrections and limitations in them remain visible.

\begin{thebibliography}{9}
\raggedright
\bibitem{v37:BJK}
J. Bronski, M. A. Johnson and T. Kapitula,
\emph{An Instability Index Theory for Quadratic Pencils and Applications},
arXiv:1207.3764 (2012), \url{https://arxiv.org/abs/1207.3764}.
This is primary-source context for spectral index methods, not an imported
index formula for the present singular-mass pencil. The reciprocal descriptor,
negative-plane certificate, count, and original-action transfers are proved
in this body.
\end{thebibliography}
% END IMMUTABLE EXACT-ACTION BRIDGE V37 BODY

% Future iterations append here.
% BEGIN IMMUTABLE EXACT-ACTION BRIDGE V38 BODY
\clearpage
\renewcommand{\thesection}{\arabic{section}}
\renewcommand{\theHsection}{bridgeV38.\arabic{section}}
\setcounter{section}{306}
\section{The moving source, rather than another frozen spectral count}
\label{v38:context}

The remaining slow observation in \cref{v37:physical-scope} cannot be read
from its frozen eigenvalue. Its leading multiplier is a constant shift,
but differentiation of the original constraint graph also differentiates
its source. This body evaluates that term. It produces a nonconstant weak
rate at order $a^{-1}$, with a certified nonzero original metric-curvature
observation, on exactly constraint-tangent slow representatives. The literal
link-curvature differential on those same representatives is only $O(a)$.
These are differential statements at the initial cut of a fixed-cutoff
family, not a nonlinear departure theorem or an assertion about spatial
refinement.

The new calculation needs the original quadratic constraint coefficient,
the already verified weak signed inverse, and the slow eigenvector. It does
not need a new fourth-order action coefficient, an unlisted second
preparation, or a bound for a higher physical-time derivative. In particular,
the leading observable is independent of the selected bounded coefficient
of the base initial rate. Repeating a frozen count or prescribing that base
rate to be exactly zero would not evaluate this observation.

\paragraph{Related-source audit.}
The four supplied terminal developments were read with their hypotheses and
remaining tasks. Native stationarity V46 uses an exact nonconstant-constraint
chart and the same kind of graded active cancellation as V33, but at its
\emph{own} V39 root and on the invariant graph selected in its V44. Its
neutral forcing has rank two and is not discarded
\cite[\src{v46:thm:rank-two}, \src{v46:sec:status}]{v38:Native}.
Neither its invariant graph nor its spectral indices are certificates for
our different V29 root. YM V79 retains the complete orientation contact
and proves a spatial floor for a specified covariance coefficient; its
original-law nonlinear remainder remains unpaid
\cite[\src{r79:unpaid}]{v38:YM}. This supports retaining the complete moving
source, not transferring that positive floor to a canonical pencil.
SM V43 proves an elapsed-time/count identifiability theorem under an
independent nondegenerate renewal clock
\cite[\src{sel43:thm:clock}]{v38:SM}; it provides no license to erase the
physical time or a duration-resolved response here. Perturbative ESM V51
retains all descendants in its field-dependent current completion
\cite[\src{thm:v51-BV}, \src{eq:v51-descendants}]{v38:ESM}; its polynomial
Noether section is not a cancellation of this finite action's moving
constraint source. These are scope-checked methodological transfers, not
additional hypotheses on the gravitational law.

\paragraph{Inherited inputs.}
We retain V29's mixed-profile root, its canonical first field $Y_*$ and
normalized first tilt $\ell_*$; V30's signed stationary chart and analytic
small-rate initial preparations; V37's simple real slow root and
constraint-tangent lift; and the original ordinary-interpolation curvature
observation of V10. The physical rate formula in V34 and the exact
radial normal in V37 are used, not claimed again as new algebraic
constructions. Background on simultaneous discrete constraints and multiplier
determination is supplied by the already cited \cite{v22:DBGP}; it
supplies none of the source coefficients below. All action-specific
conclusions in this body concern $N=3$, $h=1/3$.

\section{The coefficient needed by the moving constraint graph}
\label{v38:coefficient}

Use the fixed mean-lapse-zero multiplier complement. Let
\[
 C=H_{uX},\qquad M=H_{uu},\qquad Q=H_{XX},\qquad
 F=\mathbb J\nabla_XH,\qquad r=\lambda_t,
\]
where $u$ denotes these normalized multiplier coordinates. The original
canonical linearization after their implicit elimination is
$T_a=\mathbb J(Q-C^*M^{-1}C)$. A dot on $C$ or $M$ below is the total
\emph{physical-time} derivative along the base exact history. It is not a
partial derivative holding its canonical source fixed.

Fix any analytic initial family supplied by V30 with
\begin{equation}
 X_a=aY_*+O(a^2),\qquad \lambda_a=e+a\ell_*+O(a^2),\qquad
 r_a=O(a),\qquad P(X_a,\lambda_a)=0.
 \label{v38:base}
\end{equation}
The transport-selected family is one such choice. Write
\[
 V_*=F_1Y_*+G\ell_*,\qquad G=\mathbb J L^*,\qquad
 X_t(a,0)=aV_*+O(a^2).
\]
The estimates below are local uniformly when the finite coefficients of
these analytic families range in a bounded admitted neighborhood. No
cutoff-uniform norm is meant by this statement.

For the original homogeneous quadratic constraint covector $P_2$, define
its symmetric bilinear second derivative by
\begin{equation}
 \mathsf B_2[V,W][\nu]
 :=D_X^2P_{\nu,2}(0)[V,W]
 =\frac{P_{\nu,2}(V+W)-P_{\nu,2}(V-W)}2.
 \label{v38:B2}
\end{equation}
The argument $0$ here specifies the constant Hessian of the homogeneous
quadratic polynomial, not omission of its field dependence. All
polarizations use its complete stationary envelope, with original literal
shifts and the electric contribution where the multiplier test permits it.
This notation is different from the quadratic \emph{drift}
$\widehat B_2(Y)$.

\begin{lemma}[First moving constraint derivative and differentiated Hessian orders]
\label{v38:moving-orders}
On \eqref{v38:base}, as a row map on canonical variations,
\begin{equation}
 C=L+a\,DP_2(Y_*)+O(a^2),\qquad
 \boxed{\dot C=a\,\mathsf B_2[V_*,\cdot]+O(a^2).}
 \label{v38:Cdot}
\end{equation}
In the physical active/weak decomposition, where the weak sector includes
all three constant shifts,
\begin{equation}
 M=\begin{pmatrix}
 a^2 A(a)&a^3 L_s(a)\\a^3L_s(a)^*&a^4 C_s(a)
 \end{pmatrix},\qquad
 \dot M=\begin{pmatrix}O(a^2)&O(a^3)\\O(a^3)&O(a^4)\end{pmatrix}.
 \label{v38:Mdot}
\end{equation}
The matrices have the original limiting signed Gram, with $A(0)$
invertible and weak Schur complement
$D_0=C_s(0)-L_s(0)^*A(0)^{-1}L_s(0)$ invertible. The leading term in
\eqref{v38:Cdot} is independent of every higher initial coefficient and of
the bounded leading coefficient of $r_a/a$.
\end{lemma}
\begin{proof}
V20's flat identities give $P(0,\lambda)=0$ and
$D_XP(0,\lambda)=L$ for all admitted multipliers near $e$. Hence the
first correction to $C$ on $X=aY_*+O(a^2)$, $\lambda=e+O(a)$ is
$aDP_2(Y_*)$. Differentiating along the base history gives
$D_XC[X_t]=a\mathsf B_2[V_*,\cdot]+O(a^2)$. Its other contribution is
$D_\lambda C[r]$. Since $D_\lambda C=D_XM=O(a)$ and $r=O(a)$,
that term is $O(a^2)$, not a second term at the displayed leading order.
This proves \eqref{v38:Cdot} with the complete stationary Hamiltonian.

The active reduced stationary-source columns start at order $a$ and its
weak columns at order $a^2$, by the universal weak first-source identity
and V20's graded-source result. Thus its exact signed factorization gives
\eqref{v38:Mdot}'s first formula. The same divided coefficient maps are
smooth in $(X/a,(\lambda-e)/a)$ on the fixed first-field neighborhood.
Both normalized velocities are bounded on \eqref{v38:base}. Differentiating
those maps without differentiating the fixed family parameter $a$ gives
exactly the stated orders for $\dot M$. This is not differentiation of a
frozen matrix with its state held fixed. Its leading metric and Schur
invertibility are the inherited V30 certificates at this root.
\end{proof}

\section{An explicit leading actual rate on the admissible slow representatives}
\label{v38:rate}

Let $\gamma>0$ be V37's simple slow root and let
$v=(v_R,c)$ span $\ker\mathcal P_0(\gamma)$. Normalize it by
\begin{equation}
                  |c|=1,\qquad c_1>0.
 \label{v38:normalization}
\end{equation}
The nonzero first component needed for this convention is verified below.
Here $c$ synthesizes a genuinely constant shift, so its Euclidean norm is
also its physical spatial-mean norm. The nonconstant multiplier $v_R$ is
synthesized in the original active-plus-gradient coordinates; it is not
identified with raw point coordinates without that synthesis.
Define
\begin{equation}
 \Xi_*:=\gamma^{-1}(Gv_R+\delta_cY_*),\qquad
 h_W[\nu]:=\mathsf B_2[V_*,\Xi_*][\nu]\quad(\nu\in\mathcal W),
 \label{v38:hW}
\end{equation}
where $\mathcal W=\ker\curl_\delta$ includes constants. If $I_W$
synthesizes physical weak multipliers and $I_H$ its constant-shift subspace,
put
\begin{equation}
 \boxed{k_*:=\gamma I_Hc-I_WD_0^{-1}h_W.}
 \label{v38:k}
\end{equation}
All terms in this formula are now specified by original coefficients.
The right inverse is the original signed weak inverse, not a positive
replacement. Factors changing spatial sums to spatial means multiply
$D_0$ and $h_W$ together and cancel from their solve.

\begin{theorem}[The moving-source term determines the nonconstant leading slow rate]
\label{v38:actual-rate}
Let $(\widetilde\xi_a,\widetilde\eta_a)$ be the exactly constraint-tangent
slow representatives of \cref{v37:slow-tangent}. Realize them as derivatives,
with respect to a parameter $\varepsilon$, of local curves of original
exactly constrained initial records. For their own exact-action initial
multiplier rates,
\begin{equation}
 \boxed{a\,\partial_\varepsilon r(a,0)\longrightarrow k_*,\qquad
 \partial_\varepsilon r_A=O(1),\qquad
 \partial_\varepsilon r_W=O(a^{-1}).}
 \label{v38:rate-limit}
\end{equation}
More precisely $\partial_\varepsilon r=a^{-1}k_*+O(1)$ in the weak
sector. The vector $k_*$ has zero lapse component and is weak. Its
nonconstant part comes entirely from $-I_WD_0^{-1}h_W$, not from the
frozen spectral term. The limit is unchanged by the higher preparation
choices in \eqref{v38:base} and by its V37 radial correction.
\end{theorem}
\begin{proof}
First use the uncorrected exact canonical eigenvector and its multiplier
lift. They satisfy
$T_a\xi_a=z_a\xi_a$ and $C\xi_a+M\eta_a=0$, with
\[
 z_a=\gamma/a+O(1),\quad
 \xi_a=a^2\Xi_*+O(a^3),\quad
 \eta_a=(a v_R+O(a^2),\ c+O(a))
\]
in the nonconstant/constant order. Differentiate the original linearized
constraint along the exact base history. Since the instantaneous canonical
variation has derivative $T_a\xi_a$, this gives the exact identity
\begin{equation}
 \partial_\varepsilon r
 =z_a\eta_a-M^{-1}(\dot C\xi_a+\dot M\eta_a).
 \label{v38:exact-rate}
\end{equation}
This identity already appears in V37; the terms on its right are evaluated
here at their correct orders.

Equation \eqref{v38:Cdot} gives
$\dot C\xi_a=a^3h+O(a^4)$, with
$h=\mathsf B_2[V_*,\Xi_*]$. In active/weak order,
$\dot M\eta_a=(O(a^3),O(a^4))$. Let
$D_a=\diag(aI_A,a^2I_W)$, so $M=D_a\Gamma(a)D_a$. Exact block
inversion, or elimination of the active row, gives
\begin{equation}
 aM^{-1}(a^3h+O(a^4))
       =I_WD_0^{-1}h_W+O(a),\qquad
 aM^{-1}(O(a^3),O(a^4))=O(a).
 \label{v38:scaled-inverse}
\end{equation}
In particular the active part of $M^{-1}\dot C\xi_a$ is at most
$O(1)$, while its weak part can be $O(a^{-1})$. Multiplying
\eqref{v38:exact-rate} by $a$ now yields \eqref{v38:k} and
\eqref{v38:rate-limit}. The part $a z_a\eta_a$ tends only to
$\gamma I_Hc$, a constant shift.

For completeness retain the radial correction exactly. Put
$\alpha_a=\langle g_a,\xi_a\rangle/a=O(a^3)$, so
$\widetilde\xi_a=\xi_a-\alpha_a W_a$ and
$\widetilde\eta_a=\eta_a$. Its instantaneous canonical residual is
$e_a=(T_a-z_aI)\widetilde\xi_a=-\alpha_a(T_a-z_aI)W_a$.
By V37, $CW_a=0$, $T_aW_a=O(1)$. Consequently
$Ce_a=-\alpha_a C T_aW_a=O(a^3)$ with leading flat contribution
$-\alpha_a L F_0W_0=0$: the flat subsidiary identity and $LW_0=0$
annihilate it. Hence $Ce_a=O(a^4)$. Alternatively
$CT_aW_a=C\mathbb JQW_a=O(a)$ by these same identities.
The corrected version of \eqref{v38:exact-rate} has the additional term
$-M^{-1}Ce_a=O(1)$ in its weak part. Also
$\dot C(\widetilde\xi_a-\xi_a)=O(a^4)$. Both contributions disappear
after multiplication by $a$. This proves independence of that correction,
without treating it as an exact eigenvector.

The higher initial coefficients enter neither $V_*$ nor $\Xi_*$,
and \eqref{v38:Cdot}'s leading coefficient does not use the coefficient
of $r_a/a$. This proves the final assertion. Realization by actual initial
curves and local original histories is inherited from V37 and the open
stationary rank chart; no uniform interval for the varied histories is
inferred.
\end{proof}

The small term $\dot M\eta_a$ is retained in \eqref{v38:exact-rate}; it
is bounded only after its stronger active/weak orders are used. Replacing
both moving terms by zero would miss the entire leading nonconstant rate.
Conversely a leading nonconstant rate is not yet a time-dependent
exponential mode. Equation \eqref{v38:rate-limit} is an initial-data
differential.

\section{A certificate for the complete moving weak covector}
\label{v38:certificate}

The new computation encloses the exact root and exact eigenvector, not only
a rounded field or a plotted eigenmode. The detailed data are supplied in
\src{moving_source_certificate.json}. This section specifies the acceptance
argument and normalization so that the nonvanishing is independently
reproducible.

\begin{lemma}[Right-eigenvector inclusion from the inherited isolated reciprocal root]
\label{v38:vector-inclusion}
Let the rational similarity in V37 put its reciprocal quotient in the form
$B=B_0+E$, with $B_0$ real normal and $\|E\|_2\le\delta<10^{-8}$.
Let $\mu_0$ be its positive real reference eigenvalue. Its distance to every
other reference eigenvalue exceeds $6$, as checked by rational squared
distances. The true simple real eigenvector can be normalized to have
selected coordinate one and all other coordinates $w$ satisfying
\begin{equation}
                 \boxed{\|w\|_2\le\frac{\delta}{6-2\delta}.}
 \label{v38:eigenvector-radius}
\end{equation}
The original descriptor lift and normalization \eqref{v38:normalization}
then have enclosing integer intervals on the whole V29 root cube.
\end{lemma}
\begin{proof}
V37's isolation proof puts the true reciprocal eigenvalue $\mu$ in
$|\mu-\mu_0|\le\delta$. Remove its selected scalar coordinate.
On the complementary coordinates, normality of $B_0$ gives inverse norm
at most $(6-\delta)^{-1}$ for $B_{0,cc}-\mu I$. A Neumann estimate
for the additional $E_{cc}$ gives $(6-2\delta)^{-1}$. Thus the selected
coordinate of an eigenvector cannot be zero. Normalize it to one and
solve its other rows:
$w=-(B_{cc}-\mu I)^{-1}B_{cj}$. Since $B_{0,cj}=0$,
$\|B_{cj}\|_2\le\delta$, proving \eqref{v38:eigenvector-radius}.

To retain the descriptor's exact eliminated kernel, let $X$ be V37's
rational coordinate matrix, $S_X=X^*\mathbf SX$,
$E_X=X^*\mathbf EX$, and $C_X=-S_X^{-1}E_X$. The first 84 coordinates
are its supported quotient and the remaining 26 its exact kernel. If $u$
is the enclosed quotient eigenvector, the full vector is
$(u,\mu^{-1}(C_X)_{ko}u)$. Multiplication by $X$ restores the original
110 descriptor coordinates; their first 107 entries give $v$.
The inherited verified descriptor inverse encloses this lift. The first
constant component has a strict sign, and the squared norm of the three
constant components is separated from zero. Integer square-root bounds
therefore enclose division by their Euclidean norm. This gives the stated
normalization without assuming a numerical eigenvector is exact.
\end{proof}

\begin{theorem}[Nonzero actual slow rate and metric coefficient at the original root]
\label{v38:nonzero}
At the exact V29 root with convention \eqref{v38:normalization}, the
original coefficients satisfy
\begin{equation}
 \boxed{
 4000<\|k_*^\beta-\overline{k_*^\beta}\|_h<4600,\qquad
 6300<\|\mathfrak E_h(k_*^\beta)\|_h<7000.
 }
 \label{v38:numerical-bounds}
\end{equation}
The tensor norm counts all nine matrix entries with the original
Frobenius pairing and normalized spatial mean. These constants belong to
the inherited three-site unit-torus normalization; they are not new
physically calibrated predictions.
\end{theorem}
\begin{proof}
The root input is the original dyadic 107-vector of V29 and its cube of
radius $2^{-240}$. Extract its three profile coordinates, reconstruct $Y_*$
from the same five-field polynomial bank, and recover $\ell_*$ from its
104 point representatives by subtracting all four component means and
applying the physical $\sqrt3$ shift synthesis. The first canonical
velocity is evaluated as $V_*=(V(p),\mathcal F(U))+G\ell_*$, using
the original V5 linear canonical equations.

Use \cref{v38:vector-inclusion} to enclose $v_R,c,\gamma$, then form
$\Xi_*$ by \eqref{v38:hW}. The vector source map \src{p2vec} is the
retained complete original quadratic row, defined by \eqref{v20:P2}.
The exact homogeneous identity \eqref{v38:B2} produces its full moving
covector, and the physical weak synthesis gives all 29 entries of $h_W$.
The source is not restricted to its constant or gradient part before this
contraction. An independent direct scalar-envelope evaluation of the same
quadratic polynomial, without its adjoint source assembly, checks all 108
rows at each of two specified rational field representatives. These
regressions check the implementation; the enclosing evaluation itself is
over the entire root and eigenvector intervals.

The full leading signed Gram is formed from the inherited original higher
source enclosure. Its active inverse and weak Schur inverse are certified
with proposed rational inverse matrices using the exact residual identity
$E=I-RA$ and
\[
 A^{-1}=(I+E+E^2)R+E^3(I-E)^{-1}R.
\]
The accepted active residual is below $10^{-10}$ and the accepted weak
Schur residual below $10^{-6}$. Their tails are bounded by
$\|E\|^3\|R\|/(1-\|E\|)$; both inverses therefore concern the
actual matrices on the entire enclosure. All operations after proposing
these inverse matrices use Python integers at 128 binary places with outward
rounding, and arbitrary-precision rational comparisons. The slow-reference
error and the eigenvector inclusion use the separately rerun V37 spectral
acceptance, not a newly assumed condition number.

Synthesize $k_*^\beta$ from \eqref{v38:k}, subtract its exact spatial
mean, and enclose its squared norm entrywise. At $N=3$ the ordinary
derivative of the trigonometric interpolant is
$\partial_i^{\rm sp}=2\pi\delta_i/(3\sqrt3)$ on every retained
frequency. Thus the original V10 observation is evaluated by the exact
finite expression
\begin{equation}
 \mathfrak E_{h,ij}(b)
 =\frac12\left(\frac{2\pi}{3\sqrt3}-1\right)
                  (\delta_i b_j+\delta_j b_i).
 \label{v38:observation}
\end{equation}
The square-root bounds use integer square roots. Bounds for $\pi$ use
Machin's identity with rational alternating-series remainder bounds; no
binary floating value is used to accept the trigonometric factor.
The resulting norm intervals lie inside
$(4026,4544)$ and $(6355,6956)$, respectively. Each endpoint in the
supplied report is a rational fraction enclosing the corresponding norm.
These strict intervals imply \eqref{v38:numerical-bounds}.

The source, synthesis, inverse order and norms in this proof retain their
original common normalization. A spatial-sum matrix implementation gives
the same solved vector because its source covector and Gram share the same
factor. The final norm divides its squared spatial sum by 27. This completes
the finite nonzero certificate, without a finite-amplitude difference quotient
or nonlinear trajectory simulation.
\end{proof}

\section{The original metric curvature sees the moving-source response}
\label{v38:metric}

Let $\mathcal R_a(\varepsilon)$ denote the tensor
$(R_{0i0j}(g_{a,\varepsilon})(0))_{ij}$ of the original interpolated ADM
metric on the histories from \cref{v38:actual-rate}. The derivative below
is evaluated at $\varepsilon=0$. The same anchored coordinates and metric
construction as V10 and V37 are retained.

\begin{theorem}[Slow geometric visibility after the moving terms are included]
\label{v38:curvature}
On the exactly constraint-tangent slow representatives,
\begin{equation}
 \boxed{
 a\,\partial_\varepsilon\mathcal R_a
 \longrightarrow
 \mathfrak E_h(k_*^\beta)
 =-\mathfrak E_h\bigl((I_WD_0^{-1}h_W)^\beta\bigr)\ne0.
 }
 \label{v38:curvature-limit}
\end{equation}
In particular the derivative has size comparable to $a^{-1}$, with fixed
positive lower and upper constants. The leading coefficient is entirely
from the moving constraint source, since $\mathfrak E_h$ annihilates the
constant frozen term $\gamma I_Hc$.
\end{theorem}
\begin{proof}
The background has $X,\lambda-e,X_t,\lambda_t,X_{tt}=O(a)$, with
bounded fixed-chart spatial derivatives. On the selected tangent vectors,
\[
 \partial_\varepsilon X=O(a^2),\quad
 \partial_\varepsilon\lambda=I_Hc+O(a),\quad
 \partial_\varepsilon X_t=O(a).
\]
The last assertion follows from the original smooth canonical equation:
$D_XF$ is bounded, while $D_\lambda F\,I_Hc=O(a)$ because a
constant multiplier has zero flat canonical vector. All spatial derivatives
of the leading constant shift vanish; derivatives of its nonconstant
remainder are $O(a)$ at this fixed grid.

Differentiate the original identity
$X_{tt}=D_XF\,F+D_\lambda F\,r$. Every second derivative of $F$
is bounded in the original stationary chart. Terms differentiating its
coefficients are therefore $O(a)$ or bounded here. The only possible
$a^{-1}$ contribution is $D_\lambda F\,\partial_\varepsilon r$.
Since $D_\lambda F=G+O(a)$, \cref{v38:actual-rate} gives
\[
 a\,\partial_\varepsilon X_{tt}\longrightarrow Gk_*.
\]
Its configuration part is
$\delta_i k_{*,j}^\beta+\delta_j k_{*,i}^\beta$.

Write the original interpolated metric as
$g_{ij}=\gamma_{ij}$, $g_{0i}=\gamma_{ij}\beta^j$ and
$g_{00}=-N^2+\gamma_{ij}\beta^i\beta^j$. Then
\[
 a\,\partial_\varepsilon(g_{0i,t})\longrightarrow k_{*,i}^\beta,
 \qquad
 a\,\partial_\varepsilon(g_{ij,tt})\longrightarrow
               \delta_i k_{*,j}^\beta+\delta_j k_{*,i}^\beta.
\]
The spatial Hessian of $\partial_\varepsilon g_{00}$ contributes zero
in the scaled limit. The background first Christoffel symbols are $O(a)$;
first variations of those symbols are at most $O(a^{-1})$. Their product
variations are therefore bounded and disappear after multiplication by $a$.
Variations of the inverse metric are bounded and also give no leading term.
The exact coordinate identity
\[
 R_{0i0j}=\tfrac12(g_{0j,it}+g_{i0,tj}-g_{00,ij}-g_{ij,tt})
                          +\text{Christoffel products}
\]
now gives the difference between the ordinary spatial derivatives of the
interpolant and the original phase derivatives, exactly
$\mathfrak E_h(k_*^\beta)$. No $\lambda_{tt}$ estimate is used.
The other independent curvature blocks have zero limit after multiplying
their parameter derivatives by $a$; they contain no second physical-time
metric derivative, and their Christoffel terms obey the same bound.
Finally apply \cref{v38:nonzero}. This proves the two-sided size and
\eqref{v38:curvature-limit}.
\end{proof}

The tensor observation is the original anchored metric read. No claim is
made here about a nonzero differential of a scalar curvature invariant;
that would additionally require its contraction with the actual background
curvature coefficient. A parameter-dependent coordinate change is not silently
used to change this read. In particular a frozen eigenvalue and a constant
leading shift, by themselves, determine neither its vanishing nor its size.

\section{The literal curvature differential has a different scale}
\label{v38:links}

The preceding calculation also has a useful physical-accounting comparison.
Both the metric and literal link curvatures of each base history are $O(a)$.
Their \emph{differentials} on the present admissible vectors have different
orders, which must not be confused with their base values.

\begin{theorem}[Metric and link differentials cannot be uniformly identified on this chart]
\label{v38:link-comparison}
On the same exactly constraint-tangent slow representatives,
\begin{equation}
 \boxed{
 \|\partial_\varepsilon E_h^{\rm link}(a,0)\|
 +\|\partial_\varepsilon F_h^{{\rm sp},{\rm link}}(a,0)\|=O(a).
 }
 \label{v38:link-bound}
\end{equation}
Thus, in the original compatible frames,
\begin{equation}
 a\,\partial_\varepsilon
 \left(\mathcal R_a-\mathcal R_a^{\rm link}\right)
 \longrightarrow\mathfrak E_h(k_*^\beta)\ne0,
 \label{v38:difference}
\end{equation}
where $\mathcal R_a^{\rm link}$ denotes the corresponding retained literal
curvature components. In particular there is no amplitude-independent
bound transferring the norm of the literal-curvature differential alone
to that of the metric-curvature differential on all these initial directions.
\end{theorem}
\begin{proof}
Let $z(X,\lambda)$ be the original stationary spatial connection and
$A_0(X,\lambda)$ the temporal auxiliary connection. These maps are analytic
on the inherited flat stationary chart, before eliminating the singular
multiplier consistency equation. Their fixed derivatives are bounded there.
The original stationary response formula gives, at the base points,
\[
 D_{\lambda_A}z=O(a),\qquad D_{\lambda_W}z=O(a^2),
\]
because its reduced source columns have precisely those orders. Therefore
$\partial_\varepsilon z=O(a^2)$ on
$\widetilde\xi_a=O(a^2)$ and
$\widetilde\eta_a=I_Hc+O(a)$.

At the flat point the temporal auxiliary response is
$(-\curl_\delta\beta/2,\nabla_\delta N)$, so its derivative in a
constant shift is exactly zero. Analyticity then gives
$\partial_\varepsilon A_0=O(a)$ on these directions. This checks the
constant mode explicitly; a generic bounded auxiliary derivative alone
would not give that order.

The spatial-connection velocity is
$z_t=D_Xz\,F+D_\lambda z\,r$. In its parameter derivative all terms
which differentiate these smooth coefficients are $O(a)$ because
$F,r=O(a)$, $\widetilde\xi_a=O(a^2)$ and
$\widetilde\eta_a=O(1)$. Also $D_Xz\,\partial_\varepsilon F=O(a)$.
The potentially singular term is controlled in its original graded order:
\[
 D_\lambda z\,\partial_\varepsilon r
 =O(a)\,O(1)+O(a^2)\,O(a^{-1})=O(a).
\]
Thus $\partial_\varepsilon z_t=O(a)$. In particular the original electric
formula retained in \cref{v3:physical-pole} is
\[
 E_i=h^{-1}\{\dot U_iU_i^{-1}+A_0-U_iA_{0,+i}U_i^{-1}\},
 \qquad U_i=\exp(hz_i).
\]
It uses $z_t$ and shifted $A_0$, not a time derivative of $A_0$.
The spatial curvature is the original analytic ordered link word. At fixed
$h$, the foregoing derivative bounds therefore give \eqref{v38:link-bound}.
A change to the original coframe components has bounded variation; its
product with the $O(a)$ base curvature has order $a$ and does not change
this conclusion.

Multiply by $a$ and combine with \cref{v38:curvature} to obtain
\eqref{v38:difference}. If an amplitude-independent transfer bound from
literal to metric differentials held on this family, its right side would
be $O(a)$ by \eqref{v38:link-bound}, while its left side is bounded below
by a positive multiple of $a^{-1}$. This contradiction proves the last
assertion. No finite parameter increment is estimated by this differential
argument without an additional uniform remainder and parameter-radius bound.
\end{proof}

This result is not a claim that the selected base histories have divergent
curvature, nor that an Einstein limit under spatial refinement fails. It
locates a lack of uniformity in a particular differential comparison at the
fixed cutoff. A more restricted selected manifold might exclude these
tangent directions; the theorem does not identify such a manifold with the
whole original constraint set.

\section{What has changed in the next dynamical problem}
\label{v38:next}

The slow positive pair is no longer an eigenvalue with an unevaluated
physical source. At the original root its constraint-tangent initial
differential has the explicit leading rate \eqref{v38:k} and the nonzero
metric read \eqref{v38:curvature-limit}. The constant frozen term contributes
nothing to that read. The complete effect is instead generated by the
bilinear quadratic-constraint coefficient and the retained signed weak
inverse. Neither a fourth-order envelope nor an unlisted initial-jet choice
is a prerequisite for this calculation.

The next missing implication is still time dependent. A physical variational
history satisfies the original coupled equations with moving $C,M,Q$.
Its vector need not remain a frozen slow eigenvector or the radial-corrected
initial representative. Establishing an exponential law, suppressing this
response by selecting an invariant family, or extending V33's physical
interval requires transport of those actual moving sources and of the
selected tangent space. The unrelated invariant graph in the newly supplied
native-stationarity lineage does not close that requirement at our root.
Nor can the coefficientwise current subtraction in the ESM file be treated
as a modification of this finite gravitational action.

V33's $ca^2$ physical interval remains the longest established one in this
bridge lineage. This body adds no new lifespan exponent. Every bound remains
at the three-site cutoff, with the original physical clock and all original
constraints retained. No spatial-refinement-uniform estimate, scalar-invariant
instability, uniform finite-perturbation radius, equality of normalized
metric/link observations, or exact-action Einstein continuum limit follows
from the new differential certificate.

\section{Verification and the append-only record}
\label{v38:verification}

The supplied driver reruns the V37 original cubic-source enclosure and its
complete slow spectral acceptance, then computes the new eigenvector inclusion,
all 29 moving weak-source entries, both required signed block inverses, and
the physical curvature norm. Its final acceptance uses outward integer
intervals and rational inequalities. The preconditioners are finite rational
matrices proposed by floating arithmetic and accepted only by their exact
residual estimates. The inherited V29 root theorem and V30 higher-source
enclosure remain identified inputs; their original complete regeneration
chains are not represented as rerun here.

A separate exact checker evaluates the original scalar $P_2$ envelope
forward on all 108 tests at two rational coefficient fields, and compares
all 216 values with the adjoint vector implementation. It checks the Hessian
polarization independently. Rational small-matrix tests retain noncommuting
graded blocks and every differentiated Hessian term in a separate analytic
Hamiltonian. Those matrices test algebra, not the gravitational source or a
nonlinear spacetime. The written asymptotic proofs, including the exact
radial residual estimate, are not proved by sampling and are not a
proof-assistant formalization.

The source audit records the names, hashes, terminal ranges and transfer
boundaries of all four new companion files and the V37 predecessor.
All thirty-seven preceding marked mathematical bodies are retained byte for
byte, including their corrections and scope statements. Only the cumulative
title and abstract outside those bodies are updated. The package contains
this appended body, the predecessor, the portable assembler, proof inputs,
verification code and the reports of checks actually executed.

\begin{thebibliography}{Native46}
\small\raggedright
\bibitem[Native46]{v38:Native}
\emph{Renewal Native Regenerative Stationarity Closure Notes V46},
\src{Renewal_Native_Regenerative_Stationarity_Closure_Notes_V46(2).tex}.
\bibitem[YM79]{v38:YM}
\emph{Renewal Yang--Mills Mass-Gap Research Notes V79},
\src{Renewal_Yang_Mills_Mass_Gap_Research_Notes_V79(2).tex}.
\bibitem[SM43]{v38:SM}
\emph{Renewal SM Source Selection Research Notes V43},
\src{Renewal_SM_Source_Selection_Research_Notes_V43(3).tex}.
\bibitem[ESM51]{v38:ESM}
\emph{Renewal Perturbative ESM Continuum Research Notes V51},
\src{Renewal_Perturbative_ESM_Continuum_Research_Notes_V51(2).tex}.
\end{thebibliography}
% END IMMUTABLE EXACT-ACTION BRIDGE V38 BODY


% Future iterations append here.
% BEGIN IMMUTABLE EXACT-ACTION BRIDGE V39 BODY
\clearpage
\renewcommand{\thesection}{\arabic{section}}
\renewcommand{\theHsection}{bridgeV39.\arabic{section}}
\setcounter{section}{314}
\section{A finite perturbation chart rather than another differential test}
\label{v39:context}

The differential in \cref{v38:curvature} is not, by itself, a finite
perturbation theorem. Its realizing curves were obtained at each positive
amplitude from a multiplier graph whose inverse degenerates at the flat tip.
A uniform radius or a uniform Taylor remainder cannot be inferred from that
construction. We instead solve the original constraints for canonical field
coordinates. The row-scaled canonical submersion used in
\cref{v37:lossless-normal} stays regular as the amplitude vanishes. It provides
an actual two-parameter family realizing precisely the V38 slow tangent, with
a uniform radius in the scaled perturbation parameter.

A second issue is essential. On that family, a bounded normalized parameter
changes the first multiplier tilt by a constant shift. A generic initial
curve would then still permit an inverse-amplitude weak rate. We prove an
affine cubic-drift identity for this particular constant-tilt line and combine
it with the already evaluated tangent. This removes that pole on the entire
new finite parameter interval, not only at its centre. The actual rate and
original curvature therefore have analytic limiting profiles. The derivative
of the latter is the nonzero coefficient already certified in V38. Its
nonvanishing now gives a finite secant estimate and an explicit preparation
precision scale.

\paragraph{Scope and nonduplication.}
All action-specific statements below use $N=3$, $h=1/3$. The small variable
$a>0$ is field amplitude, not spatial mesh. The auxiliary parameter is denoted
$q$, and the original tangent parameter is $\varepsilon=aq$; neither is
physical time. The action, normalized mean lapse, stationary equations,
canonical flow, and literal link observations are unchanged. We retain V30's
analytic small-rate base preparation, V37's analytic constraint-tangent slow
representatives and field submersion, and V38's actual moving-rate coefficient
and its nonzero certificate. Their matrix extractions are not new here.
V10 already exhibited a fixed-cutoff metric/link mismatch on a different
family. The new result is a uniform nonlinear completion of the \emph{specific}
V38 direction through the presently selected small-rate base, with controlled
finite parameter differences. No new frozen spectral count is needed.

The standard distinction between initial constraints and their own
continuation is retained \cite{v39:GN}. No theorem from that reference supplies
the coefficient cancellations or amplitude-uniform parameter estimates below.

\section{The constant-tilt cancellation needed before a finite displacement}
\label{v39:cubic}

Write $H(X,\lambda)$ for the original stationary Hamiltonian, $P=H_\lambda$,
$M=H_{\lambda\lambda}$, $F=\mathbb JH_X$, and
$\mathsf B=D_XP\,F$. All vector-valued multiplier formulas in this section
initially retain the complete 108-row covector. Let $L=D_XP(0,e)$ and
$G=\mathbb JL^*$, $e=(1,0)$. Let $M_2(X)$ be the homogeneous quadratic
\emph{field} coefficient of $M(X,e)$. For a constant shift $c$, use
$c_H=(0,c)$ when needed to distinguish the multiplier from its three-vector.
The weak sector $W$ includes every curl-free shift, including constants.

\begin{lemma}[A constant direction annihilates the first multiplier curvature]
\label{v39:constant-algebra}
For arbitrary symmetric finite fields and any constant shift,
\begin{equation}
 Gc_H=0,\qquad M_2(X)c_H=0,\qquad
 D_X(M_2c_H)=0.                                      \label{v39:flat-null}
\end{equation}
Moreover $\mathsf B(0,\lambda)=0$ for every admitted $\lambda$, and the
homogeneous degree-one field coefficient of
$D_\lambda^2\mathsf B(X,e)[c_H,\mu]$ is zero for every multiplier $\mu$.
\end{lemma}
\begin{proof}
A constant shift has zero flat phase gradient, so $Gc_H=0$.
The retained leading multiplier Hessian is
$M_2(X)=J_1(X)^*K_0^{-1}J_1(X)$. The universal weak first-source identity
of \cref{v3:universal-flat} gives $J_1(X)c_H=0$ for every symmetric $X$.
This proves the other identities in \eqref{v39:flat-null}, including their
field derivatives. No sign of the higher weak Gram is changed.

V20's flat identities give $D_XP(0,\lambda)=L$ and
$F(0,\lambda)=G(\lambda-e)$. Hence $\mathsf B(0,\lambda)=0$ by
$L\mathbb JL^*=0$. In particular every multiplier derivative of this
zero-field value vanishes.

Use the exact original identity
\[
 D_\lambda\mathsf B[\mu]=\mathbb K\mu+D_X(M\mu)[F],
 \qquad\mathbb K=D_XP\,\mathbb J D_XP^*.
\]
Put $C=D_XP$. Its differentiated Poisson part is
\[
 D_\lambda\mathbb K[c_H]
 =D_X(Mc_H)\mathbb JC^*+C\mathbb J D_X(Mc_H)^*.
\]
Here $D_X(Mc_H)=O(\|X\|^2)$ by \eqref{v39:flat-null}; therefore this
term has no degree-one field coefficient. In the other term the two
contributions are
\[
 D_X\{D_\lambda M[c_H]\mu\}[F]
       +D_X(M\mu)[D_\lambda F[c_H]].
\]
At $\lambda=e$, $F=O(\|X\|)$,
$D_XD_\lambda M=O(\|X\|)$ by the flat germ, and
$D_\lambda F[c_H]=\mathbb JC^*c_H=Gc_H+O(\|X\|)=O(\|X\|)$.
Also $D_XM=O(\|X\|)$. Both products are $O(\|X\|^2)$.
This proves the claimed vanishing, with all differentiated terms retained.
Symmetry of the two multiplier derivatives gives the same conclusion when
the constant direction is the second argument.
\end{proof}

\begin{proposition}[The cubic weak coefficient is affine along the prescribed line]
\label{v39:affine-cubic}
Suppose analytic initial records, not necessarily yet constrained, have
\begin{equation}
 \begin{aligned}
 X(a,q)&=aY+a^2Y_2+a^3Y_3(q)+O(a^4),\\
 \lambda(a,q)&=e+a(\ell+qc_H)+a^2(m+q\eta_1)+O(a^3),
 \end{aligned}                                                   \label{v39:jets}
\end{equation}
where $Y,Y_2,\ell,m,c_H,\eta_1$ do not depend on $q$. Then the coefficient
$[a^3]\mathsf B_W(X(a,q),\lambda(a,q))$ is affine in $q$ and is
independent of $Y_3(q)$ and of every higher initial coefficient.
\end{proposition}
\begin{proof}
The un-tilted weak linear drift is identically zero on all fields, by
\cref{v5:prepared-B2}. Thus the cubic coefficient of
$\mathsf B_W(X(a,q),e)$ uses only $Y,Y_2$ and not $Y_3(q)$.
Expand in $u=\lambda-e$ at fixed $X$. Since
$\mathsf B(0,\lambda)=0$, terms with three multiplier increments
start at total order four. The linear multiplier term through order three
uses the degree-one and degree-two field coefficients of
$D_\lambda\mathsf B$, and is affine in the two coefficients of $u$ in
\eqref{v39:jets}. The quadratic multiplier term can contribute at order
three only through the degree-one field coefficient of
$D_\lambda^2\mathsf B$, applied twice to $\ell+qc_H$.
By \cref{v39:constant-algebra}, both its $q$ and $q^2$ contributions
vanish. Its remaining value is independent of $q$. No fourth or higher
Hamiltonian coefficient is being deleted: the constant cubic weak baseline
may contain such a coefficient, but is unchanged by this parameter line.
This proves affinity of the \emph{complete} cubic weak coefficient.
\end{proof}

\section{A uniform chart of exact original initial constraints}
\label{v39:chart}

Fix one analytic base family of \cref{v38:base}. Take V37's corrected slow
tangent with the V38 normalization and write
\begin{equation}
 \widetilde\xi_a=a^2\Xi(a),\qquad
 \widetilde\eta_a=c_H+a\eta_1+O(a^2),\qquad
 \Xi(0)=\Xi_* .                                      \label{v39:tangent}
\end{equation}
These divided functions can be chosen real analytic through $a=0$.
Indeed the original balanced spectral transfer is analytic, its real slow
root is simple, and its eigenvector is fixed by a nonsingular normalization.
The radial correction in \cref{v37:lossless-normal} is analytic with the
stated divisibility. The complete constraint tangent identity is
\begin{equation}
 C_{\rm full}(a)\widetilde\xi_a+
 M_{\rm full}(a)\widetilde\eta_a=0.                    \label{v39:tangent-full}
\end{equation}
For clarity, the two identities in \cref{v37:slow-tangent} imply this full
identity: the complementary rows vanish, and homogeneity
$\lambda_a^*M_{\rm full}=0$ together with the radial tangent condition
annihilates the remaining row. Its mean lapse coefficient is one.

Let $\Pi_{\rm n}$ select 104 independent nonconstant constraint rows and
let $\Pi_{\rm m}$ select their four means. The proven canonical submersion is
\begin{equation}
 \mathcal C_0=
 \begin{pmatrix}\Pi_{\rm n}L\\\Pi_{\rm m}DP_2(Y_*)\end{pmatrix}.
                                                               \label{v39:C0}
\end{equation}
It is onto. Choose a fixed 108-column canonical insertion $J_c$ for which
$\mathcal C_0J_c$ is invertible. This is exactly the submersion established
in \cref{v37:lossless-normal}; no multiplier inverse is used in this choice.

\begin{theorem}[Finite constraint-preserving completion on an amplitude-scaled radius]
\label{v39:constraint-chart}
There are $a_0,q_0>0$ and a real analytic vector $v(a,q)$ on a neighborhood
of $[0,a_0]\times[-q_0,q_0]$ such that
\begin{equation}
 \boxed{\begin{aligned}
 \lambda(a,q)&=\lambda_a+aq\widetilde\eta_a,\\
 X(a,q)&=X_a+aq\widetilde\xi_a+a^3J_cv(a,q),\\
 P(X(a,q),\lambda(a,q))&=0,\\
 v(a,0)&=0,\qquad \partial_qv(a,0)=0,\qquad
                         \|v(a,q)\|\le Cq^2.
 \end{aligned}}                                                  \label{v39:chart-eq}
\end{equation}
The mean lapse is exactly one. Every original mean constraint, including
all three shift means, is imposed. The \emph{values} of the shift means are
allowed to change, with
$\overline{\lambda^\beta(a,q)-\lambda_a^\beta}=aq(c+O(a))$.
Uniformly on this rectangle,
\begin{equation}
 \|X(a,q)-X_a\|\le Ca^3|q|,\qquad
 \|\lambda(a,q)-\lambda_a\|\asymp a|q|.               \label{v39:record-size}
\end{equation}
At each positive $a$ the derivative with respect to
$\varepsilon=aq$ at zero is precisely
$(\widetilde\xi_a,\widetilde\eta_a)$, not a different constraint tangent.
\end{theorem}
\begin{proof}
Insert the first two formulas of \eqref{v39:chart-eq}, with $v$ temporarily
free, into the complete constraint covector. Define its scaled residual
\begin{equation}
 \mathcal Q(a,q,v)=
 \begin{pmatrix}
 a^{-3}\Pi_{\rm n}P(X_a+aq\widetilde\xi_a+a^3J_cv,
                         \lambda_a+aq\widetilde\eta_a)\\
 a^{-4}\Pi_{\rm m}P(X_a+aq\widetilde\xi_a+a^3J_cv,
                         \lambda_a+aq\widetilde\eta_a)
 \end{pmatrix}.                                      \label{v39:divided-constraints}
\end{equation}
We verify removability of both divisions. The first two field coefficients
are unchanged. The only change in the first multiplier coefficient is
$qc_H$. The constraint Taylor series through degree three has the form
\[
 LX+P_2(X)+P_3(X)+M_2(X)(\lambda-e)
                 +\text{terms of joint degree at least four}.
\]
This follows from the complete flat germ: the field-linear constraint is
$LX$ for every multiplier, and $M=O(\|X\|^2)$. The extra degree-three
term $qM_2(Y_*)c_H$ is zero by \eqref{v39:flat-null}. The other changes
at that degree are $L(q\Xi_*+J_cv)$. They have zero spatial mean.
All smaller degrees equal those of the exactly constrained base. This
proves the divisibility in \eqref{v39:divided-constraints}. Its extension
is analytic, not just a bounded quotient.

The derivative in $v$ at $(0,0,0)$ is exactly
$\mathcal C_0J_c$. To check its lower block, a canonical change $a^3J_cv$
first affects the mean through $aDP_2(Y_*)a^3J_cv$; the flat mean row is
zero, and multiplier-dependent corrections cost another order. Hence no
mean row is missing from \eqref{v39:C0}. The analytic implicit-function
theorem supplies $v(a,q)$ on one parameter rectangle. This also follows
by a contraction with the fixed inverse $(\mathcal C_0J_c)^{-1}$ after
reducing that rectangle; the derivative in $v$ stays uniformly close to
its value at the origin.

The base solves $\mathcal Q(a,0,0)=0$, so uniqueness gives $v(a,0)=0$.
Equation \eqref{v39:tangent-full} makes the $q$ derivative of the
\emph{unscaled} residual zero at $(a,0,0)$ for every positive $a$.
Analytic extension gives the same statement at zero. Differentiating the
implicit equation therefore gives $\partial_qv(a,0)=0$ throughout.
Uniform boundedness of its second derivative on a smaller closed rectangle
proves $v(a,q)=O(q^2)$. Equation \eqref{v39:tangent} then proves the
first estimate in \eqref{v39:record-size}. The second uses
$\widetilde\eta_a\to c_H$ and $|c|=1$ in the physical mean norm.

The multiplier tangent has zero mean lapse, so that normalization is
unchanged exactly. Solving all scaled rows gives all 108 unscaled
constraint equations. Positivity of the spatial metric and lapse holds
on a smaller rectangle since $X=O(a)$ and $\lambda-e=O(a)$ uniformly.
Finally differentiating \eqref{v39:chart-eq} with respect to $aq$ proves
the claimed exact tangent realization. The interval in the \emph{original}
perturbation variable is $|\varepsilon|\le aq_0$, not a fixed unscaled
radius.
\end{proof}

\section{The original rate has a finite analytic nonlinear response profile}
\label{v39:rate}

The leading canonical field and the second canonical coefficient in
\eqref{v39:chart-eq} are independent of $q$. Its first multiplier tilt is
$\ell_*+qc_H$ and its second multiplier coefficient is affine in $q$.
The higher field corrections do not have to be tabulated.

\begin{theorem}[Removal of the potential inverse-amplitude pole on the whole chart]
\label{v39:rate-profile}
On the exact initial chart of \cref{v39:constraint-chart}, the complete
original drift obeys
\begin{equation}
          \mathsf B_A=a^3b_A(a,q),\qquad
          \mathsf B_W=a^4b_W(a,q),                   \label{v39:B-orders}
\end{equation}
with analytic divided coefficients on a common rectangle. The original
normalized rate and its limiting profile have the form
\begin{equation}
 \boxed{\begin{aligned}
 r_A(a,q)&=a\,r_{A,1}(a,q),\\
 r_W(a,q)&=\mathcal U(q)+a\,r_{W,1}(a,q),\\
 \mathcal U(0)&=0,\qquad
 I_W\mathcal U'(0)=k_* .
 \end{aligned}}                                                  \label{v39:rate-profile-eq}
\end{equation}
All displayed coefficient functions are analytic and uniformly bounded with
every fixed number of parameter derivatives on a smaller rectangle. The
physical vector $k_*$ is exactly \eqref{v38:k}. Every positive-amplitude
point has its own local normalized original-action history; no common
physical-time interval is asserted for these perturbed histories.
\end{theorem}
\begin{proof}
The complete quadratic prepared drift is
$\widehat B_2(Y_*)+\mathbb K_1(Y_*)(\ell_*+qc_H)$ by V20.
It is zero: the base tilt cancels it and the four constant multiplier
fields lie in the first Poisson kernel on compatible first fields, as
proved in V28. This removes every quadratic row, not only a weak test.
Thus the active estimate in \eqref{v39:B-orders} follows.

Let $b_{3,W}(q)=[a^3]\mathsf B_W$ on the exact chart. It is affine
by \cref{v39:affine-cubic}. At $q=0$ the actual base rate is $O(a)$.
The original graded Hessian gives
$\mathsf B_W=-M_{WA}r_A-M_{WW}r_W=O(a^4)$ there, so
$b_{3,W}(0)=0$.

To evaluate the affine slope retain the differentiated consistency equation
on the \emph{exact} parameter chart:
\[
 \partial_q\mathsf B=-(\partial_qM)r-M\partial_qr.
\]
Its tangent at zero is $a$ times the V38 tangent, so V38 proves
$\partial_qr_A(a,0)=O(a)$ and
$\partial_qr_W(a,0)=\mathcal U_0+O(a)$, with physical synthesis
$I_W\mathcal U_0=k_*$. The divided source maps are smooth on the
present normalized record chart; hence
$\partial_qM_{WA}=O(a^3)$ and $\partial_qM_{WW}=O(a^4)$ at zero.
Since $r(a,0)=O(a)$, the weak row of this differentiated identity is
$O(a^4)$. It follows that $b_{3,W}'(0)=0$. An affine function with zero
value and zero slope is identically zero. This proves the second estimate
in \eqref{v39:B-orders} on the \emph{whole} parameter interval. It is not
an assumption that a small differential stays small at a finite parameter.

By V20's graded-source identity the exact normalized Hessian is
$M=D_a\Gamma(a,q)D_a$, where
$D_a=\diag(aI_A,a^2I_W)$ and $\Gamma(0,q)=\Gamma_*$ is independent
of $q$. The first canonical field has not changed. V30's signed certificate
therefore supplies an analytic bounded inverse on a common rectangle.
From \eqref{v39:B-orders},
\begin{equation}
 \boxed{r=-\diag(aI_A,I_W)\,
                  \Gamma(a,q)^{-1}\binom{b_A(a,q)}{b_W(a,q)}.}
                                                               \label{v39:exact-rate}
\end{equation}
This is the actual signed rate, not its positive minimum substitute.
The radial equation follows from the original homogeneity and mean-lapse
normalization. Formula \eqref{v39:exact-rate} proves all analytic extension
claims. At $q=0$, $r_W=O(a)$ gives $\mathcal U(0)=0$; differentiating
there and using V38 gives $I_W\mathcal U'(0)=k_*$. Analyticity gives the
uniform remainders in \eqref{v39:rate-profile-eq}.

The full Hessian has exactly its radial null direction at every positive
$a$ in this rectangle. The original fixed-cutoff continuation theorem applies
at each point, preserving all original constraints. Uniformity of this
\emph{initial parameter} chart does not bound its physical continuation
time, since differentiated rates need not be uniform as $a\downarrow0$.
\end{proof}

\section{The original metric and links on the nonlinear parameter chart}
\label{v39:observations}

Let $\mathcal R(a,q)=(R_{0i0j}(g_h)(0))_{ij}$ be the electric block of
the original interpolated metric, evaluated with its own canonical and
multiplier rates. Let $\mathcal L(a,q)$ denote the pair of original temporal
and spatial literal-link curvatures in the same fixed compatible frames.
No varied time coordinate is used.

\begin{theorem}[Analytic finite limiting metric curvature and fixed leading link coefficient]
\label{v39:observations-theorem}
Both observations have analytic extensions to the closed parameter rectangle.
With physical synthesis of the weak profile, put
\begin{equation}
 \mathcal Q(q)=\mathfrak E_h\bigl((I_W\mathcal U(q))^\beta\bigr).
                                                               \label{v39:Qprofile}
\end{equation}
Then
\begin{equation}
 \boxed{\begin{aligned}
 \mathcal R(a,q)&=\mathcal Q(q)+a\mathcal R_1(a,q),\\
 \mathcal Q(0)&=0,\qquad
 \mathcal Q'(0)=\mathfrak E_h(k_*^\beta)=:\mathcal Q_*\ne0,\\
 \mathcal L(a,q)&=a\mathcal L_0+a^2\mathcal L_1(a,q).
 \end{aligned}}                                                  \label{v39:observation-expansions}
\end{equation}
The vector $\mathcal L_0$ is independent of $q$. The other independent metric
curvature blocks are $O(a)$ uniformly. Consequently
\begin{equation}
 \begin{aligned}
 \mathcal R(a,q)-\mathcal R(a,0)
       &=q\mathcal Q_*+O(q^2+a|q|),\\
 \|\mathcal L(a,q)-\mathcal L(a,0)\|&\le Ca^2|q|.
 \end{aligned}                                                   \label{v39:secant-upper}
\end{equation}
All norms include the original physical spatial mean and Frobenius
occurrences. These estimates also hold with any fixed spatial norm at this
finite cutoff, with its own constants.
\end{theorem}
\begin{proof}
The canonical first velocity is unchanged along the chart:
\[
 X_t=F(X,\lambda)=a(F_1Y_*+G\ell_*+qGc_H)+O(a^2)
                  =aV_*+O(a^2).
\]
The remainders are analytic uniformly in $q$. By \cref{v39:rate-profile},
$r_A=O(a)$ and $r_W=I_W\mathcal U(q)+O(a)$ in physical synthesis.
Since the canonical map is smooth on the original stationary chart,
\[
 X_{tt}=D_XF\,F+D_\lambda F\,r
               =G I_W\mathcal U(q)+O(a).
\]
Its configuration part is the phase symmetric gradient of the weak shift.
The metric and its spatial derivatives approach the flat metric as $O(a)$.
Its spatial-metric first time derivative is $O(a)$, while
$\beta_t=(I_W\mathcal U(q))^\beta+O(a)$ and $N_t=O(a)$.
Therefore
\[
 g_{0i,t}=(I_W\mathcal U(q))_i^\beta+O(a),\qquad
 g_{ij,tt}=\delta_i(I_W\mathcal U(q))_j^\beta+
                         \delta_j(I_W\mathcal U(q))_i^\beta+O(a).
\]
The spatial Hessian of $g_{00}$ is $O(a)$.

One must not bound both factors in every Christoffel product by $O(a)$:
$\Gamma^\rho_{00}$ can be $O(1)$. The precise cancellation is that
$\Gamma^\rho_{0i}$ and $\Gamma^\rho_{ij}$ are $O(a)$. In a first-kind
symbol with indices $0i$, the two occurrences of $g_{0i,t}$ cancel; it
contains only spatial derivatives of $g_{0\mu}$ and $\gamma_t$.
Thus in the electric curvature product the possibly $O(1)$ symbol is
multiplied by $\Gamma^\rho_{ij}=O(a)$. Every product is $O(a)$.
Using the original formula
\[
 R_{0i0j}=\tfrac12(g_{0j,it}+g_{i0,tj}-g_{00,ij}-g_{ij,tt})
                          +\text{Christoffel products}
\]
gives exactly \eqref{v39:Qprofile}. The remaining blocks use only the
small mixed and spatial symbols and at most $\gamma_t$; they are $O(a)$.
All these expressions are analytic in the initial fields and rates, with
bounded inverse metric, proving the first analytic expansion.

For the literal observations, the stationary spatial connection has
$z=az_1(Y_*)+O(a^2)$, independent of $q$ at first order. The first temporal
connection coefficient is also unchanged: the only changed first multiplier
coefficient is constant, whose flat auxiliary rotation, boost and symmetric
velocity are all zero. Thus $A_0=aA_{0,1}+O(a^2)$ with $A_{0,1}$
independent of $q$. The original stationary derivative obeys
$D_{\lambda_A}z=O(a)$ and $D_{\lambda_W}z=O(a^2)$, uniformly on this
chart, by its signed source factorization. Hence
\[
 z_t=D_Xz\,X_t+D_\lambda z\,r
       =a\,D_Xz(0,e)V_*+O(a^2).
\]
The $O(1)$ weak rate costs two amplitude orders here. The original temporal
link curvature uses $z_t$, $A_0$, and the shifted $A_0$, not $\dot A_0$;
its fixed-cutoff analytic expansion has the same first coefficient for every
$q$. The spatial link expression uses the same first $z$ coefficient.
This proves the final expansion of \eqref{v39:observation-expansions}.

Analyticity on a slightly larger rectangle bounds the fixed derivatives of
$\mathcal R_1,\mathcal L_1,\mathcal U$. Subtract their values at $q=0$
and apply the fundamental theorem of calculus. Taylor expansion of
$\mathcal Q$ at zero gives the first line of \eqref{v39:secant-upper},
including the factor $|q|$ on the amplitude remainder. The second follows
from the same subtraction of $\mathcal L_1$. Finally V38's certified
$6300<\|\mathcal Q_*\|<7000$ proves nonvanishing. No unknown higher
coefficient is assigned a favorable sign.
\end{proof}

\section{Finite nonlinear separation and a scalar invariant}
\label{v39:finite-separation}

\begin{corollary}[Two-sided finite secants on exactly constrained records]
\label{v39:finite-secants}
After reducing $a_0,q_0$ there are fixed $c_0,C_0>0$ such that
\begin{equation}
 \boxed{
 c_0|q|\le\|\mathcal R(a,q)-\mathcal R(a,0)\|\le C_0|q|,
 \qquad \|\mathcal L(a,q)-\mathcal L(a,0)\|\le Ca^2|q|.
 }                                                            \label{v39:two-sided}
\end{equation}
In the original perturbation variable $\varepsilon=aq$ this becomes
\begin{equation}
 \boxed{\begin{aligned}
 \|\Delta X\|&\le Ca^2|\varepsilon|,\qquad
                  \|\Delta\lambda\|\asymp|\varepsilon|,\\
 \|\Delta X_t\|&\le Ca|\varepsilon|,\\
 \|\Delta\mathcal R\|&\asymp |\varepsilon|/a,\qquad
                  \|\Delta\mathcal L\|\le Ca|\varepsilon|,
 \qquad |\varepsilon|\le aq_0.
 \end{aligned}}                                                  \label{v39:epsilon}
\end{equation}
For any fixed sufficiently small $q\ne0$, the initial fields and multipliers
converge to their flat values, the literal curvatures tend to zero, but
$\mathcal R(a,q)\to\mathcal Q(q)\ne0$. Each record has a local
original-action continuation. This is a finite parameter statement, not a
time-dependent instability or a spatial-refinement counterexample.
\end{corollary}
\begin{proof}
By \eqref{v39:secant-upper}, the error from $q\mathcal Q_*$ is at most
$C(|q|+a)|q|$. Choose the rectangle so
$C(q_0+a_0)<\|\mathcal Q_*\|/2$. The triangle and reverse triangle
inequalities give \eqref{v39:two-sided}. Equations
\eqref{v39:record-size} and $q=\varepsilon/a$ give
\eqref{v39:epsilon}; its velocity estimate follows by subtracting the
uniform analytic expansion $X_t=aV_*+O(a^2)$ from the observation proof.
The limiting assertion follows from
\eqref{v39:observation-expansions} and the same lower bound at $a=0$.
The chart selects distinct original initial records; it is not a coordinate
transformation of the same spacetime or a modification of the readout.
\end{proof}

\begin{corollary}[The finite discrepancy is not removable by an initial coordinate change]
\label{v39:invariant}
For fixed nonzero $q$ as above, in the original curvature convention,
\begin{equation}
 \begin{aligned}
 G_{00}(a,q)&\longrightarrow0,\qquad G_{0i}(a,q)\longrightarrow0,\\
 (G_{ij}(a,q))_{ij}&\longrightarrow
                  -\mathcal Q(q)+\tr(\mathcal Q(q))I_3\ne0.
 \end{aligned}                                                   \label{v39:Einstein}
\end{equation}
For the scalar $\mathscr K(g)=R_{abcd}R^{abcd}$,
\begin{equation}
 \boxed{
 \lim_{a\downarrow0}\int_{\T^3}\mathscr K(g_{a,q})(0,x)\,dx
                 =4\|\mathcal Q(q)\|_{L^2}^2>0.
 }                                                            \label{v39:Kretschmann}
\end{equation}
The fixed-cutoff trigonometric and physical spatial-mean norms agree for this
limiting tensor. The conclusion concerns the same initial events; changing
the time slice is a different comparison.
\end{corollary}
\begin{proof}
The inverse metric tends to $\operatorname{diag}(-1,1,1,1)$ and all curvature
blocks except $R_{0i0j}$ tend to zero. Contracting this algebraic Riemann tensor
gives $\operatorname{Ric}_{00}=\tr\mathcal Q$,
$\operatorname{Ric}_{ij}=-\mathcal Q_{ij}$ and
$R=-2\tr\mathcal Q$, proving \eqref{v39:Einstein}. The spatial limit
cannot be zero because
\[
 \|-\mathcal Q+(\tr\mathcal Q)I_3\|_F^2
                  =\|\mathcal Q\|_F^2+(\tr\mathcal Q)^2.
\]
There are four symmetry-related occurrences of each electric component in
the fully contracted curvature square. Each has two time indices, so raising
them introduces a positive sign. Thus its limit is $4|\mathcal Q|_F^2$.
The fixed finite Fourier setting and uniform analytic expansions permit
integration, proving \eqref{v39:Kretschmann}. This is a scalar observation,
not merely a nonzero coordinate component. The limit need not be the curvature
of the limiting zeroth-order metric; convergence in two time derivatives is
precisely absent here.
\end{proof}

\section{An amplitude-dependent preparation precision, and what remains}
\label{v39:precision}

\begin{proposition}[Sharp preparation scale along the constructed direction]
\label{v39:precision-theorem}
For any selection $\varepsilon(a)$ with $|\varepsilon(a)|\le aq_0$,
\begin{equation}
 \boxed{\begin{aligned}
 \|\mathcal R(a,\varepsilon(a)/a)\|=O(a)
       &\quad\Longleftrightarrow\quad \varepsilon(a)=O(a^2),\\
 \|\mathcal R(a,\varepsilon(a)/a)\|\longrightarrow0
       &\quad\Longleftrightarrow\quad \varepsilon(a)=o(a).
 \end{aligned}}                                                  \label{v39:precision-eq}
\end{equation}
On the entire allowed chart the original literal curvature is $O(a)$,
regardless of these stricter metric requirements. The equivalences concern
this verified one-parameter family, not every possible initial perturbation.
\end{proposition}
\begin{proof}
The base metric curvature is $O(a)$. Apply \eqref{v39:two-sided} to its
difference from the perturbed observation. If both are $O(a)$, its lower
bound gives $|q(a)|=O(a)$; conversely this bound and the upper estimate
suffice. The same two inequalities show that the perturbed curvature tends
to zero exactly when $q(a)\to0$. Substitution $q=\varepsilon/a$ proves
both equivalences. The link assertion is
\eqref{v39:observation-expansions}. No lower link bound or ratio with a
possibly zero link difference is assumed.
\end{proof}

The new acquisition is the uniform finite radius and the controlled
nonlinear remainder in \eqref{v39:epsilon}, not another differential
nonvanishing certificate. It is obtained from the original field-constraint
submersion and the affine cubic weak identity before using the signed inverse.
The construction does not solve a common-lifespan problem. For fixed nonzero
$q$, these new records generally have an $O(q)$ weak rate, rather than the
$O(a)$ rate of the specially selected histories. Their individual local
continuations are not claimed to share V33's $ca^2$ interval. That remains
the longest proved interval for its own selected family. All new conclusions
are at fixed $h=1/3$, and the changed initial shift means are retained physical
values, not erased as gauge parameters.

The next useful estimate must propagate a geometry-sensitive initial budget
such as \eqref{v39:precision-eq}, or control the actual rate component on an
invariant selected family over physical time. Link values alone do not supply
that estimate. No spatial-refinement result, longer physical interval,
normalized metric/link identification, or original-action Einstein continuum
theorem is asserted.

\section{Verification and preservation}
\label{v39:verification}

The accompanying checker verifies the new finite Taylor-jet identities in a
separate exact polynomial Hamiltonian, including the dependence of the
multiplier Hessian, the vanishing mixed constant-direction second derivative,
and the affine cubic weak coefficient. These checks test the algebra used in
the written proof, not a replacement gravitational action. Exact symbolic
checks also retain the noncommuting graded inverse, the row scaling of the
constraint solve, the Christoffel cancellation with a bounded nonzero shift
rate, and all Ricci, Einstein and curvature-square contractions of the limiting
electric tensor. Explicit analytic scalar examples check that dropping the
affinity or the uniform divided chart would not justify the finite conclusion.

The complete V38 portable verification is rerun, including its original V37
cubic coefficient enclosure, complete slow spectral acceptance, moving-source
certificate, and 216 independent original scalar-row checks. Its accepted
nonzero coefficient is inherited rather than newly relabelled. The V29
nonlinear root and the V30 higher-source enclosure remain the named inputs
of that verification, not an independently reconstructed entire theory.
No new higher-action coefficient, numerical nonlinear initial root, or
spacetime integration is reported. The new parameter radius exists by the
proved analytic submersion and derivative bounds; it is not assigned an
unverified numerical value.

Every marked V1--V38 body is preserved byte for byte. The package contains the
new body, cumulative source, predecessor, assembly script, exact checker,
executed reports, and the inherited verifier with its required data. The proofs
are not proof-assistant formalizations or independently peer-reviewed results.

\begin{thebibliography}{GN79}
\small\raggedright
\bibitem[GN79]{v39:GN}
M. J. Gotay and J. M. Nester,
\emph{Presymplectic Lagrangian systems. I: the constraint algorithm and the
equivalence theorem}, Annales de l'Institut Henri Poincar\'e, Section A,
\textbf{30} (1979), 129--142.
\url{https://www.numdam.org/item/AIHPA_1979__30_2_129_0/}.
Primary-source background for retaining the original constraint system and
its own local evolution. The parameter-chart and source-specific amplitude
estimates here are proved in this body, not imported from this reference.
\end{thebibliography}
% END IMMUTABLE EXACT-ACTION BRIDGE V39 BODY


% Future iterations append here.
% BEGIN IMMUTABLE EXACT-ACTION BRIDGE V40 BODY
\clearpage
\renewcommand{\thesection}{\arabic{section}}
\renewcommand{\theHsection}{bridgeV40.\arabic{section}}
\setcounter{section}{322}
\section{Propagate the actual preparation error, not an independently chosen mode}
\label{v40:context}

V39 supplies finite exactly constrained perturbations, but only estimates their
initial observations. V31 supplies an expanding rapid mode, but excites it by a
separately chosen initial-rate variation. Neither statement establishes that
V39's \emph{actual} slow-direction preparation error excites that mode, or that
its finite metric/link separation persists during the original evolution. We
establish that connection here. A new interval calculation proves nonzero
coupling of V38's moving-source vector to the dominant rapid left mode. An
anisotropically scaled initial-value problem then propagates the whole V39
parameter chart. A logarithmic-time estimate pays its nonlinear remainder and
turns the original initial precision requirement into a sharp time-dependent
one on this family.

All action-specific statements retain $N=3$, $h=1/3$, the same V29 root, the
original signed multiplier equation, and its original metric and link reads.
The positive variable $a$ is canonical field amplitude. The parameter $q$ is
V39's scaled initial-data displacement, with original displacement
$\varepsilon=aq$. The identity $t=a^3s$ is an analytical time coordinate;
$t$ remains the original physical time. In particular the new theorem is not
a spatial-refinement statement.

We select the analytic base preparation of \cref{v30:transport-rate}, whose own
initial rate is exactly $ar_*$. Apply \cref{v39:constraint-chart} to that base,
using the analytic slow tangent of V37--V38. This is a permitted instance of
V39, not an assumption that every $O(a)$-rate base has the transport estimate.
The first canonical coefficient and first nonconstant multiplier tilt remain
$Y_*,\ell_*$. The changed constant initial shift $aqc_H$ is retained.
No endpoint reselection, damping, or reset is made in the histories studied
below.

\paragraph{Inherited inputs and new work.}
The exact root and signed chart are V29--V30. The narrow positive rapid
eigenvalue is \cref{v31:eigen-theorem}; its original block conventions are
\eqref{v31:BK}. The moving-source vector and enclosure are V38. The finite
parameter chart, cubic cancellation and rate profile are
\cref{v39:constraint-chart,v39:rate-profile}. The analytic stationary maps,
graded-source identities and homogeneous off-constraint extension are retained
from V20 and \eqref{v30:P-extension}. We do not reprove those results or claim
an independent validation of the whole inherited theory. The new acquisitions
are an actual source/eigenmode pairing, an exact limiting propagation law for
these finite perturbations, a parameter-differentiated logarithmic remainder
estimate, and the resulting necessary-and-sufficient preparation scales.

\section{The actual slow-data response excites the dominant rapid mode}
\label{v40:coupling}

For this body order the weak coordinates as the 26 gradients
$\nabla_\delta(e_j-e_{26})$, $0\le j<26$, followed by the three constant shifts.
This is the fixed permutation used by V38's supplied interval arrays. Let
$I_A,I_W$ be the corresponding physical synthesis maps, with $I_A$ in the
mean-zero lapse/transverse-shift slice. All forms use one common spatial-sum
pairing; multiplication of all forms by $1/27$ leaves their response operator
unchanged. Set
\begin{equation}
 \begin{gathered}
 \Gamma_* =\begin{pmatrix}A&L\\L^*&C\end{pmatrix},\qquad
 D=C-L^*A^{-1}L,\\
 K=I_W^*\mathbb K_1(Y_*)I_W,\qquad
 B=-D^{-1}K,\qquad H=A^{-1}L.
 \end{gathered}                                      \label{v40:blocks}
\end{equation}
Here $A$ is positive and $D$ is the original nonsingular signed weak Schur
form. Neither is replaced by a positive surrogate. Write $k\in\R^{29}$ for
V38's actual rate coefficient in these coordinates; thus $I_Wk=k_*$ in
\eqref{v38:k}. Equivalently $k=\mathcal U'(0)$ in V39's rate profile, after
this fixed permutation.

\begin{theorem}[Certified source coupling and a dominant real mode]
\label{v40:source-mode}
Let $\gamma$ be the simple positive eigenvalue in \eqref{v31:gamma}. It is the
unique eigenvalue of $B$ with maximal real part; every other eigenvalue has
real part smaller than $3/50$. There is a unique real left eigenvector $\ell$
normalized by $\ell_{10}=1$ (zero-based indexing) and
\begin{equation}
 \boxed{\ell B=\gamma\ell,\qquad
        \ell|_{\text{constant shifts}}=0,\qquad
             -640<\ell k<-490.}                    \label{v40:nonzero-coupling}
\end{equation}
In particular the specified moving-source coefficient, not an independently
chosen forcing, has a nonzero expanding component.
\end{theorem}
\begin{proof}
We give a reproducible interval argument with its complete normalizations.
The source arrays are V38's \src{rate_coefficient_bounds.json} and
\src{weak_Schur_bounds.json}, and the original first-Poisson enclosure is
V37's \src{K_bounds.json}. Their common root and physical synthesis are those
already used in V38. The new checker forms $K=I_W^*\mathbb K_1I_W$ and
$B=-D^{-1}K$. The three constant columns of $K$ vanish by the exact original
Poisson identity; they and their skew rows are set to their exact zero values,
not classified by a numerical threshold.

Every endpoint is an integer divided by $2^{128}$. Matrix products are rounded
outward using integer center--radius arithmetic. Fixed binary-rational inverse
proposals $R_D,R_J$ give
\begin{equation}
 \|I-R_DD\|_\infty<10^{-6},\qquad
 \|I-R_JJ\|_\infty<10^{-4},                         \label{v40:inverse-tests}
\end{equation}
where $J$ is the principal 28-by-28 submatrix of $B^*-\gamma I$ obtained by
omitting coordinate 10. The second interval includes the entire inherited
interval $0.09479542352<\gamma<0.09479542359$. To enclose each inverse,
write $E=I-RA$ and retain $R+ER+E^2R$; the remaining matrix has norm at most
$\|E\|^3\|R\|/(1-\|E\|)$. This is the convergent Neumann argument,
not acceptance of a floating inverse.

At the actual eigenvalue, invertibility of $J$ implies that a nonzero left
eigenvector has nonzero coordinate 10: otherwise its other 28 coordinates
would vanish. Solve the other 28 equations with coordinate 10 equal to one.
The omitted equation holds because this is the actual eigenvalue. Also
$B I_H=0$ and $\gamma\ne0$ imply $\ell I_H=0$ exactly. The resulting
interval vector, multiplied by the full original interval $k$, gives
\begin{equation}
 -636.868<\inf(\ell k)\le\sup(\ell k)<-496.867.
                                                        \label{v40:pairing-box}
\end{equation}
Only the wider strict rational interval in
\eqref{v40:nonzero-coupling} is used. The exact endpoint fractions and all
inverse proposals are supplied in the new certificate. The inverse residuals
are about $1.132\,10^{-7}$ and $3.637\,10^{-5}$; those decimals are descriptive,
while acceptance compares the complete rational fractions.

For completeness a separate 29-dimensional similarity certificate verifies
dominance. Its fixed real binary-rational matrix $V$ and real block-diagonal
reference $\Lambda$ use one-by-one real blocks, normal two-by-two conjugate
blocks, and a five-dimensional zero reference block. The checker bounds
\[
 2\|V^{-1}(BV-V\Lambda)\|_\infty<10^{-4}.
\]
The factor two pays the exact complex diagonalization of the two-by-two
blocks. The actual accepted radius is less than $8.771\,10^{-5}$. The disk
containing the inherited $\gamma$ is isolated from every other disk. Every
other disk has real part smaller than $3/50$, verified by rational comparisons
of its center and radius. Scaling the error continuously from zero to its
actual value keeps all roots in these disk groups by row dominance; the
isolated group therefore contains exactly one simple eigenvalue. Its reality
also follows from conjugate symmetry. This identifies the same inherited
$\gamma$ and proves the claimed strict separation. No conclusion about exact
zeros or Jordan blocks in the other groups is needed.
\end{proof}

Let
\begin{equation}
 \mathcal E w=\mathfrak E_h\bigl((I_Ww)^\beta\bigr)                 \label{v40:read}
\end{equation}
be the original electric metric-curvature response to a weak rate. Its kernel
is exactly the three constant coordinates, by the full-frequency symbol of
\cref{v10:symbol}. Thus $\ell$ factors through this \emph{original} observation:
there is a bounded real functional $\Lambda_{\mathcal E}$ such that
\begin{equation}
                 \Lambda_{\mathcal E}\mathcal Ew=\ell w.
                                                        \label{v40:dual-read}
\end{equation}
Indeed define the functional on $\Ran\mathcal E$ by this formula; it is
well-defined since $\ell$ annihilates the kernel, and extend it by orthogonal
projection. All spaces are finite-dimensional with their original positive
observation norms. This is not a redefinition of curvature.

\begin{corollary}[Uniform growth of the actual limiting observation]
\label{v40:profile-growth}
After reducing the interval in V39's analytic profile $\mathcal U(q)$, there
are constants $c,C>0$ such that for all $s\ge0$ and all sufficiently small
$|q|$,
\begin{equation}
 \boxed{c e^{\gamma s}|q|\le
       \|\mathcal E e^{sB}\mathcal U(q)\|_h
          \le C e^{\gamma s}|q|.}                  \label{v40:profile-growth-eq}
\end{equation}
If $v$ is normalized by $Bv=\gamma v$, $\ell v=1$, then
\begin{equation}
 e^{-\gamma s}e^{sB}\longrightarrow v\ell
                    \quad\text{in operator norm as }s\to\infty,
 \qquad \mathcal Ev\ne0.                           \label{v40:dominant-projector}
\end{equation}
These are properties of the fixed limiting response, not yet assertions about
physical histories on unbounded scaled intervals.
\end{corollary}
\begin{proof}
Simplicity gives $\ell v\ne0$ before normalization. The other spectral real
parts are strictly below $3/50<\gamma$. A Jordan decomposition bounds their
polynomial factors by a slightly larger exponential still smaller than
$e^{\gamma s}$. This proves \eqref{v40:dominant-projector} and
$\|e^{sB}\|\le Ce^{\gamma s}$. Since $\mathcal U(0)=0$ and
$\ell\mathcal U'(0)=\ell k\ne0$, Taylor's theorem gives
$c|q|\le|\ell\mathcal U(q)|\le C|q|$ on a smaller interval. Apply
\eqref{v40:dual-read} and
$\ell e^{sB}\mathcal U(q)=e^{\gamma s}\ell\mathcal U(q)$ for the lower
bound, and the semigroup upper bound for the other. If $\mathcal Ev=0$, then
$v$ would be a constant shift and $Bv=0$, a contradiction.
\end{proof}

\section{An anisotropic initial-value problem on the unchanged action}
\label{v40:scaled-flow}

Denote V39's exact initial records by $X^0(a,q),\lambda^0(a,q)$. On the
physical mean-lapse-zero complement use
\begin{equation}
 \begin{aligned}
 t&=a^3s,\\
 X&=X^0(a,q)+a^4\xi,\\
 \lambda&=\lambda^0(a,q)+a^4I_A\eta_A+a^3I_Wy,\\
 u&=\eta_A+Hy,\qquad z=(\xi,u,y).
 \end{aligned}                                                \label{v40:scaling}
\end{equation}
The initial displacement is zero in every component. The weak displacement
has scale $a^3$, rather than the $a^4$ scale appropriate to a prescribed
$O(a)$ initial weak rate. This is necessary because V39 gives an $O(q)$ weak
rate for fixed $q$. The shear has the fixed coefficient $H=A^{-1}L$ and is
invertible with determinant one; no state-dependent differential is assumed
integrable.

Let $\alpha(q)=\lim_{a\to0}r_A(a,q)/a$ in V39's rate coordinates and put
\begin{equation}
 \rho(q)=\alpha(q)+H\mathcal U(q),\qquad
 b(q)=(V_*,\rho(q),\mathcal U(q)),\qquad
 J=\diag(0_{402},B).                              \label{v40:b-J}
\end{equation}
In particular $\mathcal U(0)=0$ and $\rho(0)=r_{*,A}$ for our chosen base.

\begin{lemma}[Exact scaled equation and polynomial derivative bounds]
\label{v40:remainder}
For small fixed $q_0,d,a_0>0$, the analytical extension of the original
equations in \eqref{v30:P-extension} takes the form
\begin{equation}
                       z'=b(q)+Jz+\mathcal R_a(q,z),\qquad z(0)=0.
                                                        \label{v40:ode}
\end{equation}
It has an analytic extension through $a=0$ on bounded $z$ sets. On
$|q|\le q_0$, $0<a\le a_0$, $a^2\|z\|\le d$, one has, with a fixed $C$,
\begin{equation}
 \begin{split}
 \|\mathcal R_a\|+\|\partial_q\mathcal R_a\|
                +\|\partial_q^2\mathcal R_a\|
                      &\le Ca(1+\|z\|)^2,\\
 \|D_z\mathcal R_a\|+\|\partial_qD_z\mathcal R_a\|
                      &\le Ca(1+\|z\|)^2.
 \end{split}                                              \label{v40:R-est}
\end{equation}
Every solution starting from the specified records satisfies $P=0$ and is
therefore a solution of the original normalized action equations. No new
physical vector field is used on that constraint surface.
\end{lemma}
\begin{proof}
Write $\overline{\mathsf B}=\mathsf B+\mathscr R P$ for the inherited
analytical completion. It has no degree-zero or degree-one joint Taylor term
at $(0,e)$; its pure-shift rows equal the original drift rows. Initially
$P=0$, and V39 gives analytic divided source coefficients
\[
 \overline{\mathsf B}_A^0=a^3 b_A(a,q),\qquad
 \overline{\mathsf B}_W^0=a^4 b_W(a,q).
\]
The leading graded source depends only on $Y_*$, not on the first constant
multiplier tilt $qc_H$, so the normalized signed matrix tends to $\Gamma_*$
uniformly in $q$.

The same divisibility holds on \eqref{v40:scaling}. The first two canonical
coefficients do not change; a weak displacement now first changes the third
multiplier coefficient. The quadratic prepared drift remains zero, and the
complete cubic weak drift is independent of that third coefficient. Its
initial value vanishes by \cref{v39:rate-profile}. For its next coefficient,
the original differential identity
$D_\lambda\mathsf B=\mathbb K+D_XM[F]$ and the flat orders give
\[
 D_\lambda\overline{\mathsf B}
       =a\mathbb K_1(Y_*)+O(a^2),\qquad
 D_X\overline{\mathsf B}=O(a)
\]
on bounded normalized charts. The extra $\mathscr R M$ is $O(a^2)$.
Consequently on bounded $z$ sets,
\begin{equation}
 \begin{aligned}
 a^{-3}\overline{\mathsf B}_A
       &=b_A(0,q)+O(a),\\
 a^{-4}\overline{\mathsf B}_W
       &=b_W(0,q)+Ky+O(a),\\
 \Gamma(a,q,z)&=\Gamma_*+O(a).
 \end{aligned}                                             \label{v40:row-expansion}
\end{equation}
These are expansions of the full analytic maps, not a polynomial substitute
for the action. In particular the first remainder in
\eqref{v40:row-expansion} includes the term $aK_{AW}y$; it is not set to zero.

Using $M=\diag(aI_A,a^2I_W)\Gamma\diag(aI_A,a^2I_W)$, the exact
normalized rate is
\[
 \binom{r_A/a}{r_W}
  =-\Gamma(a,q,z)^{-1}
        \binom{a^{-3}\overline{\mathsf B}_A}
              {a^{-4}\overline{\mathsf B}_W}.
\]
Its weak leading row is $\mathcal U(q)+By$, and its active leading row is
$\rho(q)-H(\mathcal U(q)+By)$. Since $y'=r_W$ and
$u'=r_A/a+Hy'$, the latter coupling cancels in $u'$. Also
$F(X,\lambda)/a=V_*+O(a)$ on bounded $z$ sets because $Gc_H=0$.
This proves the leading vector field \eqref{v40:ode}.

We record why these are uniform bounds for growing $z$, not just formal
fixed-$z$ expansions. In the stated domain the normalized arguments are
\[
 X/a=Y_*+O(a)+a^3\xi,\qquad
 (\lambda-e)/a=\ell_*+qc_H+O(a)+a^3I_A(u-Hy)+a^2I_Wy.
\]
They lie in a fixed compact analytic chart after reducing $d,q_0,a_0$.
The normalized source and signed inverse have bounded fixed derivatives
there; their leading signed matrix is independent of the bounded first
multiplier tilt. For the completed source, Taylor expansion about the initial
record retains the following possible orders before division: the linear
weak multiplier displacement gives $a^4Ky$; its next coefficient is
$O(a^5\|z\|)$; the canonical displacement gives $O(a^5\|z\|)$;
and the quadratic cross and multiplier terms are at most
$O(a^7\|z\|^2)$. The active rows have the same or better orders before
their smaller division by $a^3$. The initial divided-coefficient remainders
are $O(a)$. These estimates follow by applying the full second-order integral
Taylor formula: $D_X\overline{\mathsf B}=O(a)$,
$D_\lambda\overline{\mathsf B}=a\mathbb K_1+O(a^2)$,
$D_\lambda^2\overline{\mathsf B}=O(a)$, and the other second derivatives
are bounded. Displacements in the mixed term have sizes $a^4\|z\|$ and
$a^3\|z\|$. Their higher remainders are controlled on the same chart.

The inverse variation has size at most $C(a+a^2\|z\|)$; multiplying it
by a source of size $C(1+\|z\|)$ remains bounded by the first line of
\eqref{v40:R-est}. The canonical row has an additional term
$a^2G I_Wy$ and terms of order $a^3\|z\|$, which satisfy the same
bound. Differentiating these integral Taylor expressions once or twice in
$q$, or once in $z$ and then in $q$, gives the displayed conservative bounds.
The initial chart derivatives are bounded, the leading $Y_*$ and $K$ do not
vary with $q$, and all leading $q$ dependence has already been placed in
$b(q)$. No derivative of a singular unscaled inverse is bounded independently
of its explicit amplitude factors.

Finally the extension obeys the exact homogeneous equation
\eqref{v30:P-extension}. Since the initial constraint is exactly zero, its
solution remains zero. On it the completed and original sources agree,
including the radial row recovered by homogeneity and mean-lapse one.
\end{proof}

\section{A populated finite-time response of the original perturbation family}
\label{v40:finite-time}

Put $\Phi_B(s)=\int_0^s e^{uB}\,\dd u$; this definition uses no inverse of
$B$, so its exact constant kernel is retained. For each fixed finite $S$ the
limiting solution of \eqref{v40:ode} is
\begin{equation}
 z_0(s,q)=\bigl(sV_*,\ s\rho(q),\ \Phi_B(s)\mathcal U(q)\bigr).
                                                        \label{v40:limiting-flow}
\end{equation}

\begin{theorem}[Nonlinear original histories realize the rapid response]
\label{v40:bounded-time}
For every fixed $S<\infty$, all initial records in a sufficiently small fixed
$q$ interval have original normalized histories on $0\le t\le a^3S$ for
all sufficiently small $a$. In the coordinates \eqref{v40:scaling},
$z_a-z_0=O_S(a)$, with one $q$ derivative and one $s$ derivative, uniformly
on that interval. In particular
\begin{equation}
 \boxed{r_W(a,a^3s;q)=e^{sB}\mathcal U(q)+O_S(a).}     \label{v40:rate-time}
\end{equation}
Let $\mathcal R_a(s,q)$ be the original metric electric curvature block and
$\mathcal L_a(s,q)$ the original temporal/spatial link pair at that physical
time. Then
\begin{equation}
 \boxed{\begin{aligned}
 \mathcal R_a(s,q)&=\mathcal E e^{sB}\mathcal U(q)+O_S(a),\\
 \mathcal L_a(s,q)&=a\mathcal L_0+O_S(a^2),
 \end{aligned}}                                             \label{v40:obs-time}
\end{equation}
where $\mathcal L_0$ is independent of $s,q$. All other independent metric
curvature blocks are $O_S(a)$. Their remainders and the displayed remainders
have bounded first $q$ derivatives of the same orders.
\end{theorem}
\begin{proof}
On bounded $s$ intervals \eqref{v40:limiting-flow} and its $q$ derivative
are bounded. The integral equation, \eqref{v40:R-est}, and Gronwall give a
uniform $O_S(a)$ difference while the scaled solution is in a fixed larger
ball. That bound prevents exit; the original stationary and signed charts
also persist since their normalized physical displacements tend to zero.
The homogeneous constraint identity makes these actual histories. Differentiate
the finite integral equation in $q$ to obtain the corresponding parameter
estimate; insert the equations to obtain the $s$ derivative estimate. The
weak displacement has time derivative $r_W=y'$, which proves
\eqref{v40:rate-time}. The active rate is $a$ times a bounded function.

Throughout such a bounded interval, $X/a\to Y_*$ and $X_t/a\to V_*$;
$(\lambda-e)/a\to\ell_*+qc_H$. Only the constant part of this last
coefficient changes with $q$. The same original metric calculation as V39,
now evaluated on the histories, gives
$R_{0i0j}=\mathfrak E_h(r_W^\beta)+O_S(a)$. The mixed and spatial
curvature blocks use only the small spatial fields and canonical first
velocity. In the electric Christoffel products a possibly bounded
$\Gamma^\rho_{00}$ always multiplies a spatial symbol $O(a)$; the
mixed symbols have no uncancelled $\beta_t$. Thus no assumption that all
symbols are $O(a)$ is made.

The stationary spatial connection has first coefficient $z_1(Y_*)$, and the
temporal connection's first coefficient is unchanged by $qc_H$. In
$z_t=D_Xz\,F+D_\lambda z\,r$, the active derivative is $O(a)$ and the
weak derivative $O(a^2)$; the latter multiplies the bounded rate in
\eqref{v40:rate-time}. Hence the first link-curvature coefficient is the
same $\mathcal L_0$ for all $s,q$. The temporal read contains $z_t$ and
shifted $A_0$, not $\dot A_0$. This proves \eqref{v40:obs-time} in the
original frames. The finite analytic formulas and the differentiated flow
bounds give the remainder derivatives asserted.
\end{proof}

This theorem populates the response by evolving V39's exact records. It is
not a claim that a slow frozen eigenvector stays aligned with the moving
flow: its actual leading rate decomposes into the rapid modes through $k$,
whose expanding projection was computed in \cref{v40:source-mode}.

\section{Paying the remainder on a logarithmically growing scaled interval}
\label{v40:logarithmic}

For a function of the two histories, $\Delta_q$ means subtraction of its value
at $q=0$, with both values evaluated at the same $a,s$. The next theorem is
not obtained by replacing a fixed-$S$ constant in \eqref{v40:obs-time} by an
unverified uniform one.

\begin{theorem}[Parameter-differentiated logarithmic propagation]
\label{v40:log-theorem}
There are fixed $M,C,c_0>0$ with $c_0M<1/4$ such that, for every fixed
$0<c<c_0$, $S_a=c\log(1/a)$, all the original initial records with
$|q|\le q_0$ have histories through $t=a^3S_a$ for sufficiently small $a$.
They satisfy
\begin{equation}
 \boxed{\begin{aligned}
 \|\Delta_q\mathcal R_a(s)-\mathcal E e^{sB}\mathcal U(q)\|_h
                         &\le Ca e^{Ms}|q|,\\
 \|\Delta_q\mathcal L_a(s)\|_h&\le Ca^2 e^{Ms}|q|,
               \qquad 0\le s\le S_a.
 \end{aligned}}                                             \label{v40:log-errors}
\end{equation}
The reference histories obey
\begin{equation}
 \sup_{s\le S_a}(\|\mathcal R_a(s,0)\|+
                         \|\mathcal L_a(s,0)\|)\le Ca.         \label{v40:base-log}
\end{equation}
In particular all perturbed literal-link curvatures remain $O(a)$ on this
whole interval. No endpoint condition is imposed on these histories.
\end{theorem}
\begin{proof}
We detail the needed growing-interval estimate. Choose a fixed $L>0$ larger
than $1+\|J\|$, and enlarge constants so that the explicit linear solution
$z_0$ and its first two $q$ derivatives are at most $C_0e^{Ls}$.
Bootstrap $\|z\|\le2C_0e^{Ls}$. The remainder bound gives
$\|\mathcal R_a(q,z)\|\le C_1ae^{2Ls}$. Variation of constants therefore
gives $\|z-z_0\|\le C_2ae^{2Ls}$, after increasing $L$ to absorb an
irrelevant polynomial factor. For $s\le c\log(1/a)$ with $c$ sufficiently
small this closes the bootstrap, and $a^2\|z\|\to0$ uniformly.

For $v=\partial_qz$, differentiate the equation:
\[
 v'=b'(q)+Jv+\partial_q\mathcal R_a+D_z\mathcal R_a\,v,
 \qquad v(0)=0.
\]
On the same bootstrap the two coefficient remainders are at most
$Cae^{2Ls}$. If $c$ is reduced so that
$\int_0^{S_a}Cae^{2Lu}\,\dd u=o(1)$, variation of constants and a second
Gronwall bound give $\|v\|\le C e^{L_1s}$ for a fixed $L_1$, and
$\|v-\partial_qz_0\|\le Ca e^{L_2s}$ for another fixed $L_2$.
The equation yields the same type of bound for one $s$ derivative and its
$q$ derivative. All exponents are fixed before $c$ is chosen; shrinking
$c$ pays each of the finitely many exponents. The argument is in a fixed
finite-dimensional norm, so it does not hide a mesh constant.

The stationary and observation formulas can be estimated on this larger
scaled domain with the same finite-order Taylor method as in
\cref{v40:remainder}. The metric remainder beyond $\mathcal E y'$ is
$a$ times a function with at most a fixed polynomial growth in $z,z'$;
its total $q$ derivative has the same type of bound. This follows directly
from the original ADM curvature formulas: $\gamma-I,\lambda-e$ and their
spatial derivatives are $O(a)$ in the normalized domain, the canonical
velocity has leading term $aV_*$, and the only unsuppressed multiplier rate
in the electric block is the already separated phase-versus-ordinary
symmetric gradient. The other products contain a small spatial symbol.
The link remainder beyond $a\mathcal L_0$ has an extra factor $a$, by the
active and weak stationary derivative orders used in
\cref{v40:bounded-time}. Its total $q$ derivative is bounded in exactly
the same way. Bounded inverses of the metric and of the normalized signed
matrix are retained throughout. Increase $M$ to cover the resulting finite
polynomial exponents. This gives $Ca e^{Ms}$ and $Ca^2e^{Ms}$ bounds for
the \emph{total parameter derivatives} of the respective observation
remainders. Integrating that derivative from $0$ to $q$ proves
\eqref{v40:log-errors}, including its essential factor $|q|$.

The chosen base is the transport preparation, so \cref{v31:log-theorem}
provides \eqref{v40:base-log} after reducing $c_0$ if necessary. This input
cannot be replaced by initial smallness of an arbitrary base rate. Finally
choose $c_0$ smaller still so $c_0M<1/4$. Physical displacement in
\eqref{v40:scaling} stays in the original positive-lapse and signed chart,
and the exact homogeneous constraint equation prevents leaving $P=0$.
Ordinary finite-dimensional continuation then reaches $S_a$. From
\eqref{v40:log-errors}, the link difference is at most
$Ca^{2-cM}|q|=o(a)$ uniformly for bounded $q$, proving its final assertion.
\end{proof}

\section{A sharp dynamical preparation scale on the actual one-parameter family}
\label{v40:precision}

\begin{theorem}[Finite dynamical secants and necessary-and-sufficient precision]
\label{v40:precision-theorem}
For $c,S_a$ as in \cref{v40:log-theorem}, after reducing $a_0,q_0$, there
are constants $c_1,C_1>0$ such that for all $|q|\le q_0$ and
$0\le s\le S_a$,
\begin{equation}
 \boxed{c_1e^{\gamma s}|q|\le
    \|\mathcal R_a(s,q)-\mathcal R_a(s,0)\|_h
                  \le C_1e^{\gamma s}|q|.}           \label{v40:dynamic-secants}
\end{equation}
For any amplitude-dependent original preparation displacement $\varepsilon(a)$
with $|\varepsilon(a)|\le aq_0$, these imply
\begin{equation}
 \boxed{\begin{aligned}
 \sup_{0\le t\le ca^3\log(1/a)}\|\mathcal R(a,t;\varepsilon/a)\|=O(a)
     &\quad\Longleftrightarrow\quad
                \varepsilon(a)=O(a^{2+c\gamma}),\\
 \sup_{0\le t\le ca^3\log(1/a)}\|\mathcal R(a,t;\varepsilon/a)\|\to0
     &\quad\Longleftrightarrow\quad
                \varepsilon(a)=o(a^{1+c\gamma}).
 \end{aligned}}                                             \label{v40:precision-eq}
\end{equation}
The link curvatures are $O(a)$ on that entire interval for every allowed
$\varepsilon$. These equivalences concern precisely the constructed family,
not all exact Cauchy records or every microscopic error law.
\end{theorem}
\begin{proof}
Combine \cref{v40:profile-growth} with \eqref{v40:log-errors}. Relative to
the lower profile bound, the error is at most
$Ca e^{(M-\gamma)s}\le Ca^{1-cM}$, which tends to zero uniformly up to
$S_a$. Triangle and reverse triangle inequalities give
\eqref{v40:dynamic-secants}. Its upper bound controls every earlier time;
its lower bound at $S_a$ controls the supremum from below. Thus the supremum
of the difference from the base is comparable to
$|q|e^{\gamma S_a}=|q|a^{-c\gamma}$. Use the $O(a)$ reference bound
\eqref{v40:base-log} and $q=\varepsilon/a$. If both observations are
$O(a)$, their difference is $O(a)$; conversely that bound on their difference
suffices. The same argument with convergence to zero proves the second
line. This establishes necessity as well as sufficiency without a lower
bound on a possibly cancelling literal-link difference.
\end{proof}

The additional factor relative to V39 is
$e^{-\gamma S_a}=a^{c\gamma}$. It is a dynamical precision cost on a
verified branch, not a proposed rescaling of the physical field or clock.
Neither $c_0$ nor the common parameter radius is given an unverified numerical
value.

\section{Initially amplitude-small geometry need not stay amplitude-small}
\label{v40:example}

\begin{corollary}[A populated shrinking-time loss of amplitude-relative control]
\label{v40:departure}
Fix any $c$ admitted above and take the actual V39 perturbation
$\varepsilon(a)=a^2$, equivalently $q(a)=a$. Its initial metric and link
curvatures are $O(a)$, and both histories are exactly constrained. At
$t_a=ca^3\log(1/a)$,
\begin{equation}
 \boxed{\|\mathcal R(a,t_a;a)\|_h\asymp a^{1-c\gamma},\qquad
       \|\mathcal R(a,t_a;a)\|_h\to0,\qquad
       a^{-1}\|\mathcal R(a,t_a;a)\|_h\to\infty.}     \label{v40:loss}
\end{equation}
All literal-link curvatures are $O(a)$ throughout. The initial record changes
satisfy $\Delta X=O(a^4)$ and $\|\Delta\lambda\|\asymp a^2$.
More precisely the nonzero vector limit is
\begin{equation}
 a^{-1+c\gamma}\bigl(\mathcal R_a(S_a,a)-\mathcal R_a(S_a,0)\bigr)
           \longrightarrow (\ell k)\mathcal Ev\ne0.             \label{v40:vector-limit}
\end{equation}
This is a statement about actual nonlinear original-action histories, not a
claim that the slow frozen eigenvector stays invariant during the motion.
\end{corollary}
\begin{proof}
V39 gives the initial estimates at $q=a$. The dynamical secants give a
terminal difference comparable to $ae^{\gamma S_a}=a^{1-c\gamma}$.
Because $cM<1/4$ and $M$ can be taken larger than $\gamma$, the exponent
lies strictly between zero and one. The base is $O(a)$ and hence negligible
at this scale, proving \eqref{v40:loss}. For the vector limit use
$\mathcal U(a)=ak+O(a^2)$, \eqref{v40:dominant-projector}, and
\eqref{v40:log-errors}; the normalized error is at most
$Ca^{1-c(M-\gamma)}\to0$. Nonvanishing is
\cref{v40:source-mode}. The record-size bounds are exactly V39's finite
chart bounds with $\varepsilon=a^2$.
\end{proof}

There is no blow-up asserted in \eqref{v40:loss}: its curvature tends to zero.
The failure is of a uniform bound proportional to the initial field amplitude,
on an interval whose physical duration also tends to zero. This does not
contradict the separately endpoint-selected $ca^2$ histories of V33. The
present perturbations are not asserted to satisfy that selection, and no
initial value is changed after the evolution begins.

\section{What is proved, what is inherited, and the remaining selection problem}
\label{v40:status}

The actual slow-direction preparation error is now linked to a certified
rapid expanding source component. Its finite nonlinear propagation is proved
on bounded rapid intervals and on a logarithmically growing rapid interval.
The latter supplies a sharp preparation scale and an actual family that
loses amplitude-relative curvature control despite meeting the initial
$O(a)$ geometry budget. Link curvature alone still does not detect that
loss. The numerical acquisition is the new left-mode pairing and separated
dominant spectral group; the analytic acquisition is propagation of the
previous finite parameter family with a uniform parameter remainder.

A useful next positive target is the intersection of a geometry-sensitive
initial tube with the actual unstable-cancelling selection, retaining its
canonical input and all original observations. The source pairing here shows
that this extra incidence is nontrivial; a slow frozen label does not keep a
perturbation out of the rapid expanding sector. The eleven-parameter endpoint
selection of V32--V33 cannot be attached to the perturbed records without
solving its changed source equations. No invariant or attracting preparation
graph, order-one physical lifespan, refinement-uniform result, or Einstein
continuum limit is established by this body. V33 remains the longest proved
physical interval for its own selected family.

\paragraph{Verification boundary.}
The new portable checker reassembles the actual $K,D,k$ in the common
V38 coordinate convention, verifies a fresh 28-dimensional left-eigenvector
solve and complete 29-dimensional dominant spectral enclosure, and checks the
strict source-coupling interval with integer/rational arithmetic. The V38
moving-source evaluation is rerun; its V29 root and higher-source interval
arrays remain named inherited inputs. The new block, derivative, observation,
and exponent identities are checked separately in exact finite algebra. Those
algebraic examples are not replacement gravitational sources or integrations
of the nonlinear action. The written evolution proof uses the inherited
analytic stationary chart and original constraint extension. No numerical
trajectory, numerically evaluated common radius, or proof-assistant
formalization is claimed.

The thirty-nine marked preceding mathematical bodies are retained byte for
byte. The package records the predecessor and body hashes, new source data,
accepted intervals, scripts and actually executed verification stages. An
attempt to rerun the complete nested V39 verification was stopped by a runtime
limit; no fresh complete rerun is claimed. The targeted original moving-source
stage and all new acceptance stages are reported separately from archival
verification reports. The standard verification methodology remains the
already cited \cite{v29:Rump}, and the original constraint interpretation is
\cite{v39:GN}; neither reference supplies the model-specific propagation or
coupling result proved here.
% END IMMUTABLE EXACT-ACTION BRIDGE V40 BODY


% Future iterations append here.
% BEGIN IMMUTABLE EXACT-ACTION BRIDGE V41 BODY
\clearpage
\renewcommand{\thesection}{\arabic{section}}
\renewcommand{\theHsection}{bridgeV41.\arabic{section}}
\setcounter{section}{330}
\section{Retain the slow displacement while selecting its expanding response}
\label{v41:context}

The precision loss in V40 is a theorem for an uncorrected one-parameter
family. It is not a theorem that every nearby exactly constrained record
must lose its amplitude-relative geometric bound. Conversely, the endpoint
selection in V33 applies to its own initial-rate family; it cannot simply be
attached to V40's records. We construct the missing joint family here.
The first step is an anchored initial-value chart: a change $aI_Wd$ of the
\emph{actual initial weak rate} costs only $O(a^3|d|)$ in the physical initial
records. All four multiplier means are kept at their original input values.
The chart equals the input record when $d=0$. This stronger cost uses both
the weak column order of the symmetric Hessian and the triangular
cubic-mean control; a generic implicit-function estimate loses one power.

Apply that chart to V40's initially geometry-safe scale
$\varepsilon=a^2\vartheta$, equivalently $q=a\vartheta$ in its notation.
An eleven-dimensional endpoint selection then cancels the expanding part of
its actual response, including the higher nonlinear forcing. The result is a
family of original histories on $0\le t\le ca^2$, with metric and literal
curvatures $O(a)$, uniformly for $\vartheta$ in a fixed interval. Its leading
slow initial displacement is not set to zero: its multiplier means are
preserved exactly. The additional record correction is $O(a^3(|\vartheta|+a))$.
These are different initial records, not a reset or a modification of the
physical flow.

All action-specific statements have $N=3$, $h=1/3$, and the original V29
root $Y_*,\ell_*$. Here $a>0$ is field amplitude, not lattice spacing;
$t=a^3s$ is only an analytical time coordinate. The interval $ca^2$ still
shrinks. No invariant manifold, attracting basin, causal microscopic
selection rule, or spatial continuum limit is asserted.

\paragraph{Inputs and nonduplication.}
The original stationary chart and radial identities, V29's three-mean
submersion, and V30's regular signed chart are inherited. V39 supplies the
finite constrained slow-displacement family and its analytic rate profile.
V40 supplies its actual vector $k=\mathcal U'(0)$ and the nonzero pairing
with an expanding left mode. V32 supplies the full rapid splitting; V33
supplies the fixed Schur shear and mixed Green estimate. The new work is the
anchored weak-rate chart, its $a^3$ record cost with unchanged means, the
parameter-dependent application to the actual perturbation family, and
uniform original-observation secants on the selected interval. The general
truncated Lyapunov--Perron methodology is established prior work
\cite{v41:PR}; no invariant-manifold theorem or new model estimate is imported
from that reference. The named inherited theorems are not independently
reproved in this body.

\section{The actual perturbation at the initially admissible scale}
\label{v41:anchor}

Use the weak coordinates of V40: 26 mean-zero gradients followed by three
constant shifts. Let $I_W$ synthesize them, and let $I_A$ synthesize the
physical active slice. Put $H=A^{-1}L$ at the same reference Gram as
\eqref{v40:blocks}. The rapid operator is $B=-D^{-1}K$, with spectral
projections $p_{\rm u},p_{\rm s},p_{\rm c}$ from V32, expressed in this
fixed weak basis. Write
\begin{equation}
 k_{\rm u}=p_{\rm u}k,\qquad
 k_{\rm n}=(p_{\rm s}+p_{\rm c})k.
 \label{v41:k-split}
\end{equation}
The subscript n denotes the nonexpanding part; its central component is not
assumed to decay. The spaces have dimensions $11,11,7$. In particular,
$\|e^{sB}k_{\rm n}\|\le C$ for $s\ge0$. Constant shifts remain in the
central space.

Let $(X^{\rm in}(a,q),\lambda^{\rm in}(a,q))$ be exactly the V39 chart
based on the analytic transport preparation chosen in V40. Define
\begin{equation}
 X_b(a,\vartheta)=X^{\rm in}(a,a\vartheta),\qquad
 \lambda_b(a,\vartheta)=\lambda^{\rm in}(a,a\vartheta),\qquad
 r_b=\lambda_t(X_b,\lambda_b).
 \label{v41:base}
\end{equation}
The rate on the right is the original normalized rate at that record.
This does not evaluate a trajectory that has already been corrected.

\begin{lemma}[Analytic normalized input rate and unchanged low jets]
\label{v41:base-lemma}
For $|\vartheta|\le\vartheta_0$ and small $a$, the full rate $r_b/a$ has
an analytic extension
\begin{equation}
 \boxed{R_b(a,\vartheta):=a^{-1}r_b
   =r_*+I_W\vartheta k+a\,\rho_b(a,\vartheta),
       \qquad \rho_b(a,0)=0.}
 \label{v41:Rb}
\end{equation}
Its fixed parameter derivatives are bounded on a smaller closed rectangle.
The records satisfy all original constraints, and
\begin{align}
 \|X_b(a,\vartheta)-X_b(a,0)\|&\le Ca^4|\vartheta|,\notag\\
 \lambda_b(a,\vartheta)-\lambda_b(a,0)
           &=a^2\vartheta\widetilde\eta_a,
          \qquad\widetilde\eta_a=c_H+O(a),\quad |c|=1.
 \label{v41:base-size}
\end{align}
In particular the canonical coefficients of orders one and two are
independent of $\vartheta$. The lapse mean is one; the shift mean is not
forced back to zero.
\end{lemma}
\begin{proof}
V39 gives $r_{b,A}=a r_{A,1}(a,q)$ and
$r_{b,W}=\mathcal U(q)+a r_{W,1}(a,q)$, with analytic functions and
$\mathcal U(0)=0$. At $q=0$ the full rate is exactly $ar_*$. Substitute
$q=a\vartheta$. The identity
\[
 a^{-1}\mathcal U(a\vartheta)
     =\vartheta\int_0^1\mathcal U'(ta\vartheta)\,dt
     =\vartheta k+O(a\vartheta^2)
\]
gives \eqref{v41:Rb}, also through $a=0$. The other parameter differences
are $O(a|\vartheta|)$ after division by the relevant rate factor. Analyticity
on a slightly larger rectangle bounds all required derivatives.
The exact V39 chart is
$X^{\rm in}=X_a+aq\widetilde\xi_a+a^3J_cv(a,q)$ and
$\lambda^{\rm in}=\lambda_a+aq\widetilde\eta_a$, where
$\widetilde\xi_a=O(a^2)$ and $v(a,q)=O(q^2)$. Its substitution gives
\eqref{v41:base-size}; $a^3q^2=O(a^5\vartheta^2)$ is smaller than the
stated bound. Its original constraint and normalization identities are
retained.
\end{proof}

The input records are exactly those for which V40 shows loss of the $O(a)$
metric bound on a logarithmic initial layer when $\vartheta\ne0$ is fixed.
Nothing below changes that conclusion for those uncorrected records.

\section{An anchored rate chart with a three-power record cost}
\label{v41:retarget}

Let $\mathcal E_0$ be the 104-dimensional four-mean-zero multiplier space
and $\mathcal H_0$ the 107-dimensional mean-lapse-zero space. Use $P$ for
the complete constraint covector, $\mathsf B=D_XP F$ for its canonical
drift, and $M=D_\lambda P$ for the symmetric Hessian. All equations use
one common physical pairing. Weak rate changes $I_Wd$ lie in
$\mathcal H_0$, and are allowed to have constant-shift components.

\begin{theorem}[Change the actual weak rate without changing the input means]
\label{v41:retarget-theorem}
There are an analytic family of original initial records
$(\widehat X(a,\vartheta,d),\widehat\lambda(a,\vartheta,d))$, for $d$ in a
fixed small ball in the full 29-dimensional weak space, and positive constants
independent of $a,\vartheta,d$ on smaller rectangles, such that
\begin{equation}
 \boxed{\begin{gathered}
 P(\widehat X,\widehat\lambda)=0,\qquad
 \widehat\lambda_t(a,0;\vartheta,d)
                      =a\{R_b(a,\vartheta)+I_Wd\},\\
 (\widehat X,\widehat\lambda)(a,\vartheta,0)
                         =(X_b,\lambda_b)(a,\vartheta),\\
 \mathsf m\widehat\lambda=\mathsf m\lambda_b,\qquad
 \|\widehat X-X_b\|+\|\widehat\lambda-\lambda_b\|
                              \le Ca^3|d|.
 \end{gathered}}
 \label{v41:retarget-main}
\end{equation}
Here $\mathsf m$ takes all four multiplier means. The displayed rate is the
record's own normalized exact-action rate. The record differences divided by
$a^3$ extend analytically through $a=0$, vanish at $d=0$, and have bounded
fixed parameter derivatives. In particular all first and second initial
coefficients are unchanged by this correction.
\end{theorem}
\begin{proof}
We give the anchoring construction, since existence of a different
initial-rate family is not sufficient. Let $z\in\R^3$ be V29's three
selected canonical profile coordinates, with $z_*$ its certified root and
its five other coordinates held fixed. Its compatible first field $Y(z)$,
normalized tilt $\ell(z)$, and cubic-mean map $\mathcal C(z)$ are analytic;
$Y(z_*)=Y_*$, $\ell(z_*)=\ell_*$, $\mathcal C(z_*)=0$, and
$D_z\mathcal C(z_*)$ is invertible. The first Poisson matrix is invertible
on $\mathcal E_0$ throughout a smaller neighborhood.

For $m\in\mathcal E_0$, prescribe the initial multiplier
\begin{equation}
 \lambda(a,\vartheta,z,m)
   =\lambda_b(a,\vartheta)+a\{\ell(z)-\ell_*\}+a^2m.
 \label{v41:lambda-chart}
\end{equation}
It has exactly the input means. To obtain the canonical chart, take the
original initializer with free canonical data of leading coefficient $Y(z)$.
Use the remaining free coordinates of $X_b$ as parameter data. Equivalently,
start with $X_b+a\{Y(z)-Y_*\}$ and solve its original constraint rows in
the fixed normal directions, with zero correction at $z=z_*,m=0$.
The nonconstant linear inverse and the four divided mean equations are the
ones in V29--V30 and V39. They give an analytic $X(a,\vartheta,z,m)$ with
\begin{equation}
 P(X,\lambda)=0,\quad X(a,\vartheta,z_*,0)=X_b,
 \qquad D_zX=O(a),\quad D_mX=O(a^3).
 \label{v41:initializer-bounds}
\end{equation}
For the last, sharper bound, varying $m$ changes the multiplier by $a^2m$.
Since $M=O(a^2)$, the induced uncorrected constraint variation is $O(a^4)$.
The canonical normal inverse loses at most one power of $a$, in the four
mean rows only, as established by the scaled normal submersion in V37--V39.
Thus $D_mX=O(a^3)$. This is uniform on the stated finite parameter chart.
An explicit analytic justification is to multiply the four mean rows of the
normal derivative by $a^{-1}$; its extension is invertible, and the scaled
right side is divisible by $a^3$. This also proves analytic divisibility of
$D_mX$. The equations for the leading field and its four quadratic means are
unchanged. Anchoring the normal solution at $X_b$, where all residuals are
already zero, proves the middle identity of \eqref{v41:initializer-bounds}.

Write $\Pi_0$ for restriction to the four-mean-zero covectors and $R_c$ for
the three constant-shift tests. Define the analytic 107-row map
\begin{equation}
 \mathcal F(a,\vartheta,z,m,d)=
 \begin{pmatrix}R_c\\\Pi_0\end{pmatrix}
 \frac{\mathsf B(X,\lambda)+aM(X,\lambda)
                              \{R_b(a,\vartheta)+I_Wd\}}{a^3}.
 \label{v41:F}
\end{equation}
Its division is legitimate. The first tilt in \eqref{v41:lambda-chart}
is $\ell(z)$, because the input shift displacement is order $a^2$.
The complete quadratic drift vanishes identically by the original Poisson
range equation for $\ell(z)$. All remaining maps are analytic. At $a=0$
the constant-shift rows are $\mathcal C(z)$, independently of $m,d$.
The lower derivative in $m$ is $\mathbb K_1(Y(z))|_{\mathcal E_0}$.
Furthermore $M_2(Y)I_W=0$ for every first field, by the universal weak
first-source identity, so $d$ does not enter the value at $a=0$.
The input family solves \eqref{v41:F} exactly at $z=z_*,m=d=0$.
It follows that the base zero at $a=0$ is $(z_*,0)$ for every small
$\vartheta,d$, with derivative
\begin{equation}
 D_{(z,m)}\mathcal F(0,\vartheta,z_*,0,d)
   =\begin{pmatrix}D_z\mathcal C(z_*)&0\\ *&
                          \mathbb K_1(Y_*)|_{\mathcal E_0}\end{pmatrix}.
 \label{v41:triangular}
\end{equation}
Both diagonal blocks are the actual inherited nonsingular matrices.
The analytic implicit-function theorem yields $z(a,\vartheta,d)$ and
$m(a,\vartheta,d)$; uniqueness gives $z=z_*,m=0$ at $d=0$.

We retain the row weights to obtain the stronger cost. At the anchored
point, the additional target is $MI_Wd/a^2$. In the active and weak rows,
respectively, this is $O(a|d|)$ and $O(a^2|d|)$, because the original
Hessian blocks have weights $a^2,a^3,a^4$. In particular its three constant
rows are $O(a^2|d|)$. The derivative of those top rows with respect to $m$
is $O(a)$, since its value at $a=0$ is zero for every $z$. On a smaller
implicit chart, integration of the exact derivatives between the anchored
and solved points consequently gives
\begin{equation}
 |z-z_*|\le C\{a^2|d|+a|m|\},\qquad
 |m|\le C\{a|d|+|z-z_*|\}.
 \label{v41:coupled-cost}
\end{equation}
The top $z$ and lower $m$ derivative blocks remain invertible; terms linear
in $z-z_*$ in the top equation have already been absorbed in that block.
For sufficiently small $a$, substitution absorbs the remaining $a|m|$ term.
Thus
\begin{equation}
                   z-z_*=O(a^2|d|),\qquad m=O(a|d|).
 \label{v41:controls}
\end{equation}
These orders also follow as analytic divisibilities: at $a=0$ both
corrections vanish, and differentiating the top implicit equation in $a$
there gives $z_a=0$. Their quotients are analytic and vanish at $d=0$.
Parameter derivatives of the quotients are bounded on a smaller rectangle.
Equations \eqref{v41:lambda-chart} and \eqref{v41:initializer-bounds} now
give the $a^3|d|$ record bound, with all means unchanged.

The tested source-rate equation has only a possible constant-lapse residual.
The exact identities $M\lambda=0$, $\langle\lambda,\mathsf B\rangle=0$
and mean lapse one make that residual zero. The full signed chart is regular
for all these records, by V30 and their unchanged first field. Its only null
direction is radial. The difference between the actual normalized rate and
$a(R_b+I_Wd)$ is in that null line and has lapse mean zero, so it is zero.
Each positive-amplitude point therefore has its own original local history.
No target rate is imposed at later times.
\end{proof}

The cost in \eqref{v41:retarget-main} is for weak targets. For an arbitrary
active target the leading term $M_2r$ need not vanish, and the same argument
does not give this cost. The canonical correction is generally larger than
V39's original $O(a^4|\vartheta|)$ field displacement; it is nonetheless of
order three and preserves the first two canonical coefficients.

\section{The original row estimates on the joint initial-data family}
\label{v41:rows}

Restrict $d$ to $E_{\rm u}=\Ran p_{\rm u}$. For each initial record from
\cref{v41:retarget-theorem}, use the original scaling and fixed shear
\begin{equation}
 t=a^3s,\quad X=\widehat X+a^4\xi,\quad
 \lambda=\widehat\lambda+a^4(I_A\eta_A+I_Wy),\quad
 x=(\xi,\zeta_A),\quad\zeta_A=\eta_A+aHy.
 \label{v41:scaling}
\end{equation}
Here $x$ has dimension 402, $y$ dimension 29. Write
$b=(b_x,b_y)=(V_*,r_*)$ and
$J_*(x,y)=(0,C_*x+By)$ as in V33; $J_*b=0$.

\begin{lemma}[Uniform, parameter-differentiated slow and fast rows]
\label{v41:row-lemma}
The same off-constraint analytical extension as V30 has exact equations
\begin{equation}
 x'=b_x+F_s(a,\vartheta,d,x,y),\qquad
 y'=b_y+\vartheta k+d+C_*x+By+F_f(a,\vartheta,d,x,y),\qquad (x,y)(0)=0.
 \label{v41:rows-eq}
\end{equation}
On fixed small parameter neighborhoods and $a^3\|z\|\le\rho$, the bounds
\begin{align}
 \|F_s\|&\le C(a+a^2\|z\|+a^4\|z\|^2),&
 \|D_zF_s\|&\le C(a^2+a^4\|z\|),\notag\\
 \|F_f\|&\le C\{a(1+\|z\|)+a^3\|z\|^2\},&
 \|D_zF_f\|&\le C(a+a^3\|z\|),\qquad
 \|D_z^2F_f\|\le Ca^3
 \label{v41:row-bounds}
\end{align}
hold. They also hold with one parameter derivative, charging bounded leading
parameter derivatives separately; in particular
\begin{align}
 \|\partial_{\vartheta,d}F_s\|&\le
                   C(a+a^2\|z\|+a^4\|z\|^2),\notag\\
 \|\partial_{\vartheta,d}F_f\|&\le
                   C\{a(1+\|z\|)+a^3\|z\|^2\},&
 \|\partial_{\vartheta,d}D_zF_f\|&\le C(a+a^3\|z\|).
 \label{v41:param-bounds}
\end{align}
The canonical part of $D_zF_s$ is $O(a^3)$. These are estimates of the
unchanged maps with a new initial parameter, not a new constitutive writer.
\end{lemma}
\begin{proof}
By \cref{v41:base-lemma,v41:retarget-theorem}, the normalized initial
arguments are $q^0=(\widehat X/a,(\widehat\lambda-e)/a)=q_*+O(a)$,
with $D_\vartheta q^0=O(a)$ and $D_dq^0=O(a^2)$. The actual initial
rate divided by $a$ is $r_*+I_W(\vartheta k+d)+O(a)$, with bounded
fixed derivatives. Its active leading row is $r_{*,A}$.
The leading signed Gram and the leading differentiated source are unchanged,
since $Y_*,\ell_*$ are unchanged. All original coefficient derivatives stay
bounded on a smaller normalized chart.

For precision, the multiplier difference equation is exactly V33's
\eqref{v33:exact-rate}, with $r_d$ replaced by the actual
$R_b(a,\vartheta)+I_Wd$, and with its integral difference quotients centered
at $q^0$. The original moving-Hessian term against that initial rate remains.
Apply the same fixed Schur row identity
$[I\ aH(a,q)]E_a\Gamma^{-1}E_a=[a^2A(a,q)^{-1}\ 0]$.
Its error after freezing at $H$ has factor
$a\{H-H(a,q)\}=O(a^2+a^4\|z\|)$. Thus the slow source difference has
size $O(a^2\|z\|+a^4\|z\|^2)$; the initial active-rate error and shear
contribute $O(a)$. The canonical source has initial error $O(a)$ and
state differential $O(a^3)$. These are the first row of
\eqref{v41:row-bounds}.

The weak leading equation is $b_y+\vartheta k+d+C_*x+By$.
All remaining terms have the same weights as V33: initial error $O(a)$,
linear state error $O(a\|z\|)$, and second state error $O(a^3\|z\|^2)$.
Each extra state differentiation of a coefficient evaluated at
$q^0+a^3\mathcal T_a^{-1}z$ has factor $a^3$.
Differentiating the integral difference quotients proves the stated first
and second state estimates. A parameter derivative of the initial argument
has at least a factor $a$; differentiation of the bounded initial rate only
hits the explicitly suppressed initial-rate and moving-Hessian terms.
This gives \eqref{v41:param-bounds}. The leading $\vartheta k+d$ has
already been removed from the remainder. These checks justify applying the
inherited row calculus on this \emph{joint} family.

The extension has the same homogeneous constraint evolution as
\eqref{v30:P-extension}. Since its new initial records satisfy $P=0$, its
solutions agree with the original action whenever they stay in the chart.
\end{proof}

\section{A parameter-dependent selection of all expanding components}
\label{v41:selection}

Define $\Phi_B(s)=\int_0^s e^{uB}\,du$, without inverting the center.
The nonexpanding reference curve is
\begin{equation}
 z^{\rm ref}_\vartheta(s)
    =s b+\vartheta\,(0,\Phi_B(s)k_{\rm n}).
 \label{v41:reference}
\end{equation}
It solves $z'=b+(0,\vartheta k_{\rm n})+J_*z$, starts at zero, and has
$\|z\|\le C(1+s)$ and uniformly bounded first and second derivatives.
The full unstable projection is
$P_{\rm u}(x,y)=(0,p_{\rm u}y+B_{\rm u}^{-1}C_{\rm u}x)$.
It annihilates \eqref{v41:reference}; projecting only its weak coordinate
would not be valid on a general displacement.

\begin{theorem}[Joint original-action selection on the $ca^2$ interval]
\label{v41:selection-theorem}
There are constants $c,a_0,\vartheta_0,C>0$ such that, for
$0<a<a_0$, $|\vartheta|\le\vartheta_0$, one can select an initial weak-rate
control $d_a(\vartheta)\in E_{\rm u}$ with
\begin{equation}
 \boxed{d_a(\vartheta)=-\vartheta k_{\rm u}+\delta_a(\vartheta),\qquad
  \|\delta_a\|+\|\partial_\vartheta\delta_a\|\le Ca.}
 \label{v41:selected-d}
\end{equation}
The exactly constrained records in \eqref{v41:retarget-main} at this control
have original normalized histories throughout $0\le t\le ca^2$.
In \eqref{v41:scaling}, write
$z=z^{\rm ref}_\vartheta+w_a$, $w_a=(u_a,v_a)$ and $S_a=c/a$.
They satisfy
\begin{equation}
 w_a(0)=0,\qquad P_{\rm u}w_a(S_a)=0,
 \label{v41:endpoint}
\end{equation}
and, uniformly for $s\le S_a$,
\begin{equation}
 \boxed{\begin{aligned}
 \|u_a(s)\|&\le Ca(1+s),&\|u_a'(s)\|&\le Ca,\\
 \|v_a(s)\|&\le Ca(1+s)^2,&\|v_a'(s)\|&\le Ca(1+s).
 \end{aligned}}
 \label{v41:w-bounds}
\end{equation}
The same bounds hold for their first $\vartheta$ derivatives. Every original
constraint is preserved. Constants are independent of amplitude and this
parameter, but not asserted independent of spatial cutoff.
\end{theorem}
\begin{proof}
Set $d=-\vartheta k_{\rm u}+\delta$ and substitute
$z=z^{\rm ref}_\vartheta+w$ in \eqref{v41:rows-eq}. The remaining equation
is precisely
\begin{equation}
 w'=J_*w+(0,\delta)+f,\qquad
 f_j(s)=F_j(a,\vartheta,-\vartheta k_{\rm u}+\delta,
                                      z^{\rm ref}_\vartheta(s)+w(s)).
 \label{v41:fixed-equation}
\end{equation}
Its linear initial--terminal solver is the Green map in
\cref{v33:Green-theorem}. In particular it returns $\delta\in E_{\rm u}$,
$w(0)=0$, $P_{\rm u}w(S)=0$ and is bounded independently of $S\ge1$ from
\[
 \|f\|_{\mathcal F_S}=\sup\|f_s\|
   +\sup\frac{\|f_f\|}{1+s}+\sup\|f_f'\|
\]
to the weighted norm
\[
 \|\delta\|+\sup\left\{
  \frac{\|u\|}{1+s}+\|u'\|
  +\frac{\|v\|}{(1+s)^2}+\frac{\|v'\|}{1+s}\right\}.
\]
This is an inherited linear lemma, not a claim that a finite-time selected
graph is invariant. Only the expanding block is solved backward; the center
is integrated forward with its bounded semigroup, retaining its constants.

For completeness the new application pays all parameter terms. On a ball
of weighted radius $Ma$ and $a(1+S)\le c_0$, with $Mc_0$ small,
$\|z^{\rm ref}_\vartheta+w\|\le C(1+s)$ and
$\|(z^{\rm ref}_\vartheta+w)'\|\le C$, uniformly in $\vartheta$.
Equations \eqref{v41:row-bounds} give $\|f\|_{\mathcal F_S}\le C_0a$;
choose $M$ larger than twice the Green constant times $C_0$. For two inputs,
use the state and control derivative bounds and differentiate the fast
composition once in $s$. Exactly as a direct chain-rule calculation, its
Lipschitz constant is at most the Green constant times
\begin{equation}
 C\{a(1+S)+a^2(1+S)^2+a^3(1+S)^2+a^4(1+S)^3\}.
 \label{v41:budget}
\end{equation}
The time derivative of the fast forcing uses $D_z^2F_f=O(a^3)$ and
$D_dD_zF_f=O(a+a^3\|z\|)$, not an uncontrolled input acceleration.
The bounded first and second derivatives of \eqref{v41:reference} suffice;
its central rotations do not create an additional secular power.
For $S=c/a$, \eqref{v41:budget} is at most $C'(c+c^2+a)$.
Choose $c$ and then $a_0$ so it is less than $1/2$ and the previous
self-mapping estimates hold. Contraction proves the existence, uniqueness
in this weighted ball, and \eqref{v41:w-bounds}.

The same constants control dependence on $\vartheta$. At fixed $(\delta,w)$,
the derivative of $f_j$ is the sum of $\partial_\vartheta F_j$,
$-D_dF_j k_{\rm u}$, and $D_zF_j\partial_\vartheta z^{\rm ref}$.
Here $\partial_\vartheta z^{\rm ref}=O(1+s)$ and its first derivative is
bounded. With \eqref{v41:param-bounds}, their \emph{forcing norm} is
$O(a)$, including the derivative of the fast composition. The uniformly
contracting fixed-point equation is differentiable; solving its linearized
equation by a Neumann series gives weighted derivative norm at most $Ca$.
This proves the derivative part of \eqref{v41:selected-d} and the stated
parameter bounds without differentiation of an unspecified endpoint root.

Finally $d$ stays in the initial chart after reducing $\vartheta_0$.
The normalized physical displacement is
$a^3\mathcal T_a^{-1}z=O(a^2)$ throughout this interval, and the initial
normalized point is $q_*+O(a)$. Thus all stationary, positive-metric,
positive-lapse and signed-inverse margins persist. The homogeneous original
constraint evolution keeps $P=0$. The solution of the analytical extension
is consequently an original normalized exact-action history on the full
interval claimed. The endpoint condition has selected its initial record;
it is not a boundary force or a reset during the history.
\end{proof}

The control $d_a(\vartheta)$ depends on the chosen horizon. Neither its
uniqueness within the weighted ball nor the contraction proves consistency
between different horizons, a forward-invariant graph, or attraction of
arbitrary initial errors. The constant $c$ is positive by the estimates but
has not been numerically evaluated.

\section{Original observations and a uniform tube of selected records}
\label{v41:geometry}

Let a superscript $\#\vartheta$ denote the selected histories of
\cref{v41:selection-theorem}, and compare them at the same physical time.
All norms below are any fixed spatial Sobolev norms at $N=3$.

\begin{theorem}[Retained slow displacement and stable geometric secants]
\label{v41:geometry-theorem}
The selected family satisfies
\begin{equation}
 \boxed{\sup_{0\le t\le ca^2}\bigl(
 \|\operatorname{Riem}(g_a^{\#\vartheta})\|
 +\|E_a^{{\rm link},\#\vartheta}\|
 +\|F_a^{{\rm sp,link},\#\vartheta}\|\bigr)\le Ca.}
 \label{v41:curvature}
\end{equation}
The cost relative to its own uncorrected V40 input, and its retained shift
mean, are
\begin{align}
 \|\widehat X-X_b\|+\|\widehat\lambda-\lambda_b\|
                         &\le Ca^3(|\vartheta|+a),\notag\\
 \overline{\widehat\lambda^\beta(a,\vartheta)
                 -\widehat\lambda^\beta(a,0)}
                         &=a^2\vartheta\overline{\widetilde\eta_a^\beta}.
 \label{v41:cost-and-mean}
\end{align}
In particular the norm of the displayed mean difference is comparable to
$a^2|\vartheta|$; the slow perturbation has not been set to zero.
The full selected record and observation secants obey
\begin{equation}
 \boxed{\begin{aligned}
 \|\widehat X(a,\vartheta)-\widehat X(a,0)\|&\le Ca^3|\vartheta|,\\
 \|\widehat\lambda(a,\vartheta)-\widehat\lambda(a,0)\|
                                      &\asymp a^2|\vartheta|,\\
 \sup_{t\le ca^2}\|\operatorname{Riem}(g_a^{\#\vartheta})
                         -\operatorname{Riem}(g_a^{\#0})\|
                                      &\le Ca|\vartheta|,\\
 \sup_{t\le ca^2}\bigl(\|\Delta_\vartheta E_a^{\rm link}\|
                         +\|\Delta_\vartheta F_a^{\rm sp,link}\|\bigr)
                                      &\le Ca^2|\vartheta|.
 \end{aligned}}
 \label{v41:secants}
\end{equation}
No lower bound on the corrected metric difference is asserted.
\end{theorem}
\begin{proof}
The chart bound and \eqref{v41:selected-d} give the cost in
\eqref{v41:cost-and-mean}. Its exact mean identity follows from the anchored
mean preservation and \eqref{v41:base-size}, independently of the chosen
control. Its lower bound uses $\overline{\widetilde\eta_a^\beta}\to c$
and $|c|=1$. Parameter differentiation of the chart and the selected control
gives the two initial secants: its $a^3$ divided correction is $C^1$ with
bounded derivatives, while the original multiplier input varies by
$a^2\vartheta\widetilde\eta_a$.

Undo \eqref{v41:scaling}, retaining the shear. From
\eqref{v41:reference} and \eqref{v41:w-bounds},
\begin{equation}
 \begin{aligned}
 X_t&=aV_*+O(a^2),\qquad\lambda_t^A=ar_{*,A}+O(a^2),\\
 \lambda_t^W&=a\{r_{*,W}+\vartheta e^{sB}k_{\rm n}\}
                                      +O(a^2(1+s)).
 \end{aligned}
 \label{v41:physical-rates}
\end{equation}
The same remainder bounds hold after one $\vartheta$ derivative. The
nonexpanding semigroup is bounded and $a(1+S_a)$ is bounded. Thus all the
rates are $O(a)$; their parameter differences are $O(a|\vartheta|)$ in
the weak row and $O(a^2|\vartheta|)$ in the active and canonical rows.
The total physical displacements are $O(a^3)$, so $X,\lambda-e=O(a)$.
The original canonical identity
$X_{tt}=D_XF\,F+D_\lambda F\,\lambda_t$ gives $X_{tt}=O(a)$ and
$\Delta_\vartheta X_{tt}=O(a|\vartheta|)$.

More specifically, integration of the parameter bounds in the displacement
variables gives, on the whole interval,
\[
 \|\Delta_\vartheta X\|\le Ca^3|\vartheta|,\quad
 \|\Delta_\vartheta\lambda\|\le Ca^2|\vartheta|,\quad
 \|\Delta_\vartheta\lambda_A\|\le Ca^3|\vartheta|,
 \quad\|\Delta_\vartheta X_t\|\le Ca^2|\vartheta|.
\]
The first-order term of the ADM metric and all necessary spatial derivatives
have size $O(a)$. Its curvature uses $X_{tt}$ and the first multiplier time
jets, not $\lambda_{tt}$. The original formulas of V21 and V39 therefore
give \eqref{v41:curvature} and the metric secant in \eqref{v41:secants}.
All Christoffel products are retained; here every needed first metric jet
is $O(a)$.

For the literal observation, the stationary spatial connection has bounded
field derivatives, $D_{\lambda_A}z=O(a)$, and
$D_{\lambda_W}z=O(a^2)$, with their analytic differentiated versions on the
normalized chart. Thus $\Delta_\vartheta z=O(a^3|\vartheta|)$, while its
physical time derivative, computed as
$z_t=D_Xz F+D_\lambda z\,\lambda_t$, has difference
$O(a^2|\vartheta|)$. In that formula the potentially largest rate difference,
$\Delta_\vartheta\lambda_t^W=O(a|\vartheta|)$, is multiplied by
$D_{\lambda_W}z=O(a^2)$, and contributes only $O(a^3|\vartheta|)$.
The largest term is instead the $O(a^2|\vartheta|)$ canonical-velocity
difference. The stationary temporal connection has difference
$O(a^2|\vartheta|)$ by its bounded ordinary derivatives. Its original
temporal link formula uses that connection, the spatial-connection velocity
and the literal spatial shifts, not a time derivative of the temporal
connection. Therefore the full literal curvature secant is
$O(a^2|\vartheta|)$. Fixed frame changes obey the same bounds. This proves
the last line of \eqref{v41:secants} for the unchanged observations.
\end{proof}

The selected family is a one-parameter tube \emph{on a boundary-selected
initial-data graph}, not an open ball of arbitrary errors around a single
history. Its canonical correction is $O(a^3)$, whereas the original slow
canonical secant at this scale was $O(a^4)$. This price is explicit.
It leaves the first two canonical coefficients unchanged, but generally
changes the third; it must not be described as the identical V40 record.

\section{The actual source fixes the nontrivial correction}
\label{v41:source-effect}

\begin{corollary}[A nonzero compensated response, not automatic robustness]
\label{v41:nontrivial}
Let $\ell$ be V40's certified left eigenfunctional, with its retained
normalization. Uniformly for small nonzero $\vartheta$,
\begin{equation}
 \boxed{\frac{\ell(d_a(\vartheta)-d_a(0))}{\vartheta}
                 =-\ell k+O(a),\qquad 490<-\ell k<640.}
 \label{v41:nontrivial-eq}
\end{equation}
The expanding part of the difference of selected initial rates is $O(a^2
|\vartheta|)$, whereas the uncorrected input has leading expanding difference
$a\vartheta k_{\rm u}$. In particular the correction is not identically zero.
\end{corollary}
\begin{proof}
A left eigenfunctional at an expanding eigenvalue annihilates the stable and
central spectral spaces; hence $\ell p_{\rm u}=\ell$. In particular
$k_{\rm u}\ne0$ by the original source certificate. Integrate
$\partial_\vartheta\delta_a=O(a)$ in \eqref{v41:selected-d} to obtain
\[
 d_a(\vartheta)-d_a(0)=-\vartheta k_{\rm u}+O(a|\vartheta|).
\]
Apply $\ell$ and use V40's actual-source interval. The exact initial rate
identity in \eqref{v41:retarget-main}, together with \eqref{v41:Rb}, gives
\[
 \widehat\lambda_t(a,0;\vartheta)-\widehat\lambda_t(a,0;0)
   =aI_W\vartheta k_{\rm n}+O(a^2|\vartheta|).
\]
Its expanding projection has the claimed order. This concerns an exact
initial-data selection, not a subtraction from the vector field.
\end{proof}

For fixed small $\vartheta\ne0$, V40's uncorrected records have
$\varepsilon=a^2\vartheta$ and lose the $O(a)$ metric bound already by a
sufficiently short $a^3\log(1/a)$ horizon. The records selected here retain
that same leading multiplier-mean displacement, differ by at most
$Ca^3(|\vartheta|+a)$ in additional initial coefficients, and retain the
$O(a)$ metric bound up to $ca^2$. There is no contradiction: V40's
necessary-and-sufficient precision theorem held on its \emph{uncorrected}
curve. The new correction leaves that curve while preserving the specified
slow displacement and the original action.

\section{Scope, next target, and verification}
\label{v41:status}

The intersection proposed at the end of V40 has now been constructed in one
concrete sense: an amplitude-scaled, finite slow-displacement family admits
an exactly constrained higher-order correction whose own original flow stays
geometrically small on the longest existing $ca^2$ interval. The correction
uses all eleven available expanding rate controls and the actual moving-source
vector. It does not require the entire physical perturbation to be refined
from $O(a^2)$ to the smaller uncorrected logarithmic precision scale.

This is not an improved lifespan exponent. The same $ca^2$ scale is now
populated by a joint parameter family with uniform geometric secants and a
proved record-cost estimate. The selection depends on the chosen endpoint,
and is not proved compatible between different endpoints. Consequently it
is not yet an invariant or attracting manifold, nor a causal implementation
by a microscopic renewal law. No claim is made that arbitrary $O(a^3)$ errors
in the selected higher coefficients are harmless. The unknown constants in
the analytic initial chart and Green argument have not been numerically
bounded in this body.

A next substantive step is an endpoint-consistent selection graph, or a paid
error theorem for the higher-coefficient selection itself, retaining the
moving canonical source. A longer physical interval still requires the
central and slow canonical transport left open in V34--V38. All constants
remain fixed-cutoff constants at $N=3$; no refinement-uniform statement,
normalized metric/link identification or exact-action Einstein continuum
theorem follows from this construction.

\paragraph{Verification boundary.}
The new checker verifies the graded weak-target powers, block implicit costs,
full unstable projection and reference-flow formulas, original three-site
weak synthesis and geometric kernel, the differentiated contraction powers,
and the physical pullback and secant bookkeeping in exact rational algebra.
Independent finite polynomial examples are explicitly labelled as algebra
checks, not replacement gravitational actions. The original V40 source/mode
acceptance is rerun from its retained input enclosures. The same enclosure
also rechecks the inherited eleven stable and eleven unstable spectral disk
counts; boundedness of the seven-dimensional center uses the named inherited
signed-inertia theorem, not an assumption about numerical near-zero roots.
No new source extraction, nonlinear numerical initial root, spacetime
integration, numerical chart radius, or proof-assistant formalization is
claimed. Preservation and compilation do not independently validate the
inherited analytical theory. All forty previous marked bodies are retained
byte for byte, with hashes and executable checks in the accompanying package.

\begin{thebibliography}{9}
\raggedright
\bibitem{v41:PR}
C. P\"otzsche and M. Rasmussen,
\emph{Computation of integral manifolds for Carath\'eodory differential equations},
IMA Journal of Numerical Analysis \textbf{30} (2010), 401--430,
\href{https://doi.org/10.1093/imanum/drn059}{doi:10.1093/imanum/drn059}.
This is methodological background for truncated Lyapunov--Perron
constructions. No external invariant-manifold or regulator estimate is
used to replace the joint original-source construction above.
\end{thebibliography}
% END IMMUTABLE EXACT-ACTION BRIDGE V41 BODY

% Future iterations append here.
% BEGIN IMMUTABLE EXACT-ACTION BRIDGE V42 BODY
\clearpage
\renewcommand{\thesection}{\arabic{section}}
\renewcommand{\theHsection}{bridgeV42.\arabic{section}}
\setcounter{section}{338}
\section{Endpoint dependence before an invariant graph is asserted}
\label{v42:context}

The selection of \cref{v41:selection-theorem} depends on its terminal time.
Its existence and uniqueness do not make the selections at two different
terminal times identical. This body compares those \emph{actual original-action}
selections. Their initial controls differ exponentially in the shorter
rescaled horizon. Their original records and curvatures consequently agree
exponentially on an earlier portion of the common physical interval. At
horizons proportional to $a^{-1}$ this is agreement smaller than every
algebraic power of $a$, not exact endpoint independence.

The distinction has a second consequence. A nearby original history that
stays in the same proved slow--fast tube must have an exponentially accurate
expanding-rate control. A conservative exponentially small control-error
budget is also proved sufficient. These are error estimates for the existing
anchored preparation chart, not assertions about all metric-small data.
There is no additional force, projection of the physical field equation,
change of action, or reset of an initial multiplier during a history.

All action-specific conclusions retain $N=3$, $h=1/3$, the original certified
root $Y_*,\ell_*$ and the original link and ADM observations. The positive
parameter $a$ is field amplitude. The notation $t=a^3s$ describes the original
physical times under study. No longer interval than $ca^2$, new source rank,
or spatial-refinement theorem is claimed.

\paragraph{Retained inputs and the new step.}
We use the anchored chart and its divided derivative bounds in
\cref{v41:retarget-theorem}, the exact parameter-dependent rows in
\cref{v41:row-lemma}, and V41's contraction construction. The full unstable
projection and the bounded rapid center are inherited from V32--V33. Their
complete spectral/source acceptance data are unchanged. The new ingredient
is a \emph{discounted difference estimate} with a nonzero terminal mismatch,
followed by its application to two different endpoints of the same nonlinear
source. It controls parameter changes of the initial chart as well as state
changes. Omitting those parameter terms would compare different equations
as though they were the same initial-value problem. Truncated
Lyapunov--Perron approximation is established methodology
\cite{v41:PR}; every difference estimate used below is proved here for the
retained finite system. No external invariant-manifold theorem is imported.

\section{One original preparation chart for all admitted horizons}
\label{v42:horizons}

Let $J=J_*$ act on $z=(x,y)\in\R^{402}\oplus\R^{29}$ by
\[
 J(x,y)=(0,C_*x+By).
\]
Let $P_\mathrm{u}$ be its full unstable projection. With $p_\mathrm{u}$ the
weak unstable projection, its exact expression is
\begin{equation}
 P_\mathrm{u}(x,y)=
 \left(0,p_\mathrm{u}y+B_\mathrm{u}^{-1}C_\mathrm{u}x\right),
 \qquad C_\mathrm{u}=p_\mathrm{u}C_*.
 \label{v42:projection}
\end{equation}
Set $P_\mathrm{n}=I-P_\mathrm{u}$. The inclusion
$\iota:E_\mathrm{u}\longrightarrow\Ran P_\mathrm{u}$ is $d\mapsto(0,d)$;
we identify these two finite spaces when taking norms. In particular
$P_\mathrm{n}\iota=0$. On the unstable space $J_\mathrm{u}$ is invertible.
The inherited spectral margin and center estimate imply, with
$\kappa=1/500$ and a finite $C_J$,
\begin{equation}
 \|e^{-sJ_\mathrm{u}}\|\le C_J e^{-\kappa s},\qquad
 \|e^{sJ}P_\mathrm{n}\|\le C_J(1+s),\qquad s\ge0.
 \label{v42:semigroups}
\end{equation}
For clarity, the exact block exponential is
\[
 e^{sJ}(x,y)=\left(x,e^{sB}y+\int_0^s e^{uB}C_*x\,du\right).
\]
On $\Ran P_\mathrm{n}$ the growing homogeneous term cancels by
\eqref{v42:projection}; the stable integral is bounded and the central
integral grows at most linearly. This proves the second estimate from the
retained weak splitting, without a center inverse. Constants in this body
may depend on this fixed coordinate norm and the fixed cutoff.

Fix $|\vartheta|\le\vartheta_0$ and retain the reference
$z^{\rm ref}_\vartheta$ of \eqref{v41:reference}. Write
\[
 d=-\vartheta k_\mathrm{u}+\delta,\qquad
 z=z^{\rm ref}_\vartheta+w,
\]
and define the \emph{original} remainder along it by
\begin{equation}
 \mathfrak F_a(\vartheta,\delta,w;s)
 =\binom{F_s}{F_f}
 \left(a,\vartheta,-\vartheta k_\mathrm{u}+\delta,
                         z^{\rm ref}_\vartheta(s)+w\right).
 \label{v42:remainder}
\end{equation}
Thus the actual scaled initial-value problem is
\begin{equation}
 w'=Jw+\iota\delta+\mathfrak F_a(\vartheta,\delta,w;s),
 \qquad w(0)=0.
 \label{v42:exact-equation}
\end{equation}
Its physical initial record is the anchored record
$\widehat Q(a,\vartheta,d)=(\widehat X,\widehat\lambda)$ of V41.
The $d$ dependence of this record is already included in $F_s,F_f$.

For $w=(u,v)$ use the inherited weighted norm
\begin{equation}
 \|w\|_{\mathcal W_S}=
 \sup_{0\le s\le S}
 \left\{\frac{\|u(s)\|}{1+s}+\|u'(s)\|
       +\frac{\|v(s)\|}{(1+s)^2}+\frac{\|v'(s)\|}{1+s}\right\}.
 \label{v42:tube-norm}
\end{equation}

\begin{lemma}[Uniform horizon range and the actual unstable endpoint size]
\label{v42:all-horizons}
There are $c_0,a_0,M_0>0$ such that for every $0<a<a_0$ and
$1\le S\le c_0/a$ the V41 construction gives a unique selected pair in its
contraction ball, denoted $(\delta_{a,S},w_{a,S})$, with
\begin{equation}
 P_\mathrm{u}w_{a,S}(S)=0,\qquad
 |\delta_{a,S}|+\|w_{a,S}\|_{\mathcal W_S}\le M_0a.
 \label{v42:selected-family}
\end{equation}
The chart, reference, coefficients and all four initial multiplier means
are the same for every $S$. On these histories,
\begin{equation}
 \|P_\mathrm{u}w_{a,S}(s)\|\le Ca(1+s),\qquad 0\le s\le S.
 \label{v42:unstable-size}
\end{equation}
All constants are uniform in $\vartheta$ on a smaller fixed interval.
\end{lemma}
\begin{proof}
V41's self-mapping and contraction bounds require only $S\ge1$ and
smallness of its displayed budget. That budget is uniform when
$a(1+S)\le a+c_0$. Choose $c_0$ and then $a_0$ small, and use the same
ball constant and initial chart as in that construction. This proves the
horizon-indexed statement; it is not a new independent existence proof.
On the ball, $\|z^{\rm ref}+w\|\le C(1+s)$ and its derivative is bounded,
because $a(1+s)$ is bounded. The original rows give
$\|\mathfrak F_a(s)\|\le Ca(1+s)$. Project the exact equation onto its
unstable space and integrate backward from its \emph{actual} zero endpoint:
\[
 P_\mathrm{u}w(s)=-\int_s^S e^{-(u-s)J_\mathrm{u}}
              \{\iota\delta+P_\mathrm{u}\mathfrak F_a(u)\}\,du.
\]
Equation \eqref{v42:semigroups} bounds this by $Ca(1+s)$. This improves the
bound one would get from the undifferentiated fast part of
\eqref{v42:tube-norm}. No terminal condition is imposed on the nonexpanding
components.
\end{proof}

\section{A discounted Green estimate retains the full terminal mismatch}
\label{v42:Green}

Choose any fixed $0<\omega<\kappa$; $\omega=1/1000$ is allowed. Define
\[
 \|f\|_{\omega,S}=\sup_{0\le s\le S}e^{-\omega s}\|f(s)\|,
 \qquad
 \|(d,v)\|_{\omega,S}^{1}
 =|d|+\sup_{s\le S}e^{-\omega s}(\|v(s)\|+\|v'(s)\|).
\]
The factor $e^{-\omega s}$ is an analytical weight in a comparison norm,
not damping inserted into \eqref{v42:exact-equation}.

\begin{lemma}[Initial control and an inhomogeneous terminal value]
\label{v42:discounted-Green}
For $S\ge1$, continuous $f$, and $b\in\Ran P_\mathrm{u}$, the equations
\begin{equation}
 v'=Jv+\iota d+f,\qquad v(0)=0,\qquad P_\mathrm{u}v(S)=b
 \label{v42:linear-BVP}
\end{equation}
have a unique pair $(d,v)$ with $d\in E_\mathrm{u}$. A constant independent
of $S$ satisfies
\begin{equation}
 \boxed{\|(d,v)\|_{\omega,S}^{1}
       \le C_\omega\{e^{-\omega S}|b|+\|f\|_{\omega,S}\}.}
 \label{v42:Green-bound}
\end{equation}
\end{lemma}
\begin{proof}
In the identified unstable coordinates put
\[
 A_S=\int_0^S e^{-uJ_\mathrm{u}}\,du
      =J_\mathrm{u}^{-1}(I-e^{-SJ_\mathrm{u}}).
\]
No eigenvalue of $e^{-SJ_\mathrm{u}}$ is one, so $A_S$ is invertible.
Its inverse is uniformly bounded for $S\ge1$: for large $S$ use the
convergent geometric series for $I-e^{-SJ_\mathrm{u}}$, and use continuity
on the remaining compact interval. The exact formulas are
\begin{align}
 \iota d&=A_S^{-1}\left(e^{-SJ_\mathrm{u}}b
                   -\int_0^S e^{-uJ_\mathrm{u}}P_\mathrm{u}f(u)\,du\right),
                                                    \label{v42:control-formula}\\
 P_\mathrm{u}v(s)&=e^{-(S-s)J_\mathrm{u}}b
       -\int_s^S e^{-(u-s)J_\mathrm{u}}
                      (\iota d+P_\mathrm{u}f(u))\,du,\notag\\
 P_\mathrm{n}v(s)&=\int_0^s e^{(s-u)J}P_\mathrm{n}f(u)\,du.
                                                    \label{v42:Green-formula}
\end{align}
The first formula makes the unstable component at zero vanish. These formulas
verify the ODE and both boundary conditions, and also prove uniqueness.

For the control, the forcing integral is bounded by
$C(\kappa-\omega)^{-1}\|f\|_{\omega,S}$. For the terminal part of the
unstable solution,
\[
 e^{-\omega s}e^{-\kappa(S-s)}\le e^{-\omega S}.
\]
The backward forcing integral has the same $(\kappa-\omega)^{-1}$ bound.
The nonexpanding integral is bounded after weighting by
\[
 C\int_0^\infty(1+u)e^{-\omega u}\,du
          =C(\omega^{-1}+\omega^{-2}).
\]
Finally the ODE itself bounds $v'$ by $\|J\|\|v\|+|d|+\|f\|$.
Combining these estimates proves \eqref{v42:Green-bound}. In particular the
center's possible linear growth and the full projection in
\eqref{v42:projection} are both retained.
\end{proof}

\section{Exponential agreement of the actual different-endpoint selections}
\label{v42:comparison}

The following comparison is the source-specific step needed before the
linear estimate can be used.

\begin{lemma}[Original nonlinear parameter terms are small in the discounted norm]
\label{v42:nonlinear-difference}
Suppose two solutions of \eqref{v42:exact-equation}, with the same
$(a,\vartheta)$, lie in a fixed enlarged version of the tube
\eqref{v42:selected-family} on $[0,S]$, where $S\le c_0/a$. Put
$v=w_2-w_1$ and $d=\delta_2-\delta_1$. Their remainder difference obeys
\begin{equation}
 \boxed{\|\mathfrak F_a(\delta_2,w_2)
               -\mathfrak F_a(\delta_1,w_1)\|_{\omega,S}
       \le C_\omega a\bigl(|d|+\|v\|_{\omega,S}\bigr).}
 \label{v42:Lip}
\end{equation}
Consequently, with $b=P_\mathrm{u}v(S)$,
\begin{equation}
 \boxed{|d|+\sup_{s\le S}e^{-\omega s}
                  (\|v(s)\|+\|v'(s)\|)
                   \le C_\omega e^{-\omega S}|b|.}
 \label{v42:nonlinear-boundary}
\end{equation}
\end{lemma}
\begin{proof}
The segment joining the two controls and paths remains in a fixed enlarged
parameter and normalized coefficient chart. In that chart
$\|z\|\le C(1+s)$. The original state bounds in
\eqref{v41:row-bounds} therefore imply
\[
 \|D_z(F_s,F_f)\|\le Ca
 \quad\hbox{for }0\le s\le c_0/a.
\]
The parameter bounds in \eqref{v41:param-bounds}, which account for the change
of the original anchored initial record, give
\[
 \|D_d(F_s,F_f)\|\le C\{a(1+s)+a^3(1+s)^2\}.
\]
The latter is not uniformly $O(a)$ without a weight. Its discounted supremum
\emph{is} at most $C_\omega a$. Integrate these two derivative estimates
along the joining segment to obtain \eqref{v42:Lip}. This is where the actual
parameter dependence of the source is paid.

The difference satisfies \eqref{v42:linear-BVP} with this remainder as its
forcing. Apply \cref{v42:discounted-Green}, then \eqref{v42:Lip}, and reduce
$a_0$ so that the resulting $C_\omega a$ is less than $1/2$. Absorb it on
the left. This proves \eqref{v42:nonlinear-boundary} independently of $S$.
There is no assertion that the unweighted nonlinear flow is contracting.
\end{proof}

\begin{theorem}[Two original endpoint selections have exponentially close controls]
\label{v42:horizon-theorem}
Uniformly for $1\le S_1\le S_2\le c_0/a$, put
$d_j=-\vartheta k_\mathrm{u}+\delta_{a,S_j}$ and $w_j=w_{a,S_j}$.
Then
\begin{equation}
 \boxed{|d_2-d_1|+\sup_{s\le S_1}e^{-\omega s}
       \left(\|w_2(s)-w_1(s)\|+\|w_2'(s)-w_1'(s)\|\right)
       \le C_\omega a(1+S_1)e^{-\omega S_1}.}
 \label{v42:horizon-bound}
\end{equation}
The two paths solve the original action with their own anchored initial
records. No path is restarted at the shorter endpoint.
\end{theorem}
\begin{proof}
Restrict the longer selected path to $[0,S_1]$. Its difference from the shorter
path satisfies the original difference equation on that interval, with
terminal mismatch $b=P_\mathrm{u}w_2(S_1)$, since
$P_\mathrm{u}w_1(S_1)=0$. By \eqref{v42:unstable-size},
$|b|\le Ca(1+S_1)$, uniformly in the longer endpoint. Apply
\eqref{v42:nonlinear-boundary}. The two different initial records are encoded
by $d_1,d_2$ in the same coefficient functions, as required by
\cref{v42:nonlinear-difference}.
\end{proof}

The factor $a(1+S_1)$ is obtained from the actual unstable equation, not from
bounding every component by the larger fast-displacement weight. A small
terminal mismatch is not assumed; at $S_1\asymp a^{-1}$ it can be of order
one while the initial control difference is exponentially small.

\section{Agreement of the original records and geometric observations}
\label{v42:observations}

Let $Q_j^0=\widehat Q(a,\vartheta,d_j)$ and let $Q_j(t)$ be its original
history. Every comparison is at the same physical time $t=a^3s$. Write
$\mathcal R_j$ for the full metric-curvature array in the original common
coordinates and frames, and $\mathcal L_j$ for the collection of original
temporal and spatial literal-link curvatures.

\begin{lemma}[A physical comparison that keeps the canonical source]
\label{v42:readout}
For two actual histories in the preceding common tube, set
$D=|d_2-d_1|$ and
$V(s)=\|w_2(s)-w_1(s)\|+\|w_2'(s)-w_1'(s)\|$.
Then
\begin{align}
 \|Q_2^0-Q_1^0\|&\le Ca^3D,&
 \lambda_{t,2}(0)-\lambda_{t,1}(0)&=aI_W(d_2-d_1),
                                                   \label{v42:initial-readout}\\
 \|\Delta\operatorname{Riem}(a^3s)\|&\le Ca\{D+V(s)\},&
 \|\Delta\mathcal L(a^3s)\|&\le Ca^2\{D+V(s)\}.
                                                   \label{v42:readout-bound}
\end{align}
The four initial multiplier means are identical. All norms can be any fixed
spatial Sobolev norms at this fixed cutoff.
\end{lemma}
\begin{proof}
The anchored chart gives the first bound and the exact initial-rate identity.
Undo the same fixed shear for both paths. It gives
\[
 \|\Delta Q(t)\|\le C(a^3D+a^4V(s)),\qquad
 \|\Delta\lambda_t\|\le CaV(s).
\]
The canonical and sheared-active rows have the stronger estimate
\[
 \|\Delta x'(s)\|\le Ca^2\|\Delta w(s)\|+CaD.
\]
Here $\|D_dF_s\|\le Ca$ uniformly for $s\le c_0/a$ follows from
its separate slow-row bound; the fast estimate used in
\cref{v42:nonlinear-difference} would lose this factor. Hence
\[
 \|\Delta X_t\|\le C(a^3V+a^2D),\qquad
 \|\Delta\lambda_t^A\|\le Ca^2(V+D).
\]
For the second inequality retain
$\eta_A=\zeta_A-aHy$ before differentiating; it contributes
$-a^2H\Delta y'$ to the physical active rate.

The actual canonical vector field $F$ and its ordinary field and multiplier
derivatives are bounded on the stationary chart. In
$X_{tt}=D_XF\,F+D_\lambda F\,\lambda_t$, the base rates are $O(a)$.
Differencing therefore bounds $\Delta X_{tt}$ by $Ca(D+V)$.
The original ADM formula uses these second canonical jets, first multiplier
time jets, and their fixed spatial derivatives. Its inverse metrics stay
bounded. Lipschitz estimates of the exact expressions, with all Christoffel
products retained, give the metric part of \eqref{v42:readout-bound}; no
estimate of $\lambda_{tt}$ is required.

For the link part, distinguish the stationary spatial connection $z^{\rm sp}$
from the scaled displacement $z$. Its original response has
$D_Xz^{\rm sp}=O(1)$, $D_{\lambda_A}z^{\rm sp}=O(a)$ and
$D_{\lambda_W}z^{\rm sp}=O(a^2)$, together with the differentiated analytic
bounds used in V41. Thus its value difference is at most
$C(a^3D+a^4V)$. Differentiate the original map along each actual history:
\[
 z_t^{\rm sp}=D_Xz^{\rm sp}X_t+D_\lambda z^{\rm sp}\lambda_t.
\]
The canonical-velocity difference is at most $C(a^2D+a^3V)$; the active
and weak rate contributions are each at most $Ca^3(D+V)$. Differences
of the coefficient derivatives against the $O(a)$ base rates are smaller.
The temporal connection difference is at most $C(a^3D+a^4V)$ by its bounded
stationary derivatives. The original temporal-link expression uses this
connection, its spatial translates, and the spatial-connection velocity,
not a time derivative of the temporal connection. Analyticity of the literal
ordered expressions consequently gives the stated conservative link bound.
All these are estimates of original readouts after undoing the shear.
\end{proof}

\begin{theorem}[A physical terminal buffer gives exponentially agreeing observations]
\label{v42:physical-theorem}
For the two selections of \cref{v42:horizon-theorem},
\begin{align}
 \|Q_2^0-Q_1^0\|&\le C_\omega a^4(1+S_1)e^{-\omega S_1},
                                                    \label{v42:records}\\
 \|\Delta\operatorname{Riem}(a^3s)\|
   &\le C_\omega a^2(1+S_1)e^{-\omega(S_1-s)},\notag\\
 \|\Delta\mathcal L(a^3s)\|
   &\le C_\omega a^3(1+S_1)e^{-\omega(S_1-s)},
             \qquad 0\le s\le S_1.                 \label{v42:curvature}
\end{align}
In particular, fix $0<c_-\le c_+\le c_0$ and $0<\sigma<1$. If
$c_-/a\le S_1\le S_2\le c_+/a$, then
\begin{equation}
 \boxed{\begin{aligned}
 |d_2-d_1|&\le Ce^{-\omega c_-/a},\\
 \|Q_2^0-Q_1^0\|&\le Ca^3e^{-\omega c_-/a},\\
 \sup_{t\le\sigma a^3S_1}\|\Delta\operatorname{Riem}(t)\|
      &\le Ca e^{-\omega(1-\sigma)c_-/a},\\
 \sup_{t\le\sigma a^3S_1}\|\Delta\mathcal L(t)\|
      &\le Ca^2e^{-\omega(1-\sigma)c_-/a}.
 \end{aligned}}
 \label{v42:beyond-orders}
\end{equation}
Thus the endpoint ambiguity is smaller than every power of $a$ in these
specified initial records and on this buffered part of their histories.
\end{theorem}
\begin{proof}
Insert \eqref{v42:horizon-bound} in \cref{v42:readout}. Since $s\ge0$,
the initial-control term is also bounded by the displayed
$e^{-\omega(S_1-s)}$ term. The factor $a(1+S_1)$ stays bounded in the
macroscopic rescaled window. This proves all estimates. For any fixed $m$,
$a^{-m}e^{-\chi/a}\to0$ when $\chi>0$, for example by setting $x=\chi/a$
and comparing the exponential series with a polynomial of degree larger
than $m$. This proves the last statement. It does not assert that these
selected families have convergent power series at $a=0$.
\end{proof}

In physical variables the buffer exponent is
$\omega(T_1-t)/a^3$, where $T_1=a^3S_1$. At $t=T_1$ that buffer is zero:
\eqref{v42:curvature} need then give only the original $O(a)$ scale. The
theorem neither changes the endpoint time nor asserts exponentially small
differences at every point of the closed shorter interval.

\section{Changing the terminal target and testing a computed preparation}
\label{v42:terminal}

There is a related estimate for terminal error that does not postulate a new
source. Keep the initial chart and original equation fixed, but let its
analytical boundary target be $b\in\Ran P_\mathrm{u}$ instead of zero.

\begin{proposition}[Original histories for a small nonzero terminal target]
\label{v42:terminal-existence}
For sufficiently small fixed $0<c_-<c_+\le c_0$, there is a
fixed $b_0>0$ such that for $c_-/a\le S\le c_+/a$ and $|b|\le b_0$
the same anchored chart has a selected original history with
$P_\mathrm{u}w_b(S)=b$. It lies in one common enlarged $O(a)$ tube.
Relative to the zero-terminal selection,
\begin{equation}
 \boxed{|d_b-d_0|+\sup_{s\le S}e^{-\omega s}
       (\|w_b-w_0\|+\|w_b'-w_0'\|)
                         \le C_\omega e^{-\omega S}|b|.}
 \label{v42:terminal-error}
\end{equation}
All original constraints and all four input multiplier means are retained.
\end{proposition}
\begin{proof}
The homogeneous part of \eqref{v42:linear-BVP} generated by its terminal
value is, entirely in the unstable space,
\[
 v_b(s)=(e^{sJ_\mathrm{u}}-I)(e^{SJ_\mathrm{u}}-I)^{-1}b,
 \qquad \iota d_b^{\rm lin}=J_\mathrm{u}(e^{SJ_\mathrm{u}}-I)^{-1}b.
\]
Using \eqref{v42:semigroups}, its contribution to the inherited mixed norm
obeys
\begin{equation}
 |d_b^{\rm lin}|+\|v_b\|_{\mathcal W_S}
                  \le \frac{C|b|}{1+S}.           \label{v42:boundary-ball}
\end{equation}
For instance,
$e^{-\kappa(S-s)}(1+S)/(1+s)\le
(1+S-s)e^{-\kappa(S-s)}$ is bounded. The subtracted identity term and
the undifferentiated weight are estimated in the same way. The active and
canonical components of this terminal contribution are zero.

At $S\ge c_-/a$, \eqref{v42:boundary-ball} is at most $Ca|b|/c_-$.
Take a ball with a fixed margin beyond V41's zero-boundary self-mapping
radius. Choose $b_0$ to fit that margin, and then $c_+$ and $a_0$ to retain
its original strict contraction budget. The same fixed-point construction
therefore provides the claimed history. It is a solution of the original
initial-value equations at the selected initial record. Apply
\eqref{v42:nonlinear-boundary} to it and the zero-terminal selection to
obtain \eqref{v42:terminal-error}. The input chart preserves the means and
constraints identically, independently of the terminal target.
\end{proof}

\begin{corollary}[An a posteriori boundary certificate in the original chart]
\label{v42:aposteriori}
Suppose an independently obtained \emph{exact original history} from the
same anchored chart lies in the enlarged tube on $[0,S]$. Let
$\widetilde b=P_\mathrm{u}\widetilde w(S)$ be its actual terminal residual.
Its control and earlier observations satisfy \eqref{v42:terminal-error} and
\cref{v42:readout} with $b=\widetilde b$, whether or not it was constructed
by an endpoint solver.
\end{corollary}
\begin{proof}
Both histories satisfy \eqref{v42:exact-equation}. The comparison lemma
uses only that equation, the common tube, and the terminal difference.
\end{proof}

This corollary does not turn an approximate numerical trajectory into an
exact one. A discretization residual would enter the forcing norm in
\eqref{v42:Green-bound} and must be bounded separately. Nor does a small
terminal residual by itself prove that a candidate stayed in the required
coefficient chart or evolution tube.

\section{A necessary and a sufficient precision budget for the higher control}
\label{v42:precision}

The same comparison clarifies why endpoint insensitivity does not imply
robustness under arbitrary algebraic-sized errors. Fix $S=c/a$, with
$c>0$ sufficiently small, and use the selected control $d_{a,S}(\vartheta)$.
An error here is a vector $e\in E_\mathrm{u}$ in the \emph{actual weak-rate
control}: the perturbed original initial record is
$\widehat Q(a,\vartheta,d_{a,S}+e)$. Its normalized initial rate differs
exactly by $aI_We$; its record change is at most $Ca^3|e|$.

\begin{theorem}[Exponential control precision for the proved tube]
\label{v42:precision-theorem}
Fix a tube constant larger than that of the selected reference. If the
perturbed original history remains in this tube through $S=c/a$, then
\begin{equation}
                       \boxed{|e|\le Ce^{-\omega c/a}.}
 \label{v42:necessary}
\end{equation}
Conversely there are $\Lambda,c_*>0$, independent of $a$, such that
\begin{equation}
                       \boxed{|e|\le c_*a e^{-\Lambda c/a}}
 \label{v42:sufficient}
\end{equation}
is sufficient for that original history to exist through $S=c/a$ and stay
in the enlarged tube. On that interval its curvature differences from the
selected history obey
\begin{equation}
 \|\Delta\operatorname{Riem}\|\le Ca^2,\qquad
                         \|\Delta\mathcal L\|\le Ca^3.
 \label{v42:precision-geometry}
\end{equation}
The constants and exponents in the two bounds are not claimed optimal or
matching. The necessity is for this tube and chart, not for every history
whose metric curvature happens to be $O(a)$.
\end{theorem}
\begin{proof}
For necessity, its unstable equation at the endpoint gives
\[
 P_\mathrm{u}\widetilde w(S)=J_\mathrm{u}^{-1}
  \{P_\mathrm{u}\widetilde w'(S)-\iota\widetilde\delta
                           -P_\mathrm{u}\widetilde{\mathfrak F}_a(S)\}.
\]
The tube bounds its derivative by $Ca(1+S)$ and the original forcing by
$Ca(1+S)$. Also $|\widetilde\delta|\le Ca$; this is part of the tube
condition, or follows from the equation at zero and its derivative bound.
The bounded inverse $J_\mathrm{u}^{-1}$ therefore bounds the terminal
residual by $Ca(1+S)$. Apply \cref{v42:aposteriori} to obtain
$|e|\le Ca(1+S)e^{-\omega S}$. Substitution $S=c/a$ proves
\eqref{v42:necessary}. No arbitrary matrix spectral projection is substituted
for this actual trajectory comparison.

For sufficiency, compare the two exact initial-value equations up to their
first exit from the enlarged tube. On this segment their state differential
is bounded by a fixed constant, and their control differential is bounded
by $C$ for $s\le c/a$, using the original parameter row estimates. The
inhomogeneous constant control difference is exactly $\iota e$.
Gronwall's inequality, obtained by integrating the scalar differential
inequality for the difference, gives fixed $C,\Lambda$ with
\begin{equation}
 \|\Delta w(s)\|+\|\Delta w'(s)\|
                         \le C|e|e^{\Lambda s}.   \label{v42:IVP-error}
\end{equation}
Under \eqref{v42:sufficient} this is at most $Cc_*a$ on the whole interval.
Take $c_*$ sufficiently small to fit the fixed margin between the reference
and enlarged weighted tubes. A first exit is impossible. The normalized
stationary chart also stays strictly inside its margins, and the usual
finite-dimensional continuation criterion extends the solution through the
endpoint. Its constraints remain zero by the inherited homogeneous
constraint-evolution identity. Apply \cref{v42:readout} with
$D+V\le Ca$ to prove \eqref{v42:precision-geometry}.
\end{proof}

The sufficient record budget induced by this control accuracy is
$Ca^4e^{-\Lambda c/a}$. It is not a lower bound on the necessary accuracy
of every raw canonical coordinate. The map from records to the full rate is
singular in amplitude, and only the anchored control chart has been used.
The result also does not assert the computational cost of achieving the
budget or the precision supplied by any physical renewal instrument.

\section{Why agreement beyond every order is not an invariant selection}
\label{v42:boundary}

There is an elementary analytical diagnostic for the logical boundary.
It is not a gravitational source or a numerical substitute for one. Let
$\gamma>0$ and consider the autonomous augmented equations
\[
 x'=1,\qquad w'=\gamma w+d+ax,\qquad x(0)=w(0)=0,
\]
with an initial constant parameter $d$, selected by $w(S)=0$. They have
\begin{equation}
 \begin{split}
 d_S&=-\frac a\gamma+\frac{aS}{e^{\gamma S}-1},\\
 w_S(s)&=\frac{aS}{\gamma(e^{\gamma S}-1)}(e^{\gamma s}-1)
                                      -\frac{as}{\gamma}.
 \end{split}
 \label{v42:counterexample}
\end{equation}
Direct substitution verifies the original equations and both endpoint
values. Their slow remainder is zero and their fast remainder is $ax$;
these satisfy the separate orders used above. Nevertheless $d_{S_1}\ne
 d_{S_2}$ when $S_1\ne S_2$. At $S=c/a$,
\[
 d_{c/a}+a/\gamma=\frac{c}{e^{\gamma c/a}-1}
\]
is nonzero and smaller than every power of $a$. This example proves that
an estimate of the type \eqref{v42:beyond-orders} cannot by itself turn
finite-endpoint selections into identical ones.

For the original gravitational family there is a further domain issue. At
fixed positive $a$, the present estimates allow only $S\le c_0/a$. They
do not take a limit $S\to\infty$ at that amplitude. Sending $a\to0$
while $S\sim c/a$ is a different limit, on physical intervals tending to
zero. There is also no claim that a time-shifted selected state belongs to
the same anchored one-parameter preparation chart. Thus neither exact
endpoint consistency, invariance, attraction, nor a causal preparation rule
follows from the new comparison. What \emph{is} proved is a quantified
insensitivity to terminal data on earlier physical windows, with the
original observation and its precision cost both retained.

\section{Results, inherited inputs, and the remaining dynamical target}
\label{v42:status}

\paragraph{Proved in this body.}
The original finite-endpoint selections agree exponentially in their initial
expanding controls and initial records. On a fixed fractional prefix of the
shorter $ca^2$ interval their original metric and link curvatures agree
exponentially as well. The same source supports small nonzero terminal
targets and an a posteriori boundary-error comparison. Remaining within the
proved tube requires exponentially accurate higher weak-rate control; a
conservative exponential error budget is sufficient. None of these facts
was an assumption in the comparison proof.

\paragraph{Inherited, not independently re-established here.}
The original stationary chart, root and higher signed-source certificates,
anchored rate realization, homogeneous constraint propagation, rapid
splitting and the $ca^2$ evolution bounds are the named preceding results.
The targeted original source/spectral acceptance is rerun, but its underlying
large source extraction, nonlinear root construction and every earlier
analytical argument are not represented as freshly verified. The exact
checks accompanying this body validate the new algebra, scale bookkeeping,
and preservation; they are not a proof-assistant formalization of the
nonlinear theorems.

\paragraph{Still required.}
There is no exact endpoint-independent graph, no invariant selection on a
neighborhood of evolved states, and no robustness theorem for arbitrary raw
record errors of size $O(a^3)$. The next genuinely different step must address
that graph's time-translation compatibility or the full central and slower
canonical dynamics beyond the proved interval. A limit of finite-horizon
controls at fixed amplitude is not supplied by horizons bounded above by
$c_0/a$. All constants remain fixed-cutoff constants; the physical interval
still shrinks. Refinement-uniform selection, differentiated metric/link
identification and a positive common physical lifespan remain separate.
The raw-action Einstein continuum bridge is not closed.

\paragraph{Source and artifact audit.}
The full V41 terminal body and its stated dependencies were read from the
uploaded source, and its preceding marked bodies were searched for the
present horizon comparison. The newer source was not retrievable through
the indexed file-search results, which returned the distinct old local-writer
lineage; those results were not substituted for this raw-action source.
The external primary publication record confirms the retained methodological
reference \cite{v41:PR}; no external estimate proves the present result.
The companion script checks the complete coupled unstable projection,
inhomogeneous terminal Green identities, boundary-rate formula, explicit
nonidentical flat endpoint example, and physical powers using exact algebra.
Additional finite-model tests, if reported, are marked as diagnostics only.
The previous forty-one marked mathematical bodies are preserved byte for
byte. Any numerical spectral data used in this body are explicitly inherited,
not newly extracted gravitational coefficients.
% END IMMUTABLE EXACT-ACTION BRIDGE V42 BODY

% Future iterations append here.

% BEGIN IMMUTABLE EXACT-ACTION BRIDGE V43 BODY
\clearpage
\renewcommand{\thesection}{\arabic{section}}
\renewcommand{\theHsection}{bridgeV43.\arabic{section}}
\setcounter{section}{347}
\section{An invariant local selection at the same root, rather than another endpoint}
\label{v43:context}

The exponential endpoint comparison of V42 does not define an invariant
set: its horizons are bounded at each positive amplitude, and its anchored
initial curve is not a chart of evolved states. We enlarge the state chart
instead. The unchanged normalized action field has a regular amplitude
extension on a fixed neighborhood. Its eleven expanding directions admit a
locally invariant graph, and the original constrained initial-rate controls
meet that graph transversely. The resulting initial family is chosen once,
without a terminal time. It retains the slow displacement of V41 and the
$ca^2$ physical interval. On that interval it is exponentially close, away
from the terminal layer, to the endpoint-selected histories of V42.

The graph requires a smooth localization of an auxiliary field outside the
original local chart. That localization is not a new gravitational law:
every physical history asserted below stays where the auxiliary field equals
the original constraint-preserving extension, and all its original constraints
are zero. Local graphs need not be unique under different localizations.
Their selected initial records are nevertheless exponentially close on the
already populated time window. We retain both statements; local invariance
is not promoted to a canonical, global or causally implemented selection.

\paragraph{Source audit and scope.}
The source-specific inputs are the original completed constraint equation
\eqref{v30:P-extension}, the V29 root and V30 signed chart, V32's splitting,
the Schur row estimates of V33, and the joint anchored chart of
\cref{v41:retarget-theorem,v41:row-lemma}. V42 supplies the finite-horizon
comparison, not an infinite-horizon limit. The related Native Regenerative
Stationarity V46 file already proves a local invariant-graph method in its
V44 at its own V39 root, with nine expanding directions. We transfer the
method, not that root, its dimensions or its estimates. Here the full
limiting field has eleven expanding directions and a $420$-dimensional
nonexpanding complement; the new points are its actual joint constrained
intersection, the retained $ca^2$ propagation, and quantitative comparison
with the original endpoint selections. The general graph argument is standard
Lyapunov--Perron theory \cite{v41:PR,v43:LL}; it is included with its weighted
norms so that no external smoothness or invariance hypothesis replaces a
needed step. All conclusions have $N=3$, $h=1/3$. The parameter $a>0$ is
field amplitude, not lattice spacing.

\section{A regular field on a fixed normalized neighborhood}
\label{v43:regular}

Let $Q_a^0=(X_a^0,\lambda_a^0)$ be the fixed analytic transport-selected
base of V40--V41 at $\vartheta=d=0$. Multiplier differences lie in the
fixed mean-lapse-zero complement. Let $\mathcal T_a$ be the fixed Schur
shear, with the same reference matrix $H$ as V41:
\[
 \mathcal T_a(\xi,\eta_A,\eta_W)
       =(\xi,\eta_A+aH\eta_W,\eta_W).
\]
It and its inverse are bounded uniformly for small $a$. Define
\begin{equation}
 u=\mathcal T_a(Q-Q_a^0)/a,\qquad
 Q=\Psi_a(u):=Q_a^0+a\mathcal T_a^{-1}u,\qquad t=a^3s.
 \label{v43:chart}
\end{equation}
Thus $u=(x,y)\in\R^{402}\oplus\R^{29}$. The canonical part of $u$ is
unchanged by the shear. This is a chart of all normalized states, not only
of the one-parameter family of prepared initial records.

Write $\mathbb F^{\rm ext}$ for the original canonical field together with
the analytically completed multiplier equation used in
\eqref{v30:P-extension}. For $a\ne0$ set
\begin{equation}
 \mathcal G(a,u)=a^2\mathcal T_a\mathbb F^{\rm ext}(\Psi_a(u)).
 \label{v43:field}
\end{equation}
The word ``extension'' here still denotes the retained completion by the
constraint residual, not a new on-constraint equation.

\begin{lemma}[Regular amplitude field and separate row orders]
\label{v43:regular-lemma}
On a fixed neighborhood of $(a,u)=(0,0)$, \eqref{v43:field} extends real
analytically. With the original $J(x,y)=(0,C_*x+By)$ and $b=(V_*,r_*)$,
\begin{equation}
 \mathcal G(a,u)=Ju+a^3b+\mathcal R(a,u),\qquad Jb=0,
 \label{v43:G}
\end{equation}
where, on a smaller neighborhood,
\begin{align}
 |\mathcal R_s|&\le C(a^4+a^2|u|+|a||u|^2),&
 |D_u\mathcal R_s|&\le C(a^2+|a||u|),\notag\\
 |\mathcal R_f|&\le C(a^4+|a||u|+|u|^2),&
 |D_u\mathcal R_f|&\le C(|a|+|u|).
 \label{v43:row-orders}
\end{align}
In particular $\mathcal G(0,0)=0$ and its full derivative there is
$(0,J)$. The stationary and signed-regular charts remain valid on this
neighborhood for sufficiently small positive $a$. On the original constraint
zero set, \eqref{v43:field} is precisely the original physical flow in the
coordinates \eqref{v43:chart}.
\end{lemma}
\begin{proof}
Write $X_a^0=a\bar X(a)$ and $\lambda_a^0=e+a\bar\chi(a)$. At
\eqref{v43:chart} all normalized stationary coefficients are evaluated at
$(\bar X(a),\bar\chi(a))+\mathcal T_a^{-1}u$. The active and weak source
columns divide analytically by $a$ and $a^2$, respectively, by the universal
weak first-source cancellation and V30's graded chart. The signed matrix
$\Gamma(a,u)$ is invertible on a fixed neighborhood of the retained root.
The completed drift $\overline{\mathsf B}=\mathsf B+\mathscr RP$ has no
constant or linear joint coefficient at $(X,\lambda-e)=0$. Hence
$\overline{\mathsf B}/a^2$ is analytic after this substitution. In the
multiplier part of $u'$ the graded solve reads
\[
 -\operatorname{diag}(aI_A,I_W)\Gamma^{-1}
       \binom{\overline{\mathsf B}_A/a}{\overline{\mathsf B}_W/a^2}.
\]
There is no negative amplitude power. The canonical row $a^2F$ is analytic
as well. This proves regularity on a fixed normalized neighborhood, not
merely along a prepared curve.

For the row estimates, set $u=a^3z$ in the exact displacement rows of V33
or V41 at $\vartheta=d=0$. The identity $u'=a^3z'$ gives
\[
 a^3F_s(a,u/a^3)=O(a^4+a^2|u|+|a||u|^2),\quad
 a^3F_f(a,u/a^3)=O(a^4+|a||u|+|u|^2).
\]
Their state derivatives give exactly \eqref{v43:row-orders}. This
substitution is valid on $|u|\le\rho$, the full domain
$a^3|z|\le\rho$ of the original row estimates. The preceding analytic
factorization makes the formulas legitimate at $a=0$. All bounds are
reduced to a smaller fixed chart if necessary. The initial canonical and
multiplier rates give $\mathcal G(a,0)=a^3b+O(a^4)$; the fixed shear's
extra term is of order four. The factorization also retains the original
indefinite higher weak block. Finally \eqref{v30:P-extension} is homogeneous
in $P$, so initial zero constraints remain zero wherever this chart is used.
\end{proof}

\section{A smooth local graph with the full nonexpanding complement}
\label{v43:graph}

Use the full commuting projections $P_u,P_n=I-P_u$ of V42. In particular
\[
 \begin{gathered}
 P_u(x,y)=(0,p_uy+B_u^{-1}C_ux),\\
 E_u=\Ran P_u,\quad E_n=\Ran P_n,\qquad
 (\dim E_u,\dim E_n)=(11,420).
 \end{gathered}
\]
The inherited estimates are
\begin{equation}
 \|e^{-sJ_u}\|\le C_J e^{-\kappa s},\quad
 \|e^{sJ_n}\|\le C_J(1+s),\quad s\ge0,
 \qquad \kappa=1/500.
 \label{v43:gap}
\end{equation}
The stable directions and the possible zero-frequency Jordan growth are
both retained in $E_n$.

Choose a fixed smooth function $\chi(a,u)$ equal to one on a small ball
about $(0,0)$ and zero outside a slightly larger ball lying in the analytic
chart. Extend
\begin{equation}
 \widetilde{\mathcal G}(a,u)
   =Ju+\chi(a,u)(\mathcal G(a,u)-Ju)
 \label{v43:localization}
\end{equation}
to all finite coordinates. The value of $a$ is constant along every flow.
The derivative of the localized nonlinear part can be made arbitrarily
small, because that part and its full first derivative vanish at the origin.
No trajectory outside $\chi=1$ is asserted to be a gravitational history.

\begin{theorem}[Locally invariant graph at the original mixed-profile root]
\label{v43:graph-theorem}
For such a sufficiently small fixed localization, there is a $C^6$ graph
\begin{equation}
 \mathcal M_a^\chi=\{n+h^\chi(a,n):n\in E_n\},
 \qquad h^\chi(a,n)\in E_u,
 \label{v43:M}
\end{equation}
locally defined near $n=0$, with
\begin{equation}
 h^\chi(0,0)=0,\quad Dh^\chi(0,0)=0,\qquad
 |D_nh^\chi(a,n)|\le C(|a|+|n|),\quad
 |h^\chi(a,0)|\le C a^4.
 \label{v43:graph-bounds}
\end{equation}
It is invariant for the localized field and locally invariant for
$\mathcal G$ while the trajectory stays where $\chi=1$. Its characterization
for the specified localization is independent of a terminal time. It is
not asserted unique among all local invariant graphs or independent of
$\chi$.
\end{theorem}
\begin{proof}
Here are the weighted details. Fix $\nu=\kappa/20$ and, for
$1\le j\le7$, let $\mathcal B_j$ have norm
$\|v\|_j=\sup_{s\ge0}e^{-j\nu s}|v(s)|$.
Write $N(a,u)=\chi(a,u)(\mathcal G(a,u)-Ju)$. For fixed $(a,n)$ define
\begin{align}
 (\mathcal Lv)_n(s)&=e^{sJ_n}n+
          \int_0^s e^{(s-r)J_n}P_nN(a,v(r))\,dr,\notag\\
 (\mathcal Lv)_u(s)&=-\int_s^\infty
          e^{-(r-s)J_u}P_uN(a,v(r))\,dr.
 \label{v43:LP}
\end{align}
The operator norm for differences is bounded by the Lipschitz constant
of $N$ times
\[
 C_J\left((j\nu)^{-1}+(j\nu)^{-2}+(\kappa-j\nu)^{-1}\right).
\]
Every denominator is positive. Shrink the localization once so this product
is less than $1/2$ for all seven weights. Constants in $N(a,0)$ are integrable
in these norms, and the affine free term belongs to each space. The contraction
therefore gives one fixed trajectory and $h^\chi(a,n)=(\mathcal Lv)_u(0)$.
The fixed points agree across weights by uniqueness, since $\mathcal B_1$
is contained in all the others.

For clarity, the smoothness does not follow by differentiating in a single
weighted space without paying products. The localized $N$ has bounded
fixed-order derivatives. A product of variations of total order $j$ has
weight $j\nu$. At derivative order $j$ the linearized equation is
$(I-\mathcal L'_{v})v_j=$ the already determined products of lower
variations, and the inverse in $\mathcal B_j$ has norm at most two.
Taylor's integral remainder, estimated in the next weight, proves continuity
of these derivatives. Induction through $j=6$ uses $\mathcal B_7$ for the
last remainder. Evaluation of the unstable integral at zero is bounded for
all these weights. This proves the stated $C^6$ regularity in $(a,n)$.
At $(0,0)$ the nonlinear derivative vanishes, so its graph is tangent to
$\R\times E_n$. Bounded second derivatives give the slope bound.

The fixed trajectory is exactly the solution whose future is in
$\mathcal B_1$. Conversely, variation of constants for any such solution
makes its unstable terminal term vanish at infinity by \eqref{v43:gap},
giving \eqref{v43:LP}. A finite forward time shift remains in that class,
so uniqueness proves invariance. A finite backward segment of the global
localized flow can also be prefixed without changing membership; thus the
graph is invariant in both directions wherever the local description applies.
This argument shifts states on the full graph, not on the old anchored
initial curve.

The invariance equation is
\begin{equation}
 D_nh\,P_n\mathcal G(a,n+h)=P_u\mathcal G(a,n+h)
 \label{v43:invariance}
\end{equation}
inside $\chi=1$. At $n=0$, $h(a,0)=O(a^2)$ initially follows from
$Dh(0,0)=0$. Also $P_ub=0$. The right side is
$J_uh+O(a^4+(|a|+|h|)|h|)$ and the left side is bounded by
$C|a|\{a^3+(|a|+|h|)|h|\}$. The inverse of $J_u$ is bounded.
Absorbing the small coefficient times $|h|$ proves $h(a,0)=O(a^4)$.
\end{proof}

\section{The graph meets the complete original constraint leaf}
\label{v43:constraints}

The graph in \cref{v43:graph-theorem} is first constructed in all 431
normalized coordinates. Its compatibility with the original 108 independent
constraint equations requires a separate argument.

\begin{theorem}[A codimension-eleven locally invariant constrained manifold]
\label{v43:constraint-manifold}
For every sufficiently small positive $a$, near the selected initial states
defined below, the set
\begin{equation}
 \mathcal S_a^\chi=
 \Psi_a(\mathcal M_a^\chi)\cap\{P=0\}
 \label{v43:physical-manifold}
\end{equation}
is a $C^6$ manifold of dimension $312$, hence codimension eleven in the
$323$-dimensional original normalized constraint leaf. It is locally invariant
under the original normalized action flow. The four constant constraint rows
and the radial energy equation are retained.
\end{theorem}
\begin{proof}
Use the 104 independent nonconstant constraint rows and the four spatial means.
After substituting $Q_a^0+a\mathcal T_a^{-1}(n+h(a,n))$, divide the former
by $a$ and the latter by $a^2$. These maps extend with $C^6$ regularity:
the flat constraint is $LX$, independent of the nearby multiplier, and its
four means vanish identically; the next means are $P_0P_2(Y)$, by the
original flat and preparation identities. No inverse power remains after
Taylor expansion. At $(a,n)=(0,0)$ the divided derivative in canonical
field directions is exactly
\begin{equation}
                  Z\longmapsto (LZ,\ P_0DP_2(Y_*)[Z]).
 \label{v43:constraint-normal}
\end{equation}
The original range-and-four-mean calibration at this same root proves rank
108 for this map. It is not a calibration at the different root in the
related notes.

Every canonical field direction is available in $E_n$. Indeed $n$ can be
parametrized by $(x,y_n)$ as
$(x,y_n-B_u^{-1}C_ux)$, with $y_n$ in the nonexpanding weak space.
Changing $h$ adds only an unstable weak multiplier, not a canonical field.
Lift the 108 fixed canonical normal directions of the inherited initializer
in this way. Their multiplier contribution to an undivided constraint is
$O(a^3)$, because $P_\lambda=M=O(a^2)$ and a coordinate change supplies
a factor $a$. It vanishes in the displayed limiting divided derivative.
Thus the derivative restricted to these lifted directions stays invertible.
The implicit-function theorem removes 108 coordinates from $E_n$, leaving
$420-108=312$. For $a>0$ dividing rows changes neither their zeros nor
ranks. The full normalized leaf has dimension $431-108=323$.

The graph is locally invariant for the extended field. The constraint
residual satisfies \eqref{v30:P-extension}; after multiplying time by $a^3$
it is still a homogeneous equation in $P$. A trajectory starting in
\eqref{v43:physical-manifold} remains in both sets while in the chart.
On $P=0$ that field is exactly the original action flow. This proves local
invariance without solving the constraints again at later times.
\end{proof}

For later use put $v(u)=P_uu-h(a,P_nu)$. In the same local chart,
\begin{equation}
 v'=N_a(u)v,\qquad \|N_a(u)-J_u\|\le C(|a|+|u|).
 \label{v43:normal-factor}
\end{equation}
To verify this exact identity, subtract \eqref{v43:invariance} from
$P_u\mathcal G-D_nhP_n\mathcal G$ and integrate its derivative along the
segment in the unstable coordinate from the graph to $u$. The linear part
is $J_u$ because the projections commute with $J$.
The positive matrix
\[
 W_u=\int_0^\infty e^{-sJ_u^*}e^{-sJ_u}\,ds
\]
satisfies $J_u^*W_u+W_uJ_u=I$. On a smaller chart
$d\langle v,W_uv\rangle/ds\ge c\langle v,W_uv\rangle$.
Thus this manifold is normally repelling in eleven directions, not an
attracting basin. This is a local statement while both the chart and solution
exist; it is not a nonlinear blowup assertion.

\section{The actual initial-rate controls select the graph transversely}
\label{v43:selection}

Retain precisely the anchored V41 records
$\widehat Q(a,\vartheta,d)$, with $d\in E_u$, and set
\[
 U(a,\vartheta,d)=
       \mathcal T_a\{\widehat Q(a,\vartheta,d)-Q_a^0\}/a.
\]
Their previously proved divisibilities give $U=O(a)$ on fixed small
parameter sets. Each record satisfies $P=0$ and has its own exact rate
$r_b(a,\vartheta)+aI_Wd$. Consequently
\begin{equation}
 \mathcal G(a,U)=a^3\mathcal V(a,\vartheta,d),\qquad
 \mathcal V(0,\vartheta,d)=b+(0,\vartheta k+d).
 \label{v43:initial-velocity}
\end{equation}
The quotient $\mathcal V$ is analytic: the original canonical rate is
$aV_*+O(a^2)$, the multiplier rate is the stated exact one, and the fixed
shear changes only the order-$a$ remainder in this quotient.

\begin{theorem}[One initial family, with no chosen terminal time]
\label{v43:selection-theorem}
For the fixed localization $\chi$ there is a $C^5$ control
\begin{equation}
 d_a^\chi(\vartheta)=-\vartheta k_u+\delta_a^\chi(\vartheta),\qquad
 |\delta_a^\chi|+|\partial_\vartheta\delta_a^\chi|\le Ca,
 \label{v43:selected-control}
\end{equation}
on a common small parameter rectangle. The original records
$\widehat Q(a,\vartheta,d_a^\chi)$ lie in $\mathcal S_a^\chi$, preserve
all four input multiplier means, and differ from their own uncorrected
V40 inputs by at most $Ca^3(|\vartheta|+a)$. No terminal time occurs in
this selection. Nearby controls have the quantitative normal displacement
\begin{equation}
 c a^3|e|\le
 |v(U(a,\vartheta,d_a^\chi+e))|\le C a^3|e|.
 \label{v43:transversality}
\end{equation}
This is normal displacement in the normalized chart; its physical scale is
$a^4|e|$, not the $a^3|e|$ upper bound for the entire initial-record change.
\end{theorem}
\begin{proof}
At the actual initial record form the divided normal velocity
\begin{equation}
 \mathcal F(a,\vartheta,d)=
 P_u\mathcal V(a,\vartheta,d)
 -D_nh(a,P_nU)P_n\mathcal V(a,\vartheta,d).
 \label{v43:normal-equation}
\end{equation}
This is a $C^5$ function through $a=0$, without differentiating a singular
rate inverse. Since $U(0,\vartheta,d)=0$ and $Dh(0,0)=0$, its limiting
value is exactly $\iota(\vartheta k_u+d)$. Its derivative in the eleven
control coordinates is the identity inclusion. The implicit-function theorem
therefore supplies a unique local branch satisfying $\mathcal F=0$ and
\eqref{v43:selected-control}, also after a $\vartheta$ derivative.

Vanishing normal velocity at one point is not in general membership in an
invariant graph. Here it is, because \eqref{v43:normal-factor} has an
invertible factor near $J_u$. At an initial record that identity reads
\[
 N_a(U)v(U)=a^3\mathcal F(a,\vartheta,d).
\]
Thus $\mathcal F=0$ implies $v(U)=0$. The original constraints already
hold on the anchored chart, proving membership in
$\mathcal S_a^\chi$. There is no post-processing of an unconstrained
or different-action initial point.

At the selected control, differentiation gives
\[
 D_dv(U)=a^3\{J_u^{-1}\iota+O(a)\}.
\]
The same bound holds uniformly for controls in a fixed small neighborhood:
$D_dU=O(a^2)$ by the anchored cost, $N_a(U)=J_u+O(a)$,
$D_d\mathcal F=\iota+O(a)$, and the term
$(D_dN_a(U))v(U)$ has higher order. Integrate this derivative on the
straight control segment and use the smallest singular value of
$J_u^{-1}\iota$ to obtain \eqref{v43:transversality}. Finally the
anchored V41 cost and exact mean preservation give the asserted physical
record bounds. No analyticity of the selected graph or control is claimed.
\end{proof}

\section{The same invariant selection retains the original \texorpdfstring{$ca^2$}{c a squared} interval}
\label{v43:propagation}

\begin{theorem}[Original histories stay on the graph through the populated interval]
\label{v43:time-theorem}
There are $c,a_0,\vartheta_0,C>0$ such that the records selected in
\cref{v43:selection-theorem} have original normalized histories on
$0\le t\le ca^2$, uniformly for $|\vartheta|\le\vartheta_0$ and
$0<a<a_0$. They remain on $\mathcal S_a^\chi$. With their own initial
records and the original V41 displacement coordinates, write
$z=z_\vartheta^{\rm ref}+w$, $w=(u_s,v_f)$. For $0\le s\le c/a$,
\begin{align}
 |u_s(s)|&\le Ca(1+s),& |u_s'(s)|&\le Ca,\notag\\
 |v_f(s)|&\le Ca(1+s)^2,& |v_f'(s)|&\le Ca(1+s).
 \label{v43:time-bounds}
\end{align}
Consequently the original observations obey
\begin{equation}
 \boxed{\sup_{0\le t\le ca^2}
  \bigl(\|\operatorname{Riem}(g_h)\|+
       \|E_h^{\rm link}\|+\|F_h^{{\rm sp},{\rm link}}\|\bigr)\le Ca.}
 \label{v43:curvature}
\end{equation}
The same selected initial family works for every shorter endpoint. The
number $c$ is not numerically evaluated, and $ca^2$ still tends to zero.
\end{theorem}
\begin{proof}
We prove retention of the interval instead of assuming that a local graph
supplies a lifetime. Let $U_0=U(a,\vartheta,d_a^\chi)$. In the coordinates
of \cref{v41:row-lemma} the absolute normalized state is $U_0+a^3z$.
It belongs to the graph, and $U_0=O(a)$. Up to a first exit from
$|z|<R/a$, its absolute state is $O(a)$ and hence $|D_nh|\le Ca$.
Subtract the graph identity at $U_0$ from the one at $U_0+a^3z$. Its
unstable displacement is a function of $P_nz$, zero at zero, with derivative
bounded by $Ca$. In physical block coordinates this gives
\begin{equation}
 p_uy'=-B_u^{-1}C_ux'+O(a)(|x'|+|p_ny'|).
 \label{v43:graph-rate}
\end{equation}
Here $p_n=p_s+p_c$ is the weak nonexpanding projection; it is not the
full $P_n$. Both are used in their proper spaces.

The exact slow row gives $x'=b_x+O(a)$ in this stopped region. Put
$r_n=p_ny'$ and
$r_n^0=b_{y,n}+\vartheta e^{sB_n}k_n$. Its norm is uniformly bounded.
The identity $C_*b_x+Bb_y=0$ implies
$(r_n^0)'=B_nr_n^0+C_nb_x$. Differentiate the original weak row, with
$a,\vartheta,d$ fixed. Since $d$ is unstable, it disappears under $p_n$;
there is no derivative of an externally reset control. We obtain
\[
 (r_n-r_n^0)'=B_n(r_n-r_n^0)
       +C_n(x'-b_x)+p_nD_zF_f\,z',
 \qquad r_n(0)-r_n^0(0)=O(a).
\]
The forward semigroup of $B_n$ is bounded. The exact row bound gives
$|D_zF_f|\le Ca$ on $|z|\le R/a$. Equation
\eqref{v43:graph-rate} bounds $|z'|$ by $C(1+|r_n|)$. Variation of
constants and the scalar integral Gronwall estimate therefore yield
\begin{equation}
 |r_n-r_n^0|\le Ca(1+s)e^{Ca s},\qquad |z'|\le C
                         \quad(0\le s\le c/a).
 \label{v43:rate-bootstrap}
\end{equation}
Choose $c$ sufficiently small that integration of the last bound cannot
reach $R/a$. The absolute normalized state remains where $\chi=1$ for
small $a$, and the stationary and signed margins persist. This closes the
first-exit argument. The localized trajectory consequently agrees with the
original one on the whole interval; its constraints remain zero by
\eqref{v30:P-extension}.

Integrating $x'-b_x=O(a)$ gives the slow bounds. The error in $r_n$ gives
$O(a(1+s)^2)$ for the weak nonexpanding displacement and
$O(a(1+s))$ for its derivative. The graph relation gives
$|P_uz|\le Ca|P_nz|\le Ca(1+s)$ and
$|(P_uz)'|\le Ca$. After subtracting the reference and undoing its full
projection, the unstable weak error is $O(a(1+s))$, with derivative
$O(a)$. These prove \eqref{v43:time-bounds}.

Undo the fixed shear in the physical multiplier. The exact original rates
are $X_t=aV_*+O(a^2)$, $\lambda_t^A=ar_{*,A}+O(a^2)$, and
\[
 \lambda_t^W=a\{r_{*,W}+\vartheta e^{sB}k_n\}
                                     +O(a^2(1+s)).
\]
They are $O(a)$ throughout $s\le c/a$. The physical displacements are
$O(a^3)$ at the endpoint. The original canonical identity
$X_{tt}=D_XF\,F+D_\lambda F\,\lambda_t$ gives $X_{tt}=O(a)$.
As in V41, the exact ADM expression uses these jets and first multiplier
time jets; the original link expression uses the spatial connection
velocity and temporal connection, not its time derivative. Their fixed-cutoff
analytic bounds prove \eqref{v43:curvature}. No transformed state is
substituted for an original geometric observation.
\end{proof}

\section{Quantitative agreement with endpoint selection and with other localizations}
\label{v43:comparison}

The graph theorem defines a genuinely time-consistent local set for each
fixed localization. It remains to identify its relation with V42's
already constructed histories, rather than call the two selections identical.

\begin{theorem}[The invariantly selected records approximate every long endpoint choice]
\label{v43:comparison-theorem}
Choose $0<\omega<1/500$ and reduce the common $c$ if necessary. For
$1\le S\le c/a$, compare the graph-selected control $d_a^\chi$ with
the original endpoint control $d_{a,S}$ of V42 at the same $(a,\vartheta)$.
Then
\begin{equation}
 |d_a^\chi-d_{a,S}|+
 \sup_{s\le S}e^{-\omega s}
    (|w^\chi-w_S|+|(w^\chi-w_S)'|)
       \le C_\omega a(1+S)e^{-\omega S}.
 \label{v43:comparison-bound}
\end{equation}
The original record and observation differences satisfy
\begin{align}
 |\Delta Q^0|&\le C_\omega a^4(1+S)e^{-\omega S},\notag\\
 \|\Delta\operatorname{Riem}(a^3s)\|
     &\le C_\omega a^2(1+S)e^{-\omega(S-s)},\notag\\
 \|\Delta\mathcal L(a^3s)\|
     &\le C_\omega a^3(1+S)e^{-\omega(S-s)}.
 \label{v43:observation-comparison}
\end{align}
The same bounds, with changed constants, compare the selected records from
any two fixed admitted localizations on a common smaller interval.
\end{theorem}
\begin{proof}
The graph-selected path is in V42's enlarged tube by
\cref{v43:time-theorem}. Its actual unstable terminal residual is
$P_uw^\chi(S)=P_uz^\chi(S)$, because the reference is nonexpanding.
The graph slope and the proved displacement bound give
$|P_uw^\chi(S)|\le Ca(1+S)$. Apply the original nonlinear boundary
comparison \eqref{v42:nonlinear-boundary}, not uniqueness with a falsely
zero terminal value. It includes the changed initial record through the
same $d$-dependent source. This proves \eqref{v43:comparison-bound}.
The anchored cost and \cref{v42:readout} give the three physical estimates.
For two localizations compare each with the same endpoint-selected original
history and use the triangle inequality. Both localized solutions have
already been proved to remain in their original-action regions on that
interval. No assertion about either auxiliary field outside it is required.
\end{proof}

In particular, with $S=c_1/a$ and a fixed prefix $0\le s\le\sigma S$,
$0<\sigma<1$, the initial-record discrepancy is at most
$Ca^3e^{-\omega c_1/a}$, the metric discrepancy at most
$Ca e^{-\omega(1-\sigma)c_1/a}$, and the literal discrepancy at most
$Ca^2e^{-\omega(1-\sigma)c_1/a}$. Local invariance is exact for each
chosen graph. Agreement between different graphs or endpoint choices is
only the quantified statement above, not exact equality.

\section{Local invariance does not canonically determine the graph}
\label{v43:nonuniqueness}

An elementary local system illustrates the remaining distinction without
standing in for the gravitational action. Take $a'=0$ and
\[
 x'=a^3,\qquad y'=\gamma y+ax,\qquad \gamma>0.
\]
For fixed $r>0$, on $|x|<r/2$ the two graphs
\begin{equation}
 h_0(a,x)=-\frac a\gamma x-\frac{a^4}{\gamma^2},\qquad
 h_1(a,x)=h_0(a,x)+
       a^4\exp\left(-\frac{\gamma(r-x)}{a^3}\right),\quad a>0,
 \label{v43:graph-example}
\end{equation}
are both locally invariant. Extend the exponential correction by zero for
$a\le0$. Both graphs are $C^\infty$ near $(0,0)$ and have identical
Taylor coefficients there, but they are unequal at every positive $a$.
Indeed direct differentiation gives $a^3\partial_xh_j=\gamma h_j+ax$;
all derivatives of the added term tend to zero on a smaller $x$-interval,
since an exponential of $-c/a^3$ dominates any fixed inverse power of $a$.
This example concerns uniqueness of a local invariant graph, not a theorem
that two particular gravitational localizations produce different graphs.

For the actual action, a chosen graph is time-consistent on its full local
constraint manifold. A time-shifted state need not belong to the old
one-parameter anchored initial family, but it does belong to the same
$312$-dimensional constrained manifold while it stays in the chart. This
resolves that particular domain issue from V42. It does not provide a
physically preferred choice among the possible graphs, a causal instrument
realizing its initial control, or an attracting basin. Normal repulsion in
\eqref{v43:normal-factor} remains present.

\section{Proof status, nonduplication, and executable checks}
\label{v43:status}

\paragraph{Proved in this body.}
At the same original mixed-profile root, the full normalized field admits a
$C^6$ locally invariant graph with eleven excluded expanding directions.
Its complete original constraint intersection is a $312$-dimensional locally
invariant manifold. The anchored joint slow-displacement/rate family meets
it by an actual $C^5$ initial control with cost $O(a^3(|\vartheta|+a))$,
unchanged multiplier means and a quantified third-order normal incidence.
One such selection retains the original $ca^2$ histories for all shorter
endpoints, without reselecting at any endpoint or later time. Its initial
records and original observations agree exponentially with V42's endpoint
selections, and with the selections from any other fixed admitted localization
on their common interval.

\paragraph{Inherited inputs.}
The root, stationary chart, graded signed source, initial submersion,
anchored rate access, rapid splitting, physical reconstruction formulas and
V42 boundary estimate are the named preceding results. The related Native
V44 invariant-graph construction is methodological input at a different
root; none of its numerical spectral data is transferred. We independently
supply the weighted construction used here and its application to the joint
original chart, but do not claim a new abstract invariant-manifold theorem.
The executable checks validate new algebra, amplitude orders and preservation;
they do not independently formalize the inherited nonlinear theory.

\paragraph{Still required.}
The localization may affect the graph and its initial selection. The graph
is not canonically or operationally selected, is not attracting, and supplies
no uniform robustness to arbitrary raw preparation error. The positive
physical interval still shrinks as $ca^2$ at $N=3$. Its existence does not
bound higher physical-time curvature derivatives or identify the normalized
metric and link coefficients. Spatial refinement, longer central/slow
transport, and a common positive physical lifespan remain unproved. No raw
action term or rate law has been changed on any asserted physical history.
The exact-action Einstein continuum bridge is not closed.

\paragraph{Verification scope.}
The new checker uses exact symbolic and rational examples to verify the full
coupled projector, normalized-field scaling, graph invariance and normal
factor, divided control transversality, constraint row orders and dimension,
nonuniqueness example, and physical comparison exponents. These examples
are labelled algebra checks, not new gravitational source extractions. The
retained original source/spectral acceptance and V42 comparison algebra are
rerun separately with their actual inputs. The preceding forty-two marked
bodies, including their corrections and qualifications, are unchanged.
No numerical nonlinear trajectory, new root radius, new source enclosure or
proof-assistant formalization is represented as performed.

\begin{thebibliography}{9}
\raggedright
\bibitem{v43:LL}
Y. Latushkin and B. Layton, \emph{The optimal gap condition for invariant
manifolds}, Discrete and Continuous Dynamical Systems \textbf{5} (1999),
233--268, doi:10.3934/dcds.1999.5.233. Primary-source context for spectral-gap
and smooth invariant-graph methods; the finite weighted construction and
source-specific application needed here are proved above. No gravitational
estimate is imported from this reference.
\end{thebibliography}
% END IMMUTABLE EXACT-ACTION BRIDGE V43 BODY

% Future iterations append here.

% BEGIN IMMUTABLE EXACT-ACTION BRIDGE V44 BODY
\clearpage
\renewcommand{\thesection}{\arabic{section}}
\renewcommand{\theHsection}{bridgeV44.\arabic{section}}
\setcounter{section}{356}
\section{Separate the decaying rapid fibres from the retained dynamics}
\label{v44:context}

The constrained manifold of \cref{v43:constraint-manifold} suppresses eleven
expanding rapid directions, but still contains eleven decaying rapid directions.
Its dimension $312$ is therefore not the dimension of the remaining central
problem. This body separates those transients on the \emph{same} original
constraint leaf. It constructs a $301$-dimensional locally invariant central
manifold inside V43's manifold, and an eleven-dimensional strong stable
foliation. The foliation matches original histories at the same physical time;
it does not replace their equations by a projected writer.

There are two different uses of this construction. First, the original anchored
rate chart can meet the central manifold using all twenty-two hyperbolic weak
controls. This preserves every initial multiplier mean exactly. Second, an
already selected V43 history has a centre-matched original history on its stable
fibre. For the supplied perturbation family that matching costs only $O(a^4)$
in the original initial record and gives exponentially decreasing differences
in the original observations. The second matching need not preserve the
multiplier means exactly. Its smaller cost must not be attributed to the first,
mean-preserving construction.

A new interval calculation proves that the original moving-source vector $k$
has a nonzero stable component. The removal is consequently nontrivial for the
actual family, not an assumption that its only hyperbolic response is expanding.
The physical interval remains $0\le t\le ca^2$ at $N=3$. Here ``central'' refers
to the limiting rapid splitting. The slower canonical pair of V36--V38 and the
retained imaginary modes are not erased or declared stable.

\paragraph{Inherited inputs and nonduplication.}
We use the actual root and signed chart of V29--V30, the full projections and
semigroup bounds of \cref{v32:full-splitting}, the full twenty-nine-coordinate
rate access of \cref{v41:retarget-theorem}, and V43's regular field, constrained
manifold and populated histories. The comparison with an auxiliary stable
foliation is a standard invariant-manifold method; see \cite{v43:LL,v44:BLZ}.
The weighted constructions needed here are proved below. No new abstract
invariant-manifold theorem or physical symplectic quotient is claimed. The new
source-specific points are the actual stable-source enclosure, the compatible
constraint charts, the twenty-two-control incidence, and the amplitude powers
in stable matching and original observations. In particular an ambient stable
leaf is not assumed to preserve the physical constraints: we first construct
the entire foliation \emph{inside} their exact zero set.

\section{A nonzero stable component of the actual moving source}
\label{v44:source}

Write $B=-D^{-1}K$ for the original $29$-dimensional rapid weak matrix at the
V29 root, in the V38 basis consisting of twenty-six phase gradients followed
by three constant shifts. Keep $k=\mathcal U'(0)$ from V38--V40, in these same
coordinates. All of these are actual coefficient objects already enclosed in
the supplied archive. Put
\[
 \gamma_-=-\gamma_*,\qquad
 -0.09479542359<\gamma_-<-0.09479542352.
\]
Existence and simplicity of this eigenvalue follow from the retained spectral
certificate and the original $D$-skew identity, not from the new pairing alone.

\begin{theorem}[The supplied perturbation also excites a decaying rapid mode]
\label{v44:stable-source}
There is a real left eigenfunctional $\ell_-$, normalized by its coefficient
number $10$ (zero-based) being $+1$, such that
\begin{equation}
 \ell_- B=\gamma_-\ell_-,\qquad
 \ell_- I_H=0,\qquad
 \boxed{-2900<\ell_- k<-2200.}
 \label{v44:coupling}
\end{equation}
In particular $p_s k\ne0$, where $p_s$ is the original eleven-dimensional weak
stable spectral projection. This is a coefficient-space assertion, not a
numerical value for a physical curvature or a new eigenvalue enclosure.
\end{theorem}
\begin{proof}
The portable certificate uses the original enclosures for $D$, the complete
Poisson matrix before weak synthesis, and $k$. The exact constant-shift null
identity of V28--V32 is retained before inversion. Set $G=B^T-\gamma_-I$, let
$j=10$, and let $I_j$ be the other twenty-eight indices. The displayed formulas
\begin{equation}
 (\ell_-)_j=1,\qquad
 (\ell_-)_{I_j}=-G_{I_j,I_j}^{-1}G_{I_j,j}
 \label{v44:left-solve}
\end{equation}
are evaluated with outward-rounded integer intervals. Fixed dyadic inverse
proposals give the acceptance bounds
\begin{equation}
 \|I-R_DD\|_\infty<10^{-6},\qquad
 \|I-R_AG_{I_j,I_j}\|_\infty<10^{-3}.
 \label{v44:inverse-accept}
\end{equation}
These inequalities certify both inverses for every represented matrix and
spectral parameter. The second minor also proves that the genuine left
nullvector at the simple eigenvalue cannot have coefficient $j$ equal to
zero. Thus it admits the normalization in \eqref{v44:left-solve}. Its omitted
row holds because it is that eigenvector, not because an interval residual
containing zero would establish a new root.

A nonzero left eigenvalue annihilates the exact constant right kernel. We use
this identity, with an interval-containment check, to set its last three
coordinates to zero. The complete interval product with the original $k$ is
contained strictly in $(-2900,-2200)$. The certificate stores the exact rational
endpoints; their approximate values are $-2889.588$ and $-2200.708$.
Only the rational comparisons are accepted. Floating inverses propose the
fixed dyadic preconditioners and do not prove either inverse or the pairing.
Finally a left eigenfunctional at $\gamma_-$ annihilates the complementary
spectral subspaces: on each of them $B-\gamma_-I$ is invertible, whereas
$\ell_-(B-\gamma_-I)=0$. Hence $\ell_-k=\ell_-p_sk\ne0$.
\end{proof}

This calculation does not replace $k$ by an independently prescribed stable
perturbation. It reuses its actual root and moving-source enclosures, without
claiming a fresh extraction of the entire gravitational coefficient bank.

\section{A regular chart inside the original constraint leaf}
\label{v44:leaf-chart}

Fix one of the admitted localizations $\chi$ in V43. Use its original normalized
coordinates $u$, original time relation $t=a^3s$, and locally invariant manifold
$\mathcal S_a^\chi$. Let $P_u,P_s,P_c$ denote the \emph{full} spectral
projections of $J$, of ranks $11,11,409$. Their weak restrictions are denoted
$p_u,p_s,p_c$. On an arbitrary full vector $(x,y)$,
\begin{equation}
 P_s(x,y)=(0,p_sy+B_s^{-1}C_sx),\qquad
 P_u(x,y)=(0,p_uy+B_u^{-1}C_ux).
 \label{v44:full-projections}
\end{equation}
In particular replacing $P_s$ by a weak-coordinate projection would change the
matching problem. Inverses here are only on the strictly hyperbolic spaces.

Let $\mathcal B$ be the derivative at $(a,u)=(0,0)$ of the divided original
constraint map used in \cref{v43:constraint-manifold}. It is the map
\begin{equation}
 Z\longmapsto (LZ,P_0DP_2(Y_*)[Z])
 \label{v44:constraint-linear}
\end{equation}
on canonical variations, with zero multiplier part, and has rank $108$.
Therefore $\mathcal BJ=0$. The pure weak hyperbolic spaces are in its kernel.
Every canonical normal variation has a lift in the full $E_c$: the responses
of its slow coordinates in $E_s,E_u$ are cancelled using
\eqref{v44:full-projections}, leaving all weak central coordinates available.
It follows that
\begin{equation}
 E_{c0}:=E_c\cap\ker\mathcal B,\qquad \dim E_{c0}=301,
 \quad \ker(\mathcal B|_{E_s\oplus E_c})=E_s\oplus E_{c0}.
 \label{v44:central-dimension}
\end{equation}
The restriction $J_{c0}$ inherits both two-sided estimates
\begin{equation}
 \|e^{sJ_{c0}}\|\le C(1+|s|),\qquad
 \|J_{c0}e^{sJ_{c0}}\|\le C.
 \label{v44:center-linear}
\end{equation}
No quotient by a constant constraint row has been taken.

\begin{lemma}[A smooth constrained chart with the correct stable and central dimensions]
\label{v44:constrained-chart}
On a fixed neighborhood of $(a,v)=(0,0)$ there is a $C^6$ parametrization
\begin{equation}
 u=\Upsilon_a(v),\qquad v=(s_0,c_0)\in E_s\oplus E_{c0},
 \label{v44:Upsilon}
\end{equation}
of $\mathcal S_a^\chi$ in normalized coordinates. At the origin its derivative
is the inclusion of these spaces, and $\Upsilon_a(0)=O(a^4)$. The induced
field has the form
\begin{equation}
 v'=J_0v+a^3b_0+N_0(a,v),\quad
 J_0=J_s\oplus J_{c0},\quad J_0b_0=0,
 \label{v44:restricted-field}
\end{equation}
where $N_0$ is $C^6$, $N_0(a,0)=O(a^4)$ and
$D_vN_0=O(|a|+|v|)$ on a smaller fixed chart. For $a>0$ every such local
trajectory reconstructs a trajectory of the unchanged normalized action.
\end{lemma}
\begin{proof}
Choose a $108$-dimensional complement $N_c$ to $E_{c0}$ in $E_c$ on which
$\mathcal B$ is invertible. Write the nonexpanding graph argument of V43 as
$n=v+N_cz$. Substitute $n+h^\chi(a,n)$ into the original divided constraints.
Their derivative in $z$ at zero is the nonsingular restriction of
\eqref{v44:constraint-linear}. The parameter-dependent implicit-function theorem
solves $z=z(a,v)$ with $D_vz(0,0)=0$. This is precisely a chart of the actual
intersection, with the same $108$ constraints as in V43.

At $v=0$, $h^\chi(a,0)=O(a^4)$. A pure weak normalized displacement has
$O(a)$ derivative in the divided mean rows and at least that order in the
other divided rows: its physical displacement is $a$ times its normalized
value and $P_\lambda=O(a^2)$. Thus the constraint residual from this graph
height is $O(a^5)$; solving the fixed normal block gives $z(a,0)=O(a^5)$.
The asserted bound for $\Upsilon_a(0)$ follows.

The inverse chart is a fixed linear projection of $P_nu$ onto
$E_s\oplus E_{c0}$, because the normal correction is in $N_c$ and the graph
height is in $E_u$. Its derivative therefore loses no differentiability.
Projecting the original field after substitution gives
\eqref{v44:restricted-field}. The linearization is $J_0$, since the tangent
space is invariant and the nonlinear graph derivatives vanish at zero.
The fixed shift $\Upsilon_a(0)=O(a^4)$ changes its value only by $O(a^4)$.
The coefficient $b$ lies in $\ker\mathcal B$ by differentiating the exact
constraint equation at the original base; it also satisfies $Jb=0$. Hence its
projection $b_0$ is purely central and has $J_0b_0=0$. Bounded second
derivatives prove the displayed differential estimate. Local invariance and
agreement with the original action follow from V43, not from a projection of
the canonical vector field onto these coordinates.
\end{proof}

\section{An invariant central manifold and its contracting fibres}
\label{v44:center-fibres}

The following construction is performed in \eqref{v44:Upsilon}, so its states
already satisfy every original constraint. Choose a smooth localization of
\eqref{v44:restricted-field} outside a smaller chart, keeping the linear part
$J_0v$ and multiplying the remaining part by a smooth cutoff. Shrink its radius
once so that the nonlinear Lipschitz constant is as small as specified below.
Denote this fixed additional choice by $\zeta$. It is not an alteration of
any asserted physical history, all of which will remain in the common region
where the localizations equal one.

\begin{theorem}[Central reduction inside the already constrained selection]
\label{v44:center-theorem}
There is a $C^4$ locally invariant manifold
\begin{equation}
 \mathcal C_a^{\chi,\zeta}
 =\Psi_a\left(\Upsilon_a\{(h_s(a,c),c):c\in E_{c0}\}\right)
 \subset \mathcal S_a^\chi
 \label{v44:C}
\end{equation}
of dimension $301$, codimension $11$ in $\mathcal S_a^\chi$ and codimension
$22$ in the original normalized constraint leaf. Its graph satisfies
\begin{equation}
 h_s(0,0)=Dh_s(0,0)=0,\qquad
 |D_ch_s(a,c)|\le C(|a|+|c|),\qquad h_s(a,0)=O(a^4).
 \label{v44:hs}
\end{equation}
A neighborhood of this manifold in $\mathcal S_a^\chi$ has $C^2$ strong
stable leaves of dimension $11$. For each point $Q$ in that neighborhood
there is a unique matching point $\pi_aQ$ on this chosen central manifold,
with
\begin{equation}
 \pi_a\Phi_t Q=\Phi_t\pi_aQ
 \label{v44:semiconjugacy}
\end{equation}
while the compared trajectories and charts remain in their common domain.
In normalized coordinates, for any fixed $0<\beta<\kappa$ after suitable
shrinking (one may fix $\beta=\kappa/2$, $\kappa=1/500$),
\begin{equation}
 |u(s)-u_c(s)|+|u'(s)-u_c'(s)|
       \le C e^{-\beta s}|u(0)-u_c(0)|.
 \label{v44:tracking}
\end{equation}
Each fibre consists of original constrained states. The retraction is a
comparison between their histories, not an identification of their physical
records, a gauge quotient, or a modified flow.
\end{theorem}
\begin{proof}
We give the required construction, including its time directions. Let
$\widehat N(a,v)$ be the localized $a^3b_0+N_0(a,v)$. For fixed $a,c$,
solve on the whole real line
\begin{align}
 v_c(t)&=e^{tJ_{c0}}c+
       \int_0^t e^{(t-r)J_{c0}}\pi_c\widehat N(a,v(r))\,\dd r,\notag\\
 v_s(t)&=\int_{-\infty}^t
       e^{(t-r)J_s}\pi_s\widehat N(a,v(r))\,\dd r.
 \label{v44:LP-center}
\end{align}
Use weights $\sup_t e^{-j\nu|t|}|v(t)|$, $1\le j\le5$, with
$0<\nu<\kappa/100$. The difference bounds are at most the Lipschitz
constant times
\[
 C\{(j\nu)^{-1}+(j\nu)^{-2}+(\kappa-j\nu)^{-1}\}.
\]
Choose one localization making every bound less than $1/2$. Constants and the
polynomial free central term are admissible in all these norms. The contraction
defines $h_s(a,c)=v_s(0)$. Differentiation through fourth order uses products
in the corresponding successively larger weights; the fifth weight bounds the
last Taylor remainder. The localized field is $C^6$, so these operations are
available. This proves $C^4$ regularity. Time shifts of the unique two-sided
weighted trajectory give local invariance. Tangency and the slope bound follow
from the vanishing nonlinear differential at zero. At $c=0$, use the graph
invariance equation, $\pi_sb_0=0$ and invertibility of $J_s$ exactly as in the
proof of \cref{v43:graph-theorem}; absorption gives $h_s(a,0)=O(a^4)$.

For the foliation, fix a global localized central trajectory $v^c(t)$ and a
small stable vector $\eta$. Seek its forward difference $w(t)$ from another
trajectory by the equations
\begin{align}
 w_s(t)&=e^{tJ_s}\eta+\int_0^t e^{(t-r)J_s}
                      \pi_s\Delta\widehat N(r,w(r))\,\dd r,\notag\\
 w_c(t)&=-\int_t^\infty e^{(t-r)J_{c0}}
                      \pi_c\Delta\widehat N(r,w(r))\,\dd r,\label{v44:fibre-equations}\\
 \Delta\widehat N(r,w)&=\widehat N(a,v^c(r)+w)-\widehat N(a,v^c(r)).\notag
\end{align}
In the norm $\sup_{t\ge0} e^{\beta t}|w(t)|$, the difference bound is
\[
 C\operatorname{Lip}(\widehat N)
       \{(\kappa-\beta)^{-1}+\beta^{-1}+\beta^{-2}\}.
\]
It is less than $1/2$ after another fixed shrinking. Thus the unique solution
obeys $|w(t)|\le C e^{-\beta t}|\eta|$. Its differential equation and bounded
coefficient derivatives give the same bound for $w'(t)$.

We include the parameter issue in using these leaves. Derivatives of a central
trajectory with respect to its initial central parameter grow at most
$C_j e^{j\nu_1t}$ for any suitably small fixed $\nu_1>\nu$; the central
polynomial bound and the small Lipschitz constant give this by the same
weighted integral estimates. In derivatives of $\Delta\widehat N$, any
undecayed derivative of that trajectory multiplies a difference of nonlinear
derivatives evaluated at points separated by $w$. It therefore has an
additional $e^{-\beta t}$ factor. Choose $3\nu_1<\beta$ and use weights
$\beta-\nu_1,\beta-2\nu_1,\beta-3\nu_1$ successively. They are positive and
below $\kappa$. The same contractions and bounded derivatives prove at least
$C^2$ dependence of $w(0)$ on $(a,c,\eta)$.

At $(a,c,\eta)=(0,0,0)$ the initial-point map is
$(c,\eta)\mapsto(\eta,c)$ in the stable/central splitting. Its derivative is
invertible. It is a local diffeomorphism on a smaller fixed chart, so the
leaves are disjoint and have a unique base point. A finite time shift of a
decaying difference still satisfies \eqref{v44:fibre-equations} with its new
stable value. Uniqueness gives \eqref{v44:semiconjugacy}. The chart and its
inverse have bounded first two derivatives, uniformly for small $a$; transporting
the bounds gives \eqref{v44:tracking}. The construction took place entirely in
the constrained coordinates \eqref{v44:Upsilon}. No assumption that an ambient
stable leaf preserves the constraints is needed. As always, physical conclusions
stop when either localized trajectory leaves the original-action chart.
\end{proof}

This result gives attraction \emph{along stable fibres within} V43's selected
manifold. In the complete constraint leaf there are still eleven repelling
normal directions. The finite-amplitude central dynamics may also have slower
growth. No positive energy for that entire dynamics, Hamiltonian quotient, or
canonical choice of either localization is inferred.

\section{All twenty-two hyperbolic rate controls are available on original records}
\label{v44:selection}

The central manifold has now been defined on the same constraint leaf as the
initial records, but its accessibility is a separate calculation. Use the full
weak-rate chart of \cref{v41:retarget-theorem}, not just the unstable restriction
used in V43. Set
\[
 d\in E_{\rm hyp}:=E_u\oplus E_s,\quad\dim E_{\rm hyp}=22,\qquad
 k_{\rm hyp}=(p_u+p_s)k,\quad k_c=p_ck.
\]

\begin{theorem}[Mean-preserving central preparation by actual initial rates]
\label{v44:central-preparation}
For the fixed choices $\chi,\zeta$, there is a $C^3$ control
\begin{equation}
 \boxed{d_a^{\rm cen}(\vartheta)
          =-\vartheta k_{\rm hyp}+O_{C^1_\vartheta}(a)}
 \label{v44:d-center}
\end{equation}
on a common small parameter rectangle, such that the original anchored
records $\widehat Q(a,\vartheta,d_a^{\rm cen})$ lie in
$\mathcal C_a^{\chi,\zeta}$. They satisfy every original constraint, preserve
all four initial multiplier means exactly, and obey
\begin{equation}
 \boxed{\|\widehat Q-Q_b(a,\vartheta)\|
                 \le Ca^3(|\vartheta|+a),\qquad
 \lambda_t(a,0)=a\{r_*+I_W\vartheta k_c\}+O(a^2).}
 \label{v44:center-cost}
\end{equation}
The first two initial coefficients remain unchanged. A normal error $e$ in
these twenty-two control coordinates has normalized normal-coordinate distance from
the central manifold comparable to $a^3|e|$, and physical normal scale
$a^4|e|$. The statement concerns actual initial rates, not rates imposed during
a history.
\end{theorem}
\begin{proof}
Use a local chart of the complete normalized constraint leaf with coordinates
$(u_h,s_h,c)$, where the first two blocks are the limiting full hyperbolic
coordinates and $c\in E_{c0}$. It is obtained by solving the same $108$
canonical normal rows in the full center as in \cref{v44:constrained-chart},
before imposing V43's unstable graph. Their derivative is nonsingular, so
these coordinates are regular through $a=0$.

In this chart V43's manifold is a graph $u_h=g_u(a,s_h,c)$, and the new
central manifold inside it has $s_h=g_s(a,c)$. Both derivatives in the state
variables vanish at $(0,0)$. Define the two normal coordinates
\[
 v_u=u_h-g_u(a,s_h,c),\qquad v_s=s_h-g_s(a,c).
\]
Invariance of the two nested manifolds and integration of the state derivative
along normal segments give an exact triangular factorization
\begin{equation}
 \binom{v_u}{v_s}'=
 \begin{pmatrix}J_u+O(|a|+|u|)&0\\
 O(|a|+|u|)&J_s+O(|a|+|u|)\end{pmatrix}
 \binom{v_u}{v_s}.
 \label{v44:normal-system}
\end{equation}
The entries denote smooth state-dependent matrices, not frozen coefficients.
Both diagonal blocks are invertible on a smaller chart. This factorization is
on the original constraint leaf. It does not assert that the stable normal is
zero off the centre-stable manifold without the lower-left term.

For the anchored point let $U=\mathcal T_a(\widehat Q-Q_a^0)/a$. The exact
rate access gives $\mathcal G(a,U)=a^3\mathcal V(a,\vartheta,d)$, with
\[
 U=O(a),\qquad \mathcal V(0,\vartheta,d)=b+(0,\vartheta k+d).
\]
Evaluate the derivative of the two normal coordinates on this \emph{actual}
velocity and divide by $a^3$. The result is a $C^3$ finite map with limiting
value
\begin{equation}
 \binom{\vartheta p_uk+d_u}{\vartheta p_sk+d_s}.
 \label{v44:control-incidence}
\end{equation}
The canonical coupling in the full projections is included; it vanishes on
$b$ because $Jb=0$. Its derivative in $d\in E_{\rm hyp}$ is the identity.
The implicit-function theorem supplies \eqref{v44:d-center} and its first
parameter derivative. Vanishing of these normal velocities implies vanishing
of both normal positions by the invertible factor in \eqref{v44:normal-system}.
Thus the selected records are on the central manifold, rather than merely
having an arbitrarily prescribed velocity tangent at one point.

The original retargeting chart already solves $P=0$ and preserves its input
means. Its $a^3|d|$ cost proves \eqref{v44:center-cost}; inserting the exact
rate identity and $R_b=r_*+I_W\vartheta k+O(a)$ gives the last assertion.
Differentiating the normal identity at the selected point gives
$a^3\operatorname{diag}(J_u^{-1},J_s^{-1})+O(a^4)$ on the control coordinates.
The same bound holds locally, and integration along a segment gives the
stated two-sided normal transversality. Finite powers of $a$ converting the
normalized and physical charts give its physical scale. There is no inverse
on $E_c$ in this argument.
\end{proof}

\section{The central preparation retains the original physical interval}
\label{v44:time}

\begin{theorem}[Central histories with the same original curvature bounds]
\label{v44:time-theorem}
For the records of \cref{v44:central-preparation}, there are fixed positive
$c,a_0,\vartheta_0,C$ such that their original normalized histories exist on
$0\le t\le ca^2$, for $0<a<a_0$, $|\vartheta|<\vartheta_0$, and remain in
$\mathcal C_a^{\chi,\zeta}$. Put $s=t/a^3$ and use the original anchored
displacement and Schur coordinates $z=(x,y)$ from V41. With
\begin{equation}
 z^{\rm ref}_{\vartheta,c}(s)
   =sb+\vartheta\left(0,\int_0^s e^{rB_c}k_c\,\dd r\right),
 \label{v44:ref}
\end{equation}
the error $w=z-z^{\rm ref}_{\vartheta,c}=(w_s,w_f)$ satisfies
\begin{align}
 |w_s(s)|&\le Ca(1+s),& |w_s'(s)|&\le Ca,\notag\\
 |w_f(s)|&\le Ca(1+s)^2,& |w_f'(s)|&\le Ca(1+s)
                     \qquad (0\le s\le c/a).
 \label{v44:time-errors}
\end{align}
Consequently
\begin{equation}
 \boxed{\sup_{0\le t\le ca^2}
  \bigl(\|\operatorname{Riem}(g_h)\|+
       \|E_h^{\rm link}\|+\|F_h^{{\rm sp},{\rm link}}\|\bigr)\le Ca.}
 \label{v44:time-curvature}
\end{equation}
No endpoint selection, time reparametrization of a compared history, or
longer physical interval is implicit in this assertion.
\end{theorem}
\begin{proof}
The analytic divided equations establishing \cref{v41:row-lemma} use the
exact weak-target property $M_2I_W=0$, not a sign of a rapid eigenvalue.
They therefore give the same slow/fast remainder and differentiated bounds
for $d$ in a fixed small ball in $E_u\oplus E_s$. The exact initial-rate
formula changes only the leading constant weak vector. With
\eqref{v44:d-center} this vector is $\vartheta k_c+O(a)$; the $O(a)$
part is included in the remainder.

Stop at the first exit from $|z|<R/a$. The absolute normalized state is
$U_0+a^3z=O(a)$ there. The two graph slopes in \eqref{v44:normal-system}
are $O(a)$. Differentiating their zero identities along the actual trajectory,
using \eqref{v44:full-projections}, gives, for both signs $\sigma=u,s$,
\begin{equation}
 p_\sigma y'=-B_\sigma^{-1}C_\sigma x'
                  +O(a)(|x'|+|p_cy'|).
 \label{v44:hyper-rate}
\end{equation}
The small terms can be absorbed simultaneously. The original slow row gives
$x'=b_x+O(a)$ in this region. Differentiate the remaining weak central row
and subtract $r_c^0=p_cr_*+\vartheta e^{sB_c}k_c$. Since
$B_cp_cr_*+C_cb_x=0$, its error obeys an equation with homogeneous part
$B_c$, forcing $O(a)(1+|r_c-r_c^0|)$, and initial value $O(a)$.
The semigroup of $B_c$ is bounded in both directions by its proved positive
metric. Variation of constants gives
\[
 |r_c-r_c^0|\le Ca(1+s)e^{Ca s},\qquad |z'|\le C
                    \quad(0\le s\le c/a).
\]
Choose $c$ small enough that integration cannot reach $R/a$. The normalized
state remains well inside both localization regions, and the stationary and
signed charts persist. Ordinary continuation proves existence through the
endpoint; local invariance keeps the states on the central manifold.

Integrating the slow rate error gives the first line of
\eqref{v44:time-errors}. The central rate bound gives the second line after
integration. The hyperbolic graph relations give displacement error
$O(a(1+s))$ and rate error $O(a)$ in the full hyperbolic projection; subtracting
the reference introduces no hyperbolic component because that reference is
central. Undoing the fixed shear gives the same bounds on the physical weak
error. The original rates $X_t,\lambda_t$ are therefore $O(a)$. The original
canonical equation gives $X_{tt}=D_XF\,F+D_\lambda F\,\lambda_t=O(a)$.
Exactly the ADM and stationary-link arguments of V41--V43 then give
\eqref{v44:time-curvature}. All constraints were solved before this propagation
and remain zero under the same original homogeneous constraint equation.
\end{proof}

The central reference need not be constant: $e^{sB_c}k_c$ retains the actual
oscillatory and zero-frequency modes. Removing decaying transients is not a
proof that these remaining modes, or their finite-amplitude coupling, permit
a common positive physical lifespan.

\section{A fourth-order stable matching of the already selected histories}
\label{v44:matching}

The foliation also applies without reselecting the V43 family in its anchored
rate chart. Let $Q_a^{cs}(\vartheta)$ denote precisely that family. Its
centre-matched point $Q_a^{c,\pi}(\vartheta):=\pi_aQ_a^{cs}(\vartheta)$ is
an original exactly constrained initial record. It is generally different from
the mean-preserving record of \cref{v44:central-preparation}.

\begin{theorem}[Stable transients have a paid original observation bound]
\label{v44:matching-theorem}
For sufficiently small common parameters, both matched original histories
exist on $0\le t\le ca^2$. At the same physical time, with $s=t/a^3$,
\begin{align}
 \boxed{\|Q_a^{cs}(t)-Q_a^{c,\pi}(t)\|
       \le Ca^4(|\vartheta|+a)e^{-\beta s},} \label{v44:record-match}\\
 \|\partial_tQ_a^{cs}(t)-\partial_tQ_a^{c,\pi}(t)\|
       &\le Ca(|\vartheta|+a)e^{-\beta s},\label{v44:rate-match}\\
 \boxed{\|\Delta\operatorname{Riem}(t)\|
       \le Ca(|\vartheta|+a)e^{-\beta s},}\label{v44:metric-match}\\
 \boxed{\|\Delta\mathcal L(t)\|
       \le Ca^2(|\vartheta|+a)e^{-\beta s}.}\label{v44:link-match}
\end{align}
For fixed sufficiently small $\vartheta\ne0$, the initial matching cost in
\eqref{v44:record-match} is genuinely of order $a^4|\vartheta|$, not merely
an upper bound. The means of the two initial multipliers may differ by that
order. Their leading $a^2\vartheta$ slow mean displacement is preserved.
\end{theorem}
\begin{proof}
In the constrained chart of V43, let $v_s$ be the stable normal to the new
central manifold. Its exact equation on $\mathcal S_a^\chi$ is
$v_s'=N_s(a,u)v_s$, with $N_s=J_s+O(a+|u|)$, by the lower diagonal part of
\eqref{v44:normal-system}. At the V43 initial point, the actual normalized
velocity is $a^3\mathcal V$, where
\[
 \mathcal V=b+(0,\vartheta k+d_a^\chi)+O(a),\qquad
 d_a^\chi=-\vartheta p_uk+O(a).
\]
Its stable normal derivative is therefore
$a^3\vartheta p_sk+O(a^4)$. Since $u=O(a)$ and $J_s$ is invertible,
\begin{equation}
 v_s(Q_a^{cs})=a^3\vartheta B_s^{-1}p_sk+O(a^4)
 \label{v44:stable-incidence}
\end{equation}
in the weak stable synthesis. The central foliation is tangent to the stable
space at $(a,u)=(0,0)$. Its $C^2$ chart then makes the difference between
this normal vector and the stable-fibre parameter $O(a|v_s|+|v_s|^2)$.
The full initial difference in normalized coordinates is consequently
\begin{equation}
 u_a^{cs}(0)-u_a^{c,\pi}(0)
       =a^3\vartheta\,\iota_sB_s^{-1}p_sk+O(a^4).
 \label{v44:matching-leading}
\end{equation}
For the uniform estimate including $\vartheta=0$, the right side is bounded
by $Ca^3(|\vartheta|+a)$. The coefficient at fixed nonzero $\vartheta$ is
nonzero by \cref{v44:stable-source}. Multiplication by the physical chart's
factor $a$ proves the two-sided initial assertion and its leading purely weak
limit. No claim that the matching leaves the exact multiplier means unchanged
follows from this vector.

The localized stable-fibre theorem gives \eqref{v44:tracking} for both normalized
states. The V43 trajectory remains $O(a)$ away from its normalized reference
inside the original chart through $ca^2$. Its matched trajectory differs by
$O(a^3(|\vartheta|+a))$ there, so it too remains in the smaller physical
chart for small $a$. This first-exit comparison shows that both localized
trajectories are actual original-action histories throughout that interval.
Multiplying the difference estimate by $a$ proves
\eqref{v44:record-match}; multiplying its $s$ derivative by $a/a^3$ proves
\eqref{v44:rate-match}. They are compared at the same physical time, with no
phase shift in the clock.

We keep the original observations in returning to physical variables. The
canonical vector field $F$ has bounded derivatives on the stationary chart,
so the canonical velocity difference is in fact $O(\|\Delta Q\|)$.
The differentiated canonical equation and \eqref{v44:rate-match} give
$\Delta X_{tt}=O(a(|\vartheta|+a)e^{-\beta s})$. Both histories have rates
and fields $O(a)$. The exact initial-jet form of the ADM curvature, applied at
each time, uses those quantities and their finite spatial derivatives but no
second time derivative of the lapse or shift. Its difference is bounded by
\eqref{v44:metric-match} in any fixed spatial norm at this cutoff.

For the link read use the same stationary derivatives
$D_{\lambda_A}z=O(a)$ and $D_{\lambda_W}z=O(a^2)$ as in V38--V43. The
coefficients and their first derivatives are uniformly bounded after extracting
these factors. In $\Delta z_t$, the canonical-rate term is
$O(\Delta Q)$, the active-rate term is at most $a\|\Delta\lambda_t^A\|$,
and the weak-rate term at most $a^2\|\Delta\lambda_t^W\|$; differences of
coefficients multiply original rates of order $a$. The temporal connection
itself is a bounded analytic function of the state. The original temporal-link
formula contains it and $z_t$, not its own time derivative. These estimates
prove the conservative bound \eqref{v44:link-match}, retaining both original
ordered spatial and temporal observations. We do not claim its exponent
optimal. Finite spatial derivatives have equivalent norms at $N=3$.
\end{proof}

At $t=\sigma ca^2$, with fixed $0<\sigma\le1$, the differences in
\eqref{v44:metric-match}--\eqref{v44:link-match} have the extra factor
$e^{-\beta\sigma c/a}$. Unlike an endpoint comparison, this exponential
improves toward the end of the common interval: it is a forward stable-fibre
estimate. At $t=0$ it allows a genuine $O(a|\vartheta|)$ metric difference.
No disappearance of all physical transients at the initial cut is asserted.

\section{What the central reduction does and does not select}
\label{v44:boundary}

The constructions have different preservation properties:
\begin{center}\small
\begin{tabularx}{\textwidth}{@{}Ycc@{}}
\toprule
Construction & Initial record cost & Exact initial means\\
\midrule
Twenty-two-control anchored central preparation relative to its original input
 & $O(a^3(|\vartheta|+a))$ & Preserved\\[2pt]
Stable-fibre matching of the already selected V43 record
 & $O(a^4(|\vartheta|+a))$ & Not asserted\\
\bottomrule
\end{tabularx}
\end{center}
Both solve the original constraints and evolve by the original law. The second
is not substituted into the first chart to claim a better mean-preserving
cost. Its fibres are dynamical comparison classes, not physical gauge orbits.

The word central is likewise scale-specific. It removes the twenty-two
\emph{order-one} hyperbolic eigenvalues of the $s=t/a^3$ limiting matrix.
It retains the seven weak central coordinates, including the isolated imaginary
pair, and all the slower canonical directions. A slow physical rate of order
$a^{-1}$ is of order $a^2$ in $s$ time and is not excluded by this construction.
No conclusion about its excitation or nonlinear transport is obtained by
renaming it central.

The full constraint leaf is $323$-dimensional; V43's manifold has dimension
$312$; the new manifold has dimension $301$. The odd final dimension includes
the original physical flow direction and its radial energy constraint. We have
not constructed a symplectic quotient of dimension $300$, nor proved that the
localized stable foliation is such a quotient. Smooth localization is an
analytical device; the retraction may depend on both choices $\chi,\zeta$.
No canonical or operational rule selecting those choices follows.

\paragraph{Proved in this body.}
The original source has a certified nonzero stable component. On the complete
original constraint leaf there is a $301$-dimensional locally invariant central
manifold inside the V43 selection, together with eleven-dimensional stable
leaves and a time-consistent local retraction. Actual initial-rate controls
reach that manifold without changing any input multiplier mean. The resulting
histories retain the original $ca^2$ interval and original curvature bounds.
The previously selected histories also have a fourth-order stable matching to
central histories, with exponentially decaying differences in their original
metric and link observations. These results remove fast stable transients from
the unresolved propagation problem without deleting their equations or costs.

\paragraph{Inherited inputs.}
The original root, signed source, rapid splitting, analytic stationary maps,
constraint submersion, full weak-rate access and V43 time theorem are inherited
with their stated domains. New calculations verify the stable-source pairing
using their retained interval inputs; they do not extract a new root or the
entire higher signed coefficient bank. The finite checks below test algebra,
weights and incidence; the manifold and propagation arguments are the written
proofs, not proof-assistant formalizations.

\paragraph{Still required.}
The reduced central dynamics has not been shown to stay in the chart on a
longer physical interval, to be attracting in the full leaf, or to have a
positive conserved coercive energy. The remaining rapid oscillations and
slower canonical source must be controlled there with the original observations.
The graph and matching are not causally or microscopically selected. All
constants remain fixed-cutoff constants, the interval still tends to zero,
and spatial refinement, higher physical-time curvature derivatives, normalized
metric/link matching and a common positive physical lifespan remain unproved.
This body neither establishes nor disproves the raw-action Einstein limit.

\paragraph{Executable verification.}
The new source checker solves the complete twenty-eight-row left-eigenvector
system with integer intervals and checks its actual pairing, keeping the
original weak synthesis, constant null identities and input fingerprints.
Separate exact symbolic checks cover full coupled stable/unstable projections,
constraint-chart dimensions and nonconstant mean retention, two-sided and
stable-fibre Green identities, original time factors, normal transversality,
and the distinct preparation and matching costs. The inherited V43 acceptance
is rerun separately; its reports distinguish algebra checks from original
coefficient acceptance. No nonlinear spacetime trajectory or numerical graph
radius is represented as computed. All forty-three preceding marked bodies
are retained byte for byte, including their qualifications and corrections.

\begin{thebibliography}{9}
\raggedright
\bibitem{v44:BLZ}
P. W. Bates, K. Lu and C. Zeng, \emph{Invariant foliations near normally
hyperbolic invariant manifolds for semiflows}, Transactions of the American
Mathematical Society \textbf{352} (2000), 4641--4676,
doi:10.1090/S0002-9947-00-02503-4. Context for stable-fibre methods, with
bibliographic details checked against the authors' publication record.
The finite-dimensional weighted constructions and their application to the
original constraint leaf are given explicitly above; no gravitational estimate
is imported from this reference.
\end{thebibliography}
% END IMMUTABLE EXACT-ACTION BRIDGE V44 BODY

% Future iterations append here.

% BEGIN IMMUTABLE EXACT-ACTION BRIDGE V45 BODY
\clearpage
\renewcommand{\thesection}{\arabic{section}}
\renewcommand{\theHsection}{bridgeV45.\arabic{section}}
\setcounter{section}{364}
\section{Resolve the central degeneracy before estimating a longer interval}
\label{v45:context}

The central reduction in \cref{v44:center-theorem} retains all directions
whose growth is not separated from zero on the rapid scale. Its linear
semigroup was bounded by $C(1+|s|)$. That estimate was sufficient for the
existing $ca^2$ physical interval, but it did not determine whether the
linear growth was an avoidable overestimate. The five eigenvalues in the
small weak cluster had deliberately not been declared exactly zero.
This body resolves that question by an exact coefficient certificate.
The weak Poisson block has rank exactly $24$ at the same root, so its
kernel consists of three constant and two additional nonconstant directions.
The constrained central linearization has exactly two Jordan chains of
length two. Its linear growth is therefore real, not a choice of norm.

The next spectral scale can then be identified without guessing another
inverse bound. Restoring the original amplitude orders gives a positive,
skew-paired intermediate pencil of rank four. The original frozen canonical
matrix has exactly four eigenvalues on this scale, with leading frequencies
of order $a^{-2}$. They are two oscillatory pairs, not a new leading expanding
sector. Their finite-amplitude real parts are allowed to be $O(a^{-1})$.
The four nonzero characteristic roots are encoded by one explicit $4\times4$
positive-metric calculation. Its coefficients include the actual next Poisson
coefficient; no unevaluated coefficient is silently set to zero.

These are fixed-$N=3$ results. They do not improve the nonlinear interval or
identify the frozen spectrum with the Jacobian of the moving constrained
central flow. The next estimate must transport the intermediate positive
energy and its original sources. The two small-parameter limits remain
different: $a$ is field amplitude, whereas spatial refinement changes $h$.

\paragraph{Inputs and nonduplication.}
The first-field polynomial bank is the original V29 bank, in its rational
component-major convention. Its root cube and full Poisson rank are
\cref{v29:root}, with their literal labels checked in the
assembly audit. The positive active Gram and signed weak inertia are V30;
the twenty-four separated nonzero weak roots and positive weak center are
\cref{v32:spectral-theorem}. The full limiting displacement matrix is
\eqref{v32:full-J}; the constrained center is V44. We also retain the exact
canonical spectral transfer of V35, the uniform feedback cancellation
of \cref{v36:uniform-H}, and the harmonic-row orders of V36.
The new finite acquisition is a common six-to-five source contraction and
its root-cube inverse check. The algebraic and spectral consequences are
proved here. Finite-dimensional perturbation is standard background
\cite{v45:Kato}; no external theorem supplies a gravitational source bound.

\section{An exact rank certificate for the original weak Poisson block}
\label{v45:weak-certificate}

Let $e_j$, $0\le j\le26$, be scalar point indicators on the three-site torus,
in the retained lexicographic order. Let $W_g$ have the twenty-six shift
columns $\nabla_\delta(e_j-e_{26})$, $0\le j<26$, synthesized in the
original component-major order. The phase derivative is still
$3(S_i-S_i^{-1})$. The three constant shifts are retained separately.

In the exact V29 coordinates $z=(a_1,b_2,b_3)$, five other mixed-profile
coefficients are fixed at
\[
 (a_2,a_3,a_4,b_1,b_4)
 =2^{-20}(43835,43443,91967,-2807,70091).
\]
Write the complete rationalized first Poisson matrix from that source as
\begin{equation}
 \widehat K(z)=K^{[0]}+z_1K^{[1]}+z_2K^{[2]}+z_3K^{[3]}
   +\left(\frac3{16}z_1^2+\frac14z_2^2+\frac14z_3^2\right)K^{[4]}.
 \label{v45:bank}
\end{equation}
These are the supplied original $108\times108$ coefficient matrices, not
interpolation fits. Their physical congruence is the one in
\eqref{v20:rational-system}. Define
\[
 \Omega_j=W_g^T(K^{[j]})_{\beta\beta}W_g,
 \qquad \Omega(z)=W_g^T\widehat K(z)_{\beta\beta}W_g.
\]
The physical weak block differs by a fixed nonzero scalar congruence, so
its rank is the same. The last coefficient satisfies $\Omega_4=0$ exactly.

\begin{lemma}[A six-dimensional scalar sector has a five-dimensional image]
\label{v45:six-five}
There are fixed integer matrices $F\in\R^{26\times6}$ and
$Z\in\R^{26\times5}$, of ranks six and five, and rational matrices
$G_j\in\R^{5\times6}$ such that
\begin{equation}
 \boxed{\Omega_jF=ZG_j\quad(0\le j\le3),\qquad \Omega_4=0.}
 \label{v45:factor}
\end{equation}
In particular $\rank\Omega(z)\le24$ for every real $z$ in this declared
three-parameter coefficient family, independently of any numerical root.
\end{lemma}
\begin{proof}
All entries of $F,Z,G_j$ are given as integers or rational strings in
\src{data/weak_kernel_certificate.json}. This is a finite exact certificate,
not a numerical nullspace. The following small minors specify its rank
checks. With zero-based rows, the first six rows of $F$ form
\[
 \diag(280364,280364,280364,135719,135719,135719).
\]
For $Z$, rows $(1,2,3,4,5)$ have determinant
\[
 -133127586847868215689907002967883013191223967985280\ne0.
\]
The accompanying verifier reads the original five coefficient matrices,
rebuilds every entry of $W_g$ from the literal shifts, contracts all five
$26\times26$ weak matrices, and checks \eqref{v45:factor} over $\mathbb Q$.
For specificity, the coordinates $G_j$ can equivalently be reconstructed by
\[
 G_j=Z_{I,:}^{-1}(\Omega_jF)_{I,:},\qquad I=(1,2,3,4,5);
\]
the verification checks the resulting equality on all twenty-six rows,
not just on $I$. The fifth contraction is the zero matrix.
Thus $\Omega(z)$ maps a six-dimensional subspace into a fixed
five-dimensional one and has a nonzero kernel vector. It is a real
skew-symmetric matrix of even size $26$. Its rank is even: orthogonal
reduction of a skew form splits its nonzero part into nonsingular
$2\times2$ skew blocks. Consequently its rank is at most $24$.
This proves the assertion at every parameter value rather than only at
sampled members of the family.
\end{proof}

\begin{theorem}[The five weak zero modes are exact at the certified root]
\label{v45:weak-rank}
On a neighborhood of the V29 root within the declared coefficient family,
\begin{equation}
 \boxed{\rank\Omega(z)=24,\qquad \dim\ker K_{WW}(z)=5.}
 \label{v45:rank}
\end{equation}
The latter kernel contains the three original constant shifts and has exactly
two additional nonconstant dimensions. At the root the rapid matrix
$B_*=-D_*^{-1}K_{WW,*}$ has five semisimple zero eigenvalues. Its other two
central eigenvalues are the inherited isolated nonzero imaginary pair.
\end{theorem}
\begin{proof}
Let $z_0$ be the first three coordinates of the exact dyadic center in
\src{data/original_V29_candidate_point.json}. The root theorem places the
root in the cube $|z-z_0|_\infty\le r$, $r=2^{-240}$, as part of its full
107-coordinate cube. Take the principal rows and columns $J=(0,\ldots,23)$
of $\Omega(z)$ and call that matrix $A(z)$. A fixed dyadic matrix $R$,
with denominator $2^{60}$, is supplied. Rational arithmetic verifies
\begin{equation}
 \boxed{
 q:=\|I-RA(z_0)\|_\infty
       +\|R\|_\infty r\sum_{j=1}^3\|(\Omega_j)_{J,J}\|_\infty
       <10^{-12}.}
 \label{v45:minor-bound}
\end{equation}
The actual upper bound lies between $4.55\,10^{-14}$ and
$4.57\,10^{-14}$; $\|R\|_\infty<0.018$. These decimal inequalities are
accepted by rational cross multiplication. Because $\Omega_4=0$, its
variation is affine and the displayed bound covers the entire root cube.
The Neumann series makes $A(z)$ invertible throughout it, proving rank at
least $24$. Combine with \cref{v45:six-five}. Continuity of this minor
extends the rank equality to an open neighborhood of the root.

The constant rows and columns of the physical weak block vanish by the
original first-Poisson identity. The gradient synthesis is injective and
has dimension $26$, so adjoining the constants gives rank $24$ on $29$
weak coordinates. Invertibility of $D_*$ gives the same nullity for $B_*$.
V32 already isolates twenty-four nonzero eigenvalues and proves that the
remaining seven-dimensional restriction is positive for $D_*$. The
restriction of $B_*$ there is similar to a real skew-adjoint matrix and
is semisimple. Nullity five therefore identifies the entire five-member
small cluster as exactly zero. The other two central roots stay in their
already separated nonzero imaginary disks. No floating tolerance defines
any of these zeros.
\end{proof}

The integer certificate is relative to the exact source bank in
\eqref{v45:bank}. That bank's derivation from the original scalar functional
and the V29 root theorem are inherited inputs, not newly reconstructed
nonlinear roots. In particular this rank theorem is not asserted at arbitrary
first fields or under spatial refinement.

\section{Exactly two secular Jordan chains on the constrained center}
\label{v45:jordan}

Write $\mathcal H$ for the three constant weak directions and
$\mathcal Z=\ker K_{WW,*}$. On $\mathcal Z$ the original $D_*$ pairing is
positive definite. Let $p_0$ be the $D_*$-orthogonal spectral projection
onto $\mathcal Z$, and define its two-dimensional complement
\[
 \mathcal Z_g=\mathcal Z\cap\mathcal H^{\perp_{D_*}}.
\]
Elements of $\mathcal Z_g$ may have a constant component in another
coordinate convention. They are not declared physical gauge modes.

The full limiting field has the original form
\[
 J(x,y)=(0,C_*x+B_*y),\qquad \dim x=402,
\]
and its constrained central restriction is $J_{c0}$ of V44, on $301$
dimensions.

\begin{theorem}[The linear central growth is sharp]
\label{v45:jordan-theorem}
The real Jordan decomposition of $J_{c0}$ consists of one nonzero
oscillatory plane, two nilpotent Jordan blocks of size two, and $295$
one-dimensional zero blocks. In particular
\begin{equation}
 \boxed{\dim\ker J_{c0}=297,\qquad
        c|s|\le\|e^{sJ_{c0}}\|\le C(1+|s|)
        \quad(|s|\ge1).}
 \label{v45:jordan-count}
\end{equation}
No fixed equivalent norm makes this entire central semigroup bounded.
The assertion concerns the limiting constrained linearization at $a=0$;
it does not assert secular physical growth on a selected nonlinear history.
\end{theorem}
\begin{proof}
The hyperbolic and nonzero imaginary weak responses to $C_*x$ can be
removed by their bounded fixed inverses, exactly as in V32. The remaining
block has $x'=0$ and $y_0'=p_0C_*x$. We compute the rank of this last map
on the constrained tangent space, not on an arbitrary ambient residual.

For an active multiplier variation $\eta_A$, direct differentiation of
the original leading weak consistency row gives
\begin{equation}
 C_* (0,\eta_A)=-D_*^{-1}K_{WA,*}\eta_A.
 \label{v45:active-coupling}
\end{equation}
The completed source adds nothing to this coefficient: its multiplier
constraint derivative starts at the next amplitude order. This is also
the active-value part of the V20 drift derivative. All active multiplier
variations are permitted in the limiting constraint tangent because the
divided normal map \eqref{v44:constraint-linear} has no multiplier part.

For every constant weak test $c$, the un-tilted quadratic drift is zero
on all fields (V5), and
$K_1(Z)[c,\ell]=\langle\delta_c\ell,LZ\rangle_h$ (V21).
Thus the canonical part of the differentiated prepared quadratic row
also pairs to zero with $c$ on $LZ=0$. The multiplier part pairs to zero
because $K_1(Y_*)c=0$. It follows that, on the constrained tangent,
\[
 p_0C_*x\in\mathcal Z_g.
\]
Here projection is only part of a linear decomposition: for $c\in\mathcal Z$,
$\langle c,D_*p_0v\rangle=\langle c,D_*v\rangle$. No constant row is
deleted from the original source. The rank is consequently at most two.

It is at least two already on the active variations. If
$z\in\mathcal Z_g$ is $D_*$-orthogonal to their image, then
$K_{AW,*}z=0$ by \eqref{v45:active-coupling}. Also $K_{WW,*}z=0$.
Therefore the pure weak vector $(0,z)$ belongs to the kernel of the full
normalized first Poisson matrix. V29's maximal-rank theorem makes that
kernel exactly $\mathcal H$. Positivity on $\mathcal Z$ and the definition
of $\mathcal Z_g$ force $z=0$. This proves rank two.

The nilpotent part squares to zero and has rank two, hence has exactly
two length-two blocks. The only nonzero central weak plane has dimension
two. Subtracting its two dimensions and the four chain dimensions from
$301$ leaves $295$ zero blocks. Their kernels give $295+2=297$.
On a chain $Je_2=e_1$, $Je_1=0$, one has
$e^{sJ}e_2=e_2+se_1$, proving the lower bound. The upper bound is V32's
fixed block reduction. Similarity by any fixed invertible coordinate map
changes constants, not the linear order.
\end{proof}

This rules out replacing the retained $C(1+|s|)$ estimate by a bounded-center
claim at $a=0$. It does not rule out an amplitude-dependent energy that resolves
the chains into oscillations at a different scale. That is the next calculation.

\section{The original canonical transfer on the intermediate scale}
\label{v45:transfer}

Fix one of the original analytic exactly prepared initial curves through
$(Y_*,\ell_*)$ in V30. The same argument applies to a curve with the displayed
Taylor expansions and remainders. Write, in the fixed mean-lapse-zero
complement,
\[
 \begin{gathered}
 C(a)=H_{uX},\qquad Q(a)=H_{XX},\qquad
 M(a)=a^4 E_a^{-1}\Gamma(a)E_a^{-1},\\
 K(a)=C(a)\mathbb JC(a)^*=aK_1+a^2K_2+O(a^3).
 \end{gathered}
\]
where $E_a=\diag(aI_A,I_W)$, $\Gamma(a)=\Gamma_*+O(a)$, and $K_2$ is
the actual coefficient along this prepared curve. The frozen canonical
linearization is still
\[
 T_a=\mathbb J\bigl(Q-C^*M^{-1}C\bigr).
\]
The unknown second preparation is not assigned a value or omitted from $K_2$.
The original harmonic-row calculation in V36 gives
\begin{equation}
 (K_1)_{H,*}=0,\qquad (K_2)_{HH}=0.
 \label{v45:harmonic-orders}
\end{equation}
Its mixed entries $(K_2)_{HR}$ are the original prepared cubic-mean tilt
map; the latter is already specified by the original $P_2,P_{c,3}$ envelopes.

\begin{lemma}[Intermediate specialization of the inherited exact transfer]
\label{v45:transfer-lemma}
For $\tau$ in a fixed compact complex annulus separated from zero and small
positive $a$, the original canonical characteristic roots at $z=\tau/a^2$
are exactly, with multiplicity, the zeros of
\begin{equation}
 \mathcal F_a(a\tau)
 =a\tau\Gamma(a)+E_a(K_1+aK_2+O(a^2))E_a
                  +\frac{a^2}{\tau}E_a\Theta(a,a^2/\tau)E_a,
 \label{v45:exact-transfer}
\end{equation}
where $\Theta$ is uniformly bounded, with fixed $\tau$ derivatives.
This statement retains the complete canonical feedback.
\end{lemma}
\begin{proof}
The exact identity is \eqref{v35:F}, and its feedback factor is uniformly
bounded for large physical spectral parameter by \cref{v36:uniform-H}.
More explicitly, with $F_a=\mathbb JQ(a)$ and $r=1/z$, that retained
factor is $a^{-1}C(a)(I-rF_a)^{-1}F_a\mathbb JC(a)^*$.
Its bounded analytic extension follows from the already established
flat identity $LF_0^jG=0$. We use that result as an input, rather than
claiming another flat cancellation. Substitution of
$r=a^2/\tau$ and $\sigma=a^3z=a\tau$ gives
\eqref{v45:exact-transfer}. The bounded free matrix $F_a$ has no spectrum
at $z=\tau/a^2$ for small $a$ on the declared annuli. All prefactors in
\eqref{v35:det} are then nonzero. Root multiplicities therefore agree.
The bounded derivatives follow from the same inherited analytic coefficient
bounds. This is an algebraic specialization; the new reduction is the
retained-kernel Schur calculation below.
\end{proof}

\section{A positive intermediate pencil of rank four}
\label{v45:intermediate}

Choose a fixed weak complement $\mathcal U$ to $\mathcal Z$ on which
$K_{WW,*}$ is nonsingular; for instance its Euclidean orthogonal complement.
It has dimension $24$. The complementary multiplier space is
\[
 \mathcal V=\mathcal A\oplus\mathcal Z,
 \qquad \dim\mathcal V=83.
\]
All restrictions in this section are congruence restrictions in fixed bases.
Define
\begin{equation}
 \mathsf G=\Gamma_*|_{\mathcal V},\qquad
 \mathsf N=
 \begin{pmatrix}
  0& (K_1)_{AZ}\cr
  -(K_1)_{AZ}^*&(K_2)_{ZZ}
 \end{pmatrix}.
 \label{v45:GN}
\end{equation}

\begin{theorem}[Exactly two leading oscillatory pairs at the intermediate scale]
\label{v45:positive-pencil}
The metric $\mathsf G$ is positive definite. The skew form $\mathsf N$ has
rank exactly four, for every finite value of the actual coefficient $K_2$
consistent with \eqref{v45:harmonic-orders}. Its operator
$-\mathsf G^{-1}\mathsf N$ has $79$ semisimple zero eigenvalues and exactly
two nonzero imaginary pairs
\begin{equation}
                  \boxed{\pm i\nu_1,\quad\pm i\nu_2,
                         \qquad \nu_1,\nu_2>0.}
 \label{v45:frequencies}
\end{equation}
The frequencies are not claimed distinct. They may depend on the chosen
higher preparation through $K_2$.
\end{theorem}
\begin{proof}
Completing the original active Schur square gives, for $z\in\mathcal Z$,
\[
 \left\langle\binom{x}{z},\Gamma_*\binom{x}{z}\right\rangle
 =\langle x+A_*^{-1}L_*z,A_*(x+A_*^{-1}L_*z)\rangle
       +\langle z,D_*z\rangle.
\]
Both terms are positive, with equality only at zero, because the active
Gram is positive and $D_*$ is positive on its weak zero space. This proves
$\mathsf G\succ0$ without declaring the full signed Gram positive.

Order $\mathcal V$ as $(\mathcal A\oplus\mathcal H)\oplus\mathcal Z_g$,
of dimensions $81+2$. By \eqref{v45:harmonic-orders}, its skew form is
\begin{equation}
 \mathsf N=\begin{pmatrix}0_{81\times81}&U\\-U^*&V\end{pmatrix},
 \qquad
 U=\binom{(K_1)_{A Z_g}}{(K_2)_{H Z_g}},\quad
 V=(K_2)_{Z_gZ_g}=-V^*.
 \label{v45:arrow}
\end{equation}
The active part of $U$ is injective. Indeed a weak vector in
$\mathcal Z_g$ killed by it would be a nonconstant null vector of the
full $K_1$, contrary to its exact normalized kernel $\mathcal H$.
Thus $\rank U=2$. Solving $\mathsf N(x,y)=0$ first gives $Uy=0$,
hence $y=0$, and then $U^*x=0$. The kernel has dimension $81-2=79$;
the rank is four, regardless of $V$ or the harmonic part of $U$.
Finally $-\mathsf G^{-1}\mathsf N$ is similar to the real skew-symmetric
matrix $-\mathsf G^{-1/2}\mathsf N\mathsf G^{-1/2}$. Such a matrix is
orthogonally a sum of rotation blocks and zero blocks. Its rank four
proves \eqref{v45:frequencies} and all semisimplicity assertions.
\end{proof}

\begin{proposition}[Exact leading Schur reduction of the canonical transfer]
\label{v45:Schur}
Eliminating the $24$ original $\mathcal U$ rows from
\eqref{v45:exact-transfer}, and dividing the retained block by $a$, gives
\begin{equation}
            \boxed{\mathsf F_a(\tau)
                       =\tau\mathsf G+\mathsf N+O(a).}
 \label{v45:reduced-transfer}
\end{equation}
The error and its fixed $\tau$ derivatives are uniform on compact nonzero
annuli. No $\mathcal U$ equation or canonical-feedback term is discarded.
\end{proposition}
\begin{proof}
The eliminated block is $(K_1)_{UU}+O(a)$ and is invertible with bounded
inverse. The leading off-diagonal blocks are $O(a)$, since $\mathcal Z$
is the exact weak kernel and every active column has the factor $a$.
Their Schur return is therefore $O(a^2)$. The retained undivided block is
$a(\tau\mathsf G+\mathsf N)+O(a^2)$: the active--active first Poisson
term has two active factors, active--zero terms have one, and zero--zero
terms first occur at $aK_2$. The canonical feedback in
\eqref{v45:exact-transfer} is $O(a^2)$ before division. Division by $a$
proves the formula. The inverse and all operations are analytic in $\tau$
on the stated annuli, giving the derivative estimates.
\end{proof}

\section{Four original canonical roots between the two known scales}
\label{v45:canonical}

\begin{theorem}[The missing inverse-square amplitude oscillations]
\label{v45:canonical-theorem}
For the fixed original prepared curve, the frozen $324$-dimensional
canonical matrix has four intermediate-scale eigenvalues, counted with
algebraic multiplicity, satisfying
\begin{equation}
 \boxed{z_{j,\pm}(a)=\frac{\pm i\nu_j}{a^2}+O(a^{-1}),
                       \qquad j=1,2.}
 \label{v45:canonical-roots}
\end{equation}
Choose fixed $r,R$ with $0<r<\min_j\nu_j$ and $R>\max_j\nu_j$.
For sufficiently small positive $a$, these are exactly the four canonical
roots in $r\le|a^2z|\le R$. More generally all canonical roots in a fixed
compact annulus separated from zero on this scale lie within $O(a)$ in
$\tau=a^2z$ of the nonzero roots of \eqref{v45:frequencies}.
There is no new leading hyperbolicity at order $a^{-2}$.
\end{theorem}
\begin{proof}
The determinant of the limiting matrix in \eqref{v45:reduced-transfer}
vanishes at zero with multiplicity $79$ and at the four indicated nonzero
imaginary roots. Its leading coefficient is $\det\mathsf G\ne0$.
Take disjoint small circles around the distinct nonzero roots, counting
a repeated frequency with its full multiplicity. Uniform analytic
convergence of the determinants preserves the number of zeros inside each
circle by the argument principle (equivalently by a homotopy on the
nonvanishing boundary). Uniform nonvanishing on the remainder of a fixed
annulus excludes other roots there. The determinant correspondence of
\cref{v45:transfer-lemma} is exact, so these are canonical eigenvalues.

For the sharper $O(a)$ distance, conjugate the limiting pencil by
$\mathsf G^{-1/2}$. It becomes $\tau I+S$ with $S$ skew-adjoint.
Off its spectrum,
$\|(\tau I+S)^{-1}\|=1/\operatorname{dist}(\tau,-\spec S)$.
At a zero of the perturbed pencil a normalized kernel vector therefore
forces that distance to be at most $Ca$. This argument also applies to
a repeated semisimple imaginary root. Returning from $\tau$ to $z$
gives \eqref{v45:canonical-roots}. It bounds the physical real part only
by $C/a$; finite-amplitude roots need not be exactly imaginary by this
argument, and finite-amplitude semisimplicity is not asserted.
\end{proof}

The interval $ca^2$ proved for selected histories has physical length of the
same order as these oscillation periods. A Gronwall estimate charging their
frequency as a positive growth rate would lose this distinction. Conversely,
a frozen positive-metric oscillation is not by itself a nonautonomous energy
estimate. The $a^{-3}$ hyperbolic sector and the real $a^{-1}$ pair of V37
remain separate, as do the regions $a^2z\to0$ and $|a^2z|\to\infty$.

\section{A four-dimensional frequency formula and the actual missing coefficient}
\label{v45:four-dimensional}

The intermediate result need not be recomputed in $83$ dimensions. In the
order of \eqref{v45:arrow}, put
\begin{equation}
 W=\begin{pmatrix}U&0\\0&I_2\end{pmatrix},\qquad
 J_4=\begin{pmatrix}0&I_2\\-I_2&V\end{pmatrix},\qquad
 H_4=W^*\mathsf G^{-1}W.
 \label{v45:four}
\end{equation}
Here $W$ is $83\times4$, is injective, and $H_4$ is positive definite.

\begin{proposition}[A positive four-coordinate representation, with no lost rows]
\label{v45:frequency-formula}
One has $\mathsf N=WJ_4W^*$ and $\det J_4=1$. The nonzero eigenvalues
of $-\mathsf G^{-1}\mathsf N$ are those of $-J_4H_4$, with multiplicity.
In particular
\begin{equation}
 \boxed{
 \nu_1^2+\nu_2^2=-\tfrac12\tr[(J_4H_4)^2],\qquad
 \nu_1^2\nu_2^2=\det H_4>0.}
 \label{v45:frequency-invariants}
\end{equation}
The quadratic $\mathcal E(x)=\langle x,\mathsf Gx\rangle/2$ is conserved
by the leading intermediate equation $x'=-\mathsf G^{-1}\mathsf N x$.
Its semigroup is bounded in this fixed positive norm.
\end{proposition}
\begin{proof}
Direct multiplication proves the factorization and block elimination proves
$\det J_4=1$. For rectangular maps $A,B$, the identities
$\det(I-tAB)=\det(I-tBA)$ follow by taking the Schur complement in
$\left(\begin{smallmatrix}I&tA\\B&I\end{smallmatrix}\right)$;
thus their nonzero characteristic roots agree. Apply this to
$A=\mathsf G^{-1}W$, $B=J_4W^*$. Positive definiteness of $H_4$ makes
$J_4H_4$ similar to a nonsingular skew-adjoint matrix, so its eigenvalues
are $\pm i\nu_1,\pm i\nu_2$. Its squared trace and determinant give
\eqref{v45:frequency-invariants}. Finally
$\mathcal E'=-\langle x,\mathsf Nx\rangle=0$.
\end{proof}

All quantities except $V=(K_2)_{Z_gZ_g}$ in this formula are obtained from
the inherited leading signed source, the original first Poisson matrix, and
the already specified prepared cubic mixed harmonic row. Since $V$ is a
$2\times2$ skew matrix, its additional content is a single real scalar
once the basis of $\mathcal Z_g$ is fixed. That scalar is an actual
next-order Poisson coefficient, not a freely selectable gyroscopic parameter.
Its numerical value, and hence numerical values of $\nu_1,\nu_2$, have not
been extracted here. No independence from the second canonical preparation
is claimed for it. The existence and rank argument does not need its sign
or value, because the injective active block of $U$ already forces rank four.

If a coefficient path is supplied for this leading intermediate equation,
the exact variable-metric identity is
\[
 \frac{d}{dt}\frac12\langle x,\mathsf G(t)x\rangle
       =\frac12\langle x,\dot{\mathsf G}(t)x\rangle
\]
for $x_t=-a^{-2}\mathsf G(t)^{-1}\mathsf N(t)x$, with constant amplitude
$a$. It charges metric motion, not the frequency $a^{-2}$. This last
identity is an auxiliary leading-equation observation. The original
nonlinear central equation has further moving-coordinate and source terms;
none has been set to zero to claim a longer lifespan.

\section{Consequences, provenance, and the next noncircular estimate}
\label{v45:status}

The exact weak nullity closes a genuine gap in the interpretation of the
retained center. The two extra zero directions cannot be called a numerical
near-zero cluster, and the resulting two Jordan chains cannot be removed by
a fixed uniformly equivalent norm. Restoring the original amplitude orders
then resolves a missing oscillatory scale of the frozen canonical matrix:
there are precisely two leading positive-metric pairs of order $a^{-2}$,
between the established rapid $a^{-3}$ sector and the slow $a^{-1}$ sector.
The nonlinear central manifold is not reduced by deleting those oscillations.

The next useful source calculation is the single remaining skew coefficient
in \eqref{v45:four}, followed by the actual moving-coordinate and source
terms in its intermediate energy. The frequency formula does not replace
that propagation problem. In particular neither the frozen canonical
matrix nor the positive comparison energy is automatically the Jacobian or
conserved energy of V44's moving constrained central flow. The radial
constraint and original metric/link observations remain separate tests.

\paragraph{Proved in this body.}
A complete rational six-to-five contraction proves the weak rank upper bound;
a new root-cube inverse certificate proves its matching lower bound. The
limiting constrained center has exactly two secular chains. The already
proved uniform canonical transfer is specialized to the intermediate scale;
its new retained-kernel Schur limit has rank four in a positive metric. Its four
canonical roots and exact four-coordinate frequency representation follow
with all elimination returns retained.

\paragraph{Inherited inputs.}
The original V29 source bank and root theorem, the V30 signed inertia,
V32 spectral isolation, V35 determinant identity, V36 uniform feedback
cancellation and harmonic orders,
and the V43--V44 constrained charts remain named inputs. They are not
independently re-proved by the new finite calculation. The new verifier
contracts all five inherited original matrices, checks every rational
factorization entry, and recomputes the root-cube inverse acceptance. A
separate run repeats the supplied V44 algebra and stable-source acceptance.

\paragraph{Still required.}
No new numerical intermediate frequencies, full higher-action extraction,
nonlinear central energy estimate, physical-time extension, causal selection
mechanism, or spatial-refinement bound is asserted. The longest established
selected interval is still $0\le t\le ca^2$ at $N=3$; it shrinks with
amplitude. The raw-action Einstein continuum bridge remains unresolved.

\paragraph{Executable checks.}
The proof package includes the complete integer and rational certificate,
its original source-bank fingerprint, the retained root center and radius,
and reproducible contractions from the original phase-gradient synthesis.
Additional exact matrix checks verify the rank-four block, the noncommuting
positive metric, Schur orders, Jordan counts, determinant identity and
frequency invariants. They are algebra regressions, not substitute
gravitational sources or a proof-assistant formalization. Floating exploration
was used to discover coefficient subspaces and propose a dyadic inverse;
only exact rational identities and acceptance inequalities enter the proofs.
The preceding forty-four bodies are preserved byte for byte.

\begin{thebibliography}{9}
\raggedright
\bibitem{v45:Kato}
T. Kato, \emph{A Short Introduction to Perturbation Theory for Linear Operators},
Springer, 1982, Chapter II,
doi:10.1007/978-1-4612-5700-4.
Bibliographic metadata and the finite-dimensional scope were checked on the
publisher's record. The matrix reductions, rank certificates and perturbation
bounds used here are proved above; no external gravitational estimate is
imported.
\end{thebibliography}
% END IMMUTABLE EXACT-ACTION BRIDGE V45 BODY


% Future iterations append here.

% BEGIN IMMUTABLE EXACT-ACTION BRIDGE V46 BODY
\clearpage
\renewcommand{\thesection}{\arabic{section}}
\renewcommand{\theHsection}{bridgeV46.\arabic{section}}
\setcounter{section}{372}
\section{Use the Hamiltonian structure before charging an oscillation as growth}
\label{v46:context}

The intermediate theorem in \cref{v45:canonical-theorem} locates four
canonical roots within $O(a^{-1})$ of two imaginary pairs of size $a^{-2}$.
It does not determine their finite-amplitude real parts. Computing the
remaining scalar in the next Poisson coefficient would determine numerical
leading frequencies, but is not necessary to settle this more basic question.
The exact canonical transfer has additional Hermitian structure on the
imaginary axis. Its crossing form is definite on both positive-frequency
root spaces. We retain that structure, rather than treat its $O(a)$ return
as an arbitrary non-Hamiltonian matrix error.

The result is exact finite-amplitude ellipticity of this isolated cluster:
all four roots are purely imaginary and semisimple for sufficiently small
positive amplitude, even if the two limiting frequencies coincide. Their
original reduced canonical Hessian is negative definite. Its negative is a
positive conserved norm for the frozen intermediate subsystem, with its
amplitude weights explicitly retained. This sign is obtained from an exact
transfer-derivative identity; it is not inferred from the positive multiplier
metric alone.

We then differentiate that original canonical energy along the already
constructed histories. Its intrinsic variation costs at most $C/a$, not
$C/a^2$. The exact moving spectral-projection formula retains the complementary
canonical input. That input has not been estimated away. Consequently this
body does not extend the nonlinear interval, or substitute a frozen
oscillatory space for the invariant central manifold.

\paragraph{Inherited inputs and scope.}
The original normalized Hamiltonian reduction and determinant correspondence
are \cref{v35:flat,v35:transfer-theorem}; the uniform high-frequency return is
\cref{v36:uniform-H}. The exact weak rank, positive intermediate metric,
rank-four pencil, and its $C^2$ spectral convergence are
\cref{v45:weak-rank,v45:positive-pencil,v45:Schur}.
The selected constrained histories and their $O(a)$ physical rates are
\cref{v44:time-theorem}, with the same stationary coefficient bounds used
in V34--V35. These are named inputs, not independently reconstructed
higher-action coefficients. The general relation between Hermitian pencil
crossings and Hamiltonian signatures is established methodology
\cite{v46:KM}; every specialized identity and inference used here is proved.
All action-specific conclusions remain at $N=3$. The small parameter $a$ is
field amplitude, not lattice spacing. No extra action term, projection in
the physical evolution, or changed physical clock is introduced.

\section{An exact Hermitian transfer and its canonical energy sign}
\label{v46:transfer}

Use fixed orthonormal coordinates for the positive physical pairings, so
$\mathbb J^*=-\mathbb J$, $\mathbb J^2=-I$. The orthonormalization is
fixed at this spatial cutoff. Keep the original real matrices
\[
 C=H_{uX},\quad Q=H_{XX}=Q^*,\quad M=H_{uu}=M^*,\quad
 S=Q-C^*M^{-1}C,\quad T=\mathbb JS,\quad F=\mathbb JQ.
\]
Here the normalized multiplier complement has dimension $107$, and the
canonical space has dimension $324$. The radial energy equation is not
included by merely eliminating these $107$ multiplier rows. It will remain
separate in the interpretation below.

For a complex spectral parameter $z$ outside the bounded free spectrum,
define the original transfer
\begin{equation}
 \mathcal B_a(z)=M+C(zI-F)^{-1}\mathbb JC^* .
 \label{v46:B}
\end{equation}
The subscript records the original prepared curve. Stars in complex
identities denote Hermitian adjoints, not analytic transposition.

\begin{lemma}[Hermitian imaginary-axis transfer]
\label{v46:hermitian}
The exact transfer satisfies
\begin{equation}
 \mathcal B_a(z)^*=\mathcal B_a(-\overline z).
 \label{v46:reflection}
\end{equation}
In particular $\mathcal B_a(i\omega)$ is Hermitian for real $\omega$ in
its domain. For $\eta\in\ker\mathcal B_a(i\omega)$ put
\begin{equation}
 \xi_\eta=(i\omega I-F)^{-1}\mathbb JC^*\eta .
 \label{v46:lift}
\end{equation}
Then $C\xi_\eta+M\eta=0$, $T\xi_\eta=i\omega\xi_\eta$, and the
lift is injective. For any two such kernel vectors,
\begin{equation}
 \boxed{\xi_\eta^*S\xi_\zeta
   =-\omega\,\eta^*\frac{d}{d\omega}\mathcal B_a(i\omega)\zeta.}
 \label{v46:energy-transfer}
\end{equation}
\end{lemma}
\begin{proof}
Since $F\mathbb J=-\mathbb JF^*$, the free resolvent $R_z=(zI-F)^{-1}$
obeys $R_z\mathbb J=\mathbb J(zI+F^*)^{-1}$. It follows directly that
$(R_z\mathbb J)^*=R_{-\bar z}\mathbb J$, proving
\eqref{v46:reflection}. At $z=i\omega$ this also gives
$\mathbb JR_z^*=-R_z\mathbb J$. Therefore, for arbitrary multiplier
vectors at that imaginary parameter,
\[
 (R_z\mathbb JC^*\eta)^*\mathbb J(R_z\mathbb JC^*\zeta)
   =-\eta^*CR_z^2\mathbb JC^*\zeta
   =\eta^*\partial_z\mathcal B_a(z)\zeta.
\]
A transfer null vector satisfies $C\xi+M\eta=0$ and
$(zI-F)\xi=\mathbb JC^*\eta$. Substitution into the definition of $T$
gives its eigenvector equation; if $\xi=0$, invertibility of $M$ forces
$\eta=0$. Since $S\xi=-i\omega\mathbb J\xi$ on that kernel and
$\frac{d}{d\omega}\mathcal B_a(i\omega)=i\partial_z\mathcal B_a(i\omega)$,
the preceding identity proves \eqref{v46:energy-transfer}. In particular
its minus sign is fixed by the original elimination, not by a convention
for a positive comparison energy.
\end{proof}

The determinant of a transfer alone need not determine its Hamiltonian
signature. We use the complete Hermitian matrix and its derivative,
including the free canonical return, in \eqref{v46:energy-transfer}.
No statement here requires $Q$, $M$, or $S$ to be positive on their full
spaces.

\section{A definite crossing survives every original elimination return}
\label{v46:crossing}

Retain $E_a=\diag(aI_{78},I_{29})$ and
$M=a^4E_a^{-1}\Gamma E_a^{-1}$. Write $z=\tau/a^2$ and keep the
fixed decomposition $\mathcal V\oplus\mathcal U$ of V45, of dimensions
$83+24$. For real $\rho\ne0$ define the Hermitian matrix
\begin{equation}
 \mathscr B_a(\rho)=a^{-4}E_a\mathcal B_a(i\rho/a^2)E_a,
 \qquad
 \mathscr H_a(\rho)=\operatorname{Schur}_{\mathcal U}\mathscr B_a(\rho).
 \label{v46:H}
\end{equation}
The $\mathcal U$ block is invertible on each fixed compact nonzero
$\rho$ interval for small $a$. Indeed its leading term is
$(ia\rho)^{-1}(K_1)_{UU}$. The exact relations to V45's functions are
\begin{equation}
 \mathcal F_a(a\tau)=\frac{\tau}{a^3}
               E_a\mathcal B_a(\tau/a^2)E_a,
 \qquad
 \boxed{\mathscr H_a(\rho)=\frac{1}{i\rho}\mathsf F_a(i\rho).}
 \label{v46:normalization}
\end{equation}
In particular, with its fixed first two spectral derivatives,
\begin{equation}
 \mathscr H_a(\rho)=\mathsf G+\frac{\mathsf N}{i\rho}+O(a),
 \qquad \mathsf G\succ0,\quad\mathsf N^*=-\mathsf N.
 \label{v46:Hlimit}
\end{equation}
Both Schur operations in this section eliminate equations analytically;
they do not project the action or the evolved fields.

Let $\mathcal L_a(\rho):\mathcal V\to\mathcal V\oplus\mathcal U$
be the exact reconstruction map with $\mathcal V$ component $x$ and
$\mathcal U$ component
$-(\mathscr B_a)_{UU}^{-1}(\mathscr B_a)_{UV}x$.
It has $\mathcal L_a=I_{\mathcal V}+O(a)$ into the full space, with the
same fixed-derivative bounds. It satisfies
\[
 \mathcal L_a^*\mathscr B_a\mathcal L_a=\mathscr H_a,
 \qquad
 \ker\mathscr H_a\ \longleftrightarrow\ \ker\mathscr B_a.
\]
When differentiating this congruence the reconstructed component can depend
on $\rho$. Its derivative lies in $\mathcal U$, whereas
$\mathscr B_a\mathcal L_a$ has zero $\mathcal U$ row. Consequently
\begin{equation}
 \mathcal L_a^*(\partial_\rho\mathscr B_a)\mathcal L_a
                         =\partial_\rho\mathscr H_a.
 \label{v46:schur-derivative}
\end{equation}
This identity does not neglect derivatives of the eliminated block.

\begin{lemma}[Positive crossing on each limiting positive-frequency space]
\label{v46:positive-crossing}
Let $i\nu$, $\nu>0$, be a distinct nonzero eigenvalue of
$-\mathsf G^{-1}\mathsf N$, with multiplicity $m$; here $m=1$ or $2$.
At its Hermitian transfer kernel,
\begin{equation}
 \boxed{x^*\mathscr H_0'(\nu)y=\nu^{-1}x^*\mathsf Gy.}
 \label{v46:crossing-sign}
\end{equation}
After a second exact Schur reduction onto that $m$-dimensional kernel,
the resulting Hermitian function $\mathscr Q_a(\rho)$ can be normalized
so that
\begin{equation}
 \begin{split}
 \mathscr Q_a(\nu)&=O(a),\\
 \mathscr Q_a'(\rho)&=\nu^{-1}I_m+O(a+|\rho-\nu|).
 \end{split}
 \label{v46:local-crossing}
\end{equation}
In particular its derivative is positive definite throughout one fixed
small real interval around $\nu$, for all sufficiently small $a$.
\end{lemma}
\begin{proof}
At the kernel, $\mathsf Nx=-i\nu\mathsf Gx$. Direct differentiation of
\eqref{v46:Hlimit} gives $\mathscr H_0'(\nu)=i\mathsf N/\nu^2$,
which proves \eqref{v46:crossing-sign}. Choose a fixed basis of the kernel
orthonormal for $\mathsf G$, and a fixed complement on which
$\mathscr H_0(\nu)$ is invertible. The latter restriction stays invertible
near $(a,\rho)=(0,\nu)$. The exact Schur return vanishes to second order
in the off-diagonal blocks at this point. Its first derivative on the
kernel is therefore the restriction just computed. The $C^2$ convergence
in \eqref{v46:Hlimit} gives \eqref{v46:local-crossing}, including the
uniform derivative remainder. Shrink the interval so the derivative is
at least $I_m/(2\nu)$. This uses a positive crossing on the whole kernel;
no assumption of distinct frequencies or simple limiting roots is needed.
\end{proof}

\section{The four finite-amplitude canonical roots are exactly imaginary}
\label{v46:ellipticity}

\begin{theorem}[Exact intermediate ellipticity and semisimplicity]
\label{v46:exact-roots}
Fix an original prepared curve covered by V45. There is $a_0>0$ such
that for $0<a<a_0$ all four canonical roots in its intermediate annulus
are exactly
\begin{equation}
 \boxed{z_{j,\pm}(a)=\pm i\omega_j(a),\qquad
     \omega_j(a)>0,\qquad a^2\omega_j(a)=\nu_j+O(a),\quad j=1,2.}
 \label{v46:roots}
\end{equation}
Every root in this cluster is semisimple. This includes a possible
coincidence of the two positive frequencies, either at leading order or
at finite amplitude. In particular
\begin{equation}
       \operatorname{Re}z=0
 \label{v46:zero-real-part}
\end{equation}
exactly for this cluster, not merely $|\operatorname{Re}z|\le C/a$.
The threshold is for the fixed prepared coefficient family; it is not
asserted uniform over unbounded next-order preparations or spatial refinement.
\end{theorem}
\begin{proof}
Consider a positive limiting frequency $\nu$ of multiplicity $m$ and the
exact reduced function in \cref{v46:positive-crossing}. On the real interval
its derivative has a fixed positive floor. Choosing a sufficiently large
fixed $c$ therefore makes
\[
 \mathscr Q_a(\nu-ca)\prec0,\qquad
 \mathscr Q_a(\nu+ca)\succ0
\]
for small $a$. Its $m$ ordered eigenvalues are continuous and strictly
increasing: the variational characterization of an ordered Hermitian
eigenvalue and the matrix derivative floor give a lower increment
$(\rho_2-\rho_1)/(2\nu)$. Each crosses zero exactly once. Simultaneous
crossings count with their kernel dimension.

At any one of those crossings $\rho_a$, split off its exact kernel and
again eliminate the invertible complement. On a kernel of dimension $r$
the Taylor expansion is
$(\rho-\rho_a)D+O((\rho-\rho_a)^2)$, with $D\succ0$. Its determinant
has a zero of order exactly $r$. Thus the sum of these real zero
multiplicities is $m$. The analytic determinant convergence in V45 places
exactly $m$ zeros in the corresponding complex neighborhood. The real
crossings already account for all of them; there is no off-real
$\rho$ zero left in that neighborhood. Since $z=i\rho/a^2$, these
canonical roots are exactly imaginary. Real canonical coefficients give
their negative conjugate partners. The two positive limiting frequencies,
counted with multiplicity, give exactly the four roots of V45.

All transfer reconstructions are injective on the kernel, and their
analytic determinant prefactors are nonzero here. The canonical geometric
multiplicity therefore equals the kernel dimension $r$, while its algebraic
multiplicity equals the zero order $r$. This proves semisimplicity. The
$O(a)$ localization in $\rho$ also follows from the derivative floor and
$\mathscr Q_a(\nu)=O(a)$, or from V45's retained distance estimate.
The proof treats a repeated limiting frequency as one complete crossing
space; it does not divide by a splitting between $\nu_1$ and $\nu_2$.
\end{proof}

An arbitrary non-Hamiltonian error of size $O(a)$ in V45's reduced pencil
could have produced a real part of order $a^{-1}$. Its actual return is
not arbitrary: \eqref{v46:reflection} holds exactly. The theorem uses that
identity and the definite matrix crossing, rather than trying to improve
the size of the generic error bound. It says nothing about the already
proved rapid and slow real pairs, which lie in different annuli.

\section{A definite original canonical Hessian on the intermediate subspace}
\label{v46:energy}

Let $\mathcal E_a$ be the real four-dimensional spectral subspace of $T_a$
corresponding to \cref{v46:exact-roots}. It is the range of the real Riesz
projection $P_a$ enclosing precisely the two conjugate pairs. This
spectral space is not yet identified with a tangent subspace of the
radial energy hypersurface or of V44's nonlinear central manifold.

\begin{theorem}[Canonical sign, conserved norm, and its amplitude cost]
\label{v46:canonical-energy}
The restriction of the original canonical Hessian $S_a$ to $\mathcal E_a$
is negative definite. If $x\in\ker\mathscr H_a(\rho)$, $\rho>0$, and
$\xi$ is its original canonical lift with
$\eta=E_a\mathcal L_a(\rho)x$, then
\begin{equation}
 \boxed{\xi^*S_a\xi
       =-a^4\rho\,x^*\mathscr H_a'(\rho)x
       =-a^4 x^*\mathsf Gx+O(a^5)\|x\|^2.}
 \label{v46:sign-formula}
\end{equation}
The estimate is uniform on the whole root kernel. The real norm
\begin{equation}
             |\xi|_{a,\mathrm{int}}^2
                   :=-a^{-4}\xi^*S_a\xi,
                   \qquad \xi\in\mathcal E_a,
 \label{v46:norm}
\end{equation}
is positive and satisfies, with $\eta=-M^{-1}C\xi$,
\begin{equation}
 \boxed{\|\xi\|\le Ca^2|\xi|_{a,\mathrm{int}},\qquad
 \|\eta_A\|\le Ca|\xi|_{a,\mathrm{int}},\qquad
 \|\eta_W\|\le C|\xi|_{a,\mathrm{int}}.}
 \label{v46:lift-bounds}
\end{equation}
For the frozen equation on this space,
\begin{equation}
          |e^{tT_a}\xi|_{a,\mathrm{int}}=|\xi|_{a,\mathrm{int}}
                    \qquad(t\in\R).
 \label{v46:isometry}
\end{equation}
No amplitude-uniform equivalence with the unweighted original canonical
norm is asserted.
\end{theorem}
\begin{proof}
In \eqref{v46:energy-transfer}, put $\omega=\rho/a^2$ and
$\eta=E_a\mathcal L_a x$. Differentiation of \eqref{v46:H} gives
$E_a\partial_\omega\mathcal B_a(i\omega)E_a
=a^6\partial_\rho\mathscr B_a$. The exact Schur derivative
\eqref{v46:schur-derivative} therefore proves the first equality of
\eqref{v46:sign-formula}. A root is $O(a)$ from a limiting frequency,
and its retained kernel is $O(a)$ from the corresponding limiting kernel.
The latter assertion follows by solving on the fixed invertible complement
used in \cref{v46:positive-crossing}. Consequently
$\rho x^*\mathscr H_a'(\rho)x=x^*\mathsf Gx+O(a)\|x\|^2$,
proving its second equality. It is strictly positive before the displayed
minus sign, uniformly on every actual root kernel.

The original identity $T_a^*S_a+S_aT_a=0$ makes eigenspaces at different
imaginary frequencies $S_a$-orthogonal: for vectors at $i\omega$ and
$i\widetilde\omega$ the pairing is multiplied by
$i(\widetilde\omega-\omega)$. A repeated root is negative definite on
its entire kernel by the same crossing calculation. The negative-frequency
spaces are conjugates of the positive ones and have the same sign.
Semisimplicity from \cref{v46:exact-roots} now makes $S_a$ negative
on their full complex direct sum and hence on its real form $\mathcal E_a$.
This also gives a second no-Jordan check: if
$(T_a-i\omega)v=u\ne0$ with $u$ an eigenvector, skew-adjointness in
$S_a$ would imply $u^*S_au=0$, contrary to its definite sign.

For one eigencomponent, \eqref{v46:sign-formula} bounds $\|x\|$ by
$C|\xi|_{a,\mathrm{int}}$. The exact reconstruction satisfies
$\|\mathcal L_a x\|\le C\|x\|$, and multiplication by $E_a$
gives the two multiplier bounds. The free resolvent at $i\rho/a^2$
has norm at most $Ca^2$; $C$ is bounded, so \eqref{v46:lift} gives the
canonical bound. Decompose a general intermediate vector into its at most
four distinct eigencomponents. They are orthogonal in the positive norm
\eqref{v46:norm}; Cauchy--Schwarz in this finite sum proves the same
bounds with a uniform constant. No ill-conditioned Euclidean eigenbasis
estimate is used when frequencies coalesce.
Finally differentiate \eqref{v46:norm} with $S_a$ fixed and use
$T_a^*S_a+S_aT_a=0$ to prove \eqref{v46:isometry}.
\end{proof}

The sign in \eqref{v46:sign-formula} is a sign of the second variation of
the original reduced Hamiltonian in these canonical directions. It is not
a new spacetime signature, a mass-gap statement, or an assertion that the
physical energy is negative on the constrained nonlinear solution set.
The positive auxiliary multiplier form $\mathsf G$ corresponds here to
\emph{negative} canonical Hessian sign; the exact transfer derivative is
what identifies that relation.

\begin{proposition}[The full intermediate projection respects the original forms]
\label{v46:projection}
For fixed $a>0$ the isolated real spectral projection satisfies
\begin{equation}
 P_a^*S_a=S_aP_a,\qquad P_a^*\mathbb J=\mathbb JP_a.
 \label{v46:orthogonal-projection}
\end{equation}
The restricted symplectic form on $\mathcal E_a$ is nondegenerate. The
projection is an analytical spectral decomposition, not a physical
constraint projection or an identification with a gauge quotient.
\end{proposition}
\begin{proof}
Both $S_a$ and $\mathbb J$ obey $T_a^*B+BT_a=0$ for $B=S_a,\mathbb J$.
The intermediate spectrum is closed under $z\mapsto-\bar z$ and separated
from the complementary spectrum. On their invariant generalized spectral
spaces the cross-form matrix solves a Sylvester equation with disjoint
opposite spectra. It is zero: in Jordan bases one solves its entries
successively, dividing only by nonzero sums of the two eigenvalues.
This proves both identities, without requiring $S_a$ to be nonsingular
on the whole canonical space. If the restriction of $\mathbb J$ were
degenerate, an intermediate vector $v\ne0$ symplectically orthogonal to
the space would make $T_a^{-1}v$ orthogonal to it in $S_a$, since
$S_aT_a^{-1}v=-\mathbb Jv$. This contradicts the definite restriction of
$S_a$. The inverse here is only on the nonzero intermediate space.
\end{proof}

\section{An original moving-energy bound, with the complementary input retained}
\label{v46:moving}

We now use the actual selected histories, not an independent coefficient
path whose regularity is assumed without verification. On any one of the
already supplied families in \cref{v44:time-theorem}, restrict to the
proved $0\le t\le ca^2$ interval and a compact allowed preparation-parameter
set. Their physical rates are $O(a)$, their displacements from the initial
cut are $O(a^3)$, and the amplitude label is held constant during time
differentiation. These are inherited bounds; no new lifetime follows here.
They imply uniform coefficient bounds
\begin{equation}
 \dot C=O(a),\qquad\dot Q=O(a),\qquad
 \dot M=D_a\dot\Gamma D_a,\quad \dot\Gamma=O(1),
 \qquad D_a=\diag(aI_{78},a^2I_{29}).
 \label{v46:time-coefficients}
\end{equation}
Indeed the normalized state coordinates have bounded speed, and the
original stationary matrices and their fixed coefficient derivatives are
bounded on the retained chart. The flat value of $C$ is independent of the
multiplier, so its ordinary time derivative has the stated amplitude order.
The same follows for $Q$ from bounded stationary derivatives and physical
rates $O(a)$.

Their $O(a^3)$ displacement also leaves the first and second prepared
coefficients in V45 unchanged to the orders needed for the intermediate
limit. Thus \cref{v46:exact-roots,v46:canonical-energy} apply uniformly
along the existing interval. The rank-four spectral projection $P(t)$ is
$C^1$ for each fixed $a>0$, including across an internal frequency crossing;
one encloses the whole cluster, rather than differentiating individual
eigenvectors.

\begin{theorem}[No inverse-square cost in the intrinsic intermediate energy]
\label{v46:metric-motion}
Along those original histories, for every $u\in\Ran P(t)$,
\begin{equation}
 \boxed{|u^*\dot S u|\le \frac Ca\,(-u^*S u).}
 \label{v46:Sdot}
\end{equation}
This is an estimate of the derivative of the original canonical Hessian,
not only of the leading multiplier metric. Its proof retains $\dot C$,
$\dot M$, and $\dot Q$.
\end{theorem}
\begin{proof}
Put $\eta=-M^{-1}Cu$. Differentiating the exact Schur Hessian gives
\begin{equation}
 u^*\dot S u
   =u^*\dot Q u+2\operatorname{Re}(\eta^*\dot C u)
                           +\eta^*\dot M\eta.
 \label{v46:exact-Sdot}
\end{equation}
By \cref{v46:canonical-energy}, with $e=|u|_{a,\mathrm{int}}$,
$\|u\|\le Ca^2e$, $\|\eta_A\|\le Cae$, and
$\|\eta_W\|\le Ce$. The first term in \eqref{v46:exact-Sdot} is
therefore bounded by $Ca^5e^2$, the second by $Ca^3e^2$.
For the last term, \eqref{v46:time-coefficients} gives
$\|D_a\eta\|\le Ca^2e$ and hence a bound $Ca^4e^2$.
Their sum is at most $Ca^3e^2=Ca^{-1}(-u^*Su)$ for small $a$.
In particular the potentially important moving constraint source is not
silently dropped; it is the term permitting the inverse-amplitude cost.
\end{proof}

Let $\xi$ be an actual canonical variation of an original constrained
history, with its multiplier variation eliminated. V34's exact equation is
$\xi_t=T(t)\xi$. The full variation, not each spectral component
separately, retains the moving radial constraint. Put
\[
          u=P\xi,\qquad v=(I-P)\xi,\qquad
          \mathcal E_{\mathrm{int}}=-\frac1{2a^4}u^*S u.
\]
Differentiating $P^2=P$ shows $\dot P u\in\ker P$ and
$\dot P v\in\Ran P$. Commutation of $P$ and $T$ gives the exact equation
\begin{equation}
              u_t=Tu+\dot P u+\dot P v.
 \label{v46:projected-original}
\end{equation}
Using \cref{v46:projection}, the normal transport $\dot P u$ contributes
zero to the instantaneous $S$ pairing. The complete energy equation is
therefore
\begin{equation}
 \boxed{
 \frac{d}{dt}\mathcal E_{\mathrm{int}}
     =-\frac1{2a^4}u^*\dot S u
       -\frac1{a^4}\operatorname{Re}(u^*S\dot P v).}
 \label{v46:moving-identity}
\end{equation}
Consequently, with the time-dependent norm in \eqref{v46:norm},
\begin{equation}
 \boxed{|u(t)|_{a,\mathrm{int}}
 \le e^{Ct/a}|u(0)|_{a,\mathrm{int}}
       +\int_0^t e^{C(t-s)/a}|\dot P(s)v(s)|_{a,\mathrm{int}}\,ds.}
 \label{v46:forced-energy}
\end{equation}
For zeros of the norm this follows by regularizing its square root and
passing to the limit. The constants absorb the harmless factor two in the
energy definition.

The homogeneous factor on the already proved interval is $e^{O(a)}$,
rather than a factor charging the oscillation frequency $a^{-2}$.
However, \eqref{v46:forced-energy} keeps the entire complementary input.
No bound for $\dot P v$ in this weighted norm is supplied here.
Removing it would change the variational equation. In particular the
auxiliary equation $u_t=Tu+[\dot P,P]u$ preserves the instantaneous
spectral space, but is not substituted for the original physical variation.
The full projected energy identity, not that auxiliary equation, is the
claim about the original dynamics.

\section{Why exact ellipticity is not yet a physical propagation theorem}
\label{v46:limits}

A simple separate Hamiltonian calculation fixes the norm distinction. With
\[
 T_a^{\mathrm{ex}}=\begin{pmatrix}0&-a^{-1}\\a^{-3}&0\end{pmatrix},
 \qquad S_a^{\mathrm{ex}}=\diag(-a^{-3},-a^{-1}),
\]
one has $T_a^{\mathrm{ex}}=\mathbb JS_a^{\mathrm{ex}}$ and exact
semisimple roots $\pm ia^{-2}$. Its negative Hessian is conserved.
Nevertheless, at $t=\pi a^2/2$,
\[
          e^{tT_a^{\mathrm{ex}}}
               =\begin{pmatrix}0&-a\\a^{-1}&0\end{pmatrix},
\]
whose unweighted norm is at least $a^{-1}$. This is an independent algebra
example, not a new gravitational coefficient or a growth theorem for the
selected gravitational histories. It explains why the amplitude weights in
\cref{v46:canonical-energy} cannot be suppressed.

Likewise the canonical spectral space need not lie in the radial energy
tangent space. The radial corrections for the earlier rapid and slow modes
cannot be imported as an exact invariant projection for these intermediate
modes. In \eqref{v46:projected-original} the complete original tangent is
retained and only decomposed for analysis. Neither its metric-curvature
read nor the literal-link read has been replaced by the canonical Hessian
energy. A bound in this energy needs an additional reader estimate before
it gives a physical curvature bound.

The scalar $V=(K_2)_{Z_gZ_g}$ in V45 is still an actual next-order coefficient.
It remains unevaluated numerically and may depend on higher initial
preparation. This body does not tune it, assume it vanishes, or establish
uniform constants when it is unbounded. The stronger qualitative result is
that every fixed finite value admitted by the original prepared expansion
has the same definite crossing and exact ellipticity once $a$ is sufficiently
small. Numerical frequency separation and resonance calculations still
require that scalar. A mutual collision of the two isolated intermediate
pairs cannot, by itself, produce an off-axis quartet in this small-amplitude
region: they have the same canonical sign. Collisions with modes outside
the isolated cluster are not covered.

\section{Status, source audit, and verification}
\label{v46:status}

\paragraph{Proved here.}
The complete canonical transfer is Hermitian on the imaginary axis. Its
original Schur reduction has a definite crossing, forcing all four
intermediate canonical eigenvalues to remain exactly imaginary and
semisimple at finite positive amplitude. The original canonical Hessian is
negative definite on their isolated real spectral space and supplies a
positive frozen conserved norm after changing its overall sign. The
physical lifting costs are explicit. Along the already constructed
histories, motion of that full Hessian costs at most $C/a$; the exact
moving-projection identity retains the complementary source.

\paragraph{Inherited, not recomputed as new physical data.}
V45's original source-bank contraction and root-cube rank certificate,
V29's root existence, V30's signed inertia, V32's rapid splitting,
and the V35--V45 exact transfers remain identified inputs. The supplied
V45 verification chain is rerun in the package. The new finite checks
verify the noncommuting transfer, derivative signs, both Schur returns,
Hamiltonian energy and Jordan identities, moving-projection terms, and
amplitude arithmetic on independent exact matrix examples. Those examples
check algebra; they do not numerically evaluate the missing gravitational
$K_2$ scalar or establish the analytic theorem by sampling.

\paragraph{Still required.}
No new nonlinear lifespan, intermediate radial-tangent eigenmode,
metric/link read estimate for this energy, estimate of the full
$\dot P(I-P)\xi$ input, causal selection mechanism, or spatial-refinement
bound is asserted. The available selected interval remains
$0\le t\le ca^2$ at $N=3$. Its extension requires control of the moving
canonical coupling and the separate slow real sector, not another frozen
imaginary-root count. The raw-action Einstein continuum bridge remains
unresolved. All forty-five preceding marked research bodies are preserved
byte for byte, including their scope statements and corrections.

\begin{thebibliography}{9}
\raggedright
\bibitem{v46:KM}
R. Koll\'ar and P. D. Miller,
\emph{Graphical Krein signature theory and Evans--Krein functions},
SIAM Review, 2014, doi:10.1137/120891423;
arXiv:1209.3185.
\url{https://arxiv.org/abs/1209.3185}.
The primary author abstract and publisher record were checked for the
matrix-pencil/signature distinction. This is methodological background;
no external gravitational spectral or propagation bound is imported.
\end{thebibliography}
% END IMMUTABLE EXACT-ACTION BRIDGE V46 BODY


% Future iterations append here.

% BEGIN IMMUTABLE EXACT-ACTION BRIDGE V47 BODY
\clearpage
\renewcommand{\thesection}{\arabic{section}}
\renewcommand{\theHsection}{bridgeV47.\arabic{section}}
\setcounter{section}{380}
\section{Control the moving quartet in the norm of the complete linearized record}
\label{v47:context}

The exact ellipticity and moving-energy theorem in V46 leaves one specific
term unbounded: $\dot P(I-P)\xi$, where $P$ is the whole intermediate
canonical spectral projection. Its omission would replace the original
variational equation. Nor does spectral separation alone bound this term
in an unweighted canonical norm: the multiplier lift becomes singular with
amplitude, and the canonical matrix can be nonnormal. This body estimates
that term in a norm retaining both components of the actual linearized
canonical--multiplier record.

The new cancellation is in an \emph{inhomogeneous} resolvent. The constraint
relation for its input cancels the apparently leading multiplier source.
The complete intermediate contour resolvent is consequently $O(a^2)$ in
the record norm. Differentiating its two original block equations, including
the differentiated constraint, proves an $O(a^{-1})$ bound for the moving
projection in that same norm. This does not require computing an individual
eigenvector derivative, dividing by a splitting between the two intermediate
frequencies, or suppressing any complementary eigenmode.

The estimates are populated on the original selected parameter families:
the intermediate component of their actual variational histories differs
from its intrinsic spectral transport by $O(a^2)$ in the intermediate
energy norm throughout $0\le t\le ca^2$. A second consequence for stable-fibre
tangents has error $O(a^4)$. These statements compare a component of an
original variation with a specified analytical transport; they do not declare
that the latter is a separate constrained physical history. A sharper full
geometric reader estimate and a longer nonlinear lifetime remain separate.

\paragraph{Inherited inputs and scope.}
We retain the original canonical reduction and saddle transfer
\cref{v35:transfer-theorem,v46:hermitian}, the exact intermediate Schur
reduction \cref{v45:Schur}, and the definite canonical norm and projection
identities \cref{v46:canonical-energy,v46:projection}. The already constructed
histories and parameter bounds used below are
\cref{v41:selection-theorem,v41:geometry-theorem,v44:time-theorem,v44:matching-theorem}.
They are not new existence results here. The contour and transported-spectral-
space method is standard; \cite{v47:ASY} is background for the distinction
between physical evolution and analytical spectral transport. All norm
estimates used here are proved directly for the original block equations;
no self-adjoint quantum adiabatic estimate is transferred to the indefinite
canonical matrix. All action-specific constants remain at $N=3$. The
parameter $a>0$ is field amplitude, held fixed during physical time
differentiation, not spatial mesh size or a changed clock.

\section{The original lifted-record norm and a bounded contour transfer}
\label{v47:record-norm}

Use the fixed positive physical pairings and the same orthonormal canonical
and multiplier coordinates as V46. Along one of its admitted original
histories set
\begin{equation}
 \begin{gathered}
 F=\mathbb JQ,\qquad S=Q-C^*M^{-1}C,\qquad T=\mathbb JS,\\
 E_a=\diag(aI_{78},I_{29}),\qquad
 M=a^4E_a^{-1}\Gamma E_a^{-1}.
 \end{gathered}                                               \label{v47:data}
\end{equation}
The amplitude factors, spaces and pairings are unchanged. In particular $M$
is the \emph{normalized} $107$-dimensional symmetric multiplier Hessian,
not the skew Poisson matrix. It is invertible on this complement; its full
radial null direction has not been inverted. The canonical space has
dimension $324$.

For every canonical vector $v$ define its original multiplier lift and its
balanced version by
\begin{equation}
 \eta_t(v)=-M(t)^{-1}C(t)v,\qquad
 y_t(v)=E_a^{-1}\eta_t(v).
 \label{v47:lift}
\end{equation}
Thus $Cv+M\eta_t(v)=0$. A variation of an original constrained history has
precisely this multiplier component. For general $v$ this definition solves
only the normalized linearized constraints; it does not assert tangency to
the remaining radial energy equation.

Define the complete-record norm
\begin{equation}
 \boxed{\mathfrak n_{a,t}(v)^2
   =a^{-4}\|v\|^2+\|y_t(v)\|^2.}
 \label{v47:norm}
\end{equation}
Equivalently its multiplier term is
$a^{-2}\|\eta_A(v)\|^2+\|\eta_W(v)\|^2$. No component, including a
constant shift, is discarded. This is an analytical positive norm, not an
additional energy in the original action. It need not be uniformly equivalent
to the unweighted canonical norm as $a\downarrow0$.

We use the actual coefficient bounds of \cref{v46:time-coefficients}:
\begin{equation}
 \begin{gathered}
 \|F\|+\|C\|+\|Q\|+\|\Gamma\|+\|\Gamma^{-1}\|\le C_0,\\
 \|\dot F\|+\|\dot C\|+\|\dot Q\|\le C_0a,\qquad
 \|\dot\Gamma\|\le C_0.
 \end{gathered}                                               \label{v47:coefficients}
\end{equation}
They hold on the established $ca^2$ interval and a compact allowed
preparation set. The same proof applies to V41's endpoint-selected histories:
they have the same required first and second coefficient orders, $O(a)$
physical rates and $O(a^3)$ displacements. Nothing in the coefficient or
intermediate contour arguments uses the invariance of V44's manifold.

Choose a fixed finite union $\mathscr C$ of positively oriented complex
contours enclosing the four nonzero intermediate limiting roots in
$\tau=a^2z$, and excluding zero and every other limiting root. Repeated
frequencies are enclosed together. Shrink the fixed coefficient neighborhood
and the amplitude range so these same contours work throughout the existing
history interval. Put
\begin{equation}
 R_0(z)=(zI-F)^{-1},\qquad
 \widehat{\mathcal B}_a(\tau)
   =a^{-4}E_a\{M+CR_0(\tau/a^2)\mathbb JC^*\}E_a.
 \label{v47:contour-transfer}
\end{equation}

\begin{lemma}[The complete normalized transfer inverse is bounded on the contour]
\label{v47:contour-bound}
Uniformly for $\tau\in\mathscr C$ and the stated original histories,
\begin{equation}
 \|R_0(\tau/a^2)\|\le Ca^2,\qquad
 \|\widehat{\mathcal B}_a(\tau)^{-1}\|\le C.
 \label{v47:contour-inverse}
\end{equation}
\end{lemma}
\begin{proof}
The first bound follows from bounded $F$ and $|\tau|\ge r>0$. For the
second retain the exact $83+24$ decomposition $\mathcal V\oplus\mathcal U$
of V45. Its $\mathcal U$ block in \eqref{v47:contour-transfer} is
$(a\tau)^{-1}(K_1)_{UU}+O(1)$, with inverse $O(a)$. Both off-diagonal
blocks are $O(1)$. The retained Schur complement is
$\mathsf G+\mathsf N/\tau+O(a)$, by \cref{v45:Schur} and
\eqref{v46:normalization}. Its limiting matrix is invertible on the
fixed compact contour. A Neumann argument gives a bounded inverse of that
Schur complement, and the exact block inverse then gives the second bound.
This uses neither a normality assertion for $T$ nor a gap inside the
intermediate quartet. All $24$ eliminated equations are reconstructed.
\end{proof}

\begin{lemma}[Comparison with the original intermediate energy]
\label{v47:energy-equivalence}
For $u\in\Ran P(t)$, where $P(t)$ is the complete rank-four projection,
\begin{equation}
 C^{-1}|u|_{a,\mathrm{int}}
 \le \mathfrak n_{a,t}(u)
 \le C|u|_{a,\mathrm{int}}.
 \label{v47:energy-equivalent}
\end{equation}
\end{lemma}
\begin{proof}
The upper bound is exactly the three lift estimates of
\cref{v46:canonical-energy}. For the reverse inequality the original Schur
identity, valid for every lifted canonical vector, is
\[
 u^*Su=u^*Qu-\eta_t(u)^*M\eta_t(u)
       =u^*Qu-a^4y_t(u)^*\Gamma y_t(u).
\]
On $\Ran P$ its negative divided by $a^4$ is the positive squared
intermediate norm. Bounded $Q,\Gamma$ bound that square by
$C\mathfrak n_{a,t}(u)^2$. No definiteness of $\Gamma$ on its whole
space is used.
\end{proof}

\section{An inhomogeneous resolvent cancellation from the original constraints}
\label{v47:resolvent}

\begin{theorem}[Full canonical contour resolvent in the lifted norm]
\label{v47:resolvent-theorem}
For the original canonical matrix and every canonical input $v$,
\begin{equation}
 \boxed{\mathfrak n_{a,t}\bigl((zI-T)^{-1}v\bigr)
       \le Ca^2\mathfrak n_{a,t}(v),
       \qquad z\in a^{-2}\mathscr C.}
 \label{v47:resolvent-bound}
\end{equation}
Consequently
\begin{equation}
 \boxed{\mathfrak n_{a,t}(Pv)+
        \mathfrak n_{a,t}((I-P)v)\le C\mathfrak n_{a,t}(v).}
 \label{v47:projection-bound}
\end{equation}
\end{theorem}
\begin{proof}
Set $x=(zI-T)^{-1}v$ and $\eta_x=-M^{-1}Cx$. They solve exactly
\begin{equation}
 (zI-F)x-\mathbb JC^*\eta_x=v,\qquad Cx+M\eta_x=0.
 \label{v47:resolvent-system}
\end{equation}
Eliminating $x$ gives, with $y_x=E_a^{-1}\eta_x$,
\[
 y_x=-a^{-4}\widehat{\mathcal B}_a^{-1}E_aCR_0(z)v.
\]
A direct estimate of this expression would spend two unnecessary powers of
$a$. Instead use the \emph{input} lift $Cv=-M\eta_t(v)$ and the exact
resolvent identity $R_0=z^{-1}I+z^{-1}R_0F$. It gives
\begin{equation}
 \boxed{y_x=z^{-1}\widehat{\mathcal B}_a^{-1}\Gamma y_t(v)
   -z^{-1}a^{-4}\widehat{\mathcal B}_a^{-1}E_aCR_0Fv.}
 \label{v47:input-cancellation}
\end{equation}
The first term is at most $Ca^2\mathfrak n(v)$. For the second,
$\|v\|\le a^2\mathfrak n(v)$ and $\|R_0\|\le Ca^2$, so
$\|E_aCR_0Fv\|\le Ca^4\mathfrak n(v)$. Since $|z^{-1}|\le Ca^2$,
it too is at most $Ca^2\mathfrak n(v)$. Thus
\[
 \|y_x\|\le Ca^2\mathfrak n(v),\qquad
 \|x\|=\|R_0(v+\mathbb JC^*E_ay_x)\|
            \le Ca^4\mathfrak n(v).
\]
These are precisely \eqref{v47:resolvent-bound}. The norm on its right is
not a bound on the canonical component alone: the original input multiplier
lift is what justifies \eqref{v47:input-cancellation}.

The Riesz formula is
$P=(2\pi i)^{-1}\int_{a^{-2}\mathscr C}(zI-T)^{-1}\,\dd z$.
The contour length is $O(a^{-2})$, which cancels the $a^2$ in the proved
resolvent bound. The triangle inequality gives the projection bound and
then the bound for $I-P$. No extension to a contour through zero or through
a complementary spectral group is asserted.
\end{proof}

For fixed input and output time, \eqref{v47:resolvent-bound} is a genuine
operator estimate on the complete $324$-dimensional canonical space with
its lifted norm. It does not delete the rapid real groups, slow groups,
or any other complementary modes. The estimate is on a specified contour;
it is not a semigroup estimate for the full canonical matrix.

\section{The full moving-projection input costs at most one inverse amplitude}
\label{v47:projection-motion}

\begin{theorem}[Moving projection with all differentiated source terms]
\label{v47:projection-theorem}
Along the admitted original histories, for every canonical vector $v$,
\begin{equation}
 \boxed{\mathfrak n_{a,t}(\dot P(t)v)
           \le \frac Ca\,\mathfrak n_{a,t}(v).}
 \label{v47:dotP-bound}
\end{equation}
In particular the unresolved input of V46 has the bound
\begin{equation}
 \boxed{|\dot P(I-P)\xi|_{a,\mathrm{int}}
    \le \frac Ca\,\mathfrak n_{a,t}((I-P)\xi)
    \le \frac Ca\,\mathfrak n_{a,t}(\xi).}
 \label{v47:coupling-bound}
\end{equation}
The constants are uniform in amplitude and on the stated fixed-cutoff
history interval. They do not involve the separation of the two frequencies
inside the quartet.
\end{theorem}
\begin{proof}
Differentiate \eqref{v47:resolvent-system} at fixed $a,z$ and fixed
canonical input $v$. With $x=(zI-T)^{-1}v$, set
\[
 r_1=\dot F x+\mathbb J\dot C^*\eta_x,\qquad
 s_1=-\dot Cx-\dot M\eta_x.
\]
The actual differentiated equations are
\begin{equation}
 (zI-F)\dot x-\mathbb JC^*\dot\eta_x=r_1,\qquad
 C\dot x+M\dot\eta_x=s_1.
 \label{v47:differentiated-saddle}
\end{equation}
They include $\dot M$ and the inhomogeneous differentiated constraint.
In particular $\dot\eta_x$ is not the multiplier lift of $\dot x$.

Write $r=\mathfrak n_{a,t}(v)$. The preceding theorem gives
$\|x\|\le Ca^4r$ and $\|y_x\|\le Ca^2r$. The bounds
\eqref{v47:coefficients}, and the exact equality
$\dot M=a^4E_a^{-1}\dot\Gamma E_a^{-1}$, imply
\begin{equation}
 \|r_1\|\le Ca^3r,\qquad
 \|E_as_1\|\le Ca^5r.
 \label{v47:derivative-sources}
\end{equation}
For example $E_a\dot M\eta_x=a^4\dot\Gamma y_x=O(a^6r)$;
the potentially larger term $E_a\dot Cx$ is $O(a^5r)$ and is retained.
The coefficient variation of the free matrix is also retained in $r_1$.

Solving \eqref{v47:differentiated-saddle} gives the exact expressions
\begin{equation}
 \begin{split}
 E_a^{-1}\dot\eta_x
 &=a^{-4}\widehat{\mathcal B}_a^{-1}
                       E_a(s_1-CR_0r_1),\\
 \dot x&=R_0(r_1+\mathbb JC^*\dot\eta_x),\\
 y_t(\dot x)&=E_a^{-1}\dot\eta_x
                         -a^{-4}\Gamma^{-1}E_as_1.
 \end{split}                                                  \label{v47:derivative-solve}
\end{equation}
Their signs follow directly from the nonzero second row of
\eqref{v47:differentiated-saddle}. Since
$\|E_aCR_0r_1\|\le Ca^5r$, they yield
\[
 \|E_a^{-1}\dot\eta_x\|+\|y_t(\dot x)\|\le Car,
 \qquad \|\dot x\|\le Ca^3r.
\]
Consequently $\mathfrak n_{a,t}(\partial_t(zI-T)^{-1}v)\le Car$.
The contours can be held fixed throughout the physical interval. Differentiate
their Riesz formula and use their $O(a^{-2})$ length to prove
\eqref{v47:dotP-bound}.

Finally $P\dot PP=0$ and $(I-P)\dot P(I-P)=0$, by differentiating
$P^2=P$. Thus $\dot P(I-P)\xi\in\Ran P$. Apply
\cref{v47:energy-equivalence} and \eqref{v47:projection-bound} to obtain
\eqref{v47:coupling-bound}. No bound on an arbitrary derivative of an
eigenbasis is required.
\end{proof}

The proof has not replaced $\dot P(I-P)\xi$ by zero. It bounds it using
the original complete record, with its canonical and multiplier amplitude
weights stated explicitly. Bounding $\dot T$ in an unweighted canonical
norm before the saddle solve would conceal the cancellation and give a much
weaker estimate. No optimality claim is made for the exponent $a^{-1}$.

\begin{proposition}[Motion of the whole spectral space over the existing interval]
\label{v47:finite-motion}
For a fixed canonical vector $v$ and any two times in the admitted interval,
\begin{align}
 \mathfrak n_{a,t}(v)&\le e^{C|t-s|/a}\mathfrak n_{a,s}(v),\notag\\
 \mathfrak n_{a,t}((P(t)-P(s))v)
 &\le C\frac{|t-s|}{a}e^{C|t-s|/a}\mathfrak n_{a,s}(v).
 \label{v47:finite-projection}
\end{align}
In particular, over $|t-s|\le ca^2$ the latter relative bound is $O(a)$.
\end{proposition}
\begin{proof}
For fixed $v$, differentiate its lift rather than its coordinates:
\[
 \partial_ty_t(v)=-\Gamma^{-1}\dot\Gamma y_t(v)
                         -a^{-4}\Gamma^{-1}E_a\dot C v.
\]
It has norm at most $Ca^{-1}\mathfrak n_{a,t}(v)$. The canonical part
of \eqref{v47:norm} is constant, so differentiation of the squared norm,
with a regularized square root when necessary, gives its first inequality
in both time directions. Integrate $\dot P(r)v$ from $s$ to $t$; apply
\eqref{v47:dotP-bound} at $r$ and transfer each norm between $s,r,t$ by
the first inequality. This proves the second inequality with a changed
constant. It is a bound for the entire spectral projection in a specified
moving norm, not a claim that its physical subspace is exactly fixed.
\end{proof}

\section{A causal comparison for the actual projected variation}
\label{v47:transport}

For fixed $a>0$ let $\mathcal U_a(t,s)$ be the solution operator of the
analytical equation
\begin{equation}
 \partial_tu=(T+[\dot P,P])u,\qquad u(s)\in\Ran P(s).
 \label{v47:intrinsic-equation}
\end{equation}
This is used only as a comparison operator. It is not substituted for the
canonical variational equation $\xi_t=T\xi$ or for nonlinear physical
histories. The normed spectral transport and the actual variation have
different equations unless the complementary input vanishes.

\begin{theorem}[Complete forced spectral transport]
\label{v47:transport-theorem}
The operator in \eqref{v47:intrinsic-equation} maps $\Ran P(s)$ into
$\Ran P(t)$ and obeys
\begin{equation}
 |\mathcal U_a(t,s)u|_{a,\mathrm{int},t}
          \le e^{C(t-s)/a}|u|_{a,\mathrm{int},s}\qquad(t\ge s).
 \label{v47:intrinsic-bound}
\end{equation}
For an actual constrained canonical variation $\xi_t=T\xi$, put
$u=P\xi$, $v=(I-P)\xi$. Then exactly
\begin{equation}
 \boxed{u(t)=\mathcal U_a(t,0)u(0)
           +\int_0^t\mathcal U_a(t,s)\dot P(s)v(s)\,\dd s.}
 \label{v47:Duhamel}
\end{equation}
In particular,
\begin{equation}
 \boxed{|u(t)-\mathcal U_a(t,0)u(0)|_{a,\mathrm{int},t}
    \le\frac Ca\int_0^t e^{C(t-s)/a}
                           \mathfrak n_{a,s}(\xi(s))\,\dd s.}
 \label{v47:full-record-error}
\end{equation}
The whole original variation retains the radial constraint. Neither term
on the right of \eqref{v47:Duhamel} is asserted to be a separate radial-
constraint-tangent physical solution.
\end{theorem}
\begin{proof}
For $K_P=[\dot P,P]$, differentiating $P^2=P$ gives
$\dot P+[P,K_P]=0$. Together with $[P,T]=0$ this proves the
intertwining $P(t)\mathcal U_a(t,s)=\mathcal U_a(t,s)P(s)$, either
by differentiating the identity or by uniqueness of the corresponding
finite linear initial-value equation.

On $\Ran P$, $K_Pu=\dot P u\in\ker P$. The original form identity
$P^*S=SP$ therefore gives $u^*SK_Pu=0$. The Hamiltonian part also
contributes zero, since $T^*S+ST=0$. The intrinsic squared norm hence
has derivative $-a^{-4}u^*\dot S u$. Apply
\cref{v46:metric-motion} and Gronwall to obtain
\eqref{v47:intrinsic-bound}, absorbing a factor two in the constant.

The actual projected equation of V46 is
$u_t=Tu+\dot P u+\dot Pv=(T+K_P)u+\dot Pv$.
The last term belongs to $\Ran P$. Variation of constants proves
\eqref{v47:Duhamel}. Applying \eqref{v47:intrinsic-bound} and then
\eqref{v47:coupling-bound} proves \eqref{v47:full-record-error}.
This is a causal integral along the actual variational history, not an
initial--terminal reselection or a prescription of its future forcing.
\end{proof}

\section{The error estimate is populated by original selected histories}
\label{v47:population}

The preceding theorems bound all canonical inputs in the specified norm.
The following applications use \emph{proved} norms of actual original
variations, rather than introduce another condition on an unknown inverse.

\begin{theorem}[Intermediate tracking of the original selected parameter tangent]
\label{v47:selected-tangent}
Take the original $C^1$ selected family of V41 on its established common
interval and a compact smaller parameter interval. At fixed physical time
let
$\xi=\partial_\vartheta X$ and
$\eta=\partial_\vartheta\lambda$.
Then the original constrained variational equation holds, and
\begin{equation}
 \sup_{0\le t\le ca^2}\mathfrak n_{a,t}(\xi(t))\le Ca.
 \label{v47:selected-record}
\end{equation}
Its intermediate component satisfies
\begin{equation}
 \boxed{\sup_{0\le t\le ca^2}|P(t)\xi(t)|_{a,\mathrm{int}}
            \le Ca,\qquad
 \sup_{0\le t\le ca^2}
 |P(t)\xi(t)-\mathcal U_a(t,0)P(0)\xi(0)|_{a,\mathrm{int}}
            \le Ca^2.}
 \label{v47:selected-error}
\end{equation}
Writing the latter difference as $e(t)\in\Ran P(t)$, its physical lifts
obey
\begin{equation}
 \boxed{\|e(t)\|\le Ca^4,\qquad
        \|\eta_t(e(t))_A\|\le Ca^3,\qquad
        \|\eta_t(e(t))_W\|\le Ca^2.}
 \label{v47:selected-lifts}
\end{equation}
\end{theorem}
\begin{proof}
The parameter-differentiated bounds in the proof of
\cref{v41:geometry-theorem} give, on the whole original interval,
\[
 \|\xi\|\le Ca^3,\qquad \|\eta\|\le Ca^2.
\]
They follow there by differentiating the anchored record chart and the
uniformly contracting selected equation, not by differentiating a bare
undifferentiated big-$O$ assertion. Both field and multiplier constraints
hold identically in the parameter, so $C\xi+M\eta=0$ and its radial
counterpart hold. The first identity makes $\eta$ exactly the lift in
\eqref{v47:lift}; the original reduction gives $\xi_t=T\xi$.
Inserting the two bounds in \eqref{v47:norm} proves
\eqref{v47:selected-record}. The more precise active bound in V41 is
not needed for this estimate.

These base histories have $O(a)$ physical rates and $O(a^3)$
displacements and keep the same leading intermediate coefficient data.
Thus the uniform contour and coefficient proof above applies to them,
as noted after \eqref{v47:coefficients}. The projection estimate and
\cref{v47:energy-equivalence} give the first bound in
\eqref{v47:selected-error}. For the second use
\eqref{v47:full-record-error}: its integrand is bounded by $Ca$, the
interval length is $ca^2$, and $e^{Ct/a}\le e^{Cca}$ is uniformly
bounded. The result is $Ca^2$. Finally apply the three original
intermediate lifting bounds to $e(t)$ to obtain
\eqref{v47:selected-lifts}.
\end{proof}

The $O(a^2)$ comparison error is absolute in the intermediate norm. It
would be a relative $O(a)$ estimate only where the initial intermediate
component has a corresponding nonzero lower bound; none is asserted here.
The chosen endpoint in the original V41 selection is kept fixed. The
intrinsic comparison introduces no additional endpoint choice.

\begin{corollary}[Stable-fibre tangents have a still smaller integrated input]
\label{v47:stable-tangent}
Fix a base central history and its original stable fibre from V44. Normalize
a local physical initial-data tangent along that fibre to have size at most
$Ca^4$. For its actual variational history,
\begin{equation}
 \mathfrak n_{a,t}(\xi(t))\le Ca^2e^{-\beta t/a^3},
 \qquad 0\le t\le ca^2,
 \label{v47:stable-record}
\end{equation}
with a fixed admitted $\beta>0$. Its intermediate forcing contribution
has the uniform bound
\begin{equation}
 \boxed{|P(t)\xi(t)-\mathcal U_a(t,0)P(0)\xi(0)|_{a,\mathrm{int}}
          \le Ca^4.}
 \label{v47:stable-error}
\end{equation}
In fact the intermediate component itself has the sharper bound
\begin{equation}
 \boxed{\sup_{0\le t\le ca^2}|P(t)\xi(t)|_{a,\mathrm{int}}\le Ca^4.}
 \label{v47:stable-contamination}
\end{equation}
For each fixed $0<\sigma<1$, on $0\le t\le\sigma ca^2$ it obeys
$|P(t)\xi(t)|_{a,\mathrm{int}}\le C_\sigma a^4e^{-\beta t/a^3}$.
\end{corollary}
\begin{proof}
Use the $C^2$ stable-fibre charts and their exponential record tracking
estimate in V44. Divide the difference between two same-fibre original
histories by their initial parameter separation and pass to the limit.
The constants are uniform on the smaller fibre chart. This gives
$\|\xi(t)\|+\|\eta(t)\|\le Ca^4e^{-\beta t/a^3}$.
All original linearized constraints remain satisfied in that limit.
Equation \eqref{v47:norm} gives \eqref{v47:stable-record}. Now
\eqref{v47:full-record-error} is at most
\[
 \frac Ca e^{Cca}\int_0^{ca^2}a^2e^{-\beta s/a^3}\,\dd s
      \le Ca^4.
\]
For the sharper conclusion, let $T_a=ca^2$. The intrinsic transport is
invertible, and the same two-sided moving-energy inequality bounds its
backward norm by $e^{C|t-s|/a}$. Apply the exact Duhamel identity backward
from the \emph{already existing} stable history at $T_a$. Its terminal
component is at most $Ca^2e^{-\beta T_a/a^3}$, by
\eqref{v47:stable-record} and the bounded projection. The remaining integral
is at most
\[
 \frac Ca\int_t^{T_a}e^{C(s-t)/a}a^2e^{-\beta s/a^3}\,\dd s
          \le Ca^4e^{-\beta t/a^3}
\]
for small $a$, since $\beta/a^3-C/a\ge\beta/(2a^3)$. It follows that
\[
 |P(t)\xi(t)|_{a,\mathrm{int}}
 \le C\{a^4e^{-\beta t/a^3}+a^2e^{-\beta c/a}\}.
\]
The last term is $O(a^4)$, proving \eqref{v47:stable-contamination}; on
$t\le\sigma T_a$ it is absorbed into the first term with a changed
$\sigma$-dependent constant. This estimates a given stable-fibre tangent;
it imposes no new terminal condition on the physical history. It is not an
estimate for arbitrary complementary initial errors or a new attracting
basin in the whole constraint leaf.
\end{proof}

\section{Why the new norm must not be replaced by a curvature claim}
\label{v47:reader-boundary}

The coupling in V46 is now controlled as an operator in the complete-record
norm, and the integrated error is evaluated on nontrivial original
variational families. There are nevertheless two different interpretations
to keep separate.

First, \eqref{v47:selected-lifts} concerns canonical and multiplier
\emph{components} of an analytical spectral comparison. A projected vector
need not satisfy the radial energy tangent condition, and the intrinsic
transport need not coincide with a tangent of V44's nonlinear central
manifold. The whole $\xi$ is the physical tangent and retains all those
constraints. Its original observations are evaluated without deleting its
complement.

Second, the energy norm does not include the whole original rate
differential. Even for $u\in\Ran P$ the frozen bound is only
\[
       |Tu|_{a,\mathrm{int}}\le Ca^{-2}|u|_{a,\mathrm{int}}.
\]
There is no justified inference that a small componentwise canonical
comparison has an equally small curvature comparison after time
differentiation. The exact rate differential also contains the moving
constraint terms recorded in V34 and V38. The original geometric
observations are not replaced by the auxiliary quadratic norm in this
body. In particular an $O(a^2)$ spectral-transport error has not been
promoted to an $O(a^2)$ metric-curvature error.

Likewise \eqref{v47:full-record-error} is not a bootstrap bound for arbitrary
full canonical variations. The rapid and slow real groups are still present,
and no uniform bound for their complete-record norm has been proved for
arbitrary data. The selected tangent and stable-fibre applications use norms
already established for those original families. The coupling constant is
uniform in amplitude at fixed cutoff, not under $N\to\infty$.

\section{Status, reproducibility, and the next source estimate}
\label{v47:status}

\paragraph{Proved here.}
The original inhomogeneous canonical resolvent is bounded by $Ca^2$ on
the intermediate contour in the complete lifted-record norm. Its input
constraint cancels the leading multiplier source. Differentiating both
original block equations proves $\|\dot P\|\le C/a$ in that same norm,
including the correction between a differentiated multiplier and the lift
of a differentiated canonical vector. This gives a normed bound for the
entire formerly unestimated complementary input. The projection itself moves
by only $O(a)$ in that norm over the existing $ca^2$ interval. The exact
Duhamel comparison retains all complementary forcing and is populated by
original selected parameter tangents, with intermediate error $O(a^2)$,
and by stable-fibre tangents, with integrated error and total intermediate
component $O(a^4)$.

\paragraph{Inherited, not re-established by numerical sampling.}
The mixed-profile root, higher signed-source enclosure, weak rank and
intermediate ellipticity remain named V29--V46 inputs. The supplied V46
checker, including its inherited chain, is rerun. The new exact checks verify
the resolvent and differentiated saddle identities on noncommuting rational
matrices, the graph-lift correction, spectral intertwining, moving energy,
and all amplitude powers. Separate examples are algebra checks, not newly
extracted gravitational coefficient banks. No numerical trajectory or new
frequency enclosure is used to prove the uniform estimates.

\paragraph{Still required.}
The nonlinear lifespan is not extended beyond $ca^2$. The physical reader
of the complete spectral comparison, arbitrary-data control of the full
lifted norm, longer-time coupling to the separate slow real sector, and
refinement remain unresolved. A useful next calculation is the original
rate/curvature differential applied to the \emph{coupled} comparison,
retaining both spectral components and the moving source; not another
frozen imaginary-root count or an assertion that $\dot P$ vanishes.
The raw-action Einstein continuum bridge remains open. All forty-six
preceding marked mathematical bodies are retained byte for byte.

\begin{thebibliography}{9}
\raggedright
\bibitem[ASY87]{v47:ASY}
J. E. Avron, R. Seiler and L. G. Yaffe,
\emph{Adiabatic theorems and applications to the quantum Hall effect},
Communications in Mathematical Physics \textbf{110} (1987), 33--49,
doi:10.1007/BF01209015; erratum \textbf{156} (1993), 649--650,
doi:10.1007/BF02096867.
The primary institutional record and the publisher's erratum record were
checked. Cited for the established distinction between spectral parallel
transport and physical evolution, not for a theorem applied without its
hypotheses to the present nonnormal constrained system. The finite contour
and derivative estimates required here are proved explicitly above.
\end{thebibliography}
% END IMMUTABLE EXACT-ACTION BRIDGE V47 BODY


% Future iterations append here.

% BEGIN IMMUTABLE EXACT-ACTION BRIDGE V48 BODY
\clearpage
\renewcommand{\thesection}{\arabic{section}}
\renewcommand{\theHsection}{bridgeV48.\arabic{section}}
\setcounter{section}{388}
\section{Pass from spectral error to the original curvature differential}
\label{v48:context}

V47 estimates the moving intermediate projection in the norm of a complete
canonical--multiplier record. It deliberately does not interpret an error of
order $a^2$ in that norm as an error of order $a^2$ in metric curvature.
This body evaluates the missing reader, including its actual multiplier-rate
variation. The resulting loss is exactly two powers of amplitude at the
operator-norm level. It is not merely an avoidable estimate: the original
intermediate pencil has no unobserved oscillatory direction under its metric
reader over a fixed rescaled observation window.

The positive consequences are correspondingly qualified. V47's stable-fibre
intermediate component has original metric contribution $O(a^2)$ and literal
temporal-link contribution $O(a^4)$. Its selected-parameter comparison error
has a proved $O(1)$ metric bound and $O(a^2)$ temporal-link bound; these upper
bounds do not assert that the actual error has a nonzero limit. We construct
separate, exactly constraint-tangent differential witnesses with precisely the
canonical and multiplier sizes of that comparison error and nonvanishing
metric output. Thus those sizes alone cannot imply a vanishing metric error.
The remaining problem is cancellation on the \emph{actual} error, not another
unweighted imaginary-spectrum argument.

\paragraph{Inputs, conventions, and nonduplication.}
The original normalized constraint graph, canonical reduction, and rate
identity are \cref{v34:context,v38:actual-rate,v47:record-norm}. The intermediate
pencil and canonical sign are \cref{v45:positive-pencil,v46:canonical-energy};
the moving projection and populated error bounds are V47. The bounded radial
normal is \cref{v37:lossless-normal}. We retain these named inputs and do not
claim independent validation of the whole cumulative argument. The finite
observability argument is proved below; \cite{v48:Hautus} credits its standard
linear-systems context, not an additional estimate for this source.

All action-specific assertions concern the original $N=3$ regulator and its
already admitted selected histories on $0\le t\le ca^2$. The positive label
$a$ is field amplitude, held fixed in physical time. Constants may depend on
the fixed chart, the actual finite higher coefficients, and a compact smaller
preparation set. They are not uniform in spatial refinement. No numerical
value for a new constant or intermediate frequency is asserted.

\section{Define the action reader before differentiating a comparison curve}
\label{v48:readers}

Use the notation of V47:
\[
 F=\mathbb JQ,\quad S=Q-C^*M^{-1}C,\quad T=\mathbb JS,\quad
 E_a=\diag(aI_{78},I_{29}),\quad M=a^4E_a^{-1}\Gamma E_a^{-1}.
\]
Write $\mathfrak f$ for the original nonlinear canonical vector and $r$ for
the original normalized multiplier rate. These are distinct from the linear
matrix $F$. At every base point on the admitted history,
\[
 \mathfrak f,r=O(a),\quad C=L+O(a),\quad
 Q,C,\Gamma,\Gamma^{-1}=O(1),\quad
 \dot Q,\dot C=O(a),\quad \dot\Gamma=O(1).
\]
All norms of fixed-order spatial derivatives are equivalent at this fixed
cutoff. The stationary maps and their fixed derivatives are bounded on the
original chart, before the singular multiplier inverse is applied.

For any canonical vector $v$, set
\begin{equation}
 \eta(v)=-M^{-1}Cv,\qquad y(v)=E_a^{-1}\eta(v),\qquad
 \mathfrak n(v)^2=a^{-4}\|v\|^2+\|y(v)\|^2.
 \label{v48:lift}
\end{equation}
This is exactly V47's norm. The normalized constraints determine $\eta$;
the remaining radial tangent equation is imposed separately below.

The original stationary equation $H_u=0$, with $u$ the 107 normalized
multiplier coordinates, locally solves these coordinates as a function of
the canonical fields. On this graph, use the original action equations to
supply $\mathfrak f,r$, and
\[
 \partial_t\mathfrak f=F\mathfrak f+\mathbb JC^*r.
\]
Substitution into the original interpolated ADM curvature defines a smooth
finite-point reader at each nonzero $a$. The same procedure defines the
original temporal and spatial literal-link readers. Their differentials in
the direction $(v,\eta(v))$ will be denoted
\begin{equation}
 \mathscr R_t(v),\qquad \mathscr L_t^{\rm el}(v),\qquad
 \mathscr L_t^{\rm sp}(v).
 \label{v48:reader-definitions}
\end{equation}
The electric metric block of $\mathscr R$ is denoted $\mathscr R^{0}$.
The notation $\mathscr R$ includes all independent Riemann components, not
only a scalar invariant. On a radial-constraint-tangent vector these are
precisely derivatives of original local histories. On a general vector they
are the linear extensions furnished by the same normalized action graph.
A spectral summand is not declared a separate physical history.

Define the ordinary-versus-phase derivative reader retained in V10 by
\begin{equation}
 (\mathfrak E_h b)_{ij}
 =\tfrac12\{(\partial_i^{\rm sp}-\delta_i)b_j+
                  (\partial_j^{\rm sp}-\delta_j)b_i\}.
 \label{v48:E}
\end{equation}
Let $\iota_0$ insert this symmetric tensor as the electric block, with its
Riemann symmetries and zero other independent blocks. At $N=3$,
$\partial_i^{\rm sp}=2\pi\delta_i/(3\sqrt3)$; hence the kernel of
$\mathfrak E_h$ consists exactly of the three constant shifts. For a nonzero
Fourier mode this also follows from
\[
 \left|\tfrac12(d\otimes b+b\otimes d)\right|_F^2
 =\tfrac12\{|d|^2|b|^2+|d\cdot b|^2\},\qquad
 d=2\pi k-\kappa_h(k)\ne0.
\]
This identity uses the original Frobenius and spatial-mean conventions.

\begin{lemma}[Exact rate differential and derivative of the lift]
\label{v48:rate}
The original action rate differential is
\begin{equation}
 \boxed{\rho_t(v):=Dr[(v,\eta(v))]
       =\eta(Tv)+\dot\eta(v),\qquad
 \dot\eta(v)=-M^{-1}\{\dot Cv+\dot M\eta(v)\}.}
 \label{v48:rho}
\end{equation}
The dot on the linear map $\eta$ holds its canonical argument fixed.
Furthermore
\begin{equation}
             \|E_a^{-1}\dot\eta(v)\|\le Ca^{-1}\mathfrak n(v).
 \label{v48:lift-motion}
\end{equation}
\end{lemma}
\begin{proof}
Differentiate $C v+M\eta(v)=0$ along an original constrained variation,
whose canonical equation is $v_t=Tv$. Its multiplier equation is
$\eta_t=\rho_t(v)$. This gives
$M\rho_t(v)=-CTv-\dot Cv-\dot M\eta(v)$ and proves the identity.
Equivalently it follows by differentiating the implicit multiplier map on
the normalized graph, so the same linear identity extends to every $v$.
Both moving terms are retained. The exact graded formula is
\[
 E_a^{-1}\dot\eta(v)
 =-\Gamma^{-1}\dot\Gamma\,y(v)
      -a^{-4}\Gamma^{-1}E_a\dot Cv.
\]
Use $\|v\|\le a^2\mathfrak n(v)$ and $\dot C=O(a)$ to prove
\eqref{v48:lift-motion}. No derivative of the inverse is estimated in
unweighted coordinates before this cancellation.
\end{proof}

\section{The original reader has a two-power loss on the quartet}
\label{v48:upper}

Let $\mathcal E_a(t)=\Ran P(t)$ be the whole intermediate quartet and use
$|u|_{a,\rm int}^2=-a^{-4}u^*Su$ there.

\begin{lemma}[A finite-jet bound with the physical rate retained]
\label{v48:jet-bound}
For every canonical direction on the normalized graph,
\begin{align}
 \|\mathscr R(v)-\iota_0\mathfrak E_h(\rho(v)^\beta)\|
 &\le C\{\|v\|+\|\eta(v)\|+\|Tv\|+a\|\rho(v)\|\},
                                                        \label{v48:metric-jet}\\
 \|\mathscr L^{\rm sp}(v)\|
 &\le C\{\|v\|+a\|\eta_A(v)\|+a^2\|\eta_W(v)\|\},
                                                        \label{v48:spatial-jet}\\
 \|\mathscr L^{\rm el}(v)\|
 &\le C\{\|v\|+\|\eta(v)\|+\|Tv\|
                       +a\|\rho_A(v)\|+a^2\|\rho_W(v)\|\}.
                                                        \label{v48:electric-jet}
\end{align}
\end{lemma}
\begin{proof}
The original canonical first velocity varies by $Tv$. Differentiating
$T=F+\mathbb JC^*\eta_t(\cdot)$ and using \eqref{v48:rho} gives the
exact second-velocity identity
\[
 \delta(\partial_t\mathfrak f)=(\dot T+T^2)v
 =F Tv+\mathbb JC^*\rho+\dot Fv+\mathbb J\dot C^*\eta(v).
\]
The last two terms are bounded by $Ca(\|v\|+\|\eta\|)$, and
$\mathbb JC^*=G+O(a)$, with the original flat constraint-generated map
$G$. Its configuration slot is
$(G\rho)^u_{ij}=\delta_i\rho_j^\beta+\delta_j\rho_i^\beta$;
the lapse-rate slot does not enter that flat configuration velocity.

In the inherited Riemann convention the electric metric block contains
\[
 \tfrac12(g_{0j,it}+g_{i0,tj}-g_{00,ij}-g_{ij,tt})
             +\hbox{Christoffel products}.
\]
The two mixed derivatives contribute ordinary spatial derivatives of
$\rho^\beta$, whereas the last term contributes its phase derivatives.
Their difference is exactly \eqref{v48:E}. The linear lapse \emph{value}
and the remaining first-velocity terms are bounded by $C(\|\eta\|+\|Tv\|)$.
All coefficients differing from their flat values are $O(a)$. Base first
metric jets are $O(a)$, so differentiated Christoffel products are bounded
by the same remainder, including $Ca\|\rho\|$. The other independent
Riemann blocks have no flat $\rho$ term. This proves
\eqref{v48:metric-jet}. Neither $r_t$ nor a second multiplier time derivative
is used in this curvature reader.

For the stationary spatial connection $z$, retain the original orders
$z_X=O(1)$, $z_{\lambda_A}=O(a)$, and $z_{\lambda_W}=O(a^2)$.
These are the universal weak first-source cancellation and the graded
stationary response, not a formula restricted to diagonal momenta.
They give the bound on $\delta z$ in \eqref{v48:spatial-jet}. Literal spatial
links are analytic functions of $z$ and its fixed shifts, proving that line.
For temporal links differentiate
$z_t=z_X\mathfrak f+z_\lambda r$. The terms with differentiated stationary
coefficients multiply $O(a)$ base velocities; the remaining terms are
$z_X Tv+z_\lambda\rho$. The temporal connection variation is bounded by
$C(\|v\|+\|\eta\|)$. Substitute in the original expression
\[
 h^{-1}\{\dot U_iU_i^{-1}+A_0-U_iA_{0,+i}U_i^{-1}\}.
\]
There is no $\dot A_0$ in this expression. All literal spatial shifts and
fixed factors of $h=1/3$ remain. This proves \eqref{v48:electric-jet}.
\end{proof}

\begin{theorem}[Rate and original-curvature transfer on the intermediate space]
\label{v48:reader-transfer}
Uniformly on the admitted histories, for $u\in\mathcal E_a(t)$,
\begin{equation}
 \boxed{
 \|\rho_A(u)\|\le Ca^{-1}|u|_{a,\rm int},\qquad
 \|\rho_W(u)\|\le Ca^{-2}|u|_{a,\rm int}.}
 \label{v48:rate-bounds}
\end{equation}
The original readers satisfy
\begin{equation}
 \boxed{
 \|\mathscr R(u)\|\le Ca^{-2}|u|_{a,\rm int},\quad
 \|\mathscr L^{\rm el}(u)\|\le C|u|_{a,\rm int},\quad
 \|\mathscr L^{\rm sp}(u)\|\le Ca^2|u|_{a,\rm int}.}
 \label{v48:reader-bounds}
\end{equation}
More precisely,
\begin{equation}
 \boxed{a^2\mathscr R(u)
   =\iota_0\mathfrak E_h\bigl(\eta(a^2Tu)^\beta\bigr)
                       +O(a)|u|_{a,\rm int}.}
 \label{v48:leading-reader}
\end{equation}
\end{theorem}
\begin{proof}
The exact imaginary semisimple roots and the definite norm of V46 give
$|Tu|_{a,\rm int}\le Ca^{-2}|u|_{a,\rm int}$ on the entire quartet,
also across an internal frequency coincidence. Apply its original lifting
bounds to $Tu$ and apply \cref{v48:rate} to $u$. The two rate estimates
follow, with the motion of the lift one order smaller in its balanced norm.
Also $\|Tu\|\le C|u|_{a,\rm int}$, $\|u\|\le Ca^2|u|_{a,\rm int}$,
$\|\eta_A(u)\|\le Ca|u|_{a,\rm int}$, and
$\|\eta_W(u)\|\le C|u|_{a,\rm int}$.
Substitution in \cref{v48:jet-bound} proves every bound. Its metric remainder
is at most $Ca^{-1}|u|_{a,\rm int}$. Multiplication by $a^2$, followed
by \eqref{v48:rho} and \eqref{v48:lift-motion}, proves
\eqref{v48:leading-reader}. No leading moving term is silently set to zero:
it is smaller on this spectral scale by a proved bound. The slow-scale
calculation in V38 had a different balance.
\end{proof}

\section{No intermediate oscillation is invisible to the metric reader}
\label{v48:observability}

Retain V45's actual limiting pair on
$\mathcal V=(\mathcal A\oplus\mathcal H)\oplus\mathcal Z_g$:
\begin{equation}
 \mathsf G\succ0,\qquad
 \mathsf N=\begin{pmatrix}0_{81\times81}&U\\-U^*&V\end{pmatrix},
 \quad\rank U=2,\quad V^*=-V,\quad
 J_0=-\mathsf G^{-1}\mathsf N.
 \label{v48:limiting-pair}
\end{equation}
The scalar in $V$ is the actual, still unevaluated coefficient from V45,
not a chosen value. Put $\mathcal V_{\rm int}=\Ran J_0$, of dimension
four. Its intersection with $\ker J_0$ is zero, because $J_0$ is skew-adjoint
in $\mathsf G$. Let $\pi_g$ select the last two coordinates and let $I_g$
synthesize them into their physical mean-zero weak shifts. Define
\[
 B_0=\mathfrak E_h I_g\pi_gJ_0\big|_{\mathcal V_{\rm int}}.
\]

\begin{lemma}[Value and one generator derivative observe all four directions]
\label{v48:two-readers}
The map
\begin{equation}
 x\longmapsto(B_0x,\,B_0J_0x),\qquad x\in\mathcal V_{\rm int},
 \label{v48:observable-pair}
\end{equation}
is injective. For every fixed $L>0$, its finite-window Gramian
\begin{equation}
 \int_0^L e^{sJ_0^*}B_0^*B_0e^{sJ_0}\,\dd s
 \label{v48:Gramian}
\end{equation}
is positive definite on $\mathcal V_{\rm int}$. The statements include
coincident nonzero frequencies.
\end{lemma}
\begin{proof}
The map $\mathfrak E_h I_g$ is injective: its shifts are mean-zero,
whereas the kernel of \eqref{v48:E} consists of constants. Suppose both
entries of \eqref{v48:observable-pair} vanish and set $w=J_0x$. Then
$\pi_gw=0$ and $\pi_gJ_0w=0$. The first 81 rows of
$\mathsf GJ_0w=-\mathsf Nw$ now give
$\mathsf G_{tt}(J_0w)_t=0$. This principal metric block is positive
and hence invertible. Therefore $J_0w=0$. Both $w$ and $x$ belong to
$\Ran J_0$, where $J_0$ is invertible, so $w=x=0$.
If the Gramian annihilates a vector, its continuous nonnegative integrand
has $B_0e^{sJ_0}x=0$ throughout the interval. Evaluating at zero and
differentiating once gives \eqref{v48:observable-pair}, so $x=0$.
Compactness of the unit sphere gives a strictly positive lower bound.
This is a finite-dimensional argument, not an imported infinite-dimensional
observability estimate.
\end{proof}

\begin{theorem}[Sharp two-power metric-reader cost for the original quartet]
\label{v48:sharp}
At each admitted base time freeze the original matrices and the original
reader differential. For every $u\in\mathcal E_a(t)$,
\begin{equation}
 \boxed{
 c|u|_{a,\rm int}^2\le
 a^4\|\mathscr R_t^0(u)\|^2+
 a^8\|\mathscr R_t^0(Tu)\|^2
 \le C|u|_{a,\rm int}^2.}
 \label{v48:jet-observability}
\end{equation}
For each fixed $L>0$ there are positive $c_L,C_L$, independent of $a$,
such that
\begin{equation}
 \boxed{
 c_L|u|_{a,\rm int}^2\le
 \int_0^L\|a^2\mathscr R_t^0(e^{s a^2T(t)}u)\|^2\,\dd s
 \le C_L|u|_{a,\rm int}^2.}
 \label{v48:window}
\end{equation}
In particular the operator norm of $\mathscr R_t^0$ on the quartet is
bounded above and below by positive multiples of $a^{-2}$.
\end{theorem}
\begin{proof}
We first specify the finite-to-limiting identification rather than infer
observability from four approximate eigenvalues. On the exact intermediate
root spaces, V45's eliminated weak block reconstructs as $O(a)$ times
the retained multiplier vector. The latter solves
$(\tau\mathsf G+\mathsf N+O(a))x=0$ with $\tau$ in a fixed annulus.
Projection onto $\ker J_0$ shows that its component outside
$\mathcal V_{\rm int}$ is $O(a)$. V46's exact energy formula and
semisimplicity imply that the map which sends an original intermediate
vector to the $\mathsf G$-projection of its retained balanced multiplier
vector is an isomorphism onto $\mathcal V_{\rm int}$ for small $a$.
Denote its inverse by $\Lambda_a$. Its norm comparison is uniform.

On every root kernel the same retained equation gives
$\Lambda_a^{-1}a^2T\Lambda_a=J_0+O(a)$. The formula holds on the full
space: distinct eigencomponents are orthogonal in the original definite
energy, and repeated roots are treated on their whole kernels. At most four
components are summed. No ill-conditioned individual eigenbasis estimate
is used. The balanced weak multiplier of $\Lambda_ax$ has limiting value
$I_Hx_H+I_g\pi_gx$. Constant shifts are annihilated by $\mathfrak E_h$.
Consequently \eqref{v48:leading-reader} gives the operator identities
\begin{equation}
 \Lambda_a^{-1}a^2T\Lambda_a=J_0+O(a),\qquad
 a^2\mathscr R_t^0\Lambda_a=B_0+O(a).
 \label{v48:finite-identification}
\end{equation}
All comparisons are uniform on the same compact coefficient set as V47.

Apply \cref{v48:two-readers} to the stacked map and perturb by
\eqref{v48:finite-identification}. A positive lower singular value persists,
proving \eqref{v48:jet-observability}. On each bounded $s$ interval,
variation of constants for the two finite four-dimensional matrices gives
uniform convergence of their exponentials. The positive Gramian in
\eqref{v48:Gramian} and its continuous perturbation prove
\eqref{v48:window}. Finally the frozen canonical evolution is an isometry
in the intermediate energy. If the norm of $a^2\mathscr R_t^0$ tended
to zero, the integral upper bound by $L$ times that squared norm would
contradict \eqref{v48:window}. The preceding transfer theorem gives the
matching upper bound.
\end{proof}

The two observations in \eqref{v48:jet-observability} are the original
reader applied to a vector and to its \emph{frozen generator} image.
The second is not the physical derivative of the reader along a moving
history: that derivative also contains $\dot{\mathscr R}_t^0u$.
Similarly \eqref{v48:window} freezes the base coefficients; it is not a
new nonlinear time interval. The theorem need not make the instantaneous
reader injective on all four real coordinates. It proves that none of their
oscillations is permanently invisible and that the two-power operator loss
is genuine for this original source.

\section{The sharp loss persists on the full original constraint tangent}
\label{v48:radial}

\begin{lemma}[A smaller radial correction for the intermediate space]
\label{v48:radial-correction}
The bounded normal $W_{a,t}$ of V37 can be chosen uniformly along the
admitted histories so that
\[
 CW_{a,t}=0,\qquad \langle g_{a,t},W_{a,t}\rangle=a,
 \qquad \|W_{a,t}\|+\|TW_{a,t}\|\le C,
\]
where $g_{a,t}=\nabla H^{\rm red}$. For $u\in\mathcal E_a(t)$ set
\begin{equation}
 \Pi_{\rm rad}u=u-\frac{\langle g_{a,t},u\rangle}{a}W_{a,t}.
 \label{v48:radial-map}
\end{equation}
It satisfies every original tangent constraint and leaves its multiplier
lift unchanged. Moreover
\begin{equation}
 \begin{gathered}
 |\langle g_{a,t},u\rangle|\le Ca^5|u|_{a,\rm int},\
 \|\Pi_{\rm rad}u-u\|\le Ca^4|u|_{a,\rm int},\
 \mathfrak n(\Pi_{\rm rad}u-u)\le Ca^2|u|_{a,\rm int}.
 \end{gathered}
 \label{v48:radial-cost}
\end{equation}
\begin{equation}
 \begin{gathered}
 \|\mathscr R(\Pi_{\rm rad}u-u)\|\le C|u|_{a,\rm int},\
 \|\mathscr L^{\rm el}(\Pi_{\rm rad}u-u)\|\le Ca^2|u|_{a,\rm int}.
 \end{gathered}
 \label{v48:radial-output}
\end{equation}
\end{lemma}
\begin{proof}
The 108-row canonical normal matrix, with its four means divided by $a$,
retains V37's surjective limit throughout the $O(a^3)$ displacement of
these histories. Its fixed right inverse therefore remains bounded. The
same target as in \cref{v37:lossless-normal} gives the displayed exact
normal identities, without differentiating that normal in time.

Along the original autonomous reduced flow,
$\dot g=S\mathfrak f=Q\mathfrak f+C^*r=O(a)$; here
$C\mathfrak f=-Mr$ is used before estimating the inverse. Also
$g^*T=-\dot g^*$. The intermediate inverse has norm at most $Ca^2$ in
its definite energy, so its canonical lift obeys
$\|T^{-1}u\|\le Ca^4|u|_{a,\rm int}$. Thus
$\langle g,u\rangle=-\langle\dot g,T^{-1}u\rangle=O(a^5)|u|_{a,\rm int}$.
This proves the field and record costs and exact radial tangency.

For the normal itself $\eta(W)=0$, $TW=O(1)$, and
$\rho(W)=-M^{-1}\{CTW+\dot CW\}$. The graded inverse gives the
sufficient bounds $\rho_A(W)=O(a^{-3})$, $\rho_W(W)=O(a^{-4})$.
The original jet estimates hence give
$\|\mathscr R(W)\|\le Ca^{-4}$,
$\|\mathscr L^{\rm el}(W)\|\le Ca^{-2}$, and
$\|\mathscr L^{\rm sp}(W)\|\le C$.
Multiplying by the $O(a^4)|u|_{a,\rm int}$ correction coefficient
proves \eqref{v48:radial-output}. These are sufficient upper bounds,
not claimed optimal normal-reader exponents.
\end{proof}

\begin{corollary}[Admissible witnesses at the exact precision of V47's error]
\label{v48:sharp-witness}
At every sufficiently small amplitude there is a real tangent
$(v_a,\eta(v_a))$ to the full original constraint leaf such that
\begin{equation}
 \boxed{
 \|v_a\|\le Ca^4,\quad\|\eta_A(v_a)\|\le Ca^3,\quad
 \|\eta_W(v_a)\|\le Ca^2,\qquad
 c\le\|\mathscr R^0(v_a)\|\le C,
 }
 \label{v48:witness-metric}
\end{equation}
while
\begin{equation}
 \|\mathscr L^{\rm el}(v_a)\|\le Ca^2,\qquad
 \|\mathscr L^{\rm sp}(v_a)\|\le Ca^4.
 \label{v48:witness-links}
\end{equation}
Each is realized by a local curve of exactly constrained original initial
records, with its own local normalized exact-action histories.
\end{corollary}
\begin{proof}
Choose a real unit intermediate vector $u_a$ realizing the lower operator
bound in \cref{v48:sharp}, and set $v_a=a^2\Pi_{\rm rad}u_a$.
The radial correction changes its metric output by $O(a^2)$, whereas
$a^2\|\mathscr R^0(u_a)\|\ge c$. All other estimates follow from
\cref{v48:reader-transfer,v48:radial-correction}. The exact constraint leaf
is a smooth submanifold at each nonzero amplitude by the original rank chart.
Every vector tangent to it is the derivative of a local curve in that leaf.
The unchanged one-null continuation theorem supplies each sufficiently nearby
initial point's own local history. No amplitude-uniform parameter radius,
finite perturbation estimate, or common lifespan for those varied points
is asserted.
\end{proof}

This is not a lower bound on the \emph{actual} V47 tracking error. It proves
that the sizes stated for that error, even augmented by all original tangent
constraints, cannot by themselves imply vanishing metric curvature.
The witnesses are different admissible differential directions. There is no
claim about a curvature-squared scalar or finite-amplitude instability.

\section{Populated metric and link bounds for the coupled spectral comparison}
\label{v48:applications}

Take an original constrained variation $\xi_t=T\xi$, put
$u=P\xi$, $v=(I-P)\xi$, and let
$\bar u(t)=\mathcal U_a(t,0)P(0)\xi(0)$ be V47's analytical transport.
Set $e=u-\bar u\in\Ran P(t)$ and $\widehat\xi=v+\bar u=\xi-e$.
All reader expressions below use \eqref{v48:reader-definitions} at the
\emph{same original base point}.

\begin{theorem}[Original differential reader of the selected comparison]
\label{v48:selected}
For V47's actual selected parameter tangent, uniformly on $0\le t\le ca^2$,
\begin{equation}
 \boxed{
 \|\mathscr R(\xi)-\mathscr R(\widehat\xi)\|\le C,\quad
 \|\mathscr L^{\rm el}(\xi)-\mathscr L^{\rm el}(\widehat\xi)\|\le Ca^2,
 \quad
 \|\mathscr L^{\rm sp}(\xi)-\mathscr L^{\rm sp}(\widehat\xi)\|\le Ca^4.
 }
 \label{v48:selected-output}
\end{equation}
The complete actual reader remains exactly the sum of its two spectral
contributions. The $O(a)$ metric and $O(a^2)$ literal reader bounds of
V41 on that full tangent are retained, not replaced by the larger
componentwise estimates.
\end{theorem}
\begin{proof}
The reader differentials are linear, so the three differences are exactly
the corresponding readers of $e$. V47 proves $|e|_{a,\rm int}\le Ca^2$.
Apply \cref{v48:reader-transfer}. The last statement is linearity applied
to $\xi=P\xi+(I-P)\xi$, together with V41's differentiated original
observation bounds. Its two separate spectral contributions can cancel;
no estimate here rules that out or proves a nonzero comparison error.
\end{proof}

\begin{theorem}[Stable-fibre intermediate contamination is geometrically small]
\label{v48:stable}
For the normalized actual stable-fibre tangents of
\cref{v47:stable-tangent}, the \emph{entire} intermediate contribution obeys
\begin{equation}
 \boxed{
 \sup_{t\le ca^2}\|\mathscr R(P\xi)\|\le Ca^2,\quad
 \sup_{t\le ca^2}\|\mathscr L^{\rm el}(P\xi)\|\le Ca^4,\quad
 \sup_{t\le ca^2}\|\mathscr L^{\rm sp}(P\xi)\|\le Ca^6.
 }
 \label{v48:stable-global}
\end{equation}
For $0\le t\le\sigma ca^2$, any fixed $0<\sigma<1$, each right side
has the additional factor $C_\sigma e^{-\beta t/a^3}$. For the spectral
transport error $e$ itself the same global powers hold, without an asserted
decay factor.
\end{theorem}
\begin{proof}
V47 proves $|P\xi|_{a,\rm int}\le Ca^4$ on the whole interval and
$|P\xi|_{a,\rm int}\le C_\sigma a^4e^{-\beta t/a^3}$ on the earlier
portion. It also proves $|e|_{a,\rm int}\le Ca^4$. Apply the three
reader bounds separately. The full stable history and its original tangent
remain the physical objects. The projected term here is their linear
contribution to the original observation, not a new radial-constraint-tangent
history obtained by projecting the dynamics.
\end{proof}

\section{A comparison path has additional, exactly identifiable jet defects}
\label{v48:path}

An important distinction remains even after \cref{v48:selected}. Let $w(t)$
be any $C^2$ canonical comparison curve along a base history and give it the
algebraic multiplier lift $\eta_t(w)$. Define its canonical equation defect
\[
 f=w_t-Tw.
\]
Then exactly
\begin{equation}
 \boxed{
 \partial_t\eta_t(w)=\rho_t(w)+\eta_t(f),\qquad
 w_{tt}=(\dot T+T^2)w+Tf+\dot f.
 }
 \label{v48:path-defects}
\end{equation}
The first-velocity defect is $f$. These identities distinguish two operations:
applying the original \emph{action} reader at each point, and differentiating
the interpolation of a comparison curve that does not solve the variational
equation.

\begin{proposition}[Complete kinematic correction to the differential reader]
\label{v48:kinematic}
There are bounded, base-dependent finite jet maps $J_X,J_\lambda,J_{XX}$
from the original metric construction such that its path-based linear
curvature satisfies
\begin{equation}
 \mathscr R_{\rm path}(w)-\mathscr R_t(w)
 =J_X f+J_\lambda\eta_t(f)+J_{XX}(Tf+\dot f).
 \label{v48:path-reader}
\end{equation}
All spatial differentiations are included in these finite maps. The original
spatial-link reader has zero path correction; the temporal-link reader has
correction obtained by applying its bounded velocity derivative to
\begin{equation}
                         z_X f+z_\lambda\eta_t(f).
 \label{v48:path-link}
\end{equation}
\end{proposition}
\begin{proof}
Differentiate the lift at fixed base time to get the first identity of
\eqref{v48:path-defects}; differentiating $w_t=Tw+f$ gives the second.
The linearized metric reader is affine linear in the variations of the
first canonical velocity, the first multiplier velocity, and the second
canonical velocity, with the zero-order record held fixed. Subtract its
values at the two jet lists. Their differences are exactly
$f,\eta_t(f),Tf+\dot f$, proving \eqref{v48:path-reader}.
The coefficients are bounded since the base metric and its needed jets
remain in the original small chart. For links, the only time derivative
in the original electric read is that of the spatial connection. Its
linearized path velocity differs from its action velocity by
\eqref{v48:path-link}. Spatial reads have no time derivative.
\end{proof}

For V47's coupled comparison $\widehat\xi=\xi-e$, the defect is explicitly
\begin{equation}
 f_{\widehat\xi}
 =-[\dot P,P]e-\dot Pv.
 \label{v48:comparison-defect}
\end{equation}
Although V47 controls $\dot P$, it does not bound the derivative of this
entire expression. Thus \eqref{v48:selected-output} must not be advertised
as a curvature estimate for the interpolated comparison \emph{path} without
also paying \eqref{v48:path-reader}. The actual full variation has $f=0$,
so all these kinematic corrections vanish there. This provides an exact
residual formula instead of silently assigning a physical trajectory to
an analytical spectral transport.

\section{What is closed and the remaining source-specific target}
\label{v48:status}

The original rate and curvature differential on the quartet are now bounded,
and their two-power metric cost is proved sharp using this source's actual
rank-four positive pencil. The lower bound does not rely on computing the
still unevaluated scalar in $K_2$. A radial correction makes the sharpness
witnesses tangent to all original constraints while preserving the multiplier
component. On those differential directions the physical sizes of V47's
tracking error coexist with a nonvanishing metric reader and vanishing
literal-link readers. This is a new intermediate-scale statement, not a
repetition of V39's finite slow-direction perturbation theorem.

The stable-fibre contamination is now controlled in the original readers:
its metric contribution is $O(a^2)$ and its temporal-link contribution is
$O(a^4)$, with decay on each earlier fraction of the common interval.
The selected comparison has only an $O(1)$ general metric-error bound from
its current energy estimate. The sharpness witnesses do not prove that its
actual error saturates this bound. A vanishing result needs a better estimate
on the actual output, an additional cancellation in its forcing, or a more
accurate comparison. A claim about the comparison path itself must also
control the exact kinematic corrections above.

No new physical lifespan, attracting set, or spatial-refinement theorem is
obtained. The admitted interval remains $ca^2$ at $N=3$. The rapid and slow
real sectors, higher observation derivatives, normalized metric/link
identification and a common positive physical interval remain separate.
The raw-action Einstein continuum bridge is neither proved nor disproved.

\paragraph{Verification.}
The companion checker verifies the noncommuting lift/rate identities, the
complete differentiated canonical acceleration, the original flat metric
reader on every retained nonzero three-site Fourier mode, the record and
radial amplitude powers, the exact path-defect formulas, and rational
rank-four observation examples including repeated frequencies. Those examples
check the general linear algebra, not a new gravitational coefficient bank.
The all-small-amplitude bounds and the source-specific visibility theorem
are the written proofs using the named inherited action inputs. The audit
records exactly which predecessor checks were run. There is no numerical
trajectory, new intermediate-frequency enclosure, numerical observation
constant, finite-error radius, or proof-assistant formalization. All 47
preceding marked bodies are preserved byte for byte.

\begin{thebibliography}{9}
\raggedright
\bibitem{v48:Hautus}
M. L. J. Hautus, \emph{Controllability and observability conditions of linear
autonomous systems}, Indagationes Mathematicae (Proceedings) \textbf{72}(5)
(1969), 443--448. The author-institution publication record was checked.
This is background for the established finite-dimensional observability
method; the two-observation proof, original reader, and amplitude bounds
used here are supplied explicitly.
\end{thebibliography}
% END IMMUTABLE EXACT-ACTION BRIDGE V48 BODY


% Future iterations append here.
% BEGIN IMMUTABLE EXACT-ACTION BRIDGE V49 BODY
\clearpage
\renewcommand{\thesection}{\arabic{section}}
\renewcommand{\theHsection}{bridgeV49.\arabic{section}}
\setcounter{section}{396}
\section{Retain a constrained source predictor before taking the curvature reader}
\label{v49:context}

The two-power reader loss in \cref{v48:reader-transfer,v48:sharp-witness}
prevents an $O(a^2)$ intermediate-energy error from implying a vanishing
metric-curvature error. It does not say that the actual selected variation
has large curvature, nor that its spectral comparison cannot be improved.
The next step here changes the \emph{analytical comparison}, not the original
action, its initial selection, or its evolution.

The useful input from the related covector-response work is that the forcing
and its actual incidence must be retained before applying a singular response.
We do not import that work's coefficient pencil or its source-derivative
constants. Instead, on the original histories already used in V47--V48, we
construct a predictor from the initial complete tangent record and the
\emph{evolving base coefficients}. A canonical normal solve makes this
predictor tangent to all 108 original constraints with its initial multiplier
variation held fixed. It approximates the actual tangent in the complete-record
norm by $O(at)$, uniformly for $0\le t\le ca^2$. Retaining its contribution
to the quartet forcing improves the actual intermediate tracking error from
$O(a^2)$ to $O(a^4)$ at the endpoint. The original metric error is then
$O(a^2)$, rather than the former general $O(1)$ bound.

A separate twice-differentiated saddle calculation pays the path-reader
question left in V48. It gives the same metric and literal-link orders for
interpolation of the corrected comparison, not only for the instantaneous
action reader. A final canonical normal adjustment restores the radial tangent
constraint without changing the multiplier component or these orders.
The full comparison retains the \emph{actual complementary component}; it is
not an autonomous approximation of all 324 unknown canonical coordinates.
The quartet correction itself needs only the base history and the initial
tangent, and is causal along that supplied base.

\paragraph{Companion audit and nonduplication.}
The four newly supplied files were inspected at their current development
blocks and at the antecedents needed to interpret them. Their results are
retained with their original scope:
\begin{center}\small
\begin{tabularx}{\textwidth}{@{}>{\raggedright\arraybackslash}p{.22\textwidth}Y@{}}
\toprule
Companion & Transfer and boundary\\\midrule
Native stationarity V54 &
\src{v54:thm:descriptor}, \src{v54:lem:moments}, and
\src{v54:thm:lift} retain two differentiated covector coordinates and pay
both fast scales of their specified coefficient pencil. Its source is an
explicit comparison covector and its fast dimensions are different. No
root, Green constant, or derivative estimate is transferred here.\\[3pt]
Yang--Mills f17 V100 &
\src{r100:resonance-theorem} and \src{r100:no-inverse} exhibit a native
resonance across every one-link cutoff; \src{r100:spectator-return} retains
energy dependence of a Schur return. Here the whole intermediate quartet
and its complementary forcing are retained; no whole-band homological
inverse or instantaneous replacement of the response is assumed.\\[3pt]
Perturbative ESM V64 &
\src{prop:v64-marks} and \src{thm:v64-hessian} distinguish complete
noncommuting derivatives and an integrable common action from independently
specified blocks. Here both differentiated constraint terms are retained.
The predictor is not declared to be a Hessian of a modified action.\\[3pt]
SM source selection V52 &
\src{sel52:thm:prepface} and \src{sel52:prop:bank} distinguish positive
executable preparations from a signed reconstruction of their responses.
The normal projector and the quartet comparison below are analytical
operations, not newly admitted primitive physical transitions.\\
\bottomrule
\end{tabularx}
\end{center}
The detailed audit records literal source filenames, line ranges, hashes,
and the narrower claims inspected. It is not a revalidation of every proof
in the four cumulative files. The standard parallel-transport method is
already cited in \cite{v47:ASY}; its original publication record was checked.
No external adiabatic or numerical-convergence theorem supplies the estimates
below. In particular this is not a claim of a new general defect-correction
method.

\paragraph{Inherited gravitational inputs.}
We retain the same $N=3$ root and the selected parameter histories of
\cref{v41:selection-theorem,v41:geometry-theorem}; the full canonical normal
submersion of \cref{v37:lossless-normal}; the complete-record resolvent and
projection bounds of V47; and the original readers and their sharp costs in
V48. All constants are uniform in $a$ on this stated fixed-cutoff family.
The parameter $a>0$ is held fixed during physical time differentiation. No
claim below changes the established physical interval $[0,ca^2]$.

\section{Complete tangent bounds and a uniformly scaled canonical normal inverse}
\label{v49:normals}

Let $(X_a(t,\vartheta),\lambda_a(t,\vartheta))$ be the existing selected
histories. Fix a parameter in a compact smaller interval and write
\[
 \xi=\partial_\vartheta X_a,\qquad
 \eta=\partial_\vartheta\lambda_a.
\]
All multiplier variations have zero mean lapse, and are represented in the
same fixed 107-dimensional complement. Use the physical active/weak split
of dimensions $78+29$. Along the base set
\[
 F=\mathbb JQ,\qquad C=H_{uX},\qquad M=H_{uu},\qquad
 T=F-\mathbb JC^*M^{-1}C,\qquad E_a=\diag(aI_A,I_W).
\]
These letters denote exactly the matrices in V47--V48. In particular $F$
is the bounded \emph{uneliminated linearized} canonical matrix, not a new
nonlinear vector field. The actual tangent equations are
\begin{equation}
 \xi_t=F\xi+\mathbb JC^*\eta=T\xi,\qquad C\xi+M\eta=0.
 \label{v49:tangent}
\end{equation}
The radial tangent equation is retained separately when the normalized
107-row system is used.

The already proved parameter bounds imply, uniformly on $[0,ca^2]$,
\begin{equation}
 \begin{gathered}
 \|\xi\|\le Ca^3,\quad \|\eta_A\|\le Ca^3,\quad
 \|\eta_W\|\le Ca^2,\\
 \|\xi_t\|\le Ca^2,\quad
 \|\eta_{A,t}\|\le Ca^2,\quad \|\eta_{W,t}\|\le Ca.
 \end{gathered}
 \label{v49:actual-budgets}
\end{equation}
For the first active bound it would suffice to use $Ca^2$; the stronger
stated bound is in the proof of \cref{v41:geometry-theorem}. These are
bounds on the actual selected tangents, not hypotheses about arbitrary
solutions of the full variational system.

Let $\mathcal C=D_XP$ contain all 108 original constraint derivatives and
let $\mathcal M=D_\lambda P$ have its \emph{columns} restricted to the
normalized multiplier complement. Thus
\begin{equation}
                 \mathcal C\xi+\mathcal M\eta=0
 \label{v49:full-tangent}
\end{equation}
contains the radial row as well. Use the original Fourier row order with
104 nonconstant rows followed by four means, and put
$S_a^{\rm row}=\diag(I_{104},a^{-1}I_4)$.

\begin{lemma}[Canonical normal solve with its source orders retained]
\label{v49:normal-lemma}
There is a fixed canonical normal inclusion $J_c$ with 108 columns for which
\begin{equation}
 R_c(t)=J_c\{S_a^{\rm row}\mathcal C(t)J_c\}^{-1}S_a^{\rm row},
 \qquad \mathcal C R_c=I_{108}
 \label{v49:right-inverse}
\end{equation}
is defined throughout the selected interval for all sufficiently small $a$.
Set $N_c=R_c\mathcal C$ and $B_c=R_c\mathcal M$. Then
\begin{equation}
 \|N_c\|+\|\dot N_c\|\le C,
 \label{v49:normal-projection-bound}
\end{equation}
and, for every multiplier vector $b=(b_A,b_W)$,
\begin{equation}
 \boxed{\|B_cb\|+\|\dot B_cb\|
          \le C(a^2\|b_A\|+a^3\|b_W\|).}
 \label{v49:normal-mass-bound}
\end{equation}
The four mean rows are not deleted in this construction.
\end{lemma}
\begin{proof}
The row-scaled matrix has limit
$(L_{\rm nc},P_0DP_2(Y_*))$, of rank 108 by the actual mean calibration
and canonical submersion in \cref{v37:lossless-normal}. Choose its fixed
nonzero minor to define $J_c$. The selected records obey
$X/a=Y_*+O(a)$ and $(\lambda-e)/a=\ell_*+O(a)$ uniformly; their normalized
first velocities are bounded. Hence the inverse in \eqref{v49:right-inverse}
and its first physical-time derivative are bounded. The matrix
$S_a^{\rm row}\mathcal C$ and its first derivative are bounded as well.
This proves \eqref{v49:normal-projection-bound} and the right-inverse identity.

For completeness, the extra mean-row order of the mass is essential.
The original source factorization is $\mathcal M=J_{\rm full}^*K^{-1}J$,
with the same stationary tangent Hessian as in V20 and V30. Its active
columns are $O(a)$, while its weak and all four constant columns are
$O(a^2)$. The constant-lapse order also follows directly from
$\mathcal M\lambda=0$ and $\lambda-e=O(a)$. Their first time derivatives
have the same orders on the normalized chart with bounded velocity.
Consequently the nonconstant rows of $\mathcal M b$ are bounded by
$C(a^2\|b_A\|+a^3\|b_W\|)$ and its four mean rows by
$C(a^3\|b_A\|+a^4\|b_W\|)$. The same statement holds for
$\partial_t\mathcal M$. Multiplication by $S_a^{\rm row}$ and the bounded
normal inverse, including its derivative, gives \eqref{v49:normal-mass-bound}.
This is the original graded source argument with mean rows retained; a
crude unweighted inverse of $\mathcal C$ would lose one power needlessly.
\end{proof}

For later use retain the original complete-record norm
\begin{equation}
 \mathfrak n_t(v)^2=a^{-4}\|v\|^2+
             \|E_a^{-1}\eta_t(v)\|^2,
 \qquad \eta_t(v)=-M^{-1}Cv.
 \label{v49:norm}
\end{equation}
The subscript $a$ is suppressed, not changed. The projection bounds and
norm transport of V47 apply in this norm.

\section{An exactly constraint-tangent predictor from the initial complete record}
\label{v49:predictor}

Write $\xi_0=\xi(0)$ and $\eta_0=\eta(0)$ in the fixed coordinates.
Define a first integral predictor using the evolving base coefficients:
\begin{equation}
 \zeta(t)=\xi_0+\int_0^t
              \{F(s)\xi_0+\mathbb JC(s)^*\eta_0\}\,ds.
 \label{v49:primitive-predictor}
\end{equation}
It is not the solution of an altered action equation. It is a vector
integral calculated from the initial tangent and the base history.
Complete it by the original canonical normals:
\begin{equation}
 \boxed{v_{\rm p}(t)=\zeta(t)-R_c(t)
            \{\mathcal C(t)\zeta(t)+\mathcal M(t)\eta_0\}.}
 \label{v49:predictor-formula}
\end{equation}
In particular no singular multiplier inversion is used to repair its
constraint residual.

\begin{theorem}[One extra amplitude power from the complete constraint incidence]
\label{v49:predictor-theorem}
The predictor satisfies, exactly at every time,
\begin{equation}
 v_{\rm p}(0)=\xi_0,\qquad
 \mathcal C(t)v_{\rm p}(t)+\mathcal M(t)\eta_0=0.
 \label{v49:predictor-constraint}
\end{equation}
Let $r=\xi-v_{\rm p}$. On $0\le t\le ca^2$,
\begin{equation}
 \begin{gathered}
 \|r(t)\|\le C(at^2+a^4t),\qquad
 \eta_t(r)=\eta(t)-\eta_0,\\
 \boxed{\mathfrak n_t(r(t))\le Cat,\qquad
                    \mathfrak n_t(r_t(t))\le Ca.}
 \end{gathered}
 \label{v49:predictor-errors}
\end{equation}
All constants are independent of amplitude on the retained family.
\end{theorem}
\begin{proof}
The right-inverse identity gives \eqref{v49:predictor-constraint} and its
initial value, using \eqref{v49:full-tangent} at zero. Set
$\delta=\xi-\zeta$ and $\Delta\eta=\eta-\eta_0$. The actual canonical
equation gives exactly
\[
 \delta_t=F(t)(\xi(t)-\xi_0)+\mathbb JC(t)^*\Delta\eta(t).
\]
By \eqref{v49:actual-budgets},
\[
 \|\xi(t)-\xi_0\|\le Ca^2t,\quad
 \|\Delta\eta_A(t)\|\le Ca^2t,\quad
 \|\Delta\eta_W(t)\|\le Cat.
\]
Bounded $F,C$ imply $\|\delta_t\|\le Cat$ and
$\|\delta\|\le Cat^2$. Substituting the \emph{complete} actual
constraint in \eqref{v49:predictor-formula} gives the useful exact identity
\begin{equation}
                   r=(I-N_c)\delta-B_c\Delta\eta.
 \label{v49:remainder-identity}
\end{equation}
The bounds on $N_c,B_c$ give the canonical estimate in
\eqref{v49:predictor-errors}. Subtracting the two full tangent equations
gives $\mathcal Cr=-\mathcal M\Delta\eta$; its normalized rows imply
$\eta_t(r)=\Delta\eta$ exactly. Therefore
\[
 \mathfrak n_t(r)\le C(t^2/a+a^2t+at)\le Cat,
\]
since $t\le ca^2$. This is where preserving the multiplier component in the
normal correction avoids the unweighted $M^{-1}$ loss.

Differentiate \eqref{v49:remainder-identity}. The first derivative bounds of
\cref{v49:normal-lemma} and the actual first rate bounds yield
$\|r_t\|\le C(at+a^4)$. Put $y_t(v)=E_a^{-1}\eta_t(v)$.
Differentiating $y_t(r)=E_a^{-1}\Delta\eta$ and using the original
fixed-input lift-motion estimate \eqref{v48:lift-motion} gives
\[
 \|y_t(r_t)\|\le
 C a+C a^{-1}\mathfrak n_t(r)\le C(a+t)\le Ca.
\]
Also $a^{-2}\|r_t\|\le C(t/a+a^2)\le Ca$. These prove the second
complete-record estimate. No second time derivative of the actual parameter
variation is used in this argument.
\end{proof}

For comparison, keeping only the constant canonical predictor $\xi_0$
gives $\mathfrak n_t(\xi(t)-\xi_0)\le Ct$. Indeed
$\mathfrak n_t(\xi_t)\le C$ follows by differentiating its actual multiplier
lift and using \eqref{v49:actual-budgets}; integrate with V47's norm transport.
The normal-completed predictor improves that remainder by a factor $a$.
It satisfies the original tangent constraints but need not satisfy the
original variational equation. The next step pays, rather than suppresses,
that distinction.

\section{A causal four-dimensional forcing correction with vanishing original-reader error}
\label{v49:correction}

Let $P(t)$ be the same complete intermediate quartet projection as V47--V48,
and set $K_P=[\dot P,P]$. Its analytical transport is
$\mathcal U_a$, generated by $T+K_P$ on the moving quartet. Define
\begin{equation}
 \boxed{\begin{aligned}
 \dot u_{\rm p}&=(T+K_P)u_{\rm p}
                  +\dot P(I-P)v_{\rm p},\\
 u_{\rm p}(0)&=P(0)\xi_0.
 \end{aligned}}
 \label{v49:forced-correction}
\end{equation}
The forcing is in $\Ran P(t)$, and the usual intertwining identity makes
$u_{\rm p}(t)\in\Ran P(t)$. Thus this is a four-dimensional moving-subspace
response. Its input uses the base history and the initial complete tangent;
it does not require the unknown actual complementary tangent as a forcing.
All coefficient evaluations and integrals are at times no later than $t$.

\begin{theorem}[Curvature-accurate correction of the selected spectral comparison]
\label{v49:corrected-reader}
For the actual selected parameter variation put
\[
 u=P\xi,\quad e=u-u_{\rm p},\qquad
 \widehat\xi=(I-P)\xi+u_{\rm p}=\xi-e.
\]
Uniformly on $0\le t\le ca^2$,
\begin{equation}
 \boxed{|e(t)|_{a,\rm int}\le Ct^2\le Ca^4.}
 \label{v49:energy-error}
\end{equation}
The instantaneous \emph{original action} readers satisfy
\begin{equation}
 \boxed{\begin{aligned}
 \|\mathscr R_t(\xi)-\mathscr R_t(\widehat\xi)\|
                         &\le Ca^{-2}t^2\le Ca^2,\\
 \|\mathscr L_t^{\rm el}(\xi)-\mathscr L_t^{\rm el}(\widehat\xi)\|
                         &\le Ct^2\le Ca^4,\\
 \|\mathscr L_t^{\rm sp}(\xi)-\mathscr L_t^{\rm sp}(\widehat\xi)\|
                         &\le Ca^2t^2\le Ca^6.
 \end{aligned}}
 \label{v49:action-readers}
\end{equation}
\end{theorem}
\begin{proof}
Subtract \eqref{v49:forced-correction} from the exact projected variational
equation, not from a frozen eigenvector equation. This gives
\begin{equation}
 e_t=(T+K_P)e+b,\qquad e(0)=0,\qquad
 b=\dot P(I-P)r\in\Ran P.
 \label{v49:error-equation}
\end{equation}
By V47 and \cref{v49:predictor-theorem},
$|b(t)|_{a,\rm int}\le C a^{-1}\mathfrak n_t(r)\le Ct$.
Variation of constants and the original energy bound for $\mathcal U_a$
therefore give
\[
 |e(t)|_{a,\rm int}\le C\int_0^t e^{C(t-s)/a}s\,ds\le Ct^2.
\]
The exponential is uniformly bounded since $t/a\le ca$. All three
readers are linear differentials at the same original base point. Apply
\cref{v48:reader-transfer} to their difference $e$. The two-power metric
loss is fully retained; it is paid by the improved $a^4$ energy error.
\end{proof}

This is a correction of the existing \emph{coupled} comparison: its
complement $(I-P)\xi$ is still the actual one. It is not an independent
numerical solution of the complete variational dynamics. In particular no
unstable complementary propagator is bounded or replaced by a stable one.
The result changes no physical initial data, law, or time scale.

For orientation, the three choices of quartet forcing give the following
proved worst-case endpoint bounds on this actual selected family:
\begin{center}\small
\begin{tabular}{@{}lccc@{}}
\toprule
Input retained in $\dot P(I-P)$ & Energy error & Metric error & Electric-link error\\\midrule
None (V47--V48) & $O(a^2)$ & $O(1)$ & $O(a^2)$\\
Fixed initial canonical tangent $\xi_0$ & $O(a^3)$ & $O(a)$ & $O(a^3)$\\
Constraint-completed $v_{\rm p}$ & $O(a^4)$ & $O(a^2)$ & $O(a^4)$\\
\bottomrule
\end{tabular}
\end{center}
These are upper estimates for different analytical comparisons. The original
full tangent retains its $O(a)$ metric bound throughout; no lower estimate
for any of these actual errors is asserted.

\section{The second moving-projector derivative has a paid original-source bound}
\label{v49:second-motion}

To pass from the pointwise action reader to the reader of a comparison path,
one more base derivative is required. It is supplied here from the actual
scaled equations, not assumed as additional smoothness of an arbitrary
source.

\begin{lemma}[Second coefficient bounds on the existing physical interval]
\label{v49:second-coefficients}
Along the selected histories,
\begin{equation}
 \|\lambda_{tt}\|\le Ca^{-2},\qquad
 \|X_{tt}\|\le Ca,
 \label{v49:base-second}
\end{equation}
and the original matrix coefficients obey
\begin{equation}
 \boxed{\|\ddot F\|+\|\ddot Q\|\le Ca^{-2},\quad
        \|\ddot C\|\le Ca^{-1},\quad
        \|\ddot\Gamma\|\le Ca^{-3}.}
 \label{v49:coefficient-second}
\end{equation}
The same $C$ bound applies before restricting its rows.
\end{lemma}
\begin{proof}
In the actual V41 scaled equations, $z=(x,y)$ satisfies
$|z|\le C(1+s)$ and $|z'|\le C$ on $s\le c/a$ after undoing the fixed
shear. The coefficients $B,C_*$ are fixed and $D_zF_f=O(a+a^3|z|)$;
therefore $|y''|\le C$. The slow rows have no larger second derivative,
and the canonical part satisfies $|\xi''_{\rm scaled}|\le Ca^3$ by the
last assertion of \cref{v41:row-lemma}. Since physical displacements have
factor $a^4$ and $t=a^3s$, these are exactly
\eqref{v49:base-second}, including the fixed shear's contribution.
This uses time differentiation only, not a second parameter derivative of
an endpoint-selected map.

All ordinary stationary derivatives of $H$ remain bounded in the original
finite chart. Thus two time derivatives of $Q$ and $F=\mathbb JQ$ cost at
most $Ca^{-2}$. For $C=D_XP$, its multiplier derivative is
$D_\lambda C=D_XM=O(a)$ by \cref{v20:flat}. Hence the only potentially
large acceleration term is $D_\lambda C\,\lambda_{tt}=O(a^{-1})$;
all field-acceleration and quadratic-velocity terms are smaller.
The full source has the same property before restriction of the rows.

Finally the retained divided signed matrix $\Gamma$ is analytic with bounded
fixed derivatives in the normalized state coordinates
$(X/a,(\lambda-e)/a)$ on the original graded chart. Their first time
derivatives are bounded and their second derivatives are $O(a^{-3})$.
The chain rule proves the last bound in \eqref{v49:coefficient-second}.
This deliberately conservative bound requires no unproved cancellation in a
higher gravitational coefficient.
\end{proof}

\begin{theorem}[Twice differentiated full contour, including its inhomogeneous constraint]
\label{v49:ddotP}
For every canonical input $v$ held fixed during differentiation,
\begin{equation}
 \boxed{\mathfrak n_t(\ddot P(t)v)\le Ca^{-3}\mathfrak n_t(v).}
 \label{v49:ddotP-bound}
\end{equation}
\end{theorem}
\begin{proof}
Use the original saddle system on $z\in a^{-2}\mathscr C$ from V47:
\[
 (z-F)x-\mathbb JC^*\eta_x=v,\qquad Cx+M\eta_x=0.
\]
For $r=\mathfrak n_t(v)$ its original and first differentiated bounds are
\[
 \|x\|\le Ca^4r,\quad \|E_a^{-1}\eta_x\|\le Ca^2r,\quad
 \|\dot x\|\le Ca^3r,\quad
 \|E_a^{-1}\dot\eta_x\|\le Car.
\]
Differentiating twice yields exactly
\begin{align*}
 (z-F)\ddot x-\mathbb JC^*\ddot\eta_x&=r_2,\\
 C\ddot x+M\ddot\eta_x&=s_2,\\
 r_2&=\ddot F x+2\dot F\dot x+
       \mathbb J\ddot C^*\eta_x+2\mathbb J\dot C^*\dot\eta_x,\\
 s_2&=-\ddot Cx-2\dot C\dot x-\ddot M\eta_x-2\dot M\dot\eta_x.
\end{align*}
By \cref{v49:second-coefficients}, $\|r_2\|\le Car$ and
$\|E_as_2\|\le Ca^3r$. In particular
$E_a\ddot M\eta_x=a^4\ddot\Gamma E_a^{-1}\eta_x=O(a^3r)$;
that term is retained. The exact inhomogeneous solve of V47 gives
\[
 \|E_a^{-1}\ddot\eta_x\|+
 \|E_a^{-1}\eta_t(\ddot x)\|\le Ca^{-1}r,\qquad
 \|\ddot x\|\le Car.
\]
The second quantity includes the correction
$-a^{-4}\Gamma^{-1}E_as_2$; it is not obtained by identifying
$\ddot\eta_x$ with the lift of $\ddot x$. Thus
$\mathfrak n_t(\ddot x)\le Ca^{-1}r$. Integrating twice-differentiated
resolvents over the fixed contour of length $O(a^{-2})$ proves the theorem.
No individual eigenvector derivative or inverse internal-frequency gap occurs.
\end{proof}

\section{The corrected comparison also has small path-based curvature error}
\label{v49:path}

\begin{theorem}[Complete jet payment for the analytical comparison]
\label{v49:path-theorem}
For the same error $e=P\xi-u_{\rm p}$, uniformly on $[0,ca^2]$,
\begin{equation}
 \begin{gathered}
 \|e\|\le Ca^2t^2,\qquad
 \|\eta_A(e)\|\le Cat^2,\quad \|\eta_W(e)\|\le Ct^2,\\
 \|e_t\|\le Ca^2t,\qquad
 \|\partial_t\eta_A(e)\|\le Cat,\quad
 \|\partial_t\eta_W(e)\|\le Ct,\qquad
 \boxed{\|e_{tt}\|\le Ca^2.}
 \end{gathered}
 \label{v49:error-jets}
\end{equation}
Give $\widehat\xi=\xi-e$ its original algebraic multiplier lift. The
linearized metric and link readers obtained by differentiating this
comparison \emph{path} then satisfy
\begin{equation}
 \boxed{
 \|\mathscr R_{\rm path}(\xi)-\mathscr R_{\rm path}(\widehat\xi)\|
     \le Ca^2,\quad
 \|\Delta\mathscr L_{\rm path}^{\rm el}\|\le Ca^4,\quad
 \|\Delta\mathscr L_{\rm path}^{\rm sp}\|\le Ca^6.}
 \label{v49:path-readers}
\end{equation}
\end{theorem}
\begin{proof}
The first three bounds follow from \eqref{v49:energy-error} and the original
intermediate lifts. Put $b=\dot P(I-P)r$ as in
\eqref{v49:error-equation}. V47 gives $\mathfrak n_t(b)\le Ct$.
The original quartet generator is bounded by $Ca^{-2}$ in its energy norm,
so \eqref{v49:error-equation} yields
$\mathfrak n_t(e_t)\le C(t^2/a^2+t)\le Ct$.
This proves the canonical first derivative estimate. Differentiate its
multiplier lift as in \cref{v48:rate}; the two contributions are the lift
of $e_t$ and the fixed-input motion of the lift applied to $e$.
They give balanced norm $C(t+t^2/a)\le Ct$, proving the next two estimates.

No derivative of the entire comparison forcing is assumed. Its exact value is
\[
 b_t=\ddot P(I-P)r-\dot P\dot Pr+\dot P(I-P)r_t.
\]
By \cref{v49:predictor-theorem,v49:ddotP} its complete-record norm is at most
$C(t/a^2+t/a+1)\le C$. In particular $\|b_t\|\le Ca^2$.
Also $\dot K_P=[\ddot P,P]$. The physical canonical identity
\[
 \dot T e=\dot F e+\mathbb J\dot C^*\eta_t(e)
                         +\mathbb JC^*\dot\eta_t(e)
\]
and V48's fixed-input lift bound give $\|\dot T e\|\le Ct^2/a$.
For $T e_t$, use the original canonical row rather than a full unweighted
operator norm: $\|T e_t\|\le C(\|e_t\|+\|\eta_t(e_t)\|)\le Ct$.
The remaining terms satisfy
$\|\dot K_P e\|\le Ct^2/a$ and $\|K_Pe_t\|\le Cat$.
Differentiating \eqref{v49:error-equation} now gives
\[
 \|e_{tt}\|\le C(t+t^2/a+at+a^2)\le Ca^2,
\]
which proves \eqref{v49:error-jets} with every term accounted for.

The original linearized metric reader of a path uses its field and multiplier
values, its first field and multiplier time derivatives, and the configuration
part of its second field time derivative. Their coefficients, with all fixed
spatial derivatives included, are bounded on the actual base histories.
Substitution of \eqref{v49:error-jets} gives the metric estimate; in particular
no $\lambda_{tt}$ variation or unbounded derivative of a comparison defect
is silently omitted.

For the stationary spatial connection, the original bounds are
$z_X=O(1)$, $z_{\lambda_A}=O(a)$, $z_{\lambda_W}=O(a^2)$, and their
first time derivatives have respectively orders $O(a),O(a),O(a^2)$.
The last two orders follow from the same divided source on a normalized chart
with bounded first velocity. Thus its path variation obeys
$\|\delta z\|\le Ca^2t^2$ and
$\|\delta\dot z\|\le Ca^2t$. The temporal connection variation is at
most $C(\|e\|+\|\eta_t(e)\|)\le Ct^2\le Ca^2t$.
The original temporal-link expression has no derivative of the temporal
connection. Its electric difference is consequently at most $Ca^2t\le Ca^4$;
its spatial difference is at most $Ca^2t^2\le Ca^6$.
All products and frame variations are differentiated on the original base.
\end{proof}

Equivalently, the path defect is exactly
\[
 f_{\widehat\xi}=-K_Pe-b.
\]
The calculation above pays the V48 terms
$J_Xf_{\widehat\xi}+J_\lambda\eta_t(f_{\widehat\xi})+
J_{XX}(Tf_{\widehat\xi}+\dot f_{\widehat\xi})$ through their full
jet identity. These terms do not vanish individually and are not replaced
by the instantaneous action reader.

\section{The radial tangent condition can be restored without losing the estimate}
\label{v49:radial}

The predictor itself is tangent to all original constraints. Its projected
quartet correction and the exact complementary component need not separately
be radial tangent. This last pointwise issue is resolved by the original
canonical normal, not by calling spectral projection a physical constraint
reduction.

\begin{theorem}[A full-constraint-tangent, curvature-accurate comparison]
\label{v49:radial-theorem}
Let $g=\nabla\mathscr H^{\rm red}$ and take the original bounded normal
$W(t)$ with $CW=0$ and $\langle g,W\rangle=a$. Define
\begin{equation}
 \alpha_e(t)=\langle g(t),e(t)\rangle/a,\qquad
 e^\flat=e-\alpha_e W,\qquad
 \widehat\xi^\flat=\xi-e^\flat.
 \label{v49:radial-comparison}
\end{equation}
Then $(\widehat\xi^\flat,\eta_t(\widehat\xi^\flat))$ satisfies all 108
original tangent constraints at every time. Its multiplier is exactly the
one of $\widehat\xi$. Both the instantaneous action-reader estimates
\eqref{v49:action-readers} at their uniform endpoint orders and the path-reader
estimates \eqref{v49:path-readers} remain valid for
$\widehat\xi^\flat$.
\end{theorem}
\begin{proof}
The normal construction in \cref{v37:lossless-normal} extends throughout
these histories by the same fixed minor. It gives $\|W\|+\|W_t\|\le C$
and $\|W_{tt}\|\le Ca^{-2}$: the twice differentiated mean-scaled canonical
matrix costs at most $a^{-2}$ by \cref{v49:second-coefficients}.
Since $CW=0$, its normalized multiplier lift is zero. The correction
therefore leaves the multiplier unchanged and makes
$\langle g,e^\flat\rangle=0$ exactly. On the normalized graph this radial
identity is precisely the missing full tangent constraint, using original
homogeneity and mean lapse one.

V48's original radial estimate on the quartet gives
$|\alpha_e|\le Ca^4|e|_{a,\rm int}\le Ca^4t^2$. Its original action-reader
correction bounds consequently give orders smaller than or equal to
\eqref{v49:action-readers}. To check the path, use the actual identities
$g=-\mathbb JX_t$, $g_t=-\mathbb JX_{tt}$ and
$\|g_{tt}\|\le Ca^{-2}$. The latter follows by differentiating the original
analytic canonical field once more and applying
\eqref{v49:base-second}; no parameter acceleration is prescribed.
Direct differentiation gives
\[
 |\alpha_{e,t}|\le Ca^2t,\qquad
 |\alpha_{e,tt}|\le C(t^2/a+a^2t+a^2)\le Ca^2.
\]
Combining these with the bounds on $W,W_t,W_{tt}$ shows that the added
canonical correction has value $O(a^4t^2)$, first derivative $O(a^2t)$,
and second derivative $O(a^2)$. Its multiplier remains identically zero as
a path. The same original metric-jet calculation therefore preserves the
$Ca^2$ bound. Its stationary spatial-connection value is $O(a^4t^2)$ and
its time derivative is $O(a^2t)$, which preserve the literal-link bounds.
\end{proof}

At each nonzero amplitude this smooth tangent section can, if desired, be
integrated pointwise to nearby exactly constrained record paths by the
original canonical submersion. Such paths are not asserted to solve the
Euler equations. The corrected comparison itself has the explicit nonzero
variational defect
\[
 -K_Pe-b+\alpha_{e,t}W+\alpha_e(W_t-TW).
\]
The theorems concern differential comparisons along the original actual
histories. They do not construct a second exact-action history with these
projected components.

\section{Robustness of the correction and the remaining dynamical problem}
\label{v49:status}

The observation gap in V48 is closed for this refined comparison on the
actual selected parameter family. The correction retains the full changing
base matrices and a constraint-compatible predictor of the input, rather
than attempting to improve the sharp $a^{-2}$ reader norm. The final coupled
comparison is tangent to every original constraint, and its original metric,
temporal-link and spatial-link \emph{differential} errors are respectively
$O(a^2)$, $O(a^4)$ and $O(a^6)$ on the established $ca^2$ interval. The same
orders hold for its path readers after the kinematic defects are paid.

Different bounded choices of the fixed canonical normal minor produce
predictors with the same proved approximation orders. Applying the two
estimates to the same actual tangent and taking their difference shows that
their corrected readers agree to the same orders. No canonical choice of
normal projection or exact equality of comparisons is inferred.

There is also a simple explicit implementation tolerance at the action-reader
level. If an independently evaluated quartet forcing differs from
$\dot P(I-P)v_{\rm p}$ by a section $\epsilon(t)\in\Ran P(t)$, and the
initial quartet value has error $e_0$, then its additional intermediate
error is bounded by
\[
 C\bigl(|e_0|_{a,\rm int}+\|\epsilon\|_{L^1_t(\rm int)}\bigr).
\]
Thus retaining the $O(a^2)$ metric-error guarantee is ensured when that quantity is $O(a^4)$;
a sufficient uniform forcing tolerance is $O(a^2)$ on an $O(a^2)$ interval.
A claim about the path reader additionally pays the differentiated forcing,
as in \cref{v49:path-theorem}. This is a tolerance theorem for an analytical
comparison, not proof that a primitive source performs the required matrix
operations or that a numerical trajectory has passed the error test.

No change has been made to the action, physical clock, nonlinear initial
selection, or accepted evolution. The exact complementary tangent is still
retained. Consequently this is not an independent low-dimensional Cauchy
solver, a finite-perturbation shadowing theorem, or equality of the nonlinear
metric and link observations. It gives a more accurate description of a
spectral part of already supplied histories. The rapid and slow real sectors,
longer-time central transport, spatial refinement, and a common positive
physical lifespan remain unresolved. All action-specific constants remain
at $N=3$. The raw-action Einstein continuum theorem is not established.

\paragraph{Verification and preservation.}
The companion program checks the complete normal-predictor identities with
moving, noncommuting matrices; its exact tangent constraint and initial
anchoring; the differentiated inhomogeneous saddle equations and lift
corrections; the off-diagonal projection source and its derivative; radial
corrections; and every amplitude/time exponent in the proved bounds.
Independent finite matrix examples test algebra, not the original nonlinear
coefficient extraction or a new physical trajectory. The unchanged V48
standalone exact algebra suite is rerun and its report kept separate.
The four companion files, predecessor source, audit ranges and hashes are
included in the package. All 48 prior marked bodies are retained byte for
byte. No new gravitational source extraction, numerical chart radius,
nonlinear spacetime integration, proof-assistant formalization, or independent
validation of the full historical chain is claimed.
% END IMMUTABLE EXACT-ACTION BRIDGE V49 BODY


% Future iterations append here.

% BEGIN IMMUTABLE EXACT-ACTION BRIDGE V50 BODY
\clearpage
\renewcommand{\thesection}{\arabic{section}}
\renewcommand{\theHsection}{bridgeV50.\arabic{section}}
\setcounter{section}{404}
\section{Close the complementary comparison by a buffered full-record reconstruction}
\label{v50:context}

The comparison in \cref{v49:radial-theorem} has small original curvature
error, but still contains the unknown exact complementary tangent. Improving
its four-dimensional equation alone cannot remove that dependence. This
body instead reconstructs the \emph{whole} tangent record from the supplied
base history and initial tangent. It retains all 431 normalized state
coordinates and the actual coefficient coupling. The new approximation
requires no exact complementary tangent at positive time.

There is a precise price. The reference response has eleven expanding
coordinates. We prescribe their \emph{analytical} terminal value to be zero,
rather than invent the unknown terminal tangent or propagate approximate
unstable data forward. The resulting reconstruction agrees with the actual
selected tangent on every fixed earlier fraction of the supplied interval.
The terminal discrepancy is exponentially screened there. It need not be
small at the endpoint itself. The construction is therefore a buffered,
base-dependent boundary reconstruction, not a causal initial-value solver.

A singular but explicitly paid reweighting of the slow tangent coordinates
makes the entire original tangent equation a bounded reference system plus
an $O(a)$ matrix perturbation. A convergent Green--Neumann iteration then
gives any prescribed algebraic accuracy in $O(\log(1/a))$ exact integral
sweeps. A canonical normal completion makes each final comparison tangent
to \emph{all 108 original constraints}. Separate derivative estimates give
the original path readers, not only instantaneous action readers. For
example, the closed comparison attains errors $O(a^2)$, $O(a^4)$, and
$O(a^6)$ in metric, electric-link, and spatial-link curvature differentials,
respectively, without retaining the exact complementary component.

\paragraph{Inputs, scope, and nonduplication.}
The actual selected histories and tangent budgets are
\cref{v41:selection-theorem,v41:geometry-theorem} and
\eqref{v49:actual-budgets}. The regular original-action extension and its
separate row orders are \cref{v43:regular-lemma}. The weak reference
splitting, including its bounded seven-dimensional center, is
\cref{v32:spectral-theorem}. The complete canonical normal is
\cref{v49:normal-lemma}; the original readers and the distinction between
path and action jets are V48--V49. These are inherited inputs, not new
coefficient extractions. The new results are the full tangent balance, the
closed finite-sweep construction, its endpoint filtering, the whole-record
constraint completion, and its original-reader error bounds.

The use of dichotomic Green operators for well-conditioned boundary
problems is established methodology; see \cite{v50:Mattheij}. The proof
below supplies every specialized estimate, including the central interval
cost and the original amplitude factors. No external theorem asserting
stability of the present action is invoked. All results remain at $N=3$;
$a$ is field amplitude and is constant under time differentiation. A smaller
fixed $c>0$ than in V49 is allowed, but the physical scale remains $ca^2$.

\section{A complete tangent balance with a bounded central reference}
\label{v50:balance}

Use the fixed normalized chart and shear of \cref{v43:regular}:
\[
 u=\mathcal T_a(Q-Q_a^0)/a,\qquad
 \mathcal T_a(\xi,\eta_A,\eta_W)
       =(\xi,\eta_A+aH\eta_W,\eta_W),\qquad t=a^3s.
\]
Here $Q=(X,\lambda)$, the normalized multiplier complement has dimension
107, and $H$ is the fixed original Schur coefficient, not a new state
variable. The extended normalized field is
\[
 u'=\mathcal G(a,u)=Ju+a^3b+\mathcal R(a,u),\qquad
 J(x,y)=(0,C_*x+By),\quad (\dim x,\dim y)=(402,29).
\]
It agrees with the original field on the original constraint leaf. Let
$u_a(s)$ be one of the supplied selected base histories, with its parameter
in a fixed compact smaller interval, and let $(\xi,\eta)$ be its actual
parameter tangent. Define
\begin{equation}
 \mathsf D_a=\diag(aI_{402},I_{29}),\qquad
 Z_a=\mathsf D_a^{-1}\mathcal T_a(\xi,\eta)/a,\qquad
 \mathsf L_a=a\mathcal T_a^{-1}\mathsf D_a.
 \label{v50:state-weight}
\end{equation}
Thus $(\xi,\eta)=\mathsf L_a Z_a$. In components, for
$Z=(Z_X,Z_A,Z_W)$,
\begin{equation}
 \xi=a^2Z_X,\qquad
 \eta_A=a^2(Z_A-HZ_W),\qquad \eta_W=aZ_W.
 \label{v50:physical-weight}
\end{equation}
This is a norm on the complete canonical--multiplier record, not an
unweighted canonical norm. In particular $\|\mathsf D_a^{-1}\|=a^{-1}$;
no uniform equivalence to the earlier unweighted chart is asserted.

Put
\begin{equation}
 D=\diag(0_{402},B),\qquad
 V_a(s)=\mathsf D_a^{-1}D_u\mathcal G(a,u_a(s))\mathsf D_a-D.
 \label{v50:coefficient}
\end{equation}
The symbol $D$ in this body is a reference generator; it is not a multiplier
Hessian or the previously used weak Schur complement.

\begin{lemma}[The full original tangent equation is a small matrix perturbation]
\label{v50:balance-lemma}
Uniformly on $0\le s\le c/a$,
\begin{equation}
 \boxed{Z_a'=(D+V_a(s))Z_a,\quad
 \|Z_a(s)\|\le Ca,\quad
 \|V_a(s)\|\le Ca,\quad \|V_a'(s)\|\le Ca^2.}
 \label{v50:balanced-equation}
\end{equation}
In addition, its 324 canonical rows obey
\begin{equation}
 (Z_a)_X'=A_X(s)Z_a,\qquad
 \boxed{\|A_X(s)\|+\|A_X'(s)\|\le Ca^2.}
 \label{v50:canonical-row}
\end{equation}
The complete coupling $C_*$ is retained in $V_a$.
\end{lemma}
\begin{proof}
The selected initial records differ from $Q_a^0$ by $O(a^2)$, and their
physical displacements on the interval are $O(a^3)$. Hence
$\|u_a\|\le Ca$. Their original physical rates are $O(a)$, so
$\|u_a'\|\le Ca^3$. The regular normalized field has bounded derivatives;
in particular $\|u_a''\|\le Ca^3$. The actual tangent is tangent to the
constraint leaf, so differentiating the extension gives precisely its
original variational equation. Differentiating the fixed chart and applying
\eqref{v50:state-weight} proves the displayed equation.

The separate original rows give the exact block expression
\[
 V_a=\begin{pmatrix}
 \mathcal R_{s,x}&a^{-1}\mathcal R_{s,y}\\
 aC_*+a\mathcal R_{f,x}&\mathcal R_{f,y}
 \end{pmatrix}_{(a,u_a(s))}.
\]
At $u_a=O(a)$ the first derivatives of $\mathcal R_s$ are $O(a^2)$ and
those of $\mathcal R_f$ are $O(a)$. This proves $\|V_a\|\le Ca$.
Analyticity on the fixed normalized chart bounds $D_u^2\mathcal G$.
Differentiation along $u_a'=O(a^3)$, with the single $a^{-1}$ weight,
gives $\|V_a'\|\le Ca^2$. This is an estimate on the actual moving
coefficient, not a frozen approximation. The tangent budgets in
\eqref{v49:actual-budgets}, including $\eta_A=O(a^3)$, give
$\|Z_a\|\le Ca$ using \eqref{v50:physical-weight}.

There is a stronger canonical row. Write the original uneliminated tangent
row as $\xi_t=F\xi+G_A\eta_A+G_W\eta_W$, where
$F=\mathbb JQ$ and $(G_A,G_W)=\mathbb JC^*$ are bounded original matrices.
Direct substitution gives
\begin{equation}
 A_X=(a^3F,\ a^3G_A,\ a^2G_W-a^3G_AH).
 \label{v50:canonical-row-exact}
\end{equation}
Its $s$ derivative is bounded by $Ca^2$, either by the original first
coefficient bounds or by the bounded stationary derivatives and
$Q_{\rm base}'=O(a^4)$. This proves the last assertion. No canonical
constraint or fast equation was deleted to obtain this row.
\end{proof}

The reference $D$ has 11 expanding directions, 11 contracting directions,
and a bounded 409-dimensional center. Let
\[
 \Pi_+=\diag(0_{402},p_{\rm u}),\qquad
 \Pi_-=I-\Pi_+.
\]
These are \emph{reference} projections of $D$, not the full projections of
$J$ used in V42. This is legitimate because $aC_*$ is retained in $V_a$;
it is iterated with every other perturbation term. The inherited weak
splitting gives, for $r\ge0$,
\begin{equation}
 \|e^{rD}\Pi_-\|\le M_D,\qquad
 \|e^{-rD}\Pi_+\|\le M_D e^{-\kappa r},\qquad \kappa=1/500.
 \label{v50:semigroups}
\end{equation}
The center is not inverted. The new singular balance does not contradict
V45's nilpotent limiting chains in the old fixed norm; it moves their
coupling into the retained $O(a)$ matrix. Nor does it remove any physical
rapid or slow growing mode.

\section{A closed finite-sweep boundary reconstruction}
\label{v50:green}

Fix $1\le S\le c/a$. The inputs are $a$, the complete base history on
$[0,a^3S]$, and the initial tangent $Z_0=Z_a(0)$. No positive-time tangent
is an input. Set
\[
 b_-:=\Pi_-Z_0,\qquad h(s)=e^{sD}b_-.
\]
For a continuous vector function $f$, define
\begin{equation}
 (\mathcal G_S f)(s)=
 \int_0^s e^{(s-r)D}\Pi_-f(r)\,dr
 -\int_s^S e^{(s-r)D}\Pi_+f(r)\,dr.
 \label{v50:Green}
\end{equation}
It solves $w'=Dw+f$, with $\Pi_-w(0)=0$ and $\Pi_+w(S)=0$.
The minus sign in the backward integral is part of this identity.
By \eqref{v50:semigroups},
\begin{equation}
 \|\mathcal G_S f\|_\infty\le C_G(1+S)\|f\|_\infty.
 \label{v50:Green-norm}
\end{equation}
Define $\mathcal V_a w=\mathcal G_S(V_aw)$ and the finite iterates
\begin{equation}
 Z^{[0]}=h,\qquad
 Z^{[m+1]}=h+\mathcal V_a Z^{[m]}.
 \label{v50:iteration}
\end{equation}
Each iterate uses all 431 coordinates. It is not a four-dimensional spectral
solver with an exact complement hidden in its forcing.

\begin{theorem}[Unique comparison and geometric convergence of the integral sweeps]
\label{v50:iteration-theorem}
There are smaller $c,a_0>0$ such that, whenever $0<a<a_0$ and
$1\le S\le c/a$, the boundary problem
\begin{equation}
 Z_\star'=(D+V_a)Z_\star,\qquad
 \Pi_-Z_\star(0)=b_-,\qquad \Pi_+Z_\star(S)=0
 \label{v50:BVP}
\end{equation}
has a unique solution, $\|Z_\star\|_\infty\le Ca$. For $m\ge2$,
\begin{equation}
 \boxed{\max_{0\le j\le2}\sup_{s\le S}
 \|\partial_s^j(Z^{[m]}-Z_\star)\|\le Ca\,2^{-m}.}
 \label{v50:iteration-error}
\end{equation}
For $j=1,2$, its canonical rows have the stronger estimate
\begin{equation}
 \boxed{\sup_{s\le S}
 \|\partial_s^j(Z^{[m]}-Z_\star)_X\|\le Ca^3\,2^{-m}.}
 \label{v50:canonical-iteration}
\end{equation}
All 420 initial and 11 terminal reference conditions in \eqref{v50:BVP}
are satisfied exactly by every iterate.
\end{theorem}
\begin{proof}
Choose $c$ and then $a_0$ so that
$C_GC_V(a+c)\le1/2$, with the coefficient constant $C_V$ from
\cref{v50:balance-lemma}. Equations \eqref{v50:Green-norm} and
\eqref{v50:balanced-equation} imply $\|\mathcal V_a\|\le1/2$.
The fixed point is the absolutely convergent Neumann series
$Z_\star=\sum_{j\ge0}\mathcal V_a^j h$. Its norm and uniqueness follow,
as does the zeroth derivative bound. Differentiating the Green formula
verifies the ODE and its boundary conditions in both directions.

Put $d_m=Z^{[m]}-Z_\star$. For $m\ge1$,
\[
 d_m'=Dd_m+V_ad_{m-1},\qquad
 d_m''=Dd_m'+V_a'd_{m-1}+V_ad_{m-1}'.
\]
The matrices $D,V_a,V_a'$ are uniformly bounded. The zeroth derivative
bound applied successively at indices $m,m-1,m-2$ gives
\eqref{v50:iteration-error}; the fixed factors at these two earlier
indices are absorbed in $C$. The canonical row of $D$ is zero, so instead
$(d_m)_X'=A_Xd_{m-1}$ and
$(d_m)_X''=A_X'd_{m-1}+A_Xd_{m-1}'$. Equation
\eqref{v50:canonical-row} gives the extra factor $a^2$ in
\eqref{v50:canonical-iteration}. This factor is needed by the original
path curvature and is not a consequence of a generic state-norm estimate.
\end{proof}

This boundary reconstruction is not a forward solution with every original
initial tangent coordinate fixed. It uses the supplied nonexpanding initial
coordinates and a declared zero terminal condition for the expanding
reference coordinates. The next theorem proves why that comparison is
accurate on the selected class, without assuming that its actual terminal
coordinates were zero. The backward integral also means the method uses
base coefficients later than the observation time.

\section{The unknown terminal tangent is screened by a physical-time buffer}
\label{v50:buffer}

\begin{theorem}[Buffered comparison to the entire actual selected tangent]
\label{v50:buffer-theorem}
Fix $0<\omega<\kappa$. After reducing $a_0$, for every $m\ge2$ and
$0\le s\le S\le c/a$,
\begin{equation}
 \boxed{\max_{0\le j\le2}
 \|\partial_s^j(Z^{[m]}(s)-Z_a(s))\|
 \le Ca\{2^{-m}+e^{-\omega(S-s)}\}.}
 \label{v50:full-error}
\end{equation}
For $j=1,2$ the same right side acquires an additional factor $a^2$ in
the canonical rows. The constant is uniform on the compact selected
parameter range. No exact complementary tangent occurs in the construction.
\end{theorem}
\begin{proof}
Let $e=Z_\star-Z_a$. It satisfies the exact homogeneous difference
equation, zero nonexpanding initial data, and terminal data
$b_+=-\Pi_+Z_a(S)$. The actual selected bound gives $\|b_+\|\le Ca$.
Use the weighted norm
$\|e\|_{\omega,S}=\sup_s e^{\omega(S-s)}\|e(s)\|$.
The forward nonexpanding Green kernel has weighted integral at most
$M_D/\omega$; the backward expanding one has weighted integral at most
$M_D/(\kappa-\omega)$. Thus its norm is $C_\omega$, independent of $S$.
The free terminal contribution has weighted norm at most $M_D\|b_+\|$.
Since $\|V_a\|\le C_Va$, take $a_0$ so that
$C_\omega C_Va_0<1/2$. Absorption gives
\[
 \|e(s)\|\le Ca e^{-\omega(S-s)}.
\]
The difference equation and its first derivative give the same estimate
for $e'$ and $e''$. Its canonical equations are
$e_X'=A_Xe$, $e_X''=A_X'e+A_Xe'$, which give the stated factor $a^2$.
Combine these bounds with \cref{v50:iteration-theorem}. The terminal
vector is used only in the proof of the error bound; it is not input data
for \eqref{v50:iteration}.
\end{proof}

For $S=c/a$ and $0\le s\le\sigma S$, $0<\sigma<1$, the terminal error
is at most $Ca\exp[-\omega(1-\sigma)c/a]$. This is smaller than every
power of $a$. The full initial tangent is therefore recovered to arbitrary
algebraic accuracy too, although it is not imposed exactly by the boundary
conditions. At $s=S$ the bound need only be $Ca$. These distinctions avoid
claiming a stable forward initial-value inverse in a system with certified
expanding modes.

\section{Complete constraint completion without an unknown comparison component}
\label{v50:constraints}

The boundary problem is posed for the analytically extended full tangent
field. Its solution and finite iterates need not themselves satisfy all
original tangent constraints. Complete the full reconstructed record rather
than projecting a spectral component alone.

Use the original $\mathcal C,\mathcal M,R_c,N_c,B_c$ of
\cref{v49:normal-lemma}. For an arbitrary joint vector set
\begin{equation}
 \mathcal P_c(t)(v,\eta)
   =\bigl((I-N_c)v-B_c\eta,\eta\bigr),\qquad
 \mathcal C R_c=I_{108}.
 \label{v50:constraint-projector}
\end{equation}
Then $\mathcal P_c^2=\mathcal P_c$ and
$\mathcal C(\mathcal P_c(v,\eta))_X+\mathcal M\eta=0$. In particular
the radial row and the four constraint means are included. Define the
closed comparison by
\begin{equation}
 \boxed{(\widehat\xi_m,\widehat\eta_m)(a^3s)
       =\mathcal P_c(a^3s)\mathsf L_aZ^{[m]}(s).}
 \label{v50:reconstruction}
\end{equation}
Every coefficient in this formula comes from the base, and every state
argument comes from the finite sweeps. The mean-lapse variation is zero
because the original normalized multiplier complement is retained. Exact
matching of all initial shift means to the true tangent is not asserted;
their error satisfies the full-record estimate below.

\begin{lemma}[Two-jet stability of the original complete constraint completion]
\label{v50:constraint-lemma}
Let $\epsilon(s)=2^{-m}+e^{-\omega(S-s)}$. The errors
$e_X=\widehat\xi_m-\xi$, $e_A=\widehat\eta_{m,A}-\eta_A$,
$e_W=\widehat\eta_{m,W}-\eta_W$ satisfy
\begin{equation}
 \begin{gathered}
 \|e_X\|+\|e_A\|\le Ca^3\epsilon(s),\qquad
 \|e_W\|\le Ca^2\epsilon(s),\\
 \|(e_X)_t\|\le Ca^2\epsilon(s),\quad
 \|(e_A)_t\|\le C\epsilon(s),\quad
 \|(e_W)_t\|\le Ca^{-1}\epsilon(s),\\
 \boxed{\|(e_X)_{tt}\|\le Ca^{-1}\epsilon(s).}
 \end{gathered}
 \label{v50:physical-error-jets}
\end{equation}
The original multiplier component is unchanged by the constraint completion.
\end{lemma}
\begin{proof}
First do not apply $\mathcal P_c$. Equations
\eqref{v50:physical-weight} and \eqref{v50:full-error}, with
$\partial_t=a^{-3}\partial_s$, give the value and multiplier-derivative
bounds. The stronger canonical derivative estimates give
$\|(e_X)_t\|\le Ca^2\epsilon$ and
$\|(e_X)_{tt}\|\le Ca^{-1}\epsilon$. These are better than applying
two time derivatives to a generic balanced-coordinate error.

The actual tangent is fixed by $\mathcal P_c$. Thus its error after
completion is exactly
\[
 (e_X,e_\eta)\longmapsto
       ((I-N_c)e_X-B_ce_\eta,e_\eta).
\]
For clarity, the needed normal second derivatives have a uniform scaled
origin. The matrices $N_c$, $a^{-2}(B_c)_A$, and $a^{-3}(B_c)_W$ are
analytic functions of the original normalized state on the fixed minor
chart used in V49. This follows from its row-scaled full constraint and
source factorization. Along $u_a'=O(a^3)$ and $u_a''=O(a^3)$ their first
two $s$ derivatives are bounded by $Ca^3$. Consequently the conservative
physical-time bounds are
\[
 \|N_c\|+\|(N_c)_t\|\le C,\quad \|(N_c)_{tt}\|\le Ca^{-3},
\]
\[
 \|B_cb\|+\|(B_c)_tb\|\le C(a^2\|b_A\|+a^3\|b_W\|),\quad
 \|(B_c)_{tt}b\|\le C(a^{-1}\|b_A\|+\|b_W\|).
\]
These include all mean rows, rather than using an unweighted inverse bound.
The uncompleted multiplier second derivatives are bounded by
$Ca^{-3}\epsilon$ on $A$ and $Ca^{-4}\epsilon$ on $W$, directly from
\eqref{v50:full-error}. Differentiate the exact error-completion formula
twice. For example the largest mass-product second terms are
$a^2a^{-3}\epsilon$ and $a^3a^{-4}\epsilon$, both $Ca^{-1}\epsilon$.
The term $(N_c)_{tt}e_X$ is only $O(\epsilon)$.
All other product terms have no larger orders. The first derivatives and
values are checked in the same formula. This proves
\eqref{v50:physical-error-jets} and the exact multiplier assertion.
\end{proof}

The completion is applied only to the analytical comparison. It is not
inserted in the physical vector field. For example, with
$\Pi_a=\mathsf L_a^{-1}\mathcal P_c\mathsf L_a$ and
$\mathcal A_a=D+V_a$, the exact comparison defect is
\begin{equation}
 (\Pi_a Z^{[m]})'-\mathcal A_a\Pi_aZ^{[m]}
 =\Pi_a f_m+
 (\Pi_a'+\Pi_a\mathcal A_a-\mathcal A_a\Pi_a)Z^{[m]},
 \quad f_m=(Z^{[m]})'-\mathcal A_a Z^{[m]}.
 \label{v50:defect}
\end{equation}
There is no assertion that the second term vanishes. Both its jet effects
and the physical reader are estimated, not silently replaced by exact
Euler evolution.

\section{Closed full-record comparison in the original curvature readers}
\label{v50:readers}

\begin{theorem}[No exact complementary tangent is needed for algebraic reader accuracy]
\label{v50:reader-theorem}
For the reconstruction \eqref{v50:reconstruction}, both the instantaneous
original-action differential readers and the readers obtained from the
comparison \emph{path} satisfy
\begin{equation}
 \boxed{\begin{aligned}
 \|\Delta\mathscr R(a^3s)\|&\le Ca^{-1}\epsilon(s),\\
 \|\Delta\mathscr L^{\rm el}(a^3s)\|&\le Ca\epsilon(s),\\
 \|\Delta\mathscr L^{\rm sp}(a^3s)\|&\le Ca^3\epsilon(s),
 \end{aligned}\qquad
 \epsilon(s)=2^{-m}+e^{-\omega(S-s)}.}
 \label{v50:reader-estimates}
\end{equation}
Fix any $p>0$ and $0<\sigma<1$. Take $S=c/a$ and
\begin{equation}
 m(a)=2+\left\lceil(p+1)\log_2(1/a)\right\rceil.
 \label{v50:sweeps}
\end{equation}
Then the whole reconstructed record is tangent to every original constraint,
and, uniformly on $0\le t\le\sigma ca^2$,
\begin{equation}
 \boxed{\Delta\mathscr R=O(a^p),\qquad
 \Delta\mathscr L^{\rm el}=O(a^{p+2}),\qquad
 \Delta\mathscr L^{\rm sp}=O(a^{p+4}).}
 \label{v50:orders}
\end{equation}
These are differential comparisons to the supplied selected history. No
second nonlinear original-action history is asserted.
\end{theorem}
\begin{proof}
For the path reader, the original linearized metric construction uses field
and multiplier values, their first time derivatives, and the canonical
configuration second time derivative, with bounded coefficients at this
fixed spatial cutoff. Substitution of \eqref{v50:physical-error-jets}
gives the first line. Neither $\lambda_{tt}$ variation nor a path-defect
term is discarded.

For the spatial stationary connection, the retained derivatives are
$z_X=O(1)$, $z_{\lambda_A}=O(a)$ and $z_{\lambda_W}=O(a^2)$, with their
first physical-time derivatives bounded at the same respective orders
(the stronger $O(a)$ bound on $\partial_tz_X$ is available).
Thus its path difference is $O(a^3\epsilon)$ and its time derivative
is $O(a\epsilon)$. The temporal connection difference is
$O(a^2\epsilon)$. The original temporal-link expression uses this
connection and the spatial-connection velocity, not $\dot A_0$.
Its analytic ordered formula therefore gives $O(a\epsilon)$, and the
spatial link reader gives $O(a^3\epsilon)$.

For the instantaneous action reader, use the original full tangent
linearization on the completed tangent vector. In balanced coordinates its
matrix is $D+V_a$, which is bounded; its canonical rows are $O(a^2)$.
The action-generated first rates hence have the same bounds as the first
jets in \eqref{v50:physical-error-jets}. The canonical second action jet
uses $A_X'$ and $A_X(D+V_a)$ and has the same second-jet bound. Equivalently
one may insert those rate bounds in the original identities of
\cref{v48:jet-bound}. The completed vector is tangent to the full leaf,
so the derivative of the analytical extension equals that of the original
flow on this vector. This proves all three action-reader bounds separately;
it does not identify action and path readers on a nonzero-defect curve.

Finally $2^{-m(a)}\le a^{p+1}/4$. On the earlier fraction,
$e^{-\omega(S-s)}\le e^{-\omega(1-\sigma)c/a}=O(a^{p+1})$.
Inserting both estimates in \eqref{v50:reader-estimates} proves
\eqref{v50:orders}. The constants may depend on $p,\sigma$; no numerical
amplitude threshold is inferred. Increasing $p$ increases the number of
integral sweeps, not the order of an uncomputed Taylor expansion of the
action or the differentiability required of the path reader.
\end{proof}

For $p=2$ this yields exactly the $O(a^2)$, $O(a^4)$, $O(a^6)$ reader
orders of V49, now for a closed approximation of \emph{all} components.
The price is a physical observation buffer and knowledge of the supplied
base through its full horizon. The method is neither low-dimensional nor
causal in that horizon, and it need not reproduce the exact initial
expanding coordinates. The term $O(\log(1/a))$ counts exact integral sweeps,
not primitive renewal events, quadrature nodes, bit operations, or an
already verified numerical implementation.

\section{A residual certificate pays the required source derivatives}
\label{v50:residual}

The formula can also be checked without assuming exact evaluation of its
integrals. Let $W\in C^2([0,S];\R^{431})$ be an independently supplied
candidate in balanced coordinates. Evaluate the \emph{original} matrices
in \eqref{v50:coefficient} and set
\[
 r=W'-(D+V_a)W,\quad
 b_-=\Pi_-W(0)-\Pi_-Z_0,\quad b_+=\Pi_+W(S).
\]
Write $\beta_0=\|b_-\|+\|b_+\|$,
$\epsilon_0=\|r\|_\infty$, $\epsilon_1=\|r'\|_\infty$,
$\epsilon_{X,0}=\|r_X\|_\infty$ and
$\epsilon_{X,1}=\|r_X'\|_\infty$.

\begin{proposition}[A posteriori full-record and path-reader budget]
\label{v50:residual-theorem}
On the same horizon, put
$E=C\{\beta_0+(1+S)\epsilon_0\}$. Relative to $Z_\star$,
\[
 \|W-Z_\star\|_\infty\le E,\quad
 \|(W-Z_\star)'\|_\infty\le C(E+\epsilon_0),\quad
 \|(W-Z_\star)''\|_\infty\le C(E+\epsilon_0+\epsilon_1).
\]
After the same complete constraint completion, the additional reader errors
satisfy the conservative bounds
\begin{equation}
 \begin{aligned}
 \|\Delta\mathscr R\|
 &\le C\{a^{-2}(E+\epsilon_0+\epsilon_1)
                          +a^{-4}\epsilon_{X,1}\},\\
 \|\Delta\mathscr L^{\rm el}\|
 &\le C\{E+\epsilon_0+a^{-1}\epsilon_{X,0}\},\\
 \|\Delta\mathscr L^{\rm sp}\|&\le Ca^2E.
 \end{aligned}
 \label{v50:residual-reader}
\end{equation}
In particular, sufficient budgets for \eqref{v50:orders}, in addition to
the same observation buffer, are
\begin{equation}
 \boxed{\beta_0=O(a^{p+2}),\quad
 \epsilon_0=O(a^{p+3}),\quad
 \epsilon_1=O(a^{p+2}),\quad
 \epsilon_{X,1}=O(a^{p+4}).}
 \label{v50:residual-budget}
\end{equation}
The exact finite sweeps with \eqref{v50:sweeps} satisfy these budgets.
\end{proposition}
\begin{proof}
Apply the Green operator to the equation for $W-Z_\star$, retaining both
boundary defects. Its norm is bounded by a constant times
$\beta_0+(1+S)\epsilon_0$ plus at most half its own norm. Absorb that
last term. Differentiating the error equation once gives the two subsequent
bounds. In the canonical rows one has more precisely
\[
 (W-Z_\star)_X'=A_X(W-Z_\star)+r_X,
\]
\[
 (W-Z_\star)_X''=A_X'(W-Z_\star)
                  +A_X(W-Z_\star)'+r_X'.
\]
Undoing the amplitude weights and time scaling gives a canonical second
jet bounded by $C\{a^{-2}(E+\epsilon_0)+a^{-4}\epsilon_{X,1}\}$.
The multiplier first derivatives cost $a^{-1}$ and $a^{-2}$; in the normal
completion its multiplier second derivatives are multiplied by the original
$a^2$ and $a^3$ mass products. The same product calculation as in
\cref{v50:constraint-lemma} gives \eqref{v50:residual-reader}.
This proves both path and the (no larger) instantaneous action-reader bounds.

For the exact iterates, the boundary defects vanish and
\[
 r_m=V_a(Z^{[m-1]}-Z^{[m]}).
\]
The already-proved iteration and derivative estimates give
$\|r_m\|+\|r_m'\|\le Ca^2 2^{-m}$ and
$\|(r_m)_X'\|\le Ca^3 2^{-m}$, the latter from the canonical row
$A_X$. Substitution of \eqref{v50:sweeps} gives every budget.
For an independent implementation they must instead be checked against its
actual residual, including the differentiated canonical residual; a small
undifferentiated residual alone does not supply the path-reader claim.
\end{proof}

These residual budgets concern the evaluated coefficient problem. They do
not assert that the primitive source executes the integrals, that approximate
base coefficients meet the same budget, or that a numerical spacetime has
passed it. Approximate base data would require their own coefficient and
source-error estimates.

\section{What has been closed and what the construction does not imply}
\label{v50:status}

The full comparison no longer contains $(I-P)\xi$ or any other exact
positive-time tangent. Its specified inputs are the base history, its
original coefficient derivatives, the initial tangent, and a horizon.
The comparison is explicitly calculated by finitely many integral sweeps,
completed to all original tangent constraints, and verified in both kinds
of original curvature reader. This closes V49's dependence on an exact
complement in a buffered boundary-reconstruction sense. It does not make
the original nonlinear base history known independently.

A scalar calibration shows why the scope cannot be enlarged for free. For
$z'=\gamma z$, $\gamma>0$, the uniformly bounded family
$z_a(s)=a e^{-\gamma(S-s)}$ has zero nonexpanding initial data. The
zero-terminal reconstruction is identically zero. Its error is exponentially
small on every earlier fraction, but is exactly $a$ at $S$. This is an
analytical boundary example, not a gravitational counterexample. It also
shows why the present procedure is not a forward initial-value solver with
all expanding initial coordinates enforced.

All finite outputs in this body are tangent \emph{sections} along a supplied
base. Their complete constraint completion has a paid nonzero variational
defect. No second exact-action history, finite nonlinear shadowing theorem,
attracting microscopic selection, or equality of normalized metric and link
curvature has been constructed. The original $ca^2$ interval has not been
extended, and no $h\downarrow0$ statement has been added. The next substantive
step is finite-perturbation or nonlinear error control without a supplied
exact base, or the separate longer-time and spatial-refinement problem.
The raw-action Einstein continuum theorem remains unresolved.

\paragraph{Verification and preservation.}
The new script checks the singular-but-paid full-state balance, the
noncommuting reference/perturbation decomposition, the forward and backward
Green signs and boundary conditions, derivative and residual recurrences,
full-row constraint completion, physical jet powers, and residual tolerances
in exact symbolic or rational arithmetic. Separate finite examples verify
identities; they are not new gravitational source extractions. The retained
V49 standalone algebra suite is rerun. The source audit records the exact
inputs and the actual executed stages. All 49 previous marked bodies are
preserved byte for byte. No numerical original-action trajectory, new root
or spectral certificate, numerical chart radius, or proof-assistant
formalization is claimed.

\begin{thebibliography}{9}
\raggedright
\bibitem{v50:Mattheij}
R. M. M. Mattheij, \emph{Decoupling and stability of algorithms for boundary
value problems}, SIAM Review \textbf{27}(1) (1985), 1--44.
\href{https://doi.org/10.1137/1027001}{doi:10.1137/1027001}.
The original publisher record and abstract were checked. This is
methodological background for stable/unstable boundary decoupling, not an
imported stability theorem for the present action. The central Green bound,
constraint completion, and all amplitude-dependent reader estimates are
proved in this body.
\end{thebibliography}
% END IMMUTABLE EXACT-ACTION BRIDGE V50 BODY


% Future iterations append here.
% BEGIN IMMUTABLE EXACT-ACTION BRIDGE V51 BODY
\clearpage
\renewcommand{\thesection}{\arabic{section}}
\renewcommand{\theHsection}{bridgeV51.\arabic{section}}
\setcounter{section}{412}
\section{Reconstruct the nonlinear history, not another tangent along it}
\label{v51:context}

The full comparison in V50 has no exact complementary tangent as an input,
but its coefficients still require a supplied exact base history. This body
removes that particular dependence. We first solve all original constraints
in canonical normal coordinates, then perform a nonlinear boundary iteration
on the resulting 323-dimensional state chart. Its limit is a new exact
history of the unchanged normalized action. Its finite iterates are exactly
constrained nonlinear curves, not merely tangent sections. Their original
metric and ordered-link curvatures converge at quantitative amplitude orders.
No positive-time reference history or variation occurs in the iteration.

The inputs are the retained certified first-field root, the analytic
transport-selected \emph{initial record}, the original stationary/action
maps, and a horizon. The initial record is defined by the inherited finite
preparation theorem; no numerical value of its higher coefficients is
invented. Nonexpanding coordinates are fixed initially and expanding
coordinates are fixed at the endpoint, relative to an explicitly known
first-order path. Thus this is a nonlinear two-point construction, not a
forward Cauchy solver with every initial coordinate prescribed. It selects
higher initial coefficients, and is not asserted to reproduce the particular
endpoint- or invariant-graph selections in V32--V44.

\paragraph{Inputs and nonduplication.}
We retain the same V29 root and V30 full signed stationary chart,
\cref{v43:regular-lemma}, the complete canonical normal submersion
\eqref{v43:constraint-normal}, and the original source orders in
\cref{v49:normal-lemma}. The rapid matrix $B$, including its eleven stable,
eleven unstable, and seven bounded central directions, is unchanged. The
fixed Schur shear and separate canonical row are the ones used in V50.
V32--V44 already prove existence and selection of other short original
histories; another existence statement alone would not be an advance here.
The additions are a fully constrained \emph{nonlinear} integral reconstruction
without a supplied trajectory, quantitative errors of its finite sweeps in
the original nonlinear readers, and a nonlinear residual certificate.

The boundary-decoupling method and contraction principle are established
methods; the primary background \cite{v50:Mattheij} is retained rather than
claimed anew. Every action-specific scaling used below is derived from the
named original chart and source identities. The root, signed matrix, and
analytic initializer are inherited proof inputs, not independently
revalidated by this body. All conclusions have fixed $N=3$ and $h=1/3$.
The small parameter $a>0$ is field amplitude. The physical interval remains
of order $a^2$, not a cutoff-independent continuum interval.

\section{A nonlinear canonical chart with every constraint retained}
\label{v51:chart}

Let $Q_a^0=(X_a^0,\lambda_a^0)$ be the analytic initial record at
$\vartheta=d=0$ in \cref{v43:regular}. In particular
\[
 P(Q_a^0)=0,\qquad X_a^0=aY_*+O(a^2),\qquad
 \lambda_a^0=e+a\ell_*+O(a^2),\qquad
 \lambda_t(Q_a^0)=ar_*.
\]
Only the \emph{initial} record and its original vector-field value are used.
Write the fixed shear coefficient as $H=A_*^{-1}L_*$, and retain the physical
active/weak multiplier dimensions $78+29$. The mean-lapse-zero complement is
used for all multiplier differences. All 108 \emph{rows} of $P$ are retained,
ordered as 104 independent nonconstant rows and four means.

Set
\[
 \mathcal C_0 Z=(L_{\rm nc}Z,\ P_0DP_2(Y_*)[Z]).
\]
Choose a fixed 108-column canonical inclusion $J_c$ for which
$\mathcal C_0J_c$ is invertible. Choose $E:\R^{216}\to\R^{324}$ onto
$\ker\mathcal C_0$, and its coordinate projection $\Lambda$ with
$\Lambda E=I$, $\Lambda J_c=0$. Thus $[E\ J_c]$ is an invertible
canonical coordinate matrix. This is a coordinate parametrization of the
full constraint leaf, not a dynamical projection or a gauge quotient.

Write
\[
 w=(q,p_A,p_W)\in\R^{216}\oplus\R^{78}\oplus\R^{29}.
\]
We shall reconstruct physical records by
\begin{equation}
 \begin{split}
 X&=X_a^0+a^2\{Eq+J_c n(a,w)\},\\
 \lambda_A&=\lambda_{a,A}^0+a^2(p_A-Hp_W),\qquad
 \lambda_W=\lambda_{a,W}^0+ap_W.
 \end{split}
 \label{v51:Theta}
\end{equation}
Denote this map by $\Theta_a(w)$. Coordinate subscripts on $\lambda^0$
mean its representatives in the fixed affine normalized chart.

\begin{lemma}[Exact constraints and the different canonical normal costs]
\label{v51:normal}
On a fixed neighborhood of $(a,w)=(0,0)$ there is a real analytic function
$n(a,w)$ with $n(a,0)=0$ such that
\begin{equation}
                    P(\Theta_a(w))=0.             \label{v51:Pzero}
\end{equation}
For every positive sufficiently small $a$, this is a coordinate chart on the
original 323-dimensional normalized constraint leaf. Its mean lapse is one;
no additional constraint mean is omitted. Uniformly on a smaller fixed chart,
\begin{equation}
 n=O(a|q|+a|q|^2+a^2|p|),\quad
 D_q n=O(a),\quad D_p n=O(a^2),\qquad p=(p_A,p_W).
 \label{v51:normal-orders}
\end{equation}
All further fixed-order derivatives of $D_qn$ and $D_pn$ have the same
respective amplitude factors. In particular, for the canonical reconstruction
$\mathcal X_a(w)$,
\begin{equation}
 \begin{gathered}
 D_q\mathcal X_a=O(a^2),\qquad D_p\mathcal X_a=O(a^4),\\
 D_q^2\mathcal X_a=O(a^3),\qquad
 D_qD_p\mathcal X_a,\ D_p^2\mathcal X_a=O(a^4).
 \label{v51:X-jets}
 \end{gathered}
\end{equation}
These are finite coordinate estimates, not a uniform inverse of the multiplier
Hessian on unweighted records.
\end{lemma}
\begin{proof}
Substitute \eqref{v51:Theta} with $n$ an independent unknown into the original
constraint map. Subtract its zero value at $w=n=0$. Divide the nonconstant
rows by $a^2$ and the four mean rows by $a^3$; call the result
$\mathcal P(a,n,w)$. This divided map is analytic. Indeed the flat constraint
is $LX$ with zero spatial means and is independent of a nearby multiplier.
The canonical displacement is $a^2(Eq+J_cn)$. Its mean linearization at
$X_a^0=aY_*+O(a^2)$ starts at $a^3$. The quadratic canonical remainder in a
mean row is $O(a^4|Eq+J_cn|^2)$. None of these divisions removes a row.

The original source factorization supplies the sharper multiplier orders.
On this normalized stationary chart, active source columns are $O(a)$,
weak source columns are $O(a^2)$, and all four constant test columns are
$O(a^2)$. The constant-lapse assertion follows also from the retained radial
homogeneity identity. Hence the nonconstant rows of $P_\lambda b$ are
$O(a^2|b_A|+a^3|b_W|)$, and its four mean rows are
$O(a^3|b_A|+a^4|b_W|)$. These are analytic divisibilities on the normalized
chart, not only bounds along one trajectory. Inserting the multiplier
synthesis in \eqref{v51:Theta} makes every divided multiplier derivative
$O(a^2)$. The weak-to-active shear has its extra factor $a$ and obeys the
same bound. This is the mean-row accounting of \cref{v49:normal-lemma},
now before imposing a time-dependent solution.

At $a=0$ the divided equation is exactly
$\mathcal P(0,n,w)=\mathcal C_0J_cn$; its $q$ term is zero because
$\mathcal C_0E=0$, and its multiplier terms are zero. The analytic
implicit-function theorem gives $n$, with $n(0,w)=0$ and $n(a,0)=0$.
Differentiation in $p$ gives
$D_pn=-(D_n\mathcal P)^{-1}D_p\mathcal P=O(a^2)$.
The leading $q$ derivative vanishes at $a=0$, so $D_qn=O(a)$.
The same divisibilities persist after any fixed number of derivatives on a
smaller analytic chart. Taylor expansion in $q$ and integration in $p$
give \eqref{v51:normal-orders}. Differentiating \eqref{v51:Theta} proves
\eqref{v51:X-jets}. The inverse chart reads $q=\Lambda(X-X_a^0)/a^2$
and recovers $p$ from the displayed multiplier values. Thus its dimension
is $216+78+29=323$, and all original constraints are identically satisfied.
\end{proof}

The pointwise solve in this lemma is part of evaluating a coordinate chart.
It does not mean that an unconstrained physical trajectory is repeatedly
projected back onto $P=0$. Every state used below is evaluated in these
coordinates from the outset.

\section{The reduced nonlinear field has a uniformly small balanced derivative}
\label{v51:field}

Let $F(X,\lambda)=\mathbb J\nabla_X H(X,\lambda)$ and let $r(X,\lambda)$
be the original normalized multiplier rate on the verified signed chart.
Define the reduced field $\mathcal H_a$ by the exact free-coordinate rows
\begin{equation}
 \begin{split}
 (\mathcal H_a)_q(w)&=a\Lambda F(\Theta_a(w)),\\
 (\mathcal H_a)_A(w)&=a\{r_A(\Theta_a(w))+aH r_W(\Theta_a(w))\},\\
 (\mathcal H_a)_W(w)&=a^2r_W(\Theta_a(w)).
 \end{split}
 \label{v51:reduced-field}
\end{equation}
Here $t=a^3s$ and $a$ is held fixed. The vector field $r$ is never replaced
by its first Poisson coefficient or a positive surrogate inverse.

Put $b_s=(\Lambda V_*,r_{*,A})\in\R^{294}$ and $b_w=r_{*,W}\in\R^{29}$.
The reduced unbalanced linear operator is
$J_{\rm red}(x,y)=(0,C_{\rm red}x+By)$. It is obtained by restricting the
original $J$ to the flat constraint tangent; it has the same weak matrix $B$.
The inherited $Jb=0$ and constraint tangency give
\begin{equation}
                    C_{\rm red}b_s+Bb_w=0.        \label{v51:b-cancel}
\end{equation}
Set $D=\diag(0_{294},B)$, and let $\Pi_+$ be its eleven-dimensional
unstable projection and $\Pi_-=I-\Pi_+$. The latter has rank 312.

\begin{lemma}[Nonlinear balance on a shrinking tube]
\label{v51:balance}
The field \eqref{v51:reduced-field} extends analytically through $a=0$ on a
fixed $w$-neighborhood. For every fixed $M_0<\infty$, on $|w|\le M_0a$,
\begin{equation}
 \|D_w\mathcal H_a-D\|\le C_{M_0}a,\qquad
 \|D_w^j\mathcal H_a\|\le C_j\quad(2\le j\le4).
 \label{v51:balanced-derivatives}
\end{equation}
The free canonical rows have the stronger bounds
\begin{equation}
 \|D_w^j(\mathcal H_a)_q\|\le C_j a^2\quad(1\le j\le3).
 \label{v51:canonical-derivatives}
\end{equation}
For the explicit path
\begin{equation}
 w_{\rm ref}(s)=(a^2s b_s,\ a^3s b_w),\qquad
 f_a(s)=\mathcal H_a(w_{\rm ref}(s))-w_{\rm ref}'(s),
 \label{v51:reference}
\end{equation}
and $1\le S\le c/a$, one has
\begin{equation}
 \|f_a\|_\infty\le Ca^3,\qquad
 \|f_a'\|_\infty\le Ca^4,\qquad
 \|w_{\rm ref}\|_\infty\le Cca,\quad
 \|w_{\rm ref}'\|\le Ca^2.
 \label{v51:reference-residual}
\end{equation}
The constant $c$ is fixed sufficiently small. These estimates require no
positive-time exact history.
\end{lemma}
\begin{proof}
For clarity, first use unbalanced free coordinates $v=(x,y)$, with physical
canonical displacement $aEq$ before normal completion. The divided normal
solve at this scale gives an analytic canonical graph $h(a,v)$, with
$h(a,0)=0$, $D_qh=O(a+|v|)$, $D_{p_A}h=O(a^2)$ and
$D_yh=O(a^3)$. Its construction is the proof of \cref{v51:normal} with row
divisions $a,a^2$ instead of $a^2,a^3$. In particular
$n(a,w)=a^{-1}h(a,\diag(aI_{294},I_{29})w)$ has exactly the divisibilities
already established. The free projection of the original analytic normalized
field in \cref{v43:regular-lemma} has form
\[
 v'=J_{\rm red}v+a^3(b_s,b_w)+R^{\rm red}(a,v).
\]
Its slow rows retain the orders
$O(a^4+a^2|v|+a|v|^2)$ with derivative $O(a^2+a|v|)$.
Its fast rows have orders $O(a^4+a|v|+|v|^2)$ with derivative
$O(a+|v|)$. Composing the old fast linear row with the canonical normal
graph contributes $O(a|v|+|v|^2)$ and is included, not discarded.
The constant first velocity belongs to the limiting constraint tangent:
differentiate the divided constraints along the original vector field at
$Q_a^0$ and take the leading $a^3$ coefficient. This proves
\eqref{v51:b-cancel} after the free projection.

Now set $v=\diag(aI_{294},I_{29})w$. The derivative perturbation has blocks
\[
 \begin{pmatrix}
 R^{\rm red}_{s,x}&a^{-1}R^{\rm red}_{s,y}\\
 aC_{\rm red}+aR^{\rm red}_{f,x}&R^{\rm red}_{f,y}
 \end{pmatrix}.
\]
On the stated tube all entries are $O(a)$. The slow remainder has an
analytic factor $a$ on the entire normalized chart, so its apparent division
by $a$ is removable and its higher balanced derivatives are bounded.
The other rows require no division. This proves
\eqref{v51:balanced-derivatives}. Alternatively the free canonical row is
exactly $a\Lambda F(\Theta_a(w))$. The original $F$ and its derivatives are
bounded on the stationary chart, and the physical state derivatives in
\cref{v51:normal} are $O(a)$ or smaller. Differentiation proves
\eqref{v51:canonical-derivatives}; the sharper normal derivatives are retained
for the later physical jet estimates.

The path before balancing is $v_{\rm ref}=a^3s(b_s,b_w)$. Its linear terms
cancel by \eqref{v51:b-cancel}. Consequently $f_a$ is the balanced image of
$R^{\rm red}(a,v_{\rm ref})$. The slow entries are bounded by
$C(a^3+a^4s+a^6s^2)$ and the fast ones by
$C(a^4+a^4s+a^6s^2)$. These are at most $Ca^3$ for $s\le c/a$.
Differentiation along $v_{\rm ref}'=a^3(b_s,b_w)$ gives $Ca^4$ in both
blocks. This proves \eqref{v51:reference-residual} from the original row
orders, rather than applying smooth consistency to an uncontrolled path.
\end{proof}

\section{A constrained nonlinear boundary iteration constructs an exact history}
\label{v51:iteration}

The reference semigroup obeys the inherited bounds
\[
 \|e^{sD}\Pi_-\|\le M_D,\qquad
 \|e^{-sD}\Pi_+\|\le M_De^{-\kappa s}\quad(s\ge0),\qquad\kappa=1/500.
\]
The bounded reference center now has dimension $294+7=301$. The original
canonical coupling $aC_{\rm red}$ remains in the nonlinear perturbation.
Define the complete forward--backward Green operator
\begin{equation}
 (\mathcal G_S f)(s)=
 \int_0^s e^{(s-r)D}\Pi_-f(r)\,dr
 -\int_s^S e^{(s-r)D}\Pi_+f(r)\,dr .
 \label{v51:Green}
\end{equation}
Its norm is at most $C_G(1+S)$; only the center spends the interval length.
For an error curve $e$ put
\[
 N_a(s,e)=\mathcal H_a(w_{\rm ref}(s)+e)-\mathcal H_a(w_{\rm ref}(s))-De,
\]
and define finite nonlinear sweeps by
\begin{equation}
 e^{[0]}=0,\qquad
 e^{[m+1]}=\mathcal G_S\{f_a+N_a(\cdot,e^{[m]})\},\qquad
 w^{[m]}=w_{\rm ref}+e^{[m]}.
 \label{v51:sweeps}
\end{equation}
Each evaluation uses only the original maps and the current iterate.

\begin{theorem}[An exact nonlinear action history without a supplied base trajectory]
\label{v51:nonlinear-history}
There are fixed $c,a_0,K_0>0$ such that for $0<a<a_0$ and
$1\le S\le c/a$, the iteration \eqref{v51:sweeps} converges to the unique
solution in $\|w-w_{\rm ref}\|_\infty\le K_0a^2$ of
\begin{equation}
 w'=\mathcal H_a(w),\qquad
 \Pi_-(w(0)-w_{\rm ref}(0))=0,\qquad
 \Pi_+(w(S)-w_{\rm ref}(S))=0.
 \label{v51:nonlinear-BVP}
\end{equation}
Its error $e_*=w_*-w_{\rm ref}$ satisfies
\begin{equation}
 \|e_*\|_\infty\le Ca^2,\qquad
 \|e_*'\|_\infty\le Ca^3.                     \label{v51:exact-sizes}
\end{equation}
The physical curve $Q_*(t)=\Theta_a(w_*(t/a^3))$ solves the unchanged
normalized original-action equations on $[0,a^3S]$ and satisfies all 108
original constraints exactly. Every iterate $Q_m=\Theta_a(w^{[m]})$ also
satisfies all those constraints pointwise, but is not claimed to solve the
evolution equations. No exact base trajectory is an input to either construction.
\end{theorem}
\begin{proof}
On a ball $\|e\|_\infty\le K_0a^2$, the states $w_{\rm ref}+e$ lie in
$|w|\le M_0a$ for fixed $M_0$. By \cref{v51:balance},
\[
 |N_a(s,e_1)-N_a(s,e_2)|\le L_0a|e_1-e_2|.
\]
Choose $c$ and then $a_0$ so that
$q:=C_G(1+S)L_0a\le1/4$ for every admitted horizon. The image of zero has
norm at most $C_G(1+S)C_fa^3\le C_0a^2$. Taking $K_0=2C_0$ makes the
ball invariant and makes the iteration a contraction. Banach's theorem
gives the asserted solution and uniqueness in this ball. The Green equation
is equivalent to the differential and boundary equations, including the
backward minus sign. All finite sweeps remain in the stationary/signed chart.

The derivative estimate needs more than the undifferentiated contraction.
Put $F_e(s)=f_a(s)+N_a(s,e_*(s))$. Initially the equation gives a finite
$C^1$ bound. Integration by parts in \eqref{v51:Green} gives the exact identity
\begin{equation}
 e_*'(s)=e^{sD}\Pi_-F_e(0)+e^{(s-S)D}\Pi_+F_e(S)
                         +\mathcal G_S F_e'(s).
 \label{v51:Green-derivative}
\end{equation}
One has $\|F_e\|\le Ca^3$ and
\[
 \|F_e'\|\le Ca^4+L_0a\|e_*'\|.
\]
For the latter bound, differentiate $N_a$ at fixed $e$: its derivative is
$[D\mathcal H_a(w_{\rm ref}+e)-D\mathcal H_a(w_{\rm ref})]w_{\rm ref}'$,
which is at most $C a^2\|e\|\le Ca^4$. The other derivative term is
$D_eN_a\,e'$ with norm $L_0a\|e'\|$. The endpoint factors in
\eqref{v51:Green-derivative} are bounded in their indicated time directions.
Absorbing the same $q\|e_*'\|$ proves $\|e_*'\|\le Ca^3$.

The canonical and multiplier free coordinates of $\Theta_a(w_*)$ have their
original rates by \eqref{v51:reduced-field}. Both the reconstructed velocity
and the original field are tangent to $P=0$. Their difference has zero free
coordinates and hence is a canonical normal vector; the invertible normal
constraint derivative forces that difference to vanish. This recovers the
full canonical rows rather than only their free projection. Thus the limit
is an original-action history. All constraints of every iterate vanish by
\eqref{v51:Pzero}, independently of whether it solves the evolution equation.
No physical constraint-restoring force has been added.
\end{proof}

\begin{corollary}[A populated original curvature bound on the existing time scale]
\label{v51:base-curvature}
For $S=c/a$ the constructed original histories satisfy
\begin{equation}
 \sup_{0\le t\le ca^2}
 \{\|\operatorname{Riem}(g_*)\|+
   \|E_*^{\rm link}\|+\|F_*^{\rm sp,link}\|\}\le Ca.
 \label{v51:base-readers}
\end{equation}
Their total displacement from $Q_a^0$ is $O(a^3)$, and
$X_t,\lambda_t,X_{tt}=O(a)$. Their initial changes have the sharper bounds
\begin{equation}
 \|X_*(0)-X_a^0\|\le Ca^7,\quad
 \|\lambda_{*,A}(0)-\lambda_{a,A}^0\|\le Ca^5,\quad
 \|\lambda_{*,W}(0)-\lambda_{a,W}^0\|\le Ca^4.
 \label{v51:initial-cost}
\end{equation}
The first canonical field $Y_*$ and first multiplier tilt $\ell_*$ are
unchanged. The entire initial record is not prescribed to equal $Q_a^0$.
\end{corollary}
\begin{proof}
At $s=0$ the free slow, stable, and central error coordinates vanish.
The remaining initial component is
$\Pi_+e_*(0)=-\int_0^S e^{-rD}\Pi_+F_e(r)\,dr$.
The exponential kernel and $\|F_e\|\le Ca^3$ give $\|e_*(0)\|\le Ca^3$,
with only weak coordinates present. Equation \eqref{v51:Theta} and
$D_pn=O(a^2)$ give \eqref{v51:initial-cost}, including the small active
component induced by the fixed shear.

The free slow coordinates of $w_*$ are $O(a)$ and its weak coordinates are
$O(a^2)$, by the explicit reference and \eqref{v51:exact-sizes}. The normal
chart therefore gives physical displacement $O(a^3)$. Its slow scaled rates
are $O(a^2)$ and weak scaled rates $O(a^3)$. Undoing \eqref{v51:Theta} and
$t=a^3s$ gives $X_t,\lambda_t=O(a)$. The original canonical equation,
whose uneliminated derivatives are bounded, gives
$X_{tt}=D_XF\,X_t+D_\lambda F\,\lambda_t=O(a)$.
The original metric curvature needs just these jets, spatial derivatives at
this fixed cutoff, and at most the first multiplier time derivatives. The
stationary connection and its velocity have order $a$; its original ordered
link curvatures have the same order. This is the original reader argument
of V21/V33, applied to a newly constructed history, not to a comparison path.
\end{proof}

\section{Finite sweeps converge with the time jets the metric actually uses}
\label{v51:finite-jets}

\begin{theorem}[Nonlinear finite-record and path-jet accuracy]
\label{v51:finite-accuracy}
For $m\ge2$ let $\delta_m=w^{[m]}-w_*$ and put $\rho_m=2^{-m}$. Then
\begin{equation}
 \max_{0\le j\le2}\|\partial_s^j\delta_m\|_\infty\le Ca^2\rho_m,
 \qquad
 \|\partial_s^j(\delta_m)_q\|_\infty\le Ca^4\rho_m\quad(j=1,2).
 \label{v51:sweep-jets}
\end{equation}
After the exact nonlinear constraint reconstruction, differences from the
original exact history have the bounds
\begin{equation}
 \begin{array}{c|c}
 \text{physical quantity}&\text{uniform error}\\ \hline
 X&a^4\rho_m\,C\\
 \lambda_A&a^4\rho_m\,C\\
 \lambda_W&a^3\rho_m\,C\\
 X_t&a^3\rho_m\,C\\
 (\lambda_A)_t&a\rho_m\,C\\
 (\lambda_W)_t&\rho_m\,C\\
 X_{tt}&\rho_m\,C
 \end{array}
 \label{v51:physical-jets}
\end{equation}
These are actual derivatives of the finite comparison curves. In particular
no equation is assigned to a finite iterate as though it were exact.
\end{theorem}
\begin{proof}
The contraction with ratio at most $1/4$ gives
$\|\delta_m\|\le Ca^2 4^{-m}$. Differentiating the defining Green equation
shows
\[
 \delta_m'=D\delta_m+
 N_a(s,e^{[m-1]})-N_a(s,e_*).
\]
This gives $Ca^2 4^{-m+1}$ for the first derivative. A second derivative
uses the bounded second derivatives of $\mathcal H_a$,
$w_{\rm ref}'=O(a^2)$ and $e_*'=O(a^3)$; the same estimate, with a fixed
index loss, follows for $m\ge2$. Fixed factors from these two index losses
are absorbed into $C2^{-m}$. The iteration derivatives themselves remain
bounded, by the same formulas and induction.

In the canonical free rows $D_q=0$, so the full iterate satisfies exactly
\[
 (w^{[m]})_q'=(\mathcal H_a)_q(w^{[m-1]}).
\]
Thus the canonical derivative error gains $a^2$ by
\eqref{v51:canonical-derivatives}. Differentiating once more and using its
second derivative bound gives the same gain for the second derivative.
This proves \eqref{v51:sweep-jets} without presuming a small derivative of
an arbitrary approximate path.

The multiplier rows of \eqref{v51:Theta} are affine, and a physical time
derivative costs $a^{-3}$. Their entries in \eqref{v51:physical-jets} follow
immediately. For the canonical reconstruction apply the chain rule using
\eqref{v51:X-jets}. In the first scaled derivative, the free-canonical
part costs $a^2\cdot a^4\rho_m$, and the multiplier-normal part costs
$a^4\cdot a^2\rho_m$. Both are $O(a^6\rho_m)$; changes in the coefficients
are no larger. In the second scaled derivative the leading two products
have the same order. Every product involving a second chart derivative is
smaller, using $D_q^2\mathcal X=O(a^3)$ and all multiplier-containing
second derivatives $O(a^4)$, together with the established bounded first
scaled jets. Division by $a^3$ and $a^6$ proves the two canonical time
entries. The zeroth-order canonical bound follows directly from the first
chart derivatives. This accounts for the nonlinear normal correction and
its derivatives, rather than applying the tangent completion of V50 to a
finite displacement.
\end{proof}

\section{The original nonlinear metric and link readers converge on the full horizon}
\label{v51:readers}

\begin{theorem}[Arbitrary algebraic reader accuracy without a supplied base]
\label{v51:reader-theorem}
Use the original nonlinear ADM reconstruction and original ordered-link
formulas on $Q_m$, and compare them to the exact history $Q_*$. For $m\ge2$,
\begin{equation}
 \boxed{\begin{aligned}
 \|\operatorname{Riem}(g_m)-\operatorname{Riem}(g_*)\|_\infty
                  &\le C2^{-m},\\
 \|E_m^{\rm link}-E_*^{\rm link}\|_\infty
                  &\le Ca^2 2^{-m},\\
 \|F_m^{\rm sp,link}-F_*^{\rm sp,link}\|_\infty
                  &\le Ca^4 2^{-m}.
 \end{aligned}}
 \label{v51:reader-errors}
\end{equation}
The same orders hold separately if each finite record is read using the
instantaneous original-action velocity instead of the comparison path's
velocity. For every fixed $p>0$, choosing
\begin{equation}
           m(a)=3+\lceil p\log_2(1/a)\rceil        \label{v51:m-choice}
\end{equation}
gives, on the \emph{entire} interval $0\le t\le ca^2$,
\begin{equation}
 \boxed{\Delta R=O(a^p),\qquad
 \Delta E^{\rm link}=O(a^{p+2}),\qquad
 \Delta F^{\rm sp,link}=O(a^{p+4}).}
 \label{v51:reader-powers}
\end{equation}
There is no terminal observation buffer relative to this problem's own exact
boundary-selected history. No equality with an older selected history is
asserted by removing that buffer.
\end{theorem}
\begin{proof}
The metric reader is a smooth finite jet function on nondegenerate spatial
metrics and positive lapses. Its required jets are the state, first time
jets and canonical configuration second time jet; second time derivatives
of lapse or shift do not occur. The exact history has all relevant jets
$O(a)$, and \eqref{v51:physical-jets} bounds the approximate jets on a common
finite jet neighborhood (increase the fixed minimum sweep index if needed).
The mean value theorem on that neighborhood gives the first bound. This is
a finite nonlinear secant estimate; no Taylor series in the comparison
error defines a different reader.

For the stationary spatial connection $z$, the original derivative orders
are $z_X=O(1)$, $z_{\lambda_A}=O(a)$ and $z_{\lambda_W}=O(a^2)$ on these
states. Their analytic source divisibilities hold on the entire intervening
chart, not just at the exact solution. Thus the spatial connection error is
$O(a^4\rho_m)$. Its time derivative has error at most $O(a^2\rho_m)$:
the three direct contributions have sizes $a^3\rho_m$,
$a\cdot a\rho_m$, and $a^2\cdot\rho_m$. Coefficient differences multiplied
by the exact $O(a)$ rates are smaller. The temporal connection error is
$O(a^3\rho_m)$ by its bounded stationary derivatives. The original temporal
link expression uses the spatial-connection velocity and shifted temporal
connection, not a derivative of the temporal connection. Its smooth ordered
formula, with fixed $h=1/3$, gives $O(a^2\rho_m)$. The spatial ordered
formula gives $O(a^4\rho_m)$.

For the instantaneous action reader, $Q_m$ lies on the exact constraint
leaf, so its induced rate is \eqref{v51:reduced-field} at $w^{[m]}$. Its
balanced rate difference is $O(a^2\rho_m)$ and its free canonical rate
difference gains $a^2$. The original canonical equation then bounds the
canonical acceleration difference by $O(\rho_m)$, retaining the multiplier
rate contribution. The same original metric and link estimates follow.
This proves the two interpretations separately, not by declaring the finite
iterate's path defect zero. Finally \eqref{v51:m-choice} implies
$2^{-m}\le a^p/8$, proving \eqref{v51:reader-powers}.
\end{proof}

The limit has curvature $O(a)$ by \cref{v51:base-curvature}. For $p\ge1$,
the finite reconstructions therefore retain that size. This is accuracy for
each reader relative to its own exact-action value. It does not assert
that normalized metric and link curvatures agree with one another, nor that
the new exact history satisfies a continuum Einstein equation.

\section{A nonlinear residual test and the input that is no longer required}
\label{v51:residual}

Let $W\in C^2([0,S];\R^{323})$ lie in the same reconstruction ball.
Evaluate the \emph{original reduced nonlinear field}, not a field linearized
along a supplied base, and define
\[
 r=W'-\mathcal H_a(W),\quad
 \beta=\|\Pi_-(W-w_{\rm ref})(0)\|+
        \|\Pi_+(W-w_{\rm ref})(S)\|.
\]
Put $\delta_0=\|r\|_\infty$, $\delta_1=\|r'\|_\infty$,
$\delta_{q,0}=\|r_q\|_\infty$, and $\delta_{q,1}=\|r_q'\|_\infty$.

\begin{proposition}[Finite nonlinear a posteriori control]
\label{v51:residual-theorem}
Assume $\delta_0\le a^2$. Relative to $w_*$, put
$E=C\{\beta+(1+S)\delta_0\}$. Then
\[
 \|W-w_*\|\le E,\quad
 \|(W-w_*)'\|\le C(E+\delta_0),\quad
 \|(W-w_*)''\|\le C(E+\delta_0+\delta_1).
\]
After the exact normal reconstruction, the nonlinear path-reader errors are
bounded by
\begin{equation}
 \begin{aligned}
 \|\Delta R\|&\le C\{a^{-2}(E+\delta_0+\delta_1)+a^{-4}\delta_{q,1}\},\\
 \|\Delta E^{\rm link}\|&\le C\{E+\delta_0+a^{-1}\delta_{q,0}\},\\
 \|\Delta F^{\rm sp,link}\|&\le Ca^2 E.
 \end{aligned}
 \label{v51:residual-readers}
\end{equation}
In particular, sufficient error budgets for \eqref{v51:reader-powers} are
\begin{equation}
 \beta=O(a^{p+2}),\quad \delta_0=O(a^{p+3}),\quad
 \delta_1=O(a^{p+2}),\quad \delta_{q,1}=O(a^{p+4}).
 \label{v51:residual-budget}
\end{equation}
The candidate's nonlinear constraints vanish only if its normal solve has
also been performed exactly; an approximate normal solve needs its own
constraint-residual certificate and is not covered by that equality.
\end{proposition}
\begin{proof}
Subtract the two nonlinear equations and integrate the mean Jacobian along
the segment between $W$ and $w_*$. Its difference from $D$ is at most $Ca$,
and the full segment remains in the convex balanced tube. The Green estimate
and the same contraction absorption give the bound $E$. Differentiating the
error equation once gives the first two derivative estimates; the bounded
second field derivative and $W'=\mathcal H_a(W)+r=O(a^2)$ control the
coefficient-motion term. In the free canonical rows the sharper identities
are
\[
 \|(W-w_*)_q'\|\le Ca^2E+\delta_{q,0},\qquad
 \|(W-w_*)_q''\|\le Ca^2(E+\delta_0)+\delta_{q,1}.
\]
Insert these and the full derivative estimates into the nonlinear chart
chain rule as in \cref{v51:finite-accuracy}. The canonical second physical
jet costs at most
$C\{a^{-2}(E+\delta_0+\delta_1)+a^{-4}\delta_{q,1}\}$;
its first jet costs at most
$C\{a(E+\delta_0)+a^{-1}\delta_{q,0}\}$.
The multiplier first jets have costs $a^{-1}$ and $a^{-2}$. The original
connection source orders and metric formulas therefore give
\eqref{v51:residual-readers}. Since $1+S=O(a^{-1})$ and
$\delta_{q,0}\le\delta_0$, the displayed budgets imply the claimed powers.
For exact sweeps, the residual is
$N_a(s,e^{[m-1]})-N_a(s,e^{[m]})$; it and its derivative are
$O(a^3 2^{-m})$, and its differentiated canonical row is
$O(a^4 2^{-m})$. These satisfy the same budgets with
\eqref{v51:m-choice}. No numerical quadrature or coefficient error is
assumed to satisfy them without evaluation.
\end{proof}

\section{What is new, what is selected, and what remains open}
\label{v51:status}

The construction now specifies a full nonlinear history by original initial
algebraic data and the original finite action maps. The history is not an
input. Each finite iterate already lies in the complete original constraint
leaf; its nonzero dynamical residual is retained and paid in the path readers.
Its limit is an actual original-action history, and the finite iterates have
arbitrary algebraic nonlinear reader accuracy on the whole horizon of that
new boundary problem. This is the advance beyond V50's tangent sections
along a supplied trajectory.

The boundary information has not disappeared. The nonexpanding initial free
coordinates and expanding terminal displacement are prescribed relative to
$w_{\rm ref}$. The initial expanding coordinates are selected by the solve.
Mean lapse is fixed; no theorem here fixes all three initial shift means to
the old values. A given earlier selected history need not meet these new
boundary conditions and is not declared identical to the new limit. The
algorithm uses backward as well as forward integrals and is not a causal
microscopic update, an attracting law, or a claim of arbitrary-initial-data
stability. All stationary and normal-map evaluations are analytical inputs;
no quadrature, bit complexity, or numerical nonlinear trajectory has been
certified. The number $O(\log(1/a))$ counts exact integral sweeps.

The longest physical scale remains $ca^2$ at $N=3$. No positive lower bound
on lifespan as $a\to0$, no spatial-refinement family, and no normalized
metric/link identification is supplied. A productive next step is a verified
finite implementation of the nonlinear residual and normal-map budgets, or
the independent central/slow transport and refinement problem. Another
comparison that assumes the same exact base is no longer needed for this
particular reconstruction task. The raw-action Einstein continuum bridge
remains unresolved.

\paragraph{Executed verification and preservation.}
The accompanying exact script tests the constraint-coordinate chain rule,
noncommuting balanced blocks, the affine reference cancellation, forward and
backward Green identities, nonlinear residual recurrence, differentiated
canonical rows, and all displayed amplitude powers on separately labelled
finite algebra examples. It checks preservation of all fifty predecessor
bodies. The supplied V50 standalone algebra suite is rerun; its older
root and high-source extraction chain is not rerun or claimed independently
validated. The present existence and uniform estimates are the written
proofs based on the explicitly named original-action inputs, not conclusions
from those algebra examples. No new source coefficient, spectral root,
numerical chart radius, nonlinear trajectory integration, or proof-assistant
formalization is claimed.
% END IMMUTABLE EXACT-ACTION BRIDGE V51 BODY


% Future iterations append here.
% BEGIN IMMUTABLE EXACT-ACTION BRIDGE V52 BODY
\clearpage
\renewcommand{\thesection}{\arabic{section}}
\renewcommand{\theHsection}{bridgeV52.\arabic{section}}
\setcounter{section}{420}
\section{Finite time reconstruction and an honest nonlinear constraint budget}
\label{v52:context}

The nonlinear boundary problem in \cref{v51:nonlinear-history} no longer
requires an exact positive-time history. Its finite sweeps nevertheless
contain exact integrals of a nonlinear forcing, and its exact constraint
statement uses an exact implicit normal map. These are the two inputs
addressed here. We replace each forcing by a cubic Hermite polynomial,
evaluate its forward and backward exponential integrals by finite matrix
functions, and quantify both field-jet evaluation errors and an inexact
normal solve. A piecewise quintic output can be formed from finitely many
certified nodal jets without losing the second time derivatives used by the
original metric reader.

The new result is an amplitude-dependent finite information bound, not a new
Einstein limit. For any fixed requested reader order $p>0$, a uniform mesh
with $O(a^{-1-p/3})$ intervals and $O(\log(1/a))$ nonlinear sweeps suffices
for metric, temporal-link and spatial-link errors of orders
$a^p,a^{p+2},a^{p+4}$. These counts refer to evaluations of the original
reduced field and its first derivative, and to fixed-dimensional matrix
operations. They are not counts of elementary bit operations or native
renewal events. Certified evaluation of the original stationary maps and
initial algebraic record remains necessary; the theorem does not supply
numerical chart radii or silently replace those maps by a polynomial action.

\paragraph{Inherited inputs and nonduplication.}
We retain the same root, source normalization and $N=3$ stationary chart.
The exact nonlinear map $\Theta_a$, its normal derivatives, the reduced
field $\mathcal H_a$, and its boundary-selected original history $w_*$ are
\cref{v51:normal,v51:balance,v51:nonlinear-history}. The original nonlinear
reader estimate is \cref{v51:residual-theorem}. Those statements are inputs;
this body does not independently revalidate the historical root or signed
source certificates. Neither cubic interpolation nor exponential integration
is claimed as a new general method; \cite{v52:HO} is background for the
matrix-function treatment, while the boundary formulation is the retained
\cite{v50:Mattheij}. All interpolation bounds used below are proved here.
The new content is the complete-source discretization, its differentiated
residual bounds, its finite normal-map budget, and their original-reader
consequences. The earlier stable-writer Hermite constructions concern a
different evolution and are not used to identify that writer with this action.

Throughout, $a>0$ is field amplitude, $h=1/3$ is fixed, $t=a^3s$, and
$S=c/a$ with $c>0$ sufficiently small. Constants are independent of $a$,
the number of time intervals, and the iteration count. They may depend on
the inherited fixed-cutoff chart and on a fixed requested order $p$.

\section{Interpolate the complete forcing, including its derivative}
\label{v52:Hermite}

Put $L_S=1+S$ and use the $323$-dimensional coordinates of V51. Thus
$D=\diag(0_{294},B)$, the projections $\Pi_-,\Pi_+$ have ranks $312,11$,
and the original $B$ retains its stable, unstable and central sectors.
For a $C^1$ error curve $e$ define
\begin{equation}
 g[e](s)=\mathcal H_a(w_{\rm ref}(s)+e(s))-De(s)-w_{\rm ref}'(s).
 \label{v52:g}
\end{equation}
This is exactly $f_a+N_a(\cdot,e)$ from V51. Its derivative is
\begin{equation}
 g[e]'=D\mathcal H_a(w_{\rm ref}+e)(w_{\rm ref}'+e')-De'.
 \label{v52:gprime}
\end{equation}
There is no finite-difference time derivative and no omitted multiplier or
normal-map derivative in this formula.

Take a uniform partition $s_j=j\Delta$, $0\le j\le J$, with
$J\Delta=S$ and $0<\Delta\le1$. For a $C^1$ vector function $f$, let
$\mathcal I_\Delta f$ be the cubic on each cell matching $f,f'$ at both
endpoints. With $u=(s-s_j)/\Delta$, its basis is
\[
 h_{00}=2u^3-3u^2+1,\quad h_{10}=u^3-2u^2+u,\quad
 h_{01}=-2u^3+3u^2,\quad h_{11}=u^3-u^2.
\]
The slope data multiply $\Delta h_{10},\Delta h_{11}$. One shared value
and slope are used at every interior node, so the forcing is globally $C^1$.

\begin{lemma}[Mesh-independent stability and differentiated interpolation error]
\label{v52:Hermite-lemma}
There is an absolute constant $C_H$ such that
\begin{align}
 \|\mathcal I_\Delta f\|_\infty
   &\le C_H(\|f\|_\infty+\Delta\|f'\|_\infty),\\
 \|(\mathcal I_\Delta f)'\|_\infty&\le C_H\|f'\|_\infty.
 \label{v52:C1-stability}
\end{align}
On a cell where $f$ has four continuous derivatives,
\begin{equation}
 \|\mathcal I_\Delta f-f\|_\infty
       \le C_H\Delta^4\|f^{(4)}\|_\infty,\qquad
 \|(\mathcal I_\Delta f-f)'\|_\infty
       \le C_H\Delta^3\|f^{(4)}\|_\infty.
 \label{v52:interp-error}
\end{equation}
For $j=2,3$, the cellwise derivative also satisfies
$\|\partial_s^j\mathcal I_\Delta f\|_\infty
\le C_H\|f^{(j)}\|_\infty$ whenever that derivative exists.
If nodal value and slope errors are bounded by $\epsilon_0,\epsilon_1$,
the resulting polynomial error $b$ satisfies
\begin{equation}
 \|b^{(j)}\|_\infty
 \le C_H(\Delta^{-j}\epsilon_0+
                    \Delta^{1-j}\epsilon_1),\quad 0\le j\le3.
 \label{v52:data-error}
\end{equation}
\end{lemma}
\begin{proof}
Rescale to $[0,1]$ and use the four displayed polynomials. For first
 derivative stability subtract the constant $f(s_j)$; the other endpoint
value differs by at most $\Delta\|f'\|_\infty$. For stability of
$\partial_s^j$, subtract the Taylor polynomial of degree $j-1$ at the left
endpoint. Hermite interpolation reproduces it; its $j$th derivative is
zero, and the remaining endpoint values and slopes are bounded by
$C\Delta^j\|f^{(j)}\|$ and $C\Delta^{j-1}\|f^{(j)}\|$.
For \eqref{v52:interp-error}, subtract the degree-three Taylor polynomial
and use its integral remainder, also after one derivative. Differentiating
the fixed basis gives the powers of $\Delta$. Applying the same basis
directly to the four nodal errors proves \eqref{v52:data-error}. The proofs
are componentwise, or by Bochner integral remainders in the finite norm.
\end{proof}

The $\epsilon_0/\Delta$ term in the derivative error cannot be dropped.
For example, an error $\epsilon_0$ at one endpoint value and zero errors
in the other three data gives $b=\epsilon_0 h_{00}$ and
$\|b'\|_\infty=3\epsilon_0/(2\Delta)$. Accurate values alone do not
supply the differentiated residual required by the metric.

\section{The finite-mesh boundary solve is a contraction on a jet space}
\label{v52:contraction}

Let $\mathcal G_S$ be exactly \eqref{v51:Green}. Define
\begin{equation}
 e_\Delta^{[0]}=0,\qquad
 e_\Delta^{[m+1]}=\mathcal G_S\mathcal I_\Delta g[e_\Delta^{[m]}],
 \qquad W_m=w_{\rm ref}+e_\Delta^{[m]}.
 \label{v52:scheme}
\end{equation}
This differs from V51: its nonlinear forcing is evaluated only at the
$J+1$ nodes, together with \eqref{v52:gprime}. Its initial nonexpanding
and terminal expanding error coordinates vanish exactly. It is not a
forward solver fixing the expanding initial coordinates.

Use the norm
\[
 \|e\|_{1,S}=\|e\|_\infty+L_S\|e'\|_\infty.
\]
The derivative identity \eqref{v51:Green-derivative} and the reference
semigroup bounds give
\begin{equation}
 \|\mathcal G_S f\|_{1,S}
 \le C L_S(\|f\|_\infty+L_S\|f'\|_\infty).
 \label{v52:Green-C1}
\end{equation}
No inverse of $D$ occurs in this estimate.

\begin{theorem}[A complete finite-node nonlinear boundary reconstruction]
\label{v52:mesh-solve}
After reducing $c$ and the amplitude threshold, \eqref{v52:scheme} is a
contraction with factor at most $1/4$ in a ball
$\|e\|_{1,S}\le K L_Sa^3$. Its limit $e_\Delta$ is globally $C^2$,
satisfies both original boundary conditions, and solves
\begin{equation}
 e_\Delta'=De_\Delta+\mathcal I_\Delta g[e_\Delta].
 \label{v52:collocated}
\end{equation}
For all iterates with $m\ge1$, and for the limit, on every closed cell,
\begin{equation}
 \begin{gathered}
 \|e\|_\infty\le C L_Sa^3,\quad
 \|e'\|_\infty\le Ca^3,\quad
 \|e^{(j)}\|_\infty\le Ca^3\quad(2\le j\le4),\\
 \|g[e]\|_\infty\le Ca^3,\quad
 \|g[e]^{(j)}\|_\infty\le Ca^4\quad(1\le j\le4).
 \end{gathered}
 \label{v52:cell-regularity}
\end{equation}
Higher derivatives in this display are cellwise; no unnecessary global
$C^4$ assertion is made. The error obeys
\begin{equation}
 \|e_\Delta^{[m]}-e_\Delta\|_{1,S}
                       \le C L_Sa^3\,4^{-m}.
 \label{v52:convergence}
\end{equation}
\end{theorem}
\begin{proof}
On the ball, $|w_{\rm ref}+e|\le M_0a$ and
$|w_{\rm ref}'+e'|\le Ca^2$. V51's derivative estimates give
\[
 \|g[e]-g[f]\|_\infty\le Ca\|e-f\|_\infty,
\]
\[
 \|g[e]'-g[f]'\|_\infty
 \le Ca\|e'-f'\|_\infty+Ca^2\|e-f\|_\infty.
\]
Since $aL_S$ is bounded by a small constant, combining these inequalities
with \cref{v52:Hermite-lemma} and \eqref{v52:Green-C1} gives a contraction
factor at most $C aL_S$. The image of zero has norm at most $CL_Sa^3$,
using the original $f_a,f_a'$ estimates. Choose $K$, then $c,a_0$ to
make the ball invariant and the factor at most $1/4$. Banach's theorem
proves existence, uniqueness in this ball, and \eqref{v52:convergence}.

The regularity bounds are uniform and do not follow by assuming an
$h$-independent fourth time derivative. The $C^1$ norm already gives
$e'=O(a^3)$, $g=O(a^3)$ and $g'=O(a^4)$. Put
$p=\mathcal I_\Delta g$. From $e'=De+p$ one obtains
$e''=De'+p'=O(a^3)$. In turn
\[
 g''=D^2\mathcal H_a(W)[W',W']+
                         (D\mathcal H_a(W)-D)e''=O(a^4).
\]
The second-derivative stability of the interpolant gives $p''=O(a^4)$,
and hence $e'''=O(a^3)$. The formula for $g'''$ has the leading term
$(D\mathcal H_a-D)e'''$ and the bounded higher field derivatives applied
to $W',W''$; it is $O(a^4)$. Third-derivative interpolation stability then
gives $p'''=O(a^4)$ and $e^{(4)}=O(a^3)$. Finally the fourth chain rule,
using V51's bounded fourth field derivative, yields $g^{(4)}=O(a^4)$.
This argument is an induction for the iterates and holds for the limit on
 each cell. The forcing and its first derivative match at nodes, so $e,e',e''$
are continuous there. Nothing requires the third derivative to match.
\end{proof}

\section{There is no remaining nonlinear time quadrature}
\label{v52:finite-Green}

On a cell let $p(s_j+u\Delta)=\sum_{k=0}^3b_{j,k}u^k$ be its Hermite
forcing. Define the entire matrix functions
\[
 \varphi_r(A)=\sum_{n=0}^\infty\frac{A^n}{(n+r)!},\qquad r\ge1.
\]
For the nonexpanding component the node recurrence is
\begin{equation}
 e^-_{j+1}=e^{\Delta D_-}e^-_j+
          \Delta\sum_{k=0}^3 k!\varphi_{k+1}(\Delta D_-)b^-_{j,k},
 \qquad e^-_0=0.
 \label{v52:forward}
\end{equation}
If $\widetilde b_{j,k}$ are the coefficients of $p(s_{j+1}-u\Delta)$,
the expanding component is evaluated backwards:
\begin{equation}
 e^+_j=e^{-\Delta D_+}e^+_{j+1}-
          \Delta\sum_{k=0}^3 k!\varphi_{k+1}(-\Delta D_+)
                                    \widetilde b^+_{j,k},
 \qquad e^+_J=0.
 \label{v52:backward}
\end{equation}
The restrictions to the two reference invariant spaces are meant here; no
large forward expanding exponential is used. The same formulas with a
partial cell give the continuous exponential-polynomial curve and its jets.

\begin{proposition}[Finite matrix-function evaluation with the center retained]
\label{v52:finite-functions}
Equations \eqref{v52:forward}--\eqref{v52:backward} equal the original
Green integrals of the cubic forcing. They require four $\varphi$ functions
and one exponential per reference block, reusable on a uniform mesh, and
$O(J)$ fixed-dimensional operations per sweep. They remain defined on the
zero-frequency center. For a fixed bound $\|A\|\le M$, truncation after
$n=N$ in any required function has a norm remainder bounded by a fixed
multiple of
\[
 e^M\frac{M^{N+1}}{(N+1)!}.
\]
Consequently the matrix-function evaluations admit certified finite
approximations to any requested positive tolerance, without inverting $D$.
\end{proposition}
\begin{proof}
Integrating the exponential series on the compact cell gives
\[
 \int_0^\Delta e^{(\Delta-r)A} (r/\Delta)^k\,dr
       =\Delta k!\varphi_{k+1}(\Delta A).
\]
Time reversal gives \eqref{v52:backward}, including its minus sign. All
series converge absolutely in operator norm, including at $A=0$. The
usual exponential-series tail follows by factoring $M^{N+1}/(N+1)!$
and bounding the remaining positive series by $e^M$. The $\varphi$ tails
are no larger up to fixed factorial constants. No assertion about the
cost of certifying the original field map is implied by this operation count.
\end{proof}

The curve in \cref{v52:mesh-solve} has a finite expression but is not a
polynomial state update defining a new action. The nonlinear action is still
evaluated in \eqref{v52:g}; only its forcing in this analytical boundary
reconstruction is interpolated. Approximate matrix evaluation must be
certified at the output jets as specified below, rather than assumed to
preserve exact node matching in floating arithmetic.

\section{The original differentiated residual meets a finite mesh budget}
\label{v52:residual}

Let $W_m$ be an exact-data sweep with $m\ge1$. Relative to the original
nonlinear equation its defect is exactly
\begin{equation}
 r_m=W_m'-\mathcal H_a(W_m)
      =\mathcal I_\Delta g[e_\Delta^{[m-1]}]-g[e_\Delta^{[m]}].
 \label{v52:residual-identity}
\end{equation}
Thus the finite sweep is not described as an exact original trajectory.

\begin{theorem}[Finite-mesh error in the original nonlinear readers]
\label{v52:reader-theorem}
All exact-normal reconstructions $\Theta_a(W_m)$ satisfy all original
nonlinear constraints. Relative to V51's original-action boundary solution,
their path-reader errors on the complete horizon satisfy
\begin{equation}
 \begin{aligned}
 \|\Delta R\|&\le C(2^{-m}+\Delta^3),\\
 \|\Delta E^{\rm link}\|&\le C(a^2 2^{-m}+a^3\Delta^4),\\
 \|\Delta F^{\rm sp,link}\|&\le C(a^4 2^{-m}+a^5\Delta^4).
 \end{aligned}
 \label{v52:reader-bounds}
\end{equation}
The same orders hold for instantaneous original-action readers on these
exactly constrained records. The two reader interpretations are proved
separately, not identified on a curve with nonzero defect.
\end{theorem}
\begin{proof}
Split \eqref{v52:residual-identity} into its interpolation defect and its
successive-iterate difference. The $C^1$ contraction and
\eqref{v52:cell-regularity} give
\begin{equation}
 \begin{split}
 \delta_0:=\|r_m\|&\le C(a^3 2^{-m}+a^4\Delta^4),\\
 \delta_1:=\|r_m'\|&\le Ca^4(2^{-m}+\Delta^3),\\
 \delta_{q,1}:=\|(r_m)_q'\|&\le Ca^4(2^{-m}+\Delta^3).
 \end{split}
 \label{v52:residual-orders}
\end{equation}
The canonical bound is stated conservatively; no improved uncomputed
coefficient is needed. These are uniform bounds on every cell. Since $W_m$
is $C^2$ and $g[\cdot]$ is $C^1$, they also hold globally for the displayed
residual derivatives. The boundary defect is zero and, for sufficiently
small $a$, $\delta_0\le a^2$. V51's actual nonlinear residual theorem
therefore gives a state error
\[
 E\le C L_S\delta_0\le C(a^2 2^{-m}+a^3\Delta^4).
\]
Insert this, \eqref{v52:residual-orders}, and
$\delta_{q,0}\le\delta_0$ into \eqref{v51:residual-readers}. The metric
terms are bounded by $C(2^{-m}+a\Delta^4+\Delta^3)$; the last two
readers give precisely \eqref{v52:reader-bounds}.

For instantaneous readers, compare $\mathcal H_a(W_m)$ with
$\mathcal H_a(w_*)$ on the actual chart. The full rate difference is
bounded by $CE$, while the free canonical rows gain $a^2$. Their physical
synthesis and the original canonical acceleration formula cost at most
$Ca^{-2}E$ in the metric, $CE$ in the temporal link, and $Ca^2E$ in the
spatial link. These are no larger than the displayed path bounds. This
separate argument retains the original signed rate at the candidate state;
it does not replace the candidate's path derivative by that rate.
\end{proof}

For example, forcing accuracy is not measured only by the nodal mismatch.
Even at a collocation fixed point that mismatch is zero at every node,
while the interior defect in \eqref{v52:residual-identity} need not vanish.
The global interpolation bound, including its derivative, is what pays the
metric error.

\section{Field-jet errors and a globally twice differentiable finite output}
\label{v52:rounding}

At every sweep suppose the node values of $g,g'$ are enclosed and selected
with errors $\epsilon_0,\epsilon_1$. Use one common selected pair at each
node, and define
\[
 \eta=\max\{\epsilon_0/\Delta,\epsilon_1\}.
\]
The data need not be unbiased or independent, and may vary between sweeps.
Assume $\eta\le c_1a^4\Delta^2$ for a fixed bounded $c_1$; reducing the
original chart constants absorbs this allowance.

\begin{proposition}[Certified finite data do not require exact field jets]
\label{v52:inexact}
The inexact sweeps stay in the same $C^1$ ball after a fixed enlargement.
Their reader estimates become
\begin{equation}
 \begin{aligned}
 \|\Delta R\|&\le C(2^{-m}+\Delta^3+a^{-4}\eta),\\
 \|\Delta E^{\rm link}\|&\le C(a^2 2^{-m}+a^3\Delta^4+a^{-2}\eta),\\
 \|\Delta F^{\rm sp,link}\|&\le C(a^4 2^{-m}+a^5\Delta^4+\eta).
 \end{aligned}
 \label{v52:inexact-readers}
\end{equation}
These statements still use the exact normal map. Its separate finite
replacement is treated in \cref{v52:normal-section}.
\end{proposition}
\begin{proof}
The interpolated data error has value $O(\Delta\eta)$ and first derivative
$O(\eta)$. Its image under \eqref{v52:Green-C1} has norm at most
$CL_S^2\eta$. The inexact contraction inequality, summed as a geometric
series, bounds the difference from the exact fixed point by
$CL_Sa^3 4^{-m}+CL_S^2\eta$. Successive differences have the same bound
with an adjusted constant. In particular the value and first derivative
errors introduced by the data are $O(a^{-2}\eta)$ and $O(a^{-1}\eta)$.
The second and third derivatives of the interpolated error are
$O(\eta/\Delta)$ and $O(\eta/\Delta^2)$; the hypothesis makes these
$O(a^4)$. Thus the inductive cellwise regularity proof still applies.
The original residual bounds become
\[
 \delta_0\le C(a^3 2^{-m}+a^4\Delta^4+a^{-1}\eta),\qquad
 \delta_1,\delta_{q,1}\le C(a^4(2^{-m}+\Delta^3)+\eta).
\]
Apply the same original residual theorem. Its state budget is now
$E\le C(a^2 2^{-m}+a^3\Delta^4+a^{-2}\eta)$, which gives
\eqref{v52:inexact-readers}. The proof is deterministic for every error
sequence obeying the stated enclosures.
\end{proof}

There is a finite polynomial representation which also handles errors in
the matrix-function evaluation. Let $\mathcal Q_\Delta W_m$ be the
piecewise quintic polynomial matching $W_m,W_m',W_m''$ at each node.
Use one set of nodal jets, so it is globally $C^2$. With exact node data,
the canonical rows are reproduced exactly: on each cell their derivative
is the cubic canonical forcing, since the canonical row of $D$ is zero.
Thus those rows of $W_m$ have degree at most four. In the other rows,
$W_m^{(6)}=O(a^3)$ on each cell, by differentiating the constant-coefficient
Green equation after its cubic forcing has disappeared. Taylor remainder
and polynomial reproduction give
\[
 \|\partial_s^j(\mathcal Q_\Delta W_m-W_m)\|
       \le Ca^3\Delta^{6-j},\quad j=0,1,2.
\]
The physical chart consequently adds only $Ca\Delta^4$ to the metric
bound, $Ca^3\Delta^5$ to the temporal-link bound, and
$Ca^6\Delta^6$ to the spatial-link bound. Each is dominated by the mesh
terms already used in \eqref{v52:inexact-readers}.

\begin{lemma}[A nodal-jet certificate for a finite polynomial state curve]
\label{v52:quintic}
If the selected nodal jets have scaled error
\[
 \nu=\max_j\{\|\delta W_j\|,\ \Delta\|\delta W_j'\|,
                              \ \Delta^2\|\delta W_j''\|\},
\]
the quintic state error in derivative order $j\le2$ is at most
$C\nu\Delta^{-j}$. Its additional physical path-reader errors are bounded by
\begin{equation}
 C a^{-4}\nu\Delta^{-2},\qquad
 C a^{-1}\nu\Delta^{-1},\qquad C a^2\nu
 \label{v52:node-reader}
\end{equation}
for metric, temporal link and spatial link respectively, provided the
candidate remains in the same fixed chart.
\end{lemma}
\begin{proof}
On the unit interval the six quintic Hermite basis functions and their first
two derivatives are bounded. Rescaling proves the jet-error bound; one
shared triple per node proves $C^2$ matching. For the metric, the largest
linear synthesis cost is $a^2a^{-6}=a^{-4}$ on the canonical second time
jet. The nonlinear chart terms have no larger cost on curves whose scaled
first and second derivatives are $O(a^2),O(a^3)$. The first canonical time
jet costs at most $a^{-1}$ and the stationary connection's multiplier
responses gain $a,a^2$, giving the temporal bound. Its spatial value costs
at most $a^2$. Quadratic difference terms are smaller on a sufficiently
small chart. These are path-reader bounds, not a declaration that the
polynomial solves the original equation.
\end{proof}

A finite approximation of an exponential is therefore not permitted to create
uncontrolled seams. It is accepted through one shared set of certified nodal
jets and this $C^2$ polynomial output. Neither a small value residual nor an
unstable forward recurrence is used as a substitute for that certificate.

\section{Finite normal solves: residuals, not fictitious exact constraints}
\label{v52:normal-section}

Let $\mathcal P(a,n,w)$ be the original divided constraint map in the proof
of \cref{v51:normal}, and put $A_0=\mathcal C_0J_c$. At $a=0$,
$\mathcal P(0,n,w)=A_0n$. Consider the frozen-Jacobian iteration
\begin{equation}
 n^{[0]}=0,\qquad
 n^{[j+1]}=n^{[j]}-A_0^{-1}\mathcal P(a,n^{[j]},w).
 \label{v52:normal-iteration}
\end{equation}
It evaluates the original constraint, not a replaced constraint polynomial.

\begin{theorem}[A uniform differentiated normal solve]
\label{v52:normal-theorem}
On a sufficiently small fixed $(a,w)$-chart, \eqref{v52:normal-iteration}
stays on the original normal branch and
\begin{equation}
 \|n^{[j]}-n\|_{C_w^2}
 +\|\mathcal P(a,n^{[j]},\cdot)\|_{C_w^2}
                                  \le Ca\,2^{-j}.
 \label{v52:normal-error}
\end{equation}
More generally a twice differentiable candidate $\widetilde n$ on this
branch has $C^2$ error at most $C\epsilon_n$ if its divided-constraint
residual is at most $\epsilon_n$ in $C^2$ and its first two derivatives stay
bounded in the fixed chart. Reconstructing the same finite state curve with
$\widetilde n$ instead of $n$ gives the original raw residuals
\begin{equation}
 \|P_{\rm nc}\|\le Ca^2\epsilon_n,\qquad
 \|P_0P\|\le Ca^3\epsilon_n,
 \label{v52:raw-constraints}
\end{equation}
and adds at most
\begin{equation}
 C a^{-1}\epsilon_n,\qquad C a\epsilon_n,\qquad C a^2\epsilon_n
 \label{v52:normal-readers}
\end{equation}
to the original metric, temporal-link and spatial-link path errors.
\end{theorem}
\begin{proof}
The derivative of the iteration map in $n$ is
$I-A_0^{-1}\mathcal P_n$. It is zero at $a=0$ on the fixed chart.
Shrinking that chart makes its norm at most $1/4$; a fixed normal ball is
invariant because $n=O(a)$ in $C^2$. The undifferentiated error decreases
by $4^{-j}$. Differentiating the error recurrence once and twice in $w$
gives the same leading contraction plus uniformly bounded multiples of the
lower derivative errors. Summing the resulting triangular recurrences gives
$Ca(1+j)^2 4^{-j}\le Ca2^{-j}$. Composition with the bounded analytic
$\mathcal P$ gives its residual bound.

For the a posteriori assertion, write
$\mathcal P(a,\widetilde n,w)-\mathcal P(a,n,w)=A(a,w)(\widetilde n-n)$,
where $A$ is the integral of $\mathcal P_n$ along the normal segment.
It has a uniformly bounded inverse by the same $1/4$ criterion.
Differentiating this identity twice and using the stated candidate bounds
proves the $C^2$ error estimate. Undoing the original row divisions gives
\eqref{v52:raw-constraints} with every mean retained.

The physical canonical error is $a^2J_c(\widetilde n-n)$; multiplier values
are unchanged. Along the finite state curves constructed above,
$W'=O(a^2)$ and $W''=O(a^3)$. Its first physical time derivative is
therefore $O(a\epsilon_n)$ and its second is $O(a^{-1}\epsilon_n)$.
The quadratic chain-rule term is $O(\epsilon_n)$ and is covered. The
original metric and stationary-link formulas give \eqref{v52:normal-readers}.
The statement is about path readers on the actual local stationary chart;
an off-constraint candidate is not called an exact constrained history.
\end{proof}

Exact constraint satisfaction belongs to $\Theta_a(W)$, using $n$.
A finite inexact normal output has the small but generally nonzero residual
\eqref{v52:raw-constraints}. The exact-normal shadow of the same free curve
exists and is controlled by \eqref{v52:normal-readers}. Mean lapse remains
one even for the inexact normal output, because it is a multiplier-coordinate
normalization and is not solved through $n$.

\section{A finite information theorem at any fixed reader order}
\label{v52:main-section}

The next theorem assembles the estimates without treating floating arithmetic
as exact. A field-jet query means an evaluation of \eqref{v52:g} and
\eqref{v52:gprime} against the original reduced map, with a verified error.
The certified root, initial preparation, stationary equations and divided
constraint solve supply the analytic objects to be evaluated; no positive-time
exact trajectory is queried.

The final normal-output tolerance below is not automatically a sufficient
tolerance for a field-jet query. In particular, errors in a stationary or
normal solve must be propagated through the original signed inverse when
certifying $\mathcal H_a$ and $D\mathcal H_a$. An evaluator can refine those
solves further; it may not substitute an approximate normal into the rate
and infer the requested field precision solely from the smaller raw
constraint residual. The evaluation bounds always refer to the exact
original maps, independently of the final output reconstruction.

\begin{theorem}[Finite-node original-action reconstruction with quantified errors]
\label{v52:main}
Fix $p>0$. For every sufficiently small positive $a$, take
\begin{equation}
 S=c/a,\quad J=\lceil Sa^{-p/3}\rceil,\quad
 \Delta=S/J,\quad
 m=3+\lceil p\log_2(1/a)\rceil.
 \label{v52:parameters}
\end{equation}
In particular $a^{p/3}/2\le\Delta\le a^{p/3}$. Require the field-jet errors
and output nodal-jet error to obey
\begin{equation}
 \epsilon_0\le c_p a^{p+4}\Delta,\qquad
 \epsilon_1\le c_p a^{p+4},\qquad
 \nu\le c_p a^{p+4}\Delta^2.
 \label{v52:evaluation-budget}
\end{equation}
Use an inexact normal map with $\epsilon_n\le c_p a^{p+2}$ in the divided
$C^2$ residual, or obtain it by
$j\ge j_0+\lceil(p+1)\log_2(1/a)\rceil$ iterations of
\eqref{v52:normal-iteration}. Then the resulting finite $C^2$ state output,
after its stated nonlinear normal reconstruction, has path-reader errors
\begin{equation}
 \boxed{\|\Delta R\|\le C_pa^p,\quad
 \|\Delta E^{\rm link}\|\le C_pa^{p+2},\quad
 \|\Delta F^{\rm sp,link}\|\le C_pa^{p+4}}
 \label{v52:main-readers}
\end{equation}
relative to V51's exact original-action boundary history, on the whole
physical interval $[0,ca^2]$. Its original constraint residuals satisfy
\begin{equation}
 \boxed{\|P_{\rm nc}\|\le C_pa^{p+4},\qquad
              \|P_0P\|\le C_pa^{p+5}.}
 \label{v52:main-constraints}
\end{equation}
Using the exact normal map instead gives exact zero constraints.
The number of field-jet queries is
\begin{equation}
 O_p\!\left(a^{-1-p/3}\log(1/a)\right),
 \label{v52:query-count}
\end{equation}
with $O(J)$ fixed-dimensional Green operations per sweep. The final
polynomial state representation has $J$ cells of degree at most five.
\end{theorem}
\begin{proof}
The mesh relation follows from the ceiling and $S\ge1$, for sufficiently
small $a$. With the stated precision,
$\eta\le c_pa^{p+4}$ and
$\eta/\Delta^2\le4c_pa^{4+p/3}$, so the inexact regularity and invariance
conditions hold after reducing $a_0$. Also $2^{-m}\le a^p/8$.
Equation \eqref{v52:inexact-readers} gives the desired powers: the mesh
terms in the two link bounds are $a^{3+4p/3}$ and $a^{5+4p/3}$, each
smaller than the corresponding requested order for every $p>0$.
The quintic interpolation terms are smaller as well. Substitution of $\nu$
in \eqref{v52:node-reader} gives at most
$a^p,a^{p+3}\Delta,a^{p+6}\Delta^2$.
The normal-map contribution in \eqref{v52:normal-readers} is at most
$a^{p+1},a^{p+3},a^{p+4}$. Summing proves \eqref{v52:main-readers}, and
\eqref{v52:raw-constraints} gives \eqref{v52:main-constraints}.
Each sweep uses $J+1$ complete field and first-derivative evaluations, plus
fixed-dimensional linear operations, proving \eqref{v52:query-count}.
The original maps and matrix functions have local finite evaluation
procedures with refinable enclosures on the inherited nonsingular charts.
This statement requires the stated error certificates; it is not an assertion
that an arbitrary implementation attains them at a fixed machine precision.
No bit-complexity bound or numerical amplitude threshold is inferred.
\end{proof}

For the useful choice $p=2$, the mesh has $J=O(a^{-5/3})$ cells and
$m=O(\log(1/a))$ sweeps. The original reader errors are
$(a^2,a^4,a^6)$; the raw nonconstant and mean constraint residuals are
$(a^6,a^7)$. Sufficient nodal field precisions are
$\epsilon_0=O(a^{20/3})$, $\epsilon_1=O(a^6)$, and the scaled output
jet tolerance is $\nu=O(a^{22/3})$. The differentiated normal tolerance
is $O(a^4)$. These sufficient powers are not claimed optimal. They show
that exponentially growing physical initial-value modes do not require
exponentially small data tolerances in this two-point reconstruction;
its expanding boundary information has not been removed.

\section{Status, actual verification, and the next independent gate}
\label{v52:status}

\paragraph{Proved in this body.}
The nonlinear boundary reconstruction now has a finite forcing-interpolation
scheme with globally $C^2$ output, populated differentiated residual budgets,
a finite-node query count, deterministic field-jet error bounds, and a
separate nonlinear normal-solve certificate. The exact center is retained
through entire $\varphi$ functions, not an inverse at zero. The positive
comparison result concerns the complete original metric and link path
readers, each relative to its own value on the exact original-action
boundary solution. It does not identify these readers with one another.

\paragraph{Inherited, not recomputed.}
The certified initial root, signed source matrices, exact nonlinear normal
chart, reduced original field estimates, and existence of the reference
boundary solution remain the named V51 and earlier inputs. The new proof
populates V51's truncation and normal-map budgets; it does not independently
prove every historical action calculation. The supplied original maps must
be evaluated with the stated field-jet and normal certificates.

\paragraph{Not claimed.}
No numerical integration of the 323-dimensional gravitational boundary
problem, numerical chart radius, total bit complexity, native renewal
implementation, causal initial-value mechanism, new cofinal family, or longer
physical lifespan has been performed or proved here. The exact path and its
boundary information are fixed by V51. The inexact finite output generally
has nonzero constraints, with the displayed bounds, rather than fictitious
exact zeros. The benchmark program is an explicitly separate low-dimensional
nonlinear example of the time kernel; it is not gravitational evidence.

The implementation gap is now specified by finite, checkable evaluation
budgets rather than uncounted exact nonlinear integrals. A useful subsequent
step is a complete numerical certificate for the original stationary-map
and boundary evaluations, or the independent central/slow physical-time
and spatial-refinement problem. The present theorem does not turn the
shrinking $ca^2$ interval at $N=3$ into a continuum Einstein result. The
raw-action Einstein bridge remains unresolved.

\paragraph{Verification boundary.}
The accompanying checker verifies the Hermite and quintic interpolation
identities, the forward/backward entire-function recurrence including zero
modes, differentiated residual identities, finite normal corrections, all
precision exponents, and preservation of the fifty-one earlier marked bodies.
It also reruns the unchanged predecessor's standalone algebra groups.
Floating tests of the delivered time kernel use a separately declared
three-coordinate nonlinear system and are reported separately from exact
checks. They do not replace any written uniform proof. Actual executed
reports, source hashes and compilation diagnostics accompany this body;
no proof-assistant formalization is claimed.

\begin{thebibliography}{9}
\raggedright
\bibitem{v52:HO}
M. Hochbruck and A. Ostermann, \emph{Exponential integrators},
Acta Numerica \textbf{19} (2010), 209--286.
\href{https://doi.org/10.1017/S0962492910000048}{doi:10.1017/S0962492910000048}.
The primary publisher and author-institution records were checked. This is
methodological background for entire matrix functions, not an imported
uniform estimate for the present gravitational source or a new attribution
of the boundary selection mechanism.
\end{thebibliography}
% END IMMUTABLE EXACT-ACTION BRIDGE V52 BODY


% BEGIN IMMUTABLE EXACT-ACTION BRIDGE V53 BODY
\clearpage
\renewcommand{\thesection}{\arabic{section}}
\renewcommand{\theHsection}{bridgeV53.\arabic{section}}
\setcounter{section}{429}

\section{Fixed positive time is now a robustness question, not an integration question}
\label{v53:context}

V52 closes a finite-node approximation problem on the already constructed
boundary-selected branch.  It does not change the geometry of the exact
constraint leaf.  In particular, V43 proved that the original normalized
action has, at the same certified root and fixed cutoff, a codimension-eleven
locally invariant constrained manifold \(\mathcal S_a^\chi\) which is normally
repelling in eleven directions.  That conclusion is independent of the later
Hermite reconstruction.  The present body returns to this exact normal
factorization and asks a different question: how accurately must an original
constrained initial record lie on the invariant manifold in order to remain in
one fixed local chart for a prescribed \emph{physical} time?

The answer is exponentially strong in \(a^{-3}\).  This does not contradict the
existence of the exactly selected histories on \(\mathcal S_a^\chi\).  It does
show that a fixed-time open basin around them cannot have any algebraic normal
thickness.  Thus another improvement of the finite-node quadrature cannot, by
itself, supply a robust fixed-time exact-action Einstein basin.

\paragraph{Inherited inputs and nonduplication.}
We use the constrained invariant manifold and exact normal factorization
\cref{v43:constraint-manifold,v43:normal-factor,v43:selection-theorem}; the
physical/normalized scaling of the anchored controls is the one stated in
\eqref{v43:transversality}.  V42's endpoint comparison and V52's finite-node
reconstruction are retained only for comparison of scales.  None of their
boundary selections is reproved.  The new work is the quantitative lower
normal growth, its physical-time exit theorem, the no-algebraic-thickness
corollary, and the distinction between exact manifold selection and robust
nearby preparation.

\section{Quantitative normal expansion on the original constrained leaf}
\label{v53:normal-growth}

Retain V43's normalized coordinates \(u\), scaled time \(s=t/a^3\), and local
normal coordinate
\[
 v(u)=P_u u-h(a,P_nu).
\]
On the original constraint leaf, while the state remains in the verified
chart,
\begin{equation}
 v'=N_a(u)v,
 \qquad \|N_a(u)-J_u\|\le C(|a|+\|u\|).
 \label{v53:normal-equation}
\end{equation}
V43 also supplies a positive definite matrix
\[
 W_u=\int_0^\infty e^{-rJ_u^*}e^{-rJ_u}\,dr
\]
for which, after shrinking the common chart,
\begin{equation}
 \frac d{ds}\langle v,W_uv\rangle
 \ge c_*\langle v,W_uv\rangle .
 \label{v53:Lyapunov-lower}
\end{equation}
All constants below are fixed-cutoff constants, uniform for sufficiently small
positive amplitude in this chart.

\begin{theorem}[Two-sided normal growth until chart exit]
\label{v53:normal-growth-theorem}
There are constants \(\mu>0\), \(c_0,C_0>0\) and a fixed chart radius
\(\rho_*>0\) such that every original normalized exact-action history on the
constraint leaf satisfies
\begin{equation}
 c_0 e^{\mu s}\|v(0)\|
 \le \|v(s)\|
 \le C_0 e^{C_0s}\|v(0)\|
 \label{v53:two-sided-normal}
\end{equation}
for every \(s\ge0\) prior to the first time \(\|v(s)\|=\rho_*\).  The lower
bound applies to every nonzero normal displacement; no genericity assumption
on its direction is required.
\end{theorem}

\begin{proof}
Let \(m_W,M_W>0\) be the extremal eigenvalue bounds of \(W_u\).  Since the
matrix is fixed,
\[
 m_W\|v\|^2\le\langle v,W_uv\rangle\le M_W\|v\|^2.
\]
Integrating \eqref{v53:Lyapunov-lower} gives
\[
 \langle v(s),W_uv(s)\rangle
 \ge e^{c_*s}\langle v(0),W_uv(0)\rangle,
\]
so
\[
 \|v(s)\|\ge (m_W/M_W)^{1/2}e^{c_*s/2}\|v(0)\|.
\]
Take \(\mu=c_*/2\).  The upper estimate follows from
\(\|v'\|\le(\|J_u\|+C\delta_*)\|v\|\) on the fixed chart and Gronwall.
Every identity is for the original normal factor \eqref{v53:normal-equation};
there is no replacement dynamics or frozen-matrix approximation in the
conclusion.
\end{proof}

\section{Exit time and the preparation precision required by a fixed physical horizon}
\label{v53:exit}

V43's anchored exact constrained family is
\[
 \widehat Q(a,\vartheta,d_a^\chi(\vartheta)+e),
 \qquad e\in E_u,
\]
where \(e=0\) is the invariant-manifold selection.  Its exact transversality
estimate is
\begin{equation}
 c_1a^3\|e\|
 \le \|v(0)\|
 \le C_1a^3\|e\|.
 \label{v53:transversality}
\end{equation}
The physical normal displacement is one further factor \(a\) smaller, because
the normalized chart itself is \(u=\mathcal T_a(Q-Q_a^0)/a\).  In the fixed
chart norm we therefore write
\begin{equation}
 \delta_{\rm phys}(0)\asymp a\|v(0)\|\asymp a^4\|e\|.
 \label{v53:physical-normal}
\end{equation}
This is the normal component; the total physical record change may contain
larger tangential terms.

\begin{theorem}[Necessary precision for residence through a physical time]
\label{v53:fixed-time-precision}
Fix \(T>0\), independent of amplitude.  Suppose an anchored, exactly
constrained initial record with control error \(e\ne0\) has an original
action history which remains inside the verified local chart throughout
\([0,T]\).  Then
\begin{equation}
 \boxed{
 \|e\|\le C a^{-3}\exp\!\left(-\mu T/a^3\right).}
 \label{v53:e-fixed-T}
\end{equation}
Equivalently, its physical normal displacement must satisfy
\begin{equation}
 \boxed{
 \delta_{\rm phys}(0)
 \le C a\exp\!\left(-\mu T/a^3\right).}
 \label{v53:physical-fixed-T}
\end{equation}
More generally, if \(T_a\) is any positive horizon for which the history
remains in the chart, the same bounds hold with \(T\) replaced by \(T_a\).
\end{theorem}

\begin{proof}
The rescaled endpoint is \(S=T/a^3\).  Residence in the chart implies
\(\|v(S)\|<\rho_*\).  Apply the lower estimate of
\cref{v53:normal-growth-theorem} and then
\eqref{v53:transversality}:
\[
 \rho_*>c_0e^{\mu S}\|v(0)\|
 \ge c_0c_1a^3e^{\mu T/a^3}\|e\|.
\]
This proves \eqref{v53:e-fixed-T}.  Multiply by the fixed norm-equivalence
factor in \eqref{v53:physical-normal} to obtain
\eqref{v53:physical-fixed-T}.  The argument is conditional only on residence
inside the chart; if the inequality fails, the theorem concludes that the
history cannot stay in that verified neighborhood for the full horizon.  It
does not assert blowup or nonexistence after exit.
\end{proof}

\begin{corollary}[Exit time for algebraically prepared errors]
\label{v53:algebraic-exit}
For a nonzero anchored control error satisfying
\(\|e(a)\|\ge c a^m\) with fixed \(m\in\R\), the first exit from the verified
normal chart, if the exact history exists until then, obeys
\begin{equation}
 \boxed{
 t_{\rm exit}(a)
 \le \frac{a^3}{\mu}
 \left[(m+3)\log(1/a)+C_m\right]
 =O\!\left(a^3\log(1/a)\right).}
 \label{v53:algebraic-exit-time}
\end{equation}
In particular no polynomial control accuracy can certify residence for one
fixed positive physical time as \(a\downarrow0\).
\end{corollary}

\begin{proof}
At exit the lower normal-growth estimate reaches the fixed radius \(\rho_*\).
Solving
\[
 c_0c_1c\,a^{m+3}e^{\mu s}\le\rho_*
\]
for \(s\), and multiplying by \(a^3\), gives the stated bound.  If the expression in brackets is negative, the
record starts outside the smaller normal tube for which the estimate was
posed, so the assertion is vacuous rather than contradictory.
\end{proof}

\section{No algebraically thick fixed-time basin around the selected exact-action manifold}
\label{v53:no-thick-basin}

The preceding theorem can be stated intrinsically, without reference to the
particular control coordinates.

\begin{theorem}[No algebraically thick robust fixed-time tube]
\label{v53:no-algebraic-basin}
Fix \(T>0\).  Let \(\mathfrak U_a(T)\) be any tubular neighborhood of
\(\mathcal S_a^\chi\) inside the original normalized constraint leaf with the
property that every exact-action history starting in \(\mathfrak U_a(T)\)
remains inside the verified local chart through time \(T\).  If
\(r_a\) denotes its physical radius in the unstable normal directions, then
\begin{equation}
 \boxed{
 r_a\le C a\exp\!\left(-\mu T/a^3\right).}
 \label{v53:tube-radius}
\end{equation}
Consequently there is no such family with
\(r_a\ge c a^m\) for any fixed power \(m\).
\end{theorem}

\begin{proof}
Take a point of the exact constrained anchored chart whose normalized normal
coordinate has fixed direction and physical size smaller than the proposed
radius.  V43's transverse selection map realizes such points for every
sufficiently small normal coordinate.  Apply
\cref{v53:fixed-time-precision} and let the chosen size approach the tubular
radius.  The exponential upper bound follows.  Since
\(a e^{-\mu T/a^3}=o(a^m)\) for every fixed real \(m\), no algebraically thick
radius is possible.
\end{proof}

\begin{remark}[What the theorem does not rule out]
\label{v53:exact-graph-exception}
The theorem does not rule out an exact history starting on
\(\mathcal S_a^\chi\) and remaining in a chart for a fixed positive time.  Its
normal coordinate is identically zero by local invariance.  Nor does it show
that a history leaving this particular chart ceases to exist or has large
curvature.  The result is a robustness obstruction: a fixed-time open basin
around the selected manifold must have exponentially thin normal width.
Longer control of the exact manifold itself is a separate central/slow
problem.
\end{remark}

\section{Relation to the finite reconstruction and to the open stable writer}
\label{v53:relation}

The result clarifies two distinctions which had become easy to blur after the
finite reconstruction programme.

\paragraph{Finite reconstruction versus forward robustness.}
V52 proves that the two-point exact-action boundary solution on the shrinking
\(ca^2\) interval can be approximated to arbitrary algebraic reader accuracy
using finite forcing data.  That theorem remains intact.  It does not imply
that an algebraically accurate perturbation of its initial record can be used
as a forward exact-action datum for one fixed physical time.  The latter
question pays the normal amplification \eqref{v53:fixed-time-precision}.
Constraint accuracy alone does not certify membership in the repelling
invariant manifold.

For comparison, if the target physical horizon is the already established
\(T_a=ca^2\), then \eqref{v53:e-fixed-T} gives the necessary normal-control
scale
\begin{equation}
 \|e\|\le C a^{-3}e^{-\mu c/a}.
 \label{v53:ca2-precision}
\end{equation}
This is consistent with the exponential endpoint/control sensitivities found
in V42.  The present derivation is different: it follows from the exact local
normal factor and therefore applies to any exact constrained off-manifold
history that stays in the chart, not only to the earlier boundary comparison.

\paragraph{Exact action versus the stable writer.}
The stable local writer in the companion programme has a cutoff-independent
small-data ball and does not require codimension-eleven exponentially precise
selection.  The raw exact action studied here instead possesses a normally
repelling constrained manifold at this verified finite cutoff.  The two
statements concern different dynamics.  The present theorem therefore
sharpens, rather than erases, the reason that a constructive Einstein
population for the stable writer is not yet a direct universality theorem for
the raw action.

No continuum or spatial-refinement conclusion follows from this fixed-cutoff
comparison.  In particular, the normal exponent, chart radius, and invariant
manifold have not been controlled uniformly as the lattice is refined.

\section{Revised frontier after the fixed-time robustness obstruction}
\label{v53:frontier}

The programme should not spend another iteration reducing V52's quadrature
error while leaving the physical selection issue unchanged.  At this point
there are three logically distinct targets:

\begin{enumerate}[label=\textup{(\roman*)},leftmargin=2.8em]
\item \emph{Exact selected-manifold propagation.}  Prove a fixed positive
physical lifespan for histories lying exactly on the constrained invariant
manifold, including the central and slow canonical dynamics and the original
metric/link observations.
\item \emph{Physical selection.}  Derive the codimension-eleven manifold, or
an exponentially accurate approximation to it, from admissible microscopic
renewal data without future endpoint information.  The present theorem shows
that algebraic normal accuracy is insufficient for a fixed positive horizon.
\item \emph{Spatial refinement.}  Construct a cofinal family for which the
rank, invariant splitting, chart geometry, source response, and physical
observations have uniform estimates as the cutoff is removed.
\end{enumerate}

A proof of (i) alone would give existence of a special exact-action branch but
not robustness.  A proof of (ii) on the shrinking \(ca^2\) horizon would not
by itself give a fixed positive time.  A proof of (iii) without either of the
first two would not identify a dynamically selected Einstein population.
The raw-action Einstein continuum bridge therefore remains open, but its
fixed-time robustness obstruction is now quantitative.

\section{Verification and append-only preservation}
\label{v53:verification}

The new checker verifies the exponential/physical-time algebra, the control
and physical normal scalings, the \(ca^2\) specialization, the algebraic-error
exit-time exponents, and byte preservation of every preceding marked body.
These are exact symbolic and rational checks of the formulas used above.  The
normal-growth theorem itself is the written consequence of the inherited
exact factorization \eqref{v43:normal-factor} and its positive Lyapunov
inequality; it is not inferred from sampled eigenvalues.

No new gravitational coefficient extraction, numerical trajectory, numerical
chart radius, spatial-refinement computation, or proof-assistant
formalization is claimed.  The V52 time-kernel checker is rerun separately as
a regression.  The fixed-time result is local to the same \(N=3\) original
action chart, and \(a\) remains field amplitude rather than lattice spacing.
% END IMMUTABLE EXACT-ACTION BRIDGE V53 BODY


% BEGIN IMMUTABLE EXACT-ACTION BRIDGE V54 BODY
\clearpage
\renewcommand{\thesection}{\arabic{section}}
\renewcommand{\theHsection}{bridgeV54.\arabic{section}}
\setcounter{section}{435}

\section{A second selection layer inside the rapid-selected manifold}
\label{v54:context}

V53 proves that a fixed-positive-time neighborhood of the rapid-selected
constraint manifold \(\mathcal S_a^\chi\) must be exponentially thin in its
eleven rapid unstable directions, on the scale \(\exp(-cT/a^3)\).  It
explicitly leaves open the behavior of histories lying \emph{exactly} on that
manifold.  The next obstruction is already present in the inherited canonical
spectrum: V37 proves a simple real pair
\[
        z_\pm(a)=\pm \gamma/a+O(1),
        \qquad 0.1585835166<\gamma<0.1585835168,
\]
while V43 removes only the eleven expanding directions at the faster
\(a^{-3}\) scale.  The slow positive root is therefore a tangential direction
of \(\mathcal S_a^\chi\), not one of its rapid normals.

This body separates that slower instability from the rapid selection.  We
construct, inside \(\mathcal S_a^\chi\), a locally invariant codimension-one
slow center--stable manifold.  Its scalar normal grows on the physical scale
\(t/a\).  A fixed-positive-time robust tube inside the rapid-selected
manifold consequently has width at most \(\exp(-cT/a)\) in this second
normal direction.  Thus exact rapid selection alone does not produce an
algebraically thick fixed-time basin.

The construction is different from V44's 301-dimensional central manifold.
V44 removes the eleven rapid \emph{stable} directions as well as the eleven
rapid unstable directions.  The manifold below retains those contracting
rapid directions and removes only the one slow expanding direction.  It is
therefore the minimal forward center--stable selection suggested by the
known real spectrum.

\paragraph{Inherited inputs and nonduplication.}
We use the exact rapid-selected constrained manifold and its local chart from
\cref{v43:constraint-manifold,v43:normal-factor}, the complete slow-root count
and location from \cref{v37:full-count,v37:canonical-roots}, and the exact
ellipticity/definite energy of the intermediate quartet from
\cref{v46:ellipticity,v46:canonical-energy}.  The finite-dimensional
Lyapunov--Perron construction is standard background; the source-specific
splitting, amplitude scaling, normal equation, and physical-time precision
claims are proved below.  We do not import V44's different central manifold
or its stable-fibre selection as a substitute for the present graph.
All statements remain at \(N=3\), and \(a>0\) is field amplitude.

\section{The slow-time spectral gap on the rapid-selected tangent space}
\label{v54:slow-gap}

Let \(Q_a^\circ\in\mathcal S_a^\chi\) be one of the selected reference
states in the common original-action chart.  Let \(A_a\) denote the
linearization of the original normalized flow restricted to
\(T_{Q_a^\circ}\mathcal S_a^\chi\).  Introduce slow time
\begin{equation}
                         \tau=t/a.                  \label{v54:slow-time}
\end{equation}
Thus the slow-time linearization is \(\mathscr A_a=aA_a\).

\begin{lemma}[One separated slow unstable direction inside \(\mathcal S_a^\chi\)]
\label{v54:slow-splitting}
For all sufficiently small positive \(a\), there is an invariant splitting
\begin{equation}
 T_{Q_a^\circ}\mathcal S_a^\chi
       =E_{+,a}^{\rm sl}\oplus E_{0,a}^{\rm sl},
 \qquad \dim E_{+,a}^{\rm sl}=1,                    \label{v54:slow-split}
\end{equation}
with the following properties.
\begin{enumerate}[label=\textup{(\roman*)},leftmargin=2.8em]
\item On \(E_{+,a}^{\rm sl}\), \(\mathscr A_a\) has the simple real
eigenvalue
\begin{equation}
                  \lambda_+(a)=\gamma+O(a),
 \qquad \lambda_+(a)>3\gamma/4.                    \label{v54:lambda-plus}
\end{equation}
\item Every spectral value of \(\mathscr A_a|_{E_{0,a}^{\rm sl}}\) has
real part at most \(\gamma/8\).
\item There is an \(a\)-dependent norm \(\|\cdot\|_{\rm sl,a}\), polynomially
equivalent to a fixed Euclidean norm, in which the Riesz projections of
\eqref{v54:slow-split} are uniformly bounded and
\begin{equation}
 \|e^{\tau\mathscr A_a}P_{0,a}\|_{\rm sl,a}
       \le C_\delta e^{\delta|\tau|}
 \quad(\tau\ge0)                                   \label{v54:slow-center-bound}
\end{equation}
for any fixed \(\gamma/8<\delta<\gamma/4\), after reducing the amplitude
threshold.
\end{enumerate}
The contracting root \(-\gamma/a+O(1)\), the eleven rapid contracting
roots, the intermediate oscillatory quartet and all remaining center
coordinates belong to \(E_{0,a}^{\rm sl}\).
\end{lemma}

\begin{proof}
V43 constructs \(\mathcal S_a^\chi\) tangent to the complement of the
\emph{rapid} eleven-dimensional expanding cluster.  Its tangent is invariant
under the full linearization because the manifold is locally invariant.
The roots \(z_\pm(a)=\pm\gamma/a+O(1)\) of
\cref{v37:canonical-roots} satisfy \(a^3z_\pm(a)\to0\); they are not rapid
roots and hence lie in this tangent space.  Multiplication by \(a\) gives
\eqref{v54:lambda-plus} and its negative partner.

It remains to separate every other root in slow time.  The rapid stable
cluster has real part at most \(-c a^{-3}\), hence at most \(-c a^{-2}\)
after multiplication by \(a\).  The intermediate quartet is exactly
imaginary and semisimple by \cref{v46:ellipticity}.  On every fixed nonzero
slow annulus, \cref{v37:slow-annulus} bounds the real parts of all roots
outside the simple real pair by \(O(1)\); after multiplication by \(a\)
these tend to zero.  Any remaining sequence with \(az\to0\) also has
\(a\operatorname{Re}z\to0\).  Thus no second positive real part can remain
above \(\gamma/8\) for all sufficiently small \(a\).  This proves (ii).

For (iii), take the Riesz projection around the simple root
\(\lambda_+(a)\).  Its separation from the rest is fixed in slow time, so
its projection is uniformly bounded in a suitable spectral norm.  On the
intermediate quartet use the positive definite energy norm of V46.  On the
rapid stable cluster use its exponentially decaying norm.  On each remaining
finite-dimensional center block put the real Jordan form into the usual
Lyapunov norm; fixed polynomial factors in \(\tau\) are absorbed into
\(e^{\delta|\tau|}\).  All coordinate changes are assembled from the finite
analytic coefficient matrices and the already declared amplitude weights, so
their equivalence constants have at most fixed algebraic powers of \(a^{-1}\).
The direct sum of these block norms gives \eqref{v54:slow-center-bound}.
No inverse of a zero center frequency is taken.
\end{proof}

\section{A nested slow center--stable manifold}
\label{v54:nested-manifold}

Choose slow-adapted local coordinates \((w,n)\) on
\(\mathcal S_a^\chi\), with scalar \(w\in E_{+,a}^{\rm sl}\) and
\(n\in E_{0,a}^{\rm sl}\).  The coordinate is normalized so that unit
\(w\) has the V37 slow multiplier profile with constant-shift part of unit
size.  In physical record coordinates a displacement \(w\) therefore has
multiplier size \(\asymp a|w|\), active-multiplier size \(O(a^2|w|)\), and
canonical size \(O(a^3|w|)\), by the slow lift in
\cref{v37:slow-tangent}.  This normalization is used only to state the
physical normal scale below.

After the time change \eqref{v54:slow-time}, localize the restricted vector
field outside a smaller original-action chart.  The localization equals one
on every physical state used below.  In the slow norm it has the form
\begin{align}
 w_\tau&=\lambda_+(a)w+b_+(a)+N_+(a,w,n),\
 n_\tau&=A_{0,a}n+b_0(a)+N_0(a,w,n),                 \label{v54:slow-system}
\end{align}
where \(D_{(w,n)}N_\bullet(a,0,0)=0\).  By shrinking the local chart for each
small \(a\), its nonlinear Lipschitz constant is less than
\(\gamma/32\).  The chart need not have a uniform Euclidean radius; no such
claim is used in the fixed-time precision theorem.

\begin{theorem}[A second locally invariant constrained graph]
\label{v54:slow-manifold-theorem}
For every sufficiently small positive \(a\), the original normalized action
possesses, inside \(\mathcal S_a^\chi\), a locally invariant \(C^3\) manifold
\begin{equation}
 \boxed{\mathcal D_a^{\chi,\sigma}
    =\{(h_{\rm sl}(a,n),n):n\in E_{0,a}^{\rm sl}\}}
                                                               \label{v54:D}
\end{equation}
of dimension \(311\).  It is codimension one in \(\mathcal S_a^\chi\) and
codimension twelve in the \(323\)-dimensional original normalized constraint
leaf.  Its tangent at the reference point is \(E_{0,a}^{\rm sl}\).

The superscript \(\sigma\) records the harmless analytical localization used
outside the physical chart.  Every history claimed here evolves where that
localized field agrees with the unchanged original action.
\end{theorem}

\begin{proof}
Fix \(\delta,\nu\) with
\[
             \gamma/8<\delta<\nu<\gamma/2.
\]
For a small initial nonunstable coordinate \(n_0\), seek a forward trajectory
of \eqref{v54:slow-system} in the weighted norm
\(\sup_{\tau\ge0}e^{-\nu\tau}\|(w(\tau),n(\tau))\|_{\rm sl,a}\).
The Lyapunov--Perron equations are
\begin{align}
 n(\tau)&=e^{\tau A_{0,a}}n_0+
   \int_0^\tau e^{(\tau-r)A_{0,a}}
       \{b_0+N_0(w(r),n(r))\}\,\dd r,\notag\\
 w(\tau)&=-\int_\tau^\infty
    e^{\lambda_+(a)(\tau-r)}
       \{b_++N_+(w(r),n(r))\}\,\dd r.              \label{v54:LP-slow}
\end{align}
The center estimate \eqref{v54:slow-center-bound} and
\(\lambda_+(a)>3\gamma/4\) bound the two integral operators by constants
proportional to \((\nu-\delta)^{-1}\) and
\((3\gamma/4-\nu)^{-1}\), respectively.  The localized nonlinear
Lipschitz constant was chosen so their sum times that constant is below
one half.  The map is therefore a contraction.

Set \(h_{\rm sl}(a,n_0)=w(0)\).  Standard differentiated contraction
estimates give \(C^3\) regularity in the finite-dimensional coordinates.
Time translation of the unique weighted trajectory proves local invariance.
The tangent statement follows from the vanishing nonlinear differential at
the reference point.  Since \(\mathcal S_a^\chi\) already lies in every
original constraint row, so does its submanifold.  Its dimension is
\(312-1=311\).  This is the finite-dimensional center--stable construction
in the present action-specific splitting; compare the general invariant
foliation background already cited in \cite{v44:BLZ}.
\end{proof}

This new manifold is not V44's \(301\)-dimensional central manifold.  The
latter removes the eleven rapid contracting directions as well.  The present
\(\mathcal D_a^{\chi,\sigma}\) retains them because they are favorable for
forward time and removes only the additional slow expanding direction.

\section{The exact slow normal factor and its physical growth scale}
\label{v54:slow-normal}

Define the scalar slow normal coordinate on \(\mathcal S_a^\chi\) by
\begin{equation}
       \omega=w-h_{\rm sl}(a,n).                     \label{v54:omega}
\end{equation}
Subtracting the graph invariance equation from the full scalar equation and
integrating the state derivative along the segment in \(w\) gives the exact
factorization
\begin{equation}
 \boxed{
 \frac{\dd\omega}{\dd\tau}
       =\{\lambda_+(a)+r_a(w,n)\}\omega,}
 \qquad |r_a(w,n)|\le\gamma/4                       \label{v54:slow-factor}
\end{equation}
inside a smaller common slow-normal tube.  No frozen approximation appears
in this equation.

\begin{theorem}[Two-sided slow normal growth until exit]
\label{v54:slow-growth}
There are constants \(c_s,C_s>0\) and
\begin{equation}
                         \mu_s:=\gamma/2>0            \label{v54:mu-s}
\end{equation}
such that every original exact-action history lying in
\(\mathcal S_a^\chi\) satisfies
\begin{equation}
 \boxed{
 c_s e^{\mu_s t/a}|\omega(0)|
   \le |\omega(t)|
   \le C_s e^{2\gamma t/a}|\omega(0)|}              \label{v54:slow-growth-eq}
\end{equation}
for every physical time before exit from the verified slow chart.
The lower bound holds for every nonzero slow-normal displacement, with no
genericity condition on its sign.
\end{theorem}

\begin{proof}
By \eqref{v54:lambda-plus} and \eqref{v54:slow-factor}, after reducing the
amplitude threshold,
\[
       \gamma/2\le \lambda_+(a)+r_a\le 2\gamma.
\]
The scalar equation can be integrated exactly.  Replace
\(\tau\) by \(t/a\) to obtain \eqref{v54:slow-growth-eq}.  The constants
allow the fixed norm equivalences in the chosen slow coordinate.
\end{proof}

In the normalization preceding \cref{v54:slow-manifold-theorem}, the physical
slow-normal record distance is comparable to \(a|\omega|\): the constant
shift multiplier component has that size, whereas the active multiplier and
canonical components carry two and three powers of \(a\), respectively.
This statement concerns the slow normal only; tangential components of a
nearby record can be larger.

\section{A second fixed-time precision barrier}
\label{v54:slow-precision}

\begin{theorem}[Necessary slow preparation precision inside the rapid-selected manifold]
\label{v54:fixed-time-slow}
Fix \(T>0\), independent of amplitude.  Suppose an exact constrained initial
record lies in \(\mathcal S_a^\chi\), has nonzero slow normal coordinate
\(\omega(0)\), and its unchanged original-action history remains inside the
verified slow chart throughout \([0,T]\).  Then
\begin{equation}
 \boxed{
       |\omega(0)|\le C\exp(-\mu_sT/a).}             \label{v54:omega-fixed-T}
\end{equation}
Equivalently, its physical slow-normal displacement satisfies
\begin{equation}
 \boxed{
       \delta_{\rm sl,phys}(0)
          \le Ca\exp(-\mu_sT/a).}                   \label{v54:slow-physical-fixed-T}
\end{equation}
The same statements hold with a varying positive horizon \(T_a\) in place
of \(T\).
\end{theorem}

\begin{proof}
Residence in the slow chart bounds \(|\omega(T)|\) by a fixed local
normal radius.  The lower estimate in \cref{v54:slow-growth} gives
\[
   C> |\omega(T)|
       \ge c_s e^{\mu_sT/a}|\omega(0)|.
\]
This proves \eqref{v54:omega-fixed-T}.  The physical scaling from the slow
lift then gives \eqref{v54:slow-physical-fixed-T}.  If the available local
slow chart is smaller, the right side only improves; no lower bound on its
radius is required for this necessary condition.
\end{proof}

\begin{corollary}[Algebraic slow errors exit on the \(a\log(1/a)\) scale]
\label{v54:slow-algebraic-exit}
Suppose a family of exact initial records is already on
\(\mathcal S_a^\chi\) and has
\(
       |\omega(0)|\ge c a^m
\)
for some fixed real \(m\), while remaining initially inside the verified
slow tube.  Then, if the exact history exists until its first slow-chart exit,
\begin{equation}
 \boxed{
 t_{\rm exit}(a)
   \le \frac{a}{\mu_s}
       \{m\log(1/a)+C_m\}
   =O\!\left(a\log(1/a)\right).}                    \label{v54:slow-exit}
\end{equation}
Thus exact rapid selection does not make polynomial slow-normal accuracy
sufficient for a fixed positive physical horizon.
\end{corollary}

\begin{proof}
Solve the lower bound in \eqref{v54:slow-growth-eq} for the time at which the
fixed slow-chart radius is reached.  The fixed radius contributes only the
constant \(C_m\).  If the proposed algebraic displacement already lies outside
the smaller slow tube, the statement is vacuous rather than contradictory.
\end{proof}

\section{The anisotropic hierarchy of exact-action selection scales}
\label{v54:hierarchy}

We can now state the fixed-time geometry of the verified local constraint
leaf without conflating the two hyperbolic scales.  There is a nested chain
\begin{equation}
 \boxed{
 \mathcal D_a^{\chi,\sigma}
      \subset \mathcal S_a^\chi
      \subset \mathcal P_a,}                         \label{v54:nested-chain}
\end{equation}
where \(\mathcal P_a\) is the \(323\)-dimensional normalized original
constraint leaf, \(\dim\mathcal S_a^\chi=312\), and
\(\dim\mathcal D_a^{\chi,\sigma}=311\).

\begin{theorem}[Two distinct exponential preparation scales for a fixed physical time]
\label{v54:two-scale-precision}
Fix \(T>0\).  Consider any local tubular family around
\(\mathcal D_a^{\chi,\sigma}\) with the property that every exact-action
history starting in it remains in the corresponding verified charts through
time \(T\).  Then its physical thicknesses in the two proved normal groups
must satisfy
\begin{align}
 \boxed{r_a^{\rm rapid}
      \le C a\exp(-\mu_rT/a^3),}                    \label{v54:rapid-radius}\\
 \boxed{r_a^{\rm slow}
      \le C a\exp(-\mu_sT/a),}                     \label{v54:slow-radius}
\end{align}
where \(\mu_r>0\) is the V53 rapid normal constant and
\(\mu_s=\gamma/2\) is \eqref{v54:mu-s}.  In particular, neither normal
group can have algebraic fixed-time thickness.
\end{theorem}

\begin{proof}
The rapid estimate is exactly \cref{v53:no-algebraic-basin}.  Restricting
the tube to its intersection with \(\mathcal S_a^\chi\) removes the rapid
normal and leaves a tube around \(\mathcal D_a^{\chi,\sigma}\) inside the
rapid-selected manifold.  Apply \cref{v54:fixed-time-slow} to its slow-normal
radius.  The two estimates use different exact normal factorizations and
therefore retain their different physical exponents.
\end{proof}

The hierarchy has a useful interpretation.  A generic algebraic error in the
rapid normal leaves the verified tube on the earlier scale
\(O(a^3\log(1/a))\).  If the rapid normal is selected exactly but the slow
normal is only algebraically accurate, the next obstruction occurs on
\(O(a\log(1/a))\).  These are local residence bounds, not curvature blowup
statements after exit.

\section{What exact selection still leaves open}
\label{v54:frontier}

The nested manifold removes both currently certified forward expanding
normal groups without deleting the favorable rapid stable directions.  This
is a stronger structural target for an exact-action Einstein branch than the
codimension-eleven manifold of V53.

It still does \emph{not} prove a fixed-positive-time history.  A trajectory
starting exactly on \(\mathcal D_a^{\chi,\sigma}\) has both normal coordinates
identically zero while it stays in the chart, but its tangential dynamics
contains the intermediate oscillatory modes, the slow contracting mode, zero
and bounded modes, and the nonlinear metric/link observations.  Controlling
those tangential variables for one fixed physical time is a separate problem.

The result also sharpens the physical-selection issue.  At fixed time \(T\),
access to a robust local branch would require, in the currently verified
coordinates, rapid accuracy of order \(e^{-cT/a^3}\) and an additional slow
accuracy of order \(e^{-cT/a}\).  No microscopic renewal mechanism producing
those two accuracies has been proved.  The invariant graphs can depend on the
analytical localizations; they are mathematical selections, not yet native
chronological selectors.

The next substantive gates are therefore:
\begin{enumerate}[label=\textup{(\roman*)},leftmargin=2.8em]
\item fixed-positive-time propagation \emph{on}
\(\mathcal D_a^{\chi,\sigma}\), with the original metric and link reads;
\item a source-native mechanism selecting this codimension-twelve manifold,
or a proof that the required exponential accuracies arise from the finite
renewal law;
\item a cofinal spatial-refinement family carrying the two nested splittings,
source response and chart constants uniformly.
\end{enumerate}
Another improvement of the finite-node reconstruction would not address these
three gates.  The raw-action Einstein continuum bridge remains open.

\section{Verification and append-only preservation}
\label{v54:verification}

The new checker verifies the slow/physical time conversion, the inherited
numerical enclosure for \(\gamma\), the dimensions in
\eqref{v54:nested-chain}, the two exponential precision scales, the algebraic
exit-time exponents, and byte preservation of every preceding marked body.
The nested-manifold theorem and exact slow normal factor are the written
finite-dimensional Lyapunov--Perron consequences of the inherited source
splitting; they are not inferred from new numerical trajectories.

No new gravitational coefficient extraction, numerical graph radius,
spatial-refinement computation, nonlinear spacetime integration, or
proof-assistant formalization is claimed.  The V53 preservation/scaling
checker is rerun separately as a regression.  Every action-specific statement
remains local to the same \(N=3\) original-action chart, and \(a\) remains
field amplitude rather than lattice spacing.
% END IMMUTABLE EXACT-ACTION BRIDGE V54 BODY


% BEGIN IMMUTABLE EXACT-ACTION BRIDGE V55 BODY
\clearpage
\renewcommand{\thesection}{\arabic{section}}
\renewcommand{\theHsection}{bridgeV55.\arabic{section}}
\setcounter{section}{443}

\section{Upstream from the nested spectrum: the slow-time principal is state dependent}
\label{v55:context}

V54 removes the two currently certified forward expanding normal groups: the
rapid eleven-dimensional group through \(\mathcal S_a^\chi\), and the single
slow real direction through \(\mathcal D_a^{\chi,\sigma}\).  It deliberately
does not prove that a history lying exactly on the nested manifold retains the
small-rate, small-curvature regime for physical times of order \(a\).

A direct attempt to extend the V51 boundary construction from
\(t=O(a^2)\) to \(t=O(a)\) reveals why the frozen spectral information is not
by itself sufficient.  On the longer scale the canonical/active free
coordinates can be of order one in the V51 balanced variables while the weak
free coordinate is naturally only of order \(a\).  After this anisotropic
rescaling, the nonlinear weak row contains an order-\(a^{-1}\)
\emph{state-dependent principal term}.  It is on the same slow-time hierarchy
as the intermediate oscillatory response and cannot be put into a small
source without changing the problem.

This body extracts that term from the already proved analytic V51 field and
reduces the first long-time compatibility problem to a finite five-component
map.  Three components are constant-shift means and two are the nonconstant
weak-zero tangency conditions.  No new numerical gravitational coefficient is
inserted.  Instead we identify the exact five contractions that must be
computed next if the existing \(O(a)\)-rate/curvature bootstrap is to be
continued to \(t=O(a)\).

\paragraph{Inherited inputs and nonduplication.}
We use the exact constrained chart and reduced field of
\cref{v51:chart,v51:field}, the rank-24 weak block and the two secular central
chains of \cref{v45:jordan-theorem}, the transport identity
\eqref{v51:b-cancel}, and the nested spectral selection of V54.  The present
result is not another invariant-graph construction and does not alter the
original action.  In particular, the five-row map below is obtained from the
actual Taylor jets of the V51 reduced field at the same V29 root.

\section{An anisotropic chart for physical time \texorpdfstring{$t=a\tau$}{t=a tau}}
\label{v55:anisotropic}

Write the V51 free coordinate as
\[
       w=(x,y),\qquad x=(q,p_A)\in\R^{294},\qquad y=p_W\in\R^{29}.
\]
The unbalanced coordinate used in the proof of \cref{v51:balance} is
\[
       v=(a x,y).
\]
For the physical scale
\begin{equation}
                         t=a\tau,
 \qquad s=t/a^3=\tau/a^2,                           \label{v55:slow-time}
\end{equation}
introduce the further weak rescaling
\begin{equation}
                         y=a z,
 \qquad \xi=(x,z)\in\R^{294}\oplus\R^{29}.          \label{v55:xi}
\end{equation}
Thus
\begin{equation}
                         v=a\xi.                    \label{v55:vaxi}
\end{equation}
For bounded \(\xi\), the physical record reconstructed by V51 differs from
\(Q_a^0\) by \(O(a^2)\) in every canonical and multiplier component:
\begin{equation}
 \boxed{
  X-X_a^0=O_R(a^2),\qquad
  \lambda_A-\lambda_{a,A}^0=O_R(a^2),\qquad
  \lambda_W-\lambda_{a,W}^0=O_R(a^2).}             \label{v55:physical-tube}
\end{equation}
Here and below the subscript \(R\) means uniformly on \(|\xi|\le R\).
This is an anisotropic free-coordinate tube but an isotropic \(O(a^2)\)
physical-record tube.

Let \(R_s^{\rm red}\) and \(R_f^{\rm red}\) be the exact slow and weak
remainders already appearing in the proof of \cref{v51:balance}, so that in
unbalanced coordinates
\begin{equation}
 \begin{split}
 (v_s)'&=a^3b_s+R_s^{\rm red}(a,v),\\
 (v_f)'&=C_{\rm red}v_s+Bv_f+a^3b_w+R_f^{\rm red}(a,v),
 \end{split}                                        \label{v55:unbalanced-exact}
\end{equation}
where prime denotes \(s\)-time.  The inherited row estimates are
\begin{align}
 |R_s^{\rm red}(a,v)|
   &\le C\{a^4+a^2|v|+a|v|^2\},                    \label{v55:Rs-order}\\
 |R_f^{\rm red}(a,v)|
   &\le C\{a^4+a|v|+|v|^2\},                       \label{v55:Rf-order}
\end{align}
and the remainders are analytic at \((a,v)=(0,0)\).

\begin{lemma}[Exact slow-time expansion on bounded anisotropic states]
\label{v55:slow-expansion}
For every fixed \(R<\infty\) there are analytic maps
\[
 \Sigma_0:\R^{323}\to\R^{294},\qquad
 \Phi_1:\R^{323}\to\R^{29},
\]
polynomial of degree at most two, and uniformly bounded analytic remainders
\(\Sigma_1(a,\xi)\), \(\Phi_2(a,\xi)\) on
\(0\le a\le a_R\), \(|\xi|\le R\), such that the unchanged original-action
reduced equations are exactly
\begin{align}
 x_\tau
   &=b_s+\Sigma_0(\xi)+a\Sigma_1(a,\xi),             \label{v55:x-slow}\\
 z_\tau
   &=a^{-2}\Phi_0(\xi)+a^{-1}\Phi_1(\xi)
       +b_w+\Phi_2(a,\xi),                           \label{v55:z-slow}
\end{align}
where
\begin{equation}
                    \boxed{\Phi_0(x,z)=C_{\rm red}x+Bz.}       \label{v55:Phi0}
\end{equation}
Moreover
\begin{align}
 \Sigma_0(\xi)
   &=\lim_{a\downarrow0}a^{-3}R_s^{\rm red}(a,a\xi),           \label{v55:Sigma0-limit}\\
 \Phi_1(\xi)
   &=\lim_{a\downarrow0}a^{-2}R_f^{\rm red}(a,a\xi),           \label{v55:Phi1-limit}
\end{align}
and in terms of the original Taylor jets
\begin{equation}
 \boxed{
 \Phi_1(\xi)
 =D_aD_vR_f^{\rm red}(0,0)[\xi]
  +\frac12D_v^2R_f^{\rm red}(0,0)[\xi,\xi].}        \label{v55:Phi1-jets}
\end{equation}
An analogous formula for \(\Sigma_0\) uses
\(\tfrac12D_a^2D_vR_s^{\rm red}(0,0)\) and
\(\tfrac12D_aD_v^2R_s^{\rm red}(0,0)\).
\end{lemma}

\begin{proof}
Equation \eqref{v55:vaxi} and \(\tau=a^2s\) give
\(v'=a^3\xi_\tau\).  Substitution into
\eqref{v55:unbalanced-exact} therefore gives the exact identities
\[
 x_\tau=b_s+a^{-3}R_s^{\rm red}(a,a\xi),
\]
and
\[
 z_\tau=a^{-2}(C_{\rm red}x+Bz)+b_w
             +a^{-3}R_f^{\rm red}(a,a\xi).
\]
The row bounds imply the required analytic divisibilities.  More explicitly,
\eqref{v55:Rs-order} and analyticity force
\[
 R_s^{\rm red}(0,v)=0,\qquad
 D_vR_s^{\rm red}(a,0)=O(a^2),\qquad
 D_v^2R_s^{\rm red}(a,0)=O(a),
\]
whereas \eqref{v55:Rf-order} forces
\[
 R_f^{\rm red}(0,0)=0,\qquad
 D_vR_f^{\rm red}(a,0)=O(a),
\]
with \(D_v^2R_f^{\rm red}\) bounded.  Taylor's formula at \((0,0)\), after
setting \(v=a\xi\), gives
\[
 R_s^{\rm red}(a,a\xi)=a^3\Sigma_0(\xi)+a^4\Sigma_1(a,\xi)
\]
and
\[
 R_f^{\rm red}(a,a\xi)=a^2\Phi_1(\xi)+a^3\Phi_2(a,\xi).
\]
The only terms surviving the first quotient are linear and quadratic in
\(\xi\); the same is true for the second quotient.  Reading those Taylor
coefficients gives \eqref{v55:Phi1-jets} and the stated formula for
\(\Sigma_0\).  Uniform boundedness on a fixed \(\xi\)-ball follows from
analyticity on the fixed V51 chart.  No trajectory estimate is used.
\end{proof}

The important point is the middle term of \eqref{v55:z-slow}.  It is
\(a^{-1}\Phi_1(\xi)\), not an \(O(a)\) perturbation in slow time.  Thus the
state dependence of the fast weak row is principal on the proposed
\(t=O(a)\) scale.

\section{The 297-dimensional leading critical sheet}
\label{v55:critical-sheet}

Let \(\mathcal H\subset\R^{29}\) be the three-dimensional constant-shift
space and let \(\mathcal Z=\ker B\).  V45 proves
\[
       \dim\mathcal Z=5,
       \qquad \mathcal Z=\mathcal H\oplus\mathcal Z_g,
       \qquad \dim\mathcal Z_g=2,
\]
with the direct sum understood in the inherited positive \(D_*\)-pairing on
the weak center.  Denote by \(p_0\) the corresponding weak zero projection,
and by \(p_g\), \(p_H\) its projections onto \(\mathcal Z_g\) and
\(\mathcal H\).

\begin{theorem}[Rank, kernel and cokernel of the leading slow constraint]
\label{v55:critical-rank}
The map \(\Phi_0:\R^{323}\to\R^{29}\) in \eqref{v55:Phi0} satisfies
\begin{equation}
 \boxed{
 \rank\Phi_0=26,
 \qquad \dim\ker\Phi_0=297,
 \qquad \dim\operatorname{coker}\Phi_0=3.}          \label{v55:critical-count}
\end{equation}
Using the inherited positive \(D_*\)-pairing to identify weak vectors with covectors, its cokernel is exactly the constant-shift space \(\mathcal H\).  Equivalently,
\begin{equation}
 p_H\Phi_0\equiv0,
 \qquad
 p_g C_{\rm red}:\R^{294}\longrightarrow\mathcal Z_g
 \quad\hbox{is onto}.                               \label{v55:cokernel-H}
\end{equation}
The linear critical sheet
\begin{equation}
                \boxed{\mathcal K_0:=\ker\Phi_0}                 \label{v55:K0}
\end{equation}
has dimension \(297\).  Under the V44 constrained-center identification it is
precisely the kernel of the limiting central generator counted in
\eqref{v45:jordan-count}.
\end{theorem}

\begin{proof}
V45 proves \(\rank B=24\) and \(\dim\ker B=5\).  Its proof of the two
nilpotent chains proves, on the same constrained tangent represented by the
V51 \(x\)-coordinates,
\[
                  \rank(p_0C_{\rm red})=2,
 \qquad p_0C_{\rm red}x\in\mathcal Z_g.
\]
Hence the image of \(C_{\rm red}\) contributes exactly two directions not
already in \(\Ran B\).  Therefore
\(\rank[C_{\rm red}\ B]=24+2=26\), proving the first two statements of
\eqref{v55:critical-count}.  The weak codomain has dimension \(29\), so the
cokernel has dimension three.

The constant-shift tests annihilate both \(B\) and the constrained coupling
\(C_{\rm red}\) by the exact translation identities used in V45.  Thus
\(\mathcal H\) is contained in the cokernel.  Since both have dimension
three, they are equal.  Finally, the limiting constrained central field is
\(J_{c0}(x,z)=(0,\Phi_0(x,z))\) after the already eliminated noncentral
responses.  Its kernel therefore has the same \(297\) dimensions, agreeing
with \eqref{v45:jordan-count}.
\end{proof}

This count explains why the longer-time problem is not captured by V54's
single remaining unstable normal.  The nested manifold has dimension \(311\),
whereas the leading small-rate critical sheet has dimension \(297\).  The
fourteen-dimensional difference contains favorable rapid-stable directions
and bounded fast/oscillatory directions which need not be removed from the
exact dynamics, but whose amplitudes cannot be treated as order-one slow
variables if the V51 \(O(a)\)-rate bootstrap is to survive to \(t=O(a)\).

\section{The first moving-principal correction and the five-row gate}
\label{v55:five-gate}

The leading equation \(\Phi_0(\xi)=0\) is not enough.  The first correction
in \eqref{v55:z-slow} has a component in the three-dimensional cokernel of
\(\Phi_0\), and that component cannot be removed by solving for a fast weak
coordinate.

Define the three-mean map
\begin{equation}
 \boxed{
       \mathfrak m(\xi):=p_H\Phi_1(\xi)
       \in\mathcal H\simeq\R^3.}                    \label{v55:m-map}
\end{equation}
Because \(\Phi_1\) is at most quadratic, so is \(\mathfrak m\), and
\(\mathfrak m(0)=0\).

There is a second compatibility already at leading order.  If a limiting
slow path remains on \(\mathcal K_0\), its tangent must also lie in
\(\mathcal K_0\).  The slow equation \eqref{v55:x-slow} therefore yields the
two nonconstant zero-mode conditions
\begin{equation}
 \boxed{
 \mathfrak g(\xi)
   :=p_g C_{\rm red}\{b_s+\Sigma_0(\xi)\}
   \in\mathcal Z_g\simeq\R^2.}                      \label{v55:g-map}
\end{equation}
We collect them as
\begin{equation}
 \boxed{
 \mathfrak C_{\rm sl}(\xi)
      :=(\mathfrak g(\xi),\mathfrak m(\xi))
      \in\mathcal Z_g\oplus\mathcal H\simeq\R^5.} \label{v55:Cslow}
\end{equation}
Every coefficient in this map is a finite Taylor contraction of the original
V51 field at the certified root.  No spectral inverse at a zero frequency is
used.

\begin{theorem}[Necessary five-row compatibility for a bounded slow-time rate]
\label{v55:five-gate-theorem}
Fix \(T,R,M>0\).  Let \(a_j\downarrow0\) and let \(Q_{a_j}(t)\) be original
exact-action histories represented in the V51 chart on
\(0\le t\le Ta_j\).  Suppose their rescaled states \(\xi_{a_j}\) from
\eqref{v55:xi} satisfy
\begin{equation}
 \sup_{0\le\tau\le T}
   \{|\xi_{a_j}(\tau)|+|\partial_\tau\xi_{a_j}(\tau)|\}
       \le R+M.                                     \label{v55:bounded-rate-hyp}
\end{equation}
Then every \(C^0\) subsequential limit is in fact compatible with a
Lipschitz limit \(\xi_0=(x_0,z_0)\), and for almost every \(\tau\)
\begin{align}
 \Phi_0(\xi_0(\tau))&=0,                            \label{v55:limit-K0}\\
 (x_0)_\tau&=b_s+\Sigma_0(\xi_0),                   \label{v55:limit-x}\\
 \boxed{\mathfrak C_{\rm sl}(\xi_0(\tau))=0.}      \label{v55:five-zero}
\end{align}
If the rescaled histories converge in \(C^1\), the identities hold at every
\(\tau\).

Equivalently, any extension of the present V51 small-rate bootstrap to a
fixed slow-time interval must, in the zero-amplitude limit, evolve inside the
297-dimensional sheet \(\mathcal K_0\) and inside the zero set of the five
finite compatibility functions \eqref{v55:Cslow}.
\end{theorem}

\begin{proof}
The bound \eqref{v55:bounded-rate-hyp} gives equiboundedness and
equicontinuity, so Arzel\`a--Ascoli gives a uniformly convergent subsequence
with a Lipschitz limit.  Multiply \eqref{v55:z-slow} by \(a^2\):
\begin{equation}
 \Phi_0(\xi_a)+a\Phi_1(\xi_a)
   =a^2\{z_{a,\tau}-b_w-\Phi_2(a,\xi_a)\}.           \label{v55:compat-exact}
\end{equation}
The right side is uniformly \(O(a^2)\).  Passing to the limit proves
\eqref{v55:limit-K0}.  Equation \eqref{v55:x-slow} has uniformly bounded
right side and converges uniformly on the bounded state set, giving
\eqref{v55:limit-x} in the weak derivative sense.

Project \eqref{v55:compat-exact} with \(p_H\).  Since
\(p_H\Phi_0=0\), division by \(a\) gives
\[
       \mathfrak m(\xi_a)=O(a),
\]
and hence \(\mathfrak m(\xi_0)=0\).

Finally \(\Phi_0(\xi_0)=0\) for every \(\tau\).  Differentiating this
identity at almost every differentiability point of the Lipschitz limit gives
\[
 C_{\rm red}(x_0)_\tau+B(z_0)_\tau=0.
\]
Project with \(p_g\).  The \(B\)-term vanishes in the weak-zero projection,
and inserting \eqref{v55:limit-x} gives
\(\mathfrak g(\xi_0)=0\).  This proves all five rows.  Under \(C^1\)
convergence the same calculation is pointwise.
\end{proof}

\begin{corollary}[The bounded slow-time rate closes the inherited reader bootstrap]
\label{v55:reader-bootstrap}
Under the hypotheses of \cref{v55:five-gate-theorem}, the physical records
satisfy, uniformly on \(0\le t\le Ta\),
\begin{equation}
 \boxed{
   X_t=O(a),\qquad \lambda_t=O(a),\qquad X_{tt}=O(a).}            \label{v55:physical-rates}
\end{equation}
Consequently the same original reconstruction estimate used in
\cref{v51:base-curvature} gives
\begin{equation}
 \boxed{
 \|\operatorname{Riem}(g)\|+
 \|E^{\rm link}\|+\|F^{\rm sp,link}\|\le C_{R,M,T}a.}           \label{v55:reader-small}
\end{equation}
Thus the five-row condition is necessary for extending the \emph{present
proved small-rate route} to an \(O(a)\)-curvature history on the scale
\(t=O(a)\).
\end{corollary}

\begin{proof}
In the V51 chart, \(q\) and \(p_A\) are components of \(x\), whereas
\(p_W=az\).  Since \(t=a\tau\), the bounded derivative in
\eqref{v55:bounded-rate-hyp} and the normal-map estimates
\eqref{v51:normal-orders}--\eqref{v51:X-jets} give
\[
 X_t=O(a),\qquad
 (\lambda_A)_t=O(a),\qquad
 (\lambda_W)_t=a z_\tau=O(a).
\]
The unchanged canonical equation is \(X_t=F(X,\lambda)\).  Its first
stationary derivatives are bounded on the physical tube
\eqref{v55:physical-tube}; hence
\[
 X_{tt}=D_XF\,X_t+D_\lambda F\,\lambda_t=O(a).
\]
The original metric and ordered-link formulas used in V51 then give
\eqref{v55:reader-small}.  No second multiplier derivative is required by
that reconstruction estimate.
\end{proof}

The theorem and corollary concern the \emph{existing sufficient small-rate
route} to \(O(a)\) geometry.  It is not asserted that every small-curvature
history must have all multiplier rates \(O(a)\); in particular, a separate
reader-specific cancellation could in principle make a larger constant-shift
rate geometrically harmless.  V55 does not assume such a cancellation.

\section{First-order correction and the exact remaining three means}
\label{v55:first-correction}

The preceding compatibility has an equivalent first-correction formulation.
Choose once and for all a right inverse
\[
 R_0:\Ran\Phi_0\longrightarrow(\ker\Phi_0)^\perp
\]
of \(\Phi_0\) in any fixed Euclidean splitting.  It is a finite fixed
operator at the certified root.

\begin{proposition}[The three means are the only first-order algebraic obstruction]
\label{v55:first-correction-prop}
Let \(\xi_0\in\mathcal K_0\).  There exists a vector \(\xi_1\) satisfying
\begin{equation}
             \Phi_0\xi_1+\Phi_1(\xi_0)=0             \label{v55:first-correction-eq}
\end{equation}
if and only if
\begin{equation}
                         \mathfrak m(\xi_0)=0.        \label{v55:m-only}
\end{equation}
When it exists, one choice is
\begin{equation}
             \xi_1=-R_0\Phi_1(\xi_0),               \label{v55:xi1}
\end{equation}
up to an arbitrary vector in \(\mathcal K_0\).  Consequently
\begin{equation}
 \Phi_0(\xi_0+a\xi_1)+a\Phi_1(\xi_0+a\xi_1)=O_R(a^2). \label{v55:corrected-principal}
\end{equation}
\end{proposition}

\begin{proof}
By \cref{v55:critical-rank}, the cokernel of \(\Phi_0\) is precisely
\(\mathcal H\).  Thus \(-\Phi_1(\xi_0)\) lies in \(\Ran\Phi_0\) exactly
when its \(p_H\)-projection vanishes, which is
\eqref{v55:m-only}.  Applying the fixed right inverse gives
\eqref{v55:xi1}.  Taylor expansion of the quadratic map \(\Phi_1\) gives
\eqref{v55:corrected-principal}.
\end{proof}

For the explicit V51 affine reference on the longer scale,
\begin{equation}
                  \xi_{\rm aff}(\tau)=\tau b,
 \qquad b=(b_s,b_w),                                \label{v55:affine}
\end{equation}
and \eqref{v51:b-cancel} says \(b\in\mathcal K_0\).  Therefore the
singular part of its weak residual is removable through first order if and
only if
\begin{equation}
                  \mathfrak m(\tau b)=0.             \label{v55:affine-means}
\end{equation}
Since \(\mathfrak m\) is at most quadratic and vanishes at zero,
\eqref{v55:affine-means} on a nontrivial interval is equivalent to the two
finite vector conditions
\begin{equation}
 \boxed{
 D\mathfrak m(0)[b]=0,
 \qquad D^2\mathfrak m(0)[b,b]=0.}                  \label{v55:affine-six}
\end{equation}
These are six scalar contractions.  They concern only the affine reference;
a curved limiting slow path may satisfy \(\mathfrak m=0\) without satisfying
that affine ansatz.

\section{The first tangency coefficient for the actual prepared branch}
\label{v55:tangency-coefficient}

The exact initial preparation already passes the five-row gate at the initial
cut.  Indeed \(\Phi_0(0)=0\), \(\mathfrak m(0)=0\), and
\eqref{v51:b-cancel} gives
\[
        p_gC_{\rm red}b_s=0,
\]
so \(\mathfrak C_{\rm sl}(0)=0\).  A longer branch must preserve these five
conditions under its limiting slow drift.

Define the finite coefficient
\begin{equation}
 \boxed{
 \mathfrak T_{\rm sl}
    :=D\mathfrak C_{\rm sl}(0)[b]
     =\left(
       p_gC_{\rm red}D\Sigma_0(0)[b],
       \ p_HD\Phi_1(0)[b]
       \right)\in\R^5.}                             \label{v55:Tslow}
\end{equation}
Every entry is an explicit contraction of the same finite coefficient bank
used to define the original V51 field.  No nonlinear trajectory is needed to
evaluate it.

\begin{corollary}[A finite necessary tangency test at the initial cut]
\label{v55:tangency-test}
Assume a family in \cref{v55:five-gate-theorem} has
\(\xi_a(0)\to0\) and \(\partial_\tau\xi_a(0)\to b\), and assume the limiting
slow path is \(C^1\) at zero.  Then necessarily
\begin{equation}
                         \boxed{\mathfrak T_{\rm sl}=0.}          \label{v55:Tzero}
\end{equation}
Thus a single nonzero entry of \(\mathfrak T_{\rm sl}\) would rule out a
\(t=O(a)\) continuation of the present \(O(a)\)-rate branch from this
initializer, while leaving open histories with different initial slow data
or reader-specific cancellations.
\end{corollary}

\begin{proof}
The limiting path satisfies
\(\mathfrak C_{\rm sl}(\xi_0(\tau))=0\) and
\(\xi_0(0)=0\).  Its initial derivative is \(b\) by hypothesis.  Differentiate
the five identities at \(\tau=0\).
\end{proof}

This is the sharp next coefficient extraction for the present programme.
V30's transport identity and V29's initial mean preparation concern the
initial source hierarchy and are fully retained, but V55 does not identify
them with \eqref{v55:Tslow} away from the initial cut without performing the
new Taylor contractions.  In particular, no vanishing of
\(\mathfrak T_{\rm sl}\) is inferred from an earlier root calculation with a
different derivative target.

\section{Why the frozen nested spectrum alone cannot prove \texorpdfstring{$t=O(a)$}{t=O(a)}}
\label{v55:not-frozen}

The need for \(\mathfrak m\) is logical, not a weakness of the estimates.
Consider an abstract slow coordinate \(x\) and one constant weak mean \(h\)
with the same vanishing linear mean coupling at the origin,
\begin{equation}
                  x_\tau=1,
 \qquad h_\tau=a^{-1}x^2.                           \label{v55:toy}
\end{equation}
The linearization at the origin has no additional unstable eigenvalue, and
the quadratic mean term does not change any frozen spectral root there.
Nevertheless \(x(\tau)=\tau\) and
\[
                  h_\tau=\tau^2/a,
\]
so the small-rate scaling fails on every fixed positive slow-time interval as
\(a\to0\).  This is exactly the mechanism represented by a nonzero quadratic
component of \(\mathfrak m\).  The example is only a logic check; it is not a
claim that the gravitational coefficient has this sign or value.

The appropriate next step is therefore not another frozen eigenvalue count.
One must evaluate \eqref{v55:Tslow}, and if it vanishes, evaluate or solve the
five-row nonlinear system
\begin{equation}
        \Phi_0(\xi)=0,
        \qquad \mathfrak C_{\rm sl}(\xi)=0,          \label{v55:five-system}
\end{equation}
together with the regular slow drift \eqref{v55:limit-x}.  Passing this gate
would remove the order-\(a^{-1}\) moving-principal obstruction and leave a
regular finite-dimensional slow descriptor to be controlled on fixed
\(\tau\)-intervals.

\section{Consequences for the Einstein-bridge programme}
\label{v55:frontier}

V55 changes the status of the proposed \(t=O(a)\) extension in three ways.

\begin{enumerate}[label=\textup{(\roman*)},leftmargin=2.8em]
\item The longer scale has a rigorously identified anisotropic state
normalization: bounded \(\xi=(x,z)\) means an \(O(a^2)\) physical displacement
from the prepared record.
\item The leading small-rate sheet is not all of V54's 311-dimensional nested
center--stable manifold.  It is the 297-dimensional kernel
\(\mathcal K_0\), and its first moving-principal correction has exactly three
unremovable mean components.
\item Preservation of that sheet adds two nonconstant zero-mode tangency rows,
so the first nonlinear slow gate is the explicit five-component map
\(\mathfrak C_{\rm sl}\).  The initial branch passes the zeroth-order gate;
its first uncomputed test is the finite vector \(\mathfrak T_{\rm sl}\).
\end{enumerate}

If \(\mathfrak T_{\rm sl}\ne0\), the current small-rate branch cannot extend
to a fixed slow-time interval while preserving the V51 rate scaling.  If it
vanishes, that is not yet a \(t=O(a)\) theorem: one must solve the curved
five-row descriptor and bound the original metric/link readers along the
resulting exact history.  Either outcome is informative and requires only a
finite same-root coefficient calculation, not another global spectral search.

The stronger fixed-positive-physical-time problem remains beyond this scale.
All statements continue to be fixed-cutoff \(N=3\) statements with \(a\) the
field amplitude.  No cutoff-uniform refinement, native microscopic selector,
or raw-action Einstein continuum theorem is claimed.

\section{Verification and append-only preservation}
\label{v55:verification}

The V55 checker verifies the two time rescalings, the physical coordinate
powers in \eqref{v55:physical-tube}, the exact dimensions
\(24+2=26\), \(323-26=297\), and \(29-26=3\), the first-correction
solvability algebra, the affine quadratic reduction
\eqref{v55:affine-six}, all new internal references, and byte preservation of
every V1--V54 immutable body.

The checker does not assign a numerical value to
\(\mathfrak T_{\rm sl}\).  That contraction is the newly isolated source
calculation for the next iteration.  No numerical nonlinear trajectory,
new root certificate, new graph radius, spatial-refinement computation, or
proof-assistant formalization is claimed.
% END IMMUTABLE EXACT-ACTION BRIDGE V55 BODY


% Future iterations append here.

\end{document}
